\documentclass[11pt]{article}
\usepackage[T1]{fontenc}
\usepackage[utf8]{inputenc}
\usepackage[a4paper,margin=25mm]{geometry}
\usepackage{amsmath,amssymb,amsthm,mathtools}
\usepackage{longtable,booktabs,array}
\usepackage{hyperref,url}
\usepackage{fancyvrb}
\hypersetup{colorlinks=true,linkcolor=blue,citecolor=blue,urlcolor=blue}
\newtheorem{theorem}{Theorem}[section]
\newtheorem{proposition}[theorem]{Proposition}
\newtheorem{lemma}[theorem]{Lemma}
\newtheorem{corollary}[theorem]{Corollary}
\theoremstyle{remark}\newtheorem{remark}[theorem]{Remark}
\newcommand{\Q}{\mathbf Q}\newcommand{\Z}{\mathbf Z}
\newcommand{\C}{\mathbf C}\newcommand{\R}{\mathbf R}
\newcommand{\Nm}{\operatorname{Nm}}
\newcommand{\cl}{\operatorname{cl}}
\newcommand{\CH}{\operatorname{CH}}
\setlength{\emergencystretch}{3em}
\title{HODGE\_DIVISION\\Effective detectors for Weil classes}
\author{HODGE\_DIVISION research run\\GPT-6 Astra Pro; inherited effort not externally exposed}
\date{21--22 November 2026 --- active research draft}
\begin{document}
\maketitle

\begin{abstract}
This active draft records the proofs and exact certificates reached in the
HODGE\_DIVISION run. It establishes an explicit Gaussian fourfold detector,
a quadratic refinement of the division step, the exact scope and rank of
the genus-four cyclic sixfold sources, an independently reconstructed
$G_{34}$ split Eisenstein eightfold source, and exact maximal image ranks
for the specified imprimitive sources at $n=5$. It also derives an effective
bound for every Eisenstein sixfold determinant by compensation. Two further
exact reflection constructions give split sixfold bounds over
$\mathbf Q(\sqrt{-7})$ and $\mathbf Q(\sqrt{-2})$. An actual semiregular
perfect complex, with its complete normalized character and ordinary
negative-Ext vanishing, is constructed at a balanced Gaussian CM sixfold.
Its coherent variant supplies an explicit numerical bound on every
split sixfold datum over every class-number-one imaginary quadratic
field, using the supplied Perry input. The proof controls the entire
regularity and integral Chern-cycle construction. A fixed lattice map
and an effective surface contraction then give numerical bounds on
every fourfold determinant over those fields.
For the Gaussian nonsplit sixfold determinant class three, two
independent calculations establish a dominant connected $\mu_4W(E_8)$
family: an exact Hodge intersection and a rigorous interval certificate
for the nine-by-nine period differential. A reconstructed cyclic
branch-exponent argument and explicit Fermat-twelve cycle suppliers
then give an integral numerical detector bound for every fixed
polarization and order in that signed determinant class. An explicit
nonsplit lattice comparison charges the entire target isogeny cost.
Each assertion's actual check history is in \texttt{CHECK\_EVENTS.tsv};
unresolved gates and conditional uses are kept separate. No unrefereed
Severn theorem is a premise of the reconstructed eightfold proof.
\end{abstract}
\tableofcontents

\section{Scope, normalization, and the accepted inputs}

Let $(A,\theta,K)$ be a polarized abelian variety of dimension $2n$, with
imaginary quadratic action compatible with Rosati and signature $(n,n)$.
The rational Weil plane is
\[
 W_K(A)=\bigwedge_K^{2n}H^1(A,\Q)\subset H^{2n}(A,\Q).
\]
Let $e_W$ denote its algebraic character projector. All degrees in this
paper are unnormalized intersection degrees:
$\deg_\theta Z=\int_Z\theta^{\dim Z}$.
An integral effective cycle has nonnegative integer multiplicities.
A detector can always be chosen to be a reduced irreducible component
with nonzero Weil projection. We write
\[
 \delta_W(A,\theta)=\min\{\deg_\theta Y:e_W[Y]\ne0,\ Y
 \text{ an integral }n\text{-subvariety}\}.
\]
For $0\ne t\in W_K(A)$, let $\delta\ge1$ satisfy $\delta t$ integral
and put $T=(-1)^n\int_A t^2>0$. The division bound is
\[
 c_t=\min\{\max(N,\deg Z^+,\deg Z^-):
 N\ge1, Nt=[Z^+]-[Z^-],\ Z^\pm\text{ integral effective}\}.
\]
The minimum is infinity when no such presentation exists.

We use the supplied reductions, rather than reprove them: HECKE\_DIV
R21--R23 and P0, GHOSTS B06, and the attack R1--R4/Theorem W;
the attack Theorem D and Lemma Dom; and Propositions O, O$'$ and
Corollary O$''$. Their scope includes fixed polarized finite-type
families, a trivialized Weil local system after finite cover, and
bounded-degree proper cycle loci. A dense set with one degree bound
therefore gives the same bound on the whole fixed family. Positivity
is used to divide classes already supplied by cycles.

The supplied GHOSTS directory contains its registers and bibliography,
but no paper. Its R12 and R20--R24 are consequently used only with the
status supplied by the intake and attack. The quoted Eisenstein
fourfold and sixfold numbers are not silently promoted to new
whole-proof checks. Markman, Severn and 032 retain their requested
preprint or unrefereed status whenever a theorem of theirs is used.
The constructions explicitly reproved here are distinguished from
such conditional uses. Exact computational certificates use rational,
quadratic-field, or finite-field arithmetic; they are not Lean proofs.

\part{Checked constructions and deductions}
% Root research draft. Integration and whole-proof audit remain pending.
\section{Gaussian fourfolds: polarization, dominance, and integral degrees}

Throughout this section $\pi:C\to X$ is a connected etale cyclic cover of
degree four, $g(X)=3$, and $\sigma$ generates its deck group. Write
$J=J(C)$, $\Theta=\Theta_C$, $D=C/\langle\sigma^2\rangle$, and
\[
 B=\operatorname{im}(1-\sigma^2)\subset J,
 \qquad P=\ker(\operatorname{Nm}_{C/X})^0.
\]
Thus $g(C)=9$, $g(D)=5$, $\dim B=4$, and $\dim P=6$.
Intersection degrees in $B$ will be measured against its principal Prym
polarization $\xi$, whereas degrees in $J$ use $\Theta$ or the explicitly
specified very ample class $3\Theta$.

\begin{proposition}[HD-G01: the kernel is the full two-torsion]
\label{HD-G01}
The restricted Jacobian polarization on $B$ equals $2\xi$. Its type is
$(2,2,2,2)$ and its polarization kernel is $B[2]$, of order $2^8$.
\end{proposition}
\begin{proof}
The map $C\to D$ is an etale double cover and its connected norm kernel
is precisely $B$. Apply Mumford, \emph{Prym Varieties I}, Section~3,
Corollary~2, immediately after Corollary~1 on printed p.~332.
That result identifies the restriction, not just its degree: its kernel
is the full two-torsion and it is twice a principal polarization.
The dimensions above then give the asserted type and order.
This is a cited classical result applied to the intermediate cover, not
a new proof of the Prym principal-polarization theorem.
\end{proof}

\begin{proposition}[HD-G06: the two principal Gaussian lattice types]
\label{HD-G06}
The principal Prym component in HD-G04 has even Gaussian Hermitian
lattice $U\oplus U$. The balanced CM point $E^2\times E^2$ with product
principal polarization has the odd lattice $I_{2,2}$. For every
principally polarized Gaussian Weil fourfold with integral maximal-order
action, either lattice type admits the uniform bound
\[
 \delta_W\le D_0(\mathbf Q(i))
       =4D_{\mathrm{Pr}}=3\,656\,994\,324\,480.
\]
For the Prym type the sharper bound of HD-G04 applies. In HD-G05 one
may use either value of $D$ according to this stated normalization.
\end{proposition}
\begin{proof}
We give the lattice details because a principal polarization does not
by itself specify which integral Gaussian component is meant.
For an etale four-cover choose a symplectic marking of the base so
that its character is nontrivial on only one handle. For each of the
other two handles, the anti-invariant homology has the basis
\[
 a_0-a_2,\ a_1-a_3,\ b_0-b_2,\ b_1-b_3,
\]
where subscripts indicate sheets. Its Jacobian intersection form is
twice the standard symplectic form; after passing to $\xi$ it is
the standard form. The deck action sends the first difference to the
second and the second to minus the first, and does the same for the
$b$'s. In the Gaussian basis consisting of the first $a$- and
$b$-differences the Hermitian matrix is
$\bigl(\begin{smallmatrix}0&i\\-i&0\end{smallmatrix}\bigr)$, which is
isometric to $U$ by a unit change of basis. These independent handle
blocks exhaust rank eight over $\mathbf Z$. This proves the first
claim. The product CM polarization instead gives diagonal signs
$(1,1,-1,-1)$ under the balanced action.

For completeness, the elementary classification required here is as
follows. A principal polarization with integral $\mathbf Z[i]$ action
gives a free unimodular Hermitian lattice of signature $(2,2)$.
Its rational Hermitian space is split: this follows from the
signature-and-determinant classification in Schoen (1998), Section~7,
since its determinant is one. Choose a rational isotropic vector,
scale it to a primitive lattice vector $e$, and use unimodularity to
choose $f$ with $H(e,f)=1$. Their span has matrix
$\bigl(\begin{smallmatrix}0&1\\1&k\end{smallmatrix}\bigr)$,
$k\in\mathbf Z$, and is a unimodular orthogonal summand. Replacing
$f$ by $f+ce$, $c\in\mathbf Z[i]$, changes $k$ by an even integer.
An even $k$ gives $U$; an odd $k$ gives $I_{1,1}$.
The complement is again unimodular and rationally split of signature
$(1,1)$, so the same argument applies. Mixed sums of $U$ and an odd
plane are isometric to $I_{2,2}$: on
$\langle1\rangle\oplus U$, with basis $a,e,f$, the integral
unimodular basis $a+e,a-f,a+e-f$ has norms $1,1,-1$ and is orthogonal.
Thus the only possibilities are the even $U^2$ and the odd $I_{2,2}$.

Finally, with $S=\bigl(\begin{smallmatrix}1&1\\1&-1\end{smallmatrix}\bigr)$,
\[
 S^*I_{1,1}S=2U,\qquad S^*US=2I_{1,1}.
\]
On two blocks this yields a fixed Gaussian-linear isogeny from an
even-type fourfold to any prescribed odd-type point, after transporting
the complex structure by the rational similitude. Its degree is
$2^4=16$ and its pullback of the target principal polarization is
$2\xi$. A detecting surface pushes forward to a detecting integral
cycle, and one reduced component detects. Its degree is at most
$2^2D_{\mathrm{Pr}}$. This proves the all-principal bound, with no
unbounded isogeny transport and no appeal to the core index $J$.
\end{proof}

\begin{proposition}[HD-G02: published rank four of the Gaussian Prym period map]
\label{HD-G02}
On a nonempty open subset of the moduli of connected etale cyclic
four-covers of genus-three curves, the differential of the period map
to the four-dimensional Gaussian Weil family has rank exactly four.
Consequently it dominates its polarized integral-lattice component.
\end{proposition}
\begin{proof}
The exact published assertion is van Geemen (1996),
\emph{Theta functions and cycles on some abelian fourfolds},
Theorem~5.3(3), printed pp.~629--630. It proves this Gaussian variant
of Schoen's argument. The following calculation records the source
check and its integral-component interpretation; the rank-four
dominance theorem itself is not claimed as a new result of this run.
Let $\eta$ be the associated line bundle of exact order four on $X$.
The cotangent map is the multiplication map
\[
 \mu:H^0(X,K_X\eta)\otimes H^0(X,K_X\eta^{-1})
       \longrightarrow H^0(X,K_X^2).
\]
The formula is the cyclic-character restriction of Schoen (1988),
Lemma~3.11, printed pp.~29--30; see also Koike (2004), equation~(1),
printed p.~568, for the same order-four character pair.
Both factors on the left have dimension two by Riemann--Roch and
$\eta^{\pm1}\ne\mathcal O_X$.

We verify the order-four variant of the argument on Schoen's printed
p.~30; his Theorem~3.2 itself is stated for the Eisenstein field.
If $K_X\eta$ has a base point $p$, Riemann--Roch gives
$h^0(\eta^{-1}(p))=1$, so $\eta^{-1}=\mathcal O_X(q-p)$ for some
$q\ne p$. Hence $4q\sim4p$. A rational function with divisor
$4q-4p$ defines a degree-four map $X\to\mathbf P^1$ totally ramified
at $p$ and $q$. Riemann--Hurwitz gives total ramification twelve;
these two points account for six. There are therefore at most six
other branch values. The corresponding Hurwitz loci, after the two
distinguished values are placed at zero and infinity and the remaining
one-dimensional automorphism group is removed, have dimension at
most five. There are finitely many ramification types and Nielsen
classes in degree four with this ramification bound. Their images
cannot contain an open subset of the six-dimensional $\mathcal M_3$.
Thus $K_X\eta$ is base-point-free for a general pair $(X,\eta)$.

The base-point-free pencil trick now identifies $\ker\mu$ with
$H^0(X,\eta^{-2})=0$, because $\eta$ has exact order four.
Thus $\operatorname{rank}\mu=2\cdot2=4$, and the differential has the
same rank. The target component has dimension four, giving dominance.
This supplies the actual Gaussian rank calculation; the source locator
for Schoen's method and the distinct field in his printed theorem are
kept separate.
\end{proof}

\begin{proposition}[HD-G03: the norm fibre needed in the Bezout argument]
\label{HD-G03}
The fibre of $\operatorname{Nm}_{C/X}:\operatorname{Pic}^4(C)
\to\operatorname{Pic}^4(X)$ over $K_X$ is a union of translates of the
sixfold $P$, not translates of the fourfold $B$. Each connected
component $T_0$ has
\[
 \operatorname{type}(\Theta|_P)=(1,1,1,1,4,4),\qquad
 \deg_{3\Theta}T_0=3^6\,6!\,16=8\,398\,080.
\]
The full fibre $W_4(C)\cap\operatorname{Nm}^{-1}(K_X)$ has dimension
at most two.
\end{proposition}
\begin{proof}
The norm is surjective, and its connected kernel has dimension
$g(C)-g(X)=6$. The polarization type is the etale cyclic-cover
formula of Lange--Ortega, \emph{Prym varieties of cyclic coverings},
Section~4, first paragraph on the induced polarization (preprint
arXiv:0805.1020, p.~6), citing Birkenhake--Lange, Corollary~12.1.4 and
Lemma~12.3.1. In their notation the type is
$(1^{(d-2)(g-1)},d^{g-1})$; here $d=4,g=3$.
Riemann--Roch for polarized abelian varieties gives the displayed
degree. The morphism $\operatorname{Sym}^4 C\to\operatorname{Sym}^4 X$
is finite. The inverse image of $|K_X|\simeq\mathbf P^2$ therefore
has dimension two. Its image under the Abel map is exactly the
indicated intersection, which proves the dimension assertion.

In particular the four-dimensional degree
$3^4\,4!\,16$ cannot be substituted for the degree of this norm fibre.
The polarization calculation in HD-G01 does not justify that
substitution: it concerns a different abelian subvariety.
\end{proof}

\begin{lemma}[HD-G04a: an effective supplier component detects after projection]
\label{HD-G04a}
Some irreducible surface component $Y$ of
$W_4(C)\cap T_0$, for one connected norm-fibre component $T_0$,
has the property that $(1-\sigma^2)_*[Y]$ has nonzero Weil projection
in $B$.
\end{lemma}
\begin{proof}
Let $s:C^4\to J$ be the sum of Abel maps, and put
$u=(1-\sigma^2)s:C^4\to B$.
For each primitive character $\chi$ of the cyclic group, the
determinant line in the tensor product of four copies of
$H^1(C)_\chi$ is the pullback of the corresponding Weil line on $B$.
The alternating-tensor map is injective. Products of two classes of
the same primitive character in $H^1(C)$ vanish, since their product
would have nontrivial character $\chi^2$ in the trivial representation
$H^2(C)$. It follows that $u^*W$ lies in the all-$H^1$ Kunneth
summand of $H^4(C^4)$ and is primitive for the product polarization.
It is a real Hodge subspace of type $(2,2)$; the middle pairing on it
is positive definite by Hodge--Riemann.

Schoen (1988), Theorem~2.0, its unramified construction on printed
pp.~12--13, and Lemma~2.6 show that this determinant subspace is
spanned by the classes of signed combinations of surfaces supported
over $|K_X|$ in $C^4\to\operatorname{Sym}^4 X$. No ramified-cover
repair is used here. For nonzero $w\in W$, the nonzero class $u^*w$
therefore pairs nontrivially with at least one of the effective
irreducible supplier surfaces $Z$. By the projection formula
$u_*[Z]$ pairs nontrivially with $w$.

Put $Y=s(Z)$. This image must have dimension two: otherwise $u_*[Z]$
vanishes in the relevant dimension. The pushforward $s_*[Z]$ is a
positive integer times $[Y]$, so $(1-\sigma^2)_*[Y]$ still has
nonzero pairing. The support lies in the norm fibre because
$Z$ lies over $|K_X|$. By HD-G03 that fibre intersected with $W_4(C)$
has dimension at most two; hence $Y$ is an irreducible component.
The Weil projector is orthogonal for the middle pairing on this
primitive summand, so the nonzero pairing is precisely detection.
\end{proof}

\begin{theorem}[HD-G04: an explicit detector bound on the principal Prym component]
\label{HD-G04}
For the Gaussian principal Prym component $S_{\mathrm{Pr}}$, with
polarization $\xi$ and integral $\mathbf Z[i]$ action,
\[
 \delta_W\le D_{\mathrm{Pr}}=914\,248\,581\,120
\]
on every fibre.
\end{theorem}
\begin{proof}
The Poincare formula gives
\[
 \deg_{3\Theta}W_4(C)=3^4\frac{9!}{5!}=244\,944.
\]
For the detecting component $Y$ supplied by HD-G04a, the projective
Bezout bound for components of an improper intersection (Fulton,
\emph{Intersection Theory}, second edition, Example~8.4.6) gives
\[
 \deg_{3\Theta}Y\le244\,944\cdot8\,398\,080.
\]
Set $p=1-\sigma^2:J\to B$. The involution $\sigma^2$ preserves
$\Theta$ and is self-adjoint for its Rosati involution. Thus
$p^\dagger p=2(1-\sigma^2)$ has eigenvalues zero and four, so
$p^*(\Theta|_B)\le4\Theta$ as semipositive invariant forms.
Since $\Theta|_B=2\xi$, this says $p^*\xi\le2\Theta$.
The image $p(Y)$ is a surface and is still detecting; its reduced
image has no greater degree than its effective pushforward. Therefore
\[
 \deg_\xi p(Y)
 \le \int_Y(p^*\xi)^2
 \le4\deg_\Theta Y
 =\frac49\deg_{3\Theta}Y
 \le914\,248\,581\,120.
\]
All cycles here are reduced integral subvarieties or integral proper
pushforwards; no division of an effective cycle was performed.
HD-G02 supplies a dense set of fibres of the component. The supplied
detector closedness result (division attack, Theorem~D, equivalently
Lemma~Dom together with the Chow argument) extends the same integer
bound to the whole component.
\end{proof}

\section{A quadratic refinement of the division step}

\begin{theorem}[HD-WQ: fixed quadratic multiplication removes a Gram factor]
\label{HD-WQ}
Use the setting and projector data of the supplied Theorem~W, with
$R$ replaced by its absolute value and $\kappa$ its stated integral
degree mass. Put
\[
 q=\operatorname{Nm}(a)^n,\quad
 s=a^{2n}+\bar a^{2n}\in\mathbf Z,\quad H=4q^2-s^2>0.
\]
For a detector of integral degree $D$, a nonzero rational Weil class
$t$, and $\delta t$ integral, let $T=(-1)^n\int t^2>0$.
There exist integral effective positive and negative cycles representing
$Nt$, with
\[
 N\le\delta H\frac{R^2}{(n!)^2}D^2,
 \qquad
 \deg Z^\pm\le
 \frac{4\delta q\sqrt{HT}\,\kappa R}{n!}D^2.
\]
Ceilings make these into an explicit integral bound. In particular
the supplied quartic comparison may be sharpened to a quadratic
comparison; this does not by itself bound $D$ in a family.
\end{theorem}
\begin{proof}
Use, without re-proving, Steps~1--4 of Theorem~W. They give an
integral $y=R e_W[Y]\ne0$ represented with effective positive and
negative parts of degrees at most $\kappa D$, and
\[
 d=\lVert y\rVert^2=(-1)^n\int y^2\in\mathbf Z_{>0},
 \qquad \sqrt d\le RD/n!.
\]
For $A=[a]^*|_W$, put $y_1=y$ and $y_2=Ay$.
The second generator has parts of degrees at most $q\kappa D$.
The Rosati adjoint and the two eigenvalues give
\[
 A+A^\dagger=s\operatorname{id},\qquad
 A^\dagger A=q^2\operatorname{id}.
\]
Consequently the Gram matrix is exactly
\[
 G=d\begin{pmatrix}1&s/2\\s/2&q^2\end{pmatrix}.
\]
This fixed shape is the point that permits cancelling a common
factor from the original Cramer calculation.
Let $b_j=\langle\delta t,y_j\rangle\in\mathbf Z$, with the positive
middle pairing on $W$. Direct inversion gives
\[
 \delta Hd\,t=(4q^2b_1-2sb_2)y_1+(4b_2-2sb_1)y_2.
\]
Its multiplier $N=\delta Hd$ is a positive integer even when $s$ is odd.
The two coefficient functionals are pairing with
$4q^2y_1-2sy_2$ and $4y_2-2sy_1$. Their squared norms, computed
from the displayed Gram matrix, are respectively $4q^2Hd$ and
$4Hd$. Cauchy--Schwarz bounds their absolute values by
$2\delta q\sqrt{HTd}$ and $2\delta\sqrt{HTd}$.
Forming the indicated integral linear combination of the signed
cycle representatives gives each effective part degree at most
\[
 \kappa D\bigl(2\delta q\sqrt{HTd}
             +q\,2\delta\sqrt{HTd}\bigr)
 \le\frac{4\delta q\sqrt{HT}\,\kappa R}{n!}D^2.
\]
The bound for $N$ follows from the same estimate on $d$.
There is no cancellation by a rational coefficient in the cycle group;
the common factor was removed from the integer Gram formula before
the cycles were assembled.
\end{proof}

\begin{theorem}[HD-G05: an integral Gaussian fourfold division constant]
\label{HD-G05}
Let $(A,\xi,\mathbf Z[i])$ be a principally polarized Gaussian Weil
fourfold possessing a detector of degree at most the integer $D$.
For $\delta t$ integral and $T=\int t^2>0$, put $h=\lceil\sqrt T\rceil$.
Then
\[
 c_t(A)\le C_4(D;\delta,T)
 =\delta D^2\max\{36\,864,76\,800h\}.
\]
In particular $D=D_{\mathrm{Pr}}$ gives a fully explicit constant on
the component in HD-G04.
\end{theorem}
\begin{proof}
Write $I=[i]^*$ and $S=[1+i]^*$ on $H^4(A)$.
On the character pieces indexed by $j=0,1,2,3,4$, the eigenvalues
of $I$ are $(-1)^j$ and those of $S$ are
$-4,-4i,4,4i,-4$, respectively (the order of conjugate entries is
immaterial). Hence
\[
 (1+I)(4-S)=16e_W.
\]
For a detector $Y$ of degree at most $D$, the integral class
\[
 y=(4+4I-S-IS)[Y]
\]
has effective positive and negative parts of degrees at most $8D$:
$I$ preserves degree and $S$ multiplies surface degree by four.
The orthogonal-projector and comass estimates already proved in
Theorem~W give $d=\int y^2\le64D^2$.

Let $A_*= [2+i]^*|_W$. Its eigenvalues are $-7\pm24i$, and
$y_2=A_*y$ has positive and negative degrees at most $200D$.
The Gram matrix of $(y,y_2)$ is
\[
 d\begin{pmatrix}1&-7\\-7&625\end{pmatrix}.
\]
For $b_j=\langle\delta t,y_j\rangle\in\mathbf Z$ there is the exact
integral identity
\[
 576\delta d\,t=(625b_1+7b_2)y+(7b_1+b_2)y_2.
\]
The dual vectors $625y+7y_2$ and $7y+y_2$ have norms
$600\sqrt d$ and $24\sqrt d$. Thus $N=576\delta d$ is at most
$36\,864\delta D^2$, while each effective part of the displayed
representative has degree at most
\[
 8D\,(600+25\cdot24)\,\delta\sqrt{Td}
 \le76\,800\delta\sqrt T\,D^2.
\]
Replacing $\sqrt T$ by its integer ceiling proves the statement.
\end{proof}

% Sixfold contribution, GPT-6 Astra Pro /root/sixfolds, 2026-11-21.
% The source proofs below were read in full; exact finite checks are in
% sixfold_certificates.py and sixfold_division_certificates.py.
\section{Sixfolds: the exact cyclic source and explicit bounds}
\label{sec:sixfolds}

The natural genus-four cyclic construction reaches precisely the two
quadratic cyclotomic fields and the split rational Hermitian class.
The cubic and quartic period maps both have differential rank nine.
The published dominance results are cited at their exact locators below;
the character, determinant, degree, lattice and division consequences are
proved separately. Throughout, determinant means its class in
$\mathbf Q^\times/\operatorname{Nm}_{K/\mathbf Q}(K^\times)$,
without absorbing a dimension-dependent sign.

\begin{proposition}[SF-01: the cyclic character-field restriction]
\label{SF-01}
Let $C\to X$ be a connected etale cyclic cover with group $C_d$ and
$g(X)=4$. For $e>1$ dividing $d$, let $B_e$ be its rational cyclotomic
isotypic abelian factor of order $e$. Then
\[
 \dim B_e=3\varphi(e).
\]
If $B_e$ is a sixfold, then $e=3,4,$ or $6$, and its coefficient
field is respectively $\mathbf Q(\sqrt{-3}),\mathbf Q(i),$ or
$\mathbf Q(\sqrt{-3})$.
\end{proposition}
\begin{proof}
Write $L$ for the order-$d$ line bundle defining the cover. The nontrivial
holomorphic eigenspaces are $H^0(X,\omega_X\otimes L^j)$. Riemann--Roch
and $H^0(X,L^{-j})=0$ give dimension $g(X)-1=3$. There are
$\varphi(e)$ characters of exact order $e$, giving the displayed
formula. The equation $\varphi(e)=2$ has exactly the three stated
solutions. Passing from the original cover to its quotient of degree
$e$ does not change this cyclotomic constituent.
\end{proof}
This is a restriction on the prescribed rational isotypic construction.
It does not deny that an individual special factor could have accidental
additional endomorphisms. Such special points do not establish dominance
in a general Weil family with generic endomorphism algebra $K$.

\begin{proposition}[SF-02: the cyclic Hermitian space is split]
\label{SF-02}
For a connected etale cyclic cover of a genus-$(n+1)$ curve, the
order-$e$ primitive homology in the quadratic cases $e=3,4,6$ is a
split Hermitian space of $K$-rank $2n$. Its determinant is
$[(-1)^n]$. In particular every sixfold in SF-01 belongs to determinant
class $[-1]$.
\end{proposition}
\begin{proof}
A connected cover is specified by a surjection
$\chi:H_1(X,\mathbf Z)\to\mathbf Z/d$. It has a primitive integral
lift. Here is a direct justification. Choose integer representatives
$v_1,\ldots,v_{2g}$ and put $b=\gcd(v_2,\ldots,v_{2g})$.
If $b=0$, replace one zero coordinate by $d$. For primes dividing both
$b$ and $d$, surjectivity implies that $v_1$ is not divisible by that
prime. For each prime dividing $b$ but not $d$, choose $k$ so that
$v_1+kd$ is nonzero modulo that prime. The Chinese remainder theorem
makes these finitely many choices simultaneously. The resulting lift
has relatively prime coordinates. A primitive vector in a unimodular
symplectic lattice extends to a symplectic basis. Surjectivity of the
surface mapping-class representation onto $\operatorname{Sp}_{2g}(\mathbf Z)$
\cite[Theorem~6.4]{SF-FarbMargalit} therefore puts all monodromy on the
first handle.

Decompose $X$ into a one-holed torus and a genus-$n$ surface with one
boundary circle. The latter has $d$ disjoint lifts; the connected cover
of the former has genus one and $d$ boundary circles. On each unaffected
handle choose cycles $a_{j,k},b_{j,k}$, $1\le j\le n$,
$k\in\mathbf Z/d$, with intersection form
\[
 (a_{j,k},a_{l,r})=(b_{j,k},b_{l,r})=0,
 \qquad (a_{j,k},b_{l,r})=\delta_{jl}\delta_{kr}.
\]
The deck generator shifts the sheet index. These cycles give a
nondegenerate rank-$2nd$ summand. Mayer--Vietoris shows that the
complement is the rank-two deck-invariant homology of the covered
torus after its boundary circles are filled by the other pieces.
Consequently all nontrivial isotypic homology comes from
\[
 \mathbf Q[C_d]^n\oplus\mathbf Q[C_d]^n,
\]
whose two displayed summands are invariant, totally isotropic and
paired perfectly. The order-$e$ idempotent yields two $K$-subspaces of
dimension $n$. They are Hermitian-isotropic: all field multiples of a
vector stay in the corresponding symplectically isotropic subspace.
Nondegeneracy now gives a Hermitian matrix
\[
 \begin{pmatrix}0&A\\ A^*&0\end{pmatrix},
 \qquad \det=(-1)^n\operatorname{Nm}(\det A).
\]
This proves the assertion.
\end{proof}

\paragraph{The missing components.}
The union of the cubic and quartic sources has image dimension exactly
nine in the split Eisenstein and Gaussian families, by the rank checks
below. It has no image in the other determinant components as a map
preserving the prescribed $K$-action and polarization. An empty image
is not a period map of rank zero. Among the class-number-one fields,
all components for
\[
 \mathbf Q(\sqrt{-2}),\ \mathbf Q(\sqrt{-7}),\
 \mathbf Q(\sqrt{-11}),\ \mathbf Q(\sqrt{-19}),\
 \mathbf Q(\sqrt{-43}),\ \mathbf Q(\sqrt{-67}),\
 \mathbf Q(\sqrt{-163})
\]
are absent, as are all nonsplit Gaussian and Eisenstein components.
The standard class-number-one list is used only to spell out these
names \cite[``Class Number 1'']{SF-Stein}; the propositions themselves
have no class-number hypothesis. Explicit missing Hermitian forms are
\[
 \operatorname{diag}(1,1,1,-1,-1,-3)\quad(K=\mathbf Q(i)),
 \qquad
 \operatorname{diag}(1,1,1,-1,-1,-2)\quad(K=\mathbf Q(\sqrt{-3})).
\]
Their determinants are $-3$ and $-2$. The positive factors $3$ and $2$
are not norms in the respective fields: the primes are inert, so norm
valuations are even. Both signatures are $(3,3)$. These are
counterexamples to universal cyclic-source dominance, not to the Hodge
conjecture.

\subsection{Published differential ranks with exact finite checks}
\label{sec:sixfold-ranks}

\paragraph{SF-CW: the cubic rank is nine.}
Schoen \cite[Lemma~13, equations~(13.1)--(13.2), pp.~334--335]{SF-Schoen1998}
proves the dominance of the cubic sixfold source. His cover of
$\mathbf P^1$ has group $(\mathbf Z/3)^3$, inertia
$e_1,e_2,e_3,-e_1-e_2-e_3$, and etale subgroup generated by $e_1+e_2$.
The codifferential is multiplication of the two three-dimensional
character spaces. Their one-dimensional constituent characters are
\[
 \begin{split}
 U_+ &: (0,1,1),\ (1,0,1),\ (2,2,1),\\
 U_- &: (1,1,0),\ (1,1,1),\ (1,1,2).
 \end{split}
\]
The nine sums modulo three are
\[
 \begin{array}{ccc}
 (1,2,1)&(1,2,2)&(1,2,0)\\
 (2,1,1)&(2,1,2)&(2,1,0)\\
 (0,0,1)&(0,0,2)&(0,0,0).
 \end{array}
\]
They are distinct. Each product is nonzero because the curve is smooth
and integral. Thus the nine products are independent and the
codifferential has exact rank nine. Source and target both have
dimension nine. The Chevalley--Weil multiplicities and this finite
character computation were independently recomputed in
\texttt{sixfold\_certificates.py}; the dominance theorem remains a
published cited input.

\paragraph{SF-KO: the quartic rank is nine.}
Koike \cite[Theorem~2.1 and Proposition~3.1]{SF-Koike2004} resolves the
sixfold question recorded in van Geemen
\cite[Remark~5.4, p.~630]{SF-vanGeemen1996}. Koike's codifferential is
\[
 H^0(X,\omega_XL)\otimes H^0(X,\omega_XL^3)
 \longrightarrow H^0(X,\omega_X^2),\qquad 3\cdot3\longrightarrow9.
\]
Its exact locator is equation~(1), p.~568; the eigenvectors are (2),
p.~568; the relations are (6)--(7), p.~571; coefficient substitution
(8) and monomial basis (9) occur on pp.~571--572. The source consists
of invariant complete intersections in $(\mathbf P^1)^4$ with free
cyclic-four action. Three normalized product coordinates give an
identity block. In the remaining six coordinates, the polynomial
matrix $M$ specializes formally at $A=B=C=0$ to
\[
 M_0=\begin{pmatrix}
 1&0&0&0&0&0\\
 0&1&0&0&-2&0\\
 0&0&1&0&0&0\\
 0&0&0&1&0&0\\
 0&0&0&0&-1&0\\
 0&0&0&0&0&-1
 \end{pmatrix},\qquad \det M_0=1.
\]
This point is a witness that the determinant polynomial is nonzero;
it is not asserted to be a smooth curve. The nonempty rank-open set
therefore intersects the nonempty smooth/free-action open set of
Proposition~3.1 in the irreducible coefficient space. The differential
rank at such smooth points is nine.

As a separate finite check, at $A=(1,2,3)$, $B=(2,3,5)$,
$C=(3,5,7)$ the six-product determinant is $-389639$, and the minor
formed by the first six rows of the seven quadratic $b$-relations is
$227264$. All sixteen polynomial identities in (6)--(7) were expanded
exactly and verified over $\mathbf Q(i)$ in
\texttt{sixfold\_certificates.py}. This second parameter point also
serves only as a rank-polynomial witness. Koike's printed symbol
$\delta=1$ is not silently identified with our unsigned determinant
class; SF-02 establishes our convention $[-1]$.

\subsection{Integral detector degrees and the polarization}

\begin{theorem}[SF-03: the Gaussian principal Prym sixfold bound]
\label{SF-03}
Let $C\to X$ be an etale cyclic-four cover with $g(X)=4$, deck
generator $\sigma$, and
\[
 B=\operatorname{Prym}(C\to C/\langle\sigma^2\rangle),\qquad
 P=\ker(\operatorname{Nm}_{C/X})^0.
\]
Write $\Theta=\Theta_C$ and $\Theta|_B=2\Xi$, where $\Xi$ is the
integral principal Prym polarization. Then
\[
 \dim J(C)=13,\qquad \dim P=9,\qquad\dim B=6,
\]
$\Theta|_P$ has type $(1^6,4^3)$, and an integral effective detecting
threefold in $B$ has $\Xi$-degree at most
\[
 D_{G,6}=121\,993\,819\,950\,101\,299\,200.
\]
The same bound holds on every fibre of the connected principal Prym
Weil component dominated by these covers.
\end{theorem}
\begin{proof}
Use the unramified, $r=0$, supplier construction in Schoen
\cite[Theorem~2.0, pp.~12--13, Lemma~2.6, and \S3, pp.~24--25]{SF-Schoen1988}.
For $h=6$ its supplier cycles on $C^6$ are supported over
$|K_X|\subset\operatorname{Sym}^6X$. We spell out the detecting
component argument because its support determines the degree bound.

Let $a:C^6\to J(C)$ be the sum Abel map and $p=1-\sigma^2:J(C)\to B$.
The pullback $(p\circ a)^*$ maps the two Weil lines injectively to the
two alternating top tensors of the primitive $\pm i$ eigenspaces in
$H^1(C)^{\otimes6}$. Products of two forms from the same such
eigenspace on one curve factor vanish: their character is nontrivial
whereas $H^2(C)$ has trivial character. Thus the determinant pullback
has exactly this primitive Kunneth component. Its middle pairing is
nondegenerate, since the pairing of the two eigenspaces is perfect
and the pairing of their top exterior powers is its nonzero
determinant. The relevant primitive Kunneth projections of Schoen's
supplier classes span this summand; other Kunneth components pair
trivially with it. Therefore one supplier has nonzero pairing with a
nonzero pulled-back Weil vector. Upon decomposing it into integral
irreducible components, at least one effective component still has
nonzero pairing.

Its Abel image $Z$ and its image $p(Z)$ both have dimension three;
otherwise the corresponding pushforward would vanish in that
dimension. Each proper pushforward is a positive integral multiple
of the reduced image, so reduced images preserve detection. The norm
of the supplier divisor has class $K_X$. Hence
\[
 Z\subset W_6(C)\cap T_0,
\]
where $T_0$ is a connected component of the appropriate norm fibre,
a translate of the full $P$, not of $B$. Conversely an effective
divisor representing a point of this intersection has norm in
$|K_X|$. The map $\operatorname{Sym}^6C\to\operatorname{Sym}^6X$ is
finite. Its inverse image of the three-dimensional canonical system
has dimension at most three. It follows that
$\dim(W_6(C)\cap T_0)\le3$, and $Z$ is an irreducible component of
that intersection.

The standard etale cyclic Prym type gives $\Theta|_P$ type
$(1^6,4^3)$ and Euler characteristic $64$; see
\cite[\S4, polarization-type paragraph, author version p.~6]{SF-LangeOrtega2011},
which cites Birkenhake--Lange Corollary~12.1.4 and Lemma~12.3.1.
Equivalently, the complementary factor $\pi^*J(X)$ has induced type
$(1,4,4,4)$: the degree-four isogeny from $J(X)$ has cyclic kernel
and pulls the restricted form back to $4\Theta_X$. Complementarity
in the principal $J(C)$ gives the asserted type of $P$.
The intermediate double cover is etale, so the double-Prym theorem
gives $\Theta|_B=2\Xi$ with $\Xi$ principal
\cite[\S3, Corollary~2 after Corollary~1, p.~332]{SF-Mumford1974}.

In the embedding by $3\Theta$, Poincare's formula gives
\[
 \deg_{3\Theta} W_6(C)=3^6\frac{13!}{7!}=900\,694\,080,
 \qquad
 \deg_{3\Theta} T_0=3^9\,9!\,64=457\,124\,290\,560.
\]
The component Bezout bound, valid for improper intersections
\cite[Example~8.4.6]{SF-Fulton}, therefore gives
\[
 \deg_\Theta Z\le
 \frac{900\,694\,080\cdot457\,124\,290\,560}{27}.
\]
The deck action preserves $\Theta$ and has Rosati adjoint
$\sigma^{-1}$. Thus
\[
 p^\dagger p=2(1-\sigma^2)\le4\operatorname{id},
 \quad p^*(\Theta|_B)\le4\Theta,
 \quad p^*\Xi\le2\Theta
\]
in the semipositive order on invariant forms. The projection formula
for the effective threefold gives
\[
 \deg_\Xi(p_*Z)\le8\deg_\Theta Z.
\]
Removing its positive generic-degree multiplicity only decreases the
degree and leaves a detector. The resulting integer is
\[
 \frac8{27}\cdot900\,694\,080\cdot457\,124\,290\,560
 =121\,993\,819\,950\,101\,299\,200.
\]
No rational division of a cycle was used. Koike's rank-nine dominance
supplies a dense set of fibres; the authorized detector closedness
input, Theorem~D/Lemma~Dom of the division attack, extends this same
integer bound to the whole component.
\end{proof}

\paragraph{Why the full norm fibre matters.}
For a cyclic-four cover of a genus-$(n+1)$ curve the supplier norm
fibre has dimension $3n$ and polarization Euler characteristic $4^n$.
The primitive double Prym has dimension $2n$ and Euler characteristic
$2^{2n}=4^n$. The equal Euler characteristics do not permit replacing
the norm fibre's factorial $(3n)!$ and power $3^{3n}$ by those of
the smaller primitive Prym.

\paragraph{The Eisenstein source normalization.}
The supplied GHOSTS R20--R24 input uses the cubic cover
$C_{10}\to X_4$ and the full six-dimensional norm kernel $P_6$,
with restricted Jacobian polarization of type $(1^3,3^3)$. Its source
detector bound is
\[
 D_{E,6}=57\,854\,793\,024\,000
 =\frac{110\,224\,800\cdot14\,171\,760}{27},
\]
where
\[
 110\,224\,800=3^6\frac{10!}{4!},\qquad
 14\,171\,760=3^6\,6!\,27.
\]
The supplier already lies in a translate of this full $P_6$, so no
$1-\sigma$ projection cost occurs in that calculation. This
normalization is not called principal. An isogenous principal
quotient or another polarized lattice requires its own charged
lattice map. The source statement is used as authorized background;
its arithmetic was recomputed, but the earlier full proof is not
present in the supplied archive and is not claimed independently
reproved here.

\begin{proposition}[SF-04: a fixed split-polarization transfer]
\label{SF-04}
Fix a polarized Weil sixfold lattice with rational Hermitian space
isometric to the Gaussian source of SF-03 or the cubic source just
described. Choose one rational $K$-linear isometry
$F:V_{\rm source}\to V_{\rm target}$ of their alternating
polarization forms, and let $m\ge1$ clear its integral-lattice
denominators. Then on the target component
\[
 \delta_W\le m^6D_{\rm source},
 \qquad D_{\rm source}=D_{G,6}\ \hbox{or}\ D_{E,6}.
\]
\end{proposition}
\begin{proof}
The isometry identifies the period domains. Intersecting the two
lattice-stabilizer groups gives a finite-index subgroup, so after a
finite level cover $mF$ induces a family of $K$-equivariant
isogenies $f:B\to A$ with
$f^*c_1(L)=m^2c_1(L_{\rm source})$. For an effective source detector
$Y$ of dimension three,
\[
 \deg_L(f_*Y)=\int_Y f^*c_1(L)^3
 =m^6\deg_{L_{\rm source}}Y.
\]
On the rational Weil plane $f_*$ is an isomorphism and respects the
$K$-character decomposition. The reduced image is still detecting
and has no larger degree. Source dominance and detector closedness
give the bound throughout the target component. Existence of $F$
for signature $(3,3)$ and determinant $[-1]$ follows from rational
Hermitian classification, Schoen \cite[\S7, p.~332]{SF-Schoen1998};
the family isogeny is explained in \S6 on that page. Enumerating
rational $K$-matrices and testing $F^*H_{\rm target}F=H_{\rm source}$
terminates by this classification, after which $m$ is an exact
integer.
\end{proof}
The isometry and denominator are fixed for one polarized family.
This does not assert control by a varying isogeny degree or the core
index $J$ on a Hecke orbit. Rescaling a rank-six Hermitian form by a
rational $c$ multiplies its determinant by $c^6=\operatorname{Nm}(c^3)$;
it cannot change its determinant class.

\begin{proposition}[SF-05: both maximal-order Gaussian principal types]
\label{SF-05}
For an integral maximal $\mathbf Z[i]$ action compatible with a
principal polarization, a balanced rank-$2n$ Gaussian Hermitian
lattice has exactly two possible isometry types:
\[
 U^n,\quad U=\begin{pmatrix}0&1\\1&0\end{pmatrix},
 \qquad I_{n,n}=\operatorname{diag}(1^n,-1^n).
\]
The cyclic-four principal Prym has the first type. A fixed isogeny
from the first type to the second pulls back the principal
polarization to twice the source principal polarization. Thus the
sixfold bound valid on either maximal-order principal type is
\[
 D_{G,6}^{\rm principal}=8D_{G,6}
 =975\,950\,559\,600\,810\,393\,600.
\]
\end{proposition}
\begin{proof}
The Gaussian module is free. If $E$ is the principal alternating
form, its associated integral Hermitian matrix satisfies
$\det_{\mathbf Z}E=(\det_K H)^2$. Hence $H$ is unimodular and its
signature gives determinant $(-1)^n$. Rational Hermitian
classification makes its rational space split. Choose a rational
isotropic vector and divide an integral multiple by its Gaussian
content to obtain a primitive isotropic lattice vector $e$.
Unimodularity gives $f$ with $H(e,f)=1$. The span has matrix
\[
 \begin{pmatrix}0&1\\1&k\end{pmatrix},\quad k\in\mathbf Z,
\]
and determinant $-1$, hence splits integrally. Replacing $f$ by
$f-ce$ changes $k$ by $2\operatorname{Re}(c)$. An even $k$ reduces
to $U$; an odd $k$ reduces to the block with $k=1$, whose basis
$f,e-f$ has matrix $I_{1,1}$. The complement is balanced,
unimodular and rationally split, so induction gives a sum of these
two sorts of planes. If an odd block occurs, convert every remaining
$U$ using
\[
 \langle1\rangle\oplus U\simeq\langle1,1,-1\rangle:
 \quad a+e,\ a-f,\ a+e-f
\]
is an orthogonal unimodular basis with norms $1,1,-1$ when
$(a,e,f)$ has matrix $\operatorname{diag}(1,U)$. Retaining the
negative line of the original odd block gives $I_{2,2}$, and the
process continues. Parity of vector norms distinguishes the even
$U^n$ from the odd $I_{n,n}$.

For the source, the topological model of SF-02 gives on every
unaffected handle the saturated primitive integral basis
\[
 a_0-a_2,\quad a_1-a_3,\quad b_0-b_2,\quad b_1-b_3.
\]
Its Jacobian intersection form is twice symplectic and its half is
$\Xi$. The deck action takes the first difference to the second and
the second to minus the first. In the Gaussian basis
$a_0-a_2,b_0-b_2$ the Hermitian form is
$\bigl(\begin{smallmatrix}0&i\\-i&0\end{smallmatrix}\bigr)$,
which becomes $U$ by a unit change. The covered torus contributes
only invariant homology, so these blocks exhaust the primitive
lattice. The published lattice calculation is also in van Geemen
\cite[Theorem~5.3(1)--(2), p.~629]{SF-vanGeemen1996}.

Finally, with
$M=\bigl(\begin{smallmatrix}1&1\\1&-1\end{smallmatrix}\bigr)$,
\[
 M^*I_{1,1}M=2U.
\]
Direct sums give an integral Gaussian map, invertible over $K$,
which transports the complex structure and identifies the two
period domains. Its induced isogeny has degree $2^{2n}$ and
pullback polarization $2\Xi$. The degree of a dimension-$n$
effective cycle grows by $2^n$ by the projection formula. Detection
is preserved by the $K$-isogeny and reduced images lower degree.
Putting $n=3$ gives the stated factor eight.
\end{proof}
For a rational action or a nonmaximal integral order, the two-lattice
conclusion is not asserted; SF-04 retains the fixed denominator cost.

\subsection{Small integral projectors and quadratic division constants}

\begin{theorem}[SF-06: an explicit Gaussian sixfold division constant]
\label{SF-06}
Let a polarized Gaussian Weil sixfold have integral $\mathbf Z[i]$
action and an effective threefold detector of degree at most $D$.
For $0\ne t\in W$, with $\delta t$ integral and
$T=-\int t^2>0$, one has
\[
 c_t\le\delta D^2\max\{31\,719\,424,
              563\,200\,000\lceil\sqrt T\rceil\}.
\]
Here degrees use the stated integral compatible polarization;
principality is not needed for this division step.
\end{theorem}
\begin{proof}
On $H^6$ put $I=[i]^*$, $A=[2+i]^*$ and
$\bar A=[2-i]^*$. In the character decomposition indexed by
$j=0,\ldots,6$, the eigenvalue of $I$ is $(-1)^{j+1}$.
On the even indices, the eigenvalues of $A+\bar A$ are $-234$
for $j=0,6$ and $150$ for $j=2,4$. Therefore the integral
correspondence
\[
 (1-I)(150-A-\bar A)=768e_W.
\]
For a detector $Y$ set $y=768e_W[Y]$. The expansion
\[
 y=(150+IA+I\bar A)[Y]-(150I+A+\bar A)[Y]
\]
has integral effective positive and negative parts, each of degree
at most $(150+125+125)D=400D$, because the pullbacks by $2\pm i$
multiply threefold degree by $5^3=125$. The authorized orthogonal
projector and comass estimates of Theorem~W give
\[
 d=-\int y^2\in\mathbf Z_{>0},\qquad
 \sqrt d\le768D/3!=128D,\qquad d\le16384D^2.
\]
The two eigenvalues of $A$ on $W$ are
$(2\pm i)^6=-117\pm44i$. For $y_1=y$, $y_2=Ay$, the positive
middle pairing thus has Gram matrix
\[
 d\begin{pmatrix}1&-117\\-117&15625\end{pmatrix},
 \qquad 15625-117^2=1936.
\]
Write $b_j=\langle\delta t,y_j\rangle\in\mathbf Z$. Direct
matrix inversion, cancelling the common scalar before assembling
cycles, gives the integral identity
\[
 1936\delta d\,t=(15625b_1+117b_2)y_1+(117b_1+b_2)y_2.
\]
The coefficient-dual vectors $15625y_1+117y_2$ and
$117y_1+y_2$ have norms $5500\sqrt d$ and $44\sqrt d$.
The signed parts of $y_2$ have degrees at most $125\cdot400D$.
Cauchy--Schwarz therefore bounds the degree of either effective
sign by
\[
 400D\,\delta\sqrt{Td}\,(5500+125\cdot44)
 \le563\,200\,000\delta\sqrt T\,D^2.
\]
The multiplier is at most
$1936\cdot16384\delta D^2=31\,719\,424\delta D^2$.
The stated ceiling produces an integer bound. No coefficient in
the cycle group is divided.
\end{proof}

\begin{theorem}[SF-07: an explicit Eisenstein sixfold division constant]
\label{SF-07}
For a polarized Eisenstein Weil sixfold with integral
$\mathbf Z[\omega]$ action, $\omega^2+\omega+1=0$, a detector of
degree at most $D$ gives
\[
 c_t\le\delta D^2\max\{70\,858\,800,
          270\,050\,760\lceil\sqrt{3T}\rceil\},
 \qquad \delta t\text{ integral},\quad T=-\int t^2>0.
\]
This applies in particular to the $(1^3,3^3)$-polarized cubic source
with $D=D_{E,6}$; it does not relabel that polarization as principal.
\end{theorem}
\begin{proof}
Put $I=[\omega]^*$ and $S=[2+\omega]^*$ on $H^6$. The sum
$1+I+I^2$ kills all character indices except $j=0,3,6$.
The $S$-eigenvalues there are $-27,27,-27$. Hence
\[
 (1+I+I^2)(27-S)=162e_W.
\]
Its positive part on $[Y]$ has degree $3\cdot27D=81D$, and
its negative part has the same bound since $\operatorname{Nm}(2+\omega)=3$.
For $y=162e_W[Y]$ the Theorem~W estimates give
\[
 d=-\int y^2>0,\qquad \sqrt d\le27D,\qquad d\le729D^2.
\]
Take $A=[3+\omega]^*$. Its norm degree factor is $7^3=343$;
its Weil eigenvalues are
\[
 (3+\omega)^6=37+360\omega=-143+180\sqrt{-3}
\]
and its conjugate. The Gram matrix of $y_1=y,y_2=Ay$ is
\[
 d\begin{pmatrix}1&-143\\-143&117649\end{pmatrix},
 \qquad117649-143^2=97200.
\]
For $b_j=\langle\delta t,y_j\rangle$ the exact integral equation is
\[
 97200\delta d\,t=(117649b_1+143b_2)y_1+(143b_1+b_2)y_2.
\]
The coefficient-dual norms are $61740\sqrt{3d}$ and
$180\sqrt{3d}$. Signed representatives of $y_1,y_2$ have each
sign of degree at most $81D$ and $343\cdot81D$, respectively.
Thus either sign in the displayed identity has degree at most
\[
 81D\,\delta\sqrt{3Td}\,(61740+343\cdot180)
 \le270\,050\,760\delta\sqrt{3T}\,D^2.
\]
The multiplier is bounded by
$97200\cdot729\delta D^2=70\,858\,800\delta D^2$.
All coefficients are integral and the claimed ceiling is safe.
\end{proof}
All seven character eigenvalues, both projector identities, both
Gram determinants, the four dual-vector norms and the displayed
integer constants were checked by exact algebra in
\texttt{sixfold\_division\_certificates.py}. This finite computation
does not purport to verify the geometric existence of the detectors
or the supplied Theorem~W. With a smaller integral order one instead
uses HD-WQ with endomorphisms lying in that order and its explicitly
charged projector.

\subsection{Interpretation and the remaining scope}
The new detector bounds give the uniform integral form of condition A,
cote~16 scan~050, with the correction SIGN\_NEXT SN-C03/C04. Their
positivity/division step is condition I, cote~16 scan~049 and cote~44
Theorem~4. Positivity is applied only after the effective supplier
exists; it supplies no cycle in an absent component. The integral
projectors implement the Cayley--Hamilton inverse mechanism of
cote~16 scans~048--049.

The endomorphisms and character projectors are algebraic
correspondences; the resulting cycle has nonzero component in the
motive whose Betti realization is $\bigwedge_K^6H^1$. The covers,
polarizations, finite-type cycle spaces and their maps spread over a
finitely generated integral $\mathbf Z$-algebra after inverting the
finite cover, polarization and level denominators. Their degree-six
classes have weight six and Tate twist $(3)$ in etale realization.
This does not replace uniform cycle degree by pointwise algebraicity
over finite fields or by an $L$-value estimate. The coherent
$K_0$-classes of the detecting subvarieties recover the cycle classes
through the topological filtration and Chern character; this
interpretation makes no semiregularity claim.

SF-01 and SF-02 settle the proposed universal cyclic dominance in the
negative. They leave the broader bounded-division question on the
missing families to other constructions, including the independently
proved eightfold contraction where applicable. Merely adding a
compensating surface changes dimension to eight; a sixfold dominance
theorem by itself gives no theorem in that higher dimension. Any
hypothetical failure is a ghost only with the stated 032-dependent
background of GHOSTS B03. No new archive, EGA, SGA or FGA locator is
invented for these interpretations; the paper's separately verified
spread and cycle-space locators apply.

\begin{proposition}[SF-04a: an explicit Gaussian lattice denominator]
\label{SF-04a}
Let $(A,\theta)$ be a Gaussian Weil sixfold with integral maximal-order
$\mathbf Z[i]$ action, split rational Hermitian form, and
\[
 \Delta=\chi(\theta)=\int_A\theta^6/6!.
\]
Then
\[
 \delta_W(A,\theta)\le(2\Delta)^6D_{G,6}.
\]
More generally, if only the order $\mathbf Z+f\mathbf Z[i]$ acts
integrally, the safe explicit bound is
\[
 \delta_W(A,\theta)\le(2f^{12}\Delta)^6D_{G,6}.
\]
The sharper principal bound of SF-05 is retained when its hypotheses
apply.
\end{proposition}
\begin{proof}
Let $\Lambda$ be the maximal-order Hermitian lattice. Choose a
rational maximal isotropic $K$-subspace and intersect it with
$\Lambda$. The intersection is a saturated Gaussian submodule of
rank three. Since $\mathbf Z[i]$ is a principal ideal domain, it is
a direct summand and its basis extends to an integral basis of
$\Lambda$. In this basis the Hermitian matrix is
\[
 H=\begin{pmatrix}0&B\\B^*&C\end{pmatrix},\qquad
 B\in M_3(\mathbf Z[i]),\quad C=C^*.
\]
Here $B$ is invertible over $K$. Put $b=\det B$. The determinant
of the underlying integral alternating polarization is $\Delta^2$,
so the Hermitian determinant is $-\Delta$; hence
$\operatorname{Nm}(b)=\Delta$. Define
\[
 C'=(B^*)^{-1}CB^{-1}
   =\frac{(\operatorname{adj}B)^*C\operatorname{adj}B}{\Delta},
 \qquad
 T=\begin{pmatrix}1&-C'/2\\0&B^{-1}\end{pmatrix}.
\]
Direct multiplication gives
\[
 T^*HT=\begin{pmatrix}0&1\\1&0\end{pmatrix},\qquad
 T^{-1}=\begin{pmatrix}1&(B^*)^{-1}C/2\\0&B\end{pmatrix}.
\]
Since $b\bar b=\Delta$, all entries of both $T$ and $T^{-1}$
belong to $(2\Delta)^{-1}\mathbf Z[i]$. In particular
$(2\Delta)T$ maps the even unimodular Gaussian source lattice
integrally into the target lattice and pulls its polarization back
to $(2\Delta)^2\Xi$. Transport the target complex structure through
$T$ and use the corresponding source detector from SF-03. Its
threefold degree under pushforward is multiplied by at most
$(2\Delta)^6$, proving the first assertion. The size of the basis
coefficients does not affect this denominator estimate.

For the order of conductor $f$, enlarge the lattice to
$\Lambda'=\mathbf Z[i]\Lambda$. Then
$\Lambda\subset\Lambda'\subset f^{-1}\Lambda$ and
$k=[\Lambda':\Lambda]\le f^{12}$. On
$A'=V_{\mathbf R}/\Lambda'$ the form $f^2\theta$ is an integral
ample polarization $\theta'$ compatible with the maximal order,
and
\[
 \Delta'=\chi(\theta')=f^{12}\Delta/k\le f^{12}\Delta.
\]
The map $[f]:A'\to A$ induced by multiplication by $f$ is an
isogeny with $[f]^*\theta=\theta'$. Consequently pushing a
threefold detector forward has no degree cost relative to these
specified polarizations. Apply the first assertion to $\Delta'$
and use its upper bound. Scaling the Hermitian form by $f^2$ does
not change its determinant class, so splitness is preserved.
\end{proof}

\section{The split Eisenstein eightfold family: an exact \(G_{34}\) construction}
\label{sec:g34-new}

Put \(K=\mathbb Q(\omega)\), where \(\omega^2+\omega+1=0\), and
\(\mathcal O=\mathbb Z[\omega]\).  We give the branch datum, reconstruct its
parabolic monodromy, prove that the period map has rank exactly \(16\), and
supply the algebraic determinant classes needed for the transfer.  The final
bound is integral and effective.  Initially the target lattice is
\(\mathcal O\)-stable; the last paragraph treats an arbitrary fixed order.

\paragraph{Attribution and status of the inputs.}
The ten roots, eight cocycle representatives, and seventeen intervals below
are numerical data from Appendix A of Benjamin Severn's unrefereed manuscript
\cite{G34:Severn}, Hexagon \(2610.00387v2\), internal version \(1.5\), dated
21 October 2026.  The actual PDF and TeX title is \emph{Weil classes on abelian
sixfolds and eightfolds of Weil type for \(\mathbb Q(\sqrt{-3})\)}.  The portal
heading instead mentions only sixfolds.  The downloaded TeX has SHA256
\texttt{fc15af16599b399d8f111f6f0804916f}\allowbreak\texttt{dd39dc627e83adb5ff4e1788babc75ec},
agreeing with the portal's advertised prefix.  The source log preserves the
metadata discrepancy and the hashes of the actual files read.

Severn's Theorem A claims the split eightfold result; his Theorems B and C
concern sixfolds and powers.  Those theorems remain conditional if invoked as
inputs.  None is invoked here.  In particular, neither his dominance criterion,
his supplier-transfer theorem, nor his printed monodromy or Gram matrices is
a premise.  The numerical witness is checked by a separate exact reconstruction;
the arguments below use the published inputs identified at each step.
Markman and 032 are not used in this section.  The check events distinguish
exact computations, whole-proof checks, and published background.

\subsection*{D-G34-01. Group identification, branch datum, and the eightfold}

Equip \(K^6\) with \(\langle x,y\rangle=\overline{x}^{\,t}y\).  The roots
\(v_i\) are the rows of the following array, regarded as column vectors when
forming matrices:
\[
\begin{array}{c|rrrrrr}
i&\multicolumn{6}{c}{v_i^t}\\\hline
1&1&-1-\omega&1&\omega&-1-\omega&\omega\\
2&1&0&0&1+\omega&0&0\\
3&0&1&1+\omega&0&0&0\\
4&1&\omega&-1-\omega&\omega&-1-\omega&1\\
5&1&-1-\omega&-1-\omega&-1-\omega&-1-\omega&\omega\\
6&1&-1-\omega&\omega&\omega&\omega&\omega\\
7&0&0&1&0&0&-1\\
8&1&\omega&\omega&-1-\omega&1&-1-\omega\\
9&1&-\omega&0&0&0&0\\
10&0&1&0&0&0&-1
\end{array}
\]
Define
\[
 s_i=1-\frac{2v_i\overline v_i^{\,t}}{\langle v_i,v_i\rangle},\qquad
 (t_1,\ldots,t_{20})=(s_1,s_1,\ldots,s_{10},s_{10}).
\]
Each root has norm \(2\) or \(6\), and each \(s_i\) is an involutive
reflection.  Consequently \(t_1\cdots t_{20}=1\).

\textbf{Claim.} The \(s_i\) generate \(G_{34}\), of order \(39\,191\,040\).
Their natural representation is absolutely irreducible over \(K\), has
character field \(K\), and its rational simple block is \(M_6(K)\).

\begin{proof}
Consider the following set \(R\) of norm-two vectors:
\[
\begin{split}
 R_{\rm sp}&=\{\omega^ae_i-\omega^be_j:
          i\ne j,\ a,b\in\{0,1,2\}\},\\
 R_{\rm de}&=\left\{\frac{\pm1}{1+2\omega}
       (\omega^{a_1},\ldots,\omega^{a_6}):
       a_i\in\{0,1,2\},\ \sum a_i\equiv0\pmod3\right\}.
\end{split}
\]
They have cardinalities \(270\) and \(486\), respectively.  The exact
reconstruction checks that every \(s_i\) permutes \(R=R_{\rm sp}\cup R_{\rm de}\),
that their generated action is transitive on its \(756\) elements, and that
the permutation group has order \(39\,191\,040\).  Its stabilizer-chain
orbit lengths are
\[
                   756,\quad80,\quad12,\quad9,\quad6,
\]
whose product is the stated order.  The permutation action is faithful because
\(R\) spans \(K^6\).

Davis \cite[Theorem 5.1 and its proof, pp.~312--313]{G34:Davis} describes
\(G_{34}\) by the reflections in the hyperplanes
\(x_i-x_j=0\), \(x_1-\omega x_2=0\), and
\(x_1+\cdots+x_6=0\).  He gives the invariant degrees
\(6,12,18,24,30,42\), whose product is \(39\,191\,040\), and the same
\(270+486\) root set.  One root reflection in our group, together with
transitivity on \(R\), gives the reflections in every root direction.  In
particular it gives Davis's standard generators.  Thus the published
\(G_{34}\) is a subgroup of the computed group; equality of the exact orders
proves equality of the groups.  This is an identification by containment and
order, not just by a remembered order and dimension.

The roots have rank six and their nonorthogonality graph is connected, both
checked over \(K\).  If a complex invariant subspace \(W\) is given, then
\(W^\perp\) is invariant.  Since \(1-s_i\) has rank one, its root line
lies in \(W\) or \(W^\perp\).  Nonorthogonal roots cannot belong to opposite
summands.  Connectedness puts all roots in the same summand, and their span
is all of \(\mathbb C^6\).  This proves absolute irreducibility.

The exact permutation calculation also verifies that \(\omega I\) belongs
to the group.  The character values lie in \(K\), while
\(\operatorname{tr}(\omega I)=6\omega\) generates \(K\) over \(\mathbb Q\).
The character field is therefore exactly \(K\).  The displayed realization
over that field has Schur index one, giving the block \(M_6(K)\).
\end{proof}

Riemann existence now supplies a connected \(G_{34}\)-cover
\(C\to\mathbb P^1\) with this product-one generating tuple.  The corresponding
Hurwitz component is a finite cover of the ordered configuration space
\(M_{0,20}\), hence has dimension \(17\).  Its genus is
\[
 2g(C)-2=|G_{34}|\left(-2+20\left(1-\frac12\right)\right)
        =8|G_{34}|,
 \qquad g(C)=156\,764\,161.
\]
For the natural character \(\sigma\) and its conjugate, every inertia
eigenvalue list is \((1,1,1,1,1,-1)\).  Chevalley--Weil in the precise form
of \cite[Theorem 2.10]{G34:FGP} gives holomorphic multiplicity
\[
                         -6+20\cdot\frac12=4
\]
for each embedding.  A rational primitive idempotent of \(M_6(K)\) cuts out
an abelian variety \(A\) of dimension eight and Weil signature \((4,4)\).
The parabolic multiplicity space has \(K\)-dimension eight.  The following
calculation verifies the latter dimension independently at the cocycle level.

\subsection*{D-G34-02. The exact parabolic and boundary certificate}

Use left cocycles \(u(gh)=u(g)+g u(h)\).  Write
\[
 u_i=u(\gamma_i)=\lambda_i v_{\lceil i/2\rceil},\qquad
 Q_i=t_1\cdots t_{i-1}.
\]
The parabolic relation and the coboundaries are
\[
 \sum_{i=1}^{20}Q_i\lambda_i v_{\lceil i/2\rceil}=0,
 \qquad u_i=(1-t_i)a.
\]
In \(\lambda\)-coordinates the coboundary matrix has row
\(2\overline v_{\lceil i/2\rceil}^{\,t}/
\langle v_{\lceil i/2\rceil},v_{\lceil i/2\rceil}\rangle\).
The relation matrix and coboundary matrix have ranks six and six, and their
product is zero.  Thus the quotient dimension is \(20-6-6=8\).

For \(j=1,\ldots,4\), let \(b_j\) have
\(\lambda_{2j-1}=\lambda_{2j}=1\) and all other entries zero.  The last four
vectors are supported at the ordered indices
\[
                       (1,3,5,7,10,12,14,16,18,20)
\]
and have coefficient rows
\[
\begin{aligned}
 b_5:;&(1,0,0,0,1+\omega,1,0,-1-\omega,0,-1+\omega),\\
 b_6:;&(0,1,0,0,1+\omega,\tfrac23-\tfrac23\omega,-1,
                     -\tfrac23-\tfrac13\omega,0,1+2\omega),\\
 b_7:;&(0,0,1,0,-\tfrac13-\tfrac23\omega,-\tfrac23-\tfrac13\omega,
                          1+\omega,0,1+\omega,-\omega),\\
 b_8:;&(0,0,0,1,1+\omega,1,-2-\omega,-1-\omega,0,3\omega).
\end{aligned}
\]
All eight satisfy the relation.  Adjoining the six coboundary columns gives
rank fourteen, so these are a quotient basis, rather than just candidate
representatives.  Put \(I=\langle b_1,b_2,b_3,b_4\rangle_K\).

For an auxiliary pairing certificate, define
\[
 F_i(u)=\sum_{j<i}Q_j u_j,\qquad
 S(u,v)=\sum_i\left\langle Q_i^{-1}F_i(u)+\frac12u_i,v_i\right\rangle.
\]
The expression is the customary left-cocycle cup expression, comparable with
\cite[Theorem 2.5]{G34:DWdual}.  The proof below does not need a convention
conversion identifying it with the geometric cup product.  It is used only
as an explicitly verified auxiliary matrix.  Exact evaluation gives
\(S^*=-S\), \(\operatorname{rank}S=8\), zero upper-left block, and
upper-right block \(\operatorname{diag}(3,1,1,3)\).  All entries are saved
in the computation certificate.

The intervals and return powers are as follows:
\[
\begin{array}{c|rrr|r@{\qquad}c|rrr|r}
 j&a&b&k&\operatorname{rank}N_j&j&a&b&k&\operatorname{rank}N_j\\\hline
1&1&8&1&4&10&9&12&1&2\\
2&1&18&1&1&11&9&16&1&4\\
3&2&8&2&3&12&9&17&2&4\\
4&3&8&1&3&13&9&18&1&4\\
5&4&8&2&2&14&11&12&1&1\\
6&5&6&1&1&15&13&14&1&1\\
7&5&8&1&2&16&13&16&1&2\\
8&7&8&1&1&17&15&16&1&1\\
9&9&10&1&1&&&&&
\end{array}
\]
They are pairwise nested or disjoint and each gives a stable partition.
They therefore determine seventeen disjoint nonisotopic essential circles
on the twenty-punctured sphere, a maximal degeneration in
\(\overline M_{0,20}\).  For an interval \([a,b]\), put
\(c=t_a\cdots t_b\).  The listed \(k\) satisfies \(c^k=1\).

The \(k\)-th boundary twist fixes the tuple exactly.  Its pullback is computed
from
\[
\begin{split}
 u(c^k)&=(1+c+\cdots+c^{k-1})
             \sum_{i=a}^b(t_a\cdots t_{i-1})u_i,\\
 u'_i&=u_i+(1-t_i)u(c^k)\quad(a\le i\le b),\qquad
 u'_i=u_i\quad(i\notin[a,b]).
\end{split}
\]
Indeed the twist conjugates each inside loop by the boundary word
\(\gamma_a\cdots\gamma_b\); the cocycle identity gives the formula, and
\(c^k=1\) gives the returning action.  Reducing the resulting cocycles
modulo the verified coboundary space gives matrices
\[
 T_j=1+N_j,\qquad N_j=
 \begin{pmatrix}0&X_j\\0&0\end{pmatrix},\qquad
 N_iN_j=0,\quad T_j^*ST_j=S.
\]
The ranks are in the table.  These matrices have been derived from the roots,
not copied from the printed monodromies.

Let \(D=\operatorname{diag}(-3,-1,-1,-3)\).  The matrices \(-DX_j\) are
Hermitian positive semidefinite: all fifteen nonempty principal minors of
each are nonnegative rational numbers, checked exactly.  The leading principal
minors of their sum are
\[
                         96,\quad2379,\quad58722,\quad2344041.
\]
They are positive.  Every positive weighted sum of the \(N_j\) consequently
has rank four, with kernel and image both \(I\).

Flatten the sixteen matrices with indices
\[
                    1,2,3,4,5,6,7,9,10,11,12,13,14,15,16,17
\]
row by row into vectors of length sixteen over \(K\).  Their determinant is
\[
                                  1728\ne0.
\]
Thus the \(K\)-span of all \(X_j\) is \(M_4(K)\), and their complex span
under either embedding is \(M_4(\mathbb C)\).  The supplementary rational
flattening has rank sixteen and a nonzero minor \(64\); it is the
\(K\)-minor that proves complex tangent rank, so no rational-dimension
substitution is being made.

\paragraph{Reproducible finite certificate.}
The executable \texttt{g34\_reconstruct.py}, invoked with
\texttt{--group-order}, and the complete output
\texttt{g34\_reconstruction\_result.json} accompany this section in its
agent directory.  The latter contains every entry of \(S\) and all seventeen
\(X_j\), the selected minor, every principal minor, the ranks, and the group
orbit lengths.  Entries are encoded as two rational coefficients of
\(1,\omega\).  The input consists only of the ten roots, the eight
representatives, and the interval list, all printed above.  Fraction arithmetic,
Gaussian elimination, determinants, the parabolic quotient, and twist pullback
are implemented directly over \(\mathbb Q[\omega]/(\omega^2+\omega+1)\).
Sympy is used only for the exact finite permutation-group order.  The
additional \texttt{g34\_matrices.tex} prints the output matrices for a reader
who prefers a typeset certificate.  The saved checks are exact computations,
not proof-assistant formalizations.

\subsection*{D-G34-03. The period differential has rank exactly sixteen}

\textbf{Statement.} This Hurwitz family has smooth points where its period
map has rank sixteen.  Its image contains an analytic open subset of the
split signature-\((4,4)\) unitary period domain, and its algebraic moduli map
is dominant onto its connected Weil component.

\begin{proof}
Pass to a finite marking cover so that the \(G\)-cover and its multiplicity
factor are families.  Near the maximal degeneration above, the stable pointed
sphere has seventeen independent plumbing parameters.  On the punctured
polydisc the Hurwitz cover has finite degree.  Since its fundamental group
is \(\mathbb Z^{17}\), a finite Kummer base change kills its finite
permutation monodromy.  Choose the exponents to be multiples of the listed
return powers.  A further common multiple kills any finite central marking
action.  A base point outside the interval circles and a marking of its lifted
fibre give the cocycle pullback calculated above; equivalently this is the
geometric-monodromy identification in
\cite[Lemma 2.2 and Theorem 2.5]{G34:DWmon}.

Every positive plumbing loop is the positive twist about its vanishing
circle.  The passage from a twist to cohomological pullback may reverse all
signs, but does not independently reverse individual signs.  Therefore the
logarithms of the resulting unipotent monodromies are
\[
                        M_j=\varepsilon c_jN_j,
             \qquad c_j\in\mathbb Z_{>0},\quad\varepsilon\in\{1,-1\},
\]
with one common \(\varepsilon\).  Powers only multiply these logarithms,
since \(N_j^2=0\).

Write the punctured coordinates as \(z_j=e^{2\pi i w_j}\).
Schmid's multivariable nilpotent-orbit theorem
\cite[Theorem (4.12), p.~233]{G34:Schmid} gives
\[
              F(w)=\exp\!\left(\sum_j w_jM_j\right)\Psi(z),
\]
where \(\Psi\) extends holomorphically to the polydisc with values in the
compact dual.  Let \(F_0=\Psi(0)\).  Restrict to the positive diagonal
\(z_1=\cdots=z_{17}=z\).  Its logarithm is \(M=\sum_jM_j\).
The semidefinite certificate, which is purely a statement about the computed
matrices and does not require identifying the auxiliary form with the
polarization, gives
\[
              \operatorname{rank}_K M=4,\qquad
                         \operatorname{im}M=\ker M=I.
\]
On the underlying rational \(H^1\), the rank is eight.  The monodromy weight
filtration centered at one is
\[
                 W_0=W_1=I_{\mathbb Q},\qquad W_2=V_{\mathbb Q}.
\]
The diagonal restriction has the same limiting filtration \(F_0\).
Schmid's limiting mixed-Hodge-structure theorem
\cite[Theorem (6.16), p.~255]{G34:Schmid} applies.  Since
\(F_0^0=V_{\mathbb C}\) and \(F_0^2=0\), the weight-zero graded piece is
of type \((0,0)\).  More explicitly, both the filtration and its conjugate
have index-zero term equal to the whole space, so Hodge indices on a pure
graded piece are nonnegative; at weight zero they must both be zero.
Consequently
\[
                     F_0^1\cap I_{\mathbb C}=0.
\]
On either \(K\)-embedding block \(V_\sigma\), dimensions are four and four,
and hence
\[
                         V_\sigma=I_\sigma\oplus F_{0,\sigma}^1.
\]

The tangent space of the compact dual \(\operatorname{Gr}(4,V_\sigma)\)
at \(F_{0,\sigma}^1\) is
\(\operatorname{Hom}(F_{0,\sigma}^1,V_\sigma/F_{0,\sigma}^1)\).
Under the displayed decomposition, the infinitesimal action of \(M_j\)
is its induced map from \(F_{0,\sigma}^1\) to \(I_\sigma\).
Changing from the fixed quotient complement in the finite calculation to
\(F_{0,\sigma}^1\) is a fixed invertible map.  Thus these tangent maps
span the same dimension as the \(c_jX_j\), namely all sixteen dimensions
by the determinant \(1728\).

This limiting rank persists at an interior point.  Differentiate the period
map with respect to \(w_j\), and translate its tangent vectors back by
\(\exp(-\sum w_iM_i)\).  The result is
\[
                   M_j\Psi(z)+2\pi i z_j\frac{\partial\Psi}{\partial z_j}.
\]
The second term tends to zero, and the first to \(M_jF_0\).  In a holomorphic
coordinate chart near \(F_0\), take the determinant of the sixteen columns
selected in the finite certificate.  Its limit is nonzero.  It is therefore
nonzero for some sufficiently small point with all coordinates nonzero.
Translation by the group preserves rank.  The smooth period map consequently
has rank at least sixteen; the unitary period domain has dimension
\(4\cdot4=16\), so the rank is exactly sixteen.

It remains to verify that the component is split using the actual geometric
polarization.  The logarithm \(M_1\) is skew-adjoint for that polarization,
its square is zero, and its image is \(I\).  For any \(u,v\), the pairing
of \(M_1u,M_1v\) is minus the pairing of \(u,M_1^2v\), hence zero.  Thus
\(I\) is an actual \(K\)-rational totally isotropic four-space in the
rank-eight Hermitian multiplicity space.  The form is split and its
determinant class is \(1\).  This reasoning is independent of the auxiliary
cup-form convention.  The period image lies in the required component;
rank sixteen makes it analytically open locally, and its algebraic moduli
map is dominant.  The source-dimension inequality was never used as a
sufficient dominance test.
\end{proof}

\subsection*{D-G34-04. The exceptional-cohomology step in Schoen's construction}

\textbf{Statement.} In the paired-branch blowup construction of
\cite[\S2]{G34:Schoen}, the inverse image of each exceptional divisor has
zero cohomology in every faithful cyclic character.  This holds for every
cyclic order, in particular for inertia exponents \(m/2\).

\begin{proof}
Work over \(\mathbb C\).  Suppose a normalized cyclic \(m\)-cover has
normal-crossing branch divisors, and two divisors of nonzero exponents
\(a,-a\pmod m\) are blown up along their transverse intersection \(Z\).
The exceptional divisor \(p:E\to Z\) is a \(\mathbb P^1\)-bundle.
Its total exponent is zero modulo \(m\), so the normalized pullback is
generically unramified along it.  The two strict branch transforms meet
\(E\) in disjoint sections.  Every other branch divisor restricts to the
full inverse image of its intersection with \(Z\).  These statements can
also be checked in the two blowup charts of the local equation
\[
             y^m=(\text{unit})s^a t^{m-a}\prod_\ell u_\ell^{a_\ell}:
\]
remove the \(m\)-th power of the exceptional parameter.

Let \(\beta:F\to E\) be the finite inverse image and \(\chi\) a faithful
character of \(\mu_m\).  Put
\(\mathscr F_\chi=(\beta_*\mathbb C)_\chi\).
At a point with nontrivial inertia, its stalk is zero.  Indeed the fibre
representation is a permutation representation, and its \(\chi\)-part
is nonzero exactly when \(\chi\) is trivial on the stabilizer.  Away from
the branch locus it is a rank-one local system.

On a \(\mathbb P^1\)-fibre not contained in another branch divisor, its
restriction is therefore extension by zero of a rank-one local system on
\(\mathbb P^1\setminus\{0,\infty\}\) with monodromy
\(e^{2\pi i a/m}\ne1\).  On a fibre contained in another branch divisor,
the restriction is zero.  A rank-one local system on \(\mathbb C^*\) with
monodromy \(\lambda\ne1\) has zero cohomology: its cellular complex has
differential \(\lambda-1\), an isomorphism.  Duality gives zero compactly
supported cohomology as well.  Proper base change therefore gives
\(Rp_*\mathscr F_\chi=0\).  Since \(\beta\) is finite and proper,
\(R\beta_*\mathbb C=\beta_*\mathbb C\), and hence
\[
                            H^*(F,\mathbb C)_\chi=0.
\]
One may apply the exact standard forms SGA4 XII, Theorem 5.1(iii), and
XVI, Theorem 4.1, by first using finite coefficients of residue characteristic
prime to \(m\), adjoining the \(m\)-th roots of unity, and then passing
to characteristic-zero coefficients and classical comparison
\cite{G34:SGA4}.  The cellular differential remains invertible with those
coefficients.

The inverse image of the original centre \(Z\) has the same faithful-character
vanishing: its inertia contains the nontrivial subgroup generated by \(a\),
so the corresponding pushforward eigensheaf is zero at every point.
The blowup is an isomorphism on the common complement.  Comparing the two
compact-support localization sequences, whose closed terms both have zero
faithful-character cohomology, identifies the faithful-character parts
before and after the blowup.  In particular no new primitive character is
introduced.  At subsequent stages the same argument applies, since the
branch transforms retain normal crossings.  This proves the assertion
needed in Schoen's Lemma 2.8.
\end{proof}

The repair is necessary for the proof being used: on p.~18 of
\cite{G34:Schoen}, a prime \(p\mid m\) is chosen such that the degree-\(p\)
quotient remains ramified at a selected point.  For \(m=12\) and local
exponent \(6\), all prime-degree quotients are unramified there.  The direct
faithful-character argument above avoids that choice.  The remaining published
paired-branch construction is used through Theorem 2.0, Lemmas 2.2--2.6,
and Corollary 3.1 of \cite{G34:Schoen}; their hypotheses are verified next.

\subsection*{D-G34-05. Algebraic cyclic suppliers for every needed character}

\textbf{Statement.} Let \(C\to\mathbb P^1\) be a connected finite
\(G\)-cover with all inertia groups of order two.  For a subgroup
\(H'\le G\) and a linear character \(\chi:H'\to\mathbb C^*\) of order
\(m\ge3\), put \(C_\chi=C/\ker\chi\), \(X=C/H'\), and let
\(P_\chi\) be the primitive cyclotomic factor of \(J(C_\chi)\), with
\(F=\mathbb Q(\zeta_m)\).  If its \(F\)-rank \(h\) in \(H^1\) is
positive, then \(h\) is even, each primitive embedding has Hodge numbers
\((h/2,h/2)\), and \(\bigwedge_F^hH^1(P_\chi,\mathbb Q)\) is algebraic.

\begin{proof}
The morphism \(C_\chi\to X\) is a connected cyclic \(m\)-cover.  Its
inertia groups are images of intersections with conjugates of the original
order-two inertia groups.  Thus every nontrivial image is \(-1\in\mu_m\).
For odd \(m\) the cover is unramified.  For even \(m\), every nonzero local
exponent is \(m/2\).  The relation in the abelianized punctured surface
group says that the sum of the exponents is zero modulo \(m\), so their
number \(s\) is even.  They therefore pair with their negatives.

At every primitive embedding, the nonzero eigenangles are all \(1/2\).
Writing \(q=g(X)\), cyclic-cover Riemann--Roch gives
\[
                     h=2q-2+s,\qquad
       h^{1,0}_\chi=h^{0,1}_\chi=q-1+\frac{s}{2}=\frac h2.
\]
Thus the cover is simple in Schoen's exact definition on p.~11 of
\cite{G34:Schoen}.  His Theorem 2.0 constructs the determinant classes on
\(C_\chi^h\).  D-G34-04 supplies its exceptional-cohomology step for the
present exponents.  The Jacobian correspondence and the cyclotomic idempotent
then give the determinant subspace on \(P_\chi\), as in his Corollary 3.1,
pp.~24--25.  A zero-dimensional factor is omitted.  All primitive embeddings
have been covered, not only the distinguished character \(\chi\).
\end{proof}

\subsection*{D-G34-06. Brauer induction and the algebraic contraction}

\textbf{Statement.} Both Weil lines of the \(G_{34}\) eightfold are algebraic.
More generally the transfer below applies to a non-self-dual multiplicity
determinant whenever the required factor is realized by a rational
group-algebra idempotent as in D-G34-01.  We make no claim here that every
abstract irreducible character, regardless of Schur index, has such a
one-copy rational factor.

\begin{proof}
Let \(J=J(C)\), \(H=H^1(J,\mathbb Q)\), and
\(U=\operatorname{Cent}_{\operatorname{Sp}(H)}(G)\), the full algebraic
centralizer.  Decompose \(H_{\mathbb C}\) into irreducible \(G\)-types.
A non-self-dual pair \(\tau,\overline\tau\) has multiplicity spaces
\(M_\tau,M_\tau^*\), giving a general linear factor of \(U_{\mathbb C}\).
A self-dual orthogonal \(G\)-type gives a symplectic multiplicity factor;
a self-dual quaternionic type \(\eta\) gives an orthogonal multiplicity
factor.  Only types occurring in \(H\) need be retained.  Consequently
\[
 X^*(U_{\mathbb C})=
       \mathbb Z^{\{\text{non-self-dual pairs}\}}
       \oplus(\mathbb Z/2)^{\{\text{quaternionic types}\}}.
\]
The disconnected orthogonal factors must be included.

For a complex \(G\)-representation \(V\), define \(a(V)\) by the free
coordinates
\(\operatorname{mult}_\tau(V)-\operatorname{mult}_{\overline\tau}(V)\)
and torsion coordinates \(\operatorname{mult}_\eta(V)\bmod2\).
This is additive on the representation ring.  Frobenius reciprocity identifies
\(\operatorname{Hom}_{H'}(\chi,H_{\mathbb C})\) with the
\(\chi\)-eigenspace of \(H^1(C_\chi,\mathbb C)\).  The character of its
determinant under \(U\) is precisely \(a(\operatorname{Ind}_{H'}^G\chi)\):
the general linear factors contribute the indicated determinant powers,
the symplectic factors have no character, and the orthogonal factors
contribute the indicated parities.

Brauer's integral induction theorem
\cite[\S10.5, Theorems 19--20]{G34:Serre} expresses the natural irreducible
character as
\[
                          \sigma=\sum_i n_i
                                  \operatorname{Ind}_{H_i}^G\chi_i,
                             \qquad n_i\in\mathbb Z.
\]
Applying \(a\) gives
\[
                           e_\sigma=\sum_i n_i
                             a(\operatorname{Ind}_{H_i}^G\chi_i),
\]
where \(e_\sigma\) is the unit determinant character for \(M_\sigma\).
Characters of orders one and two contribute zero.  Their free coordinates
are zero because they are real.  On a quaternionic irreducible, their
eigenspaces inherit a nondegenerate alternating form: distinct eigencharacters
are orthogonal unless their product is one, and a real eigencharacter is
its own inverse.  Their dimensions are therefore even, giving zero parity
coordinates.  A term whose eigenspace in \(H_{\mathbb C}\) is zero also
contributes zero.  Every remaining term is supplied by D-G34-05.
Replacing \(\chi_i\) by \(\chi_i^{-1}\) negates its full character:
the free coordinates change sign, and the parity coordinates are their own
negatives.  This deals with every negative coefficient.

We spell out the contraction.  For each retained supplier choose a nonzero
rational algebraic class \(y_i\) in its cyclotomic determinant space.
Its components at all primitive embeddings are nonzero.  Indeed these
components are the embeddings of a rank-one \(F_i\)-space, and a nonzero
rational element has no zero embedding.  Choose the embedding determinant
line \(L_i\subset H^{h_i}(P_i,\mathbb C)\) with character
\(\operatorname{sign}(n_i)a(\operatorname{Ind}\chi_i)\).
The complementary top exterior line
\[
                     L_i^\vee\subset
                       H^{2\dim P_i-h_i}(P_i,\mathbb C)
\]
has the inverse character and pairs perfectly with \(L_i\).  This follows
also by taking the product of all other embedding determinant lines: the
full determinant of \(H^1(P_i)\) is trivial under \(U\), because \(U\)
preserves the polarization on the factor.

On \(B=\prod_iP_i^{|n_i|}\), take a pure tensor \(z\) in these complementary
lines, and let \(x\ne0\) lie in \(\det M_\sigma\subset H^8(A,\mathbb C)\).
The character identity makes \(\Gamma=z\otimes x\) on \(B\times A\)
\(U\)-invariant.  Each supplier is an algebraic isogeny factor of \(J(C)\):
use quotient-curve pullback and its rational cyclotomic projector.
The factor \(A\) is similarly given by a rational group-algebra idempotent.
These correspondences commute with \(U\).  A rational \(U\)-equivariant
splitting on \(H^1\), followed by exterior powers, lifts \(\Gamma\) to a
\(U\)-invariant class on a power of \(J\).

The Lefschetz group \(S(J)\) is contained in \(U\), since
\(\operatorname{End}^0(J)\) contains \(\mathbb Q[G]\).
Milne's theorem \cite[Theorem 3.2; Lemma 3.1]{G34:Milne} says that the
\(S(J)\)-invariants on all powers of \(J\) are generated by divisor
classes, also after extension to complex coefficients.  The lifted class,
and hence \(\Gamma\), is therefore algebraic with complex coefficients.
Contract it with the rational algebraic class
\(\boxtimes_i y_i^{\boxtimes|n_i|}\) and push forward along \(B\).
Only the chosen primitive component pairs with each complementary line.
The result is a nonzero scalar multiple of \(x\), so \(x\) is algebraic
with complex coefficients.  Repeat for its conjugate line.  Both the rational
Weil space and the rational cycle-class image are defined over \(\mathbb Q\);
inclusion after extension to \(\mathbb C\) implies inclusion over
\(\mathbb Q\).  This proves algebraicity of the whole Weil plane.

For the quantitative construction there is a useful middle-degree version.
On the curve power \(C_i^{h_i}\), the determinant line is contained in
\(H^1(C_i)^{\otimes h_i}\).  The determinant line of the conjugate
eigenspace has inverse character and is its dual under integration.  Both
have degree \(h_i\).  Thus, with
\(\ell=\sum_i|n_i|h_i\), the contracting tensor on
\((\prod_i C_i^{h_i|n_i|})\times A\) has degree \(\ell+8\).
It lifts to a \(U\)-invariant \(H^1\)-tensor on a power of \(J\) by the
same correspondences.  This is the precise interface to the supplied degree
estimate; no inverse Lefschetz operation with uncontrolled denominators
is needed.
\end{proof}

\subsection*{D-G34-07. An explicit effective bound and the order correction}

The full proofs of \cite[SQ1--SQ7]{G34:Hecke} have been read and checked in
this run by two readers.  Their geometric premises are now supplied by
D-G34-01--06, with no Severn theorem premise: a split rank-eight factor,
rank-sixteen dominance, algebraic cyclic determinant suppliers with every
primitive component nonzero, and the integral relation in the full
centralizer character group.  The arithmetic and intersection estimates of
SQ1--SQ7 are used as supplied results, with the following explicit definitions.

Set
\[
\begin{gathered}
 G=39\,191\,040,\qquad g=4G+1,\qquad
 R=(2G+1)G!G^G,\qquad L=2gR,\\
 S_0=G^{2g}(g!)^2,\qquad S_B=(L/2)!S_0^R,\qquad k=L/2+4,\\
 B_0=(40G)^k\binom{k}{4}S_B\,8!G^8.
\end{gathered}
\]
If the target lattice is \(\mathcal O\)-stable and its polarization has
elementary divisors \(d_1,\ldots,d_8\), write
\(\Delta=\prod_i d_i\) and define
\[
             N=4G^8\Delta,\qquad
             \boxed{D_8(\Delta)=N^8B_0.}
\]
These are positive integers.  The precise estimates used are: coefficient
mass at most \(R\) in the integral character relation (SQ3); supplier ranks
\(h_i\le g\) and integral orbit-sum/Ramanujan-projector mass at most
\(S_0\) (SQ4); exterior-product mass at most \(S_B\) (SQ5); globally
generated divisor representatives bounded numerically by \(40G\) times the
product polarization and the corresponding mixed-degree estimate \(B_0\)
(SQ6); and a split-lattice comparison cleared by \(4G^8\Delta\), using
\(\Delta_A\mid G^8\) for the source (SQ2, SQ7).

Consequently a dominant analytic open subset of the fixed-data Weil family
has an effective integral four-cycle of polarization degree at most
\(D_8(\Delta)\) with nonzero Weil projection.  The accepted bounded-cycle
closedness and transport results of the project extend the same bound over
the connected fixed-data family.  The accepted Theorem W then gives, in
its normalization of denominator \(\delta\) and norm bound \(T\),
\[
       \delta_W\le D_8(\Delta),\qquad
       c_t\le \left\lceil C_{4,3+\omega}(\delta,T)\right\rceil
                       D_8(\Delta)^4.
\]
If the chosen definition of \(C\) is already integral, the ceiling is
redundant.  No bound here is obtained merely from algebraicity point by point;
it is the uniform supplier/intersection construction and the proved
rank-sixteen source that give the degree bound.

\paragraph{D-G34-07b: a nonmaximal endomorphism order.}
Suppose the target lattice \(\Lambda\) is stable under
\(\mathcal O_f=\mathbb Z+f\mathcal O\).  Define
\(\Lambda'=f\mathcal O\Lambda\).  Then
\[
                  f\Lambda\subset\Lambda'\subset\Lambda,
      \qquad I=[\Lambda:\Lambda']\le f^{16},
\]
and \(\Lambda'\) is \(\mathcal O\)-stable.  The identity on the rational
vector space induces an isogeny
\(\phi:A'=V_{\mathbb R}/\Lambda'\to A\).
The pulled-back polarization \(\theta'=\phi^*\theta\) is integral,
positive, and compatible with \(\mathcal O\).  Its elementary-divisor
product is
\(\Delta'=I\Delta\), because its polarization matrix is the restriction
of the original alternating form to an index-\(I\) sublattice.  Thus
\(\Delta'\le f^{16}\Delta\).
Apply the maximal-order bound to \(A'\) and push its effective integral
four-cycle forward by \(\phi\).  The projection formula gives
\[
                 \deg_\theta(\phi_*Z)=\deg_{\theta'}Z.
\]
An isogeny induces an isomorphism on rational cohomology and carries the
Weil subspace to the Weil subspace, so the nonzero Weil projection remains
nonzero.  Since \(D_8\) is increasing in its positive integer argument,
\[
 \boxed{\delta_W(A,\theta)\le D_8(f^{16}\Delta),\qquad
 c_t\le\left\lceil C_{4,3+\omega}(\delta,T)\right\rceil
                         D_8(f^{16}\Delta)^4.}
\]
This proves the fixed-order extension without suppressing a conductor factor.
A rational \(K\)-action on a fixed integral local system has some such fixed
order; the bound is uniform with that order and polarization data fixed.

\section{A split \(\mathbb Q(\sqrt{-7})\) sixfold source from \(G_{24}\)}
\label{sec:g24-new}

Let \(K=\mathbb Q(\tau)\), where
\(\tau=-(1+\sqrt{-7})/2\), so \(\tau^2+\tau+2=0\) and
\(\mathcal O_K=\mathbb Z[\tau]\).  We construct a dominant family of split
Weil sixfolds, supply its determinant classes, and give an explicit effective
cycle bound for every fixed split integral datum over this field.

\paragraph{Provenance.}
The initial three reflection matrices are the specialization \(p=2\) of
Deraux's published matrices \cite[equation (13), author manuscript p.~24;
\S3, pp.~4--6]{G24:Deraux}.  The six roots and nine boundary intervals below
were found in this run by an exact search, beginning with those three matrices.
They are not copied from Severn or rStation.  The unrefereed rStation papers
58 and 75 suggested looking at \(G_{24}\), but no statement or numerical
witness from them is a premise.  The fixed certificate
\texttt{g24/g24\_reconstruct.py} reconstructs every matrix from the displayed
roots.  No numerical approximation, claimed period-rank routine, or dimension
inequality is used to prove dominance.

\subsection*{D-G24-01. The finite group, integral action, and branch data}

Put
\[
 H=\begin{pmatrix}
 2&-\tau&\bar\tau\\
 -\bar\tau&2&-\tau\\
 \tau&-\bar\tau&2
 \end{pmatrix},\qquad
 \langle x,y\rangle=\bar x^{\,t}Hy.
\]
Its leading principal minors are \(2,2,1\), so it is positive definite.
For a nonzero root \(v\), define
\[
 s_v=1-\frac{2v\bar v^{\,t}H}{\langle v,v\rangle}.
\]
The coordinate roots give Deraux's three involutions
\[
 R_1=\begin{pmatrix}-1&\tau&-\bar\tau\\0&1&0\\0&0&1\end{pmatrix},\quad
 R_2=\begin{pmatrix}1&0&0\\\bar\tau&-1&\tau\\0&0&1\end{pmatrix},\quad
 R_3=\begin{pmatrix}1&0&0\\0&1&0\\-\tau&\bar\tau&-1\end{pmatrix}.
\]
Their group is the rank-three reflection group \(G_{24}\), of order
\(336\).  In addition to the published identification, a direct enumeration
over \(K\) gives exactly \(336\) distinct matrices, all preserving \(H\).
Their orbit on the first coordinate vector has \(42\) vectors of norm two,
or \(21\) root lines.  These are explicit finite equality checks, not a
recognition based only on the order.

Take the following six norm-two roots:
\[
\begin{array}{c|rrr}
 i&\multicolumn{3}{c}{v_i^t}\\\hline
 1&0&-1&1+\tau\\
 2&-1&0&\tau\\
 3&-2&\tau&1+\tau\\
 4&-1-\tau&-2&\tau\\
 5&-1&0&0\\
 6&-1&-1&\tau
\end{array}
\]
Set \(s_i=s_{v_i}\) and
\[
                (t_1,\ldots,t_{12})=(s_1,s_1,\ldots,s_6,s_6).
\]
The exact enumeration verifies that the \(s_i\) generate the same \(336\)
matrices as \(R_1,R_2,R_3\).  Their root span has rank three and their
nonorthogonality graph is connected.  The orthogonal-summand argument used
in D-G34-01 therefore proves absolute irreducibility over \(\mathbb C\).
All matrices lie in \(K\), and
\[
                  \operatorname{tr}(R_1R_2R_3)=1+\tau.
\]
Thus the character field is exactly \(K\).  The displayed realization has
Schur index one, so the rational simple block is \(M_3(K)\).

The doubled tuple has product one and generates the group.  Riemann
existence gives a connected \(G_{24}\)-cover \(C\to\mathbb P^1\), with
twelve branch points of inertia order two.  Its Hurwitz component is a
finite cover of \(M_{0,12}\).  Riemann--Hurwitz gives
\[
 2g(C)-2=336\bigl(-2+12(1-1/2)\bigr)=4\cdot336,
 \qquad g(C)=673.
\]
For the natural character \(\sigma\) and its conjugate, each local
eigenvalue list is \((1,1,-1)\).  Chevalley--Weil
\cite[Theorem 2.10]{G34:FGP} gives holomorphic multiplicity
\[
                         -3+12/2=3
\]
at each embedding.  A rational primitive idempotent of \(M_3(K)\) therefore
cuts out an abelian sixfold of signature \((3,3)\).

We need its integral structure, not just its rational isogeny class.  Let
\(G=336\), let \(\chi\) be the natural character, and put
\[
 \varepsilon=\frac3G\sum_{a\in G_{24}}
       \operatorname{Tr}_{K/\mathbb Q}\bigl(\chi(a^{-1})\bigr)a,
 \qquad e=\varepsilon\frac{1-s_1}{2}.
\]
The character traces are integers.  Thus \(2Ge\in\mathbb Z[G_{24}]\).
The exact coefficient calculation also verifies that \(Ge\) itself has
denominator two; it would be incorrect to reuse the even-dimensional
\(G_{34}\) projector denominator.  The element \(e\) is a primitive
Rosati-symmetric rational idempotent.  Let \(A\subset J(C)\) be its image,
with the restriction \(\theta_A\) of the principal polarization.

There is an integral \(\mathcal O_K\)-action on this particular subvariety.
Let \(g_0=R_1R_2R_3\).  Direct multiplication gives
\[
 g_0=\begin{pmatrix}
 1+\tau&1&1\\1&1&-\tau\\-\tau&-1-\tau&-1
 \end{pmatrix},\qquad
 \langle v_1,g_0v_1\rangle=\tau,
 \qquad (1-s_1)g_0v_1=\tau v_1.
\]
The integral endomorphism \((1-s_1)g_0\) of \(J(C)\) preserves \(A\).
Indeed \(g_0\) preserves the central \(\varepsilon\)-block and \(1-s_1\)
maps that block to \(A\).  On the natural root line its restriction is
\(\tau\), and on the conjugate line it is \(\bar\tau\).  Its restriction
to \(A\) thus supplies \(\tau\in\operatorname{End}(A)\), proving
\(\mathcal O_K\)-stability of \(H_1(A,\mathbb Z)\) without saturation.

Write \(\Lambda=H_1(J(C),\mathbb Z)\) and
\(\Lambda_A=\Lambda\cap e\Lambda_{\mathbb Q}\).  The alternating
dual is \(\Lambda_A^\#=e\Lambda\): the primitive inclusion of
\(\Lambda_A\) in the unimodular lattice \(\Lambda\) lets every integral
functional extend, and self-adjointness of \(e\) identifies its representing
vector with its projection.  Hence
\[
 2G\Lambda_A^\#\subseteq\Lambda_A.
\]
Every polarization elementary divisor divides \(2G=672\).  Therefore
\begin{equation}\label{eq:g24-source-delta}
 \Delta_A\mid(2G)^6,
 \qquad \int_A\theta_A^6=6!\Delta_A\le6!(2G)^6.
\end{equation}

\subsection*{D-G24-02. Nine exact boundary logarithms}

Use left cocycles \(u(gh)=u(g)+g u(h)\), and write
\[
 u_i=\lambda_i v_{\lceil i/2\rceil},\qquad Q_i=t_1\cdots t_{i-1}.
\]
The relation is \(\sum_i Q_i u_i=0\); the coboundaries are
\(u_i=(1-t_i)a\).  The coefficient matrices have ranks three and three,
their product is zero, and the quotient dimension is \(12-3-3=6\).

Let \(b_1,b_2,b_3\) be the equal-pair cocycles supported respectively at
\((1,2),(3,4),(5,6)\), with coefficient one at each member of the pair.
The remaining representatives are supported at \((1,3,5,7,9,11)\), with
rows
\[
\begin{aligned}
 b_4:\;&(-1-\tau/2,\ 1,\ -1-\tau/2,\ 1,\ 0,\ 0),\\
 b_5:\;&(1-\tau/2,\ 1+\tau,\ -1-\tau/2,\ 0,\ 1,\ 0),\\
 b_6:\;&(-1-\tau,\ 1,\ -1,\ 0,\ 0,\ 1).
\end{aligned}
\]
All six satisfy the relation.  Adjoining the three coboundary columns gives
rank nine, so they are a quotient basis.  Set
\(I=\langle b_1,b_2,b_3\rangle_K\).

As an auxiliary algebraic certificate, set
\[
 S(u,v)=\sum_i\left\langle Q_i^{-1}\sum_{j<i}Q_j u_j+u_i/2,v_i\right\rangle.
\]
The reconstruction checks \(S^*=-S\), \(\operatorname{rank}S=6\),
and a zero upper-left \(3\)-by-\(3\) block.  Its upper-right block is
\[
 B=\begin{pmatrix}
 -1-\tau/2&1-\tau/2&-1-\tau\\
 1&1+\tau&1\\
 -1-\tau/2&-1-\tau/2&-1
 \end{pmatrix}.
\]
The proof uses this form only for its explicitly verified matrix properties;
no unproved conversion between cup-product conventions is required.

Take the nine intervals
\[
 [1,2],\ [3,4],\ [5,6],\ [7,8],\ [9,10],\ [11,12],
 \ [1,4],\ [5,8],\ [1,8].
\]
They are pairwise nested or disjoint and define distinct stable partitions.
Thus they are a maximal compatible collection of nine boundary circles on
the twelve-punctured sphere.  Their products \(c=t_a\cdots t_b\) are all
the identity.  In particular every boundary twist already returns the
monodromy tuple, without a nontrivial return power.

The twist formula is
\[
 u(c)=\sum_{i=a}^b(t_a\cdots t_{i-1})u_i,\qquad
 u_i'=u_i+(1-t_i)u(c)\quad(a\le i\le b),
\]
with unchanged outside coordinates.  It follows by applying the cocycle
identity to conjugation by the interval boundary word.  Reducing these
transformed representatives modulo coboundaries gives
\[
 T_j=1+N_j,\qquad N_j=\begin{pmatrix}0&X_j\\0&0\end{pmatrix},
 \qquad N_iN_j=0,\qquad T_j^*ST_j=S.
\]
The ranks of the \(N_j\), in the displayed interval order, are
\[
                         1,1,1,1,1,1,2,2,2.
\]
All nine \(3\)-by-\(3\) blocks and the full matrix \(S\) are printed
in the associated exact matrix appendix.

For each \(j\), let \(Q_j'\) be minus the lower-right block of \(SN_j\).
It is Hermitian positive semidefinite.  Here is the complete principal-minor
certificate, with columns ordered as
\(1,2,3,12,13,23,123\):
\[
\begin{array}{c|rrrrrrr}
 j&1&2&3&12&13&23&123\\\hline
1&2&4&4&0&0&0&0\\
2&2&4&2&0&0&0&0\\
3&2&2&2&0&0&0&0\\
4&2&0&0&0&0&0&0\\
5&0&2&0&0&0&0&0\\
6&0&0&2&0&0&0&0\\
7&3/2&5/2&2&2&1&1&0\\
8&3/2&2&2&2&2&0&0\\
9&0&2&2&0&0&2&0
\end{array}
\]
The leading principal minors of \(\sum_jQ_j'\) are
\[
                              11,\quad143,\quad732.
\]
Thus every positive weighted sum of the \(Q_j'\) is positive definite.
Since the lower-left block of \(S\) is invertible, every positive weighted
sum of the \(N_j\) has rank three and has image and kernel both \(I\).

Flatten the nine \(X_j\) row by row.  The resulting \(9\)-by-\(9\)
matrix over \(K\) has determinant
\begin{equation}\label{eq:g24-period-minor}
             -308-126\tau\ne0,
 \qquad N_{K/\mathbb Q}(-308-126\tau)=87808=2^8\,7^3.
\end{equation}
This proves complex linear independence under either embedding of \(K\),
not merely rational independence of a doubled list of coefficients.

\subsection*{D-G24-03. The analytic rank bridge and splitness}

We give the analytic step explicitly.  The nine circles determine a
zero-dimensional boundary stratum of \(\overline M_{0,12}\), with nine
independent plumbing coordinates \(z_1,\ldots,z_9\).  After a finite cover
of the Hurwitz space marking the finite monodromy, and a coordinatewise
Kummer base change, the associated polarized weight-one variation has
unipotent monodromy logarithms
\[
 M_j=\epsilon c_jN_j,\qquad c_j\in\mathbb Z_{>0},
 \qquad\epsilon\in\{1,-1\}\text{ independent of }j.
\]
The common sign reflects the global choice of pullback versus transport
and orientation; it cannot be chosen independently at different nodes.
The geometric parabolic action is the cocycle action by
\cite[Lemma 2.2 and Theorem 2.5]{G34:DWmon}.  Finite residual central
ambiguities are removed by the finite base change.  Since \(N_j^2=0\),
the logarithm is exactly the indicated scalar multiple of \(N_j\).

Put \(z_j=e^{2\pi i w_j}\).  The nilpotent orbit theorem
\cite[Theorem 4.12, p.~233]{G34:Schmid} gives
\[
 F(w)=\exp\!\left(\sum_jw_jM_j\right)\Psi(z),
 \qquad\Psi\text{ holomorphic at }0,
 \qquad F_0=\Psi(0).
\]
Along the diagonal, \(M=\sum_jM_j\) has rank three over \(K\), and
\(M^2=0\), \(\operatorname{im}M=\ker M=I\).  On the underlying
twelve-dimensional rational space the monodromy weight filtration is
\[
                    W_0=W_1=I_{\mathbb Q},\qquad W_2=V_{\mathbb Q}.
\]
Schmid's limiting mixed Hodge structure theorem
\cite[Theorem 6.16, p.~255]{G34:Schmid} applies.  The original weight-one
filtration has \(F^0=V_{\mathbb C}\), \(F^2=0\); these conditions remain
true at \(F_0\).  On \(\operatorname{Gr}_0^W\), both Hodge indices are
therefore nonnegative and sum to zero.  It has type \((0,0)\), so
\[
                            F_0^1\cap I_{\mathbb C}=0.
\]
In either embedding of \(K\), \(F_0^1\) is a three-dimensional complement
to \(I\) in a six-dimensional space.

The tangent space of \(\operatorname{Gr}(3,6)\) there is
\(\operatorname{Hom}(F_0^1,I)\).  After projection along the complement,
the tangent vectors \(M_jF_0\) are represented by \(\epsilon c_jX_j\),
up to the same invertible basis changes.  Equation
\eqref{eq:g24-period-minor} makes all nine independent.

Translate a period derivative back by
\(\exp(-\sum_jw_jM_j)\).  In any local Grassmannian chart it is
\[
              M_j\Psi(z)+2\pi i z_j\frac{\partial\Psi}{\partial z_j}(z).
\]
As every \(z_j\) tends to zero through nonzero values, its tangent value
tends to \(M_jF_0\).  The determinant of these nine translated tangent
vectors thus has a nonzero limit.  It is nonzero at sufficiently nearby
smooth points.  The actual period map has rank nine there; the unitary
period domain of signature \((3,3)\) has dimension nine.  The period
image consequently contains a nonempty analytic open subset.

Finally \(M\), as a geometric logarithm, is skew-adjoint for the genuine
alternating polarization.  Thus
\[
 E(Mu,Mv)=-E(u,M^2v)=0.
\]
Its image \(I\) is a three-dimensional isotropic \(K\)-subspace.  Hence
the Hermitian form is split.  This conclusion uses the actual polarized
variation; it does not identify the auxiliary matrix \(S\) with the
geometric cup product by fiat.

\subsection*{D-G24-04. Algebraic determinant classes and transfer}

The supplier and transfer proofs D-G34-04--06 apply to this cover.  We
specify all their hypotheses and the dimensional changes to avoid hiding
an additional theorem premise.

For a subgroup \(H'\le G_{24}\) and a linear character \(\chi\) of order
\(m\ge3\), the intermediate map
\(C/\ker\chi\to C/H'\) is cyclic.  Every nonzero inertia image is the
unique element of order two.  If \(m\) is odd the map is unramified; if
\(m\) is even, all nonzero branch exponents equal \(m/2\).  The abelianized
surface-group relation makes their number even, so they pair with their
negatives.  For base genus \(q\) and \(s\) nonzero branch points, every
primitive eigenspace has even rank \(h=2q-2+s\) and signature
\((h/2,h/2)\).  Schoen's cyclic-product theorem
\cite[Theorem 2.0 and Corollary 3.1]{G34:Schoen}, with the faithful-character
localization proof of D-G34-04 replacing the defective prime-quotient step,
supplies rational algebraic determinant classes.  That repair is independent
of \(G_{34}\), the field of the selected target character, and the target
dimension.  It uses only paired nonzero branch exponents.  Its original
center and exceptional-divisor vanishings both apply here.

The rational primitive cyclotomic determinant space has rank one over its
cyclotomic field.  A nonzero rational element has a nonzero component at
every primitive embedding, so the suppliers cover all required characters.
Let \(U=\operatorname{Cent}_{\operatorname{Sp}(H^1(C))}(G_{24})\), including
its disconnected orthogonal factors.  Brauer induction gives an integral
relation for \(\sigma\), hence for its determinant character in
\[
                       X^*(U_{\mathbb C})=\mathbb Z^a\oplus(\mathbb Z/2)^b.
\]
The orders-one-and-two terms vanish: free coordinates cancel under complex
conjugation, and quaternionic restriction multiplicities are even by the
nondegenerate alternating form on a real-character eigenspace.  Zero
cohomology terms also vanish.  The remaining terms are exactly the supplied
cyclic determinants.  Replacing a character by its inverse handles negative
coefficients in both the free and parity coordinates.

Choose one primitive determinant line for each term and its complementary
exterior line.  The integral character relation makes their product with
the desired line \(\det M_\sigma\subset H^6(A,\mathbb C)\) \(U\)-invariant.
All factors are algebraic isogeny factors of \(J(C)\) via quotient pullbacks
and rational group-algebra projectors.  A \(U\)-equivariant splitting lifts
the tensor to a power of \(J(C)\).  The Lefschetz group of \(J(C)\) is
contained in \(U\), so \cite[Theorem 3.2]{G34:Milne} makes this tensor a
divisor-generated algebraic class.  Contracting the complementary lines
against the rational suppliers and pushing forward isolates the desired
primitive determinant line with a nonzero coefficient.  Repeat for the
conjugate line and descend from complex to rational coefficients.  The
whole rational Weil plane is algebraic.

For the quantitative argument use the middle-degree version on the curve
powers: a primitive determinant in \(H^1(C_i)^{\otimes h_i}\) pairs with
the inverse-character determinant in the same degree \(h_i\).  If
\(\ell=\sum_i|n_i|h_i\), the invariant contracting tensor on the product
of these curve powers with \(A\) has degree \(\ell+6\).  Thus its divisor
degree is \(\ell/2+3\).  This is the sixfold replacement of
\(\ell/2+4\) in the eightfold construction, and involves no inverse
Lefschetz operation of uncontrolled degree.

\subsection*{D-G24-05. An explicit detector bound}

Define the following positive integers:
\[
\begin{gathered}
 G=336,\qquad g=673,\qquad M=2G=672,\\
 R=(2G+1)G!G^G,\qquad L=2gR,\qquad
 S_0=G^{2g}(g!)^2,\qquad S_B=(L/2)!S_0^R,\\
 k=L/2+3,\qquad
 B_{0,7}=(40G)^k\binom{k}{3}S_B\,6!M^6.
\end{gathered}
\]
For a fixed split \(\mathcal O_K\)-stable Weil datum with polarization
elementary-divisor product \(\Delta\), put
\begin{equation}\label{eq:g24-detector-bound}
 N_7=4M^6\Delta,\qquad
 \boxed{D_{6,7}(\Delta)=N_7^6B_{0,7}.}
\end{equation}
Every fibre of its connected fixed-data family has an effective integral
three-cycle with nonzero Weil projection and degree at most
\(D_{6,7}(\Delta)\).

\begin{proof}
We adapt the fully checked arithmetic and intersection arguments SQ3--SQ7
of \cite{G34:Hecke}; their original \(G_{34}\) geometric premises have
already been replaced by D-G24-01--04.  The integral relation has coefficient
mass at most \(R\): the character lattice has at most \(G\) rows, all
entries have absolute value at most \(G\), and appending the columns
\(2e_j\) treats the orthogonal parities.  The integer-lattice determinant
argument of SQ3 gives \((2G+1)G!G^G\), independent of the target dimension.
The cyclic primitive projectors are cleared by \(G\), since their
coefficients are Ramanujan sums divided by subgroup orders.  Supplier ranks
are at most \(g\), and the integral orbit-sum cycle mass is bounded by
\(S_0\) as in SQ4.  The product supplier on the curve product has degree
mass at most \(S_B\), with \(\ell\le gR\le L\), as in SQ5.

The full-centralizer invariant algebra is generated by degree-two tensors
coming from \(\mathbb Q[G_{24}]\), with disconnected orthogonal factors
included.  The effective divisor spanning and moving argument of SQ6
therefore applies.  Its globally generated divisors are bounded in the nef
order by \(40G(H_B+\theta_A)\).  The bound \(40G\) depends on the finite
cover degree and on the curve-to-Jacobian pullbacks, not on the representation
dimension.  Intersect the signed product supplier times \(A\) with a
divisor monomial of degree \(c=\ell/2+3\), then push the resulting
three-cycles to \(A\).  Their total degree mass is at most
\[
\begin{aligned}
 &(40G)^c\int_{|Y_B|\times A}(H_B+\theta_A)^c\theta_A^3\\
 &\qquad=(40G)^c\binom c3
       \operatorname{mass}_{H_B}(Y_B)\int_A\theta_A^6
 \le B_{0,7}.
\end{aligned}
\]
The invariant contraction is in the complex span of these divisor monomials
and has nonzero Weil projection.  Therefore at least one of the finite
collection of effective integral output cycles has nonzero Weil projection.
In fact their projected classes span both Weil eigenlines.  No bound on
the scalar coefficients of a complex divisor-polynomial presentation is
needed for this extraction.

It remains to compare integral split lattices.  For the alternating form
\(E\), define
\[
                  h(x,y)=E(x,\tau y)-\bar\tau E(x,y).
\]
This is an integral \(\mathcal O_K\)-valued Hermitian form, and the
coefficient of \(\tau\) recovers \(E\).  In an \(\mathcal O_K\)-basis,
\(\det E=(\det h)^2\); in signature \((3,3)\), \(\det h=-\Delta\).
The ring \(\mathcal O_K\) is a PID.  For completeness, it is norm
Euclidean: choose the \(\tau\)-coefficient of an approximation so that the
imaginary error has absolute value at most \(\sqrt7/4\), then choose the
integer coefficient so that the real error has absolute value at most
\(1/2\).  The error norm is at most \(7/16+1/4=11/16<1\).

A primitive maximal isotropic submodule therefore extends to an integral
basis with matrix
\[
                 \begin{pmatrix}0&B\\B^*&D\end{pmatrix},
 \qquad b=\det B,\qquad N_{K/\mathbb Q}(b)=\Delta.
\]
Writing \(X=(B^*)^{-1}D/2\), the substitutions \(w'=w-vX\) and
\(u=w'B^{-1}\) make a standard hyperbolic basis.  The forward and inverse
coordinate matrices have denominators dividing \(2\Delta\): the only
denominators are \(2\bar b,b,2b\bar b\).  Applying this to source and
target gives a comparison denominator dividing \(4\Delta_A\Delta\),
which is cleared by \(N_7=4M^6\Delta\) in view of
\eqref{eq:g24-source-delta}.  Multiplying the rational isometry by \(N_7\)
gives an isogeny with
\[
                  \phi^*\theta_{\mathrm{target}}=N_7^2\theta_A.
\]
Its pushforward multiplies the degree of a three-cycle by \(N_7^6\),
and preserves nonzero Weil projection.  D-G24-03 gives an analytic open
period image, so \eqref{eq:g24-detector-bound} holds on an analytic open
subset of the fixed target family.  The project's accepted bounded-cycle
closedness and transport theorem extends that same degree bound to every
fibre of its connected component.
\end{proof}

For a fixed nonmaximal order \(\mathcal O_f=\mathbb Z+f\mathcal O_K\),
put \(\Lambda'=f\mathcal O_K\Lambda\).  Then
\[
 f\Lambda\subset\Lambda'\subset\Lambda,\qquad
 I=[\Lambda:\Lambda']\le f^{12},\qquad \Delta'=I\Delta.
\]
The identity isogeny from the torus with lattice \(\Lambda'\) to the
original torus has pullback polarization of volume \(\Delta'\).  Pushing
an effective three-cycle forward preserves its degree for this pair of
polarizations.  Thus
\[
                     \boxed{\delta_W\le D_{6,7}(f^{12}\Delta).}
\]
This is a fixed-order assertion, with the conductor dependence displayed.
Applying the project's accepted Weil degree-to-division theorem gives the
corresponding explicit bounded-division estimate for every fixed rational
Weil class.  In particular the principal balanced \(E^6\) datum over
\(K=\mathbb Q(\sqrt{-7})\) is covered.  The argument establishes the
split sixfold component; it does not assert an all-determinant sixfold
theorem over this field.

\section{A split \(\mathbb Q(\sqrt{-2})\) sixfold source from an \(F_4\) extension}
\label{sec:gf4-new}

Put \(K=\mathbb Q(t)\), \(t^2=-2\), and \(\mathcal O_K=\mathbb Z[t]\).
We construct a finite subgroup of \(U(4)\) over \(K\), a twelve-involution
cover, and an exact rank-nine period certificate.  Two local involutions
have a two-dimensional negative eigenspace; this is handled explicitly
in the parabolic calculation.  The resulting split sixfold source gives
an explicit effective detector bound.

\paragraph{Provenance.}
The standard \(F_4\) root system and four simple roots below are the
published ones in \cite[\S3, p.~83]{GF4:Cahn}.  The extra matrix, five
chosen reflection roots, and exact boundary certificate are constructed
and checked here.  The unrefereed rStation paper75 suggested an \(F_4\)
extension over \(\mathbb Q(\sqrt{-2})\); none of its theorem statements,
matrices, or numerical witnesses is used.  The proof establishes the
finite group directly rather than relying on an unverified classification
of possible sources.

\subsection*{D-GF4-01. The group and its integral sixfold factor}

Use the standard positive definite form \(\bar x^{\,t}y\) on \(K^4\).
For a nonzero vector \(v\), write
\[
                  s_v=1-2v\bar v^{\,t}/(\bar v^{\,t}v).
\]
Let \(R\) be the usual \(48\)-root system
\[
 R=\{\pm e_i\}\ \cup\ \{\pm e_i\pm e_j:i<j\}
       \ \cup\ \{(\pm e_1\pm e_2\pm e_3\pm e_4)/2\},
\]
and \(W=\langle s_v:v\in R\rangle=W(F_4)\).  It is generated by the
four simple reflections in
\[
 e_2-e_3,\quad e_3-e_4,\quad e_4,\quad(e_1-e_2-e_3-e_4)/2.
\]

Here is a direct order proof.  The signed permutation group \(B_4\)
lies in \(W\), and a reflection in a half-integral root sends a coordinate
root to a half-integral short root.  Thus \(W\) is transitive on the
\(24\) short roots.  Its stabilizer of \(e_1\) permutes the six short
roots perpendicular to \(e_1\), namely \(\pm e_2,\pm e_3,\pm e_4\).
Restriction embeds that stabilizer in the signed permutation group
\(B_3\), of order \(48\); the entire \(B_3\) already belongs to it.
Consequently \(|W|=24\cdot48=1152\).  The fact that the displayed
reflections preserve \(R\) follows directly from their formulas, and is
also checked by the finite reconstruction.

Define
\[
 \Phi=\begin{pmatrix}
 0&-1&-1&0\\1&0&0&1\\1&0&0&-1\\0&-1&1&0
 \end{pmatrix},\qquad z=\frac t2\Phi.
\]
We have \(\Phi^t=-\Phi\), \(\Phi^2=-2I\), and therefore
\(z^*=z\), \(z^2=1\).  Its trace is zero, so the negative and positive
eigenspaces both have dimension two.  The matrix \(\Phi\) sends every
short root to a long root and every long root to twice a short root,
up to sign.  For coordinate roots this follows from its columns.  For
half-integral roots its coordinates belong to \(\{0,1,-1\}\), and its
squared norm is two, so exactly two coordinates are nonzero.  The assertion
for long roots then follows from \(\Phi^2=-2I\).  Hence \(z\) permutes
the root lines and normalizes \(W\).  As \(z\) has nonreal entries,
\[
             G_F=\langle W,z\rangle=W\sqcup Wz,\qquad |G_F|=2304.
\]
This is a group over \(K\) preserving the standard positive form.

The subgroup \(B_4\) already makes the natural representation absolutely
irreducible: the independent coordinate sign changes force an invariant
complex subspace to be a sum of coordinate lines, and the permutations
act transitively on those lines.  Put \(s_0=s_{e_1}\), let \(P_{12}\)
swap the first two coordinates, and set \(g_0=P_{12}z\).  Directly,
\[
       (1-s_0)g_0e_1=te_1,
       \qquad \operatorname{tr}(s_0g_0)=-t.
\]
Thus its character field is exactly \(K\).  The displayed realization
has Schur index one, and its rational simple block is \(M_4(K)\).

Choose five reflection roots
\[
\begin{array}{c|rrrr}
 i&\multicolumn{4}{c}{v_i^t}\\\hline
1&0&-1&-1&0\\
2&0&-1&0&1\\
3&0&0&0&-1\\
4&-1/2&1/2&-1/2&-1/2\\
5&0&0&-1&1
\end{array}
\]
and set \(s_i=s_{v_i}\).  In terms of these, the four standard simple
reflections listed above are, respectively,
\[
                 s_2s_5s_2,\quad s_5,\quad s_3,\quad
                 s_1s_2s_4s_2s_1.
\]
These short word identities are checked by direct multiplication.  Thus
the \(s_i\) generate \(W\), and the doubled tuple
\[
       (s_1,s_1,s_2,s_2,s_3,s_3,s_4,s_4,s_5,s_5,z,z)
\]
has product one and generates \(G_F\).  All twelve inertia elements
are involutions.  Ten have negative-eigenspace dimension one, and two
have negative-eigenspace dimension two.

The connected cover \(C\to\mathbb P^1\) has
\[
 2g(C)-2=2304\bigl(-2+12(1-1/2)\bigr)=4\cdot2304,
 \qquad g(C)=4609.
\]
Chevalley--Weil \cite[Theorem 2.10]{G34:FGP} gives holomorphic multiplicity
\[
                     -4+\frac{10+2\cdot2}{2}=3
\]
for each conjugate natural character.  Its multiplicity factor is a
Weil sixfold of signature \((3,3)\).  Notice that the source genus uses
the order of each inertia element, whereas Chevalley--Weil uses its
negative-eigenspace multiplicity; the two rank-two involutions are not
being counted as reflections.

Let \(G=2304\), let \(\chi\) be the natural character, and define
\[
 \varepsilon=\frac4G\sum_{a\in G_F}
       \operatorname{Tr}_{K/\mathbb Q}\bigl(\chi(a^{-1})\bigr)a,
 \qquad e=\varepsilon(1-s_0)/2.
\]
Then \(e\) is a rational primitive Rosati-symmetric idempotent and
\(Ge\in\mathbb Z[G_F]\), because \(4/2=2\) and the character traces
are integers.  Let \(A\subset J(C)\) be its image.  The integral
endomorphism \((1-s_0)P_{12}z\) preserves \(A\) and restricts to \(t\)
on its natural root line, and to \(-t\) on the conjugate line.  It therefore
supplies the integral \(\mathbb Z[t]\)-action on \(A\).  As in the
dual-lattice argument of D-G24-01, \(G\Lambda_A^\#\subseteq\Lambda_A\).
For the restricted principal polarization,
\begin{equation}\label{eq:gf4-source-delta}
        \Delta_A\mid G^6,\qquad \int_A\theta_A^6\le6!G^6.
\end{equation}

\subsection*{D-GF4-02. Parabolic coordinates with rank-two inertia}

For the first five negative eigenspaces use the column vectors
\[
\begin{gathered}
 q_1=(0,1,1,0)^t,\quad q_2=(0,1,0,-1)^t,\quad
 q_3=(0,0,0,2)^t,\\
 q_4=(1,-1,1,1)^t/2,\qquad q_5=(0,0,1,-1)^t.
\end{gathered}
For \(\operatorname{im}(1-z)\) use
\[
                   q_6=(1,-t/2,-t/2,0)^t,
                   \qquad q_7=(t/2,1,0,t/2)^t.
\]
These are obtained by taking independent columns of \(1-t_i\), and
their asserted eigenspace dimensions are verified exactly.

Write the cocycle values at the first ten branch loops as
\(u_i=\lambda_iq_{\lceil i/2\rceil}\).  At the last two put
\[
          u_{11}=\lambda_{11}q_6+\lambda_{12}q_7,
          \qquad u_{12}=\lambda_{13}q_6+\lambda_{14}q_7.
\]
Thus the parabolic coefficient space has dimension fourteen, not twelve.
With \(Q_i=t_1\cdots t_{i-1}\), impose \(\sum_iQ_iu_i=0\), and quotient
by \(u_i=(1-t_i)a\).  The relation and coboundary matrices each have
rank four, their product is zero, and the quotient has dimension
\(14-4-4=6\).

The first three quotient representatives \(b_1,b_2,b_3\) have equal
coefficients one in the first, second and third branch pair respectively.
The remaining representatives are supported in coefficient positions
\((1,3,5,7,9,11,12)\), with rows
\[
\begin{aligned}
 b_4:\;&(-1,1,1,0,1,0,0),\\
 b_5:\;&(1+t/2,-2,-1/2,-2,0,1,0),\\
 b_6:\;&(t/2,-1-t,-1/2-t/2,-t,0,0,1).
\end{aligned}
\]
They satisfy the relation, and the six columns together with the four
coboundary columns have rank ten.  Hence they form a quotient basis.
Put \(I=\langle b_1,b_2,b_3\rangle_K\).

Define the auxiliary form
\[
 S(u,v)=\sum_i\left\langle Q_i^{-1}\sum_{j<i}Q_ju_j+u_i/2,v_i\right\rangle.
\]
Here \(u_i/2\) applies to the entire negative eigenspace, including the
two-dimensional one for \(z\); no rank-one formula is substituted there.
The exact computation gives a nondegenerate anti-Hermitian form with zero
upper-left \(3\)-by-\(3\) block.  Its full matrix is in the matrix appendix.

Use the same nine compatible intervals
\[
 [1,2],\ [3,4],\ [5,6],\ [7,8],\ [9,10],\ [11,12],
 \ [1,4],\ [5,8],\ [1,8].
\]
They form a maximal boundary stratum of \(\overline M_{0,12}\), and
every interval monodromy product equals one.  Applying the cocycle identity
to the returning boundary twist gives
\[
 u_i'=u_i+(1-t_i)\sum_{j=a}^b(t_a\cdots t_{j-1})u_j
       \quad(a\le i\le b),
\]
with unchanged outside coordinates.  Coordinates in each negative
eigenspace are solved in its displayed basis, then the quotient reduction
is performed using the verified ten-column basis.  The resulting matrices
satisfy
\[
 T_j=1+N_j,\qquad N_j=\begin{pmatrix}0&X_j\\0&0\end{pmatrix},
 \qquad N_iN_j=0,\qquad T_j^*ST_j=S.
\]
Their ranks, in interval order, are
\[
                         1,1,1,1,1,2,2,2,3.
\]

Minus the lower-right block of \(SN_j\) is Hermitian positive semidefinite.
Here are all its nonempty principal minors, ordered
\(1,2,3,12,13,23,123\):
\[
\begin{array}{c|rrrrrrr}
j&1&2&3&12&13&23&123\\\hline
1&2&3&1&0&0&0&0\\
2&2&8&6&0&0&0&0\\
3&4&1&3&0&0&0&0\\
4&0&4&2&0&0&0&0\\
5&2&0&0&0&0&0&0\\
6&0&2&2&0&0&2&0\\
7&2&7&5&9/2&9/2&3/2&0\\
8&4&7&7&12&6&9/2&0\\
9&2&2&2&7/2&7/2&2&2
\end{array}
\]
The leading principal minors of their sum are
\[
                             18,\quad379,\quad1286.
\]
Every positive weighted sum of the \(N_j\) thus has rank three, image
\(I\), and kernel \(I\).  The exact determinant obtained by flattening
the nine \(X_j\) row by row is
\begin{equation}\label{eq:gf4-rank-minor}
           \frac{256+640t}{27}\ne0,
           \qquad N_{K/\mathbb Q}\left(\frac{256+640t}{27}\right)
                  =\frac{2^{15}}{3^3}.
\end{equation}
This is a determinant over \(K\), and hence a nonzero determinant under
either complex embedding.

\subsection*{D-GF4-03. Actual period rank and splitness}

The geometric braid/cocycle comparison
\cite[Lemma 2.2 and Theorem 2.5]{G34:DWmon} applies to arbitrary finite
inertia, so it applies to the two rank-two involutions as well.  Finite
marking and a coordinatewise Kummer cover of the nine-node plumbing
polydisc give geometric logarithms
\[
                  M_j=\epsilon c_jN_j,
        \qquad c_j\in\mathbb Z_{>0},\quad \epsilon\in\{1,-1\},
\]
with one common sign from the orientation/transport convention.  The
finite central ambiguity is killed by the finite cover.  All products
\(M_iM_j\) vanish.  The positive-sum calculation makes
\(M=\sum_jM_j\) have rank three over \(K\) and image and kernel \(I\).

Write the period map as
\[
       F(w)=\exp\!\left(\sum_jw_jM_j\right)\Psi(z),
       \qquad z_j=e^{2\pi i w_j},\qquad \Psi(0)=F_0,
\]
using \cite[Theorem 4.12]{G34:Schmid}.  On the underlying rational space
the weight filtration of the diagonal limit is
\(W_0=W_1=I_{\mathbb Q}\), \(W_2=V_{\mathbb Q}\).
By \cite[Theorem 6.16]{G34:Schmid} it is a limiting mixed Hodge structure.
The conditions \(F_0^0=V_{\mathbb C}\), \(F_0^2=0\) imply that
\(\operatorname{Gr}_0^W\) has type \((0,0)\).  Therefore
\(F_0^1\cap I_{\mathbb C}=0\), and in either \(K\)-embedding the
three-dimensional \(F_0^1\) is a complement to \(I\).

At that point the Grassmannian tangent space is
\(\operatorname{Hom}(F_0^1,I)\).  The tangent vectors \(M_jF_0\)
are represented, up to fixed invertible changes of basis, by the nine
nonzero scalar multiples \(\epsilon c_jX_j\).  Equation
\eqref{eq:gf4-rank-minor} makes them independent.  Translating a period
derivative by \(\exp(-\sum_jw_jM_j)\) gives
\[
                  M_j\Psi(z)+2\pi i z_j\partial_j\Psi(z).
\]
Its limit is \(M_jF_0\), so the derivative determinant remains nonzero
at sufficiently nearby smooth points.  The actual period rank is nine,
and the signature-\((3,3)\) period image contains a nonempty analytic open.
This conclusion uses the interior derivative limit, not just the nine
boundary parameters or the dimension of their span over \(\mathbb Q\).

The genuine geometric logarithm \(M\) is skew-adjoint for the polarization
and has square zero.  Thus \(E(Mu,Mv)=0\).  Its \(K\)-rank-three image
is a maximal isotropic subspace, proving that the source Hermitian form
is split.  The auxiliary form was only used to certify the rank of every
positive weighted sum, not to identify the geometric polarization.

\subsection*{D-GF4-04. Suppliers and the explicit bound}

Every inertia group of this cover has order two.  The rank of its action
in the selected four-dimensional representation does not change that
fact.  Consequently every cyclic quotient supplier has only the paired
nonzero branch exponents \(m/2\), or is unramified for odd \(m\).
The complete supplier proof D-G34-04--05 and the full-centralizer Brauer
transfer D-G34-06 therefore apply.  In particular all primitive embeddings
have balanced even rank and receive nonzero algebraic determinant classes;
the faithful-character vanishing argument repairs the original-center
and exceptional-divisor step; and the full character group includes the
orthogonal parity coordinates.  The natural representation has a rational
one-copy factor by its displayed Schur-index-one realization.  The same
algebraic contraction as D-G24-04 gives the entire Weil plane on \(A\).
No assumption that the local involutions are all reflections is needed
in this supplier or transfer argument.

Define
\[
\begin{gathered}
 G=2304,\quad g=4609,\quad
 R=(2G+1)G!G^G,\quad L=2gR,\\
 S_0=G^{2g}(g!)^2,\quad S_B=(L/2)!S_0^R,\quad
 k=L/2+3,\\
 B_{0,2}=(40G)^k\binom{k}{3}S_B\,6!G^6,
 \qquad N_2=4G^6\Delta,\\
                 \boxed{D_{6,2}(\Delta)=N_2^6B_{0,2}.}
\end{gathered}
\]
Then every fibre of a connected fixed split \(\mathcal O_K\)-stable
Weil sixfold datum of polarization volume \(\Delta\) has an effective
integral three-cycle of degree at most \(D_{6,2}(\Delta)\) with nonzero
Weil projection.

\begin{proof}
The arithmetic supplier mass, divisor spanning, and moving arguments in
D-G24-05 are independent of the imaginary quadratic field, once the
source integral action, projector denominator, and actual dominance have
been supplied.  Here the projector denominator is \(G\), and
\eqref{eq:gf4-source-delta} replaces \(\Delta_A\mid(2G)^6\).
The coefficient mass is bounded by \(R\), the individual supplier mass
by \(S_0\), and the product mass by \(S_B\), exactly as in the checked
SQ3--SQ5 of \cite{G34:Hecke}.  The degree-six contracting tensor has
divisor degree \(c=\ell/2+3\le k\).  The effective mixed-intersection
estimate is
\[
 (40G)^c\binom c3\operatorname{mass}_{H_B}(Y_B)
                         \int_A\theta_A^6\le B_{0,2}.
\]
It gives a finite spanning collection of effective integral detector
three-cycles on the source; extracting a nonzero projected member incurs
no bound on rational or complex coefficients.

For the fixed target alternating form use
\[
                         h(x,y)=E(x,ty)+tE(x,y).
\]
This integral Hermitian form recovers \(E\) as the coefficient of \(t\).
Its determinant in signature \((3,3)\) is \(-\Delta\).  The ring
\(\mathbb Z[\sqrt{-2}]\) is norm Euclidean: approximation in its
rectangular lattice gives real and imaginary errors at most \(1/2\)
and \(\sqrt2/2\), with total norm at most \(3/4<1\).  It is a PID,
so the split-lattice comparison of D-G24-05 applies verbatim.  Both
directions have denominator dividing \(2\Delta\), and the source-to-target
comparison is cleared by \(4\Delta_A\Delta\), hence by \(N_2\).
The resulting isogeny pulls target polarization back to \(N_2^2\theta_A\)
and multiplies three-cycle degree by \(N_2^6\).  D-GF4-03 supplies the
nonempty analytic open image.  The accepted bounded-cycle closedness
then proves the displayed uniform bound on the whole fixed component.
\end{proof}

For the fixed order \(\mathcal O_f=\mathbb Z+f\mathcal O_K\), the
lattice \(f\mathcal O_K\Lambda\) has index at most \(f^{12}\), and its
pullback polarization has volume at most \(f^{12}\Delta\).  The same
projection-formula argument as D-G24-05 therefore gives
\[
                       \boxed{\delta_W\le D_{6,2}(f^{12}\Delta).}
\]
The project's accepted degree-to-division theorem makes this an explicit
bounded-division estimate for every fixed rational Weil class.  The result
includes the principal balanced \(E^6\) datum over \(\mathbb Q(\sqrt{-2})\).
It proves the split component and does not assert that nonsplit sixfold
determinants over this field have been covered.

\section{Compensation gives every Eisenstein sixfold determinant}

\begin{theorem}[HD-E6: an effective bound for every integral Eisenstein sixfold datum]
\label{HD-E6}
Let $(A,L)$ be an abelian sixfold of Weil type for
$K=\mathbf Q(\omega)$, with signature $(3,3)$ and integral
$\mathbf Z[\omega]$ action preserving the polarization. Let $\Delta$
be the product of the polarization elementary divisors. Once
D-G34-01--07 are applied with their checked proofs, every member of
this fixed datum has an effective integral three-cycle detecting its
Weil space, of degree
\[
 \delta_W(A,L)\le 2D_8(\Delta^2).
\]
The same formula is uniform on its whole Weil family, without a
restriction on the determinant class. Theorem HD-WQ then gives the
stated quadratic division bound for each normalized rational Weil class.
\end{theorem}
\begin{proof}
Use the Hermitian normalization of HECKE\_DIV SQ7,
$h(x,y)=E(x,\omega y)-\bar\omega E(x,y)$.
Its matrix on an integral Eisenstein basis has determinant $-\Delta$:
its determinant has sign $(-1)^3$, and its absolute value is the
product of the elementary divisors of $E$, by the determinant identity
in the cited proof.

Let $E_\omega$ have its degree-one principal class $\vartheta$, and put
\[
 T=E_\omega\times E_\omega,\quad
 a\cdot(x,y)=(ax,\bar a y),\quad
 \lambda=\vartheta_1+\Delta\vartheta_2.
\]
This is a polarized Weil surface with Hermitian matrix
$\operatorname{diag}(1,-\Delta)$ and polarization product $\Delta$.
The product $A\times T$, with class
$H=p_A^*c_1(L)+p_T^*\lambda$, has signature $(4,4)$, polarization
product $\Delta^2$, and Hermitian determinant $\Delta^2$.
That determinant is a norm from $K$, so the rational Hermitian space
is split, by Schoen (1998), Section~7, printed p.~332.
It is an integral maximal-order datum. D-G34-07 therefore supplies an
effective integral four-cycle $Z$ on $A\times T$ with
\[
 \deg_H Z\le D_8(\Delta^2),\qquad e_{W(A\times T)}[Z]\ne0.
\]
If needed, choose a reduced component with nonzero projection.

Here is a completely effective contraction. Denote by $F_1,F_2$
the two coordinate elliptic divisors on $T$, and by $\Gamma_1$ and
$\Gamma_\omega$ the graphs of $1$ and $\omega$. Their classes span
$\operatorname{NS}(T)_\mathbf Q$: the two diagonal Hermitian matrices
and the off-diagonal entries $1,\omega$ form a rational basis of
the Hermitian $2$-by-$2$ matrices over $K$. More specifically, the
classes
\[
 [F_1]+[F_2]-[\Gamma_1],\qquad
 [F_1]+[F_2]-[\Gamma_\omega]
\]
span $W_K(T)$, since they are precisely the two off-diagonal
Hermitian classes for the balanced action. The intersection pairing
on this plane is nondegenerate. This also follows directly from its
negative-definite Hermitian-matrix intersection form, or from the
surface Hodge index theorem.

The Kunneth summand of $[Z]$ in
$H^6(A)\otimes H^2(T)$ has a nonzero projection to
$W_K(A)\otimes W_K(T)$: the nonzero product Weil component is in
this summand, and over $\mathbf C$ its two lines are the same-embedding
determinant products. Contracting this tensor against all of
$W_K(T)$ cannot give zero. Thus at least one of the four effective
divisor classes just listed has
\[
 e_{W_K(A)}p_{A*}\bigl([Z]\,p_T^*[D]\bigr)\ne0.
\]
This conclusion needs no choice of a pure Weil cycle representing $Z$.
Indeed a nonzero contraction against a difference of the displayed
effective divisors forces a nonzero contraction against at least one
of them.

Translate that elliptic divisor generally on $T$. Such translates
meet every component of $Z$ properly: apply the elementary incidence
dimension theorem to $Z\to T$ and the transitive translation action
on $T$. Their intersections have effective integral multiplicities.
The pushforward to $A$ is therefore an effective integral three-cycle;
lower-dimensional images contribute zero. It detects by the preceding
cohomological calculation, since translation does not change its class.

Every one of these four divisors satisfies the nef inequality
$[D]\le2(\vartheta_1+\vartheta_2)\le2\lambda$.
For the coordinate divisors this is immediate. For either unit graph,
its Hermitian matrix has eigenvalues $0,2$, proving the first
inequality. Hence its pushed-forward degree is at most
\[
 \int_Z p_T^*[D]p_A^*c_1(L)^3
 \le 2\int_Z p_T^*\lambda\,p_A^*c_1(L)^3
 \le2\int_Z H^4
 \le2D_8(\Delta^2).
\]
All inequalities are intersections of nef classes with an effective
cycle. Passing to a detecting reduced component only decreases degree.
This proves the claimed integral bound. One could retain the positive
coefficient four in the binomial expansion of $H^4$ to sharpen the
last estimate, but no such sharpening is needed here.

The same fixed $T$ and polarization work for the entire fixed datum
of $A$, and the eightfold bound holds on every fibre, not just on the
open Hurwitz image. Thus the conclusion is uniform, including the
possibly nonsplit sixfold determinant. This is the effective version
of the compensation argument of Schoen (1998), Proposition~10,
printed pp.~332--333, with dimensions raised by two.
\end{proof}

\begin{corollary}[HD-E6b: a fixed nonmaximal order]
\label{HD-E6b}
If the sixfold datum above has an integral action only of
$R_f=\mathbf Z+f\mathbf Z[\omega]$, the bound
\[
 \delta_W(A,L)\le2D_8(f^{24}\Delta^2)
\]
holds uniformly for this fixed conductor and polarization datum.
\end{corollary}
\begin{proof}
Write $A=V/\Lambda$ and set $\Lambda'=f\mathbf Z[\omega]\Lambda$.
Then $f\Lambda\subset\Lambda'\subset\Lambda$, because $R_f$ acts
integrally. The new lattice is maximal-order stable and has index
$I=[\Lambda:\Lambda']\le f^{12}$. The identity on $V$ induces a
$K$-linear isogeny $p:A'=V/\Lambda'\to A$. Endow $A'$ with $p^*L$;
its polarization elementary-divisor product is $I\Delta$ by the
top intersection projection formula. Apply HD-E6 and push a detecting
effective three-cycle forward. Its degree against $L$ is exactly
its degree against $p^*L$, and its Weil projection stays nonzero.
The function $D_8(x)=(4G^8x)^8B_0$ is increasing for positive $x$.
Thus $2D_8((I\Delta)^2)\le2D_8(f^{24}\Delta^2)$, as required.
No claim independent of an unbounded conductor is made.
\end{proof}

\section{Compensation over other class-number-one fields}

\begin{proposition}[HD-COMP: an effective contraction from a split family]
\label{HD-COMP}
Let $K$ be imaginary quadratic with class number one, write
$\mathcal O_K=\mathbf Z[\tau]$, and let $(A,L)$ have dimension $2n$,
signature $(n,n)$, and integral compatible $\mathcal O_K$ action.
Put $\Delta=\chi(L)>0$ and $c_K=1+\operatorname{Nm}(\tau)$.
Suppose that a bound $D_{n+1}(x)$ is proved for every split polarized
Weil $(2n+2)$-fold over $K$ with polarization product $x$.
Then, with no determinant restriction on $A$,
\[
 \delta_W(A,L)\le c_K D_{n+1}(\Delta^2).
\]
All cycles in this assertion are effective integral cycles. If the
higher-dimensional bound is numerical and explicit, so is this bound.
If it is only existential, this conclusion is only existential.
\end{proposition}
\begin{proof}
The normalized Hermitian form is
$h(x,y)=E(x,\tau y)-\bar\tau E(x,y)$, where $E$ is the alternating
polarization form. In an integral $\mathcal O_K$ basis its determinant
has absolute value $\Delta$ and sign $(-1)^n$.
Here is the normalization check for a general quadratic order basis.
Writing $\tau+\bar\tau=b$ and $\tau\bar\tau=c$, the real determinant
of a Hermitian matrix is the norm of its $K$-determinant. The map from
$h$ to $E$ is the coefficient of $\tau$ in $h$ (with the fixed choice
of conjugate-linear argument), and on a rank-one block $h=(a)$ the
integral basis $(v,\tau v)$ gives the alternating matrix
$\bigl(\begin{smallmatrix}0&a\\-a&0\end{smallmatrix}\bigr)$.
Change of $K$-basis changes both determinant formulas by the same
norm square. Thus $\det_{\mathbf Z}E=(\det_K h)^2$, without an
extra discriminant factor. The class-number-one assumption supplies
the free integral $\mathcal O_K$ basis.

Let $E_K=\mathbf C/\mathcal O_K$, with principal class $\vartheta$,
and put $T=E_K^2$ with balanced action and
$\lambda=\vartheta_1+\Delta\vartheta_2$. Its Hermitian determinant
is $-\Delta$. Consequently $(A\times T,H)$, where
$H=p_A^*c_1(L)+p_T^*\lambda$, has polarization product $\Delta^2$
and determinant $(-1)^{n+1}\Delta^2$. Rational Hermitian
classification, Schoen (1998), Section~7, printed p.~332, makes this
space split: its signed determinant is the norm of the rational
integer $\Delta$. It therefore has an effective detector $Z$ of
dimension $n+1$ and degree at most $D_{n+1}(\Delta^2)$.

On $T$, the two coordinate elliptic divisors and the graphs of
$1,\tau$ span the rational Neron--Severi group. Under the usual
Hermitian-matrix description, their off-diagonal entries span $K$
over $\mathbf Q$. Subtracting the corresponding diagonal classes
from the two graphs gives a basis of $W_K(T)$ for the balanced
action. For the graph of $\tau$ the diagonal coefficients are
$\operatorname{Nm}(\tau)$ and $1$, in an order determined by the
choice of coordinates. Its matrix has eigenvalues $0,c_K$; the
unit graph has eigenvalues $0,2$. Because $c_K\ge2$, every one of
these four effective divisors $D$ satisfies
\[
 [D]\le c_K(\vartheta_1+\vartheta_2)\le c_K\lambda
\]
in the nef order.

The nonzero Weil projection of $[Z]$ lies in the Kunneth component
$W_K(A)\otimes W_K(T)$. The intersection pairing on $W_K(T)$ is
nondegenerate. If contraction against each of the four effective
divisors were zero after projection to $W_K(A)$, contraction
against their two off-diagonal linear combinations would also be
zero, contradicting that nonzero tensor. Choose a divisor with
nonzero contraction, translate it generally to meet $Z$ properly,
and push $Z\cdot p_T^*D$ to $A$. This is an effective integral
$n$-cycle and detects $W_K(A)$. The incidence and translation
argument proving properness is exactly the one given in HD-E6;
it applies to every dimension. Its degree is at most
\[
 \int_Z p_A^*c_1(L)^n p_T^*[D]
 \le c_K\int_Z p_A^*c_1(L)^n p_T^*\lambda
 \le c_K\int_Z H^{n+1}
 \le c_K D_{n+1}(\Delta^2).
\]
Choosing a detecting reduced component only decreases this degree.
This proves the assertion without dividing a cycle or supposing
that $Z$ itself is a pure Weil cycle.
\end{proof}

\begin{corollary}[HD-COMP4: numerical fourfold bounds beyond the principal case]
\label{HD-COMP4}
For compatible integral maximal-order actions the following hold on
every fixed fourfold polarization datum, without a determinant
restriction:
\[
\begin{array}{c|c}
 K & \text{upper bound for }\delta_W(A,L)\\ \hline
 \mathbf Q(i) & 128\Delta^{12}D_{G,6}\\
 \mathbf Q(\sqrt{-7}) & 3D_{6,7}(\Delta^2)\\
 \mathbf Q(\sqrt{-2}) & 3D_{6,2}(\Delta^2).
\end{array}
\]
Here $D_{G,6}$ is the integer of SF-03, and $D_{6,7},D_{6,2}$
are the fully explicit functions of D-G24-05 and D-GF4-04.
For an integral action only of $\mathbf Z+f\mathcal O_K$, replace
$\Delta$ in these bounds by $f^8\Delta$. Thus the Gaussian
bound, for example, becomes $128f^{96}\Delta^{12}D_{G,6}$.
The much smaller principal Gaussian bound HD-G06 is retained
whenever its hypotheses hold.
\end{corollary}
\begin{proof}
For $K=\mathbf Q(i)$ choose $\tau=i$, so $c_K=2$ and apply
SF-04a to the split sixfold of polarization product $\Delta^2$:
$2(2\Delta^2)^6D_{G,6}=128\Delta^{12}D_{G,6}$.
For $K=\mathbf Q(\sqrt{-7})$ use an integral generator of norm
$2$, and for $K=\mathbf Q(\sqrt{-2})$ use $\tau=\sqrt{-2}$.
Both have $c_K=3$, so the two reconstructed sixfold bounds give
the remaining lines. For conductor $f$, put
$\Lambda'=f\mathcal O_K\Lambda$. Then
$f\Lambda\subset\Lambda'\subset\Lambda$ and
$I=[\Lambda:\Lambda']\le f^8$. The isogeny
$A'=V/\Lambda'\to A=V/\Lambda$ has pulled-back polarization
product $I\Delta$. Pushing its detector forward has exactly the
same degree relative to these polarizations, and keeps its nonzero
Weil projection. The displayed functions are increasing, so
$I\Delta\le f^8\Delta$ proves the stated bounds.
\end{proof}

The same contraction applied to D-SR-05 gives bounded division on
every fourfold determinant for each class-number-one field, using
the supplied Perry deformation input and without a Markman theorem
as a premise. At this stage its detector integer is existential for
the fields without a numerical split-sixfold bound. This is not
advertised as a numerical solution of those remaining gates.
In the project's frame these are new suppliers for the integral
condition A of cote~16 scan~050; their effective intersection uses
condition I of scan~049 and cote~44 Theorem~4 only after a higher
detector has been supplied. They are correspondences between the
weight-$2n+2$ and weight-$2n$ determinant realizations, with Tate
twists $(n+1)$ and $(n)$. The spread, $K_0$, and cycle-space
interpretations are those specified in the framework section.

% Generated by build_tex_sections.py; full proof text, no abridgement.
% Requires amsmath, amssymb, and UTF-8 support.
\section{Imprimitive sources: exact reduction and rank obstruction}

Run: HODGE\_DIVISION, agent \texttt{/root/imprimitive}. Date: 2026-11-21. Model: GPT-6 Astra Pro (session identity); effort: inherited, not externally exposed. The new proofs D-IM-01--08 have independent whole-proof checks, with the exact scope of each event recorded in \texttt{CHECK\_EVENTS.tsv}. The finite certificates, including the nine rank-one Nielsen tuples of D-IM-08, were also independently reconstructed.

\subsection{Conventions and scope}

Let \(d\) be 3, 4, or 6, \(e\) a positive divisor of \(d\), and \(m\geq 3\). Put \(K=\mathbb Q(\zeta_d)\). We use \(G(d,e,m)\) for the group of \(m\) by \(m\) monomial matrices whose nonzero entries lie in \(\mu_d\) and whose product belongs to \(\mu_{d/e}\). Thus the sum of the entry exponents is 0 modulo \(e\). Let \(\rho\) be its natural K-representation.

All claims below concern this natural representation and one multiplicity abelian factor of its rational conjugate pair. They do not concern every irreducible representation of the abstract finite group. A different representation must be checked separately.

The field and group are fixed on each Hurwitz family. A countable union of different group constructions is not asserted to have a single finite-type parameter space.

Published convention references: Lehrer--Taylor, *Unitary Reflection Groups* (Cambridge, 2009), Chapter 2, sections 1--3, pp.23--27, especially Proposition 2.10 as corrected in the authors' errata dated 5 February 2025; Lewis--Wang, *The Hurwitz action in complex reflection groups*, Combinatorial Theory 2(1) (2022), section 2.1. The needed induced-representation identification is proved below in the precise form used here. \cite{IM-LehrerTaylor,IM-LewisWang}

\subsection{D-IM-01. Every natural imprimitive factor is a cyclic Prym factor}
\label{D-IM-01}

\textbf{Statement.} Suppose \(C\) is a smooth connected projective complex curve with faithful \(G\)=\(G(d,e,m)\) action. Let \(H\) be the stabilizer of the first coordinate line, let \(\chi:H\to\mu_d\) be its action on that line, and put \(L=\ker\chi\). Then \(\chi\) is onto and \(H\)/\(L\) is cyclic of order \(d\). For

\[
Y=C/L,\qquad X=C/H,
\]

the morphism \(Y\to X\) is a connected cyclic d-cover. The abelian factor \(B\) associated with one copy of the rational conjugate pair of \(\rho\) is K-isogenous to the primitive cyclotomic factor \(P\) of \(J\)(\(Y\)). Its Hodge structure is exactly the pair of \(\chi\) and \(\bar\chi\) eigenspaces of \(H^1(Y)\).

If \(B\) is chosen as the image of the self-adjoint rational primitive idempotent described below, it is the image of \(P\) under pullback \(J\)(\(Y\))-\textgreater{}\(J\)(\(C\)). If \(\theta_C\) and \(\theta_Y\) are the canonical principal polarizations, the resulting isogeny \(u:P\to B\) satisfies

\[
u^*(\theta_C|_B)=|L|\,\theta_Y|_P.
\]

The isogeny kernel is killed by |\(L\)|, hence has order at most \(|L|^{4n}\) when \(\dim P=2n\). In particular the degree and polarization scaling are constant-bounded in a fixed Hurwitz family.

\textbf{Proof.} The coordinate-line stabilizer \(H\) has index \(m\). It contains \(\operatorname{diag}(\zeta_d,\zeta_d^{-1},1,\ldots,1)\), so \(\chi\) is surjective even when \(e= d\). The map

\[
\operatorname{Ind}_H^G\chi\longrightarrow K^m,\qquad g\otimes v\longmapsto vge_1.
\]

is well defined: replacing \(g\) by gh multiplies the image by \(\chi\)(\(h\)), as required in the tensor definition of induction. It is surjective, since permutation matrices carry \(e_1\) to every \(e_j\). Its source and target both have dimension \(m\), so it is an isomorphism.

The natural representation is absolutely irreducible. Indeed the diagonal normal subgroup has distinct coordinate characters when \(m\geq 3\) and \(d\geq 3\). (For coordinates i!=j, change the i-th diagonal entry by \(\zeta_d\) and a third entry by \(\zeta_d^{-1}\), leaving the j-th entry unchanged.) An invariant subspace is consequently a sum of coordinate lines, and transitivity of the permutation matrices makes that sum either zero or the whole space. The representation is not isomorphic to its conjugate: diagonal elements with exponent triples (1,1,d-2), padded by zeros, have nonreal trace for \(d= 3\),4,6. Thus \(\rho\) and \(\bar\rho\) are the two distinct complex constituents of a rational irreducible block \(\operatorname{Mat}_m(K)\).

For any \(K\)[\(G\)]-module \(V\), Frobenius reciprocity in this explicit induced model identifies \(\operatorname{Hom}_G(\rho,V)\) with its chi-eigenspace for \(H\). Apply this to the Hodge structure \(H^1(C,K)\). Pullback for \(p:C\to Y\) identifies \(H^1(Y,K)\) with \(H^1(C,K)^L\). To see this directly, the graph correspondences satisfy \(p_*p^*=|L|\) and \(p^*p_*=\sum_{l\in L}l^*\) on \(H^1\). The latter average is the invariant projector. The identifications preserve Hodge types and the K-action. Therefore

\[
\operatorname{Hom}_G(\rho,H^1(C,K))=H^1(Y,K)_\chi,
\]

and likewise for the conjugate character. This proves the asserted equality of multiplicity Hodge structures.

For an algebraic realization, use

\[
f=\frac1{|H|}\sum_{h\in H}\bigl(\chi(h)^{-1}+\bar\chi(h)^{-1}\bigr)h\in\mathbb Q[G].
\]

The coefficients are rational because \(K\) is quadratic. This is the sum of the two H-character projectors. Frobenius reciprocity and \(\operatorname{Ind}_H^G\chi=\rho\) show that it vanishes on all other G-irreducibles and has rank one on each of \(\rho\) and \(\bar\rho\). It is thus a rational primitive idempotent of the \(\operatorname{Mat}_m(K)\) block. It is self-adjoint for the group-invariant polarization: \(h^\dagger=h^{-1}\), and the sum of the conjugate characters is inversion invariant. Its image in \(J\)(\(C\)), understood as the image of any integral multiple of \(f\), is precisely \(p^*P\). Hence the equality of Hodge structures is induced by an algebraic isogeny, not just an abstract pointwise identification. Any other fixed primitive-idempotent realization is isogenous by fixed rational group-algebra matrix units.

For Jacobian pullback, the polarization adjoint of \(p^*\) is the norm \(p_*\), and \(p_*p^*=[|L|]\). Therefore \((p^*)^*\theta_C=|L|\theta_Y\); restriction gives the displayed equality. The kernel of \(p^*\) is annihilated by |\(L\)| by the same norm identity. On \(P[|L|]\), its order is at most \(|L|^{2\dim P}\). All correspondences and group denominators are fixed in the family. This completes the proof.

\subsection{D-IM-02. Exact Chevalley--Weil criterion in colored-cycle notation}
\label{D-IM-02}

\textbf{Statement.} Let (\(g_1\),...,\(g_r\)) be nonidentity elements of \(G(d,e,m)\) which generate \(G\) and have product one. They define connected G-covers of \(\mathbb P^1\). For each permutation cycle c of \(g_b\), write \(\ell_c\) for its length and \(a_c\) in \{0,...,d-1\} for the exponent of the product of its matrix entries. Define

\[
A=\sum_{b,c}(\ell_c-1),\qquad R=\#\{(b,c):a_c\ne0\},\qquad S=\sum_{\substack{b,c\\a_c\ne0}}\frac{a_c}{d}.
\]

The natural multiplicity abelian variety is Weil type of dimension 2n and signature (\(n\),\(n\)) if and only if

\[
A+R=2m+2n,\qquad S=R/2.
\]

The intermediate curve \(X\) of D-IM-01 has genus \(h=1-m+A/2\). The cyclic cover \(Y\to X\) is branched at precisely \(R\) points, with local character exponents the nonzero \(a_c\). Equivalently the signature criterion is

\[
R=2(n-h+1),\qquad\sum a_c=dR/2.
\]

\textbf{Proof.} A colored permutation cycle has characteristic polynomial \(t^\ell-\zeta_d^a\). Its eigenvalue angles in [0,1) are (a/\(d\)+k)/ell for k=0,...,ell-1. Their sum is (ell-1)/2+a/\(d\). The number of its eigenvalues different from one is ell-1 if a=0, and ell otherwise. Hence, writing age(\(g\)) for the sum of angles,

\[
\sum_b\operatorname{age}(g_b)=A/2+S,\qquad\sum_b\operatorname{age}(g_b^{-1})=A/2+R-S.
\]

Chevalley--Weil for the nontrivial irreducible \(\rho\) over a genus-zero base gives multiplicities

\[
-m+A/2+S,\qquad-m+A/2+R-S.
\]

in the two conjugate holomorphic multiplicity spaces, in an order depending only on the convention for monodromy versus pullback. Equality of both to \(n\) is exactly the two conditions stated. This is the formula already supplied in HODGE\_DIVISION\_ATTACK, Proposition O' and Corollary O''; it is used, not reproved.

The degree-m map \(X=C/H\to\mathbb P^1\) has local permutation monodromy g\_b's permutation. Riemann--Hurwitz therefore gives \(2h-2=-2m+A\). Over a branch value, each permutation cycle gives one point of \(X\). Going ell times around that point applies \(g_b^\ell\) to the chosen coordinate line, with scalar \(\zeta_d^a\). Thus \(Y\to X\) ramifies at that point exactly when a!=0, with the stated local exponent. Substitution of the Riemann--Hurwitz formula proves the equivalent form.

This is an exact criterion for any actual generating product-one tuple. It does not replace the product-one and generation tests or classify braid orbits of all such tuples.

\subsection{D-IM-03. The imprimitive route cannot dominate for n\textgreater{}=4}
\label{D-IM-03}

\textbf{Statement.} On every fixed Hurwitz component as above with Weil signature (\(n\),\(n\)),

\[
\operatorname{rank}(d\mathcal P)\leq\dim(\operatorname{im}\mathcal P)\leq\min(r-3,h+2n-1)\leq3n.
\]

Here rank means the rank at any smooth point where the relevant moduli maps are defined; the generic rank equals image dimension in characteristic zero. If \(e< d\), the last bound improves to 3n-1. Consequently no such natural-representation family dominates an \(n^2\text{-dimensional}\) Weil family for \(n\geq 4\). For \(n= 5\) the bounds are 15 when \(e= d\) and 14 when \(e< d\), both below 25.

\textbf{Proof.} D-IM-01 gives a family-level rational polarized-Hodge identification with a cyclic Prym factor. Forgetting the degree-m map \(X\to\mathbb P^1\) therefore factors the period map, up to a fixed isogeny correspondence, through the moduli of the cyclic cover \(Y\to X\). The local system and polarization identification is fixed; a discrete choice of finite kernel adds no tangent directions. By D-IM-02 that cyclic cover has base genus \(h\) and \(R=2(n-h+1)\) branch points. With fixed local exponents, there are finitely many cyclic covers over a fixed pointed \(X\): their root data differ by d-torsion in \(J\)(\(X\)). Their moduli dimension is consequently \(3h-3+R=h+2n-1\). For \(h= 0\) or1 the same formula holds because \(R\)=2n+2 or2n respectively, so the pointed curves are stable. For \(h\geq 2\) it is the ordinary stable pointed-curve formula; when \(R= 0\) it gives 3h-3. Also \(R\geq 0\) forces \(h\leq n+1\). This proves the bound and proves it at the differential level through factorization, without inferring rank from source dimension.

If \(e< d\) and \(R= 0\), every colored cycle of every \(g_b\) has color zero. Each \(g_b\) then has total entry product one, hence belongs to \(G(d,d,m)\), a proper subgroup of \(G(d,e,m)\). This contradicts generation. Thus \(R> 0\); since \(R=2(n-h+1)\) it is even, so \(R\geq 2\) and \(h\leq n\). The bound becomes 3n-1. The asserted non-dominance follows from \(n^2>3n\) for \(n\geq 4\).

The rank inequality in the first display should be read with dim(image) as the generic rank; at special points rank cannot exceed that generic rank. No exact rank for each arbitrary Nielsen class is asserted.

\subsection{D-IM-04. Large reflection Hurwitz spaces really exist}
\label{D-IM-04}

\textbf{Statement.} For all \(m\geq 3\) and \(n\geq 1\), there are product-one generating reflection tuples in \(G(d,d,m)\) of length 2m+2n whose natural factor has signature(\(n\),\(n\)); their source has dimension 2m+2n-3 but their period image has dimension at most3n. If \(e< d\), analogous tuples of length2m+2n exist with exactly two diagonal reflections and intermediate base genus \(n\); their image has dimension at most3n-1.

\textbf{Proof.} Let \(s_i\) exchange coordinates i and i+1 with no phases, for i=1,...,m-1. Let \(t\) exchange coordinates1,2 with opposite phases \(\zeta_d\),\(\zeta_d^{-1}\). Every \(s_i\) and \(t\) is an involutive reflection. The \(s_i\) generate \(S_m\); \(s_1\) \(t\) gives a diagonal matrix with opposite primitive phases in coordinates1,2. Its conjugates by \(S_m\) generate all diagonal matrices whose entry product is one. Hence \(s_1\),...,\(s_{m-1}\),\(t\) generate \(G(d,d,m)\).

List each of these \(m\) generators twice consecutively, and append \(n\) additional pairs (\(s_1\),\(s_1\)). The resulting tuple has product one, generates the group, and has length2m+2n. Every local eigenvalue list is (1,...,1,-1), so the two ages are both1/2 and Chevalley--Weil gives both multiplicities \(n\). Riemann existence produces a nonempty Hurwitz family of source dimension r-3. Every permutation branch point is a transposition and every cycle color is zero, so the intermediate cyclic cover is etale and \(h= n+1\).

For \(e< d\), adjoin the diagonal reflection \(u=\operatorname{diag}(\zeta_d^e,1,\ldots,1)\), which generates the quotient \(G(d,e,m)/G(d,d,m)\). Start with the 2m tuple above, add (\(u\),\(u^{-1}\)), then n-1 further pairs(\(s_1\),\(s_1\)). The product is one and the generated group is \(G(d,e,m)\). The two diagonal eigenangles sum to one; all others contribute1/2. The multiplicities again equal \(n\). The only nonzero cycle colors are \(e\) and d-e, giving \(R= 2\) and \(h= n\). This proves the statement.

At \(n= 5\),\(m= 9\) these tuples give 25-dimensional Hurwitz sources, with image at most15 or14. They are concrete counterexamples to treating the parameter count \(r-3\geq n^2\) as a sufficient dominance test.

\subsection{D-IM-03G. General monomial obstruction}
\label{D-IM-03G}

\textbf{Statement.} The bound3n also holds for an arbitrary nonreal irreducible complex representation \(\rho=\operatorname{Ind}_H^G\chi\) induced from a linear character, whenever its multiplicity factor is a Weil 2n-fold with holomorphic multiplicities(\(n\),\(n\)). The order of \(\chi\) need not be3,4,6. Thus any dominant curve-cover source for \(n\geq 4\) must use a nonmonomial irreducible representation. This is a statement about representations, not about whether the abstract group has some monomial realization.

\textbf{Proof.} Put \(L=\ker\chi\) and let \(d\) be the order of \(\chi\). Frobenius reciprocity identifies the \(\rho\) and \(\bar\rho\) multiplicity Hodge structures with the \(\chi\) and \(\bar\chi\) eigenspaces of the cyclic d-cover \(C\)/L-\textgreater{}\(C\)/\(H\). The character of \(H\)/\(L\) is faithful by construction. Adding the two cyclic Chevalley--Weil formulas therefore gives2n=2(h-1)+\(R\), where \(h\)=\(g\)(\(C\)/\(H\)) and \(R\) is the number of branch points of that cyclic cover. Every deformation preserving this cyclic cover preserves the two multiplicity Hodge filtrations; any flat rational structure on the multiplicity factor coming from the fixed G-action is unchanged. Hence the local period differential factors through deformations of the cyclic cover and has rank\textless{}=\(3h-3+R\)\textless{}=3n, exactly as in D-IM-03. This argument uses only the eigenspace variation; when phi(\(d\))\textgreater{}2, it does not assert that the entire primitive cyclotomic abelian factor of \(J\)(\(C\)/\(L\)) has dimension2n.

\subsection{D-IM-05. A certified genus-six cyclic Prym differential of rank15}
\label{D-IM-05}

\textbf{Statement.} For each \(d= 3\),4,6 there exists a smooth etale cyclic d-cover of a genus-six curve whose primitive cyclotomic Prym period map has differential of rank15. Thus the dimension bound3n is attained by the cyclic Prym family when \(n= 5\). \texttt{nodal\_rank.py} and \texttt{nodal\_rank\_certificate.json} are exact finite-field minor certificates; the proof interpreting them is given next.

\textbf{Proof.} Identify pairs (\(a_j\),\(b_j\))=(2j-1,2j), \(j= 1\),...,6, on \(\mathbb P^1\) to form an irreducible stable nodal curve \(X_0\) of arithmetic genus6. Let \(\eta\) restrict trivially to the normalization and have gluing \(\zeta=\zeta_d\) at every node. It has exact orderd. Choose residue conventions so that a section of \(\omega_{X_0}\) tensor \(\eta\) pulls back to

\[
\sum_jr_j\left(\frac1{x-a_j}-\frac{\zeta}{x-b_j}\right)dx,
\]

with \(\sum_j(1-\zeta)r_j=0\). Conversely every such differential descends. Thus the five sections \(F_j-F_6\), \(j= 1\),...,5, form a basis, where \(F_j\) is the displayed j-th elementary differential. Replacing \(\zeta\) by \(\zeta^{-1}\) gives a basis for the inverse twist.

Put \(D(x)=\prod_j(x-a_j)(x-b_j)\). After multiplying the differentials by \(D\)/dx, their numerators are the integral \(Z\)[\(\zeta\)] polynomials

\[
p_j(\zeta)=\frac D{x-a_j}-\frac{\zeta D}{x-b_j}-\frac D{x-a_6}+\frac{\zeta D}{x-b_6}.
\]

They have degree\textless{}=10. The25 products \(p_i(\zeta)p_j(\zeta^{-1})\) represent the multiplication map into \(H^0(X_0,\omega_{X_0}^2)\) as polynomials of degree\textless{}=20 over the common denominator \(D^2\). Their leading polar coefficients match at each glued pair, so they belong to that target. The target has dimension15 by Riemann--Roch and \(H^1(\omega^2)=0\); the two source spaces have dimension5 because \(\eta\) is nontrivial of degreezero and \(H^1(\omega\otimes\eta)=H^0(\eta^{-1})^*=0\).

The certificate reduces these integral coefficients modulo1009 under \(\zeta_3\mapsto374\), \(\zeta_4\mapsto469\), \(\zeta_6\mapsto375\). These elements have exact orders3,4,6. In every case the same15 columns (zero-based indices0,1,2,3,4,5,6,7,8,9,10,12,13,18,24) and rows0,...,14 have nonzero determinant. The residues are841,614,849 respectively. Hence these minors are nonzero already overQ(\(\zeta_d\)); the multiplication rank is at least15, and the target dimension makes it exactly15. This is an exact lower-rank certificate in characteristiczero, not a floating point numerical approximation.

It remains to pass to a smooth curve, rather than compute the differential on a singular moduli boundary. The torsion line bundle gives an etale mu\_d-cover of \(X_0\): take \(d\) copies of its normalization and glue above each node by the cyclic sheet shift specified by \(\zeta\). This cover is connected because the shift is primitive. Smooth each node of \(X_0\) by local plumbing uv=\(t\), and smooth all of its \(d\) preimages with the same parameter. The sheet permutations respect the plumbing. This produces a smooth connected etale cyclic d-cover over nearby nonzero parameters. Equivalently the etale torsor lifts over a deformation because [\(d\)] is etale on the relative smooth Picard space in characteristiczero. The source spaces \(H^0(\omega\otimes\eta^{\pm1})\) and target \(H^0(\omega^2)\) have constant dimensions5,5,15 near the special fiber, by the vanishings just recorded and cohomology and base change. Multiplication is therefore a morphism of vector bundles whose rank15 locus is open. A nearby smooth fiber has multiplication rank15.

For smooth etale cyclic covers, the codifferential of the primitive period map is precisely \(H^0(\omega\otimes\eta)\otimes H^0(\omega\otimes\eta^{-1})\to H^0(\omega^2)\), the primitive character-pair component of Lange--Ortega, *Prym varieties of cyclic coverings*, Geom. Dedicata150(2011)391--403, Proposition4.1. Thus the smooth period differential has rank15, proving the statement. \cite{IM-LangeOrtega}

\textbf{Scope.} The finite-field matrix verification and the independent whole-proof check of the descent conventions and smoothing/base-change argument are complete; see the recorded events. Exact image dimensions for arbitrary fixed monodromy tuples are not consequences of this one witness.

\subsection{D-IM-06. A certified two-branch-point genus-five differential of rank14}
\label{D-IM-06}

\textbf{Statement.} For \(d= 3\),4,6 and each proper divisor \(e\) of \(d\), the cyclic-cover moduli space of connected d-covers of genus-five curves with two ordered branch points and character exponents(\(e\),d-e) has a point where the primitive period differential has rank14. Its dimension is14 and its primitive factor has Weil signature(5,5).

\textbf{Proof.} Let \(X_0\) be the rational nodal curve obtained by identifying

\[
a_j=(2j)(2j+1)^d,\qquad b_j=2j,\qquad j=1,\ldots,5,
\]

and mark0 and \(\infty\). These twelve points on the normalization are distinct. Let \(\eta\) have normalization \(\mathcal O(\infty)\) and gluing multiplier

\[
q_j=\frac{\zeta_d}{(2j+1)^e},
\]

from the fiber at \(a_j\) to the fiber at \(b_j\), in the affine trivialization. Since \(q_j^d=(b_j/a_j)^e\), multiplication by \(x^{-e}\) on the normalization descends to an isomorphism

\[
\eta^d=\mathcal O_{X_0}\bigl(e[0]+(d-e)[\infty]\bigr).
\]

Indeed the ratio of the image trivializations at \(b_j\) and \(a_j\) is \(q_j^d(a_j/b_j)^e=1\). The normalized cyclic covering is \(w^d=x^e\). Its \(e\) irreducible components are indexed by \(w^{d/e}/x\) in \(\mu_e\). The gluing sends this label to \(\zeta_d^{d/e}\) times itself, so the whole nodal cover is connected. (If \(e= 1\) its normalization already is connected.)

Write \(B=[0]+[\infty]\) and \(M=\eta^{-1}(B)\). The normalization of \(M\) can be identified with \(\mathcal O(\infty)\) using the frame \(x^{-1}\eta^{-1}\). Its gluing multipliers are

\[
q_j^{\prime}=\frac{b_j/a_j}{q_j}.
\]

Sections of \(\omega\otimes\eta\) are exactly the span of

\[
F_j=\left(\frac1{x-a_j}-\frac{q_j}{x-b_j}\right)dx,\qquad j=1,\ldots,5,
\]

in the normalization trivialization; an allowed simple pole at \(\infty\) removes the total-residue constraint present in D-IM-05. Sections of \(\omega\otimes M\) similarly have basis \(F_j^{\prime}\) with \(q_j^{\prime}\). Their products, in the natural frame of \(\omega^2(B)\), are \(F_i\) \(F_j^{\prime}\)/\(x\). At the two branches of every node the quadratic leading coefficients agree, because \(q_jq_j^{\prime}=b_j/a_j\). Thus these are descended sections of \(\omega^2(B)\).

Both eigenspace dimensions are5 by Riemann--Roch (\(\eta\) and \(M\) have degree1); the target dimension is14. With \(D(x)=\prod_j(x-a_j)(x-b_j)\), the product numerators are

\[
\left(\frac D{x-a_i}-\frac{q_iD}{x-b_i}\right)\left(\frac D{x-a_j}-\frac{q_j^{\prime}D}{x-b_j}\right),
\]

over the common denominator \(xD^2\). They have degree\textless{}=18. \texttt{additional\_rank\_certificates.py} computes their coefficient matrix over the residue field \(\mathbb F_{1009}\), at the same roots of unity as D-IM-05. None of the node coordinates or denominators vanishes or collides modulo1009. A14 by14 minor is nonzero in each case:

\begin{center}
\begin{tabular}{lr}
\hline
(\(d\),\(e\)) & determinant modulo1009 \\
\hline
(3,1) & 133 \\
(4,1) & 303 \\
(4,2) & 106 \\
(6,1) & 960 \\
(6,2) & 487 \\
(6,3) & 420 \\
\hline
\end{tabular}
\end{center}

The exact rows, columns and matrices are in \texttt{additional\_rank\_certificates.json}. The rational denominators are units modulo1009, so nonzero reduction implies nonzero characteristic-zero determinant. These determinants were also recomputed by independent integer determinant arithmetic on the recorded finite-field matrices. Therefore the special multiplication map has rank14.

Equivariant plumbing at the nodes, away from0 and \(\infty\), smooths this connected cyclic cover while preserving its two local branch types. Cohomology and base change applies to the two source bundles and the target: their \(H^1\) groups vanish because the dual bundles have negative degree. Thus rank14 persists on a nearby smooth cover.

For completeness, the period-map formula in this two-exponent form follows from the same cup-product argument as Lange--Ortega Proposition4.1, but does not require their special uniform-branch-exponent convention. The tangent space to the fixed local-monodromy cyclic-cover moduli is \(H^1(X,T_X(-B))\), since the root choice is finite etale over pointed curves. Its action on each holomorphic eigenspace is the Kodaira--Spencer cup product. Serre duality identifies its dual with multiplication of the two eigendifferentials followed by trace to \(X\). The faithful conjugate characters at either branch point have complementary nonzero angles; in a local coordinate \(x=t^{d/e}\) their product descends to a quadratic differential with at most a simple pole. Thus the target is \(H^0(X,\omega_X^2(B))\), and under the eigensheaf identifications it is exactly \cite{IM-LangeOrtega}

\[
H^0(\omega\otimes\eta)\otimes H^0\bigl(\omega\otimes\eta^{-1}(B)\bigr)\longrightarrow H^0\bigl(\omega^2(B)\bigr),
\]

up to a fixed nonzero trace scalar. This is the map computed. The smooth primitive period differential therefore has rank14.

\subsection{D-IM-07. Exact maximal image dimensions for n=5}
\label{D-IM-07}

\textbf{Statement.} Fix \(d= 3\),4,6 and \(m\geq 9\). Among the natural-representation Hurwitz families in \(G(d,e,m)\) with Weil signature(5,5), the maximal image dimension is exactly15 if \(e= d\) and exactly14 if \(e< d\). In particular the D-IM-04 doubled-reflection families attain these values; their generic period differential has the exact corresponding rank. This does not assert that every possible Nielsen class has maximal rank.

\textbf{Proof, step1: connected cyclic-cover moduli.} Connected etale cyclic d-covers of genus \(h\) curves form an irreducible moduli space: their monodromy characters are the primitive vectors of \(H^1(X,\mathbb Z/d)\), on which the mapping class group acts transitively. One way to verify the finite-modulus version is to lift a primitive residue vector to a primitive integer vector, extend it to a symplectic basis, and use the integral symplectic action. Primitive lifting follows by changing one coordinate by multiples of \(d\) to avoid the finitely many prime divisors of the gcd of the remaining coordinates; if all remaining coordinates vanish, change one of them to \(d\) first. The integral symplectic action on primitive vectors is transitive. The mapping class group surjects onto \(\operatorname{Sp}_{2h}(\mathbb Z)\) (Farb--Margalit, *A Primer on Mapping Class Groups*, Theorem6.4; Proposition6.2 and the change-of-coordinates principle give the primitive-vector form). \cite{IM-FarbMargalit}

For two ordered branch points with prescribed monodromies \(e\),-e modulo \(d\), write a character of the punctured surface group as its2h handle values \(v\) in \((\mathbb Z/d)^{2h}\) and those two puncture values. It defines a connected cover precisely when \(v\) modulo \(e\) is primitive, because \(\gcd(d,e,v_1,\ldots,v_{2h})=\gcd(e,v_1,\ldots,v_{2h})\). Point-pushing the first branch point around a handle loop changes the corresponding dual handle value by \(e\) and leaves puncture values unchanged: on punctured homology it adds the puncture loop times the relevant algebraic intersection number. Independent handle loops therefore add arbitrary elements of \(e(\mathbb Z/d)^{2h}\). Modulo \(e\), the symplectic action is transitive on primitive vectors by the preceding argument. Given two connected-cover characters, first match their reductions modulo \(e\) by a mapping class, then use point-pushes to match the lifts modulo \(d\). Thus the action is transitive on the entire set of connected-cover characters. This proves connectedness, hence irreducibility, of the corresponding smooth finite cover of the moduli stack of pointed curves. The point-pushes used here are those of Farb--Margalit, Theorem4.6 and Fact4.7; the homology action just described also follows by multiplying the two Dehn twists in Fact4.7. When \(e= 1\), the point-push translations alone are transitive. \cite{IM-FarbMargalit}

\textbf{Step2: forgetting the map to \(\mathbb P^1\) is locally onto.} Take any smooth cover represented by a D-IM-04 tuple and form its tower \(Y\to X\to\mathbb P^1\). The map \(f:X\to\mathbb P^1\) is simply branched, of degree \(m\). If \(e= d\), \(X\) has genus6 and \(Y\to X\) is etale. If \(e< d\), \(X\) has genus5 and \(Y\to X\) has two branch points \(p\),\(q\) of exponents \(e\),d-e; these are unramified points of \(f\) and their images are distinct from all critical values. All these are open conditions.

Given a first-order deformation of the cyclic cover \(Y\to X\), equivalently of \(X\) (and \(p\),\(q\) in the ramified case) with its discrete root data, the obstruction to extending \(f\) lies in \(H^1(X,f^*T_{\mathbb P^1})\). This vanishes because \(\deg f^*T_{\mathbb P^1}=2m>2h-2\). More explicitly, lift \(f\) on affine patches using the smoothness of \(\mathbb P^1\). Differences of local lifts form a Cech1-cocycle in \(f^*T_{\mathbb P^1}\); vanishing of \(H^1\) allows the lifts to be corrected and glued. The same argument works over successive small extensions. Therefore the local space of such maps is smooth over the deformation space of the cyclic cover, and the forgetful map has surjective differential. The simple branch conditions remain open, and the finite monodromy tuple of the Galois closure is locally constant in this open set; hence we remain in the original full-G Hurwitz component.

It follows that each doubled-reflection Hurwitz component maps onto an open set of the relevant irreducible cyclic-cover moduli. In the etale genus6 case, D-IM-05 gives generic period rank15. In the ramified genus5 case, D-IM-06 gives generic period rank14. Surjectivity of the forgetful differential implies the same generic rank on these Hurwitz components. The global upper bounds in D-IM-03 show that these are the exact maximal image dimensions.

\textbf{A note on the domain degree assumption.} Here \(m\geq 9\) is more than needed for the forgetful deformation argument: \(m> h-1\) suffices. No Brill--Noether existence theorem is required because D-IM-04 already supplied a nonempty family and hence an initial degree-m map \(f\).

\subsection{D-IM-08. Explicit Nielsen classes with exact image dimension one}
\label{D-IM-08}

\textbf{Statement.} For every \(d= 3\),4,6 and every divisor \(e\) of \(d\), \(G(d,e,12)\) has a four-branch generating product-one tuple whose natural multiplicity factor has Weil signature (5,5). Every Hurwitz component containing the specified tuple has period image of dimension exactly one. Thus the maximal dimensions 15 and 14 in D-IM-07 are not the ranks of all Nielsen classes.

\textbf{Explicit matrices.} Write \(\zeta=\zeta_d\) and use the convention that a colored permutation (pi,c) sends \(e_j\) to zeta\textasciicircum{}(c\_j)\(e_{\pi(j)}\). Let a have permutation (1 2 ... 12). If \(e= d\), all its colors vanish; let b have permutation (1 2 ... 7), with colors +1 at domain coordinate 1, -1 at domain coordinate 7, and zero elsewhere. If \(e< d\), give a the single color \(e\) at domain coordinate 12; let b have permutation (1 2 ... 6), with colors +1 at coordinate 1, -1 at coordinate 2, and zero elsewhere. Take

\[
(g_1,g_2,g_3,g_4)=(a,a^{-1},b,b^{-1}).
\]

Every matrix belongs to \(G(d,e,12)\), every entry of the tuple is nonidentity, and its product is one.

\textbf{Generation.} First consider the uncolored permutations \(a_0\)=(1 ... 12) and \(b_0\)=(1 ... k), where k=7 or 6. The permutation \(b_0^{-1}a_0b_0a_0^{-1}\) is the three-cycle (1 k k+1). Its a\_0-conjugates have supports containing every adjacent pair on the twelve-cycle. These three-cycles generate \(A_{12}\). Here is an elementary justification that avoids a permutation-group classification: three-cycles whose support hypergraph is connected generate the alternating group on that support. Start with one three-cycle, hence \(A_3\). If a further support meets the accumulated support in two vertices, its conjugates by the existing alternating group add the new vertex and generate the larger alternating group. If it meets in one vertex and adds \(x\),\(y\), conjugate its old vertex to a second old vertex; the product of one cycle with the inverse of the other gives a three-cycle with two old vertices and \(x\). Apply the preceding step twice. Connectedness permits this induction until all vertices have been added. Since \(a_0\) is odd, \(a_0\),\(b_0\) generate \(S_{12}\).

Let \(H\) be the subgroup generated by the colored a,b. Its permutation projection is therefore \(S_{12}\). Put \(c=a^{-1}b\). Direct monomial multiplication gives, in the two cases respectively,

\[
\begin{aligned}e=d:\quad c^6&=\operatorname{diag}(\zeta^6,1,1,1,1,1,\underbrace{\zeta^{-1},\ldots,\zeta^{-1}}_{6}),\\e<d:\quad c^7&=\operatorname{diag}(\zeta^7,\zeta^{-7},1,1,1,\underbrace{\zeta^{-e},\ldots,\zeta^{-e}}_{7}).\end{aligned}
\]

The last repetitions have lengths six and seven, respectively. Conjugating a diagonal matrix by a monomial matrix depends only on its permutation. In the first case, commute \(c^6\) with any lift in \(H\) of the transposition (2 7). The result has primitive opposite colors at coordinates 2 and 7. In the second, commute \(c^7\) with a lift of (1 3); the result has opposite colors 7 and -7 at coordinates 1 and 3. Since 7 is invertible modulo \(d\), its power supplies a primitive opposite-color diagonal. Conjugating by the permutation projection gives all such diagonals. Consequently \(H\) contains the entire diagonal subgroup \(D_0\) of total color zero.

If \(e= d\), this proves \(H\)=\(G(d,d,12)\). If \(e< d\), the quotient \(G(d,e,12)\)/\(D_0\) is the direct product of the total-color group \(\mu_{d/e}\) with \(S_{12}\). The images of a,b are (zeta\textasciicircum{}\(e\),\(a_0\)) and (1,\(b_0\)), and both permutations are odd. The commutator subgroup of their image contains \{1\} times \(A_{12}\): it maps onto \(A_{12}\) and lies in that subgroup, where permutation projection is injective. Modulo it, (zeta\textasciicircum{}\(e\),-1) and (1,-1) generate \(\mu_{d/e}\) times \{+1,-1\}. Thus the image is the full quotient and \(H\)=\(G(d,e,12)\). This completes the generation proof for every listed divisor, without an enumeration of the enormous full groups.

\textbf{Signature and cyclic tower.} A colored cycle contributes one fixed vector precisely when its total color is zero. Therefore

\[
\begin{gathered}(\dim\operatorname{Fix}(a),\dim\operatorname{Fix}(b))=\begin{cases}(1,6)&e=d,\\(0,7)&e<d,\end{cases}\\n=-12+\operatorname{codim}\operatorname{Fix}(a)+\operatorname{codim}\operatorname{Fix}(b)=5.\end{gathered}
\]

The inverse pairs make the two Chevalley--Weil multiplicities equal: age(\(g\))+age(g\textasciicircum{}(-1))=codim Fix(\(g\)). The cyclic tower of D-IM-01 has \(h= 6\),\(R= 0\) in the first case and \(h= 5\),\(R= 2\) in the second, in agreement with D-IM-02. Riemann existence and the generation proof give the claimed nonempty Hurwitz curves, finite over \(M_{0,4}\).

\textbf{Exact period rank.} Degenerate the four-punctured sphere along the partition separating (a,\(a^{-1}\)) from (b,\(b^{-1}\)). The monodromy at the separating node is one. The admissible G-cover has components indexed by \(G\)/\textless{}a\textgreater{} on one side and \(G\)/\textless{}b\textgreater{} on the other, and nodes indexed by \(G\). Every component is rational: its map to the corresponding base component is a cyclic cover with just the two inverse branch points. The dual graph is connected because a,b generate \(G\). Its cellular chain sequence is

\[
0\longrightarrow H_1(\Gamma,\mathbb Q)\longrightarrow\mathbb Q[G]\longrightarrow\mathbb Q[G/\langle a\rangle]\oplus\mathbb Q[G/\langle b\rangle]\longrightarrow\mathbb Q\longrightarrow0.
\]

For the nontrivial natural representation, Frobenius reciprocity gives the multiplicity in graph homology as \(12-\dim\operatorname{Fix}(a)-\dim\operatorname{Fix}(b)=5\). A smoothing is obtained by plumbing all nodes over the base node with a common parameter. To see the monodromy rank directly, let \(\delta_f\) be the oriented core circle in the annulus smoothing the edge \(f\). The vertex-incidence relations are exactly the relations among these circles: each punctured rational component contributes the sum of its oriented boundary circles, and Mayer--Vietoris gives no others. Their span is therefore graph cohomology, of dimension equal to the first Betti number of the graph. The dual paths associated to graph loops show that the intersection pairing onto the dual of this span is surjective.

Going once around the plumbing parameter applies one positive Dehn twist at every annulus. A transverse path gains its oriented core circle, so the logarithm is the sum of the transvections \(x\) -\textgreater{} (\(x\),\(\delta_f\))\(\delta_f\). Its square is zero because the core circles are pairwise disjoint. Its image is their entire span: after choosing a real basis of that span, the associated map from its dual is the sum of the positive rank-one matrices \(\delta_f\) \(\delta_f^t\), which is positive definite because the \(\delta_f\) span. Finite return powers only replace these coefficients by positive integers. The construction is G-equivariant, so taking natural multiplicity spaces preserves surjectivity onto the corresponding graph-cohomology factor. The logarithm on that factor has rank five, and in particular is nonzero. This proves the needed weight-one degeneration statement here without assuming a separate rank theorem for the Prym map.

A constant polarized period map cannot have such monodromy. A polarization isometry preserving a fixed weight-one Hodge filtration is unitary for its positive Hodge metric; a unitary unipotent operator is the identity. The nonzero logarithm contradicts constancy. Hence the period map of the Hurwitz curve is nonconstant. Its image has dimension one and its generic differential has rank one, as asserted.

The certificate \texttt{nielsen\_rank\_one.py} checks all nine (\(d\),\(e\)) cases by exact colored-permutation arithmetic: membership, product one, the two diagonal powers, the conjugate three-cycle supports, fixed-space and age counts, intermediate genera, cyclic branch counts, and graph multiplicity. The generation and period interpretations are the full proofs just given; the script does not pretend to enumerate each full \(G(d,e,12)\).

\subsection{Remaining gate and programme implications}

Target4's dominance question for the specified natural \(G(d,e,m)\) sources is answered negatively for every \(n\geq 5\), not merely \(n= 5\). D-IM-07 gives exact maximal image dimensions 15 and 14 for \(n= 5\) and explicit families achieving them; D-IM-08 gives exact rank-one Nielsen classes for every \(d\),\(e\) at \(m\)=12. A complete rank classification of arbitrary Nielsen classes and arbitrary \(m\) is not asserted. The remaining Q4c gate is narrower: for a specified Nielsen class, compute the rank of the cyclic-Prym multiplication map after restricting its dual to the image of the forgetful deformation map. This is a finite algebraic rank problem once an explicit family and its deformation constraints are given. Fixed higher ramification can impose constraints on those deformations, so the simple-branch submersion proof of D-IM-07 does not apply automatically. The strongest next dominant curve-source gate remains a nonmonomial irreducible representation, or induction from a representation of dimension greater than one, with an actual period-rank computation. Higher-dimensional perfect complexes are a separate possible source.

One must not conclude that no curve source can dominate for \(n\geq 5\). Corollary O'' permits nonreflection generators when \(m\) exceeds its lower bound. The Shephard--Todd classification of groups generated entirely by reflections then no longer exhausts the allowed groups. Even at all-reflection tuples, the present theorem only treats the natural imprimitive representation.

In the project's frame, D-IM-03 says that this geometric attempt cannot supply the dominating effective-cycle family of Lemma Dom. It therefore supplies no new bound on \(\delta_W\) or \(c_t\). It does not challenge their equivalence or conditionA (cote16, scan050, SIGN\_NEXT SN-C03/C04). The rational group-algebra projector in D-IM-01 is an algebraic correspondence and an idempotent of the \(h^1\) motive; this is distinct from producing a codimension-n detecting cycle. Positivity/conditionI (cote16 scan049, cote44 Theorem4) cannot remove this image-dimension obstruction.

The same quotient-pullback/trace identities make sense for \(h^1\) motives and for l-adic sheaves after inverting |\(G\)| and spreading the fixed cover construction over a finitely generated Z-algebra. No assertion about new weights or L-values is needed. In K-theory or lambda-ring language the induced representation is an actual equality \([\rho]=\operatorname{Ind}[\chi]\), and applying \(\operatorname{Hom}_G\) to the \(h^1\) realization realizes that equality; it does not identify exterior determinant classes with effective algebraic cycles. Failure of these source constructions is not a ghost or a counterexample to Hodge; a ghost conclusion would require the independent conditional hypotheses already specified in GHOSTS B03.

\section{Exact period ranks for specified Nielsen components}
\label{sec:nielsen-rank-algorithm}

\paragraph{Scope and check status.}
This section gives a terminating exact algorithm, not a practical running-time
bound or a table for unbounded coefficient degree.  It constructs an algebraic
generic cover of the specified component and computes its period rank by
linear algebra in a finitely generated field.  No period integral is used.
The proof and the finite example below are new in this run.  The parent
completed an independent whole-proof review of D-NR-01--04 on
21 November 2026, including the primary Romagny--Wewers inputs, and
recorded PASS.  The finite matrix example was executed exactly; the
general generic-model search has not been executed.  The classical algebraic algorithms used are
exact arithmetic in algebraic numbers and finitely generated fields,
Gr\"obner-basis elimination, and isolation of algebraic roots.

\subsection*{D-NR-01. Input and an exact monodromy subroutine}

The input consists of a finite group $G$, an ordered generating tuple
$\boldsymbol g=(g_1,\ldots,g_r)$ of nonidentity elements with product one,
and an absolutely irreducible representation $V$ over an explicitly given
imaginary quadratic field $K$.  Its dual is not isomorphic to $V$, and its
two multiplicities in $H^0(C,\omega_C)$ are assumed to be $(n,n)$.
The desired output is the generic rank on the connected ordered Hurwitz
component containing $\boldsymbol g$, modulo simultaneous inner conjugacy.
Finite additional level structures do not change this rank.  Cases $r\le3$
have rank zero.  We henceforth take $r\ge4$.  All applications here have
$g(C)\ge2$.

We first describe a subroutine which determines the labelled monodromy of
an algebraic $G$-cover over $\overline{\mathbf Q}$.  Choose a primitive
function $y$ for its function field over $\overline{\mathbf Q}(z)$, with
separable equation
\[
                 p(z,y)=0,\qquad \deg_y p=|G|.                 \tag{NR.1}
\]
Clear denominators so that $p$ is polynomial in both variables; this adds
the zeros of the cleared denominators to the exceptional set.
The deck maps are part of the exact input.  Include among the exceptional
points the poles of coefficients, the zeros of the leading coefficient and
the discriminant of $p$, and the finitely many points needed to evaluate
the chosen rational deck-map formulas.  These are an effectively computable
finite set of algebraic points.  Standard polygon generators of the
punctured sphere can be perturbed, using rational vertices, to avoid this
set while retaining their homotopy classes relative to the actual branch
points.  Small disjoint rational boxes about the exceptional points make
this a finite planar construction.  A loop around an extra discriminant
point has trivial monodromy on the normalization; such a point is not
silently treated as a new branch point.

Here is an entirely algebraic decision procedure for lifting one polygon
edge.  Subdivide it into $2^a$ equal segments.  At every subdivision vertex,
isolate the $|G|$ distinct algebraic roots of (NR.1).  Let $\delta$ be the
minimum distance between these roots at that vertex.  Choose rational
centres within $\delta/128$ of the roots and a common rational radius in
$[\delta/32,\delta/16]$.  Isolation and comparisons of algebraic real
numbers produce these choices exactly.  The resulting closed discs are
pairwise disjoint and each initially contains exactly one root.

Write $p(z,w)=\sum_{j=0}^N a_j(z)w^j$ and
$a_j(z)=\sum_k a_{jk}z^k$, where $N=|G|$.
For the segment $[z_0,z_1]$, put $h=|z_1-z_0|$ and choose a rational
$Z\ge\max(|z_0|,|z_1|,1)$.  For a root disc with centre $c_i$ and
radius $s_i$, choose a rational $R_i\ge|c_i|+s_i$ and a rational
$\varepsilon_i>|c_i-\alpha_i|$, still smaller than $\delta/128$.
On its boundary, factorization at the algebraic endpoint gives the
strictly positive lower bound
\[
 \ell_i=|a_N(z_0)|(s_i-\varepsilon_i)
               \prod_{j\ne i}(|c_i-\alpha_j|-s_i)>0.
\]
Choose a rational $L_i$ with $\ell_i/2<L_i<\ell_i$.
The coefficient derivative estimate on the straight segment gives
\[
 |p(z,w)-p(z_0,w)|\le h D_i,\qquad
 D_i\ge\sum_{j=0}^N R_i^j\sum_{k\ge1}k|a_{jk}|Z^{k-1},
 \qquad h^2D_i^2<L_i^2,                                    \tag{NR.2}
\]
where $D_i$ is a rational upper bound, chosen within a factor two after
adding one if the exact expression is zero.  All comparisons and bounds
in this test are computable with algebraic numbers and rational intervals;
$h^2$ is rational because the polygon vertices are rational.  Test the
last strict inequality for every root disc.  If it passes, the polynomial
homotopy $p(z_0,w)+u(p(z,w)-p(z_0,w))$, $0\le u\le1$, has no zero
on the boundary.  Its winding number, hence its number of zeros in the
disc, is constant.  This is the elementary Rouch\'e argument, and it proves
that exactly one root remains in each disc throughout the segment.
Identify the endpoint root by exact disc membership.  If any inequality
fails, increase $a$ and start the edge again.

This process terminates.  On the compact edge the nonvanishing leading
coefficient has a positive minimum, the roots are bounded, and their
minimum separation is positive.  The prescribed radii and centres make
every factor in $\ell_i$ uniformly bounded below by a positive constant:
$s_i-\varepsilon_i\ge3\delta/128$ and
$|c_i-\alpha_j|-s_i>\delta/2$ for $j\ne i$.
The chosen rational $L_i$ retain at least half this lower bound.
The $R_i$, $Z$ and hence the $D_i$ are uniformly bounded above; take their
rational upper bounds within a fixed additive one when choosing them.
Dyadic subdivision sends $h$ to zero, so eventually every inequality
(NR.2) holds.  All accepted steps are finite exact rational certificates.
This proves both the lifting result and its error control without any
uncertified analytic continuation.

At the common base point choose one point of the fibre.  Its $G$-orbit
identifies the fibre with the regular $G$-set.  Monodromy commutes with the
deck action, so the computed permutations recover elements of $G$, unique
up to simultaneous inner conjugacy.  At increasing real branch values
use the standard ordered positive polygon generators, transported from
$0,1,\ldots,r-1$ through the increasing real chamber.  The same tails,
loop orientations and multiplication order are used for the input tuple.
With a left deck action, continuation acts by right multiplication on
the fibre labelled by $G$; if function pullbacks instead encode the
opposite action, invert that encoding before the tuple comparison.
In this convention the standard positive generators have ordered product
one.  This gives a
terminating exact test for equality with the input inner Nielsen tuple.

\subsection*{D-NR-02. Constructing the specified generic component}

Put $k=\overline{\mathbf Q}$ with its standard computable presentation by
minimal polynomials and isolating rectangles.  Normalize the first three
ordered branch points to $0,1,2$, and let
\[
 B=\{(t_4,\ldots,t_r):0,1,2,t_4,\ldots,t_r
                  \text{ are distinct}\}\subset\mathbf A^{r-3}_k.
                                                               \tag{NR.3}
\]
This is a dense open coordinate chart in the ordered moduli of punctured
spheres.  Write $K_0=k(t_4,\ldots,t_r)$.  The following search is an actual
enumeration of finite algebraic descriptions.

All enumerations below are dovetailed: enumerate complete finite packages
in increasing total description length, including the defining polynomial
of the coefficient field, every chart, map, and certificate.  At each length
there are only finitely many strings in a fixed effective alphabet.  The
phrase ``for each $F$'' therefore never means finishing an infinite search
over the first field before proceeding to the next.  Since every individual
test terminates, every valid finite package is eventually tested.

Enumerate monic irreducible polynomials $P(T)\in K_0[T]$ and put
$F=K_0[T]/(P)$.  Irreducibility here is an effective test: after clearing
denominators and removing the polynomial content in the $t_i$, Gauss's
lemma reduces it to factorization of a multivariate polynomial over the
computable algebraically closed field $k$.  Finite-degree coefficient
equations with nonzero leading coefficients decide that factorization.
For each $F$, enumerate projective curves, rational maps
and their finite affine-chart descriptions over $F$.  Retain a package only
when it passes all the following exact tests: the curve is smooth and
geometrically connected; its map $\pi:C\to\mathbf P^1_F$ is finite of
degree $|G|$; the specified maps give a faithful action of $G$ over
$\mathbf P^1_F$; and the branch locus is exactly the ordered divisor
\[
                    0+1+2+t_4+\cdots+t_r.                    \tag{NR.4}
\]
We may restrict to the genus prescribed by Riemann--Hurwitz for the input
tuple.  Include a primitive plane equation for $F(C)/F(z)$ and explicit
birational identifications in the package, so the monodromy subroutine has
its required presentation.

For completeness, the tests just listed are algebraic decisions, not
oracles for the desired period rank.  Polynomial identities and equality
of rational maps are ideal-membership tests.  Smoothness is the Jacobian
criterion on finitely many affine charts.  Dimensions, degrees and the
Hilbert polynomial are computed from initial monomial ideals; the constant
term of the curve's Hilbert polynomial gives its genus.  Geometric
irreducibility can be tested by absolute decomposition.  Alternatively,
in the enumerated package one can require a birational identification with
an absolutely irreducible plane equation, a smooth one-dimensional
projective curve, and a zero-dimensional complement to that identification.
Absolute irreducibility of a plane equation is decided by trying the
finitely many possible positive factor degrees and testing the coefficient
equations over an algebraic closure.  Enumerate the possible nonzero
leading coefficients of each factor, fix one scaling coefficient to one,
and impose each required nonvanishing by an auxiliary equation $uv=1$.
Gr\"obner elimination decides solvability of each such finite system.
This alternative avoids a general primeness
oracle.  The absence of base points and the identities on overlaps verify
the chart descriptions of morphisms.  A nonconstant morphism from the
smooth projective curve to $\mathbf P^1$ is finite; its degree is checked
in its primitive function-field presentation.  The group law and
distinctness of all $|G|$ automorphisms then show that it is a $G$-cover.
Finally the torsion module $\Omega^1_{C/\mathbf P^1}$ is computed from
the differentials of the chart equations.  Its support, projected to
$\mathbf P^1$, is computed by elimination and is compared with (NR.4),
including the chart at infinity.  These operations terminate over the
computable characteristic-zero field $F$.

Spread a retained package to a good family over a nonempty open in the
normalization of $B$ in $F$.  This step is effective as well: clear all
denominators and remove the images of the nonsmooth, nonfinite,
nonflat, undefined-map and wrong-branch loci.  Jacobian ideals,
Fitting ideals, affine-cover certificates and elimination compute these
loci.  Generic flatness can be effected by inverting the finitely many
leading coefficients in the relevant module Gr\"obner bases.  Generic
geometric connectedness and proper smoothness then give geometric
connectedness throughout this open.  Each excluded locus is proper,
because the generic package has passed the tests.  Equivalently one may
enumerate base localizations and their finite Jacobian, Fitting and
covering certificates; generic validity guarantees that a successful
localization occurs, and the tests themselves terminate.

We need an open of the \emph{base} over which \emph{every} sheet is good.
After clearing coefficients of $P$ and inverting its discriminant, the
monogenic algebra is already finite etale over an open of $B$.  If a
further nonzero element of this algebra has to be inverted, invert its
nonzero field norm down to $K_0$.  After finitely many such steps we obtain
explicit $D\ne0$ and
\[
 U=B\cap\{D\ne0\},\qquad
 S=\operatorname{Spec}\bigl(k[t_4,\ldots,t_r,D^{-1},T]/(P)\bigr)
       \longrightarrow U                                    \tag{NR.5}
\]
finite etale, with the entire package a smooth family of connected
$G$-covers over $S$.  Factors needed to clear coefficients and branch
diagonals are included in $D$.  Because $P$ is irreducible over $K_0$,
$S$ is integral and $S(\mathbf C)$ is connected.

Choose integers $2<t_4<\cdots<t_r$ with $D\ne0$ by enumeration.
Such integers exist: that ordered integer set is Zariski dense, as follows
by induction on the number of variables from the finiteness of the zeros
of a nonzero univariate polynomial.  Compute \emph{every} root of
$P$ at these integers.  Each is an algebraic point of $S$ and gives an
exact algebraic $G$-cover with real ordered branch points.  Apply D-NR-01
to each of these finitely many covers.  If one has the input tuple up to
simultaneous inner conjugacy, retain this generic package.  Otherwise
continue the enumeration.

We prove both the correctness and termination of this selection.
The inner Hurwitz space of complex connected $G$-covers is a finite
topological covering of branch configuration space; its fibre is exactly
the finite set of generating product-one tuples modulo inner conjugacy.
The local charts and the analytic classifying map for a family are proved
in \cite[Theorem~3.1, Propositions~3.2 and~3.4, pp.~319--321]{NR:RW}.
Ordered branch points and the normalization (NR.3) are obtained by base
change and do not alter this assertion.  The classifying map from
$S(\mathbf C)$ is a map of finite covering spaces over $U(\mathbf C)$.
It is therefore a covering onto one connected Hurwitz component over
$U$: locally it maps a sheet isomorphically to a sheet, so its image is
open; finiteness over $U$ makes the image closed.  Consequently its finite
fibre over the chosen integer branch configuration contains every inner
Nielsen point in that component.  This is why all roots of $P$, not one
arbitrary root, are tested.

Restriction from $B$ to the nonempty Zariski open $U$ does not split a
connected smooth Hurwitz component: an irreducible complex algebraic
variety remains irreducible on a nonempty Zariski open.  The algebraic
Hurwitz scheme and its agreement with this analytic covering are
\cite[Theorem~4.11, pp.~328--331, and Remark~4.15(i)]{NR:RW}.
Thus a successful tuple test identifies exactly the required component.
Since $S\to U$ is finite etale, its family dominates that component and
has the same generic period rank.

\paragraph{Grothendieck's algebraic framework.}
This construction sits in the algebraic theory of tame Hurwitz schemes
over $\mathbf Z[1/|G|]$: after passage to characteristic zero, the
branch-coordinate morphism is a finite etale cover of the appropriate
configuration space.  The construction and comparison used here are
\cite[Theorem~4.11 and Remark~4.15]{NR:RW}, together with their
finite-field-of-definition Theorem~4.3.  They permit the passage from a
topologically specified cover to a finite algebraic generic description.
The algorithm itself uses only characteristic zero and computes an exact
generic Hodge matrix over a finite extension of
$\overline{\mathbf Q}(t_4,\ldots,t_r)$.  No stronger effective theorem in
positive characteristic is being asserted.  In particular we do not
attribute a direct field-of-definition statement to an unspecified SGA
locator; Romagny--Wewers explicitly discuss that source issue.

Conversely let $\mathcal H$ be the component specified by the input.
It is defined over $k$, and its function field is a finite extension of
$K_0$.  A cover representing its geometric generic point is defined over
a finite further extension $F$ of that function field: its equations,
the cover map and the finitely many labelled group automorphisms use only
finitely many algebraic coefficients.  This also follows directly from
the finite-field-of-definition theorem
\cite[Theorem~4.3 and its proof, pp.~323--327]{NR:RW}.
The center of $G$ is allowed; we do not assume that the coarse generic
point itself carries a universal cover.  The finite further extension
removes exactly this possible field-of-moduli obstruction.
Characteristic zero supplies a primitive element for $F/K_0$.
The resulting cover, any smooth projective embedding and its finite map
descriptions occur in the enumeration.  They spread to a package of the
form (NR.5), and its tuple test succeeds.  Every earlier candidate is
decided in finitely many steps.  Therefore the search terminates.

\subsection*{D-NR-03. A finite matrix whose rank is the answer}

For the accepted generic cover set
\[
 A=\operatorname{Hom}_{G}(V_F,H^0(C,\omega_C)),\qquad
 A'=\operatorname{Hom}_{G}(V_F^*,H^0(C,\omega_C)).
                                                               \tag{NR.6}
\]
Both spaces have dimension $n$.  Choose a basis $e_1,\ldots,e_m$ of $V$
and its dual basis.  Define
\[
 \mu:A\otimes_F A'\longrightarrow H^0(C,\omega_C^2)^G,
 \qquad
 \mu(a\otimes b)=\sum_{j=1}^m a(e_j)b(e_j^*).                \tag{NR.7}
\]
The expression is independent of the dual bases and is $G$-invariant.

All parts of (NR.6)--(NR.7) can be computed exactly.  One explicit
terminating method to compute $H^0(C,\omega_C)$ is useful here.  Every
rational differential has the form $f\,dz$ with $f\in F(C)$.
Enumerate these differentials.  On a finite smooth affine cover, regularity
is membership in the finitely presented K\"ahler differential module.
After a common denominator is taken, membership of $v/b$ is the module
membership test $v\in b\Omega^1$; the differential module is torsion-free
because the curve is smooth.  Module Gr\"obner bases decide it.  Keep
$F$-linearly independent regular differentials until $g(C)$ are found.
Riemann--Roch guarantees termination.  Linear dependence over $F$ is
decided by writing elements in the basis of $F(C)/F(z)$, clearing their
$z$-denominators, and comparing polynomial coefficients.  Pullback by each
deck map is exact in the same function field.  The intertwiner equations
defining (NR.6) are consequently finite linear equations over $F$.

Write $b_1=0,b_2=1,b_3=2,b_i=t_i$ and
$Q(z)=\prod_{i=1}^r(z-b_i)$.  Invariant quadratic differentials descend as
\[
 H^0(C,\omega_C^2)^G
  =\pi^*H^0\bigl(\mathbf P^1_F,\omega_{\mathbf P^1}^2(B)\bigr),
 \qquad
 \frac{z^j(dz)^2}{Q(z)},\quad0\le j\le r-4,                 \tag{NR.8}
\]
where equality means the canonical pullback identification.  At a branch
point $z=u^e$, invariance of a holomorphic quadratic differential makes
its coefficient a multiple of $u^{e-2}$ modulo powers of $u^e$; descent
therefore allows exactly a simple pole.  At an unramified point it is
holomorphic.  At infinity the degree bound in (NR.8) is precisely the
condition of regularity.

If $\mu(a\otimes b)=f(dz)^2$, its descended coefficient can also be
computed as $|G|^{-1}\operatorname{Tr}_{F(C)/F(z)}(f)$.
The trace is an exact finite matrix trace.  Express these outputs in the
basis (NR.8) to obtain an explicit
\[
                       (r-3)\times n^2\text{ matrix }M
                       \text{ over }F.                      \tag{NR.9}
\]

We prove that its rank is the desired geometric rank.  For a smooth curve,
a Cech representative $\{\theta_{ij}\}$ of a first-order deformation
sends a holomorphic differential $\alpha$ to the class
$\{\theta_{ij}\mathbin{\lrcorner}\alpha\}\in H^1(C,\mathcal O_C)$.
Serre duality pairs this class with $\beta$ by pairing the deformation
with $\alpha\beta$.  This derives the multiplication codifferential of
the Jacobian period map.  For a $G$-invariant deformation only the group
average of a product contributes.  On $V\otimes V^*$, the group average
of a rank-one matrix unit is $m^{-1}$ times the identity, by Schur's
lemma.  Restriction to the primitive $K$-factor thus gives (NR.7), up to
a fixed nonzero scalar coming from the normalization of its polarization.
The cotangent space of the signature-$(n,n)$ unitary period domain is
$A\otimes A'$.  Polarization identifies the conjugate filtration with
the annihilator of the first; no extra independent period directions are
omitted in this description.

The tangent space of the Hurwitz component is
$H^1(C,T_C)^G$.  The local computation $z=u^e$ gives
$(\pi_*T_C)^G=T_{\mathbf P^1}(-B)$: an invariant vector field is
$u h(u^e)\partial_u$ and descends to $e z h(z)\partial_z$.
Taking cohomology and using exactness of finite-group invariants gives
\[
 H^1(C,T_C)^G=H^1(\mathbf P^1,T_{\mathbf P^1}(-B)),
 \qquad \dim=r-3.                                           \tag{NR.10}
\]
The branch-coordinate map is etale, so the family retained in D-NR-02
has this entire tangent space.  Duality identifies the dual of (NR.10)
with (NR.8).  Consequently (NR.9) is the full period codifferential on
the specified component, not a restriction to a conveniently chosen
subfamily.

Gaussian elimination in $F$ now outputs $q=\operatorname{rank}_F M$.
Field equality is decidable by reduction modulo $P$ and equality of
rational functions in the algebraically independent $t_i$.
One nonzero $q$-minor and a row-reduction identity with exactly $q$
nonzero rows are finite exact lower- and upper-rank certificates.
Base change to $\mathbf C$ preserves rank.  On a nonempty open of the
Hurwitz component this equals the complex period differential rank,
and hence the dimension of its image.  This proves a terminating exact
algorithm for every specified input tuple satisfying the hypotheses.

\subsection*{D-NR-04. The smaller trace formula for natural imprimitive groups}

For the natural representation of $G(d,e,m)$, use the actual intermediate
tower already constructed in D-IM-01:
\[
             Y=C/\ker\chi\ \xrightarrow{\ q\ }\quad
             X=C/H\ \xrightarrow{\ f\ }\ \mathbf P^1,
 \qquad V=\operatorname{Ind}_H^G\chi.                         \tag{NR.11}
\]
Choose bases $\alpha_1,\ldots,\alpha_n$ and
$\beta_1,\ldots,\beta_n$ of the two reciprocal eigenspaces of
$H^0(Y,\omega_Y)$.  Their products are invariant under the cyclic group,
so each descends to a meromorphic quadratic differential
$\gamma_{ab}$ on $X$, with at most simple poles at the cyclic branch
points.  Equivalently,
\[
 \gamma_{ab}=\frac1d\operatorname{Tr}^{(2)}_{Y/X}
                          (\alpha_a\beta_b).                \tag{NR.12}
\]
Here the superscript means trace of quadratic differentials; in a common
base coordinate it is the ordinary field trace of their coefficients.
Frobenius reciprocity identifies the coset basis of the induced module
with the conjugate pullbacks of these eigendifferentials.  The dual-basis
contraction in (NR.7) sums exactly over the $m$ cosets.  Therefore it is,
up to a fixed nonzero scalar,
\[
                     \operatorname{Tr}^{(2)}_{X/\mathbf P^1}
                                      (\gamma_{ab}).         \tag{NR.13}
\]
This also proves the formula directly at the level of variations and
their polarization pairings; it is not an equality inferred from their
dimensions.  Local trace at a ramification point, or (NR.8), shows that
(NR.13) has at most simple poles at the original branch divisor $B$.

Thus the exact rank of any specified natural-representation Nielsen
component is
\[
 \boxed{\dim_F\operatorname{span}\left\{
       \operatorname{Tr}^{(2)}_{X/\mathbf P^1}(\gamma_{ab}):
                       1\le a,b\le n\right\}.}              \tag{NR.14}
\]
The finite algebraic model search above makes this a computable formula,
including for arbitrary ramification of $X\to\mathbf P^1$.
It retains that map and imposes no submersion assumption when forgetting
it.  The universal bounds and the ranks $15$, $14$ and $1$ proved in
D-IM-03--08 remain distinct claims, not a universal replacement of
(NR.14) by one of the maxima.

For each fixed $m,d,e,n$, this also gives a finite procedure to tabulate
all answers.  Faithfulness and Chevalley--Weil give
$\sum_j\operatorname{rank}(1-g_j)=2m+2n$, so $r\le2m+2n$.
There are finitely many tuples in the finite group for these lengths.
Product one, generation, and signature are finite exact tests; run the
algorithm on each passing tuple.  Computing finite braid orbits can remove
duplicate work, but is not needed for termination.  This is not a claim
that the resulting table for every unbounded $m$ has been executed.

\paragraph{A finite check of the matrix step.}
For a cyclic cubic cover of $\mathbf P^1$ with six branch exponents one
and six branch exponents two, write
\[
 y^3=P(z)Q(z)^2,\qquad \deg P=\deg Q=6,
\]
with twelve distinct roots.  The two eigenspace bases are
$z^a dz/y$ and $z^bQ(z)dz/y^2$, $0\le a,b\le4$.
Their products are $z^{a+b}(dz)^2/(PQ)$, and the exponents $a+b$
span $0,\ldots,8$.  The exact matrix has rank nine.  The accompanying
small script checks all local valuations, the two dimensions five, and
an identity $9$-minor.  This illustrates and certifies the matrix step;
it is not an execution of the generic-model search for a large
$G(d,e,m)$, and it makes no additional claim about those components.

% Generated by build_tex_sections.py; full proof text, no abridgement.
% Requires amsmath, amssymb, and UTF-8 support.
\section{A rank-zero semiregular complex at a CM sixfold}

Run HODGE\_DIVISION; agent \texttt{/root/imprimitive}; 2026-11-21. Model GPT-6 Astra Pro, effort inherited/not exposed. \textbf{D-SR-01--06 passed independent whole-proof checks by \texttt{/root/sixfolds}; D-SR-01--05 also passed a separate whole-proof audit by \texttt{/root/g34\_eightfolds}.} See the peer reports and \texttt{CHECK\_EVENTS.tsv} for the exact scope of each event. No result below relies on a theorem of Markman. His secant construction suggested the approach. The classical trace identities, Fourier--Mukai equivalence and its trace naturality are separate cited inputs. The final deformation implication is exactly the supplied DEG\_L S-14/Perry theorem; its preprint status is retained. The later positive-rank construction SF-QC01--05 supplies an explicit numerical detector bound; the present section retains the independent rank-zero construction and its integral normalization. \cite{SR-Perry2026}

\subsection{D-SR-01. A finite construction of a disjoint cyclic orbit}
\label{D-SR-01}

Let \(C\) be a smooth nonhyperelliptic genus-three curve, \(J=J(C)\), and \(\Theta\) its principal polarization. Put \(D=C-C\) in \(J\). Then \(D\) is an irreducible divisor of class \(2\Theta\). Indeed the difference map \(C\times C\) to \(D\) is generically one-to-one: a nontrivial equality \(p-q=p^{\prime}-q^{\prime}\) would yield two different effective degree-two divisors in the same linear system, contradicting nonhyperellipticity. The Poincare formula and the Pontryagin product give \([C]*[-C]=2\Theta\). We use the usual Abel--Jacobi embedding; translating it does not alter \(D\). \cite{SR-BirkenhakeLange}

For every positive integer \(r\),

\[
\#\bigl(D\cap J[r]\bigr)\leq 108r^4.\tag{SR1.1}
\]

Here points are counted without multiplicity. To prove this, embed \(J\) by \(3\Theta\). Hyperplanes through the image of zero have base locus exactly \{0\}. Their pullbacks by [\(r\)] have finite base locus \(J[r]\). Choose two general such hyperplanes so that their pullbacks cut the surface \(D\) properly. Each point of \(D\) intersect \(J[r]\) contributes at least one to the resulting zero-dimensional intersection. Its degree is \cite{SR-BauerSzemberg}

\[
D\cdot([r]^*3\Theta)^2=(2\Theta)\cdot(3r^2\Theta)^2=18r^4\Theta^3=108r^4.
\]

The existence of proper successive cuts follows from Bertini avoidance of the finitely many positive-dimensional components; smoothness of the cuts is unnecessary.

Set

\[
P=2^{2048},\quad L=\operatorname{lcm}(1,2,\ldots,P),\quad N=4L^2+1,\quad d=N-1=4L^2.\tag{SR1.2}
\]

There is a cyclic subgroup \(H\) of \(J[N]\), of exact order \(N\), such that

\[
H\cap(C-C)=\{0\}.\tag{SR1.3}
\]

Here is a finite counting proof, with an explicit search specification. Choose \(v\) uniformly among the points of exact order \(N\) in \(J[N]=(\mathbb Z/N)^6\). For every divisor \(q> 1\) of \(N\), the point \((N/q)v\) is uniformly distributed among the exact-order-q points. Their number is

\[
J_6(q)=q^6\prod_{p\mid q}(1-p^{-6})>\frac{q^6}{2}.
\]

The inequality follows, for example, from \(\prod_i(1-a_i)\geq1-\sum_i a_i\) and \(\sum_{n\geq2}n^{-6}<1/2\). If the cyclic subgroup generated by \(v\) has a nonzero point on \(D\), one of its exact-order-q generators lies on \(D\) for some \(q\mid N\). The union bound and (SR1.1) therefore give

\[
\begin{aligned}\Pr\bigl(\langle v\rangle\cap(D\setminus\{0\})\ne\varnothing\bigr)&\leq216\sum_{\substack{q\mid N\\q>1}}\frac{\varphi(q)}{q^2}\\&\leq216\left[\exp\left(\sum_{p\mid N}\frac1p\right)-1\right].\end{aligned}\tag{SR1.4}
\]

The first estimate allows all \(\varphi(q)\) generators of the cyclic order-q subgroup and thus overcounts, harmlessly. The second uses

\[
\sum_{q\mid N}\frac{\varphi(q)}{q^2}=\prod_{p^a\parallel N}\left(1+\frac{1-p^{-a}}p\right)\leq\prod_{p\mid N}\left(1+\frac1p\right).
\]

Every prime divisor of \(N\) exceeds \(P\), because every prime \textless{}=\(P\) divides \(L\) and \(N\) is 1 modulo that prime. For \(P\) a power of two, the elementary central-binomial argument gives

\[
\log L<P\log4.
\]

For completeness, if \(\psi(x)=\log\operatorname{lcm}(1,\ldots,\lfloor x\rfloor)\), then \(\psi(2r)-\psi(r)\leq\log\binom{2r}{r}\leq2r\log2\): every prime power in (\(r\),2r] contributes its prime once to the binomial coefficient. Sum over the dyadic intervals ending at \(P\). Consequently

\[
\sum_{p\mid N}\frac1p\leq\frac{\log N}{(P+1)\log(P+1)}<\frac1{500}.
\]

The last bound follows from \(\log N\leq\log5+2\log L<\log5+4P\log2\) and \(\log P=2048\log2\). Thus (SR1.4) is strictly less than 216/499\textless{}1/2, using \(\exp(x)-1\leq x/(1-x)\) for 0\textless{}\(x\)\textless{}1. This proves existence.

To specify a finite construction, fix any algebraic enumeration of the finite group \(J[N]\), test its exact-order-N points in that order, and take the first \(v\) for which kv is outside the closed algebraic divisor \(D\) for every 1\textless{}=k\textless{}\(N\). All these tests are finite algebraic membership tests, and the proof guarantees termination. This is an explicit terminating construction, although its deliberately coarse integer \(N\) is extremely large. It is not a claim that we have enumerated \(J[N]\) in this run.

Let \(q:J\to Z=J/H\) be the quotient. By (SR1.3), the restriction \(q|_C\) is a closed embedding: it is injective and unramified, and the source is proper. Write \(\bar C=q(C)\). Its full inverse image is the disjoint union of the \(N\) translates \(C\)+\(h\).

Let \(\theta\) in \(\operatorname{NS}(Z)_{\mathbb Q}\) be the unique rational class with \(q^*\theta=\Theta\). It is ample. Uniqueness and existence follow from the rational cohomology isomorphism of an isogeny and the norm of a line bundle; explicitly \(q^*q_*\Theta=N\Theta\) because translations act trivially on cohomology. Put \(F=\mathcal I_{\bar C}\). Then

\[
\operatorname{ch}(F)=1-\frac N2\theta^2+\frac N3\theta^3,\qquad\Delta:=\int_Z\frac{\theta^3}{6}=\frac1N.\tag{SR1.5}
\]

The character follows by pulling back to \(J\): a smooth genus-three curve has \(\operatorname{ch}(\mathcal O_C)=[C]-2[\mathrm{pt}]\), since the degree of its Euler characteristic is -2 and \(\operatorname{td}(J)=1\). Pullback of the ideal is the ideal of the disjoint union, so its character is \(1-N\Theta^2/2+2N[\mathrm{pt}]\).

\textbf{Extension.} For any fixed positive integer \(m\leq 163\) the same proof applies to \(N=4mL^2+1\) and \(d=4mL^2\), with the same \(P\),\(L\). Replace log5 by log(4m+1) in the final estimate; it still gives \(\sum_{p\mid N}1/p<1/500\). The resulting field below is \(\mathbb Q(\sqrt{-m})\). This extension concerns a CM sixfold isogenous to the Fermat Jacobian square; for \(m\ne 1\) it is not being identified with a power of an elliptic curve with CM by \(\mathbb Q(\sqrt{-m})\).

\subsection{D-SR-02. The descended ideal and its external square are semiregular}
\label{D-SR-02}

In the situation above, \(F=\mathcal I_{\bar C}\) on the abelian threefold \(Z\) has

\[
\bigl(\dim\operatorname{Ext}^j(F,F)\bigr)_{j=0}^3=(1,6,6,1).\tag{SR2.1}
\]

Its trace maps \(\sigma_F^j:\operatorname{Ext}^j(F,F)\to\operatorname{HH}_{-j}(Z)\) are injective for \(j= 0\),1,2. The same statements hold for the derived dual \(F^\vee\). Consequently \(F\boxtimes F^\vee\) on \(Z\times Z\) is a gluable fully semiregular perfect complex, with \(\operatorname{Ext}^2\) dimension 48.

\textbf{Proof of (SR2.1).} Since \(q\) is etale and \(q|_C\) is an isomorphism onto \(\bar C\), the normal bundles \(N_{C/J}\) and \(N_{\bar C/Z}\) are identified. The normal sequence

\[
0\longrightarrow T_C\longrightarrow\mathcal O_C^3\longrightarrow N_{C/J}\longrightarrow0.
\]

shows \(h^0(N)=3\). Indeed the next map \(H^1(T_C)\to H^1(\mathcal O_C)^3\) is injective: its dual is the multiplication \(H^0(K_C)\otimes H^0(K_C)\to H^0(K_C^2)\), which is surjective for a smooth plane quartic. In this genus-three case this last statement is elementary: \(K_C=\mathcal O_C(1)\), and restriction of plane quadrics is an isomorphism from the six-dimensional space of quadrics onto \(H^0(C,\mathcal O_C(2))\); no nonzero quadric contains \(C\). Also deg \(N= 4\) and rank \(N= 2\), so Riemann--Roch gives \(h^1(N)=3\).

Here is a direct self-Ext calculation, avoiding any global identification of moduli spaces. Apply Hom(\(F\),-) to \(0\to F\to\mathcal O_Z\to\mathcal O_{\bar C}\to0\). The map Hom(\(F\),O\_Z)-\textgreater{}\(\operatorname{Hom}(F,\mathcal O_{\bar C})\) is zero, since every global homomorphism \(F\to\mathcal O_Z\) is a scalar multiple of the ideal inclusion. Thus there is an exact sequence

\[
0\longrightarrow H^0(N_{\bar C/Z})\longrightarrow\operatorname{Ext}^1(F,F)\longrightarrow\operatorname{Ext}^1(F,\mathcal O_Z).
\]

By threefold Serre duality, \(\operatorname{Ext}^1(F,\mathcal O_Z)=H^2(F)^*\). The restriction \(H^1(\mathcal O_Z)\to H^1(\mathcal O_{\bar C})\) is an isomorphism: it is the Abel--Jacobi isomorphism after composition with the isogeny \(q\). The cohomology sequence of the ideal therefore gives \(H^2(F)=H^2(\mathcal O_Z)\), of dimension three. Thus dim \(\operatorname{Ext}^1(F,F)\leq3+3=6\). The three embedded translation directions and the three degree-zero tensor directions are independent. Indeed the embedded ideal deformations have trivial determinant, whereas the trace of a tensor direction is the original class in \(H^1(\mathcal O_Z)\), because \(\operatorname{rank}(F)=1\). Hence equality holds and

\[
\operatorname{Ext}^1(F,F)=H^1(\mathcal O_Z)\oplus H^0(N_{\bar C/Z}).\tag{SR2.2}
\]

These assertions can also be checked on the two-generator Koszul resolution of the ideal; no stability or smoothness of a moduli space of sheaves is assumed.

\(\operatorname{End}(F)=\mathbb C\): restriction outside \(\bar C\) identifies an endomorphism with a regular function, which extends across codimension two on the normal projective variety \(Z\). Serre duality, with \(\omega_Z=\mathcal O_Z\), then yields (SR2.1).

Write \(\sigma_F^j(x)=\operatorname{tr}(\exp(\operatorname{at}(F))x)\) under HKR; the Todd correction is one for an abelian variety. For \(j= 0\) the trace of the identity is \(\operatorname{ch}(F)\), whose rank component is 1. For \(j= 1\), the infinitesimal action of \(H^1(\mathcal O_Z)\) by tensoring and of \(H^0(T_Z)\) by translation gives an isomorphism \cite{SR-BuchweitzFlenner2003,SR-BuchweitzFlenner2008}

\[
H^1(\mathcal O_Z)\oplus H^0(T_Z)\longrightarrow\operatorname{Ext}^1(F,F).
\]

The translation summand is an isomorphism onto \(H^0(N)\) by the normal sequence just used, and the tensor summand maps isomorphically to the determinant direction. Composition with \(\sigma_F^1\) sends the first summand to \(H^1(\mathcal O_Z)\) by the identity; the \(H^0(T_Z)\) summand has zero component there because \(c_1(F)=0\) and maps injectively into \(H^2(\Omega_Z^1)\) by contraction with \(-(N/2)\theta^2\). Nondegeneracy of \(\theta\) proves that contraction is injective. Thus \(\sigma_F^1\) is injective.

For \(j= 2\), apply the trace/Atiyah obstruction identity to the subspace

\[
H^2(\mathcal O_Z)\oplus H^1(T_Z)\subset\operatorname{HH}^2(Z).
\]

The \(H^2(\mathcal O_Z)\) summand maps by the rank-one identity to the \(H^2(\mathcal O_Z)\) component of \(\operatorname{HH}_{-2}\). The \(H^1(T_Z)\) summand has zero component there because \(c_1(F)=0\) and maps surjectively onto \(H^3(\Omega_Z^1)\) by contraction with \(-(N/2)\theta^2\). Here is a direct rank check. In holomorphic/antiholomorphic bases \(\theta=\sum_i\alpha_i\beta_i\), contraction by a Kodaira--Spencer matrix \(\kappa\) gives \(\kappa\) contraction \(\theta^2\)=2(\(\kappa\) contraction \(\theta\))\(\theta\). The map kappa-\textgreater{}\(\kappa\) contraction \(\theta\) is onto \(\bigwedge^2\langle\beta_i\rangle\), and multiplication by \(\theta\) is an isomorphism of that three-dimensional space onto \(\langle\alpha_i\rangle\otimes\bigwedge^3\langle\beta_i\rangle\). Hence the rank is three. Therefore the full composite has rank six. Its middle space \(\operatorname{Ext}^2(F,F)\) has dimension six, so \(\sigma_F^2\) is injective. Only the classical trace/Atiyah identity is used, not a conjectural deformation assertion.

Derived duality identifies the self-Ext groups of \(F\) and \(F^\vee\); the same rank argument applies to \(\operatorname{ch}(F^\vee)=1-(N/2)\theta^2-(N/3)\theta^3\). In particular \(\sigma_{F^\vee}^j\) is injective for \(j= 0\),1,2.

Finally the trace of an external tensor product is the external product of the traces, because its Atiyah class is \(\operatorname{at}(F)\) external\_tensor id + id external\_tensor \(\operatorname{at}(F^\vee)\), and the exponential and trace factor. The \(\operatorname{Ext}^2\) Kunneth decomposition is

\[
\begin{aligned}&\operatorname{Ext}^2(F,F)\otimes\operatorname{Ext}^0(F^\vee,F^\vee)\\{}\oplus{}&\operatorname{Ext}^1(F,F)\otimes\operatorname{Ext}^1(F^\vee,F^\vee)\\{}\oplus{}&\operatorname{Ext}^0(F,F)\otimes\operatorname{Ext}^2(F^\vee,F^\vee).\end{aligned}
\]

The dimensions are 6,36,6. Their trace images lie in the three distinct Hochschild Kunneth components \(\operatorname{HH}_{-2}\otimes\operatorname{HH}_0\), \(\operatorname{HH}_{-1}\otimes\operatorname{HH}_{-1}\), and \(\operatorname{HH}_0\otimes\operatorname{HH}_{-2}\). Each restriction is injective by the preceding checks, so there is no cancellation between the three pieces. This proves semiregularity and the dimension 48. Negative self-Ext groups vanish by the same Kunneth decomposition, since they vanish for a coherent sheaf and for its derived dual. This proves gluability.

\subsection{D-SR-03. An exact Fourier--Mukai formula with a nonzero Weil term}
\label{D-SR-03}

Let \(A=Z\times\widehat Z\). Let Poincare be the normalized Poincare bundle and \(\gamma=c_1(\mathcal P)\) on A. Let \(\widehat\theta\) be the rational dual polarization satisfying \(\phi_\theta^*\widehat\theta=\theta\), where \(\phi_\theta:Z\to\widehat Z\) is the rational polarization quasi-isogeny. Define the positive rational divisor class

\[
H=N\theta_x+\widehat\theta+\gamma.\tag{SR3.1}
\]

Let \(\mu:Z\times Z\to Z\times Z\) be \(\mu(x,y)=(x+y,y)\). Define the exact equivalence

\[
\Phi=(\operatorname{id}_Z\boxtimes\operatorname{FM}_{\mathcal P^{-1}})\circ\mu^*,\qquad E=\Phi(F\boxtimes F^\vee).\tag{SR3.2}
\]

No extra cohomological shift is included in this convention. Then \(E\) is gluable and fully semiregular, and its complete Chern character satisfies

\[
\exp(-\gamma/2)\operatorname{ch}(E)=t+2H-\frac{H^3}{12d}+\frac{H^5}{960d^2}.\tag{SR3.3}
\]

where \(t\) is a nonzero rational Weil class for a \(K=\mathbb Q(\sqrt{-d})\) action of signature (3,3) on (A,\(H\)). Thus rank(\(E\))=0 causes no problem for the exact S-14 hypothesis.

\textbf{Exact calculation.} Put \(u=\sqrt{-d}\), \(s\)=-1+\(u\), \(\bar s\)=-1-u, and

\[
a=\frac{1+1/u}{2},\qquad\bar a=\frac{1-1/u}{2}.
\]

Equation (SR1.5) is exactly

\[
\operatorname{ch}(F)=a\exp(s\theta)+\bar a\exp(\bar s\theta).\tag{SR3.4}
\]

It suffices to expand through \(\theta^3\); use \(u^2=-d\) and \(N=d+1\). The coefficient of \(\theta\) is zero, that of \(\theta^2\) is -N/2, and that of \(\theta^3\) is \(N\)/3.

Choose rational symplectic Darboux coordinates for \(\theta\). In one of its three planes denote the first-factor coordinates by \(a_i\),\(b_i\), the integrated second-factor coordinates by \(c_i\),\(d_i\), and the dual output coordinates by \(u_i\),\(v_i\). Then \(\theta_x=\sum_i a_i b_i\), \(\theta_y=\sum_i c_i d_i\), \(\widehat\theta=\sum_i u_i v_i\), and \(\gamma=\sum_i(a_i v_i-b_i u_i)\). Juxtaposition means exterior multiplication. The Poincare kernel on the integrated factors has class \(\sum_i(c_i v_i-d_i u_i)\). Direct extraction of the \(c_i\) \(d_i\) coefficient gives

\[
\begin{aligned}&\int_{c_i,d_i}\exp\bigl(s(a_i+c_i)(b_i+d_i)-t c_i d_i-(c_i v_i-d_i u_i)\bigr)\\&\quad=(s-t)-st a_i b_i-s b_i u_i-u_i v_i+s a_i v_i-s a_i b_i u_i v_i.\end{aligned}\tag{SR3.5}
\]

The integration over all three planes has the additional scalar \(\Delta\)=integral\_Z \(\theta^3/6=1/N\). This scalar is essential because \(\theta\) is rational and need not be principal on \(Z\).

When \(t\)=\(s\), (SR3.5) equals

\[
-\exp\bigl((a_i v_i-b_i u_i)/2\bigr)(s a_i-u_i)(s b_i-v_i).
\]

Consequently, putting \(Q_s=\prod_{i=1}^3(s a_i-u_i)(s b_i-v_i)\), the two equal-root terms in the normalized transform are

\[
t=-\Delta\bigl(a^2Q_s+\bar a^2Q_{\bar s}\bigr).\tag{SR3.6}
\]

When \(t\)=\(\bar s\), set \(\alpha=s-\bar s=2u\). Since \(st=N\) and \(s+t=-2\), (SR3.5) is

\[
\alpha\exp(\gamma_i/2-H_i/\alpha),\qquad H_i=N a_i b_i+u_i v_i+a_i v_i-b_i u_i.
\]

The unequal-root pair therefore contributes after normalization

\[
\Delta a\bar a\alpha^3\bigl[\exp(-H/\alpha)-\exp(H/\alpha)\bigr].
\]

Use \(a\bar a=N/(4d)\), \(\alpha^2=-4d\) and \(H^7=0\). Expanding the odd powers gives precisely the polynomial in (SR3.3). Equations (SR3.4)--(SR3.5), or a sparse exterior-algebra expansion of all twelve output generators, are exact identities; a separate script checks the displayed formulas and records the coefficients.

\textbf{The Weil structure and positivity.} The rational endomorphism \(T\) of A is defined, after the quasi-isogeny identification A\textasciitilde{}\(Z\times Z\) supplied by \(\theta\), by the matrix

\[
T=\begin{pmatrix}-2&-1\\N&0\end{pmatrix}.
\]

It satisfies \(T^2+2T+N=0\). Thus \(T+1\) defines \(\sqrt{-d}\). On \(H^1\), its s-eigenspace is exactly the six-dimensional span of the forms \(s\) a\_i-u\_i, \(s\) b\_i-v\_i. Its conjugate eigenspace is the analogous \(\bar s\) span. Since the two-by-two matrix acts by scalars on each copy of \(H^1(Z)\), each eigenspace has holomorphic dimension three and antiholomorphic dimension three. Thus \(Q_s\) and \(Q_{\bar s}\) span the two determinant lines of the rational Weil plane.

The class \(H\) is ample. Under id times \(\phi_\theta\) its pullback is

\[
N\theta_1+\theta_2+(m^*\theta-\theta_1-\theta_2)=(N-1)\theta_1+m^*\theta,
\]

which is ample because (\(x\),\(y\))-\textgreater{}(\(x\),\(x\)+\(y\)) is an isomorphism and \(N> 1\). The generator \(T+1\) is Rosati anti-self-adjoint for \(H\): in these coordinates the polarization matrix is \(S=\bigl(\begin{smallmatrix}N&1\\1&1\end{smallmatrix}\bigr)\), and \(S(T+1)=\bigl(\begin{smallmatrix}0&-d\\d&0\end{smallmatrix}\bigr)\) is skew-symmetric. Hence the \(K\) action is compatible with \(H\). The coefficients a and \(\bar a\) are nonzero, and \(Q_s\),\(Q_{\bar s}\) are independent nonzero determinant vectors, so (SR3.6) is nonzero. Galois conjugation interchanges its two summands, proving rationality. This proves the entire character identity and its nonzero Weil component.

Fourier--Mukai equivalence preserves negative self-Ext and injectivity of the categorical trace, so D-SR-02 proves the stated gluability and semiregularity of \(E\). This preservation is the functoriality in Perry Remarks 2.3 and 2.5, and follows directly from transporting the evaluation/coevaluation defining the categorical trace across an equivalence. The result is not inferred from equality of K-classes alone. \cite{SR-Perry2026}

\subsection{D-SR-04. Transport to the requested balanced CM power}
\label{D-SR-04}

Choose \(C\) to be the Fermat quartic \(x^4+y^4+z^4=0\). Its Jacobian is isogenous to \(E_i^3\), where \(E_i=\mathbb C/(\mathbb Z+i\mathbb Z)\): Beauville, *Some surfaces with maximal Picard number*, JEP 1 (2014), section 4, Example 3, \(p\).107, with Lemma 2 on \(p\).106. Thus \(A=Z\times\widehat Z\) is isogenous to \(E_i^6\). \cite{SR-Beauville2014}

The isogeny to the CM power and its polarization can be specified explicitly. The three degree-two quotients of the Fermat quartic by sign changes of its coordinates are elliptic curves with j=1728. Choose origin-preserving isomorphisms of them with \(E_i\), and let

\[
j:E_i^3\longrightarrow J.
\]

be the sum of the three pullback homomorphisms of Jacobians. Their cotangent images are the three distinct nontrivial characters of the sign-change Klein four group, so j is an isogeny. Each diagonal norm-pullback composition is multiplication by two; cross compositions vanish, since they intertwine different sign characters. Therefore

\[
j^*\Theta=2\theta_{E_i^3},\qquad\deg j=8.
\]

Put \(h=qj:E_i^3\to Z\), \(r=2L\), \(d=r^2\), \(s=-1+ri\), and define the two-by-two matrices acting on threefold pairs by

\[
R=\begin{pmatrix}1&1\\\bar s&s\end{pmatrix},\quad k=\operatorname{diag}(h,\phi_\theta h),\quad f=NkR:E_i^6\longrightarrow A.\tag{SR4.1}
\]

The map k is initially a rational homomorphism, but \(f\) is an actual isogeny: Ntheta is the integral norm-polarization class \(q_*\Theta\), so \(N\phi_\theta=\phi_{N\theta}\) is an actual homomorphism. All entries of \(R\) are Gaussian integers. Its two columns are eigenvectors of \(T+1\) with eigenvalues ri and -ri. Thus \(f\) is K-linear for the specified balanced source action \(\operatorname{diag}(i,i,i,-i,-i,-i)\). If \(\theta_{\rm prod}\) is the product principal polarization on \(E_i^6\), direct matrix multiplication gives

\[
R^\dagger\begin{pmatrix}N&1\\1&1\end{pmatrix}R=2dI,\qquad f^*H=4dN^2\theta_{\rm prod},\qquad\deg f=4096d^3N^{12}.\tag{SR4.2}
\]

For the degree calculation, \(\deg h=8N\), \(\deg\phi_\theta=\Delta^2=N^{-2}\) as a rational isogeny, so \(\deg k=64\); \(\deg R=\operatorname{Norm}(2ri)^3=(4d)^3\); multiplication by \(N\) on a sixfold has degree \(N^{12}\). Multiplying gives the stated integer.

Let \(G=\ker f\), acting honestly by translations, and \(E_0=f^*E\) with its canonical coherent G-linearization. Etale descent identifies

\[
\operatorname{Ext}^2(E_0,E_0)^G=\operatorname{Ext}^2(E,E).
\]

Under HKR, the full trace of a pullback along an etale morphism is the pullback of the trace. The map \(f^*:\operatorname{HH}_{-2}(A)\to\operatorname{HH}_{-2}(E_i^6)\) is injective; indeed each Hodge-cohomology pullback is an isomorphism over Q or \(C\) for an isogeny, with inverse \(f_*/\deg f\). Since \(E\) is fully semiregular, the trace of \(E_0\) is injective on the invariant \(\operatorname{Ext}^2\). Hence \(E_0\) is G-semiregular, and therefore weakly G-semiregular.

We check ordinary negative self-Ext, rather than only its G-invariant part. The finite etale isogeny has

\[
f_*\mathcal O_{E_i^6}=\bigoplus_{\alpha\in\ker\widehat f}\alpha,
\]

where the summands are torsion degree-zero line bundles on A. Adjunction and the projection formula yield

\[
\operatorname{Ext}^j(E_0,E_0)=\bigoplus_\alpha\operatorname{Ext}^j(E,E\otimes\alpha).
\]

Conjugating tensorization by \(\alpha\) through \(\Phi\) gives an autoequivalence of \(Z\times Z\) consisting of a translation followed by tensorization with a degree-zero line bundle. This can be seen directly from (SR3.2): Fourier--Mukai interchanges a degree-zero tensor on the dual factor with a translation on the original factor, and \(\mu\) conjugates translations to translations and pulls degree-zero line bundles back to degree-zero line bundles. Every such operation on \(Z\times Z\) is a product operation, since \(\operatorname{Pic}^0(Z\times Z)=\operatorname{Pic}^0(Z)\times\operatorname{Pic}^0(Z)\). Thus, up to an irrelevant one-dimensional scalar factor, \cite{SR-Mukai1981}

\[
\Phi^{-1}(E\otimes\alpha)=(t_x^*F\otimes L_1)\boxtimes(t_y^*F^\vee\otimes L_2).
\]

The self-Ext comparison with \(F\boxtimes F^\vee\) has no negative degrees: the first factor is Ext between coherent sheaves; by derived duality the second is also Ext between coherent sheaves in the opposite order. Their Kunneth tensor product therefore has no negative cohomology. This proves \(\operatorname{Ext}^j(E_0,E_0)=0\) for every \(j< 0\), closing the gluability requirement.

The character identity already pulls back exactly: set \(B_0=-f^*\gamma/2\), \(t_0=f^*t\), and \(H_0=f^*H\). Then

\[
\exp(B_0)\operatorname{ch}(E_0)=t_0+2H_0-\frac{H_0^3}{12d}+\frac{H_0^5}{960d^2}.
\]

By (SR4.2), \(H_0=4dN^2\theta_{\rm prod}\). We can take the fixed very ample polarization \(\theta_0=3\theta_{\rm prod}\) and rewrite the right side as \(t_0+P(\theta_0)\), with the explicitly displayed rational coefficients after substituting \(H_0=(4dN^2/3)\theta_0\). Thus this is a seed on the requested balanced CM power with its product principal polarization, not merely on an unspecified isogenous polarized variety. All of the S-14 object hypotheses have now been proved and independently checked.

\subsection{D-SR-05. Split determinant and the other imaginary quadratic fields}
\label{D-SR-05}

The K-Hermitian space in D-SR-03 is split, of rank six and Witt index three; in the raw determinant convention its determinant class is \([-1]\). This does not follow merely from its real signature. Here is a direct rational proof. Let \(E_\theta\) be the rational alternating polarization form on \(V=H_1(Z,\mathbb Q)\), and write \(J\)=\(T+1\) and \(u=\sqrt{-d}\). Define \(q_H(x,y)=E_H(x,Jy)\) and \(h(x,y)=q_H(x,y)+uE_H(x,y)\). Since \(J\) is skew-adjoint and \(J^2=-d\), \(q_H\) is symmetric, \(h\) is Hermitian, \(h(x,Jy)=u h(x,y)\), and \(h(Jx,y)=-u h(x,y)\). Thus \(h\) is conjugate-linear in the first argument and linear in the second, with \(h(v,v)=q_H(v,v)\).

Choose any rational E\_theta-Lagrangian \(L\) in \(V\), of dimension three. Under the rational identification \(H_1(A,\mathbb Q)=V\oplus V\), put \(I=L\oplus L\). The scalar two-by-two matrix \(J\) preserves I, so I has K-dimension three. The polarization matrix is \(S\otimes E_\theta\); hence \(E_H\) vanishes on I. Because I is J-stable, \(q_H\) also vanishes on I, and therefore \(h\) vanishes on I. This is a half-dimensional totally isotropic K-subspace. A nondegenerate Hermitian space with such a subspace is a sum of three hyperbolic planes: choose a dual basis for I, and subtract half of its Hermitian self-pairing matrix to make the dual subspace isotropic. Each hyperbolic plane has determinant -1 modulo K-norms. Thus \(h\) is split and has determinant \((-1)^3\) modulo K-norms. This argument concerns the actual polarization and \(K\) action and works independently of the additional CM endomorphisms.

For any positive squarefree integer \(m\), choose

\[
\begin{gathered}P_m=2^{2048+\lceil\log_2(m+1)\rceil},\qquad L_m=\operatorname{lcm}(1,\ldots,P_m),\\N_m=4mL_m^2+1,\qquad d_m=4mL_m^2.\end{gathered}
\]

The proof of D-SR-01 applies verbatim. Indeed every prime of \(N_m\) exceeds \(P_m\), log \(L_m<P_m\) log4, and

\[
\sum_{p\mid N_m}\frac1p<\frac{\log(4m+1)+4P_m\log2}{P_m\log P_m}<\frac1{500}.
\]

For the last inequality use \(m<P_m\), log \(P_m\geq2048\log2\) and \(P_m\geq2^{2048}\); the two lower-order logarithmic terms are less than \(2/P_m\). Thus the same strict union bound gives the required cyclic subgroup. The construction D-SR-01--03 yields a gluable fully semiregular complex on a CM sixfold isogenous to \(E_i^6\), with the complete character identity and Weil field \(K\)=\(\mathbb Q(\sqrt{-m})\). For the nine class-number-one fields the smaller common \(P=2^{2048}\) already suffices, as stated in D-SR-01. The underlying elliptic CM field is \(\mathbb Q(i)\); the newly constructed Weil action is \(K\). Those two actions must not be confused.

Apply supplied DEG\_L S-14 with \(G\) trivial on A and with a sufficiently divisible integral multiple of \(H\) as the relatively very ample polarization. On a finite level cover of the corresponding split Weil family, the rational Weil plane is a constant local system; take the flat section whose value at this CM point is \(t\). Formula (SR3.3) is its complete Hodge-compatible character. S-14 therefore gives bounded division for this nonzero Weil class throughout that fixed family. The supplied Theorem W then gives the corresponding detector/whole-Weil-plane conclusion. This statement uses Perry's deformation theorem as specified in S-14; it has no Markman, Severn, or 032 theorem as a premise. \cite{SR-Perry2026}

The construction settles the existence of the deformation object, including at the standard balanced Gaussian CM power in D-SR-04. It supplies an explicit finite search bound for the cyclic subgroup, an exact Chern character, an actual group action, and an exact CM isogeny degree. D-SR-06 additionally supplies an ordinary, untwisted deformation and the explicit class denominator and norm. This rank-zero argument by itself leaves the detector degree existential. The later positive-rank construction SF-QC01--05 supplies the missing effective presentation and a numerical degree bound; it has a separate full proof and audit. A positive definite form or a divisor-generated CM K-class alone does not supply that presentation.

\subsection{D-SR-06. An integral twist, ordinary deformation, and exact class normalization}
\label{D-SR-06}

Retain the Gaussian CM isogeny \(f\) of D-SR-04 and put \(c=4dN^2\), so \(H_0=c\theta_{\rm prod}\). The class \(B_0\) is integral. To check this with the correct normalization, let \(C_2=\bigl(\begin{smallmatrix}0&1\\1&0\end{smallmatrix}\bigr)\). The pullback of \(\gamma\) by \(k=\operatorname{diag}(h,\phi_\theta h)\) has matrix \(2C_2\) relative to the product principal polarization, because \(h^*\theta=2\theta_{E_i^3}\). Consequently

\[
\begin{gathered}R^\dagger C_2R=2\begin{pmatrix}-1&s\\\bar s&-1\end{pmatrix},\quad f^*\gamma=4N^2\begin{pmatrix}-1&s\\\bar s&-1\end{pmatrix}\otimes I_3,\\B_0=2N^2\begin{pmatrix}1&-s\\-\bar s&1\end{pmatrix}\otimes I_3.\end{gathered}\tag{SR6.1}
\]

The entries of the last Hermitian matrix are Gaussian integers and its diagonal entries are ordinary integers. Such a matrix represents an integral divisor class on the principally polarized Gaussian product. Indeed the map from divisor classes to symmetric endomorphisms, composed with the inverse principal polarization, identifies these classes with the integral Hermitian matrices; equivalently, its alternating Riemann form is integral on the product Gaussian lattice. By the integral (1,1) theorem it is the first Chern class of a line bundle \(L_B\). No ampleness of \(L_B\) is needed.

Put \(E_0^{\prime}=E_0\otimes L_B\). Conjugate the honest G-action by tensorization with \(L_B\). For a translation \(t_g\) this replaces the action by \(t_g^*(-)\) tensor (\(L_B\) tensor \(t_g^*L_B^{-1}\)). The additional line bundle is degree zero, and the assignment is a homomorphism into the translation/tensor group; its image is finite because \(G\) is finite. The coherent composition isomorphisms are transported by the tensor equivalence. The equivariant structure, weak semiregularity, and ordinary negative-self-Ext vanishing are therefore preserved. Now the entire ordinary character is

\[
\operatorname{ch}(E_0^{\prime})=t_0+2c\theta_{\rm prod}-\frac{16}{3}d^2N^6\theta_{\rm prod}^3+\frac{16}{15}d^3N^{10}\theta_{\rm prod}^5.\tag{SR6.2}
\]

All three polynomial coefficients are integers: \(L\) is divisible by 3 and 5, and \(d=4L^2\). This is a genuine line-bundle twist of the complex; it is not a formal normalization of a rank-zero character.

In Perry's proof of Theorem 1.2 in section 6.2, pp.26--28, one can now choose \(B_0^{\prime}=0\), hence the initial Kummer class and the associated Brauer class are zero. The band pushout used to extend the group action sends the zero gerbe class to zero. The ambient category in Theorem 6.3 is consequently Perf of the ordinary abelian family. Thus the resulting object on an etale neighborhood is an ordinary perfect complex, even if a separate presentation of the invariant category involves a gerbe. This conclusion uses the actual proof of the cited preprint, rather than strengthening the word "twisted" in its introductory statement without justification. \cite{SR-Perry2026}

The class denominator can be sharpened from the safe factorial 6 to 2. The character (SR6.2) has \(\operatorname{ch}_1=2H_0\) and \(\operatorname{ch}_2=0\), so the integral Chern classes of the perfect complex satisfy \(c_1=2H_0\) and \(c_2=2H_0^2\). Newton's formula gives \(\operatorname{ch}_3=-(2/3)H_0^3+c_3/2\). Comparing with (SR3.3) yields

\[
\begin{gathered}t_0=\frac{c_3(E_0^{\prime})}{2}+\left[\frac{16}{3}d^2N^6(1-8d)\right]\theta_{\rm prod}^3,\\\delta=2,\qquad T_0:=-\int_{E_i^6}t_0^2=32768d^4N^{12}.\end{gathered}\tag{SR6.3}
\]

The bracket is an integer because 3 divides \(L\). Hence \(2t_0\) is integral. The same calculation expresses 2t as an integral algebraic cycle on the ordinary-deformation open. The degree bound of its positive and negative parts can be chosen uniformly on a nonempty finite-type open by representing the relative Chern classes and the polarization power by cycles there. The proper relative Chow-space argument of supplied DIV1 then extends this fixed multiplier 2, with some common finite degree bound, to every fiber of the fixed family. This argument establishes a uniform multiplier; it does not calculate the degree bound.

For completeness, the exact norm in (SR6.3) is obtained before the isogeny. The two decomposable determinant vectors square to zero. In the Darboux orientation of D-SR-03 their product has top coefficient \((s-\bar s)^6\). The product of the \(\theta\) and dual-theta symplectic volumes is one. Therefore

\[
\begin{gathered}\int_A t^2=2\Delta^2(a\bar a)^2(s-\bar s)^6=-8d,\\-\int_{E_i^6}(f^*t)^2=8d\deg f=32768d^4N^{12}.\end{gathered}\tag{SR6.4}
\]

Both the twelve-generator wedge calculation and the two-by-two B-field and polarization matrices are checked by \texttt{semiregular\_gaussian\_certificate.py}, with their exact outputs in its JSON file. These arithmetic checks do not claim that the large torsion subgroup was enumerated.

There is also a fully integral normalization for every field in D-SR-05 without diagonalizing its action over \(\mathbb Q(i)\). Use the same Fermat map \(h\) and rational k, put \(M=2dN\), and form \(f_m=Mk:E_i^6\to A\). It is an actual isogeny because \(M\) is divisible by \(N\). Let \(\Theta_S\) be the integral ample class on \(E_i^6\) with Hermitian matrix \(S=\bigl(\begin{smallmatrix}N&1\\1&1\end{smallmatrix}\bigr)\) tensor \(I_3\). Then

\[
\deg f_m=64M^{12},\qquad f_m^*H=2M^2\Theta_S,\qquad B_m=-\tfrac12f_m^*\gamma=-M^2C_2\otimes I_3.\tag{SR6.5}
\]

The matrix \(B_m\) is integral. The K-action on the source is the integral scalar-block matrix \(T+1\), and its Hermitian space is split by D-SR-05. Twist \(f_m^*E\) by the line bundle of class \(B_m\), and conjugate its deck action exactly as above. The negative-self-Ext proof of D-SR-04 applies to this isogeny as well. The full ordinary character is \(t_m+2H_m-H_m^3/(12d)+H_m^5/(960d^2)\), where \(t_m=f_m^*t\) and \(H_m=2M^2\Theta_S\). The same Chern-class computation gives

\[
\begin{gathered}t_m=\frac{c_3}{2}+\left[\frac{128}{3}d^5N^6(1-8d)\right]\Theta_S^3,\\\delta_m=2,\qquad T_m:=-\int_{E_i^6}t_m^2=512dM^{12}=2^{21}d^{13}N^{12}.\end{gathered}\tag{SR6.6}
\]

The bracket is an integer since 3 divides \(d\). This gives an ordinary CM seed and explicit class normalization for every imaginary quadratic \(K\). Its underlying elliptic factor still has CM by \(\mathbb Q(i)\), while its Weil action is \(K\). The balanced standard \(E_K^6\) assertion is made only for the Gaussian case of D-SR-04. The polarization \(\Theta_S\) and the product principal polarization in (SR6.3) are kept distinct.

\subsection{Quantitative scope and framework}

The independent review of D-SR-01--04 checked the full trace tensor-product argument, Fourier signs, rational volume, and ordinary negative self-Ext. Exact source locators are recorded in \texttt{SOURCES.md}: Buchweitz--Flenner2003 Proposition4.2 pp166--167; Buchweitz--Flenner2008 author-version Proposition6.2.1 and Theorem6.4.2 pp29--31; Mukai1981 Theorem2.2 p156 and equation3.1 p158; Beauville2014 Lemma2 p106 and Example3 p107; Perry2026v2 Theorem1.2 and Remarks2.3/2.5, Definition2.6/Lemma2.7. Poincare and very ampleness are classical inputs explicitly identified in that log. \cite{SR-Perry2026}

Bounding the cycle presentation after S-14 is a stricter question than existence of the semiregular seed. The present rank-zero proof answers the existence question; the separately checked SF-QC01--05 positive-rank construction and HD-CM6 lattice transfer answer the numerical question for the split class-number-one families. In the project's language, the object is now in Perf rather than only a numerical lambda-ring class, and the Hochschild trace verifies that its actual obstruction space is controlled by the Hodge character. Its nonzero weight-six determinant component is of type (3,3); the remaining components are polarization powers. The supplied bounded-division reduction turns its deformation into the integral uniform form of Grothendieck's condition A (cote16, scan050, corrected SIGN\_NEXT SN-C03/C04). The positive ample class \(H\) supplies compatibility with condition I (scan049 and cote44 Theorem4), but the cycles come from the categorical deformation theorem, not from positivity alone. The resultant/Cayley--Hamilton division step remains the supplied Theorem W, corresponding to scans048--049. No new historical attribution of perfect complexes or Hochschild semiregularity to those archive pages is intended.

% Generated from the full note by format_quantitative_tex.py; no proof paragraphs omitted.
\section{Explicit regularity and an effective detector from the CM
construction}\label{explicit-regularity-and-an-effective-detector-from-the-cm-construction}

Author/checker: GPT-6 Astra Pro, inherited effort not exposed,
HODGE\_DIVISION /root/sixfolds. Date: 2026-11-21.

This note develops a positive-rank variant of D-SR-01--05. The
regularity argument was suggested by the supplied HECKE\_DIV M7
calculation, but that statement, Markman's output-sheaf theorem, and
Markman's generic-position hypotheses are not premises. All geometry and
regularity needed here are proved below. The deformation step uses the
authorized DEG\_L S-14 input, specifically Perry arXiv:2604.00511v2,
Theorem 1.2 and its proof. Perry remains a preprint in the source-status
convention of this project.

The finite construction below is explicit but has not been executed for
its astronomical torsion order. Its termination follows from the
displayed counting inequalities. Numerical certificates in the
accompanying script concern the symbolic character, matrices, regularity
inequalities, and final integer formula, not an enumeration of the
torsion group.

\subsection{SF-QC01. A strengthened torsion
orbit}\label{SF-QC01}

Let \(C\) be the Fermat plane quartic, \(J=J(C)\), and \(\Theta\) its
principal polarization. Let \(E_i\) be the Gaussian CM elliptic curve.
The three elliptic quotient maps give an isogeny \(j:E_i^3\to J\) with
\(j^*\Theta=2\theta_{\mathrm{prod}}\) and degree eight, as proved in
D-SR-04. Fix a positive squarefree integer \(m\). Put

\[
 P=2^{2048+\lceil\log_2(m+1)\rceil},\quad
 L=\operatorname{lcm}(1,\ldots,P),\quad
 N=4mL^2+1,\quad d=N-1=4mL^2.
\]

For the nine class-number-one radicands \(m\le 163\) one can instead use
the common \(P=2^{2048}\), as checked in D-SR-01. Every prime divisor of
\(N\) exceeds \(P\). In particular \(N\) is odd, and \(j\) induces an
isomorphism \(E_i^3[N]\to J[N]\). There exists a cyclic subgroup
\(H\subset J[N]\) of order \(N\) satisfying both:

\begin{enumerate}
\def\labelenumi{\arabic{enumi}.}
\item
  The translates \(C+h\), \(h\in H\), are pairwise disjoint.
\item
  For every \(P_{0}\in \operatorname{Pic}^{0}(J)\), \[
  H^0(J,I_{Z_H}(2\Theta)\otimes P_0)=0,
  \qquad Z_H=\coprod_{h\in H}(C+h).
  \tag{QC1.1}
  \]
\end{enumerate}

\textbf{Proof.} Sample uniformly from exact-order-N vectors
\(v'\in E_i^3[N]\), and let \(v=j(v')\). D-SR-01 proves that failure of
the first condition has probability less than 216/499. Its proof uses
the bound \(\#((C-C)\cap J[q])\le108q^4\), the lower bound
\(J_{6}(q)>q^{6}/2\) for exact-order-q points, and
\(\sum_{p\mid N}1/p<1/500\). These estimates apply unchanged to the
present m-dependent \(P\).

Identify \(\operatorname{NS}(E_i^{3})\) with integral Hermitian
three-by-three Gaussian matrices using the product principal
polarization. Consider all nonzero such matrices \(M\) with
\(0\le M\le 4I\). There are at most

\[
 B=5^3\,81^3=66\,430\,125
\]

matrices: each diagonal entry is one of 0,1,2,3,4; each off-diagonal
Gaussian entry has real and imaginary parts between \(-4\) and 4.
Positive semidefiniteness actually gives its norm at most 16, so every
nonzero entry is a unit modulo \(N\), because all prime divisors of
\(N\) exceed 16. A row containing such an entry defines a surjective map
\((Z[i]/N)^3\to Z[i]/N\). Consequently \(\operatorname{ker}(M)\)
contains at most \(N^{4}\) points of \(E_i^{3}[N]\), and the probability
that an exact-order-N \(v'\) lies in it is less than \(2/N^{2}\). The
probability of lying in any of these kernels is therefore at most
\(2B/N^{2}<1/4\). The combined failure probability is below
\(216/499+1/4<1\). Choose \(v'\) outside both bad sets, and set
\(H=\langle j(v')\rangle\).

Suppose that (QC1.1) fails. Then an effective divisor \(D\) numerically
equivalent to \(2\Theta\) contains every curve \(C+h\). Write its
distinct irreducible components as \(D_{1},\ldots ,D_k\). Since
\(\Theta ^{2}D=12\) and \(\Theta ^{2}D_i\) is a positive integer,
\(k\le 12\). For every \(h\in H\), both \(D\) and \(D+h\) contain all
the \(N\) curves. If these divisors had no common component, their
proper intersection would be an effective curve of \(\Theta\)-degree 24.
Each curve \(C+h\) has \(\Theta\)-degree three, so \(3N\le 24\), a
contradiction. Thus every \(h\) belongs to at least one set

\[
 S_{ij}=\{h\in H:D_i+h=D_j\}.
\]

Each nonempty \(S_{ij}\) is a coset of
\(H\cap \operatorname{Stab}(D_i)\). If all these subgroups were proper,
each coset would have at most \(N/p_{\min}\) points, where
\(p_{\min}>P\) is the smallest prime divisor of \(N\). The at most
\(k^{2}\le 144\) cosets could not cover \(H\). Therefore some \(D_i\) is
H-stable.

Effective divisors on an abelian variety are nef: for a curve, translate
the divisor to meet it properly and compute the nonnegative
intersection. Hence \(D_i\) and \(2\Theta -D_i\) are nef. The matrix
\(M\) corresponding to \(j^*D_i\) is a nonzero integral Gaussian
Hermitian matrix with \(0\le M\le 4I\). Since \(D_i\) is invariant under
\(v=j(v')\), its polarization homomorphism kills \(v\). Pullback gives
\(Mv'=0\), contrary to the choice of \(v'\). This proves (QC1.1).

The construction is finite: enumerate exact-order-N torsion vectors,
test avoidance of \(C-C\) and the finite displayed matrix kernels, and
take the first successful vector. No test over the continuum
\(\operatorname{Pic}^{0}(J)\) is needed, because the divisor argument
proves its vanishing consequence uniformly.

\subsection{SF-QC02. The coherent output and its complete
character}\label{SF-QC02}

Let \(A=J_x\times \widehat{J}_\alpha\), identify \(\widehat{J}\) with
\(J\) by \(\Theta\) when writing matrices, and put

\[
 x=\Theta_x,\qquad y=\Theta_\alpha,\qquad
 \gamma=c_1(\mathcal P_{x,\alpha}),\qquad
 H_d=dx+y,\qquad H_b=3(x+y),\qquad h=3H_d.
\]

Both \(H_b\) and \(h\) are very ample: the third power of an ample line
bundle on an abelian variety is very ample. Let \(Y=J_y\times A\) and
\(\pi :Y\to A\). Write \(a(y,x)=y+x\). On \(Y\) set

\[
 \mathcal L=\mathcal O(\Theta_y)\otimes a^*\mathcal O(\Theta)
                       \otimes\mathcal P_{y,\alpha}^{-1},
 \qquad D_1=\{y\in Z_H\},\quad D_2=\{y+x\in-Z_H\}.
\]

Define

\[
 \mathcal M=\mathcal L\otimes I_{D_1}I_{D_2},\qquad
 \mathcal G=R\pi_*\mathcal M,\qquad
 E=\mathcal G^\vee[-1].
 \tag{QC2.1}
\]

Then \(E\) is a coherent sheaf of rank \(r=8d\), locally of projective
dimension at most one. It is strongly equivariantly semiregular for a
coherent action of a finite subgroup of \(A[N]\times \widehat{A}[N]\).

To prove the sheaf assertion, the change of variables
\((y,x)\mapsto (y,z=y+x)\) makes \(D_{1},D_{2}\) pullbacks from
independent factors. Their ideals are Tor-independent and their
intersection equals their product. There is therefore an exact sequence
on the total space

\[
 0\longrightarrow\mathcal M\longrightarrow\mathcal L
 \longrightarrow\mathcal L|_{D_1}\oplus\mathcal L|_{D_2}
 \longrightarrow\mathcal L|_T\longrightarrow0,
 \quad T=D_1\cap D_2.
 \tag{QC2.2}
\]

There are \(N^{2}\) disjoint components of \(T\), each isomorphic to
\(C\times C\times \widehat{J}\). Its map \(f:T\to A\) is finite: after
translations and signs, its first-factor map is the addition
\(C\times C\to C^{(2)}\simeq\Theta\subset J\). For a nonhyperelliptic
genus-three curve the Abel map \(C^{(2)}\to\Theta\) is an embedding,
since an effective divisor of degree two has a unique section, and the
differential criterion gives an immersion. The quotient
\(C\times C\to C^{(2)}\) is finite flat of degree two. Thus \(\Theta\)
is smooth and each component of \(f\) factors as a finite flat
degree-two cover of a smooth divisor in A.

Set

\[
 W=\pi_*\mathcal L,\quad
 U=\pi_*(\mathcal L|_{D_1}\oplus\mathcal L|_{D_2}),\quad
 \mathcal T=f_*(\mathcal L|_T).
\]

On every \(\pi\)-fibre, \(L\) has class \(2\Theta\) and no higher
cohomology; \(W\) is locally free of rank eight. Each curve restriction
has degree six, and genus three implies \(H^{1}=0\) and \(h^{0}=4\).
Hence \(U\) is locally free of rank 8N. Since \(f\) is finite, \(T\) has
no higher direct image. Sequence (QC2.2) represents \(G\) by
\([W\to U\to T]\) in degrees 0,1,2. By (QC1.1), \(W\to U\) is fibrewise
injective, so \(Q=U/W\) is a vector bundle of rank \(8N-8=r\) and
\(G=[Q\to T]\) in degrees 1,2.

Each summand of \(T\) is the pushforward of a vector bundle on a smooth
divisor; hence it has local projective dimension one and only its first
sheaf Ext into \(O_A\) is nonzero. Finite duality identifies

\[
 \mathcal T^D:=\mathcal E xt_A^1(\mathcal T,\mathcal O_A)
 =\bigoplus_{i,j} f_{ij,*}
       ((\mathcal L|_{T_{ij}})^{-1}\otimes\omega_{T_{ij}}).
\]

Dualizing \([Q\to T]\) gives a coherent sheaf \(E\) and an exact
sequence

\[
 0\longrightarrow Q^*\longrightarrow E\longrightarrow\mathcal T^D
 \longrightarrow0.
 \tag{QC2.3}
\]

This also proves \(\operatorname{pd}(E)\le 1\). No statement about a
codimension-four nonfree locus is used.

For semiregularity, let \(q:J\to Z=J/H\) and \(F=I_{q(C)}\). D-SR-02
proves injectivity of the trace in degrees zero, one and two for \(F\)
and its transforms by inversion. Thus \(F\boxtimes (-1)^*F\) is fully
semiregular by the three Hochschild Künneth summands. Pullback under
\(q\times q\) gives \(I_{Z_H}\boxtimes I_{-Z_H}\) with its honest
\(H\times H\) descent action. Its invariant \(\operatorname{Ext}^{2}\)
is the original \(\operatorname{Ext}^{2}\), and trace pullback is
injective. Tensor by \(O(\Theta )\boxtimes O(\Theta )\), make the linear
change of variables, apply Fourier--Mukai on one factor, and finally
derived-dual and shift as in (QC2.1). These equivalences or
anti-equivalences preserve trace injectivity and transport the whole
coherent action. Translation and degree-zero tensor functors are taken
to functors of the same kind under these operations. Their action on
isomorphism classes is an automorphism of the abelian variety of such
functors, so the resulting group remains an injective finite subgroup of
\(A\times \widehat{A}\) of exponent \(N\). Coherence is transported from
the honest original descent action, not inferred from invariance of a
K-class. Since the output is a coherent sheaf, its ordinary negative
self-Ext groups vanish.

Here is the full character, independently calculated by exterior
integration. Put \(u^{2}=-d\), \(a=(1+1/u)/2\) and \(b=(1-1/u)/2\). Then

\[
 \operatorname{ch}(I_{Z_H}(\Theta))
 =1+\Theta-\frac d2\Theta^2-\frac d6\Theta^3
 =a e^{u\Theta}+b e^{-u\Theta}.
\]

In Darboux coordinates define

\[
 Q_u=\prod_{j=1}^3(u a_j-U_j)(u b_j-V_j),\qquad
 t=-ab(Q_u+Q_{-u}).
\]

With the sign convention of (QC2.1),

\[
 e^{\gamma/2}\operatorname{ch}(E)=t+P_d(H_d),
 \tag{QC2.4}
\]

where

\[
 P_d(z)=8d+2(d-1)z-z^2-\frac{d-1}{12d}z^3
     +\frac1{48d}z^4+\frac{d-1}{960d^2}z^5
     -\frac1{5760d^2}z^6.
 \tag{QC2.5}
\]

For completeness, in one symplectic plane, integrating the middle
coordinates c,e gives

\[
 \int_{c,e}\exp(s(a+c)(b+e)+tce-cV+eU)
  =(s+t)+st\,ab+s(aV-bU)-UV-s\,abUV.
\]

For opposite roots \(s=-t\) this is \(-e^{\gamma _j/2}(sa-U)(sb-V)\),
producing the two Weil determinant terms. For equal roots \(s=t\) it
equals

\[
 2s\exp\left(\frac{\gamma_j}{2}+\frac{s}{2d}(d a b+UV)\right).
\]

Multiplying the three planes, then dualizing and multiplying by \(-1\)
for the shift, gives the unequal-root term \(-ab(Q_u+Q_{-u})\) and the
equal-root polynomial

\[
 -8u^3\left(a^2e^{-uH_d/(2d)}-b^2e^{uH_d/(2d)}\right).
\]

Expanding through degree six proves (QC2.5). There is no additional
volume factor because \(\Theta\) on \(J\) is principal.

The Weil operator on \(H_{1}(J)\oplus H_{1}(J)\) is

\[
 J_K=\begin{pmatrix}0&-1\\d&0\end{pmatrix},\qquad J_K^2=-d.
\]

Its two eigenspaces on \(H^{1}\) have generators \(ua_j-U_j\),
\(ub_j-V_j\) and their conjugates, each of Hodge signature (3,3). It is
Rosati-skew for the polarization matrix \(\operatorname{diag}(d,1)\);
that product is
\(\bigl(\begin{smallmatrix}0&-d\\d&0\end{smallmatrix}\bigr)\). If
\(L_{0}\subset H_{1}(J,Q)\) is a rational symplectic Lagrangian,
\(L_{0}\oplus L_{0}\) is \(J_K\)-stable and totally isotropic for the
associated K-Hermitian form. Thus the Hermitian space is split with
determinant class \([-1]\). The two determinant terms \(Q_u,Q_{-u}\) are
independent and \(ab=N/(4d)\) is nonzero. Therefore \(t\) is a nonzero
rational primitive Weil class. In particular \(H_d t=0\).

Writing

\[
 a_d=\frac{d-1}{4d},\qquad
 D=\det E,\qquad c_1(D)=2(d-1)H_d-4d\gamma,
\]

the normalized entire character is

\[
 \kappa(E):=e^{-c_1(E)/r}\operatorname{ch}(E)
 =t+e^{-a_dH_d}P_d(H_d).
 \tag{QC2.6}
\]

All terms other than \(t\) are polarization powers. The coefficients
\(c_k\) of \(e^{-a_d z}P_d(z)\), through degree six, are

\[
\begin{array}{c|l}
0&8d\\
1&0\\
2&-(d+1)^2/(4d)\\
3&(d-1)(d+1)^2/(24d^2)\\
4&-(d-3)(d+1)^2(3d-1)/(768d^3)\\
5&(d-1)(d+1)^2(d^2-6d+1)/(3840d^4)\\
6&-(d+1)^2(d^2-10d+5)(5d^2-10d+1)/(368640d^5).
\end{array}
\tag{QC2.7}
\]

\subsection{SF-QC03. A complete explicit regularity
proof}\label{SF-QC03}

The sheaf \(E\) is 112-regular with respect to \(H_b\). More strongly,
for any integer \(s\ge 1\) and any line bundle \(M\) on A satisfying

\[
 c_1(M)-112s^2H_b\quad\hbox{is nef},
 \tag{QC3.1}
\]

one has \(H^i(A,[s]^*E\otimes M)=0\) for all \(i>0\). Arbitrary flat
factors of \(M\) are allowed.

The regularity convention is \(H^i(A,F((q-i)H_b))=0\) for every \(i>0\).
The standard Castelnuovo--Mumford lemma is applied to the pushforward
under the embedding defined by \(H_b\): a zero-regular sheaf is globally
generated, remains regular under nonnegative twists, and its section
multiplication maps are surjective. An inspected published locator is
Pareschi--Popa, M-regularity and the Fourier--Mukai transform, Theorem
1.2 (author version p.2); the publisher version is Pure Appl. Math.
Q.4(3),587--611 (2008). The relevant Buchsbaum--Rim terms are in
Raicu--Sam, Compositio158(2022), pp.41--42; their exactness in the
locally split case used here is also proved below.

\textbf{Step 1: the dual bundles.} For every integer \(q\ge 1\),

\[
 H^i(A,W^*(qH_b))=H^i(A,U^*(qH_b))=0\quad(i>0).
 \tag{QC3.2}
\]

For \(W\), relative duality identifies \(W^*\) with \(R^3\pi_*L^{-1}\).
Leray and Serre duality on the ninefold \(Y\) identify the dual of
\(H^i(A,W^*(qH_b))\) with \(H^{6-i}(Y,\mathcal L(-qH_b))\). Project
\(Y\) onto \(J_y\). On an A-fibre the line bundle has matrix
\(\operatorname{diag}(1-3q,-3q)\), hence is antiample. Its cohomology is
concentrated in degree six, so total cohomology below six vanishes.

For a summand of \(U\) supported on \(C\times A\), relative duality
first introduces the relative canonical bundle; Serre duality on the
sevenfold cancels it and again gives
\(H^{6-i}(C\times A,\mathcal L_C(-qH_b))\). Project onto \(C\). The
first curve type has the same fibre matrix. For the second type use
\(z=y+x\) as the fixed curve coordinate; its matrix is

\[
 \begin{pmatrix}1-3q&1\\1&-3q\end{pmatrix}.
\]

The largest eigenvalue of the matrix before subtracting 3qI is
\((1+\sqrt 5)/2<3\), proving antiampleness for \(q\ge 1\). Flat twists
and translations do not change this calculation. Thus \(W^*,U^*\) are
7-regular.

\textbf{Step 2: a kernel estimate.} Let \(e=h^{0}(U^*(7H_b))\).
Evaluation gives

\[
 0\to K_U\to O_A(-7H_b)^e\to U^*\to0.
\]

The kernel is a vector bundle and is 14-regular. For \(i=2,\ldots ,6\),
twist by \((14-i)H_b\) and use (QC3.2) and the positivity of
\((7-i)H_b\). For \(i=1\), section multiplication from 7-regularity
gives surjectivity at twist \(13H_b\); the multiplication by degree-six
ambient polynomials factors through \(H^{0}(A,6H_b)\), which is enough
for this evaluation map.

Compose evaluation with \(U^*\to W^*\), and let \(J_{0}\) be the
resulting kernel. There is an exact sequence

\[
 0\to K_U\to J_0\to Q^*\to0.
 \tag{QC3.3}
\]

The Buchsbaum--Rim complex resolving \(J_{0}\) has terms, beginning with
a surjection \(C_{1}\to J_{0}\),

\[
 C_k=\bigwedge^{8+k}(O_A(-7H_b)^e)
       \otimes\operatorname{Sym}^{k-1}W\otimes\det W,
 \qquad 1\le k\le e-8.
 \tag{QC3.4}
\]

Its determinant factor is required globally. To verify exactness in this
case, locally split the surjection as \(K\oplus V\to V\), with dim
\(V=8\). Splitting each wedge by the number a of factors from \(K\)
identifies the residual strands, after the determinant pairing, with the
polynomial de Rham complex of total degree \(a-1\) on \(V^*\). For
\(a\ge 2\) the Euler contraction divides by the nonzero integer \(a-1\)
and gives a contracting homotopy. The \(a=1\) strand maps isomorphically
to \(K\), and the remaining \(V\) maps isomorphically to the target.
This proves the needed local exactness in characteristic zero.

The theorem of the square gives
\(W=O_A(\Theta _x)\otimes b^*\operatorname{FM}(O_J(2\Theta ))\), where
\(b(x,\alpha )=x-\alpha\). Pulling the Poincaré family back by the
polarization isogeny {[}2{]} identifies it with
\(a^*O(2\Theta )\otimes p_{2}^*O(-2\Theta )\). Taking direct image gives
\(H^{0}(O(2\Theta ))\otimes O(-2\Theta )\). Consequently

\[
 [2]^*W\simeq Q_0^{\oplus8},\qquad
 c_1(Q_0)=4\Theta_x-2\Theta_{x-\alpha}.
 \tag{QC3.5}
\]

A different symmetric normalization adds a flat factor and affects no
vanishing assertion. The matrix of \(Q_{0}\) is
\(\bigl(\begin{smallmatrix}2&2\\2&-2\end{smallmatrix}\bigr)\), with
eigenvalues \(\pm 2\sqrt 2\), while \([2]^*H_b\) has matrix 12I. Thus
\(q[2]^*H_b+(k+7)c_{1}(Q_{0})\) is ample for \(q\ge k+7\). Kodaira
vanishing on A, and trace splitting for the finite étale map {[}2{]},
prove

\[
 \operatorname{reg}_{H_b}(\operatorname{Sym}^{k-1}W\otimes\det W)
 \le k+13.
\]

Hence reg \(C_k\le 7(8+k)+k+13=69+8k\). Only the first six terms of
(QC3.4) are needed. Successive cohomology sequences show the vanishings
required for R-regularity of \(J_{0}\) whenever

\[
 R\ge\max_{1\le k\le6}(\operatorname{reg} C_k-k+1)
 \le\max_{1\le k\le6}(70+7k)=112.
\]

The final possible cohomology is in degree seven and vanishes on the
sixfold. Sequence (QC3.3), together with 14-regularity of \(K_U\), gives
112-regularity of \(Q^*\). This bound is independent of the number \(N\)
of curves and of the number \(e\) of evaluation sections.

\textbf{Step 3: the torsion dual.} On one component
\(C\times C\times \widehat{J}\), put \(F_{1}={p}\times C\),
\(F_{2}=C\times {p}\), \(S=F_{1}+F_{2}\), and let \(\Delta _C\) be the
diagonal. The sum map \(C\times C\to J\) pulls \(\Theta\) back to

\[
 A_C=4S-\Delta_C.
\]

Indeed \(\Theta |C\) has degree three and its mixed Poincaré class is
\(S-\Delta _C\). Restriction to the diagonal gives degree twelve,
agreeing with \([2]^*\Theta |C\).

The divisor \(R_C=3S-\Delta _C\) is nef. Projection of the plane quartic
from a point of \(C\) gives a degree-three morphism to \(P^{1}\). The
residual to the diagonal in its fibre correspondence is effective of
class \(R_C\). If its monodromy is \(S_{3}\), this residual divisor is
irreducible, and its self-intersection is \(18-12-4=2\). An effective
irreducible divisor of nonnegative square on a smooth surface is nef. If
the transitive monodromy is \(A_{3}\), the cover is cyclic, and the
residual divisor is the sum of the two nonidentity automorphism graphs.
Each graph has square \(-4\). Riemann--Hurwitz gives five fixed points;
the two graphs meet transversely there because the two nontrivial
eigenvalues of the order-three local action differ. Their sum has
intersection one with each graph and nonnegative intersection with any
other irreducible curve. It is again nef.

Numerically on this component,

\[
 c_1(\mathcal L_T)=3F_1+3F_2-\mathcal P_{c_1,\alpha},\qquad
 f^*H_b=3A_C+3\Theta_\alpha.
\]

For \(q\ge 3\) decompose

\[
\begin{split}
 qf^*H_b-c_1(\mathcal L_T)
 ={}&[3qA_C-6F_1-3F_2]+(3q-1)\Theta_\alpha\\
 &+[3F_1+\Theta_\alpha+\mathcal P_{c_1,\alpha}].
\end{split}
\]

The last bracket is a pullback of \(\Theta\), hence nef. The first is
\(3qR_C+(3q-6)F_{1}+(3q-3)F_{2}\), ample on \(C\times C\). Together with
the strictly positive theta class on \(\widehat{J}\) it is ample on the
product. Thus the entire class is ample. Kodaira vanishing on the smooth
fivefold gives
\(H^i(T,\omega_T\otimes\mathcal L_T^{-1}\otimes f^*O(qH_b))=0\) for
\(i>0,q\ge 3\). Therefore \(T^D\) is 8-regular. Extension (QC2.3) proves
the asserted 112-regularity of \(E\).

\textbf{Step 4: the stronger bound (QC3.1).} Pull all exact complexes
back by {[}s{]} and tensor with \(M\). In Step 1 the two output fibre
matrices are multiplied by \(s^{2}\); subtracting \(c_{1}(M)\) makes
them negative definite under (QC3.1). Thus \([s]^*U^*\otimes M\) and
\([s]^*W^*\otimes M\) have no higher cohomology. In the pulled-back
evaluation sequence the middle term is a sum of ample line bundles
\(M(-7s^{2}H_b)\), so \(H^j([s]^*K_U\otimes M)=0\) for \(j\ge 2\). No
claim about \(H^{1}\) of this kernel is needed.

For \(k\le 6\) pull \([s]^*C_k\otimes M\) back once more by {[}2{]}. Its
summands have class

\[
 [2]^*c_1(M)-7s^2(8+k)[2]^*H_b+s^2(k+7)c_1(Q_0).
\]

After removing a nef class the matrix is \(s^{2}\) times

\[
 12[112-7(8+k)]I+(k+7)\begin{pmatrix}2&2\\2&-2\end{pmatrix}.
\]

Its least eigenvalue is positive for \(1\le k\le 6\); the worst
displayed lower bound is \(168-26\sqrt 2>0\). Kodaira vanishing, trace
splitting, the six-term truncation and (QC3.3) give the required
vanishings for \([s]^*Q^*\otimes M\).

For \(T^D\), form the finite étale base change \(T'=T\times_{A,[s]}A\)
with \(g:T'\to T\) and \(f':T'\to A\). The relevant line class is

\[
 f'^*c_1(M)-g^*c_1(\mathcal L_T)
 =g^*(3f^*H_b-c_1(\mathcal L_T))
       +f'^*(c_1(M)-3s^2H_b).
\]

The first summand is ample by Step 3 and the second is nef. Kodaira
vanishing on \(T'\), finite base change, and (QC2.3) prove (QC3.1) for
\(E\).

\subsection{SF-QC04. Ordinary normalization and coherent
equivariance}\label{SF-QC04}

Let \(r=8d\) and define the coherent sheaf

\[
 B_0=[r]^*E\otimes D^{-r},\qquad D=\det E.
 \tag{QC4.1}
\]

It has rank \(r\), projective dimension at most one, \(c_{1}(B_{0})=0\),
and complete character

\[
 \operatorname{ch}(B_0)=[r]^*\kappa(E)
       =r^6t+\sum_{k=0}^6 r^{2k}c_k H_d^k.
 \tag{QC4.2}
\]

Here \(r^{6}\) is the degree-six cohomological scaling, not the isogeny
degree. Since \(\operatorname{gcd}(r,N)=1\), this normalization
preserves a strong equivariant semiregularity structure for an actual
finite subgroup of \(A\times \widehat{A}\), as follows.

Write the original group action as \(T_g=t_{a_g}^*\otimes P_g\), with
\(a_g\in A[N]\), \(P_g\in \widehat{A}[N]\), and all associativity
constraints transported from descent. Let \(a'_g\) be the unique
N-torsion point with \(r\) \(a'_g=a_g\). Define
\(T'_g=t_{a'_g}^*\otimes[r]^*P_g\). Its coherent composition maps are
the pullbacks of those for \(T_g\), since
\(t_{a'_g}^*[r]^*=[r]^*t_{a_g}^*\). The deck group \(A[r]\) acts on
\([r]^*E\) in the canonical way; each pulled-back \(P_g\) has a
canonical deck linearization. Thus these actions commute coherently and
give an action of \(A[r]\times G\). The map to \(A\times \widehat{A}\)
is

\[
 (b,g)\longmapsto(b+a'_g,\,[r]^*P_g).
\]

It is injective: the first summands have coprime r- and N-torsion
orders, and multiplication by \(r\) is injective on N-torsion in the
dual. There is no invisible central kernel. Descent for the deck group
identifies the invariant \(\operatorname{Ext}^{2}\) with
\(\operatorname{Ext}^{2}(E,E)\), and then taking G-invariants gives the
original semiregular invariant space. Trace pullback is injective under
the isogeny. Tensor by \(D^{-r}\) conjugates this honest action by an
autoequivalence, hence preserves coherence and injectivity of its
underlying subgroup. This proves strong, and therefore weak, equivariant
semiregularity of \(B_{0}\). Ordinary negative Ext vanishes because
\(B_{0}\) is a coherent sheaf.

The full character (QC4.2) stays Hodge on the split Weil period family
with polarization \(h=3H_d\). On a suitable finite level cover the
rational Weil plane is constant, and \(r^{6}t\) extends as a flat Hodge
class. Apply Perry Theorem 1.2 with B-field zero. Its proof gives an
ordinary perfect deformation, not merely a twisted one: in §6.2 choose
\(b_{0}=0,n=1\) and the zero gerbe class. Corollary 5.17 may replace its
band by \(\mu _q\), but pushing out the zero class keeps the Brauer
class zero. Theorem 6.3 returns the object in the ambient category
\(\operatorname{Perf}(A_U,0)\). Any gerbe on an auxiliary quotient used
to geometrize the invariant category does not change this ambient
conclusion.

After shrinking the étale neighbourhood, the ordinary perfect
deformation is a flat coherent sheaf B and has fibrewise projective
dimension at most one. Here are precise algebraic details of the
openness step. At every point of the special fibre the derived
residue-field complex has amplitude in \([-1,0]\), because \(B_{0}\) has
projective dimension at most one. Openness of the Betti-number
vanishings for a perfect complex (Stacks Lemma36.31.1, Tag0BDI)
therefore gives a two-term vector-bundle presentation in degrees
\(-1,0\) on a neighbourhood of that point. On the special fibre its
differential is injective as a map of sheaves, since \(B_{0}\) is a
sheaf. The local criterion in Stacks Lemma10.99.1, Tag00ME, applied to
the local rings of the base and total space, proves that the
differential is injective near this point and that its cokernel is flat
over the base there. Openness of the flat locus (Theorem37.15.1,
Tag0399, citing EGA IV\(_3\), 11.3.1) extends this to an open
neighbourhood of the entire special fibre. Properness permits shrinking
the base so that every point of the new total space lies in this
neighbourhood. The two-term presentation gives \(pd\le 1\) on every
fibre, and the cokernel gives the desired ordinary flat coherent family.
No global boundedness of all sheaves with this Chern character is
asserted.

Put

\[
 R=224r^2+r(d+2)+6=14344d^2+16d+6.
 \tag{QC4.3}
\]

Then \(B_{0}\) is R-regular with respect to \(h\). Indeed
\(2h-H_b=3((2d-1)x+y)\) is ample. The matrix of \((d+2)h-c_{1}(D)\) is

\[
 \begin{pmatrix}d(d+8)&4d\\4d&d+8\end{pmatrix},
\]

whose determinant is \(d((d+8)^{2}-16d)=d(d^{2}+64)>0\). For
\(i=1,\ldots ,6\), the line bundle \(M_i=D^{-r}\otimes h^{R-i}\)
satisfies

\[
 c_1(M_i)-112r^2H_b
 =112r^2(2h-H_b)+r((d+2)h-c_1(D))+(6-i)h,
\]

which is ample. SF-QC03 supplies all six vanishings. Upper
semicontinuity in the ordinary flat coherent family preserves them after
shrinking. Thus every \(B_s\) on this open is R-regular, of rank \(r\),
\(c_{1}=0\) and \(pd\le 1\), with the same full character.

\subsection{SF-QC05. A numerical effective
detector}\label{SF-QC05}

Define the explicit integers

\[
 \begin{split}
 r&=8d,\qquad R=14344d^2+16d+6,\\
 V&=\int_A h^6=524880d^3,\\
 M&=729d^3r(R+r^2)^6,\\
 A_2&=\frac{1024}{9}d^3(d+1)^2.
 \end{split}
 \tag{QC5.1}
\]

\(A_{2}\) is an integer because \(d=4mL^{2}\) is divisible by nine. Set

\[
 D_{\rm CM}(m)=V\left((r+2)R A_2
       +7M\binom{r+1}{2}R^2+343\binom M3\right).
 \tag{QC5.2}
\]

This is a numerical formula involving only \(m\), powers, binomial
coefficients and the explicit finite lcm defining \(L\). On the entire
fixed split Weil family containing (A,h,K), there is an effective
integral three-cycle with nonzero Weil projection and \(h\)-degree at
most \(D_{\rm CM}(m)\), using the authorized detector closedness input
to pass from the open obtained above to its component. No Chow
denominator occurs in the detecting cycle.

In particular the simpler safe integer

\[
 \boxed{D_{\rm CM}(m)\le 2^{318}d^{51}
                  =2^{420}m^{51}L^{102}}
 \tag{QC5.2a}
\]

may be used instead. Indeed \(r+2\le 2^{4}d\), \(R<2^{14}d^{2}\),
\(A_{2}<2^{9}d^{5}\), \(V<2^{20}d^{3}\) and
\(\operatorname{binom}(r+1,2)\le 2^{6}d^{2}\). Also
\(R+r^{2}\le 14430d^{2}<2^{14}d^{2}\), so \(M<2^{97}d^{16}\). The three
summands in (QC5.2) are respectively bounded by \(2^{47}d^{11}\),
\(2^{154}d^{25}\), and \(2^{317}d^{51}\); for the last use 343
\(\operatorname{binom}(M,3)<64M^{3}\). Since \(d\ge 1\), their sum is
below \(2^{318}d^{51}\).

\textbf{Proof.} On the deformation open, \(F_s=B_s(Rh)\) is zero-regular
and globally generated. Its higher cohomology vanishes and
\(v=h^{0}(F_s)=\chi (F_s)\) is constant. The evaluation sequence

\[
 0\to K_s\to O_{A_s}^{v}\to F_s\to0
\]

has a vector-bundle kernel because \(\operatorname{pd}(F_s)\le 1\). Put
\(k=v-r\). The same section multiplication argument used above shows
that \(K_s\) is 7-regular: for cohomological degrees at least two use
the positivity of \((7-i)h\) and zero-regularity of \(F_s\), and in
degree one use surjectivity of
\(H^{0}(F_s)\otimes H^{0}(O(6h))\to H^{0}(F_s(6h))\). Hence
\(G_s=K_s(7h)\) is a globally generated vector bundle.

Let

\[
 A_3=\frac{65536}{81}d^4(d-1)(d+1)^2.
\]

From (QC2.7) and \(c_{1}(B_s)=0\),

\[
 c_2(B_s)=A_2h^2,\qquad
 c_3(B_s)=A_3h^3+2r^6t_s.
 \tag{QC5.3}
\]

The sign and all coefficients here follow from
\(\operatorname{ch}_{2}=-c_{2}\) and \(\operatorname{ch}_{3}=c_{3}/2\)
when \(c_{1}=0\). Inverting \(c(F_s)\) in the evaluation sequence and
then twisting by 7h gives

\[
\begin{split}
 c_3(G_s)= {}&-2r^6t_s+ S_kh^3,\\
 S_k={}&-A_3+(r+2)R A_2-\binom{r+2}{3}R^3\\
 &+7(k-2)\left(\binom{r+1}{2}R^2-A_2\right)
   -49rR\binom{k-1}{2}+343\binom k3.
\end{split}
\tag{QC5.4}
\]

Thus \(c_{3}(G_s)\) has nonzero Weil projection. In particular
\(k\ge 3\). Choose \(k-2\) general global sections of \(G_s\). Their
first rank-drop locus has expected codimension three, or is empty,
because the bundle is globally generated and the universal evaluation is
transverse to the matrix-rank strata. The Thom--Porteous formula gives
its class \(c_{3}(G_s)\). Nonzero Weil projection excludes the empty
case. It is therefore an integral effective detecting cycle.

It remains to bound \(k\). Write \(c_j\) for (QC2.7). Since \(H_d t=0\)
and \(td(A)=1\), Riemann--Roch gives

\[
 v=V\sum_{j=0}^6\frac{r^{2j}c_j}{3^j}
                      \frac{R^{6-j}}{(6-j)!}.
\]

For \(d\ge 1\) each coefficient of \(P_d\) has absolute value at most
r/j!. Since \(0\le a_d<1/4\), multiplying by
\(\operatorname{exp}(-a_dz)\) gives \(|c_j|\le r(5/4)^j/j!\).
Consequently

\[
 v\le\frac{Vr}{6!}(R+5r^2/12)^6
    \le729d^3r(R+r^2)^6=M.
\]

Thus \(3\le k\le M\). Since \(d>1\), both \(A_{2},A_{3}\) are positive.
Discarding the negative terms in (QC5.4) and replacing \(k\) by \(M\) in
the remaining positive terms gives \(S_k\le (r+2)RA_{2}+7M\)
\(\operatorname{binom}(r+1,2)R^{2}+343\) \(\operatorname{binom}(M,3)\).
Finally \(h^{3}t_s=0\), so the degree is \(VS_k\le D_{\rm CM}(m)\),
proving (QC5.2).

The source polarization is specified: \(h=3(dx+y)\) has type
(3,3,3,3d,3d,3d), product \(729d^{3}\). This result must not be
described as a principal-polarization bound without an additional
lattice/isogeny calculation. For any fixed target split polarized
K-lattice, a chosen rational polarization isometry whose denominators
are cleared by \(q\) gives a fixed isogeny with pullback target
polarization \(q^{2}h\); pushforward transports the detector bound with
cost \(q^{6}\). A numerical bound uniform over all target polarization
products requires bounding that \(q\), separately.

For class normalization, the exterior pairing gives
\(\int t^{2}=-8dN^{2}\): each symplectic plane contributes
\(Q_uQ_{-u}=-4d\) times its oriented volume, and
\(2(ab)^{2}=2N^{2}/(16d^{2})\). Therefore, for \(t_B=r^{6}t\),

\[
 -\int t_B^2=8dN^2r^{12}=2^{39}d^{13}N^2.
\]

Moreover \(A_{3}\) is an integer for the chosen \(d\), since 9 divides
\(d\). Formula (QC5.3) then expresses \(2t_B\) as the difference of
integral algebraic cycle classes. Thus the lattice denominator
\(\delta =2\) is valid for this explicit flat Weil class. This does not
assert that an arbitrary rational Weil vector has denominator two.

\subsection{Source and status notes}\label{source-and-status-notes}

\begin{itemize}
\item
  The finite torsion avoidance and the denominator-free effective Chern
  construction are new arguments in this run. The proof does not invoke
  an unverified generic-rank or generic-position statement.
\item
  Pareschi--Popa, M-regularity and the Fourier--Mukai transform, Pure
  Appl. Math. Q.4(3),587--611 (2008), Theorem1.2; exact basic regularity
  statement inspected in the author PDF p.2. Publisher
  DOI10.4310/PAMQ.2008.v4.n3.a1.
\item
  Raicu--Sam, Bi-graded Koszul modules, K3 carpets, and Green's
  conjecture, Compositio158(2022),33--56, pp.41--42 immediately before
  Theorem3.2. The globally correct Buchsbaum--Rim determinant factor was
  inspected in the published PDF, and the locally split exactness is
  proved here.
\item
  Bauer--Szemberg, On tensor products of ample line bundles on abelian
  varieties, Math.Z.223(1996),79--85, Theorem1.1(a),(b), author version
  p.1. The proof and publication header were inspected.
\item
  Fulton, Intersection Theory, second edition (1998), Theorem14.4,
  pp.254--255. The Thom--Porteous statement used here was checked
  against the publisher's chapter abstract and the indexed primary book
  text; a full independent audit of the entire chapter proof is not
  claimed.
\item
  Mukai1981, Theorem2.2 and equation3.1, and Buchweitz--Flenner2003
  Proposition4.2, remain the published Fourier--Mukai and trace inputs
  already checked in D-SR-01--04.
\item
  Perry2604.00511v2 remains an authorized preprint input. The zero-B
  ordinary conclusion is an inference from the full proofs of
  Corollary5.17, Theorems6.1/6.3 and Theorem1.2 in §6.2, not an extra
  theorem quoted without checking.
\item
  Exact Perry page locators for this inference: Corollary5.17 pp.20--21;
  Theorems6.1/6.3 pp.25--26; §6.2 and the proof of Theorem1.2 pp.26--28.
  All were inspected in the official versioned PDF.
\item
  Stacks Project Tags0BDI,00ME,0399 and07VK were opened and checked for,
  respectively, openness of perfect-complex Betti vanishings, the
  injectivity/flat-cokernel criterion, openness of the flat locus, and
  perfect cohomology with arbitrary base change. The last two provide
  the precise family hypotheses used for upper semicontinuity.
\item
  All regularity constants used for the numerical estimate have been
  proved above. HECKE\_DIV M7 is provenance for the architecture of that
  argument, not a premise, and Markman's output sheaf and
  general-position assertions are unnecessary here.
\end{itemize}

\section{A numerical CM bound for every split class-number-one datum}

\begin{theorem}[HD-CM6: an explicit split-sixfold bound]
\label{HD-CM6}
Let $K=\mathbf Q(\sqrt{-m})$ have class number one, and define the
integers $L,d=4mL^2$ as in SF-QC01. Let $(A,\theta)$ be a split
Weil sixfold with integral compatible maximal-order action, and
$\Delta=\int_A\theta^6/6!$. Then, using the supplied Perry
deformation input with its stated preprint status,
\[
 \boxed{\delta_W(A,\theta)\le (2d\Delta)^6D_{\rm CM}(m)
                         \le2^{324}d^{57}\Delta^6.}
\]
The same bound holds uniformly on the entire fixed polarized family.
If only $\mathbf Z+f_t\mathcal O_K$ acts integrally, replace
$\Delta$ by $f_t^{12}\Delta$, giving the safe last bound
$2^{324}d^{57}f_t^{72}\Delta^6$.
\end{theorem}
\begin{proof}
We make the source-to-hyperbolic lattice map explicit; a varying
isogeny degree is not being used as an uncontrolled parameter.
The source of SF-QC05 is $J\times\widehat J$, identified with
$J^2$ by its principal polarization, with
$h=3(d\Theta_x+\Theta_y)$ and
\[
 u=\begin{pmatrix}0&-1\\d&0\end{pmatrix},\qquad u^2=-d.
\]
Choose an integral symplectic basis $a_1,a_2,a_3,b_1,b_2,b_3$
of $H_1(J,\mathbf Z)$, and use its vectors in the first factor
as a $K$-basis of the full rational homology. The second-factor
vector corresponding to $v$ is $(u/d)v$. Thus the integral source
lattice in this $K$-basis is
\[
 \Lambda_s=(\mathbf Z+(u/d)\mathbf Z)^6.
\]
Write $\mathcal O_K=\mathbf Z[\tau]$, where
$\tau=a+u/f_s$. If $m\not\equiv3\pmod4$, take
$a=0,f_s=2L$; if $m\equiv3\pmod4$, take
$a=1/2,f_s=4L$. In both cases $f_s\mid d$.
For the normalized Hermitian form
$h_K(v,w)=E_h(v,\tau w)-\bar\tau E_h(v,w)$,
the matrix in the displayed $K$-basis is
\[
 H_s=\begin{pmatrix}0&kI_3\\-kI_3&0\end{pmatrix},
 \qquad k=3du/f_s\in\mathcal O_K,
 \qquad\bar k=-k.
\]
Indeed $u$ takes a first-factor vector into the second factor,
which is $E_h$-orthogonal to the first, while the alternating form
on the first factor is $3dE_\Theta$. This proves the formula
including its normalization factor.

Let $U^3$ denote the standard hyperbolic Hermitian lattice on
$\mathcal O_K^6$. The matrix
$F=\operatorname{diag}(I_3,kI_3)$ satisfies
$F^*U^3F=H_s$. Moreover $dF\Lambda_s\subset\mathcal O_K^6$:
the first three coordinates lie in $d\mathbf Z+u\mathbf Z$,
and the last three in $dk\mathbf Z+ku\mathbf Z$.
Consequently, after transporting the complex structure, $dF$
defines an actual $K$-linear isogeny from the source to the
standard hyperbolic polarized abelian variety. Its pullback of
the target polarization is $d^2h$. A source threefold detector
of $h$-degree at most $D_{\rm CM}(m)$ therefore pushes to an
effective detector of degree at most $d^6D_{\rm CM}(m)$.
The rational isometry identifies the full period domains;
intersecting the arithmetic lattice stabilizers gives the finite
level correspondence. Thus this assertion holds throughout the
hyperbolic component, not just at its CM point.

The target lattice is free over $\mathcal O_K$ by the
class-number-one hypothesis. Intersect a maximal rational
isotropic subspace with it and extend a basis of that saturated
rank-three submodule to an integral basis. Its Hermitian matrix
is $\bigl(\begin{smallmatrix}0&B\\B^*&C\end{smallmatrix}\bigr)$,
with $\operatorname{Nm}(\det B)=\Delta$. The explicit substitution
in SF-04a changes this matrix to $U^3$ and has all denominators
dividing $2\Delta$ in both directions. That algebraic calculation
uses only a principal ideal domain and complex conjugation, so it
applies to the present ring, including the non-Euclidean
class-number-one rings. Multiplication by $2\Delta$ gives an
isogeny from the hyperbolic datum to the target with pullback
polarization $(2\Delta)^2\theta_U$. Pushforward therefore costs
at most $(2\Delta)^6$. Combining the two maps proves the first
bound, and SF-QC05 gives the second.

For a fixed smaller order, write $A=V/\Lambda$ and put
$\Lambda'=f_t\mathcal O_K\Lambda$. The inclusions
$f_t\Lambda\subset\Lambda'\subset\Lambda$ give an index
$I\le f_t^{12}$. The isogeny $V/\Lambda'\to A$ has pulled-back
polarization product $I\Delta$. Apply the maximal-order bound
and push forward; projection formula preserves the cycle degree
relative to these specified polarizations and preserves detection.
Monotonicity in $\Delta$ gives the conductor assertion.
\end{proof}

\begin{corollary}[HD-CM4: all fourfold determinants by compensation]
\label{HD-CM4}
For every class-number-one imaginary quadratic $K$, the same
explicit $d$ and $c_K=1+\operatorname{Nm}(\tau)$ give, on every
fourfold determinant and integral maximal-order polarization datum,
\[
 \delta_W(A,\theta)\le
 c_K(2d\Delta^2)^6D_{\rm CM}(m)
 \le c_K2^{324}d^{57}\Delta^{12}.
\]
For a fixed conductor $f_t$, the last expression is multiplied
by $f_t^{96}$. Smaller previously proved bounds are retained
when available. These conclusions use Perry as in SF-QC04;
no Markman, Severn or 032 theorem is a premise.
\end{corollary}
\begin{proof}
Apply HD-COMP with the explicit split-sixfold function HD-CM6.
The conductor statement uses the fourfold lattice index
$I\le f_t^8$ in the same pushforward argument.
\end{proof}

\section{One explicit quadratic division constant for every field}

\begin{proposition}[HD-U: a uniform symmetric integral projector]
\label{HD-U}
Let $(A,\theta,K)$ have dimension $2n$ and signature $(n,n)$,
and suppose $a\in K$ acts by an integral endomorphism and
$a/\bar a$ is not a root of unity. If an effective integral
detector has $\theta$-degree at most $D$, put $Q=\operatorname{Nm}(a)$.
For every nonzero rational Weil class $t$, with $\delta t$ integral
and $T=(-1)^n\int_A t^2>0$, one has the entirely explicit bound
\[
 \boxed{c_t\le
 \left\lceil\frac{8\cdot16^n}{n!}\right\rceil
 \delta\lceil\sqrt T\rceil Q^{2n(n+1)}D^2.}
\]
The bound includes the integral multiplier and the degree of each
effective positive and negative part. In particular the scalar
constants for $n=2,3,4$ are respectively $1024,5462,21846$.
\end{proposition}
\begin{proof}
Set $q=Q^n$ and $S=[a]^*+[\bar a]^*$ on $H^{2n}(A)$.
For $0\le j\le n$ put
\[
 \alpha_j=a^j\bar a^{2n-j}+a^{2n-j}\bar a^j\in\mathbf Z.
\]
These are the eigenvalues of $S$ on the $K$-character summands,
with the complementary indices paired. Write $r=a/\bar a$.
For $1\le j\le n$,
\[
 \alpha_0-\alpha_j=\bar a^{2n}(1-r^j)(1-r^{2n-j})\ne0.
\]
Consequently, after changing an overall sign if necessary,
\[
 \prod_{j=1}^n(S-\alpha_j)=\pm R e_W,
 \qquad R=\left|\prod_{j=1}^n(\alpha_0-\alpha_j)\right|
 \in\mathbf Z_{>0}.
\]
This is an integral correspondence, not a division in the cycle
group. A pullback by $a$ or $\bar a$ multiplies the degree of an
$n$-cycle by $q$: its compatible polarization is pulled back to
$Q\theta$, while the finite map has degree $Q^{2n}$.
Expanding the displayed polynomial therefore gives integral
effective positive and negative parts, each of degree at most
$\kappa D$, where
\[
 \kappa=\prod_{j=1}^n(2q+|\alpha_j|)\le(4q)^n,
 \qquad R\le(4q)^n.
\]
An overall sign just exchanges the parts.

Apply the proved quadratic Gram cancellation HD-WQ to this
projector, using $[a]^*$ as its second basis operator. In its
notation $s=\alpha_0$ and $H=4q^2-s^2$; the non-root-of-unity
hypothesis makes $H>0$, and $H\le4q^2$. The integral multiplier
is at most
\[
 \frac{\delta HR^2}{(n!)^2}D^2
 \le \frac{4\delta\,16^n q^{2n+2}}{(n!)^2}D^2,
\]
and each effective part has degree at most
\[
 \frac{4\delta q\sqrt{HT}\,\kappa R}{n!}D^2
 \le \frac{8\delta\,16^n q^{2n+2}}{n!}\sqrt T\,D^2.
\]
The second numerical coefficient dominates the first, since
$n!\ge1$. Replacing $\sqrt T$ and the scalar coefficient by
their ceilings gives the displayed integer bound. All cycles
used by HD-WQ are integral; none of these numerical simplifications
divides an effective cycle.
\end{proof}

\paragraph{An endomorphism that always works.}
Write $K=\mathbf Q(\sqrt{-m})$ with $m>0$ squarefree. If the
integral order contains $\mathbf Z+f\mathcal O_K$, choose
$a=f(2+\sqrt{-m})$. Then $a$ acts integrally and
$Q=f^2(m+4)$. Its ratio with its conjugate is not a root of unity.
For fields other than $\mathbf Q(i),\mathbf Q(\sqrt{-3})$ this
follows because the only roots of unity are $\pm1$, while $a$
has both real and imaginary parts nonzero. In the two exceptional
fields the real part of the ratio is respectively $3/5$ or $1/7$,
which is not the real part of any root of unity in that field.
Scaling by $f$ does not change the ratio. The small specialized
projectors HD-G05 and SF-06--07 give sharper constants when
their hypotheses hold.

\begin{corollary}[HD-U6: a fully explicit split-sixfold division bound]
\label{HD-U6}
For every class-number-one field $K=\mathbf Q(\sqrt{-m})$,
with $d$ as in HD-CM6, polarization product $\Delta$, and fixed
integral order containing $\mathbf Z+f\mathcal O_K$, the split
sixfold families satisfy
\[
 \boxed{c_t\le 2^{661}\,
 \delta\lceil\sqrt T\rceil(m+4)^{24}d^{114}f^{192}\Delta^{12}.}
\]
For every fourfold determinant the compensation construction gives
\[
 c_t\le c_K^2 2^{658}\,
 \delta\lceil\sqrt T\rceil(m+4)^{12}d^{114}f^{216}\Delta^{24},
 \qquad c_K=1+\operatorname{Nm}(\tau).
\]
These are safe uniform integers, not claimed optimal. Their CM
supplier premise is exactly the checked SF-QC proof using the
authorized Perry input.
\end{corollary}
\begin{proof}
For sixfolds insert $D=2^{324}d^{57}f^{72}\Delta^6$ and
$Q=f^2(m+4)$ into HD-U with $n=3$, and use $5462<2^{13}$.
For fourfolds insert
$D=c_K2^{324}d^{57}f^{96}\Delta^{12}$ into the $n=2$ formula,
whose coefficient is $1024=2^{10}$. These substitutions give
the displayed exponents. The actual effective cycles and bounded
integral multiplier come from HD-U, so the numerical upper bounds
do not introduce rational cycles.
\end{proof}

For the independently reconstructed Eisenstein eightfold source,
the same proposition gives the useful uniform statement
\[
 c_t\le21846\,\delta\lceil\sqrt T\rceil
       (7f^2)^{40}D_8(f^{16}\Delta)^2.
\]
Here the source is split; the separate compensation HD-E6 treats
every Eisenstein sixfold determinant. These formulas make explicit
the condition-A integral numerators and multipliers. Their
projectors are the polynomial-inverse construction of cote~16
scans~048--049; condition~I controls the positive Gram form only
after a detector has been supplied.

\section{A determinant formula and two universal nonsplit-source barriers}
\label{sec:finite-source-barrier}

The splitness of the two preceding witnesses is forced by their local
geometry.  We prove a parabolic determinant formula and then show that
every dominant sixfold curve-cover family using either of these natural
representations has exactly the relevant involution inertia types.
The conclusion concerns these representations and does not classify all
finite groups with quadratic character field.

\subsection*{D-DET-01. The Seifert determinant formula}

Let \(K\) be an imaginary quadratic field, let \((V,h)\) be a positive
definite Hermitian \(K\)-space of dimension \(m\), and let
\(g_1,\ldots,g_r\in U(V,h)\) be finite-order elements with product one
and no common invariant vector.  We use right actions in this proof, to
match exactly the conventions of \cite[Theorem 2.5, p.~12]{G34:DWdual}.
Put
\[
 U_i=\operatorname{im}(g_i-1),\quad d_i=\dim_KU_i,\quad
 D=\sum_i d_i,\quad q_i=g_{i+1}\cdots g_r,
 \quad L_i=U_iq_i.
\]
Each restriction \(h_i=h|_{U_i}\) is nondegenerate.  Assume the
parabolic multiplicity space has even dimension \(2n=D-2m\).
Set
\[
                   b=\prod_i\det(1-g_i|_{U_i}).
\]
Then \(b\in\mathbb Q^*\), and the signed Hermitian discriminant of
the parabolic space is
\begin{equation}\label{eq:det-seifert}
       \boxed{\displaystyle
       \left[\frac{\prod_i\det h_i}{b}\right]
          \quad\text{in }\mathbb Q^*/N_{K/\mathbb Q}(K^*).}
\end{equation}
The determinants of the \(h_i\) may be computed in arbitrary
\(K\)-bases; changing those bases changes the expression by a norm.

\begin{proof}
Let \(\mathcal V=\bigoplus_iL_i\), with summation map
\(R:\mathcal V\to V\), \(R(w)=\sum_iw_i\).  A parabolic cocycle
\(u_i\in U_i\) has transported values \(w_i=u_iq_i\), and its relation
is \(R(w)=0\).  Its coboundary map is
\[
                    d(a)_i=a(g_i-1)q_i.
\]
The product-one relation makes \(Rd=0\).  The common-invariant-vector
hypothesis makes \(d\) injective.

Let \(G_i=q_i^{-1}g_iq_i\) act on \(L_i\), and set
\[
             A_i=(1-G_i^{-1})^{-1}|_{L_i}.
\]
These inverses exist, because \(L_i\) is the moving space.  Unitarity
gives \(A_i+A_i^*=1\).  On \(\mathcal V\) define the triangular
sesquilinear form
\[
 \mathcal S(w,z)=\sum_{j<i}h(w_j,z_i)+\sum_i h(w_iA_i,z_i).
\]
Its diagonal blocks are invertible, so it is nondegenerate.  Since
\(A_i+A_i^*=1\),
\[
                  \mathcal S+\mathcal S^*=R^*hR.
\]
We also have the crucial identity
\begin{equation}\label{eq:det-coboundary}
                     \mathcal S(d(a),z)=h(a,Rz).
\end{equation}
Indeed
\(\sum_{j<i}d(a)_j=a-a g_iq_i\).  Writing \(P_i\) for the
orthogonal projector onto \(U_i\), the diagonal term satisfies
\[
 d(a)_iA_i=a(g_i-1)(1-g_i^{-1})^{-1}q_i
                         =a g_iP_iq_i.
\]
Their sum differs from \(a\) by a vector in
\(\ker(g_i-1)q_i=L_i^\perp\).  Pairing with \(z_i\in L_i\) gives
\eqref{eq:det-coboundary}.

Because \(\mathcal S\) is nondegenerate and \(d\) is injective,
\eqref{eq:det-coboundary} also proves that \(R\) is surjective.
On \(\ker R\), \(\mathcal S\) is anti-Hermitian, with radical
exactly \(dV\).  Thus it induces a nondegenerate anti-Hermitian form
\(S_P\) on
\[
                         P=\ker R/dV.
\]

This is the geometric parabolic cup form up to one overall sign.
Here is the convention check, rather than an appeal to an unspecified
auxiliary form.  Apply Dettweiler--Wewers Theorem2.5 to the dual cocycle
obtained from \(h\).  Its formula simplifies, using
\(u_i'=u_i(g_i-1)^{-1}\), to
\[
 \sum_i h(u_i,v_i(g_i-1)^{-1})
       -\sum_{j<i}h(u_j,v_i(g_{j+1}\cdots g_i)^{-1}).
\]
Transporting both arguments by the tails \(q_i\), the second term is
\(-\sum_{j<i}h(w_j,z_i)\).  The first is
\(-\sum_i h(w_iA_i,z_i)\), by the adjoint identity for a unitary
operator.  Thus their cup pairing is \(-S_P\).  Since \(\dim_KP=2n\),
this sign has no effect on its determinant.  For involutions the diagonal
operator is exactly \(A_i=\tfrac12 I\), agreeing with the triangular
form used in the explicit certificates.

Choose a basis of \(\mathcal V\) consisting of a basis of \(dV\),
representatives of a basis of \(P\), and a complement \(C\) on which
\(R\) is an isomorphism.  In this basis the form has block shape
\[
 \begin{pmatrix}
 0&0&B\\0&S_P&F\\-B^*&-F^*&J
 \end{pmatrix},\qquad B=hR_C.
\]
Taking the Schur complement of \(S_P\) shows that its determinant is
\(N_{K/\mathbb Q}(\det B)\det S_P\).  A basis change of a
sesquilinear form multiplies its determinant by a norm.  Therefore
\([\det S_P]=[\det\mathcal S]\).

The triangular diagonal blocks give
\[
               \det\mathcal S=\frac{\prod_i\det h_i}{b}.
\]
For clarity, \(\det(h(\,\cdot\, A_i,\,\cdot\,))\)
equals \(\det h_i\,\overline{\det A_i}\), and unitarity makes
\(\overline{\det A_i}=\det(1-g_i|_{U_i})^{-1}\).
Moreover
\[
 \bar b=(-1)^D\left(\prod_i\det g_i\right)^{-1}b=b,
\]
because \(D=2m+2n\) is even and \(\prod_i g_i=1\).
Thus \(b\) and the determinant expression are rational.

Finally, multiply the anti-Hermitian form by any nonzero imaginary
\(\eta\in K\) to obtain a Hermitian form.  Its signed determinant is
\[
             (-1)^n\det(\eta S_P)
                   =N_{K/\mathbb Q}(\eta^n)\det S_P.
\]
It has the class claimed in \eqref{eq:det-seifert}.  The cup pairing on
the one-copy abelian factor differs from the multiplicity pairing only
by the rational cover-degree and root-line polarization scalars.  Such
a scalar to the power \(2n\) is a norm.  Switching the selected dual
or conjugate multiplicity convention also preserves the class (the
quotient \(\mathbb Q^*/N(K^*)\) has exponent two).  Hence the formula
applies to the actual polarized abelian factor.
\end{proof}

For involution inertia, \(b=2^D\) is a norm.  Thus
\begin{equation}\label{eq:det-involution}
           \boxed{[\operatorname{disc}_{\rm signed}P]
                    =\prod_i[\det(h|_{U_i})].}
\end{equation}
This formula uses only the local Hermitian moving spaces.  It is
independent of the ordering, Nielsen class, and Hurwitz component.

\subsection*{D-DET-02. The \(G_{24}\) barrier}

For the natural \(G_{24}\)-representation over
\(K=\mathbb Q(\sqrt{-7})\), every reflection has a norm-two root.
Since \(2=N_{K/\mathbb Q}(\tau)\), each reflection moving line has
norm determinant.  Thus every product-one generating tuple of these
involutions giving a balanced multiplicity space has split signed
discriminant.  This conclusion holds for unpaired tuples as well.

In fact every involution in this representation has norm moving-space
determinant.  Rank-one involutions are the reflections.  For a rank-two
involution, its negative plane is the orthogonal complement of the root
line of its negative, which is a reflection because \(-I\in G_{24}\).
Since \(\det H=1\), its plane determinant has class \(1/2\), again a
norm.  The remaining involution is \(-I\), whose moving space has
determinant one.  Thus \eqref{eq:det-involution} also covers arbitrary
involution types in \(G_{24}\).

There is a stronger full-source obstruction.  Let the base curve have
genus \(\gamma\), let \(r\) be the number of genuine branch points, and
suppose the natural multiplicity factor is a sixfold.  The parabolic
Euler characteristic gives
\[
             6=(2\gamma-2)\cdot3+D,\qquad D=12-6\gamma.
\]
Each nontrivial inertia element has moving-space dimension at least one,
so \(r\le D\).  For \(\gamma=1\), the Hurwitz dimension is at most
\(6\); for \(\gamma=2\), only unramified covers are possible and the
base moduli dimension is three; larger \(\gamma\) is impossible.
Thus a period map of rank nine must have \(\gamma=0\).

For \(\gamma=0\), the Hurwitz dimension is \(r-3\le9\).  Rank nine
forces \(r=12=D\), so all inertia elements have moving rank one.  Their
nontrivial eigenvalue belongs to \(K\), since it is the trace minus two.
The only roots of unity in \(K\) are \(\pm1\), so they are involutive
reflections.  The determinant formula forces splitness.

\textbf{Conclusion.} No dominant Weil-sixfold curve-cover family using
the natural \(G_{24}\)-representation can lie in a nonsplit determinant
component, even if arbitrary inertia orders and base genera were initially
allowed.  The inequalities are used only as necessary conditions for
dominance; positive rank nine was proved separately for the split witness.

\subsection*{D-DET-03. The \(F_4\)-extension barrier}

First we prove the local assertion for every outside involution, without
assuming it lies in a particular finite conjugacy class.  An element
\(g\in G_F\setminus W\) has the form \(tA\), with \(A\in M_4(\mathbb Q)\)
and \(t^2=-2\).  If \(g^2=1\), unitarity gives
\[
                       A^t=-A,\qquad A^2=-I/2.
\]
Let \(U=\ker(g+1)\).  The rational-linear map
\[
               x\longmapsto (1-tA)x:\mathbb Q^4\longrightarrow U
\]
is an isomorphism of rational vector spaces.  It lands in \(U\) by the
displayed square relation; injectivity follows by comparing rational and
\(t\)-coefficients.  Its Hermitian diagonal form is
\[
              h((1-tA)x,(1-tA)x)=2x^t x.
\]
Thus every such negative plane has the same underlying rational quadratic
form as the explicit negative plane of \(z\).

We use the following elementary trace-Witt fact.  If two \(K\)-Hermitian
forms have isometric rational diagonal forms, their difference has
hyperbolic rational diagonal form.  A nonzero rational isotropic vector
is also a \(K\)-isotropic vector.  Split off a Hermitian hyperbolic plane;
its underlying rational form is hyperbolic of dimension four.  Rational
Witt cancellation leaves a hyperbolic complement, and induction shows
the whole Hermitian difference is hyperbolic.  In particular equal-rank
forms have determinant classes differing by a norm.  This argument does
not identify a real signature with a rational Witt class.

The explicit plane of \(z\), with the basis in D-GF4-02, has Gram matrix
\[
                         \begin{pmatrix}2&t\\-t&2\end{pmatrix},
                         \qquad\det=2=N_{K/\mathbb Q}(t).
\]
Every outside involution therefore has norm moving-plane determinant.
The reflection moving lines have norms one or two, also norms from
\(K=\mathbb Q(\sqrt{-2})\).  Hence every product-one generating tuple
consisting of reflections and outside involutions has split discriminant,
whether or not its entries occur in adjacent equal pairs.

Now allow arbitrary inertia.  The sixfold dimension formula is
\[
                 6=(2\gamma-2)\cdot4+D,
                 \qquad D=14-8\gamma.
\]
For \(\gamma=1\), the Hurwitz dimension is at most six; for larger
\(\gamma\), the dimension equation is impossible.  Thus rank-nine
dominance requires a genus-zero base.

On that base, the product-one generating tuple must contain at least
two elements outside \(W\): the quotient \(G_F/W\) has order two,
so their number is positive and even.  For any outside element,
\(\bar g=-g\).  Complex conjugation therefore carries its \(+1\)
eigenspace to its \(-1\) eigenspace.  Its fixed dimension is at most
two, so its moving dimension is at least two.  Consequently
\[
                 r\le D-2=12.
\]
Dominance requires \(r\ge12\), so equality holds.  There are exactly
two outside elements, each with moving dimension two, and all ten other
elements have moving dimension one.  An outside element of moving rank
two has two eigenvalues \(+1\) and two \(-1\), so it is an involution.
The ten rank-one elements are reflections in \(W\).  The determinant
formula just proved therefore forces splitness.

\textbf{Conclusion.} No dominant Weil-sixfold curve-cover family using
this natural \(G_F\)-representation can lie in a nonsplit determinant
component, for any initially allowed inertia or base genus.

\subsection*{What a nonsplit construction would have to change}

For these two representations, changing the product-one tuple, Hurwitz
component, or ordering cannot produce a dominant nonsplit source.  Using
higher-order inertia or higher-genus bases within the same representations
loses the necessary nine-dimensional parameter space.  A new dominant
nonsplit source must therefore change the representation or the geometric
source construction.

Formula \eqref{eq:det-seifert} identifies the arithmetic alternative for
another representation: either a local moving-space determinant or the
product \(\prod_i\det(1-g_i|_{U_i})\) must contribute a nonnorm class.
For higher-order inertia, the automatic supplier argument also needs
replacement or a new check.  Cyclic quotients may then have local
exponents other than \(m/2\); their exponents need not occur in opposite
pairs, and their primitive embeddings need not all be balanced.  Schoen's
paired-exponent determinant construction and the full Brauer-character
coverage cannot simply be carried over without verifying those conditions.
This is an exact missing requirement, not a claim that all higher-order
suppliers fail.

\section{Reflection normal closures and the remaining finite-group source gate}
\label{sec:reflection-reduction}

We work with a faithful absolutely irreducible representation over an
imaginary quadratic field \(K\), whose character field is exactly \(K\).
This is the split group-algebra block \(M_m(K)\) relevant to a generic
Weil factor with endomorphism algebra \(K\). The following argument is
not a classification of representations without reflections.

\subsection*{D-FIELD-01. The reflection components over the seven fields}

Let
\[
 \mathcal D=\{2,7,11,19,43,67,163\},\qquad K_d=\mathbb Q(\sqrt{-d}).
\]
An irreducible finite complex reflection group with a model over \(K_d\)
is a rational Weyl reflection representation, with the following sole
additional possibilities:
\[
  d=2:\ G_{12}\text{ in rank two};\qquad
  d=7:\ G_{24}\text{ in rank three}.
\]

Here we use the published reflection classification, rather than claim
to reprove it. Precise inputs are Cohen's rank-two classification and
table, \cite[\S\S3.4--3.6, pp.~394--395]{FIELD:Cohen}, and his
higher-rank classification, \cite[Theorem 5.12, p.~435]{FIELD:Cohen}.
The exceptional character fields are listed in
\cite[Tables 8.1--8.2, pp.~674--675]{FIELD:Malle}.
The table omits \(G_4,G_{25},G_{28},G_{35}\), as explicitly explained
on its following page. The first two have reflections of order three
(Cohen's \S3.6 and \S4.16 tables, pp.~395 and 412), so cannot have
models over these \(K_d\); the last two are the rational Weyl groups
\(F_4,E_6\). Inspection of the other listed fields leaves precisely
the two nonreal cases displayed above. Noncrystallographic real types
\(H_3,H_4\) have character field \(\mathbb Q(\sqrt5)\), and are excluded.

For completeness, the infinite imprimitive series needs no field table.
In \(G(a,b,k)\), \(k\ge2\), the diagonal element with entries
\(\zeta_a,\zeta_a^{-1},1,\ldots,1\) shows
\(2\cos(2\pi/a)\in K_d\cap\mathbb R=\mathbb Q\). Thus
\(a\in\{1,2,3,4,6\}\): the sum is a rational algebraic integer in
\([-2,2]\), and its five possible integer values give that list.
If \(k\ge3\), multiplying this diagonal element by the transposition
of the second and third coordinates gives trace
\(\zeta_a+k-3\), so \(\zeta_a\in K_d\) and \(a\le2\).
For \(k=2,b=a\), the group is the dihedral Weyl representation of
type \(A_2,B_2,G_2\), or a reducible lower case. For \(k=2,b<a\),
the element \(\operatorname{diag}(\zeta_a^b,1)\) forces
\(a/b\le2\). The remaining possibilities \((a,b)=(4,2),(6,3)\)
have traces \(2i\) and \(\zeta_6+\zeta_6^2=\sqrt{-3}\), respectively,
so are excluded. The cases \(a\le2\) are rational Weyl cases.
Rank one similarly has only the reflection \(-1\). The real standard
representations of symmetric groups are also rational Weyl cases.

\subsection*{D-FIELD-02. Rational normalizers and induced blocks}

\begin{lemma}\label{lem:weyl-normalizer}
Let \(W\) be an irreducible rational Weyl reflection group on
\(K_d^k\), with its standard rational positive definite form \(B\).
Let \(J\) be a finite subgroup of \(U(B,K_d)\) normalizing \(W\).
If \(d\in\{7,11,19,43,67,163\}\), every element of \(J\) has rational
matrix entries. If \(d=2\) and \(k\ge3\), either the same conclusion
holds or \(W\) is of type \(F_4\) and, up to the fixed rational
root coordinates, \(J\subset G_F\), the 2304-element group of D-GF4-01.
If \(J\) contains \(W\) and has a nonrational element in this last
case, then \(J=G_F\).
\end{lemma}

\begin{proof}
For \(g\in J\), conjugation by \(g\) is an automorphism of \(W\).
The equations
\[
                  A w=(g w g^{-1})A\qquad(w\in W)
\]
have rational coefficients. Their solution space over \(K_d\) is
one-dimensional, by absolute irreducibility. Hence it has a nonzero
rational solution \(A\), which is invertible, and \(g=cA\) for
some \(c\in K_d^*\). A finite order \(N\) of \(g\) gives
\(A^N=c^{-N}I\), so \(c^N\in\mathbb Q\). Consequently
\((c/\bar c)^N=1\). These fields have only the roots of unity
\(\pm1\), so either \(c\in\mathbb Q\) or
\(c\in\sqrt{-d}\,\mathbb Q\).

Unitarity implies
\[
                    A^tBA=\nu B,\qquad N(c)\nu=1,
                    \qquad \nu\in\mathbb Q_{>0}.
\]
The rational matrix \(A\) permutes the reflection root lines. For
one rational root \(\alpha\), write \(A\alpha=q\beta\), with
\(q\in\mathbb Q^*\) and \(\beta\) another standard root. Then
\[
             \nu=q^2\frac{B(\beta,\beta)}{B(\alpha,\alpha)}.
\]
An irreducible crystallographic root system has one root length or
two root lengths whose squared ratio is two or three. Thus the positive
rational square class of \(\nu\) is among \(1,2,3\).
If \(c\) is purely imaginary, its norm has square class \(d\).
The identity \(N(c)\nu=1\) therefore excludes that case for
\(d=7,11,19,43,67,163\), proving rationality of \(g\).

For \(d=2\), a nonrational \(g\) requires a root-length interchange
with squared length ratio two. In rank at least three the only
possibilities are \(B_k\) and \(F_4\). The two root-line conjugacy
classes in \(B_k\) have sizes \(k\) and \(k(k-1)\). A normalizer
cannot interchange them when \(k\ge3\), because their sizes differ.
Hence only \(F_4\) remains.

Use the standard \(F_4\) roots in \(\mathbb Q^4\) and the matrix
\(\Phi\) of D-GF4-01. The orthogonal matrix
\(\Phi/\sqrt2\) interchanges the two root-line classes, as was
checked directly on all 48 roots. An orthogonal normalizer preserving
the two classes preserves the actual root system. Such an automorphism
belongs to \(W(F_4)\): it sends a Weyl chamber to a Weyl chamber,
so after composition with \(W(F_4)\) it fixes a chamber and induces
an automorphism of its Cartan diagram; the oriented \(F_4\) diagram
has no nonidentity automorphism. It follows that a rational similarity
normalizing \(W(F_4)\) is a rational scalar times either a Weyl
element or a Weyl element times \(\Phi\).

In the first case a finite unitary element over \(K_2\) is just a
Weyl element, since its scalar is a root of unity in \(K_2\).
In the second case write it as \(\lambda w\Phi\). Unitarity gives
\(2N(\lambda)=1\), and the preceding finite-order argument gives
that \(\lambda\) is rational or purely imaginary. The rational
case would require a rational square equal to \(1/2\); hence
\(\lambda=\pm\sqrt{-2}/2\). Thus every such element lies in
\(W(F_4)z\), and \(J\subset W(F_4)\sqcup W(F_4)z=G_F\).
If \(J\) contains \(W(F_4)\) and one element of the second coset,
it contains both whole cosets and equals \(G_F\).
\end{proof}

Now let \(G\subset\operatorname{GL}_m(K_d)\) be finite and absolutely
irreducible, with exact character field \(K_d\), and suppose it contains
a reflection. Let \(R\) be the normal subgroup generated by all its
reflections. The fixed space \(V^R\) is \(G\)-stable, hence zero.
The nonorthogonality graph of the reflection root lines gives an
orthogonal decomposition
\[
                  V=V_1\oplus\cdots\oplus V_s
\]
into the irreducible reflection components of \(R\). Each root line
is defined over \(K_d\), being the image of a rank-one matrix
\(r-1\), so each \(V_i\) is defined over \(K_d\). The group \(G\)
permutes these components transitively, by irreducibility. If \(H\)
stabilizes \(V_1\), and \(\tau\) is its representation on \(V_1\),
then
\[
                         \rho=\operatorname{Ind}_H^G\tau.
\]
The image of \(H\) is finite, contains the irreducible reflection
group \(R_1\), and normalizes it. The character of \(\tau\) cannot
be rational-valued, since otherwise its induced character would be
rational-valued.

If \(R_1\) is a rational Weyl representation, a \(K_d\)-change of
basis puts it into its rational standard model. Such a change of basis
exists because the intertwining equations are defined over \(K_d\)
and the representations are absolutely irreducible. Its invariant
Hermitian form is a positive rational scalar times the standard rational
form. Lemma~\ref{lem:weyl-normalizer} therefore applies to the image of
\(H\).

We obtain the following conclusions.
\begin{enumerate}
\item For \(d=11,19,43,67,163\), the existence of a reflection
is impossible. Every component is a rational Weyl component by
D-FIELD-01, and the lemma would make \(\tau\), hence \(\rho\),
rational-valued.
\item For \(d=7\), the component must be \(G_{24}\) of rank three.
Its finite unitary normalizing overgroups are contained in scalar multiples of
\(G_{24}\), by Cohen's corrected Proposition 5.14
\cite[p.~435]{FIELD:Cohen}\cite[p.~613]{FIELD:CohenErr}.
The unique exception in that corrected proposition is \(G_{25}\),
which is irrelevant here. A scalar multiplier belonging to \(K_7\)
and giving a finite overgroup is \(\pm1\); both scalars already
belong to \(G_{24}\). Thus the image of \(H\) is exactly \(G_{24}\).
\item For \(d=2\), either \(\dim_{K_2}V_1=2\), or the image of
\(H\) is exactly \(G_F\) in rank four. This uses D-FIELD-01 for
the nonreal reflection case \(G_{12}\), and the lemma for every
real Weyl case of rank at least three. Rank one cannot have a
nonrational finite character over \(K_2\).
\end{enumerate}

\subsection*{D-FIELD-03. Consequences for dominant sixfold curve suppliers}

Suppose a family of connected finite \(G\)-covers produces a Weil
sixfold as the multiplicity factor for \(\rho\). For the block
stabilizer above, put \(Y=C/H\). Shapiro's isomorphism identifies
the \(\rho\)-multiplicity Hodge structure in \(H^1(C)\) with the
\(\tau\)-multiplicity Hodge structure for the \(H\)-cover \(C\to Y\).
One can see this directly by assigning a vector in each induced coset
summand to the corresponding lift of a singular simplex; the boundary
maps and Hodge de Rham complexes agree. The cup pairings agree with
the sum over the sheets, using the induced invariant Hermitian form.
The originally chosen irreducible invariant form differs from it by a
positive rational scalar. Thus this is an identification of polarized
rational variations up to a fixed rational scalar, and preserves the
signed determinant class in even dimension. Quotienting the kernel
of \(\tau\) gives a connected cover with the indicated faithful
image group. In particular the period map factors through the appropriate
Hurwitz space of these covers of \(Y\).

For a rank-\(k\) local system on a base of genus \(\gamma\), with
moving-dimension sum \(D\), a sixfold multiplicity space satisfies
\[
                       6=(2\gamma-2)k+D.
\]
There are at most \(D\) genuine branch points. In rank two this
gives \(D=10-4\gamma\). The possible Hurwitz dimensions are at
most seven for \(\gamma=0\), six for \(\gamma=1\), and five for
\(\gamma=2\); higher genera are impossible. Consequently an induced
rank-two source cannot have period rank nine.

For \(K_7\), D-FIELD-02 reduces every source containing a reflection
to the natural \(G_{24}\) representation on the intermediate base
\(Y\). D-DET-02 proves that any rank-nine image of this family has
split determinant. For \(K_2\), rank two is excluded by the preceding
dimension bound; the remaining possibility is the natural \(G_F\)
representation, and D-DET-03 again forces split determinant.

\textbf{Universal reflection-source conclusion.}
Over \(\mathbb Q(\sqrt{-2})\) or \(\mathbb Q(\sqrt{-7})\), every
dominant sixfold finite-group curve supplier whose faithful absolutely
irreducible representation has exact character field \(K\) and contains
a reflection lies in the split determinant component. This statement
allows arbitrary representation dimension, finite group, initial branch
orders, and base genus within that class.

For the five fields \(K_d\), \(d=11,19,43,67,163\), such a
representation contains no reflection at all. A possible dominant
source over those fields, or a nonsplit dominant source over \(K_2,K_7\),
must therefore use a representation without reflections. If its degree
is \(m\), each nonidentity inertia has moving dimension at least two.
A genus-one base has \(D=6\) and at most three branch points, hence
cannot give rank nine. On a base of genus at least two, the dimension
equation and \(m\ge2\) leave only \(\gamma=2,m\le3\), whose
Hurwitz dimension is less than nine. Thus the base must have genus zero,
and
\[
                 12\le r\le\frac{2m+6}{2}=m+3,
                 \qquad m\ge9.
\]
More generally, if every nonidentity element has moving rank at least
\(q\), the necessary inequality is \(m\ge6q-3\).

These are necessary conditions only. They neither produce a rank-nine
period map nor exclude all representations without reflections. The
remaining finite-group gate is precisely a genus-zero representation
of degree at least nine without reflections, with an actual rank-nine
period computation and a nonnorm value in D-DET-01 for a nonsplit target.

\section{A universal monomial obstruction for nonsplit sixfold sources}
\label{sec:monomial-sixfold}

Here ``monomial'' refers to induction of a complex representation from a
one-dimensional character. It does not require a reflection group or
reflection inertia. We allow the values of the inducing character to
generate a cyclotomic field strictly larger than the quadratic character
field. The conclusion concerns primitive factors of the rational group
algebra; it does not concern an unrelated quadratic action imposed on a
factor whose character field is real.

\subsection*{D-MON-01. The cyclic factor, including its Schur index}

Let $G$ act faithfully on a connected smooth projective curve $C$, and let
$\rho$ be an irreducible complex representation with character field an
imaginary quadratic field $K$. Suppose
\[
             \rho=\operatorname{Ind}_H^G\chi,
             \qquad d=\operatorname{ord}(\chi),
             \qquad L=\mathbf Q(\zeta_d).
\]
Let $B$ be the abelian factor selected by a primitive idempotent in the
corresponding simple component of $\mathbf Q[G]$. Assume that $B$ is a
sixfold with $K$-signature $(3,3)$, with the polarization induced from
$J(C)$. The character formula gives $K\subset L$, and $d>2$ because a
real-valued linear character would induce a rational-valued character.

Write that simple component as $M_k(D)$, where $D$ is a central division
algebra over $K$, and let $s^2=\dim_KD$ be its Schur-index square. The
corner algebra of a primitive idempotent acts faithfully on $H^1(B,\mathbf Q)$.
Consequently this cohomology is a module over $D$, or its opposite according
to the cohomological convention, and
\[
                     s^2\mid\dim_KH^1(B,\mathbf Q)=6.
\]
The only possible positive integer $s$ is one. Thus $\rho$ has a
$K$-realization $V$, its component is $M_{\dim\rho}(K)$, and the
multiplicity Hodge structure
\[
     M=\operatorname{Hom}_{\mathbf Q[G]}(V_{\mathbf Q},H^1(C,\mathbf Q))
\]
is the rational Hodge structure of $B$. Here $V_{\mathbf Q}$ means
restriction of scalars from $K$; its commuting algebra is $K$.

Put $J=\ker\chi$, $Y=C/J$ and $X=C/H$. Then $Y\to X$ is a connected
cyclic cover of degree $d$. Let $P_d(Y/X)$ denote the primitive
cyclotomic factor of $J(Y)$, namely the image of the rational idempotent
for $\mathbf Q(\zeta_d)$ in $\mathbf Q[H/J]$. There is an isogeny
\begin{equation}\label{eq:mon-power}
           P_d(Y/X)\ \sim\ B^{\ell},\qquad \ell=[L:K].
\end{equation}
Here is the rational representation argument, including the multiplicity.
Regard $L$ as the rational $H$-module afforded by $\chi$. Over $\mathbf C$,
the induction of this module is the sum of $\operatorname{Ind}_H^G\chi^\sigma$
over all embeddings $\sigma:L\hookrightarrow\mathbf C$. Every summand is
$\rho$ or $\bar\rho$, according to the restriction of $\sigma$ to $K$,
and each occurs $[L:K]$ times. Since the Schur index is one, semisimplicity
over $\mathbf Q$ therefore gives
\[
             \operatorname{Ind}_H^G L\cong V_{\mathbf Q}^{\oplus\ell}.
\]
Frobenius reciprocity identifies its Hom into $H^1(C,\mathbf Q)$ with
$\operatorname{Hom}_H(L,H^1(C,\mathbf Q))$. The latter is precisely
the primitive cyclotomic Hodge structure of $Y$: first take $J$-invariants,
then the primitive $\mathbf Q[H/J]$ summand. This proves
\eqref{eq:mon-power}. These constructions use fixed rational group-algebra
maps, so in a family they give an isomorphism of rational variations and,
after clearing fixed denominators, an isogeny of abelian families.

Write $h=g(X)$ and let $R$ be the number of branch points of $Y\to X$.
The faithful character $\chi$ has no invariants at any genuine branch
point. Its parabolic Euler characteristic and its $(3,3)$ multiplicities
give
\begin{equation}\label{eq:mon-dimension}
                  6=2h-2+R,\qquad R=8-2h.
\end{equation}
One can also obtain this identity by adding the two character
Chevalley--Weil formulas: each nonzero local exponent and its conjugate
contribute one. In particular $0\le h\le4$.

For fixed cyclic branch exponents, the moduli of these cyclic covers are
finite over the moduli of the $R$-pointed base curve. Indeed the cover is
specified by a $d$th root of the appropriate branch line bundle, and roots,
when they exist, form a finite torsor under $J(X)[d]$. Using
\eqref{eq:mon-dimension}, their dimensions for $h=0,1,2,3,4$ are at most
\[
                              5,6,7,8,9,
\]
respectively. The genus-zero and genus-one formulas here use $R=8$ and
$R=6$, so no unstable pointed-curve exception occurs.

The period differential of $B$ has the same kernel as that of $B^\ell$:
on the Hodge bundles the latter is the direct sum of $\ell$ copies of the
former. Isogeny \eqref{eq:mon-power} therefore bounds its rank by the
rank of the forgetful map to cyclic-cover moduli. This is a statement
about actual variations, not a sufficiency claim based on the source
dimension. If the $B$-period map has rank nine, necessarily
\begin{equation}\label{eq:mon-equality}
                  h=4,\qquad R=0,
\end{equation}
and the forgetful image is dominant in a component of the moduli of
connected etale cyclic degree-$d$ covers of genus-four curves. After
passing to finite level structures these are algebraic varieties; a
dominant constructible image contains a nonempty Zariski open set.
No assumption about simple branching of $C\to C/G$ is used here.

\subsection*{D-MON-02. The full endomorphism algebra of a general cyclic Prym}

For a connected etale cyclic degree-$d$ cover of a genus-four curve, the
primitive integral homology, modulo the harmless choice of a finite-index
lattice, is a rank-six module over $R_d=\mathbf Z[\zeta_d]$.
Looijenga's Proposition~2.2 identifies its equivariant intersection form
with the standard skew-Hermitian hyperbolic form
\[
               \Omega=\begin{pmatrix}0&I_3\\-I_3&0\end{pmatrix}.
\]
His Theorem~2.4 says that the group of lifts of mapping classes fixing
the cyclic quotient has image
\[
 U_6^\#(R_d)=
 \{g\in\operatorname{GL}_6(R_d):g^*\Omega g=\Omega,
                       \ \det g\in\langle\zeta_d^2\rangle\}.
\]
Both statements occur on page~5 of the author's version of
\cite{IM:Looijenga1997}; the hyperbolic basis is proved on pages~8--9,
and the genus-at-least-three monodromy proof ends on page~12.
We use the published theorem, and do not claim a new proof of its full
mapping-class surjectivity. A finite level cover has monodromy of finite
index in this image, which is all that is needed below.

We first prove the precise commutant assertion. For every finite-index
subgroup $\Gamma\subset U_6^\#(R_d)$,
\begin{equation}\label{eq:mon-commutant}
               \operatorname{End}_{\mathbf Q,\Gamma}(L^6)=L.
\end{equation}
Indeed $U_6^\#(R_d)$ contains $\operatorname{Sp}_6(\mathbf Z)$.
For every integral symplectic transvection $1+N$, some positive power
belongs to $\Gamma$. As $N^2=0$, commuting with that power is equivalent
over $\mathbf Q$ to commuting with $N$. Thus a commuting endomorphism
commutes with all integral symplectic transvections. The transvections
in the coordinate vectors give the elementary matrices
$E_{i,i+3},E_{i+3,i}$, whose products give the coordinate projectors.
Transvections in sums of coordinate vectors identify the resulting
coordinate endomorphisms. Hence the commutant on $L^6$, viewed as a
rational vector space, consists of operators
$\operatorname{id}_{\mathbf Q^6}\otimes f$ with
$f\in\operatorname{End}_{\mathbf Q}(L)$.

Now let $B$ be the Hermitian three-by-three matrix with
$B_{12}=\zeta_d$, $B_{21}=\bar\zeta_d$ and all other entries zero.
The matrix
\[
                          u_B=\begin{pmatrix}I_3&B\\0&I_3\end{pmatrix}
\]
belongs to $U_6^\#(R_d)$ and is unipotent with square-zero logarithm.
Some positive power belongs to $\Gamma$. Commutation with that power
forces $f(\zeta_d x)=\zeta_d f(x)$ for every $x\in L$. Thus $f$ is
$L$-linear and is multiplication by $f(1)$. This proves
\eqref{eq:mon-commutant}, without invoking an unproved assertion about
Zariski closure or independence of Galois embeddings.

It follows that the primitive abelian variety has full rational
endomorphism algebra $L$ outside a countable union of proper analytic
loci. We give the argument because a statement only about endomorphisms
commuting with the cyclic group would not suffice for our application.
On the connected universal cover of a finite-level moduli component,
trivialize rational homology. For each of its countably many rational
endomorphisms, the condition of preserving the Hodge subbundle is closed
analytic. Outside the union of those loci which are proper, the Hodge
endomorphisms are exactly those preserving the Hodge structure everywhere
on this universal cover. This rational algebra $E$ is stable under
monodromy conjugation, and its intersection with integral endomorphisms
is a lattice. The polarization gives the positive definite Rosati norm
\[
                         f\longmapsto\operatorname{Tr}(f^\dagger f)
\]
on this lattice. Monodromy conjugation preserves the lattice and this
norm, so its action on $E$ is finite. A finite-index monodromy subgroup
therefore centralizes $E$. Equation \eqref{eq:mon-commutant} gives
$E\subset L$; the reverse inclusion is the cyclic deck action. Hence
$E=L$. The complement of the countably many proper analytic loci meets
every nonempty analytic open set, by the Baire theorem. This is the only
genericity consequence needed below; an algebraicity theorem for Hodge
loci is not needed.

\subsection*{D-MON-03. The universal monomial sixfold barrier}

Under the hypotheses of D-MON-01, a rank-nine period map forces
\begin{equation}\label{eq:mon-final}
                 L=K,\qquad d\in\{3,4,6\},
                 \qquad\text{the polarization Hermitian space is split.}
\end{equation}
In particular there is no dominant nonsplit Weil-sixfold curve-cover
source from any such monomial representation. For a quadratic field
other than $\mathbf Q(i)$ and $\mathbf Q(\sqrt{-3})$, there is no
dominant source of this kind even in the split component.

To prove this, rank nine first gives \eqref{eq:mon-equality} and a
nonempty open image in the etale cyclic genus-four moduli. Choose a
point of this open image at which D-MON-02 gives
$\operatorname{End}^0(P_d)=L$. If $\ell=[L:K]>1$, the isogeny
$P_d\sim B^\ell$ gives an embedding
\[
                       M_\ell(K)\hookrightarrow\operatorname{End}^0(P_d).
\]
The left side is noncommutative and the right side is a field, a
contradiction. Thus $L=K$, so the cyclotomic degree is two and
$d=3,4,6$.

For completeness the final splitness concerns the actual induced
polarization. When $L=K$, the rational element
\[
 p_\chi=\frac1{|H|}\sum_{h\in H}
                   \bigl(\chi(h)^{-1}+\bar\chi(h)^{-1}\bigr)h
                   \in\mathbf Q[G]
\]
is a Rosati-fixed primitive idempotent of the $\rho$ component.
Frobenius reciprocity gives multiplicity one of $\chi$ in
$\rho|_H$, and induction shows that no other complex irreducible
representation contains $\chi$. Thus this idempotent selects exactly
one $K$-copy. Its abelian image is the pullback of $P_d(Y/X)$ along
$C\to Y$; the polarization pulls back to $|J|$ times the restriction
of the Jacobian polarization on $P_d$.

Looijenga's equivariant hyperbolic basis, or the cut-surface construction
in its proof, says that the latter skew-Hermitian space is the sum of
three hyperbolic planes. Multiplication by a nonzero imaginary scalar
converts it to a Hermitian form with the same half-dimensional isotropic
subspace. The actual polarization Hermitian form is therefore split.
Any other Rosati-fixed primitive factor in the same component has the
same multiplicity Hermitian form multiplied by the positive rational
norm of a line in the positive definite representation space. This
preserves splitness. Equivalently, restrict the tensor-product
polarization on the isotypic decomposition to either primitive line.
The raw determinant class is $[-1]$ in dimension six, or the trivial
class in the signed-discriminant convention.

This conclusion allows arbitrary initial inertia orders, arbitrary base
genus for $C/G$, arbitrary Nielsen class, and every monomial induction
whose character field is $K$. Its necessary equality condition was
deduced from the actual cyclic variation. No forgetful submersion for
fixed higher ramification was assumed. The remaining finite-group
nonsplit-source gate must use a nonmonomial irreducible representation,
or must change the specified group-algebra/character-field setup.

\section{The Gaussian groups \(G_{29},G_{31}\): exact determinant obstruction}
\label{sec:g31-barrier}

The two natural rank-four Gaussian reflection representations do not
supply a dominant nonsplit Weil-sixfold family. This conclusion is based
on an exhaustive exact calculation from published matrices, followed by
D-DET-01; it is not a conclusion from a numerical period experiment.

\subsection*{D-G31-01. Published generators and exhaustive local determinants}

Put \(K=\mathbb Q(i)\), and use the standard Hermitian form on \(K^4\).
Davis \cite[Theorem 4.1, proof, p.~307]{G34:Davis} gives the four
generators of \(G_{29}\):
\[
s_1=\frac12\begin{pmatrix}
1&-1&-1&-1\\-1&1&-1&-1\\-1&-1&1&-1\\-1&-1&-1&1
\end{pmatrix},\quad
s_2=\begin{pmatrix}0&i&0&0\\-i&0&0&0\\0&0&1&0\\0&0&0&1\end{pmatrix},
\]
\[
s_3=\begin{pmatrix}0&1&0&0\\1&0&0&0\\0&0&1&0\\0&0&0&1\end{pmatrix},\qquad
s_4=\begin{pmatrix}1&0&0&0\\0&0&1&0\\0&1&0&0\\0&0&0&1\end{pmatrix}.
\]
His Theorem 4.16, p.~311, adds
\(s_5=\operatorname{diag}(1,1,-1,1)\) to generate \(G_{31}\).
All five are self-adjoint involutive reflections.

For an element \(g\), let \(U_g=\operatorname{im}(1-g)\),
\(k(g)=\dim_KU_g\), and choose a \(K\)-basis of \(U_g\). Define
\[
b(g)=\det(1-g|_{U_g}),\qquad
q_+(g)=\frac{\det(h|_{U_g})}{|b(g)|}.
\]
Since all group determinants are \(\pm1\),
\(\overline{b(g)}=(-1)^{k(g)}\det(g)^{-1}b(g)\), so \(b(g)\)
is rational or purely imaginary rational. Hence \(q_+(g)\) is a
positive rational number. Its class modulo Gaussian norms is independent
of the chosen moving-space basis and of conjugation in the group.

The complete exact enumeration gives:
\[
\begin{array}{c|r|r|r|r}
&|G|&\text{conjugacy classes}&\text{reflections}&
 \text{nonidentity involutions}\\\hline
G_{29}&7680&37&40&231\\
G_{31}&46080&59&60&511
\end{array}
\]
For \(G_{29}\), every local \(q_+\) is a Gaussian norm. For
\(G_{31}\), every class of moving rank at most three, and every
involution class, has norm \(q_+\). The three exceptions are:
\[
\begin{array}{c|r|r|c|c|c|c}
\text{certificate index}&\text{class size}&\text{order}&k(g)&b(g)&q_+(g)
 &\operatorname{age}(g)-k(g)/2\\\hline
41&960&12&4&3i&3&1/2\\
43&960&12&4&-3i&3&-1/2\\
46&960&6&4&3&3&0
\end{array}
\]
The indices are those of our enumeration, not imported character-table
labels. In the selected bases all three moving Gram determinants equal
nine. Their nonzero eigenvalue exponents are respectively
\((4,7,8,11)/12\), \((1,4,5,8)/12\), and \((1,2,4,5)/6\).

Here is a transparent specification of the exhaustive certificate.
\texttt{enumerate\_gaussian\_groups.cpp} stores twice each matrix as
32 signed integers, the real and imaginary entries in row order. Every
product divides its integer numerator by two, with exact divisibility
checked. Breadth-first multiplication by the four or five displayed
generators enumerates the entire generated group. All matrices satisfy
\(g^*g=I\) exactly. Breadth-first conjugation by the same generators
partitions it into complete conjugacy classes; the class sizes sum to
the full enumerated order. Thus the class list is exhaustive.

\texttt{check\_local\_discriminants.py} independently performs rational
Gaussian elimination on each class representative. It takes independent
columns of \(1-g\), computes the induced moving-space matrix, its Gram
determinant, and \(b(g)\). It tests the Gaussian norm criterion by
factoring the positive rational number \(q_+\): each prime congruent
to three modulo four must have even valuation. Eigenvalue multiplicities
are computed by the exact finite Fourier formula in
\(\mathbb Q[\zeta_{\operatorname{lcm}(4,\operatorname{ord}g)}]\), reducing
polynomials by the cyclotomic polynomial, not by floating-point roots.
Every multiplicity is checked to be a nonnegative integer.

The files \texttt{G29\_classes.json}, \texttt{G31\_classes.json}, and
\texttt{G29\_local\_discriminants.json},
\texttt{G31\_local\_discriminants.json} retain every representative,
class size, order, local matrix determinant, Gram determinant, spectrum
and norm test. Together with the short source code, these constitute
the full finite computation certificate used here.

\subsection*{D-G31-02. The universal nonsplit-source obstruction}

Consider a genus-zero product-one generating tuple with balanced
parabolic multiplicity space. D-DET-01 gives its signed determinant
class as
\[
                  \left[\prod_j\frac{\det(h|_{U_{g_j}})}{b(g_j)}\right].
\]
The displayed representative is positive: it differs from the determinant
of the anti-Hermitian parabolic cup matrix by a positive norm, and for
a balanced Hermitian space that anti-Hermitian determinant is positive.
Since each factor is real or purely imaginary, the positive representative
equals the product of its absolute values. Therefore its norm class is
\[
                              \left[\prod_jq_+(g_j)\right].
\]
For \(G_{29}\) this is always a norm. Thus a balanced genus-zero
\(G_{29}\) source is necessarily split.

For \(G_{31}\), a nonsplit balanced source must contain an odd number
of the three exceptional classes, in particular at least one inertia
element of moving rank four. A sixfold has total moving rank
\(D=2\cdot4+6=14\). If its branch count is \(r\), that rank-four
entry and the other nonidentity entries give \(D\ge4+(r-1)\), hence
\[
                              r\le11.
\]
Its Hurwitz dimension and period-image dimension are consequently at most
eight. This is a necessary-dimension obstruction, not a claim that the
upper bound is attained by every such tuple.

Allowing higher-genus bases does not restore dominance. For a base of
genus \(\gamma\), the same sixfold condition is
\(D=14-8\gamma\). At genus one there are at most six branch points
and at most six parameters; at larger genus the dimension equation is
impossible. Therefore neither natural representation gives a dominant
nonsplit sixfold curve source, for arbitrary initial inertia orders,
Nielsen class, ordering, or base genus.

\subsection*{D-G31-03. An explicit nonsplit tuple and its rank-two boundary}

The determinant obstruction is visible in a concrete connected family.
Besides the five published generators define
\[
u=\frac12\begin{pmatrix}
1&-i&-1&i\\i&1&i&1\\-1&-i&1&i\\-i&1&-i&1
\end{pmatrix},\qquad
c=\frac12\begin{pmatrix}
0&1-i&-1-i&0\\-1+i&0&0&1+i\\
1-i&0&0&1+i\\0&-1+i&-1-i&0
\end{pmatrix}.
\]
Exact multiplication gives
\[
          c=(s_1s_4s_5s_3s_2u)^{-1},\qquad
          c^6=1,\quad \det(1-c)=3.
\]
The matrix \(u\) is a reflection in the conjugacy orbit of \(s_5\).
The eigenvalues of \(c\) have exponents \((1,2,4,5)/6\), so its
age is two and its moving rank is four. The eleven-entry tuple
\[
       (s_1,s_1,s_2,s_2,s_1,s_4,s_5,s_3,s_2,u,c)
\]
has product one and generates \(G_{31}\), since its tail includes
all five published generators. Riemann existence gives connected covers;
their branch orders are \((2^{10},6)\) and their genus is
\[
           1+\frac{46080}{2}\left(-2+10\cdot\frac12+\frac56\right)
                       =88321.
\]
Chevalley--Weil gives signature \((3,3)\): each reflection contributes
age \(1/2\), and \(c\) and its inverse each contribute age two.
The exact parabolic quotient has relation and coboundary ranks four,
dimension six, and anti-Hermitian cup determinant \(3/8\). Hence its
signed discriminant has Gaussian norm class three, which is nonsplit.
The source has dimension eight; its exact period rank has not been
computed and is not asserted to be eight.

There is a useful degeneration of precisely the form sought for a
nonsplit sixfold. Bubble off the first two equal pairs separately.
The two new base nodes have trivial monodromy. The central normalization
has branch tuple
\((s_1,s_4,s_5,s_3,s_2,u,c)\), still generating the full group;
its natural multiplicity space has dimension two, signature \((1,1)\),
and exact cup determinant \(3/2\). It is an anisotropic Hermitian plane,
since three is not a Gaussian norm. The bubble components are rational
cyclic double covers.

The dual graph has one central vertex, bubble vertices
\(G/\langle s_1\rangle\) and \(G/\langle s_2\rangle\), and two
copies of \(G\) as edge sets. Its natural-representation graph
multiplicity is
\[
             2\dim V-\dim V^{s_1}-\dim V^{s_2}=8-3-3=2.
\]
The vanishing circles satisfy exactly the vertex-incidence relations;
this follows by Mayer--Vietoris on the punctured normalization
components. A positive smoothing applies a positive Dehn twist at
each node. The sum of their rank-one transvection forms has image the
full span of those circles, since its positive quadratic form is
nondegenerate on that span. Taking the exact natural multiplicity
functor gives a logarithm of rank two. The weight-one part of the
normalization is the anisotropic plane just computed. Thus the
degeneration has graph rank two and an anisotropic residual Weil surface,
with no fictitious maximal isotropic rank-three limit.

\texttt{find\_nonsplit\_tuple.py} reconstructs the reflection orbit,
verifies the product, generation witness, exact spectrum and both
parabolic Gram matrices, and writes \texttt{nonsplit\_tuple.json}.
This establishes the source and its determinant, not algebraicity of
its sixfold Weil classes: the higher-order inertia requires a new
supplier/character-coverage argument, and the source is too small to
dominate the nine-dimensional nonsplit moduli space.

\section{Gaussian scalar extensions of the Weyl groups \(E_6,E_7\)}
\label{sec:e6-e7-gaussian}

\subsection*{D-E67-01. Exact local data from the published class representatives}

Let \(K=\mathbb Q(i)\), and let \(W_m=W(E_m)\), for \(m=6,7\),
act in its rational simple-root basis. Its invariant positive form is
the Cartan matrix \(H_m\): diagonal entries are two, and its edges
are
\[
 (1,3),(3,4),(2,4),(4,5),(5,6),
 \quad\text{and additionally }(6,7)\text{ for }m=7.
\]
Thus \(\det H_6=3\), \(\det H_7=2\), and the simple reflections
are explicitly
\[
                   s_j=I-e_j(H_m)_{j,*}.
\]
We consider \(J_m=\mu_4W_m\). All statements in the following finite
calculation use exact rational Gaussian arithmetic.

Carter's published classification gives 25 classes for \(W_6\) and
60 for \(W_7\), with their names and sizes in
\cite[Table 9, p.~50; Table 10, pp.~51--52]{WG:Carter}.
The author's CHEVIE library supplies words in the simple reflections for
these classes. The precise data are \texttt{ClassInfo.classtext} in
\texttt{weyle6.g}, and \texttt{WordsClassRepresentatives} in
\texttt{weyle7.g}; CHEVIE's manual, \S87.1, documents these as
complete conjugacy-class representatives. We use this published
classification and documented data for completeness. We do not claim
an independent enumeration of all 2,903,040 elements of \(W_7\).

The compact retained input files contain only the Cartan matrices, class
names, words, orders and sizes, together with the URL and SHA256 of each
original source. The checker reconstructs each matrix from the displayed
simple reflections, verifies its exact order and orthogonality, and
computes each of \(w,iw,-w,-iw\). For each such element \(g\), it
takes independent columns of \(1-g\) as a basis of \(U_g\), and
computes
\[
       d_g=\det(H_m|_{U_g}),\qquad
       b_g=\det(1-g|_{U_g}),\qquad k_g=\dim_K U_g.
\]
Finite Fourier inversion modulo the appropriate cyclotomic polynomial
gives exact eigenvalue multiplicities and ages.

In the rank-seven case, \(b_g\) may be a nonzero rational multiple
of \(1+i\) or \(1-i\), as well as a real or imaginary rational
number. Write
\[
              |b_g|=c_g\,2^{\epsilon_g/2},
           \qquad c_g\in\mathbb Q_{>0},\quad\epsilon_g\in\{0,1\},
           \qquad q_g=d_g/c_g.
\]
This gives a rational local norm class. In a product-one tuple of even
total moving rank, \(\prod b_g\) is rational by D-DET-01. Its phase
then forces \(\sum\epsilon_g\) to be even. The removed factor
\(2^{\sum\epsilon_g/2}\) is a Gaussian norm. For a balanced tuple
the sign of the product is positive, also by D-DET-01. Therefore the
signed determinant class of the parabolic space is precisely
\[
                              \left[\prod q_g\right].
\]
Changing moving-space bases multiplies each \(q_g\) by a Gaussian norm.

The complete calculations give the following particularly small table.
Here ``nonreal'' means \(g\in i\mathbb R^{m\times m}\).
\[
\begin{array}{c|c|c|c|c}
m&\min k_g\text{ for nonreal }g&
 \text{underlying Carter classes attaining it}&q_g&
 \text{age}(iw)-k_g/2\\\hline
6&4&D_4(a_1)&3&-1/2\\
7&5&D_4(a_1)&8&-3/4\\
7&5&D_4(a_1)+2A_1&8&1/4\\
7&5&2A_3+A_1&4&3/4\\
7&5&D_4(a_1)+A_1&8&-1/4
\end{array}
\]
Complex conjugation changes the displayed imbalance to its negative
and leaves the local norm class unchanged. All reflections in \(J_m\)
have norm local invariant and real entries. In rank six the sole
minimal nonreal class is represented by the particularly short word
\[
                  w=s_3s_4s_3s_2s_4s_5.
\]
Its characteristic polynomial is \((x^2+1)^2(x-1)^2\).
The rank-seven words and every corresponding moving basis and Gram
matrix are retained in the exact certificate.

The reproducible files are \texttt{check\_e6\_e7.py},
\texttt{E6\_witness\_input.json}, \texttt{E7\_witness\_input.json},
\texttt{E6\_local\_invariants.json}, and
\texttt{E7\_local\_invariants.json}. The checker shares only the
elementary rational Gaussian arithmetic helper with the G31 certificate;
it imports no numerical invariant or conclusion from that calculation.

\subsection*{D-E67-02. Every dominant sixfold source is split}

First suppose the base has genus zero. A sixfold natural multiplicity
space has total moving dimension
\[
                              D=2m+6.
\]
Period rank nine requires at least twelve branch points. Every inertia
is nonidentity, so the sum of the excesses \(k_g-1\) is
\(D-r\le2m-6\).

A generating tuple for \(J_m\) contains a nonreal entry. The
homomorphism recording whether its matrix is real or purely imaginary
has values in \(\mathbb Z/2\), and the product-one relation makes
the number of nonreal entries even. Thus it contains at least two.
For \(m=6\), each contributes excess at least three; for \(m=7\),
each contributes excess at least four. In either case these two entries
already exhaust the entire permitted excess. It follows that
\[
 r=12,\quad\text{ten entries are reflections, and the other two have}
 \quad k_g=4\ (m=6),\quad k_g=5\ (m=7).
\]
The table gives local classes three and three for the latter two entries
in rank six. Their product is nine, a Gaussian norm. In rank seven each
local class is itself a Gaussian norm. The ten reflections contribute
only norms. D-DET-01, with the phase normalization explained above,
therefore proves split signed determinant in both cases. This is a
necessary-dimension obstruction combined with an exact determinant
calculation, not a claim that these branch data have maximal period rank.

For a base of genus \(\gamma\), the same dimension equation is
\(D=2m+6-2m\gamma\). If \(\gamma=1\), then \(D=6\), so there
are at most six branch points and at most six moduli parameters.
If \(\gamma\ge2\), \(D<0\) for \(m=6,7\), which is impossible.
Thus higher-genus bases cannot restore dominance.

The conclusion also applies to a finite group whose irreducible
reflection normal-closure component is of type \(E_6\) or \(E_7\)
over \(K\). Indeed its stabilizer acts by a finite unitary normalizer
of that Weyl group. The rational-intertwiner argument of D-FIELD-02
writes such a normalizer as a scalar times a rational similarity.
All roots have the same length, so after rational rescaling the similarity
is orthogonal and permutes the roots. An automorphism of the root system
is a Weyl element followed by a diagram automorphism: move its chamber
back by the simply transitive Weyl action and inspect the simple roots.
For \(E_7\) the diagram has no automorphism, and for \(E_6\) its
involution is induced by \(-w_0\), where \(w_0\) is the longest
Weyl element. Hence the normalizer lies in scalar multiples of the
Weyl group. A scalar in \(K\) giving a finite group is in \(\mu_4\).
If the stabilizer character is nonrational it contains a nonreal such
element and, since it contains the Weyl group, it contains \(iI\);
its image is exactly \(J_m\).

Induction from the block stabilizer identifies the multiplicity
variation and actual polarization on the intermediate quotient curve,
as proved in D-FIELD-03. Applying the preceding obstruction there shows
that no such induced finite-group source dominates a nonsplit Weil-sixfold
component. This result concerns \(E_6,E_7\); it makes no assertion
that they exhaust the Gaussian reflection-normal-closure possibilities.

\section{The remaining Gaussian reflection-component cases}
\label{sec:gaussian-reflection-reduction}

Throughout, the original finite-group representation is faithful,
absolutely irreducible over \(K=\mathbb Q(i)\), and has exact
character field \(K\). We study the actual primitive-factor
polarization, and a sixfold of signature \((3,3)\).

\subsection*{D-GAUSS-01. The Gaussian \(F_4\) normalizer and \(D_4\) triality}

Use the standard rational \(F_4\) roots and standard form on \(K^4\).
Let \(W=W(F_4)\), and put
\[
 \Phi=\begin{pmatrix}0&-1&-1&0\\1&0&0&1\\1&0&0&-1\\0&-1&1&0\end{pmatrix},
 \qquad z=\frac{1+i}{2}\Phi,\qquad X=\langle W,z\rangle.
\]
The matrix \(\Phi/\sqrt2\) interchanges the long and short root
lines. This was verified on the complete 48-root system in D-GF4-01.
Moreover \(\Phi^2=-2I\), so \(z^2=-iI\), and \(z\) normalizes
\(W\). Consequently
\[
       X=W\sqcup iW\sqcup zW\sqcup izW,\qquad |X|=4608.
\]

Every finite unitary normalizer of \(W\) over \(K\) is contained
in \(X\). Here are the details of the scalar issue. The rational
intertwiner argument of D-FIELD-02 writes such an element as \(g=cA\)
with \(A\) rational and \(c\in K^*\). Finite order implies
\(c^N\in\mathbb Q\), hence \(c/\bar c\in\mu_4\).
Therefore \(c\) is a rational multiple of one of
\(1,i,1+i,1-i\). The rational similarity \(A\) either preserves
the two root-length classes or interchanges them. After rational
rescaling it is respectively a Weyl element or a Weyl element times
\(\Phi\), by the chamber argument and the absence of a nontrivial
automorphism of the oriented \(F_4\) diagram. In the first case
unitarity forces the remaining scalar to be in \(\mu_4\). In the
second case it forces \(2N(\lambda)=1\); the listed four possible
rational directions then give
\(\lambda\in\mu_4(1+i)/2\). This proves the containment.

For an irreducible component of type \(D_4\), Cohen's normalizer
Proposition 2.13 gives containment in scalar multiples of \(W(F_4)\),
including the triality automorphisms
\cite[\S2.13, pp.~391--392]{FIELD:Cohen}. Its rational Weyl matrices
and the \(K\)-model force the finite scalar to lie in \(\mu_4\).
Thus every such stabilizer image is contained in \(\mu_4W\subset X\).

An exhaustive exact enumeration from the four standard simple \(F_4\)
reflections and \(z\) gives 58 conjugacy classes in \(X\).
The simple roots used are
\[
 e_2-e_3,\quad e_3-e_4,\quad e_4,\quad
                   (e_1-e_2-e_3-e_4)/2.
\]
The integer program stores twice every Gaussian matrix, verifies every
product's divisibility, enumerates the full generated group, verifies
unitarity, and partitions it by conjugation. The rational checker then
computes every moving-space Gram determinant and local invariant as in
D-E67-01, with exact cyclotomic Fourier spectra. There are just three
nonnorm classes:
\[
\begin{array}{c|c|c|c|c|c}
 \text{certificate index}&\text{size}&\text{order}&k_g&q_g&
                       \text{age}(g)-k_g/2\\\hline
11&96&12&4&3&1/2\\
31&96&12&4&3&-1/2\\
37&96&6&4&3&0
\end{array}
\]
In particular every element of moving rank at most three has norm local
invariant. All matrices, classes and invariants are retained in
\texttt{F4\_gaussian\_classes.json} and
\texttt{F4\_gaussian\_local\_invariants.json}, with the short
programs \texttt{enumerate\_f4\_gaussian.cpp} and
\texttt{check\_f4\_gaussian.py}.

A genus-zero sixfold source in this rank-four representation has
\(D=14\). A nonsplit balanced determinant requires a nonnorm local
factor by D-DET-01, hence an inertia of moving rank four. Its branch
count therefore satisfies \(14\ge4+(r-1)\), or \(r\le11\).
Its period image has dimension at most eight. For a base of genus
one, \(D=6\) bounds the dimension by six; higher genera are impossible.
This argument applies to every subgroup of \(X\), since it uses the
actual matrix local invariants and does not require that the subgroup
contain all classes. In particular it excludes dominant nonsplit
sources from the \(D_4\) and \(F_4\) reflection components, including
both triality and the Gaussian length-interchange extension.

\subsection*{D-GAUSS-02. An all-rank obstruction for the standard \(A_m\) component}

Let \(m\ge3\), \(n=m+1\), and let \(W(A_m)=S_n\) act on its
standard sum-zero subspace. The finite Gaussian unitary normalizer is
contained in \(\mu_4W(A_m)\): all roots have the same length, and
the chamber argument identifies root-system automorphisms with Weyl
elements and the diagram involution \(-w_0\). The rational-intertwiner
and scalar argument used in D-E67-02 now applies. A stabilizer image
containing \(W(A_m)\) and having nonrational character therefore
contains \(iI\) and is the whole group
\[
                J=\mu_4\times S_n,\qquad
                \rho=\chi_i\otimes\operatorname{Std}_n.
\]
The product is direct because no nonidentity scalar occurs in the
standard \(S_n\)-representation when \(n\ge3\).
This standard representation is not being assumed monomial.

Consider first a genus-zero generating branch tuple. Write its elements
as \(\lambda\pi\), with \(\lambda\in\mu_4\), \(\pi\in S_n\).
If \(\lambda\) has order four, the fixed dimension of
\(\lambda\pi\) on the permutation representation is the number
of cycles of \(\pi\) whose lengths are divisible by four. The
constant line has eigenvalue \(\lambda\ne1\), so removing it does
not change the fixed dimension. Thus
\[
      k(\lambda\pi)\ge m-\lfloor n/4\rfloor.
\]
Likewise, for \(\lambda=-1\),
\[
      k(-\pi)\ge m-\lfloor n/2\rfloor.
\]
Let \(t\) be the number of order-four scalar entries and \(u\)
the number of scalar-minus-one entries. Generation and the scalar
product-one relation imply that \(t\ge2\) is even.
Every branch is nonidentity; a rank-nine period map would require
\(r\ge12\), and hence the total moving excess would satisfy
\[
    \sum_j(k_j-1)=2m+6-r\le2m-6=2n-8.
\]
If \(t\ge2\) and \(u\ge1\), its lower bound from these three
entries already exceeds the right side: the difference is
\[
 3(n-2)-2\lfloor n/4\rfloor-\lfloor n/2\rfloor-(2n-8)
      =n+2-2\lfloor n/4\rfloor-\lfloor n/2\rfloor\ge2.
\]
If \(t\ge4\), the four order-four entries alone exceed the bound,
with difference
\[
              4(n-2-\lfloor n/4\rfloor)-(2n-8)
                  =2n-4\lfloor n/4\rfloor>0.
\]
Consequently a possible rank-nine source has
\[
                  t=2,\qquad u=0.
\]
The two nontrivial scalars must be \(i\) and \(-i\).

Let \(C\) be the connected \(J\)-cover. Its cyclic quotient
\(T=C/S_n\to\mathbb P^1\) has exactly two branch points, both
fully ramified of order four. Riemann--Hurwitz gives \(g(T)=0\).
Put
\[
 H=\mu_4\times S_{n-1},\qquad
 X=C/S_{n-1},\qquad Y=C/H.
\]
Then \(X\to Y\) is a connected cyclic degree-four cover. The
representation identity
\[
 \operatorname{Ind}_H^J(\chi_i)
       =\chi_i\otimes\operatorname{Perm}_n
       =\chi_i\oplus\rho
\]
and Frobenius reciprocity identify its faithful-character multiplicity
space with the direct sum of those for \(\chi_i\) and \(\rho\).
The first is zero, because it lies in \(H^1(T)\). Applying the same
argument to conjugate characters gives a rational Hodge isogeny
\[
                        P_4(X/Y)\sim B_\rho.
\]
It is induced by fixed group-algebra correspondences in the family.

This identification also controls the actual polarization. The
Rosati-fixed rational projector for the pair \(\chi_i,\bar\chi_i\)
on \(H\) has one primitive line in the \(\rho\)-block and a line
in the scalar-character block. The latter acts as zero on \(H^1(C)\).
Its nonzero abelian image is therefore a primitive \(\rho\)-factor,
and is the pullback of the cyclic factor from \(X\). Pullback
multiplies the Jacobian polarization by \(|S_{n-1}|\). Any other
Rosati-fixed primitive line changes the multiplicity Hermitian form by
a positive rational scalar, as in D-MON-03. Hence splitness of this
cyclic factor is splitness of the specified primitive polarization.

If \(h=g(Y)\) and \(R\) is the cyclic branch count, the faithful
parabolic dimension is now exactly six:
\[
                              6=2h-2+R.
\]
The actual cyclic variation has moduli dimension at most
\(5,6,7,8,9\) for \(h=0,1,2,3,4\). Its isogeny to the
sixfold variation preserves the period differential's kernel. Thus
rank nine forces \(h=4,R=0\). Looijenga's equivariant hyperbolic
basis, already verified and applied in D-MON-02--03, makes the actual
cyclic-factor polarization split. This excludes the desired nonsplit
source for every \(m\), without invoking monomiality of the standard
representation.

Higher-genus original bases cannot give rank nine. At genus one,
\(D=6\), so their dimension is at most six. At genus two,
\(D=6-2m\), leaving only \(m=3,D=0\) among the ranks under
consideration, with base-moduli dimension three. Larger genus is
impossible. Ranks one and two are already excluded respectively by
D-MON and the rank-two Hurwitz-dimension bound in D-FIELD-03.
Induction from a reflection-component stabilizer preserves the actual
variation and polarization as in that same statement. Therefore the
obstruction also applies to arbitrary original finite groups with an
\(A_m\) reflection component.

\subsection*{D-GAUSS-03. Exact scope of the remaining reflection-containing gate}

Let \(R\) be the normal closure of the reflections in the original
irreducible Gaussian group. As in D-FIELD-02, its nonorthogonality
components are defined over \(K\), the original group permutes them
transitively, and the original representation is induced from the
stabilizer of one irreducible reflection component.

Cohen's classification and Malle's character-field tables show that
the nonreal primitive reflection components over \(K\) are \(G_8\)
in rank two and \(G_{29},G_{31}\) in rank four
\cite[\S3.6 and Theorem 5.12]{FIELD:Cohen}
\cite[Tables 8.1--8.2, pp.~674--675]{FIELD:Malle}.
The omissions \(G_4,G_{25}\) in Malle's table have order-three
reflections and are not Gaussian; the omitted \(G_{28},G_{35}\)
are the rational Weyl groups \(F_4,E_6\).
The primitive real possibilities are rational Weyl groups, since the
noncrystallographic real character fields are not contained in \(K\).

An imprimitive component has a unique system of imprimitivity except
for the four cases in Cohen's Lemma 2.7
\cite[pp.~389--390]{FIELD:Cohen}. The rank-two exceptions are already
excluded by their dimension bound; \(G(3,3,3)\) has non-Gaussian
character field; and \(G(2,2,4)=W(D_4)\) is covered by D-GAUSS-01.
In every other imprimitive case its normalizer preserves that unique
system, so its representation is monomial. Transitivity of induction
then makes the original representation monomial, and D-MON excludes
dominant nonsplit sources even if the inducing character has a larger
cyclotomic field. This covers the usual \(B_m,D_m\) cases as well.

For \(G_{29},G_{31}\), the corrected normalizer Proposition 5.14
of Cohen puts the stabilizer image in scalar multiples of that group
\cite{FIELD:Cohen,FIELD:CohenErr}. Each already contains \(\mu_4\),
as is visible in the central classes of the exact enumeration in
D-G31-01. The scalar lies in \(K\), and finite order forces it
into \(\mu_4\). Thus the stabilizer image is the same natural
group, and D-G31-02 applies.

The remaining rational Weyl components are accounted for by
D-GAUSS-01 (\(D_4,F_4\)), D-GAUSS-02 (all \(A_m\)),
D-E67-02 (\(E_6,E_7\)), and the rank-two dimension bound.
Accordingly, among reflection-containing Gaussian sources in the stated
group-algebra setup, the sole remaining possible component for a
dominant nonsplit sixfold is
\[
                              \mu_4W(E_8).
\]
This is a reduction to a candidate component, not a dominance theorem
or an algebraic-cycle supplier. Reflection-free representations remain
outside this classification and must be considered separately.

\section{The first reflection-free sixfold source degrees}
\label{sec:bireflection-gate}

The statements here concern genus-zero finite-cover sources in the
representation setting of O$''$.  Let $K$ be imaginary quadratic,
let $V$ be an absolutely irreducible $K$-representation of dimension
$m$, and let its complex character field be exactly $K$.  Its finite
image is denoted by $G$.  We consider a generating product-one tuple
$(g_1,\ldots,g_r)$ whose multiplicity variation has $K$-dimension six
and signature $(3,3)$.  Trivial inertia in this representation is
discarded.  We assume that the finite image is reflection-free.
Write $d_j=\operatorname{rank}(1-g_j)$.  The necessary conditions for
dominance are
\[
 d_j\ge2,\qquad \sum_jd_j=2m+6,\qquad r\ge12.
\tag{BI.1}
\]
The last condition is only necessary; none of the arguments below
infers dominance from it.  The exact-character-field hypothesis is
the ordinary absolutely irreducible coefficient setting.  A
rational-character coefficient endowed subsequently with an extra
$K$-action is a separate construction; generic endomorphisms already
obstruct it as in D-E8-GRR-06.

\begin{theorem}[D-BI-01: two-dimensional inducing blocks]
\label{D-BI-01}
Suppose the complex representation is induced from an absolutely
irreducible representation $\tau$ of dimension two of a subgroup
$H$, so $V_{\mathbf C}\simeq\operatorname{Ind}_H^G\tau$.
Then the period map of the associated Weil sixfolds has rank at most
seven.  This assertion does not require the character field of
$\tau$ to equal $K$.
\end{theorem}
\begin{proof}
Let $C\to\mathbf P^1$ be the connected $G$-cover, set $X=C/H$, and
let $H_0$ be the kernel of $\tau$.  The intermediate map
\[
             Y=C/H_0\longrightarrow X
\]
is a connected finite Galois cover with group $H/H_0$.  Frobenius
reciprocity and quotient pullback give an isomorphism of complex
Hodge variations
\[
 \operatorname{Hom}_G(V_{\mathbf C},H^1(C,\mathbf C))
 \simeq \operatorname{Hom}_H(\tau,H^1(Y,\mathbf C)).
\tag{BI.2}
\]
This identification holds in families after a finite cover of the
parameter space marking the subgroup and representations.  It is
induced by fixed finite-group algebra correspondences; it is not
merely an equality of fibre dimensions.

Put $h=g(X)$.  Discard precisely the points whose local inertia is
trivial on $\tau$, and denote the remaining number of branch points
by $R$.  If $D$ is the sum of the moving dimensions of these local
operators on $\tau$, the Euler characteristic of the middle-extension
local system gives
\[
              6=2(2h-2)+D,\qquad D=10-4h.
\tag{BI.3}
\]
The zero-th and second cohomology terms vanish: $\tau$ is a nontrivial
absolutely irreducible representation and the cover has its full
image $H/H_0$.  Faithfulness of this image gives $R\le D$.
In particular $h\le2$.  For $h=0$ the six-dimensional parabolic
cohomology also implies $R\ge5$, since each moving dimension is
at most two.  For $h=1$ one has $R\ge3$.  Thus all pointed curves
in question are stable; the unpointed exceptional moduli dimensions
do not occur.

With the finite group and topological branching data fixed, the
Hurwitz space of $Y\to X$ has finite fibres over the moduli space of
the $R$-pointed genus-$h$ curve: the punctured surface group is
finitely generated and has only finitely many homomorphisms into
the fixed finite group.  Locally this is the usual finite cover of
the pointed Teichmuller space.  Consequently (BI.2) factors through
a space of dimension at most
\[
       3h-3+R\le3h-3+(10-4h)=7-h\le7.
\tag{BI.4}
\]
Forgetting the additional map $X\to\mathbf P^1$ can lose parameters,
but cannot add variation to (BI.2).

Finally, the one complex multiplicity space in (BI.2) determines the
period of the rational $K$-Hodge structure.  Indeed, its filtration
$F^1_\sigma$ determines $F^1_{\bar\sigma}$ as its annihilator under
the fixed polarization pairing between the two embeddings of $K$.
Thus the bound is a bound for the full Weil period map, even when
$\tau$ itself is defined over a larger field.  No descent of
$\tau$ to $K$ was used.
\end{proof}

\begin{lemma}[D-BI-02a: tensor products in degrees nine through fourteen]
\label{D-BI-02a}
Under (BI.1), if $9\le m\le14$, the representation is tensor
indecomposable in the sense required by Guralnick--Saxl's
Theorem~7.1.
\end{lemma}
\begin{proof}
Since $2m+6<36$, (BI.1) forces at least one $d_j=2$.
Suppose there is a factorwise invariant tensor decomposition
$V_{\mathbf C}=A\otimes B$, with $2\le a=\dim A\le b=\dim B$.
The factors of any finite-order tensor operator can be rescaled to
be diagonalizable of finite order.  If one factor is nonscalar,
its largest eigenspace has codimension at least one.  Summing
products of eigenvalue multiplicities shows that a nonscalar tensor
operator has moving dimension at least $\min(a,b)$, even after
an arbitrary scalar multiplication.  Thus the existence of a
bireflection excludes $a\ge3$ in the stated degree range.

It remains that $a=2$ and $b\in\{5,6,7\}$.  The projective image
on $A$ is irreducible, since an invariant line in $A$ would give an
invariant proper subspace of $A\otimes B$.  At least three of the
inertia operators have nontrivial image in $\operatorname{PGL}(A)$.
Otherwise their product-one relation makes that image trivial or
cyclic, and a finite cyclic projective subgroup of
$\operatorname{PGL}_2(\mathbf C)$ fixes a line.  Each of these
three operators has moving dimension at least $b$: both eigenvalue
multiplicities on $A$ are one, so its tensor product has a fixed
space of dimension at most $b$.  The other inertia operators have
moving dimension at least two.  Hence
\[
       2m+6\ge3b+2(r-3)\ge\frac{3m}{2}+18,
\]
which requires $m\ge24$, a contradiction.

The word ``tensor indecomposable'' in the cited theorem excludes
factorwise tensor decompositions.  Groups permuting equal tensor
factors, for example $\operatorname{GL}(A)\wr\mathfrak S_t$ on
$A^{\otimes t}$, are explicitly included in the theorem's
tensor-indecomposable case; see its introductory paragraph on
p.~548.  Thus no tensor-permutation case has been omitted here.
\end{proof}

\begin{theorem}[D-BI-02: primitive degrees nine through fourteen]
\label{D-BI-02}
There is no source satisfying (BI.1) with primitive complex
representation and $9\le m\le14$ whose character field is imaginary
quadratic.  This exclusion does not depend on its determinant class.
\end{theorem}
\begin{proof}
There is a bireflection and D-BI-02a applies.  We use exactly the
characteristic-zero specialization of the published
Guralnick--Saxl classification \cite[Theorem~7.1 and Remark~(3),
pp.~548--549]{BI:GS}: in these dimensions the possibilities are
scalar extensions of the deleted permutation representation of an
alternating or symmetric group, or the ten-dimensional representation
of $U_5(2)$ and its linear normalizer.  The defining-characteristic
classical cases in that theorem do not occur over $\mathbf C$.

In the permutation case put $n=m+1$.  After conjugation every group
element is $\lambda\sigma$ on the deleted subspace of
$\mathbf C^n$, with $\sigma\in\mathfrak S_n$ and $\lambda$ a root
of unity.  This description includes the normalizer of
$\mathfrak A_n$: for $n\ge10$ its automorphisms are induced by
$\mathfrak S_n$, and Schur's lemma supplies only an extra scalar.
If $\lambda\notin\{1,-1\}$ has order $e\ge3$, the
$\lambda^{-1}$-eigenspace of a permutation matrix has dimension
at most $\lfloor n/e\rfloor\le\lfloor n/3\rfloor$.
Indeed a cycle contributes at most one such eigenvector and must
have length divisible by $e$.  Deleting the trivial line does not
change this nontrivial eigenspace.  Thus
\[
 \operatorname{rank}(1-\lambda\sigma)
       \ge q_m:=m-\left\lfloor\frac{m+1}{3}\right\rfloor.
\tag{BI.5}
\]
If all inertia scalars are real, their generated representation is
real and its character, lying in $K$, is rational.  Otherwise at
least two inertia scalars are nonreal: the product of the scalars
is one, since the deleted permutation group contains no nontrivial
scalar.  Equation (BI.1) would then imply
\[
        2m+6\ge2q_m+2(r-2)\ge20+2q_m.
\tag{BI.6}
\]
The following exact values contradict this inequality:
\[
\begin{array}{c|rrrrrr}
m&9&10&11&12&13&14\\ \hline
2m+6&24&26&28&30&32&34\\
q_m&6&7&7&8&9&9\\
20+2q_m&32&34&34&36&38&38.
\end{array}
\]

For $U_5(2)$ one has $m=10$.  Its degree-ten module preserves a
nondegenerate alternating form.  This is the symplectic Weil module
in \cite[Table~II, author version p.~7; Proposition~7.10]{BI:GT}; it is also
given by explicit matrices and an alternating form of determinant
one in the author's ATLAS data \cite{BI:AtlasU52}.
Irreducibility makes the invariant form line unique, so every
element of its linear normalizer is a similitude of this form.
If the multiplier of an inertia operator is not one, its fixed
space is totally isotropic: for fixed vectors $x,y$ one has
$b(x,y)=b(gx,gy)=c(g)b(x,y)$.
Its fixed dimension is therefore at most five, and its moving
dimension is at least five.  But (BI.1) bounds every moving
dimension by
\[
              26-2(r-1)\le4.
\]
Every inertia multiplier is consequently one.  They generate $G$,
so $G$ preserves the alternating form, the representation is
self-dual, and its character is real.  Once again an imaginary
quadratic character field is impossible.
\end{proof}

\begin{proposition}[D-BI-03: the full degree-nine nonsplit exclusion]
\label{D-BI-03}
No reflection-free source of coefficient degree at most nine can
dominate a nonsplit Weil sixfold family in this representation
setting.  Thus the least numerical possibility $m=9$ in (BI.1)
does not occur.
\end{proposition}
\begin{proof}
If $m<9$, then $2m+6<24$, contrary to (BI.1).  If $m=9$, then
necessarily $r=12$ and every inertia has moving dimension two.
The primitive case is D-BI-02.  An imprimitive irreducible
nine-dimensional representation has either nine one-dimensional
blocks or three three-dimensional blocks.  The first case is
monomial and is excluded for a dominant nonsplit sixfold by
D-MON-01--03.  In the second case, an operator nontrivially
permuting the blocks has moving dimension at least three: a
cycle of length $\ell\ge2$ on blocks of dimension three has
fixed dimension at most three and moving dimension at least
$3(\ell-1)$.  Every inertia therefore fixes every block.
Since the tuple generates $G$, this contradicts irreducibility.
\end{proof}

\begin{proposition}[D-BI-04: the remaining imprimitive blocks below fifteen]
\label{D-BI-04}
Suppose (BI.1) holds with $9\le m\le14$, and suppose the source
could dominate a nonsplit sixfold family.  Apart from the excluded
monomial and two-dimensional-block cases, the only imprimitive
block sizes still permitted are
\[
            (m,t,d)=(12,4,3)\quad\hbox{or}\quad(14,2,7),
\tag{BI.7}
\]
where $m=td$.  The intermediate block quotient $X\to\mathbf P^1$
has genus zero.  Equality holds in every rank estimate: the first
case has six block-transposition inertia operators of rank three
and six block-fixed inertia operators of rank two; the second has
two block-transposition operators of rank seven and ten block-fixed
operators of rank two.  This proposition is a reduction, not a
claim that either remaining case exists or dominates.
\end{proposition}
\begin{proof}
For a transitive imprimitivity system with $t$ blocks of dimension
$d$, let $\pi_j$ be the permutation of the blocks induced by
$g_j$, and let $\operatorname{ind}(\pi_j)=t-\#\text{cycles}(\pi_j)$.
The fixed space on a block cycle has dimension at most $d$, so
\[
             d_j\ge d\operatorname{ind}(\pi_j).
\]
If $\pi_j\ne1$, then
$d_j-2\ge(d-2)\operatorname{ind}(\pi_j)$; the difference in this
estimate includes $2(\operatorname{ind}(\pi_j)-1)$.
For $\pi_j=1$ the same inequality follows from $d_j\ge2$.
If $h=g(X)$ for the block quotient, Riemann--Hurwitz gives
\[
 \sum_j\operatorname{ind}(\pi_j)=2t+2h-2.
\]
Consequently
\[
 2m-18\ge\sum_j(d_j-2)
       \ge(d-2)(2t+2h-2),
\qquad
 2t+d\ge11+(d-2)h.
\tag{BI.8}
\]
For $d\ge3$, $t\ge2$ and $9\le td\le14$, the last inequality
leaves exactly the two pairs in (BI.7).  Both satisfy $2t+d=11$,
so $h=0$ and every inequality used is an equality.  In particular
$r=12$, every nontrivial block permutation has index one and hence
is a transposition, each has moving rank exactly $d$, and every
block-fixed inertia has moving rank two.  The required number of
transpositions is $2t-2$, which gives the two stated patterns.
\end{proof}

\begin{theorem}[D-BI-05: three-dimensional inducing blocks]
\label{D-BI-05}
Let $K$ be one of
\[
 \mathbf Q(i),\quad \mathbf Q(\sqrt{-2}),\quad
 \mathbf Q(\sqrt{-7}),\quad\mathbf Q(\sqrt{-11}),\quad
 \mathbf Q(\sqrt{-19}),\quad\mathbf Q(\sqrt{-43}),\quad
 \mathbf Q(\sqrt{-67}),\quad\mathbf Q(\sqrt{-163}).
\]
In the exact-character-field setting above, an irreducible
representation induced from a three-dimensional complex representation
cannot dominate a nonsplit Weil sixfold family.  The inducing
representation need not have a model over $K$, and the original
coefficient degree is unrestricted.
\end{theorem}
\begin{proof}
Write $V_{\mathbf C}=\operatorname{Ind}_H^G\tau$ with
$\dim\tau=3$ and $L=\mathbf Q(\chi_\tau)$.  Induction of
characters immediately gives
\[
                         K=\mathbf Q(\chi_V)\subset L.
\tag{BI.9}
\]
This inclusion uses values of complex characters, so does not
presuppose a rational block system.  Form the intermediate cover
$Y=C/\ker\tau\to X=C/H$ as in D-BI-01.  Its multiplicity variation
is the same one complex period variation as the original family.
If $h=g(X)$, $R$ is its genuine branch count, and $D$ the sum of
its moving dimensions, Euler characteristic gives
\[
                   D=12-6h,\qquad R\le D.
\]
The stable pointed-curve Hurwitz dimension is at most
\[
                    3h-3+R\le9-3h.
\tag{BI.10}
\]
The potentially unstable cases do not give an exception: at genus
zero the positive six-dimensional parabolic cohomology forces at
least four punctures; at genus one $D=6$ forces branching; at genus
two the unramified base has dimension three.  Thus rank nine forces
$h=0$, $R=12$, and equality $R=D$.  Every inertia on $\tau$ is
therefore a complex reflection.  Since $X=\mathbf P^1$, these
inertia generate its full image, which is a finite irreducible
rank-three reflection group.

If $\tau$ is imprimitive, its rank being three makes it monomial;
induction then makes $V$ monomial.  D-MON-01--03 excludes a
dominant nonsplit source in this case.  Otherwise the published
rank-three reflection classification and its splitting fields give
\[
\begin{array}{c|ccccc}
\text{group}&G_{23}=W(H_3)&G_{24}&G_{25}&G_{26}&G_{27}\\ \hline
k&\mathbf Q(\sqrt5)&\mathbf Q(\sqrt{-7})&
 \mathbf Q(\omega)&\mathbf Q(\omega)&
 \mathbf Q(\omega,\sqrt5).
\end{array}
\tag{BI.11}
\]
Here the classification and the $G_{25}$ realization are
\cite[\S4.16, p.~412; Theorem~5.12, p.~435]{FIELD:Cohen};
the splitting fields are \cite[Table~8.2, p.~675, with its
omission convention on p.~676]{FIELD:Malle}.  In each case
$L\subset k$, which is all that is needed.  The quadratic
subfields of the last biquadratic field are
$\mathbf Q(\sqrt5)$, $\mathbf Q(\sqrt{-3})$, and
$\mathbf Q(\sqrt{-15})$.  Hence for the eight fields in the
statement, (BI.9) leaves only $G_{24}$ and
$K=L=\mathbf Q(\sqrt{-7})$.

In this last case the natural $G_{24}$ representation has a model
over $K$.  Both $V$ and $\operatorname{Ind}_H^G\tau$ then have
$K$-models and equal absolutely irreducible complex characters;
their intertwining linear equations over $K$ give a $K$-isomorphism.
Shapiro's identification of the parabolic complexes consequently
identifies their polarized $K$-variations up to a positive rational
scalar: use the induced positive Hermitian form and sum cup products
over cosets; the specified irreducible invariant form differs by
such a scalar by Schur's lemma.  This is the polarization argument
of D-FIELD-03, now justified after establishing $L=K$.
On a six-dimensional $K$-space a rational scalar changes the
determinant by its sixth power, a $K$-norm.  D-DET-02 proves that
the natural $G_{24}$ source in question has split signed determinant.
The original source is therefore split as well, a contradiction.
\end{proof}

\begin{proposition}[D-BI-06: the remaining two blocks of dimension seven]
\label{D-BI-06}
The $(m,t,d)=(14,2,7)$ case left by D-BI-04 cannot dominate a
nonsplit Weil sixfold family with exact imaginary quadratic
character field.
\end{proposition}
\begin{proof}
Let $\tau$ be the representation on one block.  If it is
imprimitive, its prime degree seven makes it monomial, hence the
original representation is monomial; the nonsplit case is already
excluded.  We can thus assume $\tau$ is primitive.  D-BI-04 gives
$X=\mathbf P^1$, two minimal-rank block transpositions, and ten
block-fixed rank-two inertia operators.  A minimal-rank block
transposition has the form
\[
       \begin{pmatrix}0&A\\A^{-1}&0\end{pmatrix},
\]
because its fixed space has dimension seven.  Its square is the
identity, so it induces no inertia on $\tau$ above the corresponding
ramification point of $X\to\mathbf P^1$.  Every remaining inertia
on $\tau$ comes from a block-fixed operator and has rank at most two.

The intermediate six-dimensional multiplicity has total moving
dimension $2\cdot7+6=20$.  Rank nine of its period map requires at
least twelve genuine branch points on $X$.  Thus at least one
inertia on $\tau$ is a reflection; otherwise the branch count is
at most ten.  Let $R_0$ be the normal closure of all reflections
in the finite image of $\tau$.  Primitivity makes $R_0$ irreducible:
its fixed space is invariant, and its orthogonal reflection
components would otherwise be a nontrivial imprimitivity system.

By the complex reflection classification in rank seven, $R_0$ is
real or belongs to the imprimitive series.  An imprimitive group
in that rank has a unique imprimitivity system, by
\cite[Lemma~2.7, pp.~389--390]{FIELD:Cohen}; none of the four
exceptions there has rank seven.  Its normalizer would preserve
that system, contrary to primitivity of $\tau$.  Therefore $R_0$
has a real reflection representation.  The same conclusion covers
the standard $A_7$ representation as well as the exceptional $E_7$
representation; neither is being omitted.

Choose its real positive definite form $B$.  Uniqueness of the
invariant Hermitian form for the irreducible $R_0$ shows that the
finite image of $\tau$ is unitary for $B$.  If $T$ belongs to
this finite image, then $T^{-1}\bar T$ centralizes $R_0$, since the automorphism
of $R_0$ induced by $T$ has real matrices.  Schur's lemma writes
$T=\lambda A$ with $A$ real; rescaling by a positive real number
makes $A$ orthogonal for $B$ and $|\lambda|=1$.
If $\lambda$ is nonreal, its inverse is a nonreal eigenvalue of
the real seven-dimensional $A$.  Conjugate eigenvalues have the
same multiplicity, so this multiplicity is at most three.  Hence
$\operatorname{rank}(1-T)\ge4$.  Every inertia on $\tau$ has rank
at most two, so it is real.  These inertia generate the full image
over $X=\mathbf P^1$, making $\mathbf Q(\chi_\tau)$ real.
But $\chi_V=\operatorname{Ind}_H^G\chi_\tau$ implies
$K\subset\mathbf Q(\chi_\tau)$, a contradiction.
\end{proof}

\begin{corollary}[D-BI-07: a strictly smaller remaining source gate]
\label{D-BI-07}
For the eight non-Eisenstein class-number-one fields in D-BI-05,
a dominant nonsplit finite-cover source with exact character field
$K$ and reflection-free image must have coefficient degree at least
fifteen.  It cannot be induced from a complex representation of
dimension one, two, or three.  The Eisenstein field is omitted
because its all-determinant sixfold theorem is already established
elsewhere in this paper, not because the field was discarded from
the original target.
\end{corollary}
\begin{proof}
Degrees below nine fail (BI.1).  D-BI-02 excludes the primitive
degrees nine through fourteen.  For the imprimitive degrees in
that range, one-dimensional blocks are D-MON-01--03,
two-dimensional blocks are D-BI-01, and D-BI-04 leaves only
the cases excluded by D-BI-05 and D-BI-06.  Inducing dimensions
one, two and three in arbitrary original degree are excluded by
the same three results.
\end{proof}

\paragraph{Verification and next gate.}
The two-dimensional-block factorization was independently derived
by the imprimitive agent.  The character-field inclusion and its
application to all three-dimensional inducing blocks were independently
pointed out by root.  The other arguments are new deductions from
(BI.1), the stated published finite classification, and elementary
eigenvalue estimates.  The companion exact checker verifies the
permutation and block-budget tables and the invariant forms in the
retained degree-ten ATLAS matrices.  It does not claim to prove the
classification or replace the family-level factorization proof.
Degrees at least fifteen with no inducing representation of dimension
at most three are the remaining reflection-free gate for the eight
fields in D-BI-07.  This is a source obstruction, not a proof of
algebraicity for those remaining determinant families.

\section{Scalar alternating and symmetric inducing blocks}
\begin{theorem}[D-BI-08: scalar standard inducing representations]
\label{D-BI-08}
Let an exact-$K$ Weil sixfold representation be induced from a
complex representation $\tau$ whose finite image, in its given
representation, satisfies
\[
 \mathfrak A_n\subset J\subset\mathbf C^*\mathfrak S_n,
 \qquad \tau=\chi_d\otimes(\operatorname{Std}_n\circ\pi),
 \qquad n\ge4,
\tag{BI.12}
\]
where the permutation groups act on their deleted permutation space
and $\pi:J\to\mathfrak S_n$ and $\chi_d:J\to\mu_d$ are the
permutation and scalar projections, with scalar image of order $d$.
A rank-nine period
map forces $d\in\{3,4,6\}$, $K=\mathbf Q(\zeta_d)$, and split
actual polarization.  In particular this cannot supply a dominant
nonsplit family.  There is no assumption that the inducing
cyclotomic field initially equals $K$.

This applies, in particular, to an inducing representation whose
reflection normal closure is an irreducible group of type $A_{n-1}$:
its complex linear normalizer is $\mathbf C^*\mathfrak S_n$.
It also includes scalar extensions of the alternating deleted
representation in the Guralnick--Saxl bireflection classification;
a direct-product assumption is not being imposed.
\end{theorem}
\begin{proof}
Set $k=n-1$ and $L=\mathbf Q(\zeta_d)$.  The displayed matrices
give a model of $\tau$ over $L$; induced characters give
$K\subset\mathbf Q(\chi_\tau)\subset L$, so $d>2$.
The scalar and permutation projections are well-defined because
the deleted permutation group contains no nonidentity scalar.
Quotient the kernel of $\tau$ in its
stabilizer to obtain a connected $J$-cover $C_0\to X_0$.
Frobenius reciprocity identifies its $\tau$-multiplicity variation
with the original complex sixfold multiplicity, exactly as in
D-BI-01.  If $h=g(X_0)$, Euler characteristic gives
$D=6-2k(h-1)$, where $D$ is the total moving dimension.
For $h=1$ the marked-base moduli dimension is at most six.  For
$h\ge2$ and $k\ge3$ the only nonnegative case is $h=2,k=3,D=0$,
whose base moduli has dimension three.  Therefore rank nine forces
$h=0$, at least twelve genuine branch points, and
\[
       \sum_j\bigl(\operatorname{rank}(1-g_j)-1\bigr)
                    =2k+6-r\le2n-8.
\tag{BI.13}
\]

Write $g_j=\lambda_j\pi_j$ on the deleted permutation space.
If the order of $\lambda_j\ne1$ is $e_j$, its fixed space has
dimension at most $\lfloor n/e_j\rfloor$: only permutation cycles
whose lengths are divisible by $e_j$ contribute.  Hence its
contribution to (BI.13) is at least
$n-2-\lfloor n/e_j\rfloor$.
Let $t$ count the nonreal scalar entries and $u$ the entries
with scalar $-1$.  There must be a nonreal entry, since these
entries generate the nonreal scalar quotient; their scalar
product is one, so $t\ne1$.

If $t\ge3$, three of them alone exceed (BI.13), since
\[
 3(n-2-\lfloor n/3\rfloor)-(2n-8)
                   =n+2-3\lfloor n/3\rfloor\ge2.
\]
Thus $t=2$.  If $u\ge2$, two nonreal entries and two minus-one
entries exceed (BI.13), with positive difference
\[
               2n-2\lfloor n/3\rfloor-2\lfloor n/2\rfloor>0.
\]
If $u=1$, the two nonreal scalars have product $-1$.
Their orders $e_1,e_2\ge3$ satisfy
$1/e_1+1/e_2\le1/2$: if one has order three the other has order
six, and otherwise both orders are at least four.  The three
entries therefore contribute at least $2n-6$, again larger than
$2n-8$.  Consequently
\[
             t=2,\qquad u=0,\qquad
             \lambda_2=\lambda_1^{-1}.
\tag{BI.14}
\]
Since these scalars generate $\mu_d$, both have order $d$.
The cyclic quotient $T=C_0/\ker\chi_d\to X_0=\mathbf P^1$
is branched at exactly two points, fully ramified at each, and
has genus zero.

Let $H_1$ be the stabilizer of one letter under the permutation
projection of $J$.  Its scalar image is all of $\mu_d$:
for any element of $J$, an element of the real $\mathfrak A_n$
subgroup can move its image of that letter back to the letter,
without changing the scalar.  Put
$Y=C_0/\ker(\chi_d|_{H_1})$ and $X=C_0/H_1$.  Then $Y\to X$
is a connected cyclic cover of degree $d$.  The permutation
representation on the $n$ letters is $1\oplus\operatorname{Std}_n$,
so the representation identity
\[
 \operatorname{Ind}_{H_1}^{J}(\chi_d|_{H_1})
       =\chi_d\oplus\tau
\tag{BI.15}
\]
identifies its faithful-character cohomology with the direct sum
of the scalar-character and $\tau$ multiplicities.  The former
is zero because it comes from $H^1(T)$.

We spell out the rational multiplicity issue, since the inducing
field may be larger than $K$.  For every embedding
$\gamma:L\hookrightarrow\mathbf C$, (BI.15) holds with
$\chi_d^\gamma$ and $\tau^\gamma$, and its scalar-character
summand still has zero cohomology.  Induction from the original
block stabilizer to $G$ sends $\tau^\gamma$ to the original
$\rho$ or its complex conjugate according as $\gamma|_K$ is the
chosen embedding or its conjugate.  There are $\ell=[L:K]$
embeddings of each type.  Thus rational semisimplicity and
Frobenius reciprocity identify the primitive cyclic Hodge structure
with $\ell$ copies of the original primitive factor:
\[
                     P_d(Y/X)\sim B^{\ell}.
\tag{BI.16}
\]
More explicitly, inducing the rational $L$-module underlying
$\chi_d$ gives the sum of $\ell$ copies of $V_{\mathbf Q}$ and
an additional rational representation whose Hom into $H^1(C)$
vanishes.  Taking this Hom discards exactly that zero summand.
The constructions are rational group-algebra correspondences in
families, so (BI.16) is an isogeny of variations, not just an
identity of complex dimensions.
This argument does not require $\mathbf Q(\chi_\tau)=L$:
the multiplicity $\ell$ counts the embeddings of the scalar-character
module, and each induced $\tau^\gamma$ is determined by $\gamma|_K$.

The faithful-character cohomology of $Y\to X$ has dimension six.
Writing $a=g(X)$ and $R$ for its genuine cyclic branch count gives
$6=2a-2+R$.  Its cyclic Hurwitz moduli has dimension at most
$5,6,7,8,9$ for $a=0,1,2,3,4$, respectively.  Equation (BI.16)
identifies its period differential kernel with that of $B$;
therefore rank nine forces $a=4,R=0$, and a dominant image in a
component of the etale cyclic genus-four moduli.  D-MON-02, from
Looijenga's published monodromy theorem and the explicit commutant
argument there, gives $\operatorname{End}^0(P_d)=L$ at a very
general point in that image.  If $\ell>1$, (BI.16) embeds the
noncommutative algebra $M_\ell(K)$ into this field, a contradiction.
Hence $L=K$ and $d=3,4,6$.

The splitness is that of the actual polarization.  Once $L=K$,
the rational faithful-character projector for $H_1$ is
Rosati-fixed.  The character $\chi_d|_{H_1}$ has multiplicity one
in $\tau|_{H_1}$: the subgroup $\mathfrak A_{n-1}$ is transitive
on the other $n-1$ letters, also when $n=4$, so the fixed subspace
of the restricted permutation representation has dimension two,
and that of its deleted summand has dimension one.
By (BI.15) its active image is a single primitive
$\tau$-line: the other line is precisely the zero scalar-character
summand.  After induction the same assertion holds for the original
$\rho$-factor, since any other summand has zero cohomology.
Its image is the pullback of $P_d(Y/X)$, with the Jacobian
polarization scaled by the degree of the quotient map.  Looijenga's
equivariant hyperbolic basis makes the etale cyclic primitive
polarization split.  As in D-MON-03, changing the Rosati-fixed
primitive line only multiplies the multiplicity Hermitian form
by a positive rational scalar and preserves its isotropic
three-space.  This proves the asserted splitness.
\end{proof}


\section{Essentially self-dual inducing blocks}
\label{sec:essential-selfdual-blocks}

\begin{theorem}[D-BI-09: the degree-twenty-four gate]
\label{D-BI-09}
Suppose a finite-cover Weil sixfold source of exact imaginary quadratic
character field $K$ has a period map of rank nine and nonsplit actual
polarization.  Suppose its complex coefficient representation is
$\operatorname{Ind}_H^G\tau$, where $\tau$ is irreducible and primitive
of degree $k\ge9$.  If $\tau$ is essentially self-dual, then
\[
                         k\ge24.                         \tag{BI.17}
\]
Here essentially self-dual means that there is a nondegenerate complex
bilinear form $B$ and a character $\alpha$ of the finite image $J$ of
$\tau$ such that
\[
              g^{\mathsf T}Bg=\alpha(g)B\qquad(g\in J).
                                                               \tag{BI.18}
\]
The form may be symmetric or alternating.  The inducing character field
$L=\mathbf Q(\chi_\tau)$ need not equal $K$ and no model of $\tau$
over $K$ is assumed.  In particular, the conclusion applies to every
primitive inducing block that is real up to scalar matrices.

For a factorwise tensor decomposition $\tau=A\otimes B_0$ with
$2\le a=\dim A\le b=\dim B_0$ and $k=ab$, the stronger necessary
inequality is
\[
 (a-1)b\ge 9a-6+\epsilon_k,
 \qquad \epsilon_k=\begin{cases}0&k\text{ even},\\1&k\text{ odd}.
 \end{cases}                                                 \tag{BI.19}
\]
The tensor factors in this assertion may be projective representations.
\end{theorem}
\begin{proof}
Pass to the actual intermediate connected $J$-cover $C_0\to X_0$
as in D-BI-01.  Its $\tau$-multiplicity variation is the original
complex multiplicity variation, by Frobenius reciprocity and quotient
pullback.  In particular its cohomological multiplicity has dimension
six, and its period differential has the same rank as the original
$K$-factor: the conjugate filtration is its polarization annihilator.
If $h=g(X_0)$ and $D$ is the sum of the moving dimensions of the
genuine local inertias, then
\[
                         D=6-2k(h-1).
\]
For $h=1$, the marked-base moduli has dimension at most six; for
$h\ge2$ this formula is negative because $k\ge9$.  A finite cover
of a fixed marked base has only finitely many topological possibilities.
Consequently rank nine forces $h=0$ and at least twelve genuine
branch points.  Write their nonidentity inertia matrices as
$g_1,\ldots,g_r$.  They generate $J$, have product one, and satisfy
\[
 r\ge12,\qquad \sum_{j=1}^r\operatorname{rank}(1-g_j)=2k+6.
                                                               \tag{BI.20}
\]
This is a factorization through the actual intermediate Hurwitz family;
it is not a subtraction of dimensions of unrelated parameter spaces.

Induction of characters gives $K\subset L$.  If $\alpha$ were
trivial, (BI.18) would make $\tau$ self-dual.  The character of a
finite self-dual complex representation is real, contradicting the
inclusion of the imaginary quadratic field $K$ in $L$.  Thus $\alpha$
is nontrivial.  Since the $g_j$ generate and their product is one,
at least two of them have $\alpha(g_j)\ne1$.

The transpose of $B$ has the same equivariance, so Schur's lemma
gives $B^{\mathsf T}=B$ or $-B$ after rescaling; in fact the sign
already follows by transposing twice.  If $\alpha(g)\ne1$ and
$x,y\in\ker(g-1)$, (BI.18) yields
$B(x,y)=\alpha(g)B(x,y)$, hence $B(x,y)=0$.  The fixed subspace
is totally isotropic for a nondegenerate symmetric or alternating
form and therefore has dimension at most $\lfloor k/2\rfloor$.
The two distinguished inertias consequently have moving rank at least
\[
                            p=\lceil k/2\rceil.             \tag{BI.21}
\]

First suppose $J$ is tensor indecomposable in the convention of
\cite[paragraph before Theorem~7.1]{BI:GS}.  If it contains an element
of moving rank two, the characteristic-zero classification
\cite[Theorem~7.1 and Remark~(3), pp.~548--549]{BI:GS} gives either
the scalar deleted alternating/symmetric case or the degree-ten
$U_5(2)$ case.  The first is excluded by D-BI-08, which explicitly
allows the larger inducing field and the full subgroup range
$\mathfrak A_{k+1}\subset J\subset\mathbf C^*\mathfrak S_{k+1}$.
In the second case the finite image normalizes a nondegenerate
alternating form, as in D-BI-02 and
\cite[Table~II]{BI:GT}.  Its multiplier character must again be
nontrivial, since otherwise $L$ would be real.  Two inertias move at
least five dimensions.  Every remaining nonidentity inertia moves
at least two: a symplectic semisimple transformation has positive
moving rank at least two, and a nontrivial similitude moves at least
five.  Thus the total is at least $2\cdot5+10\cdot2=30$, exceeding
$2k+6=26$.  This also excludes the exceptional case.

There are no reflections either.  Indeed, after averaging a positive
Hermitian form, a reflection's moving line is perpendicular to its
fixed hyperplane.  Irreducibility gives two conjugate reflections
$u,v$ with distinct moving lines.  If $uvx=x$, then
$(v-1)x=(u^{-1}-1)x$ lies in the intersection of those two lines,
which is zero.  Hence
$\ker(uv-1)=\ker(u-1)\cap\ker(v-1)$, of codimension two.
Their product would be the already excluded bireflection.

Every other genuine inertia now moves at least three dimensions.
Combining this with (BI.20)--(BI.21) gives
\[
              2k+6\ge 2\lceil k/2\rceil+30
                         =k+\epsilon_k+30.
                                                               \tag{BI.22}
\]
Therefore $k\ge24+\epsilon_k$, proving the assertion in this case.
Groups that permute tensor factors but do not preserve a factorwise
decomposition fall within the cited tensor-indecomposable convention;
they are not omitted here.

It remains to handle a factorwise tensor decomposition of dimensions
$2\le a\le b$.  For any tensor matrix, the fixed space has dimension
at most the largest eigenspace multiplicity in the first factor
times $b$.  Thus a matrix whose first projective factor is nontrivial
moves at least $b$ dimensions.  If only the second projective factor
is nontrivial, it moves at least $a$; a nonidentity scalar matrix
moves all $k$ dimensions.  Every nonidentity inertia therefore moves
at least $a$.

At least three inertia matrices have nontrivial first projective
factor.  Otherwise their product-one relation makes the group
generated by the first projective factors cyclic.  A cyclic projective
group is represented by powers of one diagonalizable matrix up to
scalars and has an invariant line.  Tensoring that line with the
second factor contradicts irreducibility.  This argument does not
require the projective factors to admit linear lifts as representations.

Choose two of the inertias in (BI.21).  At least one of those three
nontrivial first-projective-factor inertias remains outside this pair.
The two selected inertias contribute at least $2p$, that third one
at least $b$, and at least nine others contribute at least $a$ each.
Consequently
\[
 2ab+6\ge 2\lceil ab/2\rceil+b+9a
                 =ab+\epsilon_k+b+9a,
\]
which is precisely (BI.19).  If $ab<24$, then $a\le4$.
For $a=2$, (BI.19) requires $b\ge12$; for $a=3$ it requires
$b\ge11$; for $a=4$ it requires $b\ge10$.  Each contradicts
$a\le b$ and $ab<24$.  This completes the proof.
\end{proof}

\begin{proposition}[D-BI-10: the multiplier field and odd degrees]
\label{D-BI-10}
Let $V$ be an absolutely irreducible primitive representation defined
over an imaginary quadratic field $K$, with exact character field
$K$, and suppose it is essentially self-dual.  Its similitude
character is unique and takes values in the roots of unity of $K$.
If $K\ne\mathbf Q(i),\mathbf Q(\sqrt{-3})$, then $\dim V$ is even.
In particular, in these fields an essentially self-dual primitive
coefficient representation satisfying D-BI-09 has even degree at
least twenty-four.  This last field assertion concerns $V$ itself;
it is not imposed on a larger-field inducing representation.
\end{proposition}
\begin{proof}
First a primitive irreducible finite representation has no nontrivial
linear self-twist.  If an invertible intertwiner $T$ satisfies
$T\rho(g)=\psi(g)\rho(g)T$, where $\psi$ has order $d>1$,
then $T^d$ is scalar.  Rescale $T$ so that $T^d=1$.
Its eigenspaces are permuted by the full image $\mu_d$ of $\psi$.
All $d$ eigenvalues occur, and their nonzero eigenspaces form a
transitively permuted direct-sum decomposition, contrary to primitivity.

Two nondegenerate equivariant bilinear forms with multipliers
$\alpha_1,\alpha_2$ give, by their quotient, a self-twist with
character $\alpha_2/\alpha_1$.  The preceding argument therefore
shows that the multipliers are identical and the forms proportional.
Choose an algebraic form $B$; this is possible because its defining
linear equations have algebraic coefficients and its multiplier has
finite image.  Apply any automorphism of $\overline{\mathbf Q}$
fixing $K$ to (BI.18).  The representation matrices are unchanged,
so uniqueness implies $\alpha^\gamma=\alpha$.  Thus $\alpha$
is $K$-valued.  The character is nontrivial because otherwise the
exact character field would be real.

For the specified fields the only roots of unity are $\pm1$.
The determinant character also takes values in $\{\pm1\}$.
Taking determinants in (BI.18) yields
\[
                   \det(\rho(g))^2=\alpha(g)^{\dim V}.
\]
The left side is one.  A nontrivial $\alpha$ takes the value $-1$,
which makes the right side minus one if $\dim V$ is odd.  This
contradiction proves the parity assertion.
\end{proof}

\section{Four-dimensional inducing blocks and their remaining fields}
\label{sec:rank-four-inducing}

\begin{theorem}[D-BI-11: a rank-four field obstruction]
\label{D-BI-11}
Let a finite-cover Weil sixfold source have exact imaginary quadratic
character field $K$, period rank nine, and nonsplit actual polarization.
Suppose its complex coefficient representation is induced from an
irreducible representation $\tau$ of dimension four.  Then
\[
 K\in\{\mathbf Q(\sqrt{-d}):d=1,2,3,5,6,15\}.
                                                               \tag{BI.23}
\]
In particular no four-dimensional inducing block can occur for
\[
                    d=7,11,19,43,67,163.                       \tag{BI.24}
\]
This is an all-degree assertion about the original representation;
the inducing character field is not assumed to be quadratic.

Among the eight non-Eisenstein class-number-one fields, the only
remaining possibilities have $K=\mathbf Q(i)$ or
$\mathbf Q(\sqrt{-2})$.  Write $L=\mathbf Q(\chi_\tau)$.
In these remaining cases $L$ is a proper biquadratic extension of
$\mathbf Q$, $[L:K]=2$, and its possible values are
\[
\begin{array}{c|c}
 K& L\\\hline
 \mathbf Q(i)&\mathbf Q(\zeta_8),\ \mathbf Q(\zeta_{12}),\
                         \mathbf Q(i,\sqrt5)\\
 \mathbf Q(\sqrt{-2})&\mathbf Q(\zeta_8).
\end{array}                                                    \tag{BI.25}
\]
Every such remaining case has exactly twelve branch points on the
intermediate rational base: ten reflection inertias and two inertias
of moving rank two outside the reflection normal closure.
Their quotient by that normal closure is cyclic and its covering
curve has genus zero.  These are necessary conditions, not a claim
that one of the remaining families exists or has rank nine.
\end{theorem}
\begin{proof}
Write the original complex representation as
$\rho=\operatorname{Ind}_H^G\tau$ and let $J$ be the faithful
image of $\tau$.  If $\tau$ is imprimitive, a system of blocks in
dimension four has block dimension one or two.  Transitivity of
induction then makes $\rho$ respectively monomial, or induced
from a two-dimensional representation.  D-MON-01--03 and D-BI-01
exclude these cases.  We may assume $\tau$ primitive.

Pass to the actual connected $J$-cover $C_0\to X_0$ obtained by
quotienting the kernel in the intermediate $H$-cover.  Frobenius
reciprocity identifies the complex multiplicity variation with
the original one, including its period differential kernel, as
in D-BI-01.  If $h=g(X_0)$ and $D$ is the total moving dimension,
then $D=14-8h$.  At $h=1$ the period variation factors through
at most six marked-base parameters, and $h\ge2$ is impossible.
Thus rank nine forces $h=0$ and
\[
 r\ge12,\qquad D=14,\qquad
 \sum_{j=1}^r\bigl(\operatorname{rank}(1-g_j)-1\bigr)=14-r\le2.
                                                               \tag{BI.26}
\]
In particular at least ten of the generating inertia matrices are
reflections, and no inertia has moving rank four.

Let $R\lhd J$ be the subgroup generated by all reflections of
$J$.  It is nontrivial.  Its fixed subspace is $J$-stable and
therefore zero.  The nonorthogonality components of its reflection
root lines give an orthogonal decomposition permuted by $J$.
Primitivity implies that there is only one component, so $R$ is
an irreducible rank-four reflection group.

There are at most two inertia matrices outside $R$.  There cannot
be exactly one, since their product in $J/R$ is one.  If there are
none, the inertia generators give $J=R$.  Otherwise there are
exactly two; (BI.26) forces their moving ranks to be two, $r=12$,
and every other inertia to be a reflection.  The quotient $J/R$
is generated by the coset of either outside inertia, and the other
coset is its inverse.  Hence it is cyclic and
\[
                         J=\langle R,g\rangle               \tag{BI.27}
\]
for either outside inertia $g$.  The quotient cover
$C_0/R\to\mathbf P^1$ has exactly two branch points, both with
full cyclic inertia, and Riemann--Hurwitz gives genus zero.

We now use the published rank-four reflection classification
\cite[Theorem~5.12, p.~435]{FIELD:Cohen}.  An imprimitive $R$ has
monomial normalizer except for $G(2,2,4)=W(D_4)$, by
\cite[Proposition~2.13, pp.~391--392]{FIELD:Cohen}.  The monomial
case is already excluded.  The primitive real possibilities are
$A_4,F_4,H_4$, and the nonreal primitive possibilities are
$G_{29},G_{31},G_{32}$.  The $A_4$ normalizer is the scalar
standard $\mathfrak S_5$ group, and is excluded by D-BI-08.
The remaining normalizer models can be chosen as follows:
\[
\begin{array}{c|c|c}
 R&\text{finite normalizer model }N\text{ up to scalars}
                                      &\text{matrix field }F\\\hline
 W(D_4)&W(F_4)&\mathbf Q\\
 W(F_4)&\langle W(F_4),\Phi/\sqrt2\rangle&\mathbf Q(\sqrt2)\\
 W(H_4)&W(H_4)&\mathbf Q(\sqrt5)\\
 G_{29},G_{31}&R&\mathbf Q(i)\\
 G_{32}&R&\mathbf Q(\sqrt{-3}).
\end{array}                                                    \tag{BI.28}
\]

Here the $D_4$ assertion is exactly Cohen's Proposition~2.13.
For $F_4$ the rational simple-root matrices and $\Phi$ are those
of D-GF4-01.  The matrix $\Phi/\sqrt2$ interchanges the two
root-length classes, squares to $-I$, and normalizes $W(F_4)$.
An orthogonal normalizer preserving the two classes is a root-system
automorphism and belongs to $W(F_4)$, by the chamber argument and
the absence of a nontrivial oriented-diagram automorphism.  This
gives the displayed finite group of order $2304$ and its field.
For $H_4$ all roots have one length and its Coxeter diagram has
no nonidentity automorphism, so the same chamber argument gives
$N=W(H_4)$.  Its standard roots and reflection matrices are over
$\mathbf Q(\sqrt5)$.  These real-normalizer reductions also follow
from Cohen's real-normalizer Lemma~2.12, p.~391, followed by the
root-system argument.  For the three nonreal groups the scalar
normalizer assertion is the \emph{corrected} Proposition~5.14:
\cite[p.~435]{FIELD:Cohen}\cite[p.~613]{FIELD:CohenErr}.
The original extra rank-four exception was deleted by that erratum.
The matrix fields are supplied by the standard reflection models
and \cite[Table~8.2, pp.~675--676]{FIELD:Malle}; only their role
as fields of definition is needed at this stage.

If $J=R$, its character field is contained in the displayed field
of its reflection model.  This either gives $K=\mathbf Q(i)$ or
$\mathbf Q(\sqrt{-3})$, or is incompatible with $K$ being imaginary.
In the remaining case write the outside inertia from (BI.27) as
\[
                           g=\lambda w,\qquad w\in N.
\]
Both $g$ and $w$ have finite order, so the commuting scalar
$\lambda$ is a root of unity.  Since $g$ fixes a two-dimensional
subspace, $\lambda^{-1}$ occurs with multiplicity two as an
eigenvalue of the matrix $w$ over $F$.  Its irreducible minimal
polynomial over $F$ consequently occurs at least twice in the
degree-four characteristic polynomial of $w$.  Therefore
\[
                 [F(\lambda):F]\le2,
                 \qquad K\subset L\subset F(\lambda).
                                                               \tag{BI.29}
\]
The first field inclusion is the induced-character identity; the
second follows from (BI.27) and the matrices in (BI.28).
There is no equality assumption between $K$, $L$, and $F(\lambda)$.

Here is the entire cyclotomic calculation in (BI.29).  If $F$ is
quadratic, the order $d$ of $\lambda$ has $\varphi(d)\le4$,
so $d\in\{1,2,3,4,5,6,8,10,12\}$.  This finite list follows
directly from $p-1\mid\varphi(d)$ for each prime $p\mid d$
and the prime-power formula for $\varphi$.  The degree condition
in (BI.29) gives exactly the following upper possibilities:
\[
\begin{array}{c|c|c}
 F&d&\text{imaginary quadratic subfields of all }F(\lambda)\\\hline
 \mathbf Q&1,2,3,4,6&-1,-3\\
 \mathbf Q(\sqrt2)&1,2,3,4,6,8&-1,-2,-3,-6\\
 \mathbf Q(\sqrt5)&1,2,3,4,5,6,10&-1,-3,-5,-15\\
 \mathbf Q(i)&1,2,3,4,6,8,12&-1,-2,-3\\
 \mathbf Q(\sqrt{-3})&1,2,3,4,6,12&-1,-3.
\end{array}                                                    \tag{BI.30}
\]
The entries in the last column denote the square-free radicands.
For example $\mathbf Q(\sqrt2,i)=\mathbf Q(\zeta_8)$ has
quadratic subfields with radicands $2,-1,-2$, whereas
$\mathbf Q(\sqrt2,\sqrt{-3})$ has radicands $2,-3,-6$.
For $F=\mathbf Q(\sqrt5)$ and $d=5,10$, the field is
$\mathbf Q(\zeta_5)$, whose sole quadratic subfield is the real
field $\mathbf Q(\sqrt5)$.  Every other degree-four field in
the table is biquadratic; its three quadratic subfields follow by
multiplying the two generating square classes.  This proves
(BI.23)--(BI.24).

We finish by establishing the proper-field assertion, which is
relevant to the actual polarization.  In this paragraph restrict
to $K=\mathbf Q(i)$ or $\mathbf Q(\sqrt{-2})$, and suppose $L=K$.
Frobenius reciprocity and irreducibility of $\rho$ give
\[
           \langle\tau,\operatorname{Res}_H^G\rho\rangle=1.
\]
The central $\tau$-projector of $K[H]$, acting on the given
$K$-model of $\rho$, therefore has a four-dimensional $K$-image
whose complexification is $\tau$.  Thus $\tau$ has a $K$-model;
no Schur-index assumption was silently made.  Its induced $K$-model
is isomorphic to the original one, and polarized Shapiro identifies
the actual multiplicity forms up to a positive rational scalar,
as in D-FIELD-03.  A scalar of this kind preserves splitness.
For $K=\mathbf Q(\sqrt{-2})$, the reflection-source theorem
D-FIELD-03 forces split determinant.  For $K=\mathbf Q(i)$,
the rank-four cases D-GAUSS-01 and D-G31-02 force the same
conclusion; the only Gaussian reflection component left by
D-GAUSS-03 has rank eight, not four.  This contradicts the
assumed nonsplit actual polarization.  Therefore $L\ne K$.

For $K=\mathbf Q(i)$ or $\mathbf Q(\sqrt{-2})$, inspection of
(BI.30) now gives precisely (BI.25).  Each ambient field containing
one of these $K$ is either that quadratic field itself or one of
the listed degree-four fields; once $K\subsetneq L$, no intermediate
choice remains.  This proves the theorem, including the exact
twelve-branch profile for all remaining cases.
\end{proof}

\section{A maximal-isotropic boundary for real rank-four inducing groups}
\label{sec:rank-four-real-boundary}

\begin{theorem}[D-BI-12: removal of the real rank-four cases]
\label{D-BI-12}
In the situation of D-BI-11, the reflection normal closure $R$ of
the four-dimensional primitive inducing image cannot be a real
reflection group.  This conclusion allows a larger inducing field
$L\supsetneq K$: the maximal isotropic space constructed in the
proof is defined over the original quadratic field $K$.

Consequently, among the eight non-Eisenstein class-number-one fields,
the remaining rank-four cases of D-BI-11 have
\[
\begin{array}{c|c|c}
 K&R&L\\\hline
 \mathbf Q(i)&G_{29},G_{31}&\mathbf Q(\zeta_8)\text{ or }
                                      \mathbf Q(\zeta_{12})\\
 \mathbf Q(i)&G_{32}&\mathbf Q(\zeta_{12})\\
 \mathbf Q(\sqrt{-2})&G_{29},G_{31}&\mathbf Q(\zeta_8).
\end{array}                                                    \tag{BI.31}
\]
All rows still require the twelve-branch profile and field descent
in D-BI-11.  No existence or dominance assertion is made for them.
\end{theorem}
\begin{proof}
Suppose $R$ is real.  If $J=R$, its character field is real, which
contradicts $K\subset\mathbf Q(\chi_\tau)$.  The remaining branch
pattern in D-BI-11 therefore consists of ten reflections and two
rank-two outside inertias.  Hurwitz moves place the latter at the
end of the tuple.  Moving an outside element past a reflection
conjugates that reflection by a normalizer of $R$, so all ten
reflection entries stay real reflections in $R$.  We may work after
forgetting or reordering the branch labels; these are finite changes
and leave the period rank unchanged.

Choose real orthogonal coordinates for $R$.  Every complex unitary
normalizer is scalar times real orthogonal, by Cohen's Lemma~2.12
and uniqueness of the positive invariant form.  Write the two
outside inertia matrices as
\[
                   g_1=\lambda A_1,\qquad
                   g_2=\lambda^{-1}A_2,                      \tag{BI.32}
\]
where $A_1,A_2$ are real orthogonal and $\lambda$ is nonreal of
absolute value one.  To justify the choices, the product $g_1g_2$
belongs to $R$, so the product of their two unit scalar phases is
real, hence $1$ or $-1$; absorb the latter sign into $A_2$.
If both phases were real, every inertia and hence $J$ would have
real character, again contradicting $K\subset L$.  Thus both
phases are nonreal and (BI.32) is valid.

Put $w=(g_1g_2)^{-1}\in R$.  The ten reflections give a
factorization of $w$.  Let $\ell$ be its shortest reflection length
in the finite real reflection group $R$.  Lewis--Reiner's published
pair-extraction result applies to an arbitrary element, not only
to a Coxeter element:
\cite[Corollary~1.4, p.~1459; proof \S4.3, pp.~1480--1481]{BI:LR}.
It changes this length-ten factorization by Hurwitz moves into
\[
 (r_1,r_1,\ldots,r_q,r_q,s_1,\ldots,s_\ell),
 \qquad q=\frac{10-\ell}{2},\qquad s_1\cdots s_\ell=w,
                                                               \tag{BI.33}
\]
where the final factorization is shortest.  Carter's criterion,
in the precise real-root-system form stated in
\cite[Proposition~4.1, p.~1473]{BI:LR}, says that the roots of
$s_1,\ldots,s_\ell$ are linearly independent.  In particular
$\ell\le4$.

We record explicitly the fixed-space consequence.  If those roots
are $\alpha_1,\ldots,\alpha_\ell$, the telescoping identity
\[
 s_1\cdots s_\ell-1
       =\sum_{j=1}^{\ell}s_1\cdots s_{j-1}(s_j-1)
\]
has image terms along the roots
$\beta_j=s_1\cdots s_{j-1}\alpha_j$.
Each $\beta_j$ is $\alpha_j$ plus a combination of earlier
$\alpha_i$, so the $\beta_j$ are independent.  Applying the
identity to a fixed vector of $w$ shows that every coefficient,
and hence every $(s_j-1)x$, is zero.  Therefore
\[
 \ker(w-1)=\bigcap_{j=1}^{\ell}\ker(s_j-1),\qquad
                \dim\ker(w-1)=4-\ell.                      \tag{BI.34}
\]
The parity of the original ten-reflection factorization gives
$\ell\in\{0,2,4\}$ and $q\in\{5,4,3\}$.

Let $H_0\le J$ be generated by
$s_1,\ldots,s_\ell,g_1,g_2$, and set
\[
 E=\ker(g_1-1)\cap\ker(g_2-1),\qquad h_0=\dim_{\mathbf C}E.
\]
Every vector of $E$ is fixed by $w$, and (BI.34) shows that it
is fixed by all $s_j$.  Hence $h_0=\dim\tau^{H_0}$.
Moreover $E$ and its complex conjugate $\overline E$ are disjoint:
$A_1$ has eigenvalue $\lambda^{-1}$ on $E$ and eigenvalue
$\lambda$ on $\overline E$, and these are distinct because
$\lambda$ is nonreal.  Both subspaces lie in the real fixed space
of $w$, after complexification.  Thus
\[
                         2h_0\le4-\ell.                    \tag{BI.35}
\]

We now construct the degeneration and compute its graph module.
For every equal pair $(r_j,r_j)$ in (BI.33), bubble off a rational
base component containing those two branch points and one node.
Retain all $s_i,g_1,g_2$ on the central rational component.
Since $r_j^2=1$, the node inertia is trivial.  The local cover over
each leaf has monodromy $\langle r_j\rangle$ and is a disjoint
union of rational double covers.  The local cover over the central
component has monodromy $H_0$.  Gluing the unramified node fibres
produces the admissible $J$-cover determined by the tuple, and
smoothing the nodes recovers the given Hurwitz component.  All
base components are stable: each leaf has three special points,
and the central component has $\ell+2+q\ge3$.

Let $\Gamma$ be the dual graph of the source curve.  There is one
copy of the regular $J$-set of edges for each leaf, while its vertex
sets are $J/\langle r_j\rangle$ and $J/H_0$.  The graph is
connected because the original tuple, and hence
$H_0,r_1,\ldots,r_q$, generates $J$.  Its exact cellular sequence is
\[
 0\longrightarrow H_1(\Gamma,\mathbf C)
 \longrightarrow\mathbf C[J]^q
 \longrightarrow\mathbf C[J/H_0]\oplus
       \bigoplus_{j=1}^{q}\mathbf C[J/\langle r_j\rangle]
 \longrightarrow\mathbf C\longrightarrow0.
                                                               \tag{BI.36}
\]
Taking $\tau$-multiplicities is exact.  The regular representation
has multiplicity four, each reflection has a three-dimensional
fixed subspace, and $\tau$ has no invariant vector for $J$.
Therefore (BI.35)--(BI.36) give
\[
 \dim\operatorname{Hom}_J(\tau,H_1(\Gamma,\mathbf C))
       =4q-3q-h_0=q-h_0
       \ge\frac{10-\ell}{2}-\frac{4-\ell}{2}=3.
                                                               \tag{BI.37}
\]

For clarity, this graph dimension is the rank of one actual
monodromy logarithm.  Smooth all $q$ node parameters along a
one-parameter arc with positive integer orders.  The vanishing
cycles are disjoint.  The sum $N$ of their positively weighted
Picard--Lefschetz logarithms satisfies $N^2=0$; its rank on a
coefficient multiplicity equals the corresponding graph multiplicity.
One can see the equality directly from (BI.36): the independent
vanishing cycles form the edge-cycle space, and the monodromy
pairing on that space is the restriction of the positive diagonal
edge-length form.  It is nondegenerate, so no additional kernel
appears after taking the finite-group isotypic summand.  This is
the same semistable-curve graph calculation used in D-G34-03.
The smooth $\tau$-multiplicity has dimension six; since $N^2=0$,
its rank is at most three.  Equation (BI.37) forces
\[
                             \operatorname{rank}N=3.          \tag{BI.38}
\]

It remains to descend this conclusion to the original field $K$.
The intermediate Hurwitz component has exactly twelve branch points
and dimension nine.  The original period map factors through it
with rank nine, so the original family's classifying map to this
component is dominant.  Choose a subvariety of the original base
mapping dominantly and generically finitely onto the component;
one obtains it by choosing a closed point of the nonempty generic
fibre, whose residue field is a finite extension of the component's
function field, and taking its closure.  After shrinking a nonempty open and normalizing,
this gives a finite etale cover on which the original $K$-variation
is defined and whose chosen complex multiplicity is the intermediate
$\tau$-variation.

Choose the preceding one-parameter smoothing arc generically so
that its punctured part lies in this nonempty open.  After a finite
base change it lifts to that finite cover.  Its monodromy on the
original $K$-variation is therefore a positive power of the
intermediate monodromy.  The logarithm $N_K$ is $K$-linear,
has square zero, and its rank is three because the rank after the
chosen complex embedding is (BI.38).  This argument neither
identifies $L$ with $K$ nor asserts that an $L$-isotropic subspace
automatically descends.

The original polarization makes $N_K$ skew adjoint.  Hence
\[
 h(N_Kx,N_Ky)=-h(x,N_K^2y)=0.
\]
Its image is a three-dimensional totally isotropic $K$-subspace
of the six-dimensional original Hermitian space.  That Hermitian
form is split, contrary to the hypothesis.  This proves the
exclusion of all real reflection normal closures and leaves
exactly the nonreal rows (BI.31) from D-BI-11.
\end{proof}

\section{No reflection-containing inducing block in six quadratic fields}
\label{sec:large-prime-inducing}

\begin{theorem}[D-BI-13: a ramification obstruction for every inducing block]
\label{D-BI-13}
Let
\[
 p\in\{7,11,19,43,67,163\},\qquad K=\mathbf Q(\sqrt{-p}).
\]
Suppose a finite-cover Weil sixfold source has exact character field
$K$, period rank nine, and nonsplit actual polarization.  If its
complex coefficient representation is expressed as
\[
                       \rho=\operatorname{Ind}_H^G\tau,
\]
then the finite image of $\tau$ contains no complex reflection.
The inducing representation may have arbitrary degree and arbitrary
character field.  In particular the obstruction applies to every
inducing block, not only to the original $K$-representation or to
blocks already defined over $K$.
\end{theorem}
\begin{proof}
Suppose the image of $\tau$ contains a reflection.  Let $R$ be
the normal subgroup generated by its reflections.  Its fixed space
is stable under the irreducible inducing image and hence is zero.
The nonorthogonality components of the reflection root lines give
a direct sum of irreducible reflection spaces permuted transitively
by that image.  Stabilize one component, of dimension $e$, and
denote its stabilizer representation by $\sigma$.  Transitivity
of induction gives
\[
                    \rho=\operatorname{Ind}_{H'}^G\sigma
\]
for a subgroup $H'\le H$.  Its faithful image $J$ contains an
irreducible reflection group $R_0$ of rank $e$ and normalizes it.
No rationality of this component or its field is asserted.  The
induced-character formula nevertheless gives the unconditional
inclusion
\[
                    K\subset L=\mathbf Q(\chi_\sigma).
                                                               \tag{BI.39}
\]

Ranks one, two, three and four are excluded, respectively, by
the monomial theorem D-MON-01--03, D-BI-01, D-BI-05, and
D-BI-11.  For $p=7$, the rank-three step uses the proved split
actual polarization of the $G_{24}$ family, with the polarized
Shapiro comparison in D-BI-05; this is where the nonsplit
hypothesis remains essential.  These are all-induction results with the larger field
retained, so they apply to $\sigma$.
If $R_0$ is imprimitive of rank $e\ge5$, Cohen's normalizer
Proposition~2.13 puts $J$ in the scalar monomial group
\cite[pp.~391--392]{FIELD:Cohen}; its low-rank exceptions have
already been discarded.  Then $\sigma$, and hence $\rho$, is
monomial, another contradiction.  If $R_0$ is a real group of
type $A_e$, its scalar normalizer is the deleted permutation
normalizer and D-BI-08 applies.  This also covers every rank in
that infinite family, without assuming that the scalar and
permutation parts form a direct product.

The published reflection classification now leaves precisely
\[
\begin{array}{c|c|c|c}
 R_0&e&\text{finite normalizer model }N\text{ up to scalars}&F\\\hline
 G_{33}&5&G_{33}&\mathbf Q(\zeta_3)\\
 G_{34}&6&G_{34}&\mathbf Q(\zeta_3)\\
 W(E_6)&6&\langle W(E_6),-I\rangle&\mathbf Q\\
 W(E_7)&7&W(E_7)&\mathbf Q\\
 W(E_8)&8&W(E_8)&\mathbf Q.
\end{array}                                                    \tag{BI.40}
\]
The classification is Cohen's Theorem~5.12, p.~435, and the
model fields are in Malle's Table~8.2, p.~675, with its explanation
on p.~676.  The omitted $G_{35}$ is the rational Weyl group
$W(E_6)$.  For $G_{33},G_{34}$ the normalizer assertion is
Cohen's corrected Proposition~5.14; its sole exception is rank
three and is already irrelevant
\cite{FIELD:Cohen,FIELD:CohenErr,FIELD:Malle}.
For the real rows, a normalizer is scalar times real by
Cohen's Lemma~2.12.  A root-system automorphism is a Weyl element
times a diagram automorphism.  For $E_6$ the diagram involution
is represented by $-w_0$, so adjoining $-I$ suffices; $E_7,E_8$
have no nontrivial diagram automorphisms and already contain
$-I$.  These are finite rational matrix groups.  Thus in every
row
\[
                    J\subset\mathbf C^*N,
                    \qquad N\subset\operatorname{GL}_e(F),
                    \qquad e\le8.                          \tag{BI.41}
\]

We next verify the branch count on which the argument depends.
Use the actual intermediate connected cover for $\sigma$.
Its cohomological multiplicity is six by Frobenius reciprocity,
and its full quadratic-field period rank factors through that
cover's marked-base moduli.  If the base has genus $h$, then
$D=6-2e(h-1)$.  Genus one gives at most six parameters, and
genus at least two is impossible because $e\ge5$.  Hence rank
nine forces a rational base and a product-one generating tuple
$g_1,\ldots,g_r$ of nonidentity inertia matrices with
\[
                  r\ge12,
                  \qquad\sum_j\operatorname{rank}(1-g_j)=2e+6.
                                                               \tag{BI.42}
\]

Write $g_j=\lambda_j n_j$ with $n_j\in N$.  Since both $g_j$
and $n_j$ have finite order and the scalar commutes with $n_j$,
each $\lambda_j$ is a root of unity.  The ambiguity in its choice
is exactly a scalar in $Z_0=N\cap\mathbf C^*I$.  Such a scalar
belongs to $F$, so its order divides two or six in (BI.40), and
in particular is prime to $p$.
It follows that the $p$-primary part of $\lambda_j$ is independent
of the chosen expression.  More generally these primary parts
define a character
\[
              \beta_p:J\longrightarrow\mu_{p^a}
                                                               \tag{BI.43}
\]
for some finite $a\ge0$.  One can see the homomorphism directly
from multiplication, or from the scalar quotient
$\mathbf C^*N/N\simeq\mathbf C^*/Z_0$ and its finite cyclic
image on $J$.  No scalar projection on $J$ before quotienting
$Z_0$ is being assumed.

The character $\beta_p$ is nontrivial.  Indeed, if all generating
phases $\lambda_j$ had order prime to $p$, then all matrices of
$J$ and all its character values would belong to
\[
                    E=F(\lambda_1,\ldots,\lambda_r),
\]
a cyclotomic field, or a subfield of a cyclotomic field, with
conductor prime to $p$.  Such a field is unramified at $p$.
For an elementary justification, when $p\nmid M$ the polynomial
$X^M-1$ and its derivative have no common root modulo $p$;
the cyclotomic factors therefore define etale orders locally at
$p$, and their number fields are unramified there.  Subfields
retain that property.  Here $F$ is either $\mathbf Q$ or
$\mathbf Q(\zeta_3)$, so it introduces no ramification at $p$.
But (BI.39) would put $K$ inside $E$, whereas
$\mathbf Q_p(\sqrt{-p})/\mathbf Q_p$ is ramified quadratic by
the Eisenstein polynomial $X^2+p$.  This is impossible.

Because $g_1\cdots g_r=1$, at least two of the $g_j$ have
$\beta_p(g_j)\ne1$.  A nontrivial $p$-primary phase implies
$\mathbf Q(\zeta_p)\subset\mathbf Q(\lambda_j)$, and
\[
 [F(\lambda_j):F]\ge[F(\zeta_p):F]=p-1.
                                                               \tag{BI.44}
\]
For $F=\mathbf Q$ the equality is the usual cyclotomic degree.
For $F=\mathbf Q(\zeta_3)$, the intersection with
$\mathbf Q(\zeta_p)$ is $\mathbf Q$: otherwise the quadratic
field $F$, ramified at three, would lie in a field unramified
at three.  Thus the equality holds in this case as well.
If $g_j$ fixes a subspace of dimension $b$, then
$\lambda_j^{-1}$ is an eigenvalue of $n_j$ of multiplicity $b$.
Its irreducible polynomial over $F$, of degree at least $p-1$,
therefore occurs at least $b$ times in the degree-$e$
characteristic polynomial of $n_j$.  Consequently
\[
 \dim\ker(g_j-1)\le\left\lfloor\frac{e}{p-1}\right\rfloor,
 \qquad
 \operatorname{rank}(1-g_j)\ge
              e-\left\lfloor\frac{e}{p-1}\right\rfloor.
\]
For $p\ge11$ this is full moving rank.  For $p=7$ and $e\le8$
it is at least $e-1$.

The two distinguished inertias have the preceding lower bound,
and every other genuine inertia contributes at least one.  Hence
\[
 2e+6\ge2e-2\left\lfloor\frac{e}{p-1}\right\rfloor+(r-2)
                   \ge2e+8,
\]
a contradiction.  It rules out the original assumption that any
inducing image contains a reflection, and proves the theorem.
\end{proof}

\section{Large degree bounds for scalar groups with unramified models}
\label{sec:scalar-unramified}

\begin{theorem}[D-BI-14: the scalar-unramified degree bound]
\label{D-BI-14}
Let $K=\mathbf Q(\sqrt{-p})$ for
$p\in\{7,11,19,43,67,163\}$, and suppose a finite-cover
Weil sixfold source of exact character field $K$ has period rank
nine and nonsplit actual polarization.  Suppose an inducing
representation of degree $e$ has faithful image $J$ satisfying
\[
                 J\subset\mathbf C^*N,
                 \qquad N\subset\operatorname{GL}_e(F),
                                                               \tag{BI.45}
\]
where $N$ is finite and the number field $F$ is unramified at $p$.
The field $F$ need not be Galois, and need not contain $K$ or the
inducing character field.  Then
\[
                  e\ge7(p-1).
                                                               \tag{BI.46}
\]
The six bounds, in the displayed order of $p$, are
\[
                   42,\ 70,\ 126,\ 294,\ 462,\ 1134.
\]
If the inducing image is also primitive and tensor indecomposable
in the Guralnick--Saxl convention, then
\[
                  e\ge12(p-1),
                                                               \tag{BI.47}
\]
giving respectively $72,120,216,504,792,1944$.
More generally, if every nonidentity matrix of $J$ has moving rank
at least $q$, the same argument gives
\[
                  e\ge(5q-3)(p-1).                         \tag{BI.48}
\]
The estimate (BI.48) for $q=1$ gives the weaker bound
$2(p-1)$ without invoking D-BI-13's absence of reflections.
\end{theorem}
\begin{proof}
Let $\tau$ denote the inducing representation and
$L=\mathbf Q(\chi_\tau)$.  Induction of characters gives
$K\subset L$.  D-BI-13 says that $J$ has no reflections.
In particular every genuine inertia has moving rank at least two.
There are no degree-one cases; a faithful nontrivial image in
dimension one would consist of reflections.  On the actual
intermediate base of genus $h$, the multiplicity equation is
$D=6-2e(h-1)$.  Genus one has at most three genuine branch points
and at most three marked-base parameters.  For $h\ge2$ the
only numerical possibilities have $h=2,e\le3$ and at most
$3-e$ branch points, hence dimension at most $6-e<9$.
Consequently the base is rational and
\[
 r\ge12,\qquad \sum_j\operatorname{rank}(1-g_j)=2e+6
                                                               \tag{BI.49}
\]
for a generating product-one tuple of genuine inertias.
For the stated $q=1$ variant, one can instead discard $e=1$
by D-MON and use the weaker bound of $D$ genuine branches:
genus one gives at most six parameters, and genus at least two
gives at most five.  Thus (BI.49) still follows without the
reflection-free conclusion of D-BI-13.

We provide the field details for arbitrary $F$.  The extension
$\mathbf Q(\zeta_p)/\mathbf Q$ is totally ramified at $p$:
$\Phi_p(X+1)$ is Eisenstein at $p$, as its constant term is
$p$, all nonleading coefficients are divisible by $p$, and its
leading coefficient is one.  Every nontrivial subfield is therefore
ramified at $p$.  Since $F$ is unramified at $p$,
\[
              F\cap\mathbf Q(\zeta_p)=\mathbf Q,
              \qquad [F(\zeta_p):F]=p-1.                  \tag{BI.50}
\]
For completeness, the degree conclusion can also be seen locally:
base change to an unramified completion of $F$ preserves the
Eisenstein property of $\Phi_p(X+1)$, so it remains irreducible
over that completion and hence over $F$.
This avoids imposing a Galois hypothesis on $F$.

Every scalar in $Z_0=N\cap\mathbf C^*I$ is a root of unity in
$F$ and has order prime to $p$, by (BI.50).  As in D-BI-13,
write $g_j=\lambda_j n_j$ with $n_j\in N$ and a root of unity
$\lambda_j$.  Its $p$-primary part is independent of the choice
modulo $Z_0$, and defines a character $\beta_p$ of $J$.
It is nontrivial: if all generating phases had order prime to
$p$, the field $F(\lambda_1,\ldots,\lambda_r)$ would be
unramified at $p$ and would contain $L$, hence the ramified field
$K$, an impossibility.  Here adjoining roots of unity of order
prime to $p$ is unramified locally, and a compositum of unramified
local extensions is unramified.  The product-one relation supplies
at least two $g_j$ with nontrivial $p$-primary phase.

For either of these two entries, (BI.50) gives
$[F(\lambda_j):F]\ge p-1$.  If its fixed dimension is $b$,
the irreducible polynomial of $\lambda_j^{-1}$ over $F$
occurs at least $b$ times in the characteristic polynomial of
$n_j$.  Thus its moving rank is at least
$e-\lfloor e/(p-1)\rfloor$.  If all other nonidentity entries
have moving rank at least $q$, (BI.49) yields
\[
 2e+6\ge2e-2\left\lfloor\frac{e}{p-1}\right\rfloor+10q,
 \qquad
 \left\lfloor\frac{e}{p-1}\right\rfloor\ge5q-3.
\]
This proves (BI.48), and $q=2$ proves (BI.46).

Finally assume primitivity and tensor indecomposability.  The
already established bound gives $e\ge42$, so the degree-ten
exception in Guralnick--Saxl cannot occur.  If $J$ contained a
matrix of moving rank two, their characteristic-zero classification
\cite[Theorem~7.1 and Remark~(3)]{BI:GS} would make it a scalar
deleted alternating/symmetric representation.  D-BI-08 excludes
that inducing block, with arbitrary larger scalar field.  Thus
every nonidentity element now moves at least three dimensions.
Applying (BI.48) with $q=3$ proves (BI.47).
\end{proof}

\section{The full Guralnick--Saxl strengthening}
\label{sec:scalar-unramified-radical}

\begin{corollary}[D-BI-14a: a quadratic bound in the ramified prime]
Under the primitive, tensor-indecomposable hypotheses of D-BI-14,
put $d=p-1$.  Then
\[
                       e\ge\frac{d(25d-22)}4.              \tag{BI.50a}
\]
For $p=7,11,19,43,67,163$ these bounds are respectively
\[
              192,\ 570,\ 1926,\ 10794,\ 26862,\ 163134.
\]
\end{corollary}
\begin{proof}
We use the complete statement of Guralnick--Saxl
\cite[Theorem~7.1 and Remark~(3), pp.~548--549;
Lemma~7.4, pp.~550--551]{BI:GS}, rather
than just its bireflection consequence.  Their invariant
$\nu_J(V)$ is the minimum moving rank of $\lambda g\ne1$,
where $g\in J$ and $\lambda$ is an arbitrary scalar.
Outside the listed exceptions their theorem gives
\[
                    \nu_J(V)>\max\{2,\sqrt e/2\}.
\]
In characteristic zero and degree $e\ge42$, all the exceptional
families except the scalar deleted alternating/symmetric family
are absent: the natural finite classical cases are in defining
positive characteristic, and all other listed exceptions have
degree at most ten.  D-BI-08 excludes the scalar deleted family
also for a larger inducing field.  The theorem's tensor convention
is exactly the convention in D-BI-14; it includes a group that
permutes tensor factors, and the extra tensor-cube exception has
degree eight.  Thus every nonidentity element of $J$ has moving
rank at least
\[
                        q=\lfloor\sqrt e/2\rfloor+1.
\]
Apply (BI.48) to obtain $e\ge(5q-3)d$.  The definition of $q$
also gives $4(q-1)^2\le e<4q^2$, whence
\[
                         4q^2-5dq+3d>0.                  \tag{BI.50b}
\]
Here $d\ge6$ and $d\equiv2\pmod4$.  Put
$q_0=(5d-2)/4$, an integer.  The upward quadratic in (BI.50b)
is negative at $q=1$ and at $q=q_0-1$:
\[
 4-2d<0,\qquad
 4(q_0-1)^2-5d(q_0-1)+3d=9-\tfrac92d<0.
\]
By convexity it is negative throughout $1\le q\le q_0-1$.
Therefore $q\ge q_0$, and
\[
           e\ge(5q_0-3)d=\frac{d(25d-22)}4,
\]
as asserted.  This is only a necessary degree bound: no existence
or dominance follows at the numerical threshold.
\end{proof}

\section{A symmetry obstruction for a larger inducing field}
\label{sec:inducing-symmetry}

\begin{theorem}[D-BI-15: normalizer and unique-system obstruction]
\label{D-BI-15}
Let $C\to\mathbf P^1$ be a connected finite $G$-cover in a
Hurwitz family, and let $\rho$ be an absolutely irreducible
representation with a $K$-model and exact character field $K$.
Suppose $\rho=\operatorname{Ind}_H^G\tau$, put
$L=\mathbf Q(\chi_\tau)$, and assume $L\supsetneq K$.
Suppose the actual intermediate base $X=C/H$ is rational and
the faithful cover $C/\ker\tau\to X$ has $R\ge3$ genuine
branch points.  If there exist
$\gamma\in\operatorname{Gal}(\overline{\mathbf Q}/K)$ and
$g\in N_G(H)$ such that
\[
                 \tau^\gamma\not\simeq\tau,
                 \qquad \tau^\gamma\simeq\tau^g,
                                                               \tag{BI.51}
\]
then the period rank of its multiplicity variation is at most
\[
                         \lfloor R/2\rfloor-1.             \tag{BI.52}
\]
In particular a twelve-branch intermediate cover of this kind
cannot supply period rank nine.

Condition (BI.51) follows in either of the following two cases:
$H$ is normal in $G$; or the given system of imprimitivity of
$\rho$ is the unique system with its block dimension.  In the
second assertion uniqueness among primitive systems also suffices,
if the inducing representation is primitive.  No general
uniqueness theorem for imprimitivity systems is assumed.
\end{theorem}
\begin{proof}
Write $N=\ker\tau$.  Galois conjugation preserves this kernel.
Indeed a finite-order matrix has character equal to its dimension
exactly when it is the identity: its eigenvalues have modulus one,
and equality in the triangle inequality forces them all to equal
one.  The equality $\chi_\tau(h)=\dim\tau$ is preserved and
reflected by $\gamma$.  Thus (BI.51) implies $gNg^{-1}=N$.
Moreover $g\notin H$, since inner conjugation leaves the
isomorphism class of $\tau$ unchanged.

The element $g$ consequently induces a nontrivial automorphism
$\varphi$ of $X=C/H$, and lifts to an automorphism of $C/N$.
The action of $N_G(H)/H$ on $C/H$ is faithful, by the
Galois correspondence for the connected $G$-cover.  Hence
$\varphi$ has finite order $q\ge2$ and preserves the reduced
branch divisor of $C/N\to X$.  This statement concerns exactly
the genuine branch points, since an isomorphism of covers
preserves ramification.

Here is the dimension bound without any assumption about a
generic automorphism group.  Conjugate a finite-order automorphism
of $\mathbf P^1$ to $z\mapsto\zeta_qz$.  An invariant reduced
divisor of degree $R$ consists of $f\in\{0,1,2\}$ of the fixed
points $0,\infty$ and $s=(R-f)/q$ nonfixed orbits.  Since
$R\ge3$, one has $s\ge1$ and $q\le R$.  The $s$ orbit
parameters lie in $\mathbf C^*$; quotienting by the one-dimensional
centralizer of the automorphism leaves dimension $s-1$.
Choices of labels and of the fixed points are finite.  There are
only finitely many possible $q$ and $f$, and therefore the locus
of such branch configurations in $M_{0,R}$ has dimension at most
\[
             s-1\le\lfloor R/2\rfloor-1.
\]
For fixed finite monodromy, the intermediate Hurwitz map to its
marked-base moduli is finite over the smooth locus.  Frobenius
reciprocity identifies the original multiplicity variation with
the one on this intermediate cover.  Its period rank is thus
bounded by the dimension just computed.  This proves (BI.52).

It remains to justify the two sufficient conditions.  Choose
$\gamma$ fixing $K$ but nontrivial on $L$, so that
$\tau^\gamma\not\simeq\tau$.  If $H\triangleleft G$, then
\[
       \operatorname{Res}_H^G\operatorname{Ind}_H^G\tau
                      =\bigoplus_{G/H}\tau^g.
\]
Since $\rho^\gamma\simeq\rho$, the constituent $\tau^\gamma$
must be one of these conjugates.  Normality puts the corresponding
$g$ in $N_G(H)$, proving (BI.51).

For the uniqueness assertion, let $W$ be the chosen inducing
block, whose stabilizer is $H$.  Frobenius reciprocity gives
\[
 \dim\operatorname{Hom}_H(\tau,\operatorname{Res}_H^G\rho)
 =\dim\operatorname{Hom}_G(\operatorname{Ind}_H^G\tau,\rho)=1.
\]
Thus $W$ is the image of the central $\tau$-projector in
$\overline{\mathbf Q}[H]$ acting on the $K$-model of $\rho$;
in particular the subspace and its translates are defined over
$\overline{\mathbf Q}$.  Applying $\gamma$ to this system gives
another $G$-stable system of the same block dimension, because
all matrices of $\rho$ are fixed by $\gamma$.  If the original
inducing representation is primitive, so is its Galois conjugate:
induction from a proper subgroup commutes with Galois conjugation,
and the existence of an imprimitivity system is equivalent to
such an induction.  The stated uniqueness hypothesis therefore
implies $W^\gamma=gW$ for some $g\in G$.
But
\[
 \operatorname{Stab}_G(W^\gamma)=H,
 \qquad\operatorname{Stab}_G(gW)=gHg^{-1},
\]
so $g\in N_G(H)$.  The two representations of $H$ on this
common block are respectively $\tau^\gamma$ and a conjugate of
$\tau$ by $g$ (or by its inverse, according to the induction
convention).  Replacing $g$ by its inverse if needed establishes
(BI.51).  This completes both sufficient-condition arguments.
\end{proof}

\begin{corollary}[Necessary asymmetry of the remaining twelve-branch cases]
For a rank-nine source with a strictly larger inducing field and
a twelve-branch rational intermediate base, the inducing subgroup
is nonnormal and its specified system of imprimitivity cannot be
unique among systems of that block dimension.  For a primitive
inducing block it cannot be the unique primitive system either.
This leaves the nonnormal, nonunique cases open; it does not
exclude them by a field-descent assumption.
\end{corollary}

\section{Inducing dimensions six and seven}
\label{sec:rank-six-seven-inducing}

Throughout this section $K$ is one of the eight non-Eisenstein
imaginary quadratic fields of class number one.  A source means
a connected finite-cover source with exact character field $K$,
actual Weil sixfold signature $(3,3)$, period rank nine, and
nonsplit actual polarization.  The inducing character field is
not assumed equal to $K$.

\begin{theorem}[D-BI-16: the remaining degree-six inducing group]
\label{D-BI-16}
If the coefficient representation of such a source is induced
from an irreducible representation $\tau$ of dimension six,
its faithful image and character field necessarily satisfy
\[
 K=\mathbf Q(i),\qquad
 L=\mathbf Q(\chi_\tau)=\mathbf Q(\zeta_{12}),\qquad
 R=G_{34},\qquad J=\mu_4G_{34}.                              \tag{BI.53}
\]
Here $R$ is the reflection normal closure in $J$; it has order
$39\,191\,040$ and center $\mu_6$, so $|J|=78\,382\,080$.
The actual intermediate base is rational, has $12\le r\le14$
genuine branch points, and has exactly two inertia entries outside
$R$.  Up to interchanging them, the complete necessary list of
moving-rank profiles is
\[
\begin{array}{c|c|c}
 r&\text{two outside moving ranks}&\text{inside moving ranks}\\\hline
14&(3,3)&1^{12}\\
13&(3,3)&1^{10},2\\
13&(3,4)&1^{11}\\
12&(3,3)&1^8,2,2\ \text{or}\ 1^9,3\\
12&(3,4)&1^9,2\\
12&(3,5)&1^{10}\\
12&(4,4)&1^{10}.
\end{array}                                                  \tag{BI.54}
\]
Exponents in this table denote repetitions, not powers of group
elements.  This table is necessary and does not claim that each
row occurs in a generating product-one tuple.

An outside rank-three inertia is an involution with spectrum
$1^3,(-1)^3$.  An outside rank-four inertia has spectrum
$1^2,(-1)^2,\alpha,\beta$, where, writing
$\zeta=\zeta_{12}$, the unordered residual pair is one of
\[
 \{i,i\},\ \{-i,-i\},\ \{i,-i\},\
 \{\zeta,\zeta^5\},\ \{\zeta^7,\zeta^{11}\},\
 \{\zeta,\zeta^{11}\},\ \{\zeta^5,\zeta^7\}.             \tag{BI.55}
\]
Its order is four or twelve, and its square is a rank-two
element of order two or six, respectively.  None of these
possibilities is asserted to occur merely from the eigenvalue
list.
\end{theorem}
\begin{proof}
Let $H$ be the inducing subgroup, so that the original
representation is $\operatorname{Ind}_H^G\tau$, and put
$L=\mathbf Q(\chi_\tau)$.  Induction gives $K\subset L$.
An imprimitive degree-six $\tau$ has an inducing block of
dimension one, two, or three.  Transitivity of induction and
D-MON, D-BI-01, and D-BI-05 exclude all three possibilities.
Thus $J$ is primitive.

On the actual intermediate base of genus $h$, the multiplicity
Euler equation is $D=18-12h$.  Genus one has at most six genuine
branches and at most six marked-base parameters; higher genus
is impossible.  Frobenius reciprocity for the actual variation
therefore forces
\[
          h=0,\qquad r\ge12,\qquad
          \sum_j\operatorname{rank}(1-g_j)=18.              \tag{BI.56}
\]
At least six of the inertias are reflections, since every other
entry has moving rank at least two.  The reflection normal
closure $R$ is consequently nontrivial.  Its fixed space is
$J$-stable and zero, and its irreducible root components would
give an imprimitivity system unless there were only one.  Thus
$R$ is irreducible.  If it were imprimitive, Cohen's
Proposition~2.13 would put its normalizer in a monomial group,
as none of its exceptional dimensions is six; this is already
excluded.  The real type $A_6$ is a scalar standard case excluded
by D-BI-08.  The complete reflection classification leaves
\[
                         R=W(E_6)\quad\text{or}\quad G_{34}.
\]
These are the published classification and corrected normalizer
inputs used in D-BI-13, with all lower-dimensional exceptions
already discarded.

First take $R=W(E_6)$.  Its normalizer is contained in
$\mathbf C^*R$: after the real-normalizer reduction, the
nontrivial diagram automorphism is $-w_0$, where $w_0$ is the
longest Weyl element.  Thus all matrices are scalar multiples
of the rational matrices of $R$.  Count the branch entries with
nonreal scalar phase by $t$.  Each moves at least three
dimensions, because a nonreal eigenvalue of a real matrix has
the same multiplicity as its distinct conjugate.  Therefore
$18\ge r+2t$, so $t\le3$.  There cannot be exactly one such
entry, by the scalar product-one relation modulo $\{1,-1\}$;
there cannot be none, since then $L$ would be real.

If $t=3$, the three moving ranks are all three.  Each scalar
phase has minimal polynomial over $\mathbf Q$ of degree at
most two, because its inverse has multiplicity three in a
degree-six rational characteristic polynomial.  Its order is
three, four, or six.  Modulo real signs the phases have order
three or two, and the product of three nonidentity such phases
can be one only if all three have order three.  Hence all
matrices are defined over $\mathbf Q(\omega)$, which cannot
contain the stipulated non-Eisenstein $K$.

If $t=2$, the two moving ranks sum to at most eight.  One of
the corresponding real matrices has the inverse scalar as an
eigenvalue of multiplicity at least two.  The scalar degree over
$\mathbf Q$ is at most three, and a nonreal root of unity of
such degree has degree two.  The two phases are inverse modulo
a real sign, and the other phases are real.  Since $R\subset J$,
the scalar matrices belong to $J$ and their traces show that
$L$ is exactly the quadratic scalar field.  Thus $L=K=\mathbf Q(i)$.
Multiplicity-one Frobenius reciprocity supplies a $K$-model of
$\tau$ via its central projector on the original $K$-model.
The polarized Shapiro comparison is then a positive rational
scalar comparison of the actual Hermitian forms.  The already
proved Gaussian $E_6$ source barrier forces splitness, a
contradiction.  This eliminates the real closure.

Now take $R=G_{34}$.  It has a model over $F=\mathbf Q(\omega)$,
all its reflections have order two, and its degrees are
$6,12,18,24,30,42$.  In particular its center is $\mu_6$ and
its character field is exactly $F$: the model gives one inclusion
and the trace of $\omega I$ gives the other.  These are Cohen's
root model/classification and Malle's model-field data, already
checked in D-BI-13.  The corrected normalizer theorem gives
$R\subset J\subset\mathbf C^*R$.  Write $J=CR$, where $C$
is the finite scalar subgroup of $J$.  Since every root of unity
in $F$ lies in $\mu_6\subset R$,
\[
             J\cap\operatorname{GL}_6(F)=R,
             \qquad J/R\simeq C/(C\cap R).                 \tag{BI.57}
\]

Let $t$ count inertias outside $R$.  There are at least two:
otherwise the cyclic quotient in (BI.57) is either trivial,
in which case $L=F$, or is not generated by a product-one tuple.
For an outside $g=\lambda w$, $w\in R$, one has
$\lambda\notin F$.  If its fixed dimension is $b$, the
minimal polynomial of $\lambda^{-1}$ over $F$ occurs at least
$b$ times in the degree-six characteristic polynomial of $w$.
In particular $b\le3$, so each outside inertia moves at least
three dimensions.  Again $t\le3$.

If $t=3$, each fixes three dimensions, and
$[F(\lambda_j):F]=2$.  The roots of unity giving degree two
over $F$ have order four or twelve, and all lie in
$\mathbf Q(\zeta_{12})\setminus F$.  Each therefore has the
same nonidentity class in $\mu_{12}/\mu_6$.  Three such
classes cannot have product one.  Hence $t=2$.

The two phases are inverse modulo $\mu_6$, all other entries
are in $R$, and generation gives
\[
                   J=\langle R,\lambda I\rangle,
                   \qquad L=F(\lambda).                    \tag{BI.58}
\]
The field equality follows from scalar traces and $F\subset L$.
The two moving ranks sum to at most eight, so one fixes at least
two dimensions and $[F(\lambda):F]\le3$.  The complete list
of scalar orders outside $F$ with this property is
\[
                        4,12,9,18.                         \tag{BI.59}
\]
Indeed $\varphi(d)\le6$, so one first lists
$1,2,3,4,5,6,7,8,9,10,12,14,18$.  Orders $1,2,3,6$ lie in $F$;
orders $5,8,10$ have degree four and orders $7,14$ degree six
over $F$, since their cyclotomic fields have trivial intersection
with $F$.  The remaining four
orders give $F(i)=\mathbf Q(\zeta_{12})$ or
$\mathbf Q(\zeta_9)$.  The latter field has the unique
quadratic subfield $F$, so it is excluded.  The only
non-Eisenstein imaginary quadratic subfield of $F(i)$ is
$\mathbf Q(i)$.  Since $\mu_6\subset R$, both orders four
and twelve generate the same group $\mu_4R$.  This proves
(BI.53).

Normalize every outside entry as $g=iw$ with $w\in R$.
The nontrivial automorphism of $F(i)/F$ pairs the eigenvalues
$i$ and $-i$ of $w$ with equal multiplicity.  Thus $g$ has
the eigenvalues $1$ and $-1$ with equal multiplicity $b$,
and its moving rank is $6-b$.  The two outside ranks are at
least three and have sum at most eight; together with (BI.56),
this gives exactly the table (BI.54) by partitioning the
remaining inside moving rank.  In particular $r\le14$.

For outside rank three, $b=3$ exhausts all six eigenvalues,
so $g$ is the stated involution.  For outside rank four,
$b=2$, and the characteristic polynomial of $w$ has the form
$(X^2+1)^2Q(X)$ with $Q$ monic of degree two over $F$.
Every root $\mu$ of $Q$ is a root of unity of degree at most
two over $F$.  It is not $i$ or $-i$, which would increase
$b$.  If it were primitive of order twelve, its conjugate
over $F$ would be $-\mu$, and
$\det w=-\mu^2\notin\{1,-1\}$.  This is impossible,
since $R$ is generated by order-two reflections.  The remaining
possibility is that both roots $\mu,\nu$ of $Q$ belong to
$\mu_6$, with $\mu\nu\in\{1,-1\}$.  Hence
\[
                  \alpha=i\mu,\qquad
                  \beta=i\varepsilon\mu^{-1},
                  \qquad\varepsilon\in\{1,-1\}.
\]
Listing the six choices of $\mu$ and two choices of
$\varepsilon$ gives precisely (BI.55).  The order and square
claims follow immediately from these exponents.  No unlisted
reflection power is used to exclude a possible rank-five entry.
\end{proof}

\paragraph{Signature and descent constraints.}
In (BI.54) the total Chevalley--Weil age must equal nine at every
embedding of $L$ extending either embedding of $K$.
The rank-three outside age is $3/2$; the rank-four outside
ages in (BI.55) are $3/2,2,5/2$.  Thus the $r=14$ row is
automatically balanced, the all-reflection-inside $(3,4)$ row
requires rank-four age two, the $(3,5)$ row requires outside
rank-five age $5/2$, and the $(4,4)$ row requires the two
outside ages to sum to four.  The other rows must still satisfy
the same total-age equation for their nonreflection inside entries.
There is one useful further restriction on the $r=12$, $(3,4)$
row.  Let $a$ be the age of its unique inside rank-two entry and
$b$ the age of its outside rank-four entry.  At the original
embedding the age equation is $a+b=3$.  The automorphism of $L$
fixing $i$ and conjugating $\omega$ changes inside eigenvalues
to their complex conjugates, so that $a$ becomes $2-a$.
For an outside matrix $g=iw$ it sends $g$ to $-\overline g$.
The seven pairs in (BI.55) show that this preserves the
rank-four age $b$; it also preserves the rank-three age $3/2$.
Balance at this second embedding therefore gives $2-a+b=3$.
Consequently $a=1$ and $b=2$ in this row as well.

The natural $L$-multiplicity abelian variety has dimension twelve.
Rational induction and $[L:K]=2$ require it, along the original
family, to be isogenous to the square of the original $K$-sixfold.
For $r=12$ the original rank-nine family dominates its intermediate
Hurwitz component, imposing a generic $M_2(K)$ endomorphism
condition there.  For $r=13$ or $14$ only a subvariety of
dimension at least nine with that condition is forced; dominance
of the entire intermediate component is not asserted.
Moreover D-BI-15 bounds the period rank by at most six whenever
the inducing subgroup is normal or its specified imprimitivity
system is unique.  All remaining cases must avoid those symmetry
hypotheses.

\begin{theorem}[D-BI-17: exclusion of degree-seven inducing blocks]
\label{D-BI-17}
No source in the stated non-Eisenstein class-number-one scope
has an irreducible inducing representation of dimension seven.
\end{theorem}
\begin{proof}
An imprimitive representation of prime dimension seven is
monomial, and transitivity of induction invokes D-MON.  Thus
the faithful inducing image $J$ is primitive.  The actual
intermediate Euler equation is $D=20-14h$, so the same genus
argument gives $h=0$, $r\ge12$, and $D=20$.
There are at least four reflection inertias.  The reflection
normal closure is irreducible by primitivity.  Cohen's normalizer
theorem excludes imprimitive reflection closures, and the real
$A_7$ case is excluded by D-BI-08.  The full reflection list
therefore leaves only $R=W(E_7)$.

Its normalizer is $\mathbf C^*R$, with rational Weyl matrices
and scalar center $\{1,-1\}$.  A nonreal scalar inertia fixes
at most three dimensions and moves at least four.  If $t$
counts these entries, $20\ge r+3t$ gives $t\le2$.
The nonreal character field and the scalar product-one relation
give $t\ge2$, hence $t=2$.  The two moving ranks sum to at
most ten, so one fixes at least two dimensions.  Its scalar
phase consequently has degree at most $\lfloor7/2\rfloor=3$
over $\mathbf Q$.  A nonreal root of unity of degree at most
three has degree two.  The reciprocal-phase relation, generation,
and scalar traces show that the entire character field $L$ is
this quadratic field.  Since $K\subset L$ and $K$ is
non-Eisenstein, $K=L=\mathbf Q(i)$.

As in the degree-six real case, multiplicity-one Frobenius
reciprocity supplies a genuine $K$-model of the inducing
representation, and polarized Shapiro identifies its actual
Hermitian form with the original one up to positive rational
scale.  The already proved Gaussian $E_7$ source barrier forces
splitness.  This contradiction excludes the last case.
\end{proof}

\section{Character-ratio exclusions for normal irreducible cores}
\label{sec:character-ratio}

We use the actual intermediate cover attached to an inducing block, as in
D-BI-01.  Thus an absolutely irreducible complex representation
$\tau:H\to\operatorname{GL}(V)$ of degree $e\ge2$ induces the original
representation, and its multiplicity local system has complex rank six.
Write $J=\tau(H)$, and discard precisely the branch points at which the
local $J$-monodromy is trivial.  No equality between the character field
of $\tau$ and the original imaginary quadratic field is assumed.

For a finite-order linear transformation $x$ put
\[
 \nu(x)=e-\max_{\lambda\in\mathbf C}\dim\ker(x-\lambda),
 \qquad d(x)=\operatorname{rank}(1-x).
\]
Thus $d(x)\ge\nu(x)$, whereas a nonidentity scalar has $\nu(x)=0$
and $d(x)=e$.  Keeping these two quantities distinct is essential below.

\begin{theorem}[D-BI-18: a character-ratio gate]
\label{D-BI-18}
Suppose $N\triangleleft J$ acts irreducibly on $V$, and for some
$0\le c<1$ its character $\chi$ satisfies
\[
          |\chi(x)|\le c e\qquad(x\in N\text{ nonscalar}).       \tag{BI.60}
\]
Every nonscalar $g\in J$ satisfies
\[
                         \nu(g)\ge\frac{e(1-c)}4.              \tag{BI.61}
\]
Consequently a period map of rank nine for the six-dimensional
multiplicity variation necessarily satisfies
\[
       12\left\lceil\frac{e(1-c)}4\right\rceil\le2e+6,
       \qquad e(1-3c)\le6.                                   \tag{BI.62}
\]
If $J\subset\mathbf C^*N$, the stronger necessary bounds are
\[
       12\left\lceil\frac{e(1-c)}2\right\rceil\le2e+6,
       \qquad e(4-6c)\le6.                                   \tag{BI.63}
\]
These conclusions allow inertia outside $N$, including inertia
inducing outer automorphisms of $N$.

In particular the following exclusions hold.
\begin{enumerate}
\item If $N$ is a finite quasisimple group of Lie type over
$\mathbf F_q$, its simple quotient is not alternating, and $q\ge47$,
then $e\in\{3,4,9\}$.  Thus every inducing degree $e\ge10$ is
excluded.  For the smaller field sizes in the following table, the
displayed upper bounds are necessary:
\[
\begin{array}{c|rrrrrr}
q&29&31&32&37&41&43\\ \hline
e\text{ at most}&87&45&33&21&15&15.
\end{array}                                                   \tag{BI.64}
\]
\item In the same Lie-type situation, if $J\subset\mathbf C^*N$
and $q\ge16$, then $e\le9$.  If $q\ge29$, no rank-nine source
exists at all.
\item If $N=ZE$, where $Z$ consists of scalars and $E$ is an
extraspecial group acting irreducibly, then $e\in\{3,4\}$.
In particular every such inducing degree $e\ge5$ is excluded.
\end{enumerate}
The same inequalities apply to an irreducibly acting central product
of commuting factors when each factor has a bound of the form
(BI.60): one may take the maximum of their character-ratio constants.
\end{theorem}

\begin{proof}
We first justify the moduli bound independently of any condition on
the field of definition.  Put $X=C/H$ and $Y=C/\ker\tau$.  The cover
$Y\to X$ has full deck group $J$.  Frobenius reciprocity identifies
the original multiplicity variation with the $\tau$-multiplicity
variation of $Y$, compatibly with its Hodge filtration.  For a fixed
topological type the possible covers $Y\to X$ are finite over the
moduli of the pointed curve $X$, so the actual period map factors
locally through that finite Hurwitz space.  Let $h=g(X)$, let $r$ be
the number of genuine branch points, and let their inertias be
$g_1,\ldots,g_r$.  Since $\tau$ is nontrivial and irreducible, the
Euler-characteristic formula is
\[
       6=2e(h-1)+D,\qquad D=\sum_{j=1}^r d(g_j),\qquad r\le D.
                                                                  \tag{BI.65}
\]
If $h=1$, then $D=6$ and the pointed-curve moduli have dimension at
most six.  If $h=2$, nonnegativity forces $e\le3$ and their dimension
is at most $3+D=9-2e\le5$.  If $h\ge3$ and $e\ge2$, then
$D=6-2e(h-1)<0$, which is impossible.  Therefore rank nine forces
$h=0$.  The dimension bound $r-3\ge9$ now gives
\[
                         r\ge12,\qquad D=2e+6,
                         \qquad e\ge3.                       \tag{BI.66}
\]
This is a necessary bound on a period map that already factors
through an actual intermediate Hurwitz space.  No dominance is
deduced from the dimension of that space.

Conjugate the finite group into a unitary group.  For $x\in N$
nonscalar, choose an eigenvalue $\lambda$ of largest multiplicity.
The trace of $\lambda^{-1}x$ is the sum of $e-\nu(x)$ ones and
$\nu(x)$ complex numbers of modulus one.  The reverse triangle
inequality, followed by (BI.60), gives
\[
                    e-2\nu(x)\le|\chi(x)|\le ce.
\]
Thus $\nu(x)\ge e(1-c)/2$.

Let $g\in J$ be nonscalar.  Since $N$ acts irreducibly, Schur's
lemma supplies $h\in N$ with $[g,h]\ne1$.  Normality gives
$x=[g,h]\in N$.  If $x$ is nonscalar, the preceding paragraph gives
$d(x)\ge\nu(x)\ge e(1-c)/2$.  If $x$ is a nonidentity scalar,
then $d(x)=e$, which is stronger.  For any nonzero scalar $\lambda$,
the commutator is unchanged when $g$ is replaced by
$A=\lambda^{-1}g$.  The elementary inequality
$\operatorname{rank}(1-AB)\le\operatorname{rank}(1-A)
+\operatorname{rank}(1-B)$ therefore gives
\[
 d([g,h])=d([A,h])
 \le d(A)+d(hA^{-1}h^{-1})=2d(A).
\]
Choose $\lambda$ of largest multiplicity in $g$ to get
$d([g,h])\le2\nu(g)$ and hence (BI.61).  This argument does not
require $N$ to be perfect.  A central but nonidentity commutator is
handled using its actual moving rank, not its largest-eigenspace
codimension.

Every genuine inertia has $d(g_j)\ge\lceil e(1-c)/4\rceil$:
this follows from (BI.61) for nonscalars, and from $d(g_j)=e$ for
nonidentity scalars.  Summing over the at least twelve entries in
(BI.66) proves the first inequality in (BI.62).  Dropping ceilings
and rearranging gives its second inequality.  If
$J\subset\mathbf C^*N$, a nonscalar $g=\alpha x$ has
$\nu(g)=\nu(x)\ge e(1-c)/2$ directly.  The identical summation
proves (BI.63).  These elementary support arguments are also the
characteristic-zero form of the arguments in
\cite[Definition~2.3 and Lemma~2.4, pp.~137--138;
Lemma~3.1, p.~139]{BI:DHP}; the proof above includes the scalar
adjustment needed for the invariant $\nu$.

For the Lie-type applications we use the published character bound
\cite[Theorem~4.1(2), p.~141]{BI:DHP}:
\[
 c(q)=\max\left\{\frac1{\sqrt q-1},\frac9q\right\}
                 \qquad(q>9).                              \tag{BI.67}
\]
That theorem attributes this simplified bound to Gluck
\cite{BI:Gluck1995}.  Its hypotheses are finite quasisimplicity,
Lie type over $\mathbf F_q$, and a nonalternating simple quotient.
It bounds every nontrivial irreducible complex character outside
its scalar locus.  Since $N\subset\operatorname{GL}(V)$ acts
faithfully and irreducibly, its scalar locus equals its center, so
it applies to exactly (BI.60).  We use the published formulation
in Duyan--Halasi--Podoski; no new proof of Gluck's character estimate
is claimed.

For $q\ge47$, both terms in (BI.67) are strictly less than $1/5$:
$47>45$ and $47>36$.  The continuous inequality in (BI.62) gives
$e<15$.  Its integer form excludes $5\le e\le8$, since its left
side is at least $24>2e+6$; it excludes $10\le e\le14$, since its
left side is at least $36>2e+6$.  Together with $e\ge3$, this
leaves precisely the necessary possibilities $3,4,9$.

For $q=29,31,32,37,41,43$ one has $c(q)=9/q$.  This comparison can
be checked without approximations by
$(q+9)^2\le81q$.  The continuous inequality gives, respectively,
$e\le87,46,38,22,17,16$.  Applying
$12\lceil e(q-9)/(4q)\rceil\le2e+6$ to these finite ranges
gives (BI.64).  Equivalently, set the ceiling equal to $k$ and
intersect the two rational intervals
$4q(k-1)/(q-9)<e\le4qk/(q-9)$ and $e\ge6k-3$.
The exact finite arithmetic is retained in
\texttt{check\_character\_ratio.py} and its JSON output.

For the scalar-only assertion, $q\ge16$ gives $c(q)\le9/16$,
and (BI.63) implies $e\le48/5$, hence $e\le9$.  If $q\ge29$,
then $9/q<1/3$ and $1/(\sqrt q-1)<1/3$, so (BI.63) gives
$e<3$, contrary to (BI.66).

For the extraspecial assertion, let $x\in E\setminus Z(E)$.
The nondegenerate commutator pairing of the extraspecial group
supplies $h\in E$ with $[h,x]=z\ne1$ in $Z(E)$.
Faithfulness and irreducibility make $z$ a nonidentity scalar.
As $hxh^{-1}=zx$, equality of traces gives
$\chi(x)=z\chi(x)$, hence $\chi(x)=0$.  Multiplication by a scalar
preserves vanishing, so every nonscalar element of $ZE$ has zero
trace.  Thus $c=0$.  Formula (BI.62) first gives $e\le6$, while
its integer form excludes $e=5,6$.  Since $e\ge3$, only $3,4$
remain.  This proves the stated exclusion in every degree at
least five, without a classification of all primitive groups.

Finally suppose a normal irreducible $N$ is generated by commuting
factors $N_1,\ldots,N_s$.  The irreducible representation is a
tensor product of irreducible factor representations $V_i$;
central identifications in a central product do not alter this
statement.  Its normalized character on $x=x_1\cdots x_s$ is
the product of the normalized characters of the $x_i$ on $V_i$.
A tensor product of invertible finite-order transformations is a
scalar only if every factor is a scalar, as is seen by comparing
their eigenvalue ratios.  For a nonscalar $x$, at least one
factor therefore has ratio at most $c_i$, and every other factor
has ratio at most one.  Taking $c=\max_i c_i$ proves (BI.60)
and completes the proof.
\end{proof}

\begin{corollary}[D-BI-18a: the essentially self-dual refinement]
\label{D-BI-18a}
In D-BI-18, suppose additionally that the original character field
is imaginary quadratic and that $J$ is essentially self-dual.
Then
\[
 2\lceil e/2\rceil+
 10\left\lceil\frac{e(1-c)}4\right\rceil\le2e+6,
 \qquad e(3-5c)\le12.                                      \tag{BI.68}
\]
If the original actual polarization is nonsplit, $\tau$ is primitive,
$e\ge9$, and $N$ is quasisimple of Lie type as above, then
$q\ge19$ is impossible.  These additional self-duality hypotheses
hold automatically when the irreducible representation of $N$
admits a nondegenerate invariant bilinear form.
\end{corollary}
\begin{proof}
Write $g^{\mathsf T}Bg=\alpha(g)B$ for the similitude character of
$J$.  If $\alpha$ were trivial, the character of $J$ would be real.
Induction of characters includes the original imaginary quadratic
field in the character field of $\tau$, so this is impossible.
The genuine inertias generate $J$ and have product one, hence at
least two have $\alpha(g_j)\ne1$.  Their fixed spaces are totally
isotropic for $B$, so each moves at least $\lceil e/2\rceil$
dimensions.  At least ten other inertias remain, and D-BI-18
applies to each.  This proves the integer inequality in (BI.68);
discarding ceilings gives its continuous form.

For $q\ge19$, (BI.67) gives $c(q)<1/2$, so (BI.68) implies
$e<24$.  D-BI-09 excludes primitive essentially self-dual inducing
degrees $9\le e\le23$ for a dominant nonsplit sixfold source,
including larger inducing character fields.  Together these two
statements prove the claimed exclusion.

Finally let $B$ be an $N$-invariant nondegenerate bilinear form.
For every $g\in J$, the form $g^{\mathsf T}Bg$ is again
$N$-invariant, because $g$ normalizes $N$.  Irreducibility makes
the space of such forms one-dimensional by Schur's lemma.
Therefore $g^{\mathsf T}Bg=\alpha(g)B$ for a character $\alpha$,
which proves the last assertion.
\end{proof}

\paragraph{Scope.}
The result applies to arbitrary inducing character fields and does
not require the whole inducing group to be quasisimple.  It does
require the stated normal subgroup to act irreducibly.  The
existence of such a subgroup has not been assumed for every
primitive group.  A reducible normal core, factors without an
adequate character-ratio bound, and small defining fields remain
separate cases.  The alternating groups are deliberately excluded
from (BI.67); their character ratios have no uniform upper bound
strictly below one, as also noted in
\cite[Remark~4.2, p.~141]{BI:DHP}.  None of these arguments is a
claim of a universal quarter-dimensional moving-rank theorem for
all finite primitive complex linear groups.

\section{The normal core of a primitive tensor-indecomposable block}
\label{sec:normal-core}

Here tensor indecomposable means that there is no isomorphism
$V\simeq U\otimes W$, with $\dim U,\dim W>1$, under which every
element of the group is a product $a\otimes b$.  A group that
permutes the factors of a tensor decomposition is not excluded by
this hypothesis.  This is the convention preceding
\cite[Theorem~7.1, p.~548]{BI:GS}.

\begin{theorem}[D-BI-19: a forced irreducible quasisimple core]
\label{D-BI-19}
Suppose an inducing block of a rank-nine finite-cover sixfold source
has finite image $J\subset\operatorname{GL}(V)$, is primitive and
tensor indecomposable in the preceding sense, and has degree
$e=\dim V\ge5$.  Write $Z=J\cap\mathbf C^*$, which equals
$Z(J)$ by irreducibility.  Then the following statements hold.
\begin{enumerate}
\item Every nonscalar normal subgroup of $J$ acts irreducibly on
$V$.  Every solvable normal subgroup of $J$ is contained in $Z$.
\item The quotient $J/Z$ has a unique minimal normal subgroup
$S_1\times\cdots\times S_t$, where the $S_i$ are isomorphic
nonabelian simple groups.  Its preimage has the form $ZE$, where
\[
                        E=E_1\cdots E_t                     \tag{BI.69}
\]
is a central product of commuting isomorphic quasisimple groups.
The subgroup $E$ is normal in $J$, acts irreducibly on $V$, and
its factors are permuted transitively by $J$.
\item If the simple groups $S_i$ are of Lie type over
$\mathbf F_q$ and are not alternating, then all inequalities of
D-BI-18 apply with $c(q)$ from (BI.67).  In particular $q\ge47$
is impossible when $e\ge10$, and the exact smaller-field bounds
(BI.64) apply to the total degree $e$.
\item If moreover $e\ge9$, the actual original polarization is
nonsplit, and the representation of $E$ on $V$ is self-dual,
then $q\ge19$ is impossible.  It suffices here that the irreducible
representations on the individual tensor factors are self-dual.
\end{enumerate}
The statement allows larger inducing character fields, outer
automorphisms, and permutations of the quasisimple tensor factors.
\end{theorem}

\begin{proof}
Let $M\triangleleft J$.  Complete reducibility decomposes $V$ into
its $M$-isotypic components.  Normality makes $J$ permute those
components, and irreducibility of $J$ makes the permutation
transitive.  If there were more than one component, they would
give a nontrivial imprimitivity system.  Hence
\[
                  V\simeq W\otimes U,
\]
where $W$ is an irreducible $M$-module and $M$ acts trivially on
the multiplicity factor $U$.  If $M$ is nonscalar, then
$\dim W>1$.  The complex algebra generated by $M$ is
$\operatorname{End}(W)\otimes1_U$, by the matrix-algebra
description of an irreducible summand of the semisimple algebra
$\mathbf C[M]$.

Every element of $J$ normalizes this matrix algebra.  Its induced
automorphism of $\operatorname{End}(W)$ is inner: for example,
the images of a complete set of matrix units furnish a new basis
of the defining module, and change of basis implements the
automorphism.  Thus for each $g\in J$ there is
$a_g\in\operatorname{GL}(W)$ such that
$(a_g^{-1}\otimes1)g$ centralizes
$\operatorname{End}(W)\otimes1$.  The centralizer is
$1\otimes\operatorname{End}(U)$, so
$g=a_g\otimes b_g$ with $b_g\in\operatorname{GL}(U)$.
If $\dim U>1$, this is the forbidden factorwise tensor
decomposition.  Therefore $\dim U=1$, proving that every
nonscalar normal subgroup acts irreducibly.  This argument is
over $\mathbf C$ and uses neither a finite-field specialization
nor maximality of the group.

We next rule out an abelian minimal normal subgroup of $J/Z$.
Let $A/Z$ be such a subgroup and let $A$ be its full preimage in
$J$.  It is nontrivial in $J/Z$, so $A$ is nonscalar and acts
irreducibly by the preceding paragraph.  Since $A/Z$ is abelian,
all commutators in $A$ lie in the scalar group $Z$.  A nonscalar
$x\in A$ is not central in $A$, by Schur's lemma, so some
$h\in A$ has $hxh^{-1}=zx$ with a nonidentity scalar $z$.
Taking traces gives $\operatorname{tr}(x)=z\operatorname{tr}(x)$,
and hence every nonscalar $x\in A$ has trace zero.  D-BI-18 with
$N=A$ and $c=0$ then forces $e\in\{3,4\}$, contradicting
$e\ge5$.  This excludes the entire abelian minimal-normal case,
without needing an extraspecial-group classification.

The quotient $J/Z$ is nontrivial, since $e>1$ and $J$ is
irreducible.  Choose any minimal nontrivial normal subgroup
$\overline N$ of $J/Z$.  A minimal normal subgroup of a finite
group is characteristically simple and is therefore a direct
product of isomorphic simple groups.  For completeness, choose a
minimal nontrivial normal subgroup $T$ of $\overline N$ and take
all its distinct images under automorphisms of $\overline N$.
They generate $\overline N$ by characteristic simplicity.
Distinct such subgroups commute, since their intersection is
trivial by minimality.  Every normal subgroup of $T$ is then
normal in the generated group, so $T$ itself is simple.  If $T$
is abelian it has prime order and the full product is elementary
abelian; otherwise its center is trivial and the product is a
direct product of nonabelian simple groups.  The abelian case has
just been excluded.  We therefore have
\[
                \overline N=S_1\times\cdots\times S_t
\]
with the $S_i$ isomorphic and nonabelian simple.  The action of
$J/Z$ on these factors is transitive: the product over any orbit
is a normal subgroup inside $\overline N$, so minimality leaves
only one orbit.

Let $N$ be the preimage of $\overline N$, and let $A_i$ be the
preimage in $N$ of $S_i$.  As $A_i/Z$ is perfect,
$A_i=A_i'Z$.  Consequently
\[
           A_i'=[A_i'Z,A_i'Z]=[A_i',A_i'],
\]
so $E_i:=A_i'$ is perfect.  Its map onto $S_i$ has central
kernel $E_i\cap Z$, and its center is exactly that kernel,
because $S_i$ has trivial center.  Thus $E_i$ is quasisimple.
For $i\ne j$, the commutators $[A_i,A_j]$ are scalar.  With
$y\in A_j$ fixed, $x\mapsto[x,y]$ is then a homomorphism from
$A_i$ to the abelian group $Z$, so it vanishes on $E_i$.
It follows that $[E_i,A_j]=1$, and hence that $E_i$ and $E_j$
commute.  Since each $A_i=E_iZ$ and the $A_i/Z$ generate
$N/Z$, we obtain
\[
                        N=Z E_1\cdots E_t,
                        \qquad N'=E_1\cdots E_t=E.
\]
The intersections of distinct factors are central, proving
(BI.69).  Conjugation by $J$ permutes the $A_i$ and therefore
the $E_i$ transitively.  In particular the factors are isomorphic.
The subgroup $E=N'$ is normal in $J$, and it acts irreducibly:
$N$ is nonscalar normal and irreducible, and adding the scalar
group $Z$ does not change invariant subspaces.

To see uniqueness of $\overline N$, suppose
$\overline M$ were a different minimal normal subgroup of
$J/Z$.  Distinct minimal normal subgroups commute.  If $M$ is
its preimage, this gives $[M,E]\subset Z$.  For any fixed
$m\in M$, the map $x\mapsto[m,x]$ from the perfect group $E$
to $Z$ is a homomorphism, hence trivial.  Thus $M$ centralizes
the irreducible group $E$, so Schur's lemma implies $M\subset Z$,
a contradiction.  Similarly, a nontrivial solvable normal subgroup
of $J/Z$ would contain a minimal normal subgroup of $J/Z$ that
is elementary abelian.  That case was excluded above.  As $Z$
is solvable and central, every solvable normal subgroup of $J$
is contained in $Z$.

The representation of the commuting central product $E$ is a
tensor product of irreducible representations of the $E_i$.
Each factor representation is faithful: an element acting as the
identity there would act as the identity on $V$, and $J$ is an
actual matrix group.  If the $S_i$ are of Lie type over
$\mathbf F_q$, the published bound (BI.67) applies to each
quasisimple factor.  The central-product assertion of D-BI-18
therefore bounds the normalized character of the irreducible
normal group $E$ on $V$ by the same $c(q)$.  Its degree in
(BI.62) is the total degree $e$, not a single tensor-factor
degree.  The quoted consequences now follow from D-BI-18.

If the representation of $E$ is self-dual, its invariant
nondegenerate bilinear form is unique up to scalar, so its
normalizer $J$ is essentially self-dual as proved in D-BI-18a.
If all the tensor factors are self-dual, tensoring their
nondegenerate forms supplies such a form on $V$.  The final
assertion follows from D-BI-18a, which allows arbitrary inducing
character fields and invokes the separately proved D-BI-09.
\end{proof}

\paragraph{Relation to the published structure reduction.}
The homogeneous-restriction, tensor-normalizer, and central-product
steps have the same form as the proof of
\cite[Theorem~5.1, pp.~148--151]{BI:DHP}.  That theorem itself
assumes a maximal coprime group over a finite field.  We have not
applied that theorem outside its hypotheses: the complete
characteristic-zero argument needed here is written above.
The result imposes no field-size bound on alternating or sporadic
simple factors and leaves the indicated small Lie-type fields
open.  A primitive group preserving a nontrivial factorwise
tensor decomposition is also outside this theorem.

\section{Removing the Lie fields from twenty-nine through forty-three}
\label{sec:lie-small-fields}

\begin{theorem}[D-BI-20: the remaining Lie-field bound]
\label{D-BI-20}
Consider a finite-cover Weil sixfold source of exact imaginary
quadratic character field, with nonsplit actual polarization and a
period map of rank nine.  Let $\tau$ be an absolutely irreducible
inducing block of degree $e\ge10$, and suppose that its finite image
$J$ is primitive and tensor indecomposable in the factorwise sense
specified in D-BI-19.  Suppose the simple factors of the normal core
in D-BI-19 are nonalternating groups of Lie type over $\mathbf F_q$.
Then
\[
                              q\le27.                       \tag{BI.72}
\]
The inducing character field may be larger than the original
imaginary quadratic field.  Tensor-factor permutations and all
outer automorphisms compatible with the representation are allowed.
The conclusion concerns a hypothetical dominant source; it does not
construct sources for the remaining small fields.
\end{theorem}

\begin{proof}
By D-BI-19 there is a normal irreducibly acting central product
$E=E_1\cdots E_t$, whose quasisimple factors are permuted
transitively by $J$.  Write $S=E_i/Z(E_i)$.  Restriction to the
commuting factors expresses the representation of $E$ as a tensor
product of irreducible representations of the $E_i$.  Their degrees
are equal, say $a$, because conjugation by $J$ permutes the factors
transitively.  Thus
\[
                              e=a^t.                        \tag{BI.73}
\]
Each factor representation is faithful: a factor element acting
trivially there acts trivially on the full tensor product, and $E_i$
is an actual subgroup of $\operatorname{GL}(V)$.  In particular its
degree is at least the minimum degree of a nontrivial projective
complex representation of $S$.

The case $q\ge47$ is already excluded by D-BI-18 and D-BI-19,
because $e\ge10$.  The prime powers between $28$ and $46$ are
precisely
\[
                    \mathcal Q=\{29,31,32,37,41,43\}.
\]
For these six values, the exact integer inequality in D-BI-18 gives
the following bounds on the \emph{total} degree:
\[
\begin{array}{c|rrrrrr}
q&29&31&32&37&41&43\\ \hline
e\le U(q)&87&45&33&21&15&15.
\end{array}                                                   \tag{BI.74}
\]

We first remove all types other than $A_1$.  The published
cross-characteristic projective lower bounds for the classical
groups, as collected in \cite[Table~1, author version p.~14]{BI:Hoffman},
give, for every classical simple group over one of these fields
other than $\operatorname{PSL}_2(q)$,
\[
                        a\ge(q^2-1)/2\ge420.                \tag{BI.75}
\]
Here ``other than $A_1$'' includes the rank-two symplectic and
unitary cases.  To check the coarse bound directly from that table:
the linear minimum is at least $q^2+q-1$; the odd-characteristic
symplectic minimum is $(q^n-1)/2$, $n\ge2$; its
even-characteristic counterpart is
$(q^n-1)(q^n-q)/(2(q+1))$, $n\ge2$; the unitary minima are
$(q^{2n}-1)/(q+1)$, $n\ge2$, and
$(q^{2n+1}-q)/(q+1)$, $n\ge1$.  Each is at least the right
side of (BI.75).  The three orthogonal rows, with their stated
minimum ranks, are still larger.  All exceptional parameter pairs
in those classical rows have field size at most four, so none
occurs here.  We use only the published classical rows of this
source, not its rows attributed to an unpublished exceptional-group
manuscript.

For exceptional types we instead use the characteristic-zero
minimum-degree theorem and complete tables of
\cite[Theorem~2.1 and \S5]{BI:Luebeck}.  There are no exceptional
Schur multipliers for these six field sizes: the complete list in
\cite[Chapter~5, Table~1, numbered p.~35]{BI:HHM} has none.
Consequently the simply connected character tables used by
L\"ubeck cover every possible quasisimple central cover here.
For the types $G_2,{}^3D_4,F_4,{}^2F_4,E_6,{}^2E_6,E_7,E_8$,
the first nontrivial degree in the respective tables is at least
$q(q-1)>87$ at the parameters under consideration.  This coarse
inequality can also be checked using the first-degree expressions
retained in \texttt{check\_lie\_small\_fields.py}.
The only possible Suzuki parameter in $\mathcal Q$ is $q=32$,
where the minimum is
\[
                         (32-1)\sqrt{32/2}=124.
\]
The small Ree groups require $q$ to be an odd power of three, and
so do not occur in $\mathcal Q$.  Thus every exceptional type
also contradicts (BI.74).  These are minimum-degree inputs from
published classification tables, not results of a new enumeration
of all finite groups.

We are left with $S=\operatorname{PSL}_2(q)$.  Its projective
minimum-degree bound, from the first row of the same classical
table, is
\[
 a\ge a_0(q):=\frac{q-1}{\gcd(2,q-1)}.
\]
The relevant numbers are
\[
\begin{array}{c|rrrrrr}
q&29&31&32&37&41&43\\ \hline
a_0(q)&14&15&31&18&20&21\\
a_0(q)^2&196&225&961&324&400&441.
\end{array}                                                   \tag{BI.76}
\]
Every entry in the second row exceeds its corresponding $U(q)$.
Equations (BI.73)--(BI.76) therefore force $t=1$.  In particular
$E$ itself acts irreducibly, rather than merely one tensor factor
of a higher-degree representation.  The values $q=41,43$ are
already impossible because $a_0(q)>U(q)$.

For precision about the remaining covers, the usual universal
central cover of $\operatorname{PSL}_2(q)$ is
$\operatorname{SL}_2(q)$, except for the exceptional multipliers
at $q=4,9$.  The exceptional list and the distinction between
ordinary Lie-type covers and exceptional covers are explicitly
recorded in \cite[Introduction, numbered p.~4; Chapter~5,
Table~1, numbered p.~35]{BI:HHM}.  Hence at the parameters here
the group $E$ is $\operatorname{SL}_2(q)$ or its quotient
$\operatorname{PSL}_2(q)$; for even $q$ these coincide.  We must
therefore check inverse conjugacy and automorphisms on
$\operatorname{SL}_2(q)$, not only on its projective quotient.

At $q=37$, all irreducible degrees of $\operatorname{SL}_2(37)$
belong to
\[
                       \{1,18,19,36,37,38\}.               \tag{BI.77}
\]
This is the specialization $\alpha=q-1$, $\beta=1$, $f=1$ of
\cite[Theorem~B, p.~4]{BI:WhiteSL2}; degrees of the projective
quotient are a subset.  Directly imposing
$12\lceil e(37-9)/(4\cdot37)\rceil\le2e+6$ with
$3\le e\le21$ leaves only
\[
                         \{3,4,5,9,10,15,21\}.
\]
This set is disjoint from (BI.77), so $q=37$ is excluded by the
exact ceiling inequality, without a self-duality shortcut.

We next prove that every element of $\operatorname{SL}_2(q)$
is conjugate to its inverse when $q$ is even or $q\equiv1\pmod4$.
A noncentral semisimple matrix $x$ and $x^{-1}$ have the same
characteristic polynomial and are conjugate in
$\operatorname{GL}_2(q)$.  The determinant map on the
$\operatorname{GL}_2(q)$-centralizer of $x$ is surjective onto
$\mathbf F_q^*$: for a split torus this is immediate from the
diagonal matrices, and for a nonsplit torus it is the surjective
norm $\mathbf F_{q^2}^*\to\mathbf F_q^*$.  Multiplying a
conjugator by a suitable centralizer element therefore makes its
determinant one.  Thus the semisimple inverse conjugacy takes
place in $\operatorname{SL}_2(q)$.

For odd $q$, every remaining noncentral element is conjugate to
$\pm(I+bE_{12})$ for some $b\ne0$.  If $q\equiv1\pmod4$,
choose $u\in\mathbf F_q$ with $u^2=-1$.  Conjugation by
$\operatorname{diag}(u,u^{-1})$ replaces $b$ by $-b$, producing
the inverse.  For even $q$, a nontrivial unipotent is an
involution.  Central elements are already their own inverses.
This proves the assertion on the full special linear group, and
it descends to its central quotients.  Its class description agrees
with \cite[\S3, pp.~5--7]{BI:WhiteSL2}.

It follows that every irreducible character of $E$ is real when
$q=29$ or $32$.  Indeed $\chi(x^{-1})=\overline{\chi(x)}$
for a finite-group character, and inverse-conjugate elements have
the same character value.  A real irreducible character is
self-dual and has a nondegenerate invariant complex bilinear form,
possibly alternating.  D-BI-18a applies to the entire normalizer
$J$, since the $E$-invariant form is unique up to scalar.  Its
hypotheses hold here: $q\ge19$, $e\ge10$, the inducing block is
primitive, and the original actual polarization is nonsplit.
Therefore $q=29,32$ are excluded.

Finally let $q=31$.  The outer automorphism group of
$\operatorname{PSL}_2(31)$ has order two, generated by the
diagonal automorphism; there is no field automorphism because
$31$ is prime.  This is the specialization of
\cite[\S2, p.~3; \S3, p.~6]{BI:WhiteSL2}.  The same assertion
holds for $\operatorname{SL}_2(31)$.  To justify the latter
point, an automorphism inducing the identity on the central
quotient has the form $x\mapsto z(x)x$, with $z$ a homomorphism
to the center.  The group $\operatorname{SL}_2(31)$ is perfect,
so $z=1$.  Perfectness here can be seen directly: commutators of
$\operatorname{diag}(u,u^{-1})$ with upper and lower unipotents,
for $u^2\ne1$, produce all those unipotents, and those matrices
generate $\operatorname{SL}_2(31)$.  Thus its automorphisms
inject into those of its projective quotient, and the diagonal
automorphism supplies the nontrivial outer class.

This diagonal automorphism is class-inverting on
$\operatorname{SL}_2(31)$.  For semisimple elements it preserves
the characteristic polynomial, hence the special-linear conjugacy
class by the centralizer argument above; that class also contains
the inverse.  It fixes the central elements.  On each of the two
unipotent square classes, conjugation by
$\operatorname{diag}(v,1)$ for a nonsquare $v$ exchanges the
square classes.  Inversion has the same effect because $-1$ is a
nonsquare in $\mathbf F_{31}$.  This also treats the negatives
of the unipotents.  It proves class inversion on the cover and on
its central quotient; compare the explicit class action in
\cite[Lemma~3.1, p.~7]{BI:WhiteSL2}.

If the irreducible character $\chi$ of $E$ is real, D-BI-18a
again gives a contradiction.  If it is nonreal, the nontrivial
diagonal outer automorphism sends $\chi$ to
$\overline\chi\ne\chi$.  It cannot be induced by conjugation
by an element of $J$, since such conjugation preserves the
character of the restriction of the fixed representation on $V$.
Hence every element of $J$ induces an inner automorphism of $E$.
For $g\in J$, choose $x\in E$ inducing the same automorphism.
Then $gx^{-1}$ centralizes the irreducible group $E$, and is
scalar by Schur's lemma.  Consequently
\[
                             J\subset\mathbf C^*E.
\]
The scalar-only part of D-BI-18 excludes this case because
$q=31\ge29$.  This finishes all six fields in $\mathcal Q$ and
proves (BI.72).
\end{proof}

\paragraph{Certificates and scope.}
The file \texttt{check\_lie\_small\_fields.py} verifies the exact
ceiling sets, the total-degree comparisons, the stated exceptional
minimum-degree expressions, and the $q=37$ disjointness.  It does
not reprove the published degree tables, classify all simple groups,
or test period dominance by a dimension inequality.  The
inverse-conjugacy and automorphism arguments above are universal
two-by-two proofs, not conclusions drawn from finite samples.
The remaining Lie fields are $q\le27$; alternating cores and
other simple core types are outside this theorem's scope.

\section{Exclusion of all rank-one Lie cores}
\label{sec:a1-core}

The small-field part of this section uses complete published ATLAS
character tables.  The finite spectral statements have been checked
both by GAP and by a separate exact Ramanujan-trace implementation.
The outer-automorphism coverage requires a further argument at
$q=9,25$: a universal-cover extension missing from the table list
is not silently assumed to exist.

\subsection*{D-BI-21a. The complete small-field spectral input}

Let
\[
 \mathcal Q_0=\{4,5,7,8,9,11,13,16,17,19,23,25,27\}.
\]
The duplicate simple groups at $q=4,5$ are represented once.
For each simple group $\operatorname{PSL}_2(q)$, the following
table specifies its universal cover and all cyclic outer subgroups.
The notation after a semicolon indicates that a proper central
quotient is used for that outer subgroup, subject to D-BI-21b below.
All identifiers are exact CTblLib names.
\[
\begin{array}{c|l|l}
q&U&\text{cyclic outer extension tables}\\\hline
4,5&2.A5&2.A5.2\\
7&2.L3(2)&2.L3(2).2\\
8&L2(8)&L2(8).3\\
9&6.A6&6.A6.2\_1,\ 6.A6.2\_2;\quad3.A6.2\_3\\
11&2.L2(11)&2.L2(11).2\\
13&2.L2(13)&2.L2(13).2\\
16&L2(16)&L2(16).2,\ L2(16).4\\
17&2.L2(17)&2.L2(17).2\\
19&2.L2(19)&2.L2(19).2\\
23&2.L2(23)&2.L2(23).2\\
25&2.L2(25)&2.L2(25).2\_1,\ 2.L2(25).2\_2;
                    \quad L2(25).2\_3\\
27&2.L2(27)&2.L2(27).2,\ 2.L2(27).3,\ 2.L2(27).6
\end{array}                                                   \tag{BI.78}
\]
The two additional quotient-core tables $3.A6$ and $L2(25)$
are retained too.  Thus the full finite input has fourteen core
tables and nineteen extension tables, thirty-three in total.
Each extension is supplied with its actual stored class fusion
from the core named in the certificate.

\begin{lemma}[D-BI-21a: the finite spectral bounds]
\label{D-BI-21a}
In every nonlinear irreducible representation of a core table in
(BI.78), and in every irreducible representation of an extension
table whose restriction to its specified core is nonlinear and
irreducible, every nonscalar element satisfies
\[
                              \nu(g)\ge a/5,                 \tag{BI.79}
\]
where $a$ is the representation degree.  The stronger bound
$\nu(g)\ge a/4$ fails only for the following four entries:
\[
\begin{array}{c|r|r|r|r|r}
\text{table}&\text{row}&\text{class}&a&\operatorname{ord}(g)&\nu(g)\\\hline
6.A6.2\_1&3&20&5&2&1\\
6.A6.2\_1&4&20&5&2&1\\
6.A6.2\_1&5&21&5&4&1\\
6.A6.2\_1&6&21&5&4&1
\end{array}                                                   \tag{BI.80}
\]
The restrictions of these rows are, respectively, rows
$2,2,3,3$ of $6.A6$.  Those two core characters have degree five
and are trivial on the center.
\end{lemma}

\begin{proof}
The identifications, characters and power maps are explicit
published ATLAS inputs accessed through \cite{SPOR:CTblLib}.
Every retained table's source string identifies that origin.
The universal-cover and automorphism metadata are interpreted as
in the library's \S3.7-3, with the small multiplier and isomorphism
exceptions also listed in \cite[Chapter~5, Tables~1--2,
numbered pp.~35--36]{BI:HHM}.  No group identification is inferred
from character orthogonality alone.

For every selected character row $\chi$ and every class of an
element $g$ of order $m$, compute every eigenvalue multiplicity by
\[
 m_k=\frac1m\sum_{j=0}^{m-1}\chi(g^j)\zeta_m^{-kj}.
                                                               \tag{BI.81}
\]
All arithmetic is exact in a cyclotomic field.  The criterion
$|\chi(g)|^2=\chi(1)^2$ discards exactly the scalar classes.
An extension row is selected precisely when its restriction along
the actual core fusion has character norm one and degree greater
than one.  Hence no representation with an irreducible core
restriction is dropped by a degree cutoff or a reality condition.

The executed file \texttt{a1/export\_and\_scan.g} checks all
$654$ selected rows and $18\,922$ nonscalar entries.  It verifies
the integrality, nonnegativity and degree sum of every spectrum,
and produces exactly (BI.80) and no failure of (BI.79).  The full
input also includes every unselected character row, all class
sizes, group orders, power maps, source strings and subgroup
fusions.  Its SHA-256 is
\texttt{5ed905449e1f6f64f3eb18b27a213fc5c72e0731dbd591670166e8763ad7fc88}.

The separate program \texttt{check\_a1\_spectra.py} imports no GAP
calculation or selected-row list.  It selects the rows using exact
restriction norms and reconstructs all $18\,922$ spectra using
Ramanujan sums.  Explicitly, the normalized field trace of
$\zeta_m^j$ is $c_m(j)/\varphi(m)$, where
\[
                 c_m(j)=\sum_{d\mid\gcd(m,j)}d\mu(m/d).
\]
Since (BI.81) is a rational integer, its normalized field trace
equals itself.  This computes it independently by integer sums,
after the prime power maps and cyclotomic Galois action have
reconstructed the character values on all powers.  That program
finds the same four exceptions and verifies the stated core
restrictions and central kernels.  It also checks the complete
table-name inventory and all nineteen core fusions: orders,
identity kernels, full normal-class unions and compatibility
with every retained prime power map.  These complete finite checks
prove the lemma from the published input.  They do not reprove
the character tables or the classification of the named groups.
\end{proof}

\subsection*{D-BI-21b. The missing mixed outer extensions}

\begin{lemma}[D-BI-21b: a central obstruction]
\label{D-BI-21b}
Let $q=p^2$ with $p$ odd, and let $\alpha$ be the product of a
nontrivial diagonal outer automorphism and the field automorphism
of $\operatorname{SL}_2(q)$.  No irreducible complex representation
on which $-I$ acts as $-1$ is invariant under $\alpha$.
At $q=9$, the same assertion holds on the universal cover $6.A6$
for every representation nontrivial on its central involution.
\end{lemma}

\begin{proof}
Choose a generator $\zeta$ of $\mathbf F_q^*$ and take
\[
 \alpha(x)=D x^{(p)}D^{-1},\quad
 D=\operatorname{diag}(\zeta,1),\quad
 u=\operatorname{diag}(\zeta^{(p+1)/2},\zeta^{-(p+1)/2}).
\]
Direct matrix multiplication gives
\[
                     \alpha^2=\operatorname{inn}(u),
                     \qquad \alpha(u)=-u.                  \tag{BI.82}
\]
For the second equality, each diagonal exponent changes by
$\pm(p-1)(p+1)/2=\pm(q-1)/2$, whose power of $\zeta$ is $-1$.
These are the explicit diagonal and field automorphisms described
in \cite[\S3, pp.~5--7]{BI:WhiteSL2}.

Suppose an invertible intertwiner $T$ implements $\alpha$ in an
irreducible representation $\rho$.  The first equality in
(BI.82) and Schur's lemma imply $T^2=c\rho(u)$ for a nonzero
scalar $c$.  Since $T$ commutes with its own square, it must
commute with $\rho(u)$.  The second equality instead makes it
send $\rho(u)$ to $-\rho(u)$ when $\rho(-I)=-1$.
This contradiction proves the assertion for $\operatorname{SL}_2(q)$.

For $q=9$ let $U=6.A6$, let $U\to2.A6=\operatorname{SL}_2(9)$
be the canonical central quotient, and lift $\alpha$ uniquely
to $U$.  The uniqueness of this lift follows from the universal
property of the perfect central cover; two lifts with the same
action on $A_6$ differ by a homomorphism from the perfect group
$U$ to its center, and therefore coincide.  Choose any lift
$u'\in U$ of $u$.  Uniqueness gives
$\alpha^2=\operatorname{inn}(u')$ on $U$.
The element $z=\alpha(u')u'^{-1}$ is central, and its image in
$2.A6$ is $-I$.  Since $Z(U)$ is cyclic of order six, $z^3$ is
its central involution.  If that involution acts as $-1$, then
$\rho(z)\ne1$.  But an intertwiner would again satisfy
$T^2=c\rho(u')$ and commute with $\rho(u')$, while conjugating
it to $\rho(z)\rho(u')$.  This is impossible.
\end{proof}

\subsection*{D-BI-21c. Coverage of every linear normalizer}

\begin{lemma}[D-BI-21c: small-field normalizers]
\label{D-BI-21c}
Let $N$ be a finite quasisimple group with simple quotient
$\operatorname{PSL}_2(q)$ for $q\in\mathcal Q_0$, acting faithfully
and irreducibly on a space of dimension $a$.  Every nonscalar
element of its linear normalizer, with arbitrary scalar phase,
satisfies $\nu\ge a/5$.  It satisfies $\nu\ge a/4$ unless
the representation factors through $A_6$ and has degree five.
\end{lemma}

\begin{proof}
Lift the representation to the universal cover $U$.  The
normalizer modulo scalars embeds in $\operatorname{Aut}(S)$,
where $S=N/Z(N)$: a central automorphism of a perfect group is
trivial, and the linear centralizer is scalar.  The complete
outer group is generated by diagonal and field automorphisms,
with structure $C_{\gcd(2,q-1)}\times C_f$ for $q=p^f$;
see \cite[\S2, pp.~2--3; \S3, p.~6]{BI:WhiteSL2}.
For $A_5$ we use the equivalent $q=5$ model.  Thus (BI.78)
includes every cyclic subgroup generated by any possible outer
element, not only the largest outer group.  In particular both
orders two and four occur at $q=16$, all three subgroups of order
two occur at $q=9,25$, and orders two, three and six occur at $q=27$.

We spell out why a fixed cyclic extension table controls any
actual normalizing matrix.  Let $A$ be such a table group,
$A/U$ cyclic of order $l$, and choose $t\in A$ whose coset
generates it.  Write $\beta=\operatorname{inn}(t)|_U$ and
$u_0=t^l\in U$.  A normalizing matrix $g$, after an inner
adjustment and replacing $t$ by a suitable generating power,
induces the same automorphism of $S$ as $t$.  Unique lifting to
the universal perfect cover then gives
\[
                  g\rho(u)g^{-1}=\rho(\beta(u)).
\]
In particular the represented kernel is preserved.  Schur's
lemma makes $g^l\rho(u_0)^{-1}$ a nonzero scalar.  Rescale $g$
by an $l$th root to arrange $g^l=\rho(u_0)$.  The conjugation
and power relations now define an irreducible representation
of the actual fixed group $A$, extending $\rho$.  The original
matrix is, up to scalar and the inner adjustment, an element
of that represented extension.  Its spectrum is covered by
D-BI-21a.  This reasoning does not assume that $A/U$ splits,
nor that every core character is invariant under every outer
automorphism.

There are two exceptions to using the universal-cover extension
literally.  At $q=25$, the ATLAS subscripts $2_1,2_2,2_3$ denote
the diagonal, field, and mixed diagonal-field classes, respectively.
At $q=9$ the order of the first two is reversed: $2_1$ is the
field extension $S_6$, $2_2$ is the diagonal extension
$\operatorname{PGL}_2(9)$, and $2_3$ is the remaining mixed
extension.  The construction labels and exact subgroup fusions
are retained in the source inventory.  The missing usual cover
extensions are precisely these mixed cases.

D-BI-21b shows that an invariant core representation in either
mixed case must kill its central involution.  Hence it factors
through $L2(25)$ at $q=25$, and through $3.A6$ at $q=9$.
The actual extensions $L2(25).2_3$ and $3.A6.2_3$ and their
core fusions are present in (BI.78), so the preceding cyclic
argument applies to them.  No inference from an unavailable
universal-cover extension is made.  The resulting representations
are exactly among the rows tested in D-BI-21a.

Finally scalar multiplication does not change $\nu$, and a
normalizer element with inner outer action is scalar times an
element of $N$.  This proves the universal fifth bound and the
stated sole possible quarter-bound exception.
\end{proof}

\subsection*{D-BI-21. The rank-one core is impossible}

\begin{theorem}[D-BI-21: full $A_1$-core exclusion]
\label{D-BI-21}
Let a finite-cover Weil sixfold source of exact imaginary
quadratic character field have nonsplit actual polarization and
period rank nine.  Let an inducing block have degree $e\ge5$
and finite image $J$ that is primitive and tensor indecomposable
in the factorwise convention of D-BI-19.  Then the simple factors
of its forced normal core cannot be $\operatorname{PSL}_2(q)$
for any prime power $q$ for which this group is simple.

The statement includes the exceptional identifications $A_5$ and
$A_6$, the exceptional universal cover $6.A6$, arbitrary larger
inducing character fields, compatible outer automorphisms and
permutations of tensor factors.  For the small fields
$q\in\mathcal Q_0$, total degree $e\ge6$ is excluded even
without the nonsplit-polarization hypothesis.
\end{theorem}

\begin{proof}
D-BI-19 supplies a normal irreducibly acting central product of
quasisimple factors, permuted transitively by $J$.  Its factor
representations are faithful and irreducible, and their degrees
are equal, say $a$, so $e=a^t$.

If $q\ge29$, the published projective minimum-degree bound
$a\ge(q-1)/\gcd(2,q-1)$ in
\cite[Table~1, author p.~14]{BI:Hoffman} gives $e\ge14$.
The simple quotient is then nonalternating, so D-BI-20 applies
and excludes this case.  The remaining simple $A_1$ groups have
exactly the parameters in $\mathcal Q_0$.

For these parameters, every factorwise nonscalar normalizer
element satisfies $\nu\ge e/5$, by D-BI-21c and the tensor
eigenspace bound.  Every element permuting factors nontrivially
satisfies the stronger $\nu\ge e/4$, by the weighted cyclic-word
proof of D-SPOR-03; that proof does not depend on the factor
being sporadic.  Thus every nonscalar element of $J$ satisfies
the fifth bound.  Genuine scalar inertias have full moving rank.

The actual intermediate Hurwitz factorization gives genus zero,
at least twelve genuine branch points and total moving rank
$2e+6$, as proved in D-BI-18.  Consequently
\[
          12\lceil e/5\rceil\le2e+6,
          \qquad e\le15.                                  \tag{BI.83}
\]
This is a necessary constraint on the actual period map.
If no factor has the exceptional degree-five $A_6$ representation,
the identical argument uses the quarter bound and gives
\[
                         12\lceil e/4\rceil\le2e+6.
                                                               \tag{BI.84}
\]
For $e=5,6$ the left side is $24$; for $e\ge7$ it is at
least $3e>2e+6$.  Thus (BI.84) is impossible for every $e\ge5$.

If there is an exceptional factor, transitivity makes the common
factor degree $a=5$.  With at least two factors one has
$e=5^t\ge25$, contradicting (BI.83).  The only remaining
possibility is a single factor of degree five with simple quotient
$A_6$.  This proves already the assertion for all $e\ge6$,
without using a determinant condition.

We finish the degree-five case.  The two core rows in (BI.80)
are exactly the two degree-five ordinary characters of $A_6$.
One is the deleted natural permutation representation on six
letters: the natural action is doubly transitive, so its deleted
character is irreducible.  The two characters are interchanged
by the diagonal outer automorphism of
$A_6=\operatorname{PSL}_2(9)$; this is also the half-degree pair
in \cite[Lemma~4.4, p.~9]{BI:WhiteSL2}.  Precomposing by that
automorphism leaves the matrix image of $A_6$ unchanged up to
equivalence.  Thus either core image is the deleted natural
$A_6$ image in dimension five.

Its linear normalizer is contained in $\mathbf C^*S_6$ in that
deleted representation.  Indeed the normalizer modulo scalars
injects into $\operatorname{Aut}(A_6)$ and must stabilize the
chosen degree-five character.  Its outer stabilizer has order two:
the field outer automorphism fixes the half-degree characters,
whereas the diagonal and mixed automorphisms exchange them
\cite[Lemmas~4.4 and~4.6, pp.~9--10]{BI:WhiteSL2}.
The natural $S_6$ supplies that stabilizer.  Schur's lemma then
makes every remaining difference scalar.  Since $J$ contains
the core, it lies in the exact subgroup range
\[
                          A_6\subset J\subset\mathbf C^*S_6.
\]
The previously proved D-BI-08 applies to this deleted-standard
inducing representation, with its larger-field allowance and
actual-polarization conclusion.  It excludes a dominant nonsplit
source.  This disposes of the final case and proves the theorem.
\end{proof}

\paragraph{Scope of the computation.}
The input contains all thirty-three tables in (BI.78), not just
the four exceptional entries.  All representations tested have
irreducible core restriction; incompatible outer automorphisms
are not assumed to act.  The finite computation is only the
spectral input to the complete geometric, field and normalizer
proof.  No signature or branch-parameter equality is used to
infer dominance.  Groups with higher Lie rank and ordinary
alternating groups outside the indicated $A_5,A_6$ identifications
remain separate cases.

\section{Exclusion of Suzuki and Ree normal cores}
\label{sec:twisted-lie-cores}

This section removes the three Suzuki--Ree series, together with
the Tits group.  The argument allows all normalizing outer
automorphisms and tensor permutations.  In particular, it does
not require a character table for the missing outer extension
of ${}^2F_4(8)$.

\subsection*{D-BI-22a. Complete small-field inputs}

We use $q$ for the actual finite-field size.  Put
\[
 U=2^2.\operatorname{Sz}(8),\qquad A=U.3.
\]
The center of $U$ has order four and is not cyclic.  A nonlinear
irreducible representation of $U$ need not be faithful; retaining
all such representations is essential for covering every possible
quasisimple central quotient.

\begin{lemma}[D-BI-22a: exact small-field bounds]
\label{D-BI-22a}
Every nonscalar element in every nonlinear irreducible representation
of $U$, and in every irreducible representation of $A$ whose
restriction to $U$ is nonlinear and irreducible, satisfies
\[
                         \nu(g)\ge\frac37a\ge\frac14a,
                                                               \tag{BI.85}
\]
where $a$ is the representation degree.  The smallest such degree
is fourteen.

For every nontrivial irreducible character $\chi$ of the three
groups in the following table, and every nonscalar class,
\[
                  |\chi(g)|\le c\chi(1),                  \tag{BI.86}
\]
with the indicated constants and minimum degrees:
\[
\begin{array}{c|c|r}
S&c&\min\chi(1)\\\hline
{}^2G_2(27)&26/703&703\\
{}^2F_4(2)'&3/13&26\\
{}^2F_4(8)&449/32319&64638
\end{array}                                                   \tag{BI.87}
\]
All three constants are at most $1/4$.
\end{lemma}

\begin{proof}
The exact input consists of the complete CTblLib tables named
\texttt{2\^{}2.Sz(8)}, \texttt{2\^{}2.Sz(8).3},
\texttt{R(27)}, \texttt{2F4(2)'}, and \texttt{2F4(8)}.
Their numbers of classes are respectively $35,25,35,22,113$.
The first four have ATLAS source strings.  The last is the
complete table supplied by Gunter Malle, whose provenance is
explicitly recorded in CTblLib, rather than relabelled as an
ATLAS table.  The complete Ree character table in CHEVIE is
also explicitly documented in \cite[\S10.2, numbered p.~131,
before Proposition~10.4]{BI:HHM}; Malle's published unipotent
character computation is \cite{BI:MalleRee}.
The group identifications and character tables are these
published mathematical inputs, not consequences of numerical
orthogonality tests.

The executed program \texttt{twisted\_lie/export\_and\_scan.g}
retains every character, class size, order, prime power map,
source string and the actual $U\to A$ fusion.  It checks full
row orthogonality.  The extension rows are selected by exact
norm-one restriction, giving $34,12,34,21,112$ selected rows,
respectively.  The computation examines all $213$ selected rows
and $15\,471$ nonscalar entries.

For the first two tables it computes all eigenvalue multiplicities
by the exact Fourier formula (BI.81), checks their nonnegative
integrality and degree sums, and finds minimum ratio $3/7$.
For the remaining tables write each retained cyclotomic value as
\[
                 \chi(g)=\sum_j b_j\zeta_m^j,
                 \qquad B(\chi(g))=\sum_j|b_j|.
\]
The coefficients $b_j$ are integers and all basis choices are
retained in the input.  The triangle inequality gives
$|\chi(g)|\le B(\chi(g))$.  Exact rational maximization of
$B(\chi(g))/\chi(1)$ over all nonscalar entries gives precisely
the three constants in (BI.87).  Thus no floating comparison of
cyclotomic values or unproved cancellation is used.

The independent program \texttt{check\_twisted\_lie.py} selects
the rows again using rational normalized-trace restriction norms,
reconstructs all $1330$ Suzuki spectra by Ramanujan Fourier sums,
and recomputes all $14\,141$ remaining coefficient bounds.
It verifies the complete five-table inventory, all minimum degrees,
and the fusion's orders, identity kernel, normal-class unions and
prime power-map compatibility.  Both executions pass.
The complete input is $161615$ bytes with SHA-256
\texttt{10dd43bccf28471f0c8725086666385aba1dc9970fe9d53bd1478e9ad6a1637b}.
These finite exact checks prove the lemma from the stated inputs.
\end{proof}

\subsection*{D-BI-22b. Every single-factor normalizer}

\begin{lemma}[D-BI-22b: universal normalizer bounds]
\label{D-BI-22b}
Let $N$ be a finite quasisimple group acting faithfully and
irreducibly in degree $a$, whose simple quotient is a Suzuki
or Ree group or the Tits group.  Every nonscalar element of
its linear normalizer satisfies
\[
                         \nu(g)\ge\frac{23}{128}a.        \tag{BI.88}
\]
All such factor degrees are at least fourteen.  More specifically:
for $\operatorname{Sz}(8)$ the quarter bound holds; for the
three groups (BI.87) the bound is $3a/16$; and for actual field
size $q\ge32$ the minimum degree is at least $124$.
\end{lemma}

\begin{proof}
The only exceptional Schur multiplier in these three series is
that of $\operatorname{Sz}(8)$; see
\cite[Chapter~5, Table~1, numbered p.~35]{BI:HHM}.
The generic centers are trivial, and the Tits multiplier is trivial.
The small-group universal-cover and outer metadata are also
retained from the official CTblLib extension metadata, interpreted
as in its \S3.7-3 \cite{SPOR:CTblLib}.

For $\operatorname{Sz}(8)$ the universal cover is $U$ above and
the full outer group is cyclic of order three.  Lift the given
irreducible representation to $U$.  An inner normalizer is scalar
times a represented element of $U$, and is covered by (BI.85).
For an outer normalizer, choose the matching generating coset of
the actual extension $A$.  After an inner adjustment its action
on $U$ is the unique lift of the corresponding simple-group
automorphism.  If $t\in A$ represents that coset and $t^3=u_0$,
Schur's lemma gives $g^3=c\rho(u_0)$.  Rescale $g$ by a cube
root of $c^{-1}$.  The conjugation and power relations then give
an irreducible representation of the fixed group $A$ extending
$\rho$, exactly as in D-BI-21c.  Its restriction belongs to the
tested rows.  This proves the quarter bound without assuming a
faithful irreducible representation of the noncyclic center of $U$,
or assuming that every outer automorphism preserves every central
kernel.

For the three groups (BI.87), D-BI-18's complex commutator
argument turns $|\chi(x)|/a\le1/4$ for nonscalar core elements
into
\[
                              \nu(g)\ge3a/16
\]
for every nonscalar normalizer, including every outer coset.
This avoids any need for an outer-extension table for
${}^2F_4(8)$.  The proof of D-BI-18 also handles nonidentity
scalar commutators by their full actual moving rank.

The remaining actual field sizes are at least $32$.
The published ordinary-character bound
\cite[Theorem~4.1(2), p.~141]{BI:DHP} gives
\[
 c(q)=\max\{1/(\sqrt q-1),9/q\}\le9/32.
\]
Indeed $q\ge32$ implies $\sqrt q>5$, so the first term is
less than $1/4<9/32$, and the second is at most $9/32$.
All simple quotients here are nonalternating, as required by
that statement.  The same complex commutator argument gives
$(1-c(q))a/4\ge23a/128$.

For completeness the first-degree formulas, converted from
the square-root field parameter in
\cite[\S\S5.1,5.2,5.5, numbered pp.~9,14--15]{BI:Luebeck},
are
\[
\begin{array}{c|c}
{}^2B_2(q)&(q-1)\sqrt{q/2}\\
{}^2G_2(q)&q^2-q+1\\
{}^2F_4(q)&(q-1)\sqrt{q/2}(q+1)^2(q^2-q+1).
\end{array}                                                   \tag{BI.89}
\]
For fields in this large range their first admissible values are
$q=32,243,32$, respectively, and the corresponding minima are
$124,58807,134090748$.  Each expression is increasing thereafter.
There are no exceptional covers in this range.  At the remaining
small simple groups the minima are those in D-BI-22a, so all
factor degrees are at least fourteen.

Finally $a/4$ and $3a/16$ are both at least $23a/128$.
This proves the uniform assertion and all refinements.
\end{proof}

\subsection*{D-BI-22. The whole Suzuki--Ree source gate}

\begin{theorem}[D-BI-22: Suzuki--Ree core exclusion]
\label{D-BI-22}
Let an inducing block of a finite-cover Weil sixfold source have
finite image $J\subset\operatorname{GL}(V)$.  Suppose $J$ has
an irreducibly acting normal central product of quasisimple
groups whose simple quotients all belong to
\[
 {}^2B_2(2^{2m+1}),\quad{}^2G_2(3^{2m+1}),\quad
 {}^2F_4(2^{2m+1})\quad(m\ge1),\qquad{}^2F_4(2)'.
\]
Then this source has period rank at most eight.

The statement allows mixed factor types, all compatible outer
automorphisms, arbitrary scalar phases, and tensor permutations.
It requires neither primitive or tensor-indecomposable $J$,
essential self-duality, equality of inducing and original fields,
nor a nonsplit-polarization assumption.
\end{theorem}

\begin{proof}
Write the normal central product as $N=N_1\cdots N_t$.
Its irreducible representation has the tensor form
$V=V_1\otimes\cdots\otimes V_t$, with faithful irreducible
factor representations and degrees $a_i\ge14$.
Put $e=\prod_i a_i$.
The factors $N_i$ are intrinsic: each is the derived subgroup
of the inverse image of its simple direct factor in $N/Z(N)$.
Hence $J$ permutes them.  Their full matrix algebras
$\operatorname{End}(V_i)\otimes1$, supplied by irreducibility,
are permuted as well.  Thus a normalizing operator is a tensor
product of factor intertwiners followed by a permutation of
isomorphic factor spaces.

For a factor-preserving nonscalar operator at least one factor
intertwiner is nonscalar.  Fixing the eigenvalues in all other
factors shows that its largest eigenspace fraction cannot exceed
that in this factor.  D-BI-22b therefore gives
$\nu(g)\ge23e/128$.

If an operator permutes factors nontrivially, choose one cycle
of length $l\ge2$ on factors of dimension $a$.
After tensor changes of basis the action is a cyclic word rotation
with a diagonal twist at one place.  Such a diagonalization is
possible because the operator has finite order up to scalar.
On each orbit of basis words the weighted cyclic action has at
most one vector for each fixed eigenvalue.  There are $a$ constant
words; every other orbit has size at least two.  Its largest
eigenspace therefore has dimension at most
\[
                  a+(a^l-a)/2\le3a^l/4.
\]
Tensoring with all other cycles preserves this fraction, giving
$\nu(g)\ge e/4\ge23e/128$.  This is the same elementary
weighted-word argument as D-SPOR-03, now with no restriction on
the component types.  Genuine scalar inertia moves all $e$
dimensions.

Assume for contradiction that the actual period rank is nine.
The proved intermediate-cover factorization of D-BI-01 applies
to this inducing block, with arbitrary inducing field.  If its
intermediate base has genus $h$ and $R$ genuine branch points,
the multiplicity-six Euler identity is
\[
                    D=6+2e(1-h),\qquad R\le D.
\]
At $h=1$ the actual variation factors through a finite Hurwitz
cover of dimension at most six.  At $h\ge2$ the displayed $D$
is negative since $e\ge14$.  Hence $h=0$.  The same actual
period factorization then forces $R\ge12$, while $D=2e+6$.
Applying the uniform normalizer bound to every genuine inertia
gives the necessary integer and continuous inequalities
\[
  12\left\lceil\frac{23e}{128}\right\rceil\le2e+6,
                 \qquad e\le\frac{192}{5}.                \tag{BI.90}
\]
This uses an upper bound for the actual period map, not an
inference of dominance from parameter counting.

If $t\ge2$, then $e\ge14^2=196$, contrary to (BI.90).
For a single factor of field size at least $32$, one has
$e\ge124$, again impossible.  For a single
$\operatorname{Sz}(8)$ factor the stronger quarter bound gives
$3e\le2e+6$, so $e\le6<14$.
For each remaining single small factor, the stronger $3/16$
bound gives $9e/4\le2e+6$, so $e\le24$, whereas its degree
is at least $26$ by (BI.87).  Every case is impossible.
The target period domain has dimension nine, so its rank is at
most eight, as asserted.
\end{proof}

\paragraph{Boundary of the result.}
The nonsimple group ${}^2G_2(3)$ is not added to the displayed
simple-family list.  Its derived simple group is
$\operatorname{PSL}_2(8)$ and belongs to the separately treated
$A_1$ case.  No claim about the other classical or exceptional
Lie families is inferred from the Suzuki--Ree calculation.

\section{All linear Weil representations in the transitive normal-core gate}
\label{sec:linear-weil-core}

\paragraph{Scope and source distinction.}
Here a linear Weil representation means a nontrivial constituent of the
complex permutation representation of $\operatorname{GL}_n(q)$ on
$\mathbf F_q^n$, restricted to $\operatorname{SL}_n(q)$, with $n\ge3$.
This is the general-linear family in G\'erardin's Proposition~4.2,
pp.~82--83, including its complete function-space proof
\cite{BI:GerardinLinear}. It is different from the symplectic and
unitary Weil families. The model is published; the support estimates,
source exclusion and finite endpoint computation below are new
derivations. No non-Weil representation of a linear group is excluded
merely by this section.

As before, $J$ is the finite image of an inducing representation
$\tau$ for the six-dimensional multiplicity variation, and
$e=\dim_{\mathbf C}\tau$. Its character field may be larger than the
original imaginary quadratic field $K$; only
$K\subset\mathbf Q(\chi_\tau)$ is used. Write
$\nu(T)=\min_z\operatorname{rank}(T-z)$.

\begin{lemma}[D-BI-23a: models, irreducibility and trace bounds]
\label{D-BI-23a}
Put
\[
 D=\frac{q^n-1}{q-1},\qquad
 A=D-1=\frac{q^n-q}{q-1},\qquad
 B=\frac{q^{n-1}+q-2}{q-1}.
                                                               \tag{LW.1}
\]
For every nontrivial character $\vartheta$ of $\mathbf F_q^*$,
the space
\[
 W_\vartheta=\{f:\mathbf F_q^n\setminus\{0\}\longrightarrow\mathbf C:
                f(tx)=\vartheta(t)f(x)\}
                                                               \tag{LW.2}
\]
is irreducible under $\operatorname{SL}_n(q)$ and has degree $D$.
The deleted projective-point permutation space $W_1^0$ is irreducible
of degree $A$. For a nonscalar element of the represented core,
\[
 |\chi_\vartheta(g)|\le B\quad(\vartheta\ne1),\qquad
 |\chi_1^0(g)|\le B-1.                                       \tag{LW.3}
\]
When $q=2$ the second bound improves to
$|\chi_1^0(g)|\le2^{n-1}-2$.
Every model with $q=2$ or $q=3$ is realizable over $\mathbf R$;
$W_1^0$ is real for every $q$.
\end{lemma}

\begin{proof}
G\'erardin's model and degree calculation give (LW.2) and the deleted
projective permutation summand. We check restriction to the special
linear group, rather than assuming it preserves irreducibility.
For $n\ge3$, $\operatorname{SL}_n(q)$ is transitive on ordered
linearly independent pairs: extend the pair to a basis and adjust
the third basis vector to correct the determinant. Its orbits on
$\mathbf F_q^n\times\mathbf F_q^n$ are consequently the independent
pairs, the origin, the two types with exactly one vector zero, and
the $q-1$ types $(x,tx)$ with $x\ne0$. There are $q+3$ orbits.

The permutation representation on $\mathbf F_q^n$ decomposes as
two trivial summands, $W_1^0$, and the $q-2$ nontrivial spaces
$W_\vartheta$. Each of the latter spaces is nonzero. The squared
character norm is exactly the number $q+3$ of orbits just computed,
which is already the minimum possible contribution
$2^2+1+(q-2)$. Thus all the indicated nontrivial summands are
irreducible and pairwise inequivalent upon restriction. This also
handles $q=2$, when there are no nontrivial $\vartheta$.

For $\alpha\in\mathbf F_q^*$ let
$d_\alpha=\dim\ker(g-\alpha)$. The scalar-character projector on
the nonzero-vector permutation model gives, up to replacing
$\vartheta$ by its inverse,
\[
 \chi_\vartheta(g)=\frac1{q-1}
        \sum_{\alpha\in\mathbf F_q^*}
                  \vartheta(\alpha)(q^{d_\alpha}-1).
                                                               \tag{LW.4}
\]
For $\vartheta=1$, the same sum counts fixed projective points;
one subtracts one for the deleted summand.

If $g$ is not a scalar natural matrix, then
$d_\alpha\le n-1$ and $\sum d_\alpha\le n$.
For nonnegative integers with these constraints, convexity of
$q^d-1$ bounds their sum by
$q^{n-1}-1+(q-1)=q^{n-1}+q-2$.
For example, transferring a unit from a smaller positive part to
a larger part weakly increases the sum until the pattern is
$(n-1,1)$, or a smaller total pattern. This proves the first
bound. The deleted projective character is an integer between
$-1$ and $B-1$; here $B-1\ge1$, proving the second.
When $q=2$ there is only one possible nonzero scalar eigenvalue,
so its exact character is $2^{d_1}-2$. Since $d_1\le n-1$ and
$n\ge3$, its absolute value is at most $2^{n-1}-2$.

These estimates apply to every nonscalar element of the represented
core. Indeed scalar natural matrices act scalarly on each model.
Conversely the preimage of the projective kernel is a normal
subgroup of $\operatorname{SL}_n(q)$. Its image in the simple group
$\operatorname{PSL}_n(q)$ is trivial or the whole group. The latter
would make a nontrivial irreducible representation of degree
greater than one projectively scalar, which is impossible.
Thus the projective kernel is exactly the natural center. This
argument also shows that any central quotient afforded by the
model has the same relevant projective core.

The projective permutation model is real. For $q=3$ the only
nontrivial scalar character has values $\pm1$, and (LW.2) has
a basis of real-valued functions. These give the asserted real
forms, and averaging a real inner product supplies a nondegenerate
invariant symmetric bilinear form after complexification.
\end{proof}

\begin{lemma}[D-BI-23b: the degree-twenty-one endpoint]
\label{D-BI-23b}
In either degree-twenty-one linear Weil representation of
$U=\operatorname{SL}_3(4)=3.\operatorname{L}_3(4)$, every
nonscalar finite-order linear normalizer element satisfies
\[
                         \nu(T)\ge8.                         \tag{LW.5}
\]
All scalar phases and all compatible outer automorphisms are included.
\end{lemma}

\begin{proof}
The two models have faithful central characters of the order-three
center, and hence are faithful on $U$. The published ATLAS data give
\[
 |\operatorname{L}_3(4)|=20160,\quad
 \operatorname{Mult}(\operatorname{L}_3(4))=C_3\times C_4\times C_4,
 \quad\operatorname{Out}(\operatorname{L}_3(4))=D_{12}
          \simeq C_2\times S_3.                              \tag{LW.6}
\]
The primary ATLAS page supplies the natural three-dimensional
matrices of the triple cover \cite{BI:AtlasL34}, but also carries
a general work-in-progress warning. We therefore check the
structural assertions independently. Counting ordered bases gives
$|\operatorname{GL}_3(4)|=63\cdot60\cdot48=181440$; the
determinant quotient and special-linear center both have order
three, giving $|U|=60480$ and the displayed simple-group order.
The exceptional status is confirmed in \cite[Chapter~5,
Table~1, p.~35]{BI:HHM}; that table lists the exceptional group,
rather than purporting to specify the multiplier structure.
The full structure $C_4^2\times C_3$ is independently specified
and used in Breuer's complete Schur-cover construction,
\S2.7-4, explicitly citing the printed ATLAS p.~23
\cite{BI:BreuerL34}.

Breuer's \S2.6-4 independently records the full outer group
$D_{12}$ and its three order-two subgroup types. This is also
consistent with the standard special-linear automorphism
description in his \S11.5-7: the diagonal outer generator
$\delta$ has order three, the field involution $f$ squares
matrix entries, and the graph involution $\gamma$ is inverse
transpose. The last two commute and both invert $\delta$.
Thus $(C_3\rtimes C_2)\times C_2\simeq S_3\times C_2$
is the stated outer group. The numerical character-table
premises below are the complete printed-ATLAS library records,
not an inference from the warning-bearing web page.
In particular $U$ is the actual triple cover,
not an assumed universal cover. The prime-to-three part of the
universal center is characteristic, so every simple-group
automorphism lifts to $U$. An automorphism of the perfect group
$U$ trivial on its central quotient is trivial: its central
discrepancy is a homomorphism to an abelian group.

Every cyclic subgroup of $D_{12}$ is conjugate to the identity,
one of its three order-two types, its order-three subgroup, or its
order-six subgroup. We retain the complete six published ATLAS
tables and the five actual normal-core fusions:
\[
\begin{array}{c|r|l|r|r}
\text{table}&\text{classes}&\text{selected rows}&
       \text{nonscalar entries}&\min\nu\\ \hline
3.L_3(4)&28&17,18&50&8\\
3.L_3(4).2_1&38&27,28,29,30&140&8\\
3.L_3(4).2_2&23&\text{none}&0&\text{--}\\
3.L_3(4).2_3&23&\text{none}&0&\text{--}\\
3.L_3(4).3&60&25,\ldots,30&342&8\\
3.L_3(4).6&66&31,\ldots,42&756&8.
\end{array}                                                   \tag{LW.7}
\]
A row is selected precisely when its pointwise restriction along
the retained fusion equals one of the two degree-twenty-one core
rows. Selection is not based only on the degree. Each table record
explicitly credits the ATLAS as its origin. The input file
\texttt{linear\_weil/linear\_weil\_inputs.json} contains every
character row, class size and order, all prime power maps, the
five fusions, and the source metadata. Its SHA-256 is
\texttt{12becae3b988a84d6820b658da1ac3dcba4a2a83ff3be85c1342e6cd9bed5e76}.
This is copied published numerical input; the following reconstruction
is an exact independent arithmetic implementation of its consequences.

For a class of order $d$, its eigenvalue multiplicities are
\[
 m_k=\frac1d\sum_{r=0}^{d-1}\chi(g^r)\zeta_d^{-kr},
                  \qquad0\le k<d.                           \tag{LW.8}
\]
The prime power maps and cyclotomic Galois action determine every
power in this expression. The GAP program checks all table-row
orthogonality and computes all spectra. The separate Python program
\texttt{check\_linear\_weil.py} selects rows by polynomial
restriction equality, verifies all five normal fusions and cyclic
quotients, and reconstructs (LW.8) by normalized rational traces,
that is, Ramanujan sums. It imports no GAP scan results. Both
programs pass all twenty-four selected rows and all1288 nonscalar
entries. The minimum moving rank is eight in each selected row.

We explain why this finite calculation covers an arbitrary actual
normalizing operator; no assertion that the whole normalizer equals
one tabulated extension is necessary. Let $\alpha$ be the induced
automorphism of $U$. Its outer class generates one of the cyclic
subgroups just listed. After conjugating that subgroup and changing
$\alpha$ by an inner automorphism, use a corresponding extension
$A=\langle U,t\rangle$ from (LW.7), with $A/U$ cyclic of order $d$.
The transformed core representation is still one of the two
degree-twenty-one rows. The actual operator implementing the
automorphism differs from the chosen intertwiner only by a scalar.
The relation
\[
            T^d\rho(t^d)^{-1}\in\mathbf C^*
                                                               \tag{LW.9}
\]
follows from Schur's lemma. Rescale $T$ by a $d$th root so that the
defining relation holds exactly. Then $\rho$ extends to an actual
representation of $A$, irreducible because its restriction to $U$
is irreducible. It therefore occurs in one of the selected rows.
The original element, including its possible inner factor, is scalar
times a matrix of some element of $A$, all of whose classes have
been checked. Scalar rescaling preserves $\nu$. If the relevant
table has no selected row, no such compatible operator exists.
The case of trivial outer class is covered by the core table.
This proves (LW.5).
\end{proof}

\begin{theorem}[D-BI-23: exclusion of transitive linear Weil cores]
\label{D-BI-23}
Suppose the inducing representation of an actual finite-cover
Weil-sixfold source has an irreducibly acting normal central product
of quasisimple factors whose simple quotients are
$\operatorname{PSL}_n(q)$, with $n\ge3$. Suppose every factor
representation is a linear Weil model of D-BI-23a, and the full
inducing group permutes the factors transitively. Then its period
map has rank less than nine.

This includes the transitive forced-core situation of D-BI-19,
all compatible outer automorphisms, arbitrary scalar phases and
tensor-factor permutations. It allows any imaginary quadratic
original field, a larger inducing character field, and either
polarization determinant. It does not exclude non-Weil factor
characters or claim a unitary-family result.
\end{theorem}

\begin{proof}
Transitivity makes the factor degrees equal, say $a$, and the total
degree is $e=a^t$ for some $t\ge1$. Each $a$ is $A$ or $D$, and
$a\ge6$. A nonscalar element of the represented central product
has at least one nonscalar factor, so a factor character bound
$|\chi(g)|\le ca$ implies the same normalized bound $c$ on the
entire core. D-BI-18 therefore applies to every nonscalar element
of the full inducing group, including those permuting factors.

For completeness, the actual intermediate-cover Euler identity is
$6=2e(h-1)+\sum_j\operatorname{rank}(1-g_j)$. If $h\ge2$, its
rightmost sum would be $6+2e(1-h)<0$. If $h=1$, there are at most
six genuine branch points and at most six moduli parameters.
Rank nine consequently requires $h=0$, at least twelve genuine
inertias, and total moving rank $2e+6$. These are consequences of
the actual period factorization, not a claim that the dimension
inequality suffices for dominance.

First let $q\ge5$. In (LW.3) use the common bound
$c=B/A<1$. Since $e\ge A$, the necessary continuous inequality
of D-BI-18 yields
\[
 6\ge e(1-3c)\ge A-3B
       =\frac{(q-3)q^{n-1}-4q+6}{q-1}>6.                   \tag{LW.10}
\]
The last expression increases with $n$ for $q\ge5$. At $n=3$
it equals $q^2-2q-6$, which is nine at $q=5$ and strictly
increases thereafter. This is a contradiction.
For $q=4$ and $n\ge4$, the same expression is
$(4^{n-1}-10)/3$, at least eighteen, also a contradiction.

It remains to consider $q=2,3$, and $(q,n)=(4,3)$.
In every real factor case of D-BI-23a, the full irreducible tensor
core has a unique invariant symmetric bilinear-form line. Its
normalizer is essentially self-dual. The inclusion
$K\subset\mathbf Q(\chi_\tau)$ forces a nontrivial similitude
character: if the form were preserved exactly, the character would
be real. Since genuine inertias generate and have product one,
at least two have nontrivial multiplier, and their fixed spaces
are totally isotropic. D-BI-18a gives the necessary bound
\[
                         e(3-5c)\le12.                     \tag{LW.11}
\]
This use of the original field does not identify it with the
possibly larger inducing field.

For $q=3$ and the nontrivial scalar character, $a=D\ge13$ and
\[
 B=\frac{a+2}{3},\qquad
 a(3-5B/a)=\frac{4a-10}{3}\ge14>12.
\]
Since $e\ge a$ and $3-5B/a>0$, (LW.11) is impossible.
For the deleted summand, $a=A\ge12$ and $B-1=a/3$, so it is
already impossible because $a(3-5/3)=4a/3\ge16$.

For $q=2$, there is only the real deleted summand,
$a=2^n-2$, and the improved bound gives
$c=(a-2)/(2a)$. Hence
\[
                  e\le\frac{24a}{a+10}.                    \tag{LW.12}
\]
For $n\ge5$, $a\ge30$ and even $e=a$ violates (LW.11), since
$a(3-5c)=a/2+5>12$. For $n=3,4$, (LW.12) is respectively
$e\le9$ and $e\le14$, forcing $t=1$ and $a=6,14$.
The first core is $\operatorname{SL}_3(2)\simeq
\operatorname{PSL}_2(7)$; its degree-six full normalizer has
$\nu\ge a/4$ by D-BI-21c. The second is
$\operatorname{SL}_4(2)\simeq A_8$, and its ordinary degree-fourteen
full normalizer has the same bound by HD-ALT-02a(3).
Both are incompatible with the exact branch budget:
\[
           12\lceil6/4\rceil=24>18,
       \qquad12\lceil14/4\rceil=48>34.                      \tag{LW.13}
\]
These endpoint spectral bounds themselves have no nonsplit
polarization hypothesis.

Finally take $(q,n)=(4,3)$. For the deleted summand, $a=20$,
$B-1=5$ and $c=1/4$. It is real, and
$a(3-5c)=35>12$, excluding it and all its transitive tensor powers.
For either nontrivial scalar character, $a=21$, $B=6$ and
$c=2/7$. D-BI-18 forces $e\le42$, hence $t=1$.
D-BI-23b gives actual moving rank at least eight for every
nonscalar inertia; a nonidentity scalar moves all twenty-one
dimensions. Thus total moving rank is at least
$12\cdot8=96$, whereas $2e+6=48$. This last contradiction
completes every prime-power and every tensor-exponent case.
\end{proof}

\paragraph{Reproducibility and limits.}
Run \texttt{linear\_weil/export\_and\_scan.g} in the retained GAP
runtime and \texttt{python check\_linear\_weil.py} in this directory.
The latter also checks the displayed endpoint arithmetic. Its first
run stopped on a metadata-string expectation: the library writes
the multiplier as \texttt{(3x4\^{}2)}, not \texttt{3.4\^{}2}.
Correcting that literal expectation gave a complete PASS; no
mathematical inequality or table entry changed. Source-integrity
checks are recorded separately from TeX compilation, which was
unavailable. The transitive-core hypothesis is explicit and is
supplied in the main primitive tensor-indecomposable gate by
D-BI-19. No arbitrary mixed-type tensor classification is asserted.

\section{The unitary Weil normal-core gate}
\label{sec:unitary-weil-core}

\paragraph{Status and scope.}
This is a new author derivation from published Weil-character and
automorphism inputs, with complete retained finite endpoint tables.
The author has reread the entire proof and executed both the complete
GAP scan and a separate exact Python reconstruction; all pass.
An independent whole-proof, primary-source and complete arithmetic
audit by /root/imprimitive has passed. It concerns
ordinary unitary Weil representations, not arbitrary irreducible
representations of unitary groups. The transitivity hypothesis below
is explicit; it is the one supplied by D-BI-19 in the main primitive,
tensor-indecomposable gate. No equality of the original and inducing
character fields is assumed.

Write $\nu(T)=a-\max_\lambda\dim\ker(T-\lambda)$ on a space of
dimension $a$. A scalar has $\nu=0$; a genuine nonidentity scalar
inertia has ordinary moving rank $a$. These are different quantities.

\begin{lemma}[D-BI-24a: unitary Weil characters and stabilizers]
\label{D-BI-24a}
Let $n\ge3$, let $q=p^f$, and suppose $(n,q)\ne(3,2)$.
Put $\epsilon=(-1)^n$. The ordinary unitary Weil constituents
$W_i$, $0\le i\le q$, restrict to pairwise distinct irreducible
representations of $L=\operatorname{SU}_n(q)$. Their degrees are
\[
 a_0=\frac{q^n+q\epsilon}{q+1},\qquad
 a_i=\frac{q^n-\epsilon}{q+1}\quad(i\ne0).                  \tag{UW.1}
\]
Every nonscalar $g\in\operatorname{GU}_n(q)$ satisfies
\[
 |\chi_i(g)|\le B:=\frac{q^{n-1}+2q-1}{q+1},\qquad
 a_i\ge A:=\frac{q^n-q}{q+1}.                              \tag{UW.2}
\]
Complex conjugation sends $\chi_i|_L$ to $\chi_{-i}|_L$,
and entrywise $p$-Frobenius sends its label to $pi$ modulo $q+1$.

For $q=2,3,5$, the full linear normalizer of a nonselfdual
constituent is contained in the scalar multiples of its
$\operatorname{GU}_n(q)$-extension. The same assertion holds
for every $i\ne0$ when $(n,q)=(3,4)$.
For each such nonscalar normalizer $T$,
\[
                      \nu(T)\ge(a_i-B)/2.                  \tag{UW.3}
\]
\end{lemma}

\begin{proof}
The published input is G\'erardin's total unitary Weil representation
$\Omega$ and its exact character formula
\[
 \operatorname{Tr}\Omega(g)
       =(-1)^n(-q)^{\dim_{\mathbf F_{q^2}}\ker(g-1)};
                                                               \tag{UW.4}
\]
see \cite[Theorem4.9.2 and its complete proof, pp.~99--100]{BI:GerardinLinear}.
This input includes even characteristic: unitary \S3.1 on p.~70
starts with an arbitrary finite field, and Theorem3.3(a\textnormal{'''})
on p.~73 explicitly constructs the even-$q$ representation.
The odd-characteristic restriction in that article belongs to the
symplectic \S2 and Theorem4.9.1, not to Theorem4.9.2.
The central-character projector formula is also explicitly recorded
in \cite[\S4.3, (4.3)--(4.4), author-version p.~16]{BI:GLOST2018};
the ordinary degrees and extensions are stated in
\cite[p.~293]{BI:GMST2002}.

Here is a direct check of the restriction assertion used below.
The scalar center of $\operatorname{GU}_n(q)$ is cyclic of order
$q+1$. Projecting $\Omega$ to its characters gives $q+1$ nonzero
spaces, with degrees (UW.1), by evaluating (UW.4) at scalar matrices.
Furthermore, $|\operatorname{Tr}\Omega(g)|^2=q^{2\dim\ker(g-1)}$
is the permutation character of $L$ on its natural vector space.
This action has exactly $q+1$ orbits: zero, nonzero isotropic
vectors, and one orbit for each of the $q-1$ nonzero norm values.
Witt extension gives transitivity for the full unitary group.
To adjust the determinant to one while fixing a target vector,
use a nondegenerate complement of dimension $n-1$ in the
anisotropic case, or a nondegenerate complement of dimension
$n-2$ to a hyperbolic plane in the isotropic case. Either
complement contains an anisotropic line on which any scalar
of norm one may act. Thus the determinant adjustment is available
for $n\ge3$, proving the same orbit assertion for $L$.
Consequently $\langle\Omega|_L,\Omega|_L\rangle=q+1$.
There are already $q+1$ nonzero central summands; positivity of
character inner products forces each restriction to be irreducible
and distinct from all the others.

Fix a primitive $(q+1)$st root $\rho$ in $\mathbf F_{q^2}$ and
a corresponding complex root $\zeta$. With a choice of label sign,
the projector formula is
\[
 \chi_i(g)=\frac{(-1)^n}{q+1}
       \sum_{j=0}^{q}\zeta^{ij}(-q)^{d_j(g)},\qquad
 d_j(g)=\dim\ker(g-\rho^j).                                \tag{UW.5}
\]
For a nonscalar natural matrix, $0\le d_j\le n-1$ and
$\sum_jd_j\le n$. The convex sum $\sum_jq^{d_j}$ is at
most its value at $(n-1,1,0,\ldots,0)$: move one unit from
a smaller positive entry to a larger entry until an entry
reaches $n-1$, and then concentrate the remaining unit.
That value is $q^{n-1}+q+(q-1)$, proving (UW.2).
Formula (UW.5) also proves the two label rules directly.

For precision, the published automorphism input is
\cite[\S2.2, Lemma6, p.~114]{BI:YangWuRevinVdovin2023}:
\[
 \operatorname{Aut}(\operatorname{PSU}_n(q))
      =\operatorname{PGU}_n(q)\rtimes\langle\sigma\rangle,
 \qquad |\sigma|=2f,                                       \tag{UW.6}
\]
where $\sigma$ is induced by entrywise $p$-Frobenius.
Only the standard simple unitary groups specified here are used.
For $n\ge3$ the nonsimple pair $(3,2)$ has already been excluded.
In particular the small groups $\operatorname{PSU}_3(3)$,
$\operatorname{PSU}_4(2)$, $\operatorname{PSU}_5(2)$,
$\operatorname{PSU}_6(2)$ and $\operatorname{PSU}_3(4)$,
$\operatorname{PSU}_3(5)$ all lie within this stated domain.
If a linear normalizer induces an automorphism on the represented
central quotient of $L$, take a standard lift $\alpha$ from
(UW.6) to $L$. Its action and conjugation by the normalizer agree
projectively. Their ratio is a scalar character of $L$, hence
trivial because $L$ is perfect. This also shows that the lift
preserves the kernel of the represented constituent whenever
that normalizer exists. No unproved characteristic-kernel claim
about an exceptional Schur multiplier is required.

All inner-diagonal automorphisms stabilize the constituent,
since it extends to $\operatorname{GU}_n(q)$. For $q=2,3,5$
the nontrivial field quotient sends $i$ to $-i$; it cannot
stabilize a nonselfdual constituent, by the distinctness just
proved. When $q=4$, multiplication by $2$ has order four on
every nonzero label modulo five, so no nonidentity field
quotient stabilizes such a label. Schur's lemma now makes
every allowed normalizer a scalar multiple of the known
unitary-group operator. A nonscalar represented operator
has a nonscalar natural lift. Finally, if an eigenspace has
dimension $m$, rotating its eigenvalue to one gives
$|\operatorname{Tr}(T)|\ge2m-a_i$. This and (UW.2) give (UW.3).
\end{proof}

\begin{lemma}[D-BI-24b: the three exact small endpoints]
\label{D-BI-24b}
For the following ordinary constituents, every nonscalar
finite-order linear normalizer of the represented simple
core has the indicated support bound:
\[
\begin{array}{c|c|c}
\text{simple core}&\text{constituent degrees}&\nu(T)\text{ lower bound}\\\hline
\operatorname{PSU}_3(3)&6,7&a/4\\
\operatorname{PSU}_5(2)&10&a/5\\
\operatorname{PSU}_6(2)&22&a/4.
\end{array}                                                  \tag{UW.7}
\]
The degree-ten constituent is selfdual. The degree-twenty-two
constituent is $W_0$ and is trivial on the center of
$\operatorname{SU}_6(2)$.
\end{lemma}

\begin{proof}
These are complete finite character-table computations from
published ATLAS inputs, not inferred normalizer bounds from
samples. The retained seven tables, selected rows and exact
nonscalar-entry counts are
\[
\begin{array}{c|c|c|c}
\text{table}&\text{rows}&\text{entries}&\min\nu\\\hline
\texttt{U3(3)}&2,3,4,5&52&2,3,2,2\\
\texttt{U3(3).2}&3,4,5,6&60&2,2,3,3\\
\texttt{U5(2)}&2&46&2\\
\texttt{U5(2).2}&3,4&84&2,2\\
\texttt{U6(2)}&2&45&6\\
\texttt{U6(2).2}&3,4&128&6,6\\
\texttt{U6(2).3}&4,5,6&339&6,6,6.
\end{array}                                                  \tag{UW.8}
\]
All characters, class sizes and orders, prime power maps,
Frobenius--Schur indicators, source strings, and core fusions
are retained in \texttt{unitary\_weil\_inputs.json}.
The exporter \texttt{unitary\_weil/export\_and\_scan.g}
checks full row orthogonality. The separate exact Python
checker reselects rows by pointwise restriction equality,
checks the normal fusions and cyclic quotients, and reconstructs
every eigenspace multiplicity from all powers:
\[
 m_b=\frac1d\sum_{j=0}^{d-1}\chi(g^j)\zeta_d^{-bj}.
                                                               \tag{UW.9}
\]
It checks integrality, nonnegativity and total degree. Thus
all seventeen selected rows and all754 nonscalar spectra
are checked. The full tables retain their published status;
orthogonality is a consistency check and does not replace
the published identification of the groups and tables.
The retained input SHA-256 is
\texttt{0c84735f987b0bea3a9fea197002d4cb28af2cdae9a165bad8bce3bf6f91a117};
the Python program and actual output are
\texttt{check\_unitary\_weil.py} and
\texttt{check\_unitary\_weil.json}. The first GAP attempt
stopped on an unsupported metadata operation for an extension
table. That metadata call was restricted to the simple core;
no character data changed, and the complete rerun passed.

The simple outer groups in these cases are respectively
$C_2,C_2,S_3$. This follows either from the retained official
extension metadata or from (UW.6): for the third group the
diagonal group is $C_3$, inverted by the field involution.
The specified extensions cover every cyclic subgroup of
each outer group up to conjugacy. For an actual normalizer,
adjust by an inner operator to match a chosen generator
$t$ of one such fixed extension $E/S$. If its order modulo
$S$ is $d$, then $T^d\rho(t^d)^{-1}$ is scalar by Schur's
lemma. Rescale $T$ by a $d$th root of the inverse scalar.
The conjugation and power relations then give a representation
of the fixed $E$, irreducible because its restriction to $S$
is irreducible. It is one of the exactly selected rows.
Putting back the inner operator still gives a checked element
of $E$, and scalar rescaling does not change $\nu$.

For the degree-ten and degree-twenty-two cases, the unique
rows in (UW.8) are exactly $W_0$ by (UW.1); (UW.5) makes
their characters real. In degree22 the central projector
is the trivial one, so the representation factors through
the simple group as stated. This proves every assertion.
\end{proof}

\begin{theorem}[D-BI-24: reduction of all transitive unitary Weil cores]
\label{D-BI-24}
Consider an inducing image $J\subset\operatorname{GL}(V)$
of an actual rank-nine Weil sixfold finite-cover source.
Suppose its character field contains the original imaginary
quadratic field $K$. Suppose $J$ has an irreducibly acting
normal central product $N=N_1\cdots N_t$ with all simple
quotients isomorphic to $\operatorname{PSU}_n(q)$,
$n\ge3$, $(n,q)\ne(3,2)$, and with the factor representations
given by ordinary unitary Weil constituents. Suppose $J$
acts transitively on the factors.

Then necessarily
\[
                   (n,q)=(4,2),\qquad t=1,\qquad
                   \dim V\in\{5,6\}.                       \tag{UW.10}
\]
For the eight non-Eisenstein class-number-one fields of the
main nonsplit problem, with the actual original polarization
nonsplit, only degree five remains. It is subject to the
already proved rank-five inducing restrictions; no new
existence or dominance claim is made for it.
\end{theorem}

\begin{proof}
Write $a=\dim W_i$ for the common factor degree, so $e=\dim V=a^t$.
Transitivity makes the constituent labels automorphically
related and in particular makes selfduality uniform.
The intrinsic tensor matrix algebras of the simple factors
are permuted by $J$. A factor-preserving normalizer is a
tensor product of factor intertwiners. The maximum eigenspace
fraction of a tensor product is at most that of any nonscalar
factor. A nontrivial weighted factor cycle has the stronger
quarter support bound from D-SPOR-03. The cycle product returns
to the same constituent; this permits identifications between
different automorphic or dual labels without allowing an
incompatible field automorphism on a fixed constituent.

The actual intermediate-cover Euler identity is
$D=6+2e(1-h)$. Here $a\ge5$, so genus $h\ge2$ is impossible.
For $h=1$, at most six genuine branch points remain; the
marked elliptic-base dimension is at most six (and is one
in the unbranched case). The actual period variation factors
through that finite Hurwitz cover and cannot have rank nine.
Thus a rank-nine source necessarily has
\[
                h=0,\qquad r\ge12,\qquad D=2e+6.            \tag{UW.11}
\]
This is only a necessary consequence of the actual variation
factorization, not a claim of dominance from branch dimension.

For the generic trace estimate put $c=B/a$. On the central
product, every nonscalar element has character ratio at
most $c$. D-BI-18 therefore gives
\[
                       e(1-3B/a)\le6.                      \tag{UW.12}
\]
When all factors are selfdual, their tensor product has a
unique invariant nondegenerate bilinear form line, symmetric
or alternating. Every element of $J$ acts on it by a multiplier.
The multiplier is nontrivial: otherwise the finite character
of $J$ is real, contradicting $K\subseteq\mathbf Q(\chi_J)$.
The generating product-one inertia tuple has at least two
nontrivial multipliers; their fixed spaces are isotropic.
D-BI-18a consequently yields
\[
                       e(3-5B/a)\le12.                     \tag{UW.13}
\]
These arguments apply to the total degree and do not require
primitivity of the inducing block.

\emph{First, the large-parameter ranges.}
By (UW.2),
\[
 A-3B=\frac{(q-3)q^{n-1}-7q+3}{q+1}.                       \tag{UW.14}
\]
For $q\ge4$, $n\ge4$, this is at least
$q^2-q-6+9/(q+1)$, whose value at four is $39/5>6$.
Its forward difference is $2q-9/((q+1)(q+2))>0$.
For $n=3$, $q\ge7$, (UW.14) equals
$q^2-4q-3+6/(q+1)$, with value $75/4>6$ at seven
and positive forward difference
$2q-3-6/((q+1)(q+2))$.
Since $e\ge a\ge A$, these strict bounds contradict (UW.12).
Only $q=2,3$ and $(n,q)=(3,4),(3,5)$ remain.

\emph{Second, the nonselfdual constituents.}
For $q=3$ the degree is $a=(3^n-(-1)^n)/4$ and
$B=(3^{n-1}+5)/4\le a/2$ for all $n\ge3$.
The direct stabilizer estimate (UW.3) and the tensor-cycle
bound give $\nu(g)\ge e/4$ for every nonscalar inertia.
The same holds at $(n,q)=(3,4)$, where $(a,B)=(13,23/5)$,
and at $(n,q)=(3,5)$, where $(a,B)=(21,17/3)$.
Nonidentity scalar inertia also satisfies the ordinary
moving-rank bound. Equation (UW.11) would give $e\le6$,
contrary to the respective factor degrees, which are at least seven.

For $q=2$ and $n\ge5$,
\[
 a=(2^n-(-1)^n)/3,\quad B=(2^{n-1}+3)/3,\quad
 \gamma=(a-B)/(2a)<1/4.
\]
All nonscalar normalizers on the tensor product have support
at least $\gamma e$. The exact budget implies
\[
 e(4a-6B)/a\le6,\qquad
        4a-6B=\frac{2^n-4(-1)^n-18}{3}.                    \tag{UW.15}
\]
For $n\ge6$ the rightmost expression is at least14, so
$e\ge a$ is impossible. At $n=5$, $(a,B)=(11,19/3)$
and $\gamma=7/33$; (UW.15) forces $e\le11$, hence $t=1$.
Every genuine inertia then moves at least $\lceil7/3\rceil=3$
dimensions, giving $D\ge36>28=2e+6$. The remaining
nonselfdual $q=2$ constituents have $(n,a)=(4,5)$.

\emph{Third, the selfdual constituents.}
For $q=3$, $n\ge4$, use $a\ge A$ in
\[
                    3a-5B\ge3^{n-1}-17/2>12,
\]
contradicting (UW.13). At $(n,q)=(3,4)$ the selfdual
constituent has $a=12$ and $3a-5B=13>12$.
At $(n,q)=(3,5)$ the selfdual degrees are20 and21;
already $3\cdot20-5\cdot17/3=95/3>12$.
For $q=2$ the sole selfdual constituent is $W_0$.
For $n\ge7$,
\[
 3a_0-5B=\frac{2^{n-1}+6(-1)^n-15}{3}
                   \ge\frac{2^{n-1}-21}{3}>12,
\]
again contradicting (UW.13).

The remaining small cases are covered in full by (UW.7).
For $\operatorname{PSU}_3(3)$, every tensor normalizer has
the quarter bound, so $e\le6$. This leaves only the single
degree-six possibility; its exact integer inequality is
$12\lceil6/4\rceil=24>18$, also impossible.
For the degree22 constituent of $\operatorname{PSU}_6(2)$,
the quarter bound contradicts $e\ge22$ immediately.
For the selfdual degree10 constituent of $\operatorname{PSU}_5(2)$,
the full-normalizer fifth bound and the two nontrivial
multiplier costs give
\[
                   2e+6\ge2(e/2)+10(e/5)=3e.
\]
Thus $e\le6$, contrary to $e\ge10$.

Only $\operatorname{PSU}_4(2)$ and its degrees5 and6 remain.
The complete small-core theorem D-SPW-04b, which allows
every compatible outer and every factor permutation,
excludes their tensor products and gives $t=1$.
Its final clause removes the degree-six case in precisely
the stated non-Eisenstein nonsplit scope. Its degree-four
case cannot appear here because it is not an ordinary
unitary Weil constituent of $\operatorname{SU}_4(2)$.
The remaining degree-five case is the already isolated
rank-five inducing gate. This proves every assertion.
\end{proof}

\paragraph{Boundary of the result.}
No conclusion about non-Weil unitary representations follows.
No assertion about arbitrary mixed normal-core types or a
nontransitive family of factor labels is made. The result
eliminates an infinite family of possible sixfold suppliers
and identifies its exact existing small endpoint; it produces
no cycle and proves no additional period-map dominance.

\section{An all-rank bound for odd-characteristic symplectic Weil cores}
\label{sec:symplectic-weil-core}

\paragraph{Status and scope.}
The following new author derivation has passed an independent whole-proof,
source and arithmetic review by \texttt{/root/g34\_eightfolds}.
The root reviewer has separately passed the complete mathematical
argument and finite arithmetic; its fresh source review is recorded
separately.  The author's finite arithmetic certificate has been executed; the
unbounded-rank assertions are proved by the displayed inequalities,
not by that finite check.  No assertion about a non-Weil irreducible
representation of a symplectic group is made.  In particular this
section does not remove the general small-field Lie-core gate.

For a finite-order complex operator $T$ on a space of dimension $a$, set
\[
 \nu(T)=a-\max_{\lambda}\dim\ker(T-\lambda).
\]
This quantity is unchanged by multiplying $T$ by a scalar.  It is
bounded above by the ordinary moving rank $\operatorname{rank}(T-1)$.
We allow arbitrary scalar phases throughout.

\subsection*{D-SPW-01. The Heisenberg model and its exact trace identity}

Let $q=p^f$ be odd, $n\ge2$, and let $V$ be a $2n$-dimensional
symplectic space over $\mathbf F_q$ with form $\langle\ ,\ \rangle$.
The Heisenberg group is
\[
 H(V)=V\times\mathbf F_q,\qquad
 (v,z)(w,t)=(v+w,z+t+\tfrac12\langle v,w\rangle).
\]
Fix a nontrivial additive character $\psi$ of $\mathbf F_q$.
Choose a Lagrangian decomposition of $V$ and realize the irreducible
Heisenberg representation with central character $\psi$ on the
space $\mathcal W$ of functions on one Lagrangian.  Its dimension is
$Q=q^n$.  Write $D(v)$ for the operator of $(v,0)$.

The linear Weil representation and its Schr\"odinger model are
published inputs, specifically G\'erardin
\cite[Theorem~2.4(a), pp.~63--64, proof pp.~64--69]{SPW:Gerardin}.
The model on the unipotent radical of a Lagrangian stabilizer is
in equations (2.8)--(2.9) on p.~65.  The following exact trace
and irreducibility statements are also published there; we use
them directly rather than reprove them.

\begin{lemma}[D-SPW-01: published trace input and scalar normalization]
\label{D-SPW-01}
For the full Weil representation $U$ of a finite symplectic
space in odd characteristic,
\[
          |\operatorname{tr}U(A)|^2=|\ker(A-1)|.            \tag{SPW.1}
\]
For the space $V$ above, the parity operator
$P:F(x)\mapsto F(-x)$ has two irreducible nonisomorphic
$\operatorname{Sp}(V)$ eigenspaces, of dimensions
$(Q+1)/2$ and $(Q-1)/2$.
The same absolute-trace formula holds for any unitary
Heisenberg intertwiner implementing $A$, and that intertwiner
commutes with parity.  This last formulation permits the
trace-form specialization used for field automorphisms below.
\end{lemma}

\begin{proof}
Formula (SPW.1) is exactly G\'erardin
\cite[Theorem~4.4(c), pp.~84--85]{SPW:Gerardin}.
The parity decomposition and its irreducibility are
\cite[Corollary~4.4(a) and its proof, p.~85]{SPW:Gerardin}.
These published assertions hold over every odd finite field,
so apply both to $\mathbf F_q$ and to $\mathbf F_p$ later.

It remains only to justify the formulation for an arbitrary
unitary intertwiner.  Irreducibility of the Heisenberg
representation makes two intertwiners for the same $A$ scalar
multiples of one another.  If both are unitary the scalar has
absolute value one, so (SPW.1) is unchanged.
The operators $P$ and $U_A P U_A^{-1}$ implement the same
central transformation $-1$ of the finite symplectic space,
so they too differ by a scalar.  Both have trace one, since
zero is the unique fixed point of $x\mapsto-x$ in odd
characteristic.  That scalar is one, as required.
\end{proof}

We will also use the uniqueness in Theorem~2.4(a): on a fixed
irreducible Heisenberg space any two linear Weil extensions
differ by a one-dimensional character of the symplectic group.
For $n\ge2$ the group is perfect, so that character is trivial.
This identifies the two realizations in the field-restriction
argument below without an additional extension theorem.

\subsection*{D-SPW-02. Every scalar and outer normalizer}

\begin{theorem}[D-SPW-02: a uniform one-fifth bound]
\label{D-SPW-02}
Let $\rho$ be any complex irreducible Weil constituent of
$\operatorname{Sp}_{2n}(q)$, with $q$ odd and $n\ge2$, and let
$N=\rho(\operatorname{Sp}_{2n}(q))$.  Put
$a=\dim\rho=(q^n\pm1)/2$.
Every nonscalar finite-order element $T$ of
$N_{\operatorname{GL}_a(\mathbf C)}(N)$ satisfies
\[
                              \nu(T)\ge a/5.               \tag{SPW.2}
\]
This includes every compatible field automorphism and arbitrary
scalar phases.  The same conclusion applies to any faithful
quasisimple core representation whose projective representation
is one of these Weil representations.
\end{theorem}

\begin{proof}
We first bound inner elements.  If $g\in\operatorname{Sp}_{2n}(q)$
is noncentral, put
\[
 d_+=\dim\ker(g-1),\qquad d_-=\dim\ker(g+1)
\]
on the natural $2n$-dimensional space.  These spaces are disjoint,
so $d_++d_-\le2n$.  If either dimension is $2n-1$, the other is
zero.  To see this, the generalized eigenspace of a symplectic
matrix for the eigenvalue $1$, and likewise for $-1$, is
nondegenerate alternating: different primary spaces are paired
only with their reciprocal eigenvalues.  Its dimension is
therefore even.  An eigenspace of dimension $2n-1$ consequently
has generalized eigenspace of dimension $2n$, leaving no other
eigenvalue.

If neither dimension is $2n-1$, both are at most $2n-2$.
Convexity of $x\mapsto q^{x/2}$ on the interval with fixed sum
at most $2n$ gives
\[
 q^{d_+/2}+q^{d_-/2}
 \le\max\{q^{n-1/2}+1,\ q^{n-1}+q\}
 =q^{n-1/2}+1.                                             \tag{SPW.3}
\]
For the final equality, after factoring $\sqrt q-1$ it is
enough that $q^{n-1}\ge\sqrt q+1$, which holds for $q\ge3$
and $n\ge2$.

The trace on a parity constituent is
$(\operatorname{tr}U_g\pm\operatorname{tr}(PU_g))/2$.
Since $PU_g$ implements $-g$, (SPW.1) and (SPW.3) imply
\[
 \frac{|\operatorname{tr}\rho(g)|}{a}
 \le \frac{q^{n-1/2}+1}{q^n-1}
 =q^{-1/2}+\frac{1+q^{-1/2}}{q^n-1}.                       \tag{SPW.4}
\]
There is no assumption that the largest eigenspace has eigenvalue
one.  For any unitary operator of dimension $a$, the contribution
of an eigenspace of maximal dimension $a-\nu$ gives
\[
        |\operatorname{tr}T|\ge a-2\nu(T).                \tag{SPW.5}
\]
Thus a trace ratio at most $3/5$ implies (SPW.2).

For $q\ge5$, (SPW.4) is less than
$1/2+(3/2)/24=9/16<3/5$.  For $q=3,n\ge4$, its maximum
is at $n=4$, where
\[
 \frac{27\sqrt3+1}{80}
 <\frac{239}{400}<\frac35,
 \qquad \sqrt3<\frac{26}{15}.                              \tag{SPW.6}
\]
At $q=3,n=3$, the same bound is less than $2/3$.
The two degrees are $13,14$, so (SPW.5) and integrality give
$\nu\ge3\ge a/5$.  At $q=3,n=2$, the degrees are $4,5$,
and the elementary nonscalar inequality $\nu\ge1$ suffices.
This proves (SPW.2) for all inner elements.

We now cover the actual outer normalizers.  The published
automorphism input is that for odd $q$ all automorphisms of
$\operatorname{Sp}_{2n}(q)$ are products of inner, diagonal and
field automorphisms; there is no graph automorphism.
See Taylor \cite[\S1.5, author p.~2; \S12.5, author p.~24]
{SPW:Taylor}.  The latter passage gives the full automorphism
group, not merely a subgroup of known automorphisms.  The
universal cover of $\operatorname{PSp}_{2n}(q)$ is
$\operatorname{Sp}_{2n}(q)$ in the range at hand.  In particular
$\operatorname{PSp}_4(3)$ has no exceptional multiplier in its
characteristic-three description; see
\cite[Chapter~5, Table~1 and preceding paragraph, numbered p.~35]
{BI:HHM}.  Consequently every automorphism of a represented
central quotient lifts to this universal perfect cover.

The diagonal automorphism with nonsquare similitude multiplier
does not preserve any Weil constituent.  Here is a direct check,
including the scalar issue.  In the Schr\"odinger model a
transvection with parameter $b\ne0$ acts as
\[
                  F(x)\longmapsto\psi(bx_1^2/2)F(x).
\]
Its full trace is $q^{n-1}G_\psi(b)$, where
$G_\psi(b)=\sum_{y\in\mathbf F_q}\psi(by^2/2)$, and its
parity-twisted trace is one.  Multiplying $b$ by a nonsquare
changes the nonzero Gauss sum to its negative.  Indeed
$G_\psi(b)=\sum_y\eta(y)\psi(by/2)$ for the quadratic
character $\eta$, and substitution gives
$G_\psi(cb)=\eta(c)G_\psi(b)$; its squared absolute value is
$q$ by the elementary orthogonality calculation.
Thus the traces $(q^{n-1}G_\psi(b)\pm1)/2$ on each fixed
parity differ after this automorphism.  The two representations
cannot even be projectively equivalent: a projective equivalence
between linear representations of the perfect group would
differ by a one-dimensional character, which is trivial.

For completeness, we construct and bound field intertwiners.
Write
\[
 \psi(x)=\exp\left(\frac{2\pi i}{p}
                 \operatorname{Tr}_{\mathbf F_q/\mathbf F_p}(a x)\right),
 \qquad a\in\mathbf F_q^*.
\]
Regard $V$ as a $2nf$-dimensional $\mathbf F_p$ symplectic
space with alternating form
$B_a(v,w)=\operatorname{Tr}_{\mathbf F_q/\mathbf F_p}
(a\langle v,w\rangle)$.
This form is nondegenerate by nondegeneracy of the field trace
and of $\langle\ ,\ \rangle$.
The Heisenberg representation for this $\mathbf F_p$ space is
the same representation on the same $q^n$-dimensional function
space: the center of the $\mathbf F_q$ Heisenberg group acts
through its trace character.  The uniqueness just recalled
identifies the models.  Restricting the linear Weil representation
of $\operatorname{Sp}_{2nf}(p)$ to $\operatorname{Sp}_{2n}(q)$
gives the chosen full Weil representation, because both extend
this Heisenberg representation and the latter group is perfect.
This proves the needed restriction assertion rather than
assuming it as an additional functoriality theorem.

Let $\sigma$ be coordinatewise $p$-power Frobenius and set
\[
            t=a^{(p-1)/2},\qquad S=t\sigma.
\]
Then $t^2=a^p/a$, and
\[
 B_a(Sv,Sw)=\operatorname{Tr}(a^p\langle v,w\rangle^p)
          =B_a(v,w).
\]
Hence $S\in\operatorname{Sp}_{2nf}(p)$, and conjugation by $S$
on $\operatorname{Sp}_{2n}(q)$ is the field automorphism
$g\mapsto g^{(p)}$.  It has finite order, also directly from
$S^f=\operatorname{Nm}_{q/p}(t)I=\pm I$.
The Weil operator of $S$ commutes with parity by D-SPW-01,
so every field automorphism preserves each Weil constituent.
Together with the diagonal calculation, this shows that a
normalizing matrix can induce only an inner-field automorphism.
A diagonal-field product cannot restore a constituent that the
diagonal factor exchanges.

Consider a nontrivial field component, represented by the
semilinear map $A=gS^k$ with $0<k<f$ and
$g\in\operatorname{Sp}_{2n}(q)$.
Put $d=\gcd(f,k)$, so $d\le f/2$.
The fixed vectors of $A$ have dimension at most $2n$ over
$\mathbf F_{p^d}$.  To verify this bound, suppose fixed vectors
independent over that fixed field have a minimal nonzero linear
relation over $\mathbf F_q$.  Normalize one coefficient to one,
apply $A$, and subtract the original relation.  Minimality forces
every coefficient to be fixed by $\sigma^k$.  The original
relation is then over $\mathbf F_{p^d}$, a contradiction.
The same argument applies to $-A$.  Their fixed spaces therefore
have $\mathbf F_p$ dimensions at most $2nd$.
Using (SPW.1) over $\mathbf F_p$ and the parity trace formula gives
\[
 \frac{|\operatorname{tr}(U_A|\mathcal W^{\pm})|}{a}
 \le\frac{2p^{nd}}{p^{nf}-1}
 \le\frac{2q^{n/2}}{q^n-1}
 \le\frac{18}{80}<\frac35.                                \tag{SPW.7}
\]
The last inequality uses $q\ge9$, $n\ge2$ and the decreasing
function $2x/(x^2-1)$ for $x\ge9$.
Every other linear intertwiner for this automorphism differs
by a scalar, by irreducibility.  Thus (SPW.5) proves (SPW.2)
for every field coset as well.

Finally, the normalizer modulo scalars embeds in the automorphism
group of the simple projective image: the linear centralizer is
scalar, and a central automorphism of a perfect group is trivial.
The projective kernel of a nontrivial Weil constituent is exactly
the center, since the projective quotient is simple.  The preceding
lift and intertwiner arguments therefore cover all normalizers
of the represented image, including those of central quotients.
A faithful quasisimple representation with the same projective
representation has the same image after the unique lift from
the universal perfect cover; any remaining difference is a scalar
character and is trivial.  This proves the final assertion.
\end{proof}

\subsection*{D-SPW-03. Consequence for actual sixfold Hurwitz sources}

\begin{theorem}[D-SPW-03: symplectic Weil cores reduce to degrees four and five]
\label{D-SPW-03}
Suppose an inducing block of a finite-cover sixfold source has
an irreducibly acting normal central product of quasisimple
groups, each represented by an odd-characteristic symplectic
Weil constituent with $n\ge2$.  Allow all compatible scalar
and outer automorphisms, and arbitrary permutations of the
tensor factors.  If the actual period map has rank nine, then
there is only one factor and
\[
                 (q,n,a)=(3,2,4)\quad\hbox{or}\quad(3,2,5).
                                                               \tag{SPW.8}
\]
In particular every factor with $q\ge5$ or $n\ge3$, and every
tensor product with two or more factors, is impossible.
This assertion does not require a nonsplit target polarization
or equality of the inducing and original character fields.
\end{theorem}

\begin{proof}
Let $e$ be the total inducing degree.  For a factorwise tensor
operator, fixing all eigenvalues in the other factors shows
that its largest eigenspace has dimension at most the product
of their dimensions times the largest eigenspace of any chosen
nonscalar factor.  Thus D-SPW-02 gives $\nu\ge e/5$.
If an operator permutes tensor factors nontrivially, the weighted
cyclic-word proof of D-SPOR-03 gives $\nu\ge e/4$.
That proof only uses equal dimensions along a permutation cycle,
finite order and diagonalizability of the closing twist; it is
independent of the nature of the quasisimple factors.  Hence
every nonscalar element in the full normalizer satisfies
$\nu\ge e/5$.  A nonidentity scalar inertia has moving rank $e$,
so the same lower bound holds for every genuine inertia.

The actual intermediate-cover factorization established in
D-BI-18 applies even when the inducing character field is larger.
Rank nine forces the intermediate base to have genus zero and
at least twelve genuine branch points, while the exact
Chevalley--Weil Euler identity is
\[
             \sum_i\operatorname{rank}(g_i-1)=2e+6.
\]
Consequently
\[
                   12\lceil e/5\rceil\le2e+6.             \tag{SPW.9}
\]
It follows first that $e\le15$, and then exact integer division
leaves, for $e\ge4$, only
\[
                              e\in\{4,5,9,10,15\}.         \tag{SPW.10}
\]
Every factor has degree at least four.  Two factors already
give $e\ge16$, which is impossible.

For the single remaining factor, its degree is $(q^n\pm1)/2$.
The inequality $e\le15$ implies $q^n\le31$.
For odd prime powers $q$ and $n\ge2$, the only possibilities
are $(q,n)=(3,2),(3,3),(5,2)$.
The latter two have degrees $13,14$ and $12,13$, respectively,
none in (SPW.10).  The first gives precisely (SPW.8).
All of these are necessary restrictions on an actual period
map; no dimension equality is used to assert dominance.
\end{proof}

\paragraph{Finite check and remaining cases.}
The separate arithmetic script
\texttt{lie\_small\_fields/check\_symplectic\_weil.py} checks
the rational inequalities in (SPW.6)--(SPW.7), the exceptional
degree ceilings, and the complete finite degree list after the
proved cutoff $e\le15$.  Its actual output is retained.
It is an author certificate, not an independent peer audit.
The surviving single-factor degree-four and degree-five cases
are not asserted impossible by this section; they belong to the
separate low-dimensional inducing analyses.  Non-Weil
representations, characteristic-two symplectic groups, and
other small-field Lie families require separate arguments.

\section{The complete small symplectic core reduction}
\label{sec:sp43-core}

\paragraph{Status.}
This new complete author derivation and its exact certificate have
passed the root reviewer's whole-proof and complete-code review.
The root reviewer also freshly reran the standalone certificate:
all $59$ selected rows and $1992$ nonscalar spectra passed.
That rerun is distinguished from a separate arithmetic implementation.
The imported Fourier routine had already passed the full prior
SPOR and A1 code reviews.  This section concerns every
irreducible representation of the indicated core, not just its Weil
constituents.  Its transitive-factor hypothesis is precisely the
one supplied by D-BI-19; no statement about arbitrary mixed
tensor products is inferred.

\subsection*{D-SPW-04a. Complete scalar-normalizer spectra}

Put $S=\operatorname{PSp}_4(3)\simeq\operatorname{PSU}_4(2)$,
and let $U=\operatorname{Sp}_4(3)$ be its universal perfect cover.
The multiplier and full outer group of $S$ both have order two;
see the published-cover input in D-SPW-02 and the official
CTblLib extension metadata, interpreted as in D-SPOR-01.
The relevant complete character tables have identifiers
\texttt{2.U4(2)} and \texttt{2.U4(2).2}.  Their orders are
$51840$ and $103680$ and their class counts are $34$ and $38$.
Both retained source strings identify the published ATLAS input.

\begin{lemma}[D-SPW-04a: every compatible normalizer]
\label{D-SPW-04a}
Let $\rho$ be any nonlinear complex irreducible representation
of $U$, of degree $a$, and let $T$ be a nonscalar finite-order
linear normalizer of its image.  Then
\[
 \nu(T)\ge
 \begin{cases}
 a/4,&a=4\text{ or }a\ge7,\\
 a/5,&a=5,\\
 a/6,&a=6.
 \end{cases}                                               \tag{SPW.12}
\]
There are no nonlinear degrees below four.  There is exactly
one degree-six core character; it has a trivial central
character and Frobenius--Schur indicator $+1$.

The degree-five quarter-bound expectation is false: both
degree-five characters have, on core class \texttt{2b}, the
involution spectrum $1^1,(-1)^4$, so $\nu=1$.
\end{lemma}

\begin{proof}
The compact exact input retains every character, every class
size and order, every prime power map, all source strings, and
the actual fusion of $U$ into its specified outer extension
$A=\texttt{2.U4(2).2}$.  The GAP exporter checks full row
orthogonality and the sum-of-squared-degrees identity.  These
are consistency checks; the group and character-table
identifications remain the stated published mathematical inputs.

For an element $x$ of order $d$, and each $0\le b<d$, the
exact eigenvalue multiplicity is
\[
 m_b=\frac1d\sum_{j=0}^{d-1}
               \chi(x^j)\zeta_d^{-bj}.                    \tag{SPW.13}
\]
The prime power maps and Galois conjugation determine every
$\chi(x^j)$.  The standalone checker evaluates (SPW.13) by
integer-polynomial reduction modulo $\Phi_d$, checks that
every result is a nonnegative integer, and checks
$\sum_bm_b=a$.  It selects the extension characters again by
pointwise equality of their restriction with a full irreducible
core-character row; thus the selection does not assume
irreducibility merely from degree.

There are $33$ nonlinear core rows and $26$ extension rows
with irreducible nonlinear restriction.  All $1992$ nonscalar
entries of those $59$ rows are checked.  The only entries
violating $\nu\ge a/4$ are
\[
\begin{array}{c|c|c|c|c}
\text{table}&\text{row}&a&\text{class}&\text{multiplicities}\\\hline
\texttt{2.U4(2)}&2,3&5&\texttt{2b},\ d=2&(1,4)\\
\texttt{2.U4(2).2}&4&6&\texttt{4d},\ d=4&(5,0,1,0)\\
\texttt{2.U4(2).2}&5&6&\texttt{4d},\ d=4&(1,0,5,0).
\end{array}                                                \tag{SPW.14}
\]
All multiplicities in (SPW.14) are ordered starting with the
eigenvalue one.  The two extension rows restrict to core row
four, the unique degree-six character.  That core character
has integer values, is trivial on the center, and the exact
sum $|U|^{-1}\sum_x\chi(x^2)$ is one.  This independently
recovers its symmetric invariant form.  The displayed lower
bounds are also checked row by row; in particular no unlisted
degree or quarter-bound exception is omitted.

It remains to pass from a table element to an arbitrary linear
normalizer, without assuming an unproved description of that
normalizer.  The kernel of a nonlinear irreducible
representation of the quasisimple $U$ is central.  An
automorphism of its represented image induces an automorphism
of $S$, and lifts uniquely to the universal perfect cover.
An automorphism trivial on $S$ is trivial on the image: its
central discrepancy is a homomorphism from a perfect group to
an abelian center.  Thus a linear normalizer is, modulo
scalars, inner or induces the one nontrivial outer automorphism.

In the inner case it is a scalar multiple of a represented
element of $U$, covered by the first table.  In the outer case
choose $t\in A\setminus U$ and make an inner adjustment so
that the normalizer $T$ induces precisely conjugation by $t$.
Writing $t^2=u_0\in U$, Schur's lemma gives
$T^2=c\rho(u_0)$.  Rescaling $T$ by a square root of
$c^{-1}$ makes it satisfy the conjugation and power relations
of this fixed $A$.  We obtain an irreducible representation
of $A$ extending $\rho$, hence a row in the checked list.
This argument also covers a nonsplit outer extension and any
represented central quotient.  Scalar rescaling leaves $\nu$
unchanged, proving (SPW.12).

The retained input is \texttt{sp4\_3\_character\_inputs.json},
SHA-256
\texttt{0604b7cbfdd2011c21e46a009674adda4292f1452afbce38ab1593b38c3f388c}.
The complete exporter and checker are
\texttt{export\_sp4\_3.g} and \texttt{check\_sp4\_3.py};
the output \texttt{check\_sp4\_3.json} retains every checked
row and all extremal witnesses.  The false degree-five
quarter-bound expectation is explicitly retained as refuted,
rather than silently discarded.
\end{proof}

\subsection*{D-SPW-04b. The transitive central-product core}

\begin{theorem}[D-SPW-04b: reduction to single degrees four and five]
\label{D-SPW-04b}
Suppose an inducing image $J\subset\operatorname{GL}(V)$
of an actual rank-nine Weil sixfold source has an irreducibly
acting normal central product $N=N_1\cdots N_t$, with all
$N_i/Z(N_i)\simeq S$ and with $J$ transitive on the factors.
Assume its inducing character field contains the original
imaginary quadratic field $K$, as it does for the induced
sources considered here.  Then $t=1$ and
\[
                          \dim V\in\{4,5,6\}.             \tag{SPW.15}
\]
Every tensor case $t\ge2$ is excluded, including degree-six
factors.

If $K$ is one of the eight non-Eisenstein fields in the main
nonsplit problem and the actual original polarization is
nonsplit, the single degree-six case is also excluded by
D-BI-16.  The only remaining possibilities are therefore the
single degree-four and degree-five cases, already subject to
the $G_{32}$ and $G_{33}$ low-degree source restrictions.
No existence or dominance assertion is made for those endpoints.
\end{theorem}

\begin{proof}
As an $N$-module,
$V=V_1\otimes\cdots\otimes V_t$, with irreducible nonlinear
factor modules.  Transitivity of $J$ implies that all factor
degrees have the same value $a\ge4$.  Put $e=a^t$.
The tensor matrix algebras of the intrinsic simple factors
are permuted by $J$.  Hence each normalizing operator is a
tensor product of intertwiners followed by a factor permutation.
For a factor-preserving operator, the largest eigenspace
fraction is at most the largest such fraction of any of its
nonscalar factors.  A nontrivial weighted factor cycle has
largest eigenspace fraction at most $3/4$, by the exact
constant-word argument in D-SPOR-03.  These statements allow
arbitrary scalar phases and compatible outer automorphisms.

The actual intermediate-cover Euler identity and period
factorization give
\[
             h=0,\qquad r\ge12,\qquad D=2e+6,              \tag{SPW.16}
\]
where $r$ counts genuine branch points.  Indeed at genus one
there are at most six genuine branches and at most six
marked-base parameters; at genus at least two the identity
$D=6+2e(1-h)$ is negative since $e\ge4$.  Genuine scalar
inertia has moving rank $e$, so all bounds below apply to it
as well.

If $a=4$ or $a\ge7$, (SPW.12) and the tensor argument give
$\nu(g)\ge e/4$ for every nonscalar inertia.  Equation
(SPW.16) implies $3e\le2e+6$, so $e\le6$.  Thus $a\ge7$
is impossible, and $a=4$ forces $t=1$.
If $a=5$, the corresponding fifth bound gives
$12e/5\le2e+6$, hence $e\le15$.  Therefore $t=1$ here too.

It remains to exclude $a=6$, $t\ge2$; the sixth bound alone
would not do so.  The unique degree-six core character has
indicator $+1$.  Taking the tensor product of its symmetric
forms gives a nondegenerate $N$-invariant symmetric form $B$
on $V$.  Irreducibility makes its invariant line unique, so
normality of $N$ gives a multiplier character
\[
                  B(gv,gw)=\mu(g)B(v,w),\qquad g\in J.
\]
This conclusion does not require a chosen rational model or
an asserted full normalizer identification.

The multiplier is nontrivial.  If it were trivial, every
finite-order element of $J$ would have eigenvalues paired
as $\lambda,\lambda^{-1}$ by $B$; these are complex
conjugates.  Its trace would be real, making the inducing
character field real, contrary to its containing the given
imaginary quadratic $K$.
Because the inertia tuple generates $J$ and has product one,
at least two inertia multipliers are nontrivial.  For each
such inertia $g$, its fixed space is totally isotropic for
$B$, since
$(\mu(g)-1)B(v,w)=0$ on that space.  Its moving rank is
therefore at least $e/2$.
Every other inertia has moving rank at least $e/6$ by
(SPW.12) and the tensor-cycle bound.  Consequently
\[
             2e+6=D\ge 2(e/2)+10(e/6)=8e/3,
             \qquad e\le9.                               \tag{SPW.17}
\]
This contradicts $e=6^t\ge36$ when $t\ge2$.
Only the single degree-six representation remains, proving
(SPW.15).

Under the final nonsplit, non-Eisenstein hypotheses, D-BI-16
applies to any degree-six inducing representation, independently
of its inducing field.  Its sole allowed faithful image is
$\mu_4G_{34}$.  This has projective order
$39191040/6=6531840$, whereas the present normalizer has
projective order at most $|\operatorname{Aut}(S)|=51840$.
Thus the degree-six case here cannot be that allowed image.
Equivalently, it is the already excluded real $E_6$ case.
The degree-four and degree-five conclusions refer back to
their previously proved low-degree restrictions and assert no
additional normalizer identification.  This completes the
stated reduction.
\end{proof}

\section{Exclusion of every $G_2$ and triality normal core}
\label{sec:g2-triality-cores}

\paragraph{Status and input boundary.}
This new author derivation has passed an independent whole-proof and
source audit by \texttt{/root/g34\_eightfolds}.  That reviewer also
reconstructed every original generic entry and every small-cover spectrum
with separate exact arithmetic implementations; both passed.  The
root reviewer has separately passed the complete mathematical and code
review and rerun every certificate; its remaining source read is
recorded separately.  The complete published
character tables, their official generic-table implementation, and the
published list of exceptional Schur multipliers are mathematical inputs.
The new bounds below are proved from those inputs by rational polynomial
inequalities and exact cyclic Fourier calculations.  In particular, the
assertions for unbounded field size do not extrapolate a finite scan.
The source log distinguishes original papers actually retrieved from
published inputs accessed through their official mathematical tables.

As before, for a finite-order operator on an $a$-dimensional complex space,
write
\[
 \nu(T)=a-\max_{\lambda}\dim\ker(T-\lambda).
\]
Multiplication by a scalar does not change $\nu$.  A genuine nonidentity
scalar inertia has ordinary moving rank $a$.

\subsection*{D-G2T-01. Uniform bounds from the complete generic tables}

\begin{lemma}[D-G2T-01: all-field character bounds]
\label{D-G2T-01}
Let $S$ be $G_2(q)$ with $q\ge5$, or ${}^3D_4(q)$ with $q\ge3$.
For every nontrivial ordinary irreducible character $\chi$ of $S$ and
every $s\ne1$,
\[
                  |\chi(s)|\le\tfrac14\chi(1).             \tag{G2T.1}
\]
Moreover $\chi(1)\ge124$ in the first family and
$\chi(1)\ge219$ in the second.
\end{lemma}

\begin{proof}
The complete ordinary $G_2$ tables are the published inputs of
Enomoto--Yamada in characteristic two, Enomoto in characteristic three,
and Chang--Ree in the other characteristics
\cite{G2T:EY,G2T:Enomoto,G2T:ChangRee}.
For triality we use Deriziotis--Michler's complete table
\cite[\S4, Theorem~4.3 and Tables~3.6, 4.4]{G2T:DM}.
The official GenericCharacterTables implementation, descended from
CHEVIE, records these exact sources and supplies all the mathematical
entries \cite{G2T:GCT,G2T:CHEVIE}.
We freeze its source at commit
\texttt{87231738271ed548f62e7f1c98ff2bbfc1ed30d3}.

The following inventory is exhaustive.  The integer in the fourth
column is the number of character types and also the number of
class types; a type may contain several parameterized characters
or classes.  The last column is a proved lower bound for every
nontrivial character degree throughout the stated real range.
\[
\begin{array}{c|l|r|r|r}
\text{table}&\text{field condition}&q_0&\text{types}&\text{degree bound}\\\hline
\texttt{G2.01}&q\text{ even},\ q\equiv1\ (3)&5&27&126\\
\texttt{G2.02}&q\text{ even},\ q\equiv2\ (3)&5&27&124\\
\texttt{G2.101}&q=3^f,\ q\equiv1\ (4)&5&28&280\\
\texttt{G2.103}&q=3^f,\ q\equiv3\ (4)&5&28&280\\
\texttt{G2.111}&q\text{ odd},\ q\equiv1\ (3),\ q\equiv1\ (4)&5&31&126\\
\texttt{G2.113}&q\text{ odd},\ q\equiv1\ (3),\ q\equiv3\ (4)&5&31&126\\
\texttt{G2.121}&q\text{ odd},\ q\equiv2\ (3),\ q\equiv1\ (4)&5&31&124\\
\texttt{G2.123}&q\text{ odd},\ q\equiv2\ (3),\ q\equiv3\ (4)&5&31&124\\
\texttt{3D4.0}&q\text{ even}&3&29&219\\
\texttt{3D4.11}&q\text{ odd},\ q\equiv1\ (4)&3&33&219\\
\texttt{3D4.13}&q\text{ odd},\ q\equiv3\ (4)&3&33&219
\end{array}                                                    \tag{G2T.2}
\]
For each type, all admissible values of every class and character
parameter are retained in the published interpretation.  Our estimates
are independent of their values, as follows.

Every table entry is a finite expression
\[
                z(q)=\sum_\alpha p_\alpha(q)u_\alpha,
                \qquad p_\alpha(q)\in\mathbf Q[q],
                \qquad |u_\alpha|=1.                       \tag{G2T.3}
\]
Here each $u_\alpha$ is a product of the roots of unity occurring
in that entry.  Equal formal products may be collected, while
relations between different products need not be used.  The
original generic cyclotomics have exponents that are rational
expressions in the real integral parameters, so each phase
indeed has absolute value one for every admissible specialization.

Put $x=q-q_0\ge0$ and expand exactly
\[
 p_\alpha(q_0+x)=\sum_k c_{\alpha k}x^k,
 \qquad d(q_0+x)=\sum_k d_kx^k,
\]
where $d$ is the degree polynomial of the row.  The triangle inequality
gives the uniform bound
\[
                 |z(q)|\le
                 \sum_k\left(\sum_\alpha|c_{\alpha k}|\right)x^k.
                                                               \tag{G2T.4}
\]
Thus it suffices to check, for every nontrivial row and every
nonidentity column,
\[
                  d_k-4\sum_\alpha|c_{\alpha k}|\ge0
                         \quad\hbox{for every }k.           \tag{G2T.5}
\]
Every one of these rational inequalities is nonnegative in all
eleven tables.  Also every $d_k$ is nonnegative, and $d_0$ is at
least the number recorded in (G2T.2).  This proves the claimed
uniform degree bounds as well as (G2T.1).

Here are the complete finite verification details.
\texttt{build\_generic\_input.py} reads the eleven frozen source
files, checks their separately recorded SHA-256 hashes, and
extracts only the original character entries, declared degrees,
and declared class lengths.  It does not execute downloaded
Julia code.  The resulting offline mathematical input
\texttt{generic\_character\_entries.json} has $415124$ bytes and
SHA-256
\texttt{66ad9b13d457d141d424084336787b3bb1f9d53818475b806f7a6cc24eb3664b}.

\texttt{parse\_generic\_bounds.py} accepts only integer constants,
the indeterminate $q$, rational constant division, addition,
multiplication, nonnegative integer powers, and formal unit phases.
It checks that the trivial row consists of ones, that every first
column entry equals the separately declared degree, and that the
first class length is one while all other class lengths are
strictly greater than one throughout the relevant range.  The
last assertions follow from nonnegative translated coefficients
and a positive constant term after subtracting one.
It then constructs and checks (G2T.5) for all $318$ nontrivial
row types and all $9242$ nontrivial-row, nonidentity-column
entries.  All checks pass.

The complete rational coefficient lists, excluding run-specific timing
data, are retained in
\texttt{generic\_polynomial\_bounds.json}; their SHA-256 is
\texttt{568425a2cf62e62dbe24709a3372d9bd182946e0a53ae74e7fed45ba640dcd13}.
A second, short program, \texttt{check\_generic\_polynomials.py},
uses only Python's exact rational arithmetic to recheck every
shifted coefficient inequality and every minimum degree from
these lists.  It also passes.  This second program validates
the polynomial certificate; the first program separately
verifies its correspondence with the retained original table
expressions.  Neither is offered as a new proof of the published
character-table classification itself.
\end{proof}

\subsection*{D-G2T-02. Exceptional covers and all outer normalizers}

\begin{lemma}[D-G2T-02: normalizer bounds]
\label{D-G2T-02}
Let $N$ be a finite quasisimple group acting faithfully and
irreducibly in degree $a$, whose simple quotient is $G_2(q)$,
$q\ge3$, or ${}^3D_4(q)$, $q\ge2$.
Every nonscalar finite-order element of its complex linear
normalizer satisfies
\[
                         \nu(T)\ge3a/16.                  \tag{G2T.6}
\]
Every such $a$ is at least twelve.  If the simple quotient is
$G_2(3)$, $G_2(4)$, or ${}^3D_4(2)$, the stronger bound
$\nu(T)\ge a/4$ holds and the respective minimum degrees are
$14,12,26$.  In the other cases the degree bounds in
D-G2T-01 apply.
\end{lemma}

\begin{proof}
The generic simply connected groups of these two types have
trivial center.  The only exceptional multipliers in these
families are those of $G_2(3)$ and $G_2(4)$; the complete
published exceptional list is \cite[Chapter~5, Table~1,
numbered p.~35]{BI:HHM}.  In particular there is no omitted
triality central cover.  For the three small simple groups,
the universal covers and their full cyclic outer extensions
are the following complete official CTblLib inputs:
\[
\begin{array}{c|r|r|r|r|c}
\text{table}&\text{classes}&\text{selected rows}&
 \text{nonscalar spectra}&\min a&\min\nu/a\\\hline
\texttt{3.G2(3)}&45&44&1848&14&3/7\\
\texttt{3.G2(3).2}&39&20&740&14&3/7\\
\texttt{2.G2(4)}&55&54&2862&12&1/3\\
\texttt{2.G2(4).2}&65&48&3024&12&1/3\\
\texttt{3D4(2)}&35&34&1156&26&4/13\\
\texttt{3D4(2).3}&49&39&1872&26&4/13
\end{array}                                                    \tag{G2T.7}
\]
The extension metadata record respectively multiplier and outer
orders $(3,2),(2,2),(1,3)$, in the sense of CTblLib
\S3.7-3 \cite{SPOR:CTblLib}.  Every source string records
the ATLAS; the group identities and complete tables are
published mathematical inputs.

The GAP exporter retains every character, class size, element
order, prime power map, source string, and actual core-to-extension
class fusion.  It verifies full character-row orthogonality.
The standalone verifier selects every nonlinear core character
and every extension character whose restriction equals a
nonlinear irreducible core character, by exact equality of
all its character values.  It checks fusion orders, trivial
identity kernel, normal-class size unions, and compatibility
with all retained prime power maps.

For an element $s$ of order $o$, every eigenvalue multiplicity is
computed by the exact Fourier formula
\[
       m_b=\frac1o\sum_{j=0}^{o-1}\chi(s^j)\zeta_o^{-bj},
                      \quad 0\le b<o.                    \tag{G2T.8}
\]
Prime power maps compute every power class.  Cyclotomic
polynomials and integer coefficient lists perform all arithmetic;
every $m_b$ is checked to be a nonnegative integer and their
sum is checked to equal $a$.  Scalar classes are recognized
by $|\chi(s)|^2=a^2$, which is equivalent to scalar action for
a finite-order unitary operator.  The full computation comprises
$239$ selected rows and $11502$ nonscalar spectra, with precisely
the minima in (G2T.7).  It passes, and every ratio is at least
one quarter.

The input \texttt{g2\_triality\_small\_inputs.json} has SHA-256
\texttt{e4a8aefbf5a4a32a5b31b217143ccae03db8b9f3753cc2e91a6169d0364b6a7c}.
The author programs are \texttt{export\_g2\_triality\_small.g}
and \texttt{check\_g2\_triality\_small.py}.  The latter reuses
the fully reviewed exact cyclotomic Fourier helper from D-SPOR.

We explain why these six fixed tables cover arbitrary compatible
small-group normalizers, including scalar changes in a power
relation.  Lift the representation to the universal perfect cover
$U$.  Every automorphism of its simple quotient has a unique
lift to $U$.  An inner normalizer is scalar times a represented
element of $U$.  For an outer normalizer, choose a matching
generating coset in the recorded extension $A=U.r$,
where $r=2$ or $3$.  After an inner adjustment, its action on
$U$ agrees with that of an element $t\in A$.  Write $t^r=u_0$.
Schur's lemma gives $T^r=c\rho(u_0)$ for a nonzero scalar $c$.
After multiplying $T$ by an $r$th root of $c^{-1}$, its
conjugation and power relations give a representation of this
fixed $A$ extending $\rho$.  It is irreducible because its
restriction to $U$ is irreducible.  Thus its row is among the
tested rows.  Nonfaithful central kernels cause no loss: all
nonlinear rows of the universal cover were retained.  This proves
the small-group quarter bound with no presumption that every
outer automorphism preserves every central kernel.

For all other field sizes, D-G2T-01 gives the character ratio
$1/4$ on the centerless core.  The complex commutator argument
of D-BI-18 gives (G2T.6) for every normalizer, not merely for
inner elements.  Explicitly, if $T$ is nonscalar, there is
$u\in N$ with $[T,u]\ne1$, by irreducibility.  The rank
inequality
\[
                \operatorname{rank}([T,u]-1)\le2\nu(T)
\]
holds after subtracting a largest-eigenspace scalar from $T$.
If $[T,u]$ is nonscalar, its unitary trace and (G2T.1) imply
\[
 a-2\operatorname{rank}([T,u]-1)
       \le\operatorname{Re}\chi([T,u])\le a/4,
\]
so $\nu(T)\ge3a/16$.  A nonidentity scalar commutator instead
has full moving rank $a$, giving the stronger $\nu(T)\ge a/2$.
The smallest degrees are those already certified in (G2T.2)
and (G2T.7), proving every assertion.
\end{proof}

\subsection*{D-G2T-03. The whole finite-cover source exclusion}

\begin{theorem}[D-G2T-03: $G_2$ and triality core exclusion]
\label{D-G2T-03}
Let an inducing block of a finite-cover Weil sixfold source have
finite image $J\subset\operatorname{GL}(V)$.  Suppose $J$ has
an irreducibly acting normal central product of quasisimple
groups whose simple quotients belong to
\[
                G_2(q)\ (q\ge3),\qquad
                {}^3D_4(q)\ (q\ge2).
\]
Then the actual period rank of the source is at most eight.

Mixed factor types, arbitrary scalar phases, compatible outer
automorphisms, and factor permutations are allowed.  No
primitivity, tensor-indecomposability, essential-self-duality,
equality of inducing and original fields, or nonsplit-polarization
hypothesis is required.
\end{theorem}

\begin{proof}
Write the normal central product as $N=N_1\cdots N_t$.
Its irreducible representation has the form
$V=V_1\otimes\cdots\otimes V_t$, with faithful irreducible
factor representations of dimensions $a_i\ge12$.
Set $e=\dim V=\prod_i a_i$.  The simple direct factors of
$N/Z(N)$ are intrinsic.  Each $N_i$ is the derived subgroup
of the inverse image of its simple factor, so $J$ permutes
the $N_i$.  Their generated full matrix algebras
$\operatorname{End}(V_i)\otimes1$ are consequently permuted
as well.  Every normalizing operator therefore consists of
factor intertwiners and a permutation of the appropriate
isomorphic factor spaces.

If the operator preserves all factors and is nonscalar,
one factor intertwiner is nonscalar.  Fixing all other
eigenvalues shows that the largest eigenspace fraction of
the tensor operator is no larger than the largest such
fraction in that factor.  Formula (G2T.6) therefore gives
$\nu(g)\ge3e/16$.

For a nontrivial permutation cycle of length $l\ge2$ on
factors of equal dimension $a$, tensor changes of basis
put its action in the form of cyclic word rotation with
a single diagonal twist.  The product twist is diagonalizable:
the finite order of the whole operator implies finite order
up to scalar for its powers on each tensor cycle.
Each orbit of the resulting basis words contributes at most
one vector to a specified eigenspace.  There are $a$ constant
words, and every other orbit has size at least two.  Hence
the largest eigenspace dimension is at most
\[
                    a+(a^l-a)/2\le3a^l/4.
\]
Tensoring with the other cycles preserves this bound, so
$\nu(g)\ge e/4\ge3e/16$.  This is the weighted cyclic-word
argument of D-SPOR-03 and does not assume that the twists
have distinct eigenvalues.  Nonidentity scalar inertia has
ordinary moving rank $e$.  Thus every genuine inertia has
ordinary moving rank at least $3e/16$.

Suppose the actual period rank were nine.  Apply the proved
intermediate-cover factorization D-BI-01 to this inducing
block.  It remains valid when the inducing character field
is larger than the original quadratic field: the first
complex Hodge filtration factors through that intermediate
cover, and the conjugate filtration is its Hermitian
annihilator.  If the intermediate base has genus $h$ and
$R$ genuine branch points, the multiplicity-six Euler formula is
\[
                      D=6+2e(1-h),\qquad R\le D.           \tag{G2T.9}
\]
For $h\ge2$, its right side is negative because $e\ge12$.
For $h=1$, the actual period map factors through a finite
Hurwitz cover of dimension at most six.  Consequently rank
nine requires $h=0$ and $R\ge12$, and then $D=2e+6$.
The inertia estimate proves the necessary inequality
\[
              12\left\lceil\frac{3e}{16}\right\rceil
                   \le2e+6,
                   \qquad e\le24.                        \tag{G2T.10}
\]
This is an obstruction obtained from factorization of the
actual variation, not a sufficiency inference from a source
dimension count.

If $t\ge2$, then $e\ge12^2=144$, contradicting (G2T.10).
For a single factor outside the three small cases,
D-G2T-01 gives $e\ge124$ or $e\ge219$, also impossible.
For a single factor in those cases, the quarter bound of
D-G2T-02 gives $3e\le2e+6$, hence $e\le6$, contrary to
the respective minima $14,12,26$.  Every case is excluded.
Since the target Weil period domain has dimension nine,
the actual period rank is at most eight.
\end{proof}

\paragraph{Boundary of the result.}
The nonsimple group $G_2(2)$ is not included.  Its derived
simple group is $U_3(3)$, as recorded in
\cite[Chapter~5, numbered p.~35]{BI:HHM}, and belongs to a
separate unitary-group gate.  No assertion about the other
classical or exceptional families is inferred from this
calculation.

\section{Complete even-characteristic symplectic cores in ranks two and three}
\label{sec:even-symplectic-low-rank}

\paragraph{Status and scope.}
The author has reviewed the complete derivation and executed all exact
certificates.  A separate whole-proof and primary-source peer review
has passed, together with an independent reconstruction of all
83 generic rows, 4,549 inequalities, 104 small-table rows and
3,442 spectra.  This section uses every ordinary irreducible character
in the two indicated symplectic ranks.  The complete official generic
tables are explicit mathematical inputs, including the correction
history recorded by their authors.  No claim about higher symplectic
rank, odd characteristic, or a partial unipotent table is made.

We retain $\nu(T)=a-\max_\lambda\dim\ker(T-\lambda)$ for an
$a$-dimensional finite-order complex operator.

\subsection*{D-ESC-01. Two complete generic character bounds}

\begin{lemma}[D-ESC-01: complete generic bounds]
\label{D-ESC-01}
For every nontrivial ordinary irreducible character of the indicated
group and every nonidentity element,
\[
\begin{array}{c|c|c}
\text{group and range}&|\chi(g)|/\chi(1)&\min\chi(1)\\\hline
\operatorname{Sp}_4(q),\ q\ge8\text{ even}&\le1/4&\ge196\\
\operatorname{Sp}_6(q),\ q\ge4\text{ even}&\le17/63&\ge378.
\end{array}                                                   \tag{ESC.1}
\]
Consequently every nonscalar finite-order linear normalizer of
either represented core has
\[
                           \nu(T)\ge23a/126.              \tag{ESC.2}
\]
\end{lemma}

\begin{proof}
The first complete table originates in Enomoto
\cite[\S4, pp.~80--85, and Tables~IV-1/IV-2,
pp.~92--94]{ESC:Enomoto}.  We use the official complete
\texttt{Sp4.0} implementation, not an unchecked transcription
of the original printed entries.  The official documentation
explicitly states that the printed paper has typographical
errors and that F.~L\"ubeck recomputed the implemented table
using Deligne--Lusztig theory \cite{G2T:GCT}.
That correction history is part of the input's provenance.

The second complete table is L\"ubeck's \texttt{Sp6.0}.
His full author text explicitly states the complete ordinary
character-table result, distinguishes it from the unipotent
subtable, and explains that the complete values were supplied
as CHEVIE computer files rather than printed in full
\cite[Introduction, pp.~3--4; \S2.2, pp.~18--21;
\S9, pp.~78--79]{ESC:Luebeck}.
The construction uses all Lusztig series and the additional
nonuniform class functions, so its scope is not merely the
first twelve unipotent character types.  The official
implementation is the entire \texttt{Sp6.0} table in
\cite{G2T:GCT}, at the same frozen commit as D-G2T-01.

Here is a finite exact proof of the new bounds from these
complete mathematical inputs.  As in D-G2T-01, write each
original entry as $\sum_\alpha p_\alpha(q)u_\alpha$ with
$p_\alpha\in\mathbf Q[q]$ and $|u_\alpha|=1$.
For each table put $x=q-q_0\ge0$ and write
\[
 p_\alpha(q_0+x)=\sum_k c_{\alpha k}x^k,
 \qquad d(q_0+x)=\sum_k d_kx^k.
\]
The same coefficientwise triangle inequality proves
$|\chi(g)|\le c\chi(1)$ as soon as
\[
                 c d_k-\sum_\alpha|c_{\alpha k}|\ge0
                      \quad\hbox{for every }k.             \tag{ESC.3}
\]
All of (ESC.3) pass with the following exact input parameters:
\[
\begin{array}{c|r|c|r|r|r}
\text{table}&q_0&c&\text{character/class types}&
\text{nontrivial cells}&\min d(q_0)\\\hline
\texttt{Sp4.0}&8&1/4&19&324&196\\
\texttt{Sp6.0}&4&17/63&66&4225&378.
\end{array}                                                   \tag{ESC.4}
\]
Every translated degree polynomial has nonnegative coefficients.
Thus its constant term gives the uniform degree lower bound.
Every nonidentity class-length polynomial likewise has
nonnegative translated coefficients and constant term greater
than one.  The trivial row and identity column are checked
against the separately declared degrees and class lengths.

The program \texttt{build\_even\_symplectic.py} checks the two
frozen source hashes and extracts the original mathematical
entry strings.  It executes no downloaded Julia code.
The offline original input has $436290$ bytes and SHA-256
\texttt{49ecbe27ad97b8cfe0204f96ef463ef1e8e0626ce6987bcccb0d8b33041e4fa0}.
The safe arithmetic parser already proved and reviewed in
D-G2T-01 constructs all $83$ nontrivial character-type rows
and $4549$ inequalities.  The deterministic coefficient
certificate has SHA-256
\texttt{92e32ae7b01d49469f31464472ae044b320b59a6400fedaec5003fde2760a3b3}.
\texttt{check\_even\_symplectic\_generic.py} reparses the
retained original input, matches every coefficient list, and
rechecks (ESC.3) using exact rational arithmetic.  All checks
pass.  The estimates are uniform in every admissible class
and character parameter, since only the absolute value one
of each phase is used.

For even $q$ both generic groups have trivial center, and the
published complete exceptional-multiplier list contains no
exception in these ranges \cite[Chapter~5, Table~1,
numbered p.~35]{BI:HHM}.  Thus a nontrivial irreducible
representation is faithful on its simple core.  Applying
the complex commutator estimate D-BI-18 to $c=17/63$ gives
\[
                       (1-c)/4=23/126.
\]
The first row has the stronger $3/16$ estimate, which is at
least $23/126$.  The commutator proof includes all normalizing
outer automorphisms and nonidentity scalar commutators, so
no outer-table assumption enters (ESC.2).
\end{proof}

\subsection*{D-ESC-02. The complete small-field endpoints}

\begin{lemma}[D-ESC-02: exact small-field normalizers]
\label{D-ESC-02}
In every nontrivial irreducible representation of
$\operatorname{Sp}_4(4)$, every nonscalar finite-order linear
normalizer satisfies $\nu(T)\ge a/3$.  Every such degree is
at least eighteen.

For every nonlinear irreducible representation of the universal
cover $2.\operatorname{Sp}_6(2)$, including those factoring
through its center, the following statements hold for every
nonscalar normalizer:
\begin{enumerate}
\item the unique degree-seven row has $\nu(T)\ge a/7$ and
has a real form with invariant symmetric bilinear form;
\item every other row has $\nu(T)\ge2a/9$;
\item among the nonlinear rows of degree at most nine, the
only degrees are seven and eight; the degree-eight row has
the stronger $\nu(T)\ge a/4$.
\end{enumerate}
\end{lemma}

\begin{proof}
We retain these four complete published ATLAS tables in CTblLib:
\[
\begin{array}{c|r|r|r}
\text{table}&\text{classes}&\text{selected rows}&
\text{nonscalar spectra}\\\hline
\texttt{S4(4)}&27&26&676\\
\texttt{S4(4).2}&30&20&580\\
\texttt{S4(4).4}&30&16&464\\
\texttt{2.S6(2)}&43&42&1722.
\end{array}                                                   \tag{ESC.5}
\]
The complete original input contains every character, class
size/order, prime power map, Frobenius--Schur indicator, and
actual fusion from the core to each extension.  Its size is
$19056$ bytes and SHA-256 is
\texttt{b51fae6eaa904bb5a1bdcac2f6aa85a940c9968ddba912bf115faa750e772fe2}.
The metadata record trivial multiplier and cyclic outer group
of order four for $\operatorname{Sp}_4(4)$, and multiplier
of order two with trivial outer group for
$\operatorname{Sp}_6(2)$.  The complete exceptional-multiplier
list agrees \cite[Chapter~5, Table~1, p.~35]{BI:HHM}.
The \texttt{S4(4).4} record additionally identifies its
construction as the full automorphism group.  These are
group-theoretic mathematical inputs in the standard CTblLib
metadata convention, not conclusions from numerical
orthogonality alone.

The author GAP exporter checks full row orthogonality and
every cyclic spectrum.  The separate polynomial Fourier
verifier \texttt{check\_even\_symplectic\_small.py} selects
irreducible extension restrictions anew by exact character
equality, verifies both normal fusions and cyclic quotients,
and reconstructs every eigenvalue multiplicity using (G2T.8).
Every multiplicity is a nonnegative integer and their sum
is the degree.  All $104$ selected rows and $3442$ spectra
pass.  The three $\operatorname{Sp}_4(4)$ tables have minimum
ratio $1/3$ and minimum degree eighteen.

In the universal $\operatorname{Sp}_6(2)$ table the only
quarter-bound failures are exactly
\[
\begin{array}{r|r|c}
\text{row}&a&\min\nu/a\\\hline
2&7&1/7\\
5&21&5/21\\
6&27&2/9.
\end{array}                                                   \tag{ESC.6}
\]
The other rows satisfy the quarter bound, so all non-seven
rows satisfy $2/9$.  The degree-seven character takes rational
integer values; the verifier recomputes its indicator directly:
\[
             \frac1{|U|}\sum_{g\in U}\chi(g^2)=1,
                     \qquad U=2.\operatorname{Sp}_6(2).
\]
The standard Frobenius--Schur criterion supplies its invariant
symmetric bilinear form.  The complete degree list has exactly
the two asserted nonlinear degrees at most nine.

For $\operatorname{Sp}_4(4)$ the cyclic subgroups of the outer
group have orders one, two and four, all represented in
(ESC.5).  Given an actual normalizer, adjust its automorphism
by an inner action and choose a corresponding element $t$
of the appropriate fixed extension $A$.  If $A/U$ has order
$d$, Schur's lemma gives $T^d=c\rho(t^d)$.  Rescaling $T$
by a $d$th root of $c^{-1}$ gives an honest representation
of that fixed extension, irreducible on its core.  The
original operator is scalar times an element of a tested
row.  Thus no splitting or implicit choice of scalar power
relation is assumed.  For $\operatorname{Sp}_6(2)$ there is
no outer action.  Uniqueness of automorphism lifting to its
universal perfect cover implies that every normalizer is
scalar times a core element.  Retaining every nonlinear
universal-cover row includes every possible central kernel.
The stated finite bounds therefore hold for all actual
normalizers.
\end{proof}

\subsection*{D-ESC-03. Exact source consequences}

\begin{theorem}[D-ESC-03: characteristic-two rank-two/rank-three cores]
\label{D-ESC-03}
Let $J$ be the finite image of an inducing block in a finite-cover
Weil-sixfold source.
\begin{enumerate}
\item If $J$ has an irreducibly acting normal central product
of quasisimple factors with simple quotients
$\operatorname{Sp}_4(q)$ or $\operatorname{Sp}_6(q)$,
where all $q\ge4$ are even, its actual period rank is at
most eight.  Mixed ranks and field sizes are allowed.
\item Suppose instead that the simple quotients of all
factors are $\operatorname{Sp}_6(2)$ and $J$ permutes the
factors transitively.  If the source has period rank nine,
there is exactly one factor and the inducing degree is seven.
\item In the nonsplit problem over any of the eight
non-Eisenstein imaginary quadratic fields of class number
one, that last degree-seven case is excluded by D-BI-17.
Consequently no transitive normal core of either of these
two symplectic ranks in characteristic two remains in that
problem.
\end{enumerate}
All parts allow arbitrary scalar phases, compatible outer
automorphisms and tensor permutations, and a larger inducing
character field.  Parts (1)--(2) impose no nonsplit determinant
condition.  Part (2) is an exact remaining-case statement,
not an existence assertion for its endpoint.
\end{theorem}

\begin{proof}
For part (1), write the irreducible core representation as
$V=\bigotimes V_i$, with dimensions $a_i$, and put
$e=\prod a_i$.  The intrinsic simple factors and their full
matrix algebras are permuted by $J$.  If an operator preserves
all factors and is nonscalar, a nonscalar factor intertwiner
controls its largest eigenspace fraction.  D-ESC-01 and
D-ESC-02 give $\nu(T)\ge23e/126$.  The minimum factor degree
in the combined list is eighteen.

A nontrivial cycle of $l\ge2$ tensor factors of dimension
$a$ is conjugate to cyclic word rotation with a diagonal
twist at one position.  Its constant basis words number
$a$, and every other orbit has at least two words.  A fixed
eigenvalue occurs at most once on each weighted orbit, so
its eigenspace has dimension at most
$a+(a^l-a)/2\le3a^l/4$.  Tensoring with the other cycles
preserves this fraction.  Hence factor permutations give
the stronger $\nu(T)\ge e/4$.  Genuine nonidentity scalar
inertia has full moving rank $e$.

If the actual period rank were nine, the intermediate-cover
factorization and Euler argument of D-BI-01 give genus zero,
at least twelve genuine inertias, and total moving rank
$2e+6$.  Indeed genus one has at most six marked-base
parameters and genus at least two would make
$6+2e(1-h)$ negative.  Therefore
\[
                  12\frac{23e}{126}\le2e+6,
                        \qquad e\le63/2.                 \tag{ESC.7}
\]
For at least two factors, $e\ge18^2>63/2$.
For a single generic factor the minimum degrees $196,378$
also violate (ESC.7).  A single $\operatorname{Sp}_4(4)$
factor has the stronger one-third bound, giving
$4e\le2e+6$, or $e\le3<18$.  This proves part (1).

For part (2), transitivity makes all factor dimensions equal,
say $a$, and $e=a^t$.  First take $a\ne7$.
The factor-preserving bound is $2e/9$ and the factor-permuting
bound is $e/4\ge2e/9$.  The same actual branch budget gives
\[
                        12\frac{2e}{9}\le2e+6,
                              \qquad e\le9.               \tag{ESC.8}
\]
The full degree list in D-ESC-02 forces $t=1$ and $a=8$.
Its stronger quarter bound instead gives
$12\lceil8/4\rceil=24>22=2e+6$, a contradiction.

Now let $a=7$.  Every factor has an invariant symmetric
bilinear form, so their tensor product is symmetric and
spans the unique invariant bilinear-form line of the
irreducible normal core.  The full normalizer preserves
that line, with a similitude character $\mu$.
For the original quadratic field $K$ and the inducing
character field $L$, induction gives $K\subset L$.
If $\mu$ were trivial, every finite-order matrix would
preserve a symmetric bilinear form exactly; its eigenvalues
would pair with their inverses, and its trace would be real.
This would contradict the imaginary quadratic subfield of
$L$.  Hence $\mu$ is nontrivial.

The genuine inertias generate and have product one, so at
least two have nontrivial multiplier.  Their fixed spaces
are totally isotropic, and each moves at least $e/2$
dimensions.  All other genuine inertias move at least
$e/7$ by D-ESC-02 and the tensor-cycle bound.  Thus
\[
                    2e+6\ge e+10e/7,
                              \qquad e\le14.              \tag{ESC.9}
\]
Since $e=7^t$, this forces $t=1$ and $e=7$, proving part (2).
Part (3) is now precisely the previously proved D-BI-17
exclusion of every degree-seven inducing block in its
stated eight-field nonsplit scope.  It does not require
identifying a new polarization or assuming equality $K=L$.
\end{proof}

\paragraph{Boundary and unsuccessful stronger estimates.}
The nonsimple group $\operatorname{Sp}_4(2)$ is not included;
its derived simple group is $A_6$ and belongs to D-BI-21.
The initial attempted quarter-character bounds at threshold
$q=4$ failed at six cells for \texttt{Sp4.0} and three cells
for \texttt{Sp6.0}.  Those failures are retained in
\texttt{even\_symplectic\_exploration.json}.  The proved
statement uses the different threshold and exact ratio in
(ESC.4), together with complete small-group spectra.
The degree-seven exception is retained in part (2), and
is removed only under the precise earlier theorem's scope.

% Author: /root/g34_eightfolds, GPT-6 Astra Pro; inherited/unexposed effort.
% Run: HODGE_DIVISION, 2026-11-21.
% Status: author whole-proof read; independent whole-proof/source audits by
% /root/imprimitive and /root pass; both exact multiplier reconstructions pass.
% This section supplies structural coverage, not the character-table scan.
\subsection{Complete cyclic-normalizer coverage for projective
\texorpdfstring{$\mathrm{PSL}_3(4)$}{PSL3(4)} representations}
\label{sec:l34-normalizer-coverage}

Write $S=\mathrm{PSL}_3(4)$ and let $\widetilde S$ be its universal
perfect central extension.  The published input used below is
Breuer--Malle--O'Brien \cite[\S3.3.1 and Lemma~3.1]{D:BMO2017}.
The complete invoked argument occupies pages~5--6 of the author version
retained with this proof.  The paper identifies
\[
 Z(\widetilde S)\simeq C_4^2\times C_3,
 \qquad \operatorname{Out}(S)\simeq D_{12},
\]
and lists all thirty-six bicyclic extensions up to isoclinism.
Here a bicyclic extension has a perfect central cover of $S$ as its
derived subgroup, cyclic center in that derived subgroup, and cyclic
outer quotient.  The distinction between isomorphism and isoclinism is
handled explicitly below.  No assertion that isoclinic groups have the
same character tables or the same element orders is used.

For a semisimple operator $T$ on an $a$-dimensional complex vector space,
put
\[
 \nu(T)=a-\max_{\lambda\in\mathbf C}
                  \dim\ker(T-\lambda I).
\]
This number is unchanged under scalar multiplication or conjugacy.

\paragraph{Lemma L34--1 (cyclic extension of an invariant representation).}
Let $P\lhd B$ be finite, with $B/P=\langle bP\rangle$ cyclic of
order $d$.  Let $\rho:P\to\mathrm{GL}(W)$ be irreducible, and suppose
that an invertible $T$ satisfies
\[
 T\rho(p)T^{-1}=\rho(bpb^{-1})\quad(p\in P).
\]
Then a scalar multiple of $T$ is the image of $b$ in an irreducible
representation of $B$ restricting to $\rho$ on $P$.

\emph{Proof.}
The operator $T^d\rho(b^d)^{-1}$ commutes with $\rho(P)$, hence equals
$cI$ by Schur's lemma.  Choose $z\in\mathbf C^*$ with $z^dc=1$.
The assignments $p\mapsto\rho(p)$ and $b\mapsto zT$ satisfy the
conjugation relations and $(zT)^d=\rho(b^d)$.
Every element of $B$ has a unique expression $pb^j$ with
$p\in P$ and $0\leq j<d$, so these relations give a representation
of $B$.  It is irreducible because its restriction to $P$ is.
This argument permits any central power relation for the chosen
extension; it does not require the extension to split.\hfill$\square$

\paragraph{Lemma L34--2 (what isoclinism preserves in this setting).}
Let $A,B$ be isoclinic finite groups.  Suppose that
$Q=A'$ and $P=B'$ are perfect, that $A/Q$ and $B/P$ are cyclic,
and that these quotients act faithfully as outer automorphisms of
$Q/Z(Q)$ and $P/Z(P)$, respectively.  Let $R$ be an irreducible
representation of $A$ whose restriction to $Q$ is irreducible.
For every generator $aQ$ of $A/Q$, there exist an irreducible
representation $R_B$ of $B$ with irreducible restriction to $P$ and
an element $b\in B$ such that $R(a)$ and $R_B(b)$ differ by a scalar,
after identifying their representation spaces.  In particular their
projective eigenspace multiplicities agree.

\emph{Proof.}
An isoclinism consists of isomorphisms
\[
 \phi:A/Z(A)\longrightarrow B/Z(B),\qquad
 \psi:Q\longrightarrow P
\]
compatible with commutators.  With the convention
$[x,y]=xyx^{-1}y^{-1}$, compatibility means
$\psi([x,y])=[x',y']$ whenever $x'Z(B)=\phi(xZ(A))$ and
$y'Z(B)=\phi(yZ(A))$.
For every $q\in Q$ one has
\[
       \phi(qZ(A))=\psi(q)Z(B).                    \tag{L34.1}
\]
Indeed this holds on commutators by compatibility and the homomorphism
property of $\phi$, and every element of $Q=A'$ is a product of
commutators.

Faithfulness of the outer quotient implies $Z(A)\subset Q$ and
$Z(B)\subset P$: a central element has trivial outer action.
Equation~(L34.1) therefore induces an isomorphism $A/Q\simeq B/P$.
Choose $b$ with $bZ(B)=\phi(aZ(A))$.  Then $bP$ generates $B/P$.
For $q\in Q$, commutator compatibility gives
\[
 \psi(aqa^{-1})
   =\psi([a,q])\psi(q)
   =[b,\psi(q)]\psi(q)=b\psi(q)b^{-1}.             \tag{L34.2}
\]
Transfer $R|_Q$ to the irreducible representation
$\rho=R|_Q\circ\psi^{-1}$ of $P$.  By (L34.2), the same operator
$R(a)$ implements conjugation by $b$ on $\rho(P)$.
Lemma~L34--1 extends $\rho$ to $B$ after rescaling that operator.
Scalar multiplication preserves all eigenspace dimensions, proving
the assertion.\hfill$\square$

\paragraph{Theorem L34--3 (complete table coverage).}
Let $W$ be a nontrivial irreducible projective representation of $S$,
and let $Q\subset\mathrm{GL}(W)$ be the faithful quasisimple image
of $\widetilde S$.  Let $T$ be a finite-order operator normalizing
$Q$.  Then, up to scalar multiplication and an identification of the
representation space, $T$ occurs as an element in an irreducible
representation of one of the following thirty-six groups whose
restriction to its perfect derived subgroup is irreducible:
\[
\begin{array}{c|c|c}
\text{perfect core}&\text{cyclic outer extensions to include}
                    &\text{number including the core}\\ \hline
S&.2_1,.2_2,.2_3,.3,.6&6\\
2.S&.2_1,.2_2,.2_3&4\\
3.S&.2_1,.2_2,.2_3,.3,.6&6\\
4_1.S&.2_1,.2_2,.2_3&4\\
4_2.S&.2_1,.2_2,.2_3&4\\
6.S&.2_1,.2_2,.2_3&4\\
12_1.S&.2_1,.2_2,.2_3&4\\
12_2.S&.2_1,.2_2,.2_3&4
\end{array}                                                     \tag{L34.3}
\]
The labels are the ATLAS labels in the cited published list.  In
particular, any scalar-invariant eigenspace inequality verified on
all those irreducible restrictions applies to every compatible
normalizing operator $T$.

\emph{Proof.}
The kernel of $\widetilde S\to Q$ is central, since the projective
representation of the simple group $S$ is nontrivial and hence
faithful.  The center of $Q$ acts by scalars, and a finite scalar
subgroup is cyclic.  Thus $Q$ is a cyclic-center quotient of
$\widetilde S$.  The kernel used here is the kernel of its actual
central character; no characteristic-subgroup property is assumed.

Let $d$ be the order of the outer automorphism of $S$ induced by
$T$.  If $d=1$, choose $q\in Q$ inducing the same inner
automorphism.  Then $q^{-1}T$ centralizes the irreducible group $Q$,
so $T$ is a scalar multiple of $q$.  The core groups in (L34.3)
already cover this case by the cited classification.

Suppose $d>1$.  There is $q\in Q$ such that conjugation by $T^d$
equals conjugation by $q$ on $Q$.  To see that equality on $S$
also gives equality on $Q$, note that an automorphism of the
perfect group $Q$ inducing the identity on $Q/Z(Q)$ has the form
$x\mapsto xz(x)$ with $z:Q\to Z(Q)$ a homomorphism, hence is
the identity.  Therefore $T^dq^{-1}=cI$ by Schur's lemma.
Rescale $T$ to $T_0$ with $T_0^d=q$, and put
\[
                       A=\langle Q,T_0\rangle.
\]
This group is finite: it is the union of the $d$ cosets
$QT_0^j$, $0\leq j<d$.  Its quotient $A/Q$ is cyclic of order
exactly $d$, because its image in $\operatorname{Out}(S)$ has
that order.  Since $Q$ is perfect and $A/Q$ is abelian,
$A'=Q$.  Thus $A$ is a bicyclic extension of precisely the type
classified in \cite[Lemma~3.1]{D:BMO2017}.

Choose its isoclinism representative $B$ from (L34.3).
Both outer quotients are faithful, and the perfect derived subgroup
of $B$ is the displayed core.  Lemma~L34--2 applies to the actual
matrix representation of $A$, which is irreducible on $Q$.
It identifies the projective spectrum of $T_0$, and therefore that
of $T$, with an element in an irreducible representation of $B$
having irreducible core restriction.  This proves coverage.
It also explains exactly why one fixed table for each listed
isoclinism class suffices: central power relations are absorbed by
the scalar in Lemma~L34--1.\hfill$\square$

\paragraph{Finite multiplier check and source precision.}
The complete structural classification in Theorem~L34--3 is a
published input, not a newly proved classification.  The retained
independent program \texttt{check\_multiplier.py} checks its finite
kernel-orbit count directly from the displayed action in the source.
Write $M_4=\langle a,b\rangle\simeq C_4^2$, $c=(ab)^{-1}$.
The displayed order-six generator acts by
$a\mapsto b^3$, $b\mapsto c^3$, while the displayed $2_2$
fixes $a$ and interchanges $b,c$.  The first fixes $C_3$ and
the second inverts it.  These maps generate $D_{12}$ on
$M_4\times C_3$.
There are twenty distinct kernels of characters of this multiplier.
Their eight outer orbits have quotient orders
\[
              1,2,3,4,4,6,12,12.
\]
The order-one and order-three quotients have outer stabilizer of
order twelve; each other orbit has stabilizer of order four.
Enumerating simultaneous orbits of a character kernel and a cyclic
subgroup of its stabilizer gives exactly thirty-six pairs, with
the profile in (L34.3).  This check imports no character-table data.

There is a minor label inconsistency in the source's subsequent
centralization prose: its displayed $2_2$ action inverts
$M_4/\langle a\rangle$ and centralizes
$M_4/\langle ab^2\rangle$, whereas the following sentences describe
these centralizations oppositely.  The same wording occurs in the
official character-library example.  Our proof uses neither of
those centralization identifications: it includes both order-four
cover types, all three involution types for each, and the complete
published Lemma~3.1 list.  The independent count treats the two
order-four orbits as unordered.  This local label issue therefore
does not leave any compatible normalizer out of (L34.3).

Finally, selecting rows by irreducibility of restriction to the
perfect core is sufficient.  Each representation constructed in
Lemma~L34--1 has that property.  Including further such rows,
whether faithful or not on the core center, can only enlarge the
scan.  The numerical eigenvalue bounds and the validity of the
retained character tables are separate checks; this section makes
no independent numerical-spectrum claim.

\section{Exclusion of every rank-three linear and unitary core}
\label{sec:a2-cores}

\paragraph{Status and exact scope.}
This is a new proof from complete published character tables, their
corrected official mathematical implementation, and the published
classification of bicyclic extensions of $\operatorname{PSL}_3(4)$.
The full proof has passed the root reading and an independent whole-proof
audit by \texttt{/root/g34\_eightfolds}.  That audit independently
reconstructed every generic inequality from the original frozen expressions
and every finite spectrum from the complete retained tables.  The separate
L34 normalizer proof also passed a whole-proof audit by
\texttt{/root/imprimitive}.  The precise source and computation scopes
are recorded in their reports and the check-event register.
All field sizes $q\ge3$ are covered.  The simple group
$\operatorname{PSL}_3(2)\simeq\operatorname{PSL}_2(7)$ is expressly
outside the statement: its degree-three representation occurs in the
already successful split-sixfold construction.  The group
$\operatorname{PSU}_3(2)$ is not simple.

For a finite-order operator on a complex vector space of dimension $a$,
put $\nu(T)=a-\max_\lambda\dim\ker(T-\lambda)$.
Scalar multiplication does not change $\nu$.

\subsection*{D-A2-01. A uniform bound on every ordinary character}

\begin{lemma}[D-A2-01]
\label{D-A2-01}
Let $U=\operatorname{SL}_3(q)$, $q\ge5$, or
$U=\operatorname{SU}_3(q)$, $q\ge7$.  For every nontrivial ordinary
irreducible character $\chi$ and every $u\notin Z(U)$,
\[
                 |\chi(u)|\le\chi(1)/4.                 \tag{A2.1}
\]
The minimum nontrivial degree is at least $30$ in the linear range
and at least $42$ in the unitary range.
\end{lemma}

\begin{proof}
Use the complete Simpson--Frame tables, Tables~1a, 1b and~2
\cite{A2:SF73}, with the unitary corrections of Geck,
\S3 \cite{A2:Geck90}, as implemented in the four complete official
tables of \cite{G2T:GCT}.  We freeze the implementation at commit
\texttt{87231738271ed548f62e7f1c98ff2bbfc1ed30d3}.
The inventory is
\[
\begin{array}{l|l|r|r|r|r}
\text{table}&\text{field condition}&q_0&\text{row types}&
 \text{class types}&\min\chi(1)\\\hline
\texttt{SL3.1}&q\equiv1\pmod3&5&17&10&30\\
\texttt{SL3.n1}&q\not\equiv1\pmod3&5&8&8&30\\
\texttt{SU3.2}&q\equiv2\pmod3&7&17&10&42\\
\texttt{SU3.n2}&q\not\equiv2\pmod3&7&8&8&42
\end{array}                                                     \tag{A2.2}
\]
The first column is the entire central class type.  In the tables
with center of order three it is parameterized by $a=0,1,2$;
its value need not equal the degree until $a=0$ is substituted.
The tables of types $17\times10$ are rectangular because a type
can parameterize multiple distinct characters or classes.

Every noncentral entry has the exact form
$z(q)=\sum_\alpha p_\alpha(q)u_\alpha$ with
$p_\alpha\in\mathbf Q[q]$ and $|u_\alpha|=1$ at every admissible
specialization of its integral parameters.  Set $x=q-q_0\ge0$ and
write
\[
 p_\alpha(q_0+x)=\sum_k c_{\alpha k}x^k,
 \qquad d(q_0+x)=\sum_k d_kx^k.
\]
Then
\[
 |z(q)|\le\sum_k\left(\sum_\alpha|c_{\alpha k}|\right)x^k.
                                                                  \tag{A2.3}
\]
For every nontrivial row and every column other than the first,
the exact coefficient comparison is
\[
       d_k/4-\sum_\alpha|c_{\alpha k}|\ge0
                         \quad\text{for every }k.          \tag{A2.4}
\]
All $386$ comparisons pass across all $46$ nontrivial row types.
Every translated degree polynomial has nonnegative coefficients,
with constant term at least the stated minimum.  Every class-length
polynomial except the first has nonnegative translated coefficients
and constant term greater than one.  Thus the omitted column is
exactly the center, and (A2.3)--(A2.4) prove (A2.1) for every
admissible field size and every class and character parameter.

The finite certificate is fully retained.  The four original files
have respective SHA-256 hashes
\begin{quote}\small\ttfamily
156d7173feac9283e0c3afb74c8c91c0bb13a1110bff274bc6d19be234b59c50\\
a5d5c0006bded69a9164a727104562226a6398dbe031bed15ba99e11fa58b7ec\\
364950438f2a8dac169b1834d43ec8957aae1781df1cfe9485834febf8846acc\\
f64356e81ec3f5fd951d551122cf0880dc7a8ebf352ec298148fe70695bce9b4.
\end{quote}
The retained original expressions, degrees and class lengths are in
\texttt{research/a2\_cores/generic\_inputs.json}, SHA-256
\texttt{4df14adc254b80d4d7a03ea18e9449c1e05c1ba72f2ad8b713067c8f252b01a2}.
The source extractor verifies these original hashes and executes
no downloaded Julia code.  Its identity check uses a restricted
arithmetic parser to verify that every central phase exponent
vanishes at $a=0$ and that the resulting entry is the separately
declared degree.  The previously reviewed formal-phase parser
collects exact rational polynomial coefficients.  The deterministic
program \texttt{build\_certificate.py} records every coefficient
in (A2.4), every degree polynomial and every translated degree.
Its output \texttt{generic\_certificate.json} has SHA-256
\texttt{7a8ec14dfbec9f2dbd29127c7aadbaa0bae46749dfa2375c1ba3a064318fa4db}.

The input boundary matters.  The corrected implementation includes
all central characters, not merely the third with trivial central
character in a misprinted header of Table~1b.  It uses the full
parameter moduli $q^2+q+1$ and $q^2-q+1$, the corrected phases on
the split cuspidal rows, and the corrected unitary class representative
whose eigenvalues have determinant one.  The linear mixed-torus
phase and the parameter exclusions are also corrected.  These
corrections are documented in the frozen source metadata; the
proof uses those complete corrected tables and does not certify an
uncorrected transcription of the printed tables.  The primary
character-table classification is a cited input, not a new result
of this certificate.
\end{proof}

\subsection*{D-A2-02. Every central cover and compatible normalizer}

\begin{lemma}[D-A2-02]
\label{D-A2-02}
Let $N$ act faithfully and irreducibly in degree $a$, where $N$ is
quasisimple with simple quotient $\operatorname{PSL}_3(q)$ or
$\operatorname{PSU}_3(q)$, $q\ge3$.  Then $a\ge6$, and every
nonscalar finite-order element of its complex linear normalizer
satisfies $\nu(T)\ge3a/16$.

If the simple quotient is one of
\[
 \operatorname{PSL}_3(3),\quad\operatorname{PSL}_3(4),\quad
 \operatorname{PSU}_3(3),\quad\operatorname{PSU}_3(4),\quad
 \operatorname{PSU}_3(5),                                  \tag{A2.5}
\]
the stronger estimate $\nu(T)\ge a/4$ holds.
For every other quotient the respective minimum $30$ or $42$
of D-A2-01 applies.
\end{lemma}

\begin{proof}
The complete exceptional-multiplier list
\cite[Chapter~5, Table~1, numbered p.~35]{BI:HHM} shows that the
only exceptional multiplier in the stated range is that of
$\operatorname{PSL}_3(4)$.  Otherwise the universal perfect cover
is the ordinary special linear or unitary group.  Every
representation of a central quotient lifts to such a cover.
For the generic ranges, D-A2-01 and the commutator argument
D-BI-18 give $\nu(T)\ge3a/16$: choose $u$ with $[T,u]\ne1$,
use $\operatorname{rank}([T,u]-1)\le2\nu(T)$, and compare
the trace of this commutator with $a/4$.  If it is a nonidentity
scalar, its rank is $a$, which gives a stronger inequality.
Thus this argument covers all compatible outer automorphisms
without any enumeration of generic outer cosets.

Here is the complete small-group calculation, using the official
CTblLib tables \cite{SPOR:CTblLib}.  Apart from $\operatorname{PSL}_3(4)$
the required tables and exact outcomes are
\[
\begin{array}{l|r|r|r|c}
\text{table}&\text{selected rows}&\text{nonscalar spectra}&
 \min a&\min\nu/a\\\hline
\texttt{L3(3)}&11&121&12&1/3\\
\texttt{L3(3).2}&10&140&12&1/3\\
\texttt{U3(3)}&13&169&6&2/7\\
\texttt{U3(3).2}&10&150&6&1/3\\
\texttt{U3(4)}&21&441&12&1/3\\
\texttt{U3(4).2}&10&190&12&1/3\\
\texttt{U3(4).4}&12&252&12&1/3\\
\texttt{3.U3(5)}&39&1443&20&8/21\\
\texttt{3.U3(5).2}&14&420&20&8/21\\
\texttt{3.U3(5).3}&90&8370&20&8/21
\end{array}                                                     \tag{A2.6}
\]
Their multiplier and outer-group data are respectively
$(1,C_2),(1,C_2),(1,C_4),(C_3,S_3)$, with the first pair for
$\operatorname{PSL}_3(3)$ and the others in increasing unitary
field size.  These are retained as the official extension metadata.
Every cyclic subgroup of each outer group is represented in (A2.6).
An invariant irreducible representation extends to any fixed
cyclic extension after rescaling its implementing matrix: if a
coset generator $b$ satisfies $b^d=u_0$, Schur's lemma says
$T^d=c\rho(u_0)$, and a scalar $d$th root corrects the relation.
Hence every compatible normalizer has, up to scalar, one of the
tested spectra.  This argument permits nonsplit cyclic extensions
and arbitrary central kernels.

For $S=\operatorname{PSL}_3(4)$ the universal multiplier is
$C_4^2\times C_3$ and $\operatorname{Out}(S)\simeq D_{12}$.
The new cyclic-extension and isoclinism-transfer lemmas
L34--1--3, proved in the adjoining section, reduce every actual
normalizing matrix to the complete published list of $36$
bicyclic extensions in \cite[\S3.3.1, Lemma~3.1]{D:BMO2017}.
Their core types are
\[
 S,\quad2.S,\quad3.S,\quad4_1.S,\quad4_2.S,\quad6.S,
 \quad12_1.S,\quad12_2.S.                                \tag{A2.7}
\]
For $S$ and $3.S$ the tables include the core, all three
involution extensions, and the order-three and order-six extensions.
For each other core they include the core and all three involution
extensions.  All $36$ tables are checked.  The smallest degree is
$6$, and every selected nonscalar spectrum satisfies the stronger
bound $\nu/a\ge3/10$.  Both order-four cover types and every
involution type are retained, so the source's discrepancy between
one displayed multiplier action and its following cover-label prose
does not affect coverage.

For completeness the computed universal-cover table
\texttt{(4\^{}2x3).L3(4)} is also included: its $303$ nonlinear
rows give $77568$ nonscalar spectra, minimum degree $6$ and
minimum ratio $1/3$.  Its retained provenance is construction by
\texttt{PossibleCharacterTablesOfTypeV4G}, rather than an ATLAS
source table.  It is a redundant consistency check; every essential
degree and support bound follows already from the eight cyclic-center
ATLAS core types and their published extension list.
No stored universal-cover-to-quotient fusion was available in GAP,
and none is claimed or used: completeness is supplied by L34--1--3
and the published $36$-extension classification.

The resulting complete input has $47$ tables, $1603$ selected
rows and $141024$ nonscalar spectra.  The exporter
\texttt{export\_small.g} retains every character, class size,
element order, prime power map, source string, and the actual
core-to-extension class fusions; it verifies full character-row
orthogonality.  The exact standalone verifier
\texttt{check\_small.py} selects every nonlinear core row and
every extension row whose restriction is an irreducible core row,
using exact equality of all character values.  Its checks include
degree-square and class-size sums, all $34$ core-to-extension
fusions, preservation of orders and power maps, trivial identity
kernel, normal class-size unions, and the full cyclic quotient
order detected by linear characters.  For an element of order $o$
it reconstructs all eigenvalue multiplicities by
\[
 m_b=\frac1o\sum_{j=0}^{o-1}\chi(g^j)\zeta_o^{-bj}
       \quad(0\le b<o),                                  \tag{A2.8}
\]
using exact cyclotomic-polynomial arithmetic.  Each multiplicity
is checked to be a nonnegative integer and their sum to be the
degree.  Scalar classes are exactly those with
$|\chi(g)|^2=\chi(1)^2$.  All the stated lower bounds pass.
The frozen \texttt{small\_inputs.json} has SHA-256
\texttt{3f356063efc3a3987126ddbbb6d12b0540fd71151f93eb4e83f4c3f6c26607e8}.
The individual row minima and spectra counts are retained in
\texttt{small\_checked.json}; the full row selection is not
replaced by a sample of low degrees.  This proves the lemma.
\end{proof}

\subsection*{D-A2-03. Exclusion of the entire supplier family}

\begin{theorem}[D-A2-03]
\label{D-A2-03}
Let an inducing block of a finite-cover Weil-sixfold source have
finite image $J\subset\operatorname{GL}(V)$.  Suppose it has an
irreducibly acting normal central product of quasisimple groups
whose simple quotients are $\operatorname{PSL}_3(q)$ or
$\operatorname{PSU}_3(q)$, with $q\ge3$ in every factor.
Then the actual period rank is at most eight.

Mixed linear and unitary types and field sizes, scalar phases,
compatible outer automorphisms and factor permutations are allowed.
No essential-self-duality, primitivity or tensor-indecomposability
assumption is needed, and there is no restriction on the inducing
character field or the polarization determinant.
\end{theorem}

\begin{proof}
Write $V=\bigotimes_{i=1}^t V_i$, with faithful irreducible
factor dimensions $a_i\ge6$, and $e=\prod_i a_i$.
The simple factors and their generated full matrix algebras are
intrinsic, so $J$ acts through factor intertwiners and compatible
factor permutations.  If a nonscalar element preserves every
factor, one intertwiner is nonscalar.  The largest eigenvalue
multiplicity fraction of a tensor product is at most that of
any one nonscalar factor.  D-A2-02 gives $\nu(g)\ge3e/16$.
For a nontrivial permutation cycle, the weighted cyclic-word
argument of D-SPOR-03, spelled out in D-G2T-03, gives
$\nu(g)\ge e/4$: after basis changes there is cyclic rotation
with one diagonal twist, and each word orbit contributes at
most one vector to a specified eigenspace.  Constant words
number $a$, while every other orbit has size at least two,
so the largest eigenspace is at most $(a^l+a)/2\le3a^l/4$.
Nonidentity scalar inertia has ordinary moving rank $e$.

If the actual period rank were nine, the intermediate-cover
factorization D-BI-01 requires genus zero and at least twelve
genuine branch points.  Indeed the multiplicity-six Euler formula
gives total moving rank $D=6+2e(1-h)$; for $h\ge2$ this is
negative since $e\ge6$, while for $h=1$ there are at most six
branch points and the intermediate Hurwitz image has dimension
at most six.  For $h=0$ the source dimension condition gives
$R\ge12$ and $D=2e+6$.  Therefore
\[
       12\left\lceil\frac{3e}{16}\right\rceil\le2e+6,
                     \qquad e\le24.                     \tag{A2.9}
\]
If $t\ge2$ then $e\ge36$, a contradiction.  For a single
factor in the generic range the lower bounds $30$ or $42$
also contradict (A2.9).  For one of the small quotients (A2.5),
the quarter-support bound improves the necessary inequality to
$12\lceil e/4\rceil\le2e+6$, hence $e\le6$.
The only remaining integer degree is $e=6$, for which the
left side is $24$ and the right side is $18$, again impossible.
This excludes every case and proves the actual-rank bound.
\end{proof}

\paragraph{Geometric and programme meaning.}
The theorem removes all these possible curve suppliers from the
remaining nonsplit-sixfold gate Q2i.  It gives a rank obstruction
for their actual variations of Hodge structure and associated
isotypic $\ell$-adic local systems, hence for this attempted
motivic supplier.  It produces no detecting cycle and no new
numerical bound by itself.  In the project's language it restricts
the supplier needed for the integral uniform form of condition~A
(cote~16, scan~050; SIGN\_NEXT SN-C03/C04); positivity, condition~I
(scan~049; cote~44, Th.~4), cannot repair the absent rank-nine
supplier.  The existing $K$-theory projectors and Cayley--Hamilton
division step of scans~048--049 still apply to sources that
actually supply cycles.  No conclusion about a ghost follows
from excluding this one source class.

\section{Unconditional finite even-symplectic normal-core exclusions}
\label{sec:higher-even-symplectic-finite}

\noindent\textbf{Status.} Complete authored proof and fresh complete finite
certificates; root whole-proof/code review and fresh $756$-row,
$271785$-entry rerun pass. The complete mathematical argument and all three character tables,
including the newly added third table, have independent whole-proof and
arithmetic review by /root/g34\_eightfolds. Its separate reconstruction
checks all $756$ rows and $271785$ entries; see the retained
\texttt{HIGHER\_EVEN\_SYMPLECTIC\_FINITE\_PEER\_AUDIT.md}. This section does not use the unresolved Gluck--Magaard (1998)
proper-parabolic estimate appearing in the separate all-rank draft.

\begin{theorem}[D-HES-F: three complete core families]
In the signature-$(3,3)$ sixfold-source problem, suppose the faithful image
$J$ of an irreducible inducing representation has an irreducible normal
subgroup $E$ which is a product of simple groups chosen from
\[
          \operatorname{Sp}_8(2),\quad \operatorname{Sp}_{10}(2),
                         \quad \operatorname{Sp}_{12}(2).
\]
Then the actual period image has dimension less than nine. The conclusion
allows arbitrary normalizing scalars, every compatible permutation of
simple factors, arbitrary multiplicities of the three types, and inducing
character fields larger than the original imaginary quadratic field.
No determinant or reflection assumption is needed.
\end{theorem}

\begin{proof}
First we establish a complete finite representation-theoretic bound. The
ordinary character tables supplied by CTblLib are
\texttt{S8(2)}, \texttt{S10(2)} and \texttt{S12(2)}. The first two
provenance strings name the ATLAS~\cite{SPOR:Atlas,SPOR:CTblLib}. The third attributes the table to
Christoph K\"ohler, computed using group-theoretic and character-theoretic
methods~\cite{HES:S12Data}. We use that official complete-table input as such; we do not claim
that it is an ATLAS table or that its construction has a verified separate
refereed publication. For every nontrivial
irreducible character, and every nonidentity class, the exact certificate
establishes the following table:
\[
\begin{array}{c|r|r|c|r}
\text{group}&\text{all irreducible rows}&\text{minimum degree}&
 \max_{\chi\ne1,x\ne1}|\chi(x)|/\chi(1)&\text{attaining value}\\ \hline
\operatorname{Sp}_8(2)&81&35&3/5&-21\\
\operatorname{Sp}_{10}(2)&198&155&17/31&-85\\
\operatorname{Sp}_{12}(2)&477&651&11/21&-341
\end{array}
\tag{HESF.1}
\]
The attaining values occur in rows of the displayed minimum degrees at
involutions. Every row is real valued. The exporter verifies full row
orthogonality and reality in exact cyclotomic arithmetic. Its retained
input includes every row, every class, class sizes, orders, power maps,
Frobenius--Schur indicators, and the original provenance strings.
The independent arithmetic routine uses only this retained input. For an
entry $\sum_jc_j\zeta_N^j$, it bounds the absolute value by
$\sum_j|c_j|$, comparing rational numbers only. All $271785$ nonidentity
entries of nontrivial rows satisfy the displayed bounds. At each maximum
an actual rational entry attains the bound, proving equality, rather than
merely an upper estimate. The routine also verifies group-order formulas,
class-size and degree-square sums, all degrees, and reality by reduction
modulo the relevant cyclotomic polynomial. Thus the common bound $3/5$
is proved for all representations and all nonidentity elements of these
three groups; no sample or extrapolation is involved.

If $A$ is a finite-order matrix of degree $a$ and an eigenvalue has
multiplicity $m$, then $|\operatorname{tr}A|\ge 2m-a$ by the triangle
inequality. It follows from (HESF.1) that every nonidentity core operator
has projective support
\[
          a-\max_\lambda\dim\ker(A-\lambda I)\ge a/5.
\tag{HESF.2}
\]
The core groups are simple, centerless, have trivial Schur multiplier, and
have no outer automorphisms. The precise all-$n\ge4$ classification facts
for $\operatorname{Sp}_{2n}(2)$ are recorded in Gluck--Magaard,
\emph{Absolute Fourth Moments}, author p.~25~\cite{HES:GM2006}. The complete exceptional
multiplier list in \cite[Chapter~5, Table~1, numbered p.~35]{BI:HHM}
contains only the ranks two and three in this even symplectic family.
These known classification facts do not require the fourth-moment
hypothesis used elsewhere in that paper.

Write $V=\bigotimes_iV_i$ for the restriction of the inducing space to
$E$, with $a_i=\dim V_i$ and $e=\dim V$. Every factor module is nontrivial;
a trivial factor could not belong to the faithful image. A normalizer
operator permutes isomorphic simple factors. After that permutation is
removed, each induced automorphism is inner, because the outer groups
are trivial. Schur's lemma therefore makes the operator a scalar times a
tensor of core operators. No additional projective cover is omitted,
since the multipliers of these groups are trivial.

For factor-preserving tensors, eigenvalue multiplicity distributions
convolve on roots of unity. The largest normalized multiplicity of the
convolution is no greater than the largest normalized multiplicity of any
one factor. For a nontrivial factor-permutation cycle of length $l$, the
factor dimensions are equal, say $a$. Conjugating the tensor coordinates
puts all twists on one coordinate; diagonalize their product. On the
resulting tensor basis the operator is a weighted cyclic permutation.
There are $a$ constant words among $a^l$ basis words. Every other cyclic
orbit has length at least two, and each weighted cycle contributes at
most one copy of a specified eigenvalue. Consequently its largest
multiplicity is at most
\[
           a+(a^l-a)/2\le 3a^l/4\qquad(a\ge2,\ l\ge2).
\tag{HESF.3}
\]
Eigenvalue coincidences do not alter the orbit count. Tensoring with the
other cycles cannot increase this normalized bound. Combining (HESF.2)
and (HESF.3), every nonscalar operator in $J$ has projective support at
least $e/5$. A nonidentity scalar has actual moving rank $e$. Thus every
nonidentity inertia in $J$ has actual moving rank at least $e/5$.

The reality certified for the core characters makes each $V_i$ self-dual.
Therefore $V$ has a nondegenerate $E$-invariant bilinear form, whose line
is unique by irreducibility. The full normalizer preserves this line and
acts on it by a multiplier character $\mu:J\to\mathbf C^\times$.
This character is nontrivial here: if it were trivial, every finite-order
operator in $J$ would have reciprocal eigenvalue pairs, and hence real
trace. Its character field would be real. But the inducing character
field contains the original imaginary quadratic field, by rational
induction and Frobenius reciprocity as in D-BI-01. Only that inclusion,
not equality of the two fields, is being used.

We apply the actual intermediate-cover factorization of D-BI-01 to this
inducing representation. If $h$ is the intermediate base genus and
$g_1,\ldots,g_r$ its nonidentity effective inertias, the six-dimensional
complex multiplicity variation has Euler formula
\[
               6=2(h-1)e+\sum_{j=1}^r\operatorname{rank}(g_j-I).
\tag{HESF.4}
\]
Here $e\ge35$. Hence $h\ge2$ is impossible. If $h=1$, the sum is six,
so $r\le6$, and the moduli dimension of the pointed intermediate curve
is at most six (one if $r=0$). The finite-cover factorization bounds the
actual period image by that dimension. A nine-dimensional image would
therefore require $h=0$, $r\ge12$, and
\[
                   \sum_j\operatorname{rank}(g_j-I)=2e+6.
\tag{HESF.5}
\]
This is a necessary condition from an actual variation factorization;
no sufficiency is inferred from a Hurwitz-space dimension.

The inertia tuple generates $J$ and has product one. Since $\mu$ is
nontrivial, at least two entries have multiplier different from one.
For such an entry its $1$-eigenspace is totally isotropic for the
nondegenerate invariant bilinear form, hence has dimension at most
$e/2$. These two entries cost at least $e$ in (HESF.5). The remaining
at least ten nonidentity entries each cost at least $e/5$. Therefore
\[
                    2e+6\ge e+10e/5=3e,
                    \qquad e\le6,
\tag{HESF.6}
\]
contrary to $e\ge35$. This proves the theorem with the full stated
normalizer, tensor, scalar, and larger-field scope.
\end{proof}

\paragraph{Retained certificate and scope.}
The complete frozen input is \texttt{sp2n2\_character\_inputs.json};
its SHA-256 is \texttt{5d3539d84b4988ac34c2703cbc81369b2840e50e3d74ea3afbb2db31f974e615}.
The exporter is \texttt{export\_sp2n2.g}, and the exact verifier and output
are \texttt{check\_sp2n2.py} and \texttt{check\_sp2n2.json}.
The fresh run checks $756$ rows and $271785$ entries. The original two-table
input is retained separately as \texttt{sp2n2\_character\_inputs\_v1.json}
to identify precisely the input used by the earlier peer computation.
The earlier all-rank draft retains its separate (GM) primary-source gate.
A new elementary replacement is under review; none of its conclusions
is used in this finite theorem.

\section{The full higher even symplectic normal-core family}
\label{sec:higher-even-symplectic}

\noindent\textbf{Status.} Complete new authored proof. Root and independent peer
/root/g34\_eightfolds have each whole-reviewed the proof and reported
PASS; root also read the full author certificate and reran it successfully.
The mathematical version reviewed by both has SHA-256
\texttt{72e295707440f2e06e39658d2b3528f61d9587227e99e6f18196b075515dbc3c}. The elementary construction in D-HES-02
replaces the unverified Gluck--Magaard (1998) input in the archived
conditional draft. That unverified source claim is not used here:
the only character-ratio input is the explicitly verified published
unipotent bound (HES.3). The finite tables and their precise provenance
are retained in D-HES-F.

\subsection*{D-HES-01. Elementary degree bounds and verified inputs}

Put $S_n=\operatorname{Sp}_{2n}(2)$ for $n\ge4$. These groups are simple,
centerless, have trivial Schur multiplier and have no outer automorphisms.
The exact all-$n\ge4$ multiplier and automorphism statements occur in
\cite[author p.~25]{HES:GM2006}; the exceptional-multiplier list in
\cite[Chapter~5, Table~1, numbered p.~35]{BI:HHM} is consistent with this.
These classification facts do not require the fourth-moment hypothesis
used elsewhere in the former paper.

Every nontrivial irreducible complex representation of $S_4$ has degree
at least $35$; for $S_5$ the minimum is $155$. These are complete finite
table assertions already certified in D-HES-F. Only its two ATLAS
endpoints $S_4,S_5$ are essential here: the independently checked
K\"ohler $S_6$ table is redundant for this all-rank theorem.
For $n\ge5$, every
nontrivial irreducible representation of $S_n$ consequently has degree
\[
                     a\ge\max\{155,2^{n-1}-1\}.
\tag{HES.1}
\]
Indeed it is faithful because $S_n$ is simple. Restriction to the block
subgroup $S_5$ has a nontrivial constituent, proving the first bound.
For the second embed $\operatorname{GL}_n(2)$ in $S_n$ by
$A\mapsto\operatorname{diag}(A,A^{-\mathsf T})$. Its subgroup
\[
 U_0=\left\{\begin{pmatrix}1&v\\0&I_{n-1}\end{pmatrix}:
                                  v\in\mathbf F_2^{n-1}\right\}
\]
is elementary abelian, and its normalizer contains
$\operatorname{GL}_{n-1}(2)$ transitively permuting all its nontrivial
characters. Faithfulness makes one nontrivial weight occur, hence all
$2^{n-1}-1$ occur with equal positive multiplicity. This proves (HES.1).

We use two further known inputs with their precise scopes. Gow's theorem
states that every element of a symplectic group in characteristic two is
a product of two involutions, hence is conjugate to its inverse
\cite[theorem on p.~583; proof, Section~3, pp.~589--591]{HES:Gow1981}.
All ordinary irreducible characters of $S_n$ are therefore real valued.
Gluck's unipotent character-ratio estimate, stated with its definitions in
the published survey \cite[Theorem~2.4(i), author pp.~2--4]{HES:Liebeck2017},
gives for our groups
\[
       |\chi(u)|\le\tfrac34\chi(1)
          \qquad(1\ne u\text{ unipotent},\quad 1\ne\chi\in\operatorname{Irr}(S_n)).
\tag{HES.3}
\]
The nonunipotent part of the same source has a different bound and is not
silently strengthened. The original reference for (HES.3) is Gluck (1995),
as specified in that survey.

\subsection*{D-HES-02. An elementary parabolic argument}

For a finite-order operator $x$ on a complex space of dimension $a$, write
\[
           \nu(x)=a-\max_\lambda\dim\ker(x-\lambda I).
\]

\begin{lemma}[D-HES-02a: fixed-point-free elementary subgroup]
Let $h\in S_n$ have odd prime order $p$, and suppose that $h$ stabilizes
a nonzero totally isotropic subspace of the natural symplectic space
$W=\mathbf F_2^{2n}$. Then there is a nonzero $h$-normalized elementary
abelian $2$-subgroup $U\le S_n$ with $C_U(h)=1$ and $|U|\ge4$.
\end{lemma}

\begin{proof}
Let $R$ be such an isotropic subspace. First suppose that $h|_R$ is
nontrivial, and choose a basis $r_1,\ldots,r_s$ of $R$. Write $b$ for the
symplectic pairing. For a symmetric matrix $B=(b_{ij})$ over $\mathbf F_2$
define
\[
             N_B(x)=\sum_{i,j}b_{ij}b(r_j,x)r_i.
\]
Since $R$ is isotropic, $N_BN_C=0$ for all $B,C$. Symmetry of $B$ and
characteristic two imply
$b(N_Bx,y)+b(x,N_By)=0$, while $b(N_Bx,N_By)=0$.
Thus $1+N_B$ is symplectic, and these operators form an elementary
abelian subgroup, canonically the additive group of symmetric matrices.
If $A$ is the matrix of $h|_R$, conjugation by $h$ sends $B$ to
$ABA^{\mathsf T}$. This action is nontrivial: for a rank-one matrix
$B=vv^{\mathsf T}$, the image is $(Av)(Av)^{\mathsf T}$, and over
$\mathbf F_2$ the map $v\mapsto vv^{\mathsf T}$ is injective (its diagonal
recovers $v$). Nonidentity $A$ therefore moves one of these matrices.

Since $p$ is odd, this $\mathbf F_2\langle h\rangle$-module is semisimple.
Its sum of nontrivial simple summands is a nonzero invariant subspace
without fixed vectors. The corresponding subgroup is the desired $U$.

Now suppose $h$ fixes $R$ pointwise. Choose $0\ne e\in R$, and choose
$f_0$ with $b(e,f_0)=1$. Put
$f=\sum_{j=0}^{p-1}h^jf_0$. Then $h(f)=f$ and $b(e,f)=p=1$ in
$\mathbf F_2$. Consequently
$W=\langle e,f\rangle\perp W_0$ is an $h$-stable orthogonal decomposition,
and $h|_{W_0}$ is nontrivial, since otherwise $h=1$.
For $v\in W_0$ and $c\in\mathbf F_2$ define the symplectic operator
\[
 q(v,c)e=e,\qquad
 q(v,c)w=w+b(w,v)e\ (w\in W_0),\qquad
 q(v,c)f=f+v+ce.
\]
Direct multiplication gives
\[
       q(v,c)q(w,d)=q(v+w,c+d+b(v,w)).
\tag{HES.4}
\]
In characteristic two $b(v,w)=b(w,v)$ and $b(v,v)=0$.
Thus this group $Q$ is abelian of exponent two, of order $2^{2n-1}$.
Moreover $h q(v,c)h^{-1}=q(hv,c)$, so the action on $Q$ is nontrivial.
Applying the same semisimple-module argument to the elementary abelian
group $Q$ gives a nonzero invariant subgroup $U$ with $C_U(h)=1$.

Finally a one-dimensional vector space over $\mathbf F_2$ has trivial
automorphism group. A nonzero fixed-point-free module for an odd-order
element must have dimension at least two; hence $|U|\ge4$ in both cases.
No invariant isotropic complement, parabolic character formula, or
unverified estimate is used in this construction.
\end{proof}

\begin{lemma}[D-HES-02b: all-element support]
For $n\ge5$, every nontrivial irreducible complex representation $\rho$
of $S_n$ of degree $a$ satisfies
\[
                   \nu(\rho(x))\ge a/8\qquad(1\ne x\in S_n).
\tag{HES.5}
\]
For $n=4$, the complete finite table gives the stronger bound $a/5$.
\end{lemma}

\begin{proof}
We first treat an odd-prime-order $h$ in a proper parabolic subgroup;
equivalently it stabilizes a nonzero totally isotropic subspace. Use $U$
from D-HES-02a, and put $t=\dim V^U$. Every nonidentity element of $U$
is an involution, hence unipotent. Averaging (HES.3) gives
\[
       t=\frac1{|U|}\sum_{u\in U}\chi(u)
        \le a\left(\frac34+\frac1{4|U|}\right)
        \le\frac{13}{16}a.
\tag{HES.6}
\]
The dual action on the character group $U^*$ is also fixed-point-free
away from its identity, because $h-1$ is invertible on $U$ and on its
dual. Every nontrivial character orbit has length $p$. On the sum of
weight spaces belonging to such an orbit, $h$ cyclically permutes equal
dimensional spaces, with $h^p=1$. Each $p$-th root therefore occurs with
exactly one $p$-th of that orbit-sum dimension. On $V^U$ we need only the
trivial upper bound $t$. It follows that every $h$-eigenspace has
dimension at most $t+(a-t)/p$, so
\[
       \nu(\rho(h))\ge(1-1/p)(a-t)
                   \ge\frac23\frac3{16}a=\frac18a.
\tag{HES.7}
\]

Next suppose an odd-prime-order $h$ belongs to no proper parabolic.
Decompose its natural space $W$ into primary summands. A summand for
an irreducible polynomial not equal to its reciprocal is totally
isotropic: it pairs only with the distinct reciprocal summand. Such a
summand is forbidden. A nonzero $1$-eigenspace is nondegenerate symplectic
and contains an invariant isotropic line, so it too is forbidden.
Each remaining irreducible polynomial is self-reciprocal of degree $2d$.
On its isotypic summand set $E=\mathbf F_{2^{2d}}$ and
$E_0=\mathbf F_{2^d}$. Adjointness of the symplectic pairing induces the
involution $z\mapsto\bar z=z^{2^d}$ on $E$. The trace pairing defines a
nondegenerate Hermitian form $H$ by
\[
              b(av,w)=\operatorname{Tr}_{E/\mathbf F_2}(aH(v,w))
                         \quad(a\in E).
\]
The adjoint identity and symmetry of $b$ give
$H(av,bw)=a\bar b H(v,w)$ and $H(w,v)=\overline{H(v,w)}$.
If the $E$-dimension were at least two, diagonalization and surjectivity
of the finite-field norm $E^\times\to E_0^\times$ would give a nonzero
isotropic vector. Its entire $E$-line would be totally isotropic for $b$,
contradicting the absence of an invariant isotropic subspace.
Every such primary polynomial therefore has multiplicity one.

The centralizer is consequently a product of norm-one scalar groups of
orders $2^{d_i}+1$, with $\sum_i d_i=n$. Hence
\[
                   |C_{S_n}(h)|=\prod_i(2^{d_i}+1)\le3^n.
\tag{HES.8}
\]
Character column orthogonality and (HES.1) imply
\[
        \frac{|\chi(h)|}{a}
        \le\frac{3^{n/2}}{\max\{155,2^{n-1}-1\}}<\frac34.
\tag{HES.9}
\]
For $n=5,6$ compare directly with $155$. For $n=7$ the squared inequality
is $16\cdot3^7=34992<35721=9\cdot63^2$.
For subsequent $n$, write $x=2^{n-1}\ge64$ and observe
$(2x-1)^2-3(x-1)^2=x^2+2x-2>0$. This proves the required induction.
All nonidentity powers of $h$ generate the same cyclic group and are
again elliptic, so (HES.9) holds for every one of them. The Fourier
formula for any $p$-th-root eigenspace bounds its dimension by
\[
             \frac1p\left(a+\sum_{j=1}^{p-1}|\chi(h^j)|\right)
                 \le a\left(\frac34+\frac1{4p}\right)\le\frac56a.
\]
Thus the elliptic case has the stronger support bound $a/6$.

If $h$ has order two, (HES.3) gives directly
$\nu(\rho(h))=(a-|\chi(h)|)/2\ge a/8$.
For arbitrary $1\ne x$, choose a nonidentity prime-order power $h=x^k$.
Every $x$-eigenspace lies in an $h$-eigenspace, so its projective support
is no smaller. This proves (HES.5). Finally D-HES-F gives
$|\chi(x)|\le3a/5$ for $S_4$; if an eigenvalue has multiplicity $m$, then
$|\chi(x)|\ge2m-a$, yielding $\nu\ge a/5$.
\end{proof}

\subsection*{D-HES-03. Full normalizers and the actual period image}

\begin{theorem}[D-HES-03: unconditional all-rank exclusion]
In the signature-$(3,3)$ sixfold-source problem, an irreducible inducing
representation whose faithful image has an irreducible normal subgroup
which is a product of groups $\operatorname{Sp}_{2n_i}(2)$, $n_i\ge4$,
cannot have an actual period image of dimension nine. Arbitrary
normalizing scalars, compatible permutations of factors, and inducing
character fields larger than the original imaginary quadratic field are
allowed. No determinant or reflection assumption is needed.
\end{theorem}

\begin{proof}
Let $J\le\operatorname{GL}(V)$ be the inducing image, $e=\dim V$, and
$E$ the specified irreducible normal subgroup. Its centerless simple
factors make it an actual direct product, and
$V=\bigotimes_iV_i$ with every factor nontrivial. A normalizer operator
permutes isomorphic simple factors. After removing this permutation,
its automorphisms on each factor are inner, since their outer groups
are trivial. Schur's lemma makes it a scalar times a tensor of core
operators. The trivial Schur multipliers also rule out an omitted
projective-cover ambiguity.

For a factor-preserving tensor, eigenvalue multiplicity distributions
convolve on roots of unity, and the maximum normalized multiplicity
cannot exceed that of a nonscalar factor. A nontrivial cycle permuting
$l$ equal-dimensional factors can be conjugated so all twists lie on
one coordinate, and their product can be diagonalized. In the resulting
weighted permutation of $a^l$ basis words there are $a$ constant words
and all remaining orbits have length at least two. A specified
eigenvalue thus has multiplicity at most
$a+(a^l-a)/2\le3a^l/4$. Additional tensor factors cannot increase that
normalized bound. D-HES-02 therefore gives projective support at least
$e/8$ on every nonscalar $J$-operator. Nonidentity scalars have actual
moving rank $e$. In particular every effective nonidentity inertia has
actual moving rank at least $e/8$.

Gow's theorem makes all the factor modules self-dual. Thus $V$ has a
nondegenerate $E$-invariant bilinear form with a unique invariant line,
and $J$ acts on that line by a multiplier $\mu:J\to\mathbf C^\times$.
This multiplier is nontrivial. Otherwise every $J$-operator would have
reciprocal eigenvalue pairs and hence real trace, contradicting inclusion
of the original imaginary quadratic field in the inducing character
field. The field inclusion follows from the rational-induction argument
of D-BI-01 and does not assume equality of the fields.

The actual intermediate-cover factorization of D-BI-01 identifies the
six-dimensional multiplicity variation with the parabolic cohomology
of this inducing representation. For intermediate genus $h$ and
nonidentity effective inertias $g_1,\ldots,g_r$ it gives
\[
                6=2(h-1)e+\sum_{j=1}^r\operatorname{rank}(g_j-I).
\tag{HES.10}
\]
Here $e\ge35$, so $h\ge2$ is impossible. If $h=1$ the sum is six,
whence $r\le6$ and pointed-curve moduli dimension is at most six
(one when $r=0$). The actual finite-cover factorization then bounds
period-image dimension below nine. Thus a nine-dimensional image would
require $h=0$, $r\ge12$ and moving-rank sum $2e+6$.

The generating inertia tuple has product one, so at least two entries
have nontrivial multiplier. Their $1$-eigenspaces are totally isotropic
for a nondegenerate bilinear form and each has dimension at most $e/2$.
Those entries cost at least $e$ in total; the remaining at least ten
cost at least $e/8$ each. Consequently
\[
                  2e+6\ge e+10e/8=9e/4,
                         \qquad e\le24,
\tag{HES.11}
\]
contradicting $e\ge35$. This uses an actual variation factorization and
a necessary branch budget, not a dimension inequality as a sufficiency
test for dominance.
\end{proof}

\paragraph{Verification boundary.}
The finite subgroup constructions in D-HES-02 are explicit linear
algebra, proved above. A separate exact program checks representative
matrices, group laws, the weight-orbit averaging arithmetic, and the
symbolic all-rank inequalities; it is supplementary to these proofs.
No claim to have retrieved or verified Gluck--Magaard (1998), Theorem~3.1,
is made. The former conditional argument is retained as an archive,
and its missing source is no longer a premise of this theorem.

\section{All higher symplectic cores in even characteristic}
\label{sec:hes-all-even}

\paragraph{Status and scope.}
This is the complete author derivation D-HES-04--06, with root and
independent-peer whole-proof review both passing. The reviewed
mathematical text has SHA-256
\texttt{ab4898d2725c48b0993539ead8e56c27b9630f0433aee185514c48ab1546ae18}.
The author exact certificate passes, and the peer independently
reconstructed the finite-field identities with different arithmetic.
It extends the reviewed D-HES-01--03
from $\mathbf F_2$ to every even field. It treats all ordinary irreducible
modules and all compatible linear normalizers, including field operators,
scalars and permutations of tensor factors. The conclusion concerns the
actual nine-dimensional period image in the sixfold-source problem.
Neither a reflection hypothesis nor a determinant hypothesis is imposed.

\subsection*{D-HES-04. Core support over every even field}

Let $q\ge4$ be a power of two, $n\ge4$, and
$S=\operatorname{Sp}_{2n}(q)$. It is centerless simple. Its Schur multiplier
is trivial in this range: the ordinary symplectic multiplier is
$(2,q-1)=1$, and there is no exception here in the complete list
\cite[Chapter~5, Table~1, numbered p.~35]{BI:HHM}.
For a finite-order operator $x$ on a complex space of dimension $a$, put
\[
              \nu(x)=a-\max_{\lambda}\dim\ker(x-\lambda I).
\]
This quantity is invariant under multiplying $x$ by a scalar.

We use the following precise published unipotent estimate:
\[
 |\chi(u)|\le c(q)\chi(1),\qquad
 c(4)=\frac34,\qquad
 c(q)=\frac1{\sqrt q-1}\quad(q\ge8),
 \qquad 1\ne u\text{ unipotent},\quad1\ne\chi\in\operatorname{Irr}(S).
\tag{HES.12}
\]
This is \cite[Theorem~2.4(i), definitions and statement on author
pp.~2--4]{HES:Liebeck2017}. The subtraction of $1$ is outside the
radical. In particular $c(8)$ is not less than $1/2$. The uniform
comparison needed below is instead
\[
              c(q)\le\frac1{\sqrt8-1}<\frac9{16}
                  \quad(q\ge8),
       \qquad 625<648.
\tag{HES.13}
\]
The last integer inequality is equivalent to $25/9<\sqrt8$.

Every nontrivial irreducible $S$-module is faithful, and has degree
\[
                          a\ge q^{n-1}-1.
\tag{HES.14}
\]
Indeed embed $\operatorname{GL}_n(q)$ by
$A\mapsto\operatorname{diag}(A,A^{-\mathsf T})$, and in it take the row
subgroup
\[
 U_0=\left\{\begin{pmatrix}1&v\\0&I_{n-1}\end{pmatrix}:
                   v\in\mathbf F_q^{n-1}\right\}.
\]
Its Levi subgroup $\operatorname{GL}_{n-1}(q)$ is transitive on all
nontrivial complex additive characters of $U_0$. Faithfulness ensures
that one, hence all $q^{n-1}-1$, occur with the same positive multiplicity.
For $q=4$ we also have $a\ge378$, by restricting to the block subgroup
$\operatorname{Sp}_6(4)$: its minimum nontrivial degree $378$ is the
complete exact character-table assertion certified in D-ESC-01.
Restriction has a nontrivial constituent because the original module
is faithful.

\begin{lemma}[D-HES-04a]
Let $h\in S$ have odd prime order $p$ and stabilize a nonzero totally
isotropic subspace of its natural module. There is a nonzero
$h$-normalized elementary abelian $2$-subgroup $U\le S$ on which $h$
has no nonidentity fixed point, and $|U|\ge4$.
\end{lemma}

\begin{proof}
The construction of D-HES-02a works over $\mathbf F_q$, with its
underlying additive groups viewed as $\mathbf F_2$-spaces. We give the
details needed to avoid a field-linearity assumption on the resulting
$U$. Write $b$ for the natural symplectic form and $R$ for an invariant
isotropic space with basis $r_1,\ldots,r_s$.
If $A=h|_R\ne I$, define for every symmetric $B=(b_{ij})$ over
$\mathbf F_q$
\[
               N_B(x)=\sum_{i,j}b_{ij}b(r_j,x)r_i.
\]
Then $N_BN_C=0$, and $1+N_B$ is symplectic. These elements form an
elementary abelian $2$-group; conjugation by $h$ sends $B$ to
$ABA^{\mathsf T}$. This action is nontrivial, since
$B=vv^{\mathsf T}$ goes to $(Av)(Av)^{\mathsf T}$ and the diagonal
entries $v_i^2$ recover $v$ uniquely in a finite field of characteristic
two. Thus $A\ne I$ moves a rank-one matrix.

If instead $h$ fixes $R$ pointwise, choose $0\ne e\in R$ and $f_0$
with $b(e,f_0)=1$. The vector $f=\sum_{j=0}^{p-1}h^jf_0$ is fixed and
$b(e,f)=1$. Hence the natural module is the $h$-stable orthogonal sum
$\langle e,f\rangle\perp W_0$, with nontrivial action on $W_0$.
For $v\in W_0,c\in\mathbf F_q$ define
\[
 q(v,c)e=e,\qquad q(v,c)w=w+b(w,v)e,\qquad q(v,c)f=f+v+ce.
\]
These are symplectic and obey
\[
          q(v,c)q(w,d)=q(v+w,c+d+b(v,w)).
\tag{HES.15}
\]
Because $b$ is alternating and symmetric in characteristic two, this
is an abelian group of exponent two, of order $q^{2n-1}$. Conjugation
by $h$ is $q(v,c)\mapsto q(hv,c)$, a nontrivial action.

In either case the elementary abelian group is a semisimple
$\mathbf F_2\langle h\rangle$-module, since $p$ is odd. Take its sum of
nontrivial simple summands. This gives a nonzero invariant subgroup
$U$ with $C_U(h)=1$. It has dimension at least two over $\mathbf F_2$,
because $\operatorname{GL}_1(2)=1$. No assertion that $U$ is an
$\mathbf F_q$-subspace is necessary.
\end{proof}

\begin{lemma}[D-HES-04b]
For every nontrivial irreducible complex $S$-module of degree $a$,
\[
 \nu(\rho(x))\ge a/8\quad(q=4),\qquad
 \nu(\rho(x))>7a/32\quad(q\ge8),\qquad 1\ne x\in S.
\tag{HES.16}
\]
\end{lemma}

\begin{proof}
First take odd-prime-order $h$ stabilizing an isotropic space, and use
$U$ from D-HES-04a. All its nonidentity elements are unipotent. If
$t=\dim V^U$, averaging (HES.12) yields
\[
       t\le a\left(c(q)+\frac{1-c(q)}{|U|}\right).
\]
The dual $h$-action on the character group of $U$ has no nontrivial
fixed character; its nontrivial orbits all have length $p$. On an
orbit-sum of weight spaces the eigenvalues of $h$ occur uniformly,
because $h^p=1$. A specified eigenvalue therefore has multiplicity at
most $t+(a-t)/p$, giving
\[
 \nu(\rho(h))\ge(1-1/p)(a-t)
       \ge\frac23\frac34(1-c(q))a=\frac{1-c(q)}2a.
\tag{HES.17}
\]
For an involution, its two eigenvalue multiplicities give exactly
$\nu(\rho(h))=(a-|\chi(h)|)/2$, with the same lower bound.

It remains to treat an odd-prime-order $h$ which stabilizes no nonzero
totally isotropic subspace. Its primary decomposition has no pair of
distinct reciprocal irreducible polynomials: either member would give
an invariant totally isotropic space. Its $1$-eigenspace is zero,
since a nonzero such space is nondegenerate symplectic and contains
an invariant isotropic line. Every remaining primary polynomial is
self-reciprocal of degree $2d$.

On its primary summand set $E=\mathbf F_{q^{2d}}$ and
$E_0=\mathbf F_{q^d}$. The adjoint involution on $E$ is
$z\mapsto\bar z=z^{q^d}$. The nondegenerate trace pairing uniquely
defines $H$ by
\[
 b(av,w)=\operatorname{Tr}_{E/\mathbf F_q}(aH(v,w))
                       \quad(a\in E).
\]
Adjointness and symmetry give
$H(av,bw)=a\bar bH(v,w)$ and
$H(w,v)=\overline{H(v,w)}$, so this is a nondegenerate Hermitian form.
If the $E$-dimension were at least two, diagonalization and surjectivity
of $E^\times\to E_0^\times$ by the norm would produce an isotropic
vector. Its entire $E$-line is totally isotropic for $b$, a contradiction.
Every primary multiplicity is thus one. The centralizer is a product
of norm-one tori, so
\[
               |C_S(h)|=\prod_i(q^{d_i}+1)\le(q+1)^n,
                      \qquad\sum_i d_i=n.
\tag{HES.18}
\]
Column orthogonality and (HES.14) give
\[
          \frac{|\chi(h)|}{a}
             \le\frac{(q+1)^{n/2}}{q^{n-1}-1}<\frac12.
\tag{HES.19}
\]
To check this uniformly, at $n=4,q\ge4$ the right side is at most
$100/(63q)\le25/63<1/2$: use
$(q+1)^2\le25q^2/16$ and $q^3-1\ge63q^3/64$.
As $n$ increases its ratio decreases, since
$(q^n-1)/(q^{n-1}-1)>q>\sqrt{q+1}$.
Every nonidentity power of $h$ is again elliptic. Fourier inversion
therefore bounds every eigenvalue multiplicity by
$a(1/2+1/(2p))\le2a/3$, giving the stronger $\nu(\rho(h))\ge a/3$.

For an arbitrary nonidentity $x$, choose a nonidentity prime-order
power $h=x^k$. Each $x$-eigenspace lies in an $h$-eigenspace, so
$\nu(\rho(x))\ge\nu(\rho(h))$. Equations (HES.12)--(HES.13) and
(HES.17)--(HES.19) prove the assertion.
\end{proof}

\subsection*{D-HES-05. All compatible normalizer operators}

\begin{lemma}[D-HES-05a]
For $q\ge8$, every nonscalar finite-order linear operator $T$ normalizing
an irreducible $S$-module of degree $a$ has $\nu(T)>7a/64$.
\end{lemma}

\begin{proof}
The induced automorphism $\alpha$ of $S$ is nontrivial, since the
centralizer of an irreducible representation consists of scalars.
Choose $s\in S$ with $\alpha(s)\ne s$. Then
$[T,\rho(s)]=\rho(\alpha(s)s^{-1})$ is a nonidentity core operator;
faithfulness and centerlessness ensure it is nonscalar. For an
eigenvalue $\lambda$ of $T$ of maximum multiplicity,
\[
 \begin{split}
 \nu([T,\rho(s)])
 &\le\operatorname{rank}([T,\rho(s)]-I)\\
 &=\operatorname{rank}(T\rho(s)-\rho(s)T)\\
 &\le2\operatorname{rank}(T-\lambda I)=2\nu(T).
 \end{split}
\]
Apply D-HES-04b. This proof covers every automorphism, without a
restriction on the order or type of the outer automorphism.
\end{proof}

\begin{lemma}[D-HES-05b: the field-involution endpoint]
For $S=\operatorname{Sp}_{2n}(4)$, $n\ge4$, every nonscalar finite-order
linear normalizer operator on an irreducible degree-$a$ module satisfies
\[
                              \nu(T)\ge7a/64.
\tag{HES.20}
\]
\end{lemma}

\begin{proof}
The outer automorphism group is $C_2$, by
\cite[Section~2.2.2, Table~2.1, author p.~14]{HES:LiXia2022}.
The entrywise field automorphism $\sigma:z\mapsto z^2$ represents its
nontrivial element. To check the latter without a naming convention,
embed the determinant-one matrix
$\left(\begin{smallmatrix}\omega&1\\1&0\end{smallmatrix}\right)$
in a symplectic two-plane, where $\omega^2+\omega=1$.
Its natural trace changes from $\omega$ to $\omega^2$ under $\sigma$,
so $\sigma$ is not inner. Thus every automorphism is
$\operatorname{inn}(g)$ or $\operatorname{inn}(g)\sigma$.

In the inner case, $T$ is a scalar times $\rho(g)$ and D-HES-04b gives
$\nu(T)\ge a/8$. In the outer case put
$\alpha=\operatorname{inn}(g)\sigma$. If $g\sigma(g)\ne1$, then
$T^2$ is a scalar times the nonidentity core operator
$\rho(g\sigma(g))$, and
$\nu(T)\ge\nu(T^2)\ge a/8$. We may therefore suppose
$g\sigma(g)=1$. The natural semilinear operator $s=g\sigma$ then
satisfies $s^2=1$.

Here is an explicit symplectic descent for this last case. Write
$W_*=\operatorname{Fix}(s)$. For any natural vector $v$, both
\[
             a=v+s(v),\qquad b=\omega^2v+\omega s(v)
\]
belong to $W_*$, and $v=b+\omega a$. If $u+\omega v=0$ with
$u,v\in W_*$, applying $s$ and subtracting gives $v=0$, then $u=0$.
Consequently the underlying $\mathbf F_2$-space is
$W_*\oplus\omega W_*$, with $\dim_{\mathbf F_2}W_*=2n$; an
$\mathbf F_2$-basis of $W_*$ is an $\mathbf F_4$-basis of the natural
module. Its symplectic pairing is $\mathbf F_2$-valued, since
$b(sx,sy)=b(x,y)^2$, and is nondegenerate since $W_*$ spans over
$\mathbf F_4$. A symplectic $\mathbf F_2$-basis therefore conjugates
$s$ to the entrywise $\sigma$ by an element of $S$.
After this conjugation the operator $T$ implements $\sigma$ on $\rho(S)$.
Its square is scalar, so rescale it to $T^2=1$.

Use $U_0\cong\mathbf F_4^r$, $r=n-1\ge3$, from (HES.14).
Its characters are
\[
        \chi_a(v)=(-1)^{\operatorname{Tr}_{\mathbf F_4/\mathbf F_2}(a\cdot v)}
                         \quad(a\in\mathbf F_4^r).
\]
All nontrivial characters have the same multiplicity $m$, by the
$\operatorname{GL}_r(4)$ action. Put $N=4^r$ and let $t$ be the
trivial weight multiplicity, so $a=t+(N-1)m$. The unipotent bound gives
\[
 t\le a\left(\frac34+\frac1{4N}\right),
                 \qquad m\ge\frac{a}{4N}.
\tag{HES.21}
\]
On the dual group, $\sigma$ sends the parameter to its coordinatewise
square. Exactly $M=2^r$ parameters are fixed. The other $N-M$ weights
form pairs. On each paired sum, $T^2=1$ forces equal $+$ and $-$
multiplicities, both equal to $m$. Thus either eigenvalue of $T$ misses
at least $(N-M)m/2$ dimensions, and
\[
       \nu(T)\ge\frac{N-M}{2}m
            \ge\frac{1-2^{-r}}8a\ge\frac7{64}a.
\tag{HES.22}
\]
The argument bounds actual eigenspaces of this field operator; it does
not assume that any special eigenvalue bound survives an arbitrary
scalar convolution.
\end{proof}

\subsection*{D-HES-06. The complete even-characteristic higher-rank exclusion}

\begin{theorem}[D-HES-06]
An irreducible inducing image in the signature-$(3,3)$ sixfold-source
problem cannot have a nine-dimensional actual period image if it has
an irreducible normal subgroup whose quasisimple factors have quotients
$\operatorname{Sp}_{2n_i}(q_i)$ with $n_i\ge4$ and all $q_i$ even.
The $q_i$ may differ. Arbitrary compatible scalar extensions, outer
operators, factor permutations and larger inducing character fields
are allowed.
\end{theorem}

\begin{proof}
The trivial multipliers in this range identify the core with a direct
product of these simple groups. Write its irreducible module as
$V=\bigotimes_i V_i$, with nontrivial $V_i$ of degree $a_i$ and
$e=\prod_i a_i$. A normalizer permutes simple factors and is, up to
scalar, a tensor of intertwining maps followed by this permutation.
For a factor-preserving operator the normalized eigenvalue
multiplicities convolve; their maximum cannot exceed the maximum for
any nonscalar factor. D-HES-02--03 and D-HES-05 give the common support
bound $7e/64$.

For completeness, a nontrivial cycle of length $\ell\ge2$ on equal
degree-$a$ factors can be conjugated so that all twists occur on one
coordinate. Their product is diagonalizable: a power of the whole
operator is a tensor of these products, and finite order excludes a
nontrivial unipotent Jordan part. In an eigenbasis this is a weighted
cyclic permutation of $a^\ell$ words. There are $a$ constant words;
every other word orbit has length at least two, contributing at most
one copy of a specified eigenvalue. The maximum multiplicity is at
most $a+(a^\ell-a)/2\le3a^\ell/4$. Tensoring other cycles cannot
increase the normalized maximum. Such an operator has support at
least $e/4$, so again at least $7e/64$. Nonidentity scalars have actual
moving rank $e$.

By Gow's theorem \cite[p.~583 and Section~3]{HES:Gow1981}, all the
factor characters are real. Hence $V$ is self-dual and has a unique
line of invariant nondegenerate bilinear forms. The inducing group
$J$ acts on that line through a multiplier $\mu$. It is nontrivial:
otherwise all finite-order $J$-operators would have reciprocal
eigenvalue pairs and real trace, whereas the original imaginary
quadratic field is contained in the inducing character field. This
uses the proven field inclusion from D-BI-01, without assuming that
the two fields coincide.

The actual intermediate-cover variation has Euler formula
\[
                   6=2(h-1)e+\sum_{j=1}^r\operatorname{rank}(g_j-I).
\tag{HES.23}
\]
Every factor has degree at least $35$: use the verified $q=2$ endpoint,
the $q=4$ bound $378$, or (HES.14) for $q\ge8$.
As in D-HES-03, $h\ge2$ is impossible and $h=1$ has actual image
dimension at most six. A nine-dimensional image therefore requires
$h=0$, $r\ge12$ and total moving rank $2e+6$.
In a generating product-one tuple at least two inertias have
nontrivial multiplier. Their $1$-eigenspaces are totally isotropic
and they each have moving rank at least $e/2$. The other at least
ten have moving rank at least $7e/64$. Consequently
\[
               2e+6\ge e+10\frac7{64}e=\frac{67}{32}e,
                            \qquad e\le64.
\tag{HES.24}
\]
Two or more factors would give $e\ge35^2>64$. For a single factor,
$q=4$ gives $e\ge378$, and $q\ge8$ gives
$e\ge q^{n-1}-1\ge511$, again impossible. The remaining single
$q=2$ case is already excluded by D-HES-03, whose sharper support
bound forces $e\le24<35$. This exhausts every case.
\end{proof}

\paragraph{Verification boundary.}
The new universal arguments are the elementary $2$-subgroup construction,
the elliptic trace estimate, the normalizer commutator estimate and the
explicit $\mathbf F_4/\mathbf F_2$ symplectic descent with weight pairs.
Their proofs are written above. A separate exact finite-field program
checks the displayed identities and endpoint inequalities; it does not
replace the universal proofs. The classification and unipotent inputs
are cited with exact scopes. No unavailable Gluck--Magaard (1998)
statement is used, and the computed $\operatorname{Sp}_{12}(2)$ table
remains redundant.

\section{All odd-characteristic symplectic normal cores of rank three}
\label{sec:odd-c3-core}

\paragraph{Status and scope.}
This is a new derivation from the complete official character-table
data and the published automorphism and unipotent-character inputs
specified below. The generic data are from L\"ubeck's dissertation and
its official CHEVIE/GenericCharacterTables implementation; they are
not described as a separately refereed journal table. The new
coefficient and finite-spectrum certificates have actually passed.
The whole mathematical proof below has been reread by its author.
The parent and an independent peer have completed whole-proof,
source and code reviews, all passing; both have executed exact
certificates, and the peer has separately reconstructed all
generic inequalities, finite spectra and auxiliary field identities.
No reality, selfduality,
fixed determinant, or equality of inducing and original character
fields is assumed.

For a finite-order operator $T$ on a complex space of dimension $a$,
put
\[
 \nu(T)=a-\max_{\lambda}\dim\ker(T-\lambda I).
                                                               \tag{C3.1}
\]
This projective support is unchanged by a nonzero scalar multiple.
Ordinary moving rank is at least $\nu(T)$, and a nonidentity scalar
inertia has moving rank $a$. Two elementary inequalities used below
are
\[
 \nu(T)\ge\frac{a-|\operatorname{Tr}T|}{2},\qquad
 \nu(T)\ge\nu(T^j)\quad(j\in\mathbf Z).
                                                               \tag{C3.2}
\]
For the first, rotate an eigenvalue of largest multiplicity to one
and apply the reverse triangle inequality to the remaining
eigenvalues. The second follows by inclusion of each eigenspace of
$T$ in an eigenspace of $T^j$.

\begin{lemma}[D-C3-01: complete similitude data and restriction]
\label{D-C3-01}
Let $q\ge5$ be odd, $C=\operatorname{CSp}_6(q)$, and
$S=\operatorname{Sp}_6(q)$. Every nonlinear ordinary irreducible
character $\chi$ of $C$ satisfies
\[
 \chi(1)\ge q^3-1,\qquad
 |\chi(x)|\le\frac{\chi(1)}4\quad(x\notin Z(C)).
                                                               \tag{C3.3}
\]
Every nontrivial ordinary irreducible character $\rho$ of $S$,
of degree $a$, is a multiplicity-one constituent of an irreducible
$C$-character of degree $a$ or $2a$. Consequently
\[
 a\ge\frac{q^3-1}{2},\qquad
 |\rho(s)|\le a/4
 \quad(s\in S\setminus Z(S)\text{ semisimple}).
                                                               \tag{C3.4}
\]
If $\rho$ is invariant under the diagonal outer automorphism, it
extends to $C$ in the same degree. Every nonscalar operator on its
space implementing an inner-diagonal automorphism then has support
at least $3a/8$.
\end{lemma}

\begin{proof}
The original complete-table input is
\cite[Introduction, pp.~3--4; \S2.2, pp.~18--21;
\S9, pp.~78--79]{C3:Luebeck}, as implemented by the full
\texttt{CSp6.1} table in \cite{C3:GCT}.
Its $100$ character types include every ordinary irreducible
character, and its $90$ class types include every conjugacy class.
The first character type parametrizes the linear similitude
characters. The first class type parametrizes the entire center.
The remaining $99$ character and $89$ class types are retained in
\texttt{generic\_inputs.json}, with their original expressions.
The immutable original source hash is
\texttt{f5e68e9429054370df7376f174a095bc548cf1ecf7ce44f5446f1eaebba250dc}.

Here is the exact finite coefficient certificate for the universal
inequalities (C3.3). Every original table entry is a finite sum
$\sum_j p_j(q)\zeta^{f_j}$, with $p_j\in\mathbf Q[q]$ and each
phase a root of unity at every admissible parameter specialization.
Write $d(q)=\chi(1)$ for the degree polynomial and $q=5+x$. If
$p_j(5+x)=\sum_k c_{jk}x^k$, its absolute value for $x\ge0$ is at
most $\sum_k|c_{jk}|x^k$. For every one of the $8811$ pairs of a
nonlinear row and a noncentral column, the retained rational
coefficient list of
\[
 \frac{d(5+x)}4-\sum_{j,k}|c_{jk}|x^k
                                                               \tag{C3.5}
\]
is nonnegative. Similarly, every nonlinear degree minus $q^3-1$
has nonnegative coefficients after this substitution. Every
noncentral class-length polynomial has nonnegative coefficients
and constant term greater than one. The script
\texttt{check\_generic.py} parses the retained mathematical
expressions without executing the external Julia source and
checks all these coefficients using rational arithmetic. Its
complete lists are in \texttt{generic\_certificate.json}; this
proves (C3.3) for all admissible parameters and all odd $q\ge5$,
not only for finitely many values of $q$.

For the restriction statement, $C/S\simeq\mathbf F_q^\times$
is cyclic, and $C/(SZ(C))$ has order two: the multiplier of a
scalar matrix is its square. The inertia group of $\rho$ therefore
has index one or two in $C$. An irreducible representation of a
normal subgroup extends over a cyclic inertia quotient. Explicitly,
choose an intertwiner for a lift $t$ of its generator; its appropriate
power differs from $\rho(t^k)$ by a scalar, which can be removed by
rescaling. Inducing this extension to $C$ gives an irreducible
character whose restriction is the sum of its distinct conjugates,
each with multiplicity one. Thus its degree is $D=a$ or $2a$.
It is nonlinear because its restriction contains the nontrivial
irreducible $\rho$. The degree assertion in (C3.4) follows.

We verify carefully that semisimple $S$-classes do not split under
$C$-conjugacy. If $s\in S$ is semisimple, the multiplier map
\[
              C_C(s)\longrightarrow\mathbf F_q^\times
                                                               \tag{C3.6}
\]
is surjective. Indeed, use the primary decomposition of its
natural symplectic space. The $+1$ and $-1$ primary spaces are
nondegenerate symplectic spaces and admit a commuting similitude
of every multiplier. A pair of distinct reciprocal primary factors
gives $U\oplus U^*$, on which multiplication by $\mu$ on $U$
and by one on $U^*$ has multiplier $\mu$. For a self-reciprocal
irreducible factor of degree $2d$, the alternating pairing is the
trace of a sesquilinear form over
$\mathbf F_{q^{2d}}/\mathbf F_{q^d}$. Scalar multiplication by
an element of $\mathbf F_{q^{2d}}$ acts on the primary space;
the sesquilinear form is obtained uniquely from the original
pairing by nondegeneracy of the finite-field trace pairing.
The adjoint of this scalar action is the field involution.
Consequently scalar multiplication by
$z$ multiplies this form, and hence the original pairing, by
$z\bar z$. The finite-field norm is surjective, so choose
$z\bar z=\mu$ for any prescribed $\mu\in\mathbf F_q^\times$.
Choosing the same multiplier on all orthogonal summands proves
(C3.6), including primary spaces of any multiplicity. Given a
conjugator in $C$, compose it with a centralizer element of the
inverse multiplier to obtain a conjugator in $S$.

Therefore the one or two constituents of $\chi|_S$ have the same
value at every semisimple $s$, and
$\chi(s)=(D/a)\rho(s)$. This transfers (C3.3) unchanged to the
semisimple assertion in (C3.4). Finally, diagonal invariance makes
the inertia group all of $C$, so the extension has degree $a$.
Any other operator realizing the same automorphism differs from
the extension operator by a scalar, by irreducibility. A nonscalar
extension operator comes from a noncentral natural matrix, so
(C3.2)--(C3.3) give the support $3a/8$.
\end{proof}

\begin{lemma}[D-C3-02: all compatible normalizers]
\label{D-C3-02}
Let $q\ge3$ be odd and let $\rho$ be a nontrivial ordinary
irreducible representation of $\operatorname{Sp}_6(q)$ of degree
$a$. Every nonscalar finite-order linear operator $T$ normalizing
its represented image satisfies the following bounds:
\[
\begin{array}{c|c|c}
q&\text{lower bound for }a&\text{lower bound for }\nu(T)/a\\ \hline
3&13&1/4\\
q\ge5\text{ prime}&(q^3-1)/2&1/5\\
9&364&2/9\\
q\ge25\text{ composite}&(q^3-1)/2&3/16.
\end{array}                                                     \tag{C3.7}
\]
These bounds cover both possible central kernels of $\rho$ and
every compatible diagonal and field automorphism.
\end{lemma}

\begin{proof}
There is no exceptional Schur multiplier for
$\operatorname{PSp}_6(q)$ with odd $q$; see the complete list in
\cite[Chapter~5, Table~1, p.~35]{BI:HHM}.
Its universal perfect cover is $S=\operatorname{Sp}_6(q)$, with
center of order two. Thus every quasisimple representation in
question is already a representation of $S$, with kernel contained
in its center. Automorphisms of either represented central quotient
lift uniquely to $S$. An automorphism of $S$ trivial on its simple
quotient differs from the identity by a homomorphism from the
perfect group $S$ to its center, and hence is the identity.
The standard automorphisms are inner, diagonal, and field, and
\[
 \operatorname{Out}(\operatorname{PSp}_6(q))=C_2\times C_f,
 \qquad q=p^f;
                                                               \tag{C3.8}
\]
the exact source is \cite[\S2.2.2 and Table~2.1,
author p.~14]{HES:LiXia2022}. There is no graph automorphism in
this rank and characteristic.

We also use the published unipotent bound
\[
 \frac{|\rho(u)|}{a}\le\frac1{\sqrt q-1}
 \quad(q>4, u\ne1\text{ unipotent}),
                                                               \tag{C3.9}
\]
with the quasisimple-group conventions and exact radical in
\cite[pp.~2--4, Theorem~2.4(i)]{HES:Liebeck2017}.
This is the published restatement of Gluck's unipotent estimate.

\emph{Prime $q\ge5$.}
There are no field automorphisms. If the diagonal automorphism
stabilizes $\rho$, D-C3-01 gives the stronger support $3a/8$
on the entire compatible normalizer. If it does not, every
compatible normalizer operator is a scalar multiple of an inner
operator. A semisimple inner element has support at least $3a/8$.
An element with a nontrivial unipotent part has a power which is a
nonidentity unipotent of order $p=q$. Let $u$ be that power and
let $\chi$ be the $C$-character in D-C3-01, of degree $D\le2a$.
All $p-1$ nonidentity powers of $u$ are noncentral in $C$.
Fourier inversion and (C3.3) show that every eigenvalue of
$\chi(u)$ has multiplicity at most
\[
 \frac{D+(p-1)D/4}{p}
       =\frac{D(p+3)}{4p}
       \le\frac{a(p+3)}{2p}.
                                                               \tag{C3.10}
\]
Each eigenspace in the constituent $\rho$ is no larger.
Consequently $\nu(\rho(u))\ge a(p-3)/(2p)\ge a/5$.
Taking a power cannot increase the lower bound required for the
original element, by (C3.2). This proves the prime row of (C3.7)
without claiming that restriction is always irreducible.

\emph{Composite $q\ge25$.}
By (C3.9) every nontrivial unipotent has character ratio at most
$1/4$. The semisimple assertion of D-C3-01 and a nontrivial
unipotent power therefore give support at least $3a/8$ for every
noncentral inner element. For a nonscalar normalizer $T$, choose
$h\in S$ for which $[T,\rho(h)]\ne1$; such an $h$ exists by
Schur's lemma. If the commutator is noncentral, its support is
at least $3a/8$, while
\[
 \operatorname{rank}([T,\rho(h)]-I)\le2\nu(T).
                                                               \tag{C3.11}
\]
To check (C3.11), multiply $T$ by a scalar so that
$\operatorname{rank}(T-I)=\nu(T)$ and expand the commutator
difference as the difference of two conjugates of $T-I$.
If the commutator is a nonidentity scalar, its rank is $a$,
which is even stronger. This proves $\nu(T)\ge3a/16$ for
all compatible outer automorphisms, whatever their order.

\emph{$q=9$.}
The same inner argument, now using (C3.9) with value $1/2$,
gives support at least $a/4$ for every noncentral inner operator.
An operator with nontrivial diagonal outer class and trivial field
part, if compatible,
has support at least $3a/8$, by the same-degree extension in
D-C3-01. Inner operators retain the bound $a/4$ just proved.
It remains to handle an automorphism $\alpha$ involving the
field involution $\sigma:x\mapsto x^3$.

If $\alpha^2$ is a nontrivial inner automorphism, $T^2$ is a
scalar multiple of a noncentral inner operator, giving the
stronger bound $a/4$. Suppose instead $\alpha^2=1$. On the
natural space choose a semilinear lift $s=h\sigma$ with
$h\in C=\operatorname{CSp}_6(9)$. Unique lifting to the
universal perfect cover gives
\[
 h\sigma(h)=cI.
                                                               \tag{C3.12}
\]
Applying $\sigma$ and comparing $\sigma(h)h$ shows
$c\in\mathbf F_3^\times$. If $\mu$ is the multiplier of $h$,
then $\operatorname{Nm}_{9/3}(\mu)=c^2=1$.
The norm-one subgroup of $\mathbf F_9^\times$ is precisely
its subgroup of squares. Thus the mixed diagonal-field outer
class cannot have inner-trivial square. This observation is
not needed as a classification assumption in the descent that
follows, but checks that no exceptional mixed involution was
discarded.

Choose $z\in\mathbf F_9^\times$ with
$\operatorname{Nm}_{9/3}(z)=c^{-1}$ and put $s'=zs$.
Then $(s')^2=1$. Its semilinear multiplier is
$\beta=z^2\mu$, of norm one, so choose $d\in\mathbf F_9^\times$
with $d^2=\beta$. If $b$ is the standard symplectic form, then
$b_d=d b$ satisfies
\[
 b_d(s'x,s'y)=\sigma(b_d(x,y)).
                                                               \tag{C3.13}
\]
This follows from $d\beta=d^3$.
Let $\zeta\in\mathbf F_9$ satisfy $\zeta^2=-1$, so
$\sigma(\zeta)=-\zeta$. The fixed space of $s'$ spans the
natural space, since
\[
 x=x_0+\zeta x_1,\qquad
 x_0=\frac{x+s'x}{2},\qquad
 x_1=\frac{x-s'x}{2\zeta},
                                                               \tag{C3.14}
\]
and both $x_0,x_1$ are fixed. Moreover the fixed vectors form a
six-dimensional $\mathbf F_3$-space: a dependence among a fixed
$\mathbf F_3$-basis with coefficients in $\mathbf F_9$ splits
into two $\mathbf F_3$-dependences by (C3.14).
The restriction of $b_d$ is $\mathbf F_3$-valued by (C3.13)
and nondegenerate because this space spans. Choose a symplectic
basis of it. The resulting change of basis $P$ lies in
$\operatorname{CSp}_6(9)$, with multiplier $d^{-1}$, and
$P^{-1}s'P=\sigma$.

Relabel the core representation by this similitude conjugation.
It remains an irreducible representation of the same dimension;
no assertion of its invariance under diagonal conjugation is
needed. The operator $T$ now implements pure $\sigma$. Its
square is scalar, so rescale it to have $T^2=1$.
Inside the standard $\operatorname{GL}_3(9)$ Levi subgroup of
$S$ take
\[
 U=\left\{\begin{pmatrix}1&v\\0&I_2\end{pmatrix}:
              v\in\mathbf F_9^2\right\},\qquad |U|=N=81.
                                                               \tag{C3.15}
\]
The embedded $\operatorname{GL}_2(9)$ acts transitively on the
nonzero additive characters of $U$. Write $t$ for the dimension
of the trivial weight and $m$ for the common multiplicity of
the other $N-1$ weights. Every nonidentity element of $U$ is
unipotent. Averaging (C3.9) gives
\[
 t\le\frac a2+\frac a{2N},\qquad
              m=\frac{a-t}{N-1}\ge\frac a{2N}.
                                                               \tag{C3.16}
\]
Parametrize characters by $w\in\mathbf F_9^2$ using the
$\mathbf F_9/\mathbf F_3$ trace. Frobenius fixes exactly the
$M=9$ parameters in $\mathbf F_3^2$, and pairs the other
$72$. On each paired sum of weight spaces the involution $T$
has $+1$ and $-1$ multiplicities both equal to $m$.
Consequently
\[
 \nu(T)\ge\frac{N-M}{2}\,m
          \ge\frac{N-M}{4N}a=\frac{2a}{9}.
                                                               \tag{C3.17}
\]
The semilinear and dual-character identities are additionally
checked in \texttt{check\_f9.py}; the universal argument is
(C3.12)--(C3.17), not an extrapolation from that finite witness.

\emph{$q=3$.}
The complete retained universal-cover table
\texttt{2.S6(3)} and its cyclic extension
\texttt{2.S6(3).2} are the published ATLAS-origin tables in
\cite{C3:Small}. All $140$ nonlinear core rows and all $88$
extension rows whose restriction is irreducible are included.
Their complete nonscalar spectrum counts and extrema are
\[
\begin{array}{c|r|r|r|c}
 \text{table}&\text{rows}&\text{spectra}&\min a&\min\nu/a\\\hline
 \texttt{2.S6(3)}&140&19460&13&2/7\\
 \texttt{2.S6(3).2}&88&11968&78&8/21.
\end{array}                                                     \tag{C3.18}
\]
The GAP exporter checks full row orthogonality, retains every
original character value and prime power map, and computes all
spectra. A separate exact Python reconstruction selects extension
rows by their restriction inner products and recomputes every
spectrum by rational cyclotomic Fourier arithmetic. It also
checks the fusion, orders, power-map compatibility and complete
normal class unions. Both calculations passed all $31428$
spectra. In particular every nonscalar operator has support at
least $a/4$.

For completeness, the cyclic extension covers every compatible
outer operator, not merely the chosen abstract extension.
An outer automorphism differs from its chosen representative
by an inner automorphism. For that representative take an
intertwiner $A$. Its square differs from the represented
square of a lift in the fixed extension by a scalar, removable
by rescaling $A$. Thus the representation extends to that
fixed extension, and the original operator is scalar times an
element of it. This applies equally when the central involution
acts trivially. There are no other outer automorphisms by
(C3.8). The bounds for $q=3$, together with the other three
cases, exhaust all odd prime powers and prove (C3.7).
\end{proof}

\begin{theorem}[D-C3-03: the complete odd rank-three symplectic gate]
\label{D-C3-03}
Let a finite group $J$ act irreducibly on a complex space $V$.
Suppose it has a normal central product of quasisimple subgroups
acting irreducibly, each with simple quotient
$\operatorname{PSp}_6(q_i)$ for some odd $q_i\ge3$.
The field sizes and factor representations may differ. Every
compatible outer, scalar and factor-permutation action of $J$
is allowed. If this representation is an inducing block of a
finite-group curve-cover construction whose multiplicity
variation has complex dimension six, its period map to a
sixfold Weil domain of dimension nine has rank at most eight.
This holds for every polarization determinant and every
imaginary quadratic original field, even when the inducing
character field is larger.
\end{theorem}

\begin{proof}
Lift each quasisimple factor to its universal cover. Irreducibility
of the central product gives
$V=\bigotimes_i V_i$, with $a_i=\dim V_i\ge13$ and
$e=\dim V=\prod_i a_i$. By D-C3-02 every nonscalar compatible
factor normalizer has support at least $3a_i/16$.
For a factor-preserving normalizer on the tensor product, normalized
eigenvalue multiplicities are convolutions; their maximum cannot
exceed the maximum for any one nonscalar factor. Thus it has
support at least $3e/16$. If all factors are scalar the whole
operator is scalar. A genuine permutation of factors has the
stronger support $e/4$: on a cycle of length $\ell\ge2$ among
equal-dimensional factors, conjugate the tensor intertwiners
so that all twists except the last are identities and diagonalize
the last twist. The weighted rotation of basis words has at most
one vector of each eigenvalue per word orbit. There are $a$
constant words and every other orbit has size at least two,
so the number of orbits is at most
$(a^\ell+a)/2\le3a^\ell/4$. Tensoring with the other cycles
does not increase the normalized largest eigenspace. This also
covers permutations combined with outer actions. Every genuine
nonidentity scalar inertia moves all $e$ dimensions.

Use the actual intermediate finite cover associated to the
inducing block, as in D-BI-01. If its base has genus $h$ and
its faithful-image cover has $R$ genuine branch points, the
six-dimensional multiplicity and the absence of invariant
vectors give
\[
        6=2e(h-1)+D,\qquad
        D=\sum_{j=1}^R\operatorname{rank}(g_j-I).
                                                               \tag{C3.19}
\]
Its variation factors through a finite pointed Hurwitz cover of
$\mathcal M_{h,R}$. Since $e\ge13$, $h\ge2$ is impossible;
if $h=1$, then $D=6$, $R\le6$, and the actual period rank
is at most six. For rank nine one must therefore have
\[
                  h=0,\qquad R\ge12,\qquad D=2e+6.
                                                               \tag{C3.20}
\]
No sufficiency is inferred from this dimension inequality.
Every genuine branch entry has moving rank at least
$\lceil3e/16\rceil$, so
\[
        12\lceil3e/16\rceil\le2e+6,
             \qquad\text{hence}\qquad e\le24.
                                                               \tag{C3.21}
\]
Two or more factors would give $e\ge13^2=169$.
A single factor with $q\ge5$ gives
$e\ge(q^3-1)/2\ge62$. The only remaining case is a single
$q=3$ factor, for which the stronger quarter bound gives
$12\lceil e/4\rceil\le2e+6$, hence $e\le6$, contrary to
$e\ge13$. Thus rank nine is impossible in every stated case.
The argument used one actual complex multiplicity variation and
its intermediate Hurwitz factorization; it never identified an
inducing character field with the original quadratic field.
\end{proof}

\section{The complete transitive $A_7$-cover gate}
\label{sec:a7-cores}

\paragraph{Status.}
New root derivation. The complete proof and all source scopes were
read by the author and an independent reviewer; both checks passed.
The complete GAP spectrum computation, the separate exact Python
reconstruction and an independently written polynomial-Fourier
certificate passed. The independent audit is recorded in
\texttt{A7\_CORE\_PEER\_AUDIT.md}.
The original character tables are identified published inputs, not
newly constructed tables.

\begin{lemma}[D-A7-01: every compatible normalizer]
\label{D-A7-01}
Let $W$ be a nontrivial ordinary irreducible representation of a
quasisimple group with simple quotient $A_7$, and put $a=\dim W$.
If $W$ is not the deleted permutation representation of $A_7$,
every nonscalar finite-order operator normalizing its image satisfies
\[
                  \nu(T):=a-\max_\lambda\dim W_\lambda(T)
                            \ge a/4.                 \tag{A7.1}
\]
Every such representation has $a\ge4$. The only degrees at most six
are four and six, and in either of these nonstandard cases one has
the stronger integral bound $\nu(T)\ge2$.
\end{lemma}

\begin{proof}
The universal perfect cover is $\widetilde A_7=6.A_7$, with cyclic
center of order six. All central quotients are included by using all
ordinary characters of this cover. The complete ordinary-cover list
is stated in \cite[Chapter~4, opening paragraph, p.~31]{BI:HHM};
the ATLAS-derived metadata
\texttt{ExtensionInfoCharacterTable(CharacterTable("A7"))}
are exactly \texttt{["6","2"]}. Also
$\operatorname{Aut}(A_7)=S_7$. The exact official ATLAS table record
\texttt{A7.2} identifies both $S_7$ and
\texttt{AutomorphismGroup(AlternatingGroup(7))}; see
\cite{HD:A7Tables}. These are cited classification inputs.

Use the complete tables \texttt{6.A7} and \texttt{6.A7.2}, with
their original values, class sizes, element orders, power maps and
core fusion. Their input SHA-256 is
\texttt{185c2a90af296395c33ad8c7c5e839fadba5fffaeda3aacac2d99c8d20c906ab}.
Every nonlinear core row and every extension row with irreducible
core restriction is selected. The selection and exact spectra give
\[
\begin{array}{c|r|r|r}
\text{table}&\text{selected rows}&\text{nonscalar spectra}&\min a\\\hline
\texttt{6.A7}&39&1326&4\\
\texttt{6.A7.2}&18&558&6.
\end{array}                                                   \tag{A7.2}
\]
The GAP program checks full row orthogonality and reconstructs every
spectrum. The separate Python checker selects extension rows by the
exact norm of their restricted character, verifies the full fusion,
class-union and power-map compatibility, and recomputes all $1884$
spectra using rational cyclotomic Fourier arithmetic. In more detail,
if $x$ has order $d$, every multiplicity is the integer
\[
             m_j=\frac1d\sum_{r=0}^{d-1}
                    \chi(x^r)\exp(-2\pi i jr/d).          \tag{A7.3}
\]
The Python routine computes its rational field trace with exact
Ramanujan sums; integrality and the sum of all multiplicities are
checked. Every original nonscalar entry is retained, rather than
only entries failing a preliminary character-ratio test.

All these spectra satisfy $4(a-\max m_j)\ge a$, except precisely
extension rows three and four on class~26. Both exceptional rows
have degree six and restrict to core row two; their spectra are
$(5,1)$ and $(1,5)$ up to the choice of initial eigenvalue. Core
row two is the unique degree-six row trivial on all six central
classes. It is the deleted permutation representation: the natural
$A_7$ action is doubly transitive, so its deleted module is
irreducible of degree six. The complete core table gives no other
center-trivial row of that degree. The two degree-four rows have
minimum projective support two, and every nonstandard degree-six
row has minimum support at least two. None of them is the
irreducible restriction of an extension row. The only nonlinear
core degrees at most six are four and six. These assertions are
explicitly tested in \texttt{research/a7\_cores/check.py} and
recorded for every row in \texttt{checked.json}.

It remains to justify that these two tables cover every compatible
linear normalizer, including scalar twists and other abstract
extensions. An automorphism of a represented central quotient
lifts uniquely to its universal perfect cover. An automorphism
of that cover which is trivial on $A_7$ differs from the identity
by a homomorphism from a perfect group to its center; it is therefore
the identity. If the induced automorphism is inner, Schur's lemma
makes $T$ a scalar multiple of a represented core element.

Otherwise adjust by an inner operator to the automorphism induced
by a fixed lift $t$ in the chosen group $6.A_7.2$. A compatible
intertwiner $A$ has $A^2$ equal to $\rho(t^2)$ times a scalar.
Rescale $A$ to remove that scalar. This gives an actual extension
of $W$ to the chosen group, with irreducible core restriction.
The original $T$ is a scalar times an element of that extension.
This argument does not require the chosen extension to split and
also handles an isoclinic extension with a different scalar square.
Projective support is unchanged by the scalar. Thus the complete
table inequalities give exactly (A7.1) and its stated endpoints.
\end{proof}

\begin{lemma}[D-A7-02: the standard tensor endpoint]
\label{D-A7-02}
In the actual sixfold-source setting, an inducing block with
irreducibly acting normal core $A_7^t$ on
$\operatorname{Std}_7^{\otimes t}$, $t\ge2$, cannot have a
nine-dimensional period image. This permits all compatible outer
operators, scalars and permutations of the factors, and a larger
inducing character field.
\end{lemma}

\begin{proof}
Put $e=6^t$. The full normalizer, by $\operatorname{Aut}(A_7)=S_7$
and Schur's lemma on each factor, consists of scalar multiples
$g=\lambda T$ of the rational orthogonal tensor-wreath operators
$T\in\mathcal R_{7,t}$ defined in D-AST-01. That lemma was proved
for all $n\ge3$, and gives, for every integer $j$,
\[
                 \operatorname{tr}(T^j)\ge-e/6.
                                                               \tag{A7.4}
\]
Its elementary proof contracts the tensor indices around each
factor cycle; every resulting standard trace lies in
$\{-1,0,1,\ldots,6\}$. A negative product therefore has absolute
value at most $6^{t-1}$. This also applies to every power.

For any nonreal eigenvalue $\zeta$ of $T$ with multiplicity $m$,
its conjugate has the same multiplicity. Some power has
$\operatorname{Re}(\zeta^j)\le-1/2$. The trace of that power is
at most $e-3m$, so (A7.4) gives $m\le7e/18$.
The tensor standard form has multiplier $\mu(g)=\lambda^2$.
Consequently an element of nontrivial multiplier has moving rank
at least $11e/18$, because its fixed space is a nonreal eigenspace
of this particular rational $T$.

Every other genuine inertia has moving rank at least $e/6$.
A nonscalar factor matrix has projective support at least one
in degree six; convolution of normalized spectra preserves this
lower fraction on a factor-preserving tensor. A nontrivial factor
cycle has the stronger quarter bound, by the weighted basis-word
argument in D-SPOR-03. A nonidentity scalar moves all of $V$.

The multiplier character is nontrivial: otherwise the finite
inducing image would preserve a nondegenerate bilinear form,
have real character, and could not have a character field
containing the original imaginary quadratic field. This is the
field inclusion from D-BI-01, without assuming field equality.
The generating product-one branch tuple has at least two entries
with nontrivial multiplier.

The actual intermediate-cover Euler identity and period
factorization give a genus-zero base, at least twelve genuine
inertias and total moving rank $2e+6$. Genus at least two is
impossible for $e\ge36$; genus one has actual variation rank at
most six. Hence
\[
        2e+6\ge2\frac{11e}{18}+10\frac e6
                  =\frac{26e}{9},
                \qquad e\le\frac{27}{4}<36,             \tag{A7.5}
\]
contradicting $e=6^t\ge36$. This extends the already written
standard-tensor argument to its missing $n=7$ endpoint; it does
not assert that a nonreal-eigenspace bound survives scalar
convolution without first extracting the actual scalar phase.
\end{proof}

\begin{theorem}[D-A7-03: no transitive $A_7$-cover core in the nonsplit gate]
\label{D-A7-03}
An inducing block in the nonsplit Weil-sixfold problem cannot
have a nine-dimensional actual period image if its normal
quasisimple central-product core acts irreducibly, all simple
quotients are $A_7$, and the inducing group is transitive on those
factors. All central covers, compatible outer and scalar actions,
factor permutations and larger inducing character fields are
allowed.
\end{theorem}

\begin{proof}
Write the representation of the lifted core as
$V=\bigotimes_{i=1}^tW_i$, with $e=\prod_i\dim W_i$.
Transitivity implies that all factor characters agree up to
automorphisms of the universal cover. In particular, if one is
the ordinary deleted representation, then all are: this character
is invariant under $\operatorname{Aut}(A_7)$, and its central
kernel is intrinsic. For $t\ge2$, D-A7-02 excludes this case.
For $t=1$, the full normalizer is
$\mathbf C^*S_7$ in its deleted representation. D-BI-08 proves
that rank nine then forces the actual polarization to be split,
contrary to the present hypothesis. Its proof explicitly allows
the inducing character field to be larger initially.

Otherwise every factor is nonstandard. By D-A7-01, every
nonscalar factor-preserving normalizer on $V$ has projective
support at least $e/4$; normalized eigenvalue multiplicities
convolve. A nontrivial factor permutation has the same quarter
bound, and a nonidentity scalar moves all $e$ dimensions.
Since $e\ge4$, the same actual intermediate-cover formula forces
genus zero for rank nine, at least twelve branches and total
moving rank $2e+6$. Thus
\[
                  12\lceil e/4\rceil\le2e+6,
                              \qquad e\le6.             \tag{A7.6}
\]
If $t\ge2$, then $e\ge4^2>6$. If $t=1$, the only degrees
left are four and six. Each genuine nonscalar inertia in either
case has moving rank at least two by D-A7-01; a genuine scalar
has still larger moving rank. Thus $24\le2e+6\le18$, another
contradiction. These exhaust the transitive core cases.
\end{proof}

\paragraph{Geometry, framework and boundary.}
This removes the residual $A_7$ and its double, triple and sixfold
covers from the transitive normal-core gate forced by D-BI-19.
It does not assert a result for an arbitrary nontransitive product
of differently labelled factors. The argument is about actual
finite covers and their weight-one multiplicity variations; the
integer spectrum and tensor identities belong to the associated
representation ring. The exact scheme-spreading, SGA finite-cover,
lambda-ring, K-theory and etale-weight placements are those in the
framework section. It narrows the supplier-existence part of
condition A, cote16 scan050, without producing a cycle or changing
condition I, cote16 scan049 and cote44 Theorem4, or the
Cayley--Hamilton projector in cote16 scans048--049. No new
motivic $L$-value estimate or general negative Hodge-conjecture
conclusion follows. All existing detector constants are unchanged.

\section{An elementary Heisenberg carrier estimate}
\label{sec:heisenberg-carrier}

\paragraph{Status.}
This is a complete new elementary derivation. The author, parent,
and an independent sibling reviewer have read the entire proof and
passed it in the stated scope. Their exact checks are recorded in
CHECK\_EVENTS.tsv. No Weil-character formula or representation-degree
classification is a premise.
The statements concern ordinary complex representations.

\begin{lemma}[D-HOS-01: fixed-point-free symplectic carriers]
\label{D-HOS-01}
Let $q=p^f$ be odd, let $(W,B)$ be a symplectic
$\mathbf F_q$-space of dimension $2m\ge2$, and let
\[
 H=W\times\mathbf F_q,\qquad
 (v,z)(w,t)=(v+w,z+t+\tfrac12 B(v,w))
                                                               \tag{HOS.1}
\]
be its Heisenberg group. Let $g\in\operatorname{Sp}(W)$ have
prime order $\ell\ne p$ and satisfy $W^g=0$. Suppose that a
finite-order operator $T$ on an $a$-dimensional complex
$H$-module $V$ satisfies
\[
 T\rho(v,z)T^{-1}=\rho(gv,z).
                                                               \tag{HOS.2}
\]
Put $t=\dim V^H$ and
$\nu(T)=a-\max_\lambda\dim\ker(T-\lambda I)$. Then
\[
 \nu(T)\ge
 \left(1-\frac1\ell\right)(1-q^{-m})(a-t).
                                                               \tag{HOS.3}
\]
If moreover $|\operatorname{Tr}\rho(h)|\le ca$ for every
$h\in H\setminus\{1\}$, with $0\le c\le1$, then
\[
 \nu(T)\ge
 \left(1-\frac1\ell\right)(1-q^{-m})
       (1-q^{-(2m+1)})(1-c)a.
                                                               \tag{HOS.4}
\]
Neither $\rho$ nor its restriction to $H$ is assumed irreducible,
and $T^\ell$ is not required to be scalar on $V$.
\end{lemma}

\begin{proof}
We give the carrier construction and every multiplicity step.
Choose a symplectic decomposition $W=X\oplus Y$ identifying
$X$ and $Y$ with $\mathbf F_q^m$, so
\[
 B((x,y),(x',y'))=x\cdot y'-x'\cdot y.
\]
Fix a nontrivial additive character $\psi:\mathbf F_q\to\mathbf C^*$.
On the $N=q^m$-dimensional function space
$U_\psi=\{f:\mathbf F_q^m\to\mathbf C\}$ define
\[
 (L_{x,y}f)(u)
   =\psi(y\cdot u+\tfrac12 x\cdot y)f(u+x),\qquad
 R_\psi((x,y),z)=\psi(z)L_{x,y}.
                                                               \tag{HOS.5}
\]
Direct multiplication gives
$L_vL_w=\psi(B(v,w)/2)L_{v+w}$, so (HOS.5) is a
representation of (HOS.1). Its trace is zero at $L_v$ for
$v\ne0$: a nonzero translation has no diagonal entries,
whereas for $x=0,y\ne0$ the sum of the nontrivial additive
character $u\mapsto\psi(y\cdot u)$ is zero.
It follows that
\[
        \operatorname{Tr}(L_vL_w^{-1})=N\delta_{v,w}.
                                                               \tag{HOS.6}
\]
Thus the $q^{2m}=N^2$ operators $L_v$ form a basis of
$\operatorname{End}_{\mathbf C}(U_\psi)$, and this
representation is irreducible. It is the only irreducible
representation with central character $\psi$: the corresponding
central-character summand of $\mathbf C[H]$ has dimension
$|H|/q=q^{2m}$ and maps onto this full matrix algebra of
the same dimension. The other irreducibles, with trivial
central character, are exactly the linear characters of
$H/Z(H)=W$.

The automorphism $(v,z)\mapsto(gv,z)$ preserves $\psi$.
Uniqueness therefore supplies an intertwiner $\Omega$ on
$U_\psi$ such that
\[
                  \Omega L_v\Omega^{-1}=L_{gv}.
                                                               \tag{HOS.7}
\]
Schur's lemma makes $\Omega^\ell$ scalar; rescale $\Omega$
so that $\Omega^\ell=I$. Equation (HOS.7) is exact, not
an equality with additional phases, because it implements
the specified automorphism on the actual group elements
$(v,0)$ in (HOS.1).
For $1\le j<\ell$, the permutation of the basis
$\{L_v\}$ induced by conjugation with $\Omega^j$ has
exactly one fixed vector, $L_0$. Indeed, $W^{g^j}=W^g=0$
because $\ell$ is prime. Hence
\[
 |\operatorname{Tr}\Omega^j|^2
  =\operatorname{Tr}(\operatorname{Ad}\Omega^j)
  =|W^{g^j}|=1.
                                                               \tag{HOS.8}
\]
The first equality uses finite order: the eigenvalues of
$\Omega$ are roots of unity, so the trace of its inverse
is the complex conjugate of its trace. Fourier inversion
on the cyclic group of order $\ell$ now bounds every
eigenvalue multiplicity of $\Omega$ by
\[
                  \frac{N+\ell-1}{\ell}.
                                                               \tag{HOS.9}
\]
Its projective support is consequently at least
$(1-1/\ell)(1-N^{-1})N$.

Decompose $V$ into its $H$-isotypic spaces. Every summand
with nontrivial central character has the form
$U_\psi\otimes M_\psi$. Since $g$ fixes the center,
it is $T$-stable, and (HOS.7) plus the full matrix-algebra
commutant gives
\[
                 T=\Omega\otimes A_\psi
                   \quad\hbox{on }U_\psi\otimes M_\psi.
                                                               \tag{HOS.10}
\]
The operator $A_\psi$ has finite order after the indicated
normalization: some power divisible by $\ell$ kills both
$T$ and $\Omega$. It is therefore diagonalizable. The
normalized eigenvalue multiplicities of a tensor product
are convolutions, so their maximum is no greater than
the maximum for $\Omega$. The support fraction from
(HOS.9) holds on every such isotypic summand, with no
assumption about the eigenvalues on $M_\psi$.

For the linear characters, $g$ has no fixed nonzero dual
weight because $g-1$ is invertible on $W$. Each nonzero
dual orbit therefore has length $\ell$, with all weight
spaces in it of equal multiplicity. On its direct sum,
$T$ is a cyclic block operator. Its $\ell$th power has
the same spectrum on all blocks, and each eigenvalue
of the block product gives at most one copy of any
chosen eigenvalue of $T$ per multiplicity. Equivalently,
conjugate the cyclic block operator to one with identity
intertwiners except for a final diagonal block.
Every eigenspace has dimension at most $1/\ell$ of
that orbit summand. This support fraction $1-1/\ell$
is at least the one in (HOS.9).

The remaining trivial $H$-isotypic space has dimension
$t$ and is $T$-stable. Give it the worst possible full
eigenspace of dimension $t$. Adding the preceding
eigenspace upper bounds over all other $T$-stable
summands proves (HOS.3). Finally, the averaging projector
onto $V^H$ gives
\[
 t=\frac1{|H|}\sum_{h\in H}\operatorname{Tr}\rho(h)
       \le a\left(c+\frac{1-c}{q^{2m+1}}\right).
                                                               \tag{HOS.11}
\]
Substitution into (HOS.3) proves (HOS.4).
\end{proof}

\begin{lemma}[D-HOS-02: an elementary symplectic degree bound]
\label{D-HOS-02}
Let $q$ be odd and $n\ge4$. Every nontrivial ordinary irreducible
representation of $\operatorname{Sp}_{2n}(q)$ has degree
\[
                     a\ge\frac{q-1}{2}\,q^{n-1}.
                                                               \tag{HOS.12}
\]
\end{lemma}

\begin{proof}
The line-parabolic unipotent radical is the Heisenberg group
in (HOS.1) with $m=n-1$. Its center is the additive group
$\mathbf F_q$ of long-root unipotents. The kernel of a
nontrivial irreducible representation of the quasisimple
symplectic group is central, so this noncentral unipotent
group cannot act trivially. At least one nonzero additive
central character therefore occurs. By the explicit group
algebra argument in (HOS.5)--(HOS.6), its isotypic space
has dimension at least $q^{n-1}$.

The line-parabolic torus sending a chosen isotropic vector
to $s$ times itself and its paired dual vector to $s^{-1}$
times itself acts on the center by multiplication by $s^2$.
Parametrize the additive characters by
$\psi_w(z)=\exp(2\pi i\operatorname{Tr}_{q/p}(wz)/p)$.
The orbit of each $w\ne0$ under all squares in
$\mathbf F_q^\times$ has size $(q-1)/2$. Every character
in this orbit must occur, with the same isotypic dimension.
Summing these dimensions proves (HOS.12). The symplectic
quasisimplicity used here is the standard ordinary-group
input; no exceptional cover occurs in the stated range
by the complete list in \cite[Chapter~5, Table~1, p.~35]{BI:HHM}.
\end{proof}

\section{Complete finite odd-$C_2$ endpoints}

\paragraph{Status.}
D-OC2-F has passed complete author, root and independent peer proof
readings, together with separate exact spectrum certificates. Its limited coverage
is explicit: the full universal cover at $q=5$, and center-positive
representations only at $q=7,9$. No generic odd-$C_2$ table is used.

\begin{lemma}[D-OC2-F1: exact full-normalizer endpoint bounds]
Let $V$ be a nontrivial irreducible complex representation of one of
$\operatorname{Sp}_4(5)$, $\operatorname{PSp}_4(7)$ or
$\operatorname{PSp}_4(9)$. Every finite-order nonscalar operator
normalizing its image has projective support at least $\dim(V)/4$.
The minimum degrees are respectively $12,25,41$.
\end{lemma}

\begin{proof}
The retained complete ordinary table input consists of
\[
\begin{gathered}
\texttt{2.S4(5)},\quad\texttt{2.S4(5).2},\qquad
\texttt{S4(7)},\quad\texttt{S4(7).2},\\
\texttt{S4(9)},\quad\texttt{S4(9).2\_1},\quad
\texttt{S4(9).2\_2},\quad\texttt{S4(9).2\_3}.
\end{gathered}
These are official complete CTblLib tables \cite{OC2:CTbl}. The two $q=5$ tables have
ATLAS provenance. The $q=7,9$ tables are computations of Thomas Breuer
using Dixon's algorithm and character-theoretic methods, dated 2001
and 2004 in the retained source strings; they are not mislabeled as
ATLAS computations. The full character values, class sizes and orders,
prime power maps, normal fusions, and provenance strings are retained.
Their input SHA-256 is
\texttt{6824ca705e8619733ced3c7aceab6e145b7b45735fa56a9990ed7607eb6f7a73}.

The GAP exporter verifies full row orthogonality. A separate exact
cyclotomic-polynomial program selects every irreducible core row and,
in an extension, every row restricting irreducibly to that core. It
reconstructs every eigenvalue multiplicity by the Fourier formula
\[
       m_j(x)=\frac1d\sum_{k=0}^{d-1}
                  \chi(x^k)\exp(-2\pi i j k/d),\qquad d=|x|.
\]
Scalar entries are identified by $|\chi(x)|^2=\chi(1)^2$. Exact
arithmetic on all 485 selected rows and 37502 nonscalar entries gives
no violation of the quarter bound. Normal class unions, class orders,
group orders and fusion compatibility with all retained prime power
maps are checked as well. The minimum support ratios over the
respective core and extension tables are
\[
                     4/13,\qquad9/25,\qquad16/41.
\]
The minimum nontrivial core degrees are $12,25,41$.
The complete verifier and its previously reviewed cyclotomic helper are
retained; no floating-point approximation enters these assertions.

We explain why these tables cover the claimed normalizers. The outer
groups at $q=5,7$ are cyclic of order two, and at $q=9$ they are
$C_2\times C_2$, by the symplectic row of
\cite[Section~2.2.2, Table~2.1, author p.~14]{HES:LiXia2022}.
The latter three cyclic subgroups are exactly the field, diagonal and
diagonal-field extensions identified in the official source strings.
Every individual operator belongs, projectively, to one of these
cyclic extension cases; a representation of the whole outer group is
not needed. At $q=5$ the multiplier-two universal cover and its outer
extension are both retained, so center-negative as well as
center-positive modules are covered.

More explicitly, an intertwiner for a compatible automorphism is unique
up to scalar by Schur's lemma. Choose an element of the listed extension
inducing the same automorphism, after multiplying by a core element if
necessary. Its cyclic power relation differs from the intertwiner's
power by a scalar; rescaling the intertwiner solves that relation over
$\mathbf C$. Thus it supplies an irreducible extension character whose
restriction is the given core module. Changing the cyclic lift or its
central power relation changes only scalar factors in the spectra.
This is why one listed covering extension suffices also for the
isoclinic alternative at $q=5$. All selected rows were scanned, so the
operator has the asserted support. A core representation factoring
through its central quotient remains covered by the same input.
\end{proof}

\begin{theorem}[D-OC2-F2]
An inducing image in the signature-$(3,3)$ sixfold-source problem with
an irreducible normal core formed from the covered groups in D-OC2-F1
cannot have an actual period image of dimension nine. Arbitrary
compatible scalar/outer extensions and permutations of tensor factors
are allowed. The claim for the $q=7,9$ factors is restricted to
center-positive modules.
\end{theorem}

\begin{proof}
The core representation is a tensor product of its factor modules.
For a factor-preserving operator, normalized eigenvalue multiplicity
distributions convolve, and their maximum is at most that of any
nonscalar factor. A nontrivial cycle of length $\ell$ on degree-$a$
factors is a weighted permutation of $a^\ell$ basis words, with $a$
constant words and all other orbits of length at least two. Its largest
eigenspace is at most $a+(a^\ell-a)/2\le3a^\ell/4$.
Thus all nonscalar operators in the full normalizer have support at
least $e/4$, where $e$ is the total degree. A nonidentity scalar has
actual moving rank $e$.

The actual intermediate-cover variation has
$6=2(h-1)e+\sum_j\operatorname{rank}(g_j-I)$.
Since $e\ge12$, genus $h\ge2$ is impossible; genus one has at most
six effective branch points and hence actual period-image dimension
at most six. A nine-dimensional image would require genus zero and
at least twelve branch points. Its moving-rank sum $2e+6$ would then
be at least $12e/4=3e$, forcing $e\le6$, a contradiction.
No self-duality, determinant or equality-of-character-fields assumption
is used in this argument.
\end{proof}

\paragraph{Exact remaining coverage gate.}
The executed inventory found no stored complete universal-cover tables
for $\operatorname{Sp}_4(7)$ or $\operatorname{Sp}_4(9)$, and no stored
complete odd-$C_2$ tables at $q=11,13,17,19,23,25,27$.
The preliminary full generic odd-$C_2$ table is known to contain errors,
and the available generic unipotent table is not a complete ordinary
table \cite{OC2:GCT}. Neither is used or extrapolated here. The already reviewed Weil
module and large-field exclusions retain their separate scopes.

\section{All odd rank-two symplectic cores above the small endpoint}
\label{sec:odd-c2-all-fields}

\paragraph{Status and exact scope.}
This complete derivation and its cited input scopes have passed whole
author, root and independent peer readings. The corrected unipotent
input and finite spectra have separate exact certificates.
The core factors in the final theorem are permuted transitively;
this is the normal-core configuration furnished by D-BI-19 for the
primitive, factorwise tensor-indecomposable source gate. No assertion
about an arbitrary intransitive mixture is substituted for that
hypothesis. The original inducing problem has a nonselfdual complex
irreducible representation, in particular one with exact imaginary
quadratic character field. Its inducing field may be larger.

Write $\nu(T)=a-\max_\lambda\dim\ker(T-\lambda I)$ for projective
support on an $a$-dimensional space. Scalar multiplication preserves it,
ordinary moving rank is at least it, and $\nu(T)\ge\nu(T^j)$.
Every nontrivial irreducible representation of $\operatorname{Sp}_4(q)$,
$q\ge5$, has projective kernel exactly its center: that kernel is
normal, its image in the simple central quotient is either trivial
or the whole quotient, and the latter would make a perfect group
act by scalars and hence trivially.

\begin{lemma}[D-OC2-01: the corrected unipotent input]
\label{D-OC2-01}
For odd $q\ge7$, every nontrivial ordinary irreducible character $\chi$
of $S=\operatorname{Sp}_4(q)$, of degree $a$, satisfies
\[
 a\ge(q^2-1)/2,\qquad |\chi(u)|\le a/2
       \quad(1\ne u\in S\text{ unipotent}).                 \tag{OC2.1}
\]
Apart from the four characters of degrees $(q^2+1)/2$ and
$(q^2-1)/2$, every such degree is at least $q(q-1)^2/2$.
For $q\equiv3\pmod4$ those four characters are nonselfdual.
\end{lemma}

\begin{proof}
We use the published corrected table of values of \emph{all ordinary
characters on unipotent elements} in
\cite[pp.~620--624, Table~2 and the restriction formulas]{OC2:Dabbaghian2006}.
This is a different input from an incomplete generic full table or a
table of only unipotent characters. There are six nonidentity
unipotent classes, denoted $A_{21},A_{22},A_{31},A_{32},A_{41},A_{42}$,
and $43$ nontrivial character types. Some parameterized types can be
empty at particular small $q$, which does not impair an upper bound
checked on every type. The degree column and these six values are
retained in \texttt{unipotent\_data.json}.

Here are the mathematical inputs and the exact coefficient procedure.
Put $t^2=(-1)^{(q-1)/2}q$ and $\epsilon=(-1+t)/2$,
$\epsilon'=(-1-t)/2$. Swapping the sign of $t$ exchanges paired
character labels. The six basis characters in the published
restriction formulas are $a_1,a_2,b,c,d_1,d_2$, with $a_2,d_2$
obtained from $a_1,d_1$ by that exchange. In the column order
$(1,A_{21},A_{22},A_{31},A_{32},A_{41},A_{42})$, the other four rows are
\begin{align*}
a_1={}&\bigl((q^2+1)/2,(q+1)/2+q\epsilon,
 (q+1)/2+q\epsilon',(q+1)/2,(1-q)/2,-\epsilon',-\epsilon\bigr),\\
b={}&\bigl(q(q-1)^2/2,q(1-q)/2,q(1-q)/2,0,q,0,0\bigr),\\
c={}&\bigl(q(q^2+1)/2,q(q+1)/2,q(q+1)/2,0,q,0,0\bigr),\\
d_1={}&\bigl((q-1)^2(q^2+1)/2,
 (1-q)/2-q(q-1)\epsilon',
 (1-q)/2-q(q-1)\epsilon,
 (1-q)/2,(1-3q)/2,-\epsilon,-\epsilon'\bigr).
                                                               \tag{OC2.2}
\end{align*}
The complete integer coefficient vectors in this six-character basis
and the trivial character are retained. Four rows are instead read
\emph{directly from Table~2}, as follows, with each conjugate obtained
by exchanging $t$ and $-t$:
\begin{align*}
\xi_{41}={}&\bigl((q+1)^2(q^2+1)/2,
 (q+1)/2-q(q+1)\epsilon,(q+1)/2-q(q+1)\epsilon',
 (1+3q)/2,(q+1)/2,-\epsilon,-\epsilon'\bigr),\\
-\xi'_{41}={}&\bigl((q^4-1)/2,
 -(1-q)/2+q(q+1)\epsilon,-(1-q)/2+q(q+1)\epsilon',
 -(1-q)/2,-(1+q)/2,\epsilon,\epsilon'\bigr).
                                                               \tag{OC2.3}
\end{align*}
The indexed text of the four corresponding restriction identities
interchanges $d_1,d_2$ relative to Table~2. We do not use those
inconsistent identities. The initial failed comparison and the direct
Table~2 overrides are retained explicitly; no correction is inferred
merely by fitting finitely many values of $q$.

Every resulting entry has the form $A(q)+B(q)t$ with rational
polynomials $A,B$. For $q\ge7$,
\[
        |t|=\sqrt q\le(q-1)/2,                                \tag{OC2.4}
\]
because $(q-1)^2-4q$ has positive coefficients after $q=7+x$.
Expand $A(7+x),B(7+x)$ and bound each by the polynomial obtained by
taking the absolute values of its coefficients. Multiply the second
bound by $3+x/2$. For all $43\cdot6=258$ entries, subtracting this
upper bound from $a(7+x)/2$ gives a polynomial with nonnegative
rational coefficients. The same exact calculation checks the degree
floor $(q^2-1)/2$, and the stronger floor $q(q-1)^2/2$ away from the
four specified types. The complete coefficient lists and the safe
arithmetic program are retained in \texttt{unipotent\_checked.json}
and \texttt{check\_unipotent.py}. This proves the inequalities for
all $q\ge7$, not just for sampled fields.

As a separate transcription check, the published formulas agree on
all six unipotent classes with all $60$ rows of the complete
$\operatorname{Sp}_4(5)$ table, all $52$ rows of the complete
$\operatorname{PSp}_4(7)$ table, and all $74$ rows of the complete
$\operatorname{PSp}_4(9)$ table, using a single column permutation
for each table. This is an executed exact cyclotomic comparison;
the two latter comparisons do not purport to supply missing
center-negative tables. The source's averaging identities for all six
basis characters provide another exact check on the positions of the
zero entries in (OC2.2).

Finally the four small characters are $a_1,a_2,a_1-1,a_2-1$ on the
unipotent subgroup. When $q\equiv3\pmod4$, their $A_{21}$ values
have a nonzero imaginary coefficient of $\sqrt{-q}$, so none is
selfdual. The four small characters are the four usual symplectic
Weil characters: the two nontrivial additive-character square classes
each give the two parity constituents of degrees $(q^2\pm1)/2$,
as in the already checked D-SPW-01, and this degree list contains
exactly those four characters.
\end{proof}

\begin{lemma}[D-OC2-02: every noncentral inner operator]
\label{D-OC2-02}
For $q\ge7$ odd and $\rho$ as above, every noncentral element
$x\in\operatorname{Sp}_4(q)$ satisfies
\[
                   \nu(\rho(x))\ge \frac{q-1}{4q}\,a.
                                                               \tag{OC2.5}
\]
\end{lemma}

\begin{proof}
Take a power $h$ whose image modulo $\{\pm I\}$ has prime order
$\ell$. Its projective support is no greater than that of $x$.
If $\ell=p$, multiplication by a central sign makes $h$ a nonidentity
unipotent, and (OC2.1) gives support at least $a/4$. Suppose
$\ell\ne p$. Then $h$ is semisimple. If $\ell$ is odd, multiply
by a central sign so that $h$ has actual order $\ell$. If $\ell=2$,
then $h^2=I$ or $h^2=-I$.

We first describe a useful invariant isotropic-subspace case. Let
$R$ be a nonzero totally isotropic $h$-stable natural subspace and
suppose conjugation by $h$ acts nontrivially on the symmetric shear
group attached to $R$. For a basis $r_1,\ldots,r_d$ of $R$ and a
symmetric $d\times d$ matrix $B$, put
\[
 N_B(v)=\sum_{i,j}b_{ij}\langle r_j,v\rangle r_i,
                    u_B=I+N_B.                                \tag{OC2.6}
\]
Symmetry and isotropy give $N_BN_C=0$ and
$\langle N_Bv,w\rangle+\langle v,N_Bw\rangle=0$.
Consequently $u_B$ is symplectic and $u_Bu_C=u_{B+C}$.
If $A=h|_R$ in this basis, conjugation acts by $B\mapsto ABA^t$.
The map $B\mapsto N_B$ is injective, because the pairings with
the independent $r_j$ are independent linear forms.

This conjugation action has order $\ell$: the central square of
$h$, if any, acts trivially on symmetric matrices. Since
$\ell\ne p$, its nonzero image under $h-1$ is an $\mathbf F_q$
subspace on which $h$ has no fixed vectors. It supplies an abelian
unipotent subgroup $U$ of order $Q\ge q$. Each nontrivial dual
weight has an orbit of length $\ell$. The invariant weight has
dimension at most $a/2+a/(2Q)$ by averaging (OC2.1), and every
other orbit contributes at most $1/\ell$ of its dimension to any
chosen eigenspace of $\rho(h)$. This is also valid for a cyclic
block operator with nontrivial block product. Hence
\[
 \nu(\rho(h))\ge(1-1/\ell)(1-1/Q)a/2
                    \ge(q-1)a/(4q).                          \tag{OC2.7}
\]

Both projective involution cases admit such an $R$. If $h^2=I$
and $h\ne\pm I$, its plus and minus eigenspaces are nonzero
orthogonal symplectic planes. Choose a line in each; their sum is
an isotropic plane, and its mixed symmetric shear has eigenvalue
$-1$. If $h^2=-I$ and $-1$ is a square, either eigenplane is
isotropic and its symmetric shears have eigenvalue $-1$. If
$-1$ is not a square, regard the natural space as dimension two
over $E=\mathbf F_{q^2}$, with $h$ multiplication by a square root
of $-1$. The alternating form induces, through the nondegenerate
trace pairing, a nondegenerate hermitian or skew-hermitian form
over $E/\mathbf F_q$. Scaling makes it hermitian. Every such
two-dimensional form over a finite field is isotropic: after
diagonalization this follows by solving one scalar norm equation,
using surjectivity of $E^*\to\mathbf F_q^*$. Its isotropic $E$-line
is an isotropic $\mathbf F_q$-plane $R$. On its symmetric square,
$h$ has a nontrivial minus-one eigenspace, as is seen over $E$
from eigenvalues $i^2=-1$ and $(-i)^2=-1$.

Now let $\ell$ be odd. On a one- or two-dimensional space, a
semisimple $A$ acts trivially on the symmetric square only if
$A=\pm I$: over an algebraic closure its squared eigenvalues
and pairwise products must all equal one. Thus any invariant
isotropic $R$ on which $h$ acts nontrivially is handled by
(OC2.6)--(OC2.7), since $-I$ cannot have odd order.

If an invariant isotropic subspace exists but $h$ fixes it
pointwise, choose a fixed vector $e$ and a fixed paired vector
$f$ with $\langle e,f\rangle=1$; average a paired vector over
the order-$\ell$ group to obtain $f$. The orthogonal complement
$W=\langle e,f\rangle^\perp$ is a symplectic plane. The action
of $h$ on it is nontrivial, and has no fixed vectors; otherwise
its semisimple determinant-one action would be the identity.
The line-parabolic radical consists of the maps
\[
 e\mapsto e,\qquad
 w\mapsto w+\langle v,w\rangle e,\qquad
 f\mapsto f+v+2ze,
             \quad v\in W,\ z\in\mathbf F_q.                  \tag{OC2.8}
\]
Its multiplication is $(v,z)(v',z')=(v+v',z+z'+
\langle v,v'\rangle/2)$, and $h$ acts by $(v,z)\mapsto(hv,z)$.
Every nonidentity element is unipotent. Applying the fully proved
Heisenberg carrier lemma D-HOS-01, with $m=1$, gives
\[
 \nu(\rho(h))\ge\frac23(1-q^{-1})(1-q^{-3})a/2
                       \ge(q-1)a/(4q),                        \tag{OC2.9}
\]
since $1-q^{-3}\ge3/4$.

It remains that $h$ has no invariant isotropic subspace. Its
semisimple primary decomposition has neither a plus/minus-one
summand nor a pair of distinct reciprocal factors, since these
provide an invariant isotropic subspace. A self-reciprocal
irreducible factor of multiplicity at least two also provides
one: its associated hermitian space over the finite extension
has an isotropic line, by the norm argument just given.
Consequently the natural space consists of one self-reciprocal
factor of degree four, or two distinct self-reciprocal factors
of degree two, each with multiplicity one. Their symplectic
centralizers have respective orders $q^2+1$ and $(q+1)^2$.
Indeed their commuting natural operators are field scalars,
and preserving the form imposes the norm-one condition on
each field. Column orthogonality therefore gives
\[
 |\chi(h)|\le\sqrt{|C_S(h)|}\le q+1,
 \qquad \frac{|\chi(h)|}{a}\le\frac2{q-1}.
                                                               \tag{OC2.10}
\]
Thus $\nu\ge a(q-3)/(2(q-1))\ge a(q-1)/(4q)$ for $q\ge7$;
the latter inequality is $q^2-4q-1\ge0$. All cases are exhausted.
\end{proof}

\begin{lemma}[D-OC2-03: compatible outer operators]
\label{D-OC2-03}
For $q\ge7$ odd, every finite-order nonscalar operator normalizing
the represented image of $\rho$ satisfies
\[
 \nu(T)\ge\frac{q-1}{8q}a.                                   \tag{OC2.11}
\]
If $\rho$ is nonselfdual, the stronger bound
\[
 \nu(T)\ge\frac{q-1}{4q}a                                   \tag{OC2.12}
\]
holds, including every compatible field automorphism.
\end{lemma}

\begin{proof}
For (OC2.11), choose a core element with nonidentity commutator
with $T$, using irreducibility. A noncentral commutator has
support at least $(q-1)a/(4q)$, while a nonidentity scalar
commutator has actual moving rank $a$. The elementary rank
inequality $\operatorname{rank}([T,h]-I)\le2\nu(T)$ proves
the claim. It is valid after scalar-normalizing $T$ so that
its largest eigenspace has eigenvalue one.

For the stronger statement we use the exact ordinary automorphism
description $\operatorname{Out}(\operatorname{PSp}_4(q))
=C_2\times C_f$, $q=p^f$, from
\cite[\S2.2.2, Table~2.1, author p.~14]{HES:LiXia2022}.
For $q\ge5$ the universal perfect cover is $\operatorname{Sp}_4(q)$,
with center of order two and no exceptional multiplier; the precise
exception list is \cite[Chapter~5, Table~1, p.~35]{BI:HHM}.
Automorphisms of a represented central quotient lift uniquely to
this cover, and an automorphism trivial on its simple quotient is
trivial: its difference from the identity is a homomorphism from
a perfect group to its center.

Conjugation by a fixed multiplier-minus-one matrix sends every
symplectic conjugacy class to its inverse class. To see the exact
published input, \cite[Theorem~1, author p.~1]{OC2:Vinroot2004}
states that every symplectic element is a product of two
skew-symplectic involutions. If $g=ab$, then $aga=g^{-1}$.
All multiplier-minus-one conjugations differ by an inner
symplectic conjugation, proving the assertion. If $-1$ is a
square, scalar normalization makes this conjugation inner, so
every irreducible character is selfdual. Hence nonselfduality
implies $q\equiv3\pmod4$, $p\equiv3\pmod4$ and $f$ odd.
The nontrivial diagonal outer automorphism $\delta$ then sends
$\rho$ to $\rho^\vee$ and is incompatible with $\rho$.

In fact every compatible outer class has trivial diagonal
component. If $\delta\sigma^j$ stabilized $\rho$, raising it to
the odd order of $\sigma^j$ in $C_f$ would make $\delta$
stabilize $\rho$, a contradiction. A compatible operator is
therefore scalar times an inner operator, or has a nontrivial
field component of odd order. In the latter case take a power
whose field component has prime order $\ell\ge3$; support
cannot increase on taking this power. Write the resulting
automorphism as $\alpha$, with field part $\sigma$ of order
$\ell$. If $\alpha^\ell$ is nontrivial inner, the corresponding
power of the intertwiner is scalar times a noncentral inner
operator, and (OC2.5) proves the desired bound.

Suppose instead $\alpha^\ell=1$. Because its diagonal class is
trivial, it has a natural lift $s=h\sigma$ with
$h\in\operatorname{Sp}_4(q)$. Its $\ell$th power is a central
sign $cI$. Replacing $s$ by $cs$, which implements the same
automorphism, makes $s^\ell=1$, since $\ell$ is odd. Let
$k=\mathbf F_{q_0}$ be the fixed field of $\sigma$, so
$q=q_0^\ell$ and $q_0\ge3$.

We give descent without assuming that $\ell$ is invertible in
the field. For $z\in\mathbf F_q$ and a natural vector $v$, put
\[
             P_z(v)=\sum_{j=0}^{\ell-1}\sigma^j(z)s^j(v).
                                                               \tag{OC2.13}
\]
It is fixed by $s$. Choose a $k$-basis $z_1,\ldots,z_\ell$
of $\mathbf F_q$. The matrix $(\sigma^j(z_i))_{i,j}$ is
invertible. Indeed, a dependence of its columns would give
a nonzero $q_0$-linearized polynomial of degree at most
$q_0^{\ell-1}$ vanishing on a $k$-basis and hence on all
$q_0^\ell$ elements, an impossibility. Reordering the columns
handles any generator of the cyclic Galois group. Consequently
every $v$ is in the $\mathbf F_q$-span of the fixed vectors
$P_{z_i}(v)$. The fixed space has dimension four over $k$:
if $w_i$ are $k$-independent fixed vectors and
$\sum a_iw_i=0$, application of $P_z$ gives
$\sum\operatorname{Tr}_{q/q_0}(za_i)w_i=0$ for every $z$.
Nondegeneracy of the finite-field trace pairing makes every
$a_i=0$. This trace pairing is nondegenerate because its
linearized trace polynomial has degree less than $q$ and
is nonzero, and multiplication by any nonzero scalar is
bijective. No averaging by $1/\ell$ has been used.

The natural alternating form takes values in $k$ on this
fixed space and is nondegenerate there, because the fixed
space spans. A symplectic $k$-basis gives
$P\in\operatorname{Sp}_4(q)$ with $P^{-1}sP=\sigma$.
Relabeling the representation by this inner change, the
intertwiner implements pure Frobenius. Its $\ell$th power is
scalar, so normalize that power to one.

Inside the standard $\operatorname{GL}_2(q)$ Levi take its
one-row additive subgroup $U\simeq(\mathbf F_q,+)$.
The diagonal torus is transitive on its nonzero additive
characters. If each nonzero weight has multiplicity $m$,
averaging (OC2.1) yields $m\ge a/(2q)$. Frobenius fixes
exactly $q_0$ of the $q$ dual weights and puts every other
weight in an orbit of length $\ell$. Each such orbit
contributes at least a fraction $1-1/\ell$ of its dimension
to projective support. Therefore
\[
 \nu(T)\ge\frac12(1-1/\ell)(1-q_0/q)a
       \ge\frac8{27}a>\frac14a
       >\frac{q-1}{4q}a.                                    \tag{OC2.14}
\]
Here $\ell\ge3$ and $q_0/q=q_0^{1-\ell}\le1/9$.
This completes the all-field normalizer proof.
\end{proof}

\begin{theorem}[D-OC2-04: the all-odd-$C_2$ source exclusion]
\label{D-OC2-04}
Let $J\subset\operatorname{GL}(V)$ be a finite group acting irreducibly
on a complex space $V$, and suppose an
irreducibly acting normal central product of quasisimple groups
with simple quotient $\operatorname{PSp}_4(q)$, $q\ge5$ odd,
has its factors permuted transitively by $J$. Suppose $V$ is an
inducing block of a nonselfdual original irreducible representation
in the sixfold signature-$(3,3)$ curve-cover problem. Then the
actual period image has dimension at most eight. All compatible
scalars, outer automorphisms and factor permutations are allowed.
The assertion imposes no polarization determinant condition and
does not equate the inducing and original character fields.
\end{theorem}

\begin{proof}
For $q=5$, D-OC2-F1--F2 already proves the stronger quarter-bound
exclusion. Assume $q\ge7$. Write $V=\bigotimes_{i=1}^k V_i$,
where the transitive core action makes the factor degrees all
equal to $a\ge(q^2-1)/2$. Their selfduality status is also the
same, since automorphisms and isomorphisms preserve it.
Write $e=a^k$. Factor-preserving operators have convolved
normalized eigenspace distributions. A nontrivial factor cycle
has support at least $e/4$, by the weighted basis-word orbit
argument: at most $a$ words are constant, and every other orbit
has length at least two. These statements permit arbitrary
outer intertwiners on the factors and arbitrary multiplicity
of eigenvalues. Nonidentity scalar inertias move all $e$
dimensions.

The actual intermediate-cover factorization gives
$6=2e(h-1)+D$, where $D$ is the sum of genuine inertia moving
ranks. Since $e\ge24$, genus at least two is impossible; genus
one has at most six genuine branch points and actual period
rank at most six. Rank nine would therefore require
\[
                    h=0,\quad R\ge12,\quad D=2e+6.
                                                               \tag{OC2.15}
\]
If the factors are nonselfdual, (OC2.12) and the tensor argument
give moving rank at least $(q-1)e/(4q)$ at every genuine branch.
Consequently
\[
 2e+6\ge12\frac{q-1}{4q}e,
       \qquad e\le\frac{6q}{q-3}\le\frac{21}2,
                                                               \tag{OC2.16}
\]
contradicting $e\ge24$.

Suppose the factors are selfdual. Their tensor product has a
unique invariant nondegenerate bilinear form up to scalar.
Every element of $J$ preserves its line and hence has a
similitude multiplier $\lambda:J\to\mathbf C^*$.
This multiplier is nontrivial. Otherwise $V$ would be selfdual
as a $J$-representation, and its induction to the original
group would also be selfdual, contrary to the hypothesis.
The product-one generating inertia tuple on the genus-zero
intermediate base contains at least two entries with
$\lambda\ne1$. For each, its fixed subspace is totally
isotropic for the bilinear form and has dimension at most
$e/2$. The two entries cost at least $e$ moving dimensions;
the other ten entries cost at least $(q-1)e/(8q)$ each by
(OC2.11) and the tensor argument. Thus
\[
 2e+6\ge e+10\frac{q-1}{8q}e,
                    \qquad e\le\frac{24q}{q-5}.              \tag{OC2.17}
\]
For $q\ge11$, the degree floor $(q^2-1)/2$ is already
strictly greater than this bound: the difference after
multiplication by $2(q-5)$ is
$q^3-5q^2-49q+5$, whose coefficients at $q=11+x$ are
positive. For $q=7,9$, two factors give respectively
$e\ge24^2,40^2$, whereas (OC2.17) gives $e\le84,54$.
There is therefore only one factor.

At $q=7$, selfduality excludes the four small characters by
D-OC2-01, leaving $a\ge126>84$. At $q=9$, the stronger
degree floor is $288>54$, except for the four Weil modules
of degrees $40,41$. Their full-normalizer source exclusion
is precisely the already reviewed D-SPW-02--03; it includes
arbitrary scalar and field operators and requires no
polarization determinant assumption. These are all the
possibilities, completing the contradiction to rank nine.
\end{proof}

\paragraph{What is not inferred.}
The result excludes the stated transitive ordinary symplectic
core configuration. It does not construct a source, does not
infer dominance from a parameter count, and does not extend
the complete finite-table coverage beyond the groups explicitly
listed in D-OC2-F. The $q=3$ endpoint retains its previously
proved separate reduction to the known degree-four and
degree-five possibilities.

\paragraph{Executed new certificates.}
The separate program \texttt{check\_all\_fields.py} verifies the
universal polynomial inequalities used in (OC2.9), (OC2.10),
(OC2.16), and (OC2.17), together with the symbolic symplectic
shear and Heisenberg multiplication identities. Its finite
$\mathbf F_{27}$ descent witness has Galois order equal to the
characteristic: it checks all $27^4=531441$ natural vectors,
finds exactly $3^4=81$ fixed vectors for an explicit symplectic
semilinear cocycle, checks the Moore matrix and trace pairing,
and obtains three fixed dual weights and eight length-three
orbits. This supports the implementation and the delicate
characteristic-three signs; the preceding universal proof
establishes descent in every field. This finite witness is a
separate check. The retained output is
\texttt{all\_fields\_checked.json}.

\section{Four complete exceptional small-field core exclusions}
\label{sec:exceptional-small-cores}

\paragraph{Status.}
New complete root derivation, read in full by the author and an
independent reviewer with both mathematical audits passing. Full GAP
orthogonality and spectrum checks, a separate root exact reconstruction,
and an independent polynomial-Fourier reconstruction all passed.
The original character tables are cited inputs with their distinct
provenances; they have not been independently constructed in this run.

\begin{lemma}[D-EXS-01: complete core and outer spectra]
\label{D-EXS-01}
Let $W$ be a nontrivial ordinary irreducible complex representation of a
quasisimple group with one of the simple quotients in the first column
below. Put $a=\dim W$ and
$\nu(T)=a-\max_\lambda\dim\ker(T-\lambda I)$. Every nonscalar core
operator satisfies the indicated bound:
\[
\begin{array}{c|c|c|c}
\text{simple quotient}&\text{universal cover used}&a\text{ at least}
                                      &\nu(\rho(x))/a\text{ at least}\\\hline
F_4(2)&2.F_4(2)&52&4/13\\
F_4(3)&F_4(3)&6643&32/73\\
E_6(2)&E_6(2)&2482&512/1241\\
{}^2E_6(2)&(2^2\times3).{}^2E_6(2)&1938&8/19
\end{array}                                                    \tag{EXS.1}
\]
For the first row, every nonscalar finite-order operator normalizing
the represented image has the same bound $\nu(T)\ge4a/13$.
No faithfulness on the entire center is assumed.
\end{lemma}

\begin{proof}
The universal-cover classification is used in precisely the ordinary
complex sense. The exceptional multiplier cases here are $F_4(2)$ and
${}^2E_6(2)$; the complete list is
\cite[Chapter~5, Table~1, p.~35]{BI:HHM}. The exact CTblLib metadata
\texttt{ExtensionInfoCharacterTable} for the four simple quotients,
in the order of (EXS.1), are
\[
 [\texttt{2},\texttt{2}],\quad
 [\texttt{empty},\texttt{empty}],\quad
 [\texttt{empty},\texttt{2}],\quad
 [\texttt{(2\^{}2x3)},\texttt{3.2}].                         \tag{EXS.2}
\]
The first entry specifies the multiplier and the second the outer
automorphism group. In particular the last center is noncyclic.
We retain \emph{every} nonlinear row of its universal table, so all
central quotients and all possible central-character kernels are included.
We do not replace this table by a presumed faithful representation
of the full noncyclic center.

The exact source identifiers and provenance are in
\cite{HD:ExceptionalSmallTables}. The $F_4(2)$ cover and its outer
extension are ATLAS inputs; the retained extension uses the corrected
square map. The $F_4(3)$ table was computed by Eamonn O'Brien with
Magma in May/June~2024 and added to CTblLib~1.3.10. The $E_6(2)$
table is Fischer's table with Unger's 2016 corrections to order-91
power maps and irrationalities. The universal ${}^2E_6(2)$ table
records its construction by \texttt{PossibleCharacterTablesOfTypeV4G}.
These separate data origins are not all relabelled as printed ATLAS
tables. The actual export uses CTblLib~1.3.11.

The complete mathematical input is
\texttt{research/exceptional\_small\_inputs.json}, with SHA-256
\texttt{625618617ad0b6268b8898cc39e7e04afb6ae55f8639686347a0b9cf8d244a21}.
It retains every class size, element order, ordinary character value,
prime power map, indicator, source string, and the core-to-outer fusion.
For an element $x$ of order $d$, every eigenvalue multiplicity is
computed by the exact identity
\[
 m_j=\frac1d\sum_{u=0}^{d-1}\chi(x^u)\exp(-2\pi iju/d),
                       \qquad 0\le j<d.                        \tag{EXS.3}
\]
The reconstruction uses rational field traces and Ramanujan sums,
with exact integer arithmetic, and verifies integrality,
nonnegativity and the sum $\sum m_j=a$. It also checks the degree
sum, all power-map orders, and the normal fusion's identity kernel,
complete class unions and compatibility with prime power maps.
An entry is skipped only when $|\chi(x)|^2=a^2$, which for a
finite-order operator is exactly the scalar condition.

For the table \texttt{2.F4(2).2}, a row is selected exactly when
its restriction to the core has norm one and matches an actual
nontrivial irreducible core row. The complete counts and minima are
\[
\begin{array}{c|r|r|c}
\text{table}&\text{selected rows}&\text{nonscalar spectra}&\min\nu/a\\\hline
\texttt{2.F4(2)}&169&28392&4/13\\
\texttt{2.F4(2).2}&86&12814&4/13\\
\texttt{F4(3)}&272&73984&32/73\\
\texttt{E6(2)}&179&32041&512/1241\\
\texttt{(2\^{}2x3).2E6(2)}&933&860226&8/19
\end{array}                                                    \tag{EXS.4}
\]
Thus there are $1639$ selected rows and $1007457$ spectra.
The script \texttt{check\_exceptional\_small.py} retains the minimum
and a witnessing class for every row in
\texttt{exceptional\_small\_checked.json}. These are complete exact
finite inequalities, not numerical sampling or a calculation of only
unipotent characters. The arithmetic helper is the previously reviewed
Ramanujan implementation, reused here separately from GAP's
\texttt{EigenvaluesChar}; it is not advertised as a new third algorithm.

For completeness the outer-table assertion covers actual normalizers,
not only the chosen abstract extension. The universal cover of $F_4(2)$
has center of order two and outer group of order two. An automorphism
of a represented central quotient lifts uniquely to the universal
perfect cover. If it induces the identity on the simple quotient,
its difference from the identity is a homomorphism from a perfect
group to its center and is trivial. Inner automorphisms are implemented
by core operators up to scalar, by Schur's lemma.

For an outer automorphism, adjust by an inner operator to the
automorphism induced by a fixed lift $t$ in \texttt{2.F4(2).2}.
A compatible intertwiner $A$ satisfies $A^2=c\rho(t^2)$ for a
nonzero scalar $c$. Rescale $A$ to remove $c$. The defining relation
then gives a genuine representation of this extension with irreducible
core restriction, exactly among the selected rows. The original
normalizer is a scalar multiple of an element of that representation.
This does not assume a split extension and also handles a different
isoclinic scalar square. Since scalar multiplication preserves $\nu$,
(EXS.4) proves the asserted full-normalizer bound.
\end{proof}

\begin{lemma}[D-EXS-02: uniform support for all full normalizers]
\label{D-EXS-02}
For every representation in D-EXS-01, every nonscalar finite-order
operator normalizing the represented core has
\[
                              \nu(T)\ge a/5.                   \tag{EXS.5}
\]
This includes all compatible outer automorphisms, without assuming
selfduality or a cyclic outer group in the last row.
\end{lemma}

\begin{proof}
The first row already has the stronger bound $4a/13$.
In each of the other three rows, (EXS.1) gives core support
at least $2a/5$: indeed $5\cdot32\ge2\cdot73$,
$5\cdot512\ge2\cdot1241$, and $5\cdot8\ge2\cdot19$.

Choose a represented core element $h$ with $[T,h]\ne I$;
otherwise irreducibility and Schur's lemma would make $T$ scalar.
The commutator belongs to the represented core. If it is nonscalar,
its moving rank is at least its projective support, hence at least
$2a/5$. If it is a nonidentity scalar, its moving rank is $a$.
Normalize $T$ by a scalar so that a largest eigenspace has eigenvalue
one. This does not change its commutator, and
\[
 \operatorname{rank}([T,h]-I)
   =\operatorname{rank}(Th-hT)
   \le2\operatorname{rank}(T-I)=2\nu(T).                       \tag{EXS.6}
\]
The equality follows by multiplying by invertible matrices, and the
inequality by writing $Th-hT=(T-I)h-h(T-I)$. Thus (EXS.5) follows.
No description of the outer automorphism action on the noncyclic
center of the last universal cover is needed for this commutator
argument; every compatible normalizer is handled directly.
\end{proof}

\begin{theorem}[D-EXS-03: no mixed products of these exceptional cores]
\label{D-EXS-03}
In the actual sixfold curve-cover source problem, an inducing block
cannot have period rank nine if its irreducibly acting normal
quasisimple central-product core has all simple quotients among
\[
                  F_4(2),\quad F_4(3),\quad E_6(2),
                             \quad {}^2E_6(2).                \tag{EXS.7}
\]
Mixed types, arbitrary central quotients, compatible outer and scalar
operators, and permutations of factors are permitted. No transitivity,
selfduality, character-field equality or polarization determinant
hypothesis is imposed beyond the stated actual source setting.
\end{theorem}

\begin{proof}
Lift the core to the product of its universal covers and write
$V=\bigotimes_i W_i$, with $e=\prod_i a_i$. Every $a_i\ge52$.
For a factor-preserving nonscalar normalizer, at least one factor
operator is nonscalar. Normalized eigenvalue multiplicities convolve
under tensor product; their maximum is no greater than the maximum
for that factor. D-EXS-02 therefore gives support at least $e/5$.

For a nontrivial factor permutation, the complete weighted basis-word
argument of D-SPOR-03 gives the stronger $e/4$. Explicitly, on a
cycle of length $l\ge2$ with factor dimension $a\ge2$, diagonalize
the product of the intertwiners around the cycle. In the resulting
weighted word basis, at most $a$ of the $a^l$ words are fixed;
every other orbit has length at least two and contributes at most
one dimension to any chosen eigenspace. Thus its largest eigenspace
has dimension at most $(a^l+a)/2\le3a^l/4$. Tensoring the other
cycles preserves this upper fraction. This argument does not require
the intertwiners to be pure coordinate permutations. Finally a
nonidentity scalar inertia moves all $e$ dimensions.

The actual intermediate-cover Euler identity from D-BI-01 is
$6=2e(h-1)+D$, with $D$ the total inertia moving rank. Since
$e\ge52$, genus $h\ge2$ is impossible. At genus one, $D=6$
allows at most six genuine branches and actual period rank at most
six. Rank nine must consequently have $h=0$, at least twelve
genuine branch points, and $D=2e+6$. Applying (EXS.5) at each
branch gives
\[
                 2e+6\ge12e/5,\qquad e\le15,                 \tag{EXS.8}
\]
contrary to $e\ge52$. This proves the theorem for every allowed
mixture, without a reality or multiplier argument.
\end{proof}

\paragraph{Geometry, framework, and precise boundary.}
This is a restriction on actual finite covers and their weight-one
multiplicity variations. Its spectrum and tensor identities lie in
the representation ring; the scheme-spreading, SGA finite-cover,
lambda-ring, K-theory and etale-weight placements have the exact
locators in the framework section. It strictly narrows the supplier
gate for condition A, cote16 scan050. It produces no new effective
cycle or detector constant, and changes neither condition I at
cote16 scan049 and cote44 Theorem4 nor the Cayley--Hamilton projector
at cote16 scans048--049. No motivic $L$-value estimate or general
negative Hodge-conjecture conclusion is inferred. Larger fields and
other exceptional families are outside (EXS.7).

\section{Nine complete small orthogonal core reductions}
\label{sec:orthogonal-small-cores}

\paragraph{Status and scope.}
New root derivation. Complete author and independent peer proof and
source reviews have passed, together with separate full exact spectrum
reconstructions and the retained GAP orthogonality checks. The character tables are
cited exact inputs, not independently reconstructed original tables.
Eight simple types are excluded in arbitrary mixed products. The ninth,
$\Omega_9(3)$, has a separately stated transitive-core conclusion in the
nonselfdual-original inducing problem. No assertion about arbitrary
mixtures involving that ninth type is substituted for its hypotheses.

Write $\nu(T)=a-\max_\lambda\dim\ker(T-\lambda I)$ for
projective support on a complex space of degree $a$.

\begin{lemma}[D-OSM-01: complete ordinary spectra]
\label{D-OSM-01}
For each complete ordinary table in the following display, take every
nonlinear irreducible row. For an extension table take precisely its
rows whose restriction to the displayed universal core is irreducible.
The minimum degree and the minimum $\nu/a$ over all nonscalar elements
are exactly as follows.
\[
\begin{array}{l|r|r|r|c}
\text{table}&\text{rows}&\text{spectra}&\min a&\min\nu/a\\\hline
6.O_7(3)&179&31146&27&2/9\\
6.O_7(3).2&126&22932&78&11/39\\
O_8^-(2)&38&1444&34&16/51\\
O_8^-(2).2&52&3068&34&7/34\\
2^2.O_8^+(3)&353&123550&260&28/65\\
2.O_8^-(3)&189&35532&246&52/123\\
2.O_9(3)&325&105300&780&13/41\\
O_{10}^+(2)&96&9216&155&32/93\\
O_{10}^+(2).2&140&21560&155&7/31\\
O_{10}^-(2)&114&12996&154&64/187\\
O_{10}^-(2).2&140&22820&154&5/22\\
O_{12}^+(2)&270&72900&651&128/357\\
O_{12}^+(2).2&340&132940&651&5/21\\
O_{12}^-(2)&236&55696&650&256/715\\
O_{12}^-(2).2&340&127160&650&31/130
\end{array}                                                     \tag{OSM.1}
\]
The only nonselfdual rows of $2.O_9(3)$ are rows $224,225$,
both of degree $838656$ and minimum support ratio $1/2$.
All its other rows have Frobenius--Schur indicator $+1$.
\end{lemma}

\begin{proof}
The mathematical inputs are the complete official CTblLib tables
\cite{HD:OrthogonalSmallTables}. The exporter retains every character
value, class size and order, every prime power map, all six core
fusions, all indicators, and the original provenance strings.
The complete input has SHA-256
\texttt{5236207618c2564af7ed4871c0f92a667664255d2e60be87f037257356669b79}.
GAP verifies full row orthogonality for all fifteen tables. A GAP
spectrum computation and a separate exact rational-trace program
recover every eigenvalue multiplicity by
\[
 m_k(x)=\frac1d\sum_{j=0}^{d-1}\chi(x^j)\zeta_d^{-kj},
 \qquad d=|x|.                                                \tag{OSM.2}
\]
The separate program selects extension rows by restriction inner
product one and checks exact pointwise equality with a unique core
row. It checks full class unions, the identity kernel, element orders,
and all retained power-map compatibilities for the fusions. All
Fourier multiplicities are exact nonnegative integers with the
correct degree sum. Scalar elements are detected by
$|\chi(x)|^2=\chi(1)^2$; no inertia class is omitted.

There are exactly $2938$ selected rows and $778260$ nonscalar spectra.
The retained per-row witnesses prove all the displayed minima.
The separate program also recomputes every Frobenius--Schur indicator
using $|G|^{-1}\sum_x\chi(x^2)$, with exact rational traces of the
cyclotomic values. The original sum is a rational integer, so taking
its normalized rational trace preserves it. This verifies the two
exceptional nonselfdual rows and the indicators of every other row.

The strings $O_d^\epsilon(q)$ here denote the simple groups in the
official table notation, not the full natural orthogonal groups.
The retained multiplier/outer metadata, interpreted as in
\cite[\S3.7-3]{HD:OrthogonalSmallTables}, are respectively
\[
\begin{array}{c|c|c}
S&M(S)&\operatorname{Out}(S)\\\hline
O_7(3)&6&2\\
O_8^-(2)&1&2\\
O_8^+(3)&2^2&S_4\\
O_8^-(3)&2&2^2\\
O_9(3)&2&2\\
O_{10}^{\pm}(2)&1&2\\
O_{12}^{\pm}(2)&1&2
\end{array}                                                     \tag{OSM.3}
\]
The exceptional $\Omega_7(3)$ multiplier is also located in
\cite[Chapter~5, Table~1 and preceding paragraph, p.~35]{BI:HHM}.
Every quasisimple central quotient lifts to the stated universal cover.
Keeping all nonlinear rows therefore includes all central characters,
including their nonfaithful central kernels. Faithfulness on the
noncyclic center of $\operatorname{Spin}_8^+(3)$ is not assumed.
\end{proof}

\begin{lemma}[D-OSM-02: full normalizers]
\label{D-OSM-02}
For every quasisimple core with simple quotient among
\[
\mathcal F=\{\Omega_7(3),\Omega_8^-(2),
 \mathrm{P}\Omega_8^+(3),\mathrm{P}\Omega_8^-(3),
 \Omega_{10}^+(2),\Omega_{10}^-(2),
 \Omega_{12}^+(2),\Omega_{12}^-(2)\},                         \tag{OSM.4}
\]
every finite-order nonscalar normalizer of a nontrivial irreducible
degree-$a$ representation has support at least $a/5$.
For a core with simple quotient $\Omega_9(3)$ the bound is
$13a/82$; it improves to $a/4$ when the core representation is
nonselfdual.
\end{lemma}

\begin{proof}
First consider the six types in (OSM.4) other than the two
eight-dimensional natural types over $\mathbf F_3$.
Their outer groups have order two, and the complete universal core
and covering outer table are both in (OSM.1). An automorphism of a
represented central quotient lifts to the universal perfect cover.
Such a lift is unique: an automorphism trivial on the simple quotient
differs from the identity by a homomorphism of a perfect group to its
center, which is trivial.

Let $T$ implement a compatible outer automorphism. Choose an element
$t$ in the retained covering extension inducing the same automorphism,
after multiplication by a core element. By Schur's lemma
$T^2=c\rho(t^2)$ for a nonzero scalar $c$. Rescaling $T$ solves this
actual square relation, so $\rho$ extends to a representation of the
listed covering extension with that action. It is irreducible because
its restriction is irreducible. Every such row is retained. This
argument permits a nonsplit relation and an isoclinic alternative;
it does not assert that isoclinic groups have equal character tables.
Changing scalar phases preserves projective support.
The minimum across these six complete core/outer pairs exceeds $1/5$.

For the two remaining types in (OSM.4), every nonscalar core operator
has support at least $(2/5)a$, since $28/65,52/123>2/5$.
A nonscalar normalizer $T$ fails to commute with some represented
core element $g$, by irreducibility. If $[T,g]$ is nonscalar its
actual moving rank is at least $(2/5)a$. If it is a nonidentity
scalar its actual moving rank is $a$. Scalar-normalize $T$ at its
largest eigenspace. Then
\[
 \operatorname{rank}([T,g]-I)
 =\operatorname{rank}(Tg-gT)\le2\operatorname{rank}(T-I)
 =2\nu(T).                                                    \tag{OSM.5}
\]
This proves $\nu(T)\ge a/5$ for every compatible normalizer,
including all outer automorphisms of the noncyclic central cover.
No table for the full $S_4$ outer group is required.

The identical commutator argument applied to $2.O_9(3)$ gives
$\nu(T)\ge13a/82$. For its two nonselfdual characters, the full
core support is $a/2$, so the argument gives $\nu(T)\ge a/4$.
Every central quotient is already covered by the universal rows.
\end{proof}

\begin{theorem}[D-OSM-03: eight mixed-type source exclusions]
\label{D-OSM-03}
An irreducibly acting normal central-product core whose simple
quotients all lie in $\mathcal F$ cannot supply actual period rank
nine in the signature-$(3,3)$ sixfold-source problem. Arbitrary
mixtures of those eight types, covers, scalars, compatible outer
operators and factor permutations are allowed. There is no
selfduality or determinant hypothesis.
\end{theorem}

\begin{proof}
The core module is a tensor product of nontrivial irreducible factor
modules. A factor-preserving nonscalar operator has normalized
eigenspace multiplicities given by convolution; their maximum is at
most the maximum for any nonscalar factor. D-OSM-02 therefore gives
support at least $e/5$, where $e$ is the total degree.
For a nontrivial factor cycle of length $r$ and factor degree $a$,
conjugate the intertwiners to identities except for a diagonal closing
twist. This is a weighted cyclic permutation on $a^r$ basis words.
There are $a$ constant words, and every other orbit has length at
least two. Each orbit contributes at most one to any chosen
eigenspace, so its dimension is at most
$a+(a^r-a)/2\le3a^r/4$. Tensoring other cycles cannot increase
the normalized maximum. The resulting quarter support is stronger
than required. A genuine nonidentity scalar inertia moves all $e$
dimensions.

Every factor has degree at least $27$. The actual intermediate-cover
identity, with base genus $h$ and $R$ genuine inertia points, is
$6=2e(h-1)+D$, where $D$ is the sum of actual moving ranks.
Genus $h\ge2$ is impossible; genus one has $D=6$, hence at most six
effective branches and actual period-image dimension at most six.
Rank nine would force genus zero, $R\ge12$ and $D=2e+6$.
Consequently $2e+6\ge12e/5$, giving $e\le15<27$, a contradiction.
This uses the actual variation factorization, valid for a larger
inducing character field; it does not infer dominance from a parameter
count.
\end{proof}

\begin{theorem}[D-OSM-04: the ninth transitive core]
\label{D-OSM-04}
Suppose the original sixfold-source representation is nonselfdual,
and an inducing block has a transitive irreducibly acting normal
central-product core with simple factors $\Omega_9(3)$. It cannot
have an actual period image of dimension nine. Every compatible
normalizer, scalar, factor permutation and either determinant is allowed.
\end{theorem}

\begin{proof}
Transitivity makes all core factors either selfdual or nonselfdual.
If they are nonselfdual, D-OSM-02 and the tensor-cycle argument give
quarter support for every nonscalar inertia. Rank nine would force
$2e+6\ge3e$, hence $e\le6$, whereas every factor has degree at
least $780$.

Otherwise each factor has an invariant nondegenerate bilinear form,
unique up to scalar. Their tensor product gives the unique invariant
form line for the irreducible core. The inducing image acts on it by
a multiplier character $\mu$. This character is nontrivial: if it
were trivial, the inducing representation would be selfdual, and
its finite-group induction would make the original representation
selfdual, contrary to hypothesis. No equality of character fields is
used.

The same actual Euler/Hurwitz argument forces genus zero and at least
twelve genuine inertias, generating the inducing image and having
product one. At least two have nontrivial $\mu$. Their fixed spaces
are totally isotropic for the form and their moving ranks are at least
$e/2$. Every remaining genuine inertia costs at least $13e/82$ by
D-OSM-02 and the tensor-cycle bound. Thus
\[
 2e+6\ge e+10(13e/82),\qquad e\le41/4<780,                 \tag{OSM.6}
\]
again impossible. The selfdual/nonselfdual division is essential to
this argument; arbitrary intransitive mixed factors involving this
ninth type are not asserted to be excluded.
\end{proof}

\paragraph{Boundary and geometric placement.}
These are necessary restrictions on finite-cover weight-one variations,
their integral representation-ring operations and actual period maps.
They narrow the condition-A cycle-supplier gate at cote16, scan050,
in the EGA/SGA and realization framework stated in the paper. They
produce no cycle, positivity division, Cayley--Hamilton projector or
motivic $L$-value estimate. Other fields and orthogonal ranks, in
particular the small $\Omega_8^+(2)$ cover representations, retain
their separate gates. No Markman, Severn or 032 premise is used.

\section{Higher-rank symplectic cores in every odd field at least seven}
\label{sec:higher-odd-symplectic}

\paragraph{Status and precise scope.}
This complete new derivation has passed whole author, root and
independent peer proof reviews, with separately executed exact certificates. It uses the proved elementary carrier and degree
lemmas D-HOS-01--02. The only character-ratio input is the published
unipotent estimate stated below. No full character table, classification
of small character degrees, or Weil-character formula is a premise.
The final theorem permits mixed ranks and fields and does not require
transitivity on the factors, essential selfduality, or a polarization
determinant condition. The odd fields $3$ and $5$ are outside its scope.

For a finite-order complex operator $T$ of degree $a$, put
\[
 \nu(T)=a-\max_\lambda\dim\ker(T-\lambda I).
\]
Scalar multiplication does not change this quantity. Ordinary moving
rank is at least $\nu(T)$, and $\nu(T)\ge\nu(T^j)$ for every integer
$j$. For a unitary operator, the triangle inequality at its largest
eigenspace gives
\[
                 |\operatorname{Tr}T|\ge a-2\nu(T).       \tag{HOS.13}
\]

The exact published input is
\cite[simply connected conventions, p.~2;
Theorem~2.4(i), pp.~3--4]{HES:Liebeck2017}, the published restatement
of Gluck's unipotent estimate. For $q>4$, a nontrivial ordinary
irreducible character $\chi$ of $\operatorname{Sp}_{2n}(q)$ satisfies
\[
 |\chi(u)|\le c_q\chi(1),\qquad
 c_q=\frac1{\sqrt q-1},\qquad
                 1\ne u\text{ unipotent}.                 \tag{HOS.14}
\]
The subtraction is outside the radical. In what follows $q\ge7$,
so $0<c_q<1$. Define
\[
 \delta_q=\frac12(1-q^{-1})(1-q^{-3})(1-c_q),\qquad
 \delta_* =\frac{9234}{55223}>\frac16.                     \tag{HOS.15}
\]

\begin{lemma}[D-HOS-03: every inner operator]
\label{D-HOS-03}
Let $q\ge7$ be odd, $n\ge4$, and let $\rho$ be a nontrivial
ordinary irreducible representation of $S=\operatorname{Sp}_{2n}(q)$
of degree $a$. For every noncentral $x\in S$,
\[
                         \nu(\rho(x))\ge\delta_q a.        \tag{HOS.16}
\]
\end{lemma}

\begin{proof}
The projective kernel of $\rho$ is exactly $Z(S)$: its image in the
simple quotient is normal, and if this image were the whole quotient,
the perfect group would act by scalars and hence trivially. Take a
power $h$ of $x$ whose image in $S/Z(S)$ has prime order $\ell$.
Its support is no greater than that of $x$. If $\ell=p$, where
$q=p^f$, multiplication by a central sign makes $h$ a nonidentity
unipotent. Equations (HOS.13)--(HOS.14) give support at least
$(1-c_q)a/2\ge\delta_q a$.

Suppose $\ell\ne p$. The element $h$ is semisimple. If $\ell$
is odd, a central sign makes its actual order $\ell$. If $\ell=2$,
its square can be either $I$ or $-I$. We retain this latter
possibility throughout the argument.

Let $E$ be the natural symplectic space with alternating form
$B$. First suppose $h$ has a nonzero invariant totally isotropic
subspace $R$. Choose a basis $r_1,\ldots,r_d$ of $R$. For a
symmetric matrix $C=(c_{ij})$ set
\[
 N_C(v)=\sum_{i,j}c_{ij}B(r_j,v)r_i,\qquad u_C=I+N_C.
                                                               \tag{HOS.17}
\]
Isotropy gives $N_CN_D=0$, and symmetry gives
$B(N_Cv,w)+B(v,N_Cw)=0$. Thus the $u_C$ form an additive
unipotent subgroup of $S$. The map $C\mapsto N_C$ is injective,
because the pairing functionals of the independent $r_j$ are
independent. If $A=h|_R$, conjugation is $C\mapsto ACA^t$.
The action has order dividing $\ell$, including when $h^2=-I$.

If this action is nontrivial, its moving subspace over
$\mathbf F_q$ has no fixed vectors: an operator of prime order
coprime to the characteristic is semisimple, and its moving and
fixed subspaces are complementary. It gives a subgroup $U$ of
order $Q\ge q$. Every nonzero dual weight has an orbit of length
$\ell$. Averaging (HOS.14) bounds its trivial weight dimension by
$a(c_q+(1-c_q)/Q)$. On each other orbit a cyclic block operator
has every eigenspace at most $1/\ell$ of that orbit's dimension;
this conclusion does not require its global $\ell$th power to be
scalar. Therefore
\[
 \nu(\rho(h))\ge(1-1/\ell)(1-Q^{-1})(1-c_q)a
       \ge\tfrac12(1-q^{-1})(1-c_q)a\ge\delta_q a.         \tag{HOS.18}
\]

It remains in this isotropic case that $ACA^t=C$ for all symmetric
$C$. Since $A$ is semisimple, its eigenvalue squares and all their
pairwise products equal one. Hence $A=\lambda I$ with
$\lambda=\pm1$. In particular $h^2=-I$ is impossible in this
case. Choose a nonzero $e\in R$. There is a paired vector $f$
with $B(e,f)=1$ and $hf=\lambda f$: project a paired vector to
the $\lambda$ eigenspace by the coprime-order averaging projector.
Explicitly, if $o$ is the actual order of $h$, use
$o^{-1}\sum_{j=0}^{o-1}\lambda^j h^j f$.
The plane $\langle e,f\rangle$ is nondegenerate, and its orthogonal
complement $W$ is $h$-stable. Put $g=\lambda h|_W$. This is
symplectic of order dividing $\ell$ and is nonidentity, since
otherwise $h=\lambda I$ on all of $E$. Its moving subspace
$W_0=\operatorname{im}(g-I)$ is nondegenerate, has even dimension
$2m\ge2$, and $g$ has no fixed vectors on it.

The line-parabolic radical contains the following subgroup, with
$v\in W_0$ and $z\in\mathbf F_q$:
\[
 e\mapsto e,\qquad w\mapsto w+B(v,w)e,\qquad
                 f\mapsto f+v+2ze.                         \tag{HOS.19}
\]
Its multiplication is the Heisenberg law
$(v,z)(v',z')=(v+v',z+z'+B(v,v')/2)$.
Conjugation by $h$ is $(v,z)\mapsto(gv,z)$, since
$\lambda^2=1$. Every nonidentity element is unipotent.
D-HOS-01, applied to the full restriction of $\rho$, gives
\[
 \nu(\rho(h))\ge
 (1-1/\ell)(1-q^{-m})(1-q^{-(2m+1)})(1-c_q)a
                         \ge\delta_q a.                    \tag{HOS.20}
\]
The lemma includes both linear-weight orbits and all nontrivial
central-character carriers with arbitrary multiplicity operators.
Thus no assumption about the restriction being irreducible or
Weil is needed here.

Finally suppose that $h$ has no invariant nonzero totally isotropic
subspace. In its semisimple primary decomposition, a pair of distinct
reciprocal factors provides such a subspace. A plus- or minus-one
primary summand does too. A self-reciprocal irreducible factor of
degree $2d$ and multiplicity at least two also provides one: the
associated form on the multiplicity space is Hermitian over a finite
quadratic extension, and a Hermitian space of dimension at least two
is isotropic. More explicitly, put $K=\mathbf F_q[h]$ on that
primary space. Its adjoint involution sends $h$ to $h^{-1}$;
its fixed field is $\mathbf F_{q^d}$. Defining a $K$-valued
form by $\operatorname{Tr}_{K/\mathbf F_q}(a b(v,w))=B(av,w)$
for every $a\in K$ gives a nondegenerate skew-Hermitian form,
by the trace-pairing argument written out in (HOS.23) below.
Multiplication by a nonzero scalar negated by the involution
makes it Hermitian. To see its isotropy directly, diagonalize two coordinates
and solve their norm equation using surjectivity of the finite-field
norm. The resulting extension-field line is isotropic for the
original alternating form as well.

Consequently all primary factors are distinct self-reciprocal
irreducibles of degrees $2d_i$, each with multiplicity one, where
$\sum_i d_i=n$. Their commuting natural operators are field scalars,
and preservation of the form imposes the norm-one condition.
It follows that
\[
 |C_S(h)|=\prod_i(q^{d_i}+1)\le(q+1)^n.
\]
Column orthogonality and D-HOS-02 now give
\[
 \frac{|\chi(h)|}{a}
 \le\frac{2(q+1)^{n/2}}{(q-1)q^{n-1}}
 \le\frac{2(q+1)^2}{(q-1)q^3}<\frac1q<c_q.                 \tag{HOS.21}
\]
The middle inequality follows by increasing $n$ from four, each
increment multiplying the expression by $\sqrt{q+1}/q<1$.
For the strict bound use $q+1\le8q/7$ and $q-1\ge6$, so
$2(q+1)^2/((q-1)q^2)\le128/294<1$.
Equations (HOS.13) and (HOS.21) give at least
$(1-c_q)a/2\ge\delta_q a$. All cases have been exhausted.
\end{proof}

\begin{lemma}[D-HOS-04: diagonal projective involutions]
\label{D-HOS-04}
In the setting of D-HOS-03, let $T$ implement a diagonal outer
automorphism whose square is the identity automorphism. Then
\[
                    \nu(T)\ge\tfrac12(1-q^{-1})(1-c_q)a.
                                                               \tag{HOS.22}
\]
\end{lemma}

\begin{proof}
Take a natural lift $h\in\operatorname{GSp}_{2n}(q)$ with
nonsquare multiplier $\mu$. The identity square of the automorphism
implies $h^2=cI$ for some $c\in\mathbf F_q^*$, since the natural
centralizer of the symplectic group is scalar. Thus $\mu=\pm c$.
We construct a nonzero $h$-stable totally isotropic subspace $R$.

If $c$ is a square, $h$ has eigenspaces for the two square roots
of $c$. Since $\mu$ is nonsquare, $\mu=-c$. The form vanishes
on either eigenspace, because its multiplier there would otherwise
be $c$ instead of $-c$. Choose a nonzero eigenline as $R$.

If $c$ is not a square, put $K=\mathbf F_q[u]/(u^2-c)$ and
let $u$ act as $h$ on $E$. The natural space has dimension $n$
over $K$. For $v,w\in E$ define the unique $K$-valued form
$b(v,w)$ by
\[
       \operatorname{Tr}_{K/\mathbf F_q}(a b(v,w))=B(av,w)
                    \quad\hbox{for every }a\in K.          \tag{HOS.23}
\]
Nondegeneracy of the trace pairing gives uniqueness and existence.
The adjoint of $h$ for $B$ is $(\mu/c)h$.
If $\mu=c$, the adjoint involution on $K$ is the identity.
Equation (HOS.23) is then $K$-bilinear and alternating:
$B(av,v)=B(v,av)=-B(av,v)$, so it vanishes for every $a$.
It is nondegenerate, and any $K$-line is isotropic. This case
automatically has $n$ even.

If $\mu=-c$, the adjoint involution on $K$ is its nontrivial
Galois involution. Equation (HOS.23) is linear in its first
variable, conjugate-linear in its second, and satisfies
$b(w,v)=-\overline{b(v,w)}$. Multiplying by any nonzero
trace-zero scalar makes it Hermitian. It is nondegenerate;
as $n\ge4$, its finite-field norm equation supplies an isotropic
$K$-line. In either case such a line, viewed over $\mathbf F_q$,
is an $h$-stable isotropic plane $R$ for $B$.

Use the shear subgroup (HOS.17) for this $R$. Conjugation is
now the normalized action
\[
                         C\longmapsto\mu^{-1}ACA^t,
                              \qquad A=h|_R.                \tag{HOS.24}
\]
This follows directly by moving $h$ through the pairing in
(HOS.17), which introduces the factor $\mu^{-1}$.
The action has square one. It is not the identity: otherwise,
over an algebraic closure, every eigenvalue square of the
semisimple $A$ and every pairwise product would be $\mu$.
All eigenvalues would coincide, so $A=\alpha I$ with
$\alpha\in\mathbf F_q$ and $\alpha^2=\mu$, contradicting
the nonsquare multiplier. Its minus-one eigenspace gives an
additive unipotent subgroup $U$ of order at least $q$, on which
$T$ exchanges all nonzero dual weights in pairs.
The averaging and cyclic-block calculation in (HOS.18), with
$\ell=2$, proves (HOS.22). The proof did not require $T^2$
to be the identity as an operator, only the indicated action
on the core and its shear subgroup.
\end{proof}

\begin{theorem}[D-HOS-05: every compatible normalizer]
\label{D-HOS-05}
For $q\ge7$ odd and $n\ge4$, every finite-order nonscalar
operator normalizing the image of a nontrivial ordinary irreducible
representation of $\operatorname{Sp}_{2n}(q)$ of degree $a$ satisfies
\[
                              \nu(T)\ge\delta_* a.          \tag{HOS.25}
\]
All scalar, diagonal and field possibilities are included.
The same conclusion holds for every quasisimple central quotient
with this simple quotient.
\end{theorem}

\begin{proof}
There are no exceptional multipliers in this range, by
\cite[Chapter~5, complete Table~1 and preceding paragraph,
p.~35]{BI:HHM}. Thus the universal perfect cover is the symplectic
group. An automorphism of a represented central quotient lifts
uniquely to that cover; an automorphism trivial on the simple
quotient is trivial, since its difference from the identity is a
homomorphism from a perfect group to its center. The exact
outer group is $C_2\times C_f$ for $q=p^f$, generated by diagonal
and field classes, from \cite[\S2.2.2, Table~2.1,
author p.~14]{HES:LiXia2022}. These descriptions therefore
apply to every normalizing operator under consideration.

The function $\delta_q$ is increasing for $q\ge7$, since each
of its three positive factors is increasing. Also
$\sqrt7>37/14$, because $1372>1369$. Hence
$c_7<14/23$ and
\[
 \delta_7>\tfrac12\frac67\frac{342}{343}\frac9{23}
                  =\frac{9234}{55223}=\delta_* .            \tag{HOS.26}
\]
Thus every noncentral inner operator already satisfies (HOS.25).

If $q$ is prime, there is no field component. For a diagonal
outer operator, if the square of its induced automorphism is a
nontrivial inner automorphism, $T^2$ is scalar times a noncentral
inner operator, and D-HOS-03 applies. If the square is the
identity, D-HOS-04 gives an even stronger estimate.

For every $q$, a general normalizer has support at least
$\delta_q a/2$. Indeed, a nonscalar $T$ fails to commute with
some represented core element by Schur's lemma. A noncentral
commutator has moving rank at least $\delta_q a$ by D-HOS-03;
a nonidentity scalar commutator has moving rank $a$.
After scalar-normalizing $T$ at its largest eigenspace,
\[
             \operatorname{rank}([T,g]-I)\le2\nu(T),
\]
which proves the claim in both cases. Consequently all $q\ge25$
are covered, because
\[
 \frac{\delta_q}{2}\ge\frac{\delta_{25}}2
 =\frac9{50}\frac{15624}{15625}>\delta_* .                  \tag{HOS.27}
\]
Every nonprime odd prime power at least seven is either nine
or at least twenty-five. It remains only to handle $q=9$.

The outer group at $q=9$ is $C_2\times C_2$. If the square of
the induced automorphism is nontrivial inner, use $T^2$ as
above. A pure diagonal involution is covered by D-HOS-04.
Consider a field-containing automorphism whose square is the
identity. It has a natural lift $h\sigma$, where
$h\in\operatorname{GSp}_{2n}(9)$ and $\sigma$ is coordinatewise
cubing. Its square is scalar, say
\[
                         h\sigma(h)=cI.
\]
Associativity gives $c=\sigma(c)$, hence $c\in\mathbf F_3^*$.
Writing $\mu$ for the multiplier of $h$ gives
$\mu\sigma(\mu)=c^2=1$. In particular $\mu$ is a square
in $\mathbf F_9^*$; a mixed diagonal-field class cannot have
identity square. We nevertheless allow an arbitrary such lift
in the following descent, so no representation is assumed
invariant under a separate diagonal conjugation.

Choose $z\in\mathbf F_9^*$ with $z\sigma(z)=c^{-1}$ and set
$s=zh\sigma$. Then $s^2=I$. Its multiplier coefficient
$\mu'=z^2\mu$ has norm one. Choose $d\in\mathbf F_9^*$ with
$\sigma(d)/d=\mu'$. Both choices follow from the cyclic
multiplicative group: the norm has image $\mathbf F_3^*$,
and $d\mapsto\sigma(d)/d=d^2$ has image the norm-one subgroup.
For $B'=dB$ one has
\[
                         B'(sv,sw)=\sigma(B'(v,w)).          \tag{HOS.28}
\]
The fixed space $E_0=E^s$ is $2n$-dimensional over $\mathbf F_3$
and spans $E$ over $\mathbf F_9$. Explicitly, choose
$\alpha\in\mathbf F_9$ with $\alpha^2=-1$, so
$\sigma(\alpha)=-\alpha$. Every vector has the decomposition
\[
             v=\frac{v+sv}{2}
                 +\alpha\frac{v-sv}{2\alpha},              \tag{HOS.29}
\]
whose two displayed coefficients belong to $E_0$.
Moreover $E_0\cap\alpha E_0=0$, which proves the dimension
and independence assertion. The restriction of $B'$ to $E_0$
is an alternating form over $\mathbf F_3$ and is nondegenerate,
since $E_0$ spans. A symplectic basis supplies a matrix
$P\in\operatorname{GSp}_{2n}(9)$ with
$P^{-1}sP=\sigma$. Replace the representation by
$x\mapsto\rho(PxP^{-1})$. Its intertwiner implements pure
$\sigma$, and it has the same degree and the same unipotent
estimate. This relabeling is permissible even if $P$ has
nontrivial diagonal class.

In the standard $\operatorname{GL}_n(9)$ Levi choose the
one-row subgroup
$U=\{\left(\begin{smallmatrix}1&u\\0&I_{n-1}\end{smallmatrix}\right)
:u\in\mathbf F_9^{\,n-1}\}$.
The block subgroup $\operatorname{diag}(1,\operatorname{GL}_{n-1}(9))$
normalizes $U$ and is transitive on its nonzero dual weights,
so each has the same multiplicity $m$. Put $Q=9^{n-1}$ and
$M=3^{n-1}$.
Since $c_9=1/2$, averaging (HOS.14) gives $m\ge a/(2Q)$.
Frobenius fixes exactly $M$ of the $Q$ dual weights, under
the nondegenerate trace-pairing identification of that dual.
All the remaining weights occur in pairs. Therefore
\[
 \nu(T)\ge\frac{Q-M}{2}m
      \ge\frac14(1-3^{-(n-1)})a
      \ge\frac{13}{54}a>\delta_*a.                          \tag{HOS.30}
\]
This proves the last field case. The same lifting argument
from the universal perfect cover proves the central-quotient
assertion. Every compatible finite-order normalizer is covered.
\end{proof}

\begin{theorem}[D-HOS-06: the higher odd symplectic source exclusion]
\label{D-HOS-06}
Suppose an inducing block in the signature-$(3,3)$ sixfold
curve-cover problem has an irreducibly acting normal central product
of quasisimple groups, each with simple quotient
$\operatorname{PSp}_{2n_i}(q_i)$, where $n_i\ge4$ and
$q_i\ge7$ is odd. Then the actual period image has dimension
at most eight. This allows different ranks and fields, arbitrary
compatible scalar and outer operators, and all factor permutations.
It does not require equality of the inducing and original character
fields, selfduality of any representation, or a determinant condition.
\end{theorem}

\begin{proof}
The core representation is a tensor product of nontrivial irreducible
factor modules. Let $e$ be its total degree. For a factor-preserving
nonscalar operator, its normalized eigenspace multiplicities are
convolutions. Their maximum is bounded by the maximum for any
nonscalar factor, so D-HOS-05 gives support at least $\delta_*e$.

For clarity, the same estimate holds when factors are permuted.
On a nontrivial factor cycle of length $r$ and factor degree $a$,
conjugate the intertwiners to identities except for one closing
twist and diagonalize that twist. The resulting operator is a
weighted cyclic permutation of the $a^r$ basis words. There are
$a$ constant words, and all other word orbits have length at least
two. Each orbit contributes at most one to any chosen eigenspace.
Thus the largest eigenspace is at most
$a+(a^r-a)/2\le3a^r/4$. Tensoring with the other cycles cannot
increase the normalized maximum. The resulting support $e/4$
is stronger than $\delta_*e$. Nonidentity scalar inertias move
all $e$ dimensions.

D-HOS-02 gives for every factor
\[
             a_i\ge\frac{q_i-1}{2}q_i^{n_i-1}
                       \ge3\cdot7^3=1029,
\]
so $e\ge1029$. Apply the actual intermediate-cover factorization:
if its base has genus $h$ and $R$ genuine inertia points, then
\[
              6=2e(h-1)+D,\qquad
                  D=\sum_{j=1}^{R}\operatorname{rank}(g_j-I).
                                                               \tag{HOS.31}
\]
This is the already established polarized multiplicity variation
factorization, valid also for a larger inducing character field.
Genus at least two is impossible because $D\ge0$ and $e>3$.
In genus one, $D=6$, hence $R\le6$ and the actual pointed-base
Hurwitz variation has dimension at most six. A period image of
dimension nine would consequently force genus zero and
$R\ge12$, with $D=2e+6$.

Every genuine inertia has moving rank at least $\delta_*e$.
It follows that
\[
       2e+6\ge12\delta_*e,\qquad
       e\le\frac6{12\delta_*-2}
           =\frac{165669}{181}<916,                          \tag{HOS.32}
\]
contradicting $e\ge1029$. This is a necessary actual-period
bound and does not infer dominance from parameter counts.
\end{proof}

\paragraph{Remaining fields and source boundary.}
The argument makes no assertion for higher-rank odd symplectic factors
over fields $3$ or $5$. The low-rank $C_2$ and $C_3$ results retain
their separately checked scopes. The published inputs here are precisely
the ordinary unipotent ratio, the ordinary automorphism description,
and the complete exceptional-multiplier list. All carrier, isotropic,
quadratic-form, semilinear-descent, tensor, and numerical arguments
are supplied in full above or in D-HOS-01--02. No missing semisimple
ratio theorem, unrefereed result, or full generic character table is used.

\section{Higher symplectic cores over the field of five elements}
\label{sec:higher-symplectic-five}

\paragraph{Status and scope.}
This complete new elementary proof has passed independent whole-proof
reviews by /root and /root/imprimitive. Both also read the full author
checker; the peer separately reconstructed the exact finite arithmetic.
Its character-theoretic construction is the explicit
Heisenberg matrix algebra in D-HOS-01, extended here to automorphisms
with fixed vectors and to defining-characteristic prime order. No
unipotent character bound or full character table is used. The final
source obstruction allows mixed ranks, arbitrary compatible outer
automorphisms, scalar twists, and permutations of tensor factors.
It uses the imaginary character-field inclusion from the original
induced representation, not equality of the original and inducing fields.

For a finite-order complex operator $T$ on a space of dimension $a$, write
\[
 \nu(T)=a-\max_\lambda\dim\ker(T-\lambda I).
                                                               \tag{HOS5.1}
\]
This quantity is unchanged by scalar multiplication, is no larger than
ordinary moving rank, and satisfies $\nu(T^j)\le\nu(T)$.

\begin{lemma}[D-HOS-07: opposite line radicals]
\label{D-HOS-07}
Let $q$ be odd, $n\ge2$, and let $V$ be a nontrivial irreducible
ordinary representation of $G=\operatorname{Sp}_{2n}(q)$ of degree $a$.
If $Q$ is the unipotent radical of a line parabolic, then
\[
                         \dim V^Q\le a/2.                     \tag{HOS5.2}
\]
\end{lemma}

\begin{proof}
Let the natural space be $E=\langle e,f\rangle\perp W$, with
$B(e,f)=1$. For $w\in W$ and $z\in\mathbf F_q$ put
\[
 u_e(w,z)e=e,\quad
 u_e(w,z)x=x+B(w,x)e\ (x\in W),\quad
 u_e(w,z)f=f+w+2ze.                                          \tag{HOS5.3}
\]
Direct substitution in the alternating form shows these are symplectic.
Their multiplication is
\[
 u_e(w,z)u_e(w',z')
 =u_e(w+w',z+z'+B(w,w')/2).
\]
They form $Q_e=Q$. Define $Q_f$ by the same formula with the paired
vectors $(f,-e)$ and the same $W$. The symplectic map
$e\mapsto f$, $f\mapsto-e$, $W\mapsto W$ conjugates $Q_e$ to $Q_f$.
We prove directly that these two groups generate $G$.

For $v\in E$ let
\[
                       t_v(c)x=x+cB(v,x)v.                    \tag{HOS5.4}
\]
The centers of $Q_e,Q_f$ contain $t_e(c),t_f(c)$ for every $c$.
Conjugating $t_f(c)$ by (HOS5.3) gives every transvection in a
direction $f+w+de$, with $w\in W,d\in\mathbf F_q$. The opposite
calculation gives every direction $e+w+df$. For $0\ne w\in W$,
choose $w_0\in W$ with $B(w_0,w)=-1$. Then
\[
                         u_f(w_0,0)(f+w)=w.
\]
Consequently transvections in direction $w$ also belong to
$\langle Q_e,Q_f\rangle$. Scaling a direction scales the parameter
by its square, and the parameter already ranges over the whole field.
Every nonzero vector is either in $W$ or a scalar multiple of one of
the displayed directions. Thus this generated group contains all
symplectic transvections.

For completeness these transvections generate $G$ by an elementary
induction. If $B(x,y)\ne0$, a suitable transvection in direction
$y-x$ sends $x$ to $y$. If $B(x,y)=0$, choose a vector outside the
two orthogonal hyperplanes and use two such steps; two proper
hyperplanes cannot cover a vector space over an odd field. Hence
one can first make a given symplectic map fix $e$. Its image of $f$
then has the form $f+w+de$ and can be moved to $f$ by an element of
$Q_e$. These radical elements are themselves products of
transvections: since $B(e,w)=0$, direct multiplication gives
\[
 u_e(w,0)=t_{w+e}(1)t_w(-1)t_e(-1),\qquad
 u_e(w,z)=t_e(2z)u_e(w,0).                                   \tag{HOS5.5}
\]
After fixing both $e$ and $f$, continue on $W$. This includes the
two-dimensional base case, where the remaining orthogonal complement
is zero. It proves $\langle Q_e,Q_f\rangle=G$.

The two fixed subspaces in $V$ have equal dimension. Their intersection
is $V^G=0$, because $V$ is nontrivial and irreducible. Their dimensions
therefore sum to at most $a$, proving (HOS5.2).
\end{proof}

\begin{lemma}[D-HOS-08: carriers with a fixed subspace]
\label{D-HOS-08}
In the notation of D-HOS-07, suppose a finite-order operator $T$ on
$V$ acts on the full line radical by
\[
             T\rho(u_e(w,z))T^{-1}=\rho(u_e(gw,z)),
                                                               \tag{HOS5.6}
\]
where $g\in\operatorname{Sp}(W)$ has prime order $\ell$. Put
$r=\operatorname{rank}(g-I)>0$. The prime $\ell$ may equal the
characteristic. Then
\[
 \nu(T)\ge (1-1/\ell)(1-q^{-r/2})(a-\dim V^{Q_e})
          \ge\tfrac12(1-1/\ell)(1-q^{-r/2})a.                \tag{HOS5.7}
\]
\end{lemma}

\begin{proof}
We retain the actual group law in (HOS5.3), so no choice of projective
intertwiner changes the specified automorphism. For a nontrivial
central additive character $\psi$, D-HOS-01 constructs the unique
irreducible carrier $U_\psi$ of dimension $N=q^{n-1}$ and its
operators $L_w$. These satisfy
\[
 L_wL_{w'}=\psi(B(w,w')/2)L_{w+w'},\qquad
 \{L_w:w\in W\}\text{ is a basis of }\operatorname{End}(U_\psi).
                                                               \tag{HOS5.8}
\]
The proof there is the explicit translation-and-multiplication model
and the central-character group-algebra dimension calculation; it
does not use the order of $g$ or a fixed-point-free hypothesis.
Uniqueness gives an exact intertwiner
$\Omega L_w\Omega^{-1}=L_{gw}$. Its $\ell$th power is scalar;
rescale it so that $\Omega^\ell=I$.

Every $g^j$, $1\le j<\ell$, has the same fixed space as $g$.
Indeed $g$ is a power of $g^j$ in the cyclic group of prime order.
Conjugation by $\Omega^j$ permutes the basis in (HOS5.8), whence
\[
 |\operatorname{Tr}\Omega^j|^2
 =\operatorname{Tr}(\operatorname{Ad}\Omega^j)
 =q^{2n-2-r},\qquad
 |\operatorname{Tr}\Omega^j|/N=q^{-r/2}.                     \tag{HOS5.9}
\]
The first identity is valid because $\Omega$ has finite order.
Fourier inversion bounds any one eigenspace of $\Omega$ by
\[
              \frac{1+(\ell-1)q^{-r/2}}{\ell}\,N.
                                                               \tag{HOS5.10}
\]
On the full $\psi$-isotypic summand, the operator is
$\Omega\otimes A_\psi$ by the full matrix-algebra commutant.
The multiplicity operator is finite order after the above
normalization, since a common multiple of the orders of $T$ and
$\Omega$ kills it. Its eigenvalue multiplicities convolve with
those of $\Omega$, so (HOS5.10) remains an upper bound on the
normalized multiplicity. No scalar-power hypothesis on $T$ over
the whole of $V$ is being made.

The remaining nontrivial irreducibles of $Q_e$ are its nonzero
linear characters, parametrized by the nonzero vectors in the
dual of $W$ using the finite-field trace pairing. Their
multiplicities in $V$ are all equal: the subgroup
$\operatorname{Sp}(W)\le G$ fixing $e,f$ normalizes $Q_e$ and is
transitive on these nonzero dual vectors. Of the $q^{2n-2}-1$
weights, exactly $q^{2n-2-r}-1$ are fixed by $g$. Their proportion
is less than $q^{-r}$. Every other orbit has length $\ell$;
on such a cyclic orbit each eigenvalue has multiplicity at most
$1/\ell$ of the total orbit dimension, even with arbitrary
finite-order multiplicity blocks. Thus the support fraction on
all nonzero linear weights is at least
\[
                (1-1/\ell)(1-q^{-r})
                 \ge(1-1/\ell)(1-q^{-r/2}).
\]
Give the trivial $Q_e$-isotypic summand its worst possible full
eigenspace. Adding the preceding bounds gives the first inequality
in (HOS5.7), and D-HOS-07 gives the second. This argument includes
$\ell=p$: neither semisimplicity of the natural action nor a
coprime averaging projector was used.
\end{proof}

\begin{lemma}[D-HOS-09: every compatible normalizer over $\mathbf F_5$]
\label{D-HOS-09}
Let $n\ge4$, and let $\rho$ be any nontrivial ordinary irreducible
representation of $G=\operatorname{Sp}_{2n}(5)$ of degree $a$.
Every nonscalar finite-order operator $T$ normalizing $\rho(G)$
satisfies
\[
                          \nu(T)\ge\frac{31}{250}a.           \tag{HOS5.11}
\]
This includes every compatible outer automorphism and both possible
central quotients of $G$.
\end{lemma}

\begin{proof}
First record three consequences of D-HOS-08. If $s$ is a natural
involution with minus-one eigenspace a nondegenerate plane, choose
$e,f$ in that plane. On the associated line radical, conjugation is
$(w,z)\mapsto(-w,z)$. Thus
\[
 \nu(\rho(s))\ge\frac{1-5^{-(n-1)}}4a
                          \ge\frac{31}{125}a.               \tag{HOS5.12}
\]
If $s$ is any noncentral natural involution, choose $e,f$ in its
smaller nonzero eigenspace. Both eigenspaces are nondegenerate and
even dimensional. Writing its eigenvalue on that plane as
$\lambda=\pm1$, the radical action is $w\mapsto\lambda s|_Ww$,
which moves the whole larger eigenspace of dimension $2m$ with
$m\ge\lceil n/2\rceil\ge2$. Therefore
\[
                         \nu(\rho(s))\ge\frac6{25}a.        \tag{HOS5.13}
\]

If $u=I+N\in G$ has $N^2=0$ and $\operatorname{rank}N=2$, its
image $R$ is isotropic and $N$ is skew-adjoint. To verify this
directly, apply preservation of $B$ by $I+jN$ for $j=1,2$:
subtracting twice the first identity from the second gives
$2B(Nx,Ny)=0$, and the first then gives skew-adjointness.
It follows that $\ker N=R^\perp$. Complete $R$ to a nondegenerate
four-dimensional symplectic subspace $E_0$; its orthogonal
complement lies in $\ker N$. Choose a paired $e,f$ in this
complement, which has dimension at least four. On its line radical
the induced action is $u|_W$, of order five and moving rank two.
Consequently
\[
                         \nu(\rho(u))\ge\frac8{25}a.        \tag{HOS5.14}
\]

We now pass to an arbitrary compatible normalizer. The ordinary
universal cover and automorphism inputs are exactly
\cite[Chapter~5, Table~1, p.~35]{BI:HHM} and
\cite[\S2.2.2, Table~2.1, author p.~14]{HES:LiXia2022}.
In this range the universal cover is $G$, and all automorphisms of
the simple quotient are inner-diagonal: the prime field has no
field automorphisms and type $C_n$, $n\ge4$, has no graph
automorphism. Thus the automorphism implemented by $T$ lifts to
conjugation by an $h\in\operatorname{GSp}(E)$; put
$\mu=\mu(h)\in\mathbf F_5^*$.
This assertion also holds when $\rho$ kills the central sign.
The kernel of a nontrivial irreducible representation is central,
and automorphisms of its image induce automorphisms of the simple
quotient. The natural lift and the given automorphism coincide:
their difference, if trivial on the simple quotient, would be a
central character of a perfect group. Such a character is trivial.
Hence
\[
                     T\rho(x)T^{-1}=\rho(hxh^{-1})
                            \quad(x\in G).                  \tag{HOS5.15}
\]

For $0\ne v\in E$ and $c\in\mathbf F_5^*$ use the actual
transvection (HOS5.4). Conjugation gives
\[
 h t_v(c)h^{-1}=t_{hv}(c/\mu),\qquad
 k=[h,t_v(c)]=t_{hv}(c/\mu)t_v(-c).                         \tag{HOS5.16}
\]
On the complex representation its commutator is $\rho(k)$.
For any chosen eigenvalue $\alpha$ of $T$,
\[
 \operatorname{rank}(\rho(k)-I)
 =\operatorname{rank}(T\rho(t_v(c))-\rho(t_v(c))T)
 \le2\operatorname{rank}(T-\alpha I).
                                                               \tag{HOS5.17}
\]
The equality discards invertible right factors. Choosing the
largest eigenspace bounds this by $2\nu(T)$.

Suppose first some $v$ has $b=B(v,hv)\ne0$. The plane
$P=\langle v,hv\rangle$ is nondegenerate, $k$ is the identity
on $P^\perp$, and in the indicated ordered basis its matrix is
\[
 k|_P=
 \begin{pmatrix}1&-cb\\-cb/\mu&1+c^2b^2/\mu\end{pmatrix},
 \qquad \det(k|_P)=1,\quad
 \operatorname{Tr}(k|_P)=2+c^2b^2/\mu.                     \tag{HOS5.18}
\]
If $\mu$ is a square, choose $c$ so that $c^2b^2/\mu=4$;
the trace is then one, and Cayley--Hamilton gives $k^3=-I$ on
$P$. If $\mu$ is a nonsquare, choose $c$ so that this ratio is
three; the trace is then zero, and $k^2=-I$ on $P$. These choices
exist because the nonzero squares in $\mathbf F_5$ are $1,4$
and the nonsquares are $2,3$. In either case a power of $k$ is
the rank-two involution in (HOS5.12). Therefore
\[
 2\nu(T)\ge\operatorname{rank}(\rho(k)-I)
          \ge\nu(\rho(k))
          \ge\frac{31}{125}a.                              \tag{HOS5.19}
\]

It remains that $B(v,hv)=0$ for every $v$. Polarizing this
identity gives $B(x,hy)=B(hx,y)$, so $h$ is self-adjoint for
$B$. Its adjoint as a similitude is $\mu h^{-1}$; consequently
$h^2=\mu I$.
If $\mu$ is a square, rescale $h$ to an involution $s\in G$.
Equation (HOS5.15) and Schur's lemma make $T$ a scalar multiple
of $\rho(s)$. If $s$ were central then $T$ would be scalar;
otherwise (HOS5.13) is stronger than (HOS5.11).

If $\mu$ is nonsquare, $h$ has no eigenline over $\mathbf F_5$.
For any $v\ne0$ the vectors $v,hv$ are independent and span a
totally isotropic plane. Equation (HOS5.16) now has the form
\[
 k=I+N,\qquad
 N=\frac c\mu hv\,B(hv,-)-c v\,B(v,-).
                                                               \tag{HOS5.20}
\]
The two pairing functionals and the two image vectors are
independent. Thus $\operatorname{rank}N=2$, while isotropy
gives $N^2=0$. Equations (HOS5.14) and (HOS5.17) give
$\nu(T)\ge4a/25>31a/250$. All cases have now been covered.
\end{proof}

\begin{theorem}[D-HOS-10: no higher $\mathbf F_5$ source]
\label{D-HOS-10}
In the genus-zero finite-group source problem for an imaginary-quadratic
Weil sixfold, no inducing representation whose image has an
irreducibly acting normal central product of nontrivial quasisimple
factors with simple quotients
\[
               \operatorname{PSp}_{2n_i}(5),\qquad n_i\ge4,
                                                               \tag{HOS5.21}
\]
can have period rank nine. The factors may have different ranks;
no transitivity on them is required. This conclusion concerns either
polarization determinant and remains valid if the inducing character
field is larger than the original imaginary quadratic field.
\end{theorem}

\begin{proof}
Let the inducing image be the finite faithful irreducible group
$J\le\operatorname{GL}(V)$ and put $e=\dim V$.
Restriction to its normal irreducible central product is a tensor
product of nontrivial ordinary irreducibles of the covers
$\operatorname{Sp}_{2n_i}(5)$. The elementary D-HOS-02 bound gives
\[
                   a_i\ge2\cdot5^{n_i-1}\ge250,
                      \qquad e=\prod_i a_i\ge250.           \tag{HOS5.22}
\]

First extend (HOS5.11) to all nonscalar elements of $J$. A
factor-preserving normalizer is a scalar times a tensor product
of compatible factor normalizers, by the full matrix-algebra
commutant. Its largest normalized eigenvalue multiplicity is at
most that of any one nonscalar factor. The factor operators can
be rescaled to finite order, since a power of their tensor
product is scalar. If an element permutes factors, consider any
nontrivial cycle of length $s\ge2$ on equal-dimensional factors
of degree $b\ge2$. By factorwise changes of basis its operator
is a cyclic rotation with one final diagonal twist. In the
tensor basis, the $b$ constant words are the only fixed words;
every other word orbit has at least two elements. A weighted
cyclic orbit contributes at most one vector to any one eigenspace.
Thus the largest eigenspace on that cycle has dimension at most
\[
                 b+\frac{b^s-b}{2}\le\frac34 b^s.
\]
Convolution with the remaining cycles preserves this upper bound.
In both cases every nonscalar element of $J$ has support at least
$31e/250$.

We next prove essential selfduality, including its source-relevant
nontrivial multiplier. The published factorization theorem of
Wonenburger, stated with its full field hypothesis in
\cite[Theorem~1, author p.~1]{HOS:Vinroot2004}, says that any
$x\in\operatorname{Sp}_{2n}(5)$ is a product of two
anti-symplectic involutions. If $x=h_1h_2$, then
$h_1xh_1=x^{-1}$. For a fixed anti-symplectic involution $h_0$,
the product $h_1h_0$ is symplectic, so $h_0xh_0$ is
symplectically conjugate to $x^{-1}$. Since $-1=2^2$ in
$\mathbf F_5$, the map $2h_0$ is itself symplectic and has the
same conjugation action. Therefore every $x$ is conjugate to
$x^{-1}$ in its symplectic group. All its ordinary irreducible
characters are real-valued, and their representations are
selfdual. No assertion about their Frobenius--Schur signs is needed.

Their tensor product has a unique nondegenerate invariant bilinear
form up to scalar, by irreducibility. Every normalizer in $J$
therefore preserves that line and has a bilinear similitude
character $\eta:J\to\mathbf C^*$. Let $L$ be the field of the
inducing character. The original character is induced from it,
so its values are sums of values in $L$; consequently the original
imaginary quadratic field $K$ is contained in $L$. If $\eta$
were trivial, the whole $J$-representation would be selfdual
and its character real-valued, contradicting $K\subset L$.
Thus $\eta$ is nontrivial.

We finish with the actual variation-of-Hodge-structure factorization.
The polarized Shapiro identification of D-BI-01 realizes the
six-dimensional multiplicity variation on a finite Hurwitz family
over the intermediate quotient curve $X$. If its genus is $h$
and its genuine branch inertias in $J$ are $g_1,\ldots,g_R$,
parabolic Euler characteristic gives
\[
        \sum_{j=1}^R\operatorname{rank}(g_j-I)
                        =6+2e(1-h).                         \tag{HOS5.23}
\]
Trivial inertia is omitted, so $R$ is at most this sum. For
$h\ge2$ the right side is negative because $e\ge250$. For
$h=1$ it is six, and the actual period map factors through a
finite Hurwitz cover of $\mathcal M_{1,R}$ of dimension at
most six. Rank nine therefore forces $h=0$, and then the
same actual factorization through $\mathcal M_{0,R}$ forces
$R\ge12$. These are bounds on the moduli space through which the
actual intermediate period map factors.

The genus-zero inertias generate $J$ and have product one.
Because $\eta$ is nontrivial, at least two have nontrivial
similitude. Their fixed subspaces are totally isotropic for the
nondegenerate bilinear form: for fixed vectors $v,w$,
$B(v,w)=B(gv,gw)=\eta(g)B(v,w)$. Each such inertia has moving
rank at least $\lceil e/2\rceil$. Every remaining genuine
inertia has moving rank at least $31e/250$: this follows from
the support bound if it is nonscalar, and from its full moving
rank $e$ if it is a nonidentity scalar. Applying (HOS5.23) to
two of the former and at least ten of the latter yields
\[
 2e+6\ge2\lceil e/2\rceil+10\frac{31e}{250}
          \ge\frac{56e}{25},\qquad e\le25.                 \tag{HOS5.24}
\]
This contradicts (HOS5.22) and proves the theorem.
\end{proof}

\paragraph{Certificate scope.}
The adjacent exact script checks the transvection trace and power
identities, all nonzero field parameters, the line-radical and
opposite-radical identities, the two self-adjoint multiplier types,
the constant inequalities, and the source-budget arithmetic. The
universal generation, carrier, normalizer, and source-factorization
proofs are the written arguments above. Finite matrix checks are
not substituted for them.

\section{The selfdual higher symplectic gate over three elements}
\label{sec:higher-symplectic-three}

\paragraph{Status and limitation.}
This new proof has passed complete mathematical reviews by /root and
/root/imprimitive. Both read the entire author checker and its output;
the peer also supplied a separately authored exact arithmetic certificate.
It uses the complete
opposite-radical and fixed-subspace carrier proofs D-HOS-07--08.
The normalizer bound below applies to every ordinary irreducible
representation. Its source consequence requires selfduality of the
irreducible core representations over $\mathbf F_3$. No reality of
all symplectic characters over this field is assumed. Nonselfdual
core factors remain outside the final obstruction.

\begin{lemma}[D-HOS-11: full normalizers over $\mathbf F_3$]
\label{D-HOS-11}
Let $n\ge4$ and let $\rho$ be a nontrivial ordinary irreducible
representation of $\operatorname{Sp}_{2n}(3)$ of degree $a$.
Every nonscalar finite-order operator $T$ normalizing its image
satisfies
\[
                              \nu(T)\ge\frac{13}{108}a.
                                                               \tag{HOS3.1}
\]
Both central quotients and all compatible outer automorphisms
are included.
\end{lemma}

\begin{proof}
The same complete universal-cover and automorphism inputs as in
D-HOS-09 apply, now with $q=3$: the universal cover is the
symplectic group and all automorphisms of the simple quotient are
inner-diagonal. There is no field or graph automorphism in this
range. The perfectness argument in (HOS5.15) therefore supplies a
natural $h\in\operatorname{GSp}_{2n}(3)$ implementing $T$ exactly
on the represented group; put $\mu=\mu(h)\in\{1,-1\}$.

Let $s$ be minus the identity on a nondegenerate natural plane and
the identity on its orthogonal complement. Choosing the line
radical from that plane, D-HOS-08 with $g=-I$ on its full
quotient gives
\[
               \nu(\rho(s))\ge\frac{1-3^{-(n-1)}}4a
                                      \ge\frac{13}{54}a.
                                                               \tag{HOS3.2}
\]
For any noncentral natural involution, choose the radical plane
in its smaller nonzero eigenspace. The larger eigenspace has
dimension $2m$ with $m\ge2$. The same argument gives
\[
                            \nu(\rho(s))\ge\frac29a.         \tag{HOS3.3}
\]

Suppose some $v$ has $b=B(v,hv)\ne0$. For any nonzero
$c\in\mathbf F_3$ the commutator
\[
 k=[h,t_v(c)]=t_{hv}(c/\mu)t_v(-c)
\]
acts as the identity outside $P=\langle v,hv\rangle$. Its
determinant on $P$ is one and its trace is
$2+c^2b^2/\mu=2+\mu^{-1}$. If $\mu=1$, this trace is zero
and $k^2=-I$ on $P$. If $\mu=-1$, its characteristic polynomial
on $P$ is $(X+1)^2$, and in characteristic three $k^3=-I$ on
$P$. Thus in either case a power is the rank-two involution
in (HOS3.2). The ordinary rank-commutator inequality (HOS5.17)
then gives $2\nu(T)\ge13a/54$.

Now suppose $B(v,hv)=0$ for all $v$. Polarization again gives
self-adjointness and $h^2=\mu I$. For $\mu=1$, Schur's lemma
makes $T$ a scalar multiple of a represented natural involution.
The central case would make $T$ scalar; otherwise (HOS3.3)
is stronger than the required bound.

It remains $\mu=-1$. For any $v\ne0$, put $r_1=v,r_2=hv$.
These are independent and span a totally isotropic plane. For
nonzero $c$, the same commutator is
\[
 k=I+N,\qquad
 N=-c\{r_1B(r_1,-)+r_2B(r_2,-)\},\quad
                 N^2=0,\quad\operatorname{rank}N=2.         \tag{HOS3.4}
\]
We need the equal coefficients in this expression, rather than
the absolute-trace bound for an arbitrary rank-two shear.

Complete the image plane to a nondegenerate natural four-space
$E_0$ as in D-HOS-09, and choose a paired $e,f$ in its
orthogonal complement. For the full line radical $Q_e$, the
induced quotient action is $g=k|_W$, of order three and moving
rank two. Choose symplectic coordinates $W=X\oplus Y$ with
$r_1,r_2$ the first two coordinate vectors in $Y$. Writing
$m=n-1$, this action is
\[
 g(x,y)=\bigl(x,y+c(x_1,x_2,0,\ldots,0)\bigr).
                                                               \tag{HOS3.5}
\]
In the explicit Schr\"odinger model (HOS.5), multiplication by
\[
 (\Omega f)(u)=
   \psi\bigl(-c(u_1^2+u_2^2)/2\bigr)f(u)
                                                               \tag{HOS3.6}
\]
implements (HOS3.5) exactly. Indeed conjugating $L_{x,y}$ by
the displayed multiplication operator gives $L_{x,y+cx'}$,
where $x'=(x_1,x_2,0,\ldots,0)$; the quadratic constant term
is exactly the required half-pairing term in $L_{x,y+cx'}$.

For either nontrivial central character, the values of the
quadratic exponent on its first two coordinates have
multiplicities $1,4,4$: the sum of two squares over
$\mathbf F_3$ is zero only at $(0,0)$, is one at four
vectors, and is minus one at four vectors. Multiplication by
a nonzero coefficient only interchanges the latter values.
The other $m-2$ coordinates are free. Thus every eigenspace
of this carrier has dimension at most $4/9$ of its degree.
As in D-HOS-08, tensoring with its arbitrary multiplicity
operator preserves this upper bound.

Among the nonzero linear characters of $Q_e$, all
multiplicities are equal under the full Levi action. The
proportion fixed by $g$ is less than $1/9$, and the other
orbits have length three. Hence their total support fraction
is at least $(2/3)(8/9)=16/27>5/9$. Combining the linear and
nontrivial-central-character summands, and then using
D-HOS-07, proves
\[
              \nu(\rho(k))\ge\frac59(a-\dim V^{Q_e})
                                     \ge\frac5{18}a.
                                                               \tag{HOS3.7}
\]
Equation (HOS5.17) gives $\nu(T)\ge5a/36>13a/108$.
This exhausts the self-adjoint case and proves the lemma.
\end{proof}

\begin{theorem}[D-HOS-12: selfdual mixed three-and-five cores]
\label{D-HOS-12}
In the imaginary-quadratic Weil-sixfold source problem, an inducing
image with an irreducibly acting normal central product of factors
of the following kinds cannot have period rank nine:
\begin{enumerate}
\item simple quotient $\operatorname{PSp}_{2n}(3)$, $n\ge4$,
with a nontrivial selfdual ordinary irreducible core representation;
\item simple quotient $\operatorname{PSp}_{2n}(5)$, $n\ge4$,
with any nontrivial ordinary irreducible core representation.
\end{enumerate}
Ranks and the two field sizes may be mixed; there is no transitivity
or original/inducing-field-equality assumption. Both polarization
determinants are included.
\end{theorem}

\begin{proof}
All core representations in the statement are selfdual; for the
second kind this was proved in D-HOS-10 from the published
anti-symplectic factorization. The elementary degree bound
D-HOS-02 gives degree at least $27$ for each factor of the first
kind and at least $250$ for each of the second. Thus the total
inducing degree $e$ is at least $27$.

The factorwise convolution and weighted tensor-cycle argument
in D-HOS-10 applies unchanged. Because
$13/108<31/250<1/4$, every nonscalar element of the full
compatible normalizer has projective support at least $13e/108$.
The unique core-invariant bilinear line is preserved by the
inducing image. Its similitude character is nontrivial because
the original imaginary quadratic character field is contained
in the inducing character field. The proof of these two
assertions is precisely the bilinear-form and character-induction
argument in D-HOS-10; it does not require that each inducing
character value already lie in the original quadratic field.

The actual intermediate Hurwitz factorization forces genus zero
and at least twelve genuine inertias, since $e\ge27$ and the
multiplicity variation has dimension six. At least two inertias
have nontrivial similitude and moving rank at least
$\lceil e/2\rceil$. Each remaining genuine inertia has moving
rank at least $\lceil13e/108\rceil$, including the full-rank
nonidentity scalar case. Consequently
\[
 2e+6\ge2\lceil e/2\rceil
                     +10\lceil13e/108\rceil.                \tag{HOS3.8}
\]
Dropping ceilings first gives $e\le324/11<30$. The elementary
degree lower bound leaves only $e=27,28,29$. For all three,
$\lceil13e/108\rceil=4$, so the right side of (HOS3.8) is at
least $68$, whereas the left side is at most $64$. This is a
contradiction. Nonselfdual higher $\mathbf F_3$ factors have
not been excluded by this argument.
\end{proof}

\section{A finite natural-rank cutoff for the remaining characteristic-three symplectic core}
\label{sec:symplectic-three-cutoff}

\paragraph{Status and scope.}
This new derivation has passed whole-proof reviews by /root and
/root/imprimitive, with independent exact arithmetic recorded separately.
The first
lemma is an all-rank elementary support estimate, including every
order-three unipotent. The final cutoff additionally uses the published
Larsen--Tiep estimate with its previously checked explicit constant.
It does not classify the nonselfdual representations at the remaining
smaller ranks. Put
\[
 S=\operatorname{Sp}_{2n}(3),\qquad n\ge4,\qquad
 \nu(T)=a-\max_\lambda\dim\ker(T-\lambda I)
 \tag{HOS3C.1}
\]
for an operator on a complex space of dimension $a$.
All representations below are ordinary complex representations.

\begin{lemma}[D-HOS-13: every element of order three]
\label{D-HOS-13}
For every nontrivial irreducible representation $\rho$ of $S$ of
degree $a$, and every element $h\in S$ of order three,
\[
                         \nu(\rho(h))\ge2a/9.
 \tag{HOS3C.2}
\]
No selfduality or Weil-module hypothesis is imposed.
\end{lemma}

\begin{proof}
We first treat a transvection. Fix a decomposition of the natural
symplectic space into paired Lagrangians $X\oplus Y$, each of dimension
$n$. The upper Siegel group
\[
 U=\{u_C:(x,y)\mapsto(x+Cy,y)\mid C=C^t\in M_n(\mathbf F_3)\}
 \tag{HOS3C.3}
\]
is abelian. It and its opposite generate $S$. Here are explicit
elementary details for this assertion. They contain the transvections
in every $e_i$ and $f_i$ direction. The usual three-transvection
product in $\langle e_i,f_i\rangle$ interchanges those two directions
up to sign. Conjugating an upper symmetric off-diagonal elementary
shear by this interchange in one coordinate gives a Levi elementary
matrix $I+cE_{ij}$, paired with its inverse transpose on $Y$.
The square of a coordinate interchange gives the Levi diagonal with
one entry $-1$. Elementary matrices and these signs generate
$\operatorname{GL}_n(3)$, so the generated group contains that Levi.
For any $y\ne0$ and any $x$, a symmetric $C$ with $Cy=-x$ exists:
change the paired coordinates so $y$ is the first coordinate vector,
prescribe the first column and row of a symmetric matrix, and change
back by congruence. Thus the shears, Levi, and coordinate interchanges
are transitive on all nonzero natural vectors. Conjugating a long-root
transvection then gives all symplectic transvections. The elementary
transvection generation proof in D-HOS-07 gives the assertion.

The opposite groups are conjugate. Their fixed spaces in $\rho$
therefore have the same dimension, and their intersection is zero,
since $\rho$ is a nontrivial irreducible $S$-module. Hence
\[
                             t:=\dim\rho^U\le a/2.
 \tag{HOS3C.4}
\]
Write $\omega=\exp(2\pi i/3)$. The dual of $U$ consists of the
characters $C\mapsto\omega^{\operatorname{tr}(BC)}$, with $B$
symmetric; the trace pairing is nondegenerate since two is invertible.
The Levi permutes these weights by congruence, so their multiplicities
are constant on every congruence orbit. For the transvection
$u_{cE_{11}}$, $c\ne0$, the distribution of a nonzero weight orbit
on the three eigenvalues equals the distribution of
$c\,v^tBv$ as $v$ runs through the nonzero vectors of $\mathbf F_3^n$.
Indeed, a uniformly chosen change-of-basis matrix has its relevant
row uniformly distributed on those vectors; quotienting by the
stabilizer of $B$ gives the same uniform orbit measure.

Every nonzero symmetric form can be diagonalized. A rank-one form
has value proportions $(1,2,0)/3$, up to interchange of the nonzero
values, when all vectors are counted. For rank at least two, two
nonzero diagonal coordinates have value proportions either
$(5,2,2)/9$ or $(1,4,4)/9$, again up to interchange. Adding the
remaining independent coordinates convolves these distributions
and cannot increase their largest entry. It follows, after removing
the zero vector, that the zero-value proportion of any nonzero form
is at most $5/9$, and every nonzero-value proportion is at most
\[
                 \frac{2\cdot3^{n-1}}{3^n-1}\le\frac79.
 \tag{HOS3C.5}
\]
For the zero proportion, removal changes a numerator at most
$5\cdot3^{n-2}$ to at most $5\cdot3^{n-2}-1$, which only improves
the bound. The rank-one case is smaller. For a nonzero value,
the rank-one bound $2\cdot3^{n-1}$ also dominates the rank-at-least-two
bound, and the last inequality follows from $3^n\ge7$.

Thus the $1$-eigenspace of a transvection has dimension at most
$t+(5/9)(a-t)\le7a/9$. Either other eigenspace has dimension at
most $(7/9)(a-t)\le7a/9$. This proves (HOS3C.2) for every
transvection, including both nonzero root parameters.

We now record a full-line Heisenberg estimate allowing the phases
which occur when a fixed line has no fixed paired vector. Let $x\in S$
have prime order $\ell$, fix a nonzero natural vector $e$, and act
on $e^\perp/\langle e\rangle$ as $g$ of rank
$r_0=\operatorname{rank}(g-I)>0$. Suppose every nontrivial power
$g^j$, $1\le j<\ell$, has that same moving rank. In the full line
radical $Q$, write its elements as $(v,z)$ with
$v\in\mathbf F_3^{2n-2}$ and $z\in\mathbf F_3$, using the
Heisenberg law of D-HOS-01. Conjugation by $x$ has the form
\[
                   (v,z)\longmapsto(gv,z+\lambda(v))
 \tag{HOS3C.6}
\]
for a linear functional $\lambda$: the center is fixed, and a
center-fixing Heisenberg automorphism with given symplectic quotient
differs from the standard lift by an inner automorphism.
This last assertion is also immediate from the law: its central
correction is additive, and the nondegenerate alternating pairing
represents every additive functional over $\mathbf F_3$.

For either nontrivial central character, let $W$ be its unique
irreducible carrier, of dimension $N=3^{n-1}$. Its Weyl operators
$L_v$ form a basis of $\operatorname{End}(W)$, as proved in D-HOS-01.
Choose an intertwiner $T_W$ for (HOS3C.6). It conjugates $L_v$ to
a unit scalar times $L_{gv}$. Choose a $Q$-invariant positive
Hermitian inner product on $W$. The operator $T_W^*T_W$ centralizes
$Q$ and is a positive scalar by Schur's lemma; dividing by its
positive square root makes $T_W$ unitary. Its $\ell$th power is
a scalar of absolute value one, so a further unit scalar rescaling
gives $T_W^\ell=I$ without losing unitarity.
For $1\le j<\ell$, the trace on this Weyl basis is a sum of
$3^{2n-2-r_0}$ unit scalars. Since this trace also equals
$|\operatorname{tr}(T_W^j)|^2$, we obtain
\[
 |\operatorname{tr}(T_W^j)|\le N3^{-r_0/2},\qquad
 \nu(T_W)\ge(1-1/\ell)(1-3^{-r_0/2})N.
 \tag{HOS3C.7}
\]
The second assertion is exact Fourier inversion for the $\ell$th
roots of unity. On a full isotypic summand the operator is
$T_W\otimes T'$; $T'$ is a finite-order multiplicity operator.
Normalized eigenspace multiplicities convolve, so the same lower
support fraction holds on that entire summand.

All nonzero linear weights of $Q$ have the same multiplicity,
because the full $\operatorname{Sp}_{2n-2}(3)$ Levi is transitive
on them. Their fixed proportion under $g$ is
$(3^{2n-2-r_0}-1)/(3^{2n-2}-1)<3^{-r_0}$. Every other weight
orbit has length $\ell$. Thus their aggregate support fraction is
at least $(1-1/\ell)(1-3^{-r_0})$, stronger than (HOS3C.7).
The remaining summand is $\rho^Q$. D-HOS-07 proved that $Q$ and
its opposite generate $S$, so its dimension is at most $a/2$.
Adding the summands gives the universal estimate
\[
 \nu(\rho(x))\ge\tfrac12(1-1/\ell)(1-3^{-r_0/2})a.
 \tag{HOS3C.8}
\]
No phase-free lift or scalar action on the full multiplicity space
has been assumed.

Apply this with $x=h$ of order three, and write
$r=\operatorname{rank}(h-I)$ on the natural space. The rank-one
case is the transvection case already handled. If $r=2$ or $3$,
then $\dim\ker(h-I)=2n-r>n$, so this fixed space is not totally
isotropic. Choose a nondegenerate plane inside it and take $e$
in that plane. The quotient operator has $r_0=r\ge2$.
If $r\ge4$, choose any nonzero fixed vector $e$. Restricting
$h-I$ to $e^\perp$ loses at most one rank and then quotienting
by $\langle e\rangle$ loses at most one more; hence
$r_0\ge r-2\ge2$. In either case every nontrivial power of
$g$ has the same rank, because $g$ has order three and
$g^2-I=(g-I)(g+I)$ with $g+I$ invertible. Equation (HOS3C.8)
now gives $(1/2)(2/3)(2/3)a=2a/9$, as required.
\end{proof}

\begin{lemma}[D-HOS-14: every inner operator at large natural rank]
\label{D-HOS-14}
Put $N_0=10^{30}$. If $n\ge N_0$, then every noncentral
$x\in\operatorname{Sp}_{2n}(3)$ and every nontrivial irreducible
ordinary representation $\rho$ satisfy
\[
                            \nu(\rho(x))\ge2a/9.
 \tag{HOS3C.9}
\]
\end{lemma}

\begin{proof}
The projective kernel is exactly the center, by simplicity of the
projective quotient and perfectness of the covering group. Take a
power of $x$ of prime order $\ell$ modulo that center. Projective
support does not increase on taking powers or changing a scalar.
If $\ell$ is odd, a central sign makes this power $h$ have actual
order $\ell$. For $\ell=3$, D-HOS-13 applies. For $\ell=2$,
one has either $h^2=I$ or $h^2=-I$.

Suppose first that $h^2=I$. Its two eigenspaces are nonzero,
nondegenerate, and even-dimensional. Change its central sign so
that the smaller eigenspace is the $+1$ eigenspace, and choose a
nondegenerate fixed plane in it. On the remaining symplectic space
the moving rank is the dimension of the larger eigenspace, at least
$n\ge4$. Equation (HOS3C.8) with $\ell=2$ gives support at least
$(1/4)(1-3^{-2})a=2a/9$.

Next suppose $\ell\ge5$ and $h$ has a natural eigenvalue $1$ or
$-1$. After the indicated central sign choice, it has actual odd
order and a nonzero fixed symplectic subspace. Choose a fixed plane.
The quotient operator is nonidentity semisimple of order $\ell$,
so its moving rank is even and at least two; all its nontrivial
powers have the same fixed space. Equation (HOS3C.8) gives
$(1/2)(4/5)(2/3)a=4a/15>2a/9$.

All remaining cases are semisimple and have neither $1$ nor $-1$
as a natural eigenvalue. This includes $h^2=-I$. Over $\mathbf F_3$
there is no other nonzero scalar. Each irreducible factor of its
natural characteristic polynomial therefore has degree at least
two. Galois conjugate eigenvalues have equal multiplicity; hence
the largest eigenspace over the algebraic closure has dimension at
most $n$ in the natural dimension $2n$. In Larsen--Tiep's support
convention,
\[
                      \operatorname{supp}(h)/(2n)\ge1/2.
 \tag{HOS3C.10}
\]

The published uniform theorem, with the explicitly smaller constant
$\sigma_0=10^{-29}$ verified in D-UNW-01 and D-OL-01, gives
\[
 \frac{|\chi(h)|}{a}\le a^{-\sigma_0\operatorname{supp}(h)/(2n)}
 \le a^{-\sigma_0/2}
 \le3^{-(n-1)\sigma_0/2}<\frac13.
 \tag{HOS3C.11}
\]
Here D-HOS-02 gives $a\ge3^{n-1}$ and
$(N_0-1)/(2\cdot10^{29})>1$. The exact primary input is
\cite[Theorem~5.5 and complete proof, accepted-author pp.~25--26]
{UNW:LT2024}, which explicitly includes the symplectic group and
uses the natural dimension in the denominator. Its available constant
$1/(6507\cdot2^{21}\cdot10^{18})>7\cdot10^{-29}$ is larger
than the value used here. The elementary trace inequality
$|\chi(h)|\ge a-2\nu(\rho(h))$ now gives
$\nu(\rho(h))>a/3>2a/9$. This exhausts the projective-prime
cases and proves (HOS3C.9).
\end{proof}

\begin{theorem}[D-HOS-15: the remaining characteristic-three source has bounded natural rank]
\label{D-HOS-15}
Suppose a finite inducing image $J$ has a normal irreducibly acting
central product of quasisimple factors with simple quotients
$\operatorname{PSp}_{2n_i}(3)$, $n_i\ge10^{30}$, and every
factor module is nonselfdual. Then its actual six-dimensional
multiplicity variation has period-image dimension at most eight.
Different such ranks, scalar extensions and arbitrary compatible
factor permutations are allowed.

In the transitive characteristic-three symplectic-core case of
the main imaginary-quadratic sixfold problem, combining this with
D-HOS-12 gives the necessary bound
\[
                                  n<10^{30}.
 \tag{HOS3C.12}
\]
This is a finite natural-rank cutoff, not a classification or
existence assertion for the remaining ranks.
\end{theorem}

\begin{proof}
In this range the universal perfect cover is the symplectic group,
by the complete exceptional-multiplier list
\cite[Chapter~5, Table~1 and preceding paragraph, p.~35]{BI:HHM}.
Automorphisms of represented central quotients lift uniquely to
this perfect cover. The full outer group at $q=3$ has order two,
represented by nonsquare-multiplier symplectic similitudes; there
is no field or graph part
\cite[\S2.2.2, Table~2.1, author p.~14]{HES:LiXia2022}.

That nontrivial diagonal class is class-inverting on characters.
Indeed every symplectic element is a product of two anti-symplectic
involutions, by the theorem stated in
\cite[Theorem~1, author p.~1]{HOS:Vinroot2004}.
An anti-symplectic involution appearing in such a product conjugates
the element to its inverse. All anti-symplectic conjugations belong
to the same outer coset, so their actions on class functions agree.
Thus the diagonal automorphism sends $\chi$ to $\overline\chi$.
A nonselfdual irreducible character is not fixed by it. Consequently
every compatible normalizer of that factor module is a scalar times
an inner operator: after correcting its inner automorphism, this
is Schur's lemma. This argument includes every represented central
quotient and does not assume that all nonreal characters are Weil.

For a factor-preserving nonscalar operator on the total tensor
space, at least one factor is nonscalar. D-HOS-14 gives its support
fraction at least $2/9$. Tensoring the other operators convolves
normalized eigenspace multiplicities and cannot increase their
maximum. If an operator nontrivially permutes factors, the weighted
tensor-cycle argument already proved in D-HOS-06 gives support
fraction at least $1/4$, stronger than $2/9$. Nonidentity scalar
inertias move the entire total space. Therefore every genuine
inertia has moving rank at least $2e/9$, where $e$ is the total
inducing degree.

The elementary degree bound gives $e\ge3^{10^{30}-1}>9$.
The actual intermediate Hurwitz factorization has
$6=2e(h-1)+D$. Genus at least two is impossible, and in genus one
$D=6$ gives at most six genuine marked points and period rank at
most six. Rank nine would therefore force genus zero, at least
twelve genuine inertia points, and $D=2e+6$. But then
\[
                    2e+6\ge12(2e/9),\qquad e\le9,
 \tag{HOS3C.13}
\]
a contradiction. The argument uses the actual multiplicity variation
and allows its coefficient field to exceed the original field.

Finally, in a transitive core all factor modules either are all
selfdual or are all nonselfdual, since automorphic pullback and
module isomorphism preserve selfduality. The first case was excluded
in the main imaginary-character-field setting by D-HOS-12, with
its original-field-to-inducing-field argument and nontrivial
bilinear multiplier. The second case with $n\ge10^{30}$ has just
been excluded. This proves (HOS3C.12), with the stated residual scope.
\end{proof}

\paragraph{Proof and source boundary.}
The new Siegel-weight and phase-tolerant Heisenberg calculations
are supplied in full. Their input from D-HOS-01,02,07 is the
previously whole-reviewed elementary carrier, degree and generation
argument. The published inputs are exactly the uniform character
theorem, ordinary automorphism and cover descriptions, and
anti-symplectic factorization. No missing maximum-eigenspace
classification, unknown character field, or unproved normalizer
compatibility is assumed. The small-rank nonselfdual case remains
open. The even-rank Hermitian-form counting lead is separate and
is not used in this cutoff.


\section{A published large-support estimate sharply bounds the residual symplectic rank}
\label{sec:symplectic-three-sharp}

\paragraph{Status and relation to the preceding proof.}
This separate addendum leaves D-HOS-13--15 unchanged. Its only new
published input is the explicit estimate and support refinement of
Larsen--Shalev--Tiep. All constants and the passage through the perfect
cover are recorded below. The entire argument and invoked published proof have passed independent
whole-proof reviews by /root and /root/imprimitive; the separate exact
arithmetic also passes. It does not assert existence at any surviving rank.
As before, $\nu(T)$ is the codimension of the largest complex eigenspace
of $T$, and $n$ denotes half the natural dimension of
$S=\operatorname{Sp}_{2n}(3)$.

\begin{theorem}[D-HOS-16: a finite residual rank interval]
\label{D-HOS-16}
Put
\[
 c_*=(19/20)^9=\frac{322687697779}{512000000000},\qquad
 \delta_*:=\frac{1-c_*}{2}
          =\frac{189312302221}{1024000000000}.
 \tag{HOS3S.1}
\]
For every $n\ge284$, every nontrivial irreducible ordinary complex
$S$-module $V$ of degree $a$, and every noncentral $x\in S$,
\[
                           \nu(x|V)\ge\delta_*a.
 \tag{HOS3S.2}
\]
Consequently a finite inducing image with an irreducibly acting normal
central product of quasisimple factors whose simple quotients are
$\operatorname{PSp}_{2n_i}(3)$, $n_i\ge284$, and whose factor modules
are all nonselfdual cannot have period-image dimension nine for its
actual six-dimensional multiplicity variation. Arbitrary central
quotients, scalar extensions, different such ranks, and compatible
factor permutations are allowed.

In the transitive characteristic-three symplectic-core case of the
main imaginary-quadratic sixfold problem, D-HOS-12 and this theorem
therefore reduce the remaining case with $n\ge4$ to
\[
                                  4\le n\le283.
 \tag{HOS3S.3}
\]
\end{theorem}

\begin{proof}
We first specify the primary estimate, including the distinction
between its two support conventions. Larsen--Shalev--Tiep,
\cite[Definition~4.1.1 and (4.1.1), p.~1907; Theorem~4.3.6,
pp.~1917--1918, (4.3.2)]{HOS:LST2011}, treat a simple simply
connected classical algebraic group $\mathbf G$, its finite group
$\mathbf G(\mathbf F_q)$, and every nontrivial ordinary irreducible
character $\chi$ of that finite group. If the support of $h$ is at
least $N\ge567$, their explicit estimate is
\[
 \frac{|\chi(h)|}{\chi(1)}
       <\gamma_q^{\sqrt{N/7}},\qquad \gamma_3=19/20.
 \tag{HOS3S.4}
\]
Here the usual support is the codimension of the largest eigenspace
over the algebraic closure in the natural representation. The
radical covers $N/7$: the complete proof, part~4 on p.~1922,
explicitly chooses $\alpha=\sqrt{1/7}$ in its exponent
$\alpha\sqrt N$. Thus $N=567=7\cdot9^2$ gives the constant
$c_*$ in (HOS3S.1).

The same authors then state, in Remark~4.3.7, p.~1923, that their
conclusion also holds with support restricted to the allowed scalar
set. For the symplectic group this set is $\mu_2=\{1,-1\}$.
We may therefore use (HOS3S.4) with
\[
 s_{\pm}(h):=\min_{\lambda\in\{1,-1\}}
             \operatorname{codim}\ker(h-\lambda I).
 \tag{HOS3S.5}
\]
This is an explicit published refinement, not an identification
of $s_{\pm}$ with the usual algebraic-closure support.

Apply these statements to the algebraic group
$\mathbf G=\operatorname{Sp}_{2n}$, whose rational points are
the full perfect cover $S$, including its central involution.
The theorem does not require $\chi$ to be faithful. Thus it applies
unchanged to a character inflated from any represented central
quotient; no character-ratio assertion about an unspecified quotient
is being substituted for the published statement.

Take a power of $x$ of prime order $\ell$ modulo $Z(S)$, and
change its central sign as necessary, calling the result $h$.
The projective kernel of a nontrivial irreducible $S$-module is
exactly $Z(S)$, as in D-HOS-14. Taking powers and multiplying by
a scalar cannot increase projective support. If $\ell$ is odd,
choose the sign so that $h$ has actual order $\ell$; if $\ell=2$,
then $h^2=I$ or $-I$.

All but one family of cases have the uniform estimates already
proved in D-HOS-13 and the elementary part of D-HOS-14. For
actual order three the bound is $2a/9$. If $h^2=I$, choose its
central sign so that the smaller nonzero symplectic eigenspace is
fixed, and remove a fixed hyperbolic plane. The quotient moving
rank is at least $n\ge4$, so the phase-tolerant Heisenberg estimate
(HOS3C.8) gives at least $(1/4)(1-3^{-2})a=2a/9$.
For odd $\ell\ge5$ with a natural eigenvalue in $\{1,-1\}$,
the actual odd-order representative has a fixed symplectic plane;
the quotient moving rank is at least two. The same estimate gives
at least $(1/2)(4/5)(2/3)a=4a/15$.

The residual cases are semisimple and have neither $1$ nor $-1$
as a natural eigenvalue, including $h^2=-I$. The usual
algebraic-closure support is at least $n$, because each minimal
polynomial factor over $\mathbf F_3$ has degree at least two.
Thus the unrefined theorem alone already gives the argument for
$n\ge567$. Using the explicit published refinement (HOS3S.5),
however, we have exactly $s_{\pm}(h)=2n$. For $n\ge284$ this
is at least $568>567$, so (HOS3S.4) with $N=567$ gives
$|\chi(h)|/a<c_*$. Rotating an eigenvalue of maximum multiplicity
to $1$ and applying the triangle inequality gives
$|\chi(h)|\ge a-2\nu(h|V)$. Hence $\nu(h|V)>\delta_*a$.
The exact inequalities
\[
                 \frac16<\delta_*<\frac29<\frac14
 \tag{HOS3S.6}
\]
show that the preceding elementary cases also imply (HOS3S.2).
This proves the assertion for every noncentral $x$.

For the source consequence we repeat the normalizer passage so
its scope is explicit. In the stated rank range the universal
perfect cover of each simple factor is $\operatorname{Sp}_{2n_i}(3)$,
and its only nontrivial outer class is the nonsquare-multiplier
diagonal class, by the cover and automorphism inputs already
checked in D-HOS-15. Vinroot's anti-symplectic factorization
implies that this class sends each character to its complex
conjugate. It therefore does not stabilize a nonselfdual factor
module. Every compatible factor normalizer is consequently a
scalar times an inner operator. This remains true after passing
to a represented central quotient, since its automorphisms lift
to the perfect cover and Schur's lemma identifies the remaining
intertwiner scalar.

On the full tensor space of dimension $e$, a factor-preserving
nonscalar operator has a nonscalar factor, so (HOS3S.2) and
convolution of eigenvalue multiplicities give support at least
$\delta_*e$. An operator permuting factors nontrivially has
support at least $e/4$ by the weighted tensor-cycle lemma, which
is stronger. A nonidentity scalar moves the entire space.
Thus every genuine inertia moves at least $\delta_*e$ dimensions.

The actual intermediate-cover formula is
$6=2e(h-1)+D$. The elementary degree bound D-HOS-02 gives
$e\ge3^{283}$, so genus at least two is impossible. In genus
one, $D=6$ permits at most six genuine marks and hence period
rank at most six. Rank nine would force genus zero, at least
twelve genuine inertia points, and $D=2e+6$. It would follow that
\[
 e\le\frac{6}{12\delta_*-2}
   =\frac{1536000000000}{55936906663}<28,
 \tag{HOS3S.7}
\]
contradicting $e\ge3^{283}$. This is a bound on the actual
multiplicity variation and does not identify the inducing
character field with the original imaginary quadratic field.

Finally selfduality is invariant under automorphic pullback and
module isomorphism. In a transitive core all factors therefore
have the same selfduality status. The selfdual case has already
been excluded by D-HOS-12 in the main imaginary-character-field
setting. The nonselfdual case with $n\ge284$ is excluded above.
This proves (HOS3S.3); it makes no assertion about the existence
or exclusion of the remaining nonselfdual ranks $4$ through $283$.
\end{proof}

\paragraph{Certificate boundary.}
The accompanying rational certificate verifies the displayed
constants, support thresholds, and the source degree contradiction.
It does not enumerate groups or certify the published character
theorem. The latter is used at the exact stated primary locators,
and the full proof of Theorem~4.3.6 and its subsequent remark were
read in the published version. The new addendum is independent
of the very small explicit exponent used in D-HOS-14.

\section{The remaining characteristic-three symplectic involutions}
\label{sec:symplectic-three-involutions}

\paragraph{Status.}
This separate complete derivation has passed whole-proof reviews by
/root and /root/imprimitive, with an independent exact supplement. It leaves D-HOS-13--16 unchanged. The new result
does not exclude all of the remaining natural ranks $4\le n\le283$.
It removes their projective involutions as a possible low-support
inertia class and specifies the resulting necessary class gate.
As before $\nu(T)$ is the codimension of the largest complex eigenspace.

\begin{lemma}[D-HOS-17: symmetric forms fixed by a complex structure]
\label{D-HOS-17}
Let $N=2m\ge4$, let $A\in\operatorname{GL}_{N}(3)$ satisfy
$A^2=-I$, and let $\mathcal O$ be any nonzero congruence orbit of
the general linear group on symmetric bilinear forms on
$\mathbf F_3^N$. Then
\[
 \frac{|\{B\in\mathcal O:A^tBA=B\}|}{|\mathcal O|}\le\frac1{39}.
 \tag{HOS3I.1}
\]
For rank two, the fraction is either zero or exactly
$2/[3(3^{N-1}-1)]$. For rank at least four it is at most $1/351$.
\end{lemma}

\begin{proof}
Put $\mathbf F_9=\mathbf F_3[t]$, $t^2=-1$, and let $t$ act as
$A$. If a symmetric form $B$ satisfies the fixed equation, then
$B(Ax,y)=-B(x,Ay)$, and
\[
                   H(x,y)=B(x,y)+tB(x,Ay)                    \tag{HOS3I.2}
\]
is Hermitian, linear in its first variable. Conversely the coefficient
of $1$ in a Hermitian form gives such a $B$. Their radicals agree;
an $\mathbf F_9$-rank-$r$ Hermitian form therefore gives an
$\mathbf F_3$-rank-$2r$ symmetric form. On the nondegenerate
quotient every Hermitian form is equivalent to $I_r$. Its underlying
symmetric form is $I_{2r}$: for $x=a+bt$ and $y=c+dt$, the
coefficient of $1$ in $x\bar y$ is $ac+bd$. Consequently only
the orthogonal type $\epsilon=(-1)^r$ occurs in rank $2r$.

We use the ordinary finite-form classification and group orders,
with their proofs in Taylor, \cite[pp.~115--118 and 138--141]{HOS:Taylor1992}.
These are orders of the full unitary and orthogonal groups, not their
special or projective quotients. The numbers of nondegenerate
Hermitian forms on $\mathbf F_9^r$ and symmetric forms of type
$\epsilon=(-1)^r$ on $\mathbf F_3^{2r}$ are respectively
\[
 \begin{split}
 H_r&=\frac{|\operatorname{GL}_r(9)|}{|\operatorname U_r(3)|}
      =3^{r(r-1)/2}\prod_{j=1}^{r}(3^j+(-1)^j),\\
 S_r&=\frac{|\operatorname{GL}_{2r}(3)|}
                  {|\operatorname O_{2r}^{(-1)^r}(3)|}
      =\frac12\,3^{r^2}(3^r+(-1)^r)
                       \prod_{j=1}^{r}(3^{2j-1}-1).
 \end{split}                                                   \tag{HOS3I.3}
\]
The formulas follow by orbit--stabilizer after choosing an ordered
basis; the orders in question are obtained in the cited source by
counting orthogonal or hyperbolic bases. Thus the relevant quotient is
\[
 R_r:=\frac{H_r}{S_r}
   =2\,3^{-r(r+1)/2}
       \frac{\prod_{j=1}^{r-1}(3^j+(-1)^j)}
            {\prod_{j=1}^{r}(3^{2j-1}-1)}.                     \tag{HOS3I.4}
\]
In particular $R_1=1/3$, $R_2=1/351$, and
\[
 \frac{R_{r+1}}{R_r}
     =3^{-r-1}\frac{3^r+(-1)^r}{3^{2r+1}-1}<1.               \tag{HOS3I.5}
\]

For a fixed rank, choose the radical first, then a nondegenerate form
on the quotient. An $A$-stable radical is an $\mathbf F_9$-subspace.
Hence the fraction in (HOS3I.1), when its type can occur, is exactly
\[
       \frac{{m\brack r}_{9}}{{2m\brack2r}_{3}}\,R_r,         \tag{HOS3I.6}
\]
where the brackets denote counts of subspaces. The first quotient is
at most one, since these $\mathbf F_9$-subspaces form a subset of
all the indicated $\mathbf F_3$-subspaces. For $r\ge2$ the bound
$R_r\le1/351$ proves the assertion. For $r=1$ the elementary
subspace counts give
\[
 \frac{{m\brack1}_{9}}{{2m\brack2}_{3}}\,R_1
     =\frac{2}{3(3^{N-1}-1)}\le\frac1{39},                   \tag{HOS3I.7}
\]
as $N\ge4$. Odd ranks and the other even-rank type have no fixed
forms. This proves the complete orbit statement.
\end{proof}

\begin{theorem}[D-HOS-18: the square-minus-one support bound]
\label{D-HOS-18}
Let $S=\operatorname{Sp}_{2n}(3)$, $n\ge4$, and let $V$ be any
nontrivial ordinary irreducible complex $S$-module of degree $a$.
If $h\in S$ satisfies $h^2=-I$, then
\[
 \nu(h|V)\ge\frac{19a}{78}\quad(n\text{ even}),\qquad
 \nu(h|V)>\frac{19a}{104}\quad(n\text{ odd}).                 \tag{HOS3I.8}
\]
The statement includes characters inflated from a represented
central quotient and makes no selfduality assumption.
\end{theorem}

\begin{proof}
Let $W$ be the natural space and $\langle\ ,\ \rangle$ its
alternating form. Regard $W$ as an $n$-dimensional space over
$\mathbf F_9$ by letting $t$ act as $h$. Then
\[
 \mathcal H(x,y)=\langle x,hy\rangle-t\langle x,y\rangle       \tag{HOS3I.9}
\]
is a nondegenerate Hermitian form, linear in the first variable.
For example $\mathcal H(hx,y)=t\mathcal H(x,y)$, by
$h^2=-I$ and symplecticity; the other identities follow directly.
The finite Hermitian Witt decomposition has $\lfloor n/2\rfloor$
hyperbolic planes and an anisotropic remainder of dimension zero
or one over $\mathbf F_9$.

Choose their two opposite maximal isotropic subspaces $L,L'$.
Their $\mathbf F_3$-dimension is
\[
              N=n\quad(n\text{ even}),\qquad
              N=n-1\quad(n\text{ odd}).                     \tag{HOS3I.10}
\]
They are $h$-stable, and $W_0=L\oplus L'$ is a nondegenerate
symplectic space of dimension $2N$ over $\mathbf F_3$.
Its orthogonal complement is zero or a symplectic plane. Let
$S_0=\operatorname{Sp}(W_0)$ act identically on that complement.
Inside $S_0$ take the opposite Siegel radicals $U,U^-$ associated
to $L,L'$. They are conjugate in $S_0$ and generate $S_0$.

Here generation can also be seen directly. The two radicals contain
the transvections with centers in $L$ and in $L'$. If
$z=x+y$, $x\in L$, $y\in L'$, and $\langle x,y\rangle\ne0$,
conjugation by a suitable transvection with center $x$ gives one
with center $z$. If this pairing vanishes and $y\ne0$, choose
$x'\in L$ with $\langle x',y\rangle\ne0$. First obtain the
center $x+x'+y$, then remove $x'$ by another such conjugation.
The cases $x=0$ or $y=0$ are already included. Thus all symplectic
transvections occur; elementary symplectic elimination shows that
they generate $S_0$, as also used in D-HOS-13.

Put $b=\dim V^{S_0}$ and $v_0=\dim V^U=\dim V^{U^-}$.
The intersection of these two invariant spaces is $V^{S_0}$,
so the subspace dimension formula gives
\[
                         v_0\le\frac{a+b}{2}.                \tag{HOS3I.11}
\]
When $n$ is even, $S_0=S$ and $b=0$. When $n$ is odd, $n\ge5$
and $2N=2n-2\ge8$. This is exactly a restriction subgroup in
Larsen--Shalev--Tiep, \cite[Proposition~4.2.3, pp.~1909--1912]{HOS:LST2011}:
its natural space is a nondegenerate summand and the group fixes
the orthogonal complement pointwise. Applied to our nontrivial
irreducible representation of the full simply connected group, its
explicit estimate gives
\[
 \frac ba<3^{8/3-(2N)/2}=3^{11/3-n}
             \le3^{-4/3}<\frac14.                            \tag{HOS3I.12}
\]
The last inequality is equivalent to $81>64$. The complete cited
proof, in particular its symplectic step 4 and summation step 6,
was read; no equality of ordinary and projective group orders is
being used. Thus $v_0\le a/2$ for even $n$ and $v_0<5a/8$ for
odd $n$.

Identify $U$ with the additive symmetric matrices on $L$. Under
the nondegenerate matrix-trace pairing over $\mathbf F_3$, its
complex character group is the same symmetric-matrix space.
The Levi subgroup $\operatorname{GL}(L)\subset S_0$ acts by
congruence. Therefore all weight spaces in one such orbit have
equal multiplicity in the restriction of $V$. The element $h$
normalizes $U$; its restriction $A$ to $L$ satisfies $A^2=-I$,
and its fixed weights are precisely the forms in D-HOS-17.
Replacing $A$ by its inverse in the contragredient convention
does not change the fixed equation.

Consequently at most $1/39$ of the total nonzero weight-space
dimension is fixed by $h$. On the other weights $h$ acts in pairs,
since its square acts trivially by conjugation on $U$. For each
pair of equal-dimensional weight spaces, an eigenvector of $h$
is determined by its component in one space. Every eigenvalue
therefore occupies at most half of the pair's total dimension.
This remains valid whether the central involution acts as $1$
or $-1$ on $V$. Summing over pairs proves
\[
 \nu(h|V)\ge\frac12\frac{38}{39}(a-v_0).                    \tag{HOS3I.13}
\]
Combining (HOS3I.11)--(HOS3I.12) gives exactly (HOS3I.8).
The argument uses the actual irreducible representation and its
weight spaces, not only a character-ratio bound.
\end{proof}

\begin{corollary}[D-HOS-19: a necessary residual inertia class]
\label{D-HOS-19}
Consider the main sixfold source problem with a finite inducing
image having an irreducibly acting, transitive normal central
product of factors whose simple quotients are
$\operatorname{PSp}_{2n}(3)$, $n\ge4$. Suppose its factor modules
are nonselfdual. If its actual six-dimensional multiplicity
variation has period rank nine, then $4\le n\le283$, and at
least one genuine inertia has all of the following properties:
it preserves every tensor factor; it is not a scalar on the full
space; and in every nonscalar factor component, every noncentral
power of prime order modulo the symplectic center has odd prime
order $\ell\ge5$ and, in its odd-order lift, has no natural
$1$ or $-1$ eigenvalue. In particular the projective orders of
those noncentral factor components are coprime to six.

The final coprimality assertion concerns projective factor orders.
It does not impose that condition on an independent scalar part
of the actual inertia. No existence claim is made for these classes.
\end{corollary}

\begin{proof}
The rank interval is D-HOS-16. For nonselfdual modules the full
compatible normalizer of one factor is scalar times inner:
the only nonsquare diagonal outer class is class-inverting by
Vinroot and hence cannot stabilize the module. This is the exact
normalizer result already checked in D-HOS-15, including lifts
through represented central quotients.

Taking a power and multiplying by a central sign cannot increase
projective support. The elementary cases in D-HOS-13--14 give
support at least $2a/9$ for actual order three or $h^2=I$, and
at least $4a/15$ for odd prime order with a natural fixed plane.
D-HOS-18 now supplies the remaining projective-involution case
$h^2=-I$. Thus a nonscalar factor component with at least one
prime-order power outside the residual class in the statement
has support at least
\[
 \beta_n a,\qquad
 \beta_n=\frac29\ (n\text{ even}),\qquad
 \beta_n=\frac{19}{104}\ (n\text{ odd}).                      \tag{HOS3I.14}
\]
Convolution of tensor-factor spectra gives the same proportional
bound on the full degree $e$ whenever one component has this
property. A nontrivial tensor-factor permutation has support
at least $e/4$, and a nonidentity scalar moves all $e$ dimensions.

If no inertia had all the asserted residual properties, every
genuine inertia would therefore move at least $\beta_n e$.
The elementary degree bound D-HOS-02 gives $e\ge3^{n-1}$.
For the actual intermediate-cover variation,
$6=2e(h-1)+D$ rules out $h\ge2$, and $h=1$ permits at most
six genuine marks and period rank at most six. Rank nine thus
requires $h=0$, at least twelve genuine marks, and $D=2e+6$.
It follows that
\[
 e\le\frac6{12\beta_n-2}
   =9\ (n\text{ even}),\qquad
 e\le\frac{156}{5}\ (n\text{ odd}).                          \tag{HOS3I.15}
\]
These contradict respectively $e\ge27$ for even $n\ge4$ and
$e\ge81$ for odd $n\ge5$. Hence the required inertia exists.
If a noncentral factor's projective order were divisible by two
or three, taking a suitable power would give one of the excluded
prime-order cases. This proves the projective coprimality claim.
The argument uses the actual inducing multiplicity variation;
its coefficient field need not equal the original imaginary
quadratic field.
\end{proof}

\paragraph{Exact supplement and limits.}
The accompanying exact finite reconstruction
enumerates every symmetric form in dimensions two and four,
classifies its rank and type by exact finite-field elimination,
and compares the fixed counts with (HOS3I.6). It also checks the
Hermitian construction, radical matrices and exact rational
constants. The orbit-counting and restriction arguments above
are the proofs in arbitrary dimension. A finite sample is not
used to extrapolate the result, and the published branching
estimate is not claimed as a new theorem of this run.

\section{Closing the remaining higher symplectic source in characteristic three}
\label{sec:symplectic-three-semisimple}

\paragraph{Status and notation.}
This separate complete new derivation has passed whole-proof reviews
by /root and /root/imprimitive, with separate exact supplements. It uses D-HOS-17--18, leaving their frozen proof
unchanged. The finite checks described below supplement, rather than
replace, the universal arguments. Write $S_n=\operatorname{Sp}_{2n}(3)$
and $\nu(T)$ for the codimension of its largest complex eigenspace.

\begin{lemma}[D-HOS-20: odd-prime fixed forms]
\label{D-HOS-20}
Let $A\in\operatorname{GL}_{N}(3)$ have prime order $\ell\ge5$
and have no eigenvalue $1$. For every nonzero general-linear
congruence orbit $\mathcal O$ of symmetric bilinear forms,
\[
             |\{B\in\mathcal O:A^tBA=B\}|\le|\mathcal O|/9.
                                                               \tag{HOS3P.1}
\]
\end{lemma}

\begin{proof}
The matrix $A$ is semisimple, and it has neither $1$ nor $-1$ as
an eigenvalue. A fixed form has an $A$-stable radical. On its
nondegenerate quotient of dimension $r>0$, the induced matrix
$A_0$ still has order $\ell$ and no such eigenvalue. Over the
algebraic closure, an invariant symmetric form pairs only
eigenspaces with eigenvalues $\lambda$ and $\lambda^{-1}$.
Nondegeneracy requires equal dimensions within each inverse pair.
Thus $r$ is even. If $r=2$, the matrix would belong to
$O_2^+(3)$ or $O_2^-(3)$, whose orders are four and eight,
respectively. This is impossible for $\ell\ge5$. These orders
follow directly from the diagonal and norm-one descriptions of
the two binary orthogonal groups, or from
\cite[Theorem~11.4 and its proof, pp.~139--140]{HOS:Taylor1992}.
Hence any nonzero fixed form has quotient dimension $r\ge4$.

Write $m_j$ for the common dimensions in the inverse eigenvalue
pairs of $A_0$. The vector space of all invariant symmetric forms
on the quotient has dimension
\[
                \sum_jm_j^2\le\left(\sum_jm_j\right)^2
                         =r^2/4.                            \tag{HOS3P.2}
\]
The dimension is the same over $\mathbf F_3$ and over its
algebraic closure because it is the kernel dimension of a linear
system over $\mathbf F_3$. It therefore contains at most
$3^{r^2/4}$ nondegenerate forms of either prescribed type.

For comparison, fix a nonsingular diagonal matrix $D$ of that
type. For every upper unitriangular matrix $L$, the matrix
$L^tDL$ has the same type. These matrices are distinct: if two
are equal, their quotient is an upper unitriangular isometry of
$D$; its first column is the first standard vector, so
orthogonality forces all other entries of its first row to vanish,
and induction gives the identity. Thus the type has at least
$3^{r(r-1)/2}$ nondegenerate forms. For each fixed radical,
the fixed proportion of quotient forms is at most
\[
             3^{r^2/4-r(r-1)/2}
                =3^{-r(r-2)/4}\le1/9.                       \tag{HOS3P.3}
\]
Every form in the original orbit is specified by its radical
and a nondegenerate quotient form of the specified type. The
number of quotient forms is independent of the chosen radical.
Non-$A$-stable radicals contribute no fixed forms; for each
stable radical the preceding estimate applies. Summing proves
(HOS3P.1). No enumeration of general-linear orbits is assumed.
\end{proof}

\begin{theorem}[D-HOS-21: all odd-prime semisimple elements]
\label{D-HOS-21}
Let $n\ge4$, let $V$ be a nontrivial ordinary irreducible complex
$S_n$-module of degree $a$, and let $h\in S_n$ have prime order
$\ell\ge5$. Then
\[
                            \nu(h|V)\ge4a/15.               \tag{HOS3P.4}
\]
This includes all characters inflated from represented central
quotients and requires no selfduality assumption.
\end{theorem}

\begin{proof}
If $h$ has a natural eigenvalue $1$, its fixed space is a
nondegenerate symplectic space, hence contains a hyperbolic pair.
The phase-tolerant Heisenberg argument already proved in
D-HOS-14 gives (HOS3P.4). Since $h$ has odd order, $-1$ cannot
be a natural eigenvalue. We therefore assume that it has neither.

We first treat any case in which there are opposite $h$-stable
totally isotropic subspaces $L,L'$ of common dimension $N\ge4$
whose sum $W_0$ is nondegenerate. The subgroup
$S_0=\operatorname{Sp}(W_0)$ acts identically on its perpendicular
space. Its two opposite Siegel radicals $U,U^-$ generate $S_0$
and are conjugate, by the explicit argument in D-HOS-18.
If $S_0=S_n$, its invariant subspace in $V$ is zero. Otherwise
Larsen--Shalev--Tiep's precise restriction estimate
\cite[Proposition~4.2.3 and complete proof, pp.~1909--1912]{HOS:LST2011}
gives
\[
 \frac{\dim V^{S_0}}a<3^{8/3-N}\le3^{-4/3}<1/4.
                                                               \tag{HOS3P.5}
\]
The subspace intersection inequality for $V^U,V^{U^-}$ therefore
gives $\dim V^U<5a/8$ (with the stronger $a/2$ bound if $S_0=S_n$).

The dual of $U$ is the symmetric-form space on $L$, with the
congruence action of its full Levi subgroup $\operatorname{GL}(L)$.
Consequently weight multiplicities are constant on each orbit.
The matrix $A=h|_L$ has order $\ell$ and no eigenvalue $1$.
D-HOS-20 bounds the fixed fraction in every nonzero orbit by
$1/9$. All other weight orbits of $h$ have length $\ell$.
An eigenvector on such an orbit is determined by its component
in one weight space, so each eigenvalue occupies at most a
$1/\ell$ fraction there. Thus
\[
 \nu(h|V)\ge(1-1/\ell)\frac89\,(a-\dim V^U)
             >\frac45\frac89\frac38a=\frac4{15}a.           \tag{HOS3P.6}
\]

It remains to show that the absence of suitable $L,L'$ forces
a small centralizer. We give the primary-space argument fully.
Put $f=\operatorname{ord}_{\ell}(3)$. Every nontrivial irreducible
factor $P$ of $X^\ell-1$ over $\mathbf F_3$ has degree $f$.
The primary space $W_P$ pairs under the symplectic form only
with $W_{P^*}$, where $P^*$ is the reciprocal polynomial. This
follows either by diagonalizing over the algebraic closure, or
by observing that the adjoint of $h$ is $h^{-1}$.

If $f$ is odd, no such $P$ is self-reciprocal: the condition
$\lambda^{-1}=\lambda^{3^j}$ for a primitive $\ell$-th root
would put the order-two element $-1$ in the odd-order subgroup
generated by $3$ modulo $\ell$. Choose one primary space from
each pair to form $L$ and the other to form $L'$. They are
opposite Lagrangians, of dimension $n\ge4$, so (HOS3P.6) applies.

If $f$ is even, write $f=2d$. The unique element of order two
in the group generated by $3$ modulo $\ell$ is
$3^d=-1$, so every primary polynomial is self-reciprocal.
Moreover $d\ge2$: degree two would require the prime $\ell\ge5$
to divide $3^2-1=8$. Put $E=\mathbf F_{3^{2d}}$ and
$E_0=\mathbf F_{3^d}$. On a primary space of multiplicity $m_P$,
$h$ acts as multiplication by a root $\lambda$ generating $E$,
with adjoint involution $\sigma(\lambda)=\lambda^{-1}$.

For clarity, the associated Hermitian structure follows directly
from trace duality. There is a unique $E$-valued pairing $K$ such
that
\[
 \operatorname{Tr}_{E/\mathbf F_3}(bK(x,y))
                         =\langle bx,y\rangle\quad(b\in E).
                                                               \tag{HOS3P.7}
\]
The adjoint rule for $E$ gives
$K(bx,y)=bK(x,y)$, $K(x,by)=\sigma(b)K(x,y)$ and
$K(y,x)=-\sigma K(x,y)$. It is nondegenerate. Choose a nonzero
$c\in E$ with $\sigma(c)=-c$; then $cK$ is Hermitian.
Conversely the original alternating form is recovered by the
trace in (HOS3P.7), so an $E$-isotropic Hermitian subspace is
also symplectically isotropic over $\mathbf F_3$.

The finite Hermitian Witt decomposition has index
$\lfloor m_P/2\rfloor$, as proved in
\cite[Lemma~10.2, Corollary~10.3 and the following decomposition,
pp.~116--117]{HOS:Taylor1992}.
If any $m_P\ge2$, one Hermitian hyperbolic plane supplies
opposite $\mathbf F_3$-isotropic spaces of dimension
$2d\ge4$, whose sum is symplectically nondegenerate. Again
(HOS3P.6) applies.

The only case left is therefore $m_P=1$ for every primary
polynomial. Distinct primary spaces cannot be permuted by an
operator commuting with $h$. On each space that operator is
$E$-linear, and hence multiplication by an element $b\in E^*$.
Preserving the alternating form, or equivalently (HOS3P.7),
is exactly $b\sigma(b)=1$. Thus the centralizer in the full
symplectic group has order
\[
               |C_{S_n}(h)|=(3^d+1)^k,\qquad n=dk.          \tag{HOS3P.8}
\]
The scalar center is already included in this product; no
projective-group order is substituted.

For all integers $d\ge2$, $3^d+1\le10^{d/2}$: equality holds
at $d=2$, and the ratio on increasing $d$ is less than
$3/\sqrt{10}<1$. Hence $|C_{S_n}(h)|\le10^{n/2}$.
Character column orthogonality gives
$|\chi(h)|\le|C_{S_n}(h)|^{1/2}$, the elementary statement
\cite[Lemma~4.3.1, p.~1912]{HOS:LST2011}.
Using the proven degree bound D-HOS-02 gives
\[
 \frac{|\chi(h)|}{a}
       \le\frac{10^{n/4}}{3^{n-1}}\le\frac{10}{27},
 \qquad
       \nu(h|V)\ge\frac{1-10/27}{2}a
                         =\frac{17}{54}a>\frac4{15}a.       \tag{HOS3P.9}
\]
The middle expression decreases with $n$, since $10^{1/4}<3$,
and its value at $n=4$ is $10/27$. The support inequality uses
$|\chi(h)|\ge a-2\nu(h|V)$ and remains valid without a
faithfulness assumption on the character. The cases are
exhaustive, proving (HOS3P.4).
\end{proof}

\begin{theorem}[D-HOS-22: the higher characteristic-three source is excluded]
\label{D-HOS-22}
For every $n\ge4$, every nontrivial ordinary irreducible complex
$S_n$-module of degree $a$, and every noncentral $g\in S_n$,
\[
 \nu(g|V)\ge\frac29a\quad(n\text{ even}),\qquad
 \nu(g|V)\ge\frac{19}{104}a\quad(n\text{ odd}).               \tag{HOS3P.10}
\]

A finite inducing image with an irreducibly acting normal central
product of such factors, all with nonselfdual factor modules,
cannot supply period rank nine for its actual six-dimensional
multiplicity variation. Different natural ranks, represented
central quotients, scalar extensions and all compatible factor
permutations are allowed. Consequently the transitive ordinary
symplectic-core case in the main imaginary-quadratic sixfold
problem has no source with $q=3$ and $n\ge4$, in either
determinant class.
\end{theorem}

\begin{proof}
The projective kernel of a nontrivial irreducible representation
of $S_n$ is its scalar center. Taking a noncentral power of $g$
of prime order modulo that center, and changing its central
sign if necessary, cannot increase projective support. For odd
prime order the representative has actual prime order; for
projective order two its square is $I$ or $-I$.

Actual order three has support at least $2a/9$ by D-HOS-13.
The case $h^2=I$ has that same bound by the already proved
Heisenberg argument of D-HOS-14. D-HOS-18 treats $h^2=-I$,
with bounds $19a/78>2a/9$ when $n$ is even and
$19a/104$ when $n$ is odd. D-HOS-21 treats every remaining
odd prime, with bound $4a/15>2a/9$. These exhaust all cases,
so (HOS3P.10) follows.

For a nonselfdual factor module the only compatible factor
normalizers are scalar times inner, by the exact class-inverting
diagonal-automorphism argument in D-HOS-15. Factorwise tensor
operators retain the support fraction of any nonscalar factor
by convolution. A nontrivial factor permutation has support at
least one quarter of the full dimension. A nonidentity scalar
has actual moving rank equal to that dimension.

If every factor has even $n_i$, every genuine inertia thus moves
at least $2e/9$ dimensions. If some factor has odd $n_i$, the
uniform bound $19e/104$ applies. The factor degree bound gives
$e\ge27$ in the first case and $e\ge81$ in the second.
As before the exact intermediate-cover Euler formula
$6=2e(h-1)+D$ excludes genus at least two; genus one has at most
six genuine marks and period-image dimension at most six.
Rank nine therefore requires genus zero, at least twelve
genuine inertia points, and $D=2e+6$. It would force
\[
                 e\le9\quad\hbox{or}\quad e\le156/5,
                                                               \tag{HOS3P.11}
\]
respectively, contradicting the preceding lower bounds.
No invariant bilinear form is required for this nonselfdual
normal-core conclusion, and the inducing character field is
not identified with the original imaginary quadratic field.

In a transitive core selfduality is preserved by automorphic
pullback, so either every factor module is selfdual or none is.
The former case is already excluded in the main imaginary-field
setting by D-HOS-12. The latter is excluded above. This proves
the final assertion; the small ranks $n<4$ remain within their
separately stated earlier results.
\end{proof}

\paragraph{Primary inputs and certificate scope.}
The new semisimple primary-space and centralizer reductions, the
unitriangular count and the source argument are proved here.
The published restriction estimate is used at its exact ordinary
simply connected-group scope, as checked in D-HOS-18. The small
binary orthogonal orders and finite Hermitian decomposition are
the stated elementary primary inputs from Taylor. The accompanying
exact finite supplement tests the fixed-form argument, both kinds
of primary behavior, a toral centralizer and all constants. It
does not infer the universal theorem from those finite examples.

\section{An explicit rank cutoff for arbitrary unitary characters}
\label{sec:unitary-nonweil-cutoff}

\paragraph{Status and scope.}
This is a new author derivation from two published character estimates.
The complete proof, invoked primary sources and exact constants have
passed author, root and independent peer reviews. The peer certificate
independently reconstructs every quantitative reduction.
The new estimate concerns all ordinary non-Weil irreducible characters
of special unitary groups of sufficiently large natural dimension.
It applies to their full compatible linear normalizers, without a
self-duality hypothesis or an equality of inducing and original
character fields.  The numerical cutoff is deliberately conservative.
The theorem does not classify the characters below the cutoff.

Put
\[
 \sigma_0=10^{-29},\qquad J_0=10^{30},\qquad N_0=10^{61}.
                                                               \tag{UNW.1}
\]
Here $n$ always denotes the natural dimension of
$\operatorname{SU}_n(q)$; the degree of a complex character is denoted
by $a$.  For a finite-order complex operator $T$ on a space of
dimension $a$, write
\[
 \nu(T)=a-\max_{\lambda}\dim\ker(T-\lambda).
                                                               \tag{UNW.2}
\]
Thus $\nu$ is unchanged by multiplying $T$ by a scalar.  Ordinary
moving rank means $\operatorname{rank}(T-1)$ and is a different
quantity.

\begin{lemma}[D-UNW-01: the two quantitative published inputs]
\label{D-UNW-01}
Let $S=\operatorname{SU}_n(q)$, $n\ge3$, and let
$\chi\in\operatorname{Irr}(S)$ have degree $a$ and character level
$j$.  Then
\[
 \frac{|\chi(g)|}{a}\le
 a^{-\sigma_0\operatorname{supp}(g)/n}\qquad(g\in S),
                                                               \tag{UNW.3}
\]
where $\operatorname{supp}(g)$ is $n$ minus the largest eigenspace
dimension of $g$ on the natural module over an algebraic closure.
The level satisfies
\[
 a\ge \frac{q^{j(n-j)}}{2(q+1)},\qquad
 j\ge n/2\ \Longrightarrow\
 a>\frac{q^{n^2/4-2}}{q+1}.                                  \tag{UNW.4}
\]

If $2\le j<n/2$ and $n\ge j^2+j+1$, set
$Q=\operatorname{GU}_j(q)$ and
$R_j=\sqrt{16.52q^{j^2+j-1}}$.
There is an $\alpha\in\operatorname{Irr}(Q)$ and an extension of
$\chi$ to $\operatorname{GU}_n(q)$, of the same degree, such that
for every noncentral $g\in S$, with
$k=n-\operatorname{supp}(g)\le n-1$,
\[
\begin{split}
 a&\ge\frac{\alpha(1)}{|Q|}
 \left(q^{nj}-|Q|q^{n(j-1)}-|Q|R_jq^{n(j-2)}\right),\\
 |\chi(g)|&\le\frac{\alpha(1)}{|Q|}
 \left(q^{kj}+|Q|q^{(n-1)(j-1)+1}
                  +|Q|R_jq^{n(j-2)}\right),\\
 |Q|&\le1.5q^{j^2}.
\end{split}                                                    \tag{UNW.5}
\]
These estimates include $q=2$.
\end{lemma}

\begin{proof}
We specify the known inputs and their constants rather than reprove
the published character theorems.
Larsen--Tiep prove (UNW.3) with an explicit larger constant
\[
 \sigma=\frac1{6507\cdot2^{21}\cdot10^{18}}>7\cdot10^{-29}.
                                                               \tag{UNW.6}
\]
This is Theorem~5.5 and the last displayed computation of its
complete proof in \cite[accepted-author version, pp.~25--26]{UNW:LT2024}.
The constants supplied in their Theorem~4.4 are
$C=2^{14}10^{12}$ and $\gamma=(2^{13}10^{12})^{-1}$, and the proof
takes
$\sigma=\min\{[241(9C)^{3/2}]^{-1},\gamma/82\}$.
Consequently the smaller rational constant $\sigma_0$ in (UNW.1)
is valid.  The proof read here is the accepted author version;
the journal publication and DOI are verified separately.  No claim
is made to have retrieved its typeset journal PDF.

We use the level defined for $S$ in
\cite[Definition~2 and Lemma~7.1, pp.~51--52]{UNW:GLT2020}:
it is the level of any irreducible character of
$\operatorname{GU}_n(q)$ lying above $\chi$.
Lemma~7.1 proves independence of the choice.  Equations (UNW.4)
are precisely Theorem~7.3(i)--(ii) there, with
$\sigma_-=1/[2(q+1)]$.

For $j<n/2$, Corollary~8.6(i), including its proof on pp.~67--68,
shows that a character above $\chi$ restricts irreducibly to $S$.
Its determinant twists contain a unique character of true level
$j$.  The determinant twists are trivial on $S$.
Theorem~1.1(iii), proved in full in \S8.3 on pp.~65--67, identifies
this true-level character with $D_\alpha^\circ$ for some
$\alpha\in\operatorname{Irr}(\operatorname{GU}_j(q))$ whenever
$j^2+j\le n-1$.  This condition is sufficient in even as well as
odd characteristic.  The exact estimate (8.18), p.~66, gives
(UNW.5); the bound on $|Q|$ is recorded immediately below it.
The tail estimate used in (8.18) is (8.17), and the complete
subsequent argument proves that these characters are irreducible,
pairwise distinct and exhaust true level $j$.
We use this exact estimate, rather than its weaker numerical
corollary with coefficient $2.43$.
\end{proof}

\begin{theorem}[D-UNW-02: a uniform non-Weil character ratio]
\label{D-UNW-02}
Let $n\ge N_0$ and $q\ge2$, and let
$\chi\in\operatorname{Irr}(\operatorname{SU}_n(q))$ have level
$j\ge2$ and degree $a$.  Then every noncentral $g$ satisfies
\[
                         |\chi(g)|\le\frac4{15}a.            \tag{UNW.7}
\]
Moreover $a>q^n>30$.  Level zero is the trivial representation;
level-one characters are the ordinary unitary Weil characters
and have degree at most $q^n$.
\end{theorem}

\begin{proof}
First record a convenient degree consequence of (UNW.4).  Since
$2(q+1)\le q^3$ and $q+1\le q^2$, for $n\ge12$ we have
\[
 a\ge q^{nj/4}\quad(1\le j\le n/2),\qquad
 a\ge q^{n^2/8}\quad(j\ge n/2).                             \tag{UNW.8}
\]
For the first assertion, $j(n-j)-3\ge nj/2-3\ge nj/4$
once $nj\ge12$, as is automatic here.  For the second,
$n^2/4-4\ge n^2/8$ for $n\ge8$.

If $j>J_0$, put $j'=\min\{j,n/2\}$.  Because $n\ge N_0>2J_0$,
we have $j'\ge J_0$.  Every noncentral natural operator has support
at least one.  Equations (UNW.3) and (UNW.8) therefore give
\[
 \frac{|\chi(g)|}{a}
 \le q^{-\sigma_0 j'/4}
 \le 2^{-\sigma_0J_0/4}
 =2^{-5/2}<\frac14<\frac4{15}.                             \tag{UNW.9}
\]
The high-level argument does not assume that $\chi$ extends to
the general unitary group.

It remains to consider $2\le j\le J_0$.  Our constants satisfy
$N_0\ge J_0^2+J_0+8$, so (UNW.5) applies and $n-j^2\ge8$.
Use the deliberately weaker estimates
\[
                    |Q|\le2q^{j^2},\qquad R_j\le5q^{j^2}.
                                                               \tag{UNW.10}
\]
The second follows from $16.52<25$ and
$j^2+j-1\le2j^2$.  Define
\[
 u=q^{j^2-n}\le2^{-8},\qquad d=2u+10u^2.
\]
After dividing the lower bound in (UNW.5) by
$\alpha(1)q^{nj}/|Q|$, its right side is at least $1-d$.
The normalized numerator in the upper bound is at most
\[
 q^{-j}+2q^{j^2-n-j+2}+10q^{2j^2-2n}
 \le q^{-j}+d\le\frac14+d.
\]
Here $j\ge2$ is essential.  Exact rational arithmetic gives
\[
 d\le\frac2{256}+\frac{10}{256^2}
      =\frac{261}{32768}<\frac1{100},\qquad
 \frac{1/4+1/100}{1-1/100}=\frac{26}{99}<\frac4{15}.
                                                               \tag{UNW.11}
\]
The denominator is positive, so (UNW.7) follows for every low
level as well, including $q=2,j=2$.

For completeness the degree separation is uniform.  If
$2\le j\le n/2$, monotonicity of $j(n-j)$ on this interval and
(UNW.4) give $a\ge q^{2n-7}>q^n$ for $n\ge8$.
If $j\ge n/2$, they give $a>q^{n^2/4-4}>q^n$ for $n\ge8$.
Theorem~7.3(i) gives the upper bound $q^n$ at level one.
Definition~2 and Lemma~7.1 identify level zero with the trivial
character.  The identification at level one is also explicitly
recorded at the end of \S8.3, p.~67, of \cite{UNW:GLT2020}.
\end{proof}

\begin{theorem}[D-UNW-03: full normalizers and the finite-cover cutoff]
\label{D-UNW-03}
Let a finite group $J\subset\operatorname{GL}(V)$ normalize an
irreducibly acting central product of quasisimple groups
$E_1\cdots E_t$.  Suppose that, on the associated tensor factors,
each $E_i$ acts through a nontrivial ordinary level-at-least-two
representation of $\operatorname{SU}_{n_i}(q_i)$ with
$n_i\ge N_0$.  Factor permutations are allowed, including all
permutations compatible with the represented central quotients.
Then every nonscalar $T\in J$ satisfies
\[
                         \nu(T)\ge\frac{11}{60}\dim V.
                                                               \tag{UNW.12}
\]
Such a representation cannot be an inducing block of an actual
finite-cover Weil sixfold source with period rank nine.

In particular, in the primitive factorwise tensor-indecomposable
normal-core situation of D-BI-19, if the simple factors are
$\operatorname{PSU}_n(q)$, a dominant source must satisfy
\[
                              n<10^{61}.                    \tag{UNW.13}
\]
This conclusion combines the new non-Weil estimate with the
already proved unitary Weil theorem D-BI-24.  In the nonsplit
main-source scope of D-BI-20 the remaining general unitary gate
therefore has both $n<10^{61}$ and $q\le27$.
\end{theorem}

\begin{proof}
We first prove a full-normalizer estimate for one factor of degree
$a$.  Average a positive Hermitian form over its finite image.
If $T$ is a nonscalar normalizer, Schur's lemma gives an element
$x$ of the image such that $b=TxT^{-1}x^{-1}\ne1$.
Choose an eigenvalue $\lambda$ of $T$ with largest eigenspace.
The rank identity for a commutator gives
\[
 \operatorname{rank}(b-1)
 =\operatorname{rank}(Tx-xT)
 \le2\operatorname{rank}(T-\lambda)=2\nu(T).                \tag{UNW.14}
\]
The operator $b$ lies in the finite core image.  If it is a
nonidentity scalar, its ordinary moving rank is $a$.
Otherwise it is represented by a noncentral element of the
special unitary group and (UNW.7) applies.  Its eigenvalues have
absolute value one, so
\[
 2\operatorname{rank}(b-1)
 \ge a-\operatorname{Re}\operatorname{Tr}(b)
 \ge(1-4/15)a=11a/15.
\]
Together with (UNW.14) this proves $\nu(T)\ge11a/60$.
No assumption on which outer automorphism $T$ implements has
entered the proof.  Represented central quotients also cause no
problem: lift the character to $\operatorname{SU}_n(q)$, and
apply the estimate to any preimage of the commutator.

If an element fixes all tensor factors, it is a tensor product
of the corresponding intertwiners.  An eigenspace proportion of
a tensor product is at most the largest eigenspace proportion
on any chosen factor: this follows by summing the products of
the eigenvalue multiplicities.  If the whole operator is
nonscalar, some factor is nonscalar, giving (UNW.12).

For clarity the factor-permutation estimate does not require
any character-field assumption.  On a cycle of length
$\ell\ge2$ with factor degree $a\ge2$, changes of tensor bases
reduce the weighted cyclic permutation to a cyclic shift with
one return operator.  This return operator is diagonalizable:
a power is scalar because the original operator has finite
projective order.  In its eigenbasis, tensor words are permuted
in cyclic orbits, each of which contributes at most one copy
of a prescribed eigenvalue.  There are $a$ constant words and
every other orbit has length at least two.  The number of orbits
is at most
\[
 a+(a^\ell-a)/2\le3a^\ell/4.
\]
Hence such a cycle has $\nu\ge a^\ell/4$; convolution with the
other factors preserves the proportion.  Since $1/4>11/60$,
this proves (UNW.12) in every case.

We spell out the geometric constraint used to convert this
estimate into an obstruction.  Let $e=\dim V$ be an inducing
degree, let $X$ be the intermediate base curve of genus $h$,
and write $r$ for its nontrivial inertia points.  Polarized
Shapiro and the finite-cover Hurwitz factorization, already
proved for this source problem, identify the two chosen Hodge
multiplicities with $(3,3)$.  For the irreducible nontrivial
coefficient representation the Euler identity is
\[
 6=2e(h-1)+\sum_{i=1}^r\operatorname{rank}(g_i-1).
                                                               \tag{UNW.15}
\]
The additional finite-cover data over a pointed $X$ are discrete,
so the period map factors through a finite cover of its moduli.
For $h\ge2$ the right side is already greater than six, since
$e>30$.  For $h=1$, the sum of moving ranks is six, hence
$r\le6$ and the pointed-curve moduli dimension is at most six
(or one if $r=0$).  Such a source cannot have period rank nine.
Therefore $h=0$, $r\ge12$, and
\[
                  \sum_i\operatorname{rank}(g_i-1)=2e+6.
                                                               \tag{UNW.16}
\]
Every nonidentity scalar inertia has moving rank $e$, and every
nonscalar inertia has moving rank at least $\nu(g_i)\ge11e/60$.
Thus
\[
               2e+6\ge12\cdot\frac{11e}{60}
                  =\frac{11e}{5},\qquad e\le30.
                                                               \tag{UNW.17}
\]
Each factor already has degree greater than thirty by
D-UNW-02, a contradiction.

Finally consider the transitive core of D-BI-19.  At these large
natural dimensions there is no exceptional multiplier; the
universal perfect cover of $\operatorname{PSU}_n(q)$ is
$\operatorname{SU}_n(q)$, and all its quasisimple covers are
central quotients.  The precise finite-exception boundary is
the published list in \cite[Chapter~5, Table~1, p.~35]{BI:HHM},
with the ordinary-cover distinction on p.~4.  Thus every factor
character lifts to the special unitary group.  Transitive
conjugation makes all factor degrees equal.  By the degree
separation at the end of D-UNW-02, either they all have level
one or they all have level at least two.  The latter case was
just excluded.  The former is excluded by D-BI-24; none of
its small unitary exceptions has natural dimension at least
$N_0$.  This proves (UNW.13).  The additional field cutoff is
exactly D-BI-20, with its nonsplit and primitive-core hypotheses
retained.
\end{proof}

\paragraph{What the finite reduction does and does not decide.}
The exact integer $10^{61}$ is an explicit upper cutoff for the
remaining natural unitary rank in the stated main-source problem.
Together with the prior field bound it leaves a finite collection
of finite groups, each with finitely many ordinary characters.
It does not supply their tables, classify all lower-rank tuples,
or prove existence of any remaining source.  Nor does (UNW.12)
exclude an arbitrary mixed core containing additional uncovered
small-rank factors: its mixed statement requires every factor
to satisfy the displayed large-rank non-Weil hypotheses.

\paragraph{Certificate.}
The standalone standard-library program
\texttt{unitary\_nonweil/check\_constants.py} checks all numerical
constants using exact integers and rational numbers, together
with the polynomial identities underlying the two degree regimes.
It does not enumerate ranks up to $10^{61}$; universal monotonicity
and the two-regime proof are given above.  The retained source
retrieval manifest records the exact full-text hashes and the
unsuccessful attempts to obtain a typeset LT journal PDF.

\section{The same explicit cutoff for arbitrary linear characters}
\label{sec:linear-nonweil-cutoff}

\paragraph{Status.}
New derivation, completely read and checked by the author, root and
an independent peer, with separately executed exact constant certificates. The published inputs are the same two articles as in
D-UNW-01, with the full general-linear dual-pair proof specified below.
No new character table or representation classification is assumed.

\begin{theorem}[D-LNW-01: an explicit linear normal-core rank cutoff]
\label{D-LNW-01}
Set $N_0=10^{61}$. If $n\ge N_0$, $q\ge2$ and
$\chi\in\operatorname{Irr}(\operatorname{SL}_n(q))$ has
character level at least two and degree $a$, then
\[
            |\chi(g)|\le\frac4{15}a
           \quad(g\notin Z(\operatorname{SL}_n(q))),
           \qquad a>q^n>30.                                \tag{LNW.1}
\]
Every finite-order nonscalar normalizer of its represented image has
projective support at least $11a/60$.

An irreducibly acting normal central-product core all of whose factors
have such large-natural-dimension non-Weil linear representations
cannot supply period rank nine in the sixfold finite-cover problem.
Arbitrary mixtures among these factors, scalar and outer intertwiners,
central quotients and compatible factor permutations are allowed.

Consequently, in the transitive primitive factorwise
tensor-indecomposable normal-core situation of D-BI-19, a linear
simple quotient $\operatorname{PSL}_n(q)$ must have
\[
                            n<10^{61}.                      \tag{LNW.2}
\]
The level-one alternative is excluded by the already proved linear
Weil theorem D-BI-23. In the nonsplit main-source scope of D-BI-20
the remaining general linear gate also has $q\le27$.
\end{theorem}

\begin{proof}
Write $j$ for the character level. The definition and its independence
are Definition~2 and Lemma~7.1 of \cite[pp.~51--52]{UNW:GLT2020}.
The uniform support estimate of \cite[Theorem~5.5, accepted-author
pp.~25--26]{UNW:LT2024} applies to $\operatorname{SL}_n(q)$ as
well as the unitary group, with the same certified constant
$\sigma_0=10^{-29}$. Theorem~7.3(i)--(ii) of
\cite[p.~52]{UNW:GLT2020} gives
\[
 a\ge\frac{q^{j(n-j)}}{q-1},\qquad
 j\ge n/2\ \Longrightarrow\
 a>\frac{q^{n^2/4-2}}{q-1}.                                \tag{LNW.3}
\]
These imply the weaker common bounds (UNW.4) and hence (UNW.8).
Put $J_0=10^{30}$. If $j>J_0$, the complete argument in
(UNW.9) applies verbatim and gives $|\chi(g)|/a<1/4$.
This does not require a same-degree general-linear extension in
the high-level range.

For $2\le j\le J_0$, we have $j<n/2$. Corollary~8.6, including
its full proof on pp.~67--68 of \cite{UNW:GLT2020}, gives an
irreducible extension to $\operatorname{GL}_n(q)$ of the same
degree; determinant twisting selects true level $j$ and is trivial
on the special linear group. The general-linear dual-pair argument
in the complete \S8.1, pp.~55--61, identifies this extension with
$D_\alpha^\circ$ for an irreducible character $\alpha$ of
$Q=\operatorname{GL}_j(q)$. The condition in Theorem~1.1(ii)
is automatic for $j\le n/2$.

Here is the precise stronger estimate used, rather than the
numerical $1.76$ corollary. Let
$R_j=\sqrt{8q^{j^2+j}}$, and let $k\le n-1$ be the largest
natural eigenspace dimension of a noncentral $g$. Equation (8.11),
p.~60, gives
\[
\begin{split}
 a&\ge\frac{\alpha(1)}{|Q|}
  \left(q^{nj}-|Q|q^{n(j-1)}-|Q|R_jq^{n(j-1)}\right),\\
 |\chi(g)|&\le\frac{\alpha(1)}{|Q|}
  \left(q^{kj}+|Q|q^{(n-1)(j-1)+1}
                         +|Q|R_jq^{n(j-1)}\right).
\end{split}                                                    \tag{LNW.4}
\]
The complete derivation of the tail bound is (8.9)--(8.10),
pp.~59--60. The earlier geometric dual-pair argument proves the
irreducibility and exhaustion assertion; no nonnegative virtual
character is silently substituted for an irreducible one.

The elementary order formula gives $|Q|\le q^{j^2}$, and
$R_j\le3q^{j^2}$ because $8<9$ and $j^2+j\le2j^2$ for $j\ge1$.
Since $N_0\ge2J_0^2+9$, put
\[
                   u=q^{2j^2-n}\le2^{-9},\qquad d=4u.
\]
After division by $\alpha(1)q^{nj}/|Q|$, the lower bound in
(LNW.4) is at least $1-d$, since
\[
 |Q|(1+R_j)q^{-n}
 \le q^{j^2-n}+3q^{2j^2-n}\le4u.
\]
The normalized upper bound is at most $q^{-j}+d$: the middle
term has exponent at most $j^2-n-j+2\le j^2-n$.
Consequently
\[
 \frac{|\chi(g)|}{a}
 \le\frac{1/4+1/128}{1-1/128}
 =\frac{33}{127}<\frac{26}{99}<\frac4{15}.                 \tag{LNW.5}
\]
This includes $q=2,j=2$. The degree separation $a>q^n$ for
levels at least two follows from the same inequalities at the end
of D-UNW-02, using (LNW.3). Level-one characters have degree at
most $q^n$ by Theorem~7.3(i); they are the linear Weil characters,
as explicitly recorded on p.~61 of \cite{UNW:GLT2020}.

For a nonscalar normalizer $T$, choose a core element $x$ with
$[T,x]\ne1$ by irreducibility. A nonidentity scalar commutator
has ordinary moving rank $a$. Otherwise (LNW.1) and the unitary
eigenvalue trace inequality give ordinary moving rank at least
$(1-4/15)a/2$. In either case
$\operatorname{rank}([T,x]-1)\le2\nu(T)$, so
$\nu(T)\ge11a/60$. This treats every compatible outer
automorphism directly, without imposing a field or graph
stabilizer hypothesis.

For tensor products, normalized eigenspace multiplicities are
convolutions, and a nontrivial factor-permutation cycle has at
least one-quarter projective support by the complete weighted
word-orbit argument in D-UNW-03. The common bound is therefore
$11e/60$ on total degree $e$. The same actual intermediate-base
Euler identity (UNW.15) forces genus zero and at least twelve
genuine inertia points for period rank nine, and then
$2e+6\ge12(11e/60)$ gives $e\le30$. Each factor already has
degree greater than thirty, a contradiction.

At natural dimension $n\ge N_0$, all quasisimple groups with
simple quotient $\operatorname{PSL}_n(q)$ are central quotients
of $\operatorname{SL}_n(q)$; the complete exceptional-cover list
in \cite[Chapter~5, Table~1, p.~35]{BI:HHM} has no such parameter.
Thus all factor characters lift to this ordinary group. In the
transitive-core case their degrees are equal. Degree separation
forces all of them to have level one or all to have level at least
two. The latter case has just been excluded; the former is exactly
the range of D-BI-23, with no large-rank exception. This proves
(LNW.2). The field cutoff invokes D-BI-20 only with its original
nonsplit and primitive-core hypotheses.
\end{proof}

\paragraph{Verification boundary.}
The entire cited general-linear \S8.1, including Lemma~8.2 and
the exact estimate (8.11), was read from the complete published PDF.
The standalone \texttt{check\_linear\_constants.py} executes the
new exact low-level estimates and checks the reduction to the already
certified high-level constants. Neither this corollary nor the
certificate enumerates the finite remaining range, claims that a
remaining profile is sufficient, or admits additional uncovered
small factors in its mixed-core statement.

\section{Explicit level and spin cutoffs for orthogonal normal cores}
\label{sec:orthogonal-level-cutoff}

\paragraph{Status and scope.}
This is a new author derivation from published quantitative character
estimates. The complete proof, invoked sources and author arithmetic certificate
have been read by /root and /root/g34\_eightfolds. Both whole-proof
reviews pass. Separately authored exact arithmetic checks also pass;
their actual scopes and timings are retained in the check register.
This section gives a strict reduction of the orthogonal normal-core
gate. It does not give an upper bound on the natural dimension of
the remaining ordinary low-level orthogonal representations.

Write $d$ for the natural orthogonal dimension, $a$ for the degree
of a complex irreducible representation, and put
\[
 \sigma_0=10^{-29},\qquad J_0=10^{30},\qquad
 M_0=2\cdot10^{29}=J_0/5,\qquad D_0=2\cdot10^{30}.
                                                               \tag{OL.1}
\]
For odd $d$ we assume $q$ odd. For even $d$ either characteristic
is allowed, with the appropriate sign $\epsilon\in\{+,-\}$.
For a finite-order operator $T$ on a complex space of dimension $a$,
let
\[
                 \nu(T)=a-\max_\lambda\dim\ker(T-\lambda).
                                                               \tag{OL.2}
\]
This number is unchanged by scalar multiplication of $T$.

\begin{lemma}[D-OL-01: the quantitative orthogonal inputs]
\label{D-OL-01}
Let $d\ge10$. Let $S=\Omega_d^\epsilon(q)$, and for odd $q$
also allow $\widetilde S=\operatorname{Spin}_d^\epsilon(q)$.
The sign is suppressed when $d$ is odd. Every ordinary
irreducible character $\chi$ of either indicated group satisfies
\[
 \frac{|\chi(g)|}{\chi(1)}
 \le \chi(1)^{-\sigma_0\operatorname{supp}(g)/d}.
                                                               \tag{OL.3}
\]
For a spin element, support is computed on its image in the
natural orthogonal group.

If $\chi\in\operatorname{Irr}(S)$ has level $\ell$, then
\[
 0\le\ell\le\lfloor d/2\rfloor,\qquad
 \chi(1)\le\left(\frac{q^d-1}{q-1}\right)^\ell<q^{d\ell}
 \quad(\ell>0).
                                                               \tag{OL.4}
\]
For $d\ge50$ and $k=\lceil\ell/3\rceil$,
\[
              \chi(1)\ge q^{k(d-2k-2)}(q-1)^k.
                                                               \tag{OL.5}
\]

For odd $q$, the following two failures both imply
\[
                 a>q^{E(d)},\qquad
                 E(d)=d^2/8-d/4-2:
                                                               \tag{OL.6}
\]
an irreducible spin character of degree $a$ does not descend to
$S$; or an irreducible character of $S$ of degree $a$ does not
extend to $\operatorname{SO}_d^\epsilon(q)$.
\end{lemma}

\begin{proof}
Equation (OL.3) is Larsen--Tiep's published uniform character
theorem, with the smaller explicit constant $\sigma_0$ already
checked in D-UNW-01. Its exact source is
\cite[Theorem~5.5 and complete proof, accepted-author pp.~25--26]
{UNW:LT2024}. That theorem explicitly includes $\Omega(V)$ and,
in odd characteristic, $\operatorname{Spin}(V)$. The support
convention for the latter is also stated immediately before
Proposition~6.4, author p.~30. Its larger available constant is
$1/(6507\cdot2^{21}\cdot10^{18})>7\cdot10^{-29}$.

We use exactly the level of
\cite[Definition~3.2(ii), p.~160]{OL:GLT2024}: the least exponent
of $\tau_A+\zeta_A$ containing the character. Here
$\tau_A(g)=q^{\dim\ker(g-1)}$ and
$\zeta_A(g)=(-1)^d(-q)^{\dim\ker(g-1)}$ on the natural
quadratic space $A$. Lemma~3.3 proves that these two characters
agree on $\Omega(A)$ in all the characteristics at issue.
Lemma~3.4(iii), including its proof on pp.~160--161, gives the
range in (OL.4). Theorem~1.5(iii), with its proof at the end of
Section~3 on p.~167, gives the degree upper bound.
Theorem~3.7 and its complete proof, pp.~165--167, give (OL.5):
for $\ell>0$ we have $\ell>3(k-1)$, so its contrapositive applies;
the case $\ell=0$ is immediate. The small
exceptional pairs in that theorem are absent for $d\ge50$.

For the spin assertion, write $d=2m+1$ or $2m$.
Proposition~6.3 and its complete proof, pp.~193--194 of
\cite{OL:GLT2024}, give respectively
\[
 a>q^{m(m+1)/2}/4,\qquad a>q^{m(m-1)/2}/4.
                                                               \tag{OL.7}
\]
In the odd-dimensional case the kernel of the spin map is the
entire center of order two. A character nontrivial on that
kernel is faithful because the spin group is quasisimple, so
the faithfulness hypothesis of Proposition~6.3(i) is satisfied.
In even dimension Proposition~6.3(ii) directly has the required
noninflation hypothesis; faithfulness on the entire center is
not imposed. Since $4\le q^2$, both estimates imply (OL.6).

For the extension assertion, $S$ has index two in
$\operatorname{SO}_d^\epsilon(q)$. An invariant irreducible
character of $S$ extends across this cyclic quotient: an
intertwiner for a coset representative can be rescaled so that
its square is the prescribed image of the square of that
representative. Thus if the character does not extend, its
induction is irreducible and has degree $2a$.
Proposition~7.3 and its complete proof, pp.~198--200 of
\cite{OL:GLT2024}, apply to this induced character and give
\[
 2a>q^{m^2/2}\quad(d=2m+1),\qquad
 2a>(q-1)q^{m(m-1)/2-1}\quad(d=2m).
                                                               \tag{OL.8}
\]
Using $2\le q$ and $q-1\ge1$, these imply (OL.6).
Indeed $m^2/2-1=E(2m+1)+9/8$, while the weaker exponent
$m(m-1)/2-2$ in even dimension is exactly $E(2m)$.
All propositions used here are ordinary-character statements;
no cross-characteristic modular variant is needed.
\end{proof}

\begin{lemma}[D-OL-02: an explicit degree barrier]
\label{D-OL-02}
Suppose a finite group $J\subset\operatorname{GL}(V)$ normalizes
an irreducibly acting central product of quasisimple groups
$E_1\cdots E_t$. On each associated tensor factor, suppose the
representation lifts to an ordinary irreducible character of
$\operatorname{Spin}_{d_i}^{\epsilon_i}(q_i)$ for odd $q_i$, or
$\Omega_{d_i}^{\epsilon_i}(q_i)$ for even $q_i$, with $d_i\ge10$.
Let its degree be $a_i$, and assume for every factor that
\[
                              a_i\ge q_i^{M_0d_i}.
                                                               \tag{OL.9}
\]
Then every nonscalar element $T\in J$ satisfies
\[
                           \nu(T)\ge\frac3{16}\dim V.
                                                               \tag{OL.10}
\]
Such a representation cannot be an inducing block of an actual
finite-cover Weil sixfold source with period rank nine.

In particular, in the transitive normal-core situation of
D-BI-19 with orthogonal simple factors of natural dimension
$d\ge10$ and field size $q$, every factor degree necessarily
satisfies
\[
                                  a<q^{M_0d}.
                                                               \tag{OL.11}
\]
All compatible outer automorphisms, represented central
quotients, scalar extensions and tensor-factor permutations
are included.
\end{lemma}

\begin{proof}
For one factor, a noncentral element of the covering group has
natural support at least one. In the spin case, an element
whose natural image is scalar is central: its commutators lie
in the central kernel of the natural map, and the resulting
homomorphism from the perfect group to that kernel is trivial.
Equations (OL.3) and (OL.9) therefore imply
\[
                         |\chi(g)|\le a q^{-2}\le a/4
                         \quad(g\notin Z).
                                                               \tag{OL.12}
\]

Let $T$ be a nonscalar compatible linear normalizer of this
irreducible factor. By Schur's lemma there is a core operator
$x$ with $b=TxT^{-1}x^{-1}\ne1$. Choose an eigenvalue
$\lambda$ of $T$ with largest eigenspace. Then
\[
 \operatorname{rank}(b-1)
   =\operatorname{rank}(Tx-xT)
   \le2\operatorname{rank}(T-\lambda)=2\nu(T).
                                                               \tag{OL.13}
\]
If $b$ is a nonidentity scalar its moving rank is $a$.
Otherwise lift it to a noncentral element of the covering
group and use (OL.12). Since $b$ belongs to a finite group,
its eigenvalues have absolute value one, and
\[
 2\operatorname{rank}(b-1)
 \ge a-\operatorname{Re}\operatorname{Tr}(b)\ge3a/4.
\]
Together these give $\nu(T)\ge3a/16$. This argument does not
require a classification of the automorphism implemented by
$T$, or that its coefficient field equal the quadratic Weil
field.

An operator fixing all tensor factors is a tensor product of
the individual intertwiners. The maximal eigenspace proportion
of a tensor product is at most that of any selected factor,
by convolution of eigenvalue multiplicities. If the total
operator is nonscalar, some individual intertwiner is
nonscalar and gives (OL.10).

For an actual permutation cycle of length $s\ge2$, the factor
degrees along that cycle coincide, say they equal $a\ge2$.
Changes of tensor bases reduce the operator to a weighted
cyclic shift with one return operator. The return operator is
diagonalizable because a power is scalar, as follows from the
finite projective order of the total operator. In its
eigenbasis each orbit of tensor words contributes at most one
copy of a prescribed eigenvalue. There are $a$ constant words,
and every nonconstant orbit has size at least two. Thus every
eigenspace has dimension at most
\[
                       a+(a^s-a)/2\le3a^s/4.
\]
Consequently such a cycle gives $\nu\ge a^s/4$, and convolution
with the other factors preserves this proportion. Since
$1/4>3/16$, equation (OL.10) follows in all cases.

We recall the actual geometric constraint. For an inducing
degree $e=\dim V$, let $X$ be the intermediate curve of genus
$h$ and let $r$ be the number of nontrivial inertia points.
Polarized Shapiro and the finite-cover factorization already
proved in this source problem give
\[
 6=2e(h-1)+\sum_{j=1}^r\operatorname{rank}(g_j-1).
                                                               \tag{OL.14}
\]
The additional finite-cover data over pointed $X$ are discrete,
so the period map factors through a finite cover of its moduli.
Here $e>24$ by (OL.9). Thus $h\ge2$ is impossible from
(OL.14). If $h=1$, then $r\le6$, and pointed elliptic-curve
moduli has dimension at most six (or one when $r=0$), so cannot
support period rank nine. Hence rank nine requires $h=0$,
$r\ge12$, and total moving rank $2e+6$.
Every nonidentity scalar inertia has moving rank $e$; every
nonscalar inertia has moving rank at least $3e/16$. Therefore
\[
                   2e+6\ge12\cdot3e/16=9e/4,
                   \qquad e\le24,
                                                               \tag{OL.15}
\]
a contradiction. This is an obstruction from the actual
factorization and period rank, not a claimed sufficiency of a
dimension inequality.

Finally, at $d\ge10$ there is no exceptional Schur multiplier
for these orthogonal simple groups; the universal covers are
the indicated spin groups in odd characteristic and the
ordinary orthogonal groups in even characteristic. This is
the precise exceptional-cover list in
\cite[Chapter~5, Table~1, p.~35; ordinary-cover discussion,
p.~4]{BI:HHM}. All quasisimple factors are central quotients
of those covers. Transitive conjugation of the simple factors
makes their degrees equal, and their abstract groups have the
same orthogonal parameters at these dimensions. If one common
degree violated (OL.11), every factor would satisfy (OL.9),
and the preceding obstruction would apply. This proves
(OL.11). In particular, no assertion that character level is
invariant under an outer automorphism is needed for this step.
\end{proof}

\begin{theorem}[D-OL-03: the remaining orthogonal gate]
\label{D-OL-03}
Consider the transitive orthogonal normal-core situation of
D-OL-02 for an actual period-rank-nine Weil sixfold source.
Every factor character which descends to
$\Omega_d^\epsilon(q)$ has
\[
                                  \ell<10^{30}.
                                                               \tag{OL.16}
\]
If $d\ge2\cdot10^{30}$, every factor character descends to
$\Omega_d^\epsilon(q)$. When $q$ is odd, every such descended
character also extends to $\operatorname{SO}_d^\epsilon(q)$.
Equivalently, any genuinely spin factor, or any ordinary factor
failing this special-orthogonal extension, has
\[
                                  d<2\cdot10^{30}.
                                                               \tag{OL.17}
\]
These conclusions require neither a self-duality assumption
nor a prescribed determinant class.
\end{theorem}

\begin{proof}
If $\ell\ge J_0$, then (OL.4) forces $d\ge2J_0$, in
particular $d\ge50$. Put $k=\lceil\ell/3\rceil$.
The inequalities
\[
 \ell/3\le k\le\ell/3+2/3\le d/6+2/3
\]
give
\[
 d-2k-2\ge2d/3-10/3
             =3d/5+(d-50)/15\ge3d/5.
\]
Consequently (OL.5) implies
\[
                      a\ge q^{3dk/5}\ge q^{d\ell/5}
                        \ge q^{dJ_0/5}=q^{M_0d},
                                                               \tag{OL.18}
\]
contradicting (OL.11). This proves (OL.16).

For the remaining assertions write $d=2J_0+s$ with $s\ge0$.
An exact expansion gives
\[
 \begin{split}
 E(d)-M_0d
  &=J_0^2/10-J_0/2-2
      +(3J_0/10-1/4)s+s^2/8>0.
 \end{split}
                                                               \tag{OL.19}
\]
The constant term is positive already for every integer
$J_0\ge10$, and all remaining coefficients are positive.
Either failure in (OL.6) would therefore imply
$a>q^{M_0d}$, again contradicting (OL.11).
For even $q$ the high-dimensional universal covering group
already is the ordinary orthogonal group, so only the descent
assertion is relevant; no odd-characteristic extension
proposition is applied there.
\end{proof}

\paragraph{Precisely what remains.}
In the nonsplit primitive normal-core setting, the previous
field theorem D-BI-20 also gives $q\le27$; that additional
hypothesis is retained when using this field cutoff. The new
result leaves an ordinary orthogonal gate of bounded character
level $\ell<10^{30}$ but unbounded natural dimension $d$.
For $d\ge2\cdot10^{30}$ it consists of characters descending
to $\Omega$ and, in odd characteristic, extending to
$\operatorname{SO}$. This is a strict level/cover reduction,
not a finite classification of all orthogonal sources.
The highest-level theta constituent in
\cite[Corollary~4.7, p.~174]{OL:GLT2024} is proved nonzero;
that statement is not used as an irreducibility or exhaustion
theorem. No extension to the full orthogonal group is asserted.

\paragraph{Exact certificate and source boundary.}
The standalone program \texttt{orthogonal\_levels/check\_constants.py}
checks the rational exponent constants, the polynomial identities
in (OL.18)--(OL.19), the four spin/extension exponent reductions,
and finite boundary instances of the level inequality. It is
an arithmetic certificate for the displayed universal proof,
not an enumeration of enormous finite groups. Complete invoked
published proofs and the accepted Larsen--Tiep proof were read;
no new proof of their classification or Deligne--Lusztig inputs
is claimed. The source log preserves exact versions and hashes.


\section{An extension and reality gate for ordinary orthogonal levels}
\label{sec:orthogonal-lowlevel-reality}

\paragraph{Status and scope.}
This is a new author derivation.  The complete proof, invoked source
statements and exact author certificate have passed the author check
and complete independent mathematical reviews by /root and
/root/g34\_eightfolds. The peer also supplied a separately authored
exact supplement; all actual scopes and timings are recorded.  The new orbit argument is separate
from the published character estimates in D-OL-01.  It gives a strict
reduction of the remaining unbounded natural dimension case, without
asserting that any surviving case is a source.

Write $O(V)$ for the full isometry group of a nondegenerate quadratic
space over $\mathbf F_q$, and put $S=\Omega(V)=[O(V),O(V)]$.
For odd $q$, the natural dimension $d$ may be even or odd; for even $q$
we use a nondefective even-dimensional quadratic space.  The
orthogonal level $\ell$ is the level of Definition~3.2(ii) of
\cite{OL:GLT2024}.  In particular, Lemma~3.3 there identifies its
constituents on $S$ with those of powers of the permutation character
$\tau_V$ on the vector set $V$.

\begin{lemma}[D-OL-04: extension from a short tuple action]
\label{D-OL-04}
Assume $d\ge10$.  If an irreducible complex character $\rho$ of $S$
has level $\ell$ and
\[
                         d\ge4\ell+2,                         \tag{OL4.1}
\]
then $\rho$ extends to an irreducible character of $O(V)$ and is
self-dual.  Consequently every nonselfdual irreducible character of
$S$ satisfies
\[
                         d\le4\ell+1.                         \tag{OL4.2}
\]
The extension assertion concerns an actual representation, not just
invariance of its character under outer automorphisms.
\end{lemma}

\begin{proof}
Let $\mathcal M=\mathbf C[V^\ell]$.  We prove that $O(V)$ and $S$
have identical orbits on $V^{2\ell}$, which is the set indexing the
matrix entries of endomorphisms of $\mathcal M$.

Fix a tuple of $2\ell$ vectors, and let $R$ be its span, of dimension
$t\le2\ell$.  For the nondegenerate polar form on $V$, the radical
of $R^\perp$ is $R\cap R^\perp$.  Its dimension is at most $t$.
Thus the induced nondegenerate polar space after taking out the radical
has dimension at least $d-2t\ge d-4\ell\ge2$.  Choose a vector-space
complement $U$ to this radical inside $R^\perp$.  The polar form
restricted to $U$ is nondegenerate; in characteristic two this means
that the restricted quadratic space is nondefective.  Hence
\[
 V=U\perp U^\perp,\qquad R\subseteq U^\perp,\qquad \dim U\ge2.
                                                               \tag{OL4.3}
\]
Every element of $O(U)$ extended by the identity on $U^\perp$ fixes
the chosen tuple pointwise.

We check that this embedded $O(U)$ surjects onto $O(V)/S$.
For odd $q$, determinant and spinor norm identify this quotient with
$\{\pm1\}\times\mathbf F_q^*/(\mathbf F_q^*)^2$;
see the precise subgroup descriptions in
\cite[Section~2, pp.~41--43]{OL:KTV2023}.
A nondegenerate binary subspace of $U$ has a diagonal quadratic form
$ax^2+by^2$ with $ab\ne0$.  It represents every nonzero $c\in\mathbf F_q$:
the two subsets $a\{x^2:x\in\mathbf F_q\}$ and
$c-b\{y^2:y\in\mathbf F_q\}$ each have $(q+1)/2$ elements, so they
intersect.  Reflections in vectors of square and nonsquare nonzero
norm therefore belong to $O(U)$ and have quotient classes
$(-1,1)$ and $(-1,\eta)$, respectively, where $\eta$ is the
nontrivial square class.  These classes generate the entire quotient.
The same conclusion holds for either type of $U$ and for odd $d$.

For even $q$, the quotient $O(V)/S$ has order two
\cite[Section~2, p.~42]{OL:KTV2023}.  There is a vector $v\in U$
with $Q(v)\ne0$, since otherwise the polar form on $U$ would vanish.
The symmetry
\[
              s_v(x)=x+\frac{B(x,v)}{Q(v)}v                    \tag{OL4.4}
\]
is an isometry supported on $U$ and has $\operatorname{rank}(s_v-1)=1$.
It is outside $S$: the even-dimensional fixed-space property on $S$
used in the proof of \cite[Lemma~3.3, p.~160]{OL:GLT2024} would
otherwise give even rank.  Thus $O(U)$ surjects onto the quotient
also in characteristic two.

For any $g\in O(V)$, choose a tuple stabilizer $h\in O(U)$ with the
same class as $g$ in $O(V)/S$.  Then $gh^{-1}\in S$ sends the tuple
to exactly its image under $g$.  This proves equality of the two orbit
partitions.  Since an equivariant matrix on a permutation module is
exactly a matrix constant on the corresponding pair orbits, we obtain
an equality of algebras, not only an equality of their dimensions:
\[
       \operatorname{End}_{S}(\mathcal M)
       =\operatorname{End}_{O(V)}(\mathcal M).                 \tag{OL4.5}
\]

Decompose $\mathcal M$ into irreducible $O(V)$-summands over
$\mathbf C$.  For an irreducible summand $W$, every
$S$-endomorphism of $W$ extends by zero on an $O(V)$-stable complement
to an $S$-endomorphism of $\mathcal M$.  Equation (OL4.5) makes it
$O(V)$-equivariant, so its restriction to $W$ is scalar by Schur's
lemma.  Complete reducibility then shows that $W|_S$ is irreducible.
The same argument applied to homomorphisms between two summands
shows that distinct $O(V)$-types restrict to distinct $S$-types.

On $S$ one has $\tau_V+\zeta_V=2\tau_V$ by
\cite[Lemma~3.3, p.~160]{OL:GLT2024}.  Thus a character of level
$\ell$ occurs in $\mathcal M|_S$, and the preceding paragraph
supplies its irreducible extension to $O(V)$.
Finally the full finite orthogonal group $O(V)$ is real in both
characteristics: every element is a product of two involutions.
For odd characteristic this is Wonenburger's theorem, restated
precisely in \cite[Theorem~1, p.~44]{OL:KTV2023}; for characteristic
two the same published paper states the Ellers--Nolte and Gow theorem
explicitly on p.~40.  These known results imply that every irreducible
character of $O(V)$ is real-valued.  Its representation is therefore
isomorphic to its dual.  Restricting the extension proves the
self-duality of $\rho$.  No assertion about its Frobenius--Schur sign
is needed.  The contrapositive gives (OL4.2).
\end{proof}

\begin{theorem}[D-OL-05: the remaining unbounded orthogonal case]
\label{D-OL-05}
Retain the actual six-dimensional multiplicity variation and
rank-nine period-map setting of D-OL-02--03.  Suppose the inducing
image $J$ has a normal irreducibly acting central product of
quasisimple orthogonal factors, permuted transitively by $J$.
Let $d$ and $q$ be their common natural dimension and field size,
and let the original irreducible representation have imaginary
quadratic character field.  Put $J_0=10^{30}$ as in (OL.1).

Every nonselfdual ordinary factor character satisfies
\[
                         d\le4J_0-3.                         \tag{OL4.6}
\]
More generally, if
\[
                         d\ge4J_0-2,                         \tag{OL4.7}
\]
then all factor characters are ordinary, have level less than
$J_0$, extend irreducibly to the full orthogonal group, and are
self-dual.  The inducing image $J$ is essentially self-dual with
nontrivial multiplier.  In addition,
\[
                 q\in\{2,3,4,5,7,8,9,11,13\}.                \tag{OL4.8}
\]
In particular every transitive orthogonal-core source with
$q\ge16$ has $d\le4J_0-3$.  All conclusions allow arbitrary
compatible scalar, diagonal, graph and field normalizers, and
arbitrary transitive factor permutations.  They do not require
equality between the inducing character field and the original
imaginary quadratic field.
\end{theorem}

\begin{proof}
D-OL-03 gives $\ell<J_0$ for every ordinary factor character.
Equation (OL4.2) immediately gives (OL4.6).  Under (OL4.7), the
spin and nonextension cutoffs of D-OL-03 already force the factor
characters to be ordinary.  Since $\ell\le J_0-1$, one has
$4\ell+2\le4J_0-2\le d$.  D-OL-04 therefore gives all the stated
extension and self-duality assertions.

Let $N\triangleleft J$ be the indicated irreducible central
product and let $e$ be the degree of the inducing representation
$\tau$.  The irreducible $N$-module is a tensor product of the
irreducible factor modules.  Tensoring their nondegenerate invariant
bilinear forms yields such a form for $N$.  Its invariant-form line
is one-dimensional by Schur's lemma, and normality makes $J$
preserve that line.  Thus $J$ has a multiplier character
$\mu:J\to\mathbf C^*$, whether the resulting form is symmetric or
alternating.  This also handles elements permuting the factors;
no separate extension or compatibility assumption is needed.

If $\mu$ were trivial, $\tau$ would be self-dual and its character
would be real-valued.  But the original character is induced from
$\tau$, so the induction formula puts the original character field
inside the character field of $\tau$ at the given complex embedding.
The former is imaginary quadratic and cannot lie in $\mathbf R$.
Consequently $\mu$ is nontrivial.

Suppose now that $q\ge16$.  The published uniform character ratio
bound used in D-BI-18 is
\[
 c(q)=\max\{1/(\sqrt q-1),9/q\}\le9/16.                       \tag{OL4.9}
\]
Indeed the first term is at most $1/3$ and the second at most
$9/16$.  This is \cite[Theorem~4.1(2), p.~141]{BI:DHP}, with its
quasisimple Lie-type and nonalternating quotient hypotheses
satisfied here.  Tensoring the factor characters shows that every
nonscalar element of $N$ has character ratio at most $9/16$.
The complete commutator argument of D-BI-18, including the case
of a nonidentity scalar commutator, therefore gives
\[
             \nu(g)\ge\frac{7e}{64}\quad
                     \text{for every nonscalar }g\in J.      \tag{OL4.10}
\]
This conclusion uses only normality and irreducibility of $N$;
it includes all outer normalizers and all factor permutations.
A genuine nonidentity scalar inertia has actual moving rank $e$,
so it also satisfies the needed lower bound on moving rank.

The actual intermediate Hurwitz factorization gives genus zero,
at least twelve genuine branch points, and total moving rank
$2e+6$ as in D-BI-18.  The inertias generate $J$, have product one,
and $\mu$ is nontrivial, so at least two have multiplier different
from one.  For such an inertia the $1$-eigenspace pairs only with
the $\mu(g)$-eigenspace under the nondegenerate form, and hence
has dimension at most $e/2$.  The two corresponding moving ranks
are at least $e/2$.  The other at least ten inertias satisfy (OL4.10).
Consequently
\[
  2e+6\ge e+10\frac{7e}{64}=\frac{67e}{32},\qquad e\le64.
                                                               \tag{OL4.11}
\]
The factor characters are nontrivial.  The $k=1$ case of
\cite[Theorem~3.7, pp.~165--167]{OL:GLT2024} gives
\[
                         a\ge q^{d-4}(q-1),                  \tag{OL4.12}
\]
where $a$ is the degree of any one factor: below that threshold
the character has level zero and is trivial.  Here $d\ge10$ and
$q\ge16$, so $a\ge16^6\cdot15>64$.  Since $e$ is the product of
the factor degrees, $e\ge a$, contradicting (OL4.11).
Thus $q<16$, and listing the prime powers below sixteen gives
exactly (OL4.8).

All uses of the branch count are necessary consequences of the
actual period-map factorization.  We have not inferred dominance
from a dimension inequality, nor ruled out the ordinary self-dual
bounded-level cases in (OL4.8).
\end{proof}

\section{A uniform support bound for extendible orthogonal characters}
\label{sec:orthogonal-extendible-support}

\paragraph{Author check and review status.}
This is a new author derivation.  The complete written proof and
all invoked source passages have passed the author check.  The
exact author supplement also passed; it is separate from the
universal proof. Whole-proof reviews by /root and /root/g34\_eightfolds
pass, with independent exact arithmetic recorded separately. Actual
check scopes, hashes and times are recorded separately; no finite
experiment substitutes for an all-dimension argument below.

Let $V$ be a quadratic space of natural dimension $d\ge22$ over
$F=\mathbf F_q$, with nondegenerate polar form
$B(x,y)=Q(x+y)-Q(x)-Q(y)$.  In characteristic two $d$ is even.
Put $O=O(V)$ and $S=\Omega(V)=[O,O]$.  Let $\rho$ be a nontrivial
irreducible complex representation of $S$, of degree $a$, which
extends to an irreducible representation $R$ of $O$.  For a finite
projective-order complex operator $T$, define
\[
 \nu(T)=a-\max_{\lambda}\dim\ker(T-\lambda I).
 \tag{OS.1}
\]
Rescaling $T$ does not change this number.  A compatible normalizer
means an operator inducing an automorphism of $\rho(S)$.

\subsection*{D-OL-06. Two opposite abelian radicals and their weights}

Choose a hyperbolic pair $e,f$ and write
\[
 V=Fe\oplus W\oplus Ff,\qquad
 Q(xe+w+yf)=xy+Q_W(w),\qquad N=\dim W=d-2\ge20.
 \tag{OS.2}
\]
For $v\in W$ define
\[
 u_e(v)e=e,\quad
 u_e(v)w=w-B(v,w)e,\quad
 u_e(v)f=f+v-Q(v)e.
 \tag{OS.3}
\]
Direct substitution in (OS.2) proves that these are isometries and
that $u_e(v)u_e(v')=u_e(v+v')$.  Thus
$U_e=\{u_e(v):v\in W\}$ is the additive group of $W$.
Let $U_f$ be the opposite group obtained by interchanging $e$ and $f$.
Every nonidentity $u_e(v)$ has natural moving rank two.  In odd
characteristic it has odd prime order, and therefore belongs to $S$,
since $O/S$ is a $2$-group.  In characteristic two it belongs to $S$
by the Dickson-invariant criterion
\cite[Lemma~2.1, p.~383]{OS:TaylorVinroot2018}.  Consequently
$U_e,U_f\subseteq S$ in both characteristics.

We give a complete generation argument.  Put $K=\langle U_e,U_f\rangle$.
The stabilizer $P$ of $Fe$ in $O$ is $U_e\rtimes L$, where $L$ acts
by $e\mapsto ce$, $f\mapsto c^{-1}f$, and an arbitrary isometry on
$W$.  This description follows by first correcting the image of $f$
with a unique element of $U_e$ and then restricting to $W$.
The group $L$ normalizes both radicals, so $P$ normalizes $K$.
The isometry $w_0$ interchanging $e,f$ and fixing $W$ also normalizes
$K$.

For every nonsingular $z$, write
$r_z(x)=x-B(x,z)Q(z)^{-1}z$ for its reflection (or symmetry in
characteristic two).  If $z=xe+w+yf$ with $y=0$, then $r_z$ fixes
$e$, so lies in $P$.  If $y\ne0$, then
\[
 u_e(-w/y)z=(Q(z)/y)e+yf.
 \tag{OS.4}
\]
Reflection in the latter vector interchanges $Fe$ and $Ff$ with
nonzero scalar factors, so belongs to $Lw_0$.  Thus every reflection
normalizes $K$.  The full orthogonal group is generated by its
reflections, with the sole exception $O_4^+(2)$, which is outside
our range; this is \cite[Corollary~11.42, p.~159]{OS:Taylor1992},
including Theorems~11.39 and 11.41 and their characteristic-two
qualification.  Hence $K\triangleleft O$.  It is a noncentral
subgroup of $S$.  The perfectness of $S$ and simplicity of $S/Z(S)$,
\cite[Theorems~11.47--11.48, pp.~161--162]{OS:Taylor1992}, imply
$KZ(S)=S$, and then
$S=[S,S]=[KZ(S),KZ(S)]\subseteq K$.  Therefore
\[
                         \langle U_e,U_f\rangle=S.
 \tag{OS.5}
\]

The two radicals are conjugate inside $S$.  To see this without
assuming that $w_0\in S$, embed $O(W)$ in the pointwise stabilizer
of $e,f$.  It surjects onto $O/S$: in odd characteristic a binary
nondegenerate subspace represents both nonzero square classes and
its reflections generate both determinant/spinor classes; in even
characteristic one symmetry supplies the nontrivial Dickson class.
This is exactly the binary-complement calculation in D-OL-04.
Choose $l\in O(W)$ with $w_0l\in S$.  It conjugates $U_e$ to $U_f$.
Since $\rho$ is nontrivial irreducible, (OS.5) gives
\[
 b:=\dim \rho^{U_e}=\dim\rho^{U_f}\le a/2.
 \tag{OS.6}
\]
Indeed the two fixed spaces have zero intersection.

Fix a nontrivial additive character $\psi$ of $\mathbf F_p$.
The characters of $U_e$ are indexed by $z\in W$ through
\[
 u_e(v)\longmapsto
 \psi\bigl(\operatorname{Tr}_{F/\mathbf F_p} B(z,v)\bigr).
 \tag{OS.7}
\]
Their multiplicities in $\rho$ are constant on the orbits of the
embedded $\Omega(W)$ in the Levi subgroup.  These nonzero orbits
are precisely
\[
                  \mathcal O_t=\{z\in W\setminus\{0\}:Q(z)=t\},
                  \qquad t\in F.
 \tag{OS.8}
\]
Witt's extension theorem gives the assertion first for $O(W)$
\cite[Theorem~7.4, pp.~57--58]{OS:Taylor1992}.  The pointwise
stabilizer of a vector contains $O(A)$ for a nondegenerate
complement $A$ of dimension at least $N-2\ge2$ in its perpendicular
space.  The same binary-complement calculation makes this
stabilizer surject onto $O(W)/\Omega(W)$.  The $O(W)$-orbit and
$\Omega(W)$-orbit consequently coincide.  This proves (OS.8),
also for nonzero singular vectors.

Here are the needed orbit estimates with their full error bounds.
For any quadratic form $Q_0$ on a vector space $X$ of dimension $m$,
whose polar form has rank $s$, and any nontrivial additive character
$\eta$ of $F$, orthogonality gives
\[
 \left|\sum_{x\in X}\eta(Q_0(x))\right|^2
 =q^m\sum_{z\in\operatorname{rad}B_0}\eta(Q_0(z))
 \le q^{2m-s}.
 \tag{OS.9}
\]
The equality follows by setting the difference of the two summation
variables equal to $z$; the inner sum of $\eta(B_0(x,z))$ is zero
unless $z$ is in the radical.  Taking absolute values in the last
sum proves the inequality, including characteristic two and a
nonzero quadratic form on the polar radical.  Fourier inversion
therefore yields
\[
 \left|\#\{x\in X:Q_0(x)=t\}-q^{m-1}\right|
 \le q^{m-s/2}.
 \tag{OS.10}
\]
In particular, every orbit in (OS.8) has size at least
\[
                  L_N=q^{N-1}-q^{\lceil N/2\rceil}-1.
 \tag{OS.11}
\]
If a linear isometry $g$ of $W$ has moving rank two, then
$\dim\ker(g-1)=N-2$.  The polar form restricted to this fixed space
has radical of dimension at most two and hence rank at least $N-4$.
Applying (OS.10) gives
\[
 \frac{|\mathcal O_t\cap\ker(g-1)|}{|\mathcal O_t|}
 \le F_N(q):=
 \frac{q^{N-3}+q^{\lceil N/2\rceil}}
      {q^{N-1}-q^{\lceil N/2\rceil}-1}.
 \tag{OS.12}
\]
For $N\ge20$, division by $q^{N-1}$ gives the uniform bounds
\[
 F_N(q)\le\frac{q^{-2}+q^{-9}}{1-q^{-9}-q^{-19}}
 <\begin{cases}1/8,&q\ge3,\\1/3,&q\ge2.\end{cases}
 \tag{OS.13}
\]
For completeness the two inequalities follow by substituting
$q=3$ and $q=2$, respectively: the numerator decreases and the
positive denominator increases with $q$, and cross multiplication
uses only positive integers.  This is a universal inequality in
$N,q$, not a finite-range extrapolation.

Suppose now that an operator $A$ on the representation space
normalizes $U_e$ and induces an action of prime order $l$ on its
weight set.  For each length-$l$ weight orbit the corresponding
weight spaces have equal dimension $m$; an eigenspace of $A$ on
their direct sum has dimension at most $m$.  This follows by
successively recovering its $l-1$ other components from the first.
No assumption that $A^l$ is scalar is required.  If in each nonzero
orbit (OS.8) at most a proportion $f_0$ of the weights are fixed,
the constant multiplicities and (OS.6) give
\[
                \nu(A)\ge (1-1/l)(1-f_0)(a-b)
                         \ge\tfrac12(1-1/l)(1-f_0)a.
 \tag{OS.14}
\]
In particular, for a prime-order isometry of $V$ with moving rank
two which fixes a hyperbolic pair, (OS.13)--(OS.14) give
\[
 \nu(R(g))\ge7a/32\quad(q\text{ odd}),\qquad
 \nu(R(g))\ge2a/9\quad(q\text{ even},\ l\text{ odd}).
 \tag{OS.15}
\]
These estimates are valid for $g\in O$, not only for $g\in S$;
the irreducibility needed in (OS.6) concerns $R|_S$.

Every rank-two isometry $g$ used here does fix such a hyperbolic
pair.  Embed $\operatorname{im}(g-1)$ in a nondegenerate subspace
$A$ of dimension at most four: choose dual partners for the radical
of its restricted polar form.  Then $A$ is $g$-stable, and
$A^\perp\subseteq\ker(g-1)$ because
$\operatorname{im}(g-1)^\perp=\ker(g-1)$.
The nondegenerate space $A^\perp$ has dimension at least $d-4\ge18$
and contains a hyperbolic pair over a finite field
\cite[Theorem~11.2 and its ensuing decomposition, p.~138]{OS:Taylor1992}.

\subsection*{D-OL-07. Linear and semilinear normalizers}

We first specify completely why the normalizers below are covered
by natural semilinear similitudes.  For $d\ge22$, the orders of the
outer automorphism groups in
\cite[Section~2.2.2, Table~2.1, p.~14]{HES:LiXia2022} agree with
the following orders obtained from these similitudes.  Write
$q=p^f$, and use $\epsilon=1$ or $-1$ for the plus or minus type
when $d=2m$:
\[
\begin{array}{c|c}
\text{natural type}&|\Gamma\!GO(V)/(F^*\Omega(V))|\\ \hline
d\text{ odd},\ q\text{ odd}&2f\\
d\text{ even},\ q\text{ even}&2f\\
d=2m,\ q\text{ odd},\ q^m\not\equiv\epsilon\pmod4&4f\\
d=2m,\ q\text{ odd},\ q^m\equiv\epsilon\pmod4&8f.
\end{array}
 \tag{OS.16}
\]
Here the notation in the numerator means all semilinear similitudes,
and the quotient means its induced outer image; scalar intersections
are taken in computing the order.

We justify the calculation and its faithfulness.  One has
$|O:S|=4$ for odd $q$ and $2$ for even $q$.
For odd natural dimension a multiplier must be a square by taking
determinants; for even characteristic every multiplier is a square.
In these cases $GO=F^*O$.  For even dimension and odd $q$ the
multiplier map is surjective.  On hyperbolic planes scale the first
coordinate by the desired multiplier; on an anisotropic binary
summand use multiplication in $\mathbf F_{q^2}$ with the prescribed
norm.  The norm is surjective since finite-field multiplicative
groups are cyclic.  In odd characteristic and even dimension,
$-I\in S$ exactly when $q^m\equiv\epsilon\pmod4$: its determinant
is one and its spinor norm is the discriminant of $Q$, namely
$(-1)^m$ times a square in plus type and times a nonsquare in minus
type.  These facts give the linear factors $2,2,4,8$ in (OS.16).
The determinant/spinor description is
\cite[Section~2, pp.~41--43]{OL:KTV2023}; the reflection spinor
calculation is also \cite[p.~163]{OS:Taylor1992}.

Each field automorphism occurs in a semilinear isometry.  This
follows directly from the finite-field quadratic classification
in \cite[pp.~138--140]{OS:Taylor1992}: the split form has coefficients
in the prime field.  For odd characteristic the nonsplit binary
form may be written $x^2-\epsilon y^2$ with $\epsilon$ a nonsquare;
Frobenius is corrected by scaling the $y$ coordinate by
$\epsilon^{(p-1)/2}$.  In characteristic two write it as
$x^2+xy+\delta y^2$, where $\operatorname{Tr}(\delta)=1$.
Since $\delta^p-\delta$ has trace zero, choose $c$ with
$c^2+c=\delta^p-\delta$ and correct Frobenius by
$(x,y)\mapsto(x+cy,y)$.  In odd natural dimension rescale the form
so that its final one-dimensional coefficient is one.  Thus there
are all $f$ field components, even if a chosen lift has order
$2f$ rather than $f$.

The centralizer of $S$ in the semilinear group consists only of
linear scalars.  Indeed, an element centralizing $S$ preserves
every space $\operatorname{im}(u_e(v)-1)=\langle e,v\rangle$;
their intersection is $Fe$.  As $e$ is arbitrary, it fixes every
singular projective point.  With two mutually perpendicular singular
basis vectors $e_1,e_2$, fixing $F(e_1+e_2)$ first makes their
scalar factors equal, and fixing all $F(e_1+te_2)$ then forces its
field automorphism to fix every $t\in F$.  It is now linear.  Three
hyperbolic pairs and the singular vectors $e_i+f_j$ for $i\ne j$
force a common scalar on the hyperbolic part.  For a vector $z$ of
the anisotropic part, the singular vector
$e_1+z-Q(z)f_1$ forces the same scalar on $z$.
The same conclusion holds if the element acts trivially on
$S/Z(S)$: its commutators with $S$ are central and define a
homomorphism from the perfect group $S$ to $Z(S)$, hence are trivial.
Consequently the image of the semilinear similitudes in
$\operatorname{Aut}(S/Z(S))$ has precisely the outer order (OS.16).
Comparison with the published complete table proves that it is the
whole automorphism group.  Dimension eight triality does not occur.

Finally, a nontrivial irreducible $S$-representation has central
kernel.  The center of the natural $S$ has order at most two, so
every subgroup of it is preserved by the above automorphisms.
An automorphism of $\rho(S)$ and the induced natural lift which
agree on the simple quotient differ by a central homomorphism;
perfectness again makes that difference trivial.  Thus every
compatible normalizer $T$ is implemented on $S$ by a natural
semilinear similitude $H$.

There is an important extension issue.  The chosen extension $R$
need not itself be invariant under conjugation by $H$.  Instead
Schur's lemma gives a linear character $\kappa_H$ of $O/S$ such that
\[
 T R(x)T^{-1}=\kappa_H(x)R(HxH^{-1})\qquad(x\in O).
 \tag{OS.17}
\]
To prove this, both operators on the right and left induce the
same conjugation on $\rho(S)$; their ratio is scalar.  Multiplication
shows that these scalars form a character, and the character is
trivial on $S$.  For a reflection $r$, $\kappa_H(r)^2=1$.
For a linear $H=h$, it follows that
\[
 [T,R(r)]=\kappa_h(r)R([h,r]),\qquad
 \nu(R([h,r]))\le\operatorname{rank}([T,R(r)]-I)\le2\nu(T).
 \tag{OS.18}
\]
The last inequality follows by choosing $\lambda$ with
$\operatorname{rank}(T-\lambda I)=\nu(T)$ and writing
$TR(r)-R(r)T=(T-\lambda I)R(r)-R(r)(T-\lambda I)$.
The first inequality uses projective support, so the scalar
$\kappa_h(r)$ is harmless.  No equality of ordinary moving ranks
after discarding that scalar is asserted.

\paragraph{Linear normalizers in odd characteristic.}
If $h$ is not scalar, some reflection $r_v$ does not commute with
$h$, because reflections generate $O$ and its natural centralizer
is scalar.  Then
$k=[h,r_v]=r_{hv}r_v\ne1$ is an isometry of determinant one and
natural moving rank at most two.  A rank-one isometry in odd
characteristic is a reflection and has determinant $-1$, so $k$
has rank two.  The same argument applies to every nonidentity
power of $k$.  Choose a nonidentity prime-order power $g$ of $k$.
Its rank is two.  Equation (OS.15) and
$\nu(R(k))\ge\nu(R(g))$ imply from (OS.18) that
\[
                         \nu(T)\ge7a/64.
 \tag{OS.19}
\]
The rank-one assertion can also be checked directly: an isometry
$I+v\ell$ of rank one has $\ell=cB(v,-)$, and preservation of $Q$
gives $cQ(v)=-1$, so its determinant is $1+2cQ(v)=-1$.
The element $k$ need not belong to $S$; use of the actual extension
$R$ in (OS.15) is essential here.

\paragraph{Linear normalizers in characteristic two.}
Rescale $h$ to an isometry.  If some nonsingular $v$ satisfies
$B(v,hv)\ne0$, the plane $\langle v,hv\rangle$ is nondegenerate.
The product $k=r_{hv}r_v$ is a nonidentity rotation on this plane
and is the identity on its perpendicular space.  In a hyperbolic
plane these rotations have order dividing $q-1$; in an anisotropic
plane they are multiplication by norm-one elements of
$\mathbf F_{q^2}^*$, of order dividing $q+1$.  Hence a nonidentity
prime-order power $g$ has odd order and rank two.  Equations
(OS.15) and (OS.18) give
\[
                         \nu(T)\ge a/9>7a/64.
 \tag{OS.20}
\]

It remains to treat the case $B(v,hv)=0$ for every nonsingular $v$.
Put $F_0(v)=B(v,hv)$, a homogeneous quadratic polynomial.  We show
that it vanishes on every vector.  If $q\ge4$, for any nonzero
singular $x$ choose a nonsingular $y\in x^\perp$.  Such a $y$
exists because the polar rank of $x^\perp$ is $d-2>0$.
For every $t\ne0$, $x+ty$ is nonsingular, and
$0=F_0(x+ty)=F_0(x)+tB_{F_0}(x,y)$.
Two distinct nonzero $t$ give $F_0(x)=0$.
For $q=2$, the nondegenerate quadratic quotient
$x^\perp/Fx$ contains an anisotropic plane: this is immediate if
its anisotropic kernel is a plane, and in two hyperbolic planes
one may take $y=e_1+f_1$ and $z=e_1+e_2+f_2$.
Lift its vectors to $x^\perp$; then
$Q(y)=Q(z)=Q(y+z)=1$.  The three equations for
$x+y,x+z,x+y+z$ give
$B_{F_0}(x,y)=B_{F_0}(x,z)=B_{F_0}(x,y+z)=F_0(x)$;
bilinearity forces $F_0(x)=0$.

Polarizing the now identically zero $F_0$ gives $h^*=h$ for the
adjoint relative to $B$.  As $h$ is an isometry, $h^*=h^{-1}$,
so $h^2=I$.  Moreover $B(x,(h-I)y)$ is alternating, since its
diagonal values are $F_0(x)=0$.  Its rank, which equals
$\operatorname{rank}(h-I)$, is even.  The Dickson criterion puts
$h$ in $S$.  Unless $T$ is scalar, $h\ne I$ and is a noncentral
unipotent element.  The known bound
\cite[Theorem~2.4(i), with conventions on pp.~2--4]{HES:Liebeck2017}
gives $|\operatorname{tr}\rho(h)|\le3a/4$ for $q=2,4$ and at most
$a/(\sqrt8-1)<9a/16$ for $q=8$.
Here $S$ is the simply connected finite group in characteristic
two in these large dimensions, so the source hypotheses apply.
Since $T$ is a scalar times $\rho(h)$ by Schur's lemma,
\[
                         \nu(T)\ge a/8>7a/64.
 \tag{OS.21}
\]
The subtraction in $\sqrt8-1$ is outside the radical.  The
comparison with $9/16$ follows from $625<648$.

\paragraph{General semilinear normalizers for $q=4,8,9$.}
These are the only nonprime fields below sixteen.  Their field
degrees $f=2,3,2$ are prime.  If the field part of $H$ is trivial,
the linear argument applies.  Otherwise write its field
automorphism as $\sigma$, of order $f$, with fixed field $F_0$.
If $T^f$ is nonscalar, its natural implementing map $H^f$ is linear,
so (OS.19)--(OS.21) give the same bound for $T$.

Suppose instead that $T^f$ is scalar.  Its conjugation on $\rho(S)$
is trivial.  Projective faithfulness, perfectness, and the centralizer
calculation above imply $H^f=cI$ on $V$ for some $c\in F^*$.
Commuting $H$ with $H^f$ shows $\sigma(c)=c$.  Surjectivity of the
norm $F^*\to F_0^*$ lets us rescale $H$ so that $H^f=I$.
If $Q(Hv)=\mu\sigma(Q(v))$, then the norm of $\mu$ is one.
The image of $b\mapsto\sigma(b)/b$ has size
$(q-1)/(|F_0|-1)$ and equals the norm-one subgroup.  Choose $b$
with $\sigma(b)/b=\mu$ and replace $Q$ by $bQ$.
The isometry groups do not change, and now
\[
                      Q(Hv)=\sigma(Q(v)),\qquad H^f=I.
 \tag{OS.22}
\]

We include the vector descent proof; the case is not assumed to
be a pure field automorphism in the original coordinates.
For every $t\in F$ and $v\in V$, the vector
$\sum_{j=0}^{f-1}H^j(tv)$ is fixed by $H$.
If the fixed vectors did not span $V$ over $F$, a nonzero $F$-linear
functional $\lambda$ annihilating their span would give
\[
       \sum_{j=0}^{f-1}\sigma^j(t)\lambda(H^jv)=0
                  \quad\text{for all }t\in F,\ v\in V.
 \tag{OS.23}
\]
The distinct field automorphisms $1,\sigma,\ldots,\sigma^{f-1}$
are linearly independent as functions to $F$: evaluate a linear
relation on $1,z,\ldots,z^{f-1}$ for an element $z$ with distinct
conjugates, and use the Vandermonde determinant.  Equation (OS.23)
would therefore force $\lambda(v)=0$ for every $v$, a contradiction.
Choose an $F$-basis of fixed vectors.  The fixed vectors are then
exactly its $F_0$-linear span, and (OS.22) says that the quadratic
form on this $d$-dimensional $F_0$-space is nondegenerate.
Choose a hyperbolic pair over $F_0$ and apply (OS.2) in this basis.
The induced action on $U_e$ parameters, and on their dual weight
vectors, is $\sigma$ entrywise.

There are at most $|F_0|^N$ fixed weight vectors in all of $W$.
Consequently each orbit (OS.8) has fixed proportion at most
\[
 \frac{|F_0|^N}{L_N}
 \le\frac{q\,|F_0|^{-N(f-1)}}{1-q^{-9}-q^{-19}}<1/4
       \quad(N\ge20,\ q\in\{4,8,9\}).
 \tag{OS.24}
\]
The inequality follows already at $N=20$ in each of the three
displayed fields, and the numerator decreases with $N$ while
the exact denominator bound is fixed.  It is an exact rational
comparison.  Norm levels not fixed by $\sigma$ have no fixed
weight vectors at all.  Applying (OS.14) with $l=f$ gives
\[
                         \nu(T)\ge3a/16>7a/64.
 \tag{OS.25}
\]
This proves the desired support bound for all semilinear
normalizers, including diagonal--field and graph--field mixtures,
not merely for pure field lifts.

\begin{theorem}[D-OL-08: closing the ordinary orthogonal dimension escape]
\label{D-OL-08}
For $d\ge22$ and $q\in\{2,3,4,5,7,8,9,11,13\}$, every
nonscalar compatible normalizer of a nontrivial irreducible
$\Omega(V)$-representation which extends irreducibly to $O(V)$
satisfies $\nu(T)\ge7a/64$.

In the actual rank-nine six-dimensional multiplicity variation
setting of D-OL-02--05, suppose the inducing image has a normal
irreducibly acting central product of quasisimple orthogonal
factors, permuted transitively.  Then their natural dimension obeys
\[
                             d\le4\cdot10^{30}-3.
 \tag{OS.26}
\]
This conclusion allows all compatible scalars, outer automorphisms
and factor permutations, and allows the inducing character field
to be larger than the original imaginary quadratic field.  It is
a bounded-dimension reduction, not a classification or an existence
assertion for the surviving parameters.
\end{theorem}

\begin{proof}
The first assertion is (OS.19)--(OS.25), since the remaining fields
in the displayed set are prime.  To prove (OS.26), suppose instead
that $d\ge4\cdot10^{30}-2$.  D-OL-05 forces all factor characters
to be ordinary, to extend irreducibly to their full orthogonal
groups, to be self-dual, and to have $q$ in the displayed set.
Their common degree is $a$ and the degree of the inducing
representation is $e=a^t$.

For an element fixing all factors, the induced operator is, up to
scalar, a tensor product of compatible factor intertwiners.  A
nonscalar factor has support at least $7/64$ of its dimension;
eigenvalue multiplicities in a tensor product are convolutions,
so the same fraction bounds the support of the whole operator.
If an element moves factors in a cycle of length $l\ge2$, conjugate
the corresponding weighted cyclic permutation to one with a
single twist.  This twist has finite projective order and an
eigenbasis.  On tensor words in that basis there are at most $a$
constant words with orbits of length one; every other word orbit
has length at least two, and on an orbit of length $b$ an
eigenspace has dimension at most one.  Thus the largest eigenspace
on that cycle has dimension at most
\[
                    a+(a^l-a)/2\le3a^l/4.
 \tag{OS.27}
\]
Tensoring the remaining cycles preserves this upper fraction.
Hence every nonscalar inertia in the inducing image has projective
support at least $7e/64$.  A nonidentity scalar inertia has actual
moving rank $e$, so satisfies the required moving-rank bound too.

The normal-core invariant bilinear form makes the inducing image
essentially self-dual.  Its multiplier is nontrivial: if it were
trivial the inducing character would be real-valued, whereas its
character field contains the original imaginary quadratic field
by the induction formula.  This is the same actual form and field
argument as D-OL-05; no equality of those fields is assumed.
The actual Hurwitz factorization has genus zero, at least twelve
genuine branch points, and total moving rank $2e+6$.
At least two inertias have nontrivial multiplier because the
inertias generate and have product one.  Each then has moving
rank at least $e/2$, by pairing the $1$- and multiplier-eigenspaces.
The remaining at least ten have moving rank at least $7e/64$.
Therefore
\[
                 2e+6\ge e+10(7e/64),\qquad e\le64.
 \tag{OS.28}
\]
But the $k=1$ case of the published degree estimate
\cite[Theorem~3.7, pp.~165--167]{OL:GLT2024} gives
$a\ge q^{d-4}(q-1)\ge2^{18}>64$ for our nontrivial ordinary
factor and $d\ge22$.  Since $e\ge a$, this is a contradiction.
This proof uses the rank-nine assumption only through the already
proved actual intermediate Hurwitz factorization.  It never
infers dominance from a source dimension inequality.
\end{proof}

\section{Orthogonal support without an extension hypothesis}
\label{sec:orthogonal-nonextendible}

\paragraph{Status.}
This is a separate new author derivation.  It does not change the
reviewed D-OL-06--08 text.  The complete written proof and invoked
source interfaces have passed the author check, and the exact
author supplement passed. Whole-proof reviews by /root and
/root/g34\_eightfolds also pass, with a separate exact reconstruction.
The numerical supplement checks constants and finite witnesses;
it is not a substitute for the universal arguments below.

Keep the notation $S=\Omega(V)$, $O=O(V)$, $d=\dim V$ and
$\nu(T)=a-\max_\lambda\dim\ker(T-\lambda I)$.
In characteristic two the natural quadratic space is nondefective
and $d$ is even.  All representations in this section are ordinary
complex representations of the natural group $S$.

\subsection*{D-OL-09. A bound for every ordinary character when $q\ge3$}

\begin{theorem}[D-OL-09]
\label{D-OL-09}
Let $d\ge22$, $q\ge3$, and let $\rho$ be any nontrivial
irreducible complex representation of $S$, of degree $a$.
Every nonscalar compatible normalizer $T$ of finite projective
order satisfies
\[
                          \nu(T)\ge7a/64.                 \tag{ON.1}
\]
No extension of $\rho$ to $O$ is assumed.

Consequently, in the main imaginary-quadratic six-dimensional
multiplicity-variation setting, a normal irreducibly acting central
product of selfdual ordinary orthogonal factors with $q_i\ge3$
and $d_i\ge22$ cannot give period-image dimension nine.  The
factors may have different fields and dimensions, and all compatible
normalizers, scalars and factor permutations are allowed.
\end{theorem}

\begin{proof}
The abelian radicals $U_e,U_f$ and their generation proof in
D-OL-06 use only a nontrivial irreducible $S$-module.  In particular
\[
       \langle U_e,U_f\rangle=S,\qquad
       b:=\dim\rho^{U_e}=\dim\rho^{U_f}\le a/2.             \tag{ON.2}
\]
The descriptions of their weights and the $\Omega(W)$-orbits
$\mathcal O_t$ also require no extension.  We first enlarge the
moving-rank range in that calculation.

Suppose $g\in S$ has prime order and natural moving rank
$2\le r\le4$.  A nondegenerate subspace containing
$\operatorname{im}(g-I)$ can be chosen of dimension at most $2r$,
by adjoining dual partners to the radical of that image.  Its
perpendicular space is fixed by $g$ and has dimension at least
$d-2r\ge14$.  It contains a hyperbolic pair.  In coordinates
adapted to such a pair, put $N=d-2\ge20$ and use the line radical
$U_e$ on $W$ from D-OL-06.  The moving rank on $W$ is still $r$.
The fixed space in $W$ has dimension $N-r$, and its polar form
has radical dimension at most $r$.  The same quadratic Gauss-sum
calculation gives
\[
 \frac{|\mathcal O_t\cap\ker(g-I)|}{|\mathcal O_t|}
 \le
 \frac{q^{N-r-1}+q^{\lceil N/2\rceil}}
      {q^{N-1}-q^{\lceil N/2\rceil}-1}
 \le F_N(q)<1/8.                                           \tag{ON.3}
\]
Indeed, for the fixed space the exponent in the error is at most
$(N-r)-(N-2r)/2=N/2$.  Thus it is independent of $r$ in this
range.  The final inequality is the universal $q\ge3$ estimate
(OS.13).  Prime-order weight orbits and (ON.2) now imply
\[
                         \nu(\rho(g))\ge7a/32.              \tag{ON.4}
\]

The full natural semilinear-normalizer description in D-OL-07
holds for all $q$, not only for the small fields in its final
application: its outer-order calculation (OS.16), field lifts,
centralizer proof and represented-central-quotient argument were
all stated and proved with arbitrary $q=p^f$.  The complete
published outer table used there is
\cite[Section~2.2.2, Table~2.1, p.~14]{HES:LiXia2022}.
Accordingly a compatible normalizer is implemented on $S$ by a
natural semilinear similitude $H$.

First suppose this implementing map is linear, $H=h$.
If $T$ is nonscalar, $h$ is not a natural scalar.  Some element
$u$ of $U_e\cup U_f$ then fails to commute with $h$, because these
two radicals generate $S$ and their natural centralizer is scalar.
Set
\[
             k=huh^{-1}u^{-1}\in S\setminus\{1\}.          \tag{ON.5}
\]
Both $u$ and $huh^{-1}$ have natural moving rank two, so
$\operatorname{rank}(k-I)\le4$.  Every power of $k$ fixes its
fixed space.  In particular a nonidentity prime-order power $g$
has moving rank at most four and cannot be a nonidentity central
scalar.  Its rank is at least two: a rank-one odd-characteristic
isometry is a reflection of determinant $-1$, whereas in even
characteristic rank one is excluded by the Dickson invariant.
Both statements are the exact rank-one facts used in D-OL-07.
Thus (ON.4) applies to $g$.

Here all group operators are in $S$, so the commutator identity
is exact, with no choice of an $O$-extension:
\[
 [T,\rho(u)]=\rho(k),\qquad
 \nu(\rho(k))\le\operatorname{rank}(\rho(k)-I)\le2\nu(T).
                                                               \tag{ON.6}
\]
The last inequality follows by subtracting a scalar from $T$ in
$T\rho(u)-\rho(u)T$.  Since taking powers cannot increase
projective support, (ON.4)--(ON.6) prove (ON.1) for every linear
normalizer.  The commutator need not have rank two; retaining
the rank-four possibility is essential.

For a general $T$, first take a nonscalar power of prime order
$l$ in the projective linear group.  Its support is no larger
than $\nu(T)$, so it suffices to treat that power, still denoted
$T$.  Its field component has order either one or $l$.
The former is the linear case just proved.  In the latter case
write $F=\mathbf F_q$, $F_0=\mathbf F_{q_0}$ and
\[
                        q=q_0^l.                           \tag{ON.7}
\]
Because $T^l$ is scalar, projective faithfulness, perfectness and
the semilinear-centralizer calculation give $H^l=cI$ for some
$c\in F_0^*$.  The norm rescaling and quadratic-form rescaling
of (OS.22)--(OS.23) apply without any change for this cyclic
extension of prime degree: after rescaling $H$ and $Q$,
\[
                  H^l=I,\qquad Q(Hv)=\sigma(Q(v)).           \tag{ON.8}
\]
The averaging and independence-of-field-automorphisms proof there
gives an $F$-basis of $H$-fixed vectors and a nondegenerate
quadratic form over $F_0$.  The dimension $d\ge22$ guarantees a
hyperbolic pair over $F_0$.  In the resulting line radical, $H$
acts on weight vectors by entrywise $\sigma$.

At most $q_0^N$ weight vectors are fixed.  In any nonzero orbit
the fixed proportion is therefore at most
\[
 \frac{q_0^N}{q^{N-1}-q^{\lceil N/2\rceil}-1}
 \le\frac{q_0^{\,l-N(l-1)}}{1-q^{-9}-q^{-19}}
 \le\frac{2^{-18}}{1-4^{-9}-4^{-19}}<1/4.                   \tag{ON.9}
\]
Here $N\ge20$, $l\ge2$, $q_0\ge2$ and $q\ge4$; in
particular $l-N(l-1)\le-18$.  Norm levels not fixed by $\sigma$
have no fixed vectors.  The weight-orbit inequality gives
$\nu(T)\ge(1/2)(1-1/l)(3/4)a\ge3a/16$, stronger than
(ON.1).  This proves the first assertion for every field degree.

For the source consequence, the tensor product of the selfdual
factor modules has a unique nondegenerate invariant bilinear
form up to scalar.  The inducing image preserves its line.
Its similitude character is nontrivial: otherwise its character
would be real-valued, while the induction formula puts the
original imaginary quadratic field inside its character field.
This is exactly the already proved field and form interface
in D-OL-05 and D-OL-08.

On tensor products, factor-preserving operators retain the
support fraction $7/64$ by convolution.  Nontrivial factor cycles
have support at least $1/4$, as proved in (OS.27).  At least two
inertias have nontrivial similitude multiplier, each with actual
moving rank at least $e/2$.  The actual genus-zero rank-nine
factorization has at least twelve genuine inertias and total
moving rank $2e+6$.  Therefore
\[
                      2e+6\ge e+10(7e/64),\qquad e\le64.
                                                               \tag{ON.10}
\]
The genus-zero assertion does not precede the degree estimate
without justification: $a_i\ge q_i^{d_i-4}(q_i-1)>64$ by
\cite[Theorem~3.7, pp.~165--167]{OL:GLT2024}, so the actual Euler
formula excludes genus at least two; genus one has at most six
genuine marks and image rank at most six.  Finally $e\ge a_i>64$
contradicts (ON.10).  No equality of the inducing and original
character fields, and no sufficiency of a source-dimension
inequality, has been used.
\end{proof}

\subsection*{D-OL-10. Nonextendible characters in even characteristic}

\begin{theorem}[D-OL-10]
\label{D-OL-10}
Let $q$ be even, $d\ge1136$, and let $\rho$ be a nontrivial
ordinary irreducible $S$-module of degree $a$ which does not extend
to $O$.  Every nonscalar compatible finite-projective-order
normalizer satisfies
\[
                           \nu(T)\ge11a/64.                 \tag{ON.11}
\]
An irreducibly acting normal central product of such factor
modules, including mixed even fields and compatible factor
permutations, cannot supply image rank nine for its actual
six-dimensional multiplicity variation.  This conclusion does
not require selfduality or a determinant assumption.
\end{theorem}

\begin{proof}
We first prove a support bound for every nonidentity element of
$S$, without an extension hypothesis.  The natural group $S$ has
trivial center in these dimensions in even characteristic and is
simple.  Thus a nontrivial irreducible representation is faithful.
Taking a nonidentity prime-order power reduces the claim to an
element $g$ of prime order $l$.

Suppose first that $g$ fixes a hyperbolic pair.  Let $r\ge2$ be
its moving rank on the complementary space $W$, of dimension
$N=d-2$.  The numerator estimate in (ON.3) remains valid for
arbitrary $r$, not only $r\le4$.  If $r\le N/2$, the restricted
polar rank is at least $N-2r$ and its Gauss error is at most
$q^{N/2}$.  If $r>N/2$, the trivial bound on that Gauss error
is $q^{N-r}\le q^{N/2}$.  In both cases the fixed proportion
is at most $F_N(q)$.  Uniformly for $N\ge20$ and $q\ge2$,
\[
 F_N(q)\le\frac{q^{-2}+q^{-9}}{1-q^{-9}-q^{-19}}
 \le\frac{2^{-2}+2^{-9}}{1-2^{-9}-2^{-19}}<5/16.             \tag{ON.12}
\]
Together with (ON.2) and prime-order weight orbits this yields
$\nu(\rho(g))\ge(1/2)(11/16)(a/2)=11a/64$.

Suppose that no hyperbolic pair is fixed.  If $l$ is odd, $g$
is semisimple and its fixed space is a nondegenerate quadratic
space.  A nondegenerate quadratic space of dimension at least
three over a finite field contains a hyperbolic pair.  Thus
$\dim\ker(g-I)\le2$.
If $l=2$, put $r=\operatorname{rank}(g-I)$.  The fixed space
has dimension $d-r$, and its polar radical has dimension at most
$r$, because its perpendicular space is $\operatorname{im}(g-I)$.
Its polar rank is therefore at least $d-2r$.  Polar rank at least
four would give a four-dimensional nondegenerate quadratic
subspace, hence a hyperbolic pair.  Consequently
\[
 \begin{array}{ll}
 d-\dim\ker(g-I)\ge d-2, & l\text{ odd},\\
 r\ge d/2-1,             & l=2.
 \end{array}                                                   \tag{ON.13}
\]
In either case the left side, which is the natural support
restricted to $\mu_2=\{1\}$ in characteristic two, is at least
$567$ when $d\ge1136$.

Apply the exact published estimate
\cite[Theorem~4.3.6, (4.3.2), pp.~1917--1918; complete proof
pp.~1918--1923; Remark~4.3.7, p.~1923]{HOS:LST2011} with
$N=567=7\cdot9^2$.  It is stated for rational points of the
simple simply connected classical group, including
$\operatorname{Spin}_d$.  Inflate $\rho$ along the standard
simply connected finite cover onto $S$.  The inflated character
is still nontrivial and irreducible, and its natural action is
the same orthogonal action used in (ON.13).  The applicable
cover is the ordinary nonexceptional one in these large
even-characteristic dimensions; see the cover discussion
\cite[p.~4; Chapter~5, Table~1, p.~35]{BI:HHM}, already used in
D-OL-01.  No faithfulness premise appears in the published
character estimate, so inflation causes no change in its scope.

For $2\le q<43$ the published $\gamma_q$ is $19/20$; for larger
$q$ it is $1/(\sqrt q-1)\le19/20$.
The support convention in its remark allows precisely $\mu_2$
for the spin group, so (ON.13) is the required support, not an
unjustified algebraic-closure support substitution.  We obtain
\[
 \frac{|\chi(g)|}{a}<(19/20)^9,\qquad
 \nu(\rho(g))>
 \frac{1-(19/20)^9}{2}a>11a/64.                              \tag{ON.14}
\]
The last step uses $|\chi(g)|\ge a-2\nu(\rho(g))$.
This proves (ON.11) for every nonidentity inner operator.

We now explain why nonextendibility is exactly the needed linear
normalizer condition.  In even characteristic $O/S$ is cyclic of
order two, with an involution representative given by a symmetry.
If $\rho$ is invariant under that outer coset, choose an intertwiner
$A$ for a symmetry $r$.  Schur's lemma gives $A^2=cI$.
Rescale $A$ so $A^2=I$; assigning $r\mapsto A$ then extends
$\rho$ to the semidirect product $O=S\rtimes\langle r\rangle$.
Conversely any extension makes its character invariant.  Hence a
nonextendible character has no compatible operator from the
nontrivial $O/S$ coset.

Every linear similitude is a scalar times an isometry in even
characteristic.  The complete natural normalizer description
therefore makes a compatible linear normalizer a scalar times
an inner operator.  It satisfies the just-proved bound.
For a general compatible $T$, take a projective-prime-order power
as in D-OL-09.  If its field part is nontrivial, the complete
norm/form descent and uniform weight count (ON.7)--(ON.9) give
$\nu(T)\ge3a/16>11a/64$.  This includes arbitrary field degrees
and mixed graph--field lifts; the lift is not presumed to be a
pure Frobenius map.  The scalar and inner cases are already
covered.  Thus (ON.11) holds for every compatible normalizer.

Convolution and the weighted factor-cycle bound give support at
least $11e/64$ for every nonscalar inertia on the tensor product.
A nonidentity scalar has actual moving rank $e$.  The factor
degree bound $a_i\ge q_i^{d_i-4}(q_i-1)>96$ excludes genus at
least two in the actual Euler formula.  Genus one again gives
at most six genuine marks and image rank at most six.
A rank-nine source would therefore have genus zero and at least
twelve genuine inertia points.  It would satisfy
\[
                         2e+6\ge12(11e/64),\qquad e\le96,
                                                               \tag{ON.15}
\]
contradicting $e\ge a_i>96$.  This source contradiction needs
no invariant bilinear form.
\end{proof}

\subsection*{D-OL-11. Strict remaining ordinary orthogonal bounds}

\begin{corollary}[D-OL-11]
\label{D-OL-11}
In the main imaginary-quadratic sixfold problem, assume the
inducing image has a transitive irreducibly acting normal central
product of ordinary quasisimple orthogonal factors.  Then:
\begin{enumerate}
\item if $q\ge3$ and the factor module is selfdual, its natural
dimension satisfies $d\le21$;
\item if $q$ is even, irrespective of selfduality, its natural
dimension satisfies $d\le1134$.
\end{enumerate}
These are necessary bounds and do not assert existence at any
remaining parameter.
\end{corollary}

\begin{proof}
The first assertion is D-OL-09.  For the second, suppose instead
that $d\ge1136$.  Transitivity preserves extension to $O$:
the natural semilinear similitudes implementing automorphisms
normalize $O$, so an $O$-extension pulls back to an $O$-extension.
Thus either all factor modules extend or none does.
In the latter case D-OL-10 gives the contradiction, without a
selfduality assumption.  In the former case every factor module
is selfdual, because all irreducible $O$-characters are real-valued,
as used at the exact primary locators in D-OL-04.
For $q\ge4$, D-OL-09 excludes this case already at $d\ge22$.
For $q=2$, the reviewed extendible-character theorem D-OL-08
provides the same exclusion, also at $d\ge22$.
Finally $d$ is even in characteristic two, so $d<1136$ means
$d\le1134$.
\end{proof}

\paragraph{Scope of the new inputs and certificate.}
The only new published numerical input beyond D-OL-06--08 is
the explicitly stated Larsen--Shalev--Tiep estimate and its
restricted-scalar support remark, at the full primary locators
given above.  The generation, Gauss-sum, normalizer, descent,
degree and actual Hurwitz interfaces are the already proved
ones, with their exact hypotheses retained.  The supplement
checks the rational constants and an actual rank-four commutator;
it does not infer a universal representation theorem from finite
matrix examples.

\section{An explicit odd-dimensional nonselfdual orthogonal bound}
\label{sec:orthogonal-odd-nonselfdual}

\paragraph{Status and scope.}
This is a separate new author derivation.  The complete written
proof and exact supplement pass the author check; whole-proof
reviews by /root and /root/g34\_eightfolds have also passed, with
a separate exact peer supplement.  It leaves D-OL-06--11 unchanged.  All modules
below are ordinary complex modules of the natural orthogonal
derived group, rather than arbitrary spin modules.  The field is
of odd order.  All compatible finite-projective-order normalizers
are included, with arbitrary field components.

Put
\[
 \delta_*:=\frac{1-(19/20)^9}{2},\qquad
             \frac16<\delta_*<\frac3{16},\qquad
             \frac6{12\delta_*-2}<28.                       \tag{OO.1}
\]
The three displayed inequalities are exact rational inequalities.

\begin{theorem}[D-OL-12]
\label{D-OL-12}
Let $q\ge5$ be odd and let $V$ be a nondegenerate quadratic
space of odd dimension $d\ge1135$ over $\mathbf F_q$.
Write $S=\Omega(V)$ and $O=O(V)$.  Let $\rho$ be a nonselfdual
ordinary irreducible complex $S$-module of degree $a$.
Every nonscalar compatible normalizer $T$ of finite projective
order satisfies
\[
                         \nu(T)>\delta_*a.                  \tag{OO.2}
\]
An irreducibly acting normal central product of such modules
cannot supply period-image rank nine for its actual six-dimensional
multiplicity variation.  Mixed odd fields and odd natural
dimensions, represented central quotients, scalars and all
compatible factor permutations are allowed.  No determinant
assumption or selfdual structure on the inducing representation
is needed for this assertion.
\end{theorem}

\begin{proof}
We first prove the inner support bound, without using
nonselfduality.  The natural group $S$ has trivial center in odd
dimension and is simple in this range.  Its projective kernel
on any nontrivial irreducible complex module is therefore trivial:
that kernel is normal, and if it were all of $S$ the perfect
group $S$ would act by scalars.  It suffices to take a nonidentity
prime-order power $g$ of an inner operator, because taking powers
cannot increase $\nu$.

Suppose first that one of the natural eigenspaces of $g$ for
$\epsilon\in\{1,-1\}$ contains a hyperbolic pair $e,f$.
Write
\[
 V=Fe\oplus W\oplus Ff,\qquad
 Q(xe+w+yf)=xy+Q_W(w),\qquad N=d-2.                          \tag{OO.3}
\]
The two opposite abelian line radicals $U_e,U_f$ of D-OL-06
are conjugate inside $S$ and generate $S$.  Thus
\[
                       b:=\dim\rho^{U_e}\le a/2.           \tag{OO.4}
\]
The formula for $u_e(v)$ in (OS.3) gives the exact identity
\[
              gu_e(v)g^{-1}=u_e(Av),\qquad A=\epsilon g|_W.
                                                               \tag{OO.5}
\]
The same isometry $A$ acts on the dual weight vectors, up to
inversion, which does not alter fixed points or orbit lengths.
If $\epsilon=-1$, then $g$ has order two and $A^2=I$.
If $\epsilon=1$, $A$ has the prime order of $g$.  In either
case $A$ is nonidentity: otherwise $g=\epsilon I$ on $V$,
and neither the identity nor $-I$ is the nonidentity element
of $S$ under consideration.  Consequently the action on $W$
has a prime order $l$ and moving rank $r\ge1$.

It is important here to allow $r=1$.  For example, an element
which is minus a reflection can lie in $S$ in odd dimension;
the replacement $A=\epsilon g|_W$ is then a reflection.  No
rank-two lower bound is imposed in this argument.

Let $X=\ker(A-I)$.  It has dimension $N-r$, and its restricted
polar radical has dimension at most $r$.  For $r\le N/2$, the
quadratic Gauss error (OS.10) is at most
\[
              q^{(N-r)-(N-2r)/2}=q^{N/2}.
\]
For $r>N/2$, the trivial bound on the Gauss sum is at most
$q^{N-r}\le q^{N/2}$.  As in D-OL-06, the nonzero Levi orbits
are exactly the norm levels
$\mathcal O_t=\{x\ne0:Q_W(x)=t\}$, and weight multiplicities
are constant on every such orbit.  Hence the proportion fixed
by $A$ in any one of them is at most
\[
 \begin{split}
 \frac{q^{N-r-1}+q^{\lceil N/2\rceil}}
      {q^{N-1}-q^{\lceil N/2\rceil}-1}
 &\le \frac{q^{-1}+q^{-9}}{1-q^{-9}-q^{-19}}\\
 &\le \frac{5^{-1}+5^{-9}}{1-5^{-9}-5^{-19}}<\frac14.
 \end{split}                                                 \tag{OO.6}
\]
Here $N\ge20$; dividing by $q^{N-1}$ gives the exponent
$1-\lfloor N/2\rfloor\le-9$.  The numerator decreases and
the positive denominator increases with $q$, so the final
inequality is uniform, not a finite-field extrapolation.
Every remaining weight orbit has length $l$.  The component
recovery argument (OS.14), valid for arbitrary scalar phases,
now yields
\[
           \nu(\rho(g))\ge(1-1/l)\frac34(a-b)
                           \ge\frac3{16}a>\delta_*a.        \tag{OO.7}
\]

Suppose next that neither natural $\epsilon$-eigenspace contains
a hyperbolic pair.  If the prime order of $g$ differs from the
characteristic, $g$ is semisimple.  Its $1$- and $-1$-eigenspaces
are nondegenerate: their orthogonal complements are the other
primary spaces.  A nondegenerate quadratic space of dimension
at least three over a finite field of odd order contains a
hyperbolic pair.  Thus both eigenspaces have dimension at most
two, and
\[
            \min_{\epsilon=\pm1}\operatorname{rank}
                      (g-\epsilon I)\ge d-2.               \tag{OO.8}
\]
This also shows directly that the present case cannot occur for
an involution when $d>4$.

If $g$ has the characteristic prime order, it is unipotent.
Put $r=\operatorname{rank}(g-I)$.  The fixed space has dimension
$d-r$, and its polar radical has dimension at most $r$, because
its perpendicular space is $\operatorname{im}(g-I)$.
Its polar rank is at least $d-2r$.  Polar rank at least three
would give a nondegenerate three-dimensional subspace and hence
a fixed hyperbolic pair.  Therefore
\[
          r\ge(d-2)/2,\qquad \operatorname{rank}(g+I)=d.
                                                               \tag{OO.9}
\]
Since $d\ge1135$ is odd, the integer lower bound in (OO.9)
is at least $567$.  Equations (OO.8)--(OO.9) prove that the
natural support restricted to $\mu_2=\{1,-1\}$ is at least
$567$ in every remaining case.

Inflate $\rho$ along the ordinary simply connected finite cover
$\operatorname{Spin}(V)\longrightarrow S$.  It remains nontrivial
and irreducible, and has the same natural orthogonal action for
the support calculation.  The cover is nonexceptional in these
large dimensions; the precise cover scope is the one used in
D-OL-10.  Apply
\cite[Theorem~4.3.6, equation~(4.3.2), pp.~1917--1918;
complete proof pp.~1918--1923; Remark~4.3.7, p.~1923]{HOS:LST2011}
with the support threshold $567=7\cdot9^2$.
The remark permits exactly the restricted scalar set $\mu_2$
for the spin group.  The published constant is $19/20$ for
$2\le q<43$, and $1/(\sqrt q-1)\le19/20$ for the remaining
fields.  No faithful-character hypothesis is required.
Consequently
\[
             |\chi(g)|/a<(19/20)^9,
             \qquad \nu(\rho(g))>\delta_*a.                 \tag{OO.10}
\]
The second inequality follows from
$|\chi(g)|\ge a-2\nu(\rho(g))$.  This proves the inner
bound for every nonidentity $g\in S$.

We now use nonselfduality to determine the compatible linear
normalizers.  In odd dimension a linear similitude with multiplier
$\mu$ has $(\det h)^2=\mu^d$, so $\mu$ is a square.  Hence
\[
 GO(V)=F^*O(V),\qquad O(V)=\langle-I\rangle\times SO(V),
                    \qquad SO(V)/S\simeq C_2.               \tag{OO.11}
\]
The last quotient is the determinant/spinor-norm description
\cite[Section~2, pp.~41--43]{OL:KTV2023}.  Thus there is only
one nontrivial linear outer coset after natural scalars and
inner automorphisms have been removed.

If that coset stabilized $\rho$, then all of $O$ would stabilize
its isomorphism class: the additional factor $-I$ is central.
Take an irreducible constituent of $\operatorname{Ind}_S^O\rho$.
Clifford theory makes its restriction a positive multiple of
$\rho$, since the conjugacy orbit of $\rho$ would be a singleton.
But all irreducible characters of $O$ are real-valued, by the
published full-group reality theorem
\cite[Theorem~1, p.~44; attribution and exact domain
pp.~39--40]{OL:KTV2023}.  Its restricted character is therefore
real, forcing $\rho$ to be selfdual, a contradiction.
Every compatible linear normalizer is consequently a scalar
times an inner operator, and satisfies (OO.10).

For completeness, this argument is applied to every linear
component of the full natural semilinear-normalizer description,
not to a chosen subgroup of automorphisms.  The description and
its cardinality/centralizer proof in D-OL-07 apply for arbitrary
field degree in these dimensions; the complete published outer
table is \cite[Section~2.2.2, Table~2.1, p.~14]{HES:LiXia2022}.
For a general compatible $T$, take a nonscalar power of prime
projective order $l$.  If its field component is trivial, the
preceding paragraph applies.  Otherwise $q=q_0^l$.  The full
norm and quadratic-form descent (ON.7)--(ON.9) applies to its
possibly mixed linear--field lift: after scalar and form
rescaling, that lift has order $l$, fixes an $F$-basis, and its
quadratic form is defined over $F_{q_0}$.  A hyperbolic pair over
$F_{q_0}$ gives a line radical on which it acts by Frobenius.
The fixed-weight fraction is less than $1/4$, uniformly in
the arbitrary prime field-extension degree.  Equations
(ON.7)--(ON.9) therefore give
\[
                         \nu(T)\ge3a/16>\delta_*a.          \tag{OO.12}
\]
There is no assumption that the original lift was pure Frobenius.
Taking the power only decreased projective support, so (OO.2)
holds for the original normalizer.

On tensor products, convolution retains the support fraction
of any active nonscalar factor.  A nontrivial factor permutation
has projective support at least one quarter of the full degree,
which is greater than $\delta_*$.  A nonidentity scalar has full
actual moving rank.  Thus every genuine inertia moves more than
$\delta_*e$ dimensions.

The ordinary factor degree estimate
\cite[Theorem~3.7, pp.~165--167]{OL:GLT2024} gives
$a_i\ge q_i^{d_i-4}(q_i-1)>28$, so $e>28$.
The actual intermediate-cover Euler identity
$6=2e(h-1)+D$ excludes genus at least two.  In genus one the
moving sum is six, hence there are at most six genuine marks
and the actual period-image rank is at most six.  Rank nine
therefore requires genus zero and at least twelve genuine
inertia points.  It would force
\[
             2e+6>12\delta_*e,
             \qquad e<\frac6{12\delta_*-2}<28,              \tag{OO.13}
\]
a contradiction.  The argument uses the actual intermediate
variation and does not identify the inducing character field
with the original imaginary quadratic field.
\end{proof}

\begin{corollary}[D-OL-13]
\label{D-OL-13}
In the main imaginary-quadratic sixfold problem, a transitive
irreducibly acting normal central product of ordinary orthogonal
factors of odd natural dimension over an odd field $q\ge5$
must have
\[
                              d\le1133.                    \tag{OO.14}
\]
If the factor module is selfdual, the sharper earlier bound
$d\le21$ applies.  These are necessary bounds, not existence
claims for the remaining parameters.
\end{corollary}

\begin{proof}
Selfduality is preserved by automorphic pullback, so transitivity
makes all factor modules selfdual or all nonselfdual.  The first
case is D-OL-09.  The second is D-OL-12.  Since $d$ is odd,
the exclusion $d\ge1135$ is equivalent to (OO.14).
\end{proof}

\paragraph{Why this calculation does not include $q=3$.}
The rank-one possibility in (OO.5) cannot be discarded.  For
$Q_W=\sum_{j=1}^s x_jy_j+z^2$ and $A(x,y,z)=(x,y,-z)$,
the fixed fraction on the nonzero singular level is
\[
 \frac{q^{2s-1}+(q-1)q^{s-1}-1}{q^{2s}-1}.                  \tag{OO.15}
\]
The numerator and denominator follow by successively adding
hyperbolic planes: $xy=t$ has $2q-1$ solutions for $t=0$ and
$q-1$ for $t\ne0$, giving the displayed norm-level counts by
induction, with the zero vector removed in both counts.
At $q=3$ this exceeds $1/3$ when $s>1$, so the uniform
$1/4$ bound used here is false in that case.  Minus a reflection
has spinor norm one on a natural space of dimension $4k+1$,
as is seen by writing its action on the $2k$ hyperbolic planes
as products of two reflections with norm product $-1$.
Thus the rank-one weight action is an actual inner possibility.
This is a limitation of this particular estimate and asserts
no existence of a characteristic-three source.

\paragraph{Source and certificate scope.}
The new steps are the signed hyperbolic-pair action, its rank-one
Gauss bound, the odd-dimensional nonselfdual normalizer argument,
and their exact source consequence.  The previously reviewed
radical generation, arbitrary-field descent, degree estimate,
and restricted-support character theorem are used with their
stated scopes.  The exact supplement checks all constants and
finite signed-action/orbit-count witnesses; the universal proof
is the one written above.

\section{The odd-dimensional orthogonal endpoint over three elements}
\label{sec:orthogonal-three-odd}

\paragraph{Status and conventions.}
This is a separate new author derivation.  The complete written
proof and exact supplement pass the author check; complete
proof reviews by /root and /root/g34\_eightfolds have also passed,
with a separate exact peer supplement.  It does not
change D-OL-12--13.  Let $S_d=\Omega_d(3)$ for odd $d\ge1135$,
and consider ordinary nontrivial irreducible complex modules
of this natural group.  Set
\[
 \beta:=\frac16-\frac1{10000}=\frac{4997}{30000},\qquad
 \delta_*:=\frac{1-(19/20)^9}{2}.
                                                               \tag{O3.1}
\]
A negative reflection means the actual natural isometry
$-r_v$, where $v$ is nonsingular and
$r_v(x)=x-B(x,v)Q(v)^{-1}v$.  Only those negative reflections
which belong to $S_d$ will be used.

\begin{lemma}[D-OL-14: the only low-support inner operators]
\label{D-OL-14}
Let $\rho$ be a nontrivial ordinary irreducible $S_d$-module of
degree $a$.  Every nonidentity $g\in S_d$ satisfies
\[
                          \nu(\rho(g))\ge\beta a.           \tag{O3.2}
\]
Unless $g$ is an actual negative reflection, the stronger bound
\[
                          \nu(\rho(g))>\delta_*a            \tag{O3.3}
\]
holds.  If $\rho$ is nonselfdual, every compatible linear
normalizer is a scalar times $\rho(g)$ for some $g\in S_d$.
Thus a nonscalar compatible normalizer with support at most
$\delta_*a$ is a scalar times a negative-reflection operator.
\end{lemma}

\begin{proof}
First take $g$ of prime order $l$.  If it fixes a hyperbolic
pair, use the opposite line radicals of D-OL-06.  Their fixed
space has dimension at most $a/2$.  On the complementary
natural space $W$, of dimension $N=d-2\ge20$, the moving rank
is at least two: a rank-one isometry in odd characteristic is
a reflection and has determinant $-1$, whereas $g\in S_d$.
The arbitrary-rank Gauss calculation (ON.3), as extended in
D-OL-10, bounds the fixed fraction in every nonzero norm orbit
by
\[
 \frac{3^{-2}+3^{-9}}{1-3^{-9}-3^{-19}}<\frac18.
                                                               \tag{O3.4}
\]
Its proof applies for every moving rank: up to half the dimension
the polar-rank bound gives Gauss error at most $3^{N/2}$; above
half the dimension the trivial error is already no larger.
Prime-order weight orbits therefore give
\[
               \nu(\rho(g))\ge\frac12\frac78\frac a2
                              =\frac7{32}a>\delta_*a.       \tag{O3.5}
\]

If $l$ is odd and no hyperbolic pair is fixed, the restricted
natural $\mu_2$-support is at least $567$.  Indeed, if $l\ne3$
the fixed space is nondegenerate of dimension at most two and
there is no eigenvalue $-1$, giving support at least $d-2$.
If $l=3$, put $r=\operatorname{rank}(g-I)$.  The fixed-space
polar rank is at least $d-2r$; absence of a hyperbolic pair
forces it to be at most two, so $r\ge(d-2)/2$.  As $d$ is
odd and at least $1135$, the integer rank is at least $567$.
Again $-1$ is not an eigenvalue.  Applying the exact
Larsen--Shalev--Tiep theorem and its restricted-support remark,
with the same ordinary Spin-to-$\Omega$ inflation and threshold
$567=7\cdot9^2$ as in (OO.8)--(OO.10), gives (O3.3).
The primary source is
\cite[Theorem~4.3.6, (4.3.2), pp.~1917--1918; complete proof
pp.~1918--1923; Remark~4.3.7, p.~1923]{HOS:LST2011}.

It remains to consider a prime-order involution with no fixed
hyperbolic pair.  Its $1$-eigenspace is nondegenerate of
dimension at most two.  Since $d$ is odd and the determinant
is one, that dimension is odd.  It is therefore one.  Its
orthogonal complement is the $-1$-eigenspace, so the involution
is precisely $-r_v$ for a nonsingular vector spanning its
$1$-eigenspace.

Such a negative reflection has a hyperbolic pair in its
$-1$-eigenspace.  In the signed coordinates (OO.3)--(OO.5),
the action on line-radical weight vectors is
$A=-g|_W$, an actual rank-one reflection.  The rank-one Gauss
bound, without replacing it by the false quarter bound, gives
\[
 \frac{|\mathcal O_t\cap\ker(A-I)|}{|\mathcal O_t|}
 \le c_3:=\frac{3^{-1}+3^{-9}}{1-3^{-9}-3^{-19}}
                  <1-4\beta.
                                                               \tag{O3.6}
\]
Together with the half-dimensional radical fixed-space bound
and the two-element weight orbits, this yields
\[
              \nu(\rho(g))\ge\frac{1-c_3}{4}a>\beta a.
                                                               \tag{O3.7}
\]
Thus all prime-order cases have been treated, with (O3.3)
failing to be asserted only for actual negative reflections.

We now classify arbitrary orders, rather than merely asserting
that some prime-order power is exceptional.  If the order of
$g$ has an odd prime divisor, a power of odd prime order is
already covered by (O3.3), and taking a power cannot increase
projective support.  Otherwise $g$ has order $2^s$.  If $s=1$,
the preceding involution analysis applies.  If $s\ge2$, put
$h=g^{2^{s-1}}$.  If $h$ is not a negative reflection, it
already gives the strong bound for $g$.

In the remaining case $h=-r_v$.  Since $g$ commutes with $h$,
it preserves the positive line $Fv$ and its perpendicular
$W_0$, on which $h=-I$.  A vector in a $1$- or $-1$-eigenspace
of $g$ is fixed by the even power $g^{2^{s-1}}$, and hence
belongs to $Fv$.  Thus both natural $\mu_2$-eigenspaces of
$g$ have dimension at most one.  The restricted support of
$g$ itself is at least $d-1\ge567$.  The published character
estimate applies directly to this element, without a
prime-order hypothesis, and proves (O3.3).  This proves the
claimed classification for every order.

Finally suppose $\rho$ is nonselfdual.  The normalizer argument
of (OO.11) uses only odd characteristic, not $q\ge5$.
For clarity, at $q=3$ there is no field automorphism; all
automorphisms in this large odd natural dimension are induced
by the full natural linear similitude group, by D-OL-07 and
\cite[Section~2.2.2, Table~2.1, p.~14]{HES:LiXia2022}.
In odd dimension that group is $F^*O(V)$ and
$O(V)=\langle-I\rangle\times SO(V)$, with
$SO(V)/S_d\simeq C_2$.

If the nontrivial linear outer coset stabilized $\rho$, all
of $O(V)$ would do so.  An irreducible constituent of
$\operatorname{Ind}_{S_d}^{O(V)}\rho$ would restrict to a
positive multiple of $\rho$.  Every irreducible character of
$O(V)$ is real-valued, by
\cite[Theorem~1, p.~44; exact domain and attribution,
pp.~39--40]{OL:KTV2023}.  This would make $\rho$ selfdual.
Hence the compatible linear normalizers are exactly scalar
times inner operators, as asserted.  This argument makes no
claim that all characters of $\Omega_d(3)$ are real.
\end{proof}

\begin{theorem}[D-OL-15: the remaining large odd-dimensional endpoint]
\label{D-OL-15}
A finite inducing image with an irreducibly acting nonempty normal central
product of nonselfdual ordinary $\Omega_{d_i}(3)$-modules,
all $d_i$ odd and at least $1135$, cannot supply period-image
rank nine for its actual six-dimensional multiplicity variation.
Mixed natural dimensions, scalars and all compatible factor
permutations are allowed; transitivity is not needed for this
assertion.

Consequently, in the main imaginary-quadratic sixfold problem,
every transitive ordinary orthogonal normal core of odd natural
dimension over any odd field satisfies
\[
                              d\le1133.                    \tag{O3.8}
\]
For selfdual factor modules the earlier bound $d\le21$ remains
valid.  These are necessary parameter bounds, not existence
claims at the remaining parameters.
\end{theorem}

\begin{proof}
Write $e=\prod_i a_i$ for the full inducing degree.  The
ordinary degree estimate of D-OL-01, or directly
\cite[Theorem~3.7, pp.~165--167]{OL:GLT2024}, gives
$a_i\ge2\cdot3^{d_i-4}>360$.
By D-OL-14 and tensor convolution, every genuine inertia has
actual moving rank at least $\beta e$.  Nontrivial factor
permutations have the stronger projective support $e/4$;
a nonidentity scalar has moving rank $e$.

Call an inertia good if its actual moving rank is at least
$\delta_*e$, and bad otherwise.  We record precisely what a
bad operator can be.  It cannot permute factors, since
$\delta_*<1/4$.  Write a factor-preserving operator as a
scalar times $\bigotimes_i\rho_i(g_i)$, using the complete
compatible-normalizer assertion in D-OL-14.  If any active
$g_i$ is not a negative reflection, tensor convolution gives
support greater than $\delta_*e$.  Hence every active factor
of a bad operator is a negative reflection.

Each such involution has two eigenspaces, with proportions
$p_i$ and $1-p_i$ lying in $[\beta,1-\beta]$.  For two active
factors, each of the two tensor eigenvalues has proportion at
most
\[
  \max_{p,q\in[\beta,1-\beta]}
       \{pq+(1-p)(1-q),\ p(1-q)+(1-p)q\}
       =(1-\beta)^2+\beta^2.
                                                               \tag{O3.9}
\]
The equality follows because each expression is affine in each
variable separately, so its maximum is at a corner of the
rectangle.  Therefore two active factors give support at least
$2\beta(1-\beta)e>e/4>\delta_*e$.  Additional convolution
cannot increase a largest eigenvalue proportion.  Every bad
inertia consequently has exactly one active factor and is a
scalar times one negative-reflection operator.  That scalar
is $1$ or $-1$: otherwise neither of its two eigenvalues is
one and the actual moving rank is $e$.  No general assertion
that scalar convolution preserves a nonreal-eigenvalue bound
is used here.

The exact arithmetic needed for the source budget is
\[
 13\beta-2>\frac17,\qquad
 11\beta+\delta_*-2>\frac1{60},\qquad
                  2\beta(1-\beta)>\frac14.                 \tag{O3.10}
\]
The actual intermediate-cover Euler formula is
$6=2e(h-1)+D$.  Since $e>360$, genus at least two is impossible.
Genus one has $D=6$, at most six genuine marks, and period-image
rank at most six.  Rank nine would require genus zero, at least
twelve genuine inertia points, and $D=2e+6$.

If their number is at least thirteen, then
$2e+6\ge13\beta e$, whence $e<42$ by (O3.10), a contradiction.
Thus there are exactly twelve.  If any one is good, then
$2e+6\ge(\delta_*+11\beta)e$, whence $e<360$, another
contradiction.  Therefore all twelve would be bad.

We now use actual generation, not only the preceding dimension
and rank inequalities.  The twelve genus-zero inertia elements
generate the faithful inducing image $J$.  Since every bad
inertia preserves each factor, so does all of $J$.  Fix one
normal factor $S_{d_i}$.  The induced action of $J$ on that
factor is generated by the inner automorphisms of at most
twelve negative reflections $g_1,\ldots,g_k$, with $k\le12$;
inactive generators induce the identity.  But $J$ contains
the normal factor itself, so this induced image contains
all of $\operatorname{Inn}(S_{d_i})$.  As the natural group
$S_{d_i}$ has trivial center, the elements $g_j$ generate
$S_{d_i}$.

Write $g_j=-r_{v_j}$ in its natural quadratic space, and put
\[
              U=\bigcap_{j=1}^k v_j^\perp,
                     \qquad \dim U\ge d_i-k>0.             \tag{O3.11}
\]
Every $g_j$ restricts to $-I$ on $U$.  The group they generate
therefore acts on this nonzero space through the scalar
group $\{I,-I\}$, and at least one generator has nontrivial
image.  This is a nontrivial homomorphism from the perfect
group $S_{d_i}$ to a group of order two, a contradiction.
There is at least one active generator because the induced
image contains the nontrivial factor.  The natural dimension
bound in (O3.11) is much larger than needed, but is the same
one already used for the exact character estimate.

This excludes the nonselfdual source without using a
similitude character, a determinant class, or equality of the
inducing character field with the original imaginary field.
For the final main-source consequence, transitivity preserves
selfduality, so all factor modules are selfdual or all are
nonselfdual.  The former case is excluded at $d\ge22$ by
D-OL-09.  In the latter case $q=3$ is excluded at $d\ge1135$
by this theorem, and every odd $q\ge5$ is excluded at that
threshold by D-OL-12.  Since $d$ is odd, (O3.8) follows.
\end{proof}

\paragraph{Certificate and scope.}
The exact supplement checks the rational constants, the
two-involution convolution corners, an actual higher-order
root of a negative reflection, and a twelve-reflection
common-perpendicular model.  The all-order classification,
compatible-normalizer reduction and generation contradiction
are the complete arguments above, not extrapolations from
these examples.  Ordinary modules are essential: no claim
about arbitrary noninflated spin characters is added here.

\section{Two-plane carriers and ordinary orthogonal groups in odd characteristic}
\label{sec:orthogonal-even-odd}

\paragraph{Status and boundary.}
This is a separate new author derivation.  The complete author, root and independent /root/sixfolds
proof, source and code readings have passed, with distinct exact
supplements. The result is promoted only in its stated ordinary scope.  It concerns
ordinary representations of the natural group $\Omega(V)$.  It makes
no assertion for representations of a spin cover which do not factor
through this natural group.  Earlier canonical sections are unchanged.

We use the exact constants of D-OL-14:
\[
 \beta=\frac16-\frac1{10000}=\frac{4997}{30000},\qquad
 \delta=\frac{1-(19/20)^9}{2},\qquad
 \frac16<\delta<\frac3{16}<\frac7{32}.
 \tag{OE.1}
\]
For a complex operator of degree $a$, write
$\nu(T)=a-\max_\lambda\dim\ker(T-\lambda I)$.

\subsection*{D-OL-16. A two-plane orthogonal carrier}

\begin{lemma}[D-OL-16]
Let $F=\mathbf F_q$, with $q$ odd, and let $V$ be a nondegenerate
quadratic space of even dimension $d\ge24$.  Put $S=\Omega(V)$.
Suppose that a linear similitude $H$ satisfies
\[
       H^2=cI,\qquad Q(Hv)=\mu Q(v),\qquad
       c\notin(F^*)^2.
       \tag{OE.2}
\]
Let $T$ be a finite-projective-order compatible normalizer of a
nontrivial ordinary irreducible complex $S$-module of degree $a$,
implementing conjugation by $H$.  Then
\[
                              \nu(T)\ge7a/32.
                              \tag{OE.3}
\]
No extension of this module to a larger orthogonal group is used.
\end{lemma}

\begin{proof}
Since $\mu^2=c^2$, write $\varepsilon=\mu/c\in\{1,-1\}$.
Let $E=F[t]/(t^2-c)$, and make $V$ an $E$-space by letting $t$
act as $H$.  Write $B$ for the polar symmetric form, so
$Q(v)=B(v,v)/2$.  The adjoint identity is $H^*=\varepsilon H$.
Define
\[
 b(v,w)=\frac12\bigl(B(v,w)+t^{-1}B(Hv,w)\bigr).
 \tag{OE.4}
\]
If $\varepsilon=1$, this is a nondegenerate symmetric $E$-bilinear
form.  If $\varepsilon=-1$, it is a nondegenerate $E/F$-Hermitian
form, linear in its first variable.  Indeed multiplication by $t$
in the first variable multiplies (OE.4) by $t$; in the second
variable it multiplies it by $\varepsilon t$.  Symmetry of $B$
gives the indicated symmetry or conjugate symmetry, and
$\operatorname{Tr}_{E/F}b=B$ proves nondegeneracy.

The finite-field symmetric classification and the finite-field
Hermitian hyperbolic-pair theorem supply a hyperbolic pair for $b$:
the $E$-dimension is at least twelve, while the respective
anisotropic dimensions are at most two and one
\cite[pp.~138--140; Lemma~10.2, Corollary~10.3 and the ensuing
decomposition, pp.~116--117]{OS:Taylor1992}.
Its two $E$-lines give opposite, totally singular, $H$-stable
$F$-planes $L,L'$.  Their cross pairing is nondegenerate.  Choose
dual $F$-bases and write
\[
 V=L\oplus W\oplus L',\quad
 Q(x,w,y)=x^{\mathsf t}y+\tfrac12w^{\mathsf t}Bw,
 \quad N=\dim W=d-4\ge20.
 \tag{OE.5}
\]
Here and below $B$ also denotes the nonsingular matrix on $W$.
In these coordinates
\[
 H=\operatorname{diag}(A,h,\mu A^{-\mathsf t}),\quad
 A^2=cI_2,\quad h^2=cI_N,\quad
 h^{\mathsf t}Bh=\mu B,\quad\det A=-c.
 \tag{OE.6}
\]

The unipotent radical fixing $L$ is the following actual subgroup:
\[
 u(X,Z)=
 \begin{pmatrix}
 I_2&X&Z-\tfrac12XB^{-1}X^{\mathsf t}\\
 0&I_N&-B^{-1}X^{\mathsf t}\\
 0&0&I_2
 \end{pmatrix},\qquad
 X\in M_{2,N}(F),\quad Z^{\mathsf t}=-Z.
 \tag{OE.7}
\]
Direct multiplication gives
\[
 u(X,Z)u(Y,R)=u\left(X+Y,
 Z+R+\tfrac12(YB^{-1}X^{\mathsf t}-XB^{-1}Y^{\mathsf t})\right).
 \tag{OE.8}
\]
Every such matrix is an isometry, and this is a group of exponent
$p=\operatorname{char}F$.  Thus its determinant and spinor norm
are trivial and it lies in $S$.  The alternating pairing on the
$2N$-dimensional $X$-space given by the $(1,2)$ entry in (OE.8)
is nondegenerate.  Consequently this group $Q_L$ is the Heisenberg
group of order $q^{2N+1}$, with center $X=0$.

We justify the fixed-space estimate for this larger radical.
Let $e_1,e_2$ be the chosen basis of $L$ and $f_1,f_2$ its dual.
The radicals $Q_L,Q_{L'}$ contain the two full line radicals on
$Fe_2\oplus W\oplus Ff_2$, with parameters ranging over $W$.
D-OL-06 therefore puts
$\Omega(Fe_2\oplus W\oplus Ff_2)$ in their generated subgroup.
It fixes $e_1,f_1$ and is transitive on the nonzero singular
vectors of its natural space.  The center of $Q_L$ contains
the line-radical elements $u_{e_1}(s e_2)$, for every $s\in F$.
Conjugating by that smaller $\Omega$ gives the parameters $s f_2$;
the rows of $X$ already give every parameter in $W$.  Additivity
of the line radical now gives its full parameter space
$W\oplus Fe_2\oplus Ff_2$.  The same argument applies to $f_1$.
These two full opposite line radicals generate $S$, again by
D-OL-06.  Hence
\[
                    \langle Q_L,Q_{L'}\rangle=S.
                    \tag{OE.9}
\]
The radicals are conjugate inside $S$: exchange the two
hyperbolic pairs, and, if necessary, correct the determinant and
spinor norm in their nondegenerate perpendicular space.  That
space represents both nonzero square classes, so its reflections
give both corrections.  An invariant vector for both radicals
would be $S$-invariant.  Irreducibility and nontriviality thus give
\[
                        t_0:=\dim\rho^{Q_L}\le a/2.
                        \tag{OE.10}
\]

Conjugating (OE.7) by (OE.6) gives exactly
\[
             X\longmapsto AXh^{-1},\qquad
             Z\longmapsto\mu^{-1}AZA^{\mathsf t}.
             \tag{OE.11}
\]
The action on the one-dimensional center is the scalar
$\det A/\mu=-c/\mu=-\varepsilon$.
The action on $X$ is an involution.  Its fixed matrices are
$E$-linear maps $W\to L$, and therefore form an $F$-space of
dimension $N$.  Dually, using the perfect trace pairing on
matrices, the dual weight action on $Y\in\operatorname{Hom}_F(L,W)$
is $Y\mapsto hYA^{-1}$.  Its fixed space also has dimension $N$.
Every nonzero fixed $Y$ has rank two, since its kernel is an
$A$-stable $F$-subspace of the irreducible two-dimensional
$F[A]$-module $L$.

We next bound the fixed fraction of every nonzero linear-weight
orbit, not just the fixed fraction of the whole dual space.
The Levi subgroup $\Omega(W)$ lies in $S$ and acts on $Y$ by
left multiplication.  For rank two, its orbits consist precisely
of the ordered independent pairs of vectors of $W$ with a fixed
Gram matrix.  Witt's theorem extends an isometry between their
spans, even when that span is degenerate
\cite[Theorem~7.4 and complete proof, pp.~57--59]{OS:Taylor1992}.
An extension in $O(W)$ can be corrected into $\Omega(W)$ while
fixing the pair: its perpendicular space contains a nondegenerate
subspace of dimension at least $N-4$, whose reflections realize
all determinant/spinor-norm pairs.  This proves the orbit assertion.

Here is a uniform lower bound for each such orbit.  There are
at least
\[
                    q^{N-1}-q^{\lceil N/2\rceil}-1
\]
choices for its nonzero first vector, by the quadratic Gauss
estimate of D-OL-06.  For each choice the prescribed cross
pairing cuts out an affine hyperplane.  The restricted polar
rank there is at least $N-2$.  Diagonalizing the quadratic
polynomial on this affine hyperplane gives the same Gauss bound
also in the presence of linear terms: a nonzero linear term on
the radical makes the sum zero, and otherwise the absolute sum
is bounded by the usual rank estimate.  Thus the number of
second vectors of the prescribed norm and cross pairing,
after removing at most $q$ dependent vectors, is at least
\[
                    q^{N-2}-q^{N/2}-q.
\]
For $q\ge3$, $N\ge20$, each displayed lower bound exceeds
half its leading term.  Every rank-two orbit therefore has
more than $q^{2N-3}/4$ vectors.  Its fixed fraction under
(OE.11) is at most
\[
                        4q^{-N+3}<1/8.
                        \tag{OE.12}
\]
Rank-one weights have no fixed vector at all.  Weight
multiplicities are constant on these Levi orbits.  Hence on
the sum of nontrivial linear-character spaces of $Q_L$, every
eigenspace of $T$ has codimension at least $7/16$ of that
sum's dimension: the nonfixed weights are paired, and a
paired block contributes at most half its dimension to one
eigenvalue.  This remains true with arbitrary scalar phases
on the weight spaces.

For clarity, the nonlinear Heisenberg characters are also
treated with their multiplicities.  The explicit construction
and matrix-algebra proof in D-HOS-01 show that, for every
nontrivial additive central character $\psi$, there is a unique
irreducible carrier $U_\psi$ of dimension $q^N$, and its Weyl
operators $L_X$ form a basis of $\operatorname{End}(U_\psi)$.
If $\varepsilon=1$, (OE.11) sends $\psi$ to $\psi^{-1}$.
They are distinct because the characteristic is odd.  The
corresponding isotypic spaces are paired, again giving support
at least one half of their combined dimension.

If $\varepsilon=-1$, the central character is fixed.  Let
$\mathcal A:X\mapsto AXh^{-1}$ be the symplectic involution
on the $X$-space.  Its fixed space has dimension $N$.  The
carrier intertwiner $\Xi$, normalized to satisfy $\Xi^2=I$,
obeys the exact identity
\[
 \Xi L_X\Xi^{-1}=L_{\mathcal A X},\qquad
 |\operatorname{Tr}\Xi|^2
       =\operatorname{Tr}(\operatorname{Ad}\Xi)
       =q^N.
 \tag{OE.13}
\]
The first equality has no extra phase, since (OE.11) preserves
the actual zero-central-coordinate section in (OE.7)--(OE.8).
Thus the support fraction of $\Xi$ is at least
$(1-q^{-N/2})/2>7/16$.  On an isotypic space the operator is
$\Xi\otimes C$ by the matrix-algebra commutant.  Eigenvalue
convolution preserves this support fraction; diagonalizability
follows by rescaling the given finite-projective-order $T$ to
finite order.  No multiplicity operator has been discarded.

The trivial $Q_L$-isotypic space has dimension $t_0$.  Combining
all the preceding summands and (OE.10) proves
$\nu(T)\ge(7/16)(a-t_0)\ge7a/32$, as required.
\end{proof}

\subsection*{D-OL-17. All compatible normalizers and the exceptional reflections}

\begin{theorem}[D-OL-17]
Let $q$ be odd, let $d\ge1136$ be even, and let $\rho$ be any
nontrivial ordinary irreducible complex representation of
$S=\Omega_d^\pm(q)$, of degree $a$.  If $q\ge5$, every nonscalar
compatible finite-projective-order normalizer satisfies
$\nu(T)>\delta a$.  If $q=3$, every such normalizer satisfies
$\nu(T)\ge\beta a$, and
\[
 \nu(T)\le\delta a\quad\Longrightarrow\quad
 \text{its natural implementing projective map is a reflection}.
 \tag{OE.14}
\]
In the last case $T$ itself has projective order two.
\end{theorem}

\begin{proof}
The complete natural semilinear-normalizer description,
including represented central quotients, is D-OL-07.  Thus a
normalizer is implemented by a semilinear similitude $H$.
First suppose that $T$ has prime projective order $l$.  The
centralizer and perfectness argument in D-OL-07 gives $H^l=cI$.
If its field part is nontrivial, the full mixed norm/form descent
of (ON.7)--(ON.9) gives support at least $3a/16>\delta a$.
This includes diagonal--field and graph--field components.
We may therefore suppose that $H$ is linear, with multiplier $\mu$.

If $l$ is odd, then $\mu^l=c^2$.  The scalar
\[
                    z=c/\mu^{(l-1)/2}
                    \quad\text{satisfies}\quad
                    z^2=\mu,\quad z^l=c.
                    \tag{OE.15}
\]
Consequently $g=z^{-1}H$ is an isometry of order $l$.
Its determinant and spinor norm, being of orders dividing both
$l$ and two, are trivial.  Thus $g\in S$, and Schur's lemma
makes $T$ a scalar times $\rho(g)$.

If $g$ fixes a hyperbolic pair, use its complementary line
radical as in (OO.3)--(OO.7).  Its natural moving rank there
is at least two, because an odd-characteristic rank-one
isometry is a reflection and has determinant minus one.
The all-rank Gauss estimate therefore gives a fixed-weight
fraction less than $1/8$, and support at least $7a/32$.
If no hyperbolic pair is fixed, a semisimple fixed space has
dimension at most two.  For a characteristic-order unipotent,
its fixed space has polar rank at least $d-2r$, where
$r=\operatorname{rank}(g-I)$; absence of a hyperbolic pair
therefore gives $r\ge(d-2)/2\ge567$.  The negative eigenspace
is zero in both cases.  In every remaining case the natural
support restricted to $\mu_2=\{1,-1\}$ is at least $567$.
Inflation to the ordinary spin cover and the exact published
LST estimate yield
\[
                    |\chi(g)|/a<(19/20)^9,
                    \qquad\nu(\rho(g))>\delta a.
                    \tag{OE.16}
\]
The source is \cite[Theorem~4.3.6, (4.3.2), pp.~1917--1918;
complete proof pp.~1918--1923; Remark~4.3.7, p.~1923]{HOS:LST2011},
with threshold $567=7\cdot9^2$.  The remark permits the
restricted scalar set $\mu_2$ for spin groups.  Ordinary
inflation is irreducible and no faithful-character assumption
is imposed.  This is the same checked source interface as
(OO.8)--(OO.10), with even $d\ge1136$.

It remains to treat $l=2$.  Now $H^2=cI$ and $\mu=\pm c$.
If $c$ is nonsquare, D-OL-16 applies.  Suppose that $c=z^2$
and replace $H$ by $z^{-1}H$, so $H^2=I$ and $\mu=\pm1$.
If $\mu=1$, the two eigenspaces are nondegenerate.  At least
one contains a hyperbolic pair.  If its eigenvalue is
$\epsilon$, the induced action on the complementary line
radical is $A=\epsilon H|_W$.  It is a nonidentity involution.
If its moving rank is at least two, the preceding $1/8$
bound gives $7a/32$.  If its moving rank is one, the original
projective map is a reflection.  The rank-one Gauss bound is
\[
 \frac{q^{-1}+q^{-9}}{1-q^{-9}-q^{-19}}.
 \tag{OE.17}
\]
For $q\ge5$ it is less than $1/4$, giving $3a/16$.
For $q=3$ it is less than $1-4\beta$, giving $\beta a$.

If instead $\mu=-1$, the two eigenspaces of $H$ are opposite
totally singular spaces of equal dimension.  Choose a
hyperbolic pair with its two vectors in the different
eigenspaces.  The induced action on the line radical is
$A=-H|_W$ after choosing the eigenvalue of the first vector
to be one.  This is an involutory anti-isometry.  After
identifying dual weights using the polar form, the dual
action is $-A$.  Both $A$ and $-A$ have totally singular
fixed spaces of dimension $N/2$, where $N=d-2$.
A nonzero fixed weight can therefore occur only
on the norm-zero Levi orbit.  The fraction fixed in that
orbit is at most
\[
 \frac{q^{N/2}}{q^{N-1}-q^{\lceil N/2\rceil}-1}<\frac18
 \quad(N\ge20,q\ge3).
 \tag{OE.18}
\]
All other norm orbits have no fixed vector.  Opposite line
radicals and paired weights give $7a/32$ once more.  The
sign and scaling in this argument follow directly by
conjugating the line-radical matrix by
$\operatorname{diag}(1,H|_W,-1)$; its parameter action is
$-H|_W$.  Thus no assertion of isometric action was used
for this anti-isometric case.

These arguments prove the asserted bounds for prime
projective order.  Taking a nonscalar prime-order power
never increases projective support, so the full bound
for $q\ge5$ and the lower bound $\beta a$ for $q=3$ follow.
We must prove the stronger all-order classification (OE.14).
For $q=3$ there is no field component.  If the projective
order of $T$ has an odd prime divisor, the corresponding
odd-prime power already has the stronger bound (OE.16).
It remains to consider order $2^s$.  Its involution power
must be a projective reflection if $\nu(T)\le\delta a$.
Suppose that $s\ge2$, and let $H$ be the natural linear lift.
Then
\[
                         H^{2^{s-1}}=\eta r_v
                         \quad(\eta\in F^*).
                         \tag{OE.19}
\]
The map $H$ commutes with this power, hence preserves the
nonsingular one-dimensional $(-1)$-eigenspace $Fv$ of $r_v$.
If $Hv=\lambda v$, its similitude multiplier is $\lambda^2$.
Rescale $H$ by $\lambda^{-1}$ to an isometry.  Its even
power has determinant one.  But any scalar multiple of
$r_v$ which is an isometry is $\pm r_v$, and in even
dimension both have determinant minus one.  This
contradicts (OE.19).  Therefore $s=1$, proving the entire
all-order statement, not merely a statement about one
prime-order power.
\end{proof}

\subsection*{D-OL-18. The actual source obstruction}

\begin{theorem}[D-OL-18]
A nonempty irreducibly acting normal central product of nontrivial
ordinary orthogonal factor modules with odd $q_i$ and even
$d_i\ge1136$ cannot supply period-image rank nine for its
actual six-dimensional multiplicity variation.  Mixed fields,
mixed dimensions, represented central quotients, arbitrary
scalars and all compatible factor permutations are allowed.
No selfduality, determinant or inducing-field equality is
assumed.

In particular, in the main transitive orthogonal-core setting,
every surviving ordinary factor has natural dimension at most
$1134$.  For odd natural dimension the sharper bound is
$1133$.  The separate spin-representation boundary is retained.
\end{theorem}

\begin{proof}
Let $e=\prod_i a_i$.  The degree bound
\cite[Theorem~3.7, pp.~165--167]{OL:GLT2024} gives
$a_i\ge q_i^{d_i-4}(q_i-1)>360$.
Every nonscalar factor-preserving operator has support at
least $\beta e$ by D-OL-17 and eigenvalue convolution.
A nontrivial factor permutation has support at least $e/4$;
a nonidentity scalar has full actual moving rank $e$.

Call an actual inertia good if its moving rank is at least
$\delta e$, and bad otherwise.  A bad inertia fixes every
factor, and every active factor must have $q_i=3$ and be
implemented by a projective reflection.  Its factor
operator has projective order two, by (OE.14).
If two or more factors were active, normalize two of these
operators to involutions.  Their two eigenvalue proportions
lie in $[\beta,1-\beta]$.  The maximum proportion after
convolving them is at most
$(1-\beta)^2+\beta^2$, attained at a corner of this square.
Consequently their support is at least
\[
                       2\beta(1-\beta)e>e/4>\delta e.
                       \tag{OE.20}
\]
Further tensor factors cannot increase the maximum
eigenvalue proportion.  A bad inertia therefore acts
nontrivially on exactly one factor, there by a projective
reflection.  Its global scalar must make one of the two
eigenvalues equal to one, or its actual moving rank is $e$.

The Euler identity for the actual intermediate cover is
$6=2e(h-1)+D$.  Since $e>360$, genus at least two is
impossible.  At genus one $D=6$, so there are at most six
genuine marks and the actual period-image rank is at most
six.  Rank nine therefore requires $h=0$, at least twelve
genuine inertia points, and $D=2e+6$.
The exact inequalities
\[
                  13\beta-2>1/7,
             \qquad11\beta+\delta-2>1/60
             \tag{OE.21}
\]
show first that thirteen marks would imply $e<42$, and
then that twelve marks with one good inertia would imply
$e<360$.  Thus exactly twelve inertias would all have to
be bad.

We now use their actual generation.  They all preserve
the factors.  Fix one normal factor.  The conjugation
image of the inducing group on it contains its entire
inner automorphism group, because the normal factor is
a subgroup of the inducing image.  This conjugation image
is generated by the projections of the twelve actual
inertias.  Each nontrivial projection is represented
naturally by a reflection, and the others are identity.
The natural projective model is faithful on the simple
quotient, by the centralizer/perfectness argument of
D-OL-07.  Hence at most twelve natural reflections
$r_{v_j}$ generate a projective subgroup containing
$S_i/Z(S_i)$.

The space
\[
                         U=\bigcap_j v_j^\perp
                         \quad\text{has dimension at least }d_i-12>0.
                         \tag{OE.22}
\]
Every one of the chosen natural reflections fixes $U$
pointwise.  Consequently every $s\in S_i$ acts on $U$
by a scalar: its projective image is a word in those
reflections.  Those scalars define a homomorphism
$S_i\to F_i^*$, which is trivial since $S_i$ is perfect
\cite[Theorems~11.47--11.48 and proofs, pp.~161--162]{OS:Taylor1992}.
This would make $U$ a nonzero fixed space for the natural
$S_i$-action.  There is no such space: two full opposite
line radicals have natural fixed spaces the two distinct
singular lines, and they generate $S_i$.
This contradiction excludes the source.  If a factor
has $q_i\ne3$, all twelve projections on it are already
trivial, which yields the same contradiction even sooner.

Finally, the selfdual ordinary cases are covered by
D-OL-09; the even-field ordinary cases by D-OL-10--11;
the odd-natural odd-field cases by D-OL-12--15; and the
remaining even-natural odd-field cases by the present
theorem.  Applied in the main transitive setting, these
give the stated ordinary natural-dimension bounds.
None of these implications extends to a spin module
which does not descend to the natural orthogonal group.
\end{proof}

\section{Semilinear Clifford lifts and the even-dimensional Spin boundary}\label{SF-SPIN:semilinear-clifford-lifts-and-the-even-dimensional-spin-boundary}

\subsection{Status and exact scope}\label{SF-SPIN:status-and-exact-scope}

New author derivation, D-OL19--21. The complete root canonical proof/code reading and the independent /root/g34\_eightfolds whole-proof/source audit have passed; distinct exact Clifford implementations are retained. The field has odd characteristic. The lifting lemma is valid in either parity in dimension at least 22; the final support and source theorems concern even natural dimension at least 1136. No source exclusion in odd natural dimension, characteristic two, or small-rank triality is made here.

We use the complete natural-group arguments D-OL06--09 and D-OL12--18, retaining their exact orbit estimates, finite-field descent, and actual-source hypotheses. The new points are: a semilinear lift acting on every represented central quotient of Spin; unique lifts of the odd-order radicals; and a degree estimate that applies to every nontrivial Spin module. The proof does not assume that a module descends to the natural Omega group or extends to a full orthogonal group.

Write

\[
 \beta=\frac{4997}{30000},\qquad
 \delta=\frac{1-(19/20)^9}{2},\qquad
 \nu(T)=a-\max_\lambda\dim\ker(T-\lambda I)
 \tag{SP.1}
\]

for an operator on a complex vector space of dimension \(a\).

\subsubsection{Exact published group inputs}\label{SF-SPIN:exact-published-group-inputs}

Let \(F=\mathbf F_q\), \(q\) odd, and let \((V,Q)\) be nondegenerate. The finite Spin group has its usual Clifford realization

\[
 S=\operatorname{Spin}(V)
 =\{g\in C^0(V,Q)^*:gVg^{-1}=V,\ \tau(g)g=1\},
 \qquad
 1\longrightarrow\{\pm1\}\longrightarrow S
 \mathop{\longrightarrow}^{\varphi}\Omega(V)\longrightarrow1,
 \tag{SP.2}
\]

where \(\tau\) reverses Clifford words. In the dimensions used here \(S\) is perfect and quasisimple. Precise published locators are Larsen--Shalev--Tiep, \emph{The Waring problem for finite simple groups}, Annals 174 (2011), proof of Theorem 4.3.6, Step 1, pp.1918--1919: the second exact sequence on p.1919 is (SP.2), and the following paragraph explicitly identifies \(K_i=\operatorname{Spin}(W_i)\) as quasisimple for \(\dim W_i\ge5\). \cite[proof of Theorem~4.3.6, Step~1, pp.~1918--1919]{HOS:LST2011} This is an invoked published structural input, not an inference that an arbitrary commutator subgroup is perfect. The Clifford definition and exact sequence also appear in Liebeck--O'Brien--Shalev--Tiep, \emph{Commutators in finite quasisimple groups}, Bulletin LMS 43 (2011), Section 4; the freshly read author version has that passage on pp.10--11 \cite[Section~4, author pp.~10--11]{SF:LOST2011}.

The simple quotient is \(\overline S=P\Omega(V)\). One can identify its kernel without using a table of Spin centers. If \(g\in S\) maps into \(Z(\Omega(V))\), then \([g,s]\in\{\pm1\}\) for every \(s\in S\). The map \(s\mapsto[g,s]\) is a homomorphism because its values are central, and perfectness makes it trivial. Thus

\[
 Z(S)=\varphi^{-1}(Z(\Omega(V))),\qquad
 S/Z(S)\simeq P\Omega(V).
 \tag{SP.3}
\]

The complete natural semilinear lifting argument D-OL07 applies to this simple quotient in dimension at least 22. It uses the published outer-order table of Li--Xia, Section 2.2.2, Table 2.1, \cite[Section~2.2.2, Table~2.1]{HES:LiXia2022}, and explicitly proves that its full outer image is realized by natural semilinear similitudes, with scalar kernel. We retain that input and its triality restriction.

\subsection{D-OL19. Every compatible represented Spin normalizer has a natural lift}\label{SF-SPIN:d-ol19.-every-compatible-represented-spin-normalizer-has-a-natural-lift}

\textbf{Lemma.} Suppose \(q\) is odd and \(d=\dim V\ge22\), with either parity. Let \(\rho:S\to\mathrm{GL}(M)\) be a nontrivial irreducible complex representation, with arbitrary central kernel. Every operator normalizing \(\rho(S)\) is compatible with a natural semilinear similitude of \((V,Q)\). The induced projective orders agree. This assertion includes modules that do not descend through \(\varphi\).

\textbf{Proof.} We first give the natural lift explicitly. Let \(H\) be \(\sigma\)-semilinear and have multiplier \(\mu\in F^*\), so

\[
 Q(Hv)=\mu\,\sigma(Q(v)).
 \tag{SP.4}
\]

On even tensor words define a \(\sigma\)-semilinear map by

\[
 \Phi_H(v_1\cdots v_{2k})
 =\mu^{-k}Hv_1\cdots Hv_{2k}.
 \tag{SP.5}
\]

Here a Clifford word has no well-defined integer degree, so well-definedness must be checked before using this formula. In the tensor algebra, an even word containing a relation \(v\otimes v-Q(v)\) maps to a scalar multiple of the corresponding relation \(Hv\otimes Hv-Q(Hv)\). Indeed the two word lengths differ by two, and the extra factor \(\mu^{-1}\) cancels \(Q(Hv)=\mu\sigma Q(v)\). Thus the defining ideal in the even tensor algebra is preserved. The inverse is the construction for \(H^{-1}\). Products of even words acquire exactly the product of their scaling factors, so (SP.5) is a semilinear algebra automorphism of \(C^0(V,Q)\).

This is also the construction and complete ideal calculation in Wonenburger, \emph{The Clifford Algebra and the Group of Similitudes}, Canadian J. Math. 14 (1962), Part II, pp.48--49. \cite[Part~II, pp.~48--49]{SF:Wonenburger1962} Its use here requires no extension to the entire Clifford algebra; Wonenburger's separate Theorem 3 imposes conditions on such an extension, and those conditions are not being assumed.

For the needed vector covariance, define on odd words a semilinear map

\[
 \Psi_H(v_1\cdots v_{2k+1})
 =\mu^{-k}Hv_1\cdots Hv_{2k+1}.
 \tag{SP.6}
\]

The same two-length calculation preserves the odd part of the defining ideal. For \(a,b\in C^0\) and \(z\in C^1\), counting lengths and then passing to the quotient gives

\[
 \Psi_H(azb)=\Phi_H(a)\Psi_H(z)\Phi_H(b),
 \qquad \Psi_H|_V=H.
 \tag{SP.7}
\]

The pair \((\Phi_H,\Psi_H)\) is an algebra automorphism on the even part and a compatible map on its odd bimodule. It is not claimed to be an automorphism of the whole Clifford algebra: the product of two odd images has the additional factor \(\mu\).

Word reversal commutes with \(\Phi_H\), because reversing a word changes neither its length nor its scaling factor. If \(g\in S\), then

\[
 \tau(\Phi_H(g))\Phi_H(g)=1,\qquad
 \Phi_H(g)H(v)\Phi_H(g)^{-1}=H(gvg^{-1}).
 \tag{SP.8}
\]

Thus \(\Phi_H(g)\) is again in \(S\), and the inverse construction proves that \(\alpha_H=\Phi_H|_S\) is a group automorphism. Its natural image satisfies

\[
 \varphi(\alpha_H(g))=H\varphi(g)H^{-1}.
 \tag{SP.9}
\]

If \(H_1,H_2\) have semilinear parts \(\sigma_1,\sigma_2\), their composite multiplier is \(\mu_1\sigma_1(\mu_2)\). Substitution in (SP.5) gives \(\alpha_{H_1H_2}=\alpha_{H_1}\alpha_{H_2}\). Replacing \(H\) by a nonzero linear scalar multiple changes \(\mu\) by the scalar's square and leaves (SP.5) unchanged. Hence this lift depends only on the projective semilinear similitude.

Put \(R=\rho(S)\). Quasisimplicity shows that \(\ker\rho\subset Z(S)\). Moreover \(Z(R)=\rho(Z(S))\): if a lift \(s\in S\) has central image in \(R\), all commutators \([s,g]\) lie in the central kernel and again form a homomorphism from the perfect group \(S\); it is trivial. Therefore \(R/Z(R)\simeq\overline S\).

Let \(T\) normalize \(R\), and write \(\gamma\) for its conjugation automorphism of \(R\). Choose a natural semilinear similitude \(H\) realizing the induced automorphism of \(\overline S\), using D-OL07. The two homomorphisms \(\gamma\circ\rho\) and \(\rho\circ\alpha_H\) have the same projection to \(\overline S\). Consequently

\[
 c(s)=\gamma(\rho(s))\,\rho(\alpha_H(s))^{-1}\in Z(R)
 \tag{SP.10}
\]

is a homomorphism \(S\to Z(R)\). For example, the multiplicativity follows by moving the central factor \(c(t)\) past the two images of \(s\). Perfectness implies \(c=1\). We obtain the exact compatibility

\[
 T\rho(s)T^{-1}=\rho(\alpha_H(s))\quad(s\in S).
 \tag{SP.11}
\]

In particular \(\alpha_H\) preserves \(\ker\rho\). Preservation of a chosen central kernel has been proved from the actual represented automorphism; it was not assumed and no classification of the different Spin centers was needed.

Finally, an automorphism of the perfect group \(R\) inducing the identity on \(R/Z(R)\) differs from the identity by a central homomorphism and is the identity. Schur's lemma then says that its implementing operator is scalar. On the natural side D-OL07 proves that a semilinear similitude inducing the identity on \(P\Omega(V)\) is a linear scalar. Applying these two observations to every power proves equality of the projective orders of \(T\) and \(H\). This proves the lemma. \(\square\)

\subsection{D-OL20. Radicals and all the support estimates transfer to Spin}\label{SF-SPIN:d-ol20.-radicals-and-all-the-support-estimates-transfer-to-spin}

\textbf{Lemma.} Let \(q\) be odd, \(d\ge1136\) even, and let \(M\) be any nontrivial irreducible complex Spin module of degree \(a\). A finite-projective-order nonscalar compatible normalizer satisfies

\[
 \nu(T)\ge\beta a.
 \tag{SP.12}
\]

If \(q\ge5\), it satisfies \(\nu(T)>\delta a\). If \(q=3\) and \(\nu(T)\le\delta a\), its natural projective lift is a reflection and has projective order two. Also

\[
 a\ge q^{d-3}-q^{\lceil(d-2)/2\rceil}-1>360.
 \tag{SP.13}
\]

\textbf{Proof: unique lifts, generation, and multiplicities.} Let \(U\) be any one- or two-plane unipotent radical used in D-OL06 or D-OL16. It is a group of odd prime-power order. The preimage \(\varphi^{-1}(U)\) is a central extension with kernel \(\{\pm1\}\). A Sylow subgroup for that odd prime maps isomorphically onto \(U\). Denote it by \(\widetilde U\). Every element of the preimage is a central element times an element of \(\widetilde U\), so \(\widetilde U\) is normal and is the unique such Sylow subgroup. This proves existence, uniqueness and functoriality of the lift without a splitting choice.

If \(U,U'\) are opposite radicals that generate \(\Omega(V)\), let \(L=\langle\widetilde U,\widetilde U'\rangle\). Its natural image is all of \(\Omega(V)\), so \(S=\{\pm1\}L\). The central kernel makes \(L\) normal and the quotient abelian. Perfectness of \(S\) gives \(L=S\). Natural conjugacies between these radicals lift through \(\varphi\) and therefore conjugate their unique lifts. On a nontrivial irreducible Spin module the two fixed spaces have equal dimension and zero intersection, giving the same bound \(t_0\le a/2\).

For a natural Levi subgroup \(\Omega(W)\), use its full inverse image under \(\varphi\). This suffices for all weight-orbit arguments: every natural Levi element has a lift, and conjugation on \(\widetilde U\) is exactly its prescribed action on \(U\). Weight multiplicities are constant on each orbit. No assertion that the represented module factors through this Levi quotient is made. Likewise (SP.9) and uniqueness imply that \(\alpha_H\) acts on each invariant lifted radical exactly by the natural matrix formulas. There is no additional central character or phase from the double cover, because the odd-order radical maps isomorphically onto its image.

It follows in particular that the full two-plane carrier proof D-OL16 applies to every Spin module: the same Heisenberg group, central character sectors, Weyl basis, fixed quotient dimension, and tensor multiplicity operator occur. For a natural nonsquare projective involution its bound is \(\nu(T)\ge7a/32\).

\textbf{Proof: the normalizer cases.} We give the transfer for every case, so that the ordinary boundary is removed only after checking each use of it. Replace \(T\) by a nonscalar prime-order projective power as in D-OL17; this cannot increase \(\nu\). By D-OL19 its natural lift has the same projective order.

If the field part is nontrivial, it has that prime order. The arbitrary-field norm and form descent in D-OL09 supplies the invariant hyperbolic pair and Frobenius action used in its line-radical count. The unique lift preserves exactly this action, the full inverse-image Levi supplies all orbits, and the bound is \(\nu(T)\ge3a/16>\delta a\).

For a linear odd-prime projective order \(\ell\), normalize the natural lift exactly as before. From \(H^\ell=cI\), with multiplier \(\mu\), one has \(\mu^\ell=c^2\). Set \(z=c/\mu^{(\ell-1)/2}\). Then \(z^2=\mu\), \(z^\ell=c\), and \(g=H/z\) is an odd-order isometry. Its determinant and spinor norm are trivial, so \(g\in\Omega(V)\). Choose \(\widetilde g\in S\) above it. Because \(g\) is an isometry, \(v\mapsto gv\) extends to an automorphism of the whole Clifford algebra whose even restriction is \(\Phi_g\). Conjugation by \(\widetilde g\) agrees with that automorphism on the generating vector space and hence everywhere. Scalar invariance of (SP.5) gives \(\alpha_H=\operatorname{Int}(\widetilde g)\), so \(T\) is a scalar multiple of \(\rho(\widetilde g)\).

When \(g\) fixes a hyperbolic pair, the complementary nontrivial natural action has rank at least two; the all-rank line-radical estimate yields \(7a/32\). When no such pair is fixed, the exact fixed-space argument in D-OL17 gives restricted support against \(\{1,-1\}\) at least \((d-2)/2\ge567\). LST Theorem 4.3.6 and Remark 4.3.7 apply directly to the irreducible Spin character of \(M\). The chosen lift may have an extra central factor; this changes neither its natural action nor the absolute character bound. Thus

\[
 \frac{|\operatorname{tr}\rho(\widetilde g)|}{a}
 <(19/20)^9,
 \qquad \nu(T)>\delta a.
 \tag{SP.14}
\]

For a linear projective involution with \(H^2=cI\), the nonsquare case was covered by the transferred two-plane proof. If \(c\) is square, normalize \(H^2=I\). For multiplier one, the eigenspaces are nondegenerate and one contains a hyperbolic pair. The same rank-at-least-two count applies unless the projective natural action is a reflection. In that last case the exact rank-one fraction from D-OL17 is less than \(1/4\) for \(q\ge5\), or less than \(1-4\beta\) for \(q=3\), giving respectively \(3a/16\) or \(\beta a\). For multiplier minus one, the opposite maximal singular eigenspaces give the same line-radical action and polar-dual action as in D-OL17. The fixed weights have norm zero and their orbit fraction is less than \(1/8\), giving \(7a/32\). Every action just used is the action on the unique lifted radical, so no ordinary-module hypothesis remains.

This proves the universal bound and the stronger bound for \(q\ge5\). Over \(F_3\), suppose \(\nu(T)\le\delta a\). The projective order has no odd prime factor, because a corresponding prime-order power would contradict (SP.14). Its involution power is a natural projective reflection. If its projective order were \(2^s\) with \(s\ge2\), the natural lift \(H\) would preserve the reflection's nonsingular line. If \(H\) acts there by \(\lambda\), its multiplier is \(\lambda^2\), so \(H/\lambda\) is an isometry. Its even power is a scalar isometry projectively equal to a reflection. In even dimension such an isometry is \(r\) or \(-r\), both of determinant minus one, whereas an even power has determinant one. This contradiction proves the order-two reflection conclusion. The parity restriction is essential at this step and is retained.

\textbf{Proof: the degree bound.} Use one full line radical \(U\), whose natural parameter space has dimension \(N=d-2\). If every weight of \(\widetilde U\) on \(M\) were trivial, its opposite conjugate would also act trivially. Their generation of \(S\) would make \(M\) trivial. Thus a nonzero weight occurs. The full inverse image of the Levi \(\Omega(W)\) moves it through its entire natural nonzero vector orbit. The exact quadratic-vector count and transitivity used in D-OL06 give an orbit of size at least \(q^{N-1}-q^{\lceil N/2\rceil}-1\). Every such weight occurs with positive multiplicity, proving (SP.13). This argument applies to faithful Spin modules and uses no ordinary GLT degree theorem. \(\square\)

\subsection{D-OL21. The even-dimensional odd-characteristic Spin source is excluded}\label{SF-SPIN:d-ol21.-the-even-dimensional-odd-characteristic-spin-source-is-excluded}

\textbf{Theorem.} In the actual inducing-representation and base-cover setting of D-OL18, allow each normal factor to be any represented central quotient of an odd-characteristic Spin group of even natural dimension at least 1136, and allow every nontrivial irreducible complex module of that factor. Then there is no rank-nine sixfold source in this case. Ordinary descent to the natural Omega factors is not required.

\textbf{Proof.} Let \(e=\prod_i a_i\) be the inducing degree. The new degree estimate gives \(a_i>360\) and hence \(e>360\). The same spectral convolution and factor-permutation arguments of D-OL18 apply to these actual modules. Scalar inertia operators different from one move the entire module. Call a nonidentity actual inertia good if its moving rank is at least \(\delta e\), and bad otherwise.

A bad inertia is factor preserving. Each active factor has field \(F_3\), and its natural projective lift is a reflection of order two by D-OL20. On a normalized active factor the two eigenspace proportions lie in \([\beta,1-\beta]\). Two active factors already have projective support at least

\[
 \bigl(1-(1-\beta)^2-\beta^2\bigr)e
 =2\beta(1-\beta)e>\tfrac14e>\delta e.
 \tag{SP.15}
\]

Thus a bad inertia has exactly one active factor. Its actual scalar must make an eigenvalue equal to one; otherwise it would move all \(e\) dimensions. In every case its actual moving rank is at least \(\beta e\).

The unchanged actual Euler identity is

\[
 6=2e(h-1)+D.
 \tag{SP.16}
\]

Because \(e>360\), it gives \(h\le1\). For \(h=1\), \(D=6\) allows at most six marked points and hence period rank at most six. A rank-nine case must therefore have \(h=0\), at least twelve marked points, and \(D=2e+6\). The exact constants satisfy

\[
 13\beta-2>\tfrac17,\qquad
 11\beta+\delta-2>\tfrac1{60}.
 \tag{SP.17}
\]

Thirteen marked points would imply \(e<42\), while twelve with even one good inertia would imply \(e<360\). Consequently the only putative case has exactly twelve bad inertias.

Those actual inertia elements generate the actual inducing group on the sphere. All are factor preserving. Fix a normal factor. Its induced projection contains the full simple inner group \(P\Omega(V_i)\); this follows from the containment of the represented factor in the inducing group and does not require a choice of central splitting. If \(q_i\ne3\), all twelve bad inertias project trivially to it, an immediate contradiction. If \(q_i=3\), their projections are identities or natural projective reflections. Choose at most twelve natural reflections representing them, with nonsingular normals \(v_j\). Their common fixed space

\[
 U=\bigcap_j v_j^\perp,
 \qquad \dim U\ge d_i-12>0
 \tag{SP.18}
\]

is fixed pointwise by every word in those reflections. Every element of the natural group \(\Omega(V_i)\) is projectively such a word, so it acts by an original-field scalar on \(U\). These scalars form a homomorphism \(\Omega(V_i)\to F_3^*\); perfectness makes it trivial. This would make \(U\) a nonzero fixed space of \(\Omega(V_i)\), contradicting its opposite full line radicals and their generation. This argument uses the natural simple projection only, so possible Spin centers and represented central kernels do not create an exception. The theorem follows. \(\square\)

\subsection{Source and remaining-scope record}\label{SF-SPIN:source-and-remaining-scope-record}

The new Clifford construction uses Wonenburger's published Part II pp.48--49; its tensor-ideal proof is reproduced above. Its odd bimodule extension, norm preservation, and represented-quotient compatibility are proved here. The full Clifford-algebra extension theorem is not used. Finite Spin perfectness/quasisimplicity and its double-cover sequence are explicit published inputs at LST pp.1918--1919. The natural full-automorphism statement is the already completely reviewed D-OL07, including its precise Li--Xia source status. The character bound is the published LST Theorem 4.3.6 with Remark 4.3.7, already read in full.

No central quotient is presumed to be fixed by every natural automorphism. Only an automorphism actually compatible with the represented quotient is used, and (SP.10) proves its required kernel preservation. No Schur-multiplier classification is invoked. The theorem removes the genuine-Spin boundary only for odd \(q\), even \(d\ge1136\). Odd natural dimension and characteristic two remain separate.

\section{The odd-dimensional genuine Spin boundary}
\label{sec:spin-odd-extension}

\paragraph{Status and exact scope.}
This is a separate new author derivation, D-OL-22--24.
The complete root contributor reading and the separate
/root/g34\_eightfolds whole-proof/source/code audit have passed.
Root supplied the order-four strategy; independent origin is not claimed
for that contribution. Distinct exact carrier reconstructions are retained.  The field is odd,
the natural dimension is odd, and the module is any nontrivial
irreducible complex module of the finite Spin group, with an
arbitrary central kernel.  Descent to the natural $\Omega$ group,
extension to a full orthogonal group, selfduality and a determinant
condition are not assumed.  Earlier canonical statements are unchanged.

Let $F=\mathbf F_q$, $q$ odd, and let $(V,Q)$ be nondegenerate.
Put $S=\operatorname{Spin}(V)$, and denote its natural map by
$\varphi:S\longrightarrow\Omega(V)$.  We use
\[
 \beta=\frac{4997}{30000},\qquad
 \delta=\frac{1-(19/20)^9}{2},\qquad
 \frac16<\delta<\frac3{16}<\frac7{32},
 \qquad \nu(T)=a-\max_\lambda\dim\ker(T-\lambda I).
 \tag{OSP.1}
\]

\paragraph{The precise lift and radical interface.}
We use D-OL-19 and the parity-free opening arguments of D-OL-20.
Their full proofs give the following interface in either parity
for $d=\dim V\ge22$.  Every compatible normalizer of a represented
central quotient of $S$ is implemented by a natural semilinear
similitude $H$, and the two projective orders agree.
For clarity the lift on even Clifford words is
\[
 \Phi_H(v_1\cdots v_{2k})
       =\mu^{-k}Hv_1\cdots Hv_{2k},
 \qquad Q(Hv)=\mu\,\sigma Q(v).                              \tag{OSP.2}
\]
The two possible tensor lengths in a Clifford relation differ by
two, and the extra $\mu^{-1}$ cancels the multiplier of $Q$.
Thus this formula is well defined on the even Clifford algebra.
The corresponding odd-word map, scaled by $\mu^{-k}$ on length
$2k+1$, intertwines its two even-module actions and restricts to
$H$ on $V$.  Reversal commutes with (OSP.2), so the Spin norm
and vector covariance are preserved.  This is the explicit
semilinear Clifford construction of D-OL-19; its even-algebra
calculation is the published construction
\cite[Part II, pp.~48--49]{SF:Wonenburger1962}.
No theorem asserting an unrestricted extension to the whole
Clifford algebra is used.

We recall the represented-kernel point.  The published exact sequence
and quasisimplicity input are
\[
 1\longrightarrow\{\pm1\}\longrightarrow S
       \mathop{\longrightarrow}^{\varphi}\Omega(V)
       \longrightarrow1,                                  \tag{OSP.3}
\]
with $S$ perfect in these dimensions
\cite[proof of Theorem~4.3.6, Step~1, pp.~1918--1919]{HOS:LST2011}.
If two homomorphisms from $S$ to a represented quotient induce the
same simple-quotient automorphism, their ratio is a homomorphism
to the center; perfectness makes it trivial.  Applying this to
the actual normalizing operator and (OSP.2) proves preservation
of its actual central kernel, rather than assuming it.
The full natural semilinear image, including the scalar
centralizer and arbitrary field components, is D-OL-07 and its
complete outer-order comparison
\cite[Section~2.2.2, Table~2.1, p.~14]{HES:LiXia2022}.

Every odd-order one- or two-plane radical $U\subset\Omega(V)$
has a unique lift $\widetilde U\subset S$ isomorphic to it.
Indeed a Sylow $p$ subgroup of $\varphi^{-1}(U)$ maps
isomorphically to $U$, and the preimage is its product with
the central kernel $\{\pm1\}$.  That Sylow subgroup is
therefore normal and unique.  All natural conjugation formulas
on $U$ consequently lift exactly, without a central phase.
If opposite radicals generate $\Omega(V)$, their lifts
generate $S$: their generated subgroup maps onto $\Omega(V)$,
so its quotient in $S$ is central and abelian, and perfectness
eliminates the quotient.  Opposite conjugate radicals therefore
have equal fixed-space dimensions with zero intersection on a
nontrivial irreducible $S$-module.  Each fixed dimension is at
most $a/2$.  Using the full inverse image of a natural Levi
group gives all its natural weight orbits and the constancy of
weight multiplicities.  This does not assert that the module
factors through that Levi quotient.

\subsection*{D-OL-22. The fourth-order negative-reflection carrier}

\begin{lemma}[D-OL-22]
Let $q=3$, let $d\ge25$ be odd, and suppose
$g\in SO(V)$ satisfies $g^2=-r_v$, where $v$ is nonsingular
and $r_v$ is its reflection.  If a finite-projective-order
normalizer $T$ of a nontrivial irreducible $S$-module implements
this natural conjugation, then
\[
                              \nu(T)\ge7a/32.               \tag{OSP.4}
\]
The assertion includes $g\notin\Omega(V)$.
\end{lemma}

\begin{proof}
The element $g$ preserves $Fv$ and $W_0=v^\perp$.
On $W_0$ it satisfies $g^2=-I$.  Since $x^2+1$ is irreducible
over $\mathbf F_3$, its characteristic polynomial there is
$(x^2+1)^{(d-1)/2}$, with determinant one.
The determinant-one condition on $g$ therefore makes its
eigenvalue on $Fv$ equal to one.

Put $E=F[t]/(t^2+1)$ and make $W_0$ an $E$-space by
letting $t$ act as $g$.  If $B$ is the polar form, then
$g^*=-g$ on $W_0$, and
\[
 b(x,y)=\frac12\bigl(B(x,y)+t^{-1}B(gx,y)\bigr)              \tag{OSP.5}
\]
is a nondegenerate $E/F$-Hermitian form with trace $B$.
Multiplication by $t$ in its first argument multiplies it by
$t$, and in its second by $-t$; symmetry of $B$ gives the
Hermitian identity.  The finite-field Hermitian hyperbolic-pair
theorem applies since $\dim_EW_0=(d-1)/2\ge12$
\cite[Lemma~10.2, Corollary~10.3 and ensuing decomposition,
pp.~116--117]{OS:Taylor1992}.
Its opposite $E$-lines give opposite $g$-stable totally
singular $F$-planes $L,L'$.

Choose dual bases and write
\[
 \begin{gathered}
 V=L\oplus W\oplus L',\qquad N=\dim W=d-4\ge21,\\
 Q(x,w,y)=x^{\mathsf t}y+\tfrac12w^{\mathsf t}Bw,\qquad
 g=\operatorname{diag}(A,h,A^{-\mathsf t}),\\
 A^2=-I_2,\quad \det A=1,\quad
 W=Fv\perp W_1,\quad h|_{Fv}=1,\quad h^2|_{W_1}=-I.
 \end{gathered}                                             \tag{OSP.6}
\]
The complementary dimension $N$ is odd; no even-dimensional
assumption on $W$ is being imported.
The natural two-plane radical consists of
\[
 u(X,Z)=
 \begin{pmatrix}
 I_2&X&Z-\frac12XB^{-1}X^{\mathsf t}\\
 0&I_N&-B^{-1}X^{\mathsf t}\\
 0&0&I_2
 \end{pmatrix},
 \quad X\in M_{2,N}(F),\quad Z^{\mathsf t}=-Z.               \tag{OSP.7}
\]
Its group law has central correction
$\frac12(YB^{-1}X^{\mathsf t}-XB^{-1}Y^{\mathsf t})$.
This is the same Heisenberg group of order $q^{2N+1}$
as in D-OL-16.  The two opposite such radicals generate
$\Omega(V)$ by its complete line-radical generation argument,
which uses only a sufficiently large nondegenerate complement
and applies in either parity.  The preceding lift interface
therefore gives
\[
                     t_0:=\dim M^{\widetilde Q_L}\le a/2.   \tag{OSP.8}
\]

Conjugation by $g$ on (OSP.7) is exactly
\[
 \mathcal A:X\longmapsto AXh^{-1},\qquad
 Z\longmapsto AZA^{\mathsf t}=Z.                            \tag{OSP.9}
\]
The center is fixed.  Decompose
$\operatorname{Hom}_F(W,L)$ according to $W=Fv\oplus W_1$.
On the first summand, of dimension two, the action is $A$,
with square $-I$ and no fixed vector.  On the second summand
its square is the identity.  The fixed maps there are
$E$-linear maps $W_1\to L$, of $F$-dimension $N-1$.
It follows that the exact fixed dimensions of
$\mathcal A,\mathcal A^2,\mathcal A^3$ on the $2N$-dimensional
quotient are
\[
                              N-1,\quad 2N-2,\quad N-1.    \tag{OSP.10}
\]

The dual weight action is $Y\mapsto hYA^{-1}$ on
$\operatorname{Hom}_F(L,W)$.  A fixed map has image in $W_1$
and is $E$-linear, so the fixed space has dimension $N-1$.
Every nonzero fixed map has rank two: its kernel would otherwise
be a proper nonzero invariant subspace of the irreducible
two-dimensional $F[A]$-module $L$.

The rank-two Levi orbits are the ordered independent pairs
with prescribed Gram matrix, by Witt extension and the
determinant/spinor corrections in their large perpendicular
space.  The complete estimate of D-OL-16 applies for odd $N$
as well.  Explicitly, the first-vector and affine second-vector
counts, after excluding dependent vectors, are bounded below by
\[
 q^{N-1}-q^{\lceil N/2\rceil}-1,\qquad
 q^{N-2}-q^{N/2}-q.                                       \tag{OSP.11}
\]
The affine count uses polar rank at least $N-2$; if a
linear term is nonzero on the radical, its Gauss sum is zero.
Thus it also covers degenerate prescribed Gram matrices.
Each bound exceeds half its leading term for $q=3$, $N\ge21$.
Every rank-two orbit consequently has more than $q^{2N-3}/4$
vectors, whereas there are at most $q^{N-1}$ fixed dual maps.
Its fixed proportion is less than
\[
                              4q^{-N+2}<1/8.               \tag{OSP.12}
\]
Rank-one orbits have no fixed map.  Weight multiplicities are
constant on these full-preimage Levi orbits.
The action in (OSP.9) has order four, so every nonfixed
weight orbit has length two or four.  Each such orbit contributes
at most half its isotypic dimension to a specified eigenspace.
Thus on the sum of all nontrivial linear-character spaces of
$\widetilde Q_L$ the support is at least $7/16$ of that sum.
It is not necessary that an individual Levi orbit be invariant
under $\mathcal A$; the fixed-weight bound and weight
multiplicities give the same bound after summing all of them.

It remains to retain the nonlinear Heisenberg sectors and
their multiplicities.  For each nontrivial additive central
character $\psi$, D-HOS-01 gives its unique carrier $U_\psi$,
of degree $q^N$, with Weyl operators $L_X$ forming a basis
of $\operatorname{End}(U_\psi)$.  The center is fixed, so
choose a unitary carrier intertwiner $\Xi$ and rescale it to
$\Xi^4=I$.  The exact zero-central-coordinate action (OSP.9)
gives, for $j=1,2,3$,
\[
 \Xi^jL_X\Xi^{-j}=L_{\mathcal A^jX},\qquad
 |\operatorname{tr}\Xi^j|^2
   =\operatorname{tr}(\operatorname{Ad}\Xi^j)
   =|\ker(\mathcal A^j-I)|.                                \tag{OSP.13}
\]
There is no hidden central phase, by uniqueness of the
odd-order radical lift.  Equations (OSP.10)--(OSP.13) give
the trace magnitudes
$q^{(N-1)/2},q^{N-1},q^{(N-1)/2}$.
The four-term Fourier formula therefore bounds every
eigenspace of $\Xi$ by
\[
 \frac{q^N+q^{N-1}+2q^{(N-1)/2}}4
       <\frac9{16}q^N.                                    \tag{OSP.14}
\]
For $q=3$, $N\ge21$ the strict inequality follows directly
after division by $q^N$; the last small term is at most
$2\cdot3^{-11}$.  Thus its support fraction exceeds $7/16$.
On the full $\psi$-isotypic space, the actual operator is
$\Xi\otimes C$ by the matrix-algebra commutant.  Rescaling the
given finite-projective-order $T$ to finite order makes the
multiplicity operator diagonalizable.  Eigenvalue convolution
preserves the carrier support bound for every scalar phase.

Only the trivial-radical space of dimension $t_0$ remains
uncontrolled.  Combining all sectors with (OSP.8) proves
$\nu(T)\ge(7/16)(a-t_0)\ge7a/32$.  This proves the lemma.
\end{proof}

\subsection*{D-OL-23. All projective orders and all compatible normalizers}

\begin{theorem}[D-OL-23]
Let $q$ be odd, let $d\ge1135$ be odd, and let $M$ be any
nontrivial irreducible complex $S=\operatorname{Spin}_d(q)$-module
of degree $a$.  Every nonscalar compatible finite-projective-order
normalizer $T$ satisfies $\nu(T)\ge\beta a$.
For $q\ge5$ it satisfies $\nu(T)>\delta a$.
For $q=3$, an operator with $\nu(T)\le\delta a$ has natural
projective lift a negative reflection, and projective order two.
Moreover
\[
                   a\ge q^{d-3}-q^{\lceil(d-2)/2\rceil}-1>360.
 \tag{OSP.15}
\]
\end{theorem}

\begin{proof}
The degree estimate is the parity-free argument of D-OL-20,
which we recall to avoid an ordinary-character hypothesis.
Some nontrivial weight of a lifted full line radical occurs;
otherwise its conjugate opposite would also act trivially,
and their generation of $S$ would make $M$ trivial.
The full inverse-image Levi moves that weight through a
nonzero natural norm orbit of dimension $d-2$.  Its size is
at least the first expression in (OSP.15), by the same exact
quadratic Gauss count as in D-OL-06.  Every weight in the orbit
occurs with positive multiplicity.  This proves (OSP.15) for
every Spin character, without applying an ordinary $\Omega$
degree theorem.

Take a nonscalar power of $T$ having prime projective order
$\ell$; this cannot increase $\nu$.  The natural lift supplied
by D-OL-19 has the same projective order.  If its field component
is nontrivial, $q=q_0^\ell$.  The complete norm and quadratic-form
descent (ON.7)--(ON.9) rescales this possibly mixed linear-field
lift to order $\ell$ with an invariant hyperbolic pair over
$\mathbf F_{q_0}$.  Its action on the associated line-radical
weights is Frobenius, with fixed proportion less than $1/4$.
The unique lifted radical and full inverse-image Levi retain
this action exactly.  The resulting support is at least
$3a/16>\delta a$.  No pure-field assumption is made.

Suppose the prime-projective-order lift is linear.
In odd natural dimension a similitude multiplier is a square,
since $(\det H)^2=\mu^d$.  Rescale it to an isometry and,
if necessary, multiply by $-I$ to obtain the unique
representative $g\in SO(V)$.  The only natural scalar in
$SO(V)$ is $I$.  Thus the order of $g$ is exactly the
projective order of $T$.

If $\ell$ is odd, $g$ has trivial spinor norm and belongs to
$\Omega(V)$, since $SO(V)/\Omega(V)$ has order two.
Choose a Spin lift $\widetilde g$.  Its conjugation on the
full Clifford algebra implements $g$ on $V$, hence agrees
on the even algebra with (OSP.2).  Schur's lemma makes the
normalizing operator a scalar multiple of $\rho(\widetilde g)$.
If $g$ fixes a hyperbolic pair, the complementary natural
moving rank is at least two: a rank-one nontrivial isometry
is a determinant-minus-one reflection.  The all-rank line
Gauss count (ON.3), with the lifted radical, gives $7a/32$.

If no hyperbolic pair is fixed, the natural support restricted
to $\{1,-1\}$ is at least $567$.  For semisimple odd prime
order, the fixed space is nondegenerate of dimension at most
two and $-1$ is not an eigenvalue.  For characteristic prime
order, if $r=\operatorname{rank}(g-I)$, the fixed polar rank
is at least $d-2r$.  Absence of a hyperbolic pair forces it
to be at most two, so the integer $r\ge(d-2)/2$ is at least
$567$.  Again $-1$ is not an eigenvalue.
Apply the published estimate directly to the nontrivial
irreducible Spin character:
\[
 |\chi(\widetilde g)|/a<(19/20)^9,\qquad
                     \nu(\rho(\widetilde g))>\delta a.      \tag{OSP.16}
\]
The exact input is
\cite[Theorem~4.3.6, (4.3.2), pp.~1917--1918;
complete proof pp.~1918--1923; Remark~4.3.7, p.~1923]{HOS:LST2011}.
Its restricted scalar set for Spin is $\mu_2$ and its
threshold is $567=7\cdot9^2$.  A central choice of lift changes
neither the natural action nor the absolute character bound.
The passage to support uses
$|\chi|\ge a-2\nu$, after rotating a largest-eigenspace
eigenvalue to one.  No faithful-character hypothesis or
descent to $\Omega$ is imposed.

If $\ell=2$, the normalized $g\in SO(V)$ is an involution.
Its two eigenspaces are nondegenerate and at least one
contains a hyperbolic pair.  Choose such a pair in the
$\varepsilon$-eigenspace.  The line-radical action on its
complement $W$ is the isometry $A=\varepsilon g|_W$,
exactly as in (OO.5), and it is nonidentity.
For moving rank at least two the bound is $7a/32$.
Moving rank one means that $A$ is a reflection and
$g$ is projectively a reflection.  The unique $SO(V)$
representative in odd dimension is $g=-r_v$.
The precise rank-one Gauss fraction is bounded by
\[
 c_q=\frac{q^{-1}+q^{-9}}{1-q^{-9}-q^{-19}}.
 \tag{OSP.17}
\]
For $q\ge5$, $c_q<1/4$; for $q=3$, $c_3<1-4\beta$.
The lifted-radical fixed-space bound and the two-element
weight orbits give support at least $(1-c_q)a/4$.
This is greater than $3a/16$ in the first case and
greater than $\beta a$ in the second.
These arguments treat every prime projective order and
give the claimed universal bounds.

It remains to classify every low-support order when $q=3$.
There is no field automorphism.  Let $g\in SO(V)$ be the
unique natural representative of the original projective
normalizer; its order equals that projective order.
An odd prime divisor would yield a strong-bound power by
the preceding analysis, so the order is $2^s$.
Its involution power must be $-r_v$ if the support is at
most $\delta a$.  If $s=1$, the claimed conclusion holds.

Suppose $s\ge3$.  Then $g^2\in\Omega(V)$ and
\[
                 (g^2)^{2^{s-2}}=-r_v .
 \tag{OSP.18}
\]
The exponent is even.  Any vector with eigenvalue $1$
or $-1$ for $g^2$ is therefore fixed by $-r_v$, hence
lies in $Fv$.  Both restricted eigenspaces have dimension
at most one.  The natural $\mu_2$-support of $g^2$ is
at least $d-1\ge567$.  The operator $T^2$ is a scalar
multiple of a Spin operator above $g^2$, by the exact
inner compatibility already used.  The same published
bound (OSP.16) applies to this element without a
prime-order condition.  It gives
$\nu(T)\ge\nu(T^2)>\delta a$, a contradiction.

The remaining case is $s=2$, so $g^2=-r_v$.
D-OL-22 applies whether $g$ belongs to $\Omega(V)$
or its other $SO$ coset, and gives
$\nu(T)\ge7a/32>\delta a$, another contradiction.
Thus $s=1$ and the natural projective lift is a negative
reflection.  The actual complex operator may have a
nontrivial central scalar in its square; its \emph{projective}
order, and its two normalized eigenspaces, are exactly two.
This distinction is retained in the source argument below.
\end{proof}

\subsection*{D-OL-24. The odd-dimensional Spin source contradiction}

\begin{theorem}[D-OL-24]
Allow a nonempty irreducibly acting normal central product of
represented Spin factors of odd natural dimensions $d_i\ge1135$
over arbitrary odd fields.  Give each factor any nontrivial
irreducible complex module, with arbitrary central kernel.
All compatible scalars, outer automorphisms and factor
permutations are allowed.  Such an actual inducing image
cannot supply period-image rank nine for its six-dimensional
multiplicity variation.  Transitivity, selfduality and a
determinant hypothesis are not required for this exclusion.
\end{theorem}

\begin{proof}
Write $e=\prod_i a_i$.  The degree bound (OSP.15)
gives $e>360$.  Convolution retains the support fraction
of every active factor-preserving operator.
Nontrivial factor cycles have projective support at least
$e/4$ by the complete weighted-word argument used in
D-OL-18.  A nonidentity scalar has actual moving rank $e$.
Every genuine inertia therefore moves at least $\beta e$.

Call an inertia good if its actual moving rank is at
least $\delta e$ and bad otherwise.  A bad operator is
factor preserving.  Every active factor must have $q_i=3$
and natural projective lift a negative reflection, by
D-OL-23.  Normalize each active factor's scalar so that
its two eigenvalues are $1,-1$.  Their proportions lie
in $[\beta,1-\beta]$, regardless of the original Spin
center or of the unnormalized square of the operator.
For two active factors the largest tensor-eigenspace
proportion is at most $(1-\beta)^2+\beta^2$:
both possible convolution expressions are affine in each
proportion and are maximized at a corner of that rectangle.
Thus two active factors already have support at least
\[
                   2\beta(1-\beta)e>e/4>\delta e.            \tag{OSP.19}
\]
Additional factors cannot increase the largest
eigenvalue proportion.  Every bad inertia consequently
has exactly one active factor.  Its actual scalar has to
select a fixed eigenvalue; otherwise it moves all $e$
dimensions.  No assertion about scalar preservation of
a nonreal-eigenvalue estimate is used.

The actual intermediate-cover Euler identity is
$6=2e(h-1)+D$.  Since $e>360$, it excludes $h\ge2$.
For $h=1$ there are at most six genuine marks and the
finite-cover period image has rank at most six.
The proposed rank-nine case must therefore have a sphere
base, at least twelve genuine marks, a generating
product-one inertia tuple and $D=2e+6$.
The exact inequalities
\[
 13\beta-2>\frac17,\qquad
 11\beta+\delta-2>\frac1{60},\qquad
 2\beta(1-\beta)>\frac14                                  \tag{OSP.20}
\]
exclude thirteen marks and exclude even one good inertia
among twelve.  Indeed they would force respectively
$e<42$ or $e<360$.  Thus exactly twelve bad inertias remain.

These actual inertia elements generate the inducing image.
All preserve every factor.  Fix one normal factor; the
induced projection contains its full simple inner group
$\Omega(V_i)$, because the represented Spin factor is
contained in the inducing group and
$\Omega(V_i)$ has trivial natural center in odd dimension.
If $q_i\ne3$, all bad inertias have trivial projection
there, an immediate contradiction.  For $q_i=3$ their
projections are identities or projective reflections.
Choose at most twelve natural reflections with
nonsingular normals $v_j$ representing those projections.
Their common fixed subspace
\[
          U_0=\bigcap_jv_j^\perp,\qquad
                         \dim U_0\ge d_i-12>0              \tag{OSP.21}
\]
is fixed by every word in them.  Every element of the
natural $\Omega(V_i)$ is projectively such a word,
so it acts by an original-field scalar on $U_0$.
Those scalars form a homomorphism from the perfect group
$\Omega(V_i)$ to $F_3^*$.  They must all be one.
This would make $U_0$ a nonzero natural fixed space of
$\Omega(V_i)$, contradicting its opposite full line
radicals and their generation.  This argument uses only
the natural simple projection; no splitting of a
represented Spin center has been introduced.
The theorem follows.
\end{proof}

\paragraph{Consequences and remaining scope.}
Together with D-OL-21, the same common support and
reflection-generation argument also excludes mixed
odd-characteristic Spin factors of both parities when
each odd dimension is at least $1135$ and each even
dimension is at least $1136$.
In the transitive-core main-source problem this gives
the necessary bound $d\le1133$ in odd natural dimension
and $d\le1134$ in even natural dimension, without an
ordinary-module restriction.  These are necessary
parameter bounds and imply no existence at the
remaining dimensions.

The new argument is the fourth-order carrier and its
all-order source consequence.  The exact finite
supplement verifies natural matrices, all three fixed
dimensions, finite Weyl covariance, carrier trace
counts and the rational constants.  It does not
replace the universal Clifford lift, group generation,
finite-field Witt theory, published character estimate
or actual-cover argument.  The primary lifting and
quasisimplicity interfaces are the fully written
D-OL-19--20 proofs, with their precise source records.
No additional Schur-multiplier classification, full
orthogonal reality theorem or positive cycle supplier
is asserted here.

\section{An exact decomposition refinement for all Spin characters}
\label{sec:spin-parameter-optimization}

\paragraph{Status and separation from the withdrawn candidate.}
These are new statements D-OL-27--29. The complete author and root
whole-proof reviews and the separate /root/sixfolds whole-proof,
primary-source and exact-certificate audit have passed.  The withdrawn D-OL-25--26 is not a premise:
in particular the second line of LST (4.3.2) has exponent
$\sqrt N/6$, not $\sqrt{N/6}$, and is not used here.
The argument below instead uses the published decomposition and
restriction lemmas and proves a new constituent-counting inequality.
The scalar bound used throughout is the uniform $\gamma=19/20$;
no subsequent numerical assertion $\gamma_q\le1/7$ is needed.

All fields in this section have odd characteristic.  Write
\[
 \gamma=\frac{19}{20},\qquad
 \rho_0=\frac{99997}{150000},\qquad
 \beta_0=\frac16-\frac1{10^7},\qquad
 \delta_0=\frac16+\frac1{10^5},\qquad \rho_0=1-2\delta_0.
 \tag{SPO.1}
\]

\subsection*{D-OL-27. Counting trivial factors with their actual weights}

\begin{lemma}[D-OL-27]
Let $H=\langle K,h\rangle$ be finite, where
$K=K_1*\cdots*K_m$ is a normal central product and $h$ normalizes
each $K_i$.  Suppose the finite extensions $H_i$ and local operators
$h_i$ satisfy the action-compatibility hypotheses of LST Lemma~4.3.2,
and every $H_i$-character irreducible and nontrivial over $K_i$ has
absolute character ratio at $h_i$ at most $\gamma$, where $0<\gamma<1$.
For every complex $H$-module $A$ of dimension $a$,
\[
 \frac{|\operatorname{tr}(h|A)|}{a}
 \le\gamma^m+
       \frac{1-\gamma^m}{m}
          \sum_{i=1}^m\frac{\dim A^{K_i}}a.
 \tag{SPO.2}
\]
No irreducibility of $A|K$, independent realization of the centers,
or extension of every constituent to the full product is assumed.
\end{lemma}

\begin{proof}
Representations of a finite central product are the representations
of the direct product on which the central kernel acts trivially.
Its irreducibles therefore have tensor factors with the required
central compatibility.  For each subset $J\subseteq\{1,\ldots,m\}$,
let $A_J$ be the sum of the $K$-isotypic components whose trivial
factor set is exactly $J$.  Since $h$ normalizes every $K_i$,
it preserves whether that factor is trivial; hence $A_J$ is
an actual $H$-submodule.
LST Lemma~4.3.2, including its complete proof for arbitrary characters,
applies to the nontrivial factors outside $J$ and gives
\[
 |\operatorname{tr}(h|A_J)|
               \le\gamma^{m-|J|}\dim A_J.
 \tag{SPO.3}
\]
This remains valid when $h$ permutes irreducible $K$-constituents:
the published proof then gives zero trace on a nontrivial orbit.
When a constituent is invariant, cyclic extension and Schur's lemma
give the tensor operator up to a root of unity.  Thus neither
multiplicity spaces nor scalar phases are suppressed.
The exact primary source is
\cite[Lemma~4.3.2 and its full proof, pp.~1912--1913]{HOS:LST2011}.

The function $x\mapsto\gamma^{m-x}$ is convex.  Its chord between
$0$ and $m$ therefore bounds its value above:
\[
 \gamma^{m-k}\le\gamma^m+\frac{k}{m}(1-\gamma^m)
       \quad(0\le k\le m).
 \tag{SPO.4}
\]
For example its second derivative is
$(\log\gamma)^2\gamma^{m-x}>0$, which proves the stated chord
inequality.  On summing (SPO.3)--(SPO.4), use the exact identity
\[
 \sum_J |J|\dim A_J=\sum_{i=1}^m\dim A^{K_i}.
 \tag{SPO.5}
\]
This follows factor by factor: a nontrivial irreducible factor has
zero fixed space, while a trivial one fixes its whole component.
Equations (SPO.2)--(SPO.5) prove the lemma.
\end{proof}

\paragraph{A central-extension character bound used below.}
We record the precise cover interface rather than treating a genuine
Spin character as an ordinary orthogonal character.  If $Z$ is central
in a finite group $B$, $J=B/Z$, and $\eta\in\operatorname{Irr}(B)$,
then for every $b\in B$,
\[
                   |\eta(b)|^2\le |C_J(bZ)|.
 \tag{SPO.6}
\]
To prove this, let $\psi$ be the central character of $\eta$.
The central-character Fourier projector and ordinary column
orthogonality give
\[
 \sum_{\xi\in\operatorname{Irr}(B),\ \xi|_Z=\xi(1)\psi}
                 |\xi(b)|^2
  =\frac{|C_B(b)|}{|Z|}
        \sum_{z\in Z_b}\overline{\psi(z)},
 \quad Z_b=\{z\in Z:zb\text{ is conjugate to }b\}.
 \tag{SPO.7}
\]
The set $Z_b$ is a subgroup: it is the image of the inverse
image of $C_J(bZ)$ under the homomorphism recording the central
commutator with $b$.  Thus the last character sum is either
zero or $|Z_b|$.  Moreover
$|C_B(b)||Z_b|/|Z|=|C_J(bZ)|$.  The nonnegative summand
$|\eta(b)|^2$ in (SPO.7) proves (SPO.6).
This is a finite-group statement and has no faithfulness premise.

\subsection*{D-OL-28. A direct trace bound at restricted support 93}

\begin{theorem}[D-OL-28]
Let $q$ be odd, let $S=\operatorname{Spin}(V)$, and let
$\chi$ be any nontrivial irreducible complex character of $S$.
If $s\in S$ has natural restricted support
\[
 \operatorname{supp}_{\mu_2}(s)
      :=\min_{\varepsilon=\pm1}
             \operatorname{codim}\ker(\varphi(s)-\varepsilon I)
      \ge93,
 \tag{SPO.8}
\]
then $|\chi(s)|/\chi(1)<\rho_0$ and consequently
$\nu(\rho(s))>\delta_0\chi(1)$.
The representation may have any central kernel.
\end{theorem}

\begin{proof}
We state the exact published inputs and their domains before
choosing the new parameters.
For an embedded restriction subgroup $\operatorname{Spin}(U)$
with $\dim U=b\ge7$ and nondegenerate orthogonal complement,
the odd-characteristic orthogonal case in the complete proof
of LST Proposition~4.2.3 gives
\[
 \frac{\dim A^{\operatorname{Spin}(U)}}{\dim A}
       <\alpha_b(q):=\frac{q^{15/4-b}}{q-1}
 \tag{SPO.9}
\]
for every nontrivial irreducible complex module $A$ of the
ambient Spin group.  The precise source is
\cite[Proposition~4.2.3, pp.~1909--1912, its odd-orthogonal
case in Step~1 and the $s=1$ summation in Step~6, p.~1912]{HOS:LST2011}.
Step~6 explicitly bounds the sum by
$\sum_{n=b+1}^{\infty}q^{15/4-n}=q^{15/4-b}/(q-1)$.
This uses the published bound in that proof, not the weaker
all-classical bound $q^{8/3-b/2}$ in the proposition's statement.
For $b=7$ we use the Step~1 telescoping argument itself:
every term has index $n\ge b+1=8$, within the proved range
$n\ge7$ of the odd-orthogonal bound.  Thus this case does not
extend the proposition's displayed all-classical $N\ge8$
statement without proof.  Both ambient parities are allowed.
It is a statement about
simply connected groups, so is not restricted to ordinary
$\Omega$ characters.

The decomposition and extension inputs are LST
Proposition~4.3.3 and the complete Steps~1--3 of Theorem~4.3.6,
printed pp.~1914--1916 and pp.~1918--1922.
An admissible decomposition is an orthogonal sum of nondegenerate
$s$-stable natural subspaces with determinant-one restrictions.
The resulting subgroups $K_i=\operatorname{Spin}(U_i)$ form a
central product.  Step~1 constructs finite extensions inside
the special Clifford groups and local operators $s_i$ having
exactly the same conjugation on $K_i$ as $s$.
For $\dim U_i\ge5$, their derived groups are quasisimple and
simply connected, and the Gluck bound (4.1.1) applies to
characters irreducible and nontrivial over $K_i$.
In particular its ratio is at most $\gamma=19/20$ for every
odd $q$.  These are the exact Spin-cover and admissibility
interfaces stated and proved in Step~1, pp.~1918--1919.
Neither a natural outer-automorphism classification at dimension
eight nor an assumption excluding triality is used: only the
actual Clifford block actions constructed there are needed.

We make explicit the large-indecomposable-block estimate, including
the dimension-seven endpoint.  For an admissibly indecomposable
natural block $U$ of dimension $b\ge8$, the final paragraph of the
odd-orthogonal proof of Proposition~4.3.3 gives
\[
 |C_{O(U)}(\bar s_U)|<2q^b,\qquad
 d(\operatorname{Spin}(U))>\tfrac89q^{b-3}.
 \tag{SPO.10}
\]
Both inequalities are stated there for $b\ge7$, not just $b\ge12$.
The definition of $d$ as the minimal nontrivial complex
representation degree is p.~1908; the precise inequalities are
on p.~1916, with $L(U)=\operatorname{Spin}(U)$ defined on p.~1913.
Let $B=\langle K_U,s_U\rangle$ be the special-Clifford extension.
Its natural-map kernel is central and its image lies in $SO(U)$.
For every $\eta\in\operatorname{Irr}(B)$ irreducible and nontrivial
over $K_U$, the central-quotient bound (SPO.6) therefore gives
\[
 \frac{|\eta(s_U)|}{\eta(1)}
 <\frac{9\sqrt2}{8}q^{3-b/2}
 \le\frac{3\sqrt2}{8}<\frac9{16}\quad(b\ge8).
 \tag{SPO.11}
\]
The last strict inequality is $\sqrt2<3/2$.
This treats the genuine Spin characters without identifying them
with ordinary orthogonal characters.

Here is the additional endpoint proof when $b=7$.
Write the natural element as a commuting semisimple-unipotent
product.  The complete odd-orthogonal classification in
Proposition~4.3.3, pp.~1915--1916, shows that a non-$\pm1$
primary summand has even dimension and gives an admissible
orthogonal summand.  It cannot occupy all seven dimensions.
The minus-one primary space has even dimension because the
natural element has determinant one; both the plus and minus
restrictions have determinant one.  Admissible indecomposability
therefore forces the minus-one primary space to be zero.
The element is unipotent.  The orthogonal unipotent decomposition
in that same proof consists of blocks $J_{2k+1}$ or $2J_{2k}$,
each with determinant one.  In dimension seven, admissible
indecomposability thus forces a single regular block $J_7$.

We compute its full orthogonal centralizer directly.
The general-linear commutant is
$R=\mathbf F_q[x]/(x^7)$, with $T=1+x$ representing $J_7$.
The adjoint involution is $\sigma(T)=T^{-1}$, hence
$\sigma(x)=-x+x^2-\cdots+x^6$.
Its diagonal on the augmentation ideal $(x)$ in the basis
$x,x^2,\ldots,x^6$ is $(-1)^j$.
Since $q$ is odd, the involution is semisimple and its fixed
space on that six-dimensional ideal has dimension three.
Write $U_R=1+(x)$ for the commutative odd-order unit group,
distinguishing it from the natural block $U$.  Consequently,
\[
 |U_R|=q^6,\qquad |U_R^\sigma|=q^3,\qquad
 |\ker(U_R\longrightarrow U_R^\sigma,\ u\mapsto u\sigma(u))|=q^3.
 \tag{SPO.11a}
\]
Indeed the norm is onto: squaring is a bijection on $U_R$, and
the unique square root of any $\sigma$-fixed unit is itself
$\sigma$-fixed.  Multiplying by the constant units shows that
$|C_{O(U)}(J_7)|=2q^3$, since a constant has norm one exactly
when it is $1$ or $-1$.
The same minimum-degree bound from (SPO.10), now at $b=7$,
and (SPO.6) imply
\[
 \frac{|\eta(s_U)|}{\eta(1)}
 <\frac{9\sqrt2}{8}q^{-5/2}
 \le\frac{\sqrt{2/3}}8<\frac18.
 \tag{SPO.11b}
\]
No assertion about arbitrary outer automorphisms of dimension-seven
or dimension-eight groups is needed.

Refine the natural space into admissibly indecomposable blocks.
If some block has dimension at least seven, its fixed fraction
is bounded by (SPO.9).  The complement of the fixed space is
stable under $s$; the one-factor case of LST Lemma~4.3.2 applies
to its actual local extension.  At $b=7$,
$\alpha_7(3)<1/54$; at $b\ge8$, $\alpha_b(q)\le\alpha_8(3)<1/162$.
Thus in the two cases the full trace ratio is respectively less than
\[
 \frac18+\frac1{54}=\frac{31}{216}<\frac35,
 \qquad
 \frac9{16}+\frac1{162}=\frac{737}{1296}<\frac35<\rho_0.
 \tag{SPO.12}
\]
If the block is the entire ambient group its fixed fraction is zero.
This is Case A, with $N_1=7$.

It remains to handle the case where every indecomposable block
has dimension at most $6$.
Call a block good when its local operator is noncentral in
the corresponding extension, and bad otherwise.
The precise classification and pairing from the published
Step~3, pp.~1921--1922, applies as follows.
A bad block has dimension at most two and is either scalar
$1$ or $-1$, or is a two-dimensional rotation with eigenvalue
pair $\{\lambda,\lambda^{-1}\}$, $\lambda\ne\pm1$.
Pair different bad types, and then pair equal remaining
rotation types.  Every such pair becomes good and has
dimension at most four.  The unpaired remainder is either
at most one two-dimensional rotation block or consists
of scalar blocks of a single sign.
Consequently the total dimension $\Sigma_g$ of the good
blocks satisfies
\[
              \Sigma_g\ge\operatorname{supp}_{\mu_2}(s)-2.
 \tag{SPO.13}
\]
For a rotation remainder this follows by subtracting at most
two from the full dimension.  For a single scalar sign,
the bad subspace lies in that sign's eigenspace, so the
stronger bound without minus two holds.  This explains
directly why the restricted scalar set $\mu_2$ suffices
here, without importing an incorrectly transcribed exponent.
Every good block now has dimension in $\{3,4,5,6\}$:
in dimensions one and two the determinant-one isometry group
and its special-Clifford extension are abelian, so such a block
is bad.  The above pairings create only dimensions three or four.

We use a sharper elementary packing argument, with $m_0=N_2=8$.
Any multiset of integers in $\{3,4,5,6\}$ of total at least $91$
contains eight disjoint groups each of sum at least eight.
To prove this, choose a packing with the maximum number of such
groups, then make each group inclusion-minimal by returning
removable items to the remainder.  Each minimal group has sum
at most twelve; a sum-twelve group must be exactly $6+6$.
The remainder has sum at most seven.  If its sum is seven,
it is $3+4$.  In that case no $6+6$ group can occur, since
replacing it and the remainder by $6+3$ and $6+4$ would increase
the number of groups.  Therefore a packing with at most seven
groups has total at most
\[
 \max\{7\cdot12+6,\ 7\cdot11+7\}=90.
 \tag{SPO.14}
\]
This proves the assertion.  The threshold is sharp for this
packing problem, as fifteen items of size six have total ninety
and yield at most seven groups.

By (SPO.8)--(SPO.13), the good blocks have total at least
$93-2=91$.  Apply this packing, and append all unused blocks,
including the bad blocks, to the eighth group.  The resulting
admissible decomposition has eight subspaces of dimension at
least eight, each containing a good block and hence carrying a
noncentral local operator.  This uses the published classification
and admissibility interfaces, but the packing improvement is
proved here.  No real-valued parameter or square-root exponent
from the numerical conclusion of LST is chosen.

Apply D-OL-27 and (SPO.9) to these eight actual Spin
subgroups.  Monotonicity in their dimensions and in $q$
gives the same upper bound for each fixed fraction,
\[
 \alpha_8(q)\le\alpha_*:=\frac{3^{-17/4}}2,
 \qquad
 \frac{|\chi(s)|}{\chi(1)}
       <c+(1-c)\alpha_*,\quad c=(19/20)^8.
 \tag{SPO.15}
\]
All constants are certified by a short positive-integer
comparison.  Direct rational subtraction gives
\[
 u:=\frac{\rho_0-c}{1-c}
       =\frac{247774877}{25849310877}>0,
 \qquad
 16\cdot3^{17}\cdot247774877^4>25849310877^4.
 \tag{SPO.16}
\]
The second inequality is exactly $u^4>\alpha_*^4$.
Since both quantities are positive, $u>\alpha_*$;
(SPO.15)--(SPO.16) give the required ratio $<\rho_0$.
The finite certificate checks this integer inequality
directly.  Finally $|\chi(s)|\ge\chi(1)-2\nu(\rho(s))$
after rotating a largest-eigenspace eigenvalue, so the
support exceeds $\delta_0\chi(1)$.
This proves the theorem for every central kernel.
\end{proof}

\subsection*{D-OL-29. The resulting actual-source parameter bound}

\begin{theorem}[D-OL-29]
An actual rank-nine six-dimensional multiplicity source cannot
have a nonempty irreducibly acting normal central product of
represented odd-characteristic Spin factors for which every
odd natural dimension is at least $187$ and every even natural
dimension is at least $188$.
Every nontrivial irreducible complex factor module and every
represented central kernel is allowed, together with all
compatible scalars, outer automorphisms and factor permutations.
No selfduality or transitivity is required.
In the transitive-core problem the remaining necessary natural
dimension is therefore at most $185$ in odd dimension and
at most $186$ in even dimension.
\end{theorem}

\begin{proof}
We identify the changed quantitative interfaces in the
already complete D-OL-19--24 normalizer proof.
The Clifford lift and actual represented-kernel preservation
need only natural dimension at least $22$.
All fixed-hyperbolic-pair, quadratic-field and arbitrary
semilinear descent cases give support at least $3a/16$ or
$7a/32$, except the rank-one reflection case.
Their radical complements have dimension at least $d-4$,
so the elementary lower-dimensional conditions remain valid.
An inner prime-order case without a fixed hyperbolic pair
has restricted natural support at least
$\lceil(d-2)/2\rceil$, which is at least $93$ at both
new thresholds.  D-OL-28 replaces the previous invocation
of the first line of LST(4.3.2) there.
In the odd-dimensional all-orders reduction, the square of
a possible $2^s$-order root with $s\ge3$ lies in $\Omega$
and has restricted support at least $d-1>93$;
D-OL-28 applies to its Spin lift as well.
The remaining fourth-order case is D-OL-22 and requires
only $d\ge25$.
The complete even-dimensional determinant reduction in
D-OL-20 is unchanged.  Thus no extra natural-dimension
condition has been omitted.

Here is the sharper rank-one estimate in full, without
using the withdrawn D-OL-25 as an input.
In the nondegenerate radical dual of dimension $b\ge43$,
a nonzero norm orbit has size at least
$q^{b-1}-q^{b/2}-1$ by the elementary quadratic Fourier
sum.  The fixed hyperplane of a reflection is nondegenerate
and has at most $q^{b-2}+q^{(b-1)/2}$ vectors of that norm.
Therefore its fixed fraction is at most
\[
 c_q^*:=\frac{q^{-1}+q^{-20}}{1-q^{-20}-q^{-40}}.
 \tag{SPO.17}
\]
Indeed, after division by $q^{b-1}$ the numerator error
is $q^{-(b-1)/2}\le q^{-21}$ and the denominator errors
are $q^{1-b/2}\le q^{-20}$ and $q^{1-b}\le q^{-40}$.
The underlying Gauss bound follows by adding the $q-1$
nontrivial additive-character terms, each of magnitude
$q^{b/2-1}$; it applies also to norm zero, with the
possible zero vector removed in the denominator.
The fraction in (SPO.17) decreases with $q$.
The exact rational endpoints are
\[
           c_3^*<1-4\beta_0,\qquad c_5^*<1/4.
 \tag{SPO.18}
\]
The lifted opposite-radical fixed space has dimension
at most $a/2$, and nonfixed involution weight orbits
have length two.  Thus the rank-one support is at least
$(1-c_q^*)a/4$, greater than $\beta_0a$ at $q=3$ and
greater than $3a/16>\delta_0a$ at $q\ge5$.

It follows, with every projective order retained as in
D-OL-20 and D-OL-23, that every nonscalar compatible
normalizer has support at least $\beta_0a$.
Any operator with support at most $\delta_0a$ has
$q=3$ and natural projective lift a reflection of
projective order two.  Actual squares may still be
central scalars, as explicitly permitted in those proofs.

For every nontrivial factor the lifted-radical weight
orbit degree bound gives, with $b=d-2\ge185$,
\[
 a\ge q^{b-1}-q^{\lceil b/2\rceil}-1
      \ge(q-1)q^{b-2}-1
      \ge2\cdot3^{183}-1>10^6.
 \tag{SPO.19}
\]
This uses only a nontrivial weight and its full inverse-image
Levi orbit, so includes genuine Spin modules.
Write $e$ for the product of the factor degrees; then $e>10^6$.

Call an actual inertia good when it moves at least
$\delta_0e$ dimensions.  Every genuine inertia moves
at least $\beta_0e$; nonidentity scalars move all $e$,
and nontrivial factor cycles have support at least $e/4$.
A bad inertia is therefore factor preserving, with
every active factor a projective reflection over $\mathbf F_3$.
Two such active factors have support at least
$2\beta_0(1-\beta_0)e>e/4$, by the exact two-eigenspace
convolution calculation.  Thus a bad inertia has exactly
one active factor, independent of its scalar phase.

The actual intermediate-cover Euler and moduli arguments
force a sphere base with at least twelve genuine marks,
total moving rank $2e+6$, and a generating product-one tuple.
The exact inequalities
\[
 13\beta_0-2>1/7,\quad
 11\beta_0+\delta_0-2=89/10^7,\quad
 6\cdot10^7/89<10^6
 \tag{SPO.20}
\]
exclude thirteen marks or even one good inertia among twelve.
Exactly twelve bad inertias remain.
Their projections are trivial on factors over $q\ne3$,
which is already incompatible with generation.
On a factor over $\mathbf F_3$, at most twelve natural
reflections represent their projective images.
The common perpendicular of their normals has dimension
at least $d-12>0$ and is fixed by every reflection word.
Every element of the natural $\Omega$ is projectively
such a word and therefore acts by a ground-field scalar
on that common subspace.  Perfectness kills this scalar
character, contradicting the absence of a nonzero natural
fixed space.  In even dimension this is argued through
$P\Omega$ exactly as in D-OL-21; no splitting of either
the natural or represented Spin center is introduced.
The contradiction proves the theorem.
\end{proof}

\paragraph{Exact evidential boundary.}
The weighted trivial-factor estimate, central-extension column
bound, regular-seven-block centralizer count and sharp small-block
packing are proved explicitly above.  The last two refinements
were contributed by the root collaborator and checked here; this
is a joint derivation, not an independent-origin audit.
The decomposition and subgroup estimates are invoked with
their complete published hypotheses and exact proof locators.
The finite certificate checks all constants and packing
endpoints, and does not replace those universal arguments.
This yields necessary remaining parameter bounds and
asserts no actual source in the surviving dimensions.

\section{The exclusion of sporadic normal cores}
\label{sec:sporadic-core}

\paragraph{Status and input.}
The proofs and exact finite certificates in this section are new.
The complete mathematical input consists of the ordinary character
tables of the twenty-six universal sporadic covers and twelve outer
extensions, with their stored power maps and subgroup fusions, in
the author-maintained Character Table Library, version 1.3.11.
Every one of these thirty-eight tables identifies its origin as the
published \emph{ATLAS of Finite Groups}; the exact table identifiers,
source strings, and complete data are retained in
\texttt{character\_inputs.json} and \texttt{coverage.json}.
We use these finite tables and their identification with the named
groups as published mathematical input, not as groups reconstructed
from orthogonality alone.  The arithmetic certificates below have
been executed.  Both the parent and a sibling agent independently
read the complete proof and programs and returned PASS.  The parent
reran the standalone verifier; the sibling additionally rebuilt the
full calculation using independent cyclotomic arithmetic, restriction
inner products, and Ramanujan-sum Fourier multiplicities.  Both
reconstructions passed all thirty-eight tables, 3035 selected rows,
and 473972 nonscalar entries.  Exact review scope and source-input
limitations are retained in the check register and peer-audit report.

For a finite-order operator on an $e$-dimensional complex space put
\[
 \nu(g)=e-\max_\lambda\dim\ker(g-\lambda).
\]
This invariant is unchanged by multiplication by a nonzero scalar.
A nonidentity scalar has $\nu(g)=0$ but actual moving rank $e$.

\subsection*{D-SPOR-01. Complete finite character input}

\begin{lemma}\label{D-SPOR-01}
Let $S$ be any sporadic simple group and $U$ its universal perfect
central cover.  Every nonlinear irreducible character of $U$ has
degree at least six.  In every such representation of degree $e$,
every nonscalar element has $\nu\ge e/4$.

If $\operatorname{Out}(S)$ is nontrivial, it has order two.  Let
$A$ be the cover-by-outer-involution extension in Table~\ref{tab:sporadic-coverage}.
For every irreducible character of $A$ whose restriction to $U$ is
nonlinear and irreducible, the same bound holds for every nonscalar
element of $A$, including elements outside $U$.
\end{lemma}

\begin{table}[ht]
\centering\small
\begin{tabular}{c|c|c|c@{\qquad}c|c|c|c}
$S$&$|Z(U)|$&$|\operatorname{Out}(S)|$&$A$&
$S$&$|Z(U)|$&$|\operatorname{Out}(S)|$&$A$\\ \hline
$B$&2&1&---&$J_4$&1&1&---\\
$Co_1$&2&1&---&$Ly$&1&1&---\\
$Co_2$&1&1&---&$M$&1&1&---\\
$Co_3$&1&1&---&$M_{11}$&1&1&---\\
$Fi_{24}'$&3&2&$3.F3+.2$&$M_{12}$&2&2&$2.M12.2$\\
$Fi_{22}$&6&2&$6.Fi22.2$&$M_{22}$&12&2&$12.M22.2$\\
$Fi_{23}$&1&1&---&$M_{23}$&1&1&---\\
$HN$&1&2&$HN.2$&$M_{24}$&1&1&---\\
$HS$&2&2&$2.HS.2$&$McL$&3&2&$3.McL.2$\\
$He$&1&2&$He.2$&$O'N$&3&2&$3.ON.2$\\
$J_1$&1&1&---&$Ru$&2&1&---\\
$J_2$&2&2&$2.J2.2$&$Suz$&6&2&$6.Suz.2$\\
$J_3$&3&2&$3.J3.2$&$Th$&1&1&---
\end{tabular}
\caption{Complete sporadic coverage.  The rightmost entries in each
half are exact CTblLib identifiers.  The cover identifier is obtained
by prefixing the displayed central order when it exceeds one.
CTblLib writes $Fi_{24}'$ as $F3+$.}
\label{tab:sporadic-coverage}
\end{table}

\begin{proof}
The group and multiplier identifications in the table are the
\emph{ATLAS} data \cite{SPOR:Atlas}, accessed through
\texttt{ExtensionInfoCharacterTable} and the complete
\texttt{IsSporadicSimple} list of \cite{SPOR:CTblLib}.
The latter attribute records the Schur multiplier and outer
automorphism group, and specifies the corresponding cover and
automorphism-group identifiers; its precise manual locator is
\cite[\S3.7, entry \texttt{ExtensionInfoCharacterTable}]{SPOR:CTblLib}.
All twelve indicated extension tables and their actual index-two
subgroup fusions from $U$ are present.  No group is omitted because
its order or representation degree is large: the Baby Monster and
Monster tables are included.

Here is the finite proof of the asserted spectral bounds, including
the exact test performed on every entry.  A character value $z$ is
encoded as an integer, or as an integral polynomial in
$\zeta_n=\exp(2\pi i/n)$.  Compute $a=z\overline z$ exactly.
If $a=e^2$, the operator is scalar: equality in the triangle
inequality for a sum of $e$ unit complex eigenvalues forces them
all to coincide.  Conversely a scalar has $a=e^2$.  Such classes
are discarded only for the assertion about $\nu$.

For a nonscalar entry write $a=\sum b_j\zeta_n^j$.  The rational
integer $B=\sum|b_j|$ is an upper bound for $|z|^2$.  Whenever
$4B\le e^2$, the trace of the operator after rescaling its largest
eigenspace to eigenvalue one gives
\[
 e-2\nu\le |z|\le e/2,
 \qquad\text{and therefore}\qquad \nu\ge e/4.       \tag{SPOR.1}
\]
When that sufficient triangle bound fails, use the character values
on powers and the exact Fourier formula.  For a class represented
by an element of order $n$, the multiplicity of $\zeta_n^k$ is
\[
 m_k=\frac1n\sum_{r=0}^{n-1}\chi(g^r)\zeta_n^{-kr}.
                                                               \tag{SPOR.2}
\]
Every number in (SPOR.2) is computed in a cyclotomic field with
integer coefficients.  The certificate verifies that each $m_k$
is a nonnegative integer, that $\sum m_k=e$, and that
$4(e-\max m_k)\ge e$.  This formula is just the orthogonality
formula for the characters of the cyclic subgroup; no numerical
approximation or eigenvalue recognition is involved.

For an extension table, retain exactly the rows whose restriction
along the stored subgroup fusion has inner product one with itself
on $U$, and whose degree exceeds one.  These are precisely the
characters specified in the lemma.  The executed GAP certificate
\texttt{scan\_quarter\_bound.g} checks all twenty-six cover tables
and twelve extension tables.  It covers $3035$ selected rows and
$473972$ nonscalar entries.  Its triangle test needs (SPOR.2) only
for $24$ entries in the table $2.J2$, and all pass.  They occur in
rows $2,3$ of degree $14$ and rows $22,23$ of degree $6$; their
computed codimensions are respectively $10$ and $4$.  The minimum
nonlinear cover degree is six.  These are complete scans of the
rows and columns, not lists of favorable conjugacy classes.

The separate program \texttt{export\_inputs.g} exports all the
complete tables, not merely the selected rows or exceptional values.
It also verifies exact row orthogonality, the degree-squared sum,
and the index-two subgroup data.  The independent arithmetic
implementation \texttt{verify\_quarter\_bound.py} calls no GAP
code.  It computes in $\mathbf Z[x]/(\Phi_n(x))$, independently
selects the extension rows by equality of the restricted character
with a row of the complete cover table, and reconstructs the
Fourier multiplicities from prime power maps.  Its different
polynomial basis needs only twelve Fourier fallbacks; it verifies
the same $3035$ rows and $473972$ entries and again finds minimum
degree six and no failed inequality.  The exact input SHA-256 is
\texttt{c2268f05043ed29a9048d5fe7d55b1a096e5c8615fa049b3438cca088f703edb}.

Thus the elementary test (SPOR.1)--(SPOR.2), applied exhaustively
to the stated published input, proves the finite lemma.  The
library notes a table-automorphism ambiguity for a power map in
$3.McL$; no Fourier fallback uses that table, so this ambiguity
does not affect the proof.  The exceptional Fourier checks all
use the stored $2.J2$ power maps.
\end{proof}

\subsection*{D-SPOR-02. All scalar and outer normalizers}

\begin{proposition}\label{D-SPOR-02}
Let $N\subset\operatorname{GL}(W)$ be a finite quasisimple group
acting faithfully and irreducibly, with sporadic simple quotient,
and put $a=\dim W$.  Every nonscalar element of its linear
normalizer has largest-eigenspace codimension at least $a/4$.
The normalizing element may have any scalar phase.  No equality
of character fields or model over an imaginary quadratic field
is required.
\end{proposition}

\begin{proof}
Let $S=N/Z(N)$ and let $U\to N$ be the quotient of the universal
perfect central cover.  The given representation becomes a
nonlinear irreducible representation $\rho$ of $U$.
The centralizer of $N$ in $\operatorname{GL}(W)$ is scalar, by
Schur's lemma.  An automorphism of the perfect group $N$ that
induces the identity on $S$ is the identity: its discrepancy from
the identity is a homomorphism from $N$ to the central abelian
group $Z(N)$.  Thus the linear normalizer modulo scalars embeds
in $\operatorname{Aut}(S)$, and its quotient by the projective
image of $N$ embeds in $\operatorname{Out}(S)$.

If this outer image is trivial, every normalizer element is a
scalar multiple of an element of $N$, and D-SPOR-01 applies.
Otherwise $\operatorname{Out}(S)=C_2$.  Fix the actual finite
extension $A=U.2$ used in D-SPOR-01 and choose $t\in A\setminus U$.
Write $\alpha$ for conjugation by $t$ on $U$ and $u_0=t^2\in U$.
This argument does not require $t^2=1$ or a split extension.

Let $g$ be a normalizing matrix inducing an outer automorphism.
After multiplying $g$ by an element of $N$, its induced
automorphism of $S$ agrees with that of $\alpha$.  Automorphisms
of $S$ lift uniquely to its universal perfect cover: this follows
from the universal property by pulling back the central extension,
and uniqueness also follows because a discrepancy is a homomorphism
from the perfect group $U$ to its center.  Consequently $\alpha$
preserves the kernel of $\rho$, and
\[
             g\rho(u)g^{-1}=\rho(\alpha(u))\qquad(u\in U)
                                                               \tag{SPOR.3}
\]
after the indicated inner adjustment.  In particular the character
of $\rho$ is $\alpha$-invariant.

By (SPOR.3), $g^2\rho(u_0)^{-1}$ centralizes $\rho(U)$ and is
a nonzero scalar.  Multiply $g$ by a scalar square root of its
inverse to obtain $T$ with $T^2=\rho(u_0)$.  Equation (SPOR.3)
and this square relation give a representation of the actual
index-two extension $A$, extending $\rho$ and sending $t$ to $T$.
It is irreducible because its restriction to $U$ is irreducible.
Therefore its character is one of the extension rows selected in
D-SPOR-01.  The original normalizing element is a scalar multiple
of an element of this represented $A$, up to the inner adjustment
already made.  The bound of D-SPOR-01 is unchanged by such scalar
multiplication.  This proves the proposition.

The argument also explains why no isoclinic version of $U.2$ is
missing: an alternative square relation changes only the scalar
normalization of its intertwiner.  We use the actual fixed table
$A$ and extend to it directly.  Characters that are exchanged by
the outer automorphism cannot occur as irreducible restrictions
of a normalizing representation and are correctly omitted.
\end{proof}

\subsection*{D-SPOR-03. Tensor permutations and the sixfold exclusion}

\begin{theorem}\label{D-SPOR-03}
Suppose a complex inducing representation for a finite-cover
sixfold source has faithful finite image $J\subset\operatorname{GL}(V)$.
Suppose $J$ has a normal subgroup $E$ that acts irreducibly on $V$
and is a central product of quasisimple groups with sporadic simple
quotients.  Then the actual six-dimensional multiplicity variation
cannot have period-map rank nine.

This excludes sporadic cores in D-BI-19, including arbitrary outer
automorphisms and permutations of quasisimple tensor factors.  It
does not require a nonsplit polarization and permits an inducing
character field larger than the original quadratic field.
\end{theorem}

\begin{proof}
Write $E=E_1\cdots E_t$ with commuting quasisimple factors.
Irreducibility gives
\[
                   V=V_1\otimes\cdots\otimes V_t,
                   \qquad a_i=\dim V_i\ge6.                 \tag{SPOR.4}
\]
Each factor representation is irreducible and faithful: an element
acting as the identity on $V_i$ acts as the identity on $V$, and
$E_i$ is an actual matrix subgroup.  The factor degrees are at
least six by D-SPOR-01.

An automorphism of $E$ permutes its quasisimple factors.  Indeed,
it permutes the simple factors of $E/Z(E)$, and each $E_i$ is
the derived subgroup of the preimage of its simple factor, so
the perfect lifts are intrinsic.  The matrix algebras generated
by the $E_i$ are the full factor algebras
$\operatorname{End}(V_i)\otimes1$.  A normalizing matrix therefore
acts as a product of linear factor maps followed by the induced
permutation.  This follows by the matrix-unit proof that every
automorphism of a full matrix algebra is inner, followed by its
tensor-factor centralizer; no rational field-of-definition
assertion is needed.

First suppose $g\in J$ fixes every tensor factor.  Then
$g=A_1\otimes\cdots\otimes A_t$, with each $A_i$ normalizing
the represented $E_i$.  If $g$ is nonscalar, some $A_i$ is
nonscalar.  D-SPOR-02 bounds its largest eigenspace by $3a_i/4$.
For tensor products of diagonalizable operators, the largest
eigenspace multiplicity is at most the largest multiplicity in
one factor times the product of the other dimensions: sum over
the eigenspaces in the other factors.  Thus
\[
                              \nu(g)\ge\dim V/4.             \tag{SPOR.5}
\]
The factor operators are diagonalizable up to scalar since $g$
has finite order; alternatively D-SPOR-02 represents each one
as a scalar multiple of a finite-order matrix.

Now suppose $g$ has a tensor-factor cycle of length $l\ge2$.
All factors along that cycle have a common dimension $a$.
On their tensor product the operator has the form of cyclic
rotation composed with invertible maps between successive factors.
Transporting bases around the cycle conjugates it by a product
of changes of basis to cyclic rotation with a single remaining
twist $B$ on one factor.  This twist is diagonalizable.  To see
this without a unitarity convention, a power of the full $g$ is
the identity, so a power of each cycle operator is scalar.
Its $l$-th power is conjugate to $B^{\otimes l}$, and a power
of $B$ must therefore be scalar.  The polynomial $x^k-c$ has
distinct roots in characteristic zero when $c\ne0$.

Choose a diagonalizing basis for $B$.  The basis words of length
$l$ on an alphabet of size $a$ are then permuted by cyclic
rotation, with nonzero scalar weights.  On an orbit of length
$d$, the operator has characteristic polynomial $x^d-c$ with
$c\ne0$, so any specified eigenvalue occurs at most once in
that orbit.  There are exactly $a$ constant words; every other
word orbit has length at least two.  Therefore the largest
eigenspace on this cycle has dimension at most
\[
       a+\frac{a^l-a}{2}=\frac{a^l+a}{2}
                  \le\frac34 a^l,                         \tag{SPOR.6}
\]
because $a^{l-1}\ge2$.  Eigenvalue coincidences between different
orbits do not change this count.  Tensoring with the other cycle
operators gives (SPOR.5) for the full $g$ as well.

We now apply the bound to the actual intermediate cover attached
to the inducing representation.  Put $e=\dim V$ and let $h$ be
the genus of its base.  Its multiplicity cohomology has complex
dimension six, so with genuine inertias $g_1,\ldots,g_r$,
\[
       6=2e(h-1)+\sum_j\operatorname{rank}(1-g_j).           \tag{SPOR.7}
\]
Frobenius reciprocity identifies this complex Hodge filtration
with the original one.  Its conjugate is determined by the actual
Hermitian polarization, so the period rank factors through this
intermediate cover even for a larger inducing character field.
Finite covers of a fixed pointed curve have only finitely many
topological possibilities.  If $h=1$, (SPOR.7) gives at most six
marked points and period rank at most six; if $h\ge2$, it is
impossible because $e\ge6$.  Rank nine therefore forces
\[
               h=0,\qquad r\ge12,\qquad
               \sum_j\operatorname{rank}(1-g_j)=2e+6.
                                                               \tag{SPOR.8}
\]
For every nonscalar inertia, (SPOR.5) gives actual moving rank
at least $\lceil e/4\rceil$.  A genuine scalar inertia has moving
rank $e$, so obeys the same lower bound.  Hence
\[
                        12\lceil e/4\rceil\le2e+6.         \tag{SPOR.9}
\]
This is impossible for $e>6$, since its left side is at least
$3e$; for $e=6$ its left side is $24$ and its right side $18$.
Finally $e\ge6$ follows from (SPOR.4).  The contradiction proves
the theorem.  In particular it is not merely an obstruction to
one branch profile or one small sporadic representation.
\end{proof}

\paragraph{Reproduction and scope.}
The complete group tables are an explicit finite published input.
The two independent arithmetic implementations verify the spectral
inequality without floating-point computations.  They do not
reprove the identification of ATLAS tables with their groups or
the classification of sporadic groups.  The theorem assumes the
normal central product acts irreducibly; a reducible normal subgroup
is not silently substituted.  D-BI-19 supplies that precise hypothesis
for the primitive tensor-indecomposable blocks to which it applies.


\section{An obstruction for essentially self-dual spin cores}
\label{sec:spin-core-obstruction}

The sporadic calculation does not cover alternating simple factors.
The following separate argument removes the essentially self-dual
spin case in all alternating degrees at least eight.  It does not
remove ordinary alternating representations or nonselfdual spin
representations.

\begin{lemma}[HD-SPIN-01a: the spin-normalizer moving rank]
\label{HD-SPIN-01a}
Let $N=2.\mathfrak A_n$, $n\ge8$, act faithfully and irreducibly
on a complex space $W$ of dimension $a$.  Every nonscalar
finite-order element $g$ of its linear normalizer satisfies
\[
                  \nu(g)\ge a/8,                         \tag{SPIN.1}
\]
where $\nu$ is the codimension of a largest eigenspace.
The same conclusion holds if $g$ has finite order only after a
scalar rescaling.
\end{lemma}

\begin{proof}
We use the local character-table input appearing in the complete
proof of \cite[Lemma~2.23, pp.~621--622]{HD:GT2012}: every faithful
irreducible character $\psi$ of $2.\mathfrak A_6$ satisfies
\[
       |\psi(h)|\le\psi(1)/2\quad
       (h\notin Z(2.\mathfrak A_6)).                       \tag{SPIN.2}
\]
This input is also independently checked on the complete retained
table in the certificate described below.  We use the input (SPIN.2),
not the weaker final $7/8$ character-ratio conclusion of that lemma.

Schur's lemma identifies the linear centralizer of $N$ with the
scalars.  An automorphism of the perfect group $N$ inducing the
identity on $\mathfrak A_n$ has the form $x\mapsto z(x)x$, with
$z$ a homomorphism to the center, and is therefore the identity.
Thus the normalizer modulo scalars embeds in
$\operatorname{Aut}(\mathfrak A_n)=\mathfrak S_n$ for $n\ge8$;
the latter standard automorphism description is also used in
\cite[proof of Proposition~3.4, p.~626]{HD:GT2012}.
Let $\pi\ne1$ be the permutation induced by $g$.

There is a three-cycle $c$ not centralized by $\pi$.  Indeed, the
three-cycles generate $\mathfrak A_n$, whose centralizer in its
natural permutation group is trivial.  The last assertion can
also be checked directly: if $\pi$ commutes with every three-cycle,
it preserves the support of every three-cycle, hence every
three-element subset; the intersections of these subsets recover
each singleton for $n\ge4$.  Choose a lift $t\in N$ of $c$.
Then
\[
                   h=[g,t]\in N
\]
projects to the nonidentity product of the three-cycles
$\pi c\pi^{-1}$ and $c^{-1}$.  Its support is contained in at
most six letters.  It therefore belongs to a natural
$2.\mathfrak A_6$ subgroup of $N$ and is noncentral there.
This is precisely the natural-subgroup step in the proof of
\cite[Lemma~2.23]{HD:GT2012}; restriction of the standard double
cover to these six letters is the standard double cover.

The central involution of this subgroup acts as $-1$ on $W$.
Consequently every irreducible constituent of the restriction
to it is faithful on the center, and hence faithful, since a
proper normal subgroup of a quasisimple group is central.
Summing (SPIN.2) over the constituents gives
$|\operatorname{tr}_W(h)|\le a/2$.  If $m$ is the largest
eigenvalue multiplicity of $h$, rotate that eigenvalue to one.
The triangle inequality gives
\[
           2m-a\le|\operatorname{tr}_W(h)|\le a/2,
           \qquad \nu(h)=a-m\ge a/4.                     \tag{SPIN.3}
\]

Choose a scalar $\lambda$ with
$\operatorname{rank}(1-\lambda g)=\nu(g)$.
For any invertible matrices $A,B$,
$\operatorname{rank}(1-AB)\le
\operatorname{rank}(1-A)+\operatorname{rank}(1-B)$.
The commutator of $\lambda g$ with $t$ is still $h$, so
\[
 \nu(h)\le\operatorname{rank}(1-h)
 \le 2\operatorname{rank}(1-\lambda g)=2\nu(g).
\]
Together with (SPIN.3) this proves (SPIN.1).
\end{proof}

\begin{theorem}[HD-SPIN-01b: exclusion of essentially self-dual spin cores]
\label{HD-SPIN-01b}
Suppose a finite-cover Weil sixfold source has nonsplit actual
polarization, exact imaginary quadratic character field, and
period rank nine.  Let an irreducible primitive inducing block
have finite image $J$ and degree $e\ge9$, with arbitrary inducing
character field.  If $J$ is essentially self-dual, it cannot have
an irreducibly acting normal central product
\[
                  E=E_1\cdots E_t,\qquad
                  E_i\simeq2.\mathfrak A_{n_i},
                  \quad n_i\ge8,                        \tag{SPIN.4}
\]
whose factor representations are faithful.
In particular this removes the essentially self-dual spin-core
case from D-BI-19.  Permutations of isomorphic factors and all
compatible outer automorphisms are allowed.
\end{theorem}

\begin{proof}
Restriction to the commuting quasisimple factors gives
$V=\bigotimes_i V_i$, with $a_i=\dim V_i$.  The published
spin-degree lower bound
\[
                   a_i\ge2^{\lfloor n_i/2\rfloor-1}\ge8  \tag{SPIN.5}
\]
is recorded and used in
\cite[proof of Proposition~3.4, p.~626]{HD:GT2012}.
It concerns faithful complex spin characters of the double
alternating cover, not ordinary representations of $\mathfrak A_n$.

Every nonscalar element of $J$ has $\nu(g)\ge e/8$.
For an element preserving each tensor factor, this follows from
HD-SPIN-01a on a nonscalar factor and the tensor eigenspace bound:
a largest eigenspace multiplicity in a tensor product is at most
that in one factor times the remaining dimensions.  The factor
maps have finite order up to scalars, as follows by taking a power
of the tensor operator.
For an element permuting factors, the weighted cyclic-word
argument of D-SPOR-03, equation (SPOR.6), instead gives
$\nu(g)\ge e/4$.  That argument uses only the dimensions of
the factors, not any sporadic group hypothesis.  The intrinsic
factor-algebra normalizer argument in the same proof applies
unchanged to (SPIN.4).  Thus the weaker bound $e/8$ holds for
every nonscalar element.  Every genuine scalar inertia moves all
$e$ dimensions.

The actual intermediate-cover argument of D-BI-09 gives
$r\ge12$ genuine inertia elements with total moving rank
$2e+6$.  The essentially self-dual similitude character is
nontrivial because the inducing character field contains the
original imaginary quadratic field.  At least two inertias have
nontrivial multiplier, and each moves at least $\lceil e/2\rceil$
dimensions, by the isotropic-fixed-space proof in D-BI-09.
The other at least ten inertias move at least $\lceil e/8\rceil$.
Therefore
\[
       2\lceil e/2\rceil+10\lceil e/8\rceil\le2e+6,
       \qquad e\le24.                                   \tag{SPIN.6}
\]
The second inequality follows since the left side is at least
$9e/4$.  D-BI-09, with its complete primitive and tensor cases,
gives $e\ge24$.  Hence $e=24$.

There can be only one factor in (SPIN.4), since two factors
already give $e\ge8^2=64$.  Equation (SPIN.5) also forces
$8\le n\le11$.  The complete faithful-on-center character rows
of the published tables $2.A8,2.A9,2.A10,2.A11$ have the
following degree-twenty-four data:
\[
\begin{array}{c|c|c}
\text{table}&\text{degree-24 rows}&\text{self-duality}\\ \hline
2.A8&16,\ 17&\text{both nonreal}\\
2.A9&\varnothing&---\\
2.A10&\varnothing&---\\
2.A11&\varnothing&---
\end{array}                                             \tag{SPIN.7}
\]
Here rows are the exact CTblLib identifiers in the retained
certificate, not an assumption that all spin degrees are powers
of two.  In $2.A8$ the two rows are complex conjugates, each
with Frobenius--Schur indicator zero.  Their values include
$b_7=\zeta_7+\zeta_7^2+\zeta_7^4$ and $\overline b_7$;
the character field is $\mathbf Q(\sqrt{-7})$.

Finally, the multiplier of an essentially self-dual $J$-module
is a one-dimensional character.  Its restriction to the perfect
group $E$ is trivial.  The invariant nondegenerate bilinear form
would therefore make the irreducible $E$-module self-dual.
Both actual degree-twenty-four rows in (SPIN.7) contradict this.
This finishes the exclusion.
\end{proof}

\paragraph{Certificate and scope.}
The complete five small-cover tables, their source strings and
power maps are retained alongside the alternating-cover
certificate in \texttt{agents/imprimitive/sporadic/}.
They are published ATLAS inputs accessed through the exact
identifiers of \cite{SPOR:CTblLib}.  The numerical verification
checks (SPIN.2), the full faithful degree lists, the degree-24
exception and its nonreal witnesses; it does not claim to prove
the general spin-degree lower bound or the identification of
the tables with the named groups.  The latter premises retain
their explicit published locators above.
The theorem leaves nonselfdual spin cores, all ordinary
alternating cores, and the small covers in degrees $n<8$
outside its scope.

\section{The full spin-core exclusion}
\label{sec:spin-quarter-obstruction}

The complete local spectra sharpen HD-SPIN-01 and remove both
its self-duality and its nonsplit-polarization hypotheses.

\begin{lemma}[HD-SPIN-02a: the local half-eigenspace bound]
\label{HD-SPIN-02a}
Let $U=2.\mathfrak A_6$ and let $W$ be any complex representation
on which the central involution acts as $-1$.  For every
$h\in U\setminus Z(U)$ and every eigenvalue $\lambda$,
\[
                   \dim\ker(h-\lambda)\le\tfrac12\dim W.
                                                               \tag{SPINQ.1}
\]
\end{lemma}

\begin{proof}
It suffices to check each irreducible constituent, since the
multiplicities of a fixed eigenvalue add.  There are six
center-negative rows in the complete published table $2.A6$.
The following table gives their degrees and largest eigenvalue
multiplicities on every noncentral class.  The columns are the
exact class indices of the retained CTblLib input; the order row
prevents a class-name convention from being tacitly substituted.
\[
\begin{array}{c|r|rrrrrrrrrrr}
\text{row}&\deg&3&4&5&6&7&8&9&10&11&12&13\\\hline
\text{order}&&4&3&6&3&6&8&8&5&10&5&10\\\hline
8 &4 &2&2&2&2&2&1&1&1&1&1&1\\
9 &4 &2&2&2&2&2&1&1&1&1&1&1\\
10&8 &4&3&3&3&3&2&2&2&2&2&2\\
11&8 &4&3&3&3&3&2&2&2&2&2&2\\
12&10&5&4&4&4&4&3&3&2&2&2&2\\
13&10&5&4&4&4&4&3&3&2&2&2&2
\end{array}                                                   \tag{SPINQ.2}
\]
Every entry is at most half its row degree.  For completeness,
the exact calculation is the cyclic Fourier formula: if $h$ has
order $m$ and $\zeta_m$ is a primitive $m$th root, then the
multiplicity of $\zeta_m^k$ is
\[
       \frac1m\sum_{r=0}^{m-1}\chi(h^r)\zeta_m^{-kr}.           \tag{SPINQ.3}
\]
The retained complete character table and power maps supply all
terms.  Prime power maps determine $h^d$ for $d=\gcd(r,m)$;
the remaining coprime power is the corresponding cyclotomic
Galois automorphism on its character value.  Reducing the sum
in $\mathbf Z[x]/(\Phi_m)$ gives a constant divisible by $m$.
The exact script checks nonnegativity, integrality and sum equal
to the degree for each of the $66$ spectra, then gives (SPINQ.2).
It selects the six rows from their actual central-involution
values, rather than assuming a supplied row-selection list.
The full published table is an explicit cited input
\cite{SPOR:CTblLib}; this calculation does not purport to reconstruct
the table or identify an unknown group from numerical orthogonality.
\end{proof}

\begin{theorem}[HD-SPIN-02b: all faithful spin normal cores]
\label{HD-SPIN-02b}
Let a finite-cover sixfold source have an inducing image $J$
with an irreducibly acting normal central product
\[
                  E=E_1\cdots E_t,\qquad
                  E_i\simeq2.\mathfrak A_{n_i},\quad n_i\ge8,
                                                               \tag{SPINQ.4}
\]
where the factor representations are faithful.  Then the actual
six-dimensional multiplicity variation has period-map rank at
most eight.  This holds for either polarization determinant class,
any inducing character field, and all compatible outer
automorphisms and permutations of tensor factors.
\end{theorem}

\begin{proof}
First take a single factor $N=2.\mathfrak A_n$ on a faithful
irreducible space $W$ of dimension $a$, and a nonscalar element
$g$ of its linear normalizer with finite order up to scalar.
The complete group-theoretic argument of HD-SPIN-01a gives a
lift $t$ of a three-cycle for which $h=[g,t]$ is noncentral in
a natural $2.\mathfrak A_6$ subgroup.  The central involution
acts as $-1$ on every restriction constituent.  HD-SPIN-02a
therefore gives $\nu(h)\ge a/2$.  The same exact commutator
rank inequality, with $g$ rescaled at a largest eigenvalue, gives
\[
                  a/2\le\nu(h)\le2\nu(g),
                  \qquad \nu(g)\ge a/4.                    \tag{SPINQ.5}
\]
The restriction can be reducible: (SPINQ.1) holds for each fixed
eigenvalue across all constituents, which is precisely what is
needed here.

Write $V=\bigotimes_iV_i$ for (SPINQ.4).  Every $a_i=\dim V_i$
is at least eight by the published general spin minimum used in
(SPIN.5).  The intrinsic factor-algebra and tensor-normalizer
argument in D-SPOR-03 applies.  A factor-preserving nonscalar
$g$ satisfies $\nu(g)\ge\dim V/4$ by (SPINQ.5) and the tensor
eigenspace inequality.  A nontrivial factor permutation satisfies
the same bound by the weighted cyclic-word count (SPOR.6).
Consequently every nonscalar element satisfies this quarter bound,
and every genuine scalar inertia moves the entire space.

Put $e=\dim V\ge8$.  The intermediate-cover and period-factorization
argument in D-SPOR-03 applies without any sporadic hypothesis.
Explicitly, if its base has genus $h$, the actual six-dimensional
multiplicity cohomology gives
\[
          6=2e(h-1)+\sum_j\operatorname{rank}(1-g_j).
\]
For $h\ge2$ this is impossible.  For $h=1$ there are at most six
genuine branch points and its pointed-curve moduli, hence period
image, has dimension at most six.  For $h=0$, period rank nine
would require at least twelve genuine branch points.  The
quarter bound would then imply
\[
          3e\le12\lceil e/4\rceil
             \le\sum_j\operatorname{rank}(1-g_j)=2e+6,
\]
so $e\le6$, contrary to $e\ge8$.  Therefore rank nine is
impossible, as asserted.
\end{proof}

\paragraph{Checks and precise remaining scope.}
The complete local spectra and source-input hash are retained in
\texttt{research/spin\_quarter\_certificate.json}; its integer-polynomial
implementation is separate from the peer's Ramanujan-trace
implementation.  The theorem supersedes the ESD restriction in
HD-SPIN-01b.  Ordinary alternating representations and exceptional
central covers in alternating degrees below eight remain outside
this result.  No new detector cycle, coefficient-field descent,
or actual dominant component is inferred from this exclusion.

\section{An ordinary alternating-core shape restriction}
\label{sec:ordinary-alternating-shapes}

Spin representations have been removed by HD-SPIN-02.  Ordinary
alternating representations require a different local estimate:
restriction to the six-letter subgroup can contain trivial summands.
We retain their contribution explicitly.

\begin{lemma}[HD-ALT-01a: the six-letter fixed-space estimate]
\label{HD-ALT-01a}
Let $N=\mathfrak A_n$, $n\ge8$, act faithfully and irreducibly
on $W$, with character $\chi$ and degree $a$.  Let $c$ be a
three-cycle and put $u=\chi(c)/a$.  Every nonscalar element $g$
of the linear normalizer of $N$, of finite order up to a scalar,
satisfies
\[
                \nu(g)\ge\frac{2a(1-u)}{15}.              \tag{ALT.1}
\]
Here $u$ is real and belongs to $[-1/2,1)$.
\end{lemma}

\begin{proof}
The complete ordinary character table of $\mathfrak A_6$ gives
\[
       \dim\ker(h-\lambda)\le\frac35\dim U              \tag{ALT.2}
\]
for every nontrivial irreducible $U$, every $h\ne1$ and every
$\lambda$.  We record the entire finite scope of this input.
The six nontrivial center-positive rows of the retained table
$2.A6$, which factor through $A6$, have degrees
$5,5,8,8,9,10$.  The minima of $\nu(h)/\dim U$ across all
noncentral lift classes are respectively
\[
                   2/5,\ 2/5,\ 1/2,\ 1/2,\ 4/9,\ 2/5.
                                                               \tag{ALT.3}
\]
The exact Fourier formula (SPINQ.3) checks all $66$ spectra,
including every possible eigenvalue.  The full spectra,
nonnegative integral multiplicities and degree sums are retained
in \texttt{research/alternating\_local\_certificate.json}.
The character table and identification with the named group
are published inputs \cite{SPOR:CTblLib}; the finite calculation
proves the claimed numerical inequality on those complete inputs.

As in HD-SPIN-01a, the normalizer modulo scalars embeds in
$\operatorname{Aut}(\mathfrak A_n)=\mathfrak S_n$.
For a nonscalar $g$ choose a three-cycle $c'$ not centralized by
its induced permutation, and put $h=[g,c']$.  This is nonidentity
and supported on at most six letters.  Embed its support in a
six-letter set and let $A6$ be the corresponding natural subgroup.
All such subgroups are conjugate in $\mathfrak A_n$ for $n\ge8$:
an odd conjugating permutation can be corrected by a transposition
on two letters outside the target six-set.  Thus the dimension
$d_0=\dim W^{A6}$ is independent of this choice.

On restriction, write $W=\mathbf1^{d_0}\oplus W'$, with no
trivial constituent in $W'$.  For every fixed eigenvalue,
(ALT.2) adds across its irreducible constituents.  The trivial
part contributes at most $d_0$, whence
\[
                    \nu(h)\ge\frac25(a-d_0).
\]
The scalar-adjusted commutator inequality gives
$\nu(h)\le2\nu(g)$, and therefore
$\nu(g)\ge(a-d_0)/5$.

The $A6$-invariants are contained in the fixed space of any
three-cycle in that subgroup.  A three-cycle is conjugate to
its inverse in $\mathfrak A_n$ (use a transposition on its
support and another on two fixed letters).  Consequently
$\chi(c)$ is real, and its two nontrivial cube-root eigenvalues
have equal multiplicity.  Its fixed dimension is
$(a+2\chi(c))/3$, so
\[
              d_0\le\frac{a+2\chi(c)}3.
\]
Combining the inequalities proves (ALT.1).  Nonnegativity of the
fixed dimension gives $u\ge-1/2$, and faithfulness gives $u<1$.
\end{proof}

\begin{lemma}[HD-ALT-01b: the Young-diagram bound]
\label{HD-ALT-01b}
Let $n\ge15$ and let $\chi$ be an irreducible ordinary character
of $\mathfrak A_n$ occurring in the restriction of the Specht
character indexed by $\lambda\vdash n$.  If
\[
               \max(\lambda_1,\lambda'_1)\le\lceil n/2\rceil,
                                                               \tag{ALT.4}
\]
then its three-cycle ratio satisfies
\[
             u\le\frac{n-5}{4(n-2)}
               =\frac14-\frac3{4(n-2)}.                  \tag{ALT.5}
\]
For even $n$ the sharper numerator is $n-7$.
\end{lemma}

\begin{proof}
We use the published Jucys--Murphy content-eigenvalue theorem
\cite[Theorem~4.2; author version p.~9]{HD:CMN2022}.
For $J_i=\sum_{j<i}(j\ i)$, direct multiplication in
$\mathbf C[\mathfrak S_n]$ gives
\[
       \sum_iJ_i^2=\binom n2\,1+
                  \sum_{c\text{ a three-cycle}}c.
\]
Indeed the equal-transposition products give $i-1$ identities,
and the ordered unequal pairs give each oriented three-cycle
exactly once, indexed by its largest letter.  The content theorem
therefore gives the following short specialization:
\[
 u=\frac{3\left(\sum_{(i,j)\in\lambda}(j-i)^2-\binom n2\right)}
           {n(n-1)(n-2)}.                                \tag{ALT.6}
\]
This also applies if the Specht restriction splits.  The two
constituents have equal degrees and are interchanged by odd
conjugation.  The three-cycle class does not split in
$\mathfrak A_n$, since a transposition on two fixed letters
centralizes the cycle.  Their character values on it are equal,
so their normalized ratio is the Specht ratio in (ALT.6).
The restriction rule itself is the standard one explicitly
recorded in \cite[Section~2A, p.~604]{HD:BNOT2015}.

Put $k=\lceil n/2\rceil$ and use Frobenius coordinates
$(a_1,\ldots,a_d\mid b_1,\ldots,b_d)$ of $\lambda$.
The diagonal hooks partition its boxes, with contents
$-b_i,\ldots,-1,0,1,\ldots,a_i$ on the $i$th hook.
Condition (ALT.4) gives $0\le a_i,b_i\le k-1$ and
$\sum_i(a_i+b_i)=n-d\le n-1$.  For $0\le s\le k-1$,
\[
 \sum_{j=1}^s j^2=\frac{s(s+1)(2s+1)}6
                  \le\frac{s\,k(2k-1)}6.
\]
Hence the content-square sum is at most
$k(2k-1)(n-1)/6$.  Substitution in (ALT.6) gives
\[
                u\le\frac{k(2k-1)-3n}{2n(n-2)}.
\]
For even $n$ this equals $(n-7)/(4(n-2))$, and for odd $n$
it equals $(n-5)/(4(n-2))$.  This proves both assertions.
\end{proof}

\begin{theorem}[HD-ALT-01c: exclusion of the indicated ordinary shapes]
\label{HD-ALT-01c}
Suppose a primitive inducing block of a dominant nonsplit Weil
sixfold source is essentially self-dual and has an irreducibly
acting normal core $E=\mathfrak A_n^t$, with $n\ge15$.
Suppose every irreducible tensor factor has a Specht label
satisfying (ALT.4).  Then these data cannot occur.

In particular, for the transitive ordinary alternating core of
D-BI-19, a remaining essentially self-dual source in alternating
degree $n\ge15$ must have
$\max(\lambda_1,\lambda'_1)>\lceil n/2\rceil$.
All self-dual constituents of self-conjugate labels are excluded
in this range.  Nonselfdual constituents are not covered.
\end{theorem}

\begin{proof}
Let $e$ be the total inducing degree.  Lemmas HD-ALT-01a/b give,
for each factor-preserving nonscalar normalizer element,
\[
        \nu(g)\ge\frac{n-1}{10(n-2)}e.                   \tag{ALT.7}
\]
The same conclusion holds for nontrivial factor permutations.
Indeed the intrinsic factor-algebra argument and the weighted
cyclic-word estimate (SPOR.6) give $\nu(g)\ge e/4$ in that
case, and $(n-1)/(10(n-2))<1/4$.  Genuine scalar inertia moves
all $e$ dimensions.  This argument permits every compatible
outer automorphism and all tensor permutations.

The actual intermediate-curve and nontrivial-similitude argument
of D-BI-09 gives at least twelve genuine inertias, of which at
least two move at least $\lceil e/2\rceil$ dimensions.  Its
hypotheses hold here: the source has nonsplit actual polarization,
exact imaginary quadratic original character field, and a
primitive essentially self-dual inducing block; its inducing
character field may be larger.  The other ten inertias obey
(ALT.7).  The continuous consequence of the exact integer
budget is
\[
 e+\frac{n-1}{n-2}e\le2e+6,
                    \qquad e\le6(n-2).                  \tag{ALT.8}
\]
The factor degrees are at least $n-1$ in this range by
\cite[Lemmas~3.1--3.2, p.~607]{HD:BNOT2015}; in particular
$e\ge9$, as required in D-BI-09.

Choose any one factor character and its label $\lambda$.
The degree of the corresponding $\mathfrak S_n$ character
is either its factor degree or twice that degree, so by (ALT.8)
it is at most $12(n-2)$.  For $n\ge15$,
\[
       12(n-2)<\frac{2n^3}{27}.
\]
To see this uniformly, the difference after multiplication by
$27$ is $2n^3-324n+648$; at $n=15$ it is $2538$, and its forward
difference is $6n^2+6n-322>0$ for $n\ge15$.
The published precise shape consequence of Rasala's classification
\cite[proof of Theorem~9.7, p.~312]{HD:GKT2019}, stated there
for every $\mathfrak S_n$ character of degree
$2\le d<2n^3/27$, therefore forces $\lambda$, up to transpose,
to be one of
\[
             (n-1,1),\qquad(n-2,2),\qquad(n-2,1,1).
\]
Each of these shapes has a row or column of length at least
$n-2>\lceil n/2\rceil$, contrary to (ALT.4).

Finally, if a partition is self-conjugate then
$2\lambda_1-1\le n$, so (ALT.4) holds automatically.
The preceding conclusion thus excludes its self-dual constituents.
For the transitive core of D-BI-19, factor labels agree up to the
outer action and transpose, which preserves (ALT.4); hence the
stated necessary shape restriction applies to that gate.
\end{proof}

\paragraph{Scope.}
The new local calculation and shape inequality do not prove
an all-alternating exclusion.  Long first-row or first-column
shapes and nonselfdual split characters remain.  The already
excluded standard, exterior-square and unordered-pair families
retain their stronger independent all-degree statements.

\section{Removing self-duality from the ordinary alternating shape restriction}
\label{sec:ordinary-alternating-sharp}

\paragraph{Status.}
Complete new root derivation with author whole-proof reading and an
independent whole-proof review by /root/g34\_eightfolds, both PASS
on 21 November 2026. The independent symbolic and permutation
certificate also passed. The finite checks below do not substitute
for the all-degree argument. No new finite character table is a premise.

\begin{lemma}[HD-ALT-03a: three-cycle and double-transposition commutators]
\label{HD-ALT-03a}
Let $n\ge15$ and let $W$ be a nontrivial irreducible ordinary
$A_n$-module of degree $a$. Put
\[
 u=\frac{\chi((123))}{a},\qquad
 v=\frac{\chi((12)(34))}{a}.
\]
For every nonscalar finite-order linear normalizer $T$ of the
represented core,
\[
 \nu(T)\ge a\min\left\{\frac{1-u}{3},\frac{1-v}{4}\right\}.
                                                               \tag{ALTN.1}
\]
All scalar phases are allowed. The conclusion includes both
constituents of a split Specht restriction, whether real or nonreal.
\end{lemma}

\begin{proof}
First suppose $W$ is the irreducible restriction of a Specht
module $S^\lambda$ of $S_n$. Every normalizer is scalar times
a matrix in this symmetric extension: its projective action lies
in $\operatorname{Aut}(A_n)=S_n$, and removing the implementing
symmetric matrix leaves a scalar by Schur's lemma. Let $g\in S_n$
be the implementing permutation, which is nonidentity.

If $g$ has a cycle of length at least three, choose two consecutive
letters $b,g(b)$ in that cycle and set $x=(b\ g(b))$. Its conjugate
by $g$ shares exactly one letter with $x$, so $[g,x]$ is a
three-cycle. If all nontrivial cycles of $g$ have length two and
$g$ has a fixed letter $c$, use $x=(b\ c)$ with $g(b)\ne b$;
again the commutator is a three-cycle. The only remaining case
is a fixed-point-free involution. Choose letters $b,c$ in two
different two-cycles and set $x=(b\ c)$. Its conjugate is the
disjoint transposition $(g(b)\ g(c))$, so $[g,x]$ is a double
transposition.

For any invertible matrices $A,B$ and any scalar $z$,
\[
 \operatorname{rank}([A,B]-1)
 =\operatorname{rank}(AB-BA)
 \le2\operatorname{rank}(A-z).
                                                               \tag{ALTN.2}
\]
Indeed $AB-BA=(A-z)B-B(A-z)$. Minimize over $z$ to get a bound
by $2\nu(A)$. A three-cycle has its two nontrivial eigenvalues
with equal multiplicity, since it is conjugate to its inverse
inside $A_n$. Its actual moving rank is $2a(1-u)/3$.
The moving rank of a double transposition is $a(1-v)/2$.
Applying (ALTN.2) proves (ALTN.1) in this case.

Now let $\lambda$ be self-conjugate, so
$S^\lambda|_{A_n}=W\oplus W'$ with two distinct constituents
of equal degree. We use the standard restriction rule
\cite[\S2A, p.~604]{HD:BNOT2015}. An odd permutation exchanges
them. Thus no odd automorphism is implemented on $W$, and every
normalizer is scalar times a represented element $g\in A_n$.

Suppose first that an odd permutation centralizes $g$. Its
operator on $S^\lambda$ exchanges $W$ and $W'$ and commutes
with $g$. Hence the two complete spectra of $g$ are identical,
and the largest-eigenspace codimension on $S^\lambda$ is
exactly $2\nu(g|W)$. Both the three-cycle and double-transposition
classes have odd centralizers: for the former use two fixed
letters, and for the latter use one of its transpositions.
Their normalized characters on $S^\lambda$ therefore equal
$u,v$. The preceding symmetric-module argument, in dimension
$2a$, now gives (ALTN.1) on $W$ after dividing by two.
This comparison is of full spectra, not merely the trace of $g$.

Finally suppose $g$ has no odd centralizer. It has no even
cycle, since such a cycle itself is odd and centralizes $g$.
It also has no two cycles of the same odd length: swapping
those two cycles letter by letter is an odd centralizer.
Thus its cycle lengths are distinct and odd. Since $n\ge15$,
some cycle has length at least five. In that cycle set
\[
 b_0=b,\quad b_1=g(b),\quad b_2=g^2(b),\quad b_3=g^3(b),
 \qquad x=(b_0\ b_1\ b_2)\in A_n.
\]
Direct multiplication gives
$[g,x]=(b_1\ b_2\ b_3)(b_2\ b_1\ b_0)
=(b_0\ b_3\ b_1)$, a three-cycle. Both matrices now act on
$W$ itself, so (ALTN.2) gives the stronger first entry of
(ALTN.1). This exhausts every split-constituent case.
\end{proof}

\begin{lemma}[HD-ALT-03b: the double-transposition ratio]
\label{HD-ALT-03b}
Let $n\ge15$ and let $\lambda\vdash n$ satisfy
\[
             \max(\lambda_1,\lambda'_1)\le\lceil n/2\rceil.
                                                               \tag{ALTN.3}
\]
For either irreducible constituent of its restriction to $A_n$,
\[
 u<\frac14,\qquad v<\frac3{10},\qquad
                 \nu(T)\ge\frac7{40}a                 \tag{ALTN.4}
\]
for every nonscalar linear normalizer element.
\end{lemma}

\begin{proof}
The bound on $u$ is HD-ALT-01b. Write
\[
 t=\sum_{(i,j)\in\lambda}(j-i),\quad
 S=\sum_{(i,j)\in\lambda}(j-i)^2,\quad
 k=\lceil n/2\rceil.
\]
Let $C$ be the sum of all transpositions in $\mathbf C[S_n]$.
It is $\sum_iJ_i$ and acts on $S^\lambda$ by the scalar $t$,
by the same published content-eigenvalue theorem
\cite[Theorem~4.2; author p.~9]{HD:CMN2022} used in HD-ALT-01b.
In the group algebra,
\[
 C^2=\binom n2\,1+3K_3+2K_{22},
                                                               \tag{ALTN.5}
\]
where $K_3$ and $K_{22}$ sum all three-cycles and all double
transpositions, respectively. The coefficients count the ordered
factorizations into two transpositions. Their class sizes are
$n(n-1)(n-2)/3$ and $n(n-1)(n-2)(n-3)/8$.
Using (ALT.6) in (ALTN.5) gives the exact specialization
\[
 v=\frac{4\bigl(t^2+n(n-1)-3S\bigr)}
           {n(n-1)(n-2)(n-3)}.                            \tag{ALTN.6}
\]
The nonsplitting argument in HD-ALT-03a shows that the same ratio
holds for either constituent when the restriction splits.

The elementary row/column expression for $t$ is
\[
 t=\sum_i\binom{\lambda_i}{2}
                   -\sum_j\binom{\lambda'_j}{2}.
\]
Since every row is at most $k$, the first sum is at most
$n(k-1)/2$. For every positive integer $h$, $\binom h2\ge h-1$;
therefore the second sum is at least $n-\lambda_1\ge n-k$.
Applying the same argument to the transpose gives
\[
 |t|\le\frac{n(k-1)}2-(n-k)
                     \le\frac{(n-1)(n-2)}4.              \tag{ALTN.7}
\]
For odd $n$ the final inequality is equality of the displayed
upper bounds; for even $n$ it is strict. Dropping the nonpositive
term $-3S$ from (ALTN.6) now gives
\[
 v\le\frac{(n-1)(n-2)^2+16n}{4n(n-2)(n-3)}<\frac3{10}.
                                                               \tag{ALTN.8}
\]
For the final strict inequality, after multiplication by the
positive denominator the difference is
$2(n^3-5n^2-84n+20)$. At $n=15$ the polynomial in parentheses
is $1010$, and its forward difference is
$3n^2-7n-88>0$ for $n\ge15$.
Thus the first term in (ALTN.1) is greater than $1/4$ and the
second greater than $7/40$. The stated weak bound follows.
\end{proof}

\begin{theorem}[HD-ALT-03c: all balanced ordinary shapes are excluded]
\label{HD-ALT-03c}
Suppose an inducing block of an actual rank-nine Weil sixfold
finite-cover source has an irreducibly acting normal central-product
core of ordinary alternating groups $A_n$, with $n\ge15$,
and every factor character has a Specht label satisfying (ALTN.3).
Then these data cannot occur. The assertion includes all scalar
phases, compatible outer automorphisms and factor permutations.
It does not require essential self-duality or nonsplit polarization.

In the transitive forced-core setting D-BI-19, every remaining
ordinary alternating source with $n\ge15$ therefore has a row or
column longer than $\lceil n/2\rceil$. Its factor restriction
is irreducible and extends to a real Specht representation of $S_n$;
in particular the full inducing image is essentially self-dual.
No existence or dominance is asserted for these remaining long shapes.
\end{theorem}

\begin{proof}
Let $e$ be the total inducing degree. For a factor-preserving
nonscalar normalizer, (ALTN.4) passes to the tensor product by
fixing all other eigenvalues. A nontrivial factor permutation
has the stronger quarter bound by the weighted tensor-cycle
argument used in D-SPOR-03. Genuine scalar inertia moves all
$e$ dimensions. Thus every genuine inertia moves at least $7e/40$.

Each factor has degree at least $n^3/27\ge125$. Indeed the precise
published consequence of Rasala's classification in
\cite[proof of Theorem~9.7, p.~312]{HD:GKT2019} says that a
nonlinear Specht character of degree below $2n^3/27$ has, up to
transpose, one of the three shapes $(n-1,1),(n-2,2),(n-2,1,1)$.
None satisfies (ALTN.3). Its degree is either the factor degree
or twice that degree, proving the stated lower bound.

The actual intermediate-cover Euler identity and period
factorization give genus zero, at least twelve genuine inertias,
and total moving rank $2e+6$. Genus one has at most six branches
and period-parameter dimension at most six; genus at least two
would make $6+2e(1-h)<0$, since $e\ge125$. Consequently
\[
            12\frac{7e}{40}\le2e+6,
                         \qquad e\le60,                  \tag{ALTN.9}
\]
contrary to $e\ge125$.

For the final structural conclusion, a self-conjugate partition
has $2\lambda_1-1\le n$, so it satisfies (ALTN.3). Such partitions,
including both nonselfdual and selfdual split constituents, have
therefore been excluded. A remaining partition is not self-conjugate;
its restriction to $A_n$ is the restriction of an irreducible
real Specht representation of $S_n$. The transitive tensor factors
are thus self-dual, and the irreducible full core has a unique
invariant line of nondegenerate bilinear forms. Every normalizer
preserves that line, proving essential self-duality. This conclusion
is a necessary representation restriction, not a new geometric source.
\end{proof}

\paragraph{Verification scope.}
The short group-algebra identities, scalar-normalizer argument,
and split-constituent comparison are proved above. The retained
exact script checks the polynomial bounds, the commutator cycle
identities, and all Young-diagram content and ratio formulas in a
finite range. The all-$n$ and all-tensor-exponent statements use
the inequalities in the proof, not extrapolation. The published
content theorem, Specht restriction rule, automorphism theorem and
precise small-degree classification remain cited inputs.

\section{A three-quarter shape restriction for ordinary alternating cores}
\label{sec:ordinary-alternating-three-quarters}

\paragraph{Status.}
New root derivation, with an author whole-proof reading and an
independent whole-proof review by /root/g34\_eightfolds, both PASS
on 21 November 2026. The independent arithmetic reconstruction passed. The exact finite certificate checks the
formulas and endpoints; the proof below supplies the all-degree
statement. The published content and low-degree inputs are exactly
those already read and cited in HD-ALT-01--03.

\begin{lemma}[HD-ALT-04a: a Cauchy improvement of the content bound]
\label{HD-ALT-04a}
Let $n\ge15$ and let $\lambda\vdash n$ satisfy
\[
                k:=\max(\lambda_1,\lambda'_1)\le3n/4.
                                                               \tag{ALTQ.1}
\]
For either irreducible constituent $W$ of $S^\lambda|_{A_n}$,
of degree $a$, put $u=\chi((123))/a$ and
$v=\chi((12)(34))/a$. Then
\[
                u<9/16,\qquad v<9/16.
                                                               \tag{ALTQ.2}
\]
Every nonscalar finite-order linear normalizer of the represented
core satisfies $\nu(T)\ge7a/64$, with every scalar phase allowed.
\end{lemma}

\begin{proof}
Write $t=\sum_{(i,j)\in\lambda}(j-i)$ and
$S=\sum_{(i,j)\in\lambda}(j-i)^2$. The exact specializations
(ALT.6) and (ALTN.6), including their split-constituent justification,
give
\[
 u=\frac{3(S-\binom n2)}{n(n-1)(n-2)},\qquad
 v=\frac{4(t^2+n(n-1)-3S)}{n(n-1)(n-2)(n-3)}.
                                                               \tag{ALTQ.3}
\]
The diagonal-hook estimate in HD-ALT-01b holds for any upper
bound $k$ on both the first row and the first column:
\[
 S\le\frac{k(2k-1)(n-1)}6.
\]
Its proof uses only that every arm and leg has length at most
$k-1$ and that the sum of their lengths is at most $n-1$.
Thus, by (ALTQ.1),
\[
 u\le\frac{k(2k-1)-3n}{2n(n-2)}
    \le\frac{9n-30}{16(n-2)}<\frac9{16}.
                                                               \tag{ALTQ.4}
\]

The row/column estimate (ALTN.7), before its specialization of $k$,
and its transpose give
\[
 |t|\le\frac{n(k-1)}2-(n-k)
       \le\frac{3n(n-2)}8.
                                                               \tag{ALTQ.5}
\]
Apply Cauchy's inequality to the $n$ contents to obtain $S\ge t^2/n$.
Keeping this term, rather than discarding $-3S$ in (ALTQ.3), yields
\[
 \begin{split}
 v&\le\frac{4t^2}{n^2(n-1)(n-2)}
              +\frac4{(n-2)(n-3)}\\
  &\le\frac{9(n-2)}{16(n-1)}
              +\frac4{(n-2)(n-3)}<\frac9{16}.
 \end{split}                                                   \tag{ALTQ.6}
\]
The final strict inequality is equivalent to
$9n^2-109n+118>0$. Its value at $15$ is $508$,
and its forward difference is $18n-100>0$ for $n\ge15$.
This proves (ALTQ.2). The complete normalizer bound HD-ALT-03a
now gives
\[
 \nu(T)\ge a\min\{(1-u)/3,(1-v)/4\}\ge7a/64.
\]
No assertion about the reality of a split constituent is needed
for this lemma.
\end{proof}

\begin{theorem}[HD-ALT-04b: transitive ordinary cores require a longer shape]
\label{HD-ALT-04b}
Consider an inducing block of an actual rank-nine Weil sixfold
finite-cover source. Suppose its finite image $J$ has an irreducibly
acting normal central-product core with simple factors $A_n$,
$n\ge15$, and ordinary factor representations. Suppose that $J$
acts transitively on the factors and that its character field
contains the original imaginary quadratic field $K$.

Then every factor's Specht label satisfies
\[
                   \max(\lambda_1,\lambda'_1)>3n/4.
                                                               \tag{ALTQ.7}
\]
In particular this is a proper refinement of the remaining
ordinary alternating gate of D-BI-19 and HD-ALT-03.
It permits all scalar phases, compatible outer automorphisms
and tensor permutations. No additional nonsplit-polarization
or inducing-field-equality hypothesis is required.
\end{theorem}

\begin{proof}
Suppose (ALTQ.1) holds for one factor. Transitivity implies that
all factor labels agree up to transpose and the alternating
restriction ambiguity, so (ALTQ.1) holds for every factor.

Each factor has degree at least $n^3/27\ge125$.
Indeed the precise published consequence of Rasala's classification
in \cite[proof of Theorem~9.7, p.~312]{HD:GKT2019} says that any
nonlinear Specht character of degree below $2n^3/27$ has, up to
transpose, one of the three shapes $(n-1,1),(n-2,2),(n-2,1,1)$.
For $n\ge15$ each has a row of length at least $n-2>3n/4$.
They cannot occur here. Passing to an alternating constituent
loses a factor of at most two and proves the bound.

A self-conjugate factor label has first row and column at most
$\lceil n/2\rceil$. By transitivity all factors would then satisfy
the stronger shape hypothesis of HD-ALT-03c, which already
excludes these data. We can therefore assume that no factor
label is self-conjugate. Every factor is consequently the
irreducible restriction of a real Specht representation.
The tensor-product core has a unique invariant line of symmetric
nondegenerate bilinear forms. Every element of $J$ preserves
this line, so acts by a multiplier character
$\mu:J\longrightarrow\mathbf C^*$.

This character is nontrivial. Otherwise $J$ preserves the form,
its representation is self-dual, and its finite-group character
is real valued. Its character field would be a subfield of
$\mathbf R$, contrary to $K\subseteq\mathbf Q(\chi_J)$.

Let $e$ be the total inducing degree. The actual intermediate-cover
factorization used in D-BI-01 and HD-ALT-03c forces genus zero,
at least twelve genuine inertias, and total moving rank $2e+6$:
genus at least two is impossible for $e\ge125$, and genus one
has period-parameter dimension at most six. The genus-zero
inertias generate $J$ and have product one. Therefore at least
two have nontrivial multiplier. For either such inertia, its
fixed space is totally isotropic for the preserved form line,
so its ordinary moving rank is at least $\lceil e/2\rceil$.

Every other genuine inertia moves at least $7e/64$ dimensions.
For factor-preserving operators this follows from HD-ALT-04a
and the largest-eigenspace tensor bound. A nontrivial factor
permutation has the stronger quarter bound from the weighted
tensor-cycle argument in D-SPOR-03. A genuine scalar inertia
moves all $e$ dimensions. Thus the exact necessary budget implies
\[
 2\lceil e/2\rceil+10\lceil7e/64\rceil\le2e+6,
 \qquad e+70e/64\le2e+6,
 \qquad e\le64.
                                                               \tag{ALTQ.8}
\]
This contradicts $e\ge125$ and proves (ALTQ.7).
\end{proof}

\paragraph{Meaning and boundary.}
This excludes a strictly larger infinite family of ordinary
alternating local systems from dominant Weil sixfold suppliers.
For example the nonselfconjugate shape $(11,4)$ at $n=15$
has $11>\lceil15/2\rceil$ and $11\le45/4$; it lies outside
the earlier half-size region and inside (ALTQ.1).
The conclusion concerns the finite monodromy representation and
the actual period-map factorization. Shapes with a still longer
row or column remain open unless covered by the separate standard,
exterior-square or unordered-pair results. No new cycle is produced.
In the project's language this narrows the supplier gate for
condition A (cote~16, scan~050). It does not strengthen condition I
(scan~049 and cote~44, Theorem~4), change the Cayley--Hamilton
projector in scans~048--049, or supply a new assertion about
weights or motivic $L$-functions.

\section{A tail-content restriction for the remaining long ordinary shapes}
\label{sec:ordinary-alternating-long-row}

\paragraph{Status.}
New root proof, read in full by its author and independently by
/root/g34\_eightfolds; both whole-proof reviews PASS on 21 November
2026. Separate exact arithmetic certificates also passed. This is one further uniform shape reduction;
no claim that it settles all long shapes is made.

\begin{lemma}[HD-ALT-05a: exact tail-content estimates]
\label{HD-ALT-05a}
Let $n\ge15$ and let $\lambda=(n-k,\mu)$, where $\mu\vdash k$
and $0<k<n/4$. Put $t=\sum(j-i)$ and $S=\sum(j-i)^2$ over
the boxes of $\lambda$. Then
\begin{align}
 S&\le\sum_{j=1}^{n-k-1}j^2+\sum_{j=1}^k j^2,               \tag{ALTL.1}\\
 0<t&\le\binom n2-k(n-k+1).                                \tag{ALTL.2}
\end{align}
For the three-cycle ratio of its irreducible alternating restriction,
\[
 \frac{1-u}{3}\ge\frac{k(n-1-k)}{(n-1)(n-2)}.
                                                               \tag{ALTL.3}
\]
If also $k\ge n/7$, then its double-transposition ratio satisfies
\[
                         v<7/12.
                                                               \tag{ALTL.4}
\]
Every nonscalar finite-order linear normalizer in that range consequently
satisfies $\nu(T)\ge5a/48$, where $a=\dim S^\lambda$.
\end{lemma}

\begin{proof}
Choose a standard tableau of the tail diagram $\mu$, which orders its
boxes by successive valid Young-diagram additions. If its $r$th box
is $(i,j)$ in tail coordinates, then $i,j\le r$. Its content in
$\lambda$ is $j-i-1$, lying between $-r$ and $r-2$. Its square is
therefore at most $r^2$. Adding these bounds and the exact first-row
squares proves (ALTL.1). The same lower bound for each tail content
gives
\[
 t\ge\binom{n-k}{2}-\frac{k(k+1)}2
     =\frac{n(n-2k-1)}2>0.
\]
The tail content sum is its own ordinary content sum minus $k$.
The ordinary sum for a partition of $k$ is at most $\binom k2$:
write it as the sum of row binomials minus the sum of column
binomials, and use $\sum_i\binom{\mu_i}{2}\le\binom k2$.
Hence
\[
 t\le\binom{n-k}{2}+\binom k2-k
       =\binom n2-k(n-k+1),
\]
proving (ALTL.2). Substituting (ALTL.1) into the exact content
formula (ALTQ.3), then simplifying, gives (ALTL.3).

Now suppose $n/7\le k<n/4$. On this interval both
$k(n-1-k)$ and $k(n+1-k)$ increase with $k$. Thus
\[
 \frac{1-u}{3}\ge\frac{n(6n-7)}{49(n-1)(n-2)}
                 >\frac6{49},\qquad
 0<t\le\frac{n(37n-63)}{98}.
                                                               \tag{ALTL.5}
\]
The first strict inequality follows on subtraction because the
numerator difference is $11n-12>0$. Cauchy's inequality as in
(ALTQ.6) now gives
\[
 v\le\frac{(37n-63)^2}{2401(n-1)(n-2)}
                +\frac4{(n-2)(n-3)}<\frac7{12}.
                                                               \tag{ALTL.6}
\]
After multiplication by the positive denominator
$12\cdot2401(n-1)(n-2)(n-3)$, the final strict inequality
is equivalent to positivity of
\[
 P(n)=379n^3+4386n^2-145831n+157290.
\]
Writing $n=15+x$ gives
\[
 P(15+x)=379x^3+21441x^2+241574x+235800>0
                         \quad(x\ge0).
\]
This proves (ALTL.4) for every $n\ge15$. Since
$6/49>5/48$, the complete normalizer estimate HD-ALT-03a gives
$\nu(T)\ge a\min\{(1-u)/3,(1-v)/4\}\ge5a/48$.
The long first row prevents a self-conjugate label, so its
alternating restriction is irreducible and has degree $a$.
\end{proof}

\begin{theorem}[HD-ALT-05b: the remaining row or column exceeds six sevenths]
\label{HD-ALT-05b}
In the transitive ordinary alternating-core setting of HD-ALT-04b,
every factor label satisfies
\[
                    \max(\lambda_1,\lambda'_1)>6n/7.
                                                               \tag{ALTL.7}
\]
The same allowance for arbitrary compatible outers, scalar phases,
tensor permutations and larger inducing character fields applies.
\end{theorem}

\begin{proof}
By HD-ALT-04b any remaining label has a row or column longer
than $3n/4$. Transposing if necessary, write it as $(n-k,\mu)$,
with $0<k<n/4$. A nontrivial ordinary factor has $k>0$.
If (ALTL.7) failed, then $k\ge n/7$. Lemma HD-ALT-05a gives
the bound $5a/48$ on each nonscalar factor normalizer. It passes
to factor-preserving tensor operators, while every nontrivial
factor permutation has the stronger quarter bound. Genuine
nonidentity scalar inertia moves the whole space.

The real-factor and nontrivial-multiplier argument of HD-ALT-04b
still applies. At least two of the at least twelve genuine
inertias have ordinary moving rank at least $\lceil e/2\rceil$.
The exact total-moving-rank identity therefore implies
\[
 2\lceil e/2\rceil+10\lceil5e/48\rceil\le2e+6,
 \qquad \frac{e}{24}\le6,
 \qquad e\le144.
                                                               \tag{ALTL.8}
\]
Every factor here has degree at least $2n^3/27\ge250$.
Use the same published precise low-degree classification consequence
\cite[proof of Theorem~9.7, p.~312]{HD:GKT2019} as in HD-ALT-04.
All three exceptional shapes have a row or column of length
at least $n-2>6n/7$ for $n\ge15$, and so cannot occur.
There is no factor-of-two loss on restriction now: the row
longer than $3n/4$ ensures the partition is not self-conjugate.
Thus $e\ge250$, contradicting (ALTL.8).
\end{proof}

\paragraph{Strictness and geometric meaning.}
At $n=15$ the shape $(12,3)$ lies outside the three-quarter
region and inside the newly excluded range, so this is a strict
further restriction. The result again concerns actual finite
local systems and the factorization of the Weil period map.
It narrows the remaining supplier gate for condition A, cote~16
scan~050, without producing a detecting cycle or altering the
condition I/projector distinctions in scan~049, cote~44 Theorem~4,
and scans~048--049. No new assertion about motivic $L$-values
or weights is inferred. The still longer shapes remain open
except where the separate standard, exterior-square and
unordered-pair exclusions apply.

\section{Keeping the first-row squares sharpens the alternating gate}
\label{sec:ordinary-alternating-87}

\paragraph{Status and scope.}
New root derivation, whole-read by the author and independently
reviewed in full by /root/imprimitive. Both mathematical reviews pass;
independent exact rational and Young-diagram checks also pass.
It uses the already proved exact content identities and full normalizer
lemma of HD-ALT-03--05. No new asymptotic character estimate is assumed.
The conclusion retains precisely the transitive ordinary-core and actual
inducing-field hypotheses of HD-ALT-04b.

\begin{lemma}[HD-ALT-06a: a sharper tail estimate]
\label{HD-ALT-06a}
For $n\ge15$ and $\lambda=(n-k,\mu)$ with $\mu\vdash k$ and
$1\le k<n/4$, put
\[
 T(n,k)=\binom n2-k(n-k+1),\qquad
 S_0(n,k)=\frac{(n-k-1)(n-k)(2n-2k-1)}6,
\]
\[
 V(n,k)=\frac{4\{T(n,k)^2+n(n-1)-3S_0(n,k)\}}
                    {n(n-1)(n-2)(n-3)}.
                                                               \tag{ALT87.1}
\]
For the ordinary irreducible alternating restriction of $S^\lambda$,
with degree $a$, its three-cycle and double-transposition ratios obey
\[
 \frac{1-u}{3}\ge\frac{k(n-1-k)}{(n-1)(n-2)},\qquad
                         v\le V(n,k).                         \tag{ALT87.2}
\]
Consequently every nonscalar finite-order linear normalizer has
projective support at least $a\delta(n,k)$, where
\[
 \delta(n,k)=\min\left\{\frac14,
       \frac{k(n-1-k)}{(n-1)(n-2)},\frac{1-V(n,k)}4\right\}.
                                                               \tag{ALT87.3}
\]
On the real interval $1\le k<n/4$, the second entry of this minimum
is increasing and $V(n,k)$ is strictly decreasing.
\end{lemma}

\begin{proof}
Write $t$ and $S$ for the sum of contents and their squares.
HD-ALT-05a gives $0<t\le T(n,k)$ and the first inequality in
(ALT87.2). The first row alone contributes exactly $S_0(n,k)$
to $S$, so $S\ge S_0(n,k)$. Keeping that contribution in the
exact identity (ALTQ.3) gives
\[
 v=\frac{4\{t^2+n(n-1)-3S\}}{n(n-1)(n-2)(n-3)}\le V(n,k).
\]
The long row prevents a self-conjugate label, so no split-constituent
degree adjustment is needed. The full normalizer conclusion follows
from HD-ALT-03a, with the harmless additional minimum $1/4$.

For monotonicity write $m=n-k$ and $N=T^2+n(n-1)-3S_0$.
The first derivative in (ALT87.2) is positive because
$n-1-2k>0$. Also
\[
 \frac{\partial N}{\partial k}
       =-2T(n+1-2k)+3m^2-3m+\frac12.                      \tag{ALT87.4}
\]
For $k\ge1$, the exact expression for $T$ gives
$T\ge n(n-2k-1)/2\ge n(n-2)/4$ on this interval.
Using $n+1-2k\ge n/2$ and $m\le n$ in (ALT87.4) gives
\[
 \frac{\partial N}{\partial k}
 \le-\frac{n^2(n-14)}4+\frac12<0\qquad(n\ge15).
\]
The denominator of $V$ is positive and independent of $k$.
This proves both monotonicity assertions.
\end{proof}

\begin{theorem}[HD-ALT-06b: the remaining row or column exceeds $87n/100$]
\label{HD-ALT-06b}
In the transitive ordinary alternating-core setting of HD-ALT-04b,
every factor label satisfies
\[
                       \max(\lambda_1,\lambda'_1)>87n/100.
                                                               \tag{ALT87.5}
\]
This allows every compatible outer operator, scalar phase, tensor
permutation and larger inducing character field in that earlier theorem.
More precisely, if a remaining label is written $(n-k,\mu)$ with
$k<n/4$ and total inducing degree $e$, then the explicit necessary
inequality is
\[
                          \delta(n,k)\le\frac1{10}+\frac3{5e}.
                                                               \tag{ALT87.6}
\]
The inequality does not assert the existence of a source whenever it holds.
\end{theorem}

\begin{proof}
By HD-ALT-04b, transpose if necessary so that $\lambda=(n-k,\mu)$
with $0<k<n/4$. Transitivity gives the same label up to transpose
on every factor. The ordinary factor is real, so the entire core has
a unique nondegenerate symmetric form line. The inducing image acts
on it by a nontrivial multiplier, exactly as in HD-ALT-04b: a trivial
multiplier would make its character real, contrary to the actual
inducing field containing the original imaginary quadratic field.

The actual intermediate-cover argument forces genus zero, at least
twelve genuine inertias and total moving rank $2e+6$. At least two
inertias have nontrivial multiplier and moving rank at least $e/2$.
The bound in HD-ALT-06a passes to tensor products by convolution;
nontrivial tensor cycles have the quarter bound, already included
in (ALT87.3). A genuine nonidentity scalar moves all dimensions.
Therefore
\[
                   2e+6\ge e+10e\delta(n,k),                \tag{ALT87.7}
\]
which is (ALT87.6). This remains a necessary constraint on the actual
period map, not a sufficiency assertion based on parameter counts.

Suppose (ALT87.5) fails, so $k\ge cn$ for $c=13/100$.
Put
\[
 A=1-2c+2c^2=\frac{3869}{5000},\qquad
 \delta_* =\frac{1-A^2}4=\frac{10030839}{100000000}>\frac1{10}.
                                                               \tag{ALT87.8}
\]
Direct polynomial simplification gives
\[
 \frac{cn(n-1-cn)}{(n-1)(n-2)}-c(1-c)
       =\frac{13(161n-174)}{10000(n-1)(n-2)}>0,             \tag{ALT87.9}
\]
and
\[
 A^2-V(n,cn)=
 \frac{13(1906518n^2-6812633n+4129618)}
              {25000000(n-1)(n-2)(n-3)}>0.                \tag{ALT87.10}
\]
For completeness, the numerator polynomial inside the parentheses,
at $n=15+x$, is
\[
              1906518x^2+50382907x+330906673>0.
\]
Since $c(1-c)=1131/10000>\delta_*$, monotonicity in the lemma
shows $\delta(n,k)>\delta_*$ whenever $k\ge cn$.
For $n\ge30$, this and (ALT87.7) force
\[
                  e<\frac{6}{10\delta_*-1}
                     =\frac{60000000}{30839}<2000.         \tag{ALT87.11}
\]
But $k\ge cn$ excludes the three low-degree shapes with tails one
or two. The precise published degree consequence of Rasala used
in HD-ALT-04--05 gives
$a\ge2n^3/27\ge2000$; see
\cite[proof of Theorem~9.7, p.~312]{HD:GKT2019}.
There is no restriction factor of two, since the long-row label is
not self-conjugate. Thus $e\ge a\ge2000$, a contradiction.

For $15\le n\le29$, monotonicity allows $k$ to be replaced by
$k_0=\lceil13n/100\rceil$. The complete exact endpoint table is
\[
\begin{array}{c|c|c@{\qquad}c|c|c}
n&k_0&\delta(n,k_0)&n&k_0&\delta(n,k_0)\\\hline
15&2&4001/32760&23&3&3989/35420\\
16&3&1091/7280&24&4&1875/14168\\
17&3&341/2380&25&4&19483/151800\\
18&3&1679/12240&26&4&573/4600\\
19&3&2039/15504&27&4&2831/23400\\
20&3&2447/19380&28&4&28891/245700\\
21&3&1453/11970&29&4&10865/95004\\
22&3&3419/29260&&&
\end{array}                                                     \tag{ALT87.12}
\]
Every displayed fraction is greater than $1/9$. Equation (ALT87.7)
therefore gives $e<54$. For $n=15,k=2$, the two possible Specht
degrees are $15(15-3)/2=90$ and $(15-1)(15-2)/2=91$.
For every other case the same published degree bound gives
$a\ge2n^3/27\ge250$; for $n=15$ and a larger $k$, this is also
the applicable bound. Hence $e\ge90$, the final contradiction.
\end{proof}

\paragraph{Strictness, exact checks and placement.}
The label $(19,3)$ at $n=22$ has row length greater than $6n/7$
but at most $87n/100$, so the reduction strictly narrows HD-ALT-05b.
The standalone exact certificate checks both polynomial identities,
the derivative identity, all fifteen displayed endpoint fractions,
and the degree ceilings. It does not enumerate all alternating groups
or decide the remaining still-longer shapes.

This is an additional restriction on actual finite-monodromy
weight-one variations and the condition-A supplier gate, cote16
scan050, with the precise EGA/SGA and representation-ring placement
in the common framework section. It supplies no cycle, no new
condition-I positivity argument, no Cayley--Hamilton inverse, and
no motivic $L$-value estimate. The existing exact detector constants
and the conditional status of Markman, Severn and032 are unaffected.

\section{Removing the fixed-point-free involution loss in the alternating gate}
\label{sec:ordinary-alternating-fixed-free}

\paragraph{Status and scope.}
This complete new root derivation has passed the full root and
independent /root/g34\_eightfolds proof reviews, with separate exact
arithmetic supplements. Its new published input is Fomin--Lulov's exact
rim-hook estimate. The paper's actual theorem and complete proof in
Sections1--2, pp.~219--223, have been read; no reconstruction of its
upstream combinatorial inputs is claimed. The arguments below concern
the same transitive ordinary alternating core and original imaginary
character-field inclusion as HD-ALT-04--06.

\begin{lemma}[HD-ALT-07a: the three-cycle bound suffices for all long shapes]
\label{HD-ALT-07a}
Let $n\ge15$, let $\lambda=(n-k,\mu)$ with $\mu\vdash k$ and
$1\le k<n/4$, and let $W=S^\lambda|_{A_n}$ have degree $a$.
Every nonscalar finite-order linear normalizer $T$ of this represented
core satisfies
\[
 \nu(T)\ge a\beta(n,k),\qquad
 \beta(n,k)=\frac{k(n-1-k)}{(n-1)(n-2)}<\frac14.
 \tag{ALTFF.1}
\]
All scalar phases and all compatible outer automorphisms are allowed.
\end{lemma}

\begin{proof}
The long first row prevents $\lambda$ from being self-conjugate, so its
alternating restriction is irreducible. Every normalizer is scalar
times the operator of an element $g\in S_n$, by the full-automorphism
and Schur-lemma argument of HD-ALT-03a. A nonscalar normalizer has
$g\ne1$. Multiplication by its scalar phase preserves $\nu$.

If $k=1$, this is the deleted standard module. Every nonidentity
permutation has a nonscalar operator on it and hence support at least
one, which equals $a\beta(n,1)$ since $a=n-1$.
Now assume $k\ge2$. The precise degree consequence used in
HD-ALT-04--06 gives
\[
                         a\ge\frac{n(n-3)}2>4(n+1).
 \tag{ALTFF.2}
\]
Here the two tail-two degrees are $n(n-3)/2$ and
$(n-1)(n-2)/2$; every other nonstandard long shape has
$a\ge2n^3/27\ge n(n-3)/2$. These are the same published
low-degree classification consequences, with locator
\cite[proof of Theorem~9.7, p.~312]{HD:GKT2019}.
For the strict inequality in (ALTFF.2), twice the difference is
$n^2-11n-8$; at $n=15+x$ this is $x^2+19x+52>0$.

The commutator construction of HD-ALT-03a produces a three-cycle
unless $g$ is a fixed-point-free involution. In every case with that
three-cycle, the exact rank inequality
$\operatorname{rank}([A,B]-I)\le2\nu(A)$ and equality of the two
nontrivial three-cycle eigenmultiplicities give
\[
 \nu(T)\ge\frac{1-u}{3}a\ge\beta(n,k)a,
 \qquad u=\chi_W((123))/a,
 \tag{ALTFF.3}
\]
where the last inequality is the proved tail-content estimate
HD-ALT-05a. The transposition in the commutator acts in the actual symmetric
extension; no compatible outer action is assumed on a split constituent.

It remains that $n=2m$ and $g$ has cycle type $2^m$.
Fomin--Lulov's published Theorem1.1, with the character interpretation
immediately following it on p.~220, states
\[
 |\chi^{\lambda}(g)|\le
       \frac{m!\,2^m}{\sqrt{n!}}\sqrt a.
 \tag{ALTFF.4}
\]
See \cite[Theorem~1.1 and following character equality, p.~220;
complete proof, pp.~221--223]{HD:FominLulov1995}.
If the shape is not domino-decomposable, its value is zero by the
same rim-hook rule, so the displayed bound still holds.
The central binomial coefficient is the largest coefficient of
$(1+x)^{2m}$, hence
\[
 \binom{2m}{m}\ge\frac{4^m}{2m+1},\qquad
 \left(\frac{m!2^m}{\sqrt{n!}}\right)^2
     =\frac{4^m}{\binom{2m}{m}}\le n+1.
 \tag{ALTFF.5}
\]
Consequently (ALTFF.2) gives
\[
 \frac{|\chi^\lambda(g)|}{a}
      \le\sqrt{\frac{n+1}{a}}<\frac12,
 \qquad
 \nu(g|W)=\frac{a-|\chi^\lambda(g)|}{2}>\frac a4.
 \tag{ALTFF.6}
\]
The last equality uses the two eigenvalues of an involution and
continues to compute projective support after any scalar phase.
It applies whether $g$ is even or odd, since $W$ has its stated
symmetric extension and the underlying vector space is unchanged.

Finally $k(n-1-k)$ is increasing for $1\le k<n/4$, and
\[
 \beta(n,k)<\frac{n(3n-4)}{16(n-1)(n-2)}<\frac14.
 \tag{ALTFF.7}
\]
The final inequality is equivalent to $n^2-8n+8>0$, whose value
at $15+x$ is $x^2+22x+113$. Thus (ALTFF.6) is stronger than
(ALTFF.1) and covers the only remaining case.
\end{proof}

\begin{theorem}[HD-ALT-07b: an exact longer-row threshold]
\label{HD-ALT-07b}
In the transitive ordinary alternating-core setting of HD-ALT-04b,
for every remaining $n\ge15$ factor label one has
\[
 \max(\lambda_1,\lambda'_1)>
                \frac{1+\sqrt{3/5}}2\,n.
 \tag{ALTFF.8}
\]
More precisely, on writing the long-row label as $(n-k,\mu)$ and
letting $e$ be the actual total inducing degree, the necessary
finite-degree inequality is
\[
 \frac{k(n-1-k)}{(n-1)(n-2)}
                   \le\frac1{10}+\frac3{5e}.
 \tag{ALTFF.9}
\]
The assertion includes compatible outers, scalar phases, central scalar
extensions and all factor permutations. It does not assert the existence
of any source satisfying these inequalities.
\end{theorem}

\begin{proof}
By the already proved three-quarter bound, transpose the label if
necessary and write $\lambda=(n-k,\mu)$ with $0<k<n/4$.
Transitivity gives the same label, up to transpose, on every factor.
The factor representation is real, so the tensor core has a unique
nondegenerate symmetric-form line. Its multiplier on the actual
inducing group is nontrivial: otherwise the character would be real,
contrary to the inclusion of the original imaginary quadratic field.
The actual intermediate-cover and period-factorization arguments give
genus zero, at least twelve genuine inertias, and total moving rank
$2e+6$, as in HD-ALT-06b. At least two inertias have nontrivial form
multiplier, each moving at least $e/2$ dimensions.

The support fraction (ALTFF.1) passes to factorwise tensor operators
by spectral convolution. Every nontrivial factor-permutation operator
has the stronger quarter bound, since $\beta(n,k)<1/4$.
A genuine nonidentity scalar moves the entire space. Applying these
bounds to two multiplier inertias and ten others gives
\[
                         2e+6\ge e+10e\beta(n,k).
 \tag{ALTFF.10}
\]
This proves (ALTFF.9), with the actual inducing field and actual
multiplicity variation throughout.

Put
\[
 c=\frac{1-\sqrt{3/5}}2,
 \qquad c(1-c)=\frac1{10},\qquad \frac1{10}<c<\frac18.
 \tag{ALTFF.11}
\]
Suppose (ALTFF.8) fails, so $k\ge cn$.
Since $cn>1$ for $n\ge15$, the integer $k$ is at least two and
(ALTFF.2) applies. Monotonicity in $k$ gives
\[
 \begin{split}
 \beta(n,k)-\frac1{10}
 &\ge\beta(n,cn)-c(1-c)\\
 &=\frac{c\{(2-3c)n-2+2c\}}{(n-1)(n-2)}
   >\frac{13n-16}{80n^2}.
 \end{split}
 \tag{ALTFF.12}
\]
Indeed $c>1/10$, $2-3c>13/8$, $-2+2c>-2$, and the positive
denominator is less than $n^2$. The lower bracket $13n/8-2$ and the factor $c$ are positive
for $n\ge15$, which justifies multiplying these strict inequalities.
Furthermore
\[
 \frac{13n-16}{80n^2}
       >\frac6{5n(n-3)}\ge\frac3{5e}.
 \tag{ALTFF.13}
\]
The first inequality is equivalent to
$13n^2-151n+48>0$; at $n=15+x$ this polynomial is
$13x^2+239x+708>0$. The second follows from
$e\ge a\ge n(n-3)/2$. Equations (ALTFF.12)--(ALTFF.13)
contradict (ALTFF.9). This proves the exact threshold for all
$n\ge15$, with no asymptotic exception or finite table extrapolation.
\end{proof}

\paragraph{Strictness and framework.}
For example $(88,12)$ at $n=100$ has longest row $88>87n/100$
but $88<\tfrac12(1+\sqrt{3/5})n$. Moreover the stronger previous
finite-degree gate is also satisfied: HD-ALT-06a gives
$\delta(100,12)=1519571/15684900<1/10$, whereas the new bound is
$\beta(100,12)=58/539>1/10$. Thus this is a strict reduction even
relative to the previous explicit inequality, not only its uniform
$87/100$ consequence. Its two-row degree is
$\binom{100}{12}-\binom{100}{11}=908791246463100$.
No existence of an actual source with that label is asserted.
This is a restriction on actual finite local systems and weight-one
period maps, hence on condition-A suppliers (cote16 scan050;
SIGN\_NEXT SN-C03/C04). It does not produce cycles through condition I
(scan049 and cote44 Theorem4), change the resultant/Cayley--Hamilton
projector of scans048--049, or add an $L$-value or weight premise.
Its scheme, representation-ring, lambda/K-theory, EGA/SGA/FGA and
motivic placements are the common framework already stated.
The still longer shapes, including any standard tensor cases not
separately excluded, remain open. Markman, Severn and032 keep their
previous conditional status wherever used.

\section{The finite ordinary alternating range}
\label{sec:ordinary-alternating-small}

This section treats ordinary representations in alternating degree
$8\le n\le14$.  It does not extrapolate a finite scan to larger $n$.

\begin{lemma}[HD-ALT-02a: the complete finite spectral input]
\label{HD-ALT-02a}
Let $8\le n\le14$.  Among nontrivial ordinary irreducible representations of
$\mathfrak A_n$, call its deleted natural representation, of degree
$n-1$, standard.  The following statements hold.
\begin{enumerate}
\item Every nonstandard irreducible representation has degree at
least fourteen.  Every nonscalar core element has $\nu\ge a/5$.
\item In every nonstandard representation, every nonscalar linear
normalizer element has $\nu\ge a/7$.
\item The only nonstandard degree at most fifteen is the degree-fourteen
$A_8$ character with Specht label $(4,4)$, up to transpose.  Its full
linear normalizer satisfies $\nu\ge a/4$.
\item Every nonreal core character has degree greater than fifteen.
Its linear normalizer is scalar times the core, and therefore satisfies
$\nu\ge a/5$.
\end{enumerate}
All scalar phases are allowed in these assertions.
\end{lemma}

\begin{proof}
We use the complete generic character tables
\texttt{CharacterTable("Alternating",n)} and
\texttt{CharacterTable("Symmetric",n)} from the published GAP
character-table implementation.  Their exact meaning and parameter
conventions are documented in \cite[\S4.2 and \S4.2-1]{SPOR:CTblLib}.
These are computed generic tables, not relabelled ATLAS tables.
The retained inputs contain every character, class size and order,
all prime power maps, the seven actual natural-subgroup fusions,
the complete partition parameters and the source strings.
The input SHA-256 is
\texttt{f406cf1ad89801369196f36422c4dc2608e5cf96c763e5fab53b8a2ecfa1939c}.

The exact GAP scan checks all nonlinear alternating rows and all
symmetric rows whose alternating restriction has norm one.  The
fourteen tables give respectively
\[
\begin{array}{c|rrrrrrr}
n&8&9&10&11&12&13&14\\\hline
A_n&13&17&23&30&42&54&71\\
S_n&18&26&38&52&72&96&130
\end{array}
\]
selected rows, totaling $682$.  For all $49\,650$ nonscalar entries
the exact cyclic Fourier formula computes every eigenspace,
checks nonnegative integral multiplicities and their degree sum.
The standalone integer-polynomial reconstruction in
\texttt{research/alternating\_small\_certificate.py} selects rows
again by pointwise equality of restricted characters, checks all
seven fusions and power maps, and reconstructs every spectrum.
It matches every minimum and every complete minimizing spectrum
from the GAP output.  This is a separate arithmetic implementation,
with the generic character-table algorithm retained as an explicit
mathematical input.

The complete checks give the first three numerical assertions
on all these rows.  In particular all nonstandard alternating
rows have $\nu/a\ge1/5$, all selected nonstandard symmetric
rows have $\nu/a\ge1/7$, and the degree-fourteen rows all have
$\nu/a\ge1/4$.  The only core degrees at most fifteen are
$n-1$ and the indicated extra fourteen at $n=8$.
The exact conjugation test on every character value shows that
all nonreal alternating rows have degree greater than fifteen;
every symmetric character value in these tables is real.

To pass from those tables to all normalizers, use
$\operatorname{Aut}(A_n)=S_n$ and Schur's lemma, as in
HD-ALT-01a.  An inner normalizer is scalar times a core element.
If an odd outer action preserves a core representation, an
intertwiner for a transposition has scalar square.  Rescaling
it makes its square one, giving an actual representation of
$A_n\rtimes C_2=S_n$ with irreducible core restriction.
The original normalizer matrix, after its inner adjustment,
is therefore scalar times an element of one of the selected
represented symmetric tables.  This includes arbitrary actual
odd elements, not only transpositions.

If the core character were nonreal, such a symmetric extension
would be impossible: all symmetric character values here are
real, so their restrictions are real.  Thus its normalizer
has only inner projective action, proving the fourth assertion.
Scalar multiplication preserves $\nu$ throughout.
\end{proof}

\begin{theorem}[HD-ALT-02b: the remaining small ordinary core]
\label{HD-ALT-02b}
Let a dominant nonsplit Weil sixfold finite-cover source have exact
imaginary quadratic original character field.  Suppose an inducing
block of degree at least five is primitive and factorwise tensor
indecomposable in the convention of D-BI-19.  If its forced core
is ordinary $A_n^t$ with $8\le n\le14$, then necessarily
\[
 n\in\{11,12,13,14\},\qquad t\ge2,
 \quad V|_{A_n^t}\simeq\operatorname{Std}_n^{\otimes t}.
                                                               \tag{ALTS.1}
\]
In particular no single-factor ordinary core in this finite range
can occur.  Larger inducing character fields and all compatible
factor permutations are included.
\end{theorem}

\begin{proof}
The core factors are permuted transitively by D-BI-19.  Consequently
their representations differ only by automorphisms of $A_n$.
Being standard and being self-dual are preserved by those
automorphisms, and all factor degrees are the same, say $a$.
Write $e=a^t$.

First suppose the factor characters are nonreal.  HD-ALT-02a gives
$a>15$ and a fifth moving-rank bound on every factor normalizer.
For a factor-preserving tensor operator that bound passes to the
total degree by fixing all other eigenvalues.  For a nontrivial
factor permutation the weighted cyclic-word argument gives the
stronger quarter bound.  Genuine scalar inertia has full moving
rank.  The actual intermediate-cover genus and branch budget
therefore gives
\[
                 12\lceil e/5\rceil\le2e+6,
                 \qquad e\le15,
\]
contrary to $e\ge a>15$.

The factor characters are thus real, hence the core representation
is self-dual.  Its invariant nondegenerate bilinear form is unique
up to scalar, so the full inducing group is essentially self-dual.
Suppose its factors are nonstandard.  Their degree is at least
fourteen, so the primitive ESD bound D-BI-09 applies.  At least
two genuine inertias move at least half the total degree; the
other ten obey the seventh bound of HD-ALT-02a, including tensor
permutations because $1/4>1/7$.  Hence
\[
                   e+10e/7\le2e+6,
                   \qquad e\le14.                         \tag{ALTS.2}
\]
The finite degree list forces $t=1$, $n=8$ and $a=e=14$.
Its stronger quarter normalizer bound instead gives
$12\lceil14/4\rceil=48>34=2e+6$, a contradiction.

It remains that all factors are standard.  If $t=1$, the complete
scalar deleted-standard theorem D-BI-08 excludes this nonsplit
source.  If $t\ge2$ then $e\ge(n-1)^2\ge49$, so D-BI-09 applies.
Every nonscalar factor normalizer has $\nu\ge1=a/(n-1)$,
and factor permutations satisfy the stronger quarter bound.
The same two-half-inertia budget gives
\[
             e+\frac{10e}{n-1}\le2e+6.
\]
For $n=8,9,10$ it implies respectively $e\le14,24,54$,
whereas two factors already have degree at least $49,64,81$.
These cases are impossible.  Precisely the stated necessary
conditions (ALTS.1) remain.  No assertion of their existence or
dominance is made.
\end{proof}

\paragraph{Limits.}
This theorem concerns the ordinary core.  Spin cores in $n\ge8$
were excluded separately by HD-SPIN-02.  Exceptional covers with
$n<8$ and ordinary alternating degree $n\ge15$ retain their own
separate restrictions.  In (ALTS.1) the scalar deleted-standard
single-factor theorem is not applied to a tensor power.

\section{Excluding the remaining small standard alternating tensor cores}
\label{sec:alternating-standard-tensors}

This section closes the four necessary tensor cases left by
HD-ALT-02b.  The key estimate is for nonreal eigenvalues of the
actual rational orthogonal tensor operator before its scalar
phase is applied.  It is proved from traces of every power;
no assertion that scalar convolution preserves a nonreal
eigenspace bound is used.

\paragraph{Review status.}
This complete new author derivation and its entire checker have
passed the root reviewer's independent whole-proof reading.
The root reviewer also freshly reran the exact certificate:
$463$ partitions, $20551$ two-factor tensor spectra and $463$
swap spectra passed in $3.787$ seconds.  This rerun is
distinguished from a separate arithmetic implementation.

\subsection*{D-AST-01. Traces of all tensor-wreath powers}

Let $n\ge3$, let $W=\operatorname{Std}_n$ be the deleted
permutation module over $\mathbf Q$, and write $a=n-1$.
For $t\ge1$, let $\mathcal R_{n,t}$ be the image on
$V=W^{\otimes t}$ of the tensor-wreath action of
$S_n\wr S_t$.  It consists of tensor products of standard
permutation matrices followed by permutations of tensor
factors.  Put $e=a^t$.

\begin{lemma}[D-AST-01: lower trace bound for every power]
\label{D-AST-01}
For every $T\in\mathcal R_{n,t}$ and every integer $r$,
\[
                         \operatorname{tr}(T^r)\ge-e/a.
                                                               \tag{AST.1}
\]
Every nonreal eigenvalue of $T$ has multiplicity at most
\[
                              \frac{a+1}{3a}e.             \tag{AST.2}
\]
\end{lemma}

\begin{proof}
For a permutation $\pi\in S_n$ one has
\[
          \operatorname{tr}(\operatorname{Std}_n(\pi))
                       =\#\operatorname{Fix}(\pi)-1
                       \in\{-1,0,1,\ldots,a\}.            \tag{AST.3}
\]
For a tensor operator followed by a permutation of the factors,
contraction of the tensor-basis indices in its trace gives one
factor for each cycle of that factor permutation.  The factor
is the trace of the product of the standard matrices around
the cycle.  That product is another standard permutation
matrix.  Hence, if there are $c$ cycles, its trace is a product
of $c$ integers from (AST.3).

If this product is negative, one of its factors is $-1$;
the absolute values of the remaining $c-1$ factors are at
most $a$.  Thus the product is at least $-a^{c-1}$, which
is at least $-a^{t-1}=-e/a$.  If the product is nonnegative
the same bound is immediate.  Every power $T^r$ belongs to
the same finite tensor-wreath image, so the argument proves
(AST.1) for every integer $r$, not just for $r=1$.

The matrices in $\mathcal R_{n,t}$ are real and orthogonal
for the tensor product of the standard positive forms.
Thus all eigenvalues have absolute value one, and the
multiplicities of an eigenvalue $\zeta$ and its complex
conjugate agree.  Suppose $\zeta$ is nonreal, with order
$d\ge3$ and multiplicity $m$.  Some power satisfies
\[
                         \operatorname{Re}(\zeta^r)\le-1/2.
\]
Indeed, writing $\zeta=\exp(2\pi i b/d)$ with $(b,d)=1$,
choose $r$ so that $br\equiv\lfloor d/2\rfloor\pmod d$.
The resulting argument is in $[2\pi/3,\pi]$.
The two original eigenspaces contribute at most $-m$ to
the real trace of $T^r$, and all other $e-2m$ eigenvalues
contribute at most $e-2m$.  Consequently
\[
                      -e/a\le\operatorname{tr}(T^r)
                                  \le e-3m.
\]
This is (AST.2).  The two original eigenspaces are distinct
even when their powers coincide at the real value $-1$, so
that possibility does not affect the estimate.
\end{proof}

\subsection*{D-AST-02. Extraction of the actual scalar phase}

\begin{lemma}[D-AST-02: the full standard tensor normalizer]
\label{D-AST-02}
Let $n\ge8$.  Let $J$ be a finite linear group normalizing
the irreducibly acting ordinary core $A_n^t$ on
$\operatorname{Std}_n^{\otimes t}$.  Then every $g\in J$
has the form
\[
                     g=\lambda T,\qquad
             T\in\mathcal R_{n,t},\quad \lambda\in\mathbf C^*.
                                                               \tag{AST.4}
\]
For the tensor product standard symmetric form, the multiplier
of $g$ is $\lambda^2$.  If that multiplier is nontrivial,
\[
                 \operatorname{rank}(1-g)
                         \ge\frac{2a-1}{3a}e.              \tag{AST.5}
\]
Every other nonidentity $g$ has moving rank at least $e/a$.
\end{lemma}

\begin{proof}
The simple direct factors of the centerless core are intrinsic,
so $J$ permutes them and their full matrix algebras.  Remove
the corresponding standard permutation of tensor factors.
Each remaining factor automorphism is an element of
$\operatorname{Aut}(A_n)=S_n$, the same full-automorphism
input used in HD-ALT-01a and HD-ALT-02a.  The standard
$S_n$ matrix implements that automorphism on the factor
module.  After removing all these matrices, the residual
operator centralizes the irreducibly acting entire core and
is scalar by Schur's lemma.  This proves (AST.4), including
all outer actions and all factor permutations.  It is an
extraction of the scalar phase from the actual operator
before its nonreal eigenspaces are estimated.

The operator $T$ is rational and preserves the positive
standard tensor form, so the multiplier is $\lambda^2$.
Both $g$ and $T$ have finite order, whence $\lambda$ is a
root of unity.  If $\lambda^2\ne1$, then
$\lambda^{-1}$ is nonreal.  The fixed space of $g$ is the
$\lambda^{-1}$ eigenspace of this particular $T$.
Applying (AST.2) to it gives (AST.5).

For the remaining lower bound, a nonscalar standard factor
matrix has largest eigenspace codimension at least one.
Its fraction passes to a factor-preserving tensor operator
by fixing the eigenvalues in all other factors, giving
$\nu(g)\ge e/a$.  A nontrivial factor-permutation cycle
has the stronger bound $\nu(g)\ge e/4$ by the weighted
constant-word argument, and $a\ge7$ here.  If the actual
matrix is a nonidentity scalar, its moving rank is $e$.
This proves the last assertion in all cases.  Multiplication
by a scalar preserves the largest-eigenspace codimension
$\nu$; it is not used to preserve the nonreal-eigenvalue
bound (AST.2).
\end{proof}

\subsection*{D-AST-03. Closure of the finite alternating endpoint}

\begin{theorem}[D-AST-03: no standard tensor core for $8\le n\le14$]
\label{D-AST-03}
Let an inducing block of an actual rank-nine Weil sixfold
finite-cover source have an irreducibly acting normal ordinary
core $A_n^t$, with representation
$\operatorname{Std}_n^{\otimes t}$, where
$8\le n\le14$ and $t\ge2$.
Suppose its inducing character field contains the imaginary
quadratic original field, as it does for the induced sources
under discussion.  These data cannot occur.

In particular the four remaining cases of HD-ALT-02b are
excluded.  Combined with that theorem, no forced ordinary
alternating core of degree $8\le n\le14$ occurs in its
dominant nonsplit source problem.  The assertion allows all
scalar phases, larger inducing fields, and compatible tensor
permutations.
\end{theorem}

\begin{proof}
Put $a=n-1$ and $e=a^t$.  The unique invariant line of
symmetric forms is normalized by $J$, with multiplier
character $\mu$.  If $\mu$ were trivial, the character of
the finite self-dual representation of $J$ would be real.
That contradicts the inclusion of the original imaginary
quadratic field in its inducing character field.  Thus
$\mu$ is nontrivial.

The actual intermediate-cover Euler identity and period
factorization, already proved for arbitrary inducing fields,
force a rational intermediate base, at least twelve genuine
inertias, and
\[
                 \sum_{j=1}^r\operatorname{rank}(1-g_j)
                           =2e+6,\qquad r\ge12.             \tag{AST.6}
\]
For clarity, genus one would have at most six genuine
branches and a marked-base variation of dimension at most
six, whereas higher genus makes
$6+2e(1-h)$ negative since $e\ge49$.  Thus (AST.6) is a
constraint on the actual period map.

Because the inertias generate $J$ and have product one,
at least two have nontrivial multiplier.  Applying (AST.5)
to two such entries and the last assertion of D-AST-02
to the at least ten others yields
\[
 2e+6\ \ge\ 2\frac{2a-1}{3a}e+10\frac ea
             =\frac{4a+28}{3a}e.
\]
Since $7\le a\le13$, this implies
\[
                          e\le\frac{9a}{14-a}.              \tag{AST.7}
\]
But $t\ge2$ gives $e\ge a^2$, and
\[
              a^2>\frac{9a}{14-a}
             \quad\Longleftrightarrow\quad a(14-a)>9.
\]
On the integer interval $7\le a\le13$ the left-hand
side is at least thirteen.  This is a contradiction.
For the four previously remaining values the exact bounds
are, respectively,
\[
\begin{array}{c|rrrr}
n&11&12&13&14\\\hline
e\text{ upper bound}&45/2&33&54&117\\
e\text{ lower bound}&100&121&144&169.
\end{array}
\]
This proves the theorem.
\end{proof}

\paragraph{Certificate and scope.}
The exact program \texttt{check\_alternating\_standard\_tensor.py}
checks every endpoint inequality.  It also reconstructs the
deleted permutation spectra from all partitions for
$8\le n\le14$, and independently checks (AST.2) on all
two-factor tensor spectra and every two-factor swap spectrum.
The proof for arbitrary $t$ is (AST.1), not an extrapolation
from that finite computation.  Nothing here asserts a bound
for arbitrary shifted convolutions of real spectra or excludes
the ordinary alternating range $n\ge15$.

\section{Scalar conductors in ordinary standard alternating tensor cores}
\label{sec:alternating-standard-conductor}

\paragraph{Status and scope.}
This new root derivation has passed the complete root and independent
/root/sixfolds proof, source and code reviews, with separate exact
supplements. It uses the fully checked D-AST-01--02 normalizer and all-power
trace statements. The original imaginary quadratic field is retained
as a subfield of the actual inducing character field; those two fields
are not identified. All tensor multiplicities, scalar phases and
compatible factor permutations are included.

Let $n\ge8$, put $a=n-1$, and let
$V=\operatorname{Std}_n^{\otimes t}$, of dimension $e=a^t$.
Write $\mathcal R=\mathcal R_{n,t}$ for the actual rational orthogonal
tensor-wreath image from D-AST-01. If a finite group $J$ normalizes
the represented ordinary core $A_n^t$, D-AST-02 writes each element
as $g=\lambda T$, with $T\in\mathcal R$.

\begin{lemma}[HD-AST-04: the scalar character and its arithmetic marks]
\label{HD-AST-04}
The decomposition $g=\lambda T$ is unique, and its scalar coordinate
is a character $\lambda:J\longrightarrow\mu_M$ with surjective image.
If $F$ is the actual character field of $J$ on $V$, then
\[
                         F\subseteq\mathbf Q(\zeta_M).
 \tag{ASC.1}
\]
For a generating product-one tuple in $J$, the following conclusions hold.
If its original field $K\subseteq F$ differs from
$\mathbf Q(\sqrt{-3})$, at least two entries have
$\lambda\notin\mu_6$. If $p\mid |\Delta_K|$ is an odd prime,
at least two entries have scalar order divisible by $p$.
If $8\mid |\Delta_K|$, at least two entries have scalar order
divisible by eight.
\end{lemma}

\begin{proof}
A scalar in $\mathcal R$ centralizes every represented simple factor.
Its tensor-factor permutation is therefore trivial. Each remaining
standard $S_n$ operator must induce the identity automorphism of $A_n$;
the centralizer of $A_n$ in $S_n$ is trivial for $n\ge8$.
For example, commuting with every three-cycle fixes its support, and
the intersection of all three-element supports containing a given
letter is that letter. Thus each permutation is the identity.
Consequently $\mathcal R\cap\mathbf C^*I=\{I\}$.
Uniqueness follows. Since $\mathcal R$ is a group and scalars commute
with it, multiplication makes the scalar coordinate a homomorphism.
Its finite image is the full group $\mu_M$ for some $M$.
Rationality of $T$ gives
$\operatorname{tr}(g)=\lambda(g)\operatorname{tr}(T)
\in\mathbf Q(\zeta_M)$, proving (ASC.1).

If all scalar values belonged to $\mu_6$, (ASC.1) would put the
imaginary quadratic subfield $K$ inside
$\mathbf Q(\zeta_6)=\mathbf Q(\sqrt{-3})$, forcing equality.
Hence $\mu_M/(\mu_M\cap\mu_6)$ is nontrivial when the fields
differ. The images of the tuple generate this quotient and have
product one. Exactly one nonidentity image cannot have product one,
so at least two are nonidentity, as claimed.

The precise standard quadratic conductor criterion is
\[
 K\subseteq\mathbf Q(\zeta_M)\quad\Longrightarrow\quad
                             |\Delta_K|\mid M.
 \tag{ASC.2}
\]
We use it as a cited known result: Milne, \emph{Class Field Theory},
v4.03, Chapter V, Remark 3.8 and Examples 3.10--3.11,
printed pp.~158--159.\cite{HD:MilneCFT2020}
These complete invoked passages have been read. The general reciprocity
and existence theorems upstream are not reproved here. If
$M\equiv2\pmod4$, replace $M$ by $M/2$ in the cyclotomic field;
the resulting conductor divides the original $M$ as well.

For any prime $p\mid M$, the quotient
$\mu_M/\mu_{M/p}$ has order $p$. At least two tuple images
in it are nonidentity. If $p^v\Vert M$, every such scalar has
order divisible by $p^v$: writing it as $\zeta_M^b$, its quotient
image is nontrivial exactly when $p\nmid b$.
For odd $p\mid|\Delta_K|$ this proves the second assertion.
For $8\mid|\Delta_K|$, use $p=2$ and $v\ge3$ to prove the third.
\end{proof}

\begin{lemma}[HD-AST-05: arithmetic scalar support]
\label{HD-AST-05}
For $T\in\mathcal R$ and a root of unity $\xi$ of order $d$,
its eigenspace multiplicity $m_\xi$ obeys
\[
             \varphi(d)m_\xi\le e.
 \tag{ASC.3}
\]
If $\xi\notin\mu_6$, then
\[
                         m_\xi\le\frac{a+1}{4a}e.
 \tag{ASC.4}
\]
Thus a scalar phase outside $\mu_6$ has actual moving rank at least
$(3a-1)e/(4a)$. A scalar phase of order divisible by an odd prime
$p$ has actual moving rank at least $(1-1/(p-1))e$.
A scalar phase of order divisible by eight has actual moving rank
at least $3e/4$.
\end{lemma}

\begin{proof}
The characteristic polynomial of the actual finite-order rational
matrix $T$ is rational. The primitive $d$th roots are the conjugate
roots of the irreducible cyclotomic polynomial, and therefore have
the same multiplicity. This proves (ASC.3).
The exact cyclotomic Galois-group description used here is Milne,
Chapter V, Example3.2, printed p.~157.\cite{HD:MilneCFT2020}
No field of definition for a scalar twist is assumed.

If $d\notin\{1,2,3,4,6\}$, then $\varphi(d)\ge4$.
Indeed, an odd prime divisor $p\ge5$ contributes $p-1\ge4$;
otherwise $d=2^u3^v$, and the totient formula gives a value at
most two only for the five displayed possibilities.
Thus (ASC.3) implies (ASC.4) unless $d=4$.
In that last case the $i$ and $-i$ eigenspaces have the same
dimension $m_\xi$. Their squares contribute $-2m_\xi$ to
$\operatorname{tr}(T^2)$; all other $e-2m_\xi$ eigenvalues
contribute real part at most $e-2m_\xi$. By D-AST-01,
\[
       -e/a\le\operatorname{tr}(T^2)\le e-4m_\xi,
 \tag{ASC.5}
\]
giving (ASC.4). This is applied to the rational operator before
the scalar is introduced.

For $g=\lambda T$, its fixed space is the $\lambda^{-1}$
eigenspace of that particular $T$. If the scalar order is divisible
by $p$, its totient is at least $p-1$; if divisible by eight,
the totient is at least four. Equation (ASC.3) yields the remaining
moving-rank assertions. They also hold when that eigenvalue is absent.
\end{proof}

\begin{theorem}[HD-AST-06: field-dependent standard-tensor exclusions]
\label{HD-AST-06}
In the actual rank-nine sixfold-source setting of D-AST-03, suppose
the inducing core is $A_n^t$ on $\operatorname{Std}_n^{\otimes t}$,
where $n\ge15$ and $t\ge2$, and the original imaginary quadratic
field $K$ is contained in the actual inducing character field.
There is no such source in any of the following ranges:
\[
\begin{array}{c|c}
\text{original-field condition}&\text{excluded range}\\\hline
K\ne\mathbf Q(\sqrt{-3})&15\le n\le19\\
8\mid|\Delta_K|&15\le n\le20\\
p\mid|\Delta_K|,\ p\ge7\text{ prime}&15\le n\le5(p-1).
\end{array}
 \tag{ASC.6}
\]
In particular the class-number-one fields in the target have the
following remaining necessary lower bounds for standard tensor cores:
\[
\begin{array}{c|rrrrrrrr}
K&\mathbf Q(i)&\mathbf Q(\sqrt{-2})&\mathbf Q(\sqrt{-7})&
\mathbf Q(\sqrt{-11})&\mathbf Q(\sqrt{-19})&
\mathbf Q(\sqrt{-43})&\mathbf Q(\sqrt{-67})&\mathbf Q(\sqrt{-163})\\\hline
n\ge&20&21&31&51&91&211&331&811.
\end{array}
 \tag{ASC.7}
\]
The Eisenstein families already have their stated positive detector
bounds; no new standard-tensor exclusion for their larger $n$ is asserted.
\end{theorem}

\begin{proof}
The actual intermediate-cover and period-factorization argument gives
genus zero, at least twelve genuine inertias, and total moving rank
$2e+6$, exactly as in D-AST-03. Here $e\ge a^2\ge196$,
so the same higher-genus and genus-one exclusions apply.
Every other genuine inertia has moving rank at least $e/a$ by
D-AST-02, including scalar and factor-permuting operators.

For $K\ne\mathbf Q(\sqrt{-3})$, HD-AST-04 supplies two
scalar phases outside $\mu_6$. HD-AST-05 and the other ten marks give
\[
  2e+6\ge2\frac{3a-1}{4a}e+10\frac ea
          =\frac{3a+19}{2a}e.
 \tag{ASC.8}
\]
For $14\le a\le18$ this implies
$e\le12a/(19-a)$. But $e\ge a^2$, and the concave polynomial
$a(19-a)$ has minimum eighteen on this integer interval, which
is strictly greater than twelve. This is a contradiction.

For $8\mid|\Delta_K|$, the two specified scalar orders are divisible
by eight, and the analogous inequality is
\[
       2e+6\ge\left(\frac32+\frac{10}{a}\right)e,
       \qquad e\le\frac{12a}{20-a}\quad(a<20).
 \tag{ASC.9}
\]
On $14\le a\le19$, the minimum of $a(20-a)$ is nineteen,
again larger than twelve. This proves that row.

Finally put $h=p-1\ge6$. The two scalar orders divisible by $p$
give
\[
 2e+6\ge\left(2-\frac2h+\frac{10}{a}\right)e,
 \qquad e\le\frac{3ah}{5h-a}\quad(a<5h).
 \tag{ASC.10}
\]
For $14\le a\le5h-1$, concavity reduces the comparison
$a(5h-a)>3h$ to its endpoints. At the first endpoint the
difference is $67h-196>0$, and at the second it is
$2h-1>0$. Thus $a^2$ exceeds the upper bound throughout the
claimed range. This proves (ASC.6), and the explicit discriminants
$4,8,7,11,19,43,67,163$ give (ASC.7).
\end{proof}

\paragraph{Strictness, actual scope and framework.}
For instance the standard label at $n=15$, $t=2$, satisfies the
previous ALT07 necessary inequality: its support fraction is $1/14$,
below $1/10+3/(5\cdot196)$. The old D-AST-03 budget also
gives no contradiction at $a=14$. The new non-Eisenstein scalar
marks exclude this case, so the necessary main gate is strictly
smaller. No existence of any source in the remaining range is implied.
Single standard factors retain their previous separate splitness
exclusion; nonstandard long shapes are outside this theorem.

These are arithmetic restrictions on finite local systems and their
actual weight-one period maps, hence on condition-A suppliers
(cote16 scan050; SIGN\_NEXT SN-C03/C04). The scalar character is
a one-dimensional representation and a line in the representation
lambda-ring; its cyclotomic conductor is ordinary finite arithmetic,
not an $L$-value positivity argument. The common scheme-over-$\mathbf Z$,
EGA/SGA/FGA, $K$-theory and motivic-realization framework remains
as stated. No cycle is supplied by condition I (scan049/cote44
Theorem4), and no change is made to the resultant projector of
scans048--049 or to the conditional status of Markman, Severn or032.

\section{Conductor-sensitive shape bounds for all ordinary alternating cores}
\label{sec:ordinary-alternating-conductor}

\paragraph{Status and scope.}
This complete new root derivation has passed the full root and
independent /root/sixfolds proof and code reviews, with separate exact
supplements. The cyclotomic scalar-line footer precision is incorporated. It combines the actual scalar-conductor argument of HD-AST-04--05
with the full-normalizer support estimate of HD-ALT-07. No equality
between the original quadratic field and the inducing character field
is assumed. Its new shape restriction concerns every remaining long
ordinary Specht label, not only the standard label.

\begin{theorem}[HD-ALT-08: arithmetic long-row restriction]
\label{HD-ALT-08}
In the transitive ordinary alternating-core setting of HD-ALT-07, let
$n\ge15$, transpose the label if needed, and write
$\lambda=(n-k,\mu)$ with $1\le k<n/4$.
Suppose that the original imaginary quadratic field $K$ satisfies
one of the following conditions, with the indicated integer $h$:
\[
 h=4\quad\text{if }8\mid|\Delta_K|,
 \qquad h=p-1\quad\text{if an odd prime }p\ge5
                         \text{ divides }|\Delta_K|.
 \tag{AC.1}
\]
If several choices are available, the largest gives the strongest bound.
Every actual rank-nine source in this setting necessarily satisfies
\[
 \beta(n,k):=\frac{k(n-1-k)}{(n-1)(n-2)}
                     \le\frac1{5h}+\frac3{5e},
 \tag{AC.2}
\]
where $e$ is its actual total inducing degree.
For every remaining label after the separate single-standard exclusion,
\[
 \max(\lambda_1,\lambda'_1)>
          \frac{1+\sqrt{1-4/(5h)}}2\,n.
 \tag{AC.3}
\]
Thus the remaining sixfold supplier gate is restricted by the actual
original field, with the same full scalar, outer and tensor-permutation
scope as HD-ALT-07. No source is asserted to exist within the bound.
\end{theorem}

\begin{proof}
First verify the arithmetic interface for these general labels.
Each long Specht module has its ordinary rational $S_n$ model, and
its restriction to $A_n$ is irreducible and nontrivial.
Transitivity identifies the factor labels up to transpose; choosing
the same long-label extension on every factor gives a rational finite
tensor-wreath image $\mathcal R_\lambda$. The normalizer argument of
HD-ALT-03 and HD-ALT-07 writes every actual inducing operator as
$g=\eta T$, with $T\in\mathcal R_\lambda$.
Any different sign choice for an $S_n$ extension is absorbed by $\eta$.

There is no nonidentity scalar in $\mathcal R_\lambda$.
A scalar has trivial action by conjugation on each represented
simple factor, so its tensor-factor permutation is trivial.
For every remaining factor its $S_n$ permutation centralizes $A_n$,
because the represented $A_n$ is faithful and has trivial projective
kernel. That centralizer is trivial, as proved in HD-AST-04.
Hence all factor permutations are identities. The scalar coordinate
is consequently unique and multiplicative, with finite image $\mu_M$.
Rationality gives $F\subseteq\mathbf Q(\zeta_M)$ for the actual
inducing character field $F$. Since $K\subseteq F$, the cited
quadratic conductor criterion gives $|\Delta_K|\mid M$.
This is precisely the criterion and source scope in HD-AST-04,
Milne Chapter V, Remark3.8 and Examples3.10--3.11,
printed pp.~158--159.\cite{HD:MilneCFT2020}

For either choice in (AC.1), the cyclic quotient argument of
HD-AST-04 gives at least two genuine inertia entries whose scalar
orders have totient at least $h$. The characteristic polynomial of
the corresponding actual rational operator $T$ has equal
multiplicities on conjugate primitive roots. Its prescribed scalar
inverse can therefore have eigenspace dimension at most $e/h$.
Each of these two actual inertia entries moves at least
$(1-1/h)e$ dimensions. No trace lower bound for the general Specht
character is needed in this step.

The actual intermediate-cover/period argument supplies at least twelve
genuine inertias and total moving rank $2e+6$. For every other entry,
HD-ALT-07 supplies moving rank at least $\beta(n,k)e$, with its
stronger quarter bound for genuine factor permutations and full
moving rank for a nonidentity scalar. Hence
\[
       2e+6\ge2(1-1/h)e+10\beta(n,k)e.
 \tag{AC.4}
\]
This is exactly (AC.2).

Put
\[
 c=\frac{1-\sqrt{1-4/(5h)}}2,
 \qquad c(1-c)=\frac1{5h},\qquad 0<c<\frac1{16}.
 \tag{AC.5}
\]
For the last inequality, $h\ge4$ gives $c(1-c)\le1/20$,
whereas $(1/16)(15/16)=15/256>1/20$, and the function
$x(1-x)$ is increasing on $[0,1/2]$.
Suppose (AC.3) fails, so $k\ge cn$.

If $k=1$, then $cn\le1$. The single standard factor is already
excluded in the nonsplit source gate, so this remaining case has
at least two tensor factors and $e\ge(n-1)^2$. Therefore
\[
 \beta(n,1)-c(1-c)
 >\frac1{n-1}-\frac1n
 =\frac1{n(n-1)}
 >\frac3{5(n-1)^2}\ge\frac3{5e}.
 \tag{AC.6}
\]
The strict middle comparison is $5(n-1)>3n$, valid for $n\ge15$.
This contradicts (AC.2).

Now let $k\ge2$. The already cited low-degree classification
consequence in HD-ALT-07 gives $e\ge n(n-3)/2$.
If $cn\le1$, monotonicity in $k$ gives
\[
 \begin{split}
 \beta(n,k)-c(1-c)
 &>\frac{2(n-3)}{(n-1)(n-2)}-\frac1n\\
 &>\frac1{2n}>\frac6{5n(n-3)}\ge\frac3{5e}.
 \end{split}
 \tag{AC.7}
\]
Here $2(n-3)/((n-1)(n-2))>3/(2n)$ is equivalent to
$n^2-3n-6>0$, and the next comparison is $5(n-3)>12$.
Both hold throughout the stated range.

Finally suppose $cn\ge1$. The exact difference used in ALT07 is
\[
 \begin{split}
 \beta(n,k)-c(1-c)
 &\ge\frac{c\{(2-3c)n-2+2c\}}{(n-1)(n-2)}\\
 &>\frac{29n-32}{16n^3}
 >\frac6{5n(n-3)}\ge\frac3{5e}.
 \end{split}
 \tag{AC.8}
\]
Indeed $2-3c>29/16$, $-2+2c>-2$, $c\ge1/n$,
and $(n-1)(n-2)<n^2$. The resulting lower bracket is positive.
The final strict rational comparison is equivalent to
$49n^2-595n+480>0$; writing $n=15+x$ gives
$49x^2+875x+2580>0$. Every case contradicts (AC.2), proving (AC.3).
\end{proof}

\paragraph{Explicit original fields and strictness.}
For the class-number-one fields other than the Gaussian and Eisenstein
fields, the values in (AC.1) are respectively
\[
\begin{array}{c|rrrrrrr}
K&\mathbf Q(\sqrt{-2})&\mathbf Q(\sqrt{-7})&
\mathbf Q(\sqrt{-11})&\mathbf Q(\sqrt{-19})&
\mathbf Q(\sqrt{-43})&\mathbf Q(\sqrt{-67})&
\mathbf Q(\sqrt{-163})\\\hline
h&4&6&10&18&42&66&162.
\end{array}
 \tag{AC.9}
\]
The first row threshold is $(1+\sqrt{4/5})n/2$, about
$0.9472135955n$, and the other six are stronger.
For a strict finite-degree witness, take $K=\mathbf Q(\sqrt{-2})$,
$n=100$ and $\lambda=(94,6)$. Its degree is
$\binom{100}{6}-\binom{100}{5}=1116764880$ and its support
fraction is $31/539$. This lies below the old ALT07 threshold
$1/10+3/(5e)$ but above $1/20+3/(5e)$ already at the
single-factor degree, and hence at every tensor degree.
This label is nonstandard and is not a tail-two label; the new
restriction is therefore strict beyond the separate standard and
tail-two tensor statements as well. It is an inequality witness,
not a source construction.

The Gaussian field has no conductor divisor giving $h\ge4$ by this
argument, so its previous shape gate is retained. The Eisenstein
families keep their positive detector bounds. These are restrictions
on actual condition-A suppliers (cote16 scan050; SIGN\_NEXT SN-C03/C04)
in the finite-local-system and weight-one-variation language.
The scalar character defines a line over its cyclotomic field;
restriction of scalars gives a rational representation class.
Its arithmetic conductor does not supply an effective cycle.
Condition I (scan049/cote44 Theorem4), the Cayley--Hamilton projector
of scans048--049, and the stated scheme, EGA/SGA/FGA, $K$-theory,
realization, weight and motivic/$L$-function framework are unchanged.
No new use of Markman, Severn or032 is made.

\section{Integral lower degrees for the remaining alternating labels}
\label{sec:alternating-integral-cutoffs}

\paragraph{Status and scope.}
This is a new root corollary of the accepted exact shape inequalities
and all-count standard/two-box exclusions. The complete root proof/code
reading and the independent /root/sixfolds whole-proof audit have passed. It makes the existing necessary gate numerical; it is not
counted as an additional strict reduction beyond those exclusions.

In the main non-Eisenstein ordinary transitive alternating-core problem,
orient the partition, by transposing if necessary, as
$\lambda=(n-k,\mu)$ with $1\le k<n/4$ and $n\ge15$.
The accepted single-factor and all-count tensor exclusions now give
$k\ge3$. Put $h=2$ for the universal ALT07 bound. Where applicable,
one may instead use the larger values from ALT08: $h=4$ if eight
divides the original quadratic discriminant, or $h=p-1$ for an
odd discriminant prime $p\ge5$. Every such $h$ is an integer
at least two. The exact necessary inequality is
\[
 \beta(n,k)=\frac{k(n-1-k)}{(n-1)(n-2)}
       \le\frac1{5h}+\frac3{5e},
 \qquad e\ge\frac{n(n-3)}2.                         \tag{ALT9.1}
\]
Here $e$ is the actual inducing coefficient degree, not a chosen
upper bound; the degree floor is the reviewed nonstandard-label
bound used in ALT03/ALT08 and remains valid under tensor powers.

\begin{corollary}[HD-ALT-09]
\label{HD-ALT-09}
Under these hypotheses every surviving ordinary alternating label
satisfies the effective integral bound
\[
                              n\ge15h-1.                 \tag{ALT9.2}
\]
The integer is the smallest one allowed by the displayed numerical
conditions: at $n=15h-1$, the label $(n-3,3)$ satisfies
$\beta(n,3)<1/(5h)$. This last assertion is about necessary
inequalities only, not the existence of a curve source.
\end{corollary}

\begin{proof}
ALT07 already gives
$k<\tfrac12(1-\sqrt{3/5})n$. If $n\le26$, its right side is
at most $13(1-\sqrt{3/5})<3$: the strict last inequality follows
from $3\cdot169>5\cdot100$. Since $k\ge3$, necessarily
$n\ge27$.

On $3\le k<n/4$, the function $k(n-1-k)$ is increasing.
Consequently $\beta(n,k)\ge\beta(n,3)$.
Put $m=5h\ge10$ and suppose, for a contradiction, that
$27\le n\le3m-2$. Then $m\ge(n+2)/3$.
The desired strict inequality
\[
 \beta(n,3)-\frac1m>\frac6{5n(n-3)}                     \tag{ALT9.3}
\]
is equivalent, after multiplying positive denominators, to
\[
 F(n,m):=5n(n-3)\bigl(3m(n-4)-(n-1)(n-2)\bigr)
              -6m(n-1)(n-2)>0.                         \tag{ALT9.4}
\]
The coefficient of $m$ is
$15n(n-3)(n-4)-6(n-1)(n-2)>0$ for $n\ge27$.
For example, $(n-3)(n-4)>n$ in that range, so its first
term is greater than $15n^2$ and the subtracted term is
less than $6n^2$.
Thus it suffices to set $m=(n+2)/3$. Direct expansion gives
\[
 \begin{split}
 F\bigl(n,(n+2)/3\bigr)
  &=3n^3-63n^2+158n-8\\
  &=3(n-27)^3+180(n-27)^2+3317(n-27)+17380>0.
 \end{split}                                             \tag{ALT9.5}
\]
This proves (ALT9.3). The degree floor in (ALT9.1) gives
$3/(5e)\le6/(5n(n-3))$, so (ALT9.3) contradicts (ALT9.1).
Therefore $n\ge3m-1=15h-1$.

At the claimed numerical endpoint $n=3m-1$,
\[
 \beta(n,3)-\frac1m
             =\frac{-6}{m(n-1)(n-2)}<0.                 \tag{ALT9.6}
\]
By the same two-row hook formula used in ALT03,
the ordinary label $(n-3,3)$ has degree
$\binom n3-\binom n2=n(n-1)(n-5)/6$, at least
$n(n-3)/2$ since $(n-7)(n-2)\ge0$.
Also $3<n/4$ because $n\ge29$. It therefore satisfies
the numerical hypotheses and inequality at this endpoint,
including under any positive tensor power. This proves
the stated numerical minimality, without a source-existence claim.
\end{proof}

For the eight non-Eisenstein imaginary quadratic fields of class
number one, the resulting table is
\[
\begin{array}{c|r|r}
K&h&\text{necessary }n\text{ lower bound}\\\hline
\mathbf Q(i)&2&29\\
\mathbf Q(\sqrt{-2})&4&59\\
\mathbf Q(\sqrt{-7})&6&89\\
\mathbf Q(\sqrt{-11})&10&149\\
\mathbf Q(\sqrt{-19})&18&269\\
\mathbf Q(\sqrt{-43})&42&629\\
\mathbf Q(\sqrt{-67})&66&989\\
\mathbf Q(\sqrt{-163})&162&2429
\end{array}                                               \tag{ALT9.7}
\]
The exact tail and field inequalities remain necessary in addition
to these integers. This is an effective presentation of the remaining
condition-A local-system supplier gate (cote16, scan050;
SIGN_NEXT SN-C03/C04), with the same actual finite-cover,
representation lambda-ring and schemes-over-$\mathbf Z$ interpretation
as its antecedents. It gives no detecting cycle, does not invoke
condition-I positivity as a supplier, and adds no Markman, Severn
or032 premise.

\section{The two-row ordinary label with a three-box tail}
\label{sec:alternating-three-row-tail}

\paragraph{Status and scope.}
New root derivation HD-AT301--03; the complete root author proof,
checker and first exact output have been read and passed.
The complete separate /root/g34\_eightfolds proof/source/code
audit and its distinct actual-subset reconstruction have also passed.
Shared preceding weighted-cycle authorship is explicitly disclosed.  No result in this section is promoted merely
because its finite tests pass.  The label is $(n-3,3)$.
Throughout $n\ge15$, $W=S^{(n-3,3)}$, and
\[
 a=\frac{n(n-1)(n-5)}6,\quad
 A=\frac{n^2-8n+3}6,\quad B=\frac{n(n-5)}6,\quad
 b=\frac n2,\quad d_0=A+B.                              \tag{AT3.1}
\]
The degree $a$ is an integer. The constants $A,B,b,d_0$ are
used as rational bounds and need not be integers.
For a permutation $x$, write $I(x)=n-\#\operatorname{cycles}(x)$,
including fixed letters.  Projective support is
$\nu(T)=\dim T-\max_z\dim\ker(T-zI)$.

\paragraph{The exact known character input.}
Let $P_j$ be the permutation representation on $j$-subsets.
The published two-row character formula gives
$\chi_W=\chi_{P_3}-\chi_{P_2}$, with the degree in (AT3.1).
We use James--Liebeck, \emph{Representations and Characters of
Groups}, second edition (2001), Theorem~29.13, pp.~344--345
\cite[Theorem~29.13, pp.~344--345]{HD:JamesLiebeck2001};
the entire proof and its preceding Proposition~29.12 were freshly
read in the institutional three-page scan.  This known result is
cited, not re-proved here.  Counting the fixed two- and three-subsets
therefore gives the elementary specialization
\[
 \chi_W(x)=\frac{f(f-1)(f-5)}6+(f-1)u+v,                 \tag{AT3.2}
\]
where $f,u,v$ are the numbers of one-, two-, and three-cycles of $x$.
The long label is not self-conjugate, so its restriction to $A_n$
is the ordinary irreducible module already used in the preceding
alternating-core interfaces.  Its rational $S_n$ extension supplies
all compatible outer automorphisms.

\subsection*{HD-AT301. A uniform index bound for every eigenvalue}

\begin{lemma}[HD-AT301]
For every nonidentity $x\in S_n$ and every scalar phase,
\[
                 \nu(Wx)\ge A I(x)+B\ge d_0,
 \qquad \chi_W(x)\ge-b.                                \tag{AT3.3}
\]
The character lower bound also holds at the identity.
\end{lemma}

\begin{proof}
If $f=0$, (AT3.2) is $-u+v\ge-n/2$.  If $f=1$, it is
nonnegative.  For $f=2,3,4$, its first term is respectively
$-1,-2,-2$, and its other terms are nonnegative.  For $f\ge5$
all terms are nonnegative.  This proves the character lower bound.

Put
\[
 c=\frac{n^2-6n+3}6,\qquad
 c_f=\frac{f^2+(n-9)f+(n-2)(n-4)}6.                     \tag{AT3.4}
\]
For $f\ge2$, the constraint $2u+3v\le n-f$ implies
$(f-1)u+v\le(f-1)(n-f)/2$.  Substitution in (AT3.2)
and exact division give, for $f<n$,
\[
             \frac{a-\chi_W(x)}{n-f}\ge c_f\ge c.
                                                                  \tag{AT3.5}
\]
The last inequality follows from
$6(c_f-c)=f^2+(n-9)f+5>0$.
For $f=0$, the character is at most $n/3$, giving the same
bound $c=a/n-1/3$.  For $f=1$, it is at most $(n-1)/3$,
and the corresponding quotient exceeds $c$ by $(n-5)/6$.
Thus $a-\chi_W(x)\ge c(n-f)$ for every $x$, including the
identity, where both sides are zero.

Let $L$ be the order of a fixed permutation $x$, let $m_z$
be the multiplicity of $z$ on $W$, and let $I=I(x)$.
Averaging over the cyclic group gives
\[
 a-m_1=\frac1L\sum_{j=0}^{L-1}(a-\chi_W(x^j))\ge cI.
                                                                  \tag{AT3.6}
\]
Here the average number of fixed letters of $x^j$ is the
number of cycles of $x$, by the actual permutation action.
Since the characters are real and at least $-b$, for $z\ne1$
with $z^L=1$ the exact Fourier formula yields
\[
 \begin{split}
 m_z-m_1
 &=\frac1L\sum_j\chi_W(x^j)(\operatorname{Re}z^{-j}-1)\\
 &\le\frac bL\sum_j(1-\operatorname{Re}z^{-j})=b.
 \end{split}                                                   \tag{AT3.7}
\]
The inequality $m_z\le m_1+b$ is also true for $z=1$.
Consequently $\nu(Wx)\ge cI-b$.

If $I\ge(n-2)/2$, then $c-A=n/3$ gives
\[
 \nu(Wx)-AI\ge\frac n3 I-\frac n2
                 \ge\frac{n(n-5)}6=B.                    \tag{AT3.8}
\]
Suppose instead $1\le I<(n-2)/2$.  A nontrivial cycle of
length $l$ moves $l\le2(l-1)$ letters, so every power of $x$
fixes at least $f_0=n-2I>2$ letters.  The polynomial $c_f$
is increasing for $f\ge0$, since $n\ge15$.  Applying (AT3.5)
to every nonidentity power improves (AT3.6) to
$a-m_1\ge c_{f_0}I$.  Together with (AT3.7), the desired
inequality is equivalent to
\[
 (n-f_0)\bigl(f_0^2+(n-9)f_0+2n+5\bigr)
                               \ge2n(n-2).                \tag{AT3.9}
\]
After subtracting the right side, its left side is
\[
 p(f_0)=-f_0^3+9f_0^2+(n^2-11n-5)f_0+9n.                 \tag{AT3.10}
\]
The polynomial is concave on $3\le f_0\le n-2$.
At $f_0=3$ and $n-2$ its values are respectively
$3n^2-24n+39$ and $2(n^2-11n+27)$, both positive for
$n\ge15$.  Hence it is positive throughout that interval.
For completeness the remaining integer endpoint $f_0=2$
has value $2n^2-13n+18>0$.
This proves (AT3.9), and therefore (AT3.3), for all nonidentity
permutations and all target eigenspaces.
\end{proof}

\subsection*{HD-AT302. The one-factor source, over every original field}

\begin{theorem}[HD-AT302]
In the actual rank-nine sixfold inducing-source setup, a single
ordinary core $A_n$ acting as $S^{(n-3,3)}$, $n\ge15$, is
impossible.  Every compatible scalar and outer automorphism is
allowed, and the original imaginary quadratic field is arbitrary.
\end{theorem}

\begin{proof}
The normalizer is scalar times the rational $S_n$ image.
That image has trivial scalar intersection: an element acting
as a scalar centralizes the faithful $A_n$ image, hence has
trivial conjugation on $A_n$; its $S_n$ permutation is the
identity.  The scalar and rational coordinates are therefore
actual homomorphisms.  The latter gives the genuine transitive
degree-$n$ quotient cover.

As in the preceding actual-source theorems, the Euler identity
$6=2a(h-1)+D$ excludes $h\ge2$, and $h=1$ has at most six
genuine marks and period rank at most six.  Rank nine forces
$h=0$, $r\ge12$, and $D=2a+6$.  The generating product-one
tuple in the actual $n$-letter quotient satisfies
$\sum I(x_i)\ge2n-2$ by Riemann--Hurwitz.
Every nonscalar genuine inertia has actual moving rank at
least $A I(x_i)+B$, by (AT3.3).  A genuine scalar has rank
$a\ge B$, while its index is zero, so the same formula applies.
It follows that
\[
 \begin{split}
 D-2a&\ge2A(n-1)+12B-2a\\
     &=n^2-8n-1>6 \qquad(n\ge15),                        \tag{AT3.11}
 \end{split}
\]
contradicting the Euler identity.  No conductor restriction
or equality of the original and inducing fields is used.
\end{proof}

\subsection*{HD-AT303. Every transitive tensor count over non-Eisenstein fields}

\begin{theorem}[HD-AT303]
Let the ordinary transitive core be $A_n^t$ on $W^{\otimes t}$,
where $W=S^{(n-3,3)}$, $n\ge15$ and $t\ge2$.
If the character field of the actual inducing representation
contains an imaginary quadratic field
$K\ne\mathbf Q(\sqrt{-3})$, it cannot give the actual
period-rank-nine sixfold source.  All compatible scalar,
outer and tensor-factor permutation actions are included.
\end{theorem}

\begin{proof}
Write $e=a^t$.  The same centralizer argument gives a unique
scalar times rational tensor-wreath normalizer, with trivial
scalar intersection.  Its rational coordinate gives the actual
transitive $nt$-letter action.  The scalar character has finite
image $\mu_M$, and the inducing character field is contained
in $\mathbf Q(\zeta_M)$.  Containment of the original
non-Eisenstein quadratic field forces at least two genuine
inertia marks with scalar outside $\mu_6$, by generation
and product one.  Equality of the two fields is not assumed.

The actual-cover argument again gives $h=0$, $r\ge12$,
$D=2e+6$.  The trace of every power of a rational tensor-wreath
operator is at least $-(b/a)e$: if the cycle-trace product
is negative, at least one factor has absolute value at most
$b$, and every other factor has absolute value at most $a$.
Rational cyclotomic multiplicities and the fourth-root square
trace estimate of D-ATT-02 consequently give moving rank
at least $(3-b/a)e/4$ at each of the two scalar marks.

Every genuine inertia moves at least $(d_0/a)e$ dimensions,
by (AT3.3) and convolution if it fixes factors, and by the
weighted-word bound if it moves them.  A factor-moving entry
has support at least $(a-1)e/(2a)$.  Here
$2d_0\le a-1$, $b\le d_0$, and $d_0\ge43$ for $n\ge15$;
these follow directly from (AT3.1).  If three entries move
factors, select the two scalar marks and a third distinct
factor-moving entry.  At least nine further genuine marks give
\[
 D\ge2e+a^{t-1}\left(9d_0-\frac{b+1}{2}\right)>2e+6.
                                                                  \tag{AT3.12}
\]
Thus exactly two entries move factors.  Product one makes their
factor permutations inverse, and transitivity makes each a full
$t$-cycle.  This is precisely the elementary two-moving-branch
argument of D-ATT-07 with the new constants.

Let $H$ be the sum of the natural indices of their two cycle
products.  Let $N$ be the total natural index of all within-factor
permutations at the other $r-2\ge10$ entries.  The actual
$nt$-letter quotient has total index
$2n(t-1)+H+N\ge2nt-2$, hence
\[
                              H+N\ge2n-2.                \tag{AT3.13}
\]
At a factor-fixing entry with $s\ge1$ active factors, convolution
and (AT3.3) give support at least the maximum of the $s$ active
supports.  It is therefore at least
$a^{t-1}(A I/s+B)\ge a^{t-1}(A I/t+B)$, where $I$ is that
entry's total natural index.  A genuine scalar has full rank
and satisfies the same bound with $I=0$, since $a\ge B$.

For a weighted full $t$-cycle with product $Wx$, the complete
general weighted-cycle lemma D-ATT-08 gives moving rank at least
\[
 \left(1-\frac1t\right)e
     +\frac{a^{t-1}}t\nu(Wx)-\frac{R_t(a)}t.              \tag{AT3.14}
\]
It remains valid when the product is the identity: here only
$\nu(Wx)\ge A I(x)$ is used, with both sides zero in that case.
The exact word error is $R_2(a)=a$; for $t\ge3$ the accepted
all-integer estimate is $R_t(a)\le2a^{t-2}$.  Thus
$2R_t(a)/a^{t-1}\le2$ for every $t\ge2$, because $a\ge350$.
Summing (AT3.14) and all factor-fixing estimates, using
(AT3.13), proves
\[
 \begin{split}
 D-2e
 &\ge\frac{a^{t-1}}t
       \left(2A(n-1)-2a+10tB-2\right)\\
 &=\frac{a^{t-1}}t\left(10tB-(n-1)^2-2\right)
   >B a^{t-1}>6.                                         \tag{AT3.15}
 \end{split}
\]
For the strict step, $9tB\ge18B=3n(n-5)$ and
$3n(n-5)-(n-1)^2-2=2n^2-13n-3>0$ for $n\ge15$.
This contradicts the actual Euler identity and proves the theorem.
\end{proof}

\paragraph{Strictness and the project frame.}
The earlier exact shape/conductor gate allows the label $(97,3)$
at $n=100$ over $\mathbf Q(i)$, for every positive tensor count:
its exact three-cycle fraction is below $1/10$ and its long row
exceeds the ALT07 threshold.  It is neither a standard nor a
two-box-tail label.  HD-AT302/303 remove that whole supplier
family, rather than restating the previous numerical cutoffs.
No actual source at that profile was previously asserted.
Other three-box and larger tails remain separate.

These are obstructions in the actual finite-cover representation
lambda-ring and its associated local systems; the scalar coordinate
and quotient cover are genuine homomorphisms and morphisms.
They shrink the condition-A supplier gate (cote16scan050;
SIGN_NEXT SN-C03/C04) and provide no condition-I cycle.
The preceding schemes-over-$\mathbf Z$, EGA/SGA/FGA, projector,
K-theoretic and motivic interpretation is unchanged.
No Markman, Severn or032 premise is added.

\section{The ordinary hook with a three-box tail}
\label{sec:hook-three-source}

\paragraph{Status, attribution and scope.}
New author derivation D-HK301--03. The complete author, root and
separate /root/imprimitive whole-proof and primary-source reviews have
passed, with distinct exact spectrum reconstructions.  The cyclic-average character-defect method
was supplied by the root's HD-AT301; the new hook character
bounds and constants are derived here.  The exact weighted-cycle
and actual quotient framework is the previously checked
D-ATT07--09.  No finite computation is treated as a universal
proof.  Throughout $n\ge15$, let $U$ be the standard
$(n-1)$-dimensional rational representation of $S_n$, and put
\[
 W=\Lambda^3U\simeq S^{(n-3,1,1,1)},\qquad
 a=\binom{n-1}{3},\quad
 A=\frac{n^2-9n+8}{6},\quad
 B=\frac{n^2-2n-1}{8},\quad
 b=\frac{(n-1)^2}{8},\quad d_0=A+B.
 \tag{HK3.1}
\]
These bounds are rational numbers; no integrality of $A,B,b$
is assumed.  For a permutation $x$, set
$I(x)=n-\#\operatorname{cycles}(x)$, including fixed letters.
For a complex operator $T$ on a space of dimension $a$, write
$\nu(T)=a-\max_z\dim\ker(T-zI)$.

\paragraph{Exact published representation and character inputs.}
Fulton--Harris, Proposition~3.12 and its complete proof,
pp.~31--32, give the irreducibility of every exterior power
of the standard representation
\cite[Proposition~3.12, pp.~31--32]{HK:FultonHarris1991}.
The hook identification is also described in Exercise~4.6,
p.~48, and the branching discussion on p.~59.
For an explicit published character formula with the precise
partition label, use Garsia--Goupil, Proposition~I.1, p.~2,
its complete proof (2.1)--(2.6), pp.~9--11, and Table~5.1,
p.~31 \cite[Proposition~I.1; Table~5.1]{HD:GarsiaGoupil2009}.
Writing $f,u,v$ for the numbers of one-, two-, and three-cycles,
their formula is
\[
 \chi_W(x)=\frac{(f-1)(f-2)(f-3)}6-(f-1)u+v.
 \tag{HK3.2}
\]
It agrees directly with the exterior-power trace identity
$\chi_{\Lambda^3U}(x)=
(\chi_U(x)^3-3\chi_U(x)\chi_U(x^2)+2\chi_U(x^3))/6$.
In particular the formula identifies the exterior cube and
the labeled irreducible character without an unproved labeling
assumption.  Its value at the identity is the degree in (HK3.1).
The label is not self-conjugate for $n\ge15$, so restriction
to $A_n$ is the ordinary irreducible module of the accepted
alternating-core interface.  Its rational $S_n$ extension
supplies all compatible outer automorphisms.

\subsection*{D-HK301. Uniform support with a full positive constant}

\begin{lemma}[D-HK301]
For every nonidentity $x\in S_n$,
\[
                  \nu(Wx)\ge A I(x)+B\ge d_0.
 \tag{HK3.3}
\]
This bounds every target eigenspace and is invariant under
arbitrary scalar multiplication of $Wx$.  For every $x$,
including the identity, $\chi_W(x)\ge-b$.
\end{lemma}

\begin{proof}
For $f=0$, (HK3.2) equals $-1+u+v\ge-1\ge-b$.
For $f\ge1$, its first term is nonnegative, including
$f=1,2,3$, when it vanishes.  Since $2u\le n-f$,
\[
 \chi_W(x)\ge-\frac{(f-1)(n-f)}2
             \ge-\frac{(n-1)^2}{8}=-b.
 \tag{HK3.4}
\]
This proves the global character lower bound.

Put
\[
 c=\frac{n^2-6n+8}{6},\qquad
 c_f=\frac{f^2+(n-6)f+n^2-6n+9}{6}.
 \tag{HK3.5}
\]
For $f\ge1$ the negative $u$ term may be discarded and
$v\le(n-f)/3$.  Thus
\[
 \chi_W(x)\le\frac{(f-1)(f-2)(f-3)}6+\frac{n-f}{3}.
 \tag{HK3.6}
\]
Subtracting this upper bound from $a$ and dividing by
$n-f$ gives exactly $c_f$ when $f<n$.
Moreover $6(c_f-c)=f^2+(n-6)f+1>0$.
For $f=0$, the constraint $2u+3v\le n$ gives
$\chi_W(x)=-1+u+v\le n/2-1$, and
$a-n/2+1=nc$.  The identity has zero defect and zero
moved-letter count.  Consequently
\[
                    a-\chi_W(x)\ge c(n-f)
 \tag{HK3.7}
\]
for every permutation, with the stronger coefficient $c_f$
available whenever $f\ge1$ and $x\ne1$.

Fix $x$ of order $L$, put $I=I(x)$, and let $m_z$ be
the multiplicity of a root $z$ in $Wx$.  The average number
of fixed letters of $x^j$ is the number of cycles of $x$.
Averaging (HK3.7) therefore gives $a-m_1\ge cI$.
Every character value is real and at least $-b$.
For $z^L=1$, $z\ne1$, the exact Fourier formula gives
\[
 \begin{split}
 m_z-m_1
 &=\frac1L\sum_{j=0}^{L-1}
       \chi_W(x^j)(\operatorname{Re}z^{-j}-1)\\
 &\le\frac bL\sum_{j=0}^{L-1}(1-\operatorname{Re}z^{-j})=b.
 \end{split}
 \tag{HK3.8}
\]
The same upper bound is trivial for $z=1$.
Hence $\nu(Wx)\ge cI-b$; this is an all-eigenvalue bound,
not an identification of the largest eigenspace with the fixed space.

If $I\ge(n-2)/2$, use $c-A=n/2$ to obtain
\[
 \nu(Wx)-AI\ge\frac n2I-b
     \ge\frac{n(n-2)}4-\frac{(n-1)^2}{8}=B.
 \tag{HK3.9}
\]
If $1\le I<(n-2)/2$, put $f_0=n-2I>2$.
Every power of $x$ fixes at least $f_0$ letters: a
nontrivial cycle of length $l$ moves at most $2(l-1)$
letters, and taking a power cannot decrease the fixed set.
The coefficient $c_f$ is increasing for $f\ge0$, because
$n\ge15$.  Applying (HK3.6)--(HK3.7) to each nonidentity
power improves the cyclic average to
$a-m_1\ge c_{f_0}I$.  Combining this with (HK3.8), the
desired bound is equivalent to
\[
 (n-f_0)\bigl(f_0^2+(n-6)f_0+3n+1\bigr)
                              \ge3n(n-2).
 \tag{HK3.10}
\]
After subtracting the right side, the polynomial is
\[
 p(f_0)=-f_0^3+6f_0^2+(n^2-9n-1)f_0+7n.
 \tag{HK3.11}
\]
It is concave for $2\le f_0\le n-2$.
Its values at the endpoints are
\[
                p(2)=2n^2-11n+14,\qquad
                p(n-2)=n^2-12n+34.
 \tag{HK3.12}
\]
Both are positive for all $n\ge15$, so the whole interval
has positive values.  This proves the small-index case,
and thus the complete lemma.
\end{proof}

\subsection*{D-HK302. One factor over every original quadratic field}

\begin{theorem}[D-HK302]
In the actual rank-nine sixfold inducing-source setup, a
single ordinary core $A_n$ acting as $\Lambda^3U$, $n\ge15$,
is impossible.  Every compatible scalar and outer action
is allowed, and the original imaginary quadratic field
is arbitrary.
\end{theorem}

\begin{proof}
The normalizer is scalar times the rational $S_n$ image.
That rational image meets the scalars trivially.  Indeed a
scalar image centralizes the faithful nontrivial $A_n$ image,
so its permutation centralizes $A_n$; in $S_n$ this forces
the identity.  The scalar and permutation coordinates
are consequently actual homomorphisms.  The latter gives
the genuine transitive $n$-letter quotient cover.

For the actual intermediate base of genus $h$, the Euler
identity is $6=2a(h-1)+D$.  Since $a\ge364$, genus at least
two is impossible.  Genus one has at most six genuine
marks and period-image rank at most six.  Rank nine
therefore forces $h=0$, $r\ge12$ and $D=2a+6$.
The generating product-one tuple in the actual transitive
$n$-letter quotient has total natural index at least
$2n-2$ by Riemann--Hurwitz.

For every genuine entry with nonidentity permutation,
its actual moving rank is at least $AI+B$, by D-HK301.
A genuine scalar has index zero and moving rank $a\ge B$,
so the same inequality holds.  Summing gives
\[
 \begin{split}
 D-2a&\ge2A(n-1)+12B-2a\\
     &=\frac{n^2-6n-13}{6}>6\qquad(n\ge15),
 \end{split}
 \tag{HK3.13}
\]
contradicting the actual Euler identity.
No conductor condition or inducing-field equality is used.
\end{proof}

\subsection*{D-HK303. Every transitive tensor count over non-Eisenstein fields}

\begin{theorem}[D-HK303]
Suppose the ordinary transitive normal core is $A_n^t$
on $W^{\otimes t}$, with $W=\Lambda^3U$, $n\ge15$,
and $t\ge2$.  If the character field of the actual
inducing representation contains an imaginary quadratic
field $K\ne\mathbf Q(\sqrt{-3})$, this representation
cannot give an actual period-rank-nine sixfold source.
All compatible scalar, outer and factor-permuting
actions are included.
\end{theorem}

\begin{proof}
Write $e=a^t$.  The rational tensor-wreath normalizer
has trivial scalar intersection by its faithful conjugation
on the simple core factors.  It gives actual scalar and
rational coordinates, and an actual transitive $nt$-letter
quotient.  Let the finite scalar image be $\mu_M$.
The inducing character field lies in $\mathbf Q(\zeta_M)$.
If that image were contained in $\mu_6$, it could not
contain the stated field $K$.  The generating product-one
inertia tuple thus contains at least two genuine marks
whose scalar coordinate lies outside $\mu_6$.
This does not assume equality of the inducing and original fields.

As before, the actual base has genus zero, $r\ge12$ and
$D=2e+6$.  A rational tensor-wreath operator has trace
equal to the product of the traces of its factor-cycle
products.  If this is negative, at least one factor has
absolute value at most $b$, and all others have absolute
value at most $a$.  The same applies to every power, so
all such traces are at least $-(b/a)e$.

At a scalar mark outside $\mu_6$, a fixed eigenvalue of
the actual operator would be a root of the rational
operator of order outside $1,2,3,6$.  If its order is
not four, its rational cyclotomic multiplicity is at
most $e/4$.  For order four the two primitive roots have
equal multiplicity $m$; the real trace of the square
is at most $e-4m$ and at least $-(b/a)e$.
Thus each of the two scalar marks moves at least
\[
                           \frac{3-b/a}{4}e.
 \tag{HK3.14}
\]
This proves the marked estimate directly with every scalar phase retained.

Every genuine entry moves at least $(d_0/a)e$ by
D-HK301 and convolution when factors are fixed.
A factor-moving entry has the proved weighted-word
support bound $(a-1)e/(2a)$.
For $n\ge15$ the elementary constants satisfy
\[
                  d_0\ge40,\qquad b\le d_0,\qquad
                  2d_0\le a-1,\qquad a\ge364.
 \tag{HK3.15}
\]
For example the numerators of $d_0-40$, $d_0-b$ and
$a-1-2d_0$, after substituting $n=x+15$, are respectively
$7x^2+168x+14$, $2x^2+42x+193$ and
$2x^3+71x^2+844x+3382$, with positive denominators.
Hence the lower $(d_0/a)e$ bound also covers moving factors.

If three entries move factors, choose the two scalar
marks and a third distinct factor-moving entry.
At least nine further genuine entries remain, and
\[
 D\ge2e+a^{t-1}\left(9d_0-\frac{b+1}{2}\right)>2e+6.
 \tag{HK3.16}
\]
This is impossible.  Transitivity rules out zero
factor-moving entries, and product one rules out one.
There are therefore exactly two.  Their factor
permutations are inverse, and transitivity makes each
a full $t$-cycle.

Let $H$ be the sum of the natural indices of the two
cycle products, and $N$ the sum of all within-factor
natural indices at the other $r-2\ge10$ entries.
The actual transitive $nt$-letter cover has total index
$2n(t-1)+H+N$, so Riemann--Hurwitz gives
\[
                             H+N\ge2n-2.
 \tag{HK3.17}
\]
At a factor-fixing entry with $s\ge1$ active factors,
the maximum of their projective supports is at least
their average.  D-HK301 therefore gives moving rank
at least $a^{t-1}(AI/s+B)\ge a^{t-1}(AI/t+B)$,
where $I$ is the entry's total natural index.
The entire positive constant $B$ is retained.
A genuine scalar satisfies the same formula with $I=0$.

For each weighted full $t$-cycle the complete general
lemma D-ATT08 gives
\[
 D_i\ge\left(1-\frac1t\right)e
       +\frac{a^{t-1}}t\nu(Wx_i)-\frac{R_t(a)}t.
 \tag{HK3.18}
\]
Only $\nu(Wx_i)\ge AI(x_i)$ is used here, so an identity
cycle product is included with both sides zero.
The exact error is $R_2(a)=a$, and for $t\ge3$ the
proved estimate is $R_t(a)\le2a^{t-2}$.
Consequently $2R_t(a)/a^{t-1}\le2$ for every $t\ge2$.
Summing the two cycles and the other entries gives
\[
 \begin{split}
 D-2e
 &\ge\frac{a^{t-1}}t
       \bigl(2A(n-1)-2a+10tB-2\bigr)\\
 &=\frac{a^{t-1}}t
       \left(10tB-\frac{2(n-1)(2n-1)}3-2\right)
       >B a^{t-1}>6.
 \end{split}
 \tag{HK3.19}
\]
For the strict step, $9tB\ge18B$ and
\[
 18B-\frac{2(n-1)(2n-1)}3-2
                   =\frac{11n^2-30n-59}{12}>0
 \tag{HK3.20}
\]
for all $n\ge15$.  This contradicts $D=2e+6$
and proves the theorem.
\end{proof}

\paragraph{Strictness and boundaries.}
This removes the entire hook-three supplier profile
in the stated ranges, not just a restatement of an
earlier numerical long-row condition.  For example
$(97,1,1,1)$ at $n=100$ has a long row beyond the
earlier Gaussian cutoff and is neither a standard nor
a two-box-tail label.  No actual source at that profile
was previously asserted.
The all-tensor result retains the non-Eisenstein
condition; the one-factor theorem does not need it.
No conclusion about a different label, a nontransitive
factor system or a positive algebraic-cycle supplier
is inferred.  These are actual finite-cover/representation
obstructions in the existing condition-A source gate.
No Markman, Severn or032 premise is introduced.

\section{The mixed three-box ordinary alternating tail}
\label{sec:alternating-mixed-three-tail}

\paragraph{Status and scope.}
New root derivation HD-MT301--03. Complete root and separate
/root/g34\_eightfolds whole-proof, source and checker reviews have
passed, including a distinct exact cyclic and actual-flag reconstruction. The ordinary label is $(n-3,2,1)$,
with $n\ge15$. Put
\[
 a=\frac{n(n-2)(n-4)}3,\quad
 A=\frac{n^2-8n+8}3,\quad B=\frac{n(n-3)}3,\quad
 b=\frac n3,\quad d_0=A+B.                              \tag{MT3.1}
\]
Here $a$ is the integral degree of $W=S^{(n-3,2,1)}$;
the remaining constants are rational bounds. Its restriction
to $A_n$ is irreducible, as this long label is not self-conjugate.
The rational $S_n$ model and full compatible normalizer are
the accepted ordinary alternating interfaces.

The exact known character input is
\[
               \chi_W(x)=\frac{f(f-2)(f-4)}3-v,          \tag{MT3.2}
\]
where $f$ and $v$ count the one- and three-cycles of $x$.
It is the $(n-3,2,1)$ row of Garsia--Goupil,
\emph{Character Polynomials, their q-Analogs and the Kronecker
Product}, Table~5.1, p.~31 \cite{HD:GarsiaGoupil2009}.
The general character-polynomial statement is Proposition~I.1,
p.~2, whose complete invoked proof is Section~2, equations
(2.1)--(2.6), pp.~9--11. Those passages and the table were
freshly read. We cite this published formula instead of
reproving the known character theory.

\subsection*{HD-MT301. The all-eigenvalue index estimate}

\begin{lemma}[HD-MT301]
For every nonidentity permutation $x\in S_n$,
\[
 \nu(Wx)\ge A I(x)+B\ge d_0,
 \qquad \chi_W(x)\ge-b,                                \tag{MT3.3}
\]
where $I(x)=n-\#\operatorname{cycles}(x)$ includes fixed
letters, and $\nu$ is maximum-eigenspace codimension.
The character lower bound also holds at the identity.
\end{lemma}

\begin{proof}
The constraint $3v\le n-f$ gives
\[
 \chi_W(x)\ge\frac{f(f-2)(f-4)+f-n}3
            =\frac{f(f-3)^2-n}3\ge-\frac n3.             \tag{MT3.4}
\]
Also $v\ge0$, so for $f<n$,
\[
 \frac{a-\chi_W(x)}{n-f}\ge c_f,
 \qquad
 c_f=\frac{n^2+nf+f^2-6n-6f+8}3,
 \quad c=c_0=\frac{(n-2)(n-4)}3.                         \tag{MT3.5}
\]
This follows by exact polynomial division. Since $n\ge15$,
$c_f$ is increasing for $f\ge0$, and hence $c_f\ge c$.
At the identity the un-divided defect inequality has both
sides zero.

Fix $x$ of order $L$ and index $I$. Cyclic averaging gives
$a-m_1\ge cI$, where $m_1$ is its fixed-space multiplicity:
the average number of fixed letters equals its number of
cycles. The all-eigenvalue Fourier calculation (AT3.7),
using the just-proved real-character lower bound $-b$,
gives $m_z\le m_1+b$ for every root $z$. Explicitly,
\[
 m_z-m_1=\frac1L\sum_{j=0}^{L-1}\chi_W(x^j)
                 (\operatorname{Re}z^{-j}-1)\le b
 \quad(z\ne1),                                         \tag{MT3.6}
\]
and the inequality is immediate for $z=1$.
Thus $\nu(Wx)\ge cI-b$.

If $I\ge(n-2)/2$, then $c-A=2n/3$ implies
\[
 \nu(Wx)-AI\ge\frac{2nI-n}{3}
                         \ge\frac{n(n-3)}3=B.           \tag{MT3.7}
\]
Otherwise every power of $x$ fixes at least $f_0=n-2I>2$
letters, since the moved support of a permutation is at
most twice its index. Monotonicity of $c_f$ improves the
cyclic bound to $a-m_1\ge c_{f_0}I$.
After using (MT3.6) and clearing positive denominators,
$c_{f_0}I-b\ge AI+B$ is exactly
\[
 (n-f_0)\bigl(f_0^2+(n-6)f_0+2n\bigr)\ge2n(n-2).
                                                                  \tag{MT3.8}
\]
The difference of the two sides is
\[
 p(f_0)=-f_0^3+6f_0^2+(n^2-8n)f_0+4n.                   \tag{MT3.9}
\]
It is concave on $[2,n-2]$, because $p''(f)=12-6f\le0$.
Its endpoint values are $2(n-2)(n-4)$ and $2(n-4)^2$,
both strictly positive. It is therefore positive throughout
that interval. This proves the required uniform bound for
every nonidentity permutation and every target eigenvalue.
\end{proof}

\subsection*{HD-MT302. The one-factor source in every field}

\begin{theorem}[HD-MT302]
The actual rank-nine sixfold inducing source cannot have
a single ordinary core $A_n$ on $S^{(n-3,2,1)}$, $n\ge15$.
This permits all compatible scalar and outer actions over
every original imaginary quadratic field.
\end{theorem}

\begin{proof}
The rational $S_n$ normalizer meets the scalars trivially:
a scalar centralizes the faithful $A_n$ image, while
the centralizer of $A_n$ in $S_n$ is trivial. Extracting
the natural permutation is therefore an actual homomorphism,
and its image contains $A_n$. The associated quotient
cover is connected and has degree $n$.

The actual Euler identity is $6=2a(h-1)+D$.
Since $a\ge715$, genus $h\ge2$ is impossible, and
$h=1$ permits at most six genuine marks and period rank
at most six. Thus rank nine requires $h=0$, at least
twelve genuine inertia entries, and $D=2a+6$.
Riemann--Hurwitz for the actual degree-$n$ quotient gives
$\sum I(x_i)\ge2n-2$.
Each nonscalar entry moves at least $AI(x_i)+B$ dimensions
by (MT3.3). A genuine scalar moves all $a\ge B$ dimensions
and has $I=0$, so it satisfies the same bound.
Consequently
\[
 D-2a\ge2A(n-1)+12B-2a
       =\frac{6n^2-20n-16}{3}>6,                         \tag{MT3.10}
\]
contradicting the source Euler identity. No scalar-conductor
assumption is used in this one-factor argument.
\end{proof}

\subsection*{HD-MT303. All transitive tensor counts}

\begin{theorem}[HD-MT303]
Let the actual inducing image normalize the ordinary core
$A_n^t$ on $W^{\otimes t}$ transitively on its factors,
where $W=S^{(n-3,2,1)}$, $n\ge15$ and $t\ge2$.
If its character field contains a non-Eisenstein imaginary
quadratic field, it cannot give the rank-nine sixfold source.
All compatible outer automorphisms, scalar phases and
factor permutations are included.
\end{theorem}

\begin{proof}
Write $e=a^t$. The faithful rational tensor-wreath normalizer
has trivial scalar intersection, by its conjugation action
on the individual faithful simple factors. Thus its scalar
and rational coordinates are actual homomorphisms, and its
natural $nt$-letter action is transitive. The inducing
character field is contained in the cyclotomic field of
the scalar character. The original non-Eisenstein field
therefore forces at least two inertia phases outside $\mu_6$,
by generation and product one. It need only be contained
in the inducing field; equality is never assumed.

The actual Euler and marked-cover arguments again give
$h=0$, $r\ge12$, $D=2e+6$.
Every power of a rational tensor-wreath operator has trace
at least $-(b/a)e$, since its trace is a product of
cycle-product traces bounded below by $-b$ and in absolute
value by $a$. Rational cyclotomic multiplicities and the
fourth-root square-trace calculation of D-ATT-02 give
moving rank at least $(3-b/a)e/4$ at both scalar marks.
Every genuine inertia has support at least $(d_0/a)e$:
use (MT3.3) and spectral convolution if factors are fixed;
if they move, the weighted-word argument gives the stronger
$(a-1)e/(2a)$. Directly from (MT3.1),
$2d_0\le a-1$, $b\le d_0$, and $d_0>97$ for $n\ge15$.

If three entries move factors, choose one distinct from
both marked scalar entries. With nine further genuine marks,
\[
 D\ge2e+a^{t-1}\left(9d_0-\frac{b+1}{2}\right)>2e+6.
                                                                  \tag{MT3.11}
\]
Therefore exactly two entries move factors: zero contradicts
transitivity, and one contradicts product one. Their factor
permutations are inverse full $t$-cycles.

Let $H$ be the sum of the natural indices of their two
cycle products and $N$ the sum of all natural indices
at the factor-fixing entries. Riemann--Hurwitz for the
actual degree-$nt$ quotient gives
$2n(t-1)+H+N\ge2nt-2$, hence $H+N\ge2n-2$.
At a factor-fixing genuine entry, averaging the active
factor estimates in (MT3.3) gives support at least
$a^{t-1}(AI/t+B)$, preserving the full $B$.
A genuine scalar has full rank and satisfies this with
$I=0$. At each full factor cycle the exact D-ATT-08 bound
is
\[
 (1-1/t)e+\frac{a^{t-1}}t\nu(Wx)-\frac{R_t(a)}t.
                                                                  \tag{MT3.12}
\]
Here use only $\nu(Wx)\ge AI(x)$, which is also valid
for the identity cycle product. The exact errors satisfy
$R_2(a)=a$ and $R_t(a)\le2a^{t-2}$ for $t\ge3$.
Thus $2R_t(a)/a^{t-1}\le2$ for every $t\ge2$.
There are at least ten factor-fixing entries. Summing gives
\[
 \begin{split}
 D-2e
 &\ge\frac{a^{t-1}}t
 \left(2A(n-1)-2a+10tB-2\right)\\
 &=\frac{a^{t-1}}t
 \left(10tB-\frac{6n^2-16n+16}{3}-2\right)
 >B a^{t-1}>6.
 \end{split}                                             \tag{MT3.13}
\]
Indeed $9tB\ge18B$ and
\[
 18B-\frac{6n^2-16n+16}{3}-2
                    =\frac{12n^2-38n-22}{3}>0
 \quad(n\ge15).                                         \tag{MT3.14}
\]
This contradicts the Euler budget and proves the theorem.
\end{proof}

\paragraph{Strictness and interpretation.}
The earlier exact shape gate allows $(97,2,1)$ at $n=100$
over $\mathbf Q(i)$: its tail has size three and its
ALT07 fraction is less than $1/10$ even in the tensor limit.
This label is distinct from the already excluded standard,
two-box and two-row three-box labels. The new result
removes another entire necessary supplier family, without
asserting that an actual source previously existed there.
It constrains actual finite-cover local systems and their
representation lambda-rings and K-theory; in the project's
frame it shrinks condition-A suppliers (cote16scan050;
SIGN_NEXT SN-C03/C04). No condition-I cycle or new conditional
Markman, Severn or032 input is produced.

\section{Integral lower degrees after the three-box exclusions}
\label{sec:alternating-four-tail-integral-cutoffs}

\paragraph{Status and scope.}
Root derivation HD-ALT-10. Complete root and separate
/root/sixfolds whole-proof reviews have passed. Its three-box
premises are now promoted after their own complete peer audits. This is an effective numerical
corollary of those exclusions and the already accepted shape bounds.
It is not counted as a separate strict reduction of the main gate.

In the non-Eisenstein ordinary transitive alternating-core problem,
transpose the label if necessary and write $\lambda=(n-k,\mu)$.
The all-count standard, two-box and three-box exclusions give $k\ge4$.
The accepted ALT07/ALT08 necessary inequalities are
\[
 k<\frac{1-\sqrt{3/5}}2n,\qquad
 \frac{k(n-1-k)}{(n-1)(n-2)}
       \le\frac1{5h}+\frac3{5e},\qquad
 e\ge\frac{n(n-3)}2.                                  \tag{ALT10.1}
\]
Here $h=2$ universally, $h=4$ if eight divides the original
quadratic discriminant, or $h=p-1$ for an odd discriminant
prime $p\ge5$. The actual inducing degree is $e$; equality
of the original field and the inducing character field is
not required by the cited shape inequalities.

\begin{corollary}[HD-ALT-10]
Every surviving label in this gate satisfies
\[
                              n\ge20h-2.                \tag{ALT10.2}
\]
This integer is optimal for the displayed numerical conditions:
at $n=20h-2$, the label $(n-4,4)$ satisfies them. This
numerical endpoint asserts no existence of a curve source.
\end{corollary}

\begin{proof}
If $n\le35$, the first inequality in (ALT10.1) implies
$k<35(1-\sqrt{3/5})/2<4$; the latter strict inequality
is equivalent to $3\cdot35^2>5\cdot27^2$.
Thus $n\ge36$. On $4\le k<n/4$, the expression
$k(n-1-k)$ is increasing, so it suffices to consider $k=4$.

Put $m=5h\ge10$ and suppose $36\le n\le4m-3$.
The strict contradiction to (ALT10.1) is
\[
 \frac{4(n-5)}{(n-1)(n-2)}-\frac1m
                              >\frac6{5n(n-3)}.         \tag{ALT10.3}
\]
After clearing positive denominators its numerator is
\[
 F(n,m)=5n(n-3)\bigl(4m(n-5)-(n-1)(n-2)\bigr)
                         -6m(n-1)(n-2).                 \tag{ALT10.4}
\]
Its coefficient of $m$ is positive for $n\ge36$:
$(n-3)(n-5)>n$ makes its first contribution exceed
$20n^2$, whereas the subtracted contribution is less
than $6n^2$. Because $m\ge(n+3)/4$, it follows that
\[
 \begin{split}
 F(n,m)&\ge F\bigl(n,(n+3)/4\bigr)\\
 &=\frac72n^3-100n^2+\frac{531}2n-9\\
 &=\frac72(n-36)^3+278(n-36)^2
                 +\frac{13347}2(n-36)+43245>0.
 \end{split}                                             \tag{ALT10.5}
\]
This proves (ALT10.3), which contradicts the degree floor
$3/(5e)\le6/(5n(n-3))$. Hence $n\ge4m-2=20h-2$.

At $n=4m-2$ the exact difference is
\[
 \frac{4(n-5)}{(n-1)(n-2)}-\frac1m
                    =\frac{-12}{m(n-1)(n-2)}<0.          \tag{ALT10.6}
\]
The known two-row degree formula, in the already cited
James--Liebeck Theorem~29.13 \cite{HD:JamesLiebeck2001}, gives
\[
 a=\binom n4-\binom n3=\frac{n(n-1)(n-2)(n-7)}{24},
 \qquad
 a-\frac{n(n-3)}2
       =\frac{n(n+1)(n^2-11n+22)}{24}>0                 \tag{ALT10.7}
\]
for $n\ge38$. The strict long-row condition also holds:
at $n=38$ it is $4<19(1-\sqrt{3/5})$, equivalent
to $3\cdot19^2<5\cdot15^2$, and its right side increases
with $n$. Thus the endpoint label has all the required
numerical properties, including for every positive tensor
power. Nothing in this verification supplies a generating
tuple, period-map rank or effective detecting cycle.
\end{proof}

For the eight non-Eisenstein class-number-one fields, the
resulting necessary lower bounds are
\[
\begin{array}{c|r|r}
K&h&n\text{ lower bound}\\\hline
\mathbf Q(i)&2&38\\
\mathbf Q(\sqrt{-2})&4&78\\
\mathbf Q(\sqrt{-7})&6&118\\
\mathbf Q(\sqrt{-11})&10&198\\
\mathbf Q(\sqrt{-19})&18&358\\
\mathbf Q(\sqrt{-43})&42&838\\
\mathbf Q(\sqrt{-67})&66&1318\\
\mathbf Q(\sqrt{-163})&162&3238
\end{array}                                               \tag{ALT10.8}
\]
This presents the reduced condition-A supplier gate
(cote16, scan050; SIGN\_NEXT SN-C03/C04) in effective integers,
with its existing finite-cover, local-system, representation
lambda-ring and K-theory interpretation. It changes no
cycle-supply premise and adds no Markman, Severn or032 input.

\section{An effective obstruction for every fixed ordinary alternating tail}
\label{sec:effective-fixed-tail}

\paragraph{Status, origin and exact scope.}
New derivation D-AFT01--03, Q2ct. Complete author and root
whole-proof reviews and the separate /root/imprimitive proof,
primary-source and independent coefficient reconstruction have passed.
The cyclic character-defect method originates in the root's
HD-AT301 proof. The exact tensor-cycle and actual quotient-cover
framework is D-ATT07--09. The new assertion here is uniform over
every fixed tail partition, with an explicit computable threshold;
it does not assert a threshold uniform over all tail sizes.

Fix a partition $\mu\vdash k$, where $k\ge3$. For sufficiently
large $n$, let $W_n=S^{(n-k,\mu)}$ be the rational ordinary
$S_n$ representation. We consider an actual rank-nine sixfold
inducing source whose irreducibly acting ordinary core is $A_n^t$
on $W_n^{\otimes t}$, with the inducing image transitive on the
simple factors. All compatible scalar and outer actions and all
factor permutations are allowed. For $t=1$ the conclusion holds
over every original imaginary quadratic field. For $t\ge2$ we
assume that the inducing character field contains a non-Eisenstein
imaginary quadratic field. Equality with the original field is
not assumed.

\subsection*{D-AFT01. The coefficient constants and uniform spectrum bound}

The published character-polynomial formula of Garsia--Goupil,
Proposition~I.1, printed p.~2, with its complete proof in
equations (2.1)--(2.6), pp.~9--11, gives
\[
 \chi_{W_n}(x)=q_\mu(X_1(x),\ldots,X_k(x))
 \quad(n-k\ge\mu_1).                                  \tag{AFT.1}
\]
Its exact rational polynomial is computed by the displayed
finite umbral formula in that proposition, or its recursive
Corollary~I.1. We use this published formula, not an unverified
asymptotic character assertion \cite{HD:GarsiaGoupil2009}.

Write $f=X_1$ and expand the polynomial in ordinary monomials:
\[
 q_\mu=Cf^k+R(f,X_2,\ldots,X_k),\qquad
 C=\frac{f^\mu}{k!}>0.                                 \tag{AFT.2}
\]
Every monomial has weighted degree at most $k$, where $X_i$
has weight $i$. The displayed leading term is the only monomial
of ordinary total degree $k$. Indeed each product
$\prod_i(iX_i-1)^{a_i}$ in the published formula has weighted
degree at most $\sum ia_i=k$, and the falling-factorial
substitution only lowers this degree. Ordinary degree $k$
therefore requires the monomial $X_1^k$, arising from the
identity class of $S_k$ with coefficient $f^\mu/k!$.

Let $r_\alpha$ be the coefficients of $R$, including its
constant coefficient, and set
\[
 \begin{aligned}
 M&=1+C^{-1}\sum_\alpha|r_\alpha|\ge1,
 &F&=\lceil kM\rceil+1,\\
 N_\mu&=\left\lceil\max\left\{
 2k+2,\ 200,\ \frac{100}{C},\ 100kM,
 [100(1+M)]^5,\ [100(1+M)F^k]^2
 \right\}\right\rceil .
 \end{aligned}                                         \tag{AFT.3}
\]
These are explicit rational computations and integer ceilings
from a finite published polynomial. In particular the threshold
does not conceal an ineffective asymptotic existence statement.

For every integer $n\ge N_\mu$, define
\[
 \begin{gathered}
 a=q_\mu(n,0,\ldots,0)=\dim W_n,\qquad S=Cn^{k-1},\\
 b=S/100,\qquad c=a/n-Cn^{k-2}/100,\\
 A=a/n-101Cn^{k-2}/100,\qquad B=2S/5,
 \qquad d_0=A+B.
 \end{gathered}                                        \tag{AFT.4}
\]
The symbol $S$ in (AFT.4) is a positive real number, not a
finite group. All displayed bounds below are rational for
integer $n$.

\begin{theorem}[D-AFT01]
For every nonidentity permutation $x\in S_n$, with natural
index $I(x)=n-\#\mathrm{cycles}(x)$,
\[
 \nu(W_nx)\ge AI(x)+B.                                 \tag{AFT.5}
\]
Here $\nu$ is maximum-eigenspace codimension, so the estimate
holds for every target eigenvalue and survives every scalar
twist. For every permutation, including the identity,
\[
 \chi_{W_n}(x)\ge-b.                                   \tag{AFT.6}
\]
Moreover
\[
 \begin{gathered}
 S\ge100,\qquad \tfrac{99}{100}nS\le a\le\tfrac{101}{100}nS,
 \qquad \tfrac9{10}S\le A\le\tfrac{101}{100}S,\\
 \tfrac{13}{10}S\le d_0\le\tfrac{141}{100}S,
 \qquad B\ge40,\quad b\le d_0,\quad 2d_0\le a-1.
 \end{gathered}                                        \tag{AFT.7}
\]
\end{theorem}

\begin{proof}
Write $q_\mu=Cf^k+P(f)+Q(f,X_2,\ldots,X_k)$, where $P$
contains the lower pure-$f$ terms and each monomial of $Q$
contains at least one other variable. Their total absolute
coefficient mass is at most $CM$.
Since $|P(n)|\le CMn^{k-1}$ and $M/n\le1/100$, the degree
bounds in (AFT.7) hold. The threshold also gives $S\ge100$.

First prove the global negative-character bound. If
$0\le f\le n^{3/5}$, a monomial $f^r\prod_{i\ge2}X_i^{a_i}$,
with $s=\sum_{i\ge2}a_i$, satisfies $r+2s\le k$ and has
absolute value at most $n^{3r/5+s}\le n^{3k/5}$.
The same estimate holds for the leading term. Therefore
\[
 |q_\mu|\le C(1+M)n^{3k/5}\le S/100.                  \tag{AFT.8}
\]
The last inequality uses $k\ge3$ and
$n^{1/5}\ge100(1+M)$.
If $f\ge n^{3/5}$, each lower pure term divided by $f^k$
is at most $n^{-3/5}$, and each mixed monomial divided by
$f^k$ is at most
$(n/f^2)^s\le n^{-1/5}$. Hence
$|R|\le CMn^{-1/5}f^k<Cf^k/100$ and $q_\mu>0$.
This proves (AFT.6) in all cases.

We next prove the particularly sharp global defect estimate
\[
             a-\chi_{W_n}(x)\ge c(n-f).                \tag{AFT.9}
\]
Put $m=n-f$. Every mixed monomial has ordinary degree at
most $k-1$. Selecting one non-$f$ variable and using
$X_i\le m/i\le m$, while bounding the remaining variables
by $n$, gives absolute value at most $mn^{k-2}$.
For a pure lower monomial of degree $r\ge2$,
\[
 \left|\frac fn n^r-f^r\right|
 =f(n^{r-1}-f^{r-1})\le(k-2)mn^{k-2}.
\]
For $r=1$ the difference is zero; for $r=0$ it is $m/n$,
also bounded by $(k-2)mn^{k-2}$. Consequently
\[
 \frac fn a-q_\mu
 \ge Cfm n^{k-2}-CM(k-2)mn^{k-2}\ge0
 \quad\hbox{if }f\ge F.                               \tag{AFT.10}
\]
Here the first positive term follows from
$fn^{k-1}-f^k\ge fm n^{k-2}$. The identity $m=0$ is
included without division.

For the remaining integer values $0\le f<F$, the weighted
degree bound instead gives
\[
 |q_\mu|\le C(1+M)F^k n^{k/2}\le S/100,                \tag{AFT.11}
\]
because $\sqrt n\ge100(1+M)F^k$ and
$(k+1)/2\le k-1$. If $f\ge1$, the lower degree bound
gives $af/n\ge99S/100$, so again $q_\mu\le af/n$.
If $f=0$, (AFT.11) gives $q_\mu\le S/100$.
These cases prove (AFT.9), with exactly the $c$ in (AFT.4).
The use of integer $f\ge1$ here is explicit.

We also need a stronger defect when many letters are fixed.
For a pure term of degree $r\le k-1$,
$|n^r-f^r|\le(k-1)mn^{k-2}$; the constant cancels.
Together with the mixed-monomial estimate already proved,
this gives
\[
 a-q_\mu\ge (n-f)c_f,\qquad
 c_f=C\left(\sum_{j=0}^{k-1}n^{k-1-j}f^j
                  -(k-1)Mn^{k-2}\right).               \tag{AFT.12}
\]
The function $c_f$ is nondecreasing for $f\ge0$. The upper
degree estimate before replacing $M/n$ by $1/100$ yields
\[
                 c_f-A\ge Cn^{k-2}(f-kM).              \tag{AFT.13}
\]
Indeed retain just the $j=0,1$ terms in the sum, use
$a/n\le Cn^{k-1}+CMn^{k-2}$, and discard the positive
$101Cn^{k-2}/100$.

Let $x$ have order $L$ and index $I\ge1$. If $m_z$ is a
target-eigenvalue multiplicity, cyclic averaging in (AFT.9)
gives $a-m_1\ge cI$. The real character bound (AFT.6) gives
\[
 m_z-m_1=\frac1L\sum_{j=0}^{L-1}\chi(x^j)
                    (\operatorname{Re}z^{-j}-1)\le b
 \quad(z\ne1);                                        \tag{AFT.14}
\]
the trivial target needs no estimate. Thus $\nu\ge cI-b$.
If $I\ge9n/20$, then $c-A=Cn^{k-2}$ implies
\[
                  \nu-AI\ge(9/20-1/100)S=11S/25>B.
                                                               \tag{AFT.15}
\]

Suppose instead $I<9n/20$. Every power of $x$ fixes at least
$f_0=n-2I>n/10$ letters, because the moved support is at most
twice the natural index. Averaging (AFT.12) gives
$a-m_1\ge c_{f_0}I$. Since $kM\le n/100$, (AFT.13)
gives $c_{f_0}-A\ge9S/100$. If $I\ge5$, this again yields
$\nu-AI\ge11S/25>B$. If $1\le I\le4$, then
$f_0\ge n-8\ge9n/10$ and the stronger bound is
$c_{f_0}-A\ge89S/100$. Therefore
$\nu-AI\ge88S/100>B$. This covers all nonidentity
permutations and every target eigenvalue.

Finally $n\ge200$ gives
$A\ge(99/100-101/(100n))S\ge9S/10$; its upper bound is
immediate. Adding $B$ proves the $d_0$ bounds. Since
$S\ge100$, $n\ge200$, and $a\ge99nS/100$, all remaining
comparisons in (AFT.7) follow directly. In particular
$a\ge19800$, $d_0\ge130$, and $2d_0\le a-1$.
\end{proof}

\subsection*{D-AFT02. One ordinary core in every field}

\begin{theorem}[D-AFT02]
For $n\ge N_\mu$, a rank-nine sixfold inducing source
cannot have the single ordinary core $A_n$ on $W_n$.
This holds over every original imaginary quadratic field,
with all compatible scalar and outer actions.
\end{theorem}

\begin{proof}
The threshold $n\ge2k+2$ makes $(n-k,\mu)$ non-self-conjugate:
its first row exceeds its number of rows. The restriction
to $A_n$ is therefore irreducible. It is nontrivial and
faithful because $A_n$ is simple. Its rational $S_n$ model
has trivial scalar intersection: a scalar centralizes the
faithful $A_n$ image, while $C_{S_n}(A_n)=1$.
The accepted ordinary alternating normalizer interface
therefore extracts an actual permutation homomorphism.
Its image contains $A_n$, so the associated degree-$n$
quotient cover is connected.

The exact intermediate Euler identity is
$6=2a(h-1)+D$. Since $a\ge19800$, genus at least two is
impossible. At genus one there are at most six genuine
marks and the actual period rank is at most six. Thus
rank nine requires $h=0$, at least twelve genuine inertia
entries, and $D=2a+6$. Riemann--Hurwitz for the connected
degree-$n$ quotient gives $\sum I(x_i)\ge2n-2$.

Every nonscalar inertia moves at least $AI(x_i)+B$
dimensions by D-AFT01. A genuine scalar has full moving
rank $a\ge B$ and index zero. Hence
\[
 \begin{split}
 D-2a&\ge2A(n-1)+12B-2a\\
     &=12B-\left(\frac{2a}{n}
             +\frac{202}{100}C(n-1)n^{k-2}\right)\\
     &\ge\left(\frac{24}{5}-\frac{101}{25}\right)S
       =\frac{19}{25}S\ge76>6.
 \end{split}                                           \tag{AFT.16}
\]
This contradicts the exact source budget. No condition on
the scalar cyclotomic field is needed for this case.
\end{proof}

\subsection*{D-AFT03. Every transitive tensor count}

\begin{theorem}[D-AFT03]
For $n\ge N_\mu$ and every $t\ge2$, a finite inducing
image normalizing $A_n^t$ on $W_n^{\otimes t}$ transitively
cannot give a rank-nine sixfold source if its character
field contains a non-Eisenstein imaginary quadratic field.
All compatible scalar phases, outer actions and factor
permutations are included.
\end{theorem}

\begin{proof}
Put $e=a^t$. The rational tensor-wreath normalizer meets
the scalars trivially by its conjugation action on the
faithful simple factors. Thus every inducing operator
has actual scalar and rational coordinates, and the
natural $nt$-letter quotient is transitive. The inducing
character field is contained in the cyclotomic field of
the scalar character. If this cyclotomic field contains
a non-Eisenstein imaginary quadratic field, its scalar
image is not contained in $\mu_6$. Generation and product
one in the scalar quotient consequently force at least
two inertia phases outside $\mu_6$. Only field inclusion
is used. The actual intermediate-cover Euler and Hurwitz
arguments give $h=0$, $r\ge12$, and $D=2e+6$.

Every power of a rational tensor-wreath operator has trace
at least $-(b/a)e$. Indeed its trace is a product of traces
of the factor-cycle products. A negative product contains
one negative factor of absolute value at most $b$, and
all other factors have absolute value at most $a$.
The rational cyclotomic multiplicity bound, together with
the fourth-root square-trace argument of D-ATT02, therefore
gives moving rank at least $(3-b/a)e/4$ at the two marked
scalar phases.

Every genuine inertia has moving rank at least $(d_0/a)e$.
For an operator fixing factors, use D-AFT01 on any active
factor and spectral convolution; a genuine scalar has
full rank. For an operator moving factors, the exact
weighted-word estimate gives the stronger
$(a-1)e/(2a)$, which is at least $(d_0/a)e$ by (AFT.7).
If at least three entries move factors, one is distinct
from both scalar marks. With nine other genuine entries,
\[
 D\ge2e+a^{t-1}\left(9d_0-\frac{b+1}{2}\right)>2e+6,
                                                               \tag{AFT.17}
\]
using $b\le d_0$ and $d_0\ge130$. Thus exactly two entries
move factors: zero contradicts transitivity and one
contradicts product one. Their inverse factor permutations
are full $t$-cycles.

Let $H$ be the total natural index of their two cycle
products, and $N$ the sum of the natural factor indices
at the remaining entries. Riemann--Hurwitz for the actual
connected degree-$nt$ quotient gives
$2n(t-1)+H+N\ge2nt-2$, hence $H+N\ge2n-2$.
At each genuine factor-fixing entry, average the D-AFT01
bounds over its active factors, and then enlarge the
denominator to $t$. This gives
$a^{t-1}(AI/t+B)$ with the entire positive constant $B$.
Scalars satisfy the same estimate at index zero.

At each full factor cycle the exact D-ATT08 estimate is
\[
 (1-1/t)e+\frac{a^{t-1}}t\nu(W_nx)-\frac{R_t(a)}t,
 \qquad R_t(a)=\sum_{j=1}^{t-1}a^{\gcd(t,j)}.            \tag{AFT.18}
\]
Only $\nu(W_nx)\ge AI(x)$ is used here, so identity cycle
products are allowed. The proved errors are $R_2(a)=a$
and $R_t(a)\le2a^{t-2}$ for $t\ge3$; in particular
$2R_t(a)/a^{t-1}\le2$ for all $t\ge2$.
There are at least ten factor-fixing entries, and summation
therefore yields
\[
 D-2e\ge\frac{a^{t-1}}t
       \bigl(2A(n-1)-2a+10tB-2\bigr)>B a^{t-1}>6.
                                                               \tag{AFT.19}
\]
For the strict comparison, the parenthesis is greater
than $tB$ because
\[
 9tB-\bigl(2a-2A(n-1)\bigr)-2
 \ge\left(\frac{36}{5}-\frac{101}{25}\right)S-2
 =\frac{79}{25}S-2>0.                                  \tag{AFT.20}
\]
This contradicts the exact Euler budget and completes
the all-count proof.
\end{proof}

\paragraph{Effective remaining-tail statement and strictness.}
For any fixed integer $K\ge3$, take the maximum of the
explicit $N_\mu$ over the finitely many partitions with
$3\le|\mu|\le K$. Above this effective maximum, the
non-Eisenstein transitive sources in D-AFT03 cannot have
tail size between three and $K$; D-AFT02 gives the same
statement for a single factor over every original
imaginary quadratic field. In the project's main
\emph{nonsplit} problem, combine this with the previously
proved standard and two-box exclusions in their stated
scopes. The tail size in any remaining non-Eisenstein
transitive ordinary alternating sequence with $n\to\infty$
must then tend to infinity. The single-factor standard
case still uses its separate actual-polarization splitness
gate. Transposing a label gives the same $A_n$
representation, so either long-row convention may be
used. No finite bound on $n$ is asserted when the tail
size itself grows.

For example, $\mu=(4)$ and any $n\ge N_{(4)}$ give a
label $(n-4,4)$ not among the previously excluded tails
of size at most three. Since $N_{(4)}\ge200$, its earlier
Gaussian long-row fraction
$4(n-5)/((n-1)(n-2))$ is below $1/10$; the earlier
necessary shape gate did not remove this entire family.
The new theorem removes it and every other fixed tail
eventually, without asserting that an actual source
previously existed. It supplies a new obstruction to
condition-A finite-cover suppliers, not a new algebraic
cycle or an unrefereed geometric premise.

\subsection{HD-TW01: averaging all tensor cycles without a field restriction}
\paragraph{Status.}
Complete root and separate /root/g34\_eightfolds whole-proof reviews
have passed, first in full message form during the outage and then on the
complete frozen TeX file after recovery. Source and exact finite-check
scopes are recorded separately; the final published4000 sketch is not
claimed as an independently recomputed certificate.
\paragraph{Actual-source hypotheses.}
Let an actual connected sphere source have a genuine generating product-one
inertia tuple \(T_1,\ldots,T_r\), period differential of rank nine, coefficient
degree \(e=a^t\), and transitive ordinary alternating core \(A_n^t\), with
\(t\ge2\) and \(a\ge4\).
Its core module is \(W^{\otimes t}\), where the compatible factor normalizers
act, up to scalars, through an irreducible \(S_n\)-module \(W\) of dimension \(a\).
Both the factor action on \(t\) letters and the natural block action on \(nt\)
letters are the actual quotient actions of this source.  In particular
\[
 r\ge12,\qquad D:=\sum_i\operatorname{rank}(1-T_i)=2e+6.
 \tag{TW.1}
\]
For a finite-order operator \(U\), write
\(\nu(U)=\dim U-\max_z\dim\ker(U-z)\).
Suppose that, for every nonidentity \(x\in S_n\),
\[
 \nu(W(x))\ge A\,\operatorname{ind}(x)+B
 \tag{TW.2}
\]
and that
\[
 \begin{gathered}
 A\ge0,\quad B>0,\quad d_0=A+B\ge1,\quad 2d_0\le a-1,\\
 A(n-1)+B\le a,\quad
 \lambda=1-\frac{nA+1}{a}\ge0,\quad
 L=a-(n-1)A\ge1,\quad a(9B-L-1)>6.
 \end{gathered}\tag{TW.3}
\]
Then the source does not exist.  This statement imposes no condition on the
imaginary quadratic field or scalar conductor.

\paragraph{Proof.}
A factor-fixing noncentral entry has \(m\ge1\) active factor operators
\(W(x_j)\), up to scalars.  Tensor spectral convolution bounds its normalized
projective support below by the maximum of the active normalized supports.
Averaging those \(m\) bounds and using \(m\le t\) gives
\[
 \frac{\operatorname{rank}(1-T_i)}e
 \ge \frac{A I_i}{ta}+\frac Ba ,
 \qquad I_i=\sum_j\operatorname{ind}(x_j).
 \tag{TW.4}
\]
A genuine scalar entry has actual moving rank \(e\), and \(I_i=0\), so it too
satisfies (TW.4).  This uses actual moving rank for such a scalar; its
projective support is zero.  Every genuine factor-fixing entry also has
moving rank at least \(d_0e/a\).
The accepted weighted-cycle minimum in D-ATT08 gives every factor-moving
entry moving rank at least \((a-1)e/(2a)\), hence at least \(d_0e/a\).
If four entries move factors, choose four of them and eight other genuine
entries.  Equation (TW.1) would imply
\[
 D\ge \frac{4(a-1)e}{2a}+\frac{8d_0e}a
   =2e+(8d_0-2)e/a
   \ge2e+6e/a>2e+6.
 \tag{TW.5}
\]
Thus at most three entries move factors, and at least nine fix them.

For any entry, let \(\pi\) be its factor permutation.  Its cycles have lengths
\(\ell_j\), with represented \(S_n\)-cycle products \(x_j\).  The exact
short-word estimate of D-ATT08 gives each cycle operator normalized support
at least
\[
 s_j\ge1-\frac1{\ell_j}
       +\frac{\nu(W(x_j))}{\ell_j a}
       -\frac{R_{\ell_j}(a)}{\ell_j a^{\ell_j}}.
 \tag{TW.6}
\]
Here \(R_1=0,\ R_2=a\), and
\(R_\ell(a)\le2a^{\ell-2}\) for \(\ell\ge3\), in the accepted range \(a\ge4\).
Consequently \(R_\ell(a)/a^\ell\le1/a\) for every \(\ell\ge2\).
Convolution across cycle operators gives total normalized support at least
\(\max_j s_j\), hence at least \(\sum_j(\ell_j/t)s_j\).
There are at most \(\operatorname{ind}(\pi)\) nontrivial cycles.
Also (TW.2) implies \(\nu(W(x))\ge A\operatorname{ind}(x)\) including \(x=1\).
The actual natural block index is
\[
 I_i=n\operatorname{ind}(\pi)+\sum_j\operatorname{ind}(x_j).
 \tag{TW.7}
\]
Substitution yields
\[
 \frac{\operatorname{rank}(1-T_i)}e
 \ge\frac{A I_i}{ta}
       +\frac{\lambda\,\operatorname{ind}(\pi)}t .
 \tag{TW.8}
\]
For a factor-fixing entry the additional \(B/a\) from (TW.4) remains available.
The two connected quotient covers satisfy
\(\sum_iI_i\ge2nt-2\) and
\(\sum_i\operatorname{ind}(\pi_i)\ge2t-2\).
Since \(A,\lambda\ge0\), summation gives, with \(r_{\rm fix}\ge9\),
\[
 \begin{split}
 D-2e
 &\ge\frac ea\left[r_{\rm fix}B-2-\frac2t(L-1)\right]\\
 &\ge\frac ea(9B-L-1)
 \ge a(9B-L-1)>6.
 \end{split}\tag{TW.9}
\]
This contradicts (TW.1), proving the assertion. \(\square\)

\paragraph{Scope of the mechanism.}
The weights \(\ell_j/t\) preserve the complete natural block index while
charging at most \(1/a\) for each nontrivial cycle.  No signed character floor
or distinguished scalar phase is a premise.  The proof applies only to
actual quotient actions and genuine source entries; it makes no existence
claim from a numerical profile.

\subsection{HD-TW02: an absolute-character criterion for every tensor count}
\paragraph{Statement.}
Let \(W\) be an irreducible \(S_n\)-module of degree \(a\), and let \(f(y)\)
denote the number of fixed letters of \(y\in S_n\).
Suppose \(n\ge24,\ a\ge n^3\), and
\[
 \frac{|\chi_W(y)|}{a}\le
       \left(\frac{\max\{f(y),1\}}n\right)^3
       \qquad(y\in S_n).
 \tag{TW.10}
\]
Then, for every nonidentity \(x\in S_n\),
\[
 \nu(W(x))\ge A\operatorname{ind}(x)+B,\qquad
 A=\frac{a(n-4)}{n^2},\quad B=\frac an.
 \tag{TW.11}
\]
No actual rank-nine sphere source with the compatible transitive core module
\(W^{\otimes t}\) exists, for any \(t\ge1\), over any imaginary quadratic
field, under the actual quotient hypotheses of HD-TW01.

\paragraph{Proof of the local bound.}
Put \(I=\operatorname{ind}(x)\), \(q=\operatorname{ord}(x)\),
and \(f_j=f(x^j)\), for \(0\le j<q\).
Fourier inversion on the cyclic group, followed by the triangle inequality
and (TW.10), gives
\[
 \frac{\nu(W(x))}{a}\ge
 1-\frac1q\sum_{j=0}^{q-1}
       \left(\frac{\max\{f_j,1\}}n\right)^3,\qquad
 \frac1q\sum_jf_j=n-I .
 \tag{TW.12}
\]
The second identity counts the orbits of \(\langle x\rangle\) on the letters.
If \(I\ge n/3\), use
\((\max\{f,1\}/n)^3\le f/n+n^{-3}\) to obtain
\[
 \nu(W(x))/a-AI/a
 \ge4I/n^2-n^{-3}
 \ge4/(3n)-n^{-3}>1/n.
 \tag{TW.13}
\]
If \(I<n/3\), every power has at least \(n-2I>n/3\) fixed letters:
the moved support of \(x\) has size at most \(2I\).
Thus \(u_j=f_j/n>1/3\), and
\(1-u_j^3\ge(13/9)(1-u_j)\).
For \(I\ge3\) this gives
\[
 \nu(W(x))/a-AI/a
 \ge 4I/(9n)+4I/n^2>1/n.
 \tag{TW.14}
\]
Finally, if \(I=1\) or \(2\), then \(u_j\ge1-4/n\ge5/6\), so
\(1+u_j+u_j^2\ge91/36>5/2\).
The resulting margin is at least
\(3I/(2n)+4I/n^2>1/n\).
These three cases prove (TW.11) for every target eigenvalue.

\paragraph{Actual-source contradiction.}
For the single factor the actual \(n\)-letter quotient and \(r\ge12\) give
\[
 D-2a\ge2A(n-1)+12B-2a
       =\frac{a(2n+8)}{n^2}>6 .
 \tag{TW.15}
\]
Genuine scalar entries also satisfy the required local rank bound:
their actual rank is \(a\), their natural index is zero, and \(B\le a\).
For \(t\ge2\), all hypotheses of HD-TW01 follow from
\[
 \begin{gathered}
 d_0=\frac{a(2n-4)}{n^2},\quad
 A(n-1)+B=a(1-2/n)^2,\quad
 \lambda=4/n-1/a,\\
 L=\frac{a(5n-4)}{n^2},\qquad
 9B-L-1=\frac{4a(n+1)}{n^2}-1 .
 \end{gathered}\tag{TW.16}
\]
Indeed \(d_0\ge1,\ L\ge1,\ \lambda>0\) because \(a\ge n^3\);
\(2d_0\le a-1\) follows from
\(a(n^2-4n+8)/n^2\ge n(n^2-4n+8)>1\).
The final expression in (TW.16), multiplied by \(a\), is larger than \(6\).
This proves the tensor assertion. \(\square\)

\paragraph{Published effective input.}
For \(n\ge4000\), the additional condition \(a\ge n^{120}\) implies
(TW.10).
The input is M\"uller--Schlage-Puchta,
\emph{Character theory of symmetric groups, subgroup growth of Fuchsian
groups, and random walks}, Advances in Mathematics 213 (2007), 919--982,
Theorem 2, with the exact normalization also stated in M\"uller,
\emph{Character theory of symmetric groups, analysis of long relators,
and random walks}, S\'eminaire Lotharingien de Combinatoire 56 (2006),
Article B56a, Section 5.1 equation (24), Theorem 9, and Lemma 7 equation (30).
The original author version locates Theorem 2 on page 8 and its explicit
\(4000\) discussion on pages 16--17.  That discussion is a proof sketch;
we use the published theorem, not an independently recomputed version of
its final numerical sketch.\cite{HD:MSP2007,HD:MullerSLC2006}

\paragraph{Proof of the stated consequence.}
For a nonidentity conjugacy class \(c\), let \(t_c\) be the least positive
even number of independent uniform elements of \(c\) having no common
fixed letter with probability at least \(1-1/n\).
The published theorem gives statistical mixing time at most \(10t_c\).
Its squared relative \(L^2\) normalization and character orthogonality give,
at the even time \(10t_c\),
\[
 \sum_{\chi(1)>1}\chi(1)^2
       \left|\frac{\chi(c)}{\chi(1)}\right|^{20t_c}
       \le \frac2n<1 .
 \tag{TW.17}
\]
If the first mixing time is earlier, the same bound at \(10t_c\) follows
because every normalized character ratio has modulus at most one.
For the particular degree \(a\), (TW.17) implies
\(|\chi_W(c)|/a<a^{-1/(10t_c)}\).

If \(1\le f=f(c)<n\), the elementary union bound for a common fixed letter
is \(n(f/n)^u\) at time \(u\).  Hence the even rounding gives
\[
 t_c\le2\left\lceil\frac{\log n}{\log(n/f)}\right\rceil
 \le\frac{4\log n}{\log(n/f)} .
 \tag{TW.18}
\]
It follows that
\[
 |\chi_W(c)|/a
 \le a^{-\log(n/f)/(40\log n)}
 \le (f/n)^3 .
 \tag{TW.19}
\]
For \(f=0\), \(t_c=2\), so the ratio is less than
\(a^{-1/20}\le n^{-6}\le n^{-3}\).
The identity has ratio one.  This proves (TW.10), including fixed-point-free
powers and the identity power, and finishes the consequence. \(\square\)

\subsection{HD-TW03: one effective ordinary alternating cutoff}
\paragraph{Statement and exact main scope.}
Set
\[
 \begin{gathered}
 M_* = 1+241\cdot240!\cdot2^{239},\qquad
 F_* = 240M_*+2,\\
 \mathcal N=\bigl[100(1+M_*)F_*^{240}\bigr]^2 .
 \end{gathered}\tag{TW.20}
\]
This is an explicit positive integer.  In the main nonsplit sixfold gate,
over a non-Eisenstein original imaginary quadratic field and within
the primitive factorwise tensor-indecomposable transitive-core scope
of D-BI19, no actual transitive ordinary alternating factor \(A_n\) has \(n\ge\mathcal N\).
All tensor counts and compatible scalar phases are included.
The earlier treatment of the single standard module and of small
alternating degrees is retained with its exact scope.
In particular the previously unbounded growing-tail ordinary alternating
range becomes finite.  This is a necessary-source reduction, not a new
detector-cycle construction or an executed enumeration.

\paragraph{A coefficient bound uniform over small tails.}
Let \(\mu\vdash k\), \(k\ge3\), and write the stable character polynomial as
\(q_\mu=Cx_1^k+R_\mu\), with \(C=f^\mu/k!\).
For the coefficient \(\ell^1\)-norm, the elementary binomial polynomial
satisfies
\[
 \left\|\binom{x}{m}\right\|_1=1\qquad(m\ge0).
 \tag{TW.21}
\]
For \(m\ge1\), its coefficients have alternating signs and their absolute
sum is
\(\prod_{j=0}^{m-1}(1+j)/m!=1\); the case \(m=0\) is constant one.

The Garsia--Goupil identity used in D-AFT01, equation (2.8),\cite{HD:GarsiaGoupil2009} expands
\(s_{(n-k,\mu)}\) as
\(\sum_{r=0}^k(-1)^r h_{n-k+r}e_r^\perp s_\mu\).
Here \(e_r^\perp s_\mu\) is the sum of \(s_\nu\) over vertical strips
\(\mu/\nu\) of size \(r\).
For \(j=k-r\), the associated induced character polynomial is
\[
 \sum_{\substack{m_1,m_2,\ldots\ge0\\ \sum i m_i=j}}
  \chi^\nu(1^{m_1}2^{m_2}\cdots)
  \prod_i\binom{x_i}{m_i}.
 \tag{TW.22}
\]
This is the invariant-\(j\)-subset character formula: each fixed subset
is a union of cycles, and its coefficient is the trace of its induced
\(S_j\)-action.  Equation (TW.21) and
\(|\chi^\nu|\le f^\nu\) bound its coefficient norm by \(p(j)f^\nu\).

By the ordinary branching rule,
\(\sum_{\mu/\nu\ {\rm vertical},\,|\nu|=j}f^\nu\le f^\mu\):
these are the dimensions of the sign-isotypic \(S_r\)-multiplicity spaces
inside the restriction of the \(S_k\)-module of degree \(f^\mu\).
Consequently
\[
 \|q_\mu\|_1
 \le f^\mu\sum_{j=0}^kp(j)
 \le f^\mu(k+1)2^{k-1}.
 \tag{TW.23}
\]
The last inequality uses \(p(0)=1\) and the injection of partitions of a
positive integer \(j\) into its \(2^{j-1}\) compositions.
The AFT coefficient parameter therefore satisfies
\[
 M_\mu:=1+C^{-1}\|R_\mu\|_1
 \le 1+k!(k+1)2^{k-1}.
 \tag{TW.24}
\]
Thus, for \(3\le k\le240\),
\(M_\mu\le M_*\), \(C^{-1}\le240!\), and
\(F_\mu=\lceil kM_\mu\rceil+1\le F_*\).
Substitution into the full accepted threshold
\[
 N_\mu=\left\lceil\max\left\{
 2k+2,200,100/C,100kM_\mu,
 [100(1+M_\mu)]^5,
 [100(1+M_\mu)F_\mu^k]^2
 \right\}\right\rceil
 \tag{TW.25}
\]
shows \(N_\mu\le\mathcal N\).
For clarity, \(F_*\ge100(1+M_*)\), so the last term of (TW.20) dominates
also the fifth power and every displayed linear term.  It also exceeds
\(482^2=232324\) and \(4000\).

\paragraph{The fixed-tail estimates in every tensor count.}
For \(n\ge N_\mu\), D-AFT01--02 gives
\(S=Cn^{k-1}\ge100\),
\(.99nS\le a\le1.01nS\),
\[
 A=a/n-101S/(100n),\quad B=2S/5,\quad
 .9S\le A\le1.01S,\quad1.3S\le A+B\le1.41S.
 \tag{TW.26}
\]
These are precisely local all-eigenvalue bounds, so they are valid before
any conductor argument.
For HD-TW01,
\(\lambda=(1.01S-1)/a>0\),
\[
 L=a/n+1.01S(1-1/n),\qquad
 1\le L\le2.02S,\qquad
 9B-L-1\ge1.58S-1>0.
 \tag{TW.27}
\]
Since \(n\ge200\), one has
\(2(A+B)\le2.82S\le a-1\),
\(L\ge .99S>B\), and \(a(1.58S-1)>6\).
Thus all remaining hypotheses of HD-TW01 hold.
It follows that these same thresholds exclude every transitive tensor
count over every imaginary quadratic field.
The single-factor contradiction is already D-AFT02.
This extension does not change the separate standard-module premise.

\paragraph{The growing-tail range.}
Suppose now that an actual source in the stated main gate has
\(n\ge\mathcal N\).
The accepted HD-ALT-07/08 shape reduction, after transposing the partition
if necessary, gives an ordinary label \(\lambda=(n-k,\mu)\) with
\(k<n/4\).  The already accepted standard and two-box exclusions, together
with the three-box exclusions, leave \(k\ge4\).
If \(k\le240\), the preceding paragraph excludes the source.

If \(k\ge241\), the published degree bound of
M\"uller--Schlage-Puchta, Lemma 8(ii),\cite{HD:MSP2007} gives
\[
 a=f^\lambda\ge\binom{n-k}{k}f^\mu
              \ge\binom{n-k}{k}.
 \tag{TW.28}
\]
Let \(N_0=\lfloor3n/4\rfloor\).
Then \(n-k\ge N_0\), \(241\le k\le N_0/2\), and \(N_0\ge n/2\).
Monotonicity of binomial coefficients through the midpoint, followed by
the elementary product inequality
\(\binom{N_0}{241}\ge(N_0/241)^{241}\), gives
\[
 a\ge\binom{N_0}{241}
   \ge(n/482)^{241}
   \ge n^{120+1/2}\ge n^{120}.
 \tag{TW.29}
\]
Here \(n\ge482^2\) justifies the penultimate inequality.
The published effective mixing input and HD-TW02 exclude this case as well.
Every possible tail is covered, proving \(n<\mathcal N\). \(\square\)

\paragraph{Geometric and programme meaning.}
The indices in HD-TW01 are ramification indices of actual quotient covers
of the original Hurwitz source; the contradiction is a strengthened
Chevalley--Weil Euler budget for its weight-one isotypic local system.
The character-polynomial calculation uses induction and exterior-power
operations in the representation \(\lambda\)-ring.  It passes to the
corresponding group-algebra summands of curve cohomology and their
realizations, but asserts no additional algebraic cycle, motivic splitting,
weight argument, or special-value bound for a motivic \(L\)-function.
Together with HD-TC01--03, the now finite ordinary alternating catalogue
has a finite effective inducing-image and component-rank search.
Original-\(K\) descent and the actual rational polarization remain separate.

For Grothendieck's programme, this is a strict reduction of the possible
curve suppliers for the integral uniform form of condition A
(cote16, scan050, corrected in SIGN_NEXT SN-C03/C04).
The positivity of condition I (scan049 and cote44, Theorem4) still cannot
supply missing cycles.  The projector/division mechanism remains the
Cayley--Hamilton inverse in scans048--049 and the already accepted run
statements.  No use of Markman, Severn, or032 is a premise of HD-TW01--03.
A counterexample to the eventual uniform detector bound would still be
a ghost in the precise GHOSTS B03 sense.


\section{Two-box alternating tensor cores}
\label{sec:alternating-two-tail}

\paragraph{Status and scope.}
This complete new author derivation has passed whole-proof readings
by the root and independent /root/imprimitive reviewer, with separate
exact computational supplements. It concerns the actual tensor-core representations
with ordinary alternating labels $(n-2,2)$ and $(n-2,1,1)$, including
all compatible scalars and factor permutations. The inducing character
field may be larger than the original imaginary quadratic field.
Single factors retain the prior D-UP and D-EX exclusions; the new
source theorem concerns at least two factors.

For $n\ge15$, let $P_n$ be the permutation module on letters,
$U_n$ its deleted summand, and $P_{2,n}$ the permutation module
on unordered pairs. We use the following two rational orthogonal
modules, whose restrictions to $A_n$ are irreducible:
\[
\begin{array}{c|c|c|c}
W&\text{ordinary label}&a=\dim W&b\\\hline
W_{\mathrm{UP}}=P_{2,n}\ominus P_n&(n-2,2)&n(n-3)/2&1\\
W_{\mathrm{EX}}=\Lambda^2U_n&(n-2,1,1)&(n-1)(n-2)/2&\lfloor(n-1)/2\rfloor.
\end{array}                                                    \tag{ATT.1}
\]
The actual submodule and quotient realizations, including the
ordinary labels and irreducibility, are those in D-UP-01/03 and
D-EX-01/03. They are not virtual representations in what follows.

\begin{lemma}[D-ATT-01: exact minimum support and lower traces]
\label{D-ATT-01}
For every nonidentity $\pi\in S_n$,
\[
 \nu(W_{\mathrm{UP}}(\pi))\ge n-3,
 \qquad \nu(W_{\mathrm{EX}}(\pi))\ge n-2.                    \tag{ATT.2}
\]
Both bounds are attained by a transposition. Every permutation
has character at least $-b$ on the corresponding module in
\textup{(ATT.1)}.
\end{lemma}

\begin{proof}
Write $k=\operatorname{ind}(\pi)$ and
$D_\lambda=\operatorname{rank}(1-\lambda W(\pi))$.
The already proved all-order inequalities (UP.3)--(UP.4) give,
when $k\ge2$,
\[
 D_1\ge n-5+n/4\ge n-3,
 \qquad
 D_{\lambda\ne1}\ge n-5+(n-4)/2\ge n-3.                   \tag{ATT.3}
\]
Their hypotheses are precisely a nonidentity permutation in
the first inequality and a nonidentity target scalar in the
second; both hold here. If $k=1$, the permutation is a
transposition. Its negative eigenspace on $P_{2,n}$ has
dimension $n-2$, while that on $P_n$ has dimension one.
Thus the two eigenspace dimensions on $W_{\mathrm{UP}}$ are
$n-3$ and $a-(n-3)$, with the latter larger. This proves
the first exact minimum.

For the exterior square, (EX.4) gives
$D_1\ge n-4+n/3\ge n-2$ when $k\ge2$.
For every nonidentity target scalar, (EX.3) gives
\[
 D_{\lambda\ne1}\ge
 a-\lfloor(n-1)^2/4\rfloor\ge n-2.                       \tag{ATT.4}
\]
For the last inequality it suffices to omit the floor:
$a-(n-1)^2/4-(n-2)=(n^2-8n+11)/4>0$ for $n\ge15$.
A transposition on $U_n$ has a negative line and a positive
space of dimension $n-2$. Its exterior square has negative
space of dimension $n-2$ and a larger positive space, proving
the second exact minimum. Taking the minimum over all
$\lambda$ proves (ATT.2); targets which are not eigenvalues
have full moving rank and cause no exception.

Let $f_1$ and $f_2$ count the fixed letters and the two-cycles
of $\pi$. The actual permutation-module and exterior-square
character identities give
\[
 \chi_{\mathrm{UP}}(\pi)=\frac{f_1(f_1-3)}2+f_2,
 \qquad
 \chi_{\mathrm{EX}}(\pi)
     =\frac{(f_1-1)(f_1-2)}2-f_2.                          \tag{ATT.5}
\]
The first is at least $-1$: the only negative values of its
first summand occur at $f_1=1,2$, where it is $-1$.
For the second, if $f_1=0$ it is at least $1-\lfloor n/2\rfloor$;
if $f_1=1,2$ it is at least $-\lfloor(n-1)/2\rfloor$;
and if $f_1\ge3$, its first summand is nonnegative and
$f_2\le\lfloor(n-3)/2\rfloor$. This proves the stated floor
in every parity. The identities follow respectively by counting
fixed unordered pairs and from
$\chi_{\Lambda^2U}(g)=(\chi_U(g)^2-\chi_U(g^2))/2$,
where $\chi_U(g)=f_1-1$ and $\chi_U(g^2)=f_1+2f_2-1$.
\end{proof}

Put $V=W^{\otimes t}$, $e=a^t$, and let $\mathcal R$ be the
actual rational orthogonal tensor-wreath image of $S_n\wr S_t$.
Set
\[
 c=b/a,
 \qquad
 \mu=\begin{cases}2/n,&W=W_{\mathrm{UP}},\\
                  2/(n-1),&W=W_{\mathrm{EX}}.
       \end{cases}                                       \tag{ATT.6}
\]

\begin{lemma}[D-ATT-02: scalar extraction and marked inertias]
\label{D-ATT-02}
Let a finite linear group $J$ normalize the ordinary core $A_n^t$
on $V$. Then every $g\in J$ is uniquely $\lambda(g)T(g)$,
with $T(g)\in\mathcal R$ and a root of unity $\lambda(g)$.
The scalar coordinate is an actual character with image $\mu_M$,
and the actual character field $F$ of $J$ is contained in
$\mathbf Q(\zeta_M)$. Every genuine inertia has moving rank
at least $\mu e$.

For every $T\in\mathcal R$ and integer $r$,
\[
                      \operatorname{tr}(T^r)\ge-ce.        \tag{ATT.7}
\]
If $\lambda(g)\notin\mu_6$, its actual moving rank is at least
$(3-c)e/4$. In a generating product-one tuple, at least two
entries have this property whenever its original imaginary
quadratic field $K\subseteq F$ differs from $\mathbf Q(\sqrt{-3})$.
If $8\mid|\Delta_K|$, two scalar orders are divisible by eight
and those entries move at least $3e/4$. If an odd prime
$p\mid|\Delta_K|$, two scalar orders are divisible by $p$
and those entries move at least $(1-1/(p-1))e$.
\end{lemma}

\begin{proof}
The intrinsic simple factors and their matrix algebras are
permuted by every normalizer. Remove the corresponding tensor
permutation. Since $\operatorname{Aut}(A_n)=S_n$ for this range,
the remaining factor automorphisms are implemented by the
chosen actual $S_n$ operators on $W$. Removing them leaves a
scalar by Schur's lemma. This is the same full-automorphism
input and argument as D-AST-02, now on the two rational
irreducible modules in (ATT.1). Sign choices for extensions
to $S_n$ are absorbed in the scalar; no outer coset is omitted.

A scalar in $\mathcal R$ centralizes every represented simple
factor, so its factor permutation is trivial. Each of its
remaining $S_n$ permutations then centralizes $A_n$, since
the nontrivial $A_n$ representation is faithful. That
centralizer is trivial: a permutation commuting with every
three-cycle preserves each three-element support, and their
intersections fix every letter. Thus
$\mathcal R\cap\mathbf C^*I=\{I\}$. The decomposition is
unique; multiplication shows both coordinates are
homomorphisms, and finite order makes the scalar image
$\mu_M$. Rationality of $T$ gives
$\operatorname{tr}(g)\in\mathbf Q(\zeta_M)$ and hence the
asserted field containment. Equality $F=K$ is not used.

For a factor-preserving nonscalar tensor operator, the
largest eigenspace proportion is at most the largest such
proportion on any active factor, by convolution of its finite
spectral distributions. Lemma D-ATT-01 gives support at least
$\mu e$. For a nontrivial factor cycle of length $l\ge2$,
conjugate factorwise so that only the cycle product remains
as a twist, and diagonalize that finite-order product.
In the resulting tensor-word basis, each cyclic word orbit
contributes at most one vector to any chosen eigenspace.
There are $a$ constant words; all other word orbits have
length at least two. Thus the maximum is at most
$(a^l+a)/2\le3a^l/4$. The remaining tensor factors cannot
increase this proportion. Quarter support is stronger than
$\mu$ here. A nonidentity scalar has full moving rank,
which proves the assertion for every genuine inertia.

Trace contraction over the cycles of a tensor permutation
expresses $\operatorname{tr}(T)$ as a product of character
values of $W$ on their cycle products. Each value lies in
$[-b,a]$. If a product of $s\le t$ such values is negative,
one factor has absolute value at most $b$, and all others
at most $a$. Consequently it is at least
$-ba^{s-1}\ge-ba^{t-1}=-ce$. Every power belongs to the
same finite tensor-wreath group, proving (ATT.7).

For a root $\xi$ of order $d$, rationality of the actual
$T$ makes all primitive $d$th roots have the same multiplicity
$m_\xi$, so $\varphi(d)m_\xi\le e$. If $\xi\notin\mu_6$,
this gives $m_\xi\le e/4$ except possibly when $d=4$.
In that case the distinct $i,-i$ eigenspaces have equal
multiplicity $m_\xi$, and
\[
                 -ce\le\operatorname{tr}(T^2)
                                    \le e-4m_\xi.           \tag{ATT.8}
\]
Thus always $m_\xi\le(1+c)e/4$ for such a root. Apply this
to the actual $\lambda(g)^{-1}$ eigenspace of $T(g)$.
If the scalar order is divisible by eight or an odd prime
$p$, its totient is at least four or $p-1$, respectively.
These prove all marked-support bounds, including the case
that the requisite eigenvalue is absent.

For completeness, the arithmetic marks use the exact
argument of HD-AST-04, with all its representation premises
just verified here. If $K\ne\mathbf Q(\sqrt{-3})$, the
scalar image cannot be contained in $\mu_6$, since otherwise
$K\subseteq F\subseteq\mathbf Q(\zeta_6)$ would force equality.
The images of a generating product-one tuple in
$\mu_M/(\mu_M\cap\mu_6)$ cannot have exactly one nonidentity
entry. There are therefore at least two marked entries.

The standard quadratic conductor criterion gives
$|\Delta_K|\mid M$ from $K\subseteq\mathbf Q(\zeta_M)$.
We use Milne, Chapter V, Example3.2 p.~157, Remark3.8 p.~158,
and Examples3.10--3.11 p.~159, at exactly the primary-source
scope recorded for HD-AST-04.\cite{HD:MilneCFT2020}
The complete invoked passages were freshly read here.
For a prime $p\mid M$, the quotient
$\mu_M/\mu_{M/p}$ is nontrivial, and at least two tuple
entries have nonidentity images. If $p^v\Vert M$, each
such scalar has order divisible by $p^v$. This proves the
odd-prime and conductor-eight marks. If $M\equiv2\pmod4$,
the usual replacement of its cyclotomic conductor by $M/2$
does not affect divisibility by the original $M$.
\end{proof}

\begin{theorem}[D-ATT-03: exclusions for two-box tensor cores]
\label{D-ATT-03}
In the actual sixfold source setting, suppose the faithful
inducing image has ordinary core $A_n^t$ on either module
$W^{\otimes t}$ of (ATT.1), with $n\ge15$ and $t\ge2$.
Keep the original imaginary quadratic field only as
$K\subseteq F$. There is no period-image rank-nine source
in any of the following ranges:
\[
\begin{array}{c|c|c}
\text{field condition}&W_{\mathrm{UP}}&W_{\mathrm{EX}}\\\hline
K\ne\mathbf Q(\sqrt{-3})&15\le n\le39&15\le n\le39\\
8\mid|\Delta_K|&15\le n\le39&15\le n\le40\\
p\mid|\Delta_K|,\ p\ge7\text{ prime}
 &15\le n\le10(p-1)-1&15\le n\le10(p-1).
\end{array}                                                    \tag{ATT.9}
\]
In particular the remaining necessary lower bounds in the
eight non-Eisenstein class-number-one fields are
\[
\begin{array}{c|rrrrrrrr}
K&\mathbf Q(i)&\mathbf Q(\sqrt{-2})&\mathbf Q(\sqrt{-7})&
\mathbf Q(\sqrt{-11})&\mathbf Q(\sqrt{-19})&
\mathbf Q(\sqrt{-43})&\mathbf Q(\sqrt{-67})&\mathbf Q(\sqrt{-163})\\\hline
W_{\mathrm{UP}}:\ n\ge&40&40&60&100&180&420&660&1620\\
W_{\mathrm{EX}}:\ n\ge&40&41&61&101&181&421&661&1621.
\end{array}                                                    \tag{ATT.10}
\]
\end{theorem}

\begin{proof}
Use the actual Shapiro/Hurwitz period factorization established
in D-BI-01, not a subtraction of dimensions on the original
base. If $h$ is the intermediate genus and $D$ the sum of
actual moving ranks, multiplicity six gives
$6=2e(h-1)+D$. Here $e\ge a^2\ge8100$.
Genus at least two is impossible; genus one gives $D=6$,
at most six genuine branch marks and period-image rank at
most six. Rank nine therefore forces genus zero, at least
twelve genuine marks, and
\[
                              D=2e+6.                       \tag{ATT.11}
\]
The generating product-one inertia tuple is the actual one
of this intermediate faithful image.

For a non-Eisenstein original field, two marked entries from
D-ATT-02 and ten other genuine entries give
\[
  2e+6\ge\left(\frac{3-c}{2}+10\mu\right)e,
  \qquad \left(10\mu-\frac{1+c}{2}\right)e\le6.             \tag{ATT.12}
\]
For $W_{\mathrm{UP}}$, the coefficient times $a^2$ is
\[
 \frac{n(n-3)}8\big((40-n)(n-3)-2\big).                    \tag{ATT.13}
\]
On $15\le n\le39$, it is at least
$15\cdot12\cdot10/8>6$. Since the coefficient is positive
and $e\ge a^2$, this contradicts (ATT.12).
For $W_{\mathrm{EX}}$, use $c\le1/(n-2)$. The coefficient
times $a^2$ is at least
\[
 \frac{(n-1)(n-2)}8\big((40-n)(n-2)-1\big)
       \ge\frac{14\cdot13\cdot12}{8}>6                   \tag{ATT.14}
\]
on the same interval. This proves the first row.

For conductor divisible by eight, replace the two marked
supports by $3e/4$. The coefficient becomes $10\mu-1/2$.
The UP range already follows from the first row. In the EX
range $15\le n\le40$, its product with $a^2$ is
\[
              \frac{(n-1)(n-2)^2(41-n)}8>6,                \tag{ATT.15}
\]
so this also excludes the additional endpoint.

Finally put $h_0=p-1\ge6$. The two scalar orders divisible
by $p$ give the necessary inequality
$(10\mu-2/h_0)e\le6$. For UP its coefficient is positive
when $n<10h_0$, and its product with $a^2$ is
\[
             \frac{n(n-3)^2(10h_0-n)}{2h_0}.              \tag{ATT.16}
\]
On $15\le n\le10h_0-1$, concavity of $n(10h_0-n)$
puts its minimum at an endpoint, and that minimum is
$10h_0-1$: the other endpoint is $15(10h_0-15)$ and
their difference is $140h_0-224>0$.
Since $(n-3)^2\ge144$, expression (ATT.16) is at least
$144(10h_0-1)/(2h_0)>6$.
For EX put $x=n-1$, so $14\le x\le10h_0-1$.
The corresponding expression is
\[
             \frac{x(x-1)^2(10h_0-x)}{2h_0}.              \tag{ATT.17}
\]
Again the endpoint minimum of $x(10h_0-x)$ is $10h_0-1$,
as $14(10h_0-14)-(10h_0-1)=130h_0-195>0$.
Now $(x-1)^2\ge169$ proves the contradiction. This proves
all claimed ranges; (ATT.10) substitutes the exact quadratic
discriminants $4,8,7,11,19,43,67,163$.
\end{proof}

\paragraph{Strictness and remaining gate.}
For either label at $n=39$ and $t=2$, the previous ALT07
necessary support inequality is satisfied: its fraction is
$\beta(39,2)=36/703<1/10$. Thus the new exclusion over
$\mathbf Q(i)$ strictly narrows the previous main gate.
No claim is made that such a tuple or source existed before
the new obstruction. The conductor rows are compatible with
the separate broader long-label conductor refinement, but
use the sharper exact two-box support minima. All larger
unexcluded parameters and other long labels remain open in
this result; an all-degree tensor exclusion has not been
deduced from these inequalities.

The finite supplement checks exact cycle spectra, both
character formulas and support minima, tensor-cycle traces,
and the rational arithmetic. The all-degree proofs above,
complete compatible-normalizer extraction, original-field
containment and actual period-factorization interfaces are
not inferred from that finite calculation. The conditional
status of Markman, Severn and032 is unchanged, and no cycle
supplier is asserted by this obstruction.

\section{An all-degree obstruction for two alternating two-box factors}
\label{sec:alternating-two-factor}

\paragraph{Status and exact scope.}
This is a separate new author derivation, not a change to the frozen
D-ATT-01--03 section. The root and independent /root/imprimitive whole-proof, source and
code readings have passed, with distinct exact supplements.
Keep either actual rational module $W$ in (ATT.1), with $n\ge15$,
and put $a=\dim W$, $V=W\otimes W$, $e=a^2$.
The source theorem below assumes that the inducing image acts
transitively on the two simple ordinary factors $A_n\times A_n$.
It does not assert this transitivity for an arbitrary tensor product.

\begin{lemma}[D-ATT-04: a common index bound]
\label{D-ATT-04}
Put $A=(n-5)/2$ and $B=n/4$. For either module and every
nonidentity permutation $\pi\in S_n$,
\[
                 \nu(W(\pi))\ge A\operatorname{ind}(\pi)+B.
                                                               \tag{ATF.1}
\]
\end{lemma}

\begin{proof}
For UP, the target-one inequality (UP.3) is exactly (ATF.1).
For a nonidentity target scalar, (UP.4) has constant term
$(n-4)/2\ge n/4$, so it is at least as strong.
For EX, the target-one inequality (EX.4) is also stronger,
by $\operatorname{ind}(\pi)/2+n/12$.
Write $k=\operatorname{ind}(\pi)$ for the remaining EX targets.
If $k\ge n/2-2$, (EX.2) exceeds $Ak+B$ by
$k/2+1-n/4\ge0$. If $k<n/2-2$, (EX.3) gives the lower bound
\[
 a-\frac{(n-1)^2}{4}=rac{n^2-4n+3}{4}
 >\frac{n^2-8n+20}{4}
   =A(n/2-2)+B>Ak+B.                                      \tag{ATF.2}
\]
The strict middle inequality holds since $4n-17>0$.
Minimizing over targets proves the assertion. The identity
permutation is excluded from (ATF.1), as it must be.
\end{proof}

\begin{lemma}[D-ATT-05: actual weighted swaps]
\label{D-ATT-05}
Let $S(u\otimes v)=v\otimes u$. For $x,y\in S_n$, the
operator $T=(W(x)\otimes W(y))S$ is conjugate to
$(W(h)\otimes I)S$, where $h=xy$. For every nonzero scalar $\lambda$,
\[
 \operatorname{rank}(1-\lambda T)
       \ge\frac{a^2+a\nu(W(h))-a}{2}.                     \tag{ATF.3}
\]
In particular every scalar times a swap has support at least
$(a-1)e/(2a)$. If $h=1$ and $\lambda\notin\mu_6$, its
actual moving rank is $e$.
\end{lemma}

\begin{proof}
Conjugate by $I\otimes W(y)$. Let $\alpha_1,\ldots,\alpha_a$
be the eigenvalues of $W(h)$, with repetition. In an eigenbasis,
the tensor word $(i,i)$ has eigenvalue $\alpha_i$.
The two words $(i,j),(j,i)$, with $i\ne j$, form a weighted
two-cycle with the two distinct roots of
$z^2=\alpha_i\alpha_j$ as eigenvalues.

Fix a proposed eigenvalue $\zeta$. Let $N_\zeta$ count ordered
pairs with $\alpha_i\alpha_j=\zeta^2$, and let $d_+,d_-$
count $\alpha_i=\zeta,-\zeta$, respectively. Its exact
multiplicity is
\[
             m_\zeta=\frac{N_\zeta+d_+-d_-}{2}
                      \le\frac{N_\zeta+a}{2}.              \tag{ATF.4}
\]
For each $i$ the possible second indices lie in one eigenspace
of $W(h)$, so
$N_\zeta\le a(a-\nu(W(h)))$.
This proves (ATF.3), including off-spectrum targets, by taking
$\zeta=\lambda^{-1}$. Omitting the nonnegative support term
gives the asserted uniform swap bound.
If $h=1$, the conjugated operator is $S$, with eigenvalues
only $1,-1$. A scalar outside $\mu_6$ is neither of these,
so its actual fixed space is zero.
\end{proof}

\begin{theorem}[D-ATT-06: no transitive two-factor source in any degree]
\label{D-ATT-06}
Suppose the faithful finite inducing image $J$ contains the
ordinary irreducibly acting core $A_n\times A_n$ on $W\otimes W$,
where $n\ge15$ and $W$ is either module in (ATT.1), and
conjugation by $J$ is transitive on these two factors.
If the original imaginary quadratic field is
$K\ne\mathbf Q(\sqrt{-3})$, contained in the actual inducing
character field, its six-dimensional multiplicity variation
cannot have period-image rank nine. All actual scalar phases
and both compatible outer cosets on each factor are included.
\end{theorem}

\begin{proof}
The actual intermediate-cover argument in D-ATT-03 again gives
genus zero, a generating product-one tuple of $r\ge12$ genuine
inertias in $J$, and total moving rank $D=2e+6$.
Use D-ATT-02 to write each inertia uniquely as $\lambda T$.
Its permutation of the two factors is either the identity
or a swap. Write $P$ for the number of swaps, and put
\[
 d_0=\begin{cases}n-3,&\mathrm{UP},\\n-2,&\mathrm{EX},\end{cases}
 \qquad \mu=d_0/a,\quad c=b/a,\quad s=(a-1)/(2a).          \tag{ATF.5}
\]
Here $b$ is exactly as in (ATT.1).

Transitivity forces $P>0$, and the product-one relation
forces $P$ even. If $P\ge4$, four swap bounds and eight
other genuine marks give
\[
    D\ge(4s+8\mu)e
      =2e+a(8d_0-2)>2e+6.                                \tag{ATF.6}
\]
Thus $P=2$.

At least two scalar phases lie outside $\mu_6$, by D-ATT-02.
If any such mark were not one of these two swaps, choose
two marked entries and a swap distinct from them. The two
marked bounds, one swap bound, and nine other genuine marks
would imply
\[
 D\ge\left(2\frac{3-c}{4}+s+9\mu\right)e
      =2e+a\left(9d_0-\frac{b+1}{2}\right)>2e+6.           \tag{ATF.7}
\]
The last inequality follows from $a\ge90$, $d_0\ge12$ and
$b\le d_0$, and is valid for both parities. Therefore the
two swaps are precisely the scalar marks outside $\mu_6$.
In particular both satisfy that scalar condition.

For a non-swap inertia write
$g=\lambda(W(x)\otimes W(y))$ and
$I_g=\operatorname{ind}(x)+\operatorname{ind}(y)$.
If at least one of $x,y$ is nonidentity, convolution gives
\[
 \operatorname{rank}(1-g)
  \ge a\max\{\nu(W(x)),\nu(W(y))\}
  \ge\frac a2(AI_g+B),                                  \tag{ATF.8}
\]
since at least one factor supplies the positive constant $B$
in (ATF.1). If both are the identity, genuine inertia means
$g$ is a nonidentity scalar, so its moving rank is $e$,
which is still at least the right side of (ATF.8).

For either swap let $h=xy$ be its cycle product. If $h\ne1$,
(ATF.1) and (ATF.3) imply
\[
 \operatorname{rank}(1-g)
   \ge\frac e2+\frac a2\big(A\operatorname{ind}(h)+B\big)
                         -\frac a2.                       \tag{ATF.9}
\]
If $h=1$, the marked scalar and D-ATT-05 give moving rank
$e$. This also implies (ATF.9), because $B-1\le a$.
Thus each swap really supplies its $B$ term, even with
trivial cycle product; this uses its actual marked phase.

We now use a quotient of the actual cover. The tensor-wreath
action is faithful: its kernel would centralize both simple
factors, and the scalar-intersection argument of D-ATT-02
then makes every constituent permutation trivial. Hence
discarding the unique scalar gives an actual homomorphism
$J\longrightarrow S_n\wr S_2$. Its action on two blocks
of $n$ letters is transitive because $J$ contains $A_n^2$
and has a swap. Take the quotient of the given connected
$J$-cover by the stabilizer of one of these $2n$ letters.
This is a connected degree-$2n$ cover of the sphere.

The index of a non-swap on these letters is exactly $I_g$.
A swap with cycle product $h$ has one cycle of length $2l$
for each length-$l$ cycle of $h$, so its index is
$n+\operatorname{ind}(h)$. Riemann--Hurwitz on this actual
quotient therefore gives
\[
 E:=\sum_{\mathrm{non\!\!-\!swap}} I_g
          +\operatorname{ind}(h_1)+\operatorname{ind}(h_2)
                          \ge2n-2.                        \tag{ATF.10}
\]
No abstract generating tuple is substituted for the actual
inertia, and scalar entries simply have zero permutation index.

Summing (ATF.8)--(ATF.9) over all $r$ entries now yields
\[
 \begin{split}
 D&\ge e+\frac a2(AE+rB)-a\\
  &\ge e+\frac a2\big((n-5)(n-1)+3n\big)-a.
 \end{split}                                               \tag{ATF.11}
\]
The inner polynomial is $n^2-3n+5$. For UP it equals
$2a+5$, while for EX it equals $2a+3$. Consequently
\[
 D\ge\begin{cases}2e+3a/2,&\mathrm{UP},\\
                   2e+a/2,&\mathrm{EX},\end{cases}
                            \qquad D>2e+6.                 \tag{ATF.12}
\]
This contradicts the actual multiplicity-six Euler identity.
\end{proof}

\paragraph{Remaining scope and strict movement.}
The all-degree conclusion requires transitivity on the two
simple factors. It is the actual transitive two-factor case
of the main tensor-core gate, not a general assertion about
all factor-preserving tensor products. It removes, for example,
both labels at $n=100$, $t=2$, over $\mathbf Q(i)$, which
the prior finite ranges and the ALT07 necessary inequality
leave open. It asserts no prior existence of those sources.
Three or more factors are not treated by this addendum:
full tensor cycles of larger length may have a marked scalar
eigenvalue even when their cycle product is the identity,
so the decisive identity-product step cannot be silently
carried over. The original/inducing field distinction and
all earlier conditional inputs are unchanged.

\section{All tensor counts for the two ordinary alternating labels}
\label{sec:alternating-two-tail-all-tensors}

\paragraph{Status and exact scope.}
This is a separate new derivation, with all new proofs written below.
It leaves the reviewed D-ATT-01--06 unchanged.  The complete root and /root/imprimitive whole-proof/code
readings have passed, with a distinct matrix supplement and
explicit disclosure of the shared preceding framework.  The original imaginary quadratic field is non-Eisenstein;
it is only assumed to be contained in the character field of the
actual inducing representation.  Equality of those fields is not
assumed.  The result is an obstruction to an actual period-rank-nine
curve source, not a dimension-based construction of a source.

Fix $n\ge15$.  Let $W$ be one of the two ordinary rational
$A_n$-modules obtained by restricting
\[
 S^{(n-2,2)},\qquad S^{(n-2,1,1)}=\Lambda^2\operatorname{Std}_n.
 \tag{ATX.1}
\]
Put $a=\dim W$, and use the constants
\[
\begin{array}{c|c|c|c}
 W&a&d_0&b_0\\ \hline
 S^{(n-2,2)}&n(n-3)/2&n-3&1\\
 \Lambda^2\operatorname{Std}_n&(n-1)(n-2)/2&n-2&\lfloor(n-1)/2\rfloor
\end{array}
\qquad A=\frac{n-5}{2},\quad B=\frac n4.
\tag{ATX.2}
\]
For an operator of degree $v$ write
$\nu(T)=v-\max_z\dim\ker(T-zI)$.  For $x\in S_n$ let
$\operatorname{ind}(x)=n-\#\{\text{cycles of }x\}$, including fixed
points among its cycles.  The complete previously reviewed UP/EX
proofs, collected in D-ATT-01 and D-ATT-04, give
\[
 \nu(Wx)\ge d_0,\qquad
 \nu(Wx)\ge A\operatorname{ind}(x)+B\quad(x\ne1),
 \qquad \operatorname{tr}(Wx)\ge-b_0.
 \tag{ATX.3}
\]
The two support inequalities apply to every target eigenvalue
when $x\ne1$; neither is asserted for the identity.  The trace
bound holds for all $x$.  This distinction is retained throughout.

Let a faithful finite inducing image $J$ act on
$V=W^{\otimes t}$, of degree $e=a^t$, and normalize the ordinary
core $A_n^t$, acting transitively on its $t$ factors.  Here $t\ge3$.
Its compatible normalizer is scalar times the rational tensor-wreath
image of $S_n\wr S_t$.  Indeed $\operatorname{Aut}(A_n)=S_n$ in
this range, and after implementing its factor automorphisms and
factor permutation the remaining operator centralizes the irreducible
core and is scalar.  That rational wreath image meets the scalars
trivially: a scalar centralizes every individual $A_n$ factor, so the
factor permutation is trivial and each $S_n$ automorphism centralizes
$A_n$, hence is the identity.  Thus
\[
                  g=\lambda(g)T(g)
 \tag{ATX.4}
\]
is unique, and both factors in (ATX.4) are actual homomorphisms.
The rational image is faithful as a tensor-wreath action by the
same centralizer argument.  In particular discarding the scalar
gives the actual permutation action on $nt$ letters in $t$ blocks.

\subsection*{D-ATT-07. Two moving tensor-factor branches}

\begin{lemma}[D-ATT-07]
Suppose the actual inducing variation has period rank nine and its
character field contains an imaginary quadratic field
$K\ne\mathbf Q(\sqrt{-3})$.  Then exactly two genuine inertia
entries move the tensor factors.  Their permutations are inverse
full $t$-cycles.
\end{lemma}

\begin{proof}
The reviewed polarized Shapiro/Hurwitz factorization applies to
the actual intermediate connected finite cover, without changing
the inducing field.  Its six-dimensional multiplicity gives
$6=2e(h-1)+D$, where $h$ is the intermediate base genus and
$D$ is the sum of actual moving ranks.  If $h\ge1$, the finite
marked-cover dimension is less than nine: $h=1$ gives at most six
genuine marks, and $h\ge2$ is impossible since $e>3$.
Consequently
\[
       h=0,\qquad r\ge12,\qquad D=2e+6,
 \tag{ATX.5}
\]
for a genuine generating product-one inertia tuple.  These are
necessary conditions from the actual variation, not sufficient
conditions for dominance.

Let $\lambda(J)=\mu_M$.  Rationality of $T(g)$ gives
$\mathbf Q(\chi_J)\subseteq\mathbf Q(\zeta_M)$.  If
$\lambda(J)\subseteq\mu_6$, this field could not contain $K$.
The quotient by $\lambda(J)\cap\mu_6$ is therefore nontrivial.
Generation and product one force at least two inertia entries with
$\lambda\notin\mu_6$.

We recall the complete scalar estimate of D-ATT-02.  The trace of
each power of a rational tensor-wreath operator is a product of
cycle-product traces.  If negative, one factor is at least $-b_0$
and all other absolute values are at most $a$.  Thus, with
$c=b_0/a$, every power has trace at least $-ce$.
For a primitive root of order outside $\{1,2,3,4,6\}$, rationality
makes its multiplicity at most $e/4$ by its cyclotomic conjugates.
For a fourth root, equal multiplicities at $i,-i$ and
$\operatorname{tr}(T^2)\ge-ce$ give multiplicity at most
$(1+c)e/4$.  Therefore each of the two indicated actual scalar
marks has moving rank at least
\[
                         \frac{3-c}{4}e.
 \tag{ATX.6}
\]

Every genuine inertia has moving rank at least $\mu e$, where
$\mu=d_0/a$.  If it fixes the factors and one of its within-factor
permutations is nonidentity, this follows from (ATX.3) and convolution of
eigenspace multiplicities.  A genuine scalar has full moving rank.
If it moves factors, restrict to one nontrivial tensor cycle of
length $l$.  After diagonalizing its cycle product, the constant
words number $a$, and all other word orbits have length at least
two.  Each word orbit contributes at most one vector to a specified
eigenspace.  Its largest eigenspace fraction is at most
$(1+a^{1-l})/2\le(1+1/a)/2$.  Convolution with the remaining
cycles preserves that bound.  Thus every factor-moving inertia
has moving rank at least
\[
               s e,\qquad s=\frac{a-1}{2a}\ge\mu.
 \tag{ATX.7}
\]
The last inequality follows directly from $2d_0\le a-1$ in
(ATX.2), for $n\ge15$.

If at least three inertia entries move factors, select the two
marks in (ATX.6) and a third, distinct factor-moving entry.
Together with at least nine other genuine entries, they give
\[
 D\ge\left(2\frac{3-c}{4}+s+9\mu\right)e
   =2e+a^{t-1}\left(9d_0-\frac{b_0+1}{2}\right)>2e+6.
 \tag{ATX.8}
\]
Here $b_0\le d_0$, $d_0\ge12$ and $a\ge90$.  Thus at most two
entries move factors.  Transitivity on $t>1$ factors requires at
least one, and product one forbids exactly one.  There are exactly
two, their permutations are inverse, and their cyclic subgroup
must be transitive.  Hence both are full $t$-cycles.  This argument
does not assign a deformation or a separate geometric family to
any scalar mark.
\end{proof}

\subsection*{D-ATT-08. A weighted tensor-cycle estimate}

\begin{lemma}[D-ATT-08]
Let $X$ be any invertible diagonalizable degree-$a$ operator, and let $T$ be
a weighted tensor $t$-cycle whose cycle product is $X$.
Let $C_d(a)$ denote the number of primitive cyclic word orbits
of length $d$ on an alphabet of size $a$, and put
\[
 R_t(a)=\sum_{\substack{d\mid t\\d<t}}(t-d)C_d(a).
 \tag{ATX.9}
\]
For every nonzero scalar $\lambda$,
\[
 \operatorname{rank}(1-\lambda T)
 \ge\left(1-\frac1t\right)a^t
       +\frac{a^{t-1}}t\nu(X)-\frac{R_t(a)}t.
 \tag{ATX.10}
\]
For $a\ge4$ and $t\ge3$,
\[
                             R_t(a)\le2a^{t-2}.
 \tag{ATX.11}
\]
\end{lemma}

\begin{proof}
Successive changes of basis in the tensor factors conjugate a
weighted cycle to one with $X$ on one factor and the identity on
the others.  For example, if the cycle sends
$(v_1,\ldots,v_t)$ to $(X_1v_t,X_2v_1,\ldots,X_tv_{t-1})$,
take change-of-basis factors $I,X_2,X_3X_2,\ldots,X_t\cdots X_2$.
The cycle product is $X_1X_t\cdots X_2$, up to the harmless
choice of starting factor.  Its $t$th power applies $X$ to every
factor.

Choose an eigenbasis of $X$, with eigenvalues
$\alpha_1,\ldots,\alpha_a$, repeated with their multiplicities.
For a proposed eigenvalue $z$ of $T$, let $N_z$ count all ordered
words $(i_1,\ldots,i_t)$ satisfying
$\prod_j\alpha_{i_j}=z^t$.  Fixing the first $t-1$ indices shows
\[
                          N_z\le a^{t-1}(a-\nu(X)).
 \tag{ATX.12}
\]
A cyclic word orbit of length $d\mid t$ gives a weighted cyclic
matrix of size $d$.  Its nonzero cyclic weights make each of its
$d$ eigenvalues simple.  If it contributes to the $z$-eigenspace,
all its $d$ words are counted by $N_z$.  Its contribution to
$t\dim\ker(T-zI)-N_z$ is therefore at most $t-d$.
An orbit not contributing gives a nonpositive contribution.
Full-length orbits contribute at most zero.  Summing all shorter
orbits proves
$t\dim\ker(T-zI)\le N_z+R_t(a)$, hence (ATX.10).
Repeated eigenvalues do not merge the basis-word orbits and
cause no change in this counting argument.

For clarity, the same error has the exact fixed-rotation expression
\[
                    R_t(a)=\sum_{j=1}^{t-1}a^{\gcd(t,j)}.
 \tag{ATX.19}
\]
An orbit of length $d<t$ has $d$ words and is fixed by exactly
$t/d-1$ nonidentity rotations, so it contributes $t-d$ to both
counts.  Thus the definition includes the full short-orbit defect,
not merely the number of short words.

For $t=3$, $R_3(a)=2a$.  For $t=4$, the shorter orbits have
length one or two, so $R_4(a)=a^2+2a\le2a^2$ for $a\ge2$.
For $t\ge5$, every proper divisor is at most $\lfloor t/2\rfloor$
and $C_d(a)\le a^d/d$.  Thus
\[
 R_t(a)\le(t-1)\sum_{j=1}^{\lfloor t/2\rfloor}a^j
       <\frac{(t-1)a^{\lfloor t/2\rfloor+1}}{a-1}.
 \tag{ATX.13}
\]
The inequality
$t-1\le2(a-1)a^{\lceil t/2\rceil-3}$ holds at $t=5,6$
for $a\ge4$.  Increasing $t$ by two multiplies the right side
by $a$, while $t+1\le a(t-1)$.  Induction in each parity
therefore proves it for all $t\ge5$.  Substituting it in
(ATX.13) proves (ATX.11).  This is a universal bound; no finite
word enumeration is used to extend its range.
\end{proof}

\subsection*{D-ATT-09. Exclusion for every larger tensor count}

\begin{theorem}[D-ATT-09]
No source of period rank nine in the setting above exists for
any $t\ge3$.  Together with D-ATT-06 and the already reviewed
single-factor UP/EX exclusions, all transitive ordinary core
tensor counts for the two labels (ATX.1), $n\ge15$, are excluded
in the non-Eisenstein original-field problem.
\end{theorem}

\begin{proof}
Use the actual inertia tuple of (ATX.5).  By D-ATT-07 there are
exactly two full tensor cycles.  All other entries fix the
factors.  For such an entry, write its rational part as
$Wx_1\otimes\cdots\otimes Wx_t$, and set
$I=\sum_j\operatorname{ind}(x_j)$.  If precisely $b>0$ of the
$x_j$ are nonidentity, (ATX.3) gives
\[
 \max_j\nu(Wx_j)\ge\frac{AI+bB}{b}
                 \ge\frac{AI}{t}+B.
 \tag{ATX.14}
\]
Convolution of eigenspaces therefore bounds its actual moving
rank, with any scalar, below by
\[
                         a^{t-1}\left(\frac{AI}{t}+B\right).
 \tag{ATX.15}
\]
If all $x_j=1$, the genuine entry is a nonidentity scalar, with
moving rank $e$, and (ATX.15) still holds since $I=0$ and $B<a$.
The full $B$ term in (ATX.15), rather than $B/t$, is essential.

Let $h_1,h_2\in S_n$ be the products around the two full cycles.
For both the identity and nonidentity products, (ATX.3) gives the
valid weaker bound $\nu(Wh_i)\ge A\operatorname{ind}(h_i)$.
D-ATT-08 therefore bounds the sum of the two cycle moving ranks
below by
\[
 2\left(1-\frac1t\right)e
  +\frac{Aa^{t-1}}t\bigl(\operatorname{ind}(h_1)
                         +\operatorname{ind}(h_2)\bigr)
  -\frac{2R_t(a)}t.
 \tag{ATX.16}
\]
No positive $B$ term is demanded at an identity cycle product.
No restriction on whether the actual scalar is a $t$th root
is needed in this estimate.

The scalar-forgetting homomorphism (ATX.4) gives an actual
transitive action on $nt$ letters.  The core is transitive in
each block and the full cycle is transitive on the blocks.
Take the inverse image of a letter stabilizer in the actual
connected $J$-cover of the sphere.  Its quotient is a connected
degree-$nt$ cover.  A factor-fixing entry has permutation index
$I$, whereas a full tensor cycle has permutation index
$n(t-1)+\operatorname{ind}(h_i)$: a cycle of length $l$ of $h_i$
becomes a cycle of length $tl$ on the letters.
Riemann--Hurwitz on this actual quotient consequently gives
\[
 E:=\sum_{\text{factor-fixing entries}} I
       +\operatorname{ind}(h_1)+\operatorname{ind}(h_2)
       \ge2n-2.
 \tag{ATX.17}
\]
Scalar inertia contributes zero to this permutation index but
is retained among the genuine coefficient-cover marks.

There are at least ten factor-fixing entries.  Summing
(ATX.15)--(ATX.16), then using (ATX.11) and (ATX.17), yields
\[
\begin{split}
 D&\ge2e+\frac{a^{t-1}}t
       \left[A(2n-2)+10tB-2a-\frac4a\right]\\
  &=2e+\frac{a^{t-1}}t
       \left[\left(\frac{5t}{2}-3\right)n
                    +\kappa-\frac4a\right],
 \qquad
 \kappa=\begin{cases}5,&W=S^{(n-2,2)},\\3,&W=\Lambda^2\operatorname{Std}_n.
              \end{cases}
\end{split}
\tag{ATX.18}
\]
For $t\ge3$, $5t/2-3\ge3t/2$, and
$\kappa-4/a>0$ since $a\ge90$.  The excess in (ATX.18) is
therefore greater than $3na^{t-1}/2>6$, contradicting (ATX.5).
This proves the theorem for all $n\ge15$ and all $t\ge3$.
\end{proof}

\paragraph{Checks, sources and remaining scope.}
The new word-orbit and all-size source arguments are universal
proofs.  A separate exact supplement checks actual weighted cyclic
matrices, short-orbit errors, permutation quotient indices and
endpoint arithmetic.  It is not a tuple or Hurwitz-component
enumeration.  The scalar-field and ordinary-label inputs are the
complete reviewed D-ATT-01--04 proofs; no additional external
classification or numerical premise is introduced.  Actual
intermediate-cover factorization and Riemann--Hurwitz are retained
explicitly.  Other ordinary alternating labels, spin modules,
nontransitive mixtures and arbitrary nonalternating cores are not
decided here.  No positive supplier or algebraic cycle is claimed.

\section{Standard alternating cores with two transitive factors}
\label{sec:standard-two-factor}

\paragraph{Status and authorship.}
This is a new root derivation, with a contributed Cauchy estimate from
/root/imprimitive. The complete root author reading, the contributor verification, and
the separate /root/g34\_eightfolds whole-proof/code audit have passed.
Distinct exact supplements and the formatting-only repair are retained. The result concerns exactly two ordinary standard
alternating factors, transitive under the actual inducing image.
It allows every scalar phase and does not restrict the original
imaginary quadratic field.

Write $U$ for the deleted permutation module of $S_n$,
$a=n-1$ and $e=a^2$, with $n\ge15$. For a finite-order
operator $T$, let $\nu(T)$ be its dimension minus its largest
eigenspace dimension. Thus every scalar multiple has moving rank
at least $\nu(T)$. The permutation index is
$\operatorname{ind}(x)=n-c_x$, where $c_x$ is the number of
cycles, including fixed points.

\begin{lemma}[HD-AST-07: a uniform standard tensor index estimate]
\label{HD-AST-07}
If $x,y\in S_n$ are not both the identity, and
$k=\operatorname{ind}(x)+\operatorname{ind}(y)$, then
\[
 \nu\bigl(U(x)\otimes U(y)\bigr)
       \ge \frac{(n-3)k+n}{2}.                         \tag{AST2.1}
\]
The assertion includes every eigenvalue, not only the eigenvalue one.
\end{lemma}

\begin{proof}
Put $c=c_x$, $d=c_y$, $s=c+d=2n-k$.
First there is a bound valid at every target eigenvalue:
\[
 \nu\bigl(U(x)\otimes U(y)\bigr)
                    \ge\frac{n-2}{2}k.                \tag{AST2.2}
\]
Indeed, if $m_\alpha$ are the eigenvalue multiplicities on the
full letter module, every cycle contributes at most one to each
$m_\alpha$, so $m_\alpha\le c$ and $\sum m_\alpha=n$.
Deleting one trivial line changes its squared multiplicity sum to
\[
 \sum_\alpha u_\alpha^2
   =\sum_\alpha m_\alpha^2-2c+1
   \le(n-2)c+1.                                        \tag{AST2.3}
\]
Use a common finite cyclic group for the eigenvalues of both
operators. At a target $\zeta$, the tensor multiplicity is
$\sum_\alpha u_\alpha v_{\zeta/\alpha}$.
Cauchy--Schwarz followed by the arithmetic--geometric mean inequality
bounds this by
\[
 \sqrt{\big((n-2)c+1\big)\big((n-2)d+1\big)}
            \le\frac{n-2}{2}(c+d)+1.
                                                               \tag{AST2.4}
\]
Subtracting from $e$ proves (AST2.2). If $k\ge n$, its right
side is at least the right side of (AST2.1), completing that case.

Suppose henceforth $1\le k<n$, and put $L=n-k>0$.
The required improvement over (AST2.2) is exactly $L/2$.
All matrices are rational. For a nonreal target eigenvalue, the
complex conjugate eigenspace has the same dimension and is distinct,
so its multiplicity is at most $e/2$. The corresponding moving
rank is at least $e/2$, which exceeds the right side of (AST2.1):
the latter is increasing in $k$ and at $k=n$ is $(e-1)/2$.
An off-spectrum target has full moving rank and causes no exception.
Only the targets $1$ and $-1$ remain.

Let $l_1,\ldots,l_c$ and $m_1,\ldots,m_d$ be the cycle lengths.
The full grid permutation on $n^2$ pairs has
$C=\sum_{i,j}\gcd(l_i,m_j)$ cycles. Define
\[
 G=\frac n2(c+d)-C
   =\frac12\sum_{i,j}
        \bigl(l_i+m_j-2\gcd(l_i,m_j)\bigr)\ge0.
                                                               \tag{AST2.5}
\]
For target one, the character identity
$U\otimes U=P\otimes P-P\otimes\mathbf1
-\mathbf1\otimes P+\mathbf1$ gives multiplicity $C-c-d+1$.
Its moving rank is therefore $(n-2)k/2+G$.
Write $k_x=n-c$, $k_y=n-d$ and let $f_x,f_y$ count fixed points.
The fixed-point/nontrivial-cycle terms of (AST2.5) give
\[
 \begin{split}
 G&\ge\frac12(f_yk_x+f_xk_y)\\
  &\ge\frac12\bigl(nk-4k_xk_y\bigr)
   \ge\frac12k(n-k)\ge\frac L2.
 \end{split}                                             \tag{AST2.6}
\]
Here $f_x\ge n-2k_x$ and $f_y\ge n-2k_y$, since a moved
cycle of length $l$ has $l\le2(l-1)$. The lower bounds remain
valid if their right sides are negative. The next inequality is
$4k_xk_y\le k^2$. This proves the target-one case.

For target minus one, let $u$ and $v$ count odd cycles in $x$ and
$y$. A grid cycle is even exactly when at least one constituent
cycle is even. Thus the multiplicity on the full grid is
\[
 C_- =\sum_{\operatorname{lcm}(l_i,m_j)\ \text{even}}
                             \gcd(l_i,m_j).
                                                               \tag{AST2.7}
\]
The two full letter modules have minus-one multiplicities
$c-u$ and $d-v$. Consequently the tensor multiplicity is
$C_--(c-u)-(d-v)$, and the excess of its moving rank over
$(n-2)k/2$ is
\[
 H-u-v+1,\qquad H=\frac n2(c+d)-C_-.
                                                               \tag{AST2.8}
\]
Each odd--odd cycle pair contributes $(l_i+m_j)/2\ge1$ to
$H$. An odd--even pair contributes
$(l_i+m_j-2\gcd(l_i,m_j))/2\ge1/2$, since its numerator
is a positive odd integer. An even--even pair contributes
a nonnegative number. Therefore
\[
 H\ge uv+\frac12\bigl(u(d-v)+v(c-u)\bigr)
      =\frac12(ud+vc),
\quad
 H-u-v+1\ge1+\frac{u(d-2)+v(c-2)}2.
                                                               \tag{AST2.9}
\]
If $c,d\ge3$, the lengths of odd cycles are at least one and
those of even cycles at least two, so
$u\ge2c-n$ and $v\ge2d-n$. Hence $u+v\ge2L$ and the
last expression in (AST2.9) is at least $1+L\ge L/2$.
If $\min(c,d)=2$, both counts are at least two and
$L=c+d-n\le2$; (AST2.9) is at least $1\ge L/2$.
If $\min(c,d)=1$, positivity of $L$ forces the other count
to be $n$. One permutation is an $n$-cycle and the other is
the identity, and $L=1$. The minus-one multiplicity on the
deleted $n$-cycle is either one or zero, so its tensor
multiplicity is at most $a$. Its moving rank is at least
$e-a=a(a-1)$, exceeding the baseline $a(a-1)/2$ by at least
$1/2$. These cases exhaust all positive cycle counts and prove
the target-minus-one assertion. This finishes (AST2.1).
\end{proof}

\begin{lemma}[HD-AST-08: standard weighted swaps and the actual quotient]
\label{HD-AST-08}
Let $S$ interchange the two tensor factors, and let
$T=(U(x)\otimes U(y))S$, with $h=xy$.
For every nonzero scalar $\lambda$,
\[
 \operatorname{rank}(1-\lambda T)
   \ge\frac{e-a}{2},
 \qquad
 \operatorname{rank}(1-\lambda T)
   \ge\frac e2+\frac a2\operatorname{ind}(h)-a.
                                                               \tag{AST2.10}
\]
If an actual finite inducing image with core $A_n^2$ acts
transitively on its two standard factors, discarding its unique
scalar coordinate gives an actual transitive action on $2n$ letters.
A swap with cycle product $h$ has index
$n+\operatorname{ind}(h)$ on those letters.
\end{lemma}

\begin{proof}
The unique scalar and full compatible-normalizer description is
HD-AST-04, restricted to two factors. Its proof uses the trivial
scalar intersection of the rational tensor-wreath image and does
not impose a quadratic-conductor condition. Thus the projection
is an actual homomorphism, not a choice made separately at each
inertia. The $A_n^2$ core is transitive on each block, and a swap
joins them. A cycle of length $l$ in $h$ becomes a cycle of
length $2l$ on the two blocks, proving the index assertion.

For clarity, the weighted-word proof of D-ATT-05 applies to any
actual $a$-dimensional module, hence to $U$. Conjugation by
$I\otimes U(y)$ gives $(U(h)\otimes I)S$.
For a proposed eigenvalue $\zeta$, let $N_\zeta$ count ordered
eigenvalue pairs $\alpha_i\alpha_j=\zeta^2$, and let
$d_+,d_-$ count $\alpha_i=\zeta,-\zeta$.
The exact multiplicity is
$(N_\zeta+d_+-d_-)/2\le(a(a-\nu(U(h)))+a)/2$.
This proves moving rank at least
$(e+a\nu(U(h))-a)/2$, and hence the first bound in (AST2.10).
Every full-letter eigenspace has dimension at most the number
$c_h$ of cycles. Deleting one trivial line cannot increase it,
so $\nu(U(h))\ge a-c_h=\operatorname{ind}(h)-1$.
Substitution proves the second bound, even when $h$ is the
identity. No scalar-order assumption is used.
\end{proof}

\begin{theorem}[HD-AST-09: no standard transitive two-factor source]
\label{HD-AST-09}
Let $n\ge15$. A faithful finite inducing image containing the
normal ordinary standard core $A_n^2$ on $U\otimes U$, and transitive
on these two factors, cannot supply a six-dimensional multiplicity
variation of period-image rank nine. This holds for every original
imaginary quadratic field and includes every compatible scalar,
outer automorphism and factor permutation.
\end{theorem}

\begin{proof}
The actual intermediate-cover and period factorization used in
D-AST-03 and HD-AST-04--06 gives the following necessary data:
a connected cover of a genus-zero curve, a generating product-one
tuple of $r\ge12$ genuine inertia operators, and total moving rank
\[
                         D=2e+6.                          \tag{AST2.11}
\]
Here $e=(n-1)^2\ge196$. Intermediate genus at least two is
excluded by the multiplicity Euler identity; genus one has only
six units of moving rank and cannot give period rank nine.
These are the previously proved geometric necessities, not an
assertion that an abstract numerical tuple creates a source.

Every genuine inertia has moving rank at least $a$. Indeed a
nonidentity constituent permutation acts nonscalarly on $U$ and
has projective support at least one; convolution then gives support
at least $a$. A nonidentity scalar has full support. A factor swap
has support at least $(e-a)/2\ge a$ by (AST2.10).

Let $P$ count the swaps. Transitivity forces $P>0$, and the
product-one condition makes it even. If $P\ge4$, four swaps
and eight further genuine entries give
\[
              D\ge4\frac{e-a}{2}+8a=2e+6a>2e+6.
                                                               \tag{AST2.12}
\]
Hence $P=2$. This step uses no scalar-conductor mark.

For a non-swap write $g=\lambda(U(x)\otimes U(y))$ and
$I_g=\operatorname{ind}(x)+\operatorname{ind}(y)$.
Set $A=(n-3)/2$ and $B=n/2$.
Lemma HD-AST-07 gives moving rank at least $AI_g+B$ unless
both permutations are the identity. In that case genuine inertia
is a nonidentity scalar and has full moving rank $e\ge B$, so
the same bound still holds.
For the two swaps, with cycle products $h_1,h_2$, (AST2.10)
gives the bound $e/2+a\operatorname{ind}(h_i)/2-a$.

Use the actual $2n$-letter quotient of HD-AST-08. Its point
stabilizer defines a connected degree-$2n$ quotient of the given
cover. Riemann--Hurwitz on that actual quotient implies
\[
 N+H\ge2n-2=2a,\qquad
 N=\sum_{\rm non\text{-}swap} I_g,\quad
 H=\operatorname{ind}(h_1)+\operatorname{ind}(h_2).
                                                               \tag{AST2.13}
\]
The two swaps contribute the additional base index $2n$;
the total index is at least $2(2n-1)$. Scalar branch entries
have zero permutation index and remain among the genuine
coefficient marks.

Since $a/2-A=1$, summing the actual moving-rank bounds yields
\[
 \begin{split}
 D&\ge e+AN+\frac a2H+(r-2)B-2a\\
  &\ge e+A(N+H)+10B-2a\\
  &\ge e+(n-3)(n-1)+5n-2(n-1)\\
  &=2e+a+5>2e+6.
 \end{split}                                             \tag{AST2.14}
\]
This contradicts (AST2.11).
\end{proof}

\paragraph{Strictness, remaining scope and programme.}
For $K=\mathbf Q(i)$, the standard label at $n=100$ with two
factors satisfies the previous conductor cutoff and the ALT07
shape inequality; the theorem removes it and all larger degrees
in this exact two-factor transitive case. It asserts no prior
existence of a source. Three or more standard factors, other
ordinary labels, and nontransitive tensor-core configurations
are not removed by this theorem. No effective detecting cycle
is manufactured by a representation obstruction.

The result restricts actual weight-one local systems and the
possible dominating suppliers for condition A (cote16, scan050;
SIGN_NEXT SN-C03/C04). The rational representations and their
exterior/tensor operations live in the representation lambda-ring
and the corresponding local-system K-theory; the quotient cover
is an actual finite morphism, compatible with the common
SGA/EGA/FGA and schemes-over-$\mathbf Z$ frame. It changes
neither the condition-I positivity step nor the projector,
motivic or $L$-function implications. No new Markman, Severn
or032 premise is used.

\section{Standard alternating cores with every larger tensor count}
\label{sec:standard-all-tensors}

\paragraph{Status, contribution and exact scope.}
This is a separate new derivation.  It leaves HD-AST-07--09 unchanged.
The pair-selection argument and the all-size source summation below are
new here.  They use the jointly derived HD-AST-07 standard two-factor
estimate and the weighted-cycle argument of D-ATT-08; the complete
weighted-cycle proof is included with that attribution.  Author checks
and independent whole-proof reviews are recorded separately.  The complete root reading and /root/g34\_eightfolds whole-proof
audit have passed. Shared weighted-cycle authorship is explicit;
the separate exact implementations are not labelled independent
discoveries of that common lemma.

Let $n\ge15$, put $a=n-1$, and write $U=\operatorname{Std}_n$
for the deleted rational permutation module.  Fix $t\ge3$ and
$V=U^{\otimes t}$, of degree $e=a^t$.  A faithful finite actual
inducing image $J$ is assumed to contain the ordinary normal core
$A_n^t$ on $V$ and to act transitively on its $t$ factors.  Its
character field $F$ contains the original imaginary quadratic field
$K$, which satisfies
\[
                    K\ne\mathbf Q(\sqrt{-3}).                 \tag{AST3.1}
\]
There is no assumption that $F=K$.  For a diagonalizable operator
$T$ of degree $v$ put
$\nu(T)=v-\max_z\dim\ker(T-zI)$.  Thus every scalar multiple
has actual moving rank at least $\nu(T)$.  For $x\in S_n$ put
$k_x=\operatorname{ind}(x)=n-c_x$, where $c_x$ counts all cycles,
including fixed points.

The fully reviewed normalizer argument D-AST-01--02 and
HD-AST-04 gives a unique decomposition
\[
             g=\lambda(g)T(g),\qquad
 T(g)\in\mathcal R_{n,t},\qquad \lambda(g)\in\mu_M,              \tag{AST3.2}
\]
where $\mathcal R_{n,t}$ is the actual rational tensor-wreath image
of $S_n\wr S_t$.  Both coordinates are homomorphisms.  We recall
why the uniqueness matters here.  The automorphisms of each $A_n$
factor are implemented by $S_n$ in this range; after implementing
them and the factor permutation, Schur's lemma leaves only a scalar.
A scalar in the rational wreath image centralizes every individual
represented $A_n$ factor, so it has trivial factor permutation.
Each remaining $S_n$ element then centralizes $A_n$ and is the
identity.  The same argument proves faithfulness of the rational
wreath image.  Thus forgetting $\lambda$ gives an actual
permutation representation on $nt$ letters in $t$ blocks; it is
not an inertia-by-inertia choice of representatives.

\subsection*{HD-AST-10. Pair selection at a factor-fixing branch}

\begin{lemma}[HD-AST-10]
Let a genuine inertia $g$ fix all tensor factors.  Write
$T(g)=U(x_1)\otimes\cdots\otimes U(x_t)$ and
$I_g=\sum_{j=1}^t k_{x_j}$.  Then, for its actual scalar phase,
\[
 \operatorname{rank}(1-g)\ge
 a^{t-2}\left(\frac{n-3}{t}I_g+\frac n2\right).                \tag{AST3.3}
\]
\end{lemma}

\begin{proof}
Suppose first that $I_g>0$.  The average of $k_{x_i}+k_{x_j}$
over the $\binom t2$ unordered pairs is $2I_g/t$.  Choose a pair
whose sum is at least that average; its sum is positive.
HD-AST-07, which applies to every target eigenvalue, gives
\[
 \nu\bigl(U(x_i)\otimes U(x_j)\bigr)
 \ge\frac{(n-3)(k_{x_i}+k_{x_j})+n}{2}
 \ge\frac{n-3}{t}I_g+\frac n2.                               \tag{AST3.4}
\]
If a diagonalizable degree-$v$ operator has largest eigenspace
dimension $m$, its tensor product with a degree-$w$ operator has
each eigenspace of dimension at most $mw$: sum the first
multiplicities against the second multiplicities, whose total is
$w$.  Applying this convolution bound to the other $t-2$ factors
multiplies the support in (AST3.4) by $a^{t-2}$.  The actual
scalar can only select one of these target eigenspaces, proving
(AST3.3).

If $I_g=0$, all $x_j$ are the identity.  Since the inertia is
genuine, $g$ is a nonidentity scalar and has moving rank $e$.
The right side of (AST3.3) is then $a^{t-2}n/2\le e$, since
$(n-1)^2\ge n/2$.  This proves the remaining case.  In particular
the full constant $n/2$ is retained even if only one factor acts
nontrivially; it is not divided by the number of tensor factors.
\end{proof}

\subsection*{HD-AST-11. The exact weighted-cycle error}

\begin{lemma}[HD-AST-11, the D-ATT-08 weighted-cycle argument]
Let $X$ be an invertible diagonalizable degree-$a$ operator, and
let $T$ be a weighted tensor $t$-cycle with cycle product $X$.
Define the exact integer
\[
                 R_t(a)=\sum_{j=1}^{t-1}a^{\gcd(t,j)}.          \tag{AST3.5}
\]
For every nonzero scalar $\lambda$,
\[
 \operatorname{rank}(1-\lambda T)\ge
 \left(1-\frac1t\right)a^t+\frac{a^{t-1}}t\nu(X)
                          -\frac{R_t(a)}t.                    \tag{AST3.6}
\]
For $a\ge4$ and $t\ge3$,
\[
                          R_t(a)\le2a^{t-2}.                  \tag{AST3.7}
\]
\end{lemma}

\begin{proof}
Successive changes of basis in the tensor factors conjugate $T$
to a cyclic rotation with $X$ on one factor and identities on the
others.  Explicitly, for
$(v_1,\ldots,v_t)\mapsto(X_1v_t,X_2v_1,\ldots,X_tv_{t-1})$,
use change-of-basis factors
$I,X_2,X_3X_2,\ldots,X_t\cdots X_2$.  The resulting cycle
product is $X_1X_t\cdots X_2$.  Choosing another starting factor
conjugates that product and does not change its eigenspaces.
In this normal form $T^t=X^{\otimes t}$.

Choose an eigenbasis with eigenvalues $\alpha_1,\ldots,\alpha_a$,
repeated with their multiplicities.  For a proposed eigenvalue $z$
of $T$, let $N_z$ count the ordered words $(i_1,\ldots,i_t)$
with $\prod_j\alpha_{i_j}=z^t$.  Fixing the first $t-1$
indices leaves at most $a-\nu(X)$ possibilities for the last,
so
\[
                         N_z\le a^{t-1}(a-\nu(X)).             \tag{AST3.8}
\]
A word orbit under cyclic rotation has a length $d$ dividing $t$.
Its span is a weighted cyclic matrix of degree $d$, with nonzero
weights.  Its characteristic polynomial is $Z^d-c$ with $c\ne0$,
so each of its eigenvalues occurs exactly once.  If it contributes
to the $z$-eigenspace, all its $d$ words are counted by $N_z$.
Its contribution to $t\dim\ker(T-zI)-N_z$ is then $t-d$;
an orbit which does not contribute has nonpositive contribution.
Full-length orbits therefore require no error term.

Let $C_d(a)$ be the number of cyclic word orbits of exact length
$d$ for $d\mid t$.  The total possible short-orbit error is
\[
 \sum_{\substack{d\mid t\\d<t}}(t-d)C_d(a)=R_t(a).            \tag{AST3.9}
\]
To prove the equality, a word orbit of length $d$ has $d$ words,
each fixed by $t/d-1$ nonidentity rotations.  Its contribution to
the count of pairs consisting of a word and a nonidentity rotation
fixing it is $d(t/d-1)=t-d$.  Rotation by $j$ positions fixes
exactly $a^{\gcd(t,j)}$ words.  Counting the same pairs in these
two ways proves (AST3.9), including the factor $t-d$.
Consequently $t\dim\ker(T-zI)\le N_z+R_t(a)$, and (AST3.8)
gives (AST3.6).  Repeated eigenvalues do not merge the distinct
basis-word orbits, so the proof includes all such degeneracies.

For completeness the universal bound (AST3.7) has a short direct
proof from (AST3.5).  For $t=3,4,5$ the errors are respectively
$2a$, $a^2+2a$ and $4a$, each at most $2a^{t-2}$ for $a\ge4$.
For $t\ge6$, each $\gcd(t,j)\le\lfloor t/2\rfloor$, hence
\[
 R_t(a)\le(t-1)a^{\lfloor t/2\rfloor}.
 \qquad
 t-1\le2a^{\lceil t/2\rceil-2}.                              \tag{AST3.10}
\]
The second inequality holds at $t=6$ and $t=7$ for $a\ge4$.
Increasing $t$ by two multiplies its right side by $a$, while
$t+1\le a(t-1)$.  Induction in both parities proves it for every
$t\ge6$, and substitution proves (AST3.7).  This proof is not an
extrapolation from a finite list of tensor counts.
\end{proof}

\subsection*{HD-AST-12. Exclusion for all larger tensor counts}

\begin{theorem}[HD-AST-12]
In the setting (AST3.1)--(AST3.2), no actual curve-cover source has
six-dimensional multiplicity variation with period-image rank nine.
Thus every transitive ordinary standard core $A_n^t$ with
$n\ge15$ and $t\ge3$ is excluded for each non-Eisenstein original
imaginary quadratic field.  Together with HD-AST-09 this excludes
all such tensor counts $t\ge2$ over those fields.  HD-AST-09
also retains its stronger all-field conclusion at $t=2$.
\end{theorem}

\begin{proof}
The actual intermediate-cover factorization of D-AST-03 and
HD-AST-04--09 applies without identifying $F$ with $K$.
If $h$ is the intermediate base genus and $D$ the sum of the
actual moving ranks at its genuine marks, the multiplicity Euler
identity is $6=2e(h-1)+D$.  Since $e>6$, genus at least two
is impossible.  Genus one has $D=6$ and at most six genuine
marks; finite cover data over the pointed base have dimension
at most six, so the actual period image cannot have rank nine.
The proposed rank-nine case must consequently have
\[
 h=0,\qquad r\ge12,\qquad D=2e+6,                            \tag{AST3.11}
\]
for a genuine generating product-one inertia tuple in $J$.
The inequality $r\ge12$ follows from the actual Hurwitz family
over the $r$ pointed sphere, of dimension $r-3$.  No sufficiency
of this dimension condition is being asserted.

Rationality of $T(g)$ gives $F\subseteq\mathbf Q(\zeta_M)$.
If all scalar values lay in $\mu_6$, it would force
$K=\mathbf Q(\sqrt{-3})$, contrary to (AST3.1).  The images of
the tuple therefore generate a nontrivial quotient
$\lambda(J)/(\lambda(J)\cap\mu_6)$ and have product one.
At least two entries have $\lambda\notin\mu_6$.

We recall the exact support estimate HD-AST-05 at these marks.
Every power of a rational tensor-wreath operator has trace at
least $-e/a$.  Indeed its trace is a product of standard
cycle-product traces, each between $-1$ and $a$; a negative
product has a factor $-1$, and there are at most $t$ factors.
For a rational matrix, a primitive $d$th root eigenspace has
equal multiplicity at all its cyclotomic conjugates.  If
$d\notin\{1,2,3,4,6\}$, its multiplicity is at most $e/4$.
For $d=4$, the equal $i,-i$ multiplicities $m$ satisfy
$-e/a\le\operatorname{tr}(T^2)\le e-4m$.
Thus each of the two scalar marks above has actual moving rank
at least
\[
                           \frac{3-1/a}{4}e.                  \tag{AST3.12}
\]
This is applied to the selected eigenspace of the rational
operator before its scalar twist; it is not a claim that a
nonreal-eigenvalue bound survives arbitrary scalar convolution.

Every genuine entry has actual moving rank at least $e/a$.
For a factor-fixing nonscalar entry, an active standard factor
has projective support at least one, and convolution gives this
bound.  A genuine scalar has full moving rank.  If an entry moves
the factors, select one tensor cycle of length $l\ge2$ and put
it in the weighted normal form used above.  Of the $a^l$ basis
words, $a$ are constant and all other rotation orbits have length
at least two.  Each orbit contributes at most one vector to a
specified eigenspace.  Its largest eigenspace fraction is therefore
at most $(1+a^{1-l})/2\le(1+1/a)/2$.  Tensoring with the other
cycles preserves this bound.  Every factor-moving entry thus has
actual moving rank at least
\[
                          s e,\qquad s=\frac{a-1}{2a}\ge\frac1a.
                                                                    \tag{AST3.13}
\]

If three or more entries moved factors, choose the two scalar
marks of (AST3.12) and a third, distinct factor-moving entry.
The remaining at least nine genuine entries give
\[
 D\ge\left(2\frac{3-1/a}{4}+\frac{a-1}{2a}
                         +\frac9a\right)e
       =2e+\frac{8e}{a}>2e+6,                               \tag{AST3.14}
\]
contrary to (AST3.11).  There are at most two factor-moving
entries.  Transitivity requires at least one, and product one
forbids exactly one.  Hence there are precisely two.  Their
factor permutations are inverse, and their cyclic subgroup is
transitive, so both are full $t$-cycles.  The same estimate would
force the two scalar marks to be these two moving entries, but
no use of their scalar orders is made in the rest of the proof.

Let $h_1,h_2\in S_n$ be the products around those two full
cycles, and put $H=k_{h_1}+k_{h_2}$.  Every full-letter
eigenspace of $h_i$ has dimension at most $c_{h_i}$; deleting
one trivial line cannot increase it.  Therefore
\[
                 \nu(U(h_i))\ge a-c_{h_i}=k_{h_i}-1.         \tag{AST3.15}
\]
This intentionally weaker inequality remains valid when $h_i=1$.
Applying HD-AST-11 and (AST3.7), the sum of the two moving
ranks is at least
\[
 2\left(1-\frac1t\right)e+\frac{a^{t-1}}tH
                  -\frac{2a^{t-1}}t-\frac{4a^{t-2}}t.        \tag{AST3.16}
\]

Let $N=\sum I_g$ over the factor-fixing entries.  The actual
permutation representation in (AST3.2) is transitive on $nt$
letters: the core is transitive within each block, and the
factor action is transitive on the blocks.  The inverse image of
a letter stabilizer in the actual connected cover defines a
connected degree-$nt$ quotient cover of the same sphere.
A factor-fixing entry has index $I_g$ there.  A full tensor
cycle with product $h_i$ has index $n(t-1)+k_{h_i}$, because
each length-$l$ cycle of $h_i$ gives a length-$tl$ cycle on
the letters.  Riemann--Hurwitz on this actual quotient implies
\[
               N+H+2n(t-1)\ge2nt-2,
               \qquad N+H\ge2n-2=2a.                       \tag{AST3.17}
\]
Scalar marks have zero permutation index but remain genuine
coefficient-cover marks in the moving-rank sum.

There are at least ten factor-fixing entries.  Sum (AST3.3)
and (AST3.16), then use $a-(n-3)=2>0$ and (AST3.17):
\[
\begin{split}
 D&\ge2\left(1-\frac1t\right)e
       +\frac{a^{t-2}(n-3)}tN+\frac{a^{t-1}}tH
       +5na^{t-2}-\frac{2a^{t-1}}t-\frac{4a^{t-2}}t\\
  &\ge2\left(1-\frac1t\right)e
       +\frac{a^{t-2}(n-3)}t(2a)
       +5na^{t-2}-\frac{2a^{t-1}}t-\frac{4a^{t-2}}t\\
  &=2e+\frac{a^{t-2}}t\bigl((5t-6)n+2\bigr)>2e+6.
\end{split}                                                   \tag{AST3.18}
\]
The last excess is at least $3na^{t-2}\ge630$ for
$n\ge15$, $a\ge14$, $t\ge3$.  It contradicts (AST3.11),
proving the theorem for every tensor count in its stated range.
\end{proof}

\paragraph{Checks, sources and remaining scope.}
The all-size arguments are universal proofs.  A separate exact
supplement checks weighted cyclic spectra, the full rotation error,
actual letter-action indices, pair selection and the symbolic sum;
it does not enumerate Hurwitz components or infer an infinite
theorem from finite data.  Known normalizer, scalar-field and
actual-cover inputs are the complete reviewed D-AST-01--03 and
HD-AST-04--09.  The elementary cyclotomic input in HD-AST-05 is
Milne, \emph{Class Field Theory}, v4.03, Chapter V, Example 3.2,
printed p.~157.\cite{HD:MilneCFT2020}  The weighted-cycle proof
is credited to D-ATT-08, with its exact short-orbit defect retained.

This result restricts actual finite local systems and possible
period-rank-nine suppliers.  It supplies no algebraic cycle and
makes no new positivity or motivic assertion.  It does not decide
the Eisenstein case with $t\ge3$, nontransitive tensor configurations,
other ordinary alternating labels or faithful spin modules.  In
particular the original field is never silently replaced by the
larger inducing character field.

\section{A general tensor-count and scalar-kernel bound}
\label{sec:tensor-count-scalar-bound}

\paragraph{Status and scope.}
New root derivation HD-TC01--03. Complete root and separate
/root/g34\_eightfolds whole-proof, source and exact-matrix reviews
have passed on the refined inducing-image scope. The bounds concern an actual inducing block of
a finite-cover sixfold source with period rank nine. In its
accepted primitive normal-core gate, the irreducibly acting
perfect normal subgroup is a central product of quasisimple
groups, and the full inducing group acts transitively on its
simple factors. Their complex module dimensions are consequently
equal; write $a\ge2$, $t$ for the number of factors and $e=a^t$.
The scalar, outer and factor-permuting actions are all retained.
No equality of original and inducing fields is assumed.

The actual base-cover and marked-moduli arguments established in
D-BI-01 and the subsequent normal-core reductions give a sphere
base, a generating product-one inertia tuple with $r\ge12$
genuine entries, and the exact identity
\[
              D:=\sum_{i=1}^r\operatorname{rank}(T_i-I)
                           =2e+6.                       \tag{TC.1}
\]
The natural factor-permutation quotient is an actual connected
cover of degree $t$. Thus its inertia permutations $\pi_i$ satisfy
\[
          \sum_i\operatorname{ind}(\pi_i)\ge2t-2,
 \qquad\operatorname{ind}(\pi)=t-\#\operatorname{cycles}(\pi).
                                                               \tag{TC.2}
\]
These are hypotheses furnished by the actual source, not
properties imposed on an arbitrary list of matrices.

\subsection*{HD-TC01. A cyclic-average tensor permutation bound}

\begin{lemma}[HD-TC01]
Let $T$ be a finite-order tensor normalizer on a tensor product
of $t$ spaces of dimension $a\ge2$, with factor permutation $\pi$.
Then its maximum-eigenspace codimension satisfies
\[
 \frac{\nu(T)}{a^t}\ge
       \bigl(1-a^{-t/2}\bigr)\frac{\operatorname{ind}(\pi)}t.
                                                               \tag{TC.3}
\]
If $\pi\ne1$, it also satisfies
$\nu(T)\ge(a-1)a^{t-1}/2\ge a^t/4$.
\end{lemma}

\begin{proof}
Choose invariant positive Hermitian forms on the irreducible
finite core factors and take their tensor product. Its scalar
uniqueness follows from irreducibility. A finite-order normalizer
is unitary for this form: its positive similitude is a positive
real root of unity. Its tensor maps can therefore be rescaled
to unitary maps, with only an overall phase remaining.
The usual tensor trace identity gives, for every integer $j$,
\[
 |\operatorname{tr}(T^j)|\le
                       a^{\#\operatorname{cycles}(\pi^j)}.
                                                               \tag{TC.4}
\]
It is the product of the traces of the operators around the
cycles, each of absolute value at most $a$.

Let $f_j$ be the number of fixed letters of $\pi^j$.
Every nontrivial permutation cycle has index at least half
its length, so
\[
 a^{\#\operatorname{cycles}(\pi^j)-t}
       \le a^{-(t-f_j)/2}
       \le\frac{f_j}{t}
                  +\left(1-\frac{f_j}{t}\right)a^{-t/2}.
                                                               \tag{TC.5}
\]
The last inequality is the endpoint chord inequality for the
convex function $x\mapsto a^{-(t-x)/2}$ on $[0,t]$.
If $L$ is the order of $T$, it is a multiple of the order
of $\pi$. The cyclic orbit-counting identity gives
$L^{-1}\sum_{j=0}^{L-1}f_j=\#\operatorname{cycles}(\pi)$.
For every eigenvalue $z$ of $T$, the Fourier formula and
(TC.4)--(TC.5) therefore bound its multiplicity by
\[
 \frac{\dim\ker(T-zI)}{a^t}
 \le\frac1L\sum_j\frac{|\operatorname{tr}(T^j)|}{a^t}
 \le\frac{\#\operatorname{cycles}(\pi)}t
       +\frac{\operatorname{ind}(\pi)}t a^{-t/2}.
                                                               \tag{TC.6}
\]
This proves (TC.3) for every target eigenvalue and every phase.

For the second assertion choose one factor cycle of length
$l\ge2$. Conjugating its weighted tensor rotation puts its
cycle product on one factor; diagonalize that product.
The resulting tensor basis is permuted up to nonzero scalar
weights. There are $a$ constant words. Every other word orbit
has length at least two and contributes at most one vector
to a specified eigenspace. Hence the maximum eigenspace on
this cycle has dimension at most $(a^l+a)/2$.
Its support fraction is at least
$(1-a^{1-l})/2\ge(a-1)/(2a)$. Tensoring with the other
cycles cannot increase the largest normalized eigenvalue
multiplicity, by convolution of the spectral multiplicities.
This proves the second bound. It is also the weighted-word
argument already used in D-ATT-08, with all short orbits retained.
\end{proof}

\begin{theorem}[HD-TC02]
Every source in the stated transitive-core gate satisfies
\[
 t\le T(a):=\max\left\{4,\left\lceil\frac{5a}{4}\right\rceil-1\right\}.
                                                               \tag{TC.7}
\]
This applies over every original imaginary quadratic field,
to arbitrary nontrivial factor modules and represented central
quotients, without a character-ratio or self-duality assumption.
\end{theorem}

\begin{proof}
There is nothing to prove if $t\le4$. Otherwise $t\ge5$
and $e=a^t\ge32$. If nine inertia entries moved factors,
HD-TC01 would give $D\ge9e/4>2e+6$, since $e>24$.
At most eight entries move factors, leaving at least four
genuine factor-fixing entries.

At such an entry with a nonscalar tensor factor, that
factor has maximum-eigenspace codimension at least one.
Spectral convolution gives moving rank at least $e/a$.
If all tensor factors are scalar, the genuine entry is
a nonidentity scalar and has full moving rank $e$.
The same lower bound holds in that case.
Sum these four bounds and use HD-TC01 and (TC.2) on the
factor-moving entries. The result is
\[
 \begin{split}
 2+\frac6e
 &\ge\left(1-a^{-t/2}\right)\left(2-\frac2t\right)+\frac4a,\\
 \frac4a&\le\frac2t+2a^{-t/2}+6a^{-t}<\frac5t.
 \end{split}                                             \tag{TC.8}
\]
For the last strict inequality, $2^t>t^2$ and $2^t>6t$
for every integer $t\ge5$. The first follows from equality
at $t=4$ and induction using $2t^2>(t+1)^2$ for $t\ge3$;
the second follows from $32>30$ and immediate induction.
Thus $a^{-t/2}<1/t$ and $6a^{-t}<1/t$.
Consequently $t<5a/4$. Taking the largest admissible integer
proves (TC.7).
\end{proof}

\subsection*{HD-TC03. Scalars and finite bounded-factor catalogues}

\begin{theorem}[HD-TC03]
Let $J$ be the actual finite inducing image and let $S$ be
its irreducibly acting perfect normal central product.
Put $P=S/Z(S)$ and $M=|\operatorname{Aut}(P)|$.
Then the scalar kernel and the whole image satisfy
\[
 |J\cap\mathbf C^\times I|\mid eM,
 \qquad |J|\le eM^2.                                    \tag{TC.9}
\]
Suppose a fixed finite catalogue of allowed simple factors
and factor modules has degree bound $A\ge2$ and automorphism
order bound $B\ge1$ for its simple factors. Put
\[
 T=T(A),\qquad E=A^T,\qquad
 Q=E\bigl(B^T T!\bigr)^2.                               \tag{TC.10}
\]
Every actual source drawn from this catalogue has
\[
 e\le E,\quad |J|\le Q,\quad r\le2E+6,
 \quad g(C)\le1+Q(E+2)                                 \tag{TC.11}
\]
for its connected Galois cover $C$ of the sphere.
In particular its inducing-image, scalar and Nielsen-component
enumeration is finite and effective. Original-field descent
is retained as a separate requirement.
\end{theorem}

\begin{proof}
Conjugation gives $J\to\operatorname{Aut}(P)$.
If an element induces the identity on $P$, its conjugation
on $S$ has the form $s\mapsto z_s s$ with $z_s\in Z(S)$.
Centrality makes $s\mapsto z_s$ a homomorphism. Perfectness
of $S$ makes it trivial. The element therefore centralizes
$S$ and is scalar by irreducibility. Conversely a scalar
centralizes $S$, so the projective image of $J$ embeds in
$\operatorname{Aut}(P)$ and has order dividing $M$.

There cannot be two inertia entries of full moving rank:
they and the other at least ten genuine entries would
give $D\ge2e+10>2e+6$. Every entry except possibly one
therefore has eigenvalue one. Its $M$th power is scalar
by the projective embedding, and evaluation on its fixed
vector shows that this scalar is one. Thus its determinant
has order dividing $M$. Product one gives the same
determinant bound for the possible exceptional entry.
Since the actual inertia tuple generates $J$, the entire
determinant image is contained in $\mu_M$.

For a scalar $zI\in J$, its determinant is $z^e$; hence
$z^{eM}=1$. The finite scalar subgroup is cyclic and its
order divides $eM$. Multiplying its size by the order of
the projective image proves (TC.9). No rational scalar
section of the projective normalizer is needed.

By transitivity all simple factors are isomorphic, so
$P\simeq L^t$ and
$|\operatorname{Aut}(P)|=|\operatorname{Aut}(L)|^t t!$.
Together with HD-TC02 this gives the numerical bounds
$e\le E$ and $M\le B^T T!$, hence $|J|\le Q$.
Every genuine inertia contributes at least one to (TC.1),
so $r\le2e+6\le2E+6$. The Galois Riemann--Hurwitz formula
then gives
\[
 2g(C)-2=|J|\left(-2+\sum_i(1-1/o(T_i))\right)
              \le |J|(r-2)\le Q(2E+4),                  \tag{TC.12}
\]
which proves the genus bound.

For effectiveness, enumerate finite multiplication tables
of order at most $Q$, their irreducible characters of degree
at most $E$, and tuples of length at most $2E+6$ satisfying
generation and product one. All are finite algebraic or
finite-group operations. The original-field, Schur-index,
signature and central-core requirements are retained; they
are not inferred from the group-order bound.

There is a precise distinction in the rank step. D-NR-01--03
as stated takes an absolutely irreducible representation over
an imaginary quadratic field. In that stated domain its
terminating rank algorithm applies directly. For a nonselfdual
inducing character over a larger algebraic number field, its
complex multiplicity-variation rank is still computable by
the same construction: D-NR works over the computable field
$\overline{\mathbf Q}$; its generic cover construction uses
only the group and tuple; and the intertwiner equations and
dual-basis contraction (NR.6)--(NR.9) are finite algebraic
matrices after adjoining the representation coefficients.
The Cech/Serre-duality calculation in D-NR-03 identifies this
matrix with that complex multiplicity period codifferential,
without identifying its coefficient field with the original
quadratic field. Thus these inducing-component rank tests
are a finite computation. An actual original-field supplier
still requires its descent and polarization interface to be
identified. No enumeration of an arbitrarily larger original
inducing-overgroup, or automatic original-field descent from
the intermediate cover, is asserted. No such large finite
Nielsen enumeration has been executed here.
\end{proof}

\paragraph{Strictness and remaining boundary.}
For example the ordinary label $(96,4)$ at $n=100$ has
$a=\binom{100}{4}-\binom{100}{3}$ and satisfies the earlier
Gaussian necessary long-row and field inequalities.
Its profile with $t=2a$ was not removed by a bound on tensor
count for general remaining labels; HD-TC02 now excludes it.
This is strict removal of a necessary profile, without a
claim that an actual source existed there.
For every bounded classical-factor catalogue the result
makes the inducing-image and Nielsen-component rank enumeration
finite. Unbounded ordinary alternating natural degrees and
the stated original-field descent requirement remain separate
boundaries. No nonexistence or dominance conclusion is inferred
from finiteness alone.

These inequalities concern actual finite etale local systems
away from the branch divisor and their tensor operations in
representation lambda-rings and K-theory. They make a part of
the condition-A supplier gate effective (cote16, scan050;
SIGN\_NEXT SN-C03/C04), retaining the established schemes-over-
$\mathbf Z$, EGA/SGA/FGA and motivic framework. They produce
no condition-I positive cycle and introduce no Markman,
Severn or032 premise.

\section{A derived tensor-count bound using the actual factor support}
\label{sec:tensor-count-local-support}

\paragraph{Status and evidential classification.}
HD-TC04 is a direct quantitative consequence of accepted HD-TC01--02
and the stated factor support bound. It is recorded as derived work,
not as an independent advance for the hourly movement rule.
Complete root and separate /root/g34\_eightfolds whole-proof reviews
have passed. The exact illustrative arithmetic was checked separately;
this corollary remains derived work for the movement rule.

\begin{corollary}[HD-TC04]
In the actual transitive-core source setting of HD-TC01--03, suppose
every compatible nonscalar factor operator has maximum-eigenspace
codimension at least $\eta a$, where $0<\eta\le1$.
Then
\[
 t\le\max\left\{4,\left\lceil\frac5{4\eta}\right\rceil-1\right\}.
 \tag{TCL.1}
\]
In particular, for a remaining ordinary alternating label
$\lambda=(n-k,\mu)$ with $n\ge15$ and $0<k<n/4$,
\[
 t\le\max\left\{4,
   \left\lceil\frac{5(n-1)(n-2)}{4k(n-1-k)}\right\rceil-1\right\}.
 \tag{TCL.2}
\]
Both assertions keep all compatible scalar and outer operators and
actual factor permutations. They do not assert a source in a
numerically permitted profile.
\end{corollary}

\begin{proof}
For $t\le4$ there is nothing to show. For $t\ge5$, the complete
HD-TC02 counting argument leaves at least four genuine inertia
entries fixing every factor. At any such entry with an active
nonscalar factor, spectral convolution gives moving rank at least
$\eta e$, where $e=a^t$. If all its factors are scalar, the genuine
entry is a nonidentity scalar and has moving rank $e\ge\eta e$.

The factor-moving entries contribute the same actual quotient-cover
index sum as before. Thus, with every term justified by HD-TC01,
\[
 2+\frac6e\ge
 (1-a^{-t/2})\left(2-\frac2t\right)+4\eta,
 \qquad
 4\eta\le\frac2t+2a^{-t/2}+6a^{-t}<\frac5t.
 \tag{TCL.3}
\]
The strict last inequality is the already proved integer estimate
in (TC.8), using $a\ge2$ and $t\ge5$. It gives $t<5/(4\eta)$,
with the exact integral conversion in (TCL.1).
The published-input hypotheses needed for the alternating
substitution were checked in HD-ALT-07a: every compatible nonscalar
normalizer has support at least
$a k(n-1-k)/((n-1)(n-2))$. Substituting this positive $\eta$
proves (TCL.2); no further character theorem is invoked.
\end{proof}

\paragraph{Concrete comparison, without an existence assertion.}
For $n=100$ and $\lambda=(96,4)$, the factor degree is
\[
 a=\binom{100}{4}-\binom{100}{3}=3759525,
 \qquad \eta=\frac{190}{4851},
 \qquad \frac5{4\eta}=\frac{4851}{152}<32.
 \tag{TCL.4}
\]
The derived bound is therefore $t\le31$, whereas the general
dimension-only HD-TC02 bound is $t\le4699406$.
This compares two necessary inequalities; it neither supplies
branch data nor changes the field, period-rank or cycle requirements.

\section{Original overgroups and descent of an inducing factor}
\label{sec:original-overgroups}

\paragraph{Status and scope.}
Statements D-OG-01--03 were developed during the execution-environment
outage on 22 November 2026.  The complete message proofs received root
whole readings and an independent whole mathematical reading by
/root/g34\_eightfolds before this file was saved.  The author, root and separate peer then read this entire saved version
after recovery and passed it; the frozen-file hashes and scopes are recorded.  The new statements concern
an absolutely irreducible coefficient representation defined over its
exact imaginary quadratic character field.  They do not concern an
additional quadratic action unrelated to that coefficient representation.
The original overgroup and induction index can be unbounded.  This does
not rule out a bound on a \emph{minimal} realization of a given inducing
datum.  The determinant obstruction HD-PD-01--02 is stated separately.

\subsection*{D-OG-01. A faithful wreath representation with a fixed inducing image}

\begin{proposition}[D-OG-01]
Let $J$ be finite, let $K$ be imaginary quadratic, and let $\tau$ be a
nontrivial faithful absolutely irreducible $K$-representation of degree
$e$ and exact character field $K$.  Let
$(g_1,\ldots,g_r)$ be a generating product-one tuple of nonidentity
elements of $J$.  For every $m\ge2$ set
\[
 G_m=J^m\rtimes S_m,\qquad
 H_m=J_1\times(J^{m-1}\rtimes S_{m-1}),\qquad
 N_m=J^{m-1}\rtimes S_{m-1}.
 \tag{OG.1}
\]
Inflate $\tau$ through $H_m/N_m=J$.  Then
$\rho_m=\operatorname{Ind}_{H_m}^{G_m}\tau$ is faithful, absolutely
irreducible, defined over $K$, and has exact character field $K$.
Its degree is $me$ and $[G_m:H_m]=m$.

Put $a_j=(g_j,1,\ldots,1)$, and let $s_i=(i,i+1)$ be a pure adjacent
block swap.  The tuple
\[
 (a_1,\ldots,a_r,s_1,s_1,\ldots,s_{m-1},s_{m-1})
 \tag{OG.2}
\]
is product one and generates $G_m$.  The two occurrences of every
$s_i$ are assigned distinct branch points.
\end{proposition}

\begin{proof}
Use the block model $W=\bigoplus_{i=1}^m W_i$, with $W_i=K^e$,
where $J^m$ acts coordinatewise and $S_m$ permutes the blocks.
Over an algebraic closure, the $W_i$ are pairwise nonisomorphic
irreducible $J^m$-modules: the nontrivial representation is supported
on a different direct factor in each case.  Semisimplicity implies
that a $J^m$-stable subspace is a sum of these blocks.  Transitivity
of $S_m$ leaves only zero and $W$ as $G_m$-stable possibilities.
This proves absolute irreducibility.

An element acting identically on $W$ must have trivial block
permutation, and then every $J$-coordinate lies in $\ker\tau=1$.
Thus the representation is faithful.  Its matrices are over $K$.
Every trace is the sum of the traces of the coordinate operators
on the fixed blocks of the permutation, so belongs to $K$.
Conversely,
\[
 \operatorname{tr}\rho_m(j,1,\ldots,1)
       =\chi_\tau(j)+(m-1)e,
 \tag{OG.3}
\]
so the character field contains the exact field of $\tau$ and is $K$.

The product in (OG.2) is one.  The $a_j$ generate $J_1$, the $s_i$
generate $S_m$, and their conjugates give every coordinate copy of
$J$.  Hence the tuple generates $G_m$.  No entry is the identity.
\end{proof}

\paragraph{The Chevalley--Weil count is consistent with the construction.}
If the seed moving-rank sum is $2e+6$, then each $s_i$ has moving
rank $e$ and eigenvalue $-1$ with multiplicity $e$.  Thus
\[
 \sum \operatorname{rank}(1-\rho_m(t_j))
       =(2e+6)+(2m-2)e=2me+6.
 \tag{OG.4}
\]
At either embedding of $K$, each swap contributes age $e/2$.
The first $r$ entries have the seed nontrivial spectrum and
$(m-1)e$ extra eigenvalues one.  The two holomorphic multiplicities
therefore remain $(3,3)$.  The following geometric proof supplies
the actual family and does not treat this dimension count as sufficient.

\subsection*{D-OG-02. The rational intermediate cover, actual period rank and polarization}

\begin{proposition}[D-OG-02]
Suppose that the seed $J$-Hurwitz component over $\mathbf P^1$
with tuple $(g_1,\ldots,g_r)$ has a primitive $\tau$-factor of
signature $(3,3)$ and actual period rank nine.  For every $m\ge2$,
(OG.2) gives a connected $G_m$-Hurwitz component with a primitive
$\rho_m$-factor of signature $(3,3)$ and actual period rank nine.
Its inducing image and genuine inducing tuple are the seed $J$
and $(g_1,\ldots,g_r)$.

If $P$ is the seed primitive factor and $B_m$ is the corresponding
original primitive factor, there is an actual $K$-isogeny
$u:P\longrightarrow B_m$ for which the restricted canonical
Jacobian polarizations satisfy
\[
 u^*\theta_{B_m}=|N_m|\theta_P,\qquad
 |N_m|=|J|^{m-1}(m-1)!.
 \tag{OG.5}
\]
In particular their rank-six signed Hermitian determinant classes
in $\mathbf Q^*/\operatorname{Nm}_{K/\mathbf Q}K^*$ agree.
\end{proposition}

\begin{proof}
Riemann existence, in the finite-cover framework already checked
in D-NR-01--02, realizes (OG.2) by a connected $G_m$-cover
$C\longrightarrow\mathbf P^1$.  For $X=C/H_m$, the $a_j$
have trivial permutation, and the remaining $2m-2$ branch points
have simple transposition inertia.  Hence
\[
 2g(X)-2=-2m+(2m-2)=-2.
 \tag{OG.6}
\]
Thus $X=\mathbf P^1$, and $f:X\to\mathbf P^1$ is a simply
branched rational map with distinct critical values.

Set $Y=C/N_m$.  The quotient $Y\to X$ is a connected $J$-cover.
Take a disk containing the first $r$ base branch values and no
swap values.  The cover $X$ over this disk is $m$ disjoint disks.
The first has $J$-monodromy exactly $g_1,\ldots,g_r$, while
the other disks have trivial $J$-monodromy.  At a swap value,
the moving two-sheet orbit has return $s_i^2=1$; on any fixed
sheet the return lies in the kernel of its first-coordinate
$J$-quotient.  Thus the $J$-cover has precisely the $r$ genuine
seed marks.  The displayed disk gives its seed Nielsen tuple.
There are no additional genuine $J$-branch points.
Moreover the Galois closure of $Y\to\mathbf P^1$ is $C$:
the core of $N_m$ in $G_m$ is the kernel of the induced faithful
representation, hence is trivial.

We next prove dominance onto the seed family.  Work near this
constructed point and normalize the rational intermediate curve
by three genuine seed marks.  The seed component has $r\ge3$
(indeed rank nine requires at least twelve marks).
Retain the rational map $f_0$ in this local normalization while
deforming the marked seed cover.  Simplicity of the critical
points, distinctness of their values, distinctness of the images
of the genuine marks, and avoidance of the critical values are
open conditions.  They all hold at the constructed point.
Every sufficiently small deformation of the seed cover can
therefore be composed with this fixed $f_0$.  The ordered
topological monodromy is locally constant on this open set,
so the full wreath group persists.

Equivalently the space of pairs consisting of the seed cover
and such a map $f$ forgets $f$ submersively at this point.
Relative Galois closure supplies an algebraic family in the
finite Hurwitz framework of D-NR; it lies in one original
$G_m$-component and maps onto an open set of the seed component.
No submersion statement for arbitrary prescribed ramification
is used.  Frobenius reciprocity and quotient pullback give an
identification of the weight-one multiplicity variations:
\[
 \operatorname{Hom}_{G_m}(\rho_m,H^1(C))
 \simeq
 \operatorname{Hom}_{H_m}(\tau,H^1(C))
 \simeq
 \operatorname{Hom}_J(\tau,H^1(Y)).
 \tag{OG.7}
\]
These are fixed group-algebra maps preserving Hodge filtrations.
The original component consequently has period rank at least
nine.  Its $(3,3)$ unitary target has dimension nine, so its
rank is exactly nine.

For the polarization statement, choose the seed's rational
Rosati-fixed primitive idempotent $p$ in its component
$\operatorname{Mat}_e(K)$ of $\mathbf Q[J]$.
Such idempotents can also be obtained by orthogonal projection
onto a $K$-line for the averaged positive coefficient Hermitian
form.  The canonical quotient-algebra embedding
$\mathbf Q[J]\to e_{N_m}\mathbf Q[H_m]$ sends it to
\[
 f_p=e_{N_m}p,\qquad
 e_{N_m}=|N_m|^{-1}\sum_{n\in N_m}n.
 \tag{OG.8}
\]
Here the notation means the quotient-algebra embedding, and
does not require a choice of representatives to be multiplicative.
Frobenius reciprocity and irreducibility of
$\operatorname{Ind}\tau$ show that $f_p$ is supported only on
the rational $\rho_m$-block, where it has $K$-rank one.
Thus it is a rational primitive idempotent of $\mathbf Q[G_m]$.
Both $e_{N_m}$ and $p$ are self-adjoint under group inversion.

Let $P=pJ(Y)$ and $B_m=f_pJ(C)$ denote connected images after
clearing the rational projector denominators.  For the actual
quotient map $\pi:C\to Y$, Jacobian pullback restricts to an
isogeny $u=\pi^*|_P:P\to B_m$.  The norm is its adjoint and
$\pi_*\pi^*=[|N_m|]$.  Restricting the canonical polarizations,
or applying the intersection projection formula, gives (OG.5).
No equality of integral polarization types or uniform degree
bound is asserted.  Use the same anti-invariant element of $K$
to normalize the Hermitian forms.  On their rank-six rational
spaces the determinant scale is
\[
 |N_m|^6=\operatorname{Nm}_{K/\mathbf Q}(|N_m|^3),
 \tag{OG.9}
\]
which proves the determinant assertion.
\end{proof}

\paragraph{An actual nonsplit example for every index.}
The accepted E8 tuple and period results D-E8T and D-E8R give
$J=\mu_4W(E_8)$, $e=8$, $K=\mathbf Q(i)$, signature $(3,3)$,
actual rank nine, and nonsplit signed determinant class $3$.
D-OG-02 therefore gives actual such sources of coefficient
degree $8m$ and induction index $m$ for every $m\ge2$.
The E8 rank assertion uses the complete period certificate
D-E8R, not merely the tuple's dimension count.

\paragraph{Padding an already realized proper-field datum.}
Suppose $\rho_0=\operatorname{Ind}_{H_0}^{G_0}\tau$ already
has exact field $K$, while $\tau$ has field $L\supsetneq K$,
and the original family has actual rank nine.
Apply the preceding construction to $G_0,\rho_0$, and use
\[
 H'=H_0\times(G_0^{m-1}\rtimes S_{m-1})
       \ \le\ G_0\wr S_m.
 \tag{OG.10}
\]
Induction transitivity identifies $\operatorname{Ind}_{H'}\tau$
with the faithful wreath representation of $\rho_0$.
The faithful inducing image is still $H_0/\ker\tau$ and the
index is multiplied by $m$.  The intermediate $G_0$-cover is
the old seed family, so its further quotient by $H_0$ is the
same rational inducing base and has the same genuine inducing
datum.  This assertion is at the family level, after restriction
to the nonempty open seed locus supplied by the submersion.
It does not assert literal padding of every prescribed individual
seed point.  The new positive polarization scale is
$|G_0|^{m-1}(m-1)!$, again invisible in the determinant norm class.
No new existence of proper-field descent is inferred.

\subsection*{D-OG-03. Rational descent is an additional condition}

\begin{proposition}[D-OG-03]
Let $\rho_K$ be an absolutely irreducible $K$-model of exact
imaginary quadratic character field $K$, and assume
$\rho_{\mathbf C}=\operatorname{Ind}_H^G\tau$ with $\tau$
absolutely irreducible.  Put $L=\mathbf Q(\chi_\tau)$ and
$\ell=[L:K]$.  Then $K\subseteq L$, $\tau$ has Schur index one
over $L$, and, with its $L$-model denoted $\tau_L$,
\[
 \operatorname{Ind}_H^G
       \operatorname{Res}_{L/K}\tau_L\simeq\rho_K^{\oplus\ell}.
 \tag{OG.11}
\]
For a curve realization, let $N=\ker\tau$, $J=H/N$, and $Y=C/N$.
If the original primitive factor is the sixfold $B$, then the
rational primitive $\tau$-factor $P_\tau$ of $J(Y)$ satisfies
\[
               P_\tau\ \sim_K\ B^\ell
 \tag{OG.12}
\]
on the actual original-family image.  These isogenies are realized
by rational deck correspondences and quotient pullback.
If the original family has actual period rank nine, then very
generally on that image
$\operatorname{End}^0(P_\tau)=\operatorname{Mat}_\ell(K)$.
\end{proposition}

\begin{proof}
The induced-character formula gives $K\subseteq L$.
Complex Frobenius reciprocity gives
\[
 \dim\operatorname{Hom}_H
       (\tau,\operatorname{Res}_H^G\rho_{\mathbf C})=1.
 \tag{OG.13}
\]
Let $T_K$ be the irreducible $K[H]$-module whose absolute orbit
contains $\tau$.  The center of its simple component is
$K(\chi_\tau)=L$.  In its scalar extension the Galois-orbit
constituents occur with the common Schur-index multiplicity $s$.
If $T_K$ occurs $n$ times in $\operatorname{Res}\rho_K$, (OG.13)
gives $ns=1$.  Hence $n=s=1$, and
$T_K=\operatorname{Res}_{L/K}\tau_L$.
These standard Wedderburn and character-field facts are the
rationality results of \cite[Section 12.2]{OG:Serre1977};
the exact multiplicity, center and adjunction statements were
also freshly checked in
\cite[Propositions 4.3.2, 4.6.14 and 4.4.2]{OG:Zheng2025}.

Every automorphism fixing $K$ sends $\operatorname{Ind}\tau$
to the same absolutely irreducible $\rho$.  Thus the character
of $\operatorname{Ind}T_K$ is $\ell\chi_\rho$.
Semisimplicity over $K$ gives (OG.11).
In particular $\ell\le[G:H]$, since $T_K$ occurs in
$\operatorname{Res}\rho_K$.

To identify actual correspondences, put
$T=\operatorname{Res}_{L/\mathbf Q}\tau_L$ and
$R=\operatorname{Res}_{K/\mathbf Q}\rho_K$.
They are rational irreducibles with endomorphism fields
$L$ and $K$, and
$\operatorname{Ind}_H^G T=R^{\oplus\ell}$.
Choose a rational Rosati-fixed primitive $p_J$ in
$\operatorname{Mat}_e(L)\subset\mathbf Q[J]$ and embed it as
$p_H=e_Np_J\in\mathbf Q[H]\subset\mathbf Q[G]$.
For a rational irreducible $G$-module $U$, the space $p_HU$
is nonzero exactly when $T$ occurs in $\operatorname{Res}U$.
Rational Frobenius reciprocity therefore implies $U=R$.
Moreover
\[
 \dim_{\mathbf Q}\operatorname{Hom}_H(T,\operatorname{Res}R)
       =\ell\dim_{\mathbf Q}\operatorname{End}_G(R)
       =2\ell=\dim_{\mathbf Q}L.
 \tag{OG.14}
\]
Thus $T$ occurs once, and $p_H$ has $K$-rank $\ell$ in the
actual simple block $\operatorname{Mat}_{\dim_K\rho}(K)$.
Its matrix units split it into $\ell$ equivalent primitive
projectors.  These are elements of $\mathbf Q[G]$, hence
rational linear combinations of deck graphs on $C$.
Clearing denominators gives actual abelian homomorphisms.
The quotient pullback identifies
$p_JH^1(Y,\mathbf Q)$ with $p_HH^1(C,\mathbf Q)$.
The matrix-unit decomposition proves (OG.12).
Equivalently, rational Frobenius on weight-one Hodge structures
identifies
$\operatorname{Hom}_H(T,H^1(C))$ with
$\operatorname{Hom}_G(R^{\oplus\ell},H^1(C))$;
the preceding idempotent argument realizes these maps geometrically.

Actual rank nine makes the original period image open in its
nine-dimensional unitary domain.  The already checked
full-domain centralizer argument gives very general
$\operatorname{End}^0(B)=K$ there, whence
$\operatorname{End}^0(B^\ell)=\operatorname{Mat}_\ell(K)$.
This is a claim on the actual original-family image, not on
the whole inducing Hurwitz component.
\end{proof}

\paragraph{What remains undecided.}
When $L\supsetneq K$, the rational isogeny-power condition and its
Rosati compatibility may hold only on a proper subfamily of the
inducing period image.  The complex rank of one chosen constituent
does not establish them.  The fixed-component exact rank algorithm
D-NR and the finite inducing-image catalogue do not by themselves
decide every such rational Hodge locus or realize its points by
an original overgroup.  Wreath padding refutes a bound on arbitrary
original overgroups, but not a bound on a minimal realization.
HD-PD-01--02 supplies a different obstruction that does not need
such a minimal-index bound.

\section{A determinant obstruction for bounded inducing images}
\label{sec:bounded-inducing-determinants}

\paragraph{Status, provenance and scope.}
These are HD-PD-01--02, with question row Q2cy reserved by root.
The simultaneous-prime obstruction and its integral strengthening
originated with root.  The polarized descent, regular-trace
matrix-order proof and homothetic-lattice proof were developed
jointly with /root/imprimitive.  /root/g34\_eightfolds has
independently read the complete message derivations and then the whole
saved proof. Author, root and separate peer frozen-file readings all pass;
the shared derivation is not represented as independent authorship.
HD-PD-02 is unconditional for the stated bounded inducing-image
class.  It does not assert that a previously enumerated structural
class exhausts all finite-group sources, and does not obstruct
arbitrary algebraic cycles.

\subsection*{HD-PD-01. Polarized Frobenius reciprocity}

Fix an imaginary quadratic field $K$ and an element
$\eta\in K^*$ with $\bar\eta=-\eta$.
Use the same $\eta$ in all Hermitian normalizations.
Suppose a smooth complex curve $C$ has a finite $G$-action.
Let $\rho_K$ be an absolutely irreducible $K$-model with exact
character field $K$, and let $B$ be the six-dimensional primitive
factor selected by a rational Rosati-fixed rank-one projector in
its $\mathbf Q[G]$ component.  Its polarization is the restriction
of the canonical Jacobian polarization, and its $K$-signature
is $(3,3)$.  Write its signed determinant as
\[
 \delta_B=(-1)^3\det(h_B)
    \quad\hbox{in}\quad
 \mathbf Q^*/\operatorname{Nm}_{K/\mathbf Q}K^*.
 \tag{PD.1}
\]
Assume
$\rho_{\mathbf C}=\operatorname{Ind}_H^G\tau$ with $\tau$
irreducible, put $N=\ker\tau$, $J=H/N$, and $Y=C/N$.
The quotient map is $\pi:C\to Y$.

\begin{proposition}[HD-PD-01]
Let $E$ be a CM cyclotomic field containing $K$ and splitting
$J$, and put $F=E^+$.  The rank-six $\tau$-multiplicity Hermitian
form in $H_1(Y,E)$ has signed determinant equal to the image of
$\delta_B$ in $F^*/\operatorname{Nm}_{E/F}E^*$.
No generic endomorphism-algebra or period-rank assumption is
required.  The field $E$ need not split the original group $G$.
\end{proposition}

\begin{proof}
Choose a model $\tau_E$.  The two representations
$\operatorname{Ind}_H^G\tau_E$ and $\rho_K\otimes_KE$ have
isomorphic complexifications, which are absolutely irreducible.
They are therefore isomorphic over $E$.
We use homology and left deck action consistently.  The extension
of the alternating intersection pairing $I_C$ gives the Hermitian
pairing
\[
                  H_C(u,v)=\eta^{-1}I_C(u,\bar v).
 \tag{PD.2}
\]
The bar on vectors conjugates the coefficient field; it is not
a new action on the underlying integral homology.
On the $\rho$-isotypic summand
$V_\rho\otimes_E M_\rho$, invariance and Schur's lemma give
\[
                         H_C=h_\rho\otimes b_\rho,
 \tag{PD.3}
\]
where $h_\rho$ is a fixed positive invariant $K$-Hermitian
coefficient form.  The multiplicity space has dimension six.
The rational primitive projector selecting $B$ selects a
$K$-line of coefficient norm $a\in\mathbf Q_{>0}$.
Thus the natural Hermitian form of $B$, under the corresponding
identification, is $a b_\rho$.  Its determinant differs by
$a^6=\operatorname{Nm}_{K/\mathbf Q}(a^3)$.

Frobenius reciprocity gives
\[
 \dim_E\operatorname{Hom}_H
          (\tau_E,\operatorname{Res}\rho_E)=1.
 \tag{PD.4}
\]
It also excludes $\tau_E$ from every other $G$-constituent:
any such constituent would receive a nonzero map from
$\operatorname{Ind}\tau_E=\rho_E$.
For a nonzero intertwiner $i:V_\tau\to V_\rho$ and a fixed
positive invariant coefficient form $h_\tau$, Schur's lemma gives
\[
                       i^*h_\rho=c h_\tau,\qquad c\in F_{>0}.
 \tag{PD.5}
\]
Thus the entire $\tau$-multiplicity form inside $H_1(C,E)$ is
$c b_\rho$.  Jacobian pullback $\pi^*:J(Y)\to J(C)$ induces
the homological transfer to $N$-invariants.  It multiplies the
intersection pairing by $|N|$.  Hence, after the actual
multiplicity-space identification,
\[
                        b_{\tau,Y}=\frac{c}{|N|}\,b_\rho.
 \tag{PD.6}
\]
Changing a coefficient form, choosing another self-adjoint
primitive coefficient line, or changing the anti-invariant
normalization multiplies this form by an element of $F^*$.
Every such rank-six determinant factor is a norm:
\[
                 s^6=\operatorname{Nm}_{E/F}(s^3)
                 \qquad(s\in F^*).
 \tag{PD.7}
\]
This proves the assertion.  In particular, no assumption that
$p$ is prime to the original group order, and no generic
endomorphism calculation, is hidden in this comparison.
\end{proof}

\paragraph{The weaker finite-catalogue form.}
Suppose a finite set of inducing topological entries exhausts a
chosen source class.  Their multiplicity pairings give finitely
many fixed determinant classes.  Pass to a common cyclotomic
field containing their coefficient fields and $K$, and choose
representatives.  A prime congruent to $-1$ modulo its conductor,
outside the finite supports of those representatives, has
Frobenius equal to complex conjugation.  It splits completely
in the real fixed field and is inert in the CM quadratic
extension at every place above it.  Every representative has
valuation zero, whereas the target class $p$ has valuation one;
every local norm has even valuation.  HD-PD-01 excludes this
target class for all entries at once.  Dirichlet's theorem gives
infinitely many such primes \cite[Section 1, p.~63]{PD:Harrison2009}.
This is conditional on the stated exhaustiveness hypothesis.
The next theorem removes the representative list and support
calculation from the bounded-order case.

\subsection*{HD-PD-02. The integral obstruction depends only on the inducing order}

\begin{theorem}[HD-PD-02]
For any integer $Q\ge1$, put
\[
            M_Q=\operatorname{lcm}(8,|\operatorname{disc}K|,
                                      1,2,\ldots,Q).
 \tag{PD.8}
\]
For every prime $p\equiv-1\pmod {M_Q}$ and every naturally
polarized primitive sixfold in HD-PD-01 with $|J|\le Q$,
\[
                             v_p(\delta_B)\equiv0\pmod2.
 \tag{PD.9}
\]
Consequently the signed target class $p$ is impossible.
The original group, induction index, curve genus, quotient
genus and branch ramification are unrestricted.

There are infinitely many such pairwise distinct target
determinant classes.  A terminating prime search constructs
them from $K,Q$ without GRH or enumeration of branch tuples.
The exclusion also applies under $K$-isogeny to a target with
$\operatorname{End}^0=K$, even when its specified polarization
was not initially required to match the natural source polarization.
\end{theorem}

\begin{proof}
Set $M=M_Q$ and $E=\mathbf Q(\zeta_M)$.
The field contains $K$ and splits every $J$ of order at most $Q$,
since it contains all roots of unity of exponent dividing $|J|$.
This is the usual cyclotomic splitting-field consequence of
Brauer induction; see
\cite[Corollary 4.6.12 and its proof]{OG:Zheng2025}
and the previously checked Brauer input.

For $p\equiv-1\pmod M$, cyclotomic Frobenius is complex
conjugation.  Thus every completion $E_w/\mathbf Q_p$ is
unramified quadratic, its fixed field is $\mathbf Q_p$, and
$K_p=E_w$.  These Frobenius and tower assertions are
\cite[8.15--8.18, pp.~142--143]{PD:MilneANT}.
In particular $p$ is odd, $p\nmid |J|$, and
$\eta=\sqrt{\operatorname{disc}K}$ is an anti-invariant $p$-unit.
Let $\mathcal O$ be the integers of $E_w$.
The local field $E_w$ contains all $\exp(J)$-th roots of unity
and splits $J$.  No assertion that it splits $G$ is needed.

\paragraph{The integral group algebra is a product of matrix orders.}
For every irreducible $E_w[J]$-module $V_\alpha$, choose a
$J$-stable $\mathcal O$-lattice $\Lambda_\alpha$.  For instance,
take the $\mathcal O$-span of the finite $J$-orbit of a basis.
Under the split group-algebra isomorphism set
\[
 \mathcal A'=\prod_\alpha
       \operatorname{End}_{\mathcal O}(\Lambda_\alpha)
       \ \supseteq\ \mathcal O[J].
 \tag{PD.10}
\]
Consider the regular trace pairing
$T(x,y)=\operatorname{Tr}(L_{xy}|E_w[J])$.
On the group basis,
\[
             T(g,h)=|J|\,\mathbf 1_{\{gh=1\}}.
 \tag{PD.11}
\]
Since $|J|$ is a unit, $\mathcal O[J]$ is its own trace-dual.
For $x\in\mathcal A'$ and $y\in\mathcal O[J]$,
\[
 T(x,y)=\sum_\alpha \dim(V_\alpha)
                       \operatorname{tr}(x_\alpha y_\alpha)
                       \ \in\ \mathcal O.
 \tag{PD.12}
\]
Hence $x$ belongs to the trace-dual of $\mathcal O[J]$.
This proves
\[
                  \mathcal O[J]=\mathcal A'.
 \tag{PD.13}
\]
In particular the integral matrix-order assertion has been
proved directly, without assuming a maximal-order classification.

\paragraph{A unimodular coefficient form.}
The involution
$*:\sum a_g g\mapsto\sum\bar a_g g^{-1}$
fixes every primitive character central idempotent, because
$\overline{\chi_\alpha(g)}=\chi_\alpha(g^{-1})$.
It preserves the matrix orders in (PD.13).
Choose an invariant nondegenerate Hermitian coefficient form
$h_\alpha$; the positive averaged form of a global cyclotomic
model supplies one.  Its adjoint involution is $*$.
Consequently
\[
 \operatorname{End}_{\mathcal O}(\Lambda_\alpha^\#)
       =\operatorname{End}_{\mathcal O}(\Lambda_\alpha).
 \tag{PD.14}
\]
Here $\#$ denotes the dual lattice for $h_\alpha$.
Two full lattices with the same full matrix endomorphism
order are homothetic: in a basis of the first, the diagonal
matrix units split the second into coordinate fractional
ideals, and the off-diagonal units force all ideals equal.
Since $\mathcal O$ is a discrete valuation ring,
$\Lambda_\alpha^\#=p^k\Lambda_\alpha$ for an integer $k$.
Rescaling $h_\alpha$ by $p^k\in\mathbf Q_p$ makes
$\Lambda_\alpha$ self-dual.

\paragraph{A unimodular multiplicity form.}
The module
$U=H_1(Y,\mathbf Z_p)\otimes_{\mathbf Z_p}\mathcal O$
has the perfect Hermitian pairing
\[
                    H(u,v)=\eta^{-1}I_Y(u,\bar v).
 \tag{PD.15}
\]
The integral intersection form $I_Y$ is unimodular and
$\eta$ is a unit.  Deck transformations act unitarily.
Every self-adjoint central summand $U_\alpha=e_\alpha U$
is therefore perfect.  Morita decomposition for the full
matrix order gives
\[
 U_\alpha=\Lambda_\alpha\otimes_{\mathcal O}W_\alpha,
          \qquad H_\alpha=h_\alpha\otimes b_\alpha.
 \tag{PD.16}
\]
The second equality follows from the adjoint compatibility
with the full matrix algebra.  The modules $W_\alpha$ are
free over $\mathcal O$.

For the relevant $\tau$ constituent, Frobenius reciprocity
and the dimension of $B$ give $\operatorname{rank}W_\tau=6$.
The form $b_\tau$ is integral: a unimodular matrix of
$h_\tau$ has some unit entry, and each coefficient of $b_\tau$
is obtained from the corresponding integral coefficient of
$H_\tau$ by division by that entry.  Moreover
\[
        \det H_\tau=(\det h_\tau)^6
                            (\det b_\tau)^{\dim\tau}.
 \tag{PD.17}
\]
Both $H_\tau$ and $h_\tau$ are perfect, so $\det b_\tau$
is a unit.  Returning to any other coefficient normalization
or primitive coefficient line changes the rank-six determinant
by a sixth power of an element fixed by bar, and therefore
changes its valuation by an even integer.
Thus the intrinsic $\tau$ determinant norm class has even
valuation at $p$.

HD-PD-01 compares this class with $\delta_B$.
The transfer factor $|N|$ and all coefficient or projector
scales contribute only sixth powers.  This remains true when
$p$ divides $|N|$ or $|G|$.
Thus (PD.9) follows.  Since the norm of an element of the
unramified quadratic extension $K_p/\mathbf Q_p$ has even
valuation, the class $p$ cannot be the signed determinant
class of $B$.

\paragraph{Target existence and arbitrary polarizations.}
The $K$-Hermitian matrix
\[
                    \operatorname{diag}(1,1,1,-1,-1,-p)
 \tag{PD.18}
\]
has signature $(3,3)$ and signed determinant $p$.
A compatible complex structure and a suitable positive
integral multiple of its alternating trace form on an
$\mathcal O_K$-lattice give a polarized abelian variety.
The associated unitary period domain has dimension nine.
The previously checked full-domain centralizer argument
gives $\operatorname{End}^0=K$ for its very general member.

More generally suppose a target with endomorphism algebra
$K$ is $K$-isogenous to a natural source $B$.
Then $\operatorname{End}^0(B)=K$.
Relative to the natural polarization, the pullback of the
target polarization is a positive Rosati-symmetric
endomorphism, hence a positive element of $K^+=\mathbf Q$.
Equivalently, compare the two alternating Hodge pairings:
their ratio is a Hodge endomorphism in $K$, and alternation
makes it fixed by conjugation.  Its rank-six determinant
scale is a norm.  Thus the prescribed target polarization
has the same signed determinant class as the natural source.
This proves the stated isogeny exclusion.  No assertion is
made that a special target with larger endomorphism algebra
has only one polarization determinant type.

\paragraph{Infinitude and an effective construction.}
Dirichlet's theorem gives infinitely many primes
$p\equiv-1\pmod M$
\cite[Section 1, p.~63]{PD:Harrison2009}.
They are inert in $K$ and their norm classes are pairwise
distinct, as seen from the odd valuation at their respective
inert primes.  Compute $M$, enumerate $p=kM-1$ above any
specified bound, and test primality until one is found.
Each primality test is finite, and Dirichlet proves termination.
No GRH, numerical prime estimate, branch-tuple enumeration,
Hodge-locus decision or bound on a minimal original index
enters this algorithm.
\end{proof}

\paragraph{The resulting necessary gate.}
For every $Q$ there are infinitely many determinant types
that have no source with a faithful inducing image of order
at most $Q$.  Therefore any exact-$K$ finite-group primitive-factor
mechanism supplying every determinant type would require
unbounded minimal inducing-image orders across those types.
If a separately proved structural class has a bound $Q$,
HD-PD-02 excludes these types from that class immediately.
The current application does not assume that the
factorwise tensor-decomposable interface in BI19 has been
resolved, or that the bounded structural class is exhaustive.

\paragraph{An explicit warning about proper-field determinant fibers.}
The restriction map on determinant norm classes need not be
injective for an even relative field degree.  For
$K=\mathbf Q(i)$, $E=\mathbf Q(i,\sqrt2)$ and
$F=\mathbf Q(\sqrt2)$,
\[
                 3=\operatorname{Nm}_{E/F}(1+i\sqrt2),
 \tag{PD.19}
\]
whereas $3$ is not a $K$-norm because it has odd valuation at
the inert prime $3$.  The selected primes in HD-PD-02 avoid
this issue: at those primes $E_w=K_p$, so the local obstruction
survives every relevant coefficient-field extension.

\section{An explicit order bound in the qualified transitive-core gate}
\label{sec:explicit-inducing-order}

\paragraph{Scope and status.}
HD-FQ01--03 is a new explicit quantitative consequence of the accepted
family reductions, HD-TW03 and HD-TC01--03.  The complete author and root whole-proof readings pass. Root also read
the entire author checker, its first complete output, and the entire
source-and-status record; no fresh root execution or fresh full
reading of all previously accepted source interfaces is claimed.
This explicit quantitative synthesis is recorded as derived work
for the formal movement rule.
Fix one of the eight non-Eisenstein class-number-one imaginary quadratic
fields
\[
 K\in\{\mathbf Q(i),\mathbf Q(\sqrt{-2}),\mathbf Q(\sqrt{-7}),
 \mathbf Q(\sqrt{-11}),\mathbf Q(\sqrt{-19}),
 \mathbf Q(\sqrt{-43}),\mathbf Q(\sqrt{-67}),
 \mathbf Q(\sqrt{-163})\}.
 \tag{FQ.1}
\]
Consider an actual finite-cover Weil sixfold source of exact original
character field $K$, period rank nine, and nonsplit actual polarization.
Let $\tau$ be an absolutely irreducible inducing block of degree $e\ge10$,
and let $J$ be its finite faithful matrix image.  Assume $J$ is primitive
and tensor indecomposable in the \emph{factorwise} sense of D-BI-19.
This convention permits permutations of tensor factors.  It does not
permit a nontrivial tensor decomposition preserved factorwise by $J$.
No equality of the original and inducing character fields is assumed.

These are the precise hypotheses throughout this section.  In particular,
no result here bounds inducing blocks of degree below ten, covers the
remaining factorwise tensor-decomposable case, or bounds an arbitrarily
larger original group from which $\tau$ is induced.

\subsection*{HD-FQ01. A common simple-quotient order ceiling}

Define the following positive integers, without expanding their decimal
representations:
\[
 \begin{split}
 M_*&=1+241\cdot240!\cdot2^{239},\qquad F_*=240M_*+2,\\
 \mathcal N&=[100(1+M_*)F_*^{240}]^2,\qquad D_0=10^{62},\\
 S_0&=\max\{\mathcal N!,\ 27^{3D_0^2},\ 10^{60}\}.
 \end{split}                                                   \tag{FQ.2}
\]

\begin{lemma}[HD-FQ01]
Every simple quotient of a quasisimple factor of the D-BI-19 normal core
has order at most $S_0$.
\end{lemma}

\begin{proof}
D-BI-19 supplies a normal irreducibly acting central product
$E_1\cdots E_t$, with the quasisimple factors conjugate under $J$.
Write $S=E_i/Z(E_i)$ for their common nonabelian simple quotient.
We use the finite simple group classification in exactly its usual
alternating, Lie-type and sporadic form; a complete statement with the
classical and exceptional families is given in
\cite[\S2.2.1, author p.~13]{HES:LiXia2022}.
This proof does not reconstruct that classification.

If $S$ is alternating, the ordinary representations with $n\ge\mathcal N$
are excluded by HD-TW03 in the present non-Eisenstein nonsplit scope.
The genuine spin alternatives for $n\ge8$ are excluded by HD-SPIN-02;
the exceptional covers in smaller alternating degrees do not affect
the inequality $n<\mathcal N$.  Thus
\[
                         |S|=n!/2\le\mathcal N!.
                                                               \tag{FQ.3}
\]
The spin and ordinary alternatives here concern the actual represented
central quotient; no ordinary-character assumption is imposed on $E_i$.

If $S$ is not alternating and is of Lie type over $\mathbf F_q$,
D-BI-20 applies under the present hypotheses and gives $q\le27$.
The accepted large-rank bounds then give the following deliberately
coarse ceilings for its classical natural dimension $d$:
\[
\begin{array}{c|c|l}
\text{family}&\text{ceiling sufficient here}&\text{accepted interface}\\ \hline
\operatorname{PSL}_{d}(q)&d<10^{61}&\text{D-LNW-01 and D-BI-23}\\
\operatorname{PSU}_{d}(q)&d<10^{61}&\text{D-UNW-03 and D-BI-24}\\
\operatorname{PSp}_{d}(q)&d\le6&\text{D-HES-01--06 and D-HOS-03--22}\\
\text{orthogonal, }q\text{ even}&d\le1134&\text{D-OL-11}\\
\text{orthogonal, }q\text{ odd}&d\le1134&\text{D-OL-19--24}.
\end{array}                                                    \tag{FQ.4}
\]
The linear and unitary statements include their level-one alternatives,
as specified in the cited interfaces.  The symplectic statements cover
all represented central kernels in the large ranks; the exceptional
small-rank multipliers are left inside the displayed finite range.
In even characteristic the large-rank orthogonal simply connected cover
has trivial center and no exceptional multiplier, as in the multiplier
scope already used in D-OL-11.  In odd characteristic D-OL-19--24 include
both ordinary and genuinely spin representations and arbitrary
represented central kernels.  In particular the withdrawn stronger
threshold based on an incorrectly read second square root in LST is
\emph{not} an input.  Every entry in (FQ.4) is less than $D_0$.

For these classical groups their defining special linear, unitary,
symplectic or orthogonal matrix group surjects onto the simple group.
Its natural degree is $d$, and its matrix coefficients lie in
$\mathbf F_q$ or $\mathbf F_{q^2}$.  Counting all matrices gives
$|S|<q^{2d^2}$, which is less than $27^{3D_0^2}$.
For clarity, this estimate does not require the exact classical order
formula or a faithful linear model for the simple quotient itself.

For exceptional groups use the published matrix constructions in
\cite[\S3, author pp.~4--12; \S4, pp.~12--17]{FQ:HRT2001}.
The untwisted degrees are $7,26,27,56,248$ for
$G_2,F_4,E_6,E_7,E_8$.  The twisted constructions have degrees
$4,8,27,26,7$ for
${}^2B_2,{}^3D_4,{}^2E_6,{}^2F_4,{}^2G_2$, over extensions of
$\mathbf F_q$ of degree at most three.  Their matrix groups map
onto the simple group, or contain its Tits derived subgroup in the
stated ${}^2F_4(2)$ case.  Hence every exceptional simple group here has
\[
                        |S|<q^{3\cdot248^2}
                                 \le27^{3D_0^2}.
                                                               \tag{FQ.5}
\]
The quoted constructions include the bad-characteristic parameters;
irreducibility of the displayed matrix module is unnecessary for
this order estimate.  The defining-field convention is $q$, not the
larger matrix coefficient field.

Finally all twenty-six sporadic orders are below $10^{60}$.
The classification list and the range up to the Monster, whose order
is less than $10^{54}$, are explicitly recorded in
\cite[Lecture~1, p.~1]{FQ:Wilson2008}; the exact Monster order is also
listed in the primary ATLAS data \cite{FQ:AtlasMonster}.
This harmless envelope does not require reopening the earlier stronger
sporadic exclusion.  Combining these cases proves the lemma.
\end{proof}

\subsection*{HD-FQ02. Bounding every quasisimple central cover}

\begin{lemma}[HD-FQ02]
For a finite group $S$ of order $s$, its integral Schur multiplier satisfies
\[
                  |H_2(S,\mathbf Z)|\le4^{(s-1)^2}.
                                                               \tag{FQ.6}
\]
Consequently, if $S$ is nonabelian simple and $E\to S$ is any
quasisimple central extension, then
\[
                         |E|\le s\,4^{(s-1)^2}.
                                                               \tag{FQ.7}
\]
\end{lemma}

\begin{proof}
Use the normalized integral bar resolution, followed by trivial
coinvariants.  Its free abelian groups in degrees two and three have
bases $[g|h]$ and $[g|h|k]$ with each entry nonidentity; any symbol with
an identity entry is interpreted as zero.  The differential is
\[
 d_3[g|h|k]=[h|k]-[gh|k]+[g|hk]-[g|h].
                                                               \tag{FQ.8}
\]
The complete resolution and its use to compute integral group homology
are given in \cite[Definition~6.5.1, Theorem~6.5.3 and
Application~6.5.4, pp.~177--179]{FQ:Weibel1994}.
Each column of this matrix is a sum of four signed coordinate vectors,
so its Euclidean norm is at most four even when some terms coincide.
Its row count is $(s-1)^2$.

The chain-complex identity gives the exact sequence
\[
 0\longrightarrow H_2(S,\mathbf Z)
  \longrightarrow\operatorname{coker}(d_3)
  \longrightarrow\operatorname{im}(d_2)\longrightarrow0.
                                                               \tag{FQ.9}
\]
The last group is free abelian, being a subgroup of a free abelian
group.  The first is finite: integral positive-degree homology of a
finite group is finitely generated and killed by $s$, by
\cite[Theorem~6.5.8 and Corollaries~6.5.9--6.5.10, p.~180]{FQ:Weibel1994}.
It follows that $H_2(S,\mathbf Z)$ is exactly the torsion subgroup of
$\operatorname{coker}(d_3)$.

Let $r$ be the rational rank of $d_3$.  Smith normal form says that
the order of that torsion subgroup is the greatest common divisor
of the nonzero $r$-rowed minors.  Choose any one such minor.
Deleting rows cannot increase column norm, so Hadamard's determinant
inequality bounds its absolute determinant by $4^r$.
Since $r\le(s-1)^2$, this proves (FQ.6).
For $r=0$ the torsion group is trivial and the same conclusion holds.

A perfect group has a universal central extension with kernel
$H_2(S,\mathbf Z)$; see the complete construction and proof in
\cite[\S6.9, Theorem~6.9.5, pp.~198--201]{FQ:Weibel1994}.
Its map to any central extension $E\to S$ is surjective when $E$
is perfect: its image $U$ maps onto $S$, so $E=U Z$ for the
central kernel $Z$, and then $E=[E,E]=[U,U]\subset U$.
Every quasisimple $E$ is therefore a quotient of that universal
extension.  Equation (FQ.7) follows from (FQ.6).
No classification of Schur multipliers is used in this new estimate.
\end{proof}

\subsection*{HD-FQ03. Explicit inducing order and its arithmetic consequence}

Set
\[
 \begin{split}
 A_0&=S_0\,4^{S_0^2},\qquad T_0=2A_0,\qquad E_0=A_0^{T_0},\\
 B_0&=S_0!,\qquad
 Q_0=E_0\bigl(B_0^{T_0}T_0!\bigr)^2.
 \end{split}                                                   \tag{FQ.10}
\]
All these are explicitly specified integers; no computational
enumeration up to them is asserted.

\begin{theorem}[HD-FQ03]
Every source under the initial hypotheses has
\[
                 t\le T_0,\qquad e\le E_0,\qquad |J|\le Q_0.
                                                               \tag{FQ.11}
\]
Its connected inducing $J$-cover $Y\to X=\mathbf P^1$ has
\[
                  r\le2E_0+6,\qquad
                  g(Y)\le1+Q_0(E_0+2).
                                                               \tag{FQ.12}
\]
For each fixed $K$ in (FQ.1), set
\[
                M_0=\operatorname{lcm}(8,|\operatorname{disc}K|,
                                      1,2,\ldots,Q_0).
                                                               \tag{FQ.13}
\]
Every prime $p\equiv-1\pmod{M_0}$ is excluded as the signed
determinant class of a naturally polarized source in this qualified
class.  The same exclusion holds for a target with
$\operatorname{End}^0=K$ which is $K$-isogenous to such a source.
There are infinitely many pairwise distinct excluded norm classes.
\end{theorem}

\begin{proof}
By HD-FQ01--02 each actual quasisimple factor has order at most $A_0$.
An irreducible complex representation of a finite group has degree
at most its group order, so its factor degree satisfies $2\le a\le A_0$.
D-BI-19 gives equal factor degrees and $e=a^t$.
HD-TC02 gives
\[
 t\le\max\{4,\lceil5a/4\rceil-1\}\le2A_0=T_0,
                 \qquad e\le A_0^{T_0}=E_0.
                                                               \tag{FQ.14}
\]
An automorphism of a finite group of order $s$ permutes its underlying
set faithfully, so $|\operatorname{Aut}(S)|\le s!\le B_0$.
For $P=S^t$, the simple direct factors are the minimal normal subgroups;
therefore every automorphism permutes them and restricts to an
automorphism on each.  Thus
\[
 |\operatorname{Aut}(P)|=|\operatorname{Aut}(S)|^t t!
                                  \le B_0^{T_0}T_0!.
                                                               \tag{FQ.15}
\]
HD-TC03, which retains the entire scalar kernel, now yields
$|J|\le e|\operatorname{Aut}(P)|^2\le Q_0$.
That theorem uses the actual generating inertia tuple: all but possibly
one inertia have a fixed vector, their projective powers control their
determinants, and product one controls the last determinant as well.
Consequently the estimate does not assume a bounded scalar group
before using the source geometry.
The same source Euler identity and Galois Riemann--Hurwitz argument
in HD-TC03 give (FQ.12).

Apply HD-PD02 with its bounded inducing order $Q=Q_0$.
For every prime in (FQ.13), that theorem proves that the natural
rank-six signed Hermitian determinant has even $p$-valuation.
The local proof allows $p$ to divide the order of the original group
or of the discarded kernel; only $p\nmid|J|$ is needed, and the
congruence implies it.  Rank-six polarization rescaling contributes
a sixth power, which is a norm, so the target comparison under
$\operatorname{End}^0=K$ is exactly the one stated in HD-PD02.
The signed class $p$ has odd valuation and is therefore impossible.
Dirichlet's theorem gives infinitely many such primes.  They are inert
in $K$, and distinct such primes give different classes in
$\mathbf Q^\times/\operatorname{Nm}_{K/\mathbf Q}(K^\times)$,
since norms have even valuation at inert primes.
These are actual signature-$(3,3)$ target types, represented by
$\operatorname{diag}(1,1,1,-1,-1,-p)$.
\end{proof}

\paragraph{What the explicit bound decides.}
The multiplication-table, character, tuple and complex period-rank
enumeration in the refined HD-TC03 is now finite for the qualified
class above, with the displayed uniform integer ceilings.
Its original-$K$ descent condition remains separate, and no such
enormous enumeration has been performed.
D-OG01--02 exhibit why this cannot be promoted to a bound on the
original group or original induction index: arbitrary wreath padding
can preserve an actual primitive sixfold factor and its determinant.
The arithmetic conclusion excludes this specified source route for
infinitely many target determinant classes; it is not a claim that
the associated Hodge classes are nonalgebraic.
The remaining factorwise tensor-decomposable case and degrees below
ten are outside this theorem, so it supplies no unconditional
all-supplier obstruction.

\section{Small inducing degrees and a projective scalar bound}
\label{sec:low-degree-projective-order}

\paragraph{Status and scope.}
These are root statements HD-LQ01--02, reserved question Q2dd.
The repaired complete proof has passed root and separate independent
whole-proof readings, including the actual inducing-base and proper-field
interfaces. The initial draft, scope repair and exact source audits are retained. The projective scalar step extends the elementary
argument of HD-TC03 without assuming a perfect normal core.
All branch entries in this section are genuine nonidentity entries
of the actual faithful inducing image. Original overgroups are not
being bounded.

\subsection*{HD-LQ01. The scalar bound needs only the projective image}

\begin{lemma}[HD-LQ01]
Let a finite subgroup \(J\le\operatorname{GL}(V)\) act irreducibly,
with \(e=\dim V\). Suppose \(T_1,\ldots,T_r\) generate \(J\),
their product is one, all are nonidentity, and
\[
 r\ge9,\qquad
 D:=\sum_i\operatorname{rank}(1-T_i)=2e+6.
 \tag{LQ.1}
\]
Put \(P=J/(J\cap\mathbb C^\times I)\) and \(s=|P|\).
Then
\[
 |J\cap\mathbb C^\times I|\mid es,
 \qquad |J|\le es^2.
 \tag{LQ.2}
\]
No primitivity, perfect-core, character-field or tensor hypothesis
is needed for this assertion.
\end{lemma}

\begin{proof}
If two entries had full moving rank, the other \(r-2\ge7\)
genuine entries would give
\(D\ge2e+7\), contradicting (LQ.1).
Thus every entry except possibly one has a fixed vector.
The projective order of every \(T_i\) divides \(s\), so
\(T_i^s\) is scalar. If \(T_i\) has a fixed vector, this scalar
is one and \(\det(T_i)^s=1\).
Taking determinants in the product-one relation gives the same
identity for the possible exceptional entry. Since the tuple
generates \(J\), its determinant image is contained in \(\mu_s\).
For any scalar \(zI\in J\), we consequently have \(z^{es}=1\).
The scalar subgroup is finite cyclic, and its order divides \(es\).
Multiplication by the projective order \(s\) proves (LQ.2).
\end{proof}

\subsection*{HD-LQ02. A uniform order ceiling in inducing degrees at most nine}

Set
\[
 B_{\rm low}=10^9,\qquad
 Q_{\rm low}=9B_{\rm low}^2=9000000000000000000.
 \tag{LQ.3}
\]

\begin{theorem}[HD-LQ02]
Let \(J\le\operatorname{GL}_e(\mathbb C)\) be the faithful
primitive inducing image of an actual rank-nine sixfold source
over the sphere, with \(e\le9\).
Then
\[
 e\ge2,\qquad |J|\le Q_{\rm low},\qquad
 r\le24,\qquad g(Y)\le1+11Q_{\rm low}
 \tag{LQ.4}
\]
for the connected inducing \(J\)-cover \(Y\to\mathbb P^1\).
The estimate includes factorwise tensor-decomposable primitive
images in this degree range and arbitrary inducing character fields.

Consequently, in HD-FQ01--03 the same integer \(Q_0\) bounds the
faithful inducing image in \emph{every} inducing degree, once the
other FQ hypotheses, including primitive factorwise tensor
indecomposability, are retained. The prime determinant exclusions
therefore apply to that entire qualified class without a lower
degree cutoff.
\end{theorem}

\begin{proof}
The actual inducing sphere source has period rank nine and therefore
at least twelve genuine branch points. This is the same source
dimension interface used in HD-TC03: a family of covers of a sphere
with \(r\) genuine marks has dimension at most \(r-3\), and its
period image cannot have greater dimension. Its Chevalley--Weil
Euler identity is \(D=2e+6\).
These are properties of the actual inducing cover, not an assignment
of twelve marks to an individual tensor factor.
If \(e=1\), the identity gives \(D=8\), whereas the twelve genuine
entries have total moving rank at least twelve. Hence \(e\ge2\).

Collins' published Theorem A gives an explicit upper bound for
\([J:Z(J)]\) for every finite primitive complex linear group.
Every entry in its degree \(2\) through \(9\) table is less than
\(10^9\); see \cite[Theorem A, p.~761]{LQ:Collins2008}.
We use this coarse envelope rather than require an optimal table
constant. Irreducibility and Schur's lemma identify \(Z(J)\)
with the scalar subgroup, so \(s=|P|\le B_{\rm low}\).
HD-LQ01 gives \(|J|\le e s^2\le9B_{\rm low}^2\).

Each genuine entry contributes at least one, so
\(r\le2e+6\le24\). If its inertia order is \(o_i\), Galois
Riemann--Hurwitz gives
\[
 2g(Y)-2=|J|\left(-2+\sum_i(1-1/o_i)\right)
       \le |J|(r-2)\le22Q_{\rm low},
 \tag{LQ.5}
\]
which proves (LQ.4).

For the all-degree FQ conclusion, first dispose of degree one
before assuming that the inducing base is a sphere.
D-MON03 excludes every monomial inducing source in the nonsplit
class under consideration. Its proof explicitly permits a proper
inducing character field \(L/K\): the actual rational isotypic
factor is \(B^{[L:K]}\), and the full generic endomorphism algebra
of the dominating cyclic genus-four component forces \(L=K\),
after which its natural polarization is split.
Thus the potentially unramified genus-four case in degree one
does not enter this nonsplit conclusion.

For \(e\ge2\), let \(h\) be the genus of the actual inducing base.
The already checked actual-cover interface D-BI-18, equations
(BI.65)--(BI.66), gives
\[
 D=2e+6-2eh,\qquad r\le D.
 \tag{LQ.6}
\]
If \(h=1\), then \(D=6\) and the pointed-curve moduli have
dimension at most six. If \(h=2\), their dimension is at most
\(3+D=9-2e\le5\). If \(h\ge3\), then \(D<0\), impossible.
The multiplicity period map factors through this actual inducing
Hurwitz space, so rank nine forces \(h=0\).
This argument uses neither equality of the original and inducing
fields nor a mark-count premise on an individual tensor factor.

In the FQ notation, \(S_0\ge10^{60}\), hence
\(A_0\ge10^{60}\), \(E_0\ge A_0\) and
\(Q_0\ge E_0>Q_{\rm low}\).
HD-FQ03 already bounds the qualified inducing images for
\(e\ge10\); (LQ.4) handles \(2\le e\le9\), with no additional
structural restriction required in these small degrees.
Together with the degree-one exclusion, the unchanged FQ bound
therefore holds in all degrees.
Applying HD-PD02 with \(Q=Q_0\) proves the unchanged prime
determinant exclusions. The extension to an arbitrary specified
polarization on a \(K\)-isogenous target still requires
\(\operatorname{End}^0=K\), exactly as in HD-PD02.
\end{proof}

\paragraph{Exact input and remaining gate.}
Collins' result is a published classification-based input, not a new
classification in this run. Root read its full Theorem A statement
and table, the definitions immediately before it, and the complete
final proof in Section5, pp.~771--774. The preceding group-theoretic
and representation estimates in Sections2--4 are cited through that
published theorem; a fresh independent reconstruction of them is not
claimed. The primary Oxford deposit is the published typeset PDF.

This removes the small-degree qualification from the FQ source-route
restriction. It does not prove that any surviving small-degree
Nielsen component has the required field descent, period rank or
detector cycles. High-degree primitive factorwise tensor-decomposable
images remain outside the all-degree FQ conclusion. No bound on the
original overgroup or arbitrary induction index is inferred.

The geometric objects remain actual finite covers, their weight-one
local systems and rational deck-correspondence projectors. The result
limits suppliers for the integral uniform condition A (cote16,
scan050; SIGN\_NEXT SN-C03/C04). Condition I (scan049; cote44,
Theorem4), the Cayley--Hamilton inverse of scans048--049 and the
existing EGA/SGA/FGA spreading interfaces do not create missing
cycles. No Markman, Severn or032 premise is introduced.

\section{Finite projected-factor catalogues and a large-field tensor obstruction}
\label{sec:tensor-factor-catalogue}

\paragraph{Scope and status.}
D-TD01--03 concern an actual factorwise tensor decomposition.
They do not transfer the twelve-mark condition to a projected factor.
The finite catalogue statement is conditional only on bounds for the
projective factor images; the large-field statement below is an
unconditional exclusion under its explicit normal-core hypotheses.
These are new \texttt{/root/imprimitive} derivations, discussed with
root. Complete author and separate root whole-proof readings pass;
root also read the full source/status record, checker, output and
invoked primary DHP and Milne passages. Exact boundaries are recorded
in the retained review. No global exhaustion of all factor types is asserted.

Let \(J\subset\operatorname{GL}(V)\) be a finite faithful irreducible
matrix group of degree \(e\ge2\).  Suppose an actual connected sphere
cover has a genuine generating product-one inertia tuple
\(g_1,\ldots,g_r\in J\), satisfying
\[
 r\ge12,\qquad
 \sum_{i=1}^r\operatorname{rank}(1-g_i)=2e+6 .
 \tag{TD.1}
\]
For the tensor catalogue assume that \(J\) preserves factorwise
\[
                    V=W_1\otimes\cdots\otimes W_t,
             \qquad a_j=\dim W_j\ge2 .
                    \tag{TD.2}
\]
The resulting projective factor homomorphisms are denoted
\(J\longrightarrow P_j\subset\operatorname{PGL}(W_j)\).
They are actual homomorphisms: the factors in a tensor decomposition
of an invertible operator are unique up to scalars.
Every projective factor action is irreducible, since an invariant
proper subspace in one factor would give one in \(V\).

\subsection{D-TD01: bounded projective factors bound the whole inducing image}

\begin{theorem}[D-TD01]
Assume \(a_j\le A\) and \(|P_j|\le B\) for every \(j\), where
\(A,B\ge2\) are integers.  Define
\[
 R=5A,\qquad
 N=B^{B^2+R+1},\qquad
 M=B^N,\qquad Q=M^3 .
 \tag{TD.3}
\]
Then
\[
            r\le R,\qquad
       |J/(J\cap\mathbf C^\times I)|\le M,\qquad |J|\le Q .
       \tag{TD.4}
\]
There is no bound on \(t\) among the hypotheses.
Consequently the actual faithful inducing images, their irreducible
modules, and generating branch tuples satisfying these bounds lie in
an effectively enumerable finite list.
\end{theorem}

\begin{proof}
For a finite-order operator \(T\), write
\(\nu(T)=\dim T-\max_z\dim\ker(T-z)\).
Tensor spectral convolution gives
\[
 \frac{\operatorname{rank}(1-g_i)}e
             \ge \max_j\frac{\nu(W_j(g_i))}{a_j}.
 \tag{TD.5}
\]
Representatives of the projective operators may be chosen arbitrarily;
their scalar phases do not affect \(\nu\).
If some projected operator is nonidentity, it has \(\nu\ge1\), so
the right side is at least \(1/A\).
If all are projectively trivial, \(g_i\) is scalar; because the tuple
is genuine, its actual moving rank is \(e\), again at least \(e/A\).
Thus (TD.1) gives
\[
                  r\le A(2+6/e)\le5A=R .
                  \tag{TD.6}
\]

Pad the tuple by identities to length \(R\), and let \(F_R\) be the
free group on these labelled generators.  For a group of order
\(b\le B\), fix a labelling by \(\{1,\ldots,b\}\).
There are at most \(b^{b^2}\) binary operation tables and at most
\(b^R\) assignments of the \(R\) generators.  These deliberately
overcount all group structures and all homomorphisms to them.
In particular there are at most
\[
          \sum_{b=1}^B b^{b^2+R}
                     \le B^{B^2+R+1}=N
                     \tag{TD.7}
\]
distinct kernels among all such homomorphisms.
Each actual \(F_R\longrightarrow P_j\) belongs to this list.
Repeating a kernel has no effect on their intersection.
Select one representative for every distinct kernel occurring in
(TD.2); at most \(N\) representatives are needed.

An invertible tensor product is scalar only if each factor is scalar.
For example, conjugating \(X\otimes I\otimes\cdots\otimes I\)
by a scalar tensor product shows that its first factor commutes
with every \(X\), and therefore is scalar; repeat for all factors.
It follows that the kernel of the combined map to the \(P_j\)
is exactly the scalar kernel in \(J\).  Hence, putting
\(s=|J/(J\cap\mathbf C^\times I)|\), the selected quotient maps give
\[
                              s\le B^N=M .
                              \tag{TD.8}
\]
This bounds the order, not a claimed divisibility by \(M\).

We next bound the scalar subgroup without assuming any field.
Irreducibility implies that the span of \(J\) is
\(\operatorname{End}_{\mathbf C}(V)\), by the group-algebra
decomposition \cite[\S6.2, Proposition10, p.~48]{OG:Serre1977}.
Choose one matrix representative for each projective element.
Every matrix of \(J\) is a scalar multiple of one of these \(s\)
representatives.  Therefore
\[
                                  e^2\le s.
                                  \tag{TD.9}
\]

At most one inertia has no eigenvalue \(1\).  Otherwise two entries
cost \(2e\), and the other \(r-2\ge10\) genuine entries each cost
at least one, contradicting (TD.1).
For every other entry, \(g_i^s\) is scalar by Lagrange's theorem
in the projective image.  Since \(g_i\) has eigenvalue \(1\),
that scalar is one.  Thus \(\det(g_i)^s=1\).
If there is an exceptional entry, the determinant of the
product-one relation gives the same conclusion for it.
The tuple generates \(J\), so \(\det(J)\subset\mu_s\).
For \(zI\in J\), this implies \(z^{es}=1\), whence
\[
 |J\cap\mathbf C^\times I|\le es,\qquad
 |J|\le es^2\le s^{5/2}\le M^3 .
 \tag{TD.10}
\]
This proves (TD.4).

For effective enumeration, enumerate the finitely many multiplication
tables of order at most \(Q\), test the group axioms, and construct
their irreducible representations over a cyclotomic splitting field.
The existence of such a field is the known Serre theorem already
used in the finite-rank algorithm
\cite[\S12.3, Theorem24, p.~94]{OG:Serre1977}.
Finite semisimple-algebra computations, or the exact algebraic
enumeration in D-NR, give an effective construction; no undecidable
choice is required.  Enumerate tuples of length at most \(R\)
and test generation and product one in the tables.
It is harmless to retain a superset instead of deciding every
tensor-decomposition condition.  This finite list contains every
source under the stated hypotheses.
\end{proof}

\paragraph{Order-only version.}
If only \(|P_j|\le B\) is given, irreducibility still implies
\(a_j^2\le |P_j|\).  Indeed the inverse image of \(P_j\) in
\(\operatorname{SL}(W_j)\) is finite, its representation is
irreducible, and its matrix span is generated by one representative
for each projective element, exactly as in (TD.9).
Thus one may take \(A=B\) in (TD.3), or the sharper integer
\(A=\lfloor\sqrt B\rfloor\) whenever the hypotheses are nonvacuous.
A bound on projective factor orders alone therefore suffices.

\paragraph{Scope of the finite decision.}
The list is one of inducing images and their actual tuple data.
Within the precise projector and period-map input scope of D-NR,
its exact generic-rank procedure is available for each entry.
No claim here decides whether a proper special subfamily has an
original-\(K\) realization or a prescribed polarization descent.
D-OG shows that arbitrary original overgroups can remain unbounded.
The necessary determinant filter HD-PD02 applies to the order
ceiling \(Q\), but is not promoted into a sufficient descent test.
In a primitive \(J\), each projective factor in (TD.2) is primitive:
an imprimitivity decomposition in one factor tensors to an
imprimitivity decomposition of \(V\).
Thus any independently proved projective-order bounds for such
factors may be inserted into (TD.3).

\subsection{D-TD02: a projected orbifold bound replacing twelve marks}

\begin{lemma}[D-TD02]
Suppose an actual quotient \(P\) of the source group is a finite
nonsolvable group with \(|P|>60\).  Suppose
\(V=W\otimes U\), where \(a=\dim W\), and the action on \(W\)
is a projective representation of \(P\).
Assume that every nonidentity \(x\in P\), in a unitary scalar
normalization, satisfies
\[
                         |\operatorname{tr}W(x)|/a\le c<1 .
                         \tag{TD.11}
\]
Then every actual sphere source with moving sum \(2e+6\) satisfies
\[
                               e(1-85c)\le252 .
                               \tag{TD.12}
\]
In particular \(c\le1/170\) forces \(e\le504\).
This lemma needs no lower bound on the number of genuine projected
marks, and places no restriction on \(U\).
\end{lemma}

\begin{proof}
Remove identity images from the projected generating tuple, and
write \(o_1,\ldots,o_k\ge2\) for the exact orders of the remaining
elements.  The associated actual connected \(P\)-Galois cover
\(Z\longrightarrow\mathbf P^1\) has
\[
 2g(Z)-2=|P|\left(-2+\sum_{i=1}^k(1-1/o_i)\right).
 \tag{TD.13}
\]
We check that \(g(Z)\ge2\).  If \(k\le2\), product one makes
the generated group cyclic, contrary to the hypothesis.
If \(g(Z)=0\), then \(k=3\), since every summand is at least
\(1/2\).  Ordering the three orders, the strict inequality
\(\sum1/o_i>1\) gives exactly
\((2,2,m),(2,3,3),(2,3,4),(2,3,5)\).
In the first case two involutions generate a dihedral quotient,
which is solvable.  In the other cases (TD.13) gives respectively
\(|P|=12,24,60\).  All are excluded.
If \(g(Z)=1\), choose an origin on \(Z\).  Its automorphism group
is the translation group semidirect its origin-fixing group;
over \(\mathbf C\) the latter is cyclic of order \(2,4\), or \(6\).
These exact known statements, with proofs, are
\cite[ChapterIV, Proposition7.13 and Remark7.18, pp.~140--141]{TD:MilneEC2021}.
Consequently every finite subgroup is abelian-by-cyclic and solvable,
again a contradiction.

Put \(\Delta=\sum_i(1-1/o_i)-2>0\).
The elementary hyperbolic order estimate is
\[
                         \Delta\ge1/42 .
                         \tag{TD.14}
\]
For completeness, if \(k\ge5\) it is at least \(1/2\).
For \(k=4\), the all-two case has zero excess and every positive
case has excess at least \(1/6\).
For \(k=3\), put \(u\le v\le w\).
If \(u\ge4\) the excess is at least \(1/4\).
If \(u=3\), the smallest positive case is \((3,3,4)\),
with excess \(1/12\).
If \(u=2\), the cases \(v=2\) are negative; for \(v=3\),
positivity requires \(w\ge7\) and the minimum is \(1/42\);
for \(v=4\) it is \(1/20\), and for \(v\ge5\) it is at least
\(1/10\).  This proves (TD.14) for every possible number of marks.

For an element \(x\in P\) of order \(o\), choose a scalar
normalization of its unitary representative \(T\) with \(T^o=1\).
Every power \(T^j\), \(1\le j<o\), is projectively nonidentity.
Fourier inversion and (TD.11) therefore give
\[
 \frac{\nu(T)}a
 \ge1-\frac{1+(o-1)c}{o}
                    =(1-c)(1-1/o).
 \tag{TD.15}
\]
Tensor convolution cannot decrease this normalized lower bound
on actual moving rank, whatever the partner operator and scalar.
Summing and using (TD.14) yields
\[
             2+\frac6e\ge(1-c)\frac{85}{42}.
             \tag{TD.16}
\]
Rearrangement gives (TD.12).  The specialization \(c\le1/170\)
has \(1-85c\ge1/2\), proving \(e\le504\).
\end{proof}

\paragraph{All-power precision.}
Condition (TD.11) is an absolute character bound for the complete
actual projective image, hence for every nontrivial power.
A support estimate only on a normal subgroup, or a commutator
support estimate with no trace control, cannot be substituted for it.

\subsection{D-TD03: full normalizers of large-field Lie cores}

\begin{lemma}[Normalizer trace averaging]
Let \(N\subset\operatorname{GL}(W)\) be finite, perfect, faithful,
and irreducible, with \(a=\dim W>1\), and put \(S=N/Z(N)\).
Assume the actual core character has
\[
                  |\chi_W(x)|/a\le c<1\qquad(x\notin Z(N)),
                  \tag{TD.17}
\]
and every nontrivial ordinary irreducible representation of \(S\)
has degree at least \(b\).
Then every nonscalar unitary operator \(T\) normalizing \(N\)
satisfies
\[
 \frac{|\operatorname{tr}T|^2}{a^2}
                   \le c+\frac{1-c}{b+1}.
                   \tag{TD.18}
\]
The operator may induce an outer automorphism or permute central
factors of \(N\).
\end{lemma}

\begin{proof}
Let \(\alpha\) be the induced automorphism of \(S\).
If \(\alpha=1\), then
\(TxT^{-1}=\lambda(x)x\) for all \(x\in N\), with
\(\lambda:N\longrightarrow\mathbf C^\times\) a homomorphism.
Perfectness gives \(\lambda=1\), and Schur's lemma makes \(T\)
scalar.  Thus \(\alpha\ne1\).
Its fixed subgroup \(H=C_S(\alpha)\) is proper.
The coset permutation module \(\mathbf C[S/H]\) contains the
constants once and has a nonzero complementary representation.
Its latter dimension is \([S:H]-1\), so
\[
                 [S:H]\ge b+1,\qquad |H|/|S|\le1/(b+1).
                 \tag{TD.19}
\]
For \(x\in N\), the commutator \(TxT^{-1}x^{-1}\) is central
exactly when \(xZ(N)\in H\).
Schur averaging gives the exact identity
\[
 \frac1{|N|}\sum_{x\in N}
       \frac{\operatorname{tr}(TxT^{-1}x^{-1})}{a}
                         =\frac{|\operatorname{tr}T|^2}{a^2}.
                         \tag{TD.20}
\]
Indeed the average of \(xT^{-1}x^{-1}\) is
\(\operatorname{tr}(T^{-1})I/a\).
The central commutators have absolute character ratio one,
including their possible nontrivial scalar phases.
Every other commutator has ratio at most \(c\).
The triangle inequality, (TD.19), and (TD.20) prove (TD.18).
\end{proof}

\begin{theorem}[D-TD03]
Suppose an actual factor \(W\) of a sphere-source coefficient
module \(V=W\otimes U\) has a normal irreducibly acting central product
\[
       N=E_1\cdots E_t\subset\operatorname{GL}(W),
       \qquad S_i=E_i/Z(E_i),
       \tag{TD.21}
\]
where each \(E_i\) is a finite quasisimple group of Lie type with
nonalternating simple quotient, in the defining-field convention of
Duyan--Halasi--Podoski Theorem4.1(2).
Assume every defining field size satisfies \(q_i\ge10^{12}\).
The actual factor image may be any finite group in the full linear
normalizer of \(N\), including compatible outer automorphisms,
factor permutations and arbitrary scalar extensions.
Then no such sphere source has moving sum \(2e+6\).
There is no factorwise tensor-indecomposability hypothesis, no
transitivity requirement on the \(E_i\), no original-field restriction,
and no projected twelve-mark premise.
\end{theorem}

\begin{proof}
The precise published character estimate is
\[
 {\rm mr}(E_i)\le
       \max\{1/(\sqrt{q_i}-1),\,9/q_i\}
       \le c_0:=1/999999 .
       \tag{TD.22}
\]
Definition2.3 and Theorem4.1(2) of \cite{BI:DHP}, pp.~137,141,
apply to all nontrivial irreducible complex characters outside
their scalar loci.  The theorem has no coprime-representation
hypothesis: the coprime hypotheses belong to other results of
that paper.  For a nontrivial representation of a quasisimple
group its scalar locus is its center, since its nonabelian simple
quotient has no proper nontrivial normal subgroup.
Thus (TD.22) applies with the stated exact scope.

Here is an elementary degree consequence which avoids any table
of minimal representation degrees or group orders.
For a nonabelian simple group \(S_i\), let \(\psi\ne1\) be an
ordinary irreducible character of degree \(d\).
If \(dc_0<1\), every value \(\psi(x)\), \(x\ne1\), has all its
Galois conjugates of modulus less than one: those conjugates are
again irreducible characters of degree \(d\), to which (TD.22)
applies.  Character values are algebraic integers, so their
integer norms force \(\psi(x)=0\).
Orthogonality would then imply \(d^2=|S_i|\), contrary to
\(\sum_{\eta\in{\rm Irr}(S_i)}\eta(1)^2=|S_i|\), which also
contains the trivial character.  Therefore
\[
                              d\ge1/c_0=999999 .
                              \tag{TD.23}
\]
The algebraic-integrality, Galois-conjugacy and orthogonality
inputs here are the standard finite-group character facts in
\cite[\S\S2.3--2.4,6.4--6.5,12.1]{OG:Serre1977}.

Now \(S=N/Z(N)=\prod_iS_i\).
Every nontrivial irreducible ordinary \(S\)-module is a tensor
product with at least one nontrivial factor, so (TD.23) gives
the same lower degree \(b=999999\) for \(S\).
We have not asserted a uniform nonidentity trace bound for every
\(S\)-character: that assertion would be false for characters
trivial on a direct factor.
In contrast, the \emph{actual} irreducible \(N\)-module \(W\)
is a tensor product of nontrivial irreducible \(E_i\)-modules.
Faithfulness excludes a trivial \(E_i\)-factor.  Hence every
noncentral element of \(N\) has at least one noncentral component,
and the product character satisfies (TD.17) with \(c=c_0\).

Conjugation on \(\operatorname{End}(W)\) factors through \(S\).
Its invariants have dimension one by irreducibility, and its
nontrivial part has dimension \(a^2-1\).
This part is nonzero, since otherwise every \(N\)-matrix would
be scalar.  Equation (TD.23) consequently gives
\[
                         a^2-1\ge999999,\qquad a\ge1000 .
                         \tag{TD.24}
\]
This argument also covers genuine projective simple-group modules:
no assumption that the center acts trivially on \(W\) has entered.

Apply (TD.18) to every nonscalar normalizer operator.
With \(b=1/c_0\) it gives the exact all-power bound
\[
 \frac{|\operatorname{tr}T|^2}{a^2}
 \le c_0+\frac{1-c_0}{1+1/c_0}
     =\frac{2c_0}{1+c_0}
     =\frac1{500000}<\frac1{170^2}.
 \tag{TD.25}
\]
The same bound holds for every projectively nonidentity power,
because every such power is again a nonscalar normalizer operator.
There is no assumption on the cyclic order or on the outer coset.
An operator which only projectively normalizes the represented
core also normalizes \(N\) itself: the full scalar preimage is
\(\mathbf C^\times N\), whose derived group is \(N\) because
\(N\) is perfect.
If only a finite projective factor image is initially given, take
its inverse image in \(\operatorname{SL}(W)\).  This is a finite
central lift with kernel \(\mu_a\).  It contains \(N\), since
perfectness makes every determinant on \(N\) trivial, and it
normalizes \(N\) by the preceding derived-group argument.
Averaging over this finite lift supplies a common positive
Hermitian form, so all the unitary normalizations used above
are legitimate without a field-of-definition assumption.

Let \(P\) be the actual projective image on \(W\).
It contains the faithful copy of \(S\), and hence is nonsolvable.
Moreover (TD.23) and the ordinary degree-square identity give
\(|S_i|\ge1+999999^2>60\), so \(|P|>60\).
All hypotheses of D-TD02 hold with \(c=1/170\).
It follows that \(e\le504\), whereas
\(e\ge a\ge1000\) by (TD.24), a contradiction.
\end{proof}

\paragraph{What remains.}
D-TD03 excludes cores all of whose defining fields are at least
\(10^{12}\); it does not silently cover a product mixing a small-field
factor with a large-field factor.
The normal core must act irreducibly on the chosen tensor factor.
D-TD01 applies whenever independent arguments bound every projective
factor image, including any future small-degree catalogue.
Neither theorem makes the original inducing field equal to \(K\),
decides a proper Hodge-locus descent, or bounds the arbitrary original
overgroups of D-OG.

\section{Large-field components in a primitive inducing image}
\label{sec:primitive-component-orbit}

\paragraph{Scope and status.}
D-TD04 is a separate extension of D-TD03.  The component-orbit
proposal is due to the root agent; the homogeneous restriction
strengthening and complete interface proof below are derived by
/root/imprimitive. Complete author and separate root whole-proof
readings pass. Root also read the full source/status record, checker
and first output, and checked the precise quasisimple source-scope
repair. The shared derivation is not relabelled as independent origin.
No status or mathematical text in the accepted D-TD01--03 is changed.

Let $J\subset\operatorname{GL}(V)$ be the finite faithful irreducible
image of an actual sphere-source coefficient module.  Put
$e=\dim V$, and retain its generating product-one inertia tuple
$T_1,\ldots,T_r$, with
\[
                  \sum_{i=1}^r\operatorname{rank}(T_i-I)=2e+6.
                  \tag{TD4.1}
\]
Assume that $J$ is primitive: it preserves no nontrivial direct-sum
decomposition whose summands it permutes.  A factorwise tensor
decomposition is allowed.  No hypothesis is imposed on the original
coefficient field or on the determinant of its polarization.

\begin{theorem}[D-TD04]
Suppose $J$ has a normal subgroup
\[
             N=E_1\cdots E_t
             \tag{TD4.2}
\]
which is a central product of commuting finite quasisimple components.
If one component $E_i$ admits a quasisimple Lie-type presentation
in the exact scope of Duyan--Halasi--Podoski Theorem~4.1(2),
with nonalternating simple quotient and defining field size
$q\ge10^{12}$, then the actual source (TD4.1) is impossible.

The other components may have arbitrarily small defining fields or
belong to different simple-group families.  The action of $N$ on $V$
need not be irreducible.  Factorwise tensor indecomposability of $J$
is not assumed.  The statement concerns a normal central product
actually represented in $J$, not a merely abstract simple section.
\end{theorem}

\begin{proof}
We first justify that the components can be grouped by an actual
$J$-orbit.  The multiplication map from the direct product of the
$E_i$ onto $N$ has central kernel.  Consequently
\[
                 N/Z(N)\simeq\prod_i E_i/Z(E_i).
                 \tag{TD4.3}
\]
For completeness, if a product $x_1\cdots x_t$ is central in $N$,
then commutation with $E_i$ shows that $x_i$ centralizes $E_i$,
because every other component commutes with it.  Hence
$x_i\in Z(E_i)$.  This proves the required assertion about centers
and the direct product in (TD4.3).

The simple direct factors of (TD4.3) are exactly its minimal nontrivial
normal subgroups.  Every automorphism of $N/Z(N)$ permutes them.
The preimage of the $i$th factor in $N$ is $E_i Z(N)$, and its
derived group is exactly $E_i$, because $E_i$ is perfect.
Thus conjugation by $J$ permutes the actual subgroups $E_i$, not
only their simple quotients.

Choose the orbit $I$ of a component with the stated large-field
presentation, and put
\[
                         H=\prod_{i\in I}E_i.
                         \tag{TD4.4}
\]
Then $H\triangleleft J$ is a nontrivial perfect central product.
All its components are isomorphic as quasisimple groups to the
chosen one, since the isomorphisms are actual conjugations by
elements of $J$.  Transport the chosen Lie-type presentation
through these isomorphisms.  The same $q$ and the same published
uniform character bound therefore apply to every component in
$I$, and to every actual nontrivial irreducible component module.
No uniqueness of a finite simple group's field presentation is
being assumed, and the other orbits are not assigned this bound.

By complete reducibility, restriction of $V$ to $H$ decomposes
into its canonical isotypic summands.  Normality of $H$ makes
$J$ permute these summands.  Irreducibility of $J$ makes the
permutation transitive; otherwise the sum over one orbit would
be a proper $J$-invariant subspace.  If there were more than one
summand, they would form a nontrivial system of imprimitivity.
Primitivity therefore leaves exactly one isotypic type.
Choose an irreducible $H$-module $W$ of that type and its
multiplicity space $U=\operatorname{Hom}_H(W,V)$.  Evaluation gives
\[
              V\simeq W\otimes U,\qquad
              x|_V=\rho_W(x)\otimes I_U\quad(x\in H).
              \tag{TD4.5}
\]
The representation $\rho_W$ is faithful: an element in its kernel
would act trivially on every isotypic copy, hence on $V$, and $J$
is an actual matrix image.  The module $W$ is nontrivial and has
dimension greater than one, since a nontrivial perfect group has
no nontrivial one-dimensional representation.  Thus $H$ is a
finite faithful perfect irreducible subgroup of
$\operatorname{GL}(W)$.

We now check the whole factorwise normalizer interface.  The complex
algebra generated by $H$ is
\[
             \mathcal A=\operatorname{End}(W)\otimes I_U,
             \tag{TD4.6}
\]
by the irreducible matrix-algebra theorem
\cite[\S6.2, Proposition~10]{OG:Serre1977}.
Every $g\in J$ normalizes this algebra.  An automorphism of a full
complex matrix algebra is inner: the images of matrix units give
a defining module basis, whose change-of-basis matrix implements
the automorphism.  Choose $A_g\in\operatorname{GL}(W)$ implementing
the action on $\operatorname{End}(W)$.  The operator
$(A_g^{-1}\otimes I)g$ centralizes $\mathcal A$.  Its centralizer
is $I_W\otimes\operatorname{End}(U)$, as one sees by commuting
with the individual matrix units.  Hence
\[
                        g=A_g\otimes B_g,
                        \qquad B_g\in\operatorname{GL}(U).
                        \tag{TD4.7}
\]
The first factor is unique up to a scalar.  Indeed, if
$A\otimes B=A'\otimes B'$ with all four maps invertible,
then $A'^{-1}A\otimes I=I\otimes B'B^{-1}$ lies in the
intersection of the two full matrix algebras, which consists
of the scalar operators.  It follows that
\[
                  g\longmapsto[A_g]\quad\hbox{defines}
                  \pi:J\longrightarrow\operatorname{PGL}(W).
                  \tag{TD4.8}
\]
This is a genuine homomorphism with finite image $P$, even if
there is no coherent choice of the individual linear maps $A_g$.

For $x\in H$, conjugating (TD4.5) and using (TD4.7) gives the
exact identity
\[
          A_g\rho_W(x)A_g^{-1}=\rho_W(gxg^{-1}).
          \tag{TD4.9}
\]
Thus all representatives $A_g$ normalize the represented group
$\rho_W(H)$ itself, without an untracked scalar character.
The inverse image $\widetilde P$ of the finite group $P$ under
$\operatorname{SL}(W)\to\operatorname{PGL}(W)$ is finite, with
kernel $\mu_{\dim W}$.  Every element of $\widetilde P$ is a
scalar multiple of some $A_g$, so also normalizes $\rho_W(H)$.
Perfectness makes $\det\rho_W$ trivial, so $\widetilde P$
contains $\rho_W(H)$.  This verifies the exact finite-linear-lift
and unitary-averaging hypotheses used in D-TD03.

Finally, applying $\pi$ to the actual inertia tuple gives a
product-one generating tuple for $P$.  Omit its identity entries.
It is the monodromy tuple of the actual connected quotient of
the $J$-cover by $\ker\pi$.  No lower bound of twelve on the
number of its nonidentity entries is asserted or required.
For every actual $g=A_g\otimes B_g$, convolution of eigenvalue
multiplicities shows that the largest normalized eigenspace of
$g$ is no larger than that of $A_g$.  This holds for every scalar
normalization of $A_g$ and every choice of the partner $B_g$.
Therefore all projective support estimates on $W$ enter the
actual moving sum (TD4.1), exactly as in D-TD02--03.

All components of $\rho_W(H)$ have the transported large-field
presentation.  D-TD03 applies to $V=W\otimes U$ and yields a
contradiction.  Explicitly, its full-normalizer all-power trace
bound is at most $1/170$, its factor degree is at least $1000$,
and the actual quotient-cover orbifold estimate forces
$e\le504$, whereas $e\ge\dim W\ge1000$.
This proves the theorem.
\end{proof}

\paragraph{Exact scope of the strengthening.}
The mixed-small/large-field boundary in D-TD03 is removed here
when the components belong to a normal central product of an
actual primitive inducing image.  Its complementary components
are absorbed into the multiplicity partner in (TD4.5); no trace
bound is asserted for them.  The existence of such a normal
central product remains a hypothesis.  In particular this does
not classify arbitrary primitive groups with solvable normal
structure, replace the remaining low-field analysis, or address
arbitrary simple sections that are not normal components.

\section{Projected extraspecial factors: a prime bound and the elliptic escape}
\label{sec:projected-extraspecial}

\paragraph{Scope and review status.}
D-TD05--07 are new partial results for actual factorwise tensor quotients.
The parent suggested the scalar-times-core budget and the full normalizer
trace direction; the proofs below are the author's subsequent derivation.
Complete author and separate root whole-proof readings pass. Root also read the complete source/status record, every checker line, full first output and exact input/exposition binding; no root rerun is claimed.  No bound on the
remaining exponent $m$ at the small primes is asserted.  In particular,
the elliptic examples at the end are projected data, not sixfold sources
with the required full inertia budget or Hodge signature.

\paragraph{Actual source and projective factor.}
Let $C\to\mathbf P^1$ be a connected finite $J$-Galois cover, where
$J\subset\operatorname{GL}(V)$ is its faithful irreducible coefficient image.  Write
$e=\dim V$ and let $T_1,\ldots,T_r$ be a generating product-one inertia
tuple satisfying the sphere-source identity
\[
                 \sum_i\operatorname{rank}(T_i-I)=2e+6.
                 \tag{TDE.1}
\]
Suppose $V=W\otimes U$ is preserved factorwise by $J$, with reciprocal
scalar ambiguity allowed in its two factors.  There is then an actual
finite quotient
\[
        P=\operatorname{im}\bigl(J\longrightarrow\operatorname{PGL}(W)\bigr),
        \qquad Z=C/\ker(J\to P),
        \qquad Z/P=\mathbf P^1.
        \tag{TDE.2}
\]
Assume $P$ contains and normalizes a subgroup
\[
 Q\simeq (\mathbf F_p)^{2m},\qquad m\ge1,
 \tag{TDE.3}
\]
whose projective action on $W$ is irreducible and whose commutator pairing
on scalar lifts is nondegenerate.  Thus, for lifts $L_v,L_w$ of $v,w\in Q$,
\[
 L_vL_wL_v^{-1}L_w^{-1}=\beta(v,w)I,
 \qquad \beta:Q\times Q\longrightarrow\mu_p
 \tag{TDE.4}
\]
is a nondegenerate alternating pairing.  This is the intrinsic
projective Heisenberg hypothesis; no splitting of a normalizer extension
is included in it.  Put
$\nu(A)=\dim W-\max_\lambda\dim\ker(A-\lambda I)$.

This setup applies to a normal extraspecial subgroup of a primitive
actual image $J$.  Indeed, the isotypic summands of its restriction to
$V$ are permuted by $J$, so primitivity makes the restriction homogeneous.
Evaluation gives $V=W\otimes U$ with irreducible $W$ and faithful action
of that subgroup on $W$: an element acting trivially on $W$ would act
trivially on $V$.  Its center is scalar by Schur's lemma.  The algebra
it generates is $\operatorname{End}(W)\otimes I$, so its normalizer acts
factorwise, by the matrix-algebra normalizer argument in D-TD04.
The projective quotient in (TDE.2) is therefore genuine.  Neither
factorwise tensor indecomposability nor an equality of coefficient
fields is being assumed.

\subsection{D-TD05: exact trace of every projective normalizer}

\begin{theorem}[D-TD05]
Under (TDE.3)--(TDE.4), $a=\dim W=p^m$.  If a unitary operator $A$
normalizes $Q$ projectively, let $g$ denote its induced symplectic
automorphism of $Q$.  Then
\[
        |\operatorname{tr}A|^2\in\{0,\ p^{\dim\operatorname{Fix}(g)}\}.
        \tag{TDE.5}
\]
Every nonscalar such $A$, including one with $g=1$, satisfies
\[
                  |\operatorname{tr}A|/a\le p^{-1/2}.
                  \tag{TDE.6}
\]
Moreover, the projective centralizer is exactly
\[
                   C_{\operatorname{PGL}(W)}(Q)=Q.
                   \tag{TDE.7}
\]
Thus the kernel of the normalizer's action on $Q$ is $Q$, not merely
the scalar identity.
\end{theorem}

\begin{proof}
The inverse image of the finite group $P$ in $\operatorname{SL}(W)$
is finite, since its kernel is $\mu_a$.  Averaging a Hermitian form
unitarizes it.  The inverse image $E$ of $Q$ has scalar center $B$;
nondegeneracy of $\beta$ shows that no other element is central.
Its representation on $W$ is irreducible by hypothesis.  If $v\ne0$,
there is $w$ with $\beta(w,v)\ne1$.  Conjugating $L_v$ by $L_w$
therefore proves $\operatorname{tr}L_v=0$.
Character orthogonality for the irreducible representation of $E$
now gives
\[
 1=\frac1{|E|}\sum_{x\in E}|\operatorname{tr}x|^2
   =\frac{|B|a^2}{|E|}=\frac{a^2}{|Q|}.
 \tag{TDE.8}
\]
This proves $a=p^m$.  Choose unitary lifts $L_v$, with $L_0=I$.
They are pairwise orthogonal for the trace inner product, because
$L_vL_w^{-1}$ has zero trace when $v\ne w$.  There are $a^2$ of
them, so they are a basis of $\operatorname{End}(W)$.
The only known representation-theoretic input here is finite character
orthogonality, in \cite[Sections2.3--2.4]{OG:Serre1977}.

Conjugation by $A$ preserves commutators, so $g$ preserves $\beta$.
Write
\[
 A L_v A^{-1}=\eta_v L_{gv},\qquad |\eta_v|=1.
 \tag{TDE.9}
\]
Let $F=\operatorname{Fix}(g)$.  The restriction $\eta|_F$ is an
ordinary character of the additive group $F$.  To see this without
any phase convention, write $L_vL_w=c(v,w)L_{v+w}$.  For $v,w\in F$,
conjugating this equality gives the same factor $c(v,w)$ on both
sides, and hence $\eta_{v+w}=\eta_v\eta_w$.  Also $\eta_0=1$.
In particular this argument works at $p=2$; it neither asserts a
globally linear phase on all of $Q$ nor assumes a quadratic-refinement
convention.

Taking the trace of conjugation on the displayed basis gives
\[
 |\operatorname{tr}A|^2
 =\operatorname{tr}\bigl(\operatorname{Ad}(A)|\operatorname{End}(W)\bigr)
 =\sum_{v\in F}\eta_v.
 \tag{TDE.10}
\]
The last sum is $|F|$ for the trivial character and zero otherwise.
This proves (TDE.5).  If $g\ne1$, then $\dim F\le2m-1$, proving
(TDE.6).  If $g=1$ and $\eta$ is trivial, $A$ centralizes a basis
of the full matrix algebra, so it is scalar.  If $g=1$ and $A$
is nonscalar, $\eta$ is nontrivial and its trace is zero; this is
the other case of (TDE.6).

Finally suppose $[A]$ centralizes $Q$ projectively.  The argument for
$g=1$ makes $\eta$ a character of $Q$.  Nondegeneracy identifies
$Q$ with its character group by $u\mapsto\beta(u,{\cdot})$.
Choose the corresponding $u$.  Then $L_u^{-1}A$ centralizes every
$L_v$ and is scalar.  Conversely every element of $Q$ centralizes
$Q$ projectively, giving (TDE.7).  This last argument also applies
to a nonunitary representative, since its phases form a finite
character and the matrix-algebra conclusion is algebraic.
\end{proof}

\subsection{D-TD06: the actual-cover prime bound}

\begin{theorem}[D-TD06]
Assume (TDE.1)--(TDE.4).  If $g(Z)\ge2$, then
\[
                 e(1-85/\sqrt p)\le252,
                 \qquad p<7744.
                 \tag{TDE.11}
\]
If $m\ge2$, the genus hypothesis is automatic.  Consequently every
source with $p\ge7744$ has $m=1$ and $g(Z)=1$.

In the special scalar-times-core case $P=Q$, the necessary list is
\[
             (m,p)\in\{(1,2),(1,3),(1,5),(1,7),(2,2)\}.
             \tag{TDE.12}
\]
The list is a necessary budget condition, not a statement of existence.
\end{theorem}

\begin{proof}
Remove identity images from the projected inertia tuple of (TDE.2),
and let $o_1,\ldots,o_k\ge2$ be the exact orders of the remaining
images.  These elements generate $P$ and have product one.
For an element of order $o$, choose a unitary representative $A$
with $A^o=I$.  Every $A^j$, $1\le j<o$, is nonscalar.  Fourier
inversion of its eigenvalue multiplicities and (TDE.6) give
\[
 \frac{\nu(A)}a\ge(1-p^{-1/2})(1-1/o).
 \tag{TDE.13}
\]
This estimate applies to an arbitrary eigenvalue, so scalars cannot
affect it.  If $B$ is the operator on $U$, each eigenspace of $A\otimes B$
has dimension at most
$(a-\nu(A))\dim U$: decompose $U$ into its eigenspaces and apply
the largest-eigenspace bound to each summand.  Thus (TDE.13) is a
lower bound for the actual moving rank divided by $e$, regardless
of the tensor partner.  Identity projected images contribute a
nonnegative amount and may be discarded.

Riemann--Hurwitz for the actual cover gives
\[
       2g(Z)-2=|P|\left(\sum_{i=1}^k(1-1/o_i)-2\right).
       \tag{TDE.14}
\]
If $g(Z)\ge2$, the positive quantity in parentheses is at least
$1/42$.  Here is the elementary integer-order check.  For $k\ge5$
it is at least $1/2$.  For $k=4$ the zero case is all orders two,
and the least positive excess is $1/6$.  For $k=3$, order the
orders increasingly: first order at least four gives excess at
least $1/4$; first order three gives least positive case $(3,3,4)$,
of excess $1/12$; first order two leaves $(2,3,w)$, $w\ge7$,
$(2,4,w)$, $w\ge5$, or second order at least five, giving respectively
least excesses $1/42$, $1/20$, and $1/10$.  Fewer marks cannot
give positive excess.  Summing (TDE.13) therefore yields
\[
                  2+6/e\ge (1-p^{-1/2})85/42.
                  \tag{TDE.15}
\]
This is the first inequality of (TDE.11).  If $p\ge7744=88^2$,
then $e\ge p^m\ge p$, and
\[
 e(1-85/\sqrt p)\ge\sqrt p(\sqrt p-85)\ge88\cdot3=264>252,
 \tag{TDE.16}
\]
a contradiction.  No count of twelve nonidentity projected marks
has been used.

We next justify the genus assertions.  An elementary abelian odd
$p$-subgroup of $\operatorname{PGL}_2(\mathbf C)$ has rank at most
one: a nonidentity element has two fixed points, and every commuting
element preserves each of them.  An interchange would conjugate the
element to its inverse, which is different for odd $p$.  Such a group
is contained in the diagonal torus, whose $p$-torsion is cyclic.
For $p=2$, the subgroup fixing both points has order at most two,
and the permutation group on the pair has order at most two; the
elementary abelian rank is at most two.  Thus genus zero is impossible
when $m\ge2$, and is also impossible for $m=1$, $p$ odd.

On a complex genus-one curve choose an origin.  Its automorphism
group is translations semidirect the cyclic origin-fixing group of
order two, four, or six; see
\cite[ChapterIV, Proposition7.13 and Remark7.18, pp.~140--141]{TD:MilneEC2021}.
The kernel of the image of $Q$ in this cyclic group lies in the
$p$-torsion translations, of rank two.  The image has rank at most
one.  The full rank of $Q$ is therefore at most three, and at most
two when $p\ge5$.  The torsion assertion is
\cite[ChapterII, Proposition6.3, p.~70]{TD:MilneEC2021}.
For $m\ge2$ the rank $2m\ge4$ excludes genus one too.  Combining
these facts with (TDE.11) proves the remaining genus claims.

Finally let $P=Q$.  Every nonidentity projected element has order
$p$ and every proper nontrivial power has trace zero, so Fourier
inversion gives every eigenvalue multiplicity exactly $a/p$.
If $k$ nonzero vectors generate $\mathbf F_p^{2m}$ and sum to zero,
then $k\ge2m+1$.  Consequently
\[
      (2m+1)(1-1/p)\le 2+6/e\le 2+6/p^m.
      \tag{TDE.17}
\]
For $m\ge3$, the left side is at least $7/2$ and the right side
at most $11/4$, impossible.  For $m=2$, (TDE.17) is
$3p^2-5p-6\le0$, leaving only the prime two.  For $m=1$ it gives
$p\le9$, leaving precisely the four primes in (TDE.12).
If $J$ is primitive, this scalar-times-core case is impossible
altogether: a maximal isotropic subgroup of $Q$ has $p^m$ distinct
one-dimensional common eigenspaces on $W$ (character Fourier
inversion as above).  Explicitly, commuting lifts of a basis of that
subgroup can be scaled to have order $p$, yielding an ordinary
representation of its additive group.  Every nonzero element has
trace zero, so each character occurs once.  The group $Q$ permutes
these eigenlines.  Tensoring those lines
with $U$ gives a nontrivial system of imprimitivity for $J$.
\end{proof}

\subsection{D-TD07: exact description of the large-prime escape}

\begin{theorem}[D-TD07]
If a source in D-TD06 has $p\ge7744$, then
\[
     m=1,\quad g(Z)=1,\quad
     P=Z[p]\rtimes C_d,\qquad d\in\{2,3,4,6\}.
     \tag{TDE.18}
\]
If the actual factor action $P\subset\operatorname{PGL}(W)$ is
primitive, the only remaining cases are
\[
 \begin{array}{c|c}
 d&\text{necessary and sufficient condition for this factor action
 to be primitive}\\ \hline
 3\text{ or }6&p\equiv2\pmod3\\
 4&p\equiv3\pmod4 .
 \end{array}
 \tag{TDE.19}
\]
In particular the $d=2$ factor is imprimitive.  These conclusions
are necessary restrictions on an actual source.  They do not bound
$m$ for the primes below $7744$.
\end{theorem}

\begin{proof}
The first two assertions follow from D-TD06.  Put
$T=P\cap\operatorname{Trans}(Z)$.  As $p\ge5$, the subgroup
$Q$ acts trivially on the origin-fixing cyclic quotient, so
$Q\subset T$.  Translations commute.  Thus $T$ centralizes $Q$
projectively, and (TDE.7) gives $T=Q$.  Since both $Q$ and $Z[p]$
have order $p^2$, they coincide.  It follows that $P/Q$ is cyclic
of some order $d\in\{1,2,3,4,6\}$.  The case $d=1$ would make
$Z/P$ a quotient by translations and hence another genus-one curve,
contradicting (TDE.2).

Choose an element whose linear part generates the cyclic quotient.
On $Z=\mathbf C/\Lambda$ it has the form $z\mapsto\zeta_d z+b$.
The nonzero endomorphism $1-\zeta_d$ is surjective, so this element
has a fixed point.  Taking it as origin makes the element linear
of exact order $d$.  It supplies the asserted complement.  The
complex-torus description and the characterization of its
origin-preserving maps are
\cite[ChapterIII, Proposition3.3, p.~94, and Theorem3.10, p.~96]{TD:MilneEC2021}.

We check primitivity rather than assuming a normalizer classification.
The action of the complement on $Q=Z[p]$ is the reduction of its
action on the rank-two period lattice.  For orders $2,3,4,6$ the
characteristic polynomials are, respectively,
\[
             (X+1)^2,\quad X^2+X+1,\quad X^2+1,\quad X^2-X+1.
             \tag{TDE.20}
\]
If this action preserves a line $L\subset Q$, the lifts of $L$
have $p$ distinct common eigenlines on $W$: their nontrivial traces
vanish and $\dim W=p$.  These lines are permuted by $P$, making
the factor imprimitive.  Conversely, any nontrivial imprimitivity
system in the prime dimension $p$ consists of $p$ lines.  The
elementary abelian $p$-group $Q$ acts on them.  Its image in $S_p$
is either trivial or cyclic of order $p$.  In the trivial case
all lifts are simultaneously diagonal and commute, contrary to
(TDE.4).  In the other case the kernel is a line in $Q$ and is
$P$-invariant.  Therefore primitivity is equivalent to irreducibility
of the complement action on $Q$.  Reducing the polynomials (TDE.20)
gives exactly (TDE.19).  If the original $J$ is primitive, its
projective factor $P$ is primitive, since any block decomposition
on $W$ tensors with $U$ to a block decomposition on $V$.
\end{proof}

\paragraph{An actual projected family with unbounded primes.}
For every odd prime $p\equiv3\pmod4$ let
\[
       E:\ y^2=x^3-x,\qquad \iota(x,y)=(-x,iy),\qquad
       E\xrightarrow{[p]}E\xrightarrow{x^2}\mathbf P^1.
       \tag{TDE.21}
\]
The composite is a connected Galois cover with group
$E[p]\rtimes\langle\iota\rangle$, of order $4p^2$.
The multiplication map is unramified, and the second map is the
quotient by the order-four automorphism.  Its two order-four fixed
points are $(0,0)$ and infinity; the other nonzero two-torsion
points form one orbit with stabilizer of order two.  Thus the
inertia signature is $(2,4,4)$.

Here is a matching faithful projective Heisenberg representation.
On the basis $e_j$, $j\in\mathbf F_p$, let $\zeta$ be a primitive
$p$th root and set
\[
 Xe_j=e_{j+1},\qquad Ze_j=\zeta^j e_j,\qquad
 Fe_j=p^{-1/2}\sum_k\zeta^{jk}e_k.
 \tag{TDE.22}
\]
Then $ZX=\zeta XZ$, $FXF^{-1}=Z$, $FZF^{-1}=X^{-1}$,
$F^2e_j=e_{-j}$, and $F^4=I$.  These formulas show directly that
$\langle X,Z,F\rangle$ is finite and its projective image is
$\mathbf F_p^2\rtimes C_4$, with complement matrix
$\left(\begin{smallmatrix}0&-1\\1&0\end{smallmatrix}\right)$.
This is the same semidirect product as the cover in (TDE.21),
since the irreducible characteristic polynomial of both actions
is $X^2+1$.  The factor is primitive by the proved criterion.
There are arbitrarily large such primes by Dirichlet's theorem
\cite[Section1, p.~63]{PD:Harrison2009}.

This example establishes that the projected quotient-cover and
projective-factor hypotheses alone allow the genus-one escape
with unbounded $p$.  It does not construct a full $J$-cover and
tensor partner obeying (TDE.1), nor prove a balanced Hodge
signature, period rank nine, coefficient-field descent, or a
prescribed polarization.  Those additional source conditions,
and the exponent $m$ for the remaining small primes, remain gates.

\section{A triple-cover obstruction for large ordinary alternating tensor factors}
\label{sec:factorwise-triples}

\paragraph{Status and exact interface.}
These are new root statements HD-FT01--02, question Q2da.
The complete author/root and separate \texttt{/root/imprimitive}
whole-frozen-proof readings pass. The complete detailed independent
source audit is retained; no finite supplemental computation is
required by or claimed as completed in this integration. They concern an actual factorwise tensor decomposition, a
case explicitly outside D-BI19. They do not assume that a projected
factor inherits twelve genuine branch marks.

Let an actual connected finite-group cover of the sphere have coefficient
representation \(V\), degree \(e\), and genuine generating product-one
inertia \(T_1,\ldots,T_r\), with
\[
                  \sum_i\operatorname{rank}(1-T_i)=2e+6.
                  \tag{FT.1}
\]
Suppose there is an actual quotient \(G\longrightarrow P\), with
\(P=A_n\) or \(S_n\) in its natural action, and an actual factorization
\[
             V=W\otimes U,\qquad
             T_i=W(x_i)\otimes U_i
             \quad\hbox{up to a scalar},
             \tag{FT.2}
\]
where \(x_i\) is the image of the inertia in \(P\), \(a=\dim W\),
and \(\dim U\ge1\). The tensor partners and all scalars are unrestricted.
Only the existence of the displayed actual quotient and representation
factorization is assumed; an abstract numerical profile is insufficient.

\subsection{HD-FT01: cubic traces, actual triple covers, and no mark-count premise}

\begin{theorem}[HD-FT01]
Assume \(n>3500000\), \(a\ge n^{120}\), and, with \(f(x)\) the number
of fixed natural letters, assume
\[
 \frac{|\chi_W(x)|}{a}\le
 \begin{cases}
       (f(x)/n)^3,&f(x)>0,\\
       2n^{-6},&f(x)=0 .
 \end{cases}
 \tag{FT.3}
\]
Then no source (FT.1)--(FT.2) exists.
The theorem is independent of the number of genuine projected marks,
the degree or field of \(U\), and the original imaginary quadratic field.
\end{theorem}

\begin{proof}
For a finite-order operator \(B\), put
\(\nu(B)=\dim B-\max_z\dim\ker(B-z)\).
Tensor spectral convolution gives
\[
       \frac{\operatorname{rank}(1-T_i)}e
                   \ge\frac{\nu(W(x_i))}{a}.
       \tag{FT.4}
\]
This also allows arbitrary scalar factors on either side of (FT.2).
If \(x_i=1\), we discard its nonnegative contribution below.
Let \(r_0\) denote the number of remaining indices.

First bound \(r_0\) from the actual total budget.
Since \(2n^{-6}\le n^{-3}\), (FT.3) implies the exact cubic
all-target hypothesis in HD-TW02. Its purely local proof therefore gives
\[
 \frac{\nu(W(x))}{a}
 \ge \frac{(n-4)\operatorname{ind}(x)}{n^2}+\frac1n
 \ge \frac{2n-4}{n^2}
 \qquad(x\ne1).
 \tag{FT.5}
\]
That local proof uses neither transitivity of tensor factors nor a
branch-mark count. From (FT.1), (FT.4), and \(e\ge a\ge n^{120}\),
\[
 r_0\le
 \left(2+\frac6e\right)\frac{n^2}{2n-4}
 \le \left(2+\frac6{n^3}\right)\frac{n^2}{2n-4}<2n.
 \tag{FT.6}
\]
For the last strict inequality it suffices that
\(2n^2-8n-6/n>0\), which holds for \(n\ge5\).

Let \(\Omega=\{1,\ldots,n\}^3\), with ordered entries and repetitions
allowed. For \(x\) of order \(q\), Burnside's orbit formula gives
\[
 \operatorname{ind}(x,\Omega)
       =n^3-\frac1q\sum_{j=0}^{q-1}f(x^j)^3.
       \tag{FT.7}
\]
Cyclic Fourier inversion bounds every eigenvalue multiplicity of \(W(x)\)
by the average absolute character value. Using (FT.3), including the
identity power and every fixed-point-free power, gives
\[
 \frac{\nu(W(x))}{a}
    \ge\frac{\operatorname{ind}(x,\Omega)}{n^3}
                  -\frac2{n^6}.
 \tag{FT.8}
\]
The error can be multiplied by the actual fraction of fixed-point-free
powers, but the displayed weaker bound suffices.

The image \(P\) is three-transitive because \(n\ge5\).
Its action on \(\Omega\) has exactly five orbits: the diagonal,
three equality patterns with exactly two distinct letters, and
the ordered triples of distinct letters. These are actual connected
quotient covers of the given \(P\)-cover. Write their genera as
\(\gamma_1,\gamma_2,\gamma_2,\gamma_2,\gamma_3\).
Riemann--Hurwitz on each orbit yields the exact sum
\[
 \sum_i\operatorname{ind}(x_i,\Omega)
       =2(n^3-5)+2(\gamma_1+3\gamma_2+\gamma_3).
       \tag{FT.9}
\]
Entries with \(x_i=1\) contribute zero, so the same equality holds
over the \(r_0\) retained entries.

Consider also the actual quotient cover on unordered three-element
subsets. Its point stabilizer is
\(H_3=P\cap(S_{n-3}\times S_3)\).
The published Monderer--Neftin theorem states that for \(n>N_1\),
every subgroup not containing \(A_{n-1}\) whose fixed curve has
genus at most one must contain a conjugate of \(A_{n-2}\).
Its explicit constant is \(N_1=3500000\), using the classification
of finite simple groups; the classification-free stated alternative
is \(60000000000\).
The exact input is Theorem1.2, with Remark2.8, of
\cite{HD:MondererNeftin2022}.

Our \(H_3\) cannot contain either required alternating subgroup:
\[
 |H_3|\le6(n-3)!<
           \frac{(n-2)!}{2}=|A_{n-2}|
           \qquad(n>14).
 \tag{FT.10}
\]
Thus the unordered-three-subset curve has genus at least two.
The forget-order map from the distinct ordered-triple curve has
degree six, for both \(P=A_n\) and \(S_n\).
For \(A_n\), every permutation of the selected three letters
can be paired with a permutation of matching parity on their
complement, since \(n-3\ge2\). This proves the degree assertion
also in that case.
Riemann--Hurwitz now gives
\[
                  \gamma_3-1\ge6(2-1),\qquad\gamma_3\ge7.
                  \tag{FT.11}
\]
The other genera are nonnegative. Hence (FT.9) is at least
\(2n^3+4\).

Finally sum (FT.4) and (FT.8), and use (FT.6):
\[
 \frac1e\sum_i\operatorname{rank}(1-T_i)
 \ge2+\frac4{n^3}-\frac{2r_0}{n^6}
 >2+\frac4{n^3}-\frac4{n^5}
 >2+\frac3{n^3}.
 \tag{FT.12}
\]
Since \(e\ge n^{120}\), the last strict excess contributes more
than \(3n^{117}>6\) to the moving-rank sum, contradicting (FT.1).
This proves the theorem.
\end{proof}

\paragraph{Source-reading precision.}
The author revised version dated18 December2020 has Theorem1.2 on
page2, Theorems2.5--2.6 and Remark2.8 on pages7--8, and the complete
function-field proof, Theorem5.1, on pages16--17.
Root read those complete invoked passages, including the full
intermediate Proposition3.1--3.3 and Corollary3.4 proofs.
Publisher metadata verifies the published article, its revision date
and its2022 pagination. The publisher PDF request led to the
subscription preview, so no reading of inaccessible typeset pages
is claimed. Remark2.8 states that its constants come from inspection
of the Neftin--Zieve argument; those numerical constants have not
been independently recomputed in this run.
We invoke the published Monderer--Neftin result with that exact
source scope, not an invented effective version of a merely
asymptotic statement.

\subsection{HD-FT02: an effective growing-tail exclusion with arbitrary tensor partners}

\begin{corollary}[HD-FT02]
In (FT.1)--(FT.2), suppose \(W\) is an irreducible \(S_n\)-module
or its irreducible ordinary restriction to \(A_n\).
If \(n>3500000\) and \(a\ge n^{120}\), the source is impossible.
In particular set
\[
                         N_{\rm tail}=241^{482}.
                         \tag{FT.13}
\]
For \(n\ge N_{\rm tail}\), a surviving such factor must have a
partition label \(\lambda\) satisfying
\[
          \min\{n-\lambda_1,n-\lambda'_1\}\le240 .
          \tag{FT.14}
\]
This assertion allows arbitrary factorwise tensor partners.
It does not claim that these remaining low-tail labels occur.
\end{corollary}

\begin{proof}
The published effective MSP consequence was proved in HD-TW02:
when \(n\ge4000\) and \(a\ge n^{120}\), the absolute character ratio
is at most \((f/n)^3\) for \(f>0\).
Its separate fixed-point-free calculation gives the stronger bound
\[
             |\chi_W(x)|/a<a^{-1/20}\le n^{-6}
                         \qquad(f(x)=0).
                         \tag{FT.15}
\]
Restriction to the ordinary irreducible \(A_n\)-module preserves
both degree and trace. Thus (FT.3) holds, and HD-FT01 applies.
No claim is made here for a nonextendible constituent obtained by
splitting a self-conjugate \(S_n\)-label.

It remains to prove (FT.14) without assuming the earlier long-row
source reduction. Suppose both tails are at least241.
Transpose \(\lambda\) if necessary for the dimension calculation,
and put \(m=\max\{\lambda_1,\lambda'_1\}\).
Since the diagram fits in its first-row by first-column rectangle,
\(m\ge\sqrt n\).
There are at least241 boxes below its first row.
Remove corners from this lower partition until a partition
\(\mu\) of241 remains. Then \((m,\mu)\) is a subpartition of
\(\lambda\). The ordinary branching rule gives
\[
                 f^\lambda\ge f^{(m,\mu)}.
                 \tag{FT.16}
\]
Equivalently, restrict the irreducible successively along a chain
of corner removals; this constituent has positive multiplicity.
The exact published MSP Lemma8(ii), applied to \((m,\mu)\), gives
\[
 f^\lambda\ge\binom{m}{241}f^\mu
       \ge(m/241)^{241}
       \ge\frac{n^{120+1/2}}{241^{241}}
       \ge n^{120}
                   \qquad(n\ge241^{482}).
 \tag{FT.17}
\]
Here \(m\ge241\), so the binomial product bound is valid.
Transposition preserves dimension. Also
\(241^{482}>3500000\).
The first assertion therefore excludes this case and proves
(FT.14).
\end{proof}

\paragraph{Geometric and programme scope.}
This is a new necessary gate for actual factorwise tensor-decomposable
curve suppliers, using the genera of actual finite quotient curves and
the weight-one Chevalley--Weil budget. It neither assumes that a factor
is itself a rank-nine source nor transfers twelve marks to it.
The bound still leaves low tails in unbounded natural degree, other
projective factor types, and at this stage nonextendible split
alternating constituents.
It supplies no cycle and no detector-degree lower bound.
For condition A (cote16,scan050; SIGN_NEXT SN-C03/C04) it removes
a possible supplier range. Condition I (scan049;cote44,Theorem4),
the Cayley--Hamilton projector of scans048--049, the existing
SGA/EGA/FGA spreading and realization interfaces, and the GHOSTS B03
interpretation remain as already stated. Markman, Severn and032
are not premises of this argument.

\subsection{HD-FT03: split alternating constituents as arbitrary tensor factors}
\label{sec:split-alternating-factor}

\paragraph{Status and provenance.}
This is a new addendum by \texttt{/root/imprimitive} to the root
HD-FT01--02 argument.  The symmetric-corner construction and the
explicit inequalities below are new derivations.  The alternating
character formula is a published input, with exact locators below.
Complete author and separate root whole-proof readings pass. Root
also read the complete source/status record, bibliography, author
checker and first successful output; the exact reading and source
access boundaries are recorded in the retained root review.
No projected twelve-mark hypothesis is introduced.

\begin{theorem}[HD-FT03]
Let \(\lambda=\lambda'\vdash n\), and let \(W\) be either of the
two ordinary irreducible constituents of the restriction of the
\(S_n\)-module labelled by \(\lambda\) to \(A_n\).
Write \(a=\dim W\), and let \(f(x)\) count the fixed natural letters
of \(x\in A_n\).  If \(n\ge65536\), then
\[
 a\ge n^{120},\qquad
 \frac{|\chi_W(x)|}{a}\le
 \begin{cases}
       (f(x)/n)^3,&f(x)>0,\\
       n^{-6},&f(x)=0 .
 \end{cases}
 \tag{FTS.1}
\]
Consequently, when \(n>3500000\), such a \(W\) cannot occur in the
actual factorwise tensor configuration (FT.1)--(FT.2) of HD-FT01,
with arbitrary tensor partner and arbitrary scalar factors.
This conclusion does not depend on the original imaginary
quadratic field or on the number of nonidentity projected inertias.
\end{theorem}

\begin{proof}
Put \(F=f^\lambda\), the degree of the \(S_n\)-module, so that
\(a=F/2\).  Let \(r\) be the number of diagonal boxes of \(\lambda\),
and write its diagonal hook lengths as
\[
 d_1>d_2>\cdots>d_r>0,\qquad
 d_i=2(\lambda_i-i)+1,\qquad
 \sum_i d_i=n,\qquad r^2\le n.
 \tag{FTS.2}
\]
The hooks in (FTS.2) are distinct odd positive integers.  The sum
identity follows by partitioning the symmetric diagram into its
diagonal hooks, and their minimum possible sum is \(r^2\).

First we give a lower bound for \(F\) using standard tableaux.
Starting with any nonempty self-conjugate diagram, select a corner
by a fixed deterministic rule.  If it is off the diagonal, its
transpose is a distinct corner.  These two boxes are incomparable
in the diagram order.  Delete both, assigning to them the two
largest available labels in either order.  If the selected corner
is diagonal, delete it and assign the largest available label.
After either operation the remaining diagram is self-conjugate.
The two choices in a paired deletion do not change this remaining
diagram, and therefore do not change any subsequent selected pair.
Reversing the deletions produces standard tableaux: each deletion
is a valid corner removal, including both orders of a paired removal.
Different choices produce different labelled tableaux.
Exactly \(r\) diagonal boxes are removed singly, and all other
boxes are removed in pairs.  The usual interpretation of \(F\) as
the number of standard tableaux, equivalently repeated branching,
therefore gives
\[
                    F\ge 2^{(n-r)/2}
                       \ge2^{(n-\sqrt n)/2}.
                    \tag{FTS.3}
\]
This construction uses a subset of all tableaux and does not assert
an exact formula for their number.

We record elementary logarithmic estimates with a completely
explicit endpoint.  For every real \(x\ge2^{16}\),
\[
 1+\log_2x\le\frac{\sqrt x}{4},\qquad
 1+6\log_2x\le\frac{x}{4}.
 \tag{FTS.4}
\]
At \(x=2^{16}\) these read \(17\le64\) and \(97\le16384\).
The differences between the right and left sides have derivatives
\[
 \frac1{8\sqrt x}-\frac1{x\log2}
   \ge\frac1{8\sqrt x}-\frac2x\ge0,\qquad
 \frac14-\frac6{x\log2}\ge\frac14-\frac{12}x\ge0.
 \tag{FTS.5}
\]
Here \(\log2=\int_1^2dt/t\ge1/2\), and the displayed ranges imply
the final inequalities.  This proves (FTS.4) on the full half-line,
not merely at integer or dyadic sample points.
In particular, writing \(s=\sqrt n\ge256\), (FTS.3)--(FTS.4) give
\[
 \log_2(F/2)-120\log_2n
 \ge\frac{n-s}{2}-1-(30s-120)
 =\frac{s(s-61)}2+119>0.
 \tag{FTS.6}
\]
Thus \(a=F/2\ge n^{120}\).

We next use the exact published alternating split-character formula.
If the cycle type of \(x\) differs from \((d_1,\ldots,d_r)\), then
\[
                 \chi_W(x)=\frac12\chi^\lambda(x).
                 \tag{FTS.7}
\]
At either of the two exceptional alternating classes of that
cycle type,
\[
 \chi_W(x)=\frac12\left(\varepsilon
                \mathbin{\pm}\sqrt{\varepsilon\prod_i d_i}\right),
 \qquad \varepsilon=(-1)^{(n-r)/2}.
 \tag{FTS.8}
\]
The choice of sign depends on the class and constituent and will
not be needed.  Exact primary locators are Balcerzyk,
\cite{FTS:Balcerzyk1997}, formulas (3)--(4), pp.~148--149, and the
diagonal-hook character value on p.~151.  The same restriction and
split-value formulas are restated in the published research paper
\cite{FTS:YangStaroletov2021}, Propositions 3--4, author pp.~5--6,
with the branching rule in Proposition 5.
One may also obtain the sign in (FTS.8) by Murnaghan--Nakayama:
remove the unique rim hook of largest diagonal-hook length at
each stage; the remaining inner diagram is self-conjugate, and
the total leg length is \(\sum_i(d_i-1)/2=(n-r)/2\).

For an exceptional element, (FTS.2)--(FTS.3) yield
\[
 \frac{|\chi_W(x)|}{a}
 \le\frac{1+\sqrt{\prod_i d_i}}{F}
 \le 2^{\,1-n/2+\sqrt n(1+\log_2n)/2}.
 \tag{FTS.9}
\]
Indeed \(\prod_i d_i\le n^r\) and \(r\le\sqrt n\).
Using the two estimates in (FTS.4) in sequence, the exponent in
(FTS.9) is at most
\[
 1-\frac{3n}{8}
 =\left(1-\frac n4\right)-\frac n8
 \le-6\log_2n-\frac n8
 <-6\log_2n .
 \tag{FTS.10}
\]
Thus every exceptional value has absolute ratio strictly less
than \(n^{-6}\).  Moreover its cycle lengths are distinct odd
integers, so it has at most one fixed letter.  If \(f(x)=1\), then
\(n^{-6}\le n^{-3}=(f(x)/n)^3\); if \(f(x)=0\), the required
fixed-point-free estimate holds directly.

For every nonexceptional \(x\), (FTS.7) shows that
\(|\chi_W(x)|/a=|\chi^\lambda(x)|/F\).
The already checked effective MSP consequence used in HD-TW02
and HD-FT02 applies because \(n\ge65536>4000\) and
\(F\ge2n^{120}\).  It gives the cubic bound for positive fixed-point
count and the stronger fixed-point-free bound
\[
                  |\chi^\lambda(x)|/F<F^{-1/20}
                                      \le n^{-6}.
                  \tag{FTS.11}
\]
The identity ratio is exactly one.  This proves all of (FTS.1).

Finally the bound (FTS.1) satisfies the weaker error allowance
\(2n^{-6}\) in (FT.3).  All hypotheses of HD-FT01 now hold for the
actual quotient \(A_n\) and actual tensor factorization, so that
the source is impossible for \(n>3500000\).
This last step uses actual quotient curves and their genera, not
a claim that the projected factor has twelve branch points.
\end{proof}

\paragraph{Compatible-normalizer precision.}
For \(n\ge7\), the known automorphism description of \(A_n\) identifies
its outer automorphism with conjugation by an odd permutation.
That automorphism exchanges the two inequivalent split constituents
(the same Proposition 3(ii) cited above).
It cannot stabilize the projective class of either one:
if an intertwiner implemented it up to scalar, the scalar discrepancy
would be a linear character of the perfect group \(A_n\), hence trivial.
Schur's lemma then shows that the compatible linear normalizer is
scalar times the represented \(A_n\).  This uses the previously checked
ordinary alternating automorphism input; it asserts no new
classification theorem.  Thus the \(A_n\)-quotient interface above
also covers a factor with its full compatible scalar normalizer.
As always, an actual quotient and tensor factorization must be supplied;
they are not inferred from a numerical spectrum alone.

\paragraph{Combined range.}
HD-FT02 and HD-FT03 together treat every ordinary irreducible
alternating constituent at degree at least \(n^{120}\).
In particular all split constituents disappear already for
\(n>3500000\), and for
\(n\ge241^{482}\) every surviving ordinary symmetric or alternating
factor has a nonsplit partition label with row or column tail at
most \(240\).  Arbitrary factorwise tensor partners are still allowed.
This does not settle these finitely many tail shapes in unbounded
natural degree, and does not assert a finite catalogue of all
projective factor types.

\section{Fixed tails of size at least seven with arbitrary tensor partners}
\label{sec:factorwise-small-tail}

\paragraph{Status and exact scope.}
HD-FT04--06, reserved question Q2de, is a new derivation by
\texttt{/root/g34\_eightfolds}.  The coefficient-polynomial input is the
already checked D-AFT01, and the actual triple-cover genus argument is
the root's HD-FT01.  The new estimate controls the number of inertia
entries at which the polynomial comparison has an error; it does not
transfer twelve branch marks to a tensor factor.
Complete author and separate root whole-proof readings pass. Root also
read the complete source/status record, the entire supplementary checker
and its first successful output; no fresh root execution is claimed.

Use exactly the actual quotient and factorization of (FT.1)--(FT.2):
a connected finite-group cover of the sphere, an actual quotient onto
$P=A_n$ or $S_n$, an actual coefficient factorization $V=W\otimes U$
up to the permitted scalar factors, and genuine inertia $T_i$ with
\[
 e=\dim V,\qquad
             \sum_i\operatorname{rank}(1-T_i)=2e+6.
 \tag{FTL.1}
\]
The tensor partner $U$, its degree and coefficient field are unrestricted.
No primitivity, factorwise tensor-indecomposability, transitive equal-factor
core, or twelve-mark condition on the projected permutations is assumed.
An abstract spectrum without the actual quotient cover does not suffice.

Fix $\mu\vdash k$, $k\ge7$, and let $W=S^{(n-k,\mu)}$, or its
irreducible ordinary restriction to $A_n$.  Transposing the full partition
only twists by sign on $S_n$, which changes none of the absolute-character
or projective-support bounds below.  The threshold will imply $n>2k$,
so the indicated partition and its irreducible alternating restriction
are in the stable, non-self-conjugate range.

\subsection*{HD-FT04. A localized cubic character bound}

Write the exact Garsia--Goupil polynomial, in the notation of D-AFT01, as
\[
 \chi_W(x)=q_\mu(X_1(x),\ldots,X_k(x)),\quad
 q_\mu=C X_1^k+R,\quad C=f^\mu/k!,\quad
 L=\|R\|_1,\quad M=1+L/C.
 \tag{FTL.2}
\]
Here $\|R\|_1$ is the sum of the absolute values of its ordinary-monomial
coefficients.  The published formula and its proof give weighted degree
at most $k$, with weight $i$ on $X_i$; the sole ordinary-degree-$k$
monomial is $C X_1^k$.  Thus $R$ has ordinary degree at most $k-1$.
These are precisely the coefficient premises already proved in D-AFT01,
using \cite[Proposition~I.1 and its complete proof, pp.~9--11]{HD:GarsiaGoupil2009}.

\begin{lemma}[HD-FT04]
Suppose $k\ge7$ and
\[
                         n\ge2^{k+1}M.
 \tag{FTL.3}
\]
Put $a=\dim W=q_\mu(n,0,\ldots,0)$ and
$\varepsilon=2M/n^4$.  Then $a\ge2n^3$, and for every permutation
$x$, writing $f=X_1(x)$,
\[
 \frac{|\chi_W(x)|}{a}\le
 \begin{cases}
    (f/n)^3,& f\ge n/2,\\
    (f/n)^3+\varepsilon,& f<n/2.
 \end{cases}
 \tag{FTL.4}
\]
These bounds include the identity and all powers of every permutation.
\end{lemma}

\begin{proof}
Set $t=f/n$ and $m=n-f$.  All cycle counts are at most $n$.
The degree premise gives
\[
 |R(n,0,\ldots,0)|\le Ln^{k-1},\qquad
        a\ge Cn^k(1-M/n)\ge\tfrac12 Cn^k.
 \tag{FTL.5}
\]
We used $n\ge2^{k+1}M\ge256M$.  Also $C\ge1/k!$ and
$k!\le2^{k(k-1)/2}$, since $j\le2^{j-1}$ for each $j\ge1$.
Consequently
\[
 n^{k-3}\ge2^{(k+1)(k-3)}\ge4k!,
 \qquad
 (k+1)(k-3)-2-k(k-1)/2=(k-5)(k+2)/2>0.
 \tag{FTL.6}
\]
Equations (FTL.5)--(FTL.6) prove $a\ge2n^3$.
They also give $n>2k$, so all stable-character premises hold.

Suppose first that $f\ge n/2$.  The crude remainder bound gives
\[
 q_\mu(f,X_2,\ldots)\ge Cf^k-Ln^{k-1}>0,
 \tag{FTL.7}
\]
because $n/2^k\ge2M>L/C$.
For a pure monomial $f^j$ of $R$,
$|f^j-n^j|\le jmn^{j-1}\le k m n^{k-2}$; its constant term
has zero difference.  Every mixed monomial has ordinary degree at
most $k-1$ and contains a non-$f$ cycle count at most $m$, so its
value is at most $m n^{k-2}$ and its value at $(n,0,\ldots)$ is zero.
Summing coefficientwise yields
\[
 |R(f,X_2,\ldots)-R(n,0,\ldots)|\le k Lm n^{k-2}.
 \tag{FTL.8}
\]
Since $1-t^3\le3m/n$, the difference between the desired character
bound and its leading term is controlled as follows:
\[
 \begin{split}
 q_\mu-a t^3
 &\le-Cn^k(t^3-t^k)+(k+3)Lm n^{k-2}\\
 &\le-Cm n^{k-1}/8+(k+3)Lm n^{k-2}\le0.
 \end{split}
 \tag{FTL.9}
\]
Indeed $t\ge1/2$ and $k\ge4$ give
$t^3-t^k=t^3(1-t^{k-3})\ge(1-t)/8$.
Moreover $2^{k+1}\ge8(k+3)$ for $k\ge7$, so (FTL.3) implies
$n\ge8(k+3)L/C$.  When $m=0$ the difference is exactly zero.
Together with (FTL.7), this proves the first line of (FTL.4),
including its absolute value.

Now assume $0\le t<1/2$.  Consider any monomial of $R$ and let $j$
be its exponent of $f$, and $b$ the sum of its other exponents.
Then $j+2b\le k$ and $j+b\le k-1$.
After division by $n^k$, its absolute value is at most
$t^j n^{j+b-k}$.  If $j\ge3$, this is at most $t^3/n$.
If $j=0$, it is at most $n^{-\lceil k/2\rceil}\le n^{-4}$.
If $j=1$ or $2$, it is at most $t/n^3$ or $t^2/n^3$, respectively.
The following exact nonnegative identities control those two cases:
\[
 \begin{split}
 \frac{t^3}{n}+\frac2{n^4}-\frac{3t}{n^3}
     &=\frac{(nt-1)^2(nt+2)}{n^4}\ge0,\\
 \frac{2t^3}{n}+\frac1{n^7}-\frac{3t^2}{n^3}
     &=\frac{(n^2t-1)^2(2n^2t+1)}{n^7}\ge0.
 \end{split}
 \tag{FTL.10}
\]
Thus every normalized remainder monomial is at most
$t^3/n+n^{-4}$, and coefficientwise summation gives
\[
 \frac{|q_\mu|}{a}
 \le\frac{t^k+(M/n)t^3+M/n^4}{1-M/n}
 \le t^3+\frac{2M}{n^4}.
 \tag{FTL.11}
\]
For the last inequality, $t^k\le t^3/16$ for $k\ge7$,
and $M/n\le1/4$.  Hence the coefficient of $t^3$ is at most
$(1/16+1/4)/(3/4)=5/12<1$, while the last coefficient is at most
$4M/(3n^4)\le2M/n^4$.  This proves the second line of (FTL.4).
The argument depends only on cycle counts, so it applies to every power.
\end{proof}

\subsection*{HD-FT05. At most ten error-bearing inertia entries}

\begin{theorem}[HD-FT05]
Under the actual quotient and factorization (FTL.1), the factor $W$
described above is impossible whenever
\[
                 n>3500000,\qquad n\ge2^{k+1}M,
                 \qquad k\ge7.
 \tag{FTL.12}
\]
The conclusion allows every original imaginary quadratic field and
arbitrary tensor partners and scalar twists.
\end{theorem}

\begin{proof}
Let $x_i\in P$ be the actual projected inertia.  Call an index $i$
exceptional for this proof if $f(x_i)<n/2$, and let $b$ count these
indices.  This terminology records only where the error in (FTL.4)
can appear; it asserts no new conjugacy classification.
If an index is not exceptional, every power of $x_i$ still fixes
at least $n/2$ letters, so every power has zero error in (FTL.4).
At an exceptional index the Fourier average has error at most
$\varepsilon=2M/n^4$.

For any $x$ of finite order $h$, cyclic Fourier inversion bounds
the multiplicity of every target eigenvalue by
$h^{-1}\sum_{j=0}^{h-1}|\chi_W(x^j)|$.  Since $0\le f(x^j)/n\le1$
and its cube is at most itself, HD-FT04 gives at every exceptional index
\[
 \frac{\nu(W(x))}{a}
   \ge1-\frac1h\sum_j\frac{f(x^j)}n-\varepsilon
   =\frac{\operatorname{ind}(x)}n-\varepsilon
   >\frac14-\varepsilon>\frac15.
 \tag{FTL.13}
\]
The equality is the natural orbit-counting formula.  A permutation with
$n-f$ moving letters has index at least $(n-f)/2$; the exceptional
inequality is strict because $f<n/2$.
For the final strict inequality note that $n\ge256M$, so
\[
        \varepsilon\le\frac1{128n^3}<\frac1{20},
        \qquad
        2+\frac6e\le2+\frac3{n^3}<\frac{11}{5}.
 \tag{FTL.14}
\]
Here $e\ge a\ge2n^3$ by HD-FT04, and $n\ge256$.
Spectral convolution for the actual tensor factorization gives
$\operatorname{rank}(1-T_i)/e\ge\nu(W(x_i))/a$ for every $i$,
with any scalar twists.  Eleven exceptional entries would therefore
make (FTL.1) strictly greater than $11e/5$, contrary to (FTL.14).
Thus
\[
                              b\le10.
 \tag{FTL.15}
\]
No lower bound on the number of genuine projected marks was used.

Let $\Omega=\{1,\ldots,n\}^3$, with repetitions allowed.
The exact orbit-counting formula and (FTL.4), now retaining the cubes,
give for each inertia
\[
 \frac{\operatorname{rank}(1-T_i)}e
 \ge\frac{\nu(W(x_i))}{a}
 \ge\frac{\operatorname{ind}(x_i,\Omega)}{n^3}
                  -\varepsilon\,\mathbf1_{f(x_i)<n/2}.
 \tag{FTL.16}
\]
For $x_i=1$ both terms on the right are zero, so no issue is caused
by original marks disappearing in the quotient.

The geometric step of HD-FT01 applies to this same actual quotient.
There are five $P$-orbits on $\Omega$, of genera
$\gamma_1,\gamma_2,\gamma_2,\gamma_2,\gamma_3$.
The unordered-three-subset quotient has genus at least two for
$n>3500000$, by the published Monderer--Neftin theorem in exactly
the source scope recorded in HD-FT01.\cite[Theorem~1.2 and
Remark~2.8]{HD:MondererNeftin2022}
Indeed its stabilizer has order at most $6(n-3)!<|A_{n-2}|$ and
cannot contain either required alternating subgroup.
The actual forget-order morphism has degree six for $A_n$ as well as
$S_n$, so $\gamma_3\ge7$ by Riemann--Hurwitz.  The five actual
quotient covers therefore satisfy
\[
 \sum_i\operatorname{ind}(x_i,\Omega)
       =2(n^3-5)+2(\gamma_1+3\gamma_2+\gamma_3)
       \ge2n^3+4.
 \tag{FTL.17}
\]
This is a quotient-cover statement, not an assertion about a formal
list of permutations.

Summing (FTL.16) and using (FTL.15)--(FTL.17) yields
\[
 \frac1e\sum_i\operatorname{rank}(1-T_i)
 \ge2+\frac4{n^3}-\frac{20M}{n^4}
 >2+\frac3{n^3}\ge2+\frac6e.
 \tag{FTL.18}
\]
The strict inequality follows from $n\ge256M>20M$.
This contradicts (FTL.1) and proves the theorem.
\end{proof}

\subsection*{HD-FT06. The remaining unbounded ordinary tails have size at most six}

\begin{corollary}[HD-FT06]
In the actual factorwise tensor setting, suppose an ordinary
symmetric or alternating irreducible factor has natural degree
\[
                              n\ge241^{482}.
 \tag{FTL.19}
\]
Then its surviving partition label is extendible, and must satisfy
\[
                  \min\{n-\lambda_1,n-\lambda'_1\}\le6.
 \tag{FTL.20}
\]
The conclusion includes arbitrary factorwise partners, unequal
partner dimensions and all scalars.  It does not assert that any
of the remaining small-tail sizes occurs.
\end{corollary}

\begin{proof}
HD-FT02 removes all extendible labels with both tails at least241,
and HD-FT03 removes the split alternating constituents already for
$n>3500000$.  It remains to consider a label, transposed if necessary,
of the form $(n-k,\mu)$ with $7\le k\le240$.
The coefficient estimate (TW.23)--(TW.24) gives the uniform bound
\[
 M\le1+k!(k+1)2^{k-1}
        \le M_*:=1+241\cdot240!\cdot2^{239}.
 \tag{FTL.21}
\]
This is the accepted fixed-subset/branching calculation and does not
require enumerating all partitions of 240.
The following elementary comparison suffices:
\[
 2^{k+1}M\le2^{241}M_*
       \le2^{481}241^{241}<241^{482}.
 \tag{FTL.22}
\]
For the middle inequality use $M_*=1+241!2^{239}\le2^{240}241!$
and $241!\le241^{241}$.  For the strict final inequality use
$241^{241}>4^{241}=2^{482}>2^{481}$.
Also $241^{482}>3500000$.  Hence every such label satisfies
(FTL.12), and HD-FT05 excludes it.  This leaves exactly the necessary
small-tail range (FTL.20).
\end{proof}

\paragraph{Boundary and provenance.}
The new ingredient is the character-polynomial comparison with an error
supported on at most ten actual inertia entries.  The proof uses the
already accepted character-polynomial formula and quotient-genus theorem,
not an effective constant inferred from an asymptotic genus estimate.
It removes the previously unbounded fixed tails seven through 240
in the factorwise tensor-decomposable gate.  Tails zero through six,
other projective factor types and the actual quotient/factorization
interface remain separate.  No factor is assumed to be a sixfold
source on its own, and no independent positive cycle is constructed.

\section{A qualitative cutoff for tails of size five and six}
\label{sec:factorwise-five-six-tail}

\paragraph{Status and exact scope.}
HD-FT07--08, reserved question Q2dg, is a new derivation by
\texttt{/root/g34\_eightfolds}.  It uses the already checked character
polynomial of D-AFT01, the actual quotient interface of HD-FT01, and
the published arbitrary-fixed-genus theorem of Monderer--Neftin.
The cutoff below is qualitative: no numerical value of the published
constant \(N_g\) for the new, large values of \(g\) is asserted.
In particular it is not substituted into the effective FQ bound.
Complete author and separate root whole-proof readings pass. Root also
read the full source/status record, entire checker and first output;
the published arbitrary-genus passages were freshly reopened.
No numerical arbitrary-genus cutoff or new root execution is claimed.

Retain exactly (FT.1)--(FT.2): an actual connected finite-group cover
of the sphere, an actual quotient onto \(P=A_n\) or \(S_n\), an actual
factorization \(V=W\otimes U\) up to scalars, and genuine inertia with
\[
 e=\dim V,\qquad
 \sum_i\operatorname{rank}(1-T_i)=2e+6.
 \tag{FT56.1}
\]
The tensor partner and its field and degree are arbitrary.  No twelve-mark
condition is imposed on the projected permutations.  Let
\(\mu\vdash k\), where \(k\in\{5,6\}\), and put
\(W=S^{(n-k,\mu)}\), or its irreducible ordinary restriction to \(A_n\).
Transposition of the full partition, hence a sign twist in the symmetric
case, is allowed.  The bounds use absolute characters and projective
support, so this changes none of them.

Write the exact stable character polynomial as
\[
 q_\mu=C X_1^k+R,\qquad C=f^\mu/k!,\qquad
 L=\|R\|_1,\qquad M=1+L/C.
 \tag{FT56.2}
\]
Its weighted degree is at most \(k\), with weight \(i\) on \(X_i\);
\(R\) has ordinary degree at most \(k-1\).  These are the precise
premises of D-AFT01, from
\cite[Proposition~I.1 and complete proof, pp.~9--11]{HD:GarsiaGoupil2009}.
In particular \(M\ge1\).

\subsection*{HD-FT07. A cubic comparison with a cubic-order localized error}

\begin{lemma}[HD-FT07]
Suppose \(k=5\) or \(6\), and \(n\ge256M\).  Put
\[
 a=\dim W,\qquad
 \varepsilon=\frac{2M(M+1)}{n^3}.
 \tag{FT56.3}
\]
Then \(a\ge2n^3\), and for every permutation \(x\), with
\(f=X_1(x)\), one has
\[
 \frac{|\chi_W(x)|}{a}\le
 \begin{cases}
 (f/n)^3,&f\ge n/2,\\
 (f/n)^3+\varepsilon,&f<n/2.
 \end{cases}
 \tag{FT56.4}
\]
This includes every power of every inertia and every target eigenvalue
in the Fourier consequence used below.
\end{lemma}

\begin{proof}
Set \(t=f/n\) and \(m=n-f\).  The ordinary degree of \(R\) gives
\[
 a\ge Cn^k-Ln^{k-1}\ge Cn^k(1-M/n)\ge\frac{n^k}{2k!}.
 \tag{FT56.5}
\]
Since \(n\ge256\), \(n^{k-3}\ge4k!\) for both \(k=5,6\),
which proves \(a\ge2n^3\).  Also \(n>2k\), as required for the
stable partition and its non-self-conjugate alternating restriction.

For \(f\ge n/2\), the argument of HD-FT04 applies with these same
coefficients, and we repeat its inequalities to specify the threshold.
The character is positive, since
\(Cf^k-Ln^{k-1}>0\) follows from \(n>2^kL/C\).
For the remainder,
\[
 |R(f,X_2,\ldots)-R(n,0,\ldots)|
       \le kLm n^{k-2}.
 \tag{FT56.6}
\]
A pure \(f^j\) term satisfies the derivative estimate, and a mixed
term has a non-\(f\) cycle count at most \(m\); these observations
prove (FT56.6) term by term.  Furthermore \(1-t^3\le3m/n\).
Consequently
\[
 q_\mu-a t^3
 \le -Cn^k(t^3-t^k)+(k+3)Lm n^{k-2}\le0.
 \tag{FT56.7}
\]
Indeed \(t^3-t^k\ge(1-t)/8\) for \(1/2\le t\le1\) and \(k\ge4\),
and \(n\ge256M\ge8(k+3)L/C\).  This proves the first case of
(FT56.4), including \(f=n\).

Suppose \(f<n/2\).  For a monomial of \(R\), write \(j\) for the
exponent of \(f\) and \(b\) for the sum of the other exponents.
Then \(j+2b\le k\) and \(j+b\le k-1\).
After division by \(n^k\), its absolute value is at most
\(t^3/n\) if \(j\ge3\), at most \(t^2/n^2\) if \(j=2\),
at most \(t/n^2\) if \(j=1\), and at most \(n^{-3}\) if \(j=0\).
Thus the following deliberately coarse simultaneous bound holds:
\[
 \frac{|R(f,X_2,\ldots)|}{Cn^k}
 \le M\left(\frac{t^3}{n}+\frac{t^2}{n^2}
                   +\frac{t}{n^2}+\frac1{n^3}\right).
 \tag{FT56.8}
\]
For \(z\ge0\) and \(M\ge1\), the elementary arithmetic--geometric
mean inequality gives
\[
 3Mz\le z^3+2M^{3/2}\le z^3+2M^2.
 \tag{FT56.9}
\]
Apply this with \(z=nt\).  Because \(t^2\le t/2\) and
\(t^k\le t^3/4\) on \(0\le t<1/2\), (FT56.8)--(FT56.9) imply
\[
 \frac{|q_\mu|}{Cn^k}
 \le\left(\frac34+\frac Mn\right)t^3
                  +\frac{M(M+1)}{n^3}.
 \tag{FT56.10}
\]
Here the two linear/quadratic terms contribute at most
\(t^3/2+M^2/n^3\).  Divide by the lower bound
\(a/(Cn^k)\ge1-M/n\).
Since \(M/n\le1/256\le1/8\), the cubic coefficient is at most one,
and \((1-M/n)^{-1}\le2\).  This proves the second case of (FT56.4).
\end{proof}

\subsection*{HD-FT08. Arbitrary partners and a qualitative natural-degree cutoff}

For each \(\mu\) as above choose the integer
\[
 G_\mu=\left\lceil2M(M+1)+2\right\rceil.
 \tag{FT56.11}
\]
Let \(N_g\) be the constant in the published
\cite[Theorem~1.2]{HD:MondererNeftin2022}, at the fixed genus bound
\(g=G_\mu\), and set
\[
 N_\mu^{(56)}
   =\max\{\lceil256M\rceil,N_{G_\mu},256\}.
 \tag{FT56.12}
\]
Only existence of \(N_{G_\mu}\) is used; the known explicit constant
for genus one is not asserted to work at \(G_\mu\).

\begin{theorem}[HD-FT08]
No actual source (FT56.1) with the indicated factor \(W\) exists when
\(n>N_\mu^{(56)}\).  This holds with arbitrary scalar factors,
arbitrary coefficient fields and arbitrary factorwise tensor partners.
There are only eighteen partitions of sizes five or six, so taking
the maximum of these eighteen constants gives a finite qualitative
cutoff for both tails.
\end{theorem}

\begin{proof}
Call an inertia image \(x_i\) bad if it fixes fewer than \(n/2\)
natural letters.  If \(x_i\) is not bad, every power fixes at least
the same letters, so the first case of (FT56.4) applies to every power.
Cyclic Fourier inversion therefore gives, for every image,
\[
 \frac{\operatorname{rank}(1-T_i)}e
 \ge\frac{\nu(W(x_i))}{a}
 \ge\frac{\operatorname{ind}(x_i,\{1,\ldots,n\}^3)}{n^3}
                    -\varepsilon\,{\bf1}_{\{i\ {\rm bad}\}}.
 \tag{FT56.13}
\]
The first inequality is tensor spectral convolution; it handles
all target scalars and does not require reality of a chosen inertia.

There are at most ten bad entries.  To see this without assigning
any mark count to \(W\), use \(u^3\le u\) for \(0\le u\le1\)
in the same Fourier average.  For a bad image,
\[
 \frac{\nu(W(x_i))}{a}
 \ge\frac{\operatorname{ind}(x_i)}n-\varepsilon
 >\frac14-\varepsilon>\frac15.
 \tag{FT56.14}
\]
Indeed \(\operatorname{ind}(x_i)\ge(n-f(x_i))/2>n/4\), and
\[
 \varepsilon\le\frac{4M^2}{n^3}
       \le\frac1{16384n}<\frac1{20}.
 \tag{FT56.15}
\]
On the other hand \(e\ge a\ge2n^3\) implies
\(2+6/e\le2+3/n^3<11/5\).
Eleven bad entries would strictly violate this actual total budget.

Consider the actual quotient curve associated with the stabilizer
of an unordered three-set in \(P\).  Its stabilizer has order at
most \(6(n-3)!\); it cannot contain \(A_{n-2}\), of order
\((n-2)!/2\), because \(n\ge256>14\).
It also cannot contain \(A_{n-1}\).
Thus Monderer--Neftin's Theorem~1.2 implies that its genus
\(\widehat\gamma_3\) exceeds \(G_\mu\) when \(n>N_{G_\mu}\).
This application uses the actual \(P\)-cover and its degree-\(n\)
natural quotient, not a proposed numerical branch profile.
The ordered-distinct-triple quotient has degree six over this
unordered-triple curve.  Consequently its genus satisfies
\[
 \gamma_3\ge6\widehat\gamma_3-5\ge6G_\mu+1.
 \tag{FT56.16}
\]

As in HD-FT01, the action on ordered triples with repetitions
has five orbits, with quotient genera
\(\gamma_1,\gamma_2,\gamma_2,\gamma_2,\gamma_3\).
Riemann--Hurwitz on these actual connected covers gives
\[
 \sum_i\operatorname{ind}(x_i,\{1,\ldots,n\}^3)
 =2(n^3-5)+2(\gamma_1+3\gamma_2+\gamma_3)
 \ge2n^3+12G_\mu-8.
 \tag{FT56.17}
\]
Summing (FT56.13), using the at-most-ten bound and (FT56.11), yields
\[
 \begin{split}
 2+\frac6e
 &\ge 2+\frac{12G_\mu-8-20M(M+1)}{n^3}\\
 &\ge 2+\frac{4M(M+1)+16}{n^3}
  >2+\frac3{n^3}\ge2+\frac6e,
 \end{split}
 \tag{FT56.18}
\]
a contradiction.  This proves the theorem.
\end{proof}

\paragraph{Remaining gate and source precision.}
Combined with HD-FT01--06, this removes unbounded tails of sizes
five and six from that actual factorwise source setup.  The only
unbounded ordinary extendible alternating labels left by these
particular arguments have tail size at most four.
This is a qualitative reduction: the finitely many possible intervals
between the existing effective cutoff and \(N_\mu^{(56)}\) have
not been enumerated, and no effective replacement for those constants
has been claimed.  No conclusion about a general complex coefficient
factor without the actual quotient (FT.2) is inferred.
The published arbitrary-genus theorem, not an unrefereed numerical
genus-growth estimate, is the new source input.

\section{Tail four and the quadratic growth of the triple quotient genus}
\label{sec:factorwise-four-tail}

\paragraph{Status and exact scope.}
These are root statements HD-FT09--10, reserved question Q2dh.
The complete proof and source application have passed root and separate
independent whole-proof review. The initial text, two precision edits
and the complete independent audit are retained. No finite certificate
was needed, executed or claimed for these universal inequalities.
The local character estimate is effective. The final exclusion is qualitative,
because the cited genus-growth constants have not been evaluated.
It is not an effective detector-degree bound and is not substituted into FQ.

Use the actual quotient and tensor factorization (FT.1)--(FT.2): a connected
finite-group cover of the sphere, an actual quotient onto $P=A_n$ or $S_n$,
and an actual factor $V=W\otimes U$ up to scalars, where the genuine inertia
operators satisfy
\[
 e=\dim V,\qquad \sum_i\operatorname{rank}(1-T_i)=2e+6.
 \tag{FT4.1}
\]
The partner $U$, its degree and its field are arbitrary. Let $\mu\vdash4$,
and let $W=S^{(n-4,\mu)}$ or its irreducible ordinary restriction to $A_n$.
Transposition of the partition is permitted: its sign twist does not affect
absolute character ratios or the maximum-eigenspace support.
As in D-AFT01, write
\[
 q_\mu=CX_1^4+R,\qquad C=f^\mu/24,\qquad
 L=\|R\|_1,\qquad M=1+L/C\ge1.
 \tag{FT4.2}
\]
The exact character polynomial has weighted degree at most four, assigning
weight $j$ to $X_j$, and $R$ has ordinary degree at most three.
These are the accepted consequences of Garsia--Goupil,
Proposition I.1 and its complete proof.\cite{HD:GarsiaGoupil2009}

\subsection*{HD-FT09. A localized quadratic-order error}

\begin{lemma}[HD-FT09]
If $n\ge256M$, then $a=\dim W\ge2n^3$ and, writing $f=X_1(x)$,
\[
 \frac{|\chi_W(x)|}{a}\le
 \begin{cases}
 (f/n)^3,&f\ge n/2,\\
 (f/n)^3+4M^2/n^2,&f<n/2.
 \end{cases}
 \tag{FT4.3}
\]
The estimate holds for every permutation, hence for every power of every
projected inertia.
\end{lemma}

\begin{proof}
Put $t=f/n$ and $m=n-f$. The ordinary-degree bound gives
\[
 a\ge Cn^4(1-M/n)\ge n^4/48\ge2n^3,
 \tag{FT4.4}
\]
because $C\ge1/24$ and $n\ge256>96$. In particular $n\ge256>9$, so the
stable partition is non-self-conjugate and all character premises hold.

If $f\ge n/2$, the character is positive: the leading term is at least
$Cn^4/16$, whereas the absolute remainder is at most $Ln^3$, and
$n\ge256M>16L/C$.
For a pure monomial in $f$, the difference of its value at $f$ and $n$
is at most its degree times $mn^{\deg-1}$. Every other monomial has a
non-$f$ cycle count at most $m$. Thus
\[
 |R(f,X_2,\ldots)-R(n,0,\ldots)|\le4Lm n^2.
 \tag{FT4.5}
\]
Since $1-t^3\le3m/n$ and $t^3-t^4\ge(1-t)/8$ in this range,
\[
 q_\mu-a t^3\le-Cmn^3/8+7Lm n^2\le0,
 \tag{FT4.6}
\]
using $n\ge256M>56L/C$. This includes $m=0$, when the difference is
exactly zero, and proves the first case with its absolute value.

Suppose $0\le t<1/2$. For a remainder monomial, let $j$ be the exponent
of $f$ and $b$ the sum of all other exponents. Then $j+2b\le4$ and
$j+b\le3$. After division by $n^4$, the monomial is at most $t^3/n$
if $j\ge3$, at most $t^2/n$ if $j=2$, at most $t/n^2$ if $j=1$,
and at most $n^{-2}$ if $j=0$. Coefficientwise summation therefore gives
\[
 \frac{|R|}{Cn^4}
 \le M\left(\frac{t^3}{n}+\frac{t^2}{n}
                       +\frac{t}{n^2}+\frac1{n^2}\right).
 \tag{FT4.7}
\]
The elementary arithmetic--geometric mean inequalities
$3Az^2\le2z^3+A^3$ and $3Az\le z^3+2A^{3/2}$, applied with
$A=4M$ and $z=nt$, give
\[
 \frac{Mt^2}{n}\le\frac{t^3}{6}+\frac{16M^3}{3n^3},
 \qquad
 \frac{Mt}{n^2}\le\frac{t^3}{12}+\frac{4M^2}{3n^3}.
 \tag{FT4.8}
\]
In the second inequality we used $M^{3/2}\le M^2$.
Since $t^4\le t^3/2$, equations (FT4.7)--(FT4.8) yield
\[
 \frac{|q_\mu|}{Cn^4}
 \le\left(\frac34+\frac Mn\right)t^3
       +\frac{M}{n^2}+\frac{16M^3+4M^2}{3n^3}.
 \tag{FT4.9}
\]
Divide by $a/(Cn^4)\ge1-M/n$.
The cubic coefficient is at most one, since $M/n\le1/256<1/8$.
Moreover
\[
 M+\frac{16M^3+4M^2}{3n}
 \le M^2+\frac{20M^3}{3\cdot256M}\le2M^2,
 \qquad (1-M/n)^{-1}\le2.
 \tag{FT4.10}
\]
The remaining term is consequently at most $4M^2/n^2$.
This proves the second case of (FT4.3).
\end{proof}

\subsection*{HD-FT10. A qualitative cutoff from a published genus-growth estimate}

We state precisely the additional published input. In the proof of
Monderer--Neftin Theorem 2.5, equation (2.2), there are positive constants
$c,d$ such that, for the genera $g_j$ of the quotients by unordered
$j$-set stabilizers of an actual $A_n$ or $S_n$ sphere cover,
\[
 g_j-g_{j-1}>(cn-dj^{15})
               \frac{\binom nj}{\binom n2}\qquad(3\le j\le\lfloor n/2\rfloor).
 \tag{FT4.11}
\]
The $j\ge3$ range applies also to the exceptional ramification types
in their Table 2.1.\cite[proof of Theorem~2.5, equation~(2.2)]{HD:MondererNeftin2022}
We may include an initial threshold $n_0\ge20$ in this asymptotic
form without strengthening the source. No numerical $c,d,n_0$ is claimed.
This is the estimate expressly invoked in the published proof, read in
the author's revised 18 December 2020 text. We do not claim a fresh
proof of its cited Neftin--Zieve precursor, nor an independently computed
effective version of those constants.

For the fixed $\mu$, define the integer, depending on these unevaluated
constants,
\[
 N_\mu^{(4)}=
 \left\lceil\max\left\{n_0,256M,4,
                  \frac{2d\,3^{15}}c,\frac{80M^2}c\right\}\right\rceil.
 \tag{FT4.12}
\]

\begin{theorem}[HD-FT10]
No source in the actual factorwise setup (FT4.1) with the indicated
tail-four factor exists for $n>N_\mu^{(4)}$.
Thus all five tail-four labels have a finite qualitative natural-degree
cutoff, with arbitrary tensor partners, scalars and coefficient fields.
This statement does not specify a numerical cutoff or an effective
enumeration range.
\end{theorem}

\begin{proof}
Put $\varepsilon=4M^2/n^2$. Call an inertia image bad if it fixes fewer
than $n/2$ letters. A nonbad image has only nonbad powers. Fourier inversion,
the cube fixed-point formula and tensor spectral convolution give
\[
 \frac{\operatorname{rank}(1-T_i)}e
 \ge\frac{\operatorname{ind}(x_i,\{1,\ldots,n\}^3)}{n^3}
                -\varepsilon\mathbf1_{\{i\text{ bad}\}}.
 \tag{FT4.13}
\]
Identity quotient images cause no problem: their index and error are zero.
There are at most ten bad entries. Indeed the same Fourier bound, with
$u^3\le u$ for $0\le u\le1$, gives at a bad entry
\[
 \frac{\nu(W(x_i))}{a}
 \ge\frac{\operatorname{ind}(x_i)}n-\varepsilon
 >\frac14-\varepsilon>\frac15,
 \qquad \varepsilon\le\frac1{16384}<\frac1{20}.
 \tag{FT4.14}
\]
Here $\operatorname{ind}(x_i)\ge(n-f(x_i))/2>n/4$.
Eleven such entries would exceed $11e/5$, whereas (FT4.4) makes the
actual total budget $2+6/e\le2+3/n^3<11/5$.

Let $\widehat\gamma_3$ be the genus of the unordered-triple quotient.
Apply (FT4.11) with $j=3$ and use $g_2\ge0$. The threshold gives
\[
 \widehat\gamma_3>
   (cn-d\,3^{15})\frac{n-2}{3}
       \ge\frac{cn}{2}\frac n6=\frac{cn^2}{12}.
 \tag{FT4.15}
\]
The ordered-distinct-triple quotient has degree six over this curve, so
its genus satisfies $\gamma_3\ge6\widehat\gamma_3-5>cn^2/2-5$.
There are five $P$-orbits on the cube, with genera
$\gamma_1,\gamma_2,\gamma_2,\gamma_2,\gamma_3$.
The actual quotient-cover Riemann--Hurwitz identities therefore imply
\[
 \sum_i\operatorname{ind}(x_i,\{1,\ldots,n\}^3)
 =2(n^3-5)+2(\gamma_1+3\gamma_2+\gamma_3)
 >2n^3+cn^2-20.
 \tag{FT4.16}
\]
Sum (FT4.13) and subtract at most ten errors. Then
\[
 2+\frac6e>2+\frac cn-\frac{40M^2}{n^2}-\frac{20}{n^3}.
 \tag{FT4.17}
\]
But $cn\ge80M^2$ and $n\ge4$ show
\[
 cn^2-40M^2n-23\ge40M^2n-23>0.
 \tag{FT4.18}
\]
Thus the right side of (FT4.17) is strictly larger than
$2+3/n^3\ge2+6/e$, a contradiction.
\end{proof}

\paragraph{Gate and geometric interpretation.}
Together with HD-FT01--08, this leaves only ordinary tails of size at most
three as potentially unbounded in the stated actual factorwise setup.
The genus-growth input gives a qualitative reduction, not the effective
integral detector bound required in the main problem. Obtaining numerical
constants, deciding the remaining smaller tails, and supplying actual
field-compatible cycles are distinct outstanding tasks.
The argument concerns actual finite covers, their weight-one local systems
and rational deck projectors. It narrows supplier routes for integral
uniform condition A (cote16, scan050; SIGN\_NEXT SN-C03/C04), while the
condition I positivity step (scan049; cote44, Theorem4), the Cayley--Hamilton
inverse of scans048--049, and the accepted EGA/SGA/FGA spreading interfaces
do not produce missing cycles. No Markman, Severn or032 premise is added.

\section{All three tail-three labels with arbitrary tensor partners}
\label{sec:factorwise-three-tail}

\paragraph{Status and scope.}
HD-FT11--12, reserved question Q2dj, is a new derivation by
\texttt{/root/g34\_eightfolds}.  It uses the exact previously reviewed
three character polynomials and the published genus-growth input
isolated in HD-FT10.  The final natural-degree cutoff is qualitative:
the genus-growth constants are not evaluated. Complete author and
separate root whole-proof readings pass. Root also read the full
source/status record, entire checker and first successful output.
The retained review records the precise shared input and computation scope.

Use exactly the actual quotient and factorization (FT.1)--(FT.2):
there is an actual quotient onto \(P=A_n\) or \(S_n\), an actual
coefficient factorization \(V=W\otimes U\) up to scalars, and genuine
product-one generating inertia with
\[
 e=\dim V,\qquad \sum_i\operatorname{rank}(1-T_i)=2e+6.
 \tag{FT3.1}
\]
The partner \(U\), its degree, its field and every scalar phase are
arbitrary.  The projected permutations are not assigned twelve marks.
Let \(W\) have one of the labels
\[
 (n-3,3),\qquad(n-3,1,1,1),\qquad(n-3,2,1),
 \tag{FT3.2}
\]
or its transpose, or the irreducible restriction of one of these
ordinary representations to \(A_n\).  All natural degrees below are
at least \(64\), so the alternating restrictions are irreducible.

\subsection*{HD-FT11. Comparison with the unordered-triple permutation character}

For a permutation \(x\), put \(f=X_1(x)\), \(u=X_2(x)\), \(v=X_3(x)\)
and \(m=n-f\).  Write
\[
 F_2=\binom f2+u,\qquad F_3=\binom f3+fu+v,\qquad
 B=\binom n3,\quad C_2=\binom n2,\quad a_0=B-C_2.
 \tag{FT3.3}
\]
Thus \(F_j\) is the number of fixed unordered \(j\)-sets.
In the order of the three labels in (FT3.2), the exact characters
and dimensions are
\[
 \begin{array}{c|c|c}
 &\chi(x)&a\\ \hline
 R&F_3-F_2&a_0\\
 H&\chi_R+(f-1)(1-2u)&a_0+n-1\\
 M&2\chi_R+f-2(f-1)u-3v&2a_0+n.
 \end{array}
 \tag{FT3.4}
\]
These are precisely Garsia--Goupil Table~5.1, already used and
reviewed in HD-AT301, D-HK301 and HD-MT301.
\cite[Table~5.1, p.~31; Proposition~I.1 and proof]{HD:GarsiaGoupil2009}
The proof uses the displayed untwisted polynomials.  A transpose
multiplies character values by sign, so the absolute estimates persist.

\begin{lemma}[HD-FT11]
For every \(n\ge64\), each character in (FT3.4) has dimension
\(a\ge n^3/12\), and
\[
 \frac{|\chi(x)|}{a}\le
 \begin{cases}
 F_3/B,& f\ge n/2,\\
 F_3/B+48/n^{3/2},&f<n/2.
 \end{cases}
 \tag{FT3.5}
\]
The estimate is valid on all permutations and therefore all powers.
It bounds every target eigenspace by cyclic Fourier inversion.
\end{lemma}

\begin{proof}
The dimensions satisfy \(a\ge a_0=n(n-1)(n-5)/6\ge n^3/12\).
Also \(B\ge n^3/12\).
The key exact incidence identity is
\[
 3F_3-(n-2)F_2
  =-\frac{m f(f-1)}2+(3f-n+2)u+3v\le3v.
 \tag{FT3.6}
\]
If \(3f-n+2\le0\), the two terms before \(3v\) are nonpositive.
Otherwise \(u\le m/2\) makes their sum at most
\(m(-f^2+4f-n+2)/2\le0\), because
\(-f^2+4f-n+2=6-n-(f-2)^2\).
Since \(C_2/B=3/(n-2)\), (FT3.6) gives
\[
 \frac{\chi_R}{a_0}-\frac{F_3}{B}
 \le\frac{C_2v}{a_0B}\le\frac{24}{n^3}.
 \tag{FT3.7}
\]
Here \(v\le n/3\) and \(C_2\le n^2/2\).

Put
\[
 \delta_H=(f-1)(1-2u),\qquad
 \delta_M=f/2-(f-1)u-3v/2.
 \tag{FT3.8}
\]
For all \(f\), including \(f=0\), one has
\(\delta_H\le n-1\) and \(\delta_M\le n/2\).
Indeed for \(f\ge1\) discard their nonpositive terms, and for
\(f=0\) use \(2u\le n\).
Applying (FT3.7) to \(\chi_H=\chi_R+\delta_H\) and
\(\chi_M/2=\chi_R+\delta_M\) proves the common upper bound
\[
 \frac{\chi(x)}a\le\frac{F_3}{B}
                      +\frac{24}{n^3}+\frac{12}{n^2}
                  \le\frac{F_3}{B}+\frac{36}{n^2}.
 \tag{FT3.9}
\]

We also need negative character values.  Directly from the polynomials,
\[
 \chi_R=\frac{f(f-1)(f-5)}6+(f-1)u+v\ge-n/2,\qquad
 \chi_M=\frac{f(f-2)(f-4)}3-v\ge-n/3.
 \tag{FT3.10}
\]
For the first inequality, \(f=0\) gives \(-u+v\ge-n/2\);
the polynomial at integers \(f\ge1\) is at least \(-2\).
For the second, the only negative polynomial value at a nonnegative
integer is \(-1\) at \(f=3\), where \(v\le(n-3)/3\).

At \(f=0\), \(\chi_H=u+v-1\ge-1\).  For \(f\ge1\),
using \(u\le(n-f)/2\) gives
\[
 \chi_H\ge \frac{f^3}{6}-\frac{f^2}{2}-\frac{fn}{2}-1
       \ge\frac{f^3}{12}-\frac{fn}{2}-\frac{11}{3}
       \ge-\frac{\sqrt2}{3}n^{3/2}-\frac{11}{3}
       \ge-n^{3/2}.
 \tag{FT3.11}
\]
For the middle bounds,
\(f^3-6f^2+32=(f-4)^2(f+2)\ge0\), and minimizing
\(f^3/12-fn/2\) on \(f\ge0\) gives its value
\(-\sqrt2\,n^{3/2}/3\) at \(f=\sqrt{2n}\).
Equivalently that minimum follows by factoring
\(z^3-3z+2=(z-1)^2(z+2)\).
The last bound holds for \(n\ge4\).
Together, (FT3.9)--(FT3.11) and \(a\ge n^3/12\) give
the second case of (FT3.5), with the deliberately coarse constant \(48\).

It remains to prove the error-free first case.  Suppose \(f\ge n/2\).
The same exact incidence identity, now using \(u\le m/2\) and
\(v\le m/3\), gives
\[
 {\cal D}:=B F_2-C_2F_3
  \ge \frac{n(n-1)m}{12}\,K,\qquad
             K=f^2-4f+n-4\ge n^2/4-n-4.
 \tag{FT3.12}
\]
Every nonfixed three-set meets the moved letters, so
\[
 B-F_3\le m\binom{n-1}{2}.
 \tag{FT3.13}
\]
For \(n\ge64\),
\[
 nK-6(n-1)(n-2)
 \ge n^3/4-7n^2+14n-12>0.
 \tag{FT3.14}
\]
It follows that
\({\cal D}\ge(n-1)(B-F_3)\).
For the hook label,
\[
 B\chi_H-(a_0+n-1)F_3
 =-{\cal D}+B\delta_H-(n-1)F_3
 \le-{\cal D}+(n-1)(B-F_3)\le0.
 \tag{FT3.15}
\]
The same calculation works for \(\chi_M/2\), replacing \(n-1\)
by \(n/2\le n-1\), and for \(\chi_R\) it is immediate.

Finally the displayed untwisted characters are nonnegative in this
range.  For \(R\) this follows from \(f\ge32>5\).
For \(M\), \(f(f-2)(f-4)\ge f^3/4\ge n^3/32>n\).
For \(H\), the first lower bound in (FT3.11) is positive, because
\[
 f(f^2/6-f/2-n/2)-1
 \ge n^3/48-3n^2/8-1>0.
 \tag{FT3.16}
\]
Thus the error-free signed inequalities are also absolute inequalities.
This proves (FT3.5) completely.
\end{proof}

\subsection*{HD-FT12. A qualitative exclusion using the actual triple quotient}

Use precisely the published constants \(c,d>0\) and optional
unevaluated initial threshold \(n_0\ge20\) in (FT4.11).
For \(n\ge\max\{n_0,4,2d\,3^{15}/c\}\), the actual unordered-triple
quotient has genus
\[
                         g_3>cn^2/12.
 \tag{FT3.17}
\]
This is HD-FT10's direct consequence of Monderer--Neftin
Theorem~2.5, proof equation~(2.2), valid also for its exceptional
ramification types.  No numerical values of these constants, or
independent proof of the cited precursor, are asserted.

Define the finite but unevaluated integer
\[
 N^{(3)}=\left\lceil
 \max\left\{n_0,1024,\frac{2d\,3^{15}}c,
               \left(\frac{960}c\right)^2,\sqrt{\frac{168}c}\right\}
                  \right\rceil.
 \tag{FT3.18}
\]

\begin{theorem}[HD-FT12]
No actual source (FT3.1) with any label (FT3.2), its transpose,
or its ordinary alternating restriction exists for \(n>N^{(3)}\).
The conclusion permits arbitrary factorwise tensor partners, scalars,
and coefficient fields.  It is a qualitative natural-degree bound.
\end{theorem}

\begin{proof}
First, for every permutation and \(n\ge64\),
\[
 \frac{F_3}{B}\le \frac fn+\frac7n.
 \tag{FT3.19}
\]
Indeed \(F_3\le f^3/6+f(n-f)/2+n/3\).  Put \(t=f/n\).
The inequality
\((1-1/n)^{-1}(1-2/n)^{-1}\le1+6/n\), and \(B\ge n^3/12\),
give
\[
 \frac{F_3}{B}
 \le t^3+\frac6n(t^3+t-t^2)+\frac4{n^2}
 \le t^3+\frac7n\le t+\frac7n.
 \tag{FT3.20}
\]
Here \(t^3+t-t^2\le1\) for \(0\le t\le1\).

Call \(x_i\) bad when \(f(x_i)<n/2\), and put
\(\varepsilon=48/n^{3/2}\).  Every power of a nonbad image
is nonbad.  Fourier inversion, tensor spectral convolution and
(FT3.5) give
\[
 \frac{\operatorname{rank}(1-T_i)}e
 \ge\frac{\nu(W(x_i))}{a}
 \ge\frac{\operatorname{ind}(x_i,\binom{[n]}3)}B
                  -\varepsilon\,{\bf1}_{\{i\ {\rm bad}\}}.
 \tag{FT3.21}
\]
For a bad image, using (FT3.19) in the same average gives
\[
 \frac{\nu(W(x_i))}a
 \ge\frac{\operatorname{ind}(x_i)}n-\frac7n-\varepsilon
 >\frac14-\frac7n-\frac{48}{n^{3/2}}>\frac15.
 \tag{FT3.22}
\]
The last strict inequality follows at \(n=1024\) from
\(7/1024+48/32768=17/2048<1/20\), and improves with \(n\).
Since \(e\ge a\ge n^3/12\), the actual total is
\(2+6/e\le2+72/n^3<11/5\).
Consequently there are at most ten bad entries.
This count uses neither the number of projected marks nor a
separate source assertion for \(W\).

The action of \(P\) on unordered three-sets is transitive.
Its actual connected quotient curve has degree \(B\) and genus \(g_3\).
Riemann--Hurwitz therefore gives the exact identity
\[
 \sum_i\operatorname{ind}(x_i,\binom{[n]}3)=2B-2+2g_3.
 \tag{FT3.23}
\]
Sum (FT3.21) and use (FT3.17).  The threshold makes
\(cn^2/6-2>0\), and \(B\le n^3/6\).  Hence
\[
 \begin{split}
 2+\frac6e
 &>2+\frac{cn^2/6-2}{B}-\frac{480}{n^{3/2}}\\
 &\ge2+\frac cn-\frac{12}{n^3}-\frac{480}{n^{3/2}}
 \ge2+\frac{c}{2n}-\frac{12}{n^3}
 >2+\frac{72}{n^3}\ge2+\frac6e.
 \end{split}
 \tag{FT3.24}
\]
The penultimate bounds use
\(\sqrt n>960/c\) and \(n^2>168/c\).
This contradiction proves the theorem.
\end{proof}

\paragraph{What is and is not bounded.}
Together with HD-FT01--10, this leaves only ordinary tail sizes
at most two as potentially unbounded within the stated actual
quotient and factorwise tensor setup.  It does not give numerical
values for the finitely many remaining degree intervals, nor
enlarge the effective FQ bound without evaluating its additional
genus constants.  Actual quotient maps and tensor factorizations
remain necessary; this result does not infer them from an abstract
character or a proposed branch spectrum.

\section{An effective triple-genus reduction by exact ramification}
\label{sec:projected-triple-ramification}

\paragraph{Status and scope.}
HD-GC01--02 are new root derivations, written 22 November 2026.
Complete author and separate independent whole-proof readings pass. Root has read the entire independent audit report; its distinct orbit enumeration is recorded with its exact scope. All covers below are actual
connected smooth complex covers, with their indicated quotient maps.
No numerical constant from an asymptotic theorem is used. This result
does not assert a numerical bound in the remaining low-genus case.

\paragraph{Notation.}
Let an actual Galois cover of $\mathbf P^1$ have group $A_n$ or $S_n$,
where $n\geq8$. Let $Y_t$ be its quotient for the action on ordered
$t$-tuples of distinct letters, and $X_t$ its quotient for the action on
unordered $t$-sets, for $t=1,2,3$. Thus $Y_1=X_1$ and the natural maps
$\pi_t:Y_t\longrightarrow X_t$ have degree $t!$. All these covers are
connected. Set
\[
 g=g(Y_1),\quad h=g(Y_2),\quad b=g(X_2),\quad c=g(X_3).
\]
For a finite map $f$ of smooth curves, $R_f$ denotes its total
ramification degree $\sum_Q(e_f(Q)-1)$. For an inertia permutation $x$,
write its individual cycle lengths as $r_1,\ldots,r_k$, including
fixed-point cycles, and put
\[
 I(x)=n-k,\qquad
 A(x)=\sum_{i,j=1}^k\bigl(r_i-\gcd(r_i,r_j)\bigr).
\]
Here the index $i$ distinguishes equal-length cycles. For a positive
integer $r$, let $v_2(r)$ be its exponent of $2$.

\paragraph{HD-GC01 (exact local triple ramification).}
The contribution to $R_{Y_2/Y_1}$ above the branch point with inertia
$x$ is $A(x)$. The contribution to $R_{\pi_3}$ is
\[
 T(x)=3\!\sum_{\substack{1\leq i,j\leq k\\v_2(r_i)>v_2(r_j)}}
           \gcd(r_i,r_j)
       +4\#\{i:3\mid r_i\}.
\]
In particular
\[
                    T(x)\leq3A(x)+2I(x).
\]
These assertions also apply to the identity permutation, with all
quantities zero.

\begin{proof}
The number of $x$-orbits on ordered distinct pairs is
$\sum_{i\ne j}\gcd(r_i,r_j)+\sum_i(r_i-1)$: between two different
cycles the orbit length is their least common multiple, and within a
single cycle the nonzero displacement specifies one orbit. The
ramification chain rule for $Y_2\to Y_1\to\mathbf P^1$ therefore gives
\[
 (n-1)k-\#\operatorname{Orb}_{Y_2}(x)
 =nk-\sum_{i,j}\gcd(r_i,r_j)=A(x).
\]
This is also the $t=3$ instance of the elementary cycle-count formula
in Neftin--Zieve, Lemma 5.4; the direct calculation here fixes the
normalization without relying on any of its asymptotic constants.

Consider an $x$-orbit of unordered triples of length $L$. The
permutation $x^L$ on its three letters generates a cyclic subgroup of
$S_3$, of order $m\in\{1,2,3\}$. Above the corresponding point of
$X_3$ there are $6/m$ points of $Y_3$, each with ramification index
$m$. This orbit therefore contributes $6-6/m$ to $R_{\pi_3}$.

If $m=2$, the triple consists of an antipodal pair in an even cycle of
length $r_i$, together with one letter in another cycle of length
$r_j$. Existence of a power that interchanges the pair and fixes the
third letter is equivalent to $v_2(r_i)>v_2(r_j)$. Conversely every
such triple has $m=2$. There are $(r_i/2)r_j$ triples of this form and
their orbit length is $\operatorname{lcm}(r_i/2,r_j)$, so their number
of orbits is
\[
 \frac{(r_i/2)r_j}{\operatorname{lcm}(r_i/2,r_j)}
 =\gcd(r_i/2,r_j)=\gcd(r_i,r_j).
\]
Each contributes $3$. This counts each such triple once: its
antipodal pair, and hence the two cycles in their indicated order,
are uniquely determined by its nontrivial stabilizer on the triple.

If $m=3$, the three letters lie in one cycle of length divisible by
$3$ and are equally spaced in that cycle. Each such cycle gives
exactly one orbit of triples, contributing $4$. These alternatives
exhaust the nontrivial cyclic subgroups of $S_3$, proving the formula.

For each pair in the first sum, $\gcd(r_i,r_j)\leq r_i/2$, whence
$\gcd(r_i,r_j)\leq r_i-\gcd(r_i,r_j)$. All other summands in $A(x)$
are nonnegative. For each cycle divisible by $3$, one has
$4\leq2(r_i-1)$. Adding these two estimates proves the inequality.
\end{proof}

\paragraph{HD-GC02 (numerical genus bounds and the residual case).}
In the notation above,
\[
 12(c-b)\geq
 2(n-8)h+(6n-10)g-12n+26.                         \tag{GC1}
\]
Consequently each of the following hypotheses implies
$g(X_3)>n^2/12$:
\[
 \begin{array}{ll}
  g(Y_1)\geq2,& n\geq16,\\
  g(Y_2)\geq n,& n\geq28.
 \end{array}                                      \tag{GC2}
\]
Thus, for $n\geq28$, the task of proving a uniform numerical
quadratic lower bound for $g(X_3)$ reduces to actual covers with
$g(Y_1)\leq1$ and $g(Y_2)<n$.

\begin{proof}
Put $A=R_{Y_2/Y_1}$ and $T=R_{\pi_3}$. Summing HD-GC01 over the
actual branch points and applying Riemann--Hurwitz to
$Y_1\to\mathbf P^1$ gives
\[
       T\leq3A+4(n+g-1),\qquad
       A=2(h-1)-2(n-1)(g-1).
\]
Riemann--Hurwitz for the three actual maps
$Y_3\to Y_2$, $Y_2\to X_2$ and $Y_3\to X_3$, of respective
degrees $n-2$, $2$ and $6$, gives exactly
\[
 12(c-b)=2(n-5)(h-1)+R_{Y_3/Y_2}+3R_{\pi_2}-T.
\]
In particular
\[
 \begin{split}
 12(c-b)&\geq2(n-5)(n-1)(g-1)
                    +(n-8)A-4(n+g-1)\\
         &=2(n-8)h+(6n-10)g-12n+26,
 \end{split}
\]
which proves (GC1). No sign condition on $h-1$ was assumed when
discarding the nonnegative ramification degrees.

If $g\geq2$ and $n\geq16$, the coefficient of $g$ in the first
expression is positive, and $A\geq0$ with $n-8\geq0$. Hence
\[
       12(c-b)\geq2n^2-16n+6>n^2.
\]
The strict inequality follows from $n^2-16n+6\geq6$ at $n=16$
and its increase thereafter. If $h\geq n$ and $n\geq28$, (GC1),
$g\geq0$ and $6n-10>0$ give
\[
       12(c-b)\geq2n^2-28n+26>n^2.
\]
Here $n^2-28n+26\geq26$ at $n=28$ and increases thereafter.
In either case $b\geq0$ proves (GC2).
\end{proof}

\paragraph{Effectivity and geometric meaning.}
In either case (GC2), any required inequality $g(X_3)>G$ follows
from the additional explicit condition $n^2\geq12G$. Likewise the
symbolic triple-genus input $g(X_3)>cn^2/12$ in HD-FT09--12 can use
$c=1$ and threshold $28$ in these cases. The full surviving
low-genus class has not been bounded here. The result is a constraint
on actual quotient covers, hence on their weight-one cohomology and
possible curve suppliers. It does not construct integral cycles or
bound a detector on a new Weil component. The existing condition-A,
condition-I, projector and spreading interfaces in the paper retain
their previously recorded exact locators and hypotheses.

\section{A necessary pair-genus and partner-support gate for tail two}
\label{sec:factorwise-two-tail-gate}

\paragraph{Status and scope.}
HD-FT13--15, reserved question Q2dm, is a new derivation by
\texttt{/root/g34\_eightfolds}.  The elementary estimates are effective
for \(n\ge64\).  The final branch-list consequence uses the published
but unevaluated constant \(N_{14}\) of Monderer--Neftin; no numerical
value is assigned to it.  This is a necessary gate, not a complete
exclusion or a claim of an exhaustive catalogue of tensor partners.
Complete author and separate root whole-proof readings pass. Root also read the complete checker, first output and source/status record, the full invoked published proof algorithm, and visually inspected the entire branch table. No new root execution or numerical N14 is claimed.

Use the actual quotient and factorization of HD-FT01.  In particular,
there is a connected finite cover of the sphere with an actual quotient
onto \(P=A_n\) or \(S_n\), and an actual coefficient factorization
\(V=W\otimes U\), up to scalar operators.  Write \(x_i\in P\) for
the projected permutations of the genuine inertia operators \(T_i\).
The hypotheses used below are precisely
\[
 a=\dim W,\qquad b=\dim U,\qquad e=ab,\qquad r\ge12,
 \qquad \sum_{i=1}^r\operatorname{rank}(1-T_i)=2e+6.
 \tag{FT2P.1}
\]
Every \(T_i\ne1\), the tuple has product one and generates the actual
covering group.  The inequality \(r\ge12\) belongs to the original
rank-nine inducing sphere source.  It is not assigned to the projected
\(A_n\)- or \(S_n\)-factor.  The degree, coefficient field and all
scalar phases of \(U\) are arbitrary.

Let \(W\) have label \((n-2,2)\) or \((n-2,1,1)\), its transpose,
or the irreducible ordinary restriction of one of these representations
to \(A_n\).  The restrictions are irreducible for \(n\ge64\).
All comparisons below are absolute-character or projective-support
comparisons, so sign twists have exactly the same bounds.

\subsection*{HD-FT13. An all-eigenvalue comparison with unordered pairs}

For a permutation \(x\), put \(f=X_1(x)\), \(u=X_2(x)\), \(m=n-f\)
and
\[
 B=\binom n2,\qquad F=\binom f2+u.
 \tag{FT2P.2}
\]
Thus \(F\) is the number of fixed unordered pairs.  The exact two
characters, in the order of the two labels above, are
\[
 \begin{array}{c|c|c}
 &\chi(x)&a\\ \hline
 R&F-f=f(f-3)/2+u&B-n,\\
 H&F-f+1-2u=(f-1)(f-2)/2-u&B-n+1.
 \end{array}
 \tag{FT2P.3}
\]
These are the tail-two rows of the already reviewed character-polynomial
table.\cite[Table~5.1, p.~31; Proposition~I.1 and proof]{HD:GarsiaGoupil2009}
They also follow from the unordered-pair permutation representation and
the exterior square of the deleted permutation representation.

\begin{lemma}[HD-FT13]
For every \(n\ge64\) and every permutation,
\[
 a\ge n^2/4,\qquad
 \frac{|\chi(x)|}{a}\le \frac FB+
                      \frac8{n^2}\,\mathbf1_{f\le2}.
 \tag{FT2P.4}
\]
Let \(I_2(x)\) be the index of \(x\) on unordered pairs and
\(I(x)\) its natural index.  If
\(\nu(Wx)=a-\max_\zeta\dim\ker(Wx-\zeta)\), then
\[
 \frac{\nu(Wx)}a\ge\frac{I_2(x)}B-
                  \frac8{n^2}\,\mathbf1_{f\le2},\qquad
 \frac{\nu(Wx)}a\ge\frac{I(x)}n-
                  \frac1{n-1}-\frac8{n^2}\,\mathbf1_{f\le2}.
 \tag{FT2P.5}
\]
Both conclusions concern every target eigenvalue, including arbitrary
scalar twists of \(Wx\).
\end{lemma}

\begin{proof}
The dimension bounds and \(B\ge n^2/4\) follow immediately from
(FT2P.3).  Also \(2u\le m\).  For the first character,
\[
 B\chi_R-(B-n)F
       =\frac n2(-fm+2u).
 \tag{FT2P.6}
\]
If \(f\ge1\), the last parenthesis is at most \(m(1-f)\le0\).
If \(f=0\), the normalized difference is
\(nu/((B-n)B)\le8/n^2\).
The only possible negative values of \(\chi_R\) occur at \(f=1,2\),
and are at least \(-1\).  Their absolute normalized value is at most
\(1/a\le4/n^2\).  This proves (FT2P.4) for \(R\), with no error
when \(f\ge3\).

For the hook character the exact signed comparison is
\[
 B\chi_H-(B-n+1)F
   =-\frac{n-1}{2}\bigl(m(f-1)+2(n-1)u\bigr).
 \tag{FT2P.7}
\]
It is nonpositive for \(f\ge1\).  For \(f=0\) its numerator is
at most \(B\), so its normalized positive error is at most
\(1/a\le4/n^2\).  Negative character values are handled separately,
using \(-\chi_H=F-(f-1)^2\):
\[
 B(-\chi_H)-(B-n+1)F
   =(n-1)F-B(f-1)^2
   \le\frac{(n-1)^2}{2}f(2-f).
 \tag{FT2P.8}
\]
Here the inequality uses \(2u\le n-f\).  The right side is
nonpositive at \(f=0\) and every \(f\ge2\).  At \(f=1\), division
by \(aB\ge n^4/16\) bounds the error by \(8/n^2\).
This proves the asserted absolute comparison in all cases.

Every power of \(x\) fixes all the letters fixed by \(x\).  Thus
if \(f(x)\ge3\), all powers have the error-free estimate.
For \(d=\operatorname{ord}(x)\), cyclic Fourier inversion gives,
for every \(\zeta\),
\[
 \frac{\dim\ker(Wx-\zeta)}a
 \le\frac1d\sum_{j=0}^{d-1}\frac{|\chi(x^j)|}a
 \le\frac1d\sum_{j=0}^{d-1}\frac{F(x^j)}B
                       +\frac8{n^2}\mathbf1_{f(x)\le2}.
 \tag{FT2P.9}
\]
The average of \(F(x^j)\) is the number of cycles on unordered
pairs, proving the first support estimate.
Finally,
\[
 \frac FB\le\frac fn+\frac1{n-1}.
 \tag{FT2P.10}
\]
Indeed for \(f\ge1\), the stronger inequality \(F/B\le f/n\)
follows from
\(2F-(n-1)f\le(f-1)(f-n)\le0\); for \(f=0\) use
\(F/B\le1/(n-1)\).  Averaging (FT2P.10) over powers and using
\(d^{-1}\sum_j f(x^j)=n-I(x)\) proves the second estimate.
\end{proof}

Spectral convolution with any finite-order partner, and with any
scalar phase, gives
\[
 \frac{\operatorname{rank}(1-T_i)}e\ge\frac{\nu(Wx_i)}a.
 \tag{FT2P.11}
\]
For example, each partner eigenspace can contribute at most the
largest \(W\)-eigenspace to a specified product eigenvalue.  This
explains why no reality or coefficient-field equality is needed.

\subsection*{HD-FT14. At most four error-bearing entries and pair genus at most fourteen}

\begin{proposition}[HD-FT14]
In the actual source (FT2P.1), at most four projected inertia
permutations satisfy \(f(x_i)\le2\).  The connected quotient cover
on unordered pairs has genus \(g_2\le14\).
\end{proposition}

\begin{proof}
Call an entry with \(f(x_i)\le2\) bad.  Its natural index satisfies
\(I(x_i)\ge(n-2)/2\), since every moved cycle has length at least two.
Equations (FT2P.5) and (FT2P.11) imply
\[
 \frac{\operatorname{rank}(1-T_i)}e
 \ge\frac12-\frac1n-\frac1{n-1}-\frac8{n^2}
 >\frac9{20}\qquad(n\ge64).
 \tag{FT2P.12}
\]
For the strict inequality it suffices to check
\(1/64+1/63+1/512<1/20\), since each subtracted term decreases.
On the other hand,
\[
 \frac{2e+6}{e}\le2+\frac{24}{n^2}<\frac94.
 \tag{FT2P.13}
\]
Five bad entries would contradict these two strict inequalities.

The unordered-pair action of \(P\) is transitive.  Hence its actual
degree-\(B\) quotient is connected and Riemann--Hurwitz gives
\(\sum_i I_2(x_i)=2B-2+2g_2\), including zero contributions from
projected identity entries.  Summing the first inequality of
(FT2P.5), with at most four errors, yields
\[
 2+\frac6e\ge2+\frac{2g_2-2}{B}-\frac{32}{n^2},
 \qquad
 2g_2-2\le\frac{56B}{n^2}
             =28\left(1-\frac1n\right)<28.
 \tag{FT2P.14}
\]
The integral genus is therefore at most fourteen.
\end{proof}

\subsection*{HD-FT15. The actual branch list forces a small-support partner operator}

Let \(N_{14}\) denote the constant in Monderer--Neftin Theorem~1.2
for genus bound fourteen.  This is an existing published qualitative
constant; its value is not evaluated here.
\cite[Theorems~1.2 and 5.1; Proposition~4.1 and Table~2.1]{HD:MondererNeftin2022}

\begin{theorem}[HD-FT15]
Suppose \(n\ge64\) and \(n>N_{14}\).  In every source (FT2P.1),
the natural and unordered-pair quotient covers both have genus zero.
At most five \(x_i\) are nonidentity, so at least seven genuine
original inertia operators have \(x_i=1\).  They satisfy
\[
 \sum_{x_i=1}\frac{\operatorname{rank}(1-T_i)}e
       \le\frac{61}{n^2}.
 \tag{FT2P.15}
\]
In particular at least one such genuine operator has
\[
 0<\frac{\operatorname{rank}(1-T_i)}e\le\frac{61}{7n^2},
 \qquad b\ge\frac{7n^2}{61}.
 \tag{FT2P.16}
\]
If these kernel-projected genuine operators all have normalized moving
rank at least a specified \(\eta>0\), then necessarily
\[
                      7\eta\le\frac{61}{n^2}.
 \tag{FT2P.17}
\]
This last assertion is conditional on the stated actual support floor;
no such uniform floor for an unbounded partner family is assumed.
\end{theorem}

\begin{proof}
Apply Monderer--Neftin Theorem~1.2 to the actual natural degree-\(n\)
cover and its pair-stabilizer subgroup \(H=P\cap(S_{n-2}\times S_2)\).
The subgroup \(H\) cannot contain \(A_{n-1}\), already by its order
when \(n\ge64\).  Proposition HD-FT14 supplies the genus bound
fourteen required by that theorem.  It follows that the natural
ramification is among the types in Proposition~4.1 and Table~2.1.
Proposition~4.1(1), for either \(A_n\) or \(S_n\), gives genus zero
for the point and pair stabilizers.  Every row of Table~2.1 has at
most five genuine branch points: only rows F1.1 and F1.2 have five,
and all others have at most four.  Multiplicities such as ``twice''
and ``thrice'' count distinct branch points.  These are the projected
nonidentity entries, not the entire original tuple.  Since \(r\ge12\),
at least seven original genuine entries project to the identity.

The support comparison and nonnegative pair genus give
\[
 \sum_{x_i\ne1}\frac{\nu(Wx_i)}a
 \ge2-\frac2B-\frac{32}{n^2}
 \ge2-\frac{37}{n^2}.
 \tag{FT2P.18}
\]
The last inequality uses \(2/B=4/(n(n-1))\le5/n^2\), valid for
\(n\ge5\).  Subtract (FT2P.18), via (FT2P.11), from the actual
budget (FT2P.1), and use \(6/e\le24/n^2\).  This proves
(FT2P.15) and its seven-entry average in (FT2P.16).

If \(x_i=1\), the actual operator is \(1_W\otimes V_i\) after
absorbing its scalar in \(V_i\).  Since \(T_i\ne1\), one has
\(V_i\ne1\) and
\(\operatorname{rank}(1-T_i)/e
=\operatorname{rank}(1-V_i)/b\ge1/b\).
Summing over the seven entries proves the partner-dimension bound.
The same sum proves (FT2P.17) under its explicit hypothesis.
\end{proof}

\paragraph{Retained boundary.}
The result does not exclude an unbounded tensor partner that supplies
such small-support genuine operators.  In particular it does not
identify a maximal-degree ordinary factor, classify factor permutations,
or establish that all alternative partner images are in a bounded
catalogue.  Any consequence requiring those statements must prove them
separately.  The effective estimates through HD-FT14 and the qualitative
\(N_{14}\)-dependent conclusion of HD-FT15 are deliberately distinct;
the numerical-source question Q2dk remains open.

\section{A numerical triple-genus bound for every alternating or symmetric cover}
\label{sec:effective-triple-genus}

\paragraph{Status.}
HD-GC03--07 give a new proof of an explicit numerical form of the
known qualitative triple-genus growth. The proof was written by root
on 22 November 2026. Complete author and separate independent whole-proof readings pass. Root has also read the full independent audit, its entire exact checker and first successful output, and both targeted derivation reports. The two recorded wording repairs change no mathematical display.
The cycle-mode argument is a numerical specialization of the method
in Neftin--Zieve, Proposition~7.1 and Corollary~7.3; all constants
and the complete argument needed here are proved below. No statement
from that preprint with unspecified constants is used as a premise.
HD-GC01--02 and the standard finite automorphism description of a
complex elliptic curve are the previously checked inputs.

Use the actual covers and notation of HD-GC01--02:
\[
 g=g(Y_1),\qquad h=g(Y_2),\qquad b=g(X_2),\qquad c=g(X_3),
 \qquad A=R_{Y_2/Y_1}.
\]
The base is $\mathbf P^1$, the Galois group is $A_n$ or $S_n$, and
all quotient curves and maps are the actual connected ones. In this
section $n\geq10^6$. Identity inertia may always be omitted.

\subsection*{HD-GC03: an explicit cycle-mode lemma}

Let a permutation on $n$ letters have individual cycle lengths $r_i$,
and write $A(x)=\sum_{i,j}(r_i-\gcd(r_i,r_j))$ as in HD-GC01.
If $A(x)<4n$, either there is a $k\in\{1,\ldots,6\}$ such that
all cycles other than the $k$-cycles have total length at most
$4k(k+1)$, or each cycle length $k\leq6$ occurs at most eight times.
More generally, if $A(x)<6n$, the corresponding alternatives hold
with $6k(k+1)$ and twelve times, respectively.

\begin{proof}
We prove both statements simultaneously. Let $q=8$ in the first
case and $q=12$ in the second, so $A(x)<qn/2$. If some length at
most six has more than $q$ cycles, let $k$ be the least such length
and let $m>q$ be its multiplicity. The sum $R$ of lengths smaller
than $k$ is at most $qk(k-1)/2$. Comparing the $m$ cycles of length
$k$ with cycles of length greater than $k$ gives
\[
       \frac m2(n-km-R)\leq A(x)<\frac q2 n,
       \qquad m(n-km-R)<qn.                         \tag{EG1}
\]
Here $r-\gcd(r,k)\geq r/2$ for every $r>k$.

First suppose $m\leq n/(2k)$. Since $R\leq180$, (EG1) gives
$m<qn/(n/2-180)<2q+1$, hence $m\leq2q$. Then
$km+R\leq144+180=324$ gives
$m<qn/(n-324)<q+1$, a contradiction. All the strict numerical
comparisons hold for both $q$ at $n\geq10^6$.
Thus $m>n/(2k)$. Put $E=n-km$. Then
\[
 E-R<2qk,\qquad
 E<R+2qk\leq\frac q2 k(k+3)\leq324.
\]
Using $m=(n-E)/k$ once more in (EG1) gives
$E-R<qkn/(n-324)<qk+1$. Since $E-R$ is an integer,
\[
             E\leq R+qk\leq\frac q2 k(k+1).
\]
If no such least length exists, the second alternative holds by
definition. This proves both versions.
\end{proof}

Assign a permutation its finite mode $k$ by the construction just
given, and assign mode $\infty$ in the second alternative. For the
$4n$ version, a finite mode has exceptional length $E\leq168$;
mode one moves at most eight letters. A mode-$\infty$ permutation
has at most $48$ cycles of length at most six, and therefore at most
$n/7+48$ cycles in total. Consequently its natural index is at least
$6n/7-48$. For finite $k>1$, its index differs from $n(1-1/k)$ by
at most $168$, and is at least $n/2-168$.

\subsection*{HD-GC04: the low ordered-pair-genus case}

Assume $g\leq1$ and $h<n$. Then
\[
        A=2(h-1)-2(n-1)(g-1)<4n.                    \tag{EG2}
\]
Every local $A(x)$ is nonnegative, so HD-GC03 applies to every
branch point. Call modes greater than one large, including mode
$\infty$, and call mode one small. Let $J$ be the sum of the natural
indices of the small branch permutations.

For a small permutation with moved support $E\leq8$ and natural
index $j$, its fixed/nontrivial cycle pairs contribute exactly
$(n-E)j$ to $A(x)$. The remaining contribution is nonnegative.
Thus
\[
       A\geq(n-8)J,\qquad J<\frac{4n}{n-8}<5,
       \qquad J\leq4.                              \tag{EG3}
\]
There are at most four large branches: five would have total natural
index at least $5(n/2-168)>2n$, whereas natural Riemann--Hurwitz gives
total index $2n+2g-2\leq2n$.

The multiset of large modes is one of
\[
 (3,3,3),\ (2,4,4),\ (2,3,6),\ (2,2,2,2),\
 (\infty,2,2),\ (\infty,\infty).                    \tag{EG4}
\]
Here and below all constants are numerical. To verify the list, first
suppose every large mode is finite, and let
$S=\sum(1-1/k)$. Natural Riemann--Hurwitz and (EG3) give
$|S-2|\leq678/n<1/60$. As $S$ has denominator dividing $60$, $S=2$.
For three terms the reciprocal sum is one, giving exactly
$(3,3,3),(2,4,4),(2,3,6)$; four terms must all be two. At most two
finite terms cannot sum to two.

If there is one infinite mode and the finite modes have weight sum
$S$, its index is $(2-S)n+\epsilon$ with $|\epsilon|\leq510$.
The index bounds $6n/7-48\leq I\leq n$ give
\[
           1-510/n\leq S\leq8/7+558/n.
\]
One finite term is at most $5/6$, three are at least $3/2$,
and two terms are either $1$ or at least $7/6$.
Since $7/6-8/7=1/42>558/n$, the only possibility is two modes two.
Two infinite modes cannot be accompanied by a finite mode, since
$2(6n/7-48)+(n/2-168)>2n$; three infinite modes are also impossible.
This proves (EG4).

We claim that, apart from the branch pattern in HD-GC05 below,
\[
                              h\geq n/24.           \tag{EG5}
\]
There are three cases.

\paragraph{A finite even mode.}
Such a mode has $k\in\{2,4,6\}$ and contains at least
$n/k-4(k+1)\geq n/6-28$ even cycles. Each even cycle contributes
one to the ramification of $Y_2\to X_2$: its antipodal pairs form
one unordered-pair orbit whose ordered double cover is ramified.
All other pair orbits are unramified in that double cover. Hence
\[
       h=2b-1+\tfrac12R_{\pi_2}
          \geq n/12-15\geq n/24.                   \tag{EG6}
\]
This treats every list in (EG4) except $(3,3,3)$ and
$(\infty,\infty)$.

\paragraph{Three modes three.}
For each of the three permutations let its exceptional cycles have
total length $E_i\leq48$, number $k_i$, and number $N_i$ whose
length is not divisible by three. Put
\[
 S_i=\sum_{\rm exceptional}(r-3),\qquad
 Q_i=S_i+4N_i,
 \qquad N=N_1+N_2+N_3.
\]
Pairing an exceptional cycle of length $r$ with a main three-cycle
in both orders contributes $r+3-2\gcd(r,3)$, which is $r-3$ if
$3\mid r$ and $r+1$ otherwise. Thus its exact pair-ramification
contribution is
\[
 A_i=\frac{n-E_i}{3}Q_i+A_{{\rm exc},i},
 \qquad 0\leq Q_i\leq2E_i,\quad
 0\leq A_{{\rm exc},i}\leq E_i^2.
\]
It follows that $A_i-nQ_i/3$ lies between $-1536$ and $2304$.
For a small branch with moved support $E\leq8$, number $k$ of
nontrivial cycles and index $j$, the remaining internal contribution
is at most $(k-1)j\leq Ej$. Thus its total contribution differs
from $nj$ by a number in $[-8j,0]$. Summing and using (EG3),
the total error lies in $[-4640,6912]$, and in particular in the
slightly wider interval $[-9248,7168]$ used below.

The natural indices of the three large permutations are
$2n/3+S_i/3$. Therefore natural Riemann--Hurwitz gives
$\sum_i S_i+3J=6g-6$. The $nJ$ term from the small branches
cancels exactly, leaving
\[
 \begin{split}
 A&=n(2g-2+4N/3)+\epsilon,
           \qquad -9248\leq\epsilon\leq7168,\\
 h&=n(2g-2+2N/3)+(2-g)+\epsilon/2.                 \tag{EG7}
 \end{split}
\]
If $g=0$ and $N\leq2$, its upper bound makes $h<0$, impossible.
If $g=0,N=3$, the equality $\sum_iS_i+3J=-6$ forces exactly three
exceptional fixed-point cycles, no other exceptional cycles, and
$J=0$: each of these three non-three-divisible cycles contributes
at least $-2$, and all remaining summands are nonnegative.
If $g=1,N=0$, the same equality forces no exceptional cycles and
$J=0$. In either case all three genuine branch permutations have
order three. The actual Galois-cover Riemann--Hurwitz formula then
gives genus one for the Galois cover. A finite group of automorphisms
of a complex genus-one curve is an extension of an abelian
translation group by a cyclic group, hence solvable, contradicting
$A_n$ or $S_n$ for $n\geq5$.
\cite[Chapter IV, Proposition 7.13 and Remark 7.18, pp.~140--141]{TD:MilneEC2021}
Thus $N\geq4$ if $g=0$, and $N\geq1$ if $g=1$. Equation (EG7)
gives $h\geq(2/3)n-4623\geq n/2$, proving (EG5) here.

\paragraph{Two infinite modes.}
Let $u,v$ be the numbers of cycles in the two large permutations.
The exact natural Riemann--Hurwitz identity is
\[
                         J=u+v+2g-2.               \tag{EG8}
\]
If $g=1$, then $J\geq2$, so
$h=1+A/2\geq1+(n-8)J/2\geq n-7$.
If $g=0$ and $J\geq3$, then
$h=2-n+A/2\geq2-n+(n-8)J/2\geq n/2-10$.
Both imply (EG5).

It remains that $g=0$ and $J\leq2$. Since $u,v\geq1$, if $J=0$
then $u=v=1$ and $A=0$, giving $h=2-n<0$; this is impossible.
If $J=1$, the two cycle numbers are one and two, and the small
branch data consist of one transposition. This is exactly the
special pattern treated in HD-GC05.

If $J=2$, then $u+v=4$. An unequal partition into $u$ lengths
has $A(x)\geq n/(2u)$: compare a largest and a smallest length,
whose gcd is at most half the largest. Thus if either of the two
large partitions is unequal, $u,v\leq3$ and
\[
       h=2-n+A/2\geq2-n+(2(n-8)+n/6)/2
                         =n/12-6\geq n/24.
\]
If both partitions are equal-length partitions, their lengths are
$n/u$ and $n/v$. Here $(u,v)$ is $(1,3),(2,2)$ or $(3,1)$.
Choose a coordinate on the genus-zero natural curve and send the
two corresponding base values to zero and infinity. The defining
rational function has every zero and pole multiplicity divisible
by $d=n/\operatorname{lcm}(u,v)\geq n/3$. Factoring its zeros and
poles over $\mathbf C$ writes it as a $d$-th power of a rational
function of degree $\operatorname{lcm}(u,v)\in\{2,3\}$.
Both factors have degree greater than one. Their intermediate
fibres give a nontrivial block system on the $n$ letters,
contradicting the primitive natural action of $A_n$ or $S_n$.
This completes the proof of (EG5), with only the stated special
pattern excluded from that intermediate conclusion.

\subsection*{HD-GC05: the exceptional three-branch pattern}

Suppose the natural genus is zero and the branch cycles have types
\[
                   [n],\quad[a,n-a],\quad[2,1^{n-2}].
\]
Put $a'=n-a$ and $d=\gcd(a,a')$. The pair ramification degrees
of these three permutations are $0,n-2d,n-2$, respectively.
Therefore $h=1-d$. Connectedness of the actual ordered-pair
quotient forces $d=1$ and $h=0$.

Put $D=n(n-1)(n-2)$. The $n$-cycle has $(n-1)(n-2)$ orbits on
ordered distinct triples. The split permutation has
\[
 (a-1)(a-2)+(a'-1)(a'-2)+3(a-1)+3(a'-1)
                         =a^2+(a')^2-2
\]
such orbits: mixed triples have orbit length $aa'$ because
$\gcd(a,a')=1$. The transposition has index $3(n-2)^2$ on
ordered triples, obtained by subtracting its
$(n-2)(n-3)(n-4)$ fixed triples from $D$ and dividing by two.
Riemann--Hurwitz consequently gives the exact identity
\[
 2(g(Y_3)-1)=(n-5)(n-2)+2(a-1)(a'-1).              \tag{EG9}
\]
The local formula HD-GC01 bounds the three contributions to
$R_{\pi_3}$ by $4,11,3(n-2)$, respectively. In the middle term
there is at most one ordered pair with unequal two-adic valuations,
whose gcd is one, and at most two cycles divisible by three.
Thus $R_{\pi_3}\leq3n+9$. From (EG9) and the degree-six map
$Y_3\to X_3$ we obtain
\[
                  12c\geq n^2-10n+13>n^2/2
                         \qquad(n\geq20).          \tag{EG10}
\]
No classification theorem or realization of arbitrary branch lists
is inferred: this is a calculation for an actual cover with this list.

\subsection*{HD-GC06: the uniform numerical result}

For every actual connected $A_n$- or $S_n$-Galois cover of the
sphere and $n\geq10^6$, its unordered-triple quotient satisfies
\[
                         \boxed{\ g(X_3)>n^2/288\ }. \tag{EG11}
\]
Indeed HD-GC02 proves the stronger $c>n^2/12$ if $g\geq2$ or
$h\geq n$. Otherwise HD-GC04 gives $h\geq n/24$, apart from
the special case settled more strongly by (EG10). Substituting
$h\geq n/24$ and $g\geq0$ in (GC1) yields
\[
 12(c-b)\geq n^2/12-(38/3)n+26>n^2/24,
\]
where the last inequality holds already for $n\geq304$.
Since $b\geq0$, (EG11) follows.

In particular every fixed genus bound $G\geq0$ is exceeded by
the triple quotient once
\[
                 n\geq\max\{10^6,\lceil\sqrt{288G}\rceil\}.
\]
This supplies the numerical triple-genus input in HD-FT07--12:
one can replace their symbolic bound $g(X_3)>c_0n^2/12$ by
$c_0=1/24$ with threshold $10^6$. It does not supply a numerical
version of every subgroup or pair-genus classification theorem.

\subsection*{HD-GC07: an effective branch-count alternative for tail two}

Suppose the unordered-pair quotient has $b\leq14$. Then, for
$n\geq10^6$, the natural quotient has at most eight genuine
branch points. This consequence does not use $N_{14}$.

First $g\leq b$. To see this directly, the permutation module on
letters embeds in the module on unordered pairs by sending a
letter $i$ to the sum of pairs containing it. The map is injective:
vanishing of all pair coefficients means $a_i+a_j=0$ for all
$i\ne j$, whence every $a_i=0$ for $n\geq3$.
Taking multiplicities in the actual Galois cover's holomorphic
differentials identifies the two invariant dimensions with $g$
and $b$, respectively. Semisimplicity over $\mathbf C$ therefore
gives $g\leq b\leq14$.

Every even cycle contributes one to $R_{\pi_2}$ and at least one
to the natural index, so $R_{\pi_2}\leq2n+2g-2$. Consequently
\[
 h=2b-1+R_{\pi_2}/2\leq n+40,\qquad
 A=2(h-1)-2(n-1)(g-1)\leq4n+76<6n.
\]
Use the $6n$ version of HD-GC03. Its finite modes have exceptional
length at most $252$; mode one moves at most twelve letters;
mode infinity has index at least $6n/7-72$. Each large branch
therefore has index at least $n/2-252$. Five would exceed the
total natural index $2n+2g-2\leq2n+26$, so there are at most four.
The small total index satisfies
\[
                (n-12)J\leq A\leq4n+76,
                \qquad J<5,
\]
hence there are at most four small genuine branches. This proves
the eight-branch bound without classifying their exact types.

For an actual tail-two source satisfying (FT2P.1), HD-FT14 gives
$b\leq14$. Therefore at least four of its at least twelve genuine
original inertia operators project to the identity in $A_n$ or
$S_n$. The effective support estimate (FT2P.18), whose proof uses
only $b\geq0$, leaves total normalized moving rank at most $61/n^2$
for those entries. If $B$ denotes the tensor-partner dimension,
each such genuine entry has normalized moving rank at least $1/B$, so
\[
             B\geq4n^2/61,
 \qquad
 0<\frac{\operatorname{rank}(1-T_i)}{\dim V}
                       \leq\frac{61}{4n^2}
 \quad\hbox{for at least one partner-only entry}.  \tag{EG12}
\]
This is weaker than the seven-entry consequence above an unknown
$N_{14}$, but it is fully numerical. The precise pair-genus
classification cutoff, an unbounded-partner exclusion, and source
descent or detector construction are not asserted here.

\paragraph{Project frame.}
These effective inequalities constrain actual curve suppliers and
their weight-one cohomology. Their use in the earlier integral
projector and cycle-transfer results retains the exact condition-A,
condition-I and spreading locators already recorded there. They
do not turn a support inequality into an integral detecting cycle,
do not bound arbitrary original overgroups, and invoke none of
Markman, Severn or032. The determinant and Hodge-family gates
remain as stated until an actual source and its cycles are supplied.

\section{Derived numerical cutoffs for ordinary alternating tensor factors}
\label{sec:effective-factorwise-tail-cutoffs}

\paragraph{Status and exact scope.}
HD-FT16--17, reserved question Q2dq, is a derived quantitative
corollary by \texttt{/root/g34\_eightfolds}.  It introduces no new
character or genus theorem.  The initial conditional proof and its first
exact certificate are retained unchanged.  HD-GC06, equation (EG11),
was subsequently independently reviewed and accepted at 02:38:42 UTC
on 22 November 2026; its accepted numerical conclusion is now used.
Complete author and separate root whole-proof readings pass, together with the complete source/status, checker and first-output readings. The first conditional certificate and subsequent status-only update are both retained; no duplicate arithmetic execution is claimed.  All character estimates and
actual quotient-cover interfaces used here are the already reviewed
HD-FT01--12.  No accepted file is changed by this addendum.

Throughout retain the actual factorwise setup (FT.1)--(FT.2): a
connected finite cover of the sphere, an actual quotient onto \(A_n\)
or \(S_n\), an actual factorization \(V=W\otimes U\) up to scalars,
and the genuine generating product-one inertia budget
\[
 e=\dim V\ge a=\dim W,\qquad
                \sum_i\operatorname{rank}(1-T_i)=2e+6.
 \tag{FTE.1}
\]
There is no lower bound on the number of genuine projected marks.
The tensor partner, its degree, its field and every scalar are arbitrary.
The genus premise being substituted is precisely
\[
 g(X_3)>n^2/288\qquad(n\ge10^6),
 \tag{FTE.2}
\]
where \(X_3\) is the actual connected quotient on unordered triples.
This is the stated conclusion of HD-GC06, not a proposed numerical
interpretation of an unspecified published asymptotic constant.

\subsection*{HD-FT16. Effective small-tail cutoffs with every error retained}

For a partition \(\mu\vdash k\), write the reviewed character
polynomial as \(q_\mu=C f^k+R\), with \(C>0\), and put
\[
 M=1+\|R\|_1/C,\qquad
 G=\left\lceil2M(M+1)+2\right\rceil.
 \tag{FTE.3}
\]
The norm is the sum of absolute monomial coefficients.  The integer
values below are recomputed from the retained checked polynomials;
they are not estimates of genus constants.

\begin{proposition}[HD-FT16]
Assume (FTE.2).  The following ordinary \(S_n\)-labels, their
transposes and their irreducible \(A_n\)-restrictions cannot occur
in (FTE.1):
\[
 \begin{array}{c|c}
 k&\text{sufficient integer natural-degree threshold}\\ \hline
 3& n\ge530841600,\\
 4& n\ge\max\{10^6,\lceil256M\rceil,\lceil1920M^2\rceil\},\\
 5,6&n\ge\max\{10^6,\lceil256M\rceil,
                              \lceil\sqrt{288G}\rceil\}.
 \end{array}
 \tag{FTE.4}
\]
In particular the single threshold \(530841600\) excludes every
tail of size three through six in this actual factorwise setup.
\end{proposition}

\begin{proof}
For tails five and six, HD-FT07--08 supplies an error-free cubic
absolute-character bound in the high-fixed-point region and error
\(2M(M+1)/n^3\) on at most ten remaining entries.  Its proof uses
\(n\ge256M\), \(a\ge2n^3\), and only the condition
\(g(X_3)>G\) in its geometric step.  Equation (FTE.2) supplies that
condition once \(n^2\ge288G\), including equality because the genus
bound is strict.  This gives the last row of (FTE.4), with the entire
\(6/e\) term already retained in HD-FT08.

For tail four, use the local estimates of HD-FT09: for \(n\ge256M\)
one has \(a\ge2n^3\), and at most ten entries have absolute-character
error \(4M^2/n^2\) compared with the natural cube.  The ordered
distinct-triple quotient maps to \(X_3\) with degree six.  Thus its
genus \(\gamma_3\) satisfies
\(\gamma_3-1\ge6(g(X_3)-1)\), whence
\(\gamma_3>n^2/48-5\).  The actual five-orbit decomposition of the
full cube gives
\[
 \sum_i\operatorname{ind}(x_i,[n]^3)
   >2n^3+n^2/24-20.
 \tag{FTE.5}
\]
After subtracting the ten errors and dividing by \(n^3\),
\[
 \frac1e\sum_i\operatorname{rank}(1-T_i)
 >2+\frac1{24n}-\frac{40M^2}{n^2}-\frac{20}{n^3}
 \ge2+\frac1{48n}-\frac{20}{n^3}
 >2+\frac3{n^3}\ge2+\frac6e.
 \tag{FTE.6}
\]
The weak inequality uses \(n\ge1920M^2\).  The last strict one
uses \(n^2>1104\), which follows from \(n\ge10^6\).
This proves the middle row without any unevaluated \(c,d,n_0\).

For tail three, HD-FT11 gives \(a\ge n^3/12\) and an error
\(48/n^{3/2}\) relative to the unordered-triple permutation
character, on at most ten entries once \(n\ge1024\).
Put \(B=\binom n3\).  Exact Riemann--Hurwitz on \(X_3\), together
with (FTE.2) and \(B\le n^3/6\), yields
\[
 \begin{split}
 \frac1e\sum_i\operatorname{rank}(1-T_i)
 &>2+\frac1{24n}-\frac{12}{n^3}-\frac{480}{n^{3/2}}\\
 &\ge2+\frac1{48n}-\frac{12}{n^3}
 >2+\frac{72}{n^3}\ge2+\frac6e.
 \end{split}
 \tag{FTE.7}
\]
For the first line keep the positive numerator together:
\((2g(X_3)-2)/B>(n^2/144-2)/B
\ge1/(24n)-12/n^3\), because \(n\ge10^6\) and
\(B\le n^3/6\).  No separate upper bound on \(2/B\) is used.
The second line uses
\(n\ge23040^2=530841600\), so \(480/\sqrt n\le1/48\).
Its strict inequality follows from \(n^2>4032\).
This proves the first row.  All estimates concern every target
eigenvalue and survive arbitrary tensor spectral convolution,
exactly as in the cited local lemmas.
\end{proof}

\paragraph{Exact per-label arithmetic.}
For tail three the values of \(M\) are \(30,18,42\) for
\(\mu=(3),(2,1),(1,1,1)\), respectively.  The sharper direct
tail-three argument uses the common threshold in (FTE.4), rather
than these coefficient constants.  For the other small tails the
complete values and thresholds are
\[
\begin{array}{c|r|r}
\mu\vdash4&M&N_\mu\\\hline
(4)&192&70778880\\
(3,1)&96&17694720\\
(2,2)&120&27648000\\
(2,1,1)&104&20766720\\
(1,1,1,1)&288&159252480
\end{array}
\qquad
\begin{array}{c|r|r|r}
\mu\vdash5&M&G&N_\mu\\\hline
(5)&1320&3487442&1000000\\
(4,1)&570&650942&1000000\\
(3,2)&504&509042&1000000\\
(3,1,1)&460&424122&1000000\\
(2,2,1)&600&721202&1000000\\
(2,1,1,1)&690&953582&1000000\\
(1,1,1,1,1)&2280&10401362&1000000
\end{array}
 \tag{FTE.8}
\]
\[
\begin{array}{c|r|r|r}
\mu\vdash6&M&G&N_\mu\\\hline
(6)&11760&276618722&3010560\\
(5,1)&4464&39863522&1142784\\
(4,2)&3040&18489282&1000000\\
(4,1,1)&3312&21945314&1000000\\
(3,3)&4224&35692802&1081344\\
(3,2,1)&2340&10955882&1000000\\
(3,1,1,1)&3360&22585922&1000000\\
(2,2,2)&5040&50813282&1290240\\
(2,2,1,1)&4000&32008002&1024000\\
(2,1,1,1,1)&6048&73168706&1548288\\
(1,1,1,1,1,1)&21600&933163202&5529600
\end{array}
 \tag{FTE.9}
\]
The maximal \(M\) at tail sizes \(3,4,5,6\) is respectively
\(42,288,2280,21600\); the maximal threshold is respectively
\(530841600,159252480,1000000,5529600\).
All bounds are sufficient thresholds, not claims of optimality.

\subsection*{HD-FT17. A numerical all-ordinary-tail conclusion}

\begin{corollary}[HD-FT17]
Assume HD-GC06.  In the actual factorwise setup (FTE.1), any
ordinary symmetric or alternating irreducible factor with
\[
                         n\ge N_*:=241^{482}
 \tag{FTE.10}
\]
must be extendible from \(S_n\), and its partition label satisfies
\[
                \min\{n-\lambda_1,n-\lambda'_1\}\le2.
 \tag{FTE.11}
\]
The conclusion includes arbitrary factorwise tensor partners and
all coefficient fields.  It does not assert existence of the
remaining sizes zero, one or two, or a bounded partner catalogue.
\end{corollary}

\begin{proof}
HD-FT02 removes all extendible labels with both tails at least241,
and HD-FT03 removes every split alternating constituent when
\(n>3500000\).  HD-FT06 already removes tails seven through240 at
\(n\ge241^{482}\), with its explicit uniform coefficient comparison.
Only sizes three through six remain to be addressed.  They are
excluded by HD-FT16 because
\(241^{482}>200^4=1600000000>530841600\).
This proves the corollary.  The actual quotient/factorization
hypothesis is retained throughout; no abstract normal component is
silently substituted for it.
\end{proof}

\paragraph{Boundary of the numerical conclusion.}
This addendum removes the unspecified triple-genus constants from
the stated ordinary-tail exclusions using the accepted HD-GC06.
It does not evaluate the exact Monderer--Neftin pair-classification
constant \(N_{14}\), replace the seven-mark conclusion of HD-FT15
by an unproved effective version, bound arbitrary original overgroups,
or construct a source or a detecting cycle.  HD-GC07's separate
eight-branch consequence has its own weaker effective
four-partner-entry conclusion.  The two assertions must not be conflated.
The bound \(N_*\) is not substituted for FQ or TW in a core whose
quasisimple factors are permuted: the required stable-factor quotient
and tensor factorization would first have to be established there.

\section{The entire tensor partner and a maximal-natural-degree tail-two gate}
\label{sec:entire-partner-support}

\paragraph{Status and exact scope.}
HD-FT18--19, reserved question Q2dr, is a derived consequence of the
accepted local support bounds and the actual tail-two gate.  The
tensor-normalizer argument was supplied independently in working
messages by \texttt{/root/imprimitive}; the full synthesis and the
maximal-natural-degree consequence are written here by
\texttt{/root/g34\_eightfolds}.  Complete author and separate root whole-proof readings pass, with the source/provenance/scope record read in full. No independent origin of the shared tensor-normalizer ingredient is claimed.  No new finite
certificate is needed or claimed for this universal argument.
No complete classification of possible
tensor partners is asserted.

For an operator \(S\) on a space \(E\), of finite order up to scalar, write
\[
 \nu(S)=\min_{z\in\mathbf C^*}\operatorname{rank}(S-zI)
       =\dim E-\max_z\dim\ker(S-zI).
 \tag{FTP.1}
\]
The quantity is unaffected by rescaling the operator.  Every operator
in a finite projective group can be rescaled to finite order, so is
semisimple and has the finite spectral distribution used below.

\subsection*{HD-FT18. A support floor independent of the number of partner factors}

Suppose the \emph{whole} partner space has a specified factorization
\[
 U=\bigotimes_{j=1}^s U_j,\qquad d_j=\dim U_j\ge2,
 \qquad b=\dim U=\prod_{j=1}^s d_j,
 \tag{FTP.2}
\]
and its actual finite projective action permutes the embedded full
matrix algebras \(\operatorname{End}(U_j)\otimes I\).
In particular there is no additional unbounded multiplicity factor
outside (FTP.2).  Fix positive numbers \(\eta_j\) such that every
compatible nonscalar operator on the \(j\)-th factor, whenever the
global operator fixes all factor algebras, has support at least
\(\eta_jd_j\).  Define
\[
                 \eta=\min\{1/4,\eta_1,\ldots,\eta_s\}.
 \tag{FTP.3}
\]

\begin{lemma}[HD-FT18]
Every nonscalar operator \(S\) in this actual projective
tensor-normalizer action satisfies \(\nu(S)\ge\eta b\).
Consequently, if all \(d_j\le A\), with \(A\ge2\), one may use
\[
                         \eta=1/\max\{A,4\}.
 \tag{FTP.4}
\]
The number of factors is arbitrary.  A genuine nonidentity scalar
on the whole source has actual moving rank equal to its full degree;
it is not assigned the projective support in (FTP.1).
\end{lemma}

\begin{proof}
An automorphism of a full complex matrix algebra is conjugation by
an invertible matrix.  Thus, after removing the actual permutation
of the factor algebras, an implementing operator is a tensor product
of factor operators, up to a scalar.  This is the matrix-algebra
normalizer argument already proved in D-BI19 and used in D-ATT02.
Factors in a permutation cycle have equal dimension.

If the operator fixes every factor algebra and is nonscalar, at
least one factor operator is nonscalar.  The maximum atom of a
convolution of finite probability distributions is at most the
maximum atom of each individual distribution.  Applying this to
the normalized spectral multiplicities gives
\(\nu(S)/b\ge\eta_j\) at that active factor.

Suppose instead there is a cycle of length \(\ell\ge2\), on
factors of degree \(d\ge2\).  Conjugate by factorwise operators so
that the cycle has only its cycle product as a twist.  This product
has finite projective order: a scalar power of the global operator,
after a multiple of the cycle length, makes the corresponding power
of every factor product scalar.  This follows because a scalar tensor
operator acts trivially by conjugation on every factor matrix algebra.
The cycle product can therefore be diagonalized after a scalar
normalization.  In a tensor-word basis, the resulting operator is
a weighted cyclic permutation.  On each word orbit an eigenvalue
has multiplicity at most one, since the successive nonzero weights
determine an eigenvector from its first coordinate.

There are exactly \(d\) constant words.  Every other word orbit
has length at least two, so every eigenspace on this cycle has
dimension at most
\[
 d+\frac{d^\ell-d}{2}
       =\frac{d^\ell+d}{2}\le\frac34d^\ell.
 \tag{FTP.5}
\]
The last inequality uses \(d^{\ell-1}\ge2\).  Convolution with
the remaining tensor cycles cannot increase the largest normalized
eigenspace.  Hence \(\nu(S)/b\ge1/4\).

Finally every nonscalar semisimple operator in degree \(d_j\)
has an eigenspace of codimension at least one, so
\(\eta_j=1/d_j\ge1/A\) is always available in the bounded-degree
case.  This proves (FTP.4) and the lemma.  For a genuine scalar,
the asserted full actual moving rank is immediate.
\end{proof}

Both local constants have actual operator witnesses.  A factorwise
diagonal involution with one negative line has normalized support
\(1/d\); a swap on \(\mathbf C^2\otimes\mathbf C^2\) has
eigenspace dimensions three and one, hence support \(1/4\).
These are operator statements, not asserted sphere-cover sources.

\subsection*{HD-FT19. Consequences for the actual tail-two source}

Retain exactly the actual source of HD-FT13--15: a connected sphere
cover, an actual quotient onto \(P=A_n\) or \(S_n\), and an actual
stable factorization
\[
 V=W\otimes U\quad\hbox{up to scalars},\qquad
 W=S^{(n-2,2)}\ \hbox{or}\ S^{(n-2,1,1)},\qquad
 e=\dim V,\qquad r\ge12,\qquad
 \sum_{i=1}^r\operatorname{rank}(1-T_i)=2e+6.
 \tag{FTP.6}
\]
The transposed labels and irreducible ordinary alternating
restrictions are permitted.  Every original inertia is genuine;
the projected tuple is not assigned twelve genuine marks.
The whole partner has the actual tensor-normalizer structure (FTP.2).

For \(n\ge10^6\), accepted HD-GC07 and HD-FT15 give at least four
genuine entries with identity projection to \(P\), and their total
normalized actual moving rank is at most \(61/n^2\).  If instead
one uses the qualitative Monderer--Neftin pair classification at
\(n>N_{14}\), with \(n\ge64\), the same conclusion has at least
seven such entries.  Put \(q=4\) in the effective case and \(q=7\)
in that separately scoped qualitative case.

\begin{theorem}[HD-FT19]
Under these actual hypotheses,
\[
                             q\eta\le\frac{61}{n^2}.
 \tag{FTP.7}
\]
In particular, if the entire partner consists of factors of degree
at most \(A\), independently of their number, then
\[
                     qn^2\le61\max\{A,4\}.
 \tag{FTP.8}
\]

More generally, suppose every partner factor is of one of the
following two kinds:
\begin{itemize}
\item a factor of degree at most \(A\), where \(A\ge2\);
\item a standard or ordinary tail-two representation of \(A_m\),
where \(15\le m\le n\), whose actual simple-core action is
normalized by every compatible local factor operator.
\end{itemize}
The actual image must still permute the full factor matrix algebras
and preserve the indicated simple-core actions when it fixes those
factors.  Then
\[
                  qn^2\le61\max\{4,A,n-1\}.
 \tag{FTP.9}
\]
For the effective case \(q=4\), \(n\ge10^6\), this forces
\[
                             A\ge\frac{4n^2}{61}.
 \tag{FTP.10}
\]
In particular there is no source in which the distinguished actual
stable tail-two factor has maximal natural degree and the entire
partner consists only of standard and tail-two factors, regardless
of the number of those factors.  Bounded additional factors are
also excluded once \(n^2>61A/4\).

Every conclusion also permits additional partner factors, of any
degree, for which the entire compatible local normalizer has a
proved normalized projective support floor of at least \(1/4\).
This is an explicit support hypothesis, not an assertion about an
unclassified residual factor.
\end{theorem}

\begin{proof}
At an identity-projection entry, \(T_i\) is a scalar multiple of
\(I_W\otimes S_i\).  If \(S_i\) is nonscalar, its actual rank is
at least \(e\nu(S_i)/b\ge\eta e\), with every scalar phase
allowed.  If \(S_i\) is scalar, genuineness of \(T_i\) instead
gives full rank \(e\ge\eta e\).  Summing over the at least \(q\)
entries and using the accepted upper budget proves (FTP.7).
Equation (FTP.8) follows from (FTP.4).

For the ordinary factors, \(\operatorname{Aut}(A_m)=S_m\) at
\(m\ge15\).  Faithfulness, irreducibility, and the chosen
ordinary extension show that every compatible local normalizer
is a scalar times the corresponding \(S_m\) operator, as in
D-ATT02 and D-AST02.  On the standard module, projective
faithfulness and integral codimension give a floor \(1/(m-1)\).
The exact D-ATT01 minima give respectively
\[
 \frac{\nu(S^{(m-2,2)}(x))}{m(m-3)/2}\ge\frac2m,
 \qquad
 \frac{\nu(S^{(m-2,1,1)}(x))}{(m-1)(m-2)/2}
                                  \ge\frac2{m-1}
 \tag{FTP.11}
\]
for every nonidentity \(x\in S_m\).  Both exceed or equal
\(1/(m-1)\ge1/(n-1)\).  The identity permutation is a scalar
local operator and is already handled separately.  Sign choices
of extension do not affect projective support.

Thus HD-FT18 gives
\(\eta\ge1/\max\{4,A,n-1\}\), proving (FTP.9).
Additional factors with the stated quarter-support floor do not
decrease this value, by the definition (FTP.3).
For \(n\ge10^6\), one has \(4n^2>61(n-1)\) and
\(4n^2>244\), so neither \(n-1\) nor four can satisfy the
maximum in (FTP.9) with \(q=4\).  This forces (FTP.10).
When no bounded-degree pieces are present, one may put \(A=2\);
the resulting contradiction proves the stated maximal-natural-degree
case.  If \(\dim U=1\), there are no nonscalar partner factors,
but all \(q\) entries have full rank, immediately contradicting
\(q\le61/n^2\).  This also covers the empty tensor product.
\end{proof}

\paragraph{Strictness and retained boundary.}
The earlier partner-dimension gate alone says
\(\dim U\ge4n^2/61\) in the effective case.  It does not prevent
\(\dim U=2^s\) with arbitrarily many two-dimensional factors.
Equations (FTP.7)--(FTP.8) exclude that entire actual tensor-normalizer
class at \(n\ge10^6\), since its floor is \(1/4\), independently
of \(s\).  The displayed dimensions can satisfy the old numerical
inequality, but no cover or branch tuple is inferred from them.
This is a strict narrowing of that previously unrestricted
large-partner family under explicit decomposition hypotheses.

The theorem does not assert that every possible partner admits such
a complete factorization, that all remaining factors have bounded
degree, or that an unbounded residual multiplicity space can be
discarded.  It also does not choose a distinguished stable factor
from a group that permutes it with the rest: the actual quotient
and factorization (FTP.6) remain indispensable.  The exact
Monderer--Neftin threshold \(N_{14}\) remains unevaluated; the
effective and qualitative counts retain their separate ranges.

\section{Five-dimensional inducing blocks: one remaining Gaussian group}
\label{sec:rank-five-inducing}

\paragraph{Scope and check status.}
This is a new reduction, recorded on 21 November 2026.  It is a
necessary-condition theorem, not a construction of a dominant family.
The scalar-field and Chevalley--Weil calculations have an exact finite
certificate.  The parent independently read the complete proof and
reran the finite certificate, recording PASS for both on 21 November
2026.  The G34 agent independently read D-R5-01--03, the cited
Cohen rank-five proof, and a separate field/spectrum certificate;
these checks also passed.  The later D-R5-04 computation and
geometric conclusion also passed an independent whole-proof review,
an independent all-eighty-class local-form calculation, and fresh
exact group enumeration by the G34 agent.  The parent separately
reran its exact group and local-form certificates.
The published classification input is Cohen's rank-five theorem and
normalizer results, with the 1978 correction; the real-boundary argument
uses Lewis--Reiner's published equal-pair extraction theorem.

\subsection*{D-R5-01. The reflection closure and the real boundary}

Suppose a finite-cover source has an exact imaginary quadratic character
field $K$, actual Weil sixfold signature $(3,3)$, period rank nine,
and nonsplit actual polarization.  Suppose its complex coefficient
representation is induced from an irreducible representation $\tau$
of dimension five.  Write $J$ for the faithful image of $\tau$ and
$L=\mathbf Q(\chi_\tau)$.  We prove that
\[
                         K\in\{\mathbf Q(i),\mathbf Q(\sqrt{-3})\}.
                                                               \tag{R5.1}
\]
Moreover, if $K\ne\mathbf Q(\sqrt{-3})$, then necessarily
\[
 \begin{gathered}
 K=\mathbf Q(i),\qquad L=\mathbf Q(\zeta_{12}),\qquad
 R=G_{33},\qquad J=\mu_4G_{33},\\
 |J|=103680,
 \end{gathered}                                                  \tag{R5.2}
\]
where $R$ is the normal subgroup generated by all reflections of $J$.
The intermediate rational base has exactly twelve branch points,
ten with reflection inertia and two with inertia of order four and
moving rank three.  The latter two spectra are respectively
\[
                 (1,1,-1,-1,i),\qquad(1,1,-1,-1,-i).
                                                               \tag{R5.3}
\]
These assertions allow arbitrary original representation degree and
do not assume that the inducing representation has a $K$-model.

If $\tau$ is imprimitive, its prime dimension makes its blocks
one-dimensional.  Transitivity of induction makes the original
representation monomial, and D-MON excludes a dominant nonsplit
source.  Thus $\tau$ is primitive.
Pass to the intermediate connected $J$-cover $C_0\to X$.
Frobenius reciprocity identifies its complex multiplicity variation
with the original one; the fixed polarization determines the
conjugate filtration, so the full original period rank factors
through this intermediate-cover moduli, as in D-BI-01.
If $h=g(X)$, the Euler equation is
\[
                            D=16-10h.
\]
For $h=1$ there are at most six branch points and at most six
marked-base parameters.  Larger $h$ is impossible.  Thus rank nine
forces $h=0$ and
\[
 r\ge12,\qquad D=16,\qquad
                  \sum_j\bigl(\operatorname{rank}(1-g_j)-1\bigr)
                       =16-r\le4.                           \tag{R5.4}
\]
In particular there are at least eight reflection inertias.

The reflection normal closure $R$ is nontrivial and irreducible.
Indeed its fixed space is $J$-stable and must vanish; the
nonorthogonality components of its reflection lines give an
orthogonal direct-sum decomposition permuted by $J$, which
primitivity forbids unless there is only one component.
If $R$ is imprimitive, Cohen's normalizer theorem
\cite[Proposition~2.13 and its proof, pp.~391--392]{FIELD:Cohen}
puts $J$ in a monomial group: none of its three exceptional
dimensions is five.  This again contradicts the already excluded
monomial case.

We next exclude every real $R$, without needing a list of real
rank-five Coxeter types.  In real orthogonal coordinates for $R$,
every unitary normalizer is scalar times a real orthogonal matrix,
by \cite[Lemma~2.12 and its proof, p.~391]{FIELD:Cohen} and
uniqueness of the positive invariant form.
There must be at least two inertia matrices with nonreal scalar
phase.  Otherwise all inertia matrices are real, or their
product is a nonreal scalar times a real matrix; the first case
would make $\chi_\tau$ and its induced character real, and the
second cannot have product one.
Such an inertia has fixed dimension at most two: the eigenspaces
of the corresponding real orthogonal matrix for its two distinct
conjugate nonreal eigenvalues have equal dimension.
Its moving rank is therefore at least three.
Equation (R5.4) now forces exactly twelve branch points, with ten
real reflections and two nonreal scalar inertias of rank three.
Hurwitz moves put the latter two at the end.  Write them as
\[
                 g_1=\lambda A_1,
                 \qquad g_2=\lambda^{-1}A_2,
                                                               \tag{R5.5}
\]
where $A_i$ are real orthogonal and $\lambda$ is nonreal.
Their original phases have real product $\pm1$, since
$g_1g_2\in R$; absorb the sign into $A_2$ to obtain (R5.5).

Put $w=(g_1g_2)^{-1}\in R$.  Lewis--Reiner
\cite[Corollary~1.4, p.~1459; proof \S4.3, pp.~1480--1481]{BI:LR}
changes the length-ten reflection factorization of $w$ by Hurwitz
moves into
\[
              (r_1,r_1,\ldots,r_q,r_q,s_1,\ldots,s_\ell),
              \qquad q=(10-\ell)/2,                         \tag{R5.6}
\]
with a shortest final factorization of length $\ell$.
Carter's independent-root criterion in
\cite[Proposition~4.1, p.~1473]{BI:LR} gives $\ell\le5$;
parity gives $\ell\in\{0,2,4\}$.
The telescoping product identity and independence of the successive
transformed roots show
\[
             \ker(w-1)=\bigcap_j\ker(s_j-1),
             \qquad \dim\ker(w-1)=5-\ell.
\]
Let $H_0$ be generated by the $s_j,g_1,g_2$, and let
$E=\ker(g_1-1)\cap\ker(g_2-1)$, of dimension $h_0$.
The preceding fixed-space identity gives $E=\tau^{H_0}$.
Both $E$ and $\overline E$ lie in the real fixed space of $w$;
they are disjoint because $A_1$ has the distinct eigenvalues
$\lambda^{-1}$ and $\lambda$ on them.  Consequently
\[
                         2h_0\le5-\ell.
                                                               \tag{R5.7}
\]

Bubble off the $q$ equal reflection pairs on rational leaves of
the base, each attached by a node of trivial inertia.  The source
dual graph is connected because the full tuple generates $J$.
Its edge module is $\mathbf C[J]^q$, and its vertex module is
$\mathbf C[J/H_0]\oplus\bigoplus_j\mathbf C[J/\langle r_j\rangle]$.
Taking $\tau$-multiplicities in the connected graph's cellular
sequence gives
\[
 \dim\operatorname{Hom}_J(\tau,H_1(\Gamma,\mathbf C))
       =5q-4q-h_0=q-h_0
       \ge\frac{10-\ell}{2}
                      -\left\lfloor\frac{5-\ell}{2}\right\rfloor
       =3.                                                    \tag{R5.8}
\]
Smooth all nodes with positive integer orders.  The positive
diagonal edge-length pairing is nondegenerate on the graph cycle
space, so the resulting single Picard--Lefschetz logarithm $N$
has exactly this rank on the coefficient multiplicity.  It has
$N^2=0$ on a six-dimensional space; hence its rank is at most
three, and (R5.8) forces rank three.

This is maximal isotropy over the original field $K$ as well.
The intermediate Hurwitz component has dimension nine, and the
original period rank nine makes its classifying map dominant.
Take a closed point of the generic fibre of that classifying map
and normalize its closure.  After shrinking it gives a finite
etale cover of a dense open of the intermediate component carrying
the original $K$-variation.  Choose the node-smoothing arc generically
so its punctured part lies in that open; after a finite base change
it lifts to the finite cover.  The logarithm on the original
$K$-variation is thus a positive integer multiple of $N$ after
the chosen complex embedding.  It is $K$-linear, square zero,
and has $K$-rank three.  Skew adjointness for the actual Hermitian
polarization gives
$h(Nx,Ny)=-h(x,N^2y)=0$.
Its image is a three-dimensional isotropic $K$-space, proving
splitness and contradicting the hypothesis.
This is the graph and finite-multisection argument of D-BI-12
with the rank-five numbers written out; it does not descend an
arbitrary $L$-isotropic subspace by assertion.

Therefore $R$ is complex primitive of rank five.  Cohen's complete
rank-five result identifies it with
\[
                             R=G_{33}=W(K_5)
                                                               \tag{R5.9}
\]
\cite[\S5.10 and proof, pp.~430--432; Table~4.16, p.~412]{FIELD:Cohen}.
Its reflection model is over $F=\mathbf Q(\omega)$, all its
reflections have order two, its degrees are $4,6,10,12,18$, and
its order is $51840$.  The model and these data appear in
\cite[\S4.15, pp.~409--410; \S4.16, pp.~411--413]{FIELD:Cohen}.
Its exact character field is $F$: it is contained in $F$, and a
complex reflection group of this type has nonreal character by
\cite[\S1.4(vi) and proof, p.~382]{FIELD:Cohen}.
Its scalar center is $\mu_2$, by the degree-gcd formula
\cite[\S1.4(ii), p.~381]{FIELD:Cohen}.

\subsection*{D-R5-02. Scalar-coupled subgroups and the exact branch profile}

The corrected scalar-normalizer result gives
\[
                       R\subset J\subset\mathbf C^*R
\]
\cite[Proposition~5.14, p.~435]{FIELD:Cohen}
\cite[p.~613]{FIELD:CohenErr}.
There is no unexamined subgroup-extension assumption here.
If $g=\lambda w\in J$, with $w\in R$, then
$\lambda I=gw^{-1}\in J$ because $R\subset J$.
Consequently $J=CR$, where $C$ is its finite cyclic scalar subgroup.
In the fixed $F$-model of $R$, put
\[
                 J_F=J\cap\operatorname{GL}_5(F)
                     =(C\cap F^*)R.
                                                               \tag{R5.10}
\]
The quotient $J/J_F$ is cyclic.  Membership in $J_F$ is independent
of the decomposition $g=\lambda w$, since its ambiguity is a scalar
in $R$, namely an element of $\mu_2\subset F$.

Induction of characters gives $K\subset L$.  Since $R\subset J$,
one also has $F\subset L$.  If $J=J_F$, its matrices are over $F$
and hence $L=F$, which forces $K=F$.
Otherwise at least two branch entries are outside $J_F$: they
generate a nontrivial quotient and their product there is one.
For such an entry write $g=\lambda w$, $\lambda\notin F$.
If it fixed at least three dimensions, the minimal polynomial of
$\lambda^{-1}$ over $F$ would occur at least three times in the
degree-five characteristic polynomial of $w$.  Its degree would
be at most one, a contradiction.  Thus each such entry moves at
least three dimensions.
Equation (R5.4) now forces
\[
 r=12,\quad \text{ten reflection entries},\quad
 \text{two entries outside }J_F\text{, each of moving rank three}.
                                                               \tag{R5.11}
\]
All other entries are in $R$.  Moving the two exceptional entries
to the end, their product lies in $R$; hence either one, say
$g=\lambda w$, generates $J$ together with $R$.
It follows that
\[
                         J=\langle R,\lambda I\rangle,
                         \qquad L=F(\lambda).
                                                               \tag{R5.12}
\]
The equality of fields uses both directions: all matrices are
over $F(\lambda)$, while $F\subset L$ and the scalar trace
$\operatorname{Tr}(\lambda I)=5\lambda$ puts $\lambda$ in $L$.

Now $\lambda^{-1}$ has multiplicity exactly two on $w$.
The same repeated-minimal-polynomial argument gives
$[F(\lambda):F]=2$.
If $d$ is the order of $\lambda$, then $\varphi(d)\le4$ and
\[
                         d\in\{1,2,3,4,5,6,8,10,12\}.
\]
The values $1,2,3,6$ lie in $F$ and are excluded.
For $d=5,10$, the sole quadratic subfield of $\mathbf Q(\zeta_d)$
is $\mathbf Q(\sqrt5)$; for $d=8$ its quadratic subfields are
$\mathbf Q(i),\mathbf Q(\sqrt2),\mathbf Q(\sqrt{-2})$.
None is $F$, so these three values would give extension degree
four over $F$ and are excluded.  Thus
\[
                    d\in\{4,12\},\qquad
                    L=F(i)=\mathbf Q(\zeta_{12}).
                                                               \tag{R5.13}
\]
Its only imaginary quadratic subfields are $\mathbf Q(i)$ and
$F$, proving (R5.1).

We can sharpen (R5.13).  The nontrivial automorphism of
$F(\lambda)/F$ sends $\lambda$ to $-\lambda$ in both remaining
cases.  Thus $w$ has eigenvalues
\[
             \lambda^{-1},\lambda^{-1},
             -\lambda^{-1},-\lambda^{-1},\mu,
             \qquad\mu\in F.
\]
Every element of $R$ has determinant $\varepsilon\in\{1,-1\}$,
because its generating reflections have order two.
Taking the product of these five eigenvalues gives
$\mu=\varepsilon\lambda^4$.
Therefore
\[
             \operatorname{Spec}(g)
                    =(1,1,-1,-1,\varepsilon\lambda^5).
                                                               \tag{R5.14}
\]
If $d=12$, then $g^2$ is a reflection whose nonidentity eigenvalue
$\lambda^{10}$ has order six.  It belongs to $J$ and hence to its
reflection normal closure $R$, contradicting the order-two
reflection list of $G_{33}$.  Thus $d=4$.
Equations (R5.12)--(R5.14) give $J=\mu_4R$ and an outside inertia
spectrum $(1,1,-1,-1,\pm i)$ of order four.  Its square is a
reflection in $R$.
The ten reflection determinants have product one, so the determinants
of the two outside inertias are reciprocal.  They must have the
opposite signs in (R5.3).

Since $R\cap\mu_4=\mu_2$, one has $J/R\simeq C_2$ and
$|J|=2\cdot51840=103680$.
The quotient cover $C_0/R\to\mathbf P^1$ has exactly two branch
points, both with full order-two inertia, and hence genus zero.
The local ages in (R5.3) are $5/4$ and $7/4$; the ten reflections
have age $1/2$ each.  Total age is $8=5+3$, as is the total dual
age, so this necessary profile has the required Chevalley--Weil
multiplicities $(3,3)$.
This verifies the signatures of a possible profile; it does not
construct a product-one generating tuple.

\subsection*{D-R5-03. What the Gaussian case still requires}

For $K=\mathbf Q(i)$, the preceding conclusions freeze the remaining
five-dimensional gate to (R5.2)--(R5.3).
The natural $L$-factor of a generic $J$-cover in a relevant
twelve-branch component is an abelian twelvefold $A_\tau$, with
$H^1(A_\tau)$ of rank six over $L$.
If it actually arises by induction from an original $K$-sixfold
$B$, the rational representation identity is
\[
 \operatorname{Ind}_H^G\operatorname{Res}_{L/\mathbf Q}\tau
       \simeq
       \left(\operatorname{Res}_{K/\mathbf Q}\rho\right)^{\oplus2}.
                                                               \tag{R5.15}
\]
Indeed scalar extension of the $K$-model of $\rho$ to $L$ is
$\operatorname{Ind}\tau$, and restriction back to $\mathbf Q$
has multiplicity $[L:K]=2$.  Frobenius reciprocity on rational
Hodge structures therefore gives, in the original family,
\[
                              A_\tau\sim B^2.
                                                               \tag{R5.16}
\]
This is an isogeny of variations after finite choices of primitive
projectors, not just a comparison of complex dimensions.
Since the original rank-nine family dominates the nine-dimensional
intermediate Hurwitz component, this would force its generic
$A_\tau$, after finite base change, to have an endomorphism algebra
containing $M_2(\mathbf Q(i))$.

No theorem proving or disproving that generic extra endomorphism
condition for every remaining component is asserted here.
Likewise the actual original $K$-polarization has not been shown
split in this complex $G_{33}$ case.  A proof of generic
$\operatorname{End}^0(A_\tau)=L$, or a suitable maximal-rank
$K$-linear boundary logarithm, would exclude the remaining case;
an actual source construction would instead have to verify the
isogeny/field descent and the period rank.  Neither follows from
the branch-count equality or the Chevalley--Weil profile alone.

\subsection*{D-R5-04. Exact local classes, a norm determinant, and one excluded component}

\paragraph{Status.}
This additional proof and its two exact computations were completed
on 21 November 2026.  The G34 agent independently checked the full
proof, all eighty local forms, and a fresh exact enumeration of the
51840-element group and its forty conjugacy classes, recording PASS.
The parent independently reran the group and local-form certificates.
It constructs a natural $L$-cover, not an original $K$-sixfold source.

Use the coordinates $(x_1,x_2,x_3,x_4,t)$ for
$(e_5-e_6)^\perp\subset\mathbf C^6$, where the last two original
coordinates both equal $t$.  The positive Hermitian form is
\[
                         H_0=\operatorname{diag}(1,1,1,1,2).
\]
After removing irrelevant scalar multiples from Cohen's five roots
in \cite[\S4.15, p.~409]{FIELD:Cohen}, choose the norm-two roots
\[
 \begin{split}
 v_1&=(0,0,0,0,1),\\
 v_2&=\tfrac12(-1,1+2\omega,1,1,1),\\
 v_3&=(1,-1,0,0,0),\quad v_4=(0,1,-1,0,0),\quad
 v_5=(0,0,1,-1,0).
 \end{split}
\]
Their Gram matrix and determinant are
\[
 H=\begin{pmatrix}
 2&1&0&0&0\\
 1&2&\omega&\omega^2&0\\
 0&\omega^2&2&-1&0\\
 0&\omega&-1&2&-1\\
 0&0&0&-1&2
 \end{pmatrix},\qquad \det H=2.
                                                               \tag{R5.17}
\]
In this basis the generating reflections are
$s_j=I-e_j H_{j,*}$; their entries lie in $\mathbf Z[\omega]$.
The exact program \texttt{g33/enumerate\_g33.cpp} enumerates their
matrix closure by right multiplication, checks closure and order
$51840$, and partitions it by conjugation by the same generators.
All arithmetic is integral in $\mathbf Z[\omega]$.
It finds forty conjugacy classes and exactly forty-five reflections.
Characteristic polynomials are computed by exact Newton identities;
the second program independently checks them by six exact determinant
evaluations per class.

There are exactly two classes with characteristic polynomials
\[
                 (x-1)(x^2+1)^2,
                 \qquad (x+1)(x^2+1)^2,
                                                               \tag{R5.18}
\]
each of size $540$ and order four.  A representative of the first is
the explicit word
\[
               w=s_1s_2s_1s_3s_2s_1s_4s_2s_5s_4.
                                                               \tag{R5.19}
\]
The existence is also visible without enumeration: Cohen's root list
contains the $D_4$ roots in the first four coordinates
\cite[pp.~409--410 and the $D_4$ stabilizer argument, p.~413]{FIELD:Cohen}.
The even signed permutation
\[
        j(x_1,x_2,x_3,x_4,t)=(x_2,-x_1,x_4,-x_3,t)
\]
lies in that $W(D_4)$ subgroup and has the first characteristic
polynomial in (R5.18).  Uniqueness of the enumerated class therefore
identifies its class with that of (R5.19).

Since adjoining the central scalar $i$ does not change conjugation
on $R$, the two cosets $R$ and $iR$ give exactly eighty conjugacy
classes in $J=\mu_4R$.  The program \texttt{g33/local\_forms.py}
works in the exact field $L=\mathbf Q(\omega,i)$ and, for every
class, computes a column basis $C$ of $\operatorname{im}(1-g)$,
its actual Hermitian Gram matrix $C^*HC$, and
$b(g)=\det(1-g|_{\operatorname{im}(1-g)})$.
It checks all eighty classes and finds exactly two outside classes
of moving rank three, namely those in (R5.18) multiplied by $i$.
There are no outside reflections.

For the two forced outside types, use the representatives $ij$ and
$-ij$.  Moving bases for them are, respectively,
\[
 (1,i,0,0,0),\ (0,0,1,i,0),\ (0,0,0,0,1)
\]
and the same vectors with $i$ replaced by $-i$.
Both moving Gram matrices are $\operatorname{diag}(2,2,2)$.
Consequently the determinant classes and denominators are
\[
 \det H_{\mathrm{mov}}=8,\qquad
 b(ij)=4(1-i),\qquad b(-ij)=4(1+i).
                                                               \tag{R5.20}
\]
The coordinate-column bases selected by the exact program give
determinants $16$ and $8(2+\sqrt3)$ instead; these agree with $8$
modulo norms because
\[
 N_{L/L^+}(1+i)=2,\qquad
 N_{L/L^+}(1-i\omega)=2+\sqrt3,
 \qquad L^+=\mathbf Q(\sqrt3).
\]
These norm equalities are checked in the field arithmetic as well.

The proof of D-DET-01 works verbatim for a CM quadratic extension
$L/L^+$: every occurrence of $\mathbf Q$ in its Hermitian determinant
argument is replaced by $L^+$, while complex conjugation is the CM
involution.  No geometric or linear-algebra step uses $L^+=\mathbf Q$.
Every reflection moving line has determinant class $2$, and its
$b$-factor is $2$.  For any product-one tuple with the forced
ten-plus-two profile, (R5.20) therefore gives signed parabolic
discriminant
\[
 \left[
 \frac{2^{10}\,8^2}
      {2^{10}\,4(1-i)\,4(1+i)}
 \right]
                  =[2]=[1]
             \quad\text{in }(L^+)^*/N_{L/L^+}(L^*).
                                                               \tag{R5.21}
\]
Thus no obstruction to rational determinant descent is obtained:
the determinant already has a rational representative which is a
Gaussian norm.

This conclusion does not descend splitness to $K=\mathbf Q(i)$.
For a concrete warning about that inference, the Gaussian Hermitian
space $U^2\perp\langle1,-3\rangle$ has signature $(3,3)$ and
nonsplit signed discriminant $3$, since $3$ is not a Gaussian norm.
After adjoining $\sqrt3$ it is split, because the last plane has
isotropic vector $(\sqrt3,1)$.  Hence even a split Hermitian space
over $L$ can be obtained from a nonsplit Gaussian space.

There is nevertheless an explicit natural $L$-cover and an exact
component obstruction.  With the word $w$ in (R5.19), set $g=iw$.
The twelve-entry tuple
\[
             (s_1,s_1,\ldots,s_5,s_5,g,g^{-1})
                                                               \tag{R5.22}
\]
has product one and generates $J$: its reflection entries generate
$R$, and $gw^{-1}=iI$ supplies the scalar.  The exact program
verifies the word, product and scalar extraction.
It has the required local spectra and coefficient multiplicities
$(3,3)$; its connected $J$-cover has genus
\[
 1+\frac{103680}{2}
        \left(-2+10\left(1-\frac12\right)
                    +2\left(1-\frac14\right)\right)
                 =233281.
\]
This does not produce a Gaussian sixfold: its natural rational
factor is the $L$-abelian twelvefold $A_\tau$ of D-R5-03.

Finally bubble off the five doubled reflection pairs in (R5.22).
The central component has inertia group $\langle g\rangle$, whose
fixed space on $\tau$ has dimension two.  The connected source
graph therefore has coefficient multiplicity
\[
                      5\cdot5-5\cdot4-2=3.
                                                               \tag{R5.23}
\]
Its positive node smoothing gives a square-zero logarithm of rank
three.  If a rank-nine original Gaussian source dominated this
particular intermediate Hurwitz component, the finite-multisection
argument in D-R5-01 would lift that logarithm to the original
$K$-variation and make its actual Hermitian space split.
Thus the component containing (R5.22) cannot supply a dominant
nonsplit original Gaussian sixfold.  No transitivity assertion
identifying every possible Nielsen component with this one is
made; that, or an independent generic-endomorphism calculation,
remains the narrower residual gate.

\section{The exact larger-field profiles for an inducing $E_8$ block}
\label{sec:e8-inducing-fields}

\begin{theorem}[HD-E8IF1]
Let a nonsplit Weil sixfold source of exact imaginary quadratic
character field $K$ have period rank nine. Suppose its coefficient
representation is induced from an irreducible eight-dimensional
representation $\tau$ whose reflection normal closure is $W(E_8)$.
Assume that $K$ is one of the eight non-Eisenstein class-number-one
fields and that the inducing character field
$L=\Q(\chi_\tau)$ is strictly larger than $K$.
Then the complete list of necessary field and branch profiles is
\[
\begin{array}{c|c|c|c}
 K&L&\text{faithful inducing image }J&
     \text{real class of either outside factor}\\\hline
 \Q(i)\text{ or }\Q(\sqrt{-2})&\Q(\zeta_8)&
       \mu_8W(E_8)&D_8(a_3),\quad\Phi_8^2\\
 \Q(i)&\Q(\zeta_{12})&\mu_{12}W(E_8)&
       E_8(a_3),\quad\Phi_{12}^2.
\end{array}                                                    \tag{E8IF.1}
\]
In both rows the intermediate base is rational and has exactly
twelve branch points: ten real reflections and two outside inertias
of moving rank six. Write the outside elements as
$g_1=\lambda w_1$, $g_2=\varepsilon\lambda^{-1}w_2$,
where $w_i\in W(E_8)$ and $\varepsilon\in\{1,-1\}$.
In the first row $\lambda$ is primitive of order eight and each
outside spectrum is
\[
                  1^2,i^2,(-1)^2,(-i)^2.                    \tag{E8IF.2}
\]
In the second row $\lambda$ is primitive of order twelve and,
up to exchanging the two entries, their spectra are
\begin{align*}
 &(1,1,\zeta_6,\zeta_6,-1,-1,\zeta_3^2,\zeta_3^2),\\
 &(1,1,\zeta_6^{-1},\zeta_6^{-1},-1,-1,\zeta_3,\zeta_3).
\end{align*}
Equivalently the distinct exponent
fractions are $\{0,1/6,1/2,2/3\}$ and
$\{0,1/3,1/2,5/6\}$, each repeated twice.
These necessary profiles have Chevalley--Weil signature $(3,3)$.
They do not assert the existence of a generating tuple, a dominant
component, or a nonsplit actual polarization.
\end{theorem}

\begin{proof}
The normalizer of the real $E_8$ reflection group is scalar times
$W(E_8)$: reduce to a real orthogonal normalizer by Cohen's
Lemma~2.12, then use the chamber argument and the absence of
a diagram automorphism, as in D-BI-13. The group already contains
$-I$ and is defined over $\Q$. Consequently
$J\subset\C^*W(E_8)$ and $W(E_8)\subset J$.
Every finite-order element of $J$ is $\lambda w$ with $w\in W(E_8)$
and $\lambda$ a root of unity. The scalar is determined modulo
$\{1,-1\}$, and this class defines a character of $J$.
Indeed the ambiguity is exactly the scalar center of $W(E_8)$.

Frobenius reciprocity identifies the original complex multiplicity
variation with the actual intermediate-cover variation for $\tau$.
Its Euler equation is $6=16(h-1)+D$. Genus at least two is impossible;
at genus one the marked-base dimension is at most six. Thus rank nine
forces genus zero and a genuine product-one generating inertia tuple
with
\[
             r\ge12,\qquad\sum_{j=1}^r\operatorname{rank}(1-g_j)=22.
                                                               \tag{E8IF.4}
\]
Let $t$ count the inertias outside $W(E_8)$. Such an inertia has a
nonreal scalar phase. Since $w$ is real orthogonal, a nonreal
eigenvalue and its conjugate have equal multiplicities, at most four.
Each outside inertia therefore has moving rank at least four.
Every other genuine inertia has moving rank at least one, so
$22\ge r+3t$, and hence $t\le3$.
The scalar product-one relation excludes $t=1$; $t=0$ would give
$J=W(E_8)$ and a real character field, contrary to $K\subset L$.

Suppose $t=3$. The three outside moving ranks sum to at most thirteen;
each is at least four, so each is at most five. Each corresponding
$w_j$ has $\lambda_j^{-1}$ as an eigenvalue with multiplicity at
least three. Since $w_j$ is rational, the minimal polynomial of
$\lambda_j$ occurs at least three times in its degree-eight
characteristic polynomial. Thus $[\Q(\lambda_j):\Q]\le2$.
The nonreal phases have order three, four or six, and their classes
modulo $\{1,-1\}$ have order three or two. In the product of three
nonidentity such classes, the number of order-two entries must be even.
Two order-two entries cancel but leave a nonidentity order-three entry;
zero order-two entries is the only possibility. All three phases
therefore have order three or six. The group and all character values
are defined over $\Q(\zeta_3)$, contradicting the non-Eisenstein
quadratic subfield $K$. This excludes $t=3$.

We have $t=2$. The two scalar phases multiply to $\varepsilon=\pm1$,
and their moving ranks have sum at most twelve. At least one fixes a
space of dimension at least two. Its rational real factor shows that
$[\Q(\lambda):\Q]\le4$. Because the tuple generates $J$ and
$W(E_8)\subset J$, its scalar matrices generate the same cyclotomic
field as $\lambda$. Their traces are $8\lambda^a$, so
\[
                         L=\Q(\lambda).                    \tag{E8IF.5}
\]
The orders of a nonreal root of unity of degree at most four are
$3,4,5,6,8,10,12$. Orders three, four and six give a quadratic
field and cannot satisfy $K\subsetneq L$. Orders five and ten
give $\Q(\zeta_5)$, whose only quadratic subfield is real.
The three quadratic subfields of $\Q(\zeta_8)$ have radicands
$2,-1,-2$, and those of $\Q(\zeta_{12})$ have radicands
$3,-1,-3$. This gives exactly the two field rows in (E8IF.1).
These elementary subfield calculations are also certified in D-BI-11.

In either remaining row the phase has degree four. On either real
factor its multiplicity as an eigenvalue is at most two, so both
outside moving ranks are at least six. Equation (E8IF.4) now forces
$r=12$, both ranks equal to six, and all ten other entries to be
reflections. Each real factor has characteristic polynomial
$\Phi_d^2$, where $d=8$ or $12$; all eight dimensions are exhausted
by the four Galois-conjugate eigenvalues, each with multiplicity two.
Since $J$ contains $W(E_8)$ and the primitive scalar $\lambda I$,
it is the asserted full scalar extension.

For $d=8$, adding a primitive eighth-root exponent to
$\{1,3,5,7\}/8$ yields $\{0,1,2,3\}/4$; each occurs twice.
The moving rank is six and the age is three. For $d=12$, the
primitive exponents are $\{1,5,7,11\}/12$. Adding the scalar
exponent gives one of the two displayed exponent sets; the reciprocal phase,
even after adding $1/2$ for $\varepsilon=-1$, gives the other set.
The ages are $8/3$ and $10/3$. Thus in either row the two outside
ages sum to six, while ten reflections contribute five. Subtracting
the coefficient degree eight in Chevalley--Weil gives three, and
the conjugate calculation gives three as well.

For the class names in (E8IF.1), the complete $E_8$ class list is
the published Carter Table~11, pp.~52--54, with the documented
CHEVIE class representatives already checked in D-GAUSS
\cite{WG:Carter,WG:CHEVIE}. The present exact certificate reconstructs
the two word matrices from the retained rational Cartan matrix,
checks their orders and characteristic polynomials, and verifies that
they are the only classes in that complete list with $\Phi_8^2$
or $\Phi_{12}^2$. The class sizes are respectively $3\,628\,800$
and $2\,419\,200$. No product-one existence assertion is inferred
from this local class calculation.
\end{proof}

\section{A uniform obstruction for scalar exterior-square permutation blocks}
\label{sec:exterior-square-obstruction}

\paragraph{Scope and check status.}
This is a new elementary proof, initially recorded on 21 November 2026.
It rules out a dominant sixfold source before solving the product-one
equations: even its necessary branch profiles do not exist.  The exact
partition calculation accompanying the proof checks the spectral formulas
for every partition through degree forty at scalar orders three, four,
and six (conjugate phases have equal ranks), and checks them against direct
exterior-square eigenvalues through degree twelve.  That finite calculation
is a check of the formulas, not a substitute for the uniform proof below.
The parent and the independent G34 agent have each completed a whole-proof
review, including the projection-descent and automatic-transitivity
arguments, and recorded PASS.  The parent reran the full finite certificate;
the G34 agent independently rebuilt the fractional spectra and signed pair
orbits through degree twelve without importing this certificate's code.
The only general geometric inputs are the standard finite-cover Euler
formula, Riemann--Hurwitz, and the previously proved multiplicity-variation
factorization in D-BI-01.  No classification theorem is used here.

\subsection*{D-EX-01. Three uniform spectral inequalities}

Let $n\ge4$, let $U$ be the deleted permutation representation of
$\mathfrak S_n$, and put
\[
        W=\bigwedge^2 U,\qquad m=\dim W=\binom{n-1}{2}.
                                                               \tag{EX.1}
\]
For a permutation $\pi$, write $c$ for its number of cycles, including
fixed points, and $k=n-c=\operatorname{ind}(\pi)$.  For a complex root
of unity $\lambda$, put
$D_\lambda(\pi)=\operatorname{rank}(1-\lambda W(\pi))$.
We prove
\begin{align}
 \lambda\ne1&\quad\Longrightarrow\quad
 D_\lambda(\pi)\ge \frac{n-4}{2}k+1,                  \tag{EX.2}\\
 \lambda\ne1&\quad\Longrightarrow\quad
 D_\lambda(\pi)\ge
 A_n:=m-\left\lfloor\frac{(n-1)^2}{4}\right\rfloor,  \tag{EX.3}\\
 \pi\ne1&\quad\Longrightarrow\quad
 D_1(\pi)\ge\frac{n-4}{2}k+\frac n3.                 \tag{EX.4}
\end{align}
These estimates hold for all scalar orders, not only orders three, four,
and six.

\paragraph{Proof of (EX.2).}
Put $t=\lambda^{-1}\ne1$.  Every eigenvalue of the permutation
representation has multiplicity at most $c$.  On $U$, the eigenvalue
$1$ has multiplicity $a=c-1$, and every other eigenvalue still has
multiplicity at most $c$.  Its eigenvalue-$1$ complement has dimension
$k$.  Let $b$ be the multiplicity of $t$ on that complement.
The fixed dimension in question is the multiplicity of $t$ on
$\bigwedge^2U$.

First suppose $2\le c\le n/2$.  Pair eigenvalue labels by the
involution $z\mapsto t/z$.  A two-element pair with multiplicities
$x,y\le c$ contributes $xy\le c(x+y)/2$; a fixed label contributes
$\binom{x}{2}\le cx/2$.  Treat the pair $1,t$ separately: it contributes
$ab$.  Thus the desired multiplicity $F$ satisfies
\[
 F\le ab+\frac c2(n-1-a-b)
   \le\frac c2(n-2).
\]
For the last inequality, use $a=c-1$, $b\le c$, and
\[
 \frac c2(a+b)-ab
       =\frac{c(c-1)-(c-2)b}{2}\ge\frac c2.
\]
Consequently
\[
 D_\lambda(\pi)-\left(\frac{n-4}{2}k+1\right)
       \ge k-\frac n2\ge0.
\]
If $c=1$, all eigenvalues on $U$ are distinct.  Pairing their labels
shows $F\le\lfloor(n-1)/2\rfloor\le n-2$, which is precisely enough
for (EX.2), since $k=n-1$.

It remains to treat $c>n/2$.  Here $a=c-1\ge k$.  In the complement
to the eigenvalue-$1$ space, an eigenvector with eigenvalue $t$ cannot
participate in an exterior pair of product $t$: its partner would
have eigenvalue $1$.  It follows that
\[
 F\le ab+\binom{k-b}{2}\le ak.
\]
Indeed the middle expression is a convex quadratic in $b\in[0,k]$,
and its endpoint values are $\binom k2$ and $ak$, the latter being
larger.  Since $n\ge4$ gives $c\ge3$ in this case,
\[
 D_\lambda(\pi)-\left(\frac{n-4}{2}k+1\right)
  \ge \frac{(c-k)(c-3)}2\ge0.
\]
This proves (EX.2).

\paragraph{Proof of (EX.3).}
The representation $W$ is real.  If $\lambda$ is nonreal, the
eigenvalues $\lambda^{-1}$ and $\overline\lambda^{-1}$ on $W$ have
equal multiplicity, so each has multiplicity at most $m/2$, and hence
at most $\lfloor(n-1)^2/4\rfloor$.
The only remaining nonidentity real scalar is $\lambda=-1$.
On $U$ pair eigenvalue labels by $z\mapsto-1/z$.
Each nonfixed pair of multiplicities $x,y$ contributes $xy$ to
the eigenvalue-$(-1)$ space in $\bigwedge^2U$, at most
$(x+y)^2/4$.  The two fixed labels are $i,-i$; their multiplicities
are equal, say $b$, because $U$ is real.  Taken together they
contribute $2\binom b2=b(b-1)\le(2b)^2/4$.
These blocks partition a space of total dimension $n-1$.
Summing their square bounds gives
\[
 \dim\ker(1+W(\pi))
       \le\frac14\sum_{\text{blocks}}(\dim\text{block})^2
       \le\frac{(n-1)^2}{4}.
\]
Taking the integer part proves (EX.3).

\paragraph{Proof of (EX.4).}
Let $l_1,\ldots,l_c$ be the cycle lengths, and let $e$ count the
even lengths.  On the exterior square of the full permutation space,
two different cycles contribute $\gcd(l_i,l_j)$ fixed vectors.
Inside one cycle of length $l$, the contribution is
$\lfloor(l-1)/2\rfloor$: the opposite-pair orbit for even $l$ has
negative return sign and contributes no fixed vector.  Since
$\bigwedge^2(1\oplus U)=U\oplus\bigwedge^2U$, subtraction of the
$c-1$ fixed vectors of $U$ gives the exact identity
\[
 D_1(\pi)=\frac{(n-3)k+E}{2},\qquad
 E=e+\sum_{i<j}\bigl(l_i+l_j-2\gcd(l_i,l_j)\bigr).
                                                               \tag{EX.5}
\]
All summands of $E$ are nonnegative.  We claim $k+E\ge2n/3$ for
every nonidentity permutation.
If there are no fixed points, let $o$ count the odd cycles.
They all have length at least three, so $o\le n/3$, and
$k+E\ge n-c+e=n-o\ge2n/3$.
If there are $f\ge1$ fixed points, the cross terms involving them
give $E\ge fk$.  Also $k\ge1$ and $n\le f+2k$, since each
nontrivial cycle satisfies $l\le2(l-1)$.  Hence
\[
 3(k+E)\ge3(f+1)k\ge2f+4k\ge2n;
\]
the middle inequality is $(3f-1)k\ge2f$, valid for $f,k\ge1$.
Combining this claim with (EX.5) proves (EX.4).

\subsection*{D-EX-02. No twelve-branch signature-$(3,3)$ profile}

Suppose $n\ge5$ and a product-one tuple of genuine branch elements
$g_j=\lambda_j W(\pi_j)$ has transitive permutation projection.
Assume its sixfold multiplicity satisfies the necessary genus-zero
Euler equation
\[
             \sum_jD_{\lambda_j}(\pi_j)=2m+6.            \tag{EX.6}
\]
For completeness, the general Euler formula follows by extending the
rank-$m$ local system over each puncture with its inertia invariants.
Its Euler characteristic on a genus-$h$ curve is
$m(2-2h-r)+\sum_j\dim\ker(g_j-1)$.
For a nontrivial irreducible local system its degree-zero and degree-two
cohomology vanish, so its first cohomology has dimension
$2m(h-1)+\sum_jD_j$.  Setting this dimension equal to six and $h=0$
gives (EX.6).
The transitive permutation cover of degree $n$ has nonnegative genus,
so Riemann--Hurwitz gives
\[
                    \sum_j\operatorname{ind}(\pi_j)\ge2n-2.
                                                               \tag{EX.7}
\]
We prove that the tuple has at most eleven entries.

Suppose instead that $r\ge12$, and let $s$ count the entries
with $\lambda_j\ne1$.  For $n\ge6$, (EX.3) and (EX.6) imply
$s\le6$.  Here is the full numerical check, valid for every $n$.
Writing $n-1=2v+1$, with $v\ge2$, gives $A_n=v^2$ and
\[
             7A_n-(2m+6)=3v^2-2v-6>0.
\]
Writing $n-1=2v$, with $v\ge3$, gives $A_n=v(v-1)$ and
\[
             7A_n-(2m+6)=3v^2-5v-6>0.
\]
Every scalar-$1$ entry has $\pi_j\ne1$, because branch entries are
genuine.  Thus (EX.2), (EX.4), and (EX.7) yield
\begin{align*}
 \sum_jD_{\lambda_j}(\pi_j)
 &\ge\frac{n-4}{2}\sum_j\operatorname{ind}(\pi_j)
            +\frac n3(r-s)+s\\
 &\ge(n-4)(n-1)+\frac n3(12-s)+s\\
 &\ge(n-4)(n-1)+2n+6
       =n^2-3n+10=2m+8.
\end{align*}
The last inequality uses $n\ge6$ and $s\le6$.
This contradicts (EX.6).

For $n=5$, one has $m=6$, $A_n=2$, and $s\le9$.
The quantity $D_1(\pi)-k/2$ belongs to $\frac12\mathbf Z$;
(EX.4) bounds it below by $5/3$, and hence by $2$.
Using (EX.2) for the other entries gives
\[
 \sum_jD_{\lambda_j}(\pi_j)
       \ge\frac12(2n-2)+2(r-s)+s
       \ge28-s\ge19>18=2m+6.
\]
This handles $n=5$ as well.

For $n=4$, the same conclusion holds whenever the scalar quotient
is nonreal.  In this case $m=3$ and (EX.6) is $\sum_jD_j=12$.
There must be at least two nonreal scalar entries: there is at least
one because the entries generate the nonreal scalar image, and the
product of one nonreal scalar with any collection of $\pm1$ cannot
be one.  Since $W$ is real of dimension three, each of these entries
has moving dimension at least two.  Every other genuine entry has
moving dimension at least one.  Therefore
$\sum_jD_j\ge r+2\ge14$, again a contradiction.

\subsection*{D-EX-03. Consequence for inducing blocks and the exact spectrum}

Let an exact-imaginary-quadratic-field sixfold representation be
induced from an absolutely irreducible representation $\tau$ of
dimension $m=\binom{n-1}{2}$, whose finite image in this representation
has the form
\[
      \tau(j)=\chi(j)\,\bigwedge^2\operatorname{Std}_n(\pi(j)),
      \qquad n\ge4,
                                                               \tag{EX.8}
\]
There is no restriction on the order or coefficient field of $\chi$;
the permutation image need not be all of $\mathfrak S_n$.
Then its period rank is at most eight; in particular it supplies
no dominant Weil sixfold family, of either determinant class.

The permutation image is automatically transitive.  Otherwise the
deleted permutation space contains a nonzero invariant vector.
By averaging a Hermitian form, split off its invariant line as
$U=1\oplus U'$, where $\dim U'=n-2\ge2$.  Then
$\bigwedge^2U=U'\oplus\bigwedge^2U'$ is a direct sum of two
nonzero invariant subspaces.  A scalar twist preserves these
subspaces, contrary to irreducibility of $\tau$.

Both the scalar and permutation maps descend when the kernel of
$\tau$ is quotiented.  Indeed $\bigwedge^2\operatorname{Std}_n$
is projectively faithful for $n\ge4$: if $\bigwedge^2T$ is scalar
for a semisimple $T$ on a space of dimension at least three, the
equalities between all pairwise eigenvalue products force every
eigenvalue of $T$ to be the same.  Thus $T$ itself is scalar.
The deleted permutation representation contains no nonidentity
scalar.  One way to see the latter assertion is to apply the
permutation to every difference $e_i-e_j$: scalar action preserves
every two-element support $\{i,j\}$, which forces the permutation
to be the identity for $n\ge3$.  Therefore an element of
$\ker\tau$ has $\pi=1$ and then $\chi=1$.

To justify passage to the intermediate cover, quotient the kernel
of $\tau$ in its inducing subgroup.  Frobenius reciprocity identifies
the complex six-dimensional multiplicity variation with the
$\tau$-multiplicity on the resulting finite cover of a curve $X$.
As in D-BI-01, the fixed Hermitian polarization determines the
conjugate Hodge subspace from the chosen one.  Thus the full
quadratic-field period rank factors through this intermediate-cover
moduli even when $\chi$ takes values in a larger cyclotomic field.
All finite-cover and inducing data over a fixed pointed $X$ are
discrete.

If $h=g(X)$, the Euler equation is
\[
                        6=2m(h-1)+\sum_jD_j.
                                                               \tag{EX.9}
\]
For $h=1$ there are at most six genuine marked points, and the
marked-base moduli dimension is at most six.  For $h\ge2$, since
$m\ge3$, the only possible case is $h=2,m=3$ with no ramification;
its base moduli has dimension three.  Rank nine would therefore
require $h=0$ and at least twelve branch points.  This is excluded
by D-EX-02.  In the $n=4$ case the scalar image is necessarily
nonreal: if it were contained in $\{1,-1\}$, $\tau$ would have
real character and so would its induced representation, contrary
to its exact imaginary quadratic character field.  This proves
the period-rank assertion.

For clarity, the complete finite partition spectrum used for checking
these estimates is also explicit.  For a partition
$l_1+\cdots+l_c=n$, the standard-representation eigenvalues are the
fractions
\[
 \coprod_i\{0,1/l_i,\ldots,(l_i-1)/l_i\}
                  \quad\text{with one zero removed}.
\]
If $\lambda=\exp(2\pi i a/d)$, add $a/d$ to every unordered sum
of two distinct members of this multiset, and reduce modulo one.
The number of nonzero resulting fractions is $D_\lambda$; their
sum is the age.  Signature $(3,3)$ additionally requires total
age $m+3$, as well as (EX.6).  Since (EX.6)--(EX.7) and $r\ge12$
already contradict the inequalities, imposing these extra
Chevalley--Weil, scalar-product, and tuple-generation constraints
cannot restore a profile.

An equivalent integer formula for the fixed dimension is useful
for the certificate.  Let $e$ be the order of $\lambda$.  On
$\bigwedge^2$ of the full permutation space, a pair of different
cycles contributes
\[
       \gcd(l_i,l_j)\,\mathbf1_{e\mid\operatorname{lcm}(l_i,l_j)}.
\]
A single cycle of length $l$ contributes
\[
  \left\lfloor\frac{l-1}{2}\right\rfloor\mathbf1_{e\mid l}
       +\mathbf1_{\,2\mid l\,,\ \lambda^{l/2}=-1}.
\]
Subtract $\sum_i\mathbf1_{e\mid l_i}$, and add $1$ if
$\lambda=1$, to pass to $\bigwedge^2U$.
This formula counts signed pair orbits, while the direct fraction
calculation counts eigenvectors; the executable check compares the
two without floating-point arithmetic.

\section{The scalar unordered-pair representation is also excluded}
\label{sec:unordered-pair-obstruction}

\paragraph{Scope and check status.}
This is a new elementary proof, recorded on 21 November 2026.
It concerns the representation usually denoted $S^{(n-2,2)}$,
of dimension $n(n-3)/2$, and scalar twists of its restrictions.
No hypothesis on the scalar character field is imposed.
The finite check for $5\le n\le10$ has run exactly: all 127
partitions and all 751 matching scalar phases pass, with a separate
check for every off-spectrum scalar.  Two different exact spectral
constructions agree.  The parent completed an independent whole-proof
review of D-UP-01--03, including the all-degree inequalities, arbitrary-phase
exhaustiveness, projection descent, and the degree-four exception, and
recorded PASS.  A second independent whole-proof review by the G34
agent also passed; its separate permutation-cycle implementation
reconstructed all 127 partitions and 751 matching phases exactly.

\subsection*{D-UP-01. Representation and spectral estimates}

Let $P_n$ be the permutation space on $n$ letters, let
$U_n=\operatorname{Std}_n$ be its deleted summand, and let $P_{2,n}$
be the permutation space on unordered pairs of distinct letters.
The incidence map
\[
        e_i\longmapsto\sum_{j\ne i}e_{\{i,j\}}
\]
embeds $P_n$ into $P_{2,n}$ for $n\ge3$: a relation between its
columns would have coefficients satisfying $a_i+a_j=0$ for all
distinct $i,j$, and hence all $a_i=0$.
Let $W_{2,n}$ be its orthogonal complement.  Thus
\[
  P_{2,n}=P_n\oplus W_{2,n},\qquad
  m=\dim W_{2,n}=\frac{n(n-3)}2.                       \tag{UP.1}
\]
Everything is defined over $\mathbf Q$ and is real.
The identity
$\operatorname{Sym}^2P_n=P_n\oplus P_{2,n}$, together with
$P_n=1\oplus U_n$, also gives
\[
          \operatorname{Sym}^2U_n\simeq P_{2,n}.
                                                               \tag{UP.2}
\]
These identities follow by decomposing diagonal and off-diagonal
monomials and cancelling the common summand in characteristic zero.

Fix $n\ge5$.  For a permutation $\pi$, write $c$ for its number
of cycles and $k=n-c$.  For a root of unity $\lambda$, set
$D_\lambda=\operatorname{rank}(1-\lambda W_{2,n}(\pi))$.
We prove
\begin{align}
 \pi\ne1&\quad\Longrightarrow\quad
 D_1\ge\frac{n-5}{2}k+\frac n4,                       \tag{UP.3}\\
 \lambda\ne1&\quad\Longrightarrow\quad
 D_\lambda\ge\frac{n-5}{2}k+\frac{n-4}{2}.             \tag{UP.4}
\end{align}

\paragraph{Proof of (UP.3).}
Let $l_1,\ldots,l_c$ be the cycle lengths, let $e$ count the even
ones, and put
\[
    Q=\sum_{i<j}\bigl(l_i+l_j-2\gcd(l_i,l_j)\bigr)\ge0.
\]
Two different cycles contribute $\gcd(l_i,l_j)$ fixed unordered-pair
orbits.  One cycle contributes $\lfloor l_i/2\rfloor$, including
its opposite-pair orbit when $l_i$ is even.  Subtracting the $c$
fixed vectors of $P_n$ in (UP.1) gives
\[
                     D_1=\frac{(n-3)k-e+Q}{2}.         \tag{UP.5}
\]
Equivalently it is the exterior-square moving rank of (EX.5) minus
the number of even cycles.
We claim $2k-e+Q\ge n/2$ whenever $\pi\ne1$.
If there are no fixed points, then $k\ge n/2$ and $e\le k$,
which proves the claim.  If there are $f\ge1$ fixed points, then
$Q\ge fk$, $e\le k$, $k\ge1$, and $n\le f+2k$.
Consequently
\[
       2k-e+Q\ge(f+1)k\ge\frac{f+2k}{2}\ge\frac n2.
\]
The middle inequality is $f(2k-1)\ge0$.
Subtracting $(n-5)k/2$ from (UP.5) now proves (UP.3).

\paragraph{Proof of (UP.4).}
Put $t=\lambda^{-1}\ne1$.  On $U_n$, the eigenvalue $1$ has
multiplicity $a=c-1$, the complementary space has dimension $k$,
and each eigenvalue has multiplicity at most $c$.
Let $b$ be the multiplicity of $t$ on that complement.
In $\operatorname{Sym}^2U_n$, the pair of eigenvalue spaces $1,t$
contributes $ab$.  An eigenvector with eigenvalue $t$ cannot
contribute to any other symmetric pair of product $t$.
The remaining relevant space therefore has dimension $k-b$.
On it, two distinct eigenvalue spaces of dimensions $x,y\le c$
contribute at most $c(x+y)/2$, and a self-paired eigenspace of
dimension $x$ contributes $\binom{x+1}{2}\le(c+1)x/2$.
By (UP.1)--(UP.2) one must also subtract the $b$ eigenvectors of
$P_n$ with eigenvalue $t$.  Hence the fixed dimension $F$ of
$\lambda W_{2,n}(\pi)$ is bounded by
\[
 F\le ab-b+\frac{c+1}{2}(k-b)
       =\frac{(c+1)k+(c-5)b}{2}.                       \tag{UP.6}
\]
If $c\ge5$, use $b\le c$.  Substituting in $m-F$ gives
\[
            D_\lambda-\frac{n-5}{2}k\ge\frac{n+c}{2}.
\]
If $1\le c\le4$, use $b\ge0$ and $c-5<0$ instead.  It gives
\[
 D_\lambda-\frac{n-5}{2}k
       \ge\frac{n+c(c-4)}2\ge\frac{n-4}{2},
\]
because $c(c-4)\ge-4$ for these four integers.  This proves (UP.4).

Both right-hand sides are positive in their stated nonidentity cases.
It follows in particular that $W_{2,n}(\mathfrak S_n)$ intersects
the scalar matrices only in the identity for $n\ge5$.
Thus, for a representation $\chi\otimes(W_{2,n}\circ\pi)$,
both $\pi$ and $\chi$ descend to its faithful image.
This conclusion is not asserted at $n=4$.

\subsection*{D-UP-02. Uniform obstruction and the finite degrees}

Consider an irreducible scalar twist
\[
           \tau(j)=\chi(j)W_{2,n}(\pi(j)),\qquad n\ge5,
                                                               \tag{UP.7}
\]
whose character is nonreal.  Suppose a genus-zero branch tuple in
its faithful image has product one, generates that image, and has
six-dimensional cohomological multiplicity.  Then
\[
                    \sum_jD_j=2m+6.                    \tag{UP.8}
\]
The permutation image is transitive; the verification from
irreducibility is given in D-UP-03 below.  Because the projections
descend by D-UP-01, the product-one tuple gives an actual transitive
degree-$n$ permutation cover of the sphere.  Therefore
\[
                    \sum_j\operatorname{ind}(\pi_j)\ge2n-2.
                                                               \tag{UP.9}
\]
Its scalar quotient is nonreal.  There are at least two nonreal
scalar entries, since one nonreal scalar cannot have product one
with any number of real roots of unity.  In particular the number
$s$ of entries with $\lambda_j\ne1$ satisfies $s\ge2$.

Suppose $n\ge11$ and $r\ge12$.  Applying (UP.3)--(UP.4) gives
\begin{align*}
 \sum_jD_j
 &\ge \frac{n-5}{2}\sum_j\operatorname{ind}(\pi_j)
          +\frac n4(r-s)+\frac{n-4}{2}s\\
 &\ge (n-5)(n-1)+3n+\frac{n-8}{4}s\\
 &\ge n^2-3n+5+\frac{n-8}{2}\\
 &=2m+6+\frac{n-10}{2}>2m+6.
\end{align*}
This contradicts (UP.8).  The strict inequality already holds for
$n=11$; no limiting or asymptotic argument is used.

For $5\le n\le10$, the exact minimum of $D_\lambda(\pi)$ over
all genuine scalar-permutation matrices, allowing arbitrary scalar
orders, is $n-3$.  The complete finite certificate is as follows.
For a partition $l_1+\cdots+l_c=n$, form a multiset of eigenvalue
fractions for $P_{2,n}$:
\begin{itemize}
 \item two different cycles of lengths $a,b$ contribute
 $\gcd(a,b)$ copies of every fraction $j/\operatorname{lcm}(a,b)$;
 \item one cycle of length $l$ contributes
 $\lfloor(l-1)/2\rfloor$ copies of every fraction $j/l$;
 \item if $l$ is even, add one copy of every fraction $j/(l/2)$
 for its opposite-pair orbit.
\end{itemize}
Subtract one copy of every fraction $j/l_i$ for each cycle to
remove $P_n$.  All resulting multiplicities are nonnegative and
sum to $m$, by (UP.1).  If $q$ occurs with multiplicity $b_q$,
the scalar $\exp(-2\pi iq)$ gives moving dimension $m-b_q$.
Every scalar not of this form gives moving dimension $m$.
Thus this finite list tests every possible scalar order.
Omit only $\pi=1$, $q=0$, which is the identity branch matrix.

The accompanying program enumerates the partitions recursively by
their nondecreasing parts, uses exact fractions, and compares this
orbit multiset to the independently constructed multiset of all
unordered sums with repetition of standard-representation
eigenvalue fractions, followed by subtraction of $P_n$.
The latter construction is (UP.2).
The complete results are
\[
\begin{array}{c|rrrrrr}
 n&5&6&7&8&9&10\\ \hline
 m&5&9&14&20&27&35\\
 \#\text{partitions}&7&11&15&22&30&42\\
 \min D_\lambda&2&3&4&5&6&7\\
 12\min D_\lambda&24&36&48&60&72&84\\
 2m+6&16&24&34&46&60&76
\end{array}                                                     \tag{UP.10}
\]
Every minimum is attained by a scalar-$1$ transposition.
The strict inequalities in the last two rows exclude $r\ge12$.
The certificate executed 127 partitions and 751 matching scalar
checks; off-spectrum scalars were treated as just explained.
This is an exhaustive finite proof of the small-degree inequalities,
not a numerical sampling of partitions or phases.

\subsection*{D-UP-03. Inducing representations and the exceptional degree four}

Let an exact-imaginary-quadratic-field Weil sixfold representation
be induced from an irreducible representation of the form (UP.7),
now allowing $n\ge4$.  Then its period rank is at most eight.
There is no assumption that the inducing character field or scalar
field equals the original quadratic field.  In particular none of
these blocks supplies a dominant family, split or nonsplit.

First let $n\ge5$.  We verify the transitivity used above.
If the permutation image preserves a subset of size $a$ with
$2\le a\le n-2$, it lies in
$\mathfrak S_a\times\mathfrak S_{n-a}$.
This full subgroup has three orbits on unordered pairs and two
orbits on letters.  Hence (UP.1) gives a one-dimensional invariant
subspace of $W_{2,n}$, a proper subspace since $m\ge5$.
It stays invariant under the smaller permutation image and under
the scalar twist, contrary to irreducibility.
If no such subset exists but the image is intransitive, it fixes
one letter and is transitive on the remaining $n-1$ letters.
Restricting (UP.1) to its containing $\mathfrak S_{n-1}$ gives
\[
 W_{2,n}|_{\mathfrak S_{n-1}}
       \simeq U_{n-1}\oplus W_{2,n-1}.                  \tag{UP.11}
\]
Indeed $P_{2,n}$ restricts as $P_{2,n-1}\oplus P_{n-1}$,
whereas $P_n$ restricts as $1\oplus P_{n-1}$.
Both summands in (UP.11) have positive dimension for $n\ge5$.
They remain invariant after restriction and scalar twist, again
contradicting irreducibility.  Thus the permutation image is
transitive.

Pass to the intermediate connected cover for the faithful image
of the inducing representation.  Its complex multiplicity
variation is the original one by Frobenius reciprocity and
quotient pullback, and the fixed Hermitian polarization determines
the conjugate filtration.  The full quadratic-field period rank
therefore factors through this intermediate-cover moduli, as in
D-BI-01 and D-EX-03.  The scalar and permutation maps descend by
(UP.3)--(UP.4), so no projected branch data have been lost.
If $h$ is the genus of the intermediate base, its Euler equation is
\[
                       6=2m(h-1)+\sum_jD_j.
\]
For $h=1$ the marked-base moduli has dimension at most six.
For $h\ge2$ the right-hand contribution $2m(h-1)$ already
exceeds six because $m\ge5$.
Thus rank nine would require $h=0$ and $r\ge12$.
The inducing character is nonreal: otherwise its induced character
would be real, contrary to the original exact imaginary quadratic
field.  Because $W_{2,n}$ is rational, a scalar image contained in
$\{1,-1\}$ would give a real inducing character.  Its scalar image
is therefore nonreal, and D-UP-02 supplies the contradiction.

When $n=4$, the representation $W_{2,4}$ has dimension two and
need not retain a faithful degree-four permutation projection.
We use no such projection in this case.
For any irreducible inducing representation of dimension two with
cohomological multiplicity six, the Euler equation gives
$D=10-4h$.  Since each genuine branch contributes at least one,
the dimensions of its marked-base moduli are at most seven for
$h=0$, six for $h=1$, and five for $h=2$; larger $h$ is impossible.
Thus the period rank is at most seven.  This treats degree four
without an incorrect projective-faithfulness assertion.

These obstructions use only necessary local spectra, the Euler
equation, and Riemann--Hurwitz.  No feasible aggregate profile is
being promoted to a product-one generating tuple, and no source
dimension inequality is being promoted to dominance.

\section{The exact remaining branch datum for the Gaussian scalar $E_8$ source}
\label{sec:e8-source-gate}

\subsection*{HD-E8-01. Local determinants in $\mu_4W(E_8)$}
Let $W=W(E_8)$ in its rational reflection representation, endowed
with the positive definite root form $H$. Use the simple-root basis
whose Cartan matrix has diagonal entries two and off-diagonal entries
minus one on the edges
\[
 (1,3),\ (3,4),\ (2,4),\ (4,5),\ (5,6),\ (6,7),\ (7,8).
\]
The determinant of $H$ is one. Let $s_j=I-e_jH_{j,*}$ and put
$G=\langle W,iI\rangle\subset U(H,\mathbf Q(i))$.
Since $-I\in W$, this group has two cosets $W$ and $iW$,
and order $2\cdot696729600$.

For $a\in G$ let $U_a=\operatorname{im}(1-a)$ and define, as in
D-G31-01,
\[
 k(a)=\dim U_a,\qquad
 q_+(a)=\frac{\det(H|_{U_a})}{|\det(1-a|_{U_a})|}
       \quad\bmod\operatorname{Nm}_{\mathbf Q(i)/\mathbf Q}\mathbf Q(i)^*.
\]
Every class in $W$ has norm $q_+$. Every class in $iW$ has moving
rank at least four. Its unique rank-four class is $i$ times Carter's
$2D_4(a_1)$, has age two, and has norm $q_+$. The complete list of
nonnorm classes in $iW$ is the following. The first column uses
the CHEVIE author's numbering, reproduced by the accompanying script.
All six classes have moving rank eight and local norm class three.
\[
\begin{array}{r|l|r|r|r}
\text{index}&\text{class of }w\text{ in }a=iw&
 \det(1-a)&q_+(a)&\operatorname{age}(a)-4\\\hline
55&E_6+A_2&-3&3&1\\
56&E_8(a_7)&-3&3&-1\\
57&E_6&6&6&0\\
58&E_7(a_2)+A_1&6&6&0\\
101&E_6+A_1&6i&6&1/2\\
102&E_7(a_2)&-6i&6&-1/2
\end{array}
\]
Here $6$ and $3$ have the same Gaussian norm class since $2$ is a norm.

\paragraph{Exact computation and its external input.}
The complete class list is the published Weyl-group classification
of Carter, \emph{Conjugacy classes in the Weyl group}, Compositio
Math.~25 (1972), 1--59, Proposition~21/Table~3, pp.~22--23, and
Table~11, pp.~52--54. We use the CHEVIE author's explicit complete
class representatives, documented in the GAP3 manual, section~87.1,
\emph{ChevieClassInfo}, and retained as mathematical words in
\texttt{research/e8\_class\_exploration.json}. Their retrieval URL
and SHA-256 occur in that file. Completeness of this known class
classification is an input, not a new exhaustive enumeration of
the $696729600$ group elements.

\texttt{scripts/examine\_e8\_classes.py} reconstructs each representative
as a product of the displayed integral simple reflections. It checks
all 112 matrix orders, orthogonality, and characteristic polynomials;
the supplied class sizes sum to the group order.
\texttt{scripts/e8\_gaussian\_local.py} computes a basis of every
moving space by exact Gaussian elimination, restricts both the
Hermitian form and the operator, and takes their determinants.
Its norm test factors rational numbers and tests the parity of
valuations at primes congruent to three modulo four. Ages are
computed from exact cyclotomic factors of the real characteristic
polynomial, with every exponent shifted by $1/4$ modulo one.
\texttt{scripts/e8\_real\_local.py} separately checks all real classes.
The complete rational outputs, including spectra, are retained.
This proves the displayed finite arithmetic assertions from the
specified complete representative list.

\subsection*{HD-E8-02. A necessary twelve-branch pattern}
Suppose the natural representation of $G$ supplies a dominant
nonsplit Weil-sixfold period map. Then the base is genus zero,
its signed Hermitian discriminant has Gaussian norm class three,
and its inertia tuple consists of exactly
\[
 \boxed{\text{ten real reflections, one rank-four element of }iW,
        \text{ and one rank-eight element }iw\text{ of class57 or58}.}
\]
The rank-four element is an involution, and the rank-eight element
has order twelve. In particular its branch orders are $(2^{11},12)$.
This is a necessary pattern. It does not assert product-one
generation, actual rank nine, or existence of detecting cycles.

\begin{proof}
For a base of genus $\gamma$, the dimension equation is
$6=(2\gamma-2)8+D$, where $D$ is the sum of local moving ranks.
For $\gamma\ge2$ this is impossible. At genus one $D=6$ and
the Hurwitz dimension is at most six. Dominance therefore requires
genus zero, $D=22$, and at least twelve branch points.

The signed determinant formula D-DET-01 applies. Since the space
is balanced, its determinant representative is positive; as in
D-G31-02 its norm class is the product of the positive local factors
$q_+$. Thus nonsplitness requires one of the six exceptional entries
in HD-E8-01, of moving rank eight. A product-one tuple has an even
number of entries in the nonreal coset $iW$, and a generating tuple
has at least one. Thus another nonreal entry has moving rank at
least four. If there are $r$ entries, their total moving rank is at
least $8+4+(r-2)=r+10$. With $D=22$ and $r\ge12$, equality
is forced: $r=12$, the second outside entry has rank four, and
every other entry is a reflection. Such reflections belong to $W$.

Each real reflection has age $1/2$, and the rank-four outside
involution has age two. Balance requires total age $11$, hence the
rank-eight element must have age four. The table leaves precisely
classes57 and58. Its local norm class is three, all other entries
have norm local factors, and the stated discriminant follows.
The spectra in the certificate give order twelve in both cases.
\end{proof}

\paragraph{Source URLs.}
\url{https://www.numdam.org/item/CM_1972__25_1_1_0.pdf};
\url{https://webusers.imj-prg.fr/~jean.michel/gap3/htm/chap087.htm};
\url{https://raw.githubusercontent.com/jmichel7/gap3-jm/master/pkg/chevie/tbl/weyle8.g}.
No entire third-party library file is redistributed.

\section{An exact nonsplit Gaussian $E_8$ twelve-branch source}
\label{sec:im-e8-tuple}

\paragraph{Status and scope.}
The assertions in D-E8T-01--03 below are established by the displayed
arguments and the exact certificate
\texttt{agents/imprimitive/certify\_e8\_tuple.py}, with full output in
\texttt{agents/imprimitive/e8\_search/tuple\_certificate.json}.
They establish product-one generation, the actual Weil signature and
polarization class, and a boundary with toric rank two. They do not
assert rank nine of the period map or algebraicity of its Weil classes.
The tuple was found independently by an exhaustive conjugacy-orbit
search; no tuple, cup matrix, or rank theorem from the unrefereed
Gaussian $E_8$ manuscript is a premise.

\subsection*{D-E8T-01. An explicit generating tuple}
Use the simple-root basis and Cartan matrix $H$ of HD-E8-01. Thus
$s_a=I-e_aH_{a,*}$, and the edges of the Dynkin diagram are
$(1,3),(3,4),(2,4),(4,5),(5,6),(6,7),(7,8)$.
Let $j_0$ be the following integral matrix:
\[
j_0=\begin{pmatrix}
0&1&1&0&-1&0&-1&0\\
1&2&0&0&-1&0&-1&-1\\
0&3&1&0&-2&0&-1&-1\\
1&4&1&0&-2&-1&-1&-2\\
1&3&1&0&-2&0&-1&-2\\
1&2&0&1&-2&0&-1&-1\\
1&2&0&0&-1&0&-1&0\\
0&1&0&0&0&0&-1&0
\end{pmatrix}.
\]
This is the representative of $2D_4(a_1)$ reconstructed as an explicit
simple-reflection word in HD-E8-01; its word, matrix, and provenance
are retained in \texttt{research/e8\_class\_exploration.json}.
Put
\[
u=s_6s_7s_8s_1,\qquad j=uj_0u^{-1},\qquad
c=s_1s_2s_3s_4s_5s_6,\qquad w=jc^{-1},\qquad h=ij,\quad g=iw.
\]
The twelve-entry tuple is
\[
\boldsymbol t=(s_1,s_1,s_2,s_2,s_1,s_2,s_3,s_4,s_5,s_6,h,g).
\tag{E8T}
\]
It has product one, generates $G=\mu_4W(E_8)$, and has branch orders
$(2^{11},12)$.

\begin{proof}
The matrices satisfy $j^2=-I$ and $j^tHj=H$. Consequently
$h^2=I$ and
\[
hg=-jw=-j^2c^{-1}=c^{-1}.
\]
The two initial pairs cancel, so the product of (E8T) is $chg=I$.
The exact characteristic polynomial of $w$ is
\[
(x-1)^2\Phi_3(x)\Phi_{12}(x)
=x^8-x^7-x^6+2x^4-x^2-x+1.
\]
The eight eigenangle fractions of $g=iw$, in increasing order, are
\[
\frac16,\ \frac14,\ \frac14,\ \frac13,\ \frac7{12},\
\frac23,\ \frac56,\ \frac{11}{12}.
\tag{E8T-spectrum}
\]
They are all nonzero, sum to four, and give order twelve.
The certificate independently checks $g^k\ne I$ for $1\le k<12$
and $g^{12}=I$.

Here is a certificate of generation which does not enumerate the
$1393459200$ elements of $G$. Starting with the six roots
$e_1,\ldots,e_6$, close under the seven transformations
$s_1,\ldots,s_6,j$. The resulting set has 240 distinct vectors, each
of $H$-norm two, and equals the set obtained by closing all eight
simple roots under all eight simple reflections. For every new root
the output records its parent root and the transformation applied;
these are finite exact integer equalities. The latter set is the full
$E_8$ root system. In the group generated by the tail
$(s_1,\ldots,s_6,h,g)$, conjugation by $h=ij$ carries the reflection
in a root $v$ to the reflection in $jv$: the scalar $i$ cancels in
conjugation. Hence the root-closure certificate implies that this
group contains every $E_8$ root reflection and therefore contains
$W(E_8)$. It contains $j$, and hence $hj^{-1}=iI$ as well. Thus even
the eight-entry tail generates all of $G$.

For reproducibility, the search first closed $j_0$ under conjugation
by all eight simple reflections and found precisely 15120 distinct
matrices, verifying $j^2=-I$ for each. For $c=s_1\cdots s_6$, exactly
1080 of these gave the characteristic polynomial of class 57 or 58
for $jc^{-1}$. The selected matrix is orbit index 179 and is obtained
by successive conjugations $[1,8,7,6]$. It is the first such match
examined for which the stated root closure gives 240 roots.
\end{proof}

\subsection*{D-E8T-02. The actual sixfold and its pairing}
A connected $G$-cover of $\mathbf P^1$ with tuple (E8T) has genus
$3077222401$. Its natural-representation primitive factor is an
abelian sixfold of Weil type for $K=\mathbf Q(i)$, of signature
$(3,3)$, with raw Hermitian determinant class $[-3]$. Its Hurwitz
component has dimension nine.

\begin{proof}
The $E_8$ reflection representation is absolutely irreducible: the
simple roots span, and their connected nonorthogonality graph forces
any orthogonal decomposition into invariant complex subspaces to
place all roots in one summand. The matrices are defined over $K$,
and the scalar $iI$ has trace $8i$, so the character field is exactly
$K$. The associated rational simple block is $\operatorname{Mat}_8(K)$.
A rational Rosati-fixed primitive projector in this block therefore
cuts out one copy of the multiplicity factor as an actual polarized
abelian subvariety of the covering Jacobian. Equivalently, project
orthogonally to a root line inside the central natural block. The
central scalar $iI$ preserves its integral homology lattice.

There are ten rank-one reflection entries of age $1/2$, one involution
$h$ with four eigenvalues $-1$ and four eigenvalues $1$, hence age two,
and the full-moving-rank entry $g$ with age four by
(E8T-spectrum). Their total moving rank is $10+4+8=22$ and their
total age is eleven. Chevalley--Weil gives
\[
\dim_K H^1_{\mathrm{par}}=-16+22=6,\qquad
h^{1,0}_{\sigma}=h^{1,0}_{\bar\sigma}=11-8=3.
\]
Thus the primitive factor has dimension six and the asserted
signature. Riemann existence and the usual ordered Hurwitz space
supply a finite etale covering of an open part of $M_{0,12}$, so its
component has dimension nine. Riemann--Hurwitz gives
\[
g(C)=1+\frac{1393459200}{2}
\left(-2+11\frac12+\frac{11}{12}\right)=3077222401.
\]

For the actual polarization, use the parabolic cup construction of
D-DET-01. To make the verification independently reproducible, for
each inertia matrix $t_a$ take the first independent columns $U_a$
of $I-t_a$, let $T_a$ be its restriction to that moving space, and put
$Q_a=t_1\cdots t_{a-1}$. The relation matrix has blocks $Q_aU_a$.
The coboundary matrix has blocks consisting of the coordinates of
$(I-t_a)x$ in $U_a$. The upper triangular Seifert matrix has blocks
\[
S_{ab}=(Q_aU_a)^*H(Q_bU_b)\quad(a<b),\qquad
S_{aa}=(I-T_a)^{-*}U_a^*HU_a.
\]
It satisfies, exactly,
\[
S+S^*=R^*HR,\qquad B^*S=HR,\qquad RB=0,
\qquad \operatorname{rank}R=\operatorname{rank}B=8.
\]
Thus its restriction to $\ker R$ descends to the anti-Hermitian cup
form on the quotient, up to the common cup-convention sign, which
does not affect its determinant in even dimension. The output saves
$R,B,S$ and every column of the quotient basis. The resulting matrix is
\[
S_{\mathrm{par}}=\begin{pmatrix}
0&0&-4&0&2-4i&2\\
0&0&0&-4&2-2i&0\\
4&0&0&0&-2+4i&-2\\
0&4&0&0&-2+2i&0\\
-2-4i&-2-2i&2+4i&2+2i&-6i&3-3i\\
-2&0&2&0&-3-3i&-2i
\end{pmatrix},
\qquad \det S_{\mathrm{par}}=1536.
\]
Converting to a Hermitian form by multiplying by $i$ changes the
six-dimensional determinant to $-1536$. Since $1536=3\cdot2^9$
and $2$ is a Gaussian norm, its raw determinant class is $[-3]$.
This is not the split class $[-1]$, since the valuation at the
inert prime three is odd. This cup form is the multiplicity
polarization; changing the chosen primitive root-line factor scales
it by a positive rational number, whose sixth power is a norm.
\end{proof}

\subsection*{D-E8T-03. A rank-two toric boundary and its anisotropic residual factor}
Degenerate the two initial equal-reflection pairs in (E8T) onto
separate rational bubbles. The natural multiplicity space of the
dual graph has $K$-dimension two. The remaining polarized
weight-one factor has signature $(1,1)$, with anti-Hermitian matrix
\[
S_{\mathrm{res}}=
\begin{pmatrix}-6i&3-3i\\-3-3i&-2i\end{pmatrix},
\qquad \det S_{\mathrm{res}}=6.
\]
It is anisotropic over $K$. The positive sum of the two geometric
node logarithms has rank two on the natural multiplicity space.
In particular the actual six-dimensional $K$-Hermitian space has
Witt index exactly two.

\begin{proof}
Each bubble has inertia tuple $(s_a,s_a,1)$; its connected covering
components are rational double covers of $\mathbf P^1$. The central
component is connected, because its tail tuple
$(s_1,\ldots,s_6,h,g)$ generates $G$ by D-E8T-01. The attachment
inertia is trivial. Hence the two sets of edges have permutation
modules $K[G]$, the bubble vertices have modules
$K[G/\langle s_1\rangle]$ and $K[G/\langle s_2\rangle]$, and the
central vertex contributes the trivial module. Exactness of the
connected graph's incidence sequence gives, in its representation
ring,
\[
H_1(\Gamma,K)=2K[G]-K[G/\langle s_1\rangle]
                         -K[G/\langle s_2\rangle].
\]
The natural representation has multiplicity eight in $K[G]$ and
multiplicity seven in each induced trivial module. Its graph
multiplicity is therefore $16-7-7=2$.

The central curve has the eight-entry tail tuple. Its total moving
rank is $6+4+8=18$, so its parabolic multiplicity dimension is
$18-16=2$; its two holomorphic multiplicities are $3+2+4-8=1$.
Its genus is
\[
1+\frac{1393459200}{2}
 \left(-2+7\frac12+\frac{11}{12}\right)=1683763201.
\]
The same exact cup computation as above gives the displayed residual
matrix. The Hermitian form $iS_{\mathrm{res}}$ has first diagonal
entry six and determinant $-6$; completing the square diagonalizes
it over $K$ to $\langle6,-1\rangle$. An isotropic vector would make
six a norm from $K$, impossible because its three-adic valuation is
odd. Thus the residual plane is anisotropic.

For an equivariant semistable smoothing, the vanishing cycles of
the nodes have precisely the vertex-incidence relations; their
natural multiplicity has dimension two by the graph calculation.
A positive smoothing logarithm is a sum of positive multiples of
the Picard--Lefschetz rank-one tensors of these cycles. The positive
coefficients ensure that its rank is the dimension of their span,
namely two over $K$. Its image is isotropic for the actual
polarization. The weight-one graded quotient is the central
curve's residual factor just computed. It follows that the Witt
decomposition has two hyperbolic planes and this anisotropic plane.

The two node parameters alone do not establish rank nine. In the
unitary domain, a boundary with isotropic dimension two has four
toric-block tangent directions, four mixed extension directions,
and one residual-period direction. A proof of dominance must
control the actual image in all nine directions, or give a separate
full-rank codifferential calculation at a smooth cover. This remains
separate from the exact tuple and pairing results proved here.
\end{proof}

\section{A published-GRR proof of rank nine for the nonsplit Gaussian $E_8$ datum}
\label{sec:e8-grr}

\paragraph{Scope and provenance.}
This section proves that a Hurwitz component for the Gaussian scalar
$E_8$ branch datum has a dominant Weil-sixfold period map.  The proof
uses an exact positive Hodge intersection, not a comparison of source
and target dimensions.  It is independent of the separate analytic
Riemann--Hilbert certificate.  The two exact numerical answers and the
proposal to use a rooted tree recurrence were leads in the unrefereed
rStation manuscript \cite{E8GRR:lead66}.  No theorem, character packet,
or executable from that manuscript is an input here.  The character
data are reconstructed from the CHEVIE author's table, the recurrence
is proved below, and the implementation is new.

An important source distinction is necessary.  Jarvis--Kimura's
\emph{A representation-valued relative Riemann--Hurwitz theorem and
the Hurwitz--Hodge bundle}, arXiv:0810.2488v3, supplies a useful
comparison formula, but no journal publication of that manuscript
has been verified.  It is not the published GRR premise of this proof.
The premise is Tseng's published formula \cite[\S7.3, especially
(7.3.6.1) and (7.3.7.2), pp.~60--64]{E8GRR:Tseng}.  Its normalization
is explained explicitly below.  Positivity is the published theorem
\cite[Lemma~6.17, with its proof and Lemma~6.15]{E8GRR:BBT}.
The references to these lemma numbers use the 48-page author version,
pp.~42--44.  No Severn, Markman, or 032 theorem is assumed.

\subsection*{D-E8-GRR-01. The fixed local system on every counted component}

Let $K=\mathbb Q(i)$.  In the simple-root basis let $H$ be the $E_8$
Cartan matrix: its diagonal entries are two and its entries on the
edges
\[
 (1,3),(3,4),(2,4),(4,5),(5,6),(6,7),(7,8)
\]
are minus one; other off-diagonal entries are zero.  Put
$s_j=1-e_jH_{j,*}$, $W=\langle s_1,\ldots,s_8\rangle$, and
$G=\langle W,iI\rangle\subset U(H,K)$.  Thus
\[
 |W|=696729600,\qquad |G|=1393459200,\qquad G=W\sqcup iW.
\]
Here $-I\in W$.  We use the natural representation $\rho$ on $K^8$.
It is absolutely irreducible and its rational simple group-algebra
factor is $M_8(K)$.

Write $R$ for the reflection class, $C_2$ for $i$ times the Weyl class
$2D_4(a_1)$, and $C_1$ for $i$ times the Weyl class $E_6$.  In the
CHEVIE class order these are respectively real class66, imaginary
class6, and imaginary class57.  The branch classes, with all markings
ordered, are
\[
                 (C_1,C_2,R,R,R,R,R,R,R,R,R,R).                 \tag{E8.1}
\]
The complete class words and sizes used here are retained in
\texttt{e8\_grr/character\_input.json}.  In particular, the word for
the real class57 is $s_1s_4s_2s_3s_6s_5$.  The orders and spectral
data required for (E8.1) are
\[
\begin{array}{c|r|r|l|r}
\text{class}&\text{size}&\text{order}&\text{eigenvalue exponents in }[0,1)
 &\text{age}\\\hline
R&120&2&0^7,(1/2)&1/2\\
C_2&15120&2&0^4,(1/2)^4&2\\
C_1&4838400&12&1/6,1/4,1/4,1/3,7/12,2/3,5/6,11/12&4
\end{array}                                                     \tag{E8.2}
\]
The exact matrix reconstruction described in D-E8-GRR-04 checks these
data.  Class57 is carried to class58 by multiplication by $-I$.

Let $\mathcal H$ be the stack of smooth principal $G$-covers of the
punctured genus-zero base, completed at the twelve marked points
with local classes (E8.1).  Disconnected total curves are allowed.
Equivalently its fiber over an ordered configuration is the groupoid
of product-one tuples in the specified classes, modulo simultaneous
conjugation by $G$.  Its mass is therefore the number of such tuples
divided by $|G|$.  Let $\overline{\mathcal H}$ be its proper
admissible-cover, or balanced twisted-cover, compactification.  We
use the normalization with this ordinary groupoid mass.  Its coarse
moduli map to $\overline M_{0,12}$ is finite, and it has pure
dimension nine.  It may have
several components, and the argument below initially sums all of them.

For a tuple in (E8.1), let $J\subset G$ be its generated subgroup.
The first inertia has no eigenvalue one, so $V^J=0$.  The moving-rank
sum is
\[
                 D=8+4+10=22.
\]
The parabolic-cohomology dimension is consequently $D-2\cdot8=6$
over $K$.  Chevalley--Weil gives the two Hodge dimensions three:
the age sum for $V$ is $4+2+10/2=11$, and the age sum for its dual
is also $22-11=11$, so each dimension is $11-8=3$.
These statements hold for every tuple, whether or not $J=G$.

The signed Hermitian determinant formula D-DET-01 applies since
$V^J=0$; it does not require that the image be all of $G$.
For a reflection the local positive determinant quotient is one.
For $C_2$ a column basis of $\operatorname{im}(1-g)$ has Gram
determinant16 and $\det(1-g)=16$, again giving one.  For $C_1$ the
moving space is all of $V$, with $\det H=1$ and $\det(1-g)=6$.
Thus its quotient is $1/6$, whose Gaussian norm class is three.
It follows that all the counted components have signed determinant
class three and signature $(3,3)$.

There is an actual abelian sixfold on each component after a finite
level base change.  The rational $M_8(K)$-isotypic summand of the
Jacobian of the total curve has the form $B^8$ up to isogeny; a
primitive self-adjoint projector gives $B$.  One can use the
rank-one projector $(1-s_j)/2$ inside this simple group-algebra
factor.  It is self-adjoint for the positive $G$-invariant root form,
and hence for the induced Rosati involution.  The center supplies
the $K$-action on $B$, and the restriction of the Jacobian
polarization is positive.  The cohomology of $B$ is the multiplicity
space just computed, of rational dimension12.  For a disconnected
curve its Jacobian means the product of the Jacobians of its connected
components; all these assertions continue to hold.  Thus counting
disconnected covers introduces no virtual abelian variety.

\subsection*{D-E8-GRR-02. The degree-one Hodge class and its boundary normalization}

For any finite-order element $g$ put
\[
 a(g)=\frac12\sum_{j=0}^{r-1}m_{g,j}\frac jr\left(1-\frac jr\right),
 \qquad r=\operatorname{ord}(g),                               \tag{E8.3}
\]
where $m_{g,j}$ is the multiplicity of $e^{2\pi i j/r}$ in $\rho$.
This is a class function, and $a(g^{-1})=a(g)$.  If
$\mathcal F_\rho=\operatorname{Hom}_G(V_\rho,\pi_*\omega)$, write
$\lambda_\rho=c_1(\det\mathcal F_\rho)$ on
$\overline{\mathcal H}$.  The rank of this bundle is three.

For a reduced boundary component $D_\alpha$, denote by $g_\alpha$
its node holonomy and by $r_\alpha$ its order.  Classes $\psi_j$
below are cotangent classes of the \emph{coarse base curve}.
Then
\[
 \boxed{\lambda_\rho=\sum_{j=1}^{12}a(C_j)\psi_j
          -\sum_\alpha a(g_\alpha)r_\alpha[D_\alpha].}          \tag{E8.4}
\]
When a boundary stratum has multiple branches, the sum is taken on
its normalizations with the usual pushforwards.  Formula (E8.4)
therefore has an intrinsic interpretation on the stack.

\begin{proof}
Apply Tseng's published formula (7.3.7.2) to $X=BG$ and the vector
bundle associated to $V_\rho$.  Here $TBG=0$, so the virtual class is
the actual fundamental class.  All positive-degree rational Chern
characters of a vector bundle on $BG$ vanish.  Put
\[
 A_2(g)=\sum_jm_{g,j}B_2(j/r),\qquad
 B_2(x)=x^2-x+\frac16.
\]
Thus $A_2(g)/2=8/12-a(g)$.  In degree one, the main term of Tseng's
formula is $8\kappa_1/12$, and the marking term is
$-\sum_j A_2(C_j)\psi_j/2$.  The node term in that formula is
$r^2 A_2(g)/4$ times the pushforward from the node gerbe with a
choice of branch.  The node gerbe has degree $1/r$ over the reduced
boundary, and the choice of branch has degree two.  Its contribution
is therefore
\[
                       \frac r2 A_2(g)[D].                    \tag{E8.5}
\]
This calculation includes the factor $r^2$ displayed in Tseng's
formula; dropping that factor would give a different intersection
problem.

Tseng uses the stack with sections to the marking gerbes.  By his
Remark2.4.2, p.~19, this is the fiber product of the marking gerbes
over the stack without those sections.  Its degree is
$\prod_jr_j^{-1}$.  The weighted class in \S2.5.1, p.~21, multiplies
by $\prod_jr_j$, and hence gives exactly the ordinary-cover mass
specified in D-E8-GRR-01.  Formula (E8.5) is unchanged on descending
to that normalization.  This also checks that no marking factor is
missing from the final answer.

The associated orbibundle has
$H^q(\mathcal C,V_\rho)=(H^q(C,\mathcal O_C)\otimes V_\rho)^G$.
Its $H^0$ vanishes by $V^J=0$.  Serre duality identifies its $H^1$
with $\mathcal F_\rho^*$, so the degree-one Chern character of its
derived pushforward is precisely $\lambda_\rho$.  Combining the
three terms just computed yields
\[
 \lambda_\rho=\frac8{12}
  \left(\kappa_1-\sum_j\psi_j+\sum_\alpha r_\alpha[D_\alpha]\right)
  +\sum_ja(C_j)\psi_j
  -\sum_\alpha a(g_\alpha)r_\alpha[D_\alpha].                 \tag{E8.6}
\]
The expression in parentheses is zero.  For example, apply the same
published GRR formula to the trivial representation: the derived
pushforward on a genus-zero base is the trivial line, its first
Chern class is zero, and $A_2=1/6$ gives exactly this relation.
Equivalently it is the pullback of the genus-zero Mumford relation,
since a node smoothing has the local form $t=s^{r_\alpha}$.
This proves (E8.4) from the published formula.
\end{proof}

For later use, the normal class of a reduced boundary is
\[
 [D_\alpha]|_{D_\alpha}
       =-\frac{\psi_++\psi_-}{r_\alpha}.                     \tag{E8.7}
\]
Indeed the coarse smoothing parameter is $t=s^{r_\alpha}$ and its
normal line is the tensor product of the two coarse tangent lines.
The normal line of $s=0$ is its $r_\alpha$-th root.  The local
equation also proves that the ramification-weighted class
$r_\alpha[D_\alpha]$ has the ordinary smooth-tuple mass over the
corresponding base boundary: a reduced boundary stratum has that
mass divided by $r_\alpha$.  At several nodes the factors multiply.
These facts can also be checked against Jarvis--Kimura's
Proposition3.3, equations(30)--(38), pp.~14--16, but that preprint is
not required for the proof.

\subsection*{D-E8-GRR-03. A complete rooted-tree recurrence}

We prove the finite recurrence used for the ninth power of (E8.4).
The argument applies to any finite group and any inversion-invariant
weight $a$.  Let $\mathcal A=Z(\mathbb Q[G])$, with class basis
$K_c=\sum_{g\in c}g$, ordinary central multiplication, and trace
\[
                  \tau(x)=\frac{[1]x}{|G|}.
\]
The trace pairing satisfies
\[
 \tau(K_cK_d)=\begin{cases}|c|/|G|,&d=c^{-1},\\0,&d\ne c^{-1}.
 \end{cases}                                                   \tag{E8.8}
\]
It is the gluing pairing for covers.  To see this directly, choose
trivializations of the torsor at one base point of each component
of a nodal genus-zero tree.  Each vertex imposes that the product
of its outgoing holonomies is one.  Identifying the two fibers at
an edge matches inverse holonomies.  Changes of trivialization
on the vertices can be eliminated along a spanning tree, leaving
one simultaneous conjugation by $G$.  Thus gluing yields the mass
of product-one tuples divided by $|G|$.  In the class basis the
inverse pairing is $|G|/|c|$; contracting an edge cancels one
vertex normalization and is exactly ordinary multiplication in
$\mathcal A$.  This proves the assertion without a connectedness
assumption on the torsor.

Write $\mathcal W$ for the diagonal linear operator
$\mathcal W(K_c)=a(c)K_c$.  Fix one marking as a root.  Let
$\mathbf c=(c_1,\ldots,c_q)$ count the other marked leaves by their
classes and put $k=\sum c_j$.  The following quantities are central
vectors whose coefficients are exponential generating coefficients
in the leaf counts, namely divided by $\prod_jc_j!$.

The vector $T[\mathbf c,a]$ sums rooted stable subtrees with a bare
root flag carrying an insertion $\psi_{\rm root}^{a}$; the root
itself has no marking exponential and no attached edge.  Every
other marking carries its exponential, and every internal edge
carries the boundary contribution derived below.  The possible
root exponent is $0\le a\le k-2$.

For a child attached to a parent flag carrying exponent $b$, let
$L[\mathbf c,b]$ include the additional factor $1/b!$ from the
parent vertex's psi integral.  A single marked leaf of class $c$
then has
\[
              L[\mathbf e_c,b]=\frac{a(c)^b}{(b!)^2}K_c.       \tag{E8.9}
\]
One factorial comes from its exponential and one from that psi
integral.  For a child subtree with $k\ge2$ the rule is
\[
 L[\mathbf c,b]=-
    \sum_{a=0}^{k-2}
       \frac{\mathcal W^{a+b+1}T[\mathbf c,a]}
            {(a+b+1)a!(b!)^2}.                               \tag{E8.10}
\]
Let $F_m[\mathbf c,b]$ be the coefficient of leaf count
$\mathbf c$ and total parent exponent $b$ in $L^m/m!$, with
multiplication in $\mathcal A$.  Then
\[
 T[\mathbf c,a]=
       \sum_{m=a+2}^{k}\frac{(m-2)!}{a!}
                    F_m[\mathbf c,m-2-a].                   \tag{E8.11}
\]
Finally, if the root marking has class $c_0$, the top intersection
on the $(k-2)$-dimensional Hurwitz stack is
\[
 (k-2)!\prod_jc_j!\;
 \tau\left(K_{c_0}\sum_{a=0}^{k-2}
                     \frac{a(c_0)^a}{a!}T[\mathbf c,a]\right).
                                                                  \tag{E8.12}
\]

\begin{proof}
Expand $\exp(t\lambda)$ using (E8.4).  Distinct compatible boundary
divisors give a stable tree; incompatible boundary divisors have
empty intersection.  Consider an edge whose reduced divisor is
$D$, whose inertia order is $r$, and whose weight is $a(c)$.
If this divisor occurs $\ell\ge1$ times in the expansion, its
coefficient, before the reduced-stratum mass is taken, is
\[
 \frac{(-t a(c)r)^\ell}{\ell!}
       \left(-\frac{\psi_++\psi_-}{r}\right)^{\ell-1}
 =-r\,\frac{t^\ell a(c)^\ell}{\ell!}
                 (\psi_++\psi_-)^{\ell-1}.                   \tag{E8.13}
\]
The reduced-stratum mass has the compensating factor $1/r$ by
(E8.7).  Orienting the edge toward the root chooses one of its
two branches, so it also removes the two-to-one cut-edge convention
in GRR.  The edge series in the central-algebra state sum is
therefore exactly
\[
       -\frac{\exp(t a(c)(\psi_++\psi_-))-1}
                   {\psi_++\psi_-}.                         \tag{E8.14}
\]
The coefficient of $\psi_+^a\psi_-^b$ is
\[
       -\frac{t^{a+b+1}a(c)^{a+b+1}}{(a+b+1)a!b!}.           \tag{E8.15}
\]
This proves (E8.10), including its sign and both endpoint
factorials; the extra $1/b!$ is the parent-vertex factorial.
Inversion of the child holonomy is immaterial because $a(c)$ is
inversion invariant.

For completeness, the genus-zero psi integral at a vertex of
valence $m+1$ is
\[
 \int_{\overline M_{0,m+1}}\psi_0^a\prod_{j=1}^m\psi_j^{b_j}
   =\frac{(m-2)!}{a!\prod_jb_j!},\qquad
                  a+\sum_jb_j=m-2.                           \tag{E8.16}
\]
This follows by induction from the three-point integral one and
the string equation.  In fact at least one exponent is zero;
forget that marking and apply the string equation.  The sum of
the resulting numerators is the sum of the remaining exponents,
which is $m-2$, giving (E8.16).  The string equation itself follows
from $\psi_j=f^*\psi_j+D_{j,\mathrm{forgotten}}$; the exceptional
sections are disjoint and their normal classes are the negative
cotangent classes.  Thus (E8.16) needs no assumption about the
Hurwitz cover.

A stable root vertex has $m\ge2$ children.  Its children already
include their factors $1/b_j!$.  Applying (E8.16), and summing
unordered children by the exponential construction $L^m/m!$,
gives (E8.11).  Restoring the exponential of the actual root
marking adds $a(c_0)^a/a!$, explaining the second root factorial
in (E8.12).  Restoring the labels of same-class leaves multiplies
by $\prod_jc_j!$.  Since we expanded an exponential, extracting
the top power rather than its exponential coefficient multiplies
by $(k-2)!$.  The last gluing normalization is the trace $\tau$.
This proves (E8.12).

There is no infinite series issue.  A term of $T[\mathbf c,a]$
has Hodge degree $k-2-a$, and a term of $L[\mathbf c,b]$ has
degree $k-1+b$.  For total dimension $d$, one can discard
$b>d-k+1$.  Every child in (E8.11) has fewer leaves than its
parent, so the recurrence is acyclic and terminates.  These
observations also prove the truncations in the supplied program.
\end{proof}

\subsection*{D-E8-GRR-04. Exact evaluation from the published character table}

For the classes in (E8.1), (E8.3) gives
\[
                     a(C_1)=\frac{17}{24},\quad
                     a(C_2)=\frac12,\quad a(R)=\frac18.       \tag{E8.17}
\]
The finite character table used in the evaluation is the complete
112-class table for $W(E_8)$ in the CHEVIE author's
\texttt{weyle8.g}.  Its primary URL and immutable downloaded hash are
\begin{quote}\small
\url{https://raw.githubusercontent.com/jmichel7/gap3-jm/master/pkg/chevie/tbl/weyle8.g}\\
SHA256 \texttt{7dff885669b7b86c45acb4b3b144ce673f}\allowbreak
\texttt{b0440736bd81b221d963694048e254}.
\end{quote}
The published package and the author's manual are cited in
\cite{E8GRR:Michel,E8GRR:CHEVIE}.  The complete class classification
is also Carter's published Table11, pp.~52--54
\cite{E8GRR:Carter}.  Completeness and the character status of the
CHEVIE table are standard finite-group inputs; orthogonality alone
is not being claimed as a proof that an arbitrary array consists
of genuine characters.

\texttt{build\_character\_input.py} specializes the published
polynomial Hecke-character data at $v=q=1$.  Omitted alternate rows
in that table are recovered by the documented sign twist, using
$(-1)^{\ell(w)}$ on each class word.  The retained file contains
the resulting integer character table, the degrees, the class
sizes and words, the central-translation permutation, and both
sets of Hodge weights.  It redistributes mathematical data, not
the entire third-party library.

\texttt{check\_character\_input.py} can recheck the retained
input without fetching the external library.  It reconstructs
all112 class matrices from the displayed Cartan matrix and words;
checks their exact orders and orthogonality; factors every exact
characteristic polynomial into cyclotomic polynomials; shifts
all exponents by $1/4$ for the $iW$ coset; and recomputes all224
weights by (E8.3).  It checks the trace of the natural character
on every reconstructed matrix, the complete character
orthogonality identity, the sum of class sizes and of squared
degrees, and multiplication by the central negative identity.
The natural character is row68.

Only rational arithmetic in the $W$-class algebra is needed.
Every subtree has at most one $C_2$ leaf and all its other leaves
are real reflections.  Thus it belongs to either $\mathbb Q[W]$
or $i\mathbb Q[W]$, and no multiplication inside a subtree
multiplies two imaginary factors.  For $x=\sum_cx_cK_c$ in the
$W$-class algebra its Fourier coordinates and their inverse are
\[
 \widehat x_r=\sum_cx_c\frac{|c|\chi_r(c)}{\chi_r(1)},\qquad
 x_c=\frac1{|W|}\sum_r\chi_r(1)\chi_r(c)\widehat x_r.          \tag{E8.18}
\]
All these Weyl characters are real.  The program checks that
these two matrices multiply to the identity, and multiplies
central vectors coordinatewise in Fourier coordinates.  Applying
$\mathcal W$ in (E8.10) is done in the class coordinates, using the
weight for the real or imaginary coset as appropriate.

At the final root multiplication there are exactly two imaginary
factors.  If $z=-I$, then $(iw)(iv)=zwv$.  Therefore the final
identity coefficient is the coefficient at the negative of the
root Weyl class, multiplied by that class size.  For root class57
this is class58.  The final normalization is $|G|=2|W|$, not
$|W|$.  This explains the central-product reduction used by the
program without treating the224-character table as an opaque
input.

Running \texttt{exact\_hodge\_volume.py} with ten reflections,
one $C_2$ leaf and root $C_1$ evaluates (E8.9)--(E8.12) exactly.
The result is
\[
 \boxed{\int_{\overline{\mathcal H}}\lambda_\rho^9
                       =\frac{3009825}{64}>0.}                \tag{E8.19}
\]
Independently of the Hodge weights, the ordinary tuple mass is
\[
 \boxed{\tau(K_{C_1}K_{C_2}K_R^{10})=90296156160000.}          \tag{E8.20}
\]
The retained output records exact rational arithmetic, the checked
transform identity, and the input hash.  The numerical values
from the unrefereed lead are compared only after the calculation;
they are not input to it.  The program also permits exchanging
$C_1$ and $C_2$ as the root, giving a differently rooted evaluation
of the same intersection.  The supplied verification record
distinguishes this comparison from an independent recurrence.

\subsection*{D-E8-GRR-05. Canonical extension and actual rank nine}

At least one component of $\mathcal H$ has a Weil-sixfold period
map of generic rank nine.  Consequently it is dominant onto the
unitary moduli component determined by its rational polarized
$K$-space, of signature $(3,3)$ and signed determinant three.

\begin{proof}
Let $\mathcal H_\beta$ be any irreducible component.  Pass to a
finite level cover, take a projective resolution of its
compactification, and resolve the boundary to obtain a smooth
projective variety $S$ with a normal crossings divisor $D$ and
a generically finite map to $\overline{\mathcal H}_\beta$.
After a further finite cover if necessary, local monodromy is
unipotent.  The admissible curve family is semistable; semistable
resolution after the base change does not change its Hodge
bundle.  Marking connected components of a disconnected family
requires only a further finite cover.

We check that the GRR extension used above pulls back to the
canonical Hodge extension required by the positivity theorem.
For a one-parameter semistable family of curves the extension of
$F^1$ is $\pi_*\omega$.  This published identification is recalled
in the proof of Theorem2.1 of \cite[p.~242]{E8GRR:CD}, with the
canonical extension characterized by nilpotent residues.  Locally
at a node $xy=t$, the relative logarithmic generator is
$dx/x=-dy/y$, which is precisely the dualizing generator.  Thus
semistable resolution and base change add no divisor to the
determinant line.  Apply this one-parameter identification to a
transverse disk at the generic point of every prime boundary
divisor of $S$, including its exceptional divisors.  The Hodge
bundles and canonical extensions are locally free; the
identification over $S\setminus D$ therefore extends over $S$
because it and its inverse have no codimension-one poles on the
normal variety $S$.  This proves the equality of extensions,
not just equality over the open family.  Constant rational
group-algebra projectors commute with this identification.
The direct-sum argument applies to disconnected curves as well.

Let $\mathbb V_{\rm full}$ be the rational full isotypic
variation for $\rho\oplus\overline\rho$.  It has rational rank96,
and its $F^1$ has rank48.  The fixed representation factors have
complex dimension eight, so its extended determinant line $L$
has
\[
 c_1(L)=8(\lambda_\rho+\lambda_{\overline\rho})
                         =16\lambda_\rho.                    \tag{E8.21}
\]
The last equality holds in rational first Chern classes on the
compactification: the weights (E8.3), and hence (E8.4), are
unchanged on dualizing the representation.  This is stronger
than an equality of ranks or of open Chern forms.

Apply Bakker--Brunebarbe--Tsimerman Lemma6.17 to
$\mathbb V_{\rm full}$ on $S\setminus D$.  Its hypotheses hold:
the variation is polarized and real, $S$ is smooth projective,
the boundary has normal crossings, and its local monodromies
are unipotent.  The lemma says that $L$ is nef, and that it is
big exactly when the Hodge differential
\[
 T_{S\setminus D}\longrightarrow
       \operatorname{Hom}(F^1,\mathbb V_{\rm full,\mathcal O}/F^1)
                                                                  \tag{E8.22}
\]
is injective as a map of sheaves.  In weight one this is the
period differential.  Sheaf injectivity is equivalent to
injectivity at the generic point.  On a smooth projective
ninefold a nef line bundle is big exactly when its ninth
self-intersection is positive, as also used in the published
proof of Lemma6.15.

It follows that the contribution
$\int_{\overline{\mathcal H}_\beta}\lambda_\rho^9$ of every
component is nonnegative: its pullback is $16^{-9}c_1(L)^9$,
and generically finite alterations multiply intersections by
their positive degree.  Formula (E8.19) is the sum of these
nonnegative contributions, so at least one is strictly
positive.  On that component (E8.22) is generically injective.
The full variation and its primitive abelian sixfold have the
same period differential kernel, since the former is a fixed
multiplicity of the latter.  The primitive period map therefore
has rank nine.

Finally, a polarized $K$-space of signature $(3,3)$ has a
unitary period domain of complex dimension nine.  The
algebraic family of sixfolds just constructed gives a map to
its moduli space after level structure.  A rank-nine point
gives an analytic open image, hence Zariski-dense image in
that moduli component.  This proves dominance.  The source
and target dimensions enter only at this final step, after
the full period rank has been established.
\end{proof}

\subsection*{D-E8-GRR-06. The dominant component consists of connected covers}

Using the Gaussian reflection-normal-closure reduction D-GAUSS-03
proved earlier in this work, the dominant component obtained in
D-E8-GRR-05 has full generated group $G$.  Thus its total curves
are connected.  This conclusion does not claim that every component
with branch datum (E8.1) is dominant or connected.

\begin{proof}
We first justify the generic endomorphism assertion needed here.
On the full unitary period domain of signature $(3,3)$, the only
rational endomorphisms which are Hodge for every point are the
given scalars $K$.  Here is a direct verification.  Identify the
real vector space with $\mathbb C^3\oplus\mathbb C^3$ and the
Hermitian form with $\operatorname{diag}(I_3,-I_3)$.  One of its
compatible complex structures is
$J_0=\operatorname{diag}(iI_3,-iI_3)$.  Write a real endomorphism
as a complex-linear matrix $A$ plus a complex-antilinear matrix
$B$ followed by conjugation.  Commutation with $J_0$ forces
$A$ to be block diagonal and $B$ to be block off-diagonal.
The tangent directions to compatible complex structures are,
up to a nonzero real factor,
\[
       \begin{pmatrix}0&-iZ\\ i\overline Z^{,t}&0\end{pmatrix},
                   \qquad Z\in\operatorname{Mat}_3(\mathbb C).
\]
Commutation with all these directions implies
$A_+Z=ZA_-$ for every $Z$, hence $A_+=A_-=aI_3$.
For the off-diagonal blocks of $B$ it implies
$B_{12}Z^t=ZB_{21}$ for every $Z$.  Taking $Z=I$ gives
$B_{12}=B_{21}=B_0$; taking diagonal $Z$ makes $B_0$ diagonal;
taking the off-diagonal matrix units then makes every diagonal
entry zero.  Thus $B=0$.  The real commutant is the complex
scalars, and its intersection with rational endomorphisms is
exactly $K$.

There are only countably many rational endomorphisms.  For each
one outside $K$ its Hodge locus in the domain is a proper closed
analytic subset by the calculation just made.  A nonempty open
subset cannot be a countable union of such subsets.  The
dominant period map in D-E8-GRR-05 therefore has a very general
fiber $B$ with $\operatorname{End}^0(B)=K$.  In particular $B$
is simple.  This argument uses the actual open period image and
does not assume a generic Mumford--Tate theorem.

Pass to a finite cover of the Hurwitz component to identify the
generated subgroup $J\subset G$ up to a fixed conjugation, and
decompose the coefficient representation $V|_J$ over $K$.
The semisimple algebra
\[
 \operatorname{End}_{KJ}(V)=\prod_\alpha
                         M_{m_\alpha}(D_\alpha)
\]
acts on its parabolic multiplicity Hodge structure $P$ (with
the opposite-algebra convention if the multiplicity is written
as a Hom).  Every summand projector is a Hodge projector.
Since the full rational Hodge endomorphism algebra of $P$ is
$K$, exactly one factor acts nontrivially.  A nonzero
homomorphism from a simple algebra to a field is injective,
so that factor must have $m_\alpha=1$ and $D_\alpha=K$.
Consequently there is a unique active absolutely irreducible
$K$-summand $\tau$ of $V|_J$, of dimension at most eight,
occurring once, and
\[
       P=H^1_{\rm par}(\mathbb P^1\setminus\{p_j\},\tau).
                                                                  \tag{E8.23}
\]
All the other summands have zero parabolic cohomology.
The decomposition can be made orthogonal for the positive
finite-group form.  Thus (E8.23) identifies the actual polarized
variation, not merely its Hodge numbers.

The character field of $\tau$ is $K$.  Indeed its matrices are
over $K$, so the only possible proper character subfield is
$\mathbb Q$.  If its character were rational, $\tau$ and its
conjugate would be isomorphic over $K$: an intertwiner is
obtained by solving the intertwining linear equations over
$K$, and any nonzero intertwiner between these absolutely
irreducible representations is invertible.  Composing this
intertwiner with scalar conjugation gives an invertible
$K$-antilinear, rational $J$-equivariant coefficient map.
It induces an invertible rational Hodge endomorphism of (E8.23)
which does not commute with $i$.  This contradicts
$\operatorname{End}_{\rm Hdg}(P)=K$.  This argument also excludes
a rational-character quaternionic representation; no implicit
Schur-index-one assertion over $\mathbb Q$ is being made.

Each of the ten reflection inertias in (E8.1) acts on $\tau$
either trivially or as an involutive reflection, because its
moving space on all of $V$ has dimension one.  At least one
acts nontrivially.  Otherwise the local system in (E8.23)
would have only two possible nontrivial monodromies, which
are inverse to each other.  The first has no fixed vectors
even on $\tau$, as it has none on $V$.  If $a=\dim_K\tau$,
the moving-rank sum would then be $2a$, and the genus-zero
parabolic dimension would be $2a-2a=0$, contrary to
$\dim_KP=6$.

The finite image of $J$ on $\tau$ is therefore absolutely
irreducible, has character field $K$, and contains reflections.
Its actual parabolic variation is the dominant nonsplit
sixfold variation (E8.23).  D-GAUSS-03 applies, including its
induced-representation cases.  It says that every irreducible
component of its reflection normal closure which can remain
in this situation is the natural $E_8$ representation.
Such a component has dimension eight.  Since
$\dim_K\tau\le8$, equality is forced, $\tau=V|_J$, and the
reflection normal closure in $J$ contains a copy of $W(E_8)$
in the natural representation.

All reflections in $G$ belong to its real subgroup $W$.
The reflection normal closure just found is consequently a
subgroup of the original $W$.  Its order is the full order
$696729600$ of $W(E_8)$, so it is $W$.  The tuple also
contains its specified element of $iW$, hence $J=G$.
The generated subgroup is constant up to conjugacy on a
Hurwitz component.  Its principal cover is connected exactly
when that subgroup is $G$.  This proves the assertion.
\end{proof}

\paragraph{Precise conclusion.}
There exists a dominant family of connected $G$-covers whose
actual primitive abelian factors have $K=\mathbb Q(i)$,
signature $(3,3)$, and signed determinant class three.
D-E8-GRR-01--05 establish dominance without a connectedness
premise; D-E8-GRR-06 proves connectedness using the separately
proved reflection-normal-closure barriers.  The positive total
intersection does not identify a particular explicit Nielsen
tuple as lying in the dominant component.  The construction of
algebraic detecting cycles is a separate assertion.

\section{A directly certified dominant Gaussian $E_8$ component}
\label{sec:im-e8-period}

\paragraph{Provenance and precise assertion.}
This section proves that at least one twelve-branch component for
$\mu_4W(E_8)$ has period rank nine and raw Hermitian determinant
class $[-3]$. The existence proof uses a candidate point from Benjamin
Severn's unrefereed manuscript
\emph{Weil classes on abelian sixfolds of Weil type for
$\mathbf Q(i)$ with discriminant $-3$}, draft 0.7, 21 October 2026,
Hexagon 2610.00477v1, Appendix B, and its numerical archive
Zenodo 23549514. A supplied approximate point is not a theorem input.
The exact target matrices, 87 rational candidate coordinates, and
nine fixed gauge coordinates are retained as mathematical input;
the contraction, its side conditions, and the period minor are
recomputed. The source's approximate continuation, its claimed
outputs, and its main theorems are not premises.

The programs \texttt{e8\_ball/verify\_M.py} and
\texttt{e8\_ball/verify\_L.py} are explicitly acknowledged adaptations
of the source algorithms, whose complete implementations and error
bounds were read and checked. They replace the unavailable continuation
input by the saved candidate, use safe JSON data, and recompute every
certificate inequality. The target-group and path program
\texttt{check\_group\_and\_paths.py} was independently written.
The proof below supplies the mathematical content of the certification.
It proves dominance for the component containing this smooth cover;
it does not identify that component with the two-cancelling-pair
component D-E8T-01.

\subsection*{D-E8R-01. Exact target monodromy and the covering construction}
Let $H$ be the Cartan matrix of D-E8T-01. In the fundamental-weight
basis, whose Gram matrix is $H^{-1}$, use
\[
a=\begin{pmatrix}
0&0&1&0&0&0&0&0\\
1&1&1&1&1&1&0&0\\
1&1&1&2&2&2&2&1\\
-1&-1&-1&-1&-1&-1&-1&0\\
-1&-1&-2&-3&-3&-2&-1&-1\\
0&0&0&0&1&0&0&0\\
1&1&2&3&2&2&1&0\\
1&2&2&3&2&1&1&1
\end{pmatrix},
\]
\[
b=\begin{pmatrix}
-1&-2&-3&-5&-4&-3&-2&-1\\
2&3&3&5&4&3&2&1\\
2&2&4&5&4&3&2&1\\
-2&-2&-3&-4&-4&-3&-2&-1\\
0&0&0&0&1&0&0&0\\
0&0&0&0&0&1&0&0\\
0&0&0&0&0&0&1&0\\
0&0&0&0&0&0&0&1
\end{pmatrix},\qquad
c=I-e_1e_1^t+(-e_1+e_3)e_1^t.
\]
Thus $c$ is the simple reflection in the first root in this basis.
The target input \texttt{mtuple.json} also records $d=(abc)^{-1}$,
although identifying the order of a product of based loops will not
be needed for the analytic argument.

The three matrices belong to $W(E_8)$ and generate it. Their exact
membership certificate successively right-multiplies
$H^{-1}aH,H^{-1}bH,H^{-1}cH$ by simple reflections until the identity
is reached, retaining all word indices. Starting with the first
simple root and closing under these three matrices in the simple-root
basis gives all 240 roots. Since $c$ is a root reflection, conjugation
then supplies every root reflection, proving generation. The same
certificate checks $\chi_a=\Phi_{20}$ and that $b,c$ are involutions
with one-dimensional moving spaces. Thus no order-only identification
of the generated group is involved.

We will certify a Fuchsian system
\[
\frac{dY}{dz}=
 \left(\frac{B_0}{z}+\frac{B_1}{z-1}
                    +\frac{B_2}{z-(2+i)}\right)Y,
\qquad
B_0=\operatorname{diag}(1,3,7,9,11,13,17,19)/20,
\tag{R1}
\]
whose three based monodromies are simultaneously conjugate to
$a,b,c$. Here $B_1=u_1v_1^t$, $B_2=u_2v_2^t$ have rank one and
trace $1/2$. Pull (R1) back by $z=s^5$ and tensor its solution with
$s^{1/4}$. Its monodromy is the full group $G=\mu_4W(E_8)$.

To prove the latter statement without assuming a particular formula
for a twelve-entry based tuple, let $\Pi$ be the fundamental group
of the four-punctured $z$-sphere and let $\Pi'$ be the subgroup for
$s^5$. It is the kernel of the epimorphism to $\mathbf Z/5$ which
sends the loop around zero to one, those around $1,2+i$ to zero, and
that around infinity to minus one. The image of $\Pi'$ in $W(E_8)$
is normal with cyclic quotient of order dividing five. Since a group
generated by involutions has no nontrivial quotient of order five,
this image is all of $W(E_8)$. After the scalar twist the projective
image remains $W(E_8)/\{\pm I\}$. At each of the five preimages of
one, the twisted system still has a real root reflection as local
monodromy. Conjugating one such reflection by the projective image
supplies all root reflections, hence all of $W(E_8)$. The monodromy
at zero is conjugate to $ia^5$, which is outside $W(E_8)$. Therefore
the full image is $\mu_4W(E_8)$.

The finite monodromy kernel supplies, by Riemann existence, a smooth
connected algebraic $G$-cover of the $s$-sphere. Its twelve branch
points are zero, the fifth roots of one, the fifth roots of $2+i$,
and infinity. The exact local exponents and the Weil signature are
verified below after the existence certificate.

\subsection*{D-E8R-02. Rigorous integration and a finite contraction test}
The candidate parameters are the entries of
\[
x=(v_1,\ (u_{2,j})_{j\ne j_0},\ v_2,\ Q)\in\mathbf C^{87}.
\]
The eight entries of $u_1$, the index $j_0=1$ in zero-based notation,
and the coordinate $u_{2,j_0}$ are fixed dyadic rational complex
numbers, recorded in \texttt{center\_g.json}. The 87 candidate
coordinates are likewise specified exactly as pairs of integers
$(m,e)$ representing $m2^e$ at precision 256. The decimal source
strings are retained only to make extraction auditable. The program
asserts that their 256-bit midpoints equal the recorded dyadics.

We use the maximum absolute-value norm on vectors and its induced
maximum-row-sum matrix norm. A complex interval is a rectangular
product of two real Arb intervals. Each binary floating-point number
selected by NumPy is thereafter treated as its exact dyadic rational
value. Floating linear algebra chooses pivots or preconditioners;
it never certifies invertibility or decides whether a period minor
is nonzero. All such decisions use enclosing Arb calculations.

\paragraph{Transport error bound.}
For a regular expansion center $c_0$ of an equation
$Y'=A(z)Y$, $A(z)=\sum_k A_k/(z-p_k)$, suppose
\[
d\le\min_k|c_0-p_k|,\qquad S\ge\sum_k\|A_k\|_\infty.
\]
The Taylor coefficients $A_n$ at $c_0$ satisfy
$\|A_n\|_\infty\le S/d^{n+1}$. If
$T(z)=\sum_{n\ge0}T_n(z-c_0)^n$ is the transport with $T_0=I$,
its recursion is
\[
(n+1)T_{n+1}=\sum_{j=0}^n A_jT_{n-j}.
\]
Coefficientwise induction compares its norms with the scalar solution
of $f'=S(d-t)^{-1}f$, $f(0)=1$. Consequently
\[
\|T_n\|_\infty\le\frac{(S)_n}{n!d^n}.
\]
For a step $h$, put $q\ge|h|/d$, with $q<1/2$. After degree $N$,
the absolute tail in each matrix entry is at most
\[
\epsilon_N=
\frac{(S)_{N+1}q^{N+1}/(N+1)!}
 {1-q(1+S/(N+2))},
\tag{R2}
\]
provided the denominator is positive. Indeed the ratio of consecutive
terms after the first omitted term is at most
$q(1+S/(N+2))$. This argument is uniform when the residues range
over an input interval box. Adding $\epsilon_N$ to both coordinate
radii encloses the omitted complex entry.

The code computes either this Taylor recursion or the equivalent
polynomial recurrence obtained by clearing
$\prod_k(z-p_k)$. It evaluates the truncated transport by Horner's
rule and multiplies successive interval matrices. The two recurrences
solve the same differential equation and use exactly the same bound
(R2); the direct Taylor version is used for wider input boxes to
reduce interval inflation. Bounds exported to binary floats use the
appropriate Arb endpoint followed by one \texttt{nextafter} step
outward. Products and sums used in certificate inequalities are
returned to Arb before they are compared. No printed rounded decimal
is used as an inequality test.

The three polygonal loops start at $1/2-i$, travel to the near side
of the radius-$0.45$ circle about the relevant pole, traverse a
counterclockwise twelve-gon, and return. Their vertices are exact
binary rationals after selection. The independently written path
check verifies winding numbers $(1,0,0),(0,1,0),(0,0,1)$ by exact
rational cross products; after shrinking the small circles to radius $1/10$, their straight
access paths are disjoint except at the base point and avoid the other
small discs, giving the usual free generating loops. The exact minimum
distance from any other pole to such a straight access segment is
$1/5$; the polygon vertices and their access segments are homotopic
to these loops within the corresponding puncture complements. Adaptive steps
are dyadic and are checked to have $|h|/d<1/2$. All endpoint
additions use at most 56 numerator bits, hence are exact at precision
256. The actual three paths use 31, 31, and 53 steps. The final step
is the exact difference to the endpoint; an explicit assertion rules
out premature termination caused by a rounded endpoint.
There is also an independent exact rational comparison of the rounded
adaptive polyline with each intended polygon edge, implemented in
\texttt{e8\_path\_tube\_certificate.py}. Every adaptive vertex is
within $10^{-12}$ of that edge, while every intended edge is at
distance greater than $1/10$ from every pole. The closed
$10^{-12}$-neighbourhood of a line segment is convex and contains
no pole. Both its adaptive subpath and its intended segment therefore
lie in this neighbourhood and are homotopic relative to their common
endpoints. Concatenation gives the claimed based-loop homotopy.
Thus the computed matrices are monodromies of the specified closed
generating loops, not merely of paths with the same winding numbers.

\paragraph{The 87 equations and their side conditions.}
Let $M_0,M_1,M_2$ be these monodromies as holomorphic functions of
the residue parameters. For the eight distinct eigenvalues
$\lambda_j=\exp(2\pi i m_j/20)$ of $a$, construct exact algebraic
eigenvectors $e_j$ using a seven-row solve and one coordinate equal
to one. Ball inversion encloses those eigenvectors. Write
$b=I-x_by_b^t$ and $c=I-x_cy_c^t$ exactly, with
$y_b^tx_b=y_c^tx_c=2$, and choose rational bases of
$\ker y_b^t$ and $\ker y_c^t$.
The equations $F(x)=0$ are the following:
\begin{enumerate}
\item $v_1^tu_1-1/2=0$ and $v_2^tu_2-1/2=0$;
\item seven selected rows of $(M_0-\lambda_jI)Qe_j=0$ for each $j$;
\item seven selected rows of $(M_1+I)Qx_b=0$, and one fixed row
of $(M_1-I)Qw=0$ for each of seven basis vectors $w\in\ker y_b^t$;
\item the analogous fourteen equations for $M_2,c$;
\item $\det Q-1=0$.
\end{enumerate}
Their number is $2+56+14+14+1=87$. The dropped rows for the eight
$M_0$ equations are $0,1,\ldots,7$. The negative-eigenvector dropped
row and positive-eigenvector pivot row are $(4,3)$ for $M_1$ and
$(3,1)$ for $M_2$, in zero-based notation.

At a zero, the rank-one residues have trace $1/2$, so their local
monodromies are reflections. The fixed residue $B_0$ has distinct
nonresonant exponents and its monodromy has the stated eight simple
eigenvalues. The following additional inequalities are certified
throughout the small solution ball: the seven retained rows of
$M_0-\lambda_jI$ have rank seven; the seven retained rows of each
$M_k+I$, $k=1,2$, have rank seven; and the selected row of $M_k-I$
is nonzero. Each rank assertion is witnessed by a $7\times7$
minor excluding zero. Hence the seven-row kernels in the equations
are the full kernels. For $M_k-I$, its rank is one at a solution,
so the kernel of a nonzero row is its full kernel. Thus a zero
satisfies
\[
M_0Q=Qa,\qquad M_1Q=Qb,\qquad M_2Q=Qc,\qquad \det Q=1.
\tag{R3}
\]
This justifies the reduction to 87 equations; none of the omitted
rows is inferred merely from a small numerical residual.

\paragraph{Proof of existence.}
The map $F$ is holomorphic in all 87 parameters: fundamental-solution
transport along fixed nonsingular paths depends holomorphically on
the residues. Let $x_0$ be the recorded dyadic candidate, $n=87$,
$R=2^{-40}$, and $h=2^{-80}$. Compute an interval enclosure of
$F(x_0)$ and of every forward difference
\[
D_{ij}=\frac{F_i(x_0+he_j)-F_i(x_0)}h.
\]
For the 64 variables belonging only to $Q$, the same center monodromy
matrices can be reused exactly; this removes redundant ODE solves
without changing the forward differences. Let $Y_0$ be the exact
dyadic matrix obtained from a floating approximate inverse of the
midpoint of $D$. Enclose $F$ on the complex polydisc of radius $R$
by evaluating it on the containing real-coordinate interval box,
and write $K_i$ for an upper bound for $|F_i|$ there.

Cauchy's estimate, applied on the $j$th coordinate disc, gives
\[
|D_{ij}-\partial_jF_i(x_0)|
 \le\frac{K_i h}{R(R-h)}.
\tag{R4}
\]
For completeness, the coefficients of the one-variable Taylor
series are bounded by $K_i/R^k$; subtracting its first derivative
from the forward difference leaves
$\sum_{k\ge2}K_i h^{k-1}/R^k$, which is (R4).
For $\|x-x_0\|_\infty\le\rho<R$, twice applying Cauchy's estimate
and integrating along the segment from $x_0$ to $x$ gives
\[
|\partial_jF_i(x)-\partial_jF_i(x_0)|
 \le\frac{2n\rho K_i}{(R-\rho)^2}.
\tag{R5}
\]
The factor two includes the equal-coordinate second derivative;
for distinct coordinates the smaller factor one would suffice.
Summing a row and using submultiplicativity therefore gives
\[
\begin{split}
\eta&=\|Y_0F(x_0)\|_\infty,\\
E&=\max_i\frac{nK_i h}{R(R-h)},\qquad
C=\max_i\frac{2n^2\rho K_i}{(R-\rho)^2},\\
\kappa&=\|I-Y_0D\|_\infty+\|Y_0\|_\infty(E+C).
\end{split}
\tag{R6}
\]
In the actual computation $\rho$ is the outward-rounded value
$10\eta$. The fresh run gives, with these figures rounded only for
display,
\[
\eta\simeq1.052812\,10^{-35},\quad
\rho\simeq1.052812\,10^{-34},\quad
\kappa\simeq0.005589859013,
\]
\[
\max_i K_i<1.172\,10^{-6},\qquad
\|I-Y_0D\|_\infty<2.216\,10^{-9},\qquad
\|Y_0\|_\infty<54.87.
\]
Arb verifies $\kappa<1$ and $\eta+\kappa\rho<\rho$.
The map $x\mapsto x-Y_0F(x)$ is consequently a strict contraction
of the closed maximum-norm ball of radius $\rho$ about $x_0$ into
itself. It has a unique fixed point. Moreover
$\|I-Y_0D\|_\infty<1$ makes $Y_0D$ invertible and therefore makes
$Y_0$ invertible, so a fixed point has $F(x)=0$. The side-minor tests
are performed on the entire enclosing radius-$\rho$ input box.
This proves existence of (R1) with the exact finite target monodromy
(R3), not merely proximity to finite monodromy.

\subsection*{D-E8R-03. Exact exponents and the polarized Weil factor}
Put $B_\infty=-B_0-B_1-B_2$. A numerical eigenbasis of its midpoint
is treated as an exact dyadic invertible matrix. Interval inversion
and multiplication then give Gershgorin discs, each contained within
$10^{-10}$ of the following rational labels, counted with
multiplicity:
\[
-1/12,\ -1,\ -11/12,\ -2/3,\ -7/12,\ -1/3,\ -5/12,\ -1.
\tag{R7}
\]
This localization is recomputed in Arb in both verification stages;
the displayed floating disc radii in the diagnostic log are not used
as a proof. The disjoint clusters of discs count one eigenvalue at
each single label and two at $-1$.

The labels are exact. By the residue--monodromy theorem, the
exponential of $2\pi i$ times a residue eigenvalue is a local
monodromy eigenvalue; see \cite[Corollary 16.20, author extract
p.~274]{IM:IlyashenkoYakovenko}. Every monodromy here is conjugate
to an element of the rational eight-dimensional representation of
$W(E_8)$. An eigenvalue of primitive order $e$ therefore has
$\varphi(e)\le8$. The complete possible list is
\[
1,2,3,4,5,6,7,8,9,10,12,14,15,16,18,20,24,30.
\]
One can verify completeness without a class table: each prime
$p\mid e$ has $p-1\le8$, hence $p\in\{2,3,5,7\}$; their possible
powers divide $2^4,3^2,5,7$, and checking their products gives exactly
the displayed list. Different candidates have rational angles
separated by far more than $10^{-10}$ from the labels (R7).
Thus the residue eigenvalues are precisely (R7), even without
fixing an ordering convention for the based-loop product at infinity.

For later rational differential calculations use the integer scalar
gauge $s^{-3}$ in addition to the scalar twist. The local exponents
of the pulled-back system are
\[
5m_j/20+1/4-3\quad\text{at }s=0,
\qquad (0^7,1/2)\quad\text{at each of the ten reflection points},
\]
\[
5e_j-1/4+3\quad\text{at }s=\infty,
\qquad e_j\text{ as in (R7)}.
\tag{R8}
\]
The gauge changes no monodromy. At zero these exponents modulo one
are four zeros and four halves. At infinity they are
\[
1/12,\ 1/6,\ 1/3,\ 5/12,\ 2/3,\ 3/4,\ 3/4,\ 5/6.
\]
Their sum is four. Thus the moving-rank and age sums are 22 and 11,
respectively. Chevalley--Weil, in the form
\cite[Theorem 2.10, author manuscript p.~8]{IM:FredianiGhigiPenegini},
gives the two holomorphic multiplicities $11-8=3$.
The natural character has field $K=\mathbf Q(i)$ and is realizable
over $K$, so a rational Rosati-fixed primitive projector gives an
actual polarized abelian sixfold $A_0$, with signature $(3,3)$.
This is the same group-algebra construction used in D-E8T-02.

Its determinant class is $[-3]$. Every real reflection has local
factor one in the determinant formula D-DET-01, using a root of
norm two. For the element $ia^5$, the first independent columns of
$I-ia^5$ give a moving-space Gram determinant 16 and determinant
$\det(I-ia^5|_{\mathrm{mov}})=16$, so its factor is one. At infinity
all eight directions move, the full Gram determinant is one, and
the displayed spectrum gives $\det(I-g_\infty)=6$; its local factor
is $1/6$. The signed six-dimensional determinant class is therefore
$[1/6]=[3]$, and the raw determinant class is $[-3]$. These are
computations for the actual invariant positive form, not an arbitrary
auxiliary Hermitian form.

\subsection*{D-E8R-04. Rational descriptions of the Hodge eigenspaces}
Change coordinates to $t=1/(s-s_*)$, where $s_*$ is the exact
binary rational complex value represented by $0.37+0.21i$. Then
$t=\infty$ is a regular point. The system has twelve finite simple
poles $p$, with residues $A_p$, whose exponents are (R8), and total
residue zero. Let $Y$ be a fundamental solution. Let $D_+$ and
$D_-$ denote the two three-dimensional natural multiplicity spaces
in $H^0(C_0,\omega_{C_0})$.

There are canonical identifications, up to interchanging the two
embeddings of $K$,
\[
D_+=\{r\in\mathbf C(t)^8:Y^{-1}r\,dt
                        \text{ is holomorphic on }C_0\},
\]
\[
D_-=\{q\in\mathbf C(t)^8:Y^tq\,dt
                        \text{ is holomorphic on }C_0\}.
\tag{R9}
\]
Indeed $Y$ becomes single-valued after pullback to the finite
monodromy cover. Its entries and those of its inverse have moderate
growth at the branch points, and hence extend meromorphically on
that cover. If $r$ is rational, $Y^{-1}r\,dt$ has the appropriate
contragredient transformation law. Conversely, for a vector of
holomorphic forms with that transformation law, multiplying by $Y$
and dividing by $dt$ gives a $G$-invariant meromorphic vector on the
cover. It descends to a meromorphic vector on the sphere and is
therefore rational. The argument for $Y^tq$ is identical with the
dual representation. These are spaces of algebraic holomorphic
forms on the actual compact algebraic curve, not spaces defined only
by a numerical differential equation.

Choose a residue eigenbasis $P_p$ and write a local fundamental
solution in the form
\[
Y=P_pH_p(t-p)(t-p)^{\Lambda_p},\qquad H_p(0)=I,
\qquad \Lambda_p=\operatorname{diag}(\lambda_{p,j}).
\tag{R10}
\]
Such a holomorphic frame exists also at resonant poles in this
calculation. One justification from the published local theory is
as follows. The holomorphic Poincare--Dulac--Levelt form has
monodromy $\exp(2\pi i\Lambda)\exp(2\pi iN)$, with commuting
semisimple and unipotent factors, where $N$ is nilpotent
\cite[Theorems 16.15--16.16 and Corollary 16.19,
author extract pp.~270--274]{IM:IlyashenkoYakovenko}.
Finite monodromy forces $\exp(2\pi iN)=I$, hence $N=0$.
The resonant coefficients and the residue's nilpotent part occupy
disjoint matrix-entry positions according to the eigenvalue
difference, so each vanishes. The normal form is the Euler system,
which gives (R10). Alternatively, before the fifth pullback all
four residue spectra lie in intervals of length less than one;
nonresonant Euler frames pull back to (R10).

The holomorphy conditions in (R9) are exactly
\[
\operatorname{ord}_p[(P_pH_p)^{-1}r]_j\ge\lfloor\lambda_{p,j}\rfloor,
\qquad
\operatorname{ord}_p[(P_pH_p)^tq]_j\ge\lfloor-\lambda_{p,j}\rfloor,
\tag{R11}
\]
together with regularity away from these poles and vanishing to
order at least two at infinity. To verify the floor precisely, let
$t-p=y^e$ on the cover. The first form's $j$th component has order
\[
e\operatorname{ord}_p[(P_pH_p)^{-1}r]_j
                    -e\lambda_{p,j}+e-1.
\]
As $e\lambda_{p,j}$ is an integer, this is nonnegative exactly when
the first inequality (R11) holds. The second follows by replacing
$\lambda$ by $-\lambda$. At infinity $dt$ has a pole of order two
and the system is regular, giving the stated vanishing condition.

Here is why the resonant recursion used in the certificate computes
a valid frame. Write $H_p=I+\sum_{m\ge1}H_mt^m$. Its $(i,j)$
coefficient at order $m$ has denominator
$m+\lambda_j-\lambda_i$. At a zero denominator choose the free
coefficient $(H_m)_{ij}$ to be zero. Starting from any frame (R10),
right multiplication of its fundamental matrix by a constant
unipotent matrix changes $H$ to
$H(t)t^\Lambda C t^{-\Lambda}$. An entry $C_{ij}$ for
$\lambda_i-\lambda_j=m>0$ first changes the resonant coefficient
of degree $m$. Choosing these entries in increasing order of $m$
sets all such coefficients to zero without altering lower ones.
Hence this normalization exists. The remaining recursion coefficients
are uniquely forced and converge by the holomorphic classification
theorem. The exact resonant right-hand sides are zero for this
geometric reason; the numerical test that their intervals contain
zero is only a consistency check, not a proof of vanishing.

\subsection*{D-E8R-05. Finite linear algebra for the actual codifferential}
For the first space in (R9), put
$M_p=\max(0,-\min_j\lfloor\lambda_{p,j}\rfloor)$, and use the
analogous definition with $-\lambda$ for the second. A rational
vector in either space has the unique partial-fraction expression
\[
r(t)=\sum_p\sum_{m=1}^{M_p}c_{p,m}(t-p)^{-m}.
\tag{R12}
\]
This follows because $P_pH_p$ and its inverse are holomorphic,
so these are upper bounds on all pole orders; vanishing at infinity
eliminates a polynomial part. Its order-two vanishing at infinity
adds the eight linear equations $\sum_p c_{p,1}=0$.

Expanding (R11) imposes finitely many linear equations on the
coefficients (R12). For the present exponents they give matrices
\[
\Phi_+:\mathbf C^{48}\longrightarrow\mathbf C^{45},\qquad
\Phi_-:\mathbf C^{120}\longrightarrow\mathbf C^{117}.
\]
The code computes the local frames through order twelve. Every
coefficient required by (R11) is within this range, as is checked
from the exact floors and pole orders; no infinite series tail is
used for this step. The eigenbases themselves are enclosed by
solving exact-eigenvalue kernel equations using square pivot blocks.
Their true kernel dimensions are known from the already proved
semisimplicity and exact exponent multiplicities.

Select 45 dependent columns in $\Phi_+$ and 117 in $\Phi_-$.
The corresponding square blocks $S_+,S_-$ have certified inverse
enclosures. More explicitly, using exact midpoint inverse
preconditioners $R_+,R_-$, the fresh run bounds
\[
\|I-R_+S_+\|_\infty<1.650\,10^{-24},\qquad
\|I-R_-S_-\|_\infty<3.385\,10^{-26}.
\]
Therefore both blocks are invertible. Set each of the remaining
three free coordinates in turn to one, the others to zero, and solve
the selected rows. These vectors form the kernel basis of those
rows. They are the complete kernel bases of $\Phi_\pm$: by
Chevalley--Weil the true kernels have dimension three, they lie
inside the three-dimensional selected-row kernels, and hence equal
them. Checking that all unused-row intervals contain zero is a
consistency test; the dimension argument, not that interval test,
proves that all rows vanish exactly.

We next prove that the resulting matrix is the period
codifferential. For a deformation of a curve, represent its
Kodaira--Spencer class by vector fields $v_{ab}$ on overlaps of
coordinate charts. Differentiating the gluing of a holomorphic
one-form $\omega$ gives $\mathcal L_{v_{ab}}\omega
=d(\iota_{v_{ab}}\omega)$. Thus the infinitesimal variation class
is represented in $H^1(\mathcal O_C)$ by
$\iota_{v_{ab}}\omega$. Pairing with a second holomorphic one-form
$\eta$ by Serre duality pairs the Kodaira--Spencer class with the
quadratic differential $\omega\eta$. This is the multiplication
description of the Torelli codifferential, derived here on charts.

For fixed-$G$ deformations, take invariants. Locally a cover has
$t=y^e$. An invariant vector field upstairs is
$y f(y^e)\partial_y$, whose image is
$e t f(t)\partial_t$. Hence
$(\pi_*T_C)^G=T_{\mathbf P^1}(-B)$, and the Hurwitz tangent space
is $H^1(T_{\mathbf P^1}(-B))$. Dually, invariant quadratic
differentials upstairs identify with
$H^0(\mathbf P^1,\omega_{\mathbf P^1}^{\otimes2}(B))$.
Indeed a possible pole downstairs has pullback order
$e\operatorname{ord}(q)+2e-2$, which for $e\ge2$ permits at most
a simple pole. On the natural nonreal representation pair, Schur's
lemma identifies the invariant contraction with the tensor product
of the two multiplicity spaces. Restricting the multiplication map
to this block is precisely the codifferential of the projected
abelian factor: the projector is an algebraic flat endomorphism of
the Jacobian variation. In the coordinates (R9) the contraction is
\[
\mu:D_+\otimes D_-\longrightarrow
 H^0(\mathbf P^1,\omega_{\mathbf P^1}^{\otimes2}(B)),\qquad
\mu(r\otimes q)=q^tr\,dt^2.
\tag{R13}
\]
Indeed $(Y^tq)^t(Y^{-1}r)=q^tr$. Conventions may multiply this
identification by a fixed nonzero scalar; they cannot change its rank.

All twelve branch points are finite in the $t$ coordinate. A
quadratic differential in the target is uniquely determined by its
simple-pole coefficients $\rho_p$, and these satisfy exactly
\[
\sum_p\rho_p=\sum_pp\rho_p=\sum_pp^2\rho_p=0.
\tag{R14}
\]
To see this, its rational coefficient is
$\sum_p\rho_p/(t-p)$, and regularity at infinity requires decay
$O(t^{-4})$. Expanding at infinity gives precisely (R14).
Thus the target has dimension $12-3=9$. For the three certified
bases on each side of (R13), form the nine rational contractions
$q_b^tr_a$, expand by the finite partial fractions (R12), and take
their twelve residues. The $9\times12$ residue matrix has the same
rank as the true codifferential.

\subsection*{D-E8R-06. A nonzero certified minor and dominance}
Using the freshly established solution enclosure, the minor on
zero-based residue columns $3,4,5,6,7,8,9,10,11$ is contained in
\[
[0.10557441259506563\ \mathord\pm\ 6.51\,10^{-18}]
+i[-0.04409991368479836\ \mathord\pm\ 3.88\,10^{-18}].
\tag{R15}
\]
The enclosure excludes zero, and its absolute-value lower bound is
greater than $0.11441485472524$. Matrix determinants, basis solves,
local frame coefficients, branch points, and residues here are all
evaluated in enclosing 256-bit arithmetic. The approximation of the
residues of (R1) is included from the outset as an input interval of
radius $\rho$ about the proven exact center. Consequently (R15)
contains the determinant at the exact finite-monodromy solution,
not just at its numerical midpoint. The three residue relations
(R14) are also checked as interval consistency conditions.

It follows that the differential of the actual polarized natural
factor period map at $C_0$ has rank nine. Its source Hurwitz
component has dimension nine, and the target signature-$(3,3)$
unitary domain has dimension nine. Passing to a finite marking and
level cover gives an actual family of abelian factors. The holomorphic
inverse-function theorem makes the local period image open in that
domain. The algebraic map to the fixed polarization and endomorphism
moduli component is therefore dominant. By D-E8R-03 that component
has raw determinant class $[-3]$.

This conclusion proves existence of a dominant nonsplit Gaussian
$E_8$ natural source. It gives neither a classification of its
Nielsen components nor a proof that every component, or the
particular two-pair component D-E8T-01, has rank nine. A supply of
algebraic Weil classes and a uniform bound for their detecting
cycles remain separate questions.

\paragraph{Reproduction and source limits.}
Run, in \texttt{agents/imprimitive/e8\_ball}, first
\texttt{python check\_group\_and\_paths.py}, then
\texttt{python verify\_M.py}, then \texttt{python verify\_L.py}.
The first step is exact rational algebra. The second is the
existence and side-condition proof; the third reads its newly
written enclosure and proves (R15). The successful run used
Python 3.12, python-flint 0.9.0, and 256-bit Arb arithmetic
\cite{IM:JohanssonArb}; NumPy and SciPy serve only for candidate
pivots and approximate preconditioners. The original continuation
solution is unnecessary and is not silently assumed to exist.
The compact input/provenance manifest records source and extracted
hashes. The full third-party archive is not part of the deliverable.

\section{Determinant suppliers controlled by the branch exponent}
\label{sec:branch-exponent}

\paragraph{Provenance and logical status.}
The Abel-fibre strategy originates in Schoen's published construction
\cite[\S2, pp.~12--17]{G34:Schoen}.  The arbitrary-branching version was
located in the unrefereed OpenAI manuscript \cite[Lemma~3.4]{Root:Gale},
and the improvement from the character order to the common inertia
exponent was located in \cite[\S2]{Root:Rstation66}.  These are credited
as leads, not assumed as theorems.  The proof below reconstructs the
geometric argument, including the direction of the returning
correspondence, from the stated published inputs.  In particular it
does not use the global CM Hodge conjecture, 032, or a Severn theorem.
The quantitative degree of this construction is a separate assertion;
no numerical detector bound follows just from the word ``supplied.''

Here a sub-Hodge structure of the cohomology of a smooth projective
variety is \emph{supplied by} a smooth projective variety $F$ if it is
contained in the sum of the images of finitely many algebraic
correspondences from powers of $F$, with Tate twists.  Coefficients may
first be extended to a number field.  The rational correspondence span
is stable under scalar conjugation, tensor products, algebraic
projectors, and composition.  These assertions follow directly by
external product and composition of correspondences.

\begin{theorem}[HD-FS1: the inertia-exponent cyclic supplier]
\label{HD-FS1}
Let $q:C\to B$ be a connected cyclic cover of smooth projective
complex curves, of degree $m\ge2$, and let $\chi$ be a faithful character
of its deck group.  Suppose every inertia order divides the positive
integer $M$.  Then
\[
 D_\chi=\det H^1(C,\mathbf Q(\zeta_m))_\chi
\]
is supplied by the Jacobian of the degree-$M$ Fermat curve.  Unramified
covers need only point suppliers.  The assertion includes a
zero-dimensional eigenspace, whose determinant is the unit.
\end{theorem}
\begin{proof}
Write $b=g(B)$ and let $D=t_1+\cdots+t_r$ be the reduced branch divisor.
The relation in the punctured surface group, followed by its cyclic
quotient, shows that $r\ne1$.  Put
\[
              v=\dim H^1(C,\mathbf Q(\zeta_m))_\chi=2b-2+r.
\]
Indeed the character part of $q_*\mathbf Q(\zeta_m)$ is a rank-one
local system on $B\setminus D$, extended by zero at $D$; its Euler
characteristic is $2-2b-r$.  Its degrees zero and two vanish because
the cover is connected and the character is nontrivial.  If $v=0$
there is nothing to prove.  Assume $v>0$.

\emph{The alternating line as quotient cohomology.}
Let $\Delta$ be the cyclic deck group, and let $N$ be the kernel of
the product homomorphism
\[
          \Delta^v\rtimes\mathfrak S_v\longrightarrow\Delta.
\]
The map is a homomorphism since $\Delta$ is abelian.  Set
$X=C^v$ and $Z=X/N$; there is a finite map
$\pi:Z\to\operatorname{Sym}^v B$ with residual group $\Delta$.
Finite quotient pullback identifies rational cohomology with the
invariants: proper base change kills higher direct images, and
averaging makes invariants exact.  In particular $H^v(Z)$ is pure,
since it injects into $H^v(X)$.
The $\chi$-eigenspace of $H^v(Z)$ maps to $D_\chi$ in $H^v(X)$.
To verify this, each curve factor must have degree one and character
$\chi$; the geometric permutation of degree-one classes contributes
its Koszul sign, so invariance under permutations is ordinary
alternation.  The resulting exterior power is a line.

On $U=\operatorname{Sym}^v(B\setminus D)$, the residual cover is a
$\Delta$-torsor, including at divisors with repeated points.  In a
small disk, sheet labels modulo $N$ remember only their product,
and elementary symmetric functions give coordinates on the symmetric
power.  At a divisor containing a branch point the character stalk
is zero, since its stabilizer contains nontrivial inertia.  Thus
$(\pi_*\mathbf Q(\zeta_m))_\chi$ is extension by zero from a
rank-one local system on $U$.  Its monodromy on a moving divisor
is the product of the original cover monodromies along the moving
point paths.  This description is valid across diagonals by the
local torsor calculation.

\emph{Localization at one Abel fibre.}
First suppose $r>0$, so $v\ge2b$.  Let
\[
 a:\operatorname{Sym}^v B\longrightarrow\operatorname{Pic}^v B,
 \qquad L_0=K_B(D),\qquad
 T=(\pi^{-1}(|L_0|))_{\rm red}.
\]
For $L\ne L_0$ the evaluation map
$H^0(B,L)\to\bigoplus_i L|_{t_i}$ is surjective.  Its obstruction
is dual to the section space of the nontrivial degree-zero bundle
$K_B\otimes L^{-1}(D)$.  Its kernel has dimension $b-1$.
Consequently
\[
                |L|\cap U\simeq\mathbf A^{b-1}\times
                                      (\mathbf C^*)^{r-1}.
\]
The rank-one local system on this product has a nontrivial circle
monodromy.  A transverse loop about an evaluation hyperplane winds
the product of the nearby zero coordinates once around the relevant
branch point, so its monodromy is that inertia character.  This also
handles multiple zeros, by the product of roots in a local monic
equation.  A nontrivial rank-one local system on a circle has cellular
complex $k\xrightarrow{\alpha-1}k$ with $\alpha\ne1$; tensoring
these complexes proves vanishing on the torus, and duality gives
compact-support vanishing.  Proper base change and Leray therefore
give the restriction isomorphism
\[
                    H^v(Z)_\chi\xrightarrow{\sim}H^v(T)_\chi.
\tag{FS.1}
\]
If $b=0$, the Picard variety is a point and the same assertion is
immediate.

The residue sequence identifies
\[
 H^0(B,K_B(D))\simeq W\oplus R,
 \quad W=H^0(B,K_B),\quad
 R=\{(w_1,\ldots,w_r):\textstyle\sum_iw_i=0\},
\]
after choosing a vector-space splitting.  The evaluation hyperplanes
are $w_i=0$ in $\mathbf P(W\oplus R)$.  For $r\ge3$ they are
distinct.  Over its function field the residual torsor has a Kummer
equation
\[
                    z^m=c h^m\prod_{i=2}^r(w_i/w_1)^{a_i}.
\tag{FS.2}
\]
Here $c\in\mathbf C^*$ and $h$ is a rational function.  To justify
the assertion globally, its valuations away from the hyperplanes
are divisible by $m$.  Choose the integers $a_i$ to match the
hyperplane valuations modulo $m$; the degree-zero condition matches
the remaining hyperplane.  The residual divisor is $mE$ and the
torsion-free Picard group of projective space makes $E$ principal.
The local transverse-loop calculation above shows that the order
of $\exp(2\pi i a_i/m)$ divides $M$.  Equivalently,
\[
                             m\mid Ma_i.
\tag{FS.3}
\]
This is the step that replaces the character order by the inertia
exponent.

Let $F_M^{r-2}\subset\mathbf P^{r-1}$ be the Fermat hypersurface
$\sum u_i^M=0$ and use its finite coordinate-power map
$s:F_M^{r-2}\to\mathbf P(R)$, $[u_i]\mapsto[u_i^M]$.
The smooth projective bundle
\[
        Y=\mathbf P_{F_M^{r-2}}
              ((W\otimes\mathcal O)\oplus s^*\mathcal O(-1))
\tag{FS.4}
\]
maps generically finitely onto $\mathbf P(W\oplus R)$.
After this pullback each factor in (FS.2) is an $m$th power by
(FS.3).  The generic torsor splits.  Rational lifts to its components
extend to morphisms: a rational lift to a finite scheme over the
base is regular on a normal source by integral closure of every
local ring.  Taking finitely many copies of components of $Y$ gives
a proper surjection $Y'\to T$.  Every component of $T$ dominates the
base: its translates cover every fibre, which is one $\Delta$-orbit.
Thus this surjection misses no component.  For $r=2$ the evaluation
hyperplanes coincide, the Kummer expression has no nonconstant
factor, and $T$ is instead dominated by a finite union of projective
spaces.  This is also the dimension-zero Fermat case.

Shioda--Katsura's published Theorem~1.7 and Proposition~2.4
\cite{Root:SK} give algebraic correspondence supply of the full
cohomology of $F_M^d$ by powers of $J(F_M^1)$, for every $d$.
Indeed their degree-$M$ rational map from a product of lower
Fermat hypersurfaces is resolved by their displayed smooth blowup;
the blowup formula, induction, and push-pull by a generically finite
map give surjectivity on cohomology.  Projective-bundle cohomology
is a direct sum of Tate twists of the base via powers of its
tautological class.  Hence the entire cohomology of $Y'$ has the
required supply.

For completeness, if $r=0$, then $v=2b-2$ and $b\ge2$.  The
nontrivial cover character defines a rank-one local system $E$ on
$\operatorname{Pic}^vB$, and its pullback by $a$ is the character
local system on the symmetric power.  The Abel fibres are
$\mathbf P^{b-2}$ except over $K_B$, where the fibre is
$\mathbf P^{b-1}$.  The powers of a relatively ample class show
that the even direct images up to degree $2b-4$ are constant;
the remaining direct image is a rank-one skyscraper in degree
$2b-2$ at $K_B$.  Tensor with $E$ and use the same torus
acyclicity on the Picard variety.  This proves (FS.1), now with
$T$ a split finite cover of the simply connected $|K_B|$.
It is a union of projective spaces and is supplied by points.

\emph{Returning from detection to supply.}
Deligne's published Proposition~8.2.5, p.~39, together with
Theorem~8.2.4(iii) \cite{Root:DeligneIII}, gives
\[
 \ker\bigl(H^v(T,\mathbf Q)\longrightarrow H^v(Y',\mathbf Q)\bigr)
                       =W_{v-1}H^v(T,\mathbf Q).
\tag{FS.5}
\]
The pure line from $H^v(Z)$ has nonzero restriction by (FS.1)
and cannot meet this kernel, by strictness of mixed Hodge structures.
Choose a component $Y_0$, of dimension $d$, on which a nonzero
$\alpha\in H^v(Z)_\chi$ survives.  Poincare duality gives
$\beta\in H^{2d-v}(Y_0)$ with
$\int_{Y_0}p^*\alpha\,\beta\ne0$, where $p:Y_0\to Z$.
In the branched case $2d-v=r-2\ge0$; in the unramified case it is
zero.  The full-cohomology supply just proved applies to $\beta$.

Take a component of $Y_0\times_Z X$ dominating $Y_0$ and resolve
it projectively to $\widetilde Y$, with maps
$h:\widetilde Y\to Y_0$ and $l:\widetilde Y\to X$.
Its cycle $(h,l)_*[\widetilde Y]$ defines the correspondence
$l_*h^*$ and
\[
       \int_X Q^*\alpha\, l_*h^*\beta
           =\deg(h)\int_{Y_0}p^*\alpha\,\beta\ne0,
\tag{FS.6}
\]
where $Q:X\to Z$.  No supply of the resolution's cohomology is
required.  Apply the inverse-character projector for
$\Delta^v\rtimes\mathfrak S_v$, a finite linear combination of
algebraic automorphism graphs.  In degree $v$ its image is precisely
$D_{\chi^{-1}}$, by the earlier Koszul calculation, and the
nonzero pairing (FS.6) is unchanged.  Thus that line is supplied.
Conjugating coefficient scalars supplies $D_\chi$ as well.
All maps are algebraic correspondences, and the only nontrivial
fixed source is $J(F_M^1)$.
\end{proof}

\begin{theorem}[HD-FS2: finite-group transfer with one Fermat exponent]
\label{HD-FS2}
Let a finite group $G$ act on a smooth projective curve, allowing a
disjoint union of connected curves permuted transitively.  Suppose
every inertia order in its quotient cover divides $M$.  Over a
splitting field, the determinant of every irreducible multiplicity
space in $H^1$ is supplied by $J(F_M^1)$.  If all powers of this
Jacobian have the rational Hodge conjecture, every balanced rational
sum of these determinant lines is algebraic.
\end{theorem}
\begin{proof}
First take a connected cover.  Integral Brauer induction
\cite[\S10.5, Theorems~19--20]{G34:Serre} gives an identity
\[
 V\oplus R_-\simeq R_+,
 \qquad R_\pm=\bigoplus_j
        (\operatorname{Ind}_{H_j}^G\chi_j)^{\oplus n_{j,\pm}},
 \quad n_{j,\pm}\ge0,
\tag{FS.7}
\]
with linear characters $\chi_j$.  Frobenius reciprocity and quotient
pullback identify their multiplicity spaces with the corresponding
faithful character spaces for
$C/\ker\chi_j\to C/H_j$.  The inertia generators of each of these
intermediate covers are powers of conjugates of the original
generators.  Their orders therefore divide $M$.  HD-FS1 applies.
The trivial character contributes the determinant of $H^1(C/H_j)$,
which is a Tate line supplied by a top polarization class.

Apply $\operatorname{Hom}_G(-,H^1(C))$ to (FS.7) and take
determinants.  It gives
$D_+=D_V\otimes D_-$.  Tensor with the determinant for the
conjugate of the minus representation and contract the two minus
determinants by the nondegenerate polarization pairing.  This
contraction is algebraic: the $H^1$-Kunneth projector and the
diagonal of a pointed curve realize the pairing and its inverse,
and the finite group algebra and matrix units realize the required
summands.  Consequently the tensor of the supplied plus and
conjugate-minus lines maps nontrivially onto $D_V$, with a Tate
shift.  This cancels an actual dual factor; it does not take a root
of a cycle class.

For a disconnected $G$-cover, choose one connected component,
with stabilizer $H$.  Its full cohomology induces to that of the
disjoint union.  Restrict $V$ to $H$, decompose it into irreducibles,
and apply the connected assertion and determinant tensor products.
Again the inertia orders divide $M$.

Finally a rational Hodge class in an algebraic image of a sum of
polarizable pure Hodge structures has a rational Hodge preimage:
take the orthogonal complement of the kernel with respect to a
polarization.  Under the stated Hodge-conjecture hypothesis this
preimage is algebraic on a power of the fixed Jacobian.  Its
algebraic-correspondence image is algebraic.  Extending coefficients
and then descending the rational cycle-class image deals with the
individual character lines.  No assertion is made that an
unbalanced character determinant is itself a Hodge class.
\end{proof}

\paragraph{The $E_8$ application and its remaining numerical step.}
Every branch order in HD-E8-02 and D-E8T-01 divides $12$.
Thus the fixed source in HD-FS2 is $J(F_{12})$, irrespective of
the much larger exponent of $\mu_4W(E_8)$ or of the orders of
Brauer's inducing characters.  A controlled elliptic decomposition
of this fixed Jacobian and explicit bounds for the correspondence
degrees will be recorded separately.  The present theorem proves
the supply statement, not a numerical value of $\delta_W$.

\section{Explicit Fermat-twelve isogeny and bounded Hodge-cycle
suppliers}\label{explicit-fermat-twelve-isogeny-and-bounded-hodge-cycle-suppliers}

Author/checker: GPT-6 Astra Pro; inherited effort not exposed;
HODGE\_DIVISION /root/sixfolds. Date: 2026-11-21. New quantitative
proof; independent whole-proof review PASS reported by root
on2026-11-21, with the same model and inherited-effort metadata.

Write \(C=F_{12}\), \(J=J(C)\), and let \(\Theta\) be its canonical
principal polarization. Put \(N=144\). The explicit conclusions are

\[
 \phi:E_i^{27}\times E_\omega^{28}\longrightarrow J,\qquad
 \deg\phi\le144^{55},
 \tag{F12.1}
\]

and, after choosing the elliptic coordinates described below,

\[
 \frac14\Theta_0\ \le\ \phi^*\Theta\
       \le4\cdot144^2\Theta_0
 \tag{F12.2}
\]

in the cone of nef rational divisor classes. Here \(\Theta_0\) is the
product principal polarization.

For every \(s\ge1\), put \(g_s=55s\) and let
\(\Theta_s=\sum_{a=1}^s\operatorname{pr}_a^*\Theta\) on \(J^s\). There
is an explicitly indexed finite set of effective integral cycles of
codimension \(k\), for each \(0\le k\le g_s\), whose classes span

\[
 H^{2k}(J^s,\mathbf Q)\cap H^{k,k}(J^s)
\]

and whose degrees satisfy

\[
 \deg_{\Theta_s}Z
 \le 8^k(55s)!\,144^{55s}.
 \tag{F12.3}
\]

The statement is rational spanning by actual effective integral cycles.
It does not assert that these classes form a saturated basis of the
integral Hodge lattice. Such spanning is sufficient to choose a
bounded-degree cycle pairing nontrivially with any specified nonzero
rational Hodge class of complementary degree.

\subsection{SF-F1201. Integral projectors and the exact elliptic
census}\label{SF-F1201}

Use the affine equation \(x^{12}+y^{12}=1\). The diagonal automorphisms

\[
 \sigma_1(x,y)=(\zeta_{12}x,y),\qquad
 \sigma_2(x,y)=(x,\zeta_{12}y)
\]

generate \(G\simeq(\mathbf Z/12)^2\). Adjunction for this smooth plane
curve gives the holomorphic differential basis

\[
 \eta_{a,b}=x^{a-1}y^{b-12}\,dx,\qquad
 1\le a,b,\quad a+b<12.
 \tag{F12.4}
\]

The orders at \(x=0\) and \(y=0\) are respectively \(a-1\) and \(b-1\);
for the latter use \(dx=-y^{11}x^{-11}dy\). At infinity put \(z=1/x\)
and \(v=y/x\), so \(1+v^{12}=z^{12}\); the order is \(11-a-b\). These
orders are nonnegative. The fifty-five forms are independent because
their characters are distinct, and they exhaust the genus
\((12-1)(12-2)/2=55\). Their \(G\)-characters are \((a,b)\), meaning
\(\sigma_1^*\eta=\zeta_{12}^a\eta\) and
\(\sigma_2^*\eta=\zeta_{12}^b\eta\). The entire \(H^1(J,\mathbf C)\)
therefore has the 110 characters

\[
 \mathcal X=\{(a,b):1\le a,b\le11,\ a+b\not\equiv0\pmod{12}\},
\]

each with multiplicity one.

Let \(O\) run through the orbits of \((\mathbf Z/12)^*=\{1,5,7,11\}\) on
\(\mathcal X\), acting by simultaneous multiplication. Define the
rational idempotent

\[
 e_O=\frac1{144}\sum_{x,y\bmod12}
       \left(\sum_{(a,b)\in O}\zeta_{12}^{-ax-by}\right)
       \sigma_1^x\sigma_2^y.
 \tag{F12.5}
\]

The inner coefficient is a rational algebraic integer, hence an integer.
Thus \(144e_O\in\mathbf Z[G]\). Character orthogonality shows that these
projectors are mutually orthogonal and sum to the identity on
\(H^1(J)\). The corresponding identities hold as rational endomorphisms
of \(J\), since an endomorphism acting trivially on rational first
homology is zero.

Put \(A_O=\operatorname{im}(144e_O)\), with its connected image as an
abelian subvariety of \(J\). Its dimension is \(|O|/2\), which is one or
two. Let

\[
 O^+=\{(a,b)\in O:a+b<12\}.
\]

The following table gives a complete orbit census. A representative
stands for its orbit under the four units; the rows contain disjoint
orbits.

\begin{longtable}[]{@{}
  >{\raggedright\arraybackslash}p{(\columnwidth - 6\tabcolsep) * \real{0.2143}}
  >{\raggedleft\arraybackslash}p{(\columnwidth - 6\tabcolsep) * \real{0.2857}}
  >{\raggedleft\arraybackslash}p{(\columnwidth - 6\tabcolsep) * \real{0.2857}}
  >{\raggedright\arraybackslash}p{(\columnwidth - 6\tabcolsep) * \real{0.2143}}@{}}
\toprule\noalign{}
\begin{minipage}[b]{\linewidth}\raggedright
Character conductor
\end{minipage} & \begin{minipage}[b]{\linewidth}\raggedleft
Pure CM field
\end{minipage} & \begin{minipage}[b]{\linewidth}\raggedleft
Factor dimension
\end{minipage} & \begin{minipage}[b]{\linewidth}\raggedright
Number of factors
\end{minipage} \\
\midrule\noalign{}
\endhead
\bottomrule\noalign{}
\endlastfoot
\(12\) & \(\mathbf Q(i)\) & \(2\) & \(12\) \\
\(12\) & \(\mathbf Q(\omega)\) & \(2\) & \(9\) \\
\(6\) & \(\mathbf Q(\omega)\) & \(1\) & \(9\) \\
\(4\) & \(\mathbf Q(i)\) & \(1\) & \(3\) \\
\(3\) & \(\mathbf Q(\omega)\) & \(1\) & \(1\) \\
\end{longtable}

Complete representative lists, in the same row order, are:

\begin{itemize}
\item
  Conductor12, Gaussian: \((1,1)\), \((1,3)\), \((1,5)\), \((1,6)\),
  \((1,8)\), \((1,10)\), \((2,5)\), \((3,1)\), \((3,4)\), \((4,3)\),
  \((4,5)\), \((6,1)\).
\item
  Conductor12, Eisenstein: \((1,2)\), \((1,4)\), \((1,7)\), \((1,9)\),
  \((2,1)\), \((2,3)\), \((3,2)\), \((3,7)\), \((4,1)\).
\item
  Conductor6: \((2,2)\), \((2,4)\), \((2,6)\), \((2,8)\), \((4,2)\),
  \((4,6)\), \((4,10)\), \((6,2)\), \((6,4)\).
\item
  Conductor4: \((3,3)\), \((3,6)\), \((6,3)\).
\item
  Conductor3: \((4,4)\).
\end{itemize}

Here the conductor is \(12/\gcd(12,a,b)\). The field label has the
following explicit integral meaning. For each Gaussian row, one of
\(\sigma_1^3,\sigma_2^3\) acts by the same primitive fourth root on all
the characters of \(O^+\). For each Eisenstein row, one of
\(\sigma_1^4,\sigma_2^4\) acts by the same primitive third root on all
the characters of \(O^+\). This is verified by reducing the displayed
orbit coordinates modulo four or three, respectively. The certificate
records the chosen automorphism and scalar for every orbit.

Consequently \(A_O\) admits the full maximal order \(\mathbf Z[i]\) or
\(\mathbf Z[\omega]\) as integral endomorphisms, with a single complex
embedding on its entire tangent space. This is stronger than merely
identifying a rational CM field.

Both orders are Euclidean for the usual norm: nearest Gaussian and
Eisenstein lattice points have squared distances at most \(1/2\) and
\(1/3\), respectively. In particular both are principal ideal domains.
If \(A_O=V/\Lambda\) has dimension \(d\), its lattice \(\Lambda\) is a
torsion-free module of rank \(d\) over the indicated order, hence free.
An order basis of \(\Lambda\) is a complex basis of \(V\), because its
complex span contains the full lattice and therefore all of \(V\). It
follows that

\[
 A_O\simeq E_i^d\quad\hbox{or}\quad A_O\simeq E_\omega^d
 \tag{F12.6}
\]

as unpolarized abelian varieties. The maps are analytic homomorphisms
between projective complex tori and therefore algebraic. This argument
introduces no further isogeny degree.

The table has thirty-four factors. Summing their dimensions gives

\[
 12\cdot2+3=27,\qquad
 9\cdot2+9+1=28.
\]

This proves the asserted elliptic type by a complete finite computation
and the integral lattice argument, without a height estimate.

\subsection{SF-F1202. Polarization exponents and the isogeny
degree}\label{SF-F1202}

Every \(g\in G\) preserves the Jacobian polarization, so its Rosati
adjoint is \(g^{-1}\). Each orbit \(O\) is stable under negation. Thus
(F12.5) gives \(e_O^\dagger=e_O\). Distinct factor spaces are orthogonal
for the polarization, and the addition map

\[
 \phi:\prod_O A_O\longrightarrow J
 \tag{F12.7}
\]

is an isogeny with

\[
 \phi^*\Theta=\boxplus_O\Theta_O,\qquad
 \Theta_O=\Theta|_{A_O}.
 \tag{F12.8}
\]

There is also an integral reverse map \(J\to\prod_O A_O\), with
coordinates \(144e_O\); both composites are multiplication by 144. This
already gives a degree bound, but the polarizations give the sharper
bound used here.

For a fixed factor \(A=A_O\), its polarization complement \(A^\perp\) is
the connected kernel of the surjection \(J\to\widehat A\) obtained from
the principal polarization and dual inclusion. Hence

\[
 \ker\lambda_{\Theta_O}=A\cap A^\perp.
\]

The endomorphism \(144e_O\) is multiplication by 144 on \(A\) and zero
on \(A^\perp\). It follows that this finite polarization kernel is
killed by 144. Every elementary divisor in the polarization type
therefore divides 144. If \(d=\dim A_O\),

\[
 \chi(\Theta_O)\le144^d.
 \tag{F12.9}
\]

Since \(\Theta\) is principal, the top-intersection formula for a finite
pullback and the product formula in (F12.8) give

\[
 \deg\phi=\chi(\phi^*\Theta)
       =\prod_O\chi(\Theta_O)\le144^{55}.
 \tag{F12.10}
\]

Combining (F12.6) and (F12.7) proves (F12.1).

\subsection{SF-F1203. A uniform comparison with the product
polarization}\label{SF-F1203}

The one-dimensional factors have degree between one and 144, so their
polarizations already lie between the product principal polarization and
144 times that polarization.

Consider a two-dimensional factor \(E_K^2\), with \(K=\mathbf Q(i)\) or
\(\mathbf Q(\omega)\). In an order basis, its polarization is
represented by a positive definite integral Hermitian matrix

\[
 H=\begin{pmatrix}a&b\\\bar b&c\end{pmatrix},
 \qquad a,c\in\mathbf Z,\quad b\in\mathcal O_K,\quad
 \mathcal D=ac-|b|^2=\chi(\Theta_O)\le144^2.
 \tag{F12.11}
\]

There is no extra factor in the last equality: relative to the product
principal polarization, the polarization homomorphism is the Hermitian
endomorphism \(H\). Its degree is
\(\operatorname{Nm}_{K/\mathbf Q}\det H=\mathcal D^2\), while the degree
of a polarization homomorphism is its Euler characteristic squared.

Choose a nonzero vector of smallest \(H\)-norm, denoting that norm by
\(a\). It is primitive over \(\mathcal O_K\), since division by a
nonunit common factor would strictly decrease its positive norm. Extend
it to an order basis. Replace the second vector by an order multiple of
the first so that

\[
 |b|^2\le\rho^2a^2,\qquad
 \rho^2=\frac12\ \text{or}\ \frac13.
\]

Minimality gives \(c\ge a\), and consequently

\[
 \mathcal D\ge(1-\rho^2)a^2\ge\frac{a^2}{2}.
\]

Since \(a\) is a positive integer,

\[
 a\le\sqrt{2\mathcal D}\le2\mathcal D,\qquad
 c=\frac{\mathcal D+|b|^2}{a}
    \le\mathcal D+\frac a2\le2\mathcal D.
 \tag{F12.12}
\]

Thus \(\operatorname{tr}H\le4\mathcal D\le4\cdot144^2\). Also

\[
 \operatorname{tr}H^{-1}
   =\frac{\operatorname{tr}H}{\mathcal D}\le4.
\]

A positive Hermitian matrix has maximum eigenvalue at most its trace. It
follows that

\[
 \frac14I\le H\le4\cdot144^2I.
 \tag{F12.13}
\]

Taking the direct sum over all thirty-four factors proves (F12.2). The
inequalities mean that the differences are represented by semipositive
translation-invariant forms; equivalently they are nef rational divisor
classes. Every coordinate elliptic curve in this choice of coordinates
has \(\Theta\)-degree at most \(2\cdot144^2\).

\subsection{SF-F1204. Finite effective generators on every
power}\label{SF-F1204}

Set \(B_s=E_i^{27s}\times E_\omega^{28s}\), let \(\Theta_{0,s}\) be its
product principal polarization, and let

\[
 \Phi_s=\phi^s:B_s\longrightarrow J^s,\qquad
 L_s=\Phi_s^*\Theta_s.
\]

Then

\[
 \Theta_{0,s}\le4L_s,\qquad
 \chi(L_s)=(\deg\phi)^s\le144^{55s}.
 \tag{F12.14}
\]

For each of the two CM blocks, take the following effective divisors:

\begin{enumerate}
\def\labelenumi{\arabic{enumi}.}
\item
  the coordinate kernels \(x_j=0\);
\item
  for every \(j<k\), the kernels \(x_j-x_k=0\);
\item
  for every \(j<k\), the kernels \(x_j-\tau_Kx_k=0\), with \(\tau_i=i\)
  and \(\tau_\omega=\omega\).
\end{enumerate}

There are

\[
 (27s)^2+(28s)^2=1513s^2
 \tag{F12.15}
\]

such divisors. Their Hermitian matrices are the coordinate rank-one
matrices and the matrices \(vv^\dagger\) with \(v\) supported on two
coordinates of unit absolute value. Thus every generator \(D_j\)
satisfies

\[
 0\le D_j\le2\Theta_{0,s}\le8L_s.
 \tag{F12.16}
\]

Their classes span the rational divisor space: the coordinate classes
supply the diagonal Hermitian entries, while the two indicated
difference kernels supply the two independent rational off-diagonal
entries in \(K=\mathbf Q+\mathbf Q\tau_K\).

Here is a direct proof that these divisors generate the entire rational
Hodge ring. Over \(F_0=\mathbf Q(i,\omega)\), decompose \(H^1(B_s)\)
according to the two embeddings of each imaginary quadratic field, and
take the resulting exterior multidegrees. A monomial has separate
imbalances \(d_i,d_\omega\), the numbers of holomorphic factors minus
antiholomorphic factors in the two blocks. Being of Hodge type \((k,k)\)
forces \(d_i+d_\omega=0\). For a rational Hodge class, all Galois
conjugate multidegree components are also Hodge components: their
projections are tensor and exterior constructions from the field-action
projectors on \(H^1\), so preserve Hodge types, and Galois conjugation
replaces each projector by its conjugate. The two quadratic fields are
distinct and linearly disjoint, so an automorphism conjugating just
\(\mathbf Q(i)\) gives \(-d_i+d_\omega=0\). Thus both imbalances are
zero.

Every balanced monomial is a product of degree-two pairings within the
individual CM blocks. Those degree-two pairings are spanned by the
divisor classes just listed. Descending the resulting spanning statement
from \(F_0\) to \(\mathbf Q\) proves the claim.

For each \(k\)-element subset of the list, form the simultaneous kernel
of its \(k\) elliptic homomorphisms. If those homomorphisms are linearly
independent in their respective CM blocks, their combined map to a
product of \(k\) elliptic curves is surjective. Its zero fibre is a
smooth, possibly disconnected, abelian kernel of codimension \(k\), and
is an effective integral cycle \(Y\) with class \(\prod D_j\). If the
maps are dependent, that product class is zero: their rank-one forms
have dependent exterior support. Discard that subset. Repeated
generators likewise give zero, so the remaining subsets span all
monomials of degree \(k\).

The finite set is explicit: test the corresponding matrices over
\(\mathbf Q(i)\) and \(\mathbf Q(\omega)\), and retain the independent
subsets. There are at most \(\binom{1513s^2}{k}\) cycles. By the
semipositive inequalities (F12.16),

\[
 \deg_{L_s}Y
 =\int_{B_s}D_{j_1}\cdots D_{j_k}L_s^{g_s-k}
 \le8^k\int_{B_s}L_s^{g_s}
 \le8^k(55s)!\,144^{55s}.
 \tag{F12.17}
\]

The pushforward \(Z=(\Phi_s)_*Y\) is effective with integral
multiplicities, and the projection formula identifies its
\(\Theta_s\)-degree with the left side. Pushforward by an isogeny is an
isomorphism on rational cohomology up to a nonzero scalar inverse to
pullback. It therefore transports the spanning statement to \(J^s\),
proving (F12.3).

If a single irreducible detector is wanted, a nonzero pairing with such
a \(Z\) implies a nonzero pairing with at least one of its irreducible
components. That component has degree at most the displayed bound. No
coefficient or saturation estimate is hidden in this last selection.

\subsection{Published source, new proof, and actual
verification}\label{published-source-new-proof-and-actual-verification}

The qualitative orbit and CM-type description is classical. A published
primary locator is Koblitz--Rohrlich, \emph{Simple Factors in the
Jacobian of a Fermat Curve}, Canadian Journal of
Mathematics30(6),1183--1205 (1978), introductory lattice decomposition
and CM-type stabilizer discussion on printed pp1183--1184, DOI
https://doi.org/10.4153/CJM-1978-099-6 . The official publisher PDF was
opened and those exact pages read. Its later theorems for prime-to-six
levels are not invoked at level twelve.

The integral projector, exponent144, degree144\^{}55, binary Hermitian
comparison, explicit finite cycle family and numerical bound above are
proved here. The exact script checks every orbit, its chosen integral
scalar action, all projector coefficients, their Fourier actions on
all144 characters, and the stated census. It does not claim to have
explicitly enumerated kernels for arbitrarily large powers or to
saturate the integral Hodge lattice.

\section{Quantitative Shioda--Katsura suppliers for all products of
Fermat-twelve
hypersurfaces}\label{sec:sf-fermat12-hypersurfaces}

Author/checker: GPT-6 Astra Pro; inherited effort not exposed;
HODGE\_DIVISION /root/sixfolds. Date: 2026-11-21. New full proof; root whole-proof independent audit PASS on2026-11-21,
including an actual rerun of all423 finite degree cases. Reviewer:
GPT-6 Astra Pro; inherited effort not exposed; HODGE\_DIVISION /root. This extends the independently checked
SF-F1201--SF-F1204.

Write \(F^n=F_{12}^n\subset\mathbf P^{n+1}\), with \(\sum x_i^{12}=0\),
and give every product its Segre hyperplane class, the sum of the
hyperplane classes of its factors. A zero-dimensional Fermat factor is
treated componentwise; it consists of twelve reduced points.

For \(T\ge1\), define the following explicit integers: \[
 \begin{aligned}
 P_T&=12^T T!,\\
 A_T&=2^{2T}P_T^2,\\
 M_T&=2(2T+1)\,3^T A_T^{\,2T+2},\\
 U_T&=440^T T!\,144^{55T},\\
 \mathcal H_{12}(T)&=
 U_T\prod_{j=2}^{T}M_j^{\,j-1}.
 \end{aligned}
 \tag{SKQ.1}
\] Put \(\mathcal H_{12}(0)=1\); an empty product is one.

\textbf{Theorem.} For every product \(X=\prod_\nu F^{d_\nu}\), with
\(d_\nu\ge1\) and total dimension \(T=\sum d_\nu\), and every
codimension, the rational Hodge classes of \(X\) are spanned by a finite
set of effective integral algebraic cycles, each having Segre degree at
most \(\mathcal H_{12}(T)\).

This is a numerical bound for actual effective cycles, not a bound for
formal Hodge tensors alone. The assertion concerns rational spanning. No
saturation of the integral Hodge lattice is claimed.

\subsection{SF-SKQ01. A bounded effective supplier on every product of
curves}\label{SF-SKQ01}

Let \(C=F^1\), with plane hyperplane class \(h\), and let
\(a:C\hookrightarrow J(C)\) be an Abel embedding. Under the Abel
identification on \(H^1\), the canonical theta class is
\(\sum_{j=1}^{55}\alpha_j\beta_j\) for a symplectic integral basis with
\(\int_C\alpha_j\beta_j=1\). Its pullback therefore has degree55, while
the plane hyperplane class has degree12: \[
 a^*\Theta=55[\mathrm{pt}],\qquad h=12[\mathrm{pt}]
 \quad\hbox{in }H^2(C,\mathbf Q).
\] For \(a_T:C^T\to J(C)^T\), let \(H_C=\sum\operatorname{pr}_j^*h\).
Then \[
 a_T^*\Theta_T=\frac{55}{12}H_C,\qquad
 \int_{C^T}H_C^T=T!\,12^T.
 \tag{SKQ.2}
\]

The pullback \(a_T^*:H^*(J(C)^T,\mathbf Q)\to H^*(C^T,\mathbf Q)\) is
surjective as a graded Hodge structure and admits a rational Hodge
section. On a single factor, use the inverse isomorphism in degree one,
the identity in degree zero, and the map
\([\mathrm{pt}]\mapsto\Theta/55\) in degree two. Tensor these sections
over all factors. Consequently every rational Hodge class of \(C^T\) has
a rational Hodge lift to \(J(C)^T\).

Use the effective cycle suppliers \(Z=(\Phi_T)_*Y\) of SF-F1204 in
codimension \(k\le T\). The cycles \(Y\) are simultaneous kernels of
elliptic homomorphisms on \(E_i^{27T}\times E_\omega^{28T}\). They are
finite unions of abelian subvarieties translated by torsion points.
Their finite-isogeny images \(Z\) have the same form, with positive
integral multiplicities.

A general translation of every component of \(Z\) meets \(a_T(C^T)\)
properly or not at all. Here no general moving bound is needed. For an
abelian subvariety \(A\) of codimension \(k\), compose \(a_T\) with the
quotient \(J(C)^T\to J(C)^T/A\). A general fibre either is empty, when
the map is not dominant, or has dimension \(T-k\), by generic
smoothness. The same choice works for the finitely many components. Thus
the pullback of a sufficiently general translate of \(Z\) is an
effective integral cycle \(W\) of codimension \(k\), possibly zero,
representing \(a_T^*[Z]\).

There is also a degree bound substantially smaller than applying a
general intersection estimate in the Jacobian. On the elliptic source of
\(\Phi_T\), each generating divisor satisfies \(D_j\le8L_T\), as
semipositive translation-invariant forms. Expanding one factor at a time
gives \[
 D_{j_1}\cdots D_{j_k}\le8^kL_T^k
\] in the cone of strongly positive \((k,k)\)-forms. Pushforward by the
finite isogeny preserves this inequality, and \[
 (\Phi_T)_*L_T^k=(\deg\Phi_T)\Theta_T^k,\qquad
 \deg\Phi_T\le144^{55T}.
\] Pulling back to the curve product and integrating against the
positive form \(H_C^{T-k}\) therefore yields \[
 \begin{aligned}
 \deg_{H_C}W
 &\le8^k144^{55T}
       \int_{C^T}(a_T^*\Theta_T)^kH_C^{T-k}\\
 &=8^k55^k12^{T-k}T!\,144^{55T}\\
 &\le440^T T!\,144^{55T}=U_T.
 \end{aligned}
 \tag{SKQ.3}
\] Translations do not alter these cohomology classes. Since the
original Jacobian classes span and \(a_T^*\) has a Hodge section, the
resulting effective curve-product cycles span every rational Hodge
class.

The degree-zero case uses the whole curve product; the top-codimension
case is allowed. Zero pullbacks are discarded. This proves the base of
the induction in all codimensions simultaneously.

\subsection{SF-SKQ02. The two explicit correspondences in one induction
step}\label{SF-SKQ02}

Suppose \[
 X=F^n\times R,\qquad n\ge2,\qquad \dim X=T,
\] where \(R\) is a product of positive-dimensional Fermat
hypersurfaces. Put \[
 Z=F^{n-1}\times C\times R,\qquad
 B=F^{n-2}\times F^0\times R.
\] Here \(B\subset Z\) is cut out by the last homogeneous coordinates
\(x_n=0\), \(y_2=0\), so it is smooth of codimension two, possibly
disconnected.

Fix \(\epsilon^{12}=-1\). The rational map on the first two factors is
\[
 ((x_0:\cdots:x_n),(y_0:y_1:y_2))
 \longmapsto
 (x_0y_2:\cdots:x_{n-1}y_2:
       \epsilon x_ny_0:\epsilon x_ny_1).
 \tag{SKQ.4}
\] Its coordinates have bidegree \((1,1)\), its base ideal is the ideal
of \(B\), and its generic degree is12. For the last assertion, on a
target point where the sum of the first \(n\) twelfth powers is nonzero,
normalizing \(y_2=1\) reconstructs \(x_n\) from a nonzero prescribed
twelfth power, giving twelve distinct preimages.

Let \(\pi:\widetilde Z=\operatorname{Bl}_B Z\to Z\). Locally the base
ideal is generated by the two regular parameters \(x_n,y_2\); the blowup
therefore resolves (SKQ.4), giving a morphism \(f:\widetilde Z\to X\),
of generic degree12, and the identity on \(R\). This is the concrete
Shioda--Katsura construction: Lemmas1.1--1.2, pp98--99, and Theorem1.7,
pp101--102, of their published1979 paper.

Write \(E\) for the exceptional divisor,
\(j:E\hookrightarrow\widetilde Z\), and \(p:E\to B\). For a
codimension-two smooth blowup, the rational Hodge decomposition is \[
 H^m(\widetilde Z,\mathbf Q)
 =
 \pi^*H^m(Z,\mathbf Q)
 \oplus j_*p^*H^{m-2}(B,\mathbf Q)(-1).
 \tag{SKQ.5}
\] The maps displayed here are the actual algebraic pullback and Gysin
maps. The inverse components are \(\pi_*\) and \(-p_*j^*\): one uses
\(j^*j_*\eta=c_1(\mathcal O_E(-1))\eta\) and the projective-line bundle
formula. The blowup formula follows by excision and the cohomology of
the exceptional projective bundle; the same geometric formula is the
input of Shioda--Katsura Lemma2.1 and Proposition2.4, pp103--104. Over
\(\mathbf C\) these algebraic maps respect the stated Hodge structures.

For every rational Hodge class \(\alpha\) on \(X\), \[
 12\alpha=f_*f^*\alpha.
\] Decompose \(f^*\alpha\) by (SKQ.5). Thus the entire Hodge space of
\(X\) is spanned by the images of Hodge classes on \(Z\) and on the
connected components \(B_\lambda\) under the correspondences \[
 \Gamma_0=(\pi,f)_*[\widetilde Z]\subset Z\times X,\qquad
 \Gamma_\lambda=(p,fj)_*[E|_{B_\lambda}]
                  \subset B_\lambda\times X.
 \tag{SKQ.6}
\] The first acts by \(f_*\pi^*\); the second acts by \(f_*j_*p^*\),
increasing codimension by one. This includes every exceptional and Tate
contribution. No claim about a primitive summand is substituted for the
full decomposition.

\subsection{SF-SKQ03. Degrees of the main and exceptional
correspondences}\label{SF-SKQ03}

Every product of positive-dimensional degree12 hypersurfaces of total
dimension \(T\) has Segre degree at most \[
 P_T=12^T T!.
 \tag{SKQ.7}
\] Indeed its exact degree is \(T!\prod_\nu12/d_\nu!\), and there are at
most \(T\) factors. The same bound applies to \(X\) and \(Z\).

Let \(H_Z,H_X\) be their Segre hyperplane classes. The resolved linear
system gives \[
 f^*H_X=\pi^*H_Z-E.
\] Hence the hyperplane class from the product embedding \(Z\times X\),
pulled back to \(\widetilde Z\), is \[
 G=2\pi^*H_Z-E.
\] Both \(G\) and \(2\pi^*H_Z\) are nef. Their difference is the
effective divisor \(E\). Expanding their difference of \(T\)-th powers
shows \[
 \int_{\widetilde Z}G^T\le
 2^T\int_{\widetilde Z}(\pi^*H_Z)^T.
\] Each term in the difference is an effective divisor intersected with
nef classes, hence has nonnegative degree. Since \((\pi,f)\) has degree
one onto its graph, \[
 \deg\Gamma_0\le2^T\deg Z\le2^TP_T.
 \tag{SKQ.8}
\]

For completeness the exceptional degree is controlled by its actual
projective bundle, without choosing a resolution of an unspecified
graph. On a connected component \(B_\lambda\), the \(F^0\) coordinate
\(y\) is a fixed point. If \(n=2\), the other \(F^0\) coordinate is
fixed as well. Write \(h\) for the remaining \(F^{n-2}\) hyperplane
class, or zero when that factor is a point, and \(H_R\) for the
hyperplane class of \(R\). Then \[
 H_{B_\lambda}=h+H_R.
\] The exceptional fibres map to the lines joining the point represented
by \(x\) and the fixed point represented by \(y\). Thus their incidence
variety is explicitly \[
 E_\lambda=
 \mathbf P_{\rm lines}\bigl(\mathcal O_{B_\lambda}(-h)
                              \oplus\mathcal O_{B_\lambda}\bigr).
\] Let \(\xi=c_1(\mathcal O_{E_\lambda}(1))\). The join map has
hyperplane class \(\xi\); in particular \(\xi\) is nef. Including the
identity map on \(R\), the pullback of the product hyperplane class from
\(B_\lambda\times X\) is \[
 \xi+h+2H_R\le\xi+2H_{B_\lambda}.
\] The projective bundle relation and pushforward formulas are \[
 \xi(\xi-h)=0,\qquad
 p_*\xi^j=h^{j-1}\quad(j\ge1),\qquad p_*1=0.
\] Therefore \[
 \begin{aligned}
 \deg\Gamma_\lambda
 &\le\int_{E_\lambda}(\xi+2H_{B_\lambda})^{T-1}\\
 &=\sum_{j=1}^{T-1}
       \binom{T-1}{j}2^{T-1-j}
       \int_{B_\lambda}h^{j-1}H_{B_\lambda}^{T-1-j}\\
 &\le(3^{T-1}-2^{T-1})\deg B_\lambda\\
 &\le3^TP_T.
 \end{aligned}
 \tag{SKQ.9}
\] The third line uses \(0\le h\le H_{B_\lambda}\), with both classes
nef. The component \(B_\lambda\) is a product of total dimension
\(T-2\), with any zero-dimensional factors replaced by points, so its
degree is at most \(P_{T-2}\le P_T\). The correspondence map has degree
one onto its incidence image; even the weaker inequality for its
pushforward would suffice.

Thus every correspondence in (SKQ.6) has degree at most \(3^TP_T\) in
the relevant product embedding.

\subsection{SF-SKQ04. The quantitative correspondence-action
bound}\label{SF-SKQ04}

The published intersection estimate used here is Truong, \emph{Relative
dynamical degrees of correspondences over a field of arbitrary
characteristic}, Crelle758(2020),139--182, Lemma2.2 in the author
version, p13. Its full proof uses Roberts' moving construction and
bounds an intersection of effective cycles by a difference of two
effective cycles, each of degree at most \[
 C(V)\,\deg A\,\deg B.
\] The proof gives \(C(V)=m\deg(V)^m\), for \(m=\dim V\). To avoid any
convention about collecting positive and negative terms, use throughout
the larger integer \[
 C_{\rm safe}(V)=2(m+1)\deg(V)^{m+1}.
 \tag{SKQ.10}
\] The estimate extends from irreducible to effective cycles by
bilinearity and summing positive multiplicities. Rational equivalence is
enough here, and therefore the resulting equality also holds in rational
cohomology.

Let \(Y\) be either \(Z\) or a component \(B_\lambda\), and let
\(\Gamma\) be the corresponding cycle in \(Y\times X\). The ambient
product \(V=Y\times X\) is smooth and projective, has dimension at most
\(2T\), and has degree at most \[
 \deg V
 =\binom{\dim Y+T}{\dim Y}\deg Y\deg X
 \le2^{2T}P_T^2=A_T.
\] For an effective input cycle \(W\subset Y\), \[
 \deg(W\times X)\le2^{2T}P_T\deg W.
\] Apply (SKQ.10) to \(\Gamma\cdot(W\times X)\). Since
\(\deg\Gamma\le3^TP_T\), it can be represented as \(Q_+-Q_-\), with both
effective and \[
 \begin{aligned}
 \deg Q_\pm
 &\le
 2(2T+1)A_T^{2T+1}
      (3^TP_T)(2^{2T}P_T\deg W)\\
 &=2(2T+1)3^TA_T^{2T+2}\deg W\\
 &=M_T\deg W.
 \end{aligned}
 \tag{SKQ.11}
\] Push these cycles to \(X\). A component whose image has smaller
dimension pushes to zero. On every other component the projection
formula and the nef inequality \(p_X^*H_X\le H_V\) show that degree does
not increase. Consequently \[
 \Gamma_*[W]=[W_+]-[W_-],\qquad
 W_\pm\ \text{effective integral},\qquad
 \deg W_\pm\le M_T\deg W.
 \tag{SKQ.12}
\] In particular, if bounded-degree effective input classes span the
source Hodge space, then effective outputs of degree at most \(M_T\)
times that bound span the correspondence image. It is harmless that both
signs are retained as separate generators.

No proper-intersection assumption is imposed on the original input and
correspondence. That is exactly what the cited moving lemma supplies,
with the displayed numerical cost.

\subsection{SF-SKQ05. Induction, an explicit recurrence, and the final
bound}\label{SF-SKQ05}

For a product of total dimension \(T\), define its complexity \[
 \mu=\sum_\nu(d_\nu-1)=T-\#\{\hbox{positive-dimensional factors}\}.
\] Then \(0\le\mu\le T-1\). Complexity zero means \(C^T\), already
handled in SF-SKQ01.

Replacing one factor \(F^n\), \(n\ge2\), by \(F^{n-1}\times C\)
preserves \(T\) and decreases \(\mu\) by one. Every exceptional
component has total dimension \(T-2\). Thus (SKQ.6) and (SKQ.12) give
the valid simultaneous recurrence \[
 B(T,0)=U_T,\qquad
 B(T,\mu)\le
 M_T\max\{B(T,\mu-1),\,\mathcal H_{12}(T-2)\},
 \tag{SKQ.13}
\] with the point value \(\mathcal H_{12}(0)=1\). All codimensions are
included at once. Terms whose codimension lies outside the source range
vanish and can simply be omitted.

Iterating at most \(T-1\) times yields \[
 B(T,\mu)\le
 M_T^{T-1}\max\{U_T,\mathcal H_{12}(T-2)\}.
 \tag{SKQ.14}
\] The explicit quantity in (SKQ.1) dominates this expression. Indeed
\(U_T\) increases with \(T\), every \(M_j\ge1\), and \[
 \frac{\mathcal H_{12}(T)}
      {M_T^{T-1}\mathcal H_{12}(T-2)}
 =
 \frac{U_T}{U_{T-2}}\,M_{T-1}^{T-2}\ge1
 \quad(T\ge3);
\] the case \(T=2\) follows directly. This proves the theorem by
induction on total dimension and then complexity.

Only finitely many effective generators are introduced at each stage:
the previous finite generating lists are passed through two
correspondence types, with finitely many exceptional components, and
each intersection is replaced by two effective cycles. The construction
never invokes an uncontrolled resolution, an unbounded regularity
estimate, or an unspecified isogeny degree. Every factor in (SKQ.1) has
been accounted for.

For a nonzero pairing functional on a Hodge space, at least one
generator pairs nontrivially. One may then choose one irreducible
component with nonzero pairing, still of degree at most
\(\mathcal H_{12}(T)\). This is the intended use as a numerical supplier
in the subsequent cyclic-cover and Brauer-correspondence transfer.

\subsection{Published-source and verification
record}\label{published-source-and-verification-record}

Shioda--Katsura is published in Tohoku Mathematical
Journal31(1979),97--115, DOI10.2748/tmj/1178229881. The official J-STAGE
PDF was opened; the rational map and blowup construction on pp98--102
and cohomological formulas on pp103--104 were inspected. The degree
bounds, projective-bundle calculation and full-product recurrence above
are new.

Truong's paper is published in Crelle758(2020),139--182,
DOI10.1515/crelle-2017-0052. Publication was checked at the publisher.
The complete relevant quantitative proof was read in the author copy
arXiv:1605.05049, Lemma2.2, p13. No unverified printed lemma-page
locator is asserted. The larger safe constant (SKQ.10) is used, so the
recurrence does not depend on a sharp collection-of-signs convention.

\section{An explicit lattice comparison for Gaussian determinant three}
\label{sec:gaussian-nonsplit-lattice}

Put $\mathcal O=\mathbb Z[i]$, $K=\mathbb Q(i)$, and
\[
 H_0=\begin{pmatrix}0&I_2&0\\I_2&0&0\\0&0&J_3\end{pmatrix},
             \qquad J_3=\operatorname{diag}(1,-3).
\]
We use conjugate-linear first arguments.  A Hermitian form of
signature $(3,3)$ has signed determinant $-\det h$ modulo
$\operatorname{Nm}_{K/\mathbb Q}(K^*)$.

\subsection*{D-E8-LAT-01. Two-sided denominators bounded by the polarization volume}

Let $\Lambda$ be a free rank-six $\mathcal O$-lattice with an
integral nondegenerate Hermitian form $h$ of signature $(3,3)$
and signed determinant class three.  Write
$\Delta=|\det h|\in\mathbb Z_{>0}$.  There are an
$\mathcal O$-basis of $\Lambda$, a matrix $T\in\operatorname{GL}_6(K)$,
and an integer $q>0$ such that
\[
 T^*hT=H_0,\qquad qT,\ qT^{-1}\in M_6(\mathcal O),\qquad
                        q\le69984\Delta^3.                  \tag{LAT.1}
\]
Both inclusion directions are part of the assertion.  The number
$69984\Delta^3$ is an upper bound for a clearing integer; it is
not claimed to be itself divisible by every possible clearing
integer.

\begin{proof}
Landherr's published classification identifies the rational
Hermitian space with $H_0$: see \cite[p.~246 for the invariants,
first paragraph of p.~247 for their completeness]{E8LAT:Landherr}.
Only this existence statement is used; no denominator bound is
deduced from it.  In particular the rational Witt index is two.
The binary plane $J_3$ is anisotropic since three is not a
Gaussian norm.

\emph{Step1: split the two hyperbolic planes with controlled
denominators.}  Intersect a rational maximal isotropic two-plane
with $\Lambda$.  The intersection is primitive.  Since
$\mathcal O$ is a PID, choose its basis $v=(v_1,v_2)$ and extend
it to a basis of $\Lambda$.  The pairing with the other four
basis vectors has rank two.  Smith normal form over $\mathcal O$
allows further integral basis changes so that, in blocks
$(v,w,z)$ of size two, the Gram matrix is
\[
 \begin{pmatrix}
 0&B&0\\ B^*&F&E\\0&E^*&A
 \end{pmatrix},\qquad
             B=\operatorname{diag}(b_1,b_2),\quad b_1b_2\ne0.
                                                                  \tag{LAT.2}
\]
All its entries are integral, and $F,A$ are Hermitian.  Put
$b=\det B$, $n_b=b\overline b$, and $d=|\det A|$.  Taking
determinants of (LAT.2) gives
\[
                   \Delta=n_b d.                            \tag{LAT.3}
\]
The following change of rational basis is explicit:
\[
 \begin{split}
 u&=\left(w-v(B^*)^{-1}F/2\right)B^{-1},\\
 z'&=z-v(B^*)^{-1}E.
 \end{split}                                                   \tag{LAT.4}
\]
It gives Gram matrix $\begin{psmallmatrix}0&I_2&0\\I_2&0&0\\0&0&A\end{psmallmatrix}$.
Indeed it kills the $w$-self-pairing and its pairing with $z'$,
and does not alter the $z'$-self-pairing.  The inverse change is
\[
             w=uB+v(B^*)^{-1}F/2,
             \qquad z=z'+v(B^*)^{-1}E.                        \tag{LAT.5}
\]
The denominators of both (LAT.4) and (LAT.5) divide $2n_b$:
the only products of inverse determinants are
$b^{-1}\overline b^{-1}=n_b^{-1}$, and either individual inverse
is also cleared by $n_b$.  Thus this step has a two-sided
clearing integer $2n_b$, not an uncontrolled basis-height
factor.  The residual integral binary form $A$ has signature
$(1,1)$ and signed determinant class three, hence is anisotropic.

\emph{Step2: reduce the binary coefficients using the Gaussian
Euclidean algorithm.}  Choose a nonzero integral vector on
which $|A(x,x)|$ is the smallest possible positive integer.
Anisotropy ensures these values are nonzero.  This vector is
primitive: division by a nonunit common divisor of its two
coordinates would reduce its value by that divisor's norm.
Extend it to an integral basis.  On replacing $A$ temporarily
by $-A$ if necessary, its matrix is
\[
                    A=\begin{pmatrix}a&z\\\overline z&c\end{pmatrix},
                       \qquad a>0.
\]
Shear the second vector by a Gaussian integral multiple of the
first, rounding both coordinates of $z/a$ to nearest integers.
This gives
\[
       |z|^2\le a^2/2,\qquad |c|\ge a.                       \tag{LAT.6}
\]
The second inequality is minimality; $c=0$ is excluded by
anisotropy.  Since $ac-|z|^2=-d<0$, the case $c>0$ would give
$|z|^2=ac+d>a^2$, contradicting (LAT.6).  Therefore $c<0$ and
\[
 d=a|c|+|z|^2,\qquad a\le\sqrt d,\quad |c|\le d,
 \qquad |\Re z|+|\Im z|\le a\le\sqrt d.                      \tag{LAT.7}
\]
These integral basis changes and their inverses introduce no
new denominators.

\emph{Step3: a complete small-zero argument in five variables.}
Consider the integral quadratic form
\[
 Q(X_1,X_2,X_3,X_4,t)
       =A(X_1+iX_2,X_3+iX_4)-t^2.
\]
Write $Q(x)=x^t\mathcal Qx$ with an integral symmetric matrix,
and put $S=\sum_{j,k}|\mathcal Q_{jk}|$.  Formula (LAT.7) gives
\[
 S=2a+2|c|+4(|\Re z|+|\Im z|)+1
                    \le2d+6\sqrt d+1\le9d.                  \tag{LAT.8}
\]
The form $Q$ is isotropic: by Landherr's classification $A$ is
isometric to $J_3$, which represents one.  It has rational Witt
index exactly one.  If it had a totally isotropic two-plane,
that plane would meet the hyperplane $t=0$ nontrivially,
contradicting the anisotropy of $A$.

Here is the full elementary height argument, so no unverified
version of a small-zero constant is needed.  For an integral
isotropic quadratic form in five variables of Witt index one,
choose a nonzero integral zero $x$ with minimal sup-norm
$M=\max_j|x_j|$.  It is primitive.  Reorder coordinates and
change its sign so that $x_5=M$.  Set
\[
                       Q_0=\lfloor\sqrt{3S}\rfloor+1.
\]
Suppose $M>Q_0^4$.  Among the $Q_0^4+1$ points
\[
 \left(\{j x_1/M\},\ldots,\{j x_4/M\}\right),
                         \quad 0\le j\le Q_0^4,
\]
two lie in the same one of the $Q_0^4$ half-open boxes of
side $1/Q_0$.  Subtract their indices to obtain
$1\le k\le Q_0^4<M$, and subtract their integer parts to
obtain $y\in\mathbb Z^5$ satisfying
\[
 y_5=k,\qquad y=(k/M)x+e,\qquad e_5=0,
                         |e_j|<1/Q_0.
\]
The vectors $x,y$ are not proportional: since $x$ is primitive,
an integral rational multiple of it has integral multiplier,
whereas $0<y_5/x_5=k/M<1$.
Let $\mathcal B$ denote the bilinear form with matrix
$\mathcal Q$.  The reflection numerator
\[
             z_0=Q(y)x-2\mathcal B(x,y)y
                 =Q(e)x-2\mathcal B(x,e)e                     \tag{LAT.9}
\]
is an integral zero of $Q$.  It is nonzero.  If it vanished,
either $y$ would be proportional to $x$, or
$Q(y)=\mathcal B(x,y)=0$; the latter would make their span
a totally isotropic two-plane.  Both alternatives have just
been excluded.  But
\[
 \|z_0\|_\infty
       \le\frac{3SM}{Q_0^2}<M,
\]
contradicting the choice of $x$.  It follows that
\[
 M\le Q_0^4\le(2\sqrt{3S})^4=144S^2.
\]
Applied to our $Q$, this gives a nonzero zero with
\[
              0<|t|\le11664d^2.                             \tag{LAT.10}
\]
The inequality $t\ne0$ follows from the anisotropy of $A$.
The selected maximal coordinate in the height argument need
not be this $t$ coordinate; (LAT.10) follows because every
coordinate is bounded by $M$.

\emph{Step4: an explicit binary isometry and its inverse.}
Put $p=(X_1+iX_2,X_3+iX_4)^t$ and $v=p/t$, so that
$v^*Av=1$.  Write $(\ell_1,\ell_2)=v^*A$ and
$u=(-\ell_2,\ell_1)^t$.  Direct multiplication gives
\[
          v^*Au=0,\qquad u^*Au=\det A=-d,
                    \det(v,u)=1.                             \tag{LAT.11}
\]
The coordinates of both $v$ and $u$ have denominators dividing
$|t|$.

Because $d/3$ is a rational Gaussian norm, the positive integer
$3d$ is a sum of two integral squares.  To justify the integral
assertion, factor $3d$ in $\mathbb Z$: at every prime
$p\equiv3\pmod4$ its valuation is even, as follows from it being
a rational norm.  The Gaussian Euclidean factorization then
constructs $\beta\in\mathcal O$ with
$\beta\overline\beta=3d$.  Set $\alpha=\beta/3$.  The matrix
\[
                  T_A=(v,u/\alpha)
\]
satisfies $T_A^*AT_A=J_3$.  Its entries are cleared by
$d|t|$, since $\alpha^{-1}=\overline\beta/d$.
Its determinant is $\alpha^{-1}$ by (LAT.11), and hence
$T_A^{-1}=\alpha\operatorname{adj}(T_A)$.  The entries in
the first row of this inverse have denominators dividing
$|t|$, and those in the second row have denominators dividing
$3|t|$.  In particular the single integer
\[
                            3d|t|                            \tag{LAT.12}
\]
clears both $T_A$ and $T_A^{-1}$.

If Step2 used $-A$, its construction instead gives Gram
$-J_3$ for the original $A$.  The integral determinant-one
matrix
\[
        U=\begin{pmatrix}1+i&-3\\1&-1+i\end{pmatrix}
\]
satisfies $U^*(-J_3)U=J_3$.  Multiplying by $U$ fixes the
sign without changing either denominator bound.

Combine the binary change with (LAT.4).  Both the total
change and its inverse are cleared by
\[
          q=(2n_b)(3d|t|)=6\Delta|t|
                 \le69984\Delta d^2\le69984\Delta^3.
\]
This proves (LAT.1), including its two inclusion directions.
\end{proof}

\subsection*{D-E8-LAT-02. Fixed-lattice isogeny and cycle-degree transport}

For two such lattices of volumes $\Delta_S,\Delta_T$ there is
a rational Hermitian isometry $P$ from the first to the second
and an integer $N$ such that
\[
 NP\Lambda_S\subset\Lambda_T,\quad
 NP^{-1}\Lambda_T\subset\Lambda_S,\qquad
          N\le69984^2(\Delta_S\Delta_T)^3.                    \tag{LAT.13}
\]
In particular, if the source lattice has volume at most $M^6$,
then
\[
                       N\le69984^2M^{18}\Delta_T^3.           \tag{LAT.14}
\]
For the $E_8$ source one may take conservatively
$M=2|G|=2786918400$, the primitive-projector denominator.

\begin{proof}
Choose $T_S,T_T$ and clearing integers $q_S,q_T$ from
D-E8-LAT-01, and take
$P=T_TT_S^{-1}$ and $N=q_Sq_T$.  Then
$P^*h_TP=h_S$.  Since both $q_TT_T$ and $q_ST_S^{-1}$
are integral, $NP$ is integral.  The other two integral
matrices give the same assertion for $NP^{-1}$.  This proves
(LAT.13); substituting $\Delta_S\le M^6$ gives (LAT.14).

For a compatible integral alternating polarization $E$ on an
$\mathcal O$-stable lattice use
\[
                 h(x,y)=E(x,iy)+iE(x,y).                     \tag{LAT.15}
\]
It is an integral Hermitian form, and $E$ is its coefficient
of $i$.  The determinant of the alternating matrix is
$(\det h)^2$, so $|\det h|$ is precisely the product of the
polarization elementary divisors.  Therefore a rational
Hermitian isometry is a rational symplectic isometry for the
specified polarizations, with no unspecified positive scalar.

For completeness, the stated source bound uses the actual
primitive projector, not just the group order as a heuristic.
Let $\chi$ be the degree-eight character of the natural
representation and put
\[
 \epsilon=\frac8{|G|}\sum_{g\in G}
       (\chi(g^{-1})+\overline{\chi(g^{-1})})g,\qquad
                   e=\epsilon(1-s)/2,
\]
where $s$ is a reflection.  The character trace in parentheses
is integral, so $(2|G|)e\in\mathbb Z[G]$.  This is a primitive
self-adjoint idempotent in the rational $M_8(K)$ factor.
If $L=H_1(J(C),\mathbb Z)$ is the unimodular Jacobian lattice,
the lattice of its image is $\Lambda_S=eL_{\mathbb Q}\cap L$.
Its polarization dual is $eL$: an integral functional on the
primitive sublattice $\Lambda_S$ extends to $L$, and
unimodularity followed by the self-adjoint projection gives
the representing vector in $eL$.  Therefore
$M\Lambda_S^\#\subset\Lambda_S$ for $M=2|G|$.
Each of its six polarization elementary divisors divides $M$,
proving $\Delta_S\le M^6$.  The central group element $iI$
preserves the image abelian variety and its integral lattice,
so this source lattice is indeed $\mathbb Z[i]$-stable.

Fix the two lattices and identify their period domains by $P$.
At corresponding period points, $NP$ is a $K$-linear isogeny
$\phi:B_S\to B_T$ and
\[
                \phi^*\theta_T=N^2\theta_S,
 \qquad \deg\phi=N^{12}\Delta_S/\Delta_T.                    \tag{LAT.16}
\]
The integer $N$ and the rational isometry can be fixed on a
marked polydisc, so this is a family comparison, not an
isogeny chosen with an uncontrolled denominator at each fiber.
For an effective codimension-three source cycle $Z$,
the projection formula gives
\[
                  \deg_{\theta_T}(\phi_*Z)
                      =N^6\deg_{\theta_S}(Z).                \tag{LAT.17}
\]
Thus the numerical upper bound in (LAT.14) may safely be used
to bound transported cycle degrees, even though that upper
bound is not itself claimed to clear every isometry.  Weil
projection is preserved by a rational $K$-isogeny, so a
nonzero detecting projection remains nonzero.
\end{proof}

For a target lattice stable only under the fixed order
$\mathbb Z+f\mathbb Z[i]$, replace it first by
$\Lambda'=f\mathcal O\Lambda\subset\Lambda$.  As proved in
D-G34-07b, $f\Lambda\subset\Lambda'\subset\Lambda$ and
$[\Lambda:\Lambda']\le f^{12}$ in dimension six.  Its pulled-back
polarization has volume at most $f^{12}\Delta_T$, so (LAT.14)
becomes
\[
                   N\le69984^2M^{18}f^{36}\Delta_T^3.         \tag{LAT.18}
\]
The subsequent identity isogeny to the original target pulls
back its polarization exactly, and adds no factor to (LAT.17).
The scope is a fixed conductor $f$, not a bound independent
of all orders.

\section{Bounded full Brauer identities and explicit contraction diagrams}
\label{sec:e8-transfer-diagrams}

\paragraph{Scope and provenance.}
This section proves the representation-theoretic and correspondence-count
steps needed for the quantitative branch-exponent construction.  The
integer column-lattice argument is reconstructed from the proof inside
the supplied SQ3.  The printed statement of SQ3 concerns determinant
characters; it is \emph{not} silently used as a bounded identity of
representations.  The latter identity is proved below.  The geometric
supplier theorem HD-FS1--02 and the separately bounded Fermat-product
Hodge cycles remain explicit inputs.  No arbitrary endomorphism of a
multiplicity Hodge structure is asserted to be algebraic.

Let \(\mathcal G\) be the connected-cover group, write \(G=|\mathcal G|\),
and let \(C\) be its smooth projective curve of genus \(g\geq1\).
Let \(F\subset\mathbb C\) be a finite Galois splitting field for
\(\mathcal G\) containing all \(G\)-th roots of unity, and put
\(A=F[\mathcal G]\).  This ensures that every subgroup's linear
characters used below are defined over \(F\); splitting
\(\mathcal G\) alone would not suffice for that assertion.
The coefficient representation \(V\) is absolutely irreducible.
In the application \(G=1393459200\), \(g=1+53G/24=3077222401\),
and the multiplicity space of \(V\) in \(H^1(C,F)\) has dimension six.

\subsection*{D-E8-XFER-01. A bounded identity of honest representations}

There is an integral Brauer identity
\[
 [V]=\sum_j n_j[\operatorname{Ind}_{H_j}^{\mathcal G}\chi_j],
 \qquad
 \sum_j|n_j|\leq (G+1)G!G^G
 \leq R:=(2G+1)G!G^G.                                  \tag{XD.1}
\]
Here each \(\chi_j\) is a linear character and trivial characters are
allowed.  Thus, on moving negative terms to the other side, (XD.1)
is an isomorphism
\[
 V\oplus R_-\simeq R_+                                   \tag{XD.2}
\]
of actual \(F[\mathcal G]\)-modules, with at most \(R\) induced
summands counted with their repetitions.

\begin{proof}
Use all irreducible representation multiplicities as rows of an
integer matrix.  For every pair \((H,\chi)\), take the multiplicity
vector of \(\operatorname{Ind}_H^{\mathcal G}\chi\) as a column.
There are finitely many columns, at most \(G\) rows, and every
entry lies between zero and \(G\): it is at most the dimension of
the induced representation, which is \([\mathcal G:H]\leq G\).
Integral Brauer induction puts the unit vector for \(V\) in the
integral column lattice.

For clarity, here is the needed integer bound.  Suppose an integer
column matrix has rank \(r\), entries bounded by \(M\geq1\), and a
target \(b\), with \(\|b\|_\infty\leq1\), in its integral column
lattice.  Choose \(r\) independent rows and \(r\) independent
columns forming an invertible square matrix \(B\).  Projection to
these rows is injective on the rational column span.  Set
\[
                    d=|\det B|\leq r!M^r.
\]
The subgroup generated by the remaining columns in
\(\mathbb Z^r/B\mathbb Z^r\) has at most \(d\) elements.  A shortest
path from zero to the class of \(b\), using these columns and
their negatives, has at most \(d-1\) edges.  Subtract the corresponding
signed columns.  The residual vector \(v\in B\mathbb Z^r\) satisfies
\[
                  \|v\|_\infty\leq1+(d-1)M\leq dM.
\]
Every cofactor is bounded by \((r-1)!M^{r-1}\).  Cramer's rule,
with its denominator \(d\), therefore gives
\[
              |(B^{-1}v)_j|\leq r!M^r.
\]
The total absolute coefficient sum is at most
\((r+1)r!M^r\).  The identity in the selected rows is an identity
in all rows by injectivity on the column span.  Taking \(r\leq G\)
and \(M=G\) proves (XD.1).  Semisimplicity in characteristic zero
turns equality of representation classes into (XD.2).
\end{proof}

Write \(U=H^1(C,F)\), let \(M_\pm=\operatorname{Hom}_{\mathcal G}(R_\pm,U)\),
and suppose \(\dim_F\operatorname{Hom}_{\mathcal G}(V,U)=6\).
Put \(k=\dim_F M_-\).  Each induced summand contributes a subspace
of \(U\), of dimension at most \(2g\).  Applying the contravariant
multiplicity functor to (XD.2) gives
\[
 M_+\simeq M_V\oplus M_-,\qquad
 \ell:=\dim_F M_++\dim_F M_-=6+2k\leq2gR=:L.                \tag{XD.3}
\]
Zero-dimensional summands can be omitted \emph{after} applying this
functor; no claim is made that deleting them preserves (XD.2) as a
representation identity.  In particular the actual curve-slot count
\(\ell\) is even and at least six.

\subsection*{D-E8-XFER-02. Every required map is a finite sum of deck diagrams}

Fix a very ample divisor \(H_C\) of degree \(6g\).  The homomorphism
of determinant lines obtained from (XD.3) and the polarization,
\[
             \det M_+\otimes\det\overline{M_-}
                  \longrightarrow\det M_V                 \tag{XD.4}
\]
is nonzero.  After embedding all its factors in tensor powers of
\(U\), it is a finite \(F\)-linear combination of the following maps:
apply deck transformations to individual inputs, pair \(k\) disjoint
pairs of inputs by the cup-product pairing, and send the six remaining
inputs, in some order, to six output factors.  Every such summand is
realized by an integral effective correspondence
\[
                     Z\subset C^\ell\times C^6
\]
which is, after permuting coordinates, a product of \(k+6\) deck
graphs.  Its degree for the sum of the curve polarizations is
\[
                 \deg Z=(k+6)!(12g)^{k+6}.                  \tag{XD.5}
\]
Precomposing with the faithful cyclic-character and determinant
projectors introduces only further sums of the same kind.

\begin{proof}
This is an assertion about maps induced by representations, not
about all Hodge maps on \(U\).  For a linear character define
\[
 e_{H,\chi}=\frac1{|H|}\sum_{h\in H}\chi(h)^{-1}h\in A.
\]
The left ideal \(Ae_{H,\chi}\) is the induced module
\(\operatorname{Ind}_H^{\mathcal G}\chi\), and
\[
 \operatorname{Hom}_A(Ae,U)\simeq eU,\qquad
                     f\longmapsto f(e).                    \tag{XD.6}
\]
Choose a primitive idempotent \(e_V\) with \(Ae_V\simeq V\).
In the \(E_8\) application choose
\(e_V=\epsilon_V(1-s_1)/2\), where \(\epsilon_V\) is the
individual central idempotent for the chosen embedding of the
natural character.  Then \(e_V+\overline{e_V}\) is precisely
the rational root-line projector \(p\) used in the last map.
If another primitive idempotent is first used, its fixed
module isomorphism with this one is also an entry of the
group algebra by the next argument.
The isomorphism (XD.2) is therefore an isomorphism between direct
sums of left ideals of the semisimple algebra \(A\).

Every \(A\)-linear map \(Ae\to Af\) is determined by the image of
\(e\), an element \(t\in eAf\); explicitly it sends \(ae\) to \(at\).
After applying (XD.6), its contravariant action is \(u\mapsto tu\).
Consequently the induced maps in (XD.3) are matrices whose entries
are elements of \(A\).  Each entry is a finite \(F\)-linear
combination of actual deck transformations of \(C\), and is
realized by the corresponding combination of their graphs.
This gives the required algebraic maps without a density assertion
about unknown multiplicity-Hodge endomorphisms.

The cup product is a nondegenerate alternating form \(Q\) on \(U\).
The transpose of \(e_{H,\chi}\) for this form is
\(e_{H,\chi^{-1}}\).  Thus it pairs \(e_{H,\chi}U\)
nondegenerately with \(e_{H,\chi^{-1}}U\), and the direct-sum
pairing identifies \(\overline{M_-}\) with \(M_-^*\), up to the
usual Tate twist.  Take the top exterior power of (XD.3) and then
contract these two determinant factors.  The isomorphism and the
nondegenerate determinant pairing prove that (XD.4) is nonzero.

More explicitly, write the two components of
\(M_+\to M_V\oplus M_-\) as \(a\) and \(b\).  For each choice of
six of the \(6+k\) plus inputs, send those six through \(a\);
send the others through \(b\) and pair them with the \(k\) minus
inputs.  Alternate over the choices, the six outputs and the
minus-input matching.  This is the usual exterior-power
isomorphism followed by the determinant pairing.  Expanding
each matrix entry in \(F[\mathcal G]\) gives precisely the
asserted maps.  The factors from direct-sum labels are just
scalar coefficients and choices of the corresponding entries.

For paired slots, group invariance gives
\(Q(gx,hy)=Q(x,g^{-1}hy)\); hence one deck graph on the two
input curves realizes each pairing.  A retained input and its
output use one deck graph as well.  The \(\ell+6=2(k+6)\)
coordinates are thereby partitioned into disjoint pairs.
The graph on each pair is isomorphic to \(C\) and has degree
\(2\deg H_C=12g\).  The product polarization gives (XD.5).
This also verifies the correspondence's degree shift:
its codimension is \((\ell+6)/2\), so it sends \(H^\ell(C^\ell)\)
to \(H^6(C^6)\).

The cyclic-character projectors act coordinatewise by group-algebra
elements.  Determinant alternation acts by permutations, with
the Koszul sign when it is written in cohomological tensors.
Precomposition by either operation merely changes the deck words,
the input pairing, the output ordering, or their scalar coefficients.
It remains a sum of the same diagrams.
\end{proof}

There is one further elementary projector needed when cycles have
components in other Kunneth degrees.  For a chosen point \(o\in C\),
\[
 \pi_1=[\Delta_C]-[\{o\}\times C]-[C\times\{o\}]
\]
projects onto \(H^1(C)\).  Every term of \(\pi_1^{\boxtimes\ell}\)
is a product of \(\ell\) diagonals or point-times-curve factors.
It is an integral effective cycle of dimension \(\ell\) and degree
at most
\[
                         \ell!(12g)^\ell.                   \tag{XD.7}
\]
Thus the full tensor Kunneth projection is one correspondence
operation, expanded into finitely many terms.  The number or sizes
of their scalar coefficients are not a degree cost for the
existence conclusion below.

\subsection*{D-E8-XFER-03. Trivial characters and bounded component selection}

A trivial inducing character contributes \(H^1(B,F)\) for
\(B=C/H\), with \(b=g(B)\).  Its determinant is represented,
after Kunneth projection and alternation, by a nonzero scalar
multiple of the product of \(b\) pullback diagonals
\[
                  T_H=(C\times_B C)^b\subset C^{2b}.
\]
As fundamental cycles \(C\times_B C=\sum_{h\in H}\Gamma_h\):
the equality holds over the free locus and both sides have
pure dimension one.  Thus
\[
                   \deg T_H=b!(12g|H|)^b
                              \leq b!(12gG)^b.              \tag{XD.8}
\]
Nonvanishing follows from the nondegenerate symplectic pairing
on \(H^1(B)\): its inverse-pairing tensor, repeated \(b\) times
and alternated, is its top symplectic volume.  Quotient pullback
only multiplies the pairing by the nonzero integer \(|H|\).

These trivial factors must be external factors of the returning
correspondence.  They are \emph{not} described as finite graphs over
a point.  If \(\ell_{\rm non}\) slots belong to the nontrivial
characters and \(\ell_{\rm triv}\) to trivial characters, the
combined correspondence is
\[
                 \Gamma_{\rm non}\times T_{\rm triv}
                         \subset Y\times C^\ell.            \tag{XD.9}
\]
If \(d=\dim Y\), its dimension is \(d+\ell_{\rm triv}/2\).
It sends a source class of degree \(2d-\ell_{\rm non}\) to a
class of degree \(\ell\), as required.  The factors with \(b=0\)
are units and are omitted.

For clarity, the following statement isolates the quantitative
selection step.  Suppose every ambient product used here is smooth
projective of dimension at most \(2L\), with very ample product
polarization of degree at most \(A\), and every required returning
correspondence (XD.9) has degree at most \(A\).
Suppose the relevant rational Hodge classes on each \(Y\) are spanned
by integral effective cycles of degree at most \(S\).
Take
\[
 A=(4L)!\,[1000g(G+1)(L+1)]^{4L},\qquad
                         B=(4L+4)A^{2L+4}.                 \tag{XD.10}
\]
The elementary degree estimates (XD.5), (XD.7) and (XD.8) are all
less than \(A\), since \(6\leq\ell\leq L\) and
\(k+6=(\ell+6)/2\leq\ell\).

The returning operation, the full tensor Kunneth projector, the
diagram operation of D-E8-XFER-02, and a final coordinatewise
primitive projector constitute at most four correspondence
operations.  All cyclic-character projectors, exterior-power
intertwiners and determinant alternations have already been
expanded into the third operation.  The final primitive projector
is redundant if included there, but keeping it costs at most one
additional operation.

Here is the degree cost for each operation, including improper
intersections.  If \(Z\) has degree \(E\), taking its product with
the target has degree at most \(AE\): the relevant binomial
coefficient times the target degree is bounded by the ambient
product degree.  Intersecting it with a correspondence of degree
at most \(A\), the effective positive-minus-negative version of
Chow moving in Truong's Lemma 2.2 has total mass at most
\[
        2(2L+1)A^{2L+1}(AE)A
             =2(2L+1)A^{2L+3}E\leq BE.                     \tag{XD.11}
\]
The displayed factor is deliberately larger than the published
constant \(kD^k\) for each sign.  Proper projection does not
increase degree for the sum polarization: the omitted source
hyperplane class is nef.  The same mass estimate therefore
holds after pushforward.

The nonzero determinant map (XD.4), the supplier theorem and the
nonzero pairing in its returning correspondence show that the
composite supplies the individual \(F\)-determinant line.
To use rational Hodge cycles, first expand all group-algebra and
permutation coefficients and let \(I\) be the sum of the images
of the resulting \emph{rational} correspondence maps from the
supplier cohomologies, with the required Tate shifts.
This is a rational Hodge substructure.  Its scalar extension
\(I_F\) contains the determinant line, hence also all its
\(\operatorname{Gal}(F/\mathbb Q)\)-conjugates, because \(I\)
is rational.  The conjugate pair spans the rational Weil
plane after scalar extension, so \(I\) contains that Weil
plane over \(\mathbb Q\).
Semisimplicity of polarizable rational Hodge structures now
gives a \emph{rational Hodge preimage} of a nonzero rational
Weil class in the direct sum of the supplier cohomologies.
The assumed rational spanning cycles express this preimage
as a finite linear combination.  Linearity implies that
some single initial spanning cycle and some choices of diagram
terms have nonzero final Weil projection.  At each moving
intersection write the result as integral positive minus negative
cycles using (XD.11).  Expanding these finite sums again, at least
one individual integral effective component has nonzero final
Weil projection.  Here the projection may be evaluated after the
fixed last map to \(B\) described below; it is not asserted that
an individual term still belongs to every intermediate projector
image.  Its degree on \(C^6\) is at most
\[
                              B^4S.                         \tag{XD.12}
\]
This argument remains valid after extending the coefficient field:
a nonzero linear combination over \(F\) cannot be formed from
vectors all having zero Weil projection.  No denominator clearing,
integral saturation, or bound on a scalar coefficient is used.

Finally let \(p\) be the Rosati-fixed primitive projector and
\(M=2G\) an integer clearing it.  The homomorphism \(Mp:J(C)\to B\),
composed with Abel maps and addition, gives \(C^6\to B\).
The pullback of the conjugate Weil determinant has, on the six
curve factors, the nonzero alternating tensor of the six-dimensional
conjugate character eigenspace.  Its cup pairing with the output
determinant in (XD.4) is nonzero by the perfect determinant pairing.
Projection formula therefore proves that this last pushforward
is nonzero on the determinant output.  Component selection above
can consequently be made using the resulting Weil projection on
\(B\).  Since
\[
 (Mp)^*\theta_B\leq M^2\Theta_{J(C)},\qquad
 \deg(\Theta_{J(C)}|_C)=g,
\]
the Cauchy inequality for the sum of six maps and
\(\deg H_C=6g\) give pullback polarization at most \(M^2H_{C^6}\).
A surviving three-dimensional cycle consequently has degree at most
\[
                         M^6B^4S
                       \leq M^6B^{4L+21}S.                 \tag{XD.13}
\]
The latter is the deliberately conservative exponent used in the
main numerical theorem.  Its source-degree conclusion is conditional
only on the separately verified incidence bound \(A\) and supplier
spanning bound \(S\), not on an uncounted sequence of tensor
intertwiners.

\section{A numerical detector bound for Gaussian determinant three}
\label{sec:e8-detector-degree}

\paragraph{Check status.}
The entire proof, including the incidence-component and rational
image-span arguments, has been read by the author and two independent
readers. Their actual events are recorded in CHECK\_EVENTS.tsv.
The geometric, Fermat and lattice inputs were checked separately.
D-E8-XFER01--03 proves the tensor-diagram lemma.

\subsection*{HD-E8Q1. Constants and the bound}

All degrees below are intersection degrees for the stated integral
polarizations. Put
\[
\begin{gathered}
 G=1393459200,\quad g=1+53G/24=3077222401,\quad M=2G,\\
 R=(2G+1)G!G^G,\qquad L=2gR,\\
 A=(4L)!\,[1000g(G+1)(L+1)]^{4L},\qquad
 B=(4L+4)A^{2L+4},\\
 S=[26(L+1)]^L\mathcal H_{12}(L),\qquad
 D_{\rm src}=M^6B^{4L+21}S.
\end{gathered}                                                    \tag{E8Q.1}
\]
Here the explicit integer \(\mathcal H_{12}\) is defined in
SF-SKQ01--05. Every constant uses only factorials, products and
powers of positive integers. The constants are deliberately generous.

For a maximal-order polarized Gaussian sixfold of signature
\((3,3)\), signed determinant class three, and polarization volume
\(\Delta\), define
\[
 \boxed{D_3(\Delta)=69984^{12}M^{114}\Delta^{18}B^{4L+21}S.}       \tag{E8Q.2}
\]
For the fixed order \(\mathbb Z+f\mathbb Z[i]\), replace
\(\Delta\) by \(f^{12}\Delta\), giving a factor \(f^{216}\).
The assertions to prove are
\[
 \delta_W\le D_3(f^{12}\Delta),\qquad
 c_t\le5462\,\delta\lceil\sqrt T\rceil
              (5f^2)^{24}D_3(f^{12}\Delta)^2.                   \tag{E8Q.3}
\]
The normalization of \(\delta,T,c_t\) is fixed in the introduction.
The second estimate uses HD-U with \(a=f(2+i)\), after the first
has supplied a cycle.

\subsection*{HD-E8Q2. The full Brauer identity and its slots}

We need an identity of representations, rather than only an identity
of determinant characters. D-E8-XFER01 applies the integer
column-lattice argument from HECKE\_DIV SQ3 to all irreducible
representation multiplicities. It gives
\[
 \rho\oplus R_-\simeq R_+,\qquad
 R_\pm=\bigoplus_j
       (\operatorname{Ind}_{H_j}^G\chi_j)^{n_{j,\pm}},\qquad
 \sum_j(n_{j,+}+n_{j,-})\le R.                                 \tag{E8Q.4}
\]
The \(\chi_j\) include trivial characters when necessary. This
uses Brauer's integral induction theorem and the proved integer
minor/Cayley-graph bound, not an inference from determinant
characters to a representation identity.

Let \(V_j\) be the associated character multiplicity in \(H^1(C)\),
of dimension \(v_j\le2g\), repeated according to its coefficient.
Replace every minus factor by its conjugate character space.
The number of degree-one curve slots is
\[
 l=\sum_jv_j\le2gR=L,\qquad
 l=6+2\dim\operatorname{Hom}_G(R_-,H^1(C)).                     \tag{E8Q.5}
\]
Thus \(l\) is even and at least six. Zero-dimensional character
spaces have unit determinant and are omitted. The determinant
tensor of these suppliers contracts nontrivially to the primitive
sixfold determinant by (E8Q.4) and polarization duality.
D-E8-XFER02--03 realizes this contraction by bounded algebraic
tensor diagrams. It never takes a root of a determinant cycle.

\subsection*{HD-E8Q3. Degrees of the returning cyclic correspondences}

Choose a line bundle \(H_C\) of degree \(6g\) on \(C\).
It is very ample, since \(6g\ge2g+1\). Give every curve product
the sum polarization; its numerical class is invariant under
the deck automorphisms.

For a nontrivial character in (E8Q.4), let
\(C/\ker\chi_j\to B_j=C/H_j\) be its faithful cyclic cover.
Put \(b_j=g(B_j)\), let \(r_j\) be its number of branch points,
and write \(v_j=2b_j-2+r_j\). Every intermediate inertia order
divides twelve. When \(r_j\ge3\), HD-FS1 uses
\[
 Y_j=\mathbf P_{F_{12}^{r_j-2}}
       (\mathcal O^{b_j}\oplus\mathcal O(-12)),\qquad
 d_j=\dim Y_j=b_j+r_j-2\le v_j.                               \tag{E8Q.6}
\]
Use the convention of lines, write
\(\xi_j=c_1(\mathcal O_{Y_j}(1))\), let \(x_j\) be the
Fermat hyperplane class, and put \(H_j=\xi_j+x_j\).
This is very ample: its section bundle on the base is
\(\mathcal O(1)^{b_j}\oplus\mathcal O(13)\), a sum of
very ample line bundles. Their sections separate base points
and base tangent vectors using any nonzero fibre coordinate;
the full linear system on each projective fibre separates its
points and tangent vectors. This proves the required embedding.
Both \(\xi_j\) and \(x_j\) are globally generated and nef.

For \(r_j=2\), take one component \(Y_j=\mathbf P^{b_j}\).
For \(r_j=0\), take \(Y_j=\mathbf P^{b_j-1}\). Use their
hyperplane classes. These are the split projective suppliers
of HD-FS1, and \(d_j\le v_j\) still holds. The case \(v_j=0\)
was discarded. A zero-dimensional Fermat scheme is treated
componentwise, so its number of points does not multiply an
individual degree bound.

Put \(L_j=K_{B_j}(D_j)\) in the branched case and
\(L_j=K_{B_j}\) in the unramified case. Then
\(\deg L_j=v_j\), and the morphism \(Y_j\to|L_j|\)
pulls back the hyperplane class to \(\xi_j\), or to the
specified projective-space hyperplane class.

In \(C^{v_j}\times Y_j\), require the section represented by
\(y\in Y_j\) to vanish at the image of each curve coordinate
in \(B_j\). There are \(v_j\) Cartier conditions, with classes
\[
 \operatorname{pr}_{C,a}^*q_j^*c_1(L_j)
                     +\operatorname{pr}_{Y_j}^*\xi_j,          \tag{E8Q.7}
\]
where \(q_j:C\to B_j\) has degree at most \(G\).
Thus \(\deg q_j^*L_j\le Gv_j\le2Gg\). Every class in
(E8Q.7) is bounded numerically by \(G+1\) times the total
product polarization, with nef difference.

This incidence is finite flat over \(Y_j\). The nonzero
universal section on the smooth curve has a zero divisor of
constant degree \(v_j\); its pullback to \(C\) also has
constant degree. The relative Cartier divisor is therefore
finite flat, and the independent coordinate fibre products
preserve this property. Repeated zeros produce finite
multiplicities, without excess-dimensional components.

The returning correspondence in HD-FS1 is a component of
this incidence: its ordered images form exactly the divisor
of the chosen section, with the chosen residual cyclic label.
A component assertion can be checked over a dense open set.
The system \(|L_j|\) is basepoint-free: in the branched
case its degree is at least \(2b_j\), and in the unramified
case it is the canonical system for \(b_j\ge2\).
Its general divisor is reduced and avoids the finitely many
branch points of \(C\to B_j\). The supplier map to
\(|L_j|\) is dominant, so these conditions define a dense
open in \(Y_j\). There the incidence is finite etale.
Requiring all projected zeros to be distinct and the
residual label to equal the chosen lift selects an
open-and-closed union of its components. Choosing one
component and taking its closure gives the desired
irreducible returning correspondence on the entire source.
A resolution pushes to its fundamental cycle, so its degree
is bounded by that of the whole incidence. This controls the
actual returning cycle without assigning a degree to an
unspecified resolution.

The bundle relation gives
\[
 \pi_*\xi_j^{b_j+k}=12^kx_j^k\quad(k\ge0),\qquad
 \deg_{H_j}Y_j
  =12\sum_{a=0}^{r_j-2}\binom{d_j}{a}12^{r_j-2-a}
  \le12\,13^{d_j}.                                            \tag{E8Q.8}
\]
For projective spaces the degree is one. A product of these
spaces of total dimension \(d\le L\) consequently has degree
at most \(d!\,12^L13^L\). A product with at most \(L\)
curve slots has dimension at most \(2L\) and degree at most
\((2L)!(6g)^L12^L13^L\). Intersecting the incidence equations
multiplies this upper bound by at most \((G+1)^L\).

For a trivial character, let \(b=g(C/H)\). Its determinant
has \(v=2b\) slots and is Tate. A product of \(b\) diagonals
on \((C/H)^{2b}\), followed by alternation and the
degree-one Kunneth projector, has nonzero top exterior
component: it is the \(b\)-th exterior power of the
nondegenerate symplectic tensor. Pull it to \(C^{2b}\).
Each factor is \(C\times_{C/H}C\), an effective cycle
of degree at most \(12gG\). The product degree is at most
\(b!(12gG)^b\). It is an external Tate cycle, not a
finite correspondence over a point.

Take the product of the nontrivial returning correspondences
and these Tate cycles. If the trivial factors occupy \(l_0\)
slots, the correspondence from \(Y=\prod Y_j\) to \(C^l\)
has dimension \(d+l_0/2\). The input cohomological degree
is \(2d-(l-l_0)\), as in HD-FS1; it is even because
\(l\) and \(l_0\) are even. All factorials from products
and incidence degrees are bounded by
\((2L)!^2\le(4L)!\); all the displayed degree powers are
bounded by \([1000g(G+1)(L+1)]^{4L}\). Hence each returning
correspondence has degree at most \(A\). This bound also
dominates every ambient product degree used here.

\subsection*{HD-E8Q4. Effective Hodge generators on the bundle suppliers}

Let \(F\) be the product of the positive-dimensional Fermat
bases of \(Y\), of total dimension \(t\le d\le L\).
Projective-space factors are bundles over a point. The
projective bundle formula decomposes the full cohomology as
\[
 \xi_1^{j_1}\cdots\xi_k^{j_k}\pi^*H^*(F),
                    \qquad 0\le j_i\le b_i,                  \tag{E8Q.9}
\]
where \(b_i\) is the projective fibre dimension. These
algebraic maps respect rational Hodge structures.
SF-SKQ gives finite effective integral spanning cycles
\(W\) on \(F\), each of degree at most
\(\mathcal H_{12}(t)\le\mathcal H_{12}(L)\).
For a point use its degree-one class.

Flat pullback of \(W\), intersected with general members
of the globally generated \(\xi_i\) systems to the powers
in (E8Q.9), is an effective integral cycle, possibly zero,
with the required class. Give \(Y\) its sum polarization
\(H_Y=\sum_i(\xi_i+x_i)\), omitting \(x_i\) for projective
spaces. To bound the degree, expand the required power of
\(H_Y\) and push to \(F\). A surviving term uses
\(\pi_*\xi_i^{b_i+a}=12^a x_i^a\); for projective spaces
only \(a=0\) survives. Since \(j_i\le b_i\), the factor
\(12^a\) is bounded by assigning a factor twelve to each
\(\xi_i\) selected from this expansion. The sum of all
nonnegative coefficients is therefore at most
\((13k)^d\), with \(k\le L\), or one for a point.
Each remaining monomial in the nef \(x_i\), of total
degree \(\dim W\), has intersection with \(W\) at most
\(\deg_{\sum x_i}W\), by expansion of that total power.
Consequently all these generators have degree at most
\[
                   [26(L+1)]^L\mathcal H_{12}(L)=S.           \tag{E8Q.10}
\]
They span every rational Hodge space of \(Y\), in all
codimensions, by (E8Q.9).

The Brauer supplier tensor has type \((l/2,l/2)\):
the sixfold determinant has type \((3,3)\), and the
minus determinant together with its conjugate contributes
a Tate factor. To justify rational descent, use a finite
Galois splitting field containing all \(G\)-th roots of
unity and the exact primitive root-line projector as in
D-E8-XFER02. Expand all coefficient-field projectors
and transfer maps into their integral correspondence terms.
Let \(I\) be the sum of the rational cohomological images
of those terms, before or after the fixed final map to the
sixfold. These are maps of pure Hodge structures with the
required Tate shifts, so \(I\) is a rational Hodge substructure.
The extended image contains the nonzero determinant line
for the natural character; Galois invariance then includes
its conjugate. Their rational span is the Weil plane.
In particular \(I\) contains a nonzero rational Weil class.
Semisimplicity of polarizable pure Hodge structures gives
a rational Hodge preimage in the appropriate source
cohomology. The effective generators (E8Q.10) span that
entire Hodge space, so some combination has nonzero final
Weil projection. No individual unbalanced character
determinant is asserted to be a Hodge class.

\subsection*{HD-E8Q5. Degree costs and an integral detecting component}

Use the published moving estimate
\cite[Lemma~2.2, author-version p.~13]{Truong2020}.
On a smooth projective ambient \(V\) of dimension \(m\),
the safe factor \(2(m+1)\deg(V)^{m+1}\) represents an
intersection of two effective integral cycles by a
difference of effective integral cycles, each with degree
at most that factor times the two input degrees.

Suppose a correspondence \(\Gamma\subset X\times X'\)
and its smooth product ambient have degrees at most \(A\),
and the ambient dimension is at most \(2L\). If an effective
cycle \(W\subset X\) has degree \(E\), then
\(\deg(W\times X')\le AE\). Indeed the Segre binomial
factor for \(W\times X'\) is no larger than that for
\(X\times X'\), and \(\deg X\ge1\).
The moving estimate, intersection with \(\Gamma\), and
proper projection express \(\Gamma_*[W]\) as a difference
of effective integral cycles with individual degrees at most
\[
                  2(2L+1)A^{2L+3}E\le BE.                   \tag{E8Q.11}
\]
Projection does not increase degree: the target hyperplane
pullback is nef and bounded by the total product hyperplane.
Components whose images drop dimension push to zero.
This estimate allows improper intersections at the outset.

All required correspondences meet these hypotheses.
HD-E8Q3 bounds the returning operation. The degree-one
Kunneth projector on \(C\) is
\(\Delta_C-[c]\times C-C\times[c]\); each term has degree
at most \(12g\). Its \(l\)-fold tensor has terms of degree
at most \(l!(12g)^l<A\).
By D-E8-XFER02 the determinant transfer is a sum of products
of \((l-6)/2+6\) deck graphs in \(C^{l+6}\), each of degree
\[
           ((l-6)/2+6)!(12g)^{(l-6)/2+6}<A.                  \tag{E8Q.12}
\]
Its cyclic-character projectors, representation matrix units,
and determinant permutations are already expanded into these
diagrams. The last primitive projector similarly expands
into coordinatewise deck graphs. All curve-product ambients
have dimension at most \(2L\) and degree at most \(A\).

There are at most four grouped correspondence operations:
the returning map; the full tensor Kunneth projector; the
transfer diagram; and the final primitive projector.
Thus (E8Q.11) bounds individual effective generators at
the end by \(B^4S\le B^{4L+21}S\).
The large exponent in (E8Q.1) is only a convenient
conservative bound; no uncounted tensor operation is needed.

Here is the precise coefficient and effectivity argument.
At every moving step retain the two effective parts as
separate generators. Their rational spans contain the
image of the previous span. Expand every coefficient-field
linear combination into its individual algebraic diagrams.
The final span has nonzero Weil projection after the fixed
last map to the sixfold. If all its integral effective
generators had zero such projection, so would the entire
span. Thus one generator, and then one irreducible component,
has nonzero projection. Taking that component does not
increase degree. Intermediate individual terms need not
belong to the intermediate projector images.
No integral saturation or division of an effective cycle
by a rational coefficient is claimed.

\subsection*{HD-E8Q6. The sixfold map and the target bound}

Let \(p\) be the self-adjoint primitive projector of
D-E8-LAT02. Its image \(B_C=pJ(C)\) is the abelian sixfold,
polarized by \(\theta_C=\Theta_{J(C)}|_{B_C}\).
The map \(f=Mp:J(C)\to B_C\) is an actual homomorphism and
\[
                    f^*\theta_C\le M^2\Theta_{J(C)}          \tag{E8Q.13}
\]
as semipositive forms, since \(p\) is the orthogonal
projection for the Rosati metric. Compose six copies of
\(f\) with Abel maps and addition to obtain
\(a_6:C^6\to B_C\). Cauchy's inequality for addition
of six vectors gives
\[
 a_6^*\theta_C\le6M^2\sum_{j=1}^6a_j^*\Theta_{J(C)}
                     =M^2\sum_{j=1}^6H_{C,j}                \tag{E8Q.14}
\]
numerically: the Abel theta degree is \(g\), whereas
\(\deg H_C=6g\). Positive representatives for this
numerical equality are furnished by the pulled-back
theta forms.

On the tensor product of six degree-one factors, addition
induces exterior multiplication, nonzero on the primitive
determinant. Equivalently, pullback of the conjugate Weil
determinant has its nonzero alternating tensor as the
component with one degree-one class on each curve.
The injective map \(f^*:H^1(B_C)\to H^1(C)\) has
nondegenerate restricted pairing. Its determinant pairing
with the output of D-E8-XFER02 is nonzero. Projection
formula therefore proves nonzero Weil projection after
\(a_{6*}\), as used in the component selection above.
By (E8Q.14), a surviving three-cycle of curve-product
degree at most \(B^{4L+21}S\) pushes to degree at most
\[
                    M^6B^{4L+21}S=D_{\rm src}.              \tag{E8Q.15}
\]

This bound is independent of branch points.
D-E8-GRR and D-E8R each provide a connected component
with period rank nine, so the natural sixfolds cover an
analytic open subset of their fixed-data period domain.
Their polarization volume is at most \(M^6\), by the
actual integral projector calculation in D-E8-LAT02.

For a specified maximal-order target lattice of signed
determinant class three, D-E8-LAT01--02 supplies one fixed
rational Hermitian isometry and a clearing integer
\[
                      N\le69984^2M^{18}\Delta^3.
\]
On the identified marked period domains its isogeny
pulls the target polarization back to \(N^2\theta_C\).
Pushing a detector multiplies its degree by \(N^6\).
Thus an analytic open subset of the target family has
detectors bounded by
\[
 N^6D_{\rm src}
       \le69984^{12}M^{114}\Delta^{18}B^{4L+21}S=D_3(\Delta).
\]
The accepted bounded-cycle closedness and fixed-data
spreading input extends this integer bound to the entire
connected family. Source lattice data are locally constant
on the chosen Hurwitz component, and the same estimate
works for every one of them.

For a fixed conductor \(f\), use
\(f\mathbb Z[i]\Lambda\subset\Lambda\) as in D-E8-LAT02.
Its pulled-back polarization volume is at most
\(f^{12}\Delta\); the last identity isogeny pulls the
target polarization back exactly. This proves the first
inequality in (E8Q.3). Applying HD-U with \(n=3\) and
\(a=f(2+i)\), of norm \(5f^2\), proves the second.
The resulting division bound has conductor exponent
\(432+48=480\) and polarization-volume exponent \(36\).
Cycles are supplied before positivity and the integral
polynomial projector of HD-U are applied.

\section{New restrictions on proposed split-eightfold CM
inputs}\label{new-restrictions-on-proposed-split-eightfold-cm-inputs}

Author/checker: GPT-6 Astra Pro; inherited effort not exposed;
HODGE\_DIVISION /root/sixfolds. Date: 2026-11-21.

These results address the remaining nonsplit sixfold gate through a
possible split-eightfold supplier. They take DEG\_L S-21--S-25 and
S-30--S-34 as established background and do not repeat the excluded
genus-four ample imaginary-secant square. They do not construct the
missing eightfold. The first restriction excludes direct-sum repairs of
the imaginary-secant character. The second excludes a different natural
source: the fully semiregular product of two surface point ideals, even
after derived equivalence, isogeny and rational divisor normalization.

\subsection{SF-E801. Intrinsic exterior
kernels}\label{SF-E801}

Let \(X\) be an abelian fourfold and put

\[
 W_X=H^0(X,T_X)\oplus H^1(X,\mathcal O_X)
     =\operatorname{Lie}(X\times\widehat X).
\]

The Hochschild cohomology algebra is the exterior algebra of \(W_X\).
Choose a nonzero holomorphic four-form \(\Omega\). Contraction by the
first summand and exterior multiplication by the second identify the
full Hodge cohomology with the free exterior module generated by
\(\Omega\). In particular, a diagonal Hodge character
\(v\in\bigoplus_pH^p(\Omega_X^p)\) corresponds to a four-vector
\(\widetilde v\in\bigwedge^4W_X\), and its degree-two contraction map
becomes

\[
 c_v^2:\bigwedge^2W_X\longrightarrow\bigwedge^6W_X,
 \qquad z\longmapsto z\wedge\widetilde v.
 \tag{E801.1}
\]

Changing \(\Omega\) multiplies every displayed four-vector by a nonzero
scalar and changes none of the kernels below.

\subsubsection{The secant kernel}\label{the-secant-kernel}

Let \(W=L\oplus M\), with both summands of dimension four, and let

\[
 v=\lambda\operatorname{vol}_L+\mu\operatorname{vol}_M,
 \qquad\lambda\mu\ne0.
\]

Then

\[
 \ker(z\mapsto z\wedge v)=L\otimes M,
 \qquad\operatorname{rank}c_v^2=12.
 \tag{E801.2}
\]

Indeed, the two pure exterior-square summands map injectively into
distinct bidegrees, while the mixed summand annihilates both volume
forms. Moreover,

\[
 \{w\in\bigwedge^4W:(L\otimes M)\wedge w=0\}
       =\bigwedge^4L\oplus\bigwedge^4M.
 \tag{E801.3}
\]

To prove this last assertion, decompose \(w\) by its \(L,M\) bidegree
((p,4-p)). For \(1\le p\le3\), exterior multiplication gives an
injective map

\[
 \bigwedge^pL\longrightarrow
       \operatorname{Hom}(L,\bigwedge^{p+1}L).
\]

The corresponding map for \(M\) is also injective. Their tensor product
says that a nonzero component of bidegree ((p,4-p)) cannot be
annihilated by every mixed bivector. The two extreme components are
annihilated. This proves (E801.3).

\subsubsection{The product kernel and recovery of its two
factors}\label{the-product-kernel-and-recovery-of-its-two-factors}

Suppose instead that \(W=U_1\oplus U_2\), each \(U_i\) has dimension
four, and \(\omega_i\in\bigwedge^2U_i\) is nondegenerate. Put
\(v=\omega_1\wedge\omega_2\). Then

\[
 \ker c_v^2=K_1\oplus K_2,
 \qquad K_i=\{z\in\bigwedge^2U_i:z\wedge\omega_i=0\},
 \qquad\dim K_i=5,
 \tag{E801.4}
\]

and \(c_v^2\) has rank eighteen. On the two pure exterior-square
summands its image has dimension one each, in bidegrees ((4,2)) and
((2,4)). On \(U_1\otimes U_2\), it is the tensor product of the
isomorphisms \(u\mapsto u\wedge\omega_i\), so has rank sixteen in
bidegree ((3,3)).

The decomposable nonzero bivectors in the projective kernel form exactly
two disjoint smooth quadric threefolds, one in each \(\mathbf P(K_i)\).
To see that no mixed point occurs, write \(z=z_1+z_2\in K_1\oplus K_2\).
If \(z\) is decomposable, then \(z^2=0\). Its bidegree ((2,2)) part is
\(2z_1\otimes z_2\), forcing one \(z_i\) to vanish. Within
\(\bigwedge^2U_i\), decomposability is the Plücker quadratic equation
\(z_i^2=0\). Its restriction to the hyperplane \(K_i=\omega_i^\perp\) is
nondegenerate, because the exterior-square pairing is nondegenerate and
\(\omega_i^2\ne0\). It is therefore a smooth irreducible quadric in
\(\mathbf P^4\).

The two original four-dimensional summands are recovered intrinsically
from these two components: for each component, take the span in \(W\) of
the two-planes represented by its points. This span is \(U_i\). Every
nonzero vector in a symplectic four-space lies in a Lagrangian
two-plane, which verifies the last claim directly. Thus the unordered
pair \({U_1,U_2}\) is determined by \(ker c_v^2\), without a chosen
presentation of \(v\).

The accompanying exact script verifies the ranks (12,18), the
two-dimensional common annihilator in (E801.3), and the rank-five
Plücker quadratic forms. The geometric component and recovery arguments
are the proofs above, not inferences from numerical sampling.

One further common-annihilator identity is useful:

\[
 \{w\in\bigwedge^4W:(K_1\oplus K_2)\wedge w=0\}
       =\mathbf C(\omega_1\wedge\omega_2).
 \tag{E801.5}
\]

For a bidegree \((p,4-p)\) component, the equations for \(K_1\) kill
\(p=0,1\), and those for \(K_2\) kill \(p=3,4\). At \(p=1\), the
bivectors annihilating a fixed nonzero vector form a three-dimensional
space and cannot contain the five-dimensional \(K_1\); applying
coefficient functionals gives the tensor-valued assertion. At \(p=2\),
orthogonality to the primitive hyperplane makes the first factor
proportional to \(\omega_1\), and the second condition makes the second
factor proportional to \(\omega_2\). This proves (E801.5), also checked
by the exact script.

\subsection{SF-E802. Direct sums cannot repair the imaginary-secant
character}\label{SF-E802}

Let \(E\) be a gluable fully semiregular perfect object on an abelian
fourfold. Suppose that, for a rational algebraic divisor class \(B\), a
rational ample class \(\theta\), and \(s^2=-d<0\),

\[
 e^B\operatorname{ch}(E)
 =a\frac{e^{s\theta}+e^{-s\theta}}2
  +b\frac{e^{s\theta}-e^{-s\theta}}{2s},
 \qquad (a,b)\in\mathbf Q^2\setminus\{(0,0)\}.
 \tag{E802.1}
\]

In any decomposition \(E=\bigoplus_jE_j\), at most one summand has
nonzero rational Chern character. In particular, if \(E\) is a coherent
sheaf, it is indecomposable. The same assertion holds for a
decomposition in an effective finite equivariant category of the kind in
DEG\_L S-21.

\textbf{Proof.} The B-field acts by an invertible shear of \(W_X\), so
it can be applied to all the characters at once. The two exponential
characters become decomposable four-vectors with transverse
four-dimensional supports. Transversality follows from nondegeneracy of
\(2s\theta\). Their coefficients \((a\pm b/s)/2\) are both nonzero,
since \(a,b\) are rational and \(s\) is nonreal. Thus (E801.2) applies.

The CY4 adjunction used in S-21 says that injectivity of the full trace
\(\sigma_E^2\) makes Hochschild evaluation onto, so

\[
 \operatorname{im}\sigma_E^2=\operatorname{im}c_E^2,
 \qquad\dim\operatorname{Ext}^2(E,E)=12.
\]

For a direct sum, the Atiyah class is block diagonal. The trace is zero
on off-diagonal self-Ext blocks and is the usual full trace on each
diagonal block. Injectivity therefore makes the diagonal trace images
linearly independent. For each summand,

\[
 \operatorname{im}c_{E_j}^2\subset
 \operatorname{im}\sigma_{E_j}^2\subset\operatorname{im}c_E^2.
\]

The wedge map in (E801.1), paired with another bivector to obtain a top
form, is a symmetric bilinear form on \(\bigwedge^2W_X\). Hence the
image inclusion implies

\[
 \ker c_E^2\subset\ker c_{E_j}^2.
\]

Equation (E801.3) now forces each normalized summand character into the
same two-dimensional secant plane. Rationality makes it
\(a_j\alpha+b_j\beta\) with \(a_j,b_j\in\mathbf Q\): the rank determines
\(a_j\), and the degree-two cohomology component determines
\(b_j\theta\). If this character is nonzero, both pure coefficients are
again nonzero, so its degree-two contraction rank is twelve. Two such
summands would give two independent trace images of dimension at least
twelve inside a twelve-dimensional space. This is impossible.

Every nonzero coherent sheaf has nonzero rational Chern character: its
first nonzero component is the effective cycle of its
maximal-dimensional support with positive generic lengths. A coherent
direct-sum decomposition therefore has only one nonzero summand. For an
equivariant decomposition, S-21's identity-sector Hochschild isomorphism
and invariant CY4 adjunction give exactly the same argument; the common
nonzero scalar normalization by the group order changes no kernel or
image. ∎

This does not exclude an indecomposable nonsimple complex or sheaf.
Together with S-21, it says that a coherent solution in the
imaginary-secant branch must be indecomposable and have endomorphism
dimension greater than two. The nonsimple control
\(\mathcal O_X\oplus\mathcal O_p\) from S-23 cannot itself be turned
into such a character by derived equivalence, since its two summands
have nonzero characters and the argument is categorical.

The same direct-sum conclusion holds for the balanced RM character of
S-22: use (E801.5) in place of (E801.3). Every nonzero summand character
is then proportional to the whole balanced character and has contraction
rank eighteen, whereas the trace image of the sum has dimension
eighteen. Thus a coherent fully semiregular balanced-RM realization is
necessarily indecomposable, with no simplicity or rank hypothesis.

\subsection{SF-E803. Surface point-ideal products cannot become the
genuine RM
character}\label{SF-E803}

Let \(S_1,S_2\) be abelian surfaces, let \(p_i\in S_i\), and put

\[
 X=S_1\times S_2,
 \qquad F=I_{p_1}\boxtimes I_{p_2}.
\]

Although \(F\) is a simple gluable fully semiregular object, no sequence
of derived equivalences of abelian varieties, isogeny pullbacks or
pushforwards, tensor products by line bundles, shifts, and nonzero
scalar multiplication of characters can give a character \(w\) on an
abelian fourfold \(Y\) with

\[
 e^B w=N\left(1-\frac q2\theta_1^2\right)
          \left(1-\frac q2\theta_2^2\right),
 \qquad Nq\ne0,
 \tag{E803.1}
\]

where \(B\) is a rational algebraic divisor and \(\theta_1,\theta_2\)
are the two eigenspace components of a compatible polarization for a
genuine real quadratic field acting with multiplicities ((2,2)). The
obstruction does not depend on the output rank or on an integral
normalizer.

\textbf{Proof of the input properties.} On an abelian surface, \(I_p\)
is simple of rank one and has no negative self-Ext. Serre duality gives
\(\operatorname{ext}^2=1\), and its degree-two trace is nonzero because
the trace of the identity in degree zero is its rank. Riemann--Roch
gives \(\chi(I_p,I_p)=-2\), hence \(\operatorname{ext}^1=4\). Its first
contraction rank is four: the translation directions detect the point
and the degree-zero tensor directions detect the rank, in separate Hodge
components. Thus the degree-one full trace is injective as well. The
product trace is injective on the three Künneth summands in degree two,
lying in distinct pairs of Hodge differences. This proves the assertions
for \(F\), with \(\operatorname{ext}^1(F,F)=8\),
\(\operatorname{ext}^2(F,F)=18\), and \(\chi(F,F)=4\).

\textbf{Proof of the obstruction.} For \(X\), the Fock form of
\(1-[p_i]\) is a nondegenerate two-vector \(\omega_i\) on

\[
 U_i=\operatorname{Lie}(S_i\times\widehat S_i)
       \subset W_X.
\]

Indeed its two nonzero terms are the volume forms of the two-dimensional
translation and tensor directions; their coefficients are nonzero. The
product character is therefore \(\omega_1\wedge\omega_2\). By SF-E801,
its degree-two contraction kernel intrinsically recovers the unordered
pair \(U_1,U_2\). Both recovered spaces are Lie algebras of abelian
subvarieties of \(X\times\widehat X\).

This property is preserved by every listed operation. A derived
equivalence induces an isomorphism of the abelian varieties
\(X\times\widehat X\) and \(Y\times\widehat Y\), whose differential is
the induced map on \(HH^1\); naturality of Hochschild cap product
carries the contraction kernel with it. This is the exact
translation--tensor conjugation assertion of Orlov Theorem2.10 and
Corollary2.13, not an arbitrary complex linear change of coordinates. An
isogeny \(f:Y\to X\) transports \(W_X\) to \(W_Y\) by
\((df^{-1},d\widehat f)\), the differential of a rational isogeny. On
rational cohomology, pushforward differs from inverse pullback by the
nonzero degree scalar. Thus both preserve the property that a recovered
subspace is the Lie algebra of an abelian subvariety. A rational
algebraic B-field is the differential of the rational shear

\[
 (x,\alpha)\longmapsto(x,\alpha+\phi_B(x)),
\]

where \(\phi_B\in\operatorname{Hom}(X,\widehat X)\otimes\mathbf Q\); it
has inverse obtained by replacing \(B\) with (-B). Clearing its
denominator shows that it also preserves algebraic subgroup directions.
Shifts and nonzero character scalars do not change the contraction
kernel.

For (E803.1), on the other hand, SF-E801 recovers the two spaces

\[
 W_\pm=H^0(T_Y)_\pm\oplus H^1(\mathcal O_Y)_\pm,
\]

the eigenspaces of the diagonal real-field action on
\(Y\times\widehat Y\). The polarization identifies the action on the
dual with the same real embedding, since it is Rosati invariant. Neither
\(W_+\) nor \(W_-\) is the Lie algebra of a positive-dimensional abelian
subvariety. Choose a nonrational real field element \(u\). If such a
subvariety existed, its underlying real tangent space and rational first
homology would lie in an eigenspace of the rational operator \(u\) with
irrational real eigenvalue \(\sigma(u)\). A nonzero rational vector
\(z\) cannot satisfy \(uz=\sigma(u)z\): any nonzero rational coordinate
of \(z\) would make \(\sigma(u)\) rational. This is a contradiction.

Thus the intrinsic recovered pair would have to be both algebraic and
nonalgebraic, proving the obstruction. ∎

The same proof applies to an external product of any two surface
characters whose Fock two-vectors are nondegenerate. It excludes a
concrete higher-rank construction strategy beyond the rank-one and
low-rank tests S-30/S-31. It does not exclude a new indecomposable
fourfold object with the required genuine RM character. SF-E805
strengthens the stated point-ideal-product obstruction to allow
interposed flat Hodge deformations; it does not assume that an inverse
equivalence deforms.

\subsection{SF-E804. A necessary Euler inequality for mixed-constituent
extensions}\label{SF-E804}

Let \(\mathcal H\) be a bounded heart in the ordinary or effective
finite equivariant fourfold category. Suppose \(E\in\mathcal H\) has a
finite filtration with constituents \(F_1,\ldots,F_s\), occurring with
multiplicities \(m_i\), such that

\[
 \operatorname{Hom}(F_i,F_j)=0\ (i\ne j),
 \qquad\operatorname{End}(F_i)=\mathbf C.
\]

If \(E\) has a nonzero normalized imaginary-secant character and is
fully semiregular, then

\[
 \chi(E,E)\le2\sum_i m_i^2-4.
 \tag{E804.1}
\]

Here \(\chi\) is the Euler pairing in the category in which the
filtration is taken. In particular, it is impossible that every
\(\chi(F_i,F_i)\ge2\) and every \(\chi(F_i,F_j)\ge0\) for \(i\ne j\).

\textbf{Proof.} Applying left exact Hom successively to both filtrations
gives

\[
 \dim\operatorname{Hom}(E,E)
       \le\sum_{i,j}m_im_j\dim\operatorname{Hom}(F_i,F_j)
       =\sum_i m_i^2.
\]

S-21 gives \(e_2=12\) and \(e_1\ge8\). CY4 duality and gluability give
\(\chi(E,E)=2e_0-2e_1+12\le2e_0-4\), proving (E804.1). Euler bilinearity
gives the final contradiction under the stated diagonal and off-diagonal
conditions. ∎

This extends the numerical heart-filtration test of S-24 to distinct
mutually Hom-orthogonal constituents. A remaining mixed-constituent
proposal must fail at least one of these explicit Euler inequalities or
the Hom-orthogonality hypothesis; this is a necessary gate, not an
existence assertion.

\subsection{SF-E805. The rational stabilizer algebra closes the
deformation
escape}\label{SF-E805}

For a rational even class \(v\) on an abelian fourfold, let

\[
 H=H_1(X,\mathbf Q),\qquad
 V=H\oplus H^*,\qquad S=\bigwedge^*H^*,\qquad
 \mathfrak g_v=\{a\in\mathfrak{so}(V):a\cdot v=0\}.
\]

Here \(x+\xi\in V\) acts on \(S\) by contraction with \(x\) plus wedge
with \(\xi\), and the orthogonal Lie algebra has its resulting spin
action. Define the rational algebra

\[
 \mathcal C(v)=\operatorname{End}_{\mathfrak g_v}(V).
 \tag{E805.1}
\]

For the point-ideal product character \(v=(1-[p_1])(1-[p_2])\),

\[
 \mathcal C(v)\simeq\mathbf Q\times\mathbf Q.
 \tag{E805.2}
\]

For the genuine real-quadratic balanced character in (E803.1), with
field \(F\),

\[
 \mathcal C(e^{-B}Nv_q)\simeq F.
 \tag{E805.3}
\]

These algebras are preserved up to isomorphism by the operations in
SF-E803 and by parallel transport of the full rational character in a
smooth family of abelian varieties. Consequently SF-E803 remains valid
with any number of such flat Hodge deformations interposed among the
listed operations.

\textbf{Explicit stabilizer calculation.} Work first over any
characteristic-zero field \(k\). Let \(H_i\) have dimension four, let
\(\nu_i\in\bigwedge^4H_i^*\) be nonzero, and put

\[
 v_i=1-c_i\nu_i,\qquad c_i\in k^*,\qquad
 V_i=H_i\oplus H_i^*.
\]

The three blocks of \(\mathfrak{so}(V_i)\) are
\(\operatorname{End}(H_i)\), \(\bigwedge^2H_i\), and
\(\bigwedge^2H_i^*\). For \(A\in\operatorname{End}(H_i)\), the spin
action is \(-A^*+\operatorname{tr}(A)/2\), where \(A^*\) denotes the
induced derivation on the exterior algebra. Its degree-zero and
degree-four components on \(v_i\) are

\[
 \frac{\operatorname{tr}(A)}2(1+c_i\nu_i).
\]

The other two blocks act in degree two: contraction by
\(C\in\bigwedge^2H_i\) and wedge by \(b\in\bigwedge^2H_i^*\) give
\(b-c_i\iota_C\nu_i\). Therefore

\[
 \mathfrak g_{v_i}
 =\mathfrak{sl}(H_i)\oplus
  \{C+b:\ b=c_i\iota_C\nu_i\},\qquad
 \dim\mathfrak g_{v_i}=15+6=21.
 \tag{E805.4}
\]

The displayed sum is a vector-space decomposition inside the orthogonal
Lie algebra; its two summands are not being asserted to be ideals.

Its commutant on \(V_i\) is exactly \(k\). The standard and dual
four-dimensional modules for \(\mathfrak{sl}(H_i)\) are nonisomorphic
irreducible modules. Hence an endomorphism commuting with this
subalgebra is scalar \(a\) on \(H_i\) and scalar \(b\) on \(H_i^*\). A
nonzero wedge/contraction pair in (E805.4) has an off-diagonal block, so
commutation forces \(a=b\). These assertions can also be verified
against the elementary matrices of \(\mathfrak{sl}_4\), without a
classification of spin representations.

The Clifford map

\[
 V_i\longrightarrow S_i^{\rm odd},\qquad
 (x,\xi)\longmapsto\xi-c_i\iota_x\nu_i
 \tag{E805.5}
\]

is an isomorphism: its degree-one and degree-three pieces are separate
isomorphisms. For \(V=V_1\oplus V_2\) and \(v=v_1\otimes v_2\), the
mixed block of the orthogonal Lie algebra is \(V_1\otimes V_2\). Its
action on \(v\), up to a fixed nonzero sign, is the tensor product of
(E805.5). It is therefore an isomorphism onto
\(S_1^{\rm odd}\otimes S_2^{\rm odd}\). It cannot cancel the diagonal
blocks, whose images lie in \(S_1^{\rm even}\otimes S_2^{\rm even}\).

For the diagonal blocks, the equation

\[
 a_1v_1\otimes v_2+v_1\otimes a_2v_2=0
\]

forces \(a_iv_i=\lambda_i v_i\), with \(\lambda_1+\lambda_2=0\).
Comparing degree zero and degree four in the preceding formula forces
\(\lambda_i=0\): the coefficients of \(v_i\) have opposite signs, while
those of \(A\cdot v_i\) have the same sign. The degree-two component
then gives (E805.4). Thus

\[
 \mathfrak g_v=\mathfrak g_{v_1}\oplus\mathfrak g_{v_2}.
 \tag{E805.6}
\]

There are no intertwiners between \(V_1\) and \(V_2\): the first Lie
algebra factor acts without a fixed vector or trivial quotient on
\(V_1\) and trivially on \(V_2\), and conversely. The scalar commutant
on each block consequently gives \(\mathcal C(v)=k\oplus k\).

For a product of surfaces these blocks are rational, proving (E805.2).
For a genuine RM character, extend scalars from \(\mathbf Q\) to its
real quadratic field \(F\). The two eigenspaces of \(H\) have dimension
four over \(F\), and \(q\theta_i^2/2\) is a nonzero multiple of a volume
form on each. The calculation gives
\(\mathcal C(v)\otimes F=F\oplus F\). The nontrivial automorphism of
\(F/\mathbf Q\) exchanges the two eigenspaces and their two scalar
idempotents. The fixed algebra is \(\{(a,\bar a):a\in F\}\), isomorphic
to \(F\). This proves (E805.3) before B-normalization; a rational
B-shear conjugates the algebra.

\textbf{Invariance under transformations and deformation.} A nonzero
scalar multiple of \(v\) has the same stabilizer. A rational isometry of
\(V\), together with an intertwiner of its Clifford modules, conjugates
the stabilizer and its commutant. Line-bundle twists and rational
B-fields have the explicit exponential Clifford action. Rational isogeny
pullback has the exterior action of the induced rational linear map on
\(H^*\), and pushforward is a scalar times its inverse.

For derived autoequivalences, the required Clifford compatibility is
Golyshev--Lunts--Orlov, Proposition4.3.7 and Lemma4.3.7.1; the complete
proof was checked in the author version pp17--19. Equivalences between
different abelian varieties follow from the same argument with source
and target distinct, or by this explicit reduction. For
\(\Phi:D^b(X)\to D^b(Y)\), apply the autoequivalence result on
\(X\times Y\) to \((\Phi^{-1}\boxtimes\Phi)\circ\mathrm{swap}\). On
cohomology this swaps the tensor factors through \(\Phi^H\) and its
inverse. Restricting the resulting Clifford identity to the source
factor gives the intertwining identity for \(\Phi^H\). This uses an
actual equivalence at the fibre where it is applied.

In a smooth family, rational homology, its dual, the exterior algebra,
and the Clifford action are local systems. A flat full character
therefore gives a fixed stabilizer and commutant under parallel
transport. No inverse equivalence has to extend across the deformation.
Finally, \(\mathbf Q\times\mathbf Q\) has a nontrivial idempotent, while
\(F\) is a field, so the two algebras cannot be isomorphic. This proves
the obstruction through interposed flat deformations. ∎

For a direct untransformed deformation, there is a shorter check
suggested by the parent audit. The degree-four class \([p_1]+[p_2]\) is
the sum of two transverse rational volume forms on \(H^1(X,\mathbf Q)\).
Its mixed wedge kernel recovers the unordered pair of rational
four-dimensional summands. The connected Hodge circle preserves each
summand. They remain rational Hodge substructures and hence give an
isogeny decomposition into abelian surfaces. The Clifford calculation
handles the transformed setting, where an individual degree-four
component need not stay separately recognizable.

The same deformation obstruction holds for a product of any two
nonzero-rank surface characters with nonzero self-Euler pairings.
Dividing each character by its rank and multiplying by the rational
exponential of minus its slope puts it in the form \(1-c_i\nu_i\). The
self-Euler hypothesis is precisely \(c_i\ne0\). The slope is a rational
algebraic divisor, so this is one of the rational Clifford shears
already treated. Thus using different positive-rank surface constituents
does not evade SF-E805.

\subsection{SF-E806. An all-rank integral obstruction on a principal
real-multiplication
datum}\label{SF-E806}

Let \(A=E_i^4\) with its product principal polarization \(\theta\), and
let

\[
 U=\operatorname{diag}\left(
 \begin{pmatrix}1&1\\1&-1\end{pmatrix},
 \begin{pmatrix}1&1\\1&-1\end{pmatrix}\right).
 \tag{E806.1}
\]

This integral symmetric endomorphism satisfies \(U^2=2I\), so it gives a
Rosati-invariant maximal-order action of \(F=\mathbf Q(\sqrt2)\), of
multiplicity two at each real embedding on the tangent space. Write
\(\theta_\pm\) for the two real eigencomponents of \(\theta\), and put

\[
 C=\frac{\theta_+^2+\theta_-^2}{2}.
\]

\textbf{The integral assertion.} For every positive integer \(r\) and
every integral class \(\ell\in H^2(A,\mathbf Z)\),

\[
 \frac{\ell^2}{2r}-C\notin H^4(A,\mathbf Z).
 \tag{E806.2}
\]

This quantifies over all integral degree-two classes, not only algebraic
classes or the RM divisor plane. The same assertion holds on any
deformation preserving this integral polarized RM datum.

Consequently there is no positive-rank perfect object \(E\), with any
rational divisor normalization \(B\), satisfying

\[
 e^B\operatorname{ch}(E)
 =r\left(1-\frac{\theta_+^2}{2r}\right)
     \left(1-\frac{\theta_-^2}{2r}\right).
 \tag{E806.3}
\]

In particular, the datum has no ordinary simple gluable fully
semiregular object of the balanced RM character in any positive rank.
Indeed S-22 gives \(4q^2r^2\in\{2,4\}\) for such an object. Since \(q\)
is positive rational, the value two would make \((qr)^2=1/2\),
impossible in \(\mathbf Q\). Thus \(q=1/r\), which is (E806.3). This
last implication uses the supplied S-22; the integral assertion below is
independent of semiregularity and simplicity.

\subsubsection{The integral four-form and its obstruction
tensor}\label{the-integral-four-form-and-its-obstruction-tensor}

Choose an integral basis \(a_0,b_0,\ldots,a_3,b_3\) of
\(H^1(A,\mathbf Z)\), with \(\theta=\sum_i a_i b_i\); juxtaposition
means exterior product. Set

\[
 \theta_U=\sum_{i,j}U_{ij}a_i b_j,\qquad
 \theta_\pm=\frac{\theta\pm\theta_U/\sqrt2}{2},\qquad
 Q=4C=\theta^2+\frac{\theta_U^2}{2}.
 \tag{E806.4}
\]

The divided square \(x^{[2]}=x^2/2\) of an integral two-form is
integral. Thus \(Q\) is integral. Make the unimodular integral change of
basis

\[
 \begin{aligned}
 s_0&=a_0+a_1,&s_1&=b_0+b_1,&t_0&=a_1,&t_1&=b_1,\\
 s_2&=a_2+a_3,&s_3&=b_2+b_3,&t_2&=a_3,&t_3&=b_3 .
 \end{aligned}
\]

Let \(S=\langle s_0,\ldots,s_3\rangle\),
\(T=\langle t_0,\ldots,t_3\rangle\), and

\[
 \omega=s_0s_1+s_2s_3,\quad
 \tau=t_0t_1+t_2t_3,\quad
 m=s_0t_1+t_0s_1+s_2t_3+t_2s_3.
\]

Direct substitution gives

\[
 \theta=\omega-m+2\tau,\qquad
 \theta_U=\omega-2\tau.
 \tag{E806.5}
\]

Consequently

\[
 Q\bmod2=\omega^{[2]}=s_0s_1s_2s_3,\qquad
 R:=\frac{Q_{(2,2)}}2\bmod2
    =m^{[2]}+\omega\otimes\tau .
 \tag{E806.6}
\]

The subscript denotes the \(S,T\) bidegree. The mixed two-form \(m\) is
a tensor of rank four in \(S\otimes T\). Its divided square is the
exterior-square tensor of the corresponding isomorphism, and therefore
has rank six in

\[
 \bigwedge^2(S/2S)\otimes\bigwedge^2(T/2T).
\]

Adding \(\omega\otimes\tau\), a rank-one tensor, leaves
\(\operatorname{rank}R\ge5\). This is the only numerical rank needed
below. The exact certificate computes rank six. The first identity in
(E806.6) also shows that \(Q\) is primitive at two.

\subsubsection{A two-adic divided-square lifting
lemma}\label{a-two-adic-divided-square-lifting-lemma}

Let \(L=S\oplus T\) be a free \(\mathbf Z_2\)-module with both summands
of rank four. Suppose \(Q\in\bigwedge^4L\) has reduction equal to a
nonzero volume form of \(S/2S\), and suppose the tensor

\[
 R=\frac{Q_{(2,2)}}2\bmod2
 \in\bigwedge^2(S/2S)\otimes\bigwedge^2(T/2T)
\]

has rank at least two. Then, for no positive integer \(r\) and no
\(\ell\in\bigwedge^2L\), is

\[
 \ell^{[2]}-\frac r4Q\in r\bigwedge^4L.
 \tag{E806.7}
\]

The reduction assumption makes \(Q_{(2,2)}\) divisible by two, so the
definition of \(R\) is meaningful.

\textbf{Proof.} Integrality of the left side forces \(4\mid r\), because
\(Q\) has a unit coefficient. Write \(r=2^e k\), with \(k\) odd and
\(e\ge2\). The alleged congruence is

\[
 \ell^{[2]}\equiv 2^{e-2}kQ\pmod {2^e}.
 \tag{E806.8}
\]

Its left side has minimum coefficient valuation exactly \(e-2\). If
\(\ell=2^a\ell'\) with \(\ell'\) primitive, then \(2a\le e-2\). Dividing
by \(2^{2a}\) reduces to the same congruence for \(\ell'\) and exponent
\(e'=e-2a\ge2\). It therefore suffices to consider primitive \(\ell\).

If \(e=2\), then \(\ell^{[2]}\bmod2\) is the nonzero decomposable volume
form of \(S/2S\). An alternating two-form over a field of characteristic
two whose divided square is a nonzero decomposable four-form has rank
four and support equal to the support of that four-form. Here is a
direct check. In a symplectic basis it is a sum of disjoint pairs. If it
has at least two pairs, every vector in its support occurs nontrivially
in a contraction of its divided square, so the support of the divided
square equals its own support. The nonzero decomposable divided square
has support dimension four. Thus

\[
 \ell=\ell_0+2D,\qquad
 \ell_0\in\bigwedge^2S,
\]

with \(\ell_0\bmod2\) nondegenerate. In bidegree \((2,2)\), modulo four,
its divided square is

\[
 \ell^{[2]}_{(2,2)}
   \equiv 2\ell_0\otimes D_{(0,2)}\pmod4.
\]

After division by two and reduction modulo two, this has tensor rank at
most one. The right side of (E806.8) has tensor rank
\(\operatorname{rank}R\ge2\). This is a contradiction.

Now suppose \(e\ge3\). The alternating elementary-divisor algorithm over
\(\mathbf Z_2\) puts \(\ell\) in the form of a sum of disjoint pairs
with coefficients \(2^{a_1},\ldots,2^{a_j}\), absorbing units into basis
vectors. This algorithm follows by choosing an entry of minimum
valuation, clearing its row and column using divisibility by that entry,
and continuing on the complementary summand. Since \(\ell\) is
primitive, \(a_1=0\). The minimum coefficient valuation of
\(\ell^{[2]}\) is \(a_1+a_2=e-2\), so \(a_2=e-2\) and every subsequent
coefficient has at least that valuation. In the original \(S,T\)
coordinates we may consequently write

\[
 \ell=\alpha+2^{e-2}\beta,
\]

where \(\alpha\) is primitive and decomposable and \(\beta\) is
integral. Dividing (E806.8) by \(2^{e-2}\) gives

\[
 \alpha\beta+2^{e-2}\beta^{[2]}\equiv kQ\pmod4.
 \tag{E806.9}
\]

Let \(P\) be the two-dimensional support of \(\bar\alpha=\alpha\bmod2\).
Reducing (E806.9) modulo two gives

\[
 \bar\alpha\,\bar\beta=\operatorname{vol}_{S/2S}.
 \tag{E806.10}
\]

Every vector in \(P\) wedges to zero with the left side, hence with the
right side. The wedge-annihilator of this nonzero decomposable volume
form is \(S/2S\), so \(P\subset S/2S\). Comparing the \((2,2)\) and
\((3,1)\) components in (E806.10) now gives

\[
 \bar\beta_{(0,2)}=0,\qquad
 \bar\beta_{(1,1)}\in P\otimes(T/2T).
 \tag{E806.11}
\]

Indeed wedge multiplication by the nonzero \(\bar\alpha\in\bigwedge^2P\)
is injective on \(\bigwedge^2(T/2T)\); on \(S/2S\) its kernel is exactly
\(P\).

Write \(\alpha=a b\), with \(\bar a,\bar b\) a basis of \(P\). The
\(T\)-components of \(a,b\) are divisible by two. Thus
\(\alpha_{(0,2)}\) is divisible by four, \(\alpha_{(1,1)}\) is divisible
by two, and

\[
 (\alpha_{(1,1)}/2)\bmod2\in P\otimes(T/2T).
\]

Consider the \((2,2)\) component of \(\alpha\beta\), divide by two, and
reduce modulo two. Its three contributions are

\[
 \alpha_{(2,0)}\beta_{(0,2)},\quad
 \alpha_{(1,1)}\beta_{(1,1)},\quad
 \alpha_{(0,2)}\beta_{(2,0)}.
\]

The first, after division, has its \(S\)-factor in the line
\(\mathbf F_2\bar\alpha\). The second has both of its \(S\)-factors in
\(P\) by (E806.11), so its \(S\)-factor lies in the same line
\(\bigwedge^2P\). The third vanishes after division and reduction.

If \(e\ge4\), the correction \(2^{e-2}\beta^{[2]}\) in (E806.9) vanishes
modulo four. If \(e=3\), dividing its \((2,2)\) component by two gives

\[
 (\bar\beta^{[2]})_{(2,2)}
  =\bar\beta_{(2,0)}\bar\beta_{(0,2)}
     +(\bar\beta_{(1,1)})^{[2]}.
\]

The first term is zero by (E806.11); the second again has its
\(S\)-factor in \(\bigwedge^2P\). In every case the \((2,2)\) component
of the left side of (E806.9), divided by two modulo two, belongs to

\[
 \mathbf F_2\bar\alpha\otimes
       \bigwedge^2(T/2T),
\]

and has tensor rank at most one. The right side gives \(kR=R\), of rank
at least two. This final contradiction proves the lemma. ∎

\subsubsection{Application to Chern
characters}\label{application-to-chern-characters}

Apply the lemma to (E806.6), where \(\operatorname{rank}R=6\). This
proves (E806.2). If an object satisfied (E806.3), its degree-two part
would force \(B=-\ell/r\), where \(\ell=c_1(E)\). Comparing degree four
then gives

\[
 \operatorname{ch}_2(E)=\frac{\ell^2}{2r}-C.
 \tag{E806.12}
\]

The left side is integral on a complex torus: the formula
\(\operatorname{ch}_2(E)=\ell^2/2-c_2(E)\) holds for perfect K-classes,
the second Chern class is integral, and the divided square of every
integral two-form is integral in the exterior cohomology lattice. This
contradicts (E806.2).

The obstruction is topological in the displayed integral RM datum. It
therefore persists when the point is deformed to a simple or non-CM
fourfold with the same datum. It is stronger than the
minimal-divisor-space hypothesis in supplied S-32 for this datum,
because every integral \(\ell\) was allowed. The following local
argument extends the assertion to other integral data without imposing a
hypothesis on the Néron--Severi space.

\subsection{SF-E807. Maximal-order extension at
two}\label{SF-E807}

Let \(F=\mathbf Q(\sqrt{\Delta})\), where \(\Delta>1\) is even and
squarefree. Suppose a polarized abelian fourfold \((X,\theta)\) has
Rosati-invariant RM by \(F\), with tangent multiplicities \((2,2)\), and
suppose its integral homology at two is stable under the maximal local
order

\[
 R=\mathbf Z_2[u],\qquad u^2=\Delta.
\]

Assume that \(D=\int_X\theta^4/24\) is odd. For every
\(\ell\in H^2(X,\mathbf Z_2)\), every positive integer \(r\), and every
\(z\in\mathbf Q_2\) of valuation zero,

\[
 \frac{\ell^2}{2r}-zC\notin H^4(X,\mathbf Z_2),\qquad
 C=\frac{\theta_+^2+\theta_-^2}{2}.
 \tag{E807.1}
\]

Thus no perfect object, with any rational divisor normalization, has
balanced normalized character \(r v_q\) when \(rq\) is a two-adic unit.
No assumption on \(\operatorname{NS}(X)_{\mathbf Q}\) is needed. In
particular every simple gluable fully semiregular lift for an effective
finite coherent target action \(G\) as in S-22 satisfies

\[
 2\mid D|G|.
 \tag{E807.2}
\]

\textbf{Proof of the local lattice normalization.} Let
\(L=H_1(X,\mathbf Z_2)\) and let \(E\) be its alternating polarization
pairing. Since \(D\) is odd, \(E\) is perfect over \(\mathbf Z_2\). The
polynomial \(t^2-\Delta\) is Eisenstein at two. More explicitly, for
\(a,b\in\mathbf Z_2\), the valuation of \(a+bu\) in
\(F\otimes\mathbf Q_2\), normalized by \(v(u)=1\), is

\[
 v(a+bu)=\min\{2v_2(a),\,2v_2(b)+1\};
\]

the two possible valuations have different parity. This gives the
discrete valuation ring \(R\), whose uniformizer is \(u\). The
torsion-free finite \(R\)-module \(L\) is free of rank four, by the
elementary basis algorithm over a discrete valuation ring.

Define an \(R\)-valued form

\[
 \Psi(x,y)=E(x,uy)+uE(x,y).
 \tag{E807.3}
\]

The Rosati-invariance identity \(E(ux,y)=E(x,uy)\) shows directly that
\(\Psi\) is \(R\)-bilinear. It is alternating:
\(E(x,ux)=E(ux,x)=-E(x,ux)\), hence \(E(x,ux)=0\), as the target has
characteristic zero.

Put \(\kappa=1/(2u)\). For \(a+bu\in R\), one has

\[
 \operatorname{Tr}(\kappa(a+bu))=b,\qquad
 \operatorname{Tr}(\kappa u(a+bu))=a.
\]

The trace pairing \((v,w)\mapsto\operatorname{Tr}(\kappa vw)\) is
therefore a perfect \(\mathbf Z_2\)-pairing on \(R\), and

\[
 E(x,y)=\operatorname{Tr}(\kappa\Psi(x,y)).
\]

It follows that \(\Psi:L\longrightarrow\operatorname{Hom}_R(L,R)\) is an
isomorphism: composing with the displayed perfect trace identification
is exactly the given perfect \(\mathbf Z_2\)-duality for \(E\). Choose a
symplectic \(R\)-basis \(x_1,y_1,x_2,y_2\), with
\(\Psi(x_i,y_j)=\delta_{ij}\). Such a basis follows by choosing a
primitive vector \(x_1\), finding \(y_1\) with pairing one using
perfectness, splitting off their span by the usual orthogonal
projection, and repeating.

Let \(\alpha_i,\beta_i,\gamma_i,\delta_i\) be the \(\mathbf Z_2\)-dual
basis to \(x_i,ux_i,y_i,uy_i\). Directly evaluating \(E\) and
\(E(u\,\cdot,\cdot)\) gives

\[
 \theta=m:=\sum_{i=1}^2(\alpha_i\delta_i+\beta_i\gamma_i),
 \qquad
 \theta_u=\omega+\Delta\tau,
 \tag{E807.4}
\]

where

\[
 \omega=\alpha_1\gamma_1+\alpha_2\gamma_2,\qquad
 \tau=\beta_1\delta_1+\beta_2\delta_2.
\]

Take \(S=\langle\alpha_1,\gamma_1,\alpha_2,\gamma_2\rangle\) and
\(T=\langle\beta_1,\delta_1,\beta_2,\delta_2\rangle\). The mixed tensor
\(m\) is again nondegenerate. Write \(\Delta=2h\), with \(h\) odd. The
eigencomponent formula gives

\[
 Q=4C=\theta^2+\frac{\theta_u^2}{\Delta}
   =2m^{[2]}+\frac1h\omega^{[2]}
          +2\omega\tau+4h\tau^{[2]}.
 \tag{E807.5}
\]

This is integral over \(\mathbf Z_2\), its reduction is
\(\omega^{[2]}=\operatorname{vol}_{S/2S}\), and

\[
 (Q_{(2,2)}/2)\bmod2=m^{[2]}+\omega\otimes\tau
\]

has rank six by the same exterior calculation as before. Multiplication
of \(Q\) by the unit \(z\) leaves these two mod-two properties
unchanged. The divided-square lifting lemma applied to \(zQ\) proves
(E807.1).

For the perfect-object assertion, degree-two normalization again forces
\(B=-c_1(E)/r\), and the degree-four comparison gives
\(\operatorname{ch}_2(E)=c_1(E)^2/(2r)-(rq)C\). This contradicts
(E807.1) when \(rq\) is a unit.

Finally suppose \(D|G|\) were odd. S-22 gives

\[
 \frac{4D(rq)^2}{|G|}=k,\qquad k\in\{2,4\}.
\]

For \(k=2\), the valuation of the rational square \((rq)^2\) would be
\(-1\), impossible. For \(k=4\), \(rq\) is a two-adic unit, contrary to
the integral assertion. This proves (E807.2). ∎

The extension allows every integral degree-two class, including extra
classes at CM points. It covers every maximal-at-two lattice for even
squarefree \(\Delta\), and every odd-degree compatible polarization.
SF-E808 below treats the odd squarefree radicands ramified at two. A
nonmaximal order at two and an even value of \(D|G|\) remain outside the
combined obstruction. No eightfold supplier is constructed.

\subsection{SF-E808. Odd radicands ramified at two: an unrestricted
Chern-class
obstruction}\label{SF-E808}

Let \(F=\mathbf Q(\sqrt{\Delta})\), with \(\Delta>1\) squarefree and
\(\Delta\equiv3\pmod4\). Suppose a polarized abelian fourfold
\((X,\theta)\) has Rosati-invariant RM by \(F\), with tangent
multiplicities \((2,2)\), and that \(L=H_1(X,\mathbf Z_2)\) is stable
under the maximal local order \(R=\mathbf Z_2[u]\), \(u^2=\Delta\).
Suppose \(D=\int_X\theta^4/24\) is odd. Define

\[
 C=\frac{\theta_+^2+\theta_-^2}{2}.
\]

\textbf{Assertion.} For every integral class
\(\ell\in H^2(X,\mathbf Z_2)\), every positive integer \(r\), and every
two-adic unit \(z\in\mathbf Q_2\),

\[
 \frac{\ell^2}{2r}-zC\notin H^4(X,\mathbf Z_2).
 \tag{E808.1}
\]

Consequently no positive-rank perfect object, with any rational divisor
normalization, has normalized balanced character \(r v_q\) when \(rq\)
is a two-adic unit. A simple gluable fully semiregular lift with an
effective finite coherent target action \(G\) as in supplied S-22
therefore satisfies \(2\mid D|G|\).

Together with SF-E807, this removes the minimal-divisor-space hypothesis
from supplied S-32 for all real quadratic fields ramified at two.
Maximal-order stability at two and oddness of the relevant degree
product are retained.

\subsubsection{The local polarized RM
lattice}\label{the-local-polarized-rm-lattice}

The element \(u-1\) satisfies

\[
 t^2+2t+1-\Delta=0.
\]

This polynomial is Eisenstein at two, because \(1-\Delta\) has valuation
one. Thus \(R\) is a discrete valuation ring with uniformizer
\(\pi=u-1\). Concretely \(R=\mathbf Z_2+\mathbf Z_2\pi\), and the
valuation of \(a+b\pi\) is \(\min\{2v_2(a),2v_2(b)+1\}\); the two terms
have distinct valuation parity. This identifies \(R\) with the valuation
ring. The finite torsion-free \(R\)-module \(L\) is free of rank four.

Let \(E\) be the polarization pairing on \(L\). Its determinant is odd,
so it is perfect over \(\mathbf Z_2\). Exactly as in the full local
proof of SF-E807, the form

\[
 \Psi(x,y)=E(x,uy)+uE(x,y)
\]

is alternating and \(R\)-bilinear, and it is perfect because

\[
 E(x,y)=\operatorname{Tr}_{(F\otimes\mathbf Q_2)/\mathbf Q_2}
              \left(\frac{\Psi(x,y)}{2u}\right).
\]

For completeness, the trace pairing on \(R\) is perfect: it sends
\(a+bu\) to \(b\), and multiplication first by \(u\) sends it to \(a\).
This identifies the \(R\)-linear dual with the \(\mathbf Z_2\)-linear
dual. A symplectic \(R\)-basis \(x_i,y_i\), \(i=1,2\), follows by
splitting off pairs with pairing one.

Let \(\alpha_i,\beta_i,\gamma_i,\delta_i\) be the dual integral basis to
\(x_i,ux_i,y_i,uy_i\). Put

\[
 \begin{aligned}
 S&=\langle\alpha_1,\gamma_1,\alpha_2,\gamma_2\rangle,
 &T&=\langle\beta_1,\delta_1,\beta_2,\delta_2\rangle,\\
 \omega&=\alpha_1\gamma_1+\alpha_2\gamma_2,
 &\tau&=\beta_1\delta_1+\beta_2\delta_2,\\
 m&=\alpha_1\delta_1+\beta_1\gamma_1
       +\alpha_2\delta_2+\beta_2\gamma_2.
 \end{aligned}
\]

Evaluation of the pairings gives \(\theta=m\) and
\(\theta_u=\omega+\Delta\tau\). Hence

\[
 Q:=2C=m^{[2]}+\frac1\Delta\omega^{[2]}
            +\omega\tau+\Delta\tau^{[2]}.
 \tag{E808.2}
\]

This is an integral two-adic four-form. Its reduction is independent of
the odd radicand:

\[
 \bar Q=m^{[2]}+\omega^{[2]}+\omega\tau+\tau^{[2]}.
 \tag{E808.3}
\]

Here the displayed operations are divided powers and exterior products
reduced from the integral lattice; signs disappear modulo two.

\subsubsection{A contraction with no wedge
annihilator}\label{a-contraction-with-no-wedge-annihilator}

Write \(s_0=\alpha_1,s_1=\gamma_1,s_2=\alpha_2,s_3=\gamma_2\), and
\(t_0=\beta_1,t_1=\delta_1,t_2=\beta_2,t_3=\delta_2\). Let \(\lambda\)
be the dual functional selecting \(s_0\). Direct contraction of (E808.3)
gives

\[
 \chi=\iota_\lambda\bar Q
 =s_1s_2s_3+s_1t_2t_3+s_2t_1t_3+s_3t_1t_2.
 \tag{E808.4}
\]

The map

\[
 \bar L^*\longrightarrow\bigwedge^4\bar L^*,\qquad
 v\longmapsto v\chi
 \tag{E808.5}
\]

is injective. This is a direct monomial calculation, not a generic-rank
assertion. Rename the six variables appearing in \(\chi\) as
\(a,b,c,d,e,f\), so \(\chi=abc+aef+bdf+cde\). Their six exterior
products are

\[
 \begin{array}{c|l}
 a&abdf+acde\\
 b&abef+bcde\\
 c&acef+bcdf\\
 d&abcd+adef\\
 e&abce+bdef\\
 f&abcf+cdef .
 \end{array}
\]

The twelve monomials are distinct, proving independence of these six
images. Exterior multiplication by either of the remaining variables
\(s_0,t_0\) gives a nonzero form containing that variable. These two
images and the preceding six belong to separate summands according to
their \(s_0,t_0\) content. This proves injectivity on the entire
eight-dimensional space.

In particular \(\bar Q\ne0\), so \(Q\) is primitive at two.

\subsubsection{A divided-square lemma with denominator
two}\label{a-divided-square-lemma-with-denominator-two}

Let \(M\) be a free \(\mathbf Z_2\)-module, let \(Q\in\bigwedge^4M\),
and suppose some functional \(\lambda\in(M/2M)^*\) has the property that
exterior multiplication by \(\chi=\iota_\lambda\bar Q\) is injective on
\(M/2M\). Then there are no \(\ell\in\bigwedge^2M\) and positive integer
\(r\) with

\[
 \ell^{[2]}-\frac r2Q\in r\bigwedge^4M.
 \tag{E808.6}
\]

\textbf{Proof.} The hypothesis makes \(Q\) primitive, so integrality
forces \(r\) even. Write \(r=2^e k\), with \(k\) odd and \(e\ge1\). The
alleged congruence is

\[
 \ell^{[2]}\equiv 2^{e-1}kQ\pmod{2^e}.
 \tag{E808.7}
\]

The minimum valuation of its left side is exactly \(e-1\). Removing the
largest common factor \(2^a\) from \(\ell\) reduces the congruence to
the same one for a primitive two-form and exponent \(e'=e-2a\ge1\). We
may assume that \(\ell\) is primitive.

For \(e=1\), the congruence says \(\bar\ell^{[2]}=\bar Q\). The
contraction identity for divided squares gives

\[
 \chi=(\iota_\lambda\bar\ell)\bar\ell.
\]

If \(\iota_\lambda\bar\ell=0\), this contradicts \(\chi\ne0\). If it is
nonzero, it wedges to zero with \(\chi\), contradicting the assumed
injectivity. Thus \(e=1\) is impossible.

For \(e\ge2\), the alternating elementary-divisor algorithm gives

\[
 \ell=\alpha+2^{e-1}\beta,
\]

where \(\alpha\) is primitive and decomposable. Indeed primitivity makes
the first elementary divisor a unit, while the valuation of the divided
square makes the second elementary divisor have valuation \(e-1\); all
subsequent divisors have at least that valuation. Dividing (E808.7) by
\(2^{e-1}\) and reducing modulo two gives

\[
 \bar Q=\bar\alpha\bar\beta.
\]

The two-dimensional support \(P\) of \(\bar\alpha\) therefore wedges to
zero with \(\bar Q\). There is a nonzero \(v\in P\cap\ker\lambda\).
Contracting \(v\bar Q=0\) by \(\lambda\) gives

\[
 0=\lambda(v)\bar Q+v\iota_\lambda\bar Q=v\chi,
\]

which again contradicts injectivity. Both cases are impossible. ∎

\subsubsection{Chern-character consequence and
scope}\label{chern-character-consequence-and-scope}

Multiplication by the unit \(z\) in (E808.1) leaves \(\bar Q\) and its
contraction unchanged. The lemma applied to \(zQ\), with \(Q=2C\),
proves (E808.1). For a normalized balanced character, degree two forces
\(B=-c_1(E)/r\), while degree four gives

\[
 \operatorname{ch}_2(E)=\frac{c_1(E)^2}{2r}-(rq)C.
\]

The Chern character on the left is integral on a complex torus, since
\(\operatorname{ch}_2=c_1^2/2-c_2\) and every integral two-form has an
integral divided square. Hence \(rq\) cannot be a two-adic unit.

If a simple gluable fully semiregular equivariant lift had \(D|G|\) odd,
supplied S-22 would give \(4D(rq)^2/|G|\in\{2,4\}\). The value two has
impossible odd valuation for a rational square; the value four makes
\(rq\) a two-adic unit. Thus \(2\mid D|G|\).

No restriction to the RM divisor plane was made: every integral
\(H^2\)-class is included, so extra CM divisor classes cannot evade the
conclusion. This does not exclude even-degree polarizations, even-order
effective finite actions, nonmaximal local orders, or other proposed
characters. It constructs no eightfold supplier.

\subsubsection{Source and check status}\label{source-and-check-status}

The proof uses only the supplied S-22 Euler restriction for the final
semiregularity corollary. The integral obstruction, local pairing
calculation, contraction matrix, and lifting argument are proved in full
above. S-32 is a comparison result, not an input to this proof. The
finite certificate checks exact normal-form substitutions and the
contraction rank; it does not replace the proof for every rank or every
radicand.

\subsection{Sources, checks, and remaining
scope}\label{sources-checks-and-remaining-scope}

\begin{itemize}
\item
  DEG\_L S-21--S-25 and S-30--S-34 are supplied background. Their CY4
  adjunction, identity-sector assertion and previously proved
  obstructions are retained with their existing status.
\item
  Orlov, \emph{Derived categories of coherent sheaves on Abelian
  varieties and equivalences between them}, Izv. Math.66(3),569--594
  (2002), Theorem2.10 and Corollary2.13. Publication verified at
  https://www.mathnet.ru/eng/im389 . Full theorem proof and corollary
  proof read in the author manuscript
  https://homepage.mi-ras.ru/\textasciitilde orlov/papers/abelnew.pdf ,
  pp9--12. The publisher PDF was located but its subsequent fetch
  failed; no unverified printed-page locator is asserted. This published
  result supplies the abelian-group transport, while the kernel
  reconstruction above is a new proof.
\item
  \texttt{eightfold\_kernel\_certificate.py} is an exact rational
  exterior-algebra check with its scope stated in its docstring and
  JSON. It is not a proof assistant or a sheaf-construction certificate.
\item
  Golyshev--Lunts--Orlov, \emph{Mirror symmetry for abelian varieties},
  J. Algebraic Geom.10(3),433--496 (2001), Proposition4.3.7 and
  Lemma4.3.7.1. Published work; complete relevant proof inspected in
  https://homepage.mi-ras.ru/\textasciitilde orlov/papers/GLO.pdf ,
  author-version pp17--19. The exact cohomological Clifford intertwining
  is the input used in SF-E805. Printed theorem pagination was not
  independently verified.
\item
  The new integral assertions SF-E806--SF-E808 are proved by elementary
  exterior algebra and the local alternating-module algorithms
  explicitly described above. Their semiregularity corollaries use
  supplied DEG\_L S-22, equations (S22.2)--(S22.3). The comparison with
  S-32 respects that theorem's stated minimal-divisor-space hypothesis.
\item
  The exact arithmetic certificate for SF-E806 checks the integral
  change of coordinates, the primitive four-form, the rank-six
  obstruction tensor, and all twenty-eight mod-two nondegenerate
  alternating forms in the first lifting case. The proof for arbitrarily
  large ranks is the two-adic argument above, not a bounded computation.
\item
  No split-eightfold CM input is constructed here. The original nonsplit
  sixfold gate remains open, narrowed by SF-E802--E808. In particular,
  the success of SF-QC01--05 for split sixfolds does not imply this
  remaining eightfold supplier.
\end{itemize}

\section{The complete-intersection plus points secant: an actual-object
obstruction}\label{the-complete-intersection-plus-points-secant-an-actual-object-obstruction}

Author/checker: GPT-6 Astra Pro; inherited effort not exposed;
HODGE\_DIVISION \texttt{/root/sixfolds}. Date: 2026-11-21.

\textbf{Status.} Independently whole-proof reviewed by \texttt{/root}
and \texttt{/root/imprimitive} on 2026-11-21: both GPT-6 Astra Pro,
inherited effort not exposed; PASS for SF-E809--E810, including the
Koszul and full trace ranks, invariant kernel and same-character Euler
equality. The exact computation in \texttt{ci\_secant\_certificate.py}
has actually run and passed. The sibling read every script line but did
not rerun it, and no sibling computation is claimed. It checks the
finite exterior algebra and the character/Euler identities; it does not
substitute for any geometric or categorical argument below. Canonical
deliverables are \texttt{ci\_secant\_obstruction\_proofs.md} and
\texttt{ci\_secant\_obstruction\_section.tex}; this draft-path copy is
retained so earlier review links remain usable.

These results concern a genuine rank-one ideal sheaf whose Chern
character is a real quadratic two-exponential secant. This character is
\textbf{not} the balanced real-multiplication product character of
DEG\_L S-22: their degree-two Hochschild contraction ranks are
respectively 12 and 18. No genuine RM eightfold seed is asserted or
presumed.

\subsection{SF-E809. Exact semiregularity kernel of the proposed
ideal}\label{SF-E809}

Let \((B,\theta)\) be a principally polarized complex abelian fourfold.
Let \(0<a<b\) be integers and suppose that divisors \(D_a\in|a\theta|\),
\(D_b\in|b\theta|\) meet transversely in a smooth surface

\[
 S=D_a\cap D_b.
\]

Let \(Z\) be a reduced set of \(N>0\) points, disjoint from \(S\), and
put \(F=I_{S\sqcup Z}\). The assertion is conditional only on this
stated smoothness; for example, general divisors in
\(|2\theta|,|4\theta|\) have the required property. Write

\[
 \sigma_F^j:\operatorname{Ext}^j(F,F)\longrightarrow HH_{-j}(B)
\]

for the full Atiyah--Chern trace. Then

\[
 \begin{aligned}
 \dim\operatorname{Ext}^1(F,F)
   &=a^4+b^4-(b-a)^4+10+4N,\\
 \dim\operatorname{Ext}^2(F,F)&=18+6N,\\
 \operatorname{rank}\sigma_F^2&=18,\\
 \dim\ker\sigma_F^2&=6N.
 \end{aligned}
 \tag{E809.1}
\]

More strongly, if a finite subgroup \(G\subset B\times\widehat B\)
stabilizes \(F\), with any coherent equivariant structure, and if \(r\)
is the number of its orbits on \(Z\), then

\[
 \dim\ker\left(\sigma_F^2|_{\operatorname{Ext}^2(F,F)^G}\right)=6r>0.
 \tag{E809.2}
\]

Consequently this ideal is not weakly \(G\)-semiregular for any such
finite group. Its external product with its derived dual fails the same
condition for every finite translation/tensor stabilizer on
\(B\times B\), including correlated actions on the two factors. A
Fourier--Mukai equivalence, line-bundle twist or shift cannot remove
this invariant trace kernel.

\subsubsection{1. The complete-intersection self-Ext
groups}\label{the-complete-intersection-self-ext-groups}

The two-term Koszul resolution of \(I_S\) is

\[
 0\longrightarrow\mathcal O_B(-(a+b)\theta)
 \longrightarrow\mathcal O_B(-a\theta)\oplus\mathcal O_B(-b\theta)
 \longrightarrow I_S\longrightarrow0.
 \tag{E809.3}
\]

Its sheaf endomorphism complex has terms

\[
 \begin{array}{c|c}
 \text{degree}&\text{term}\\ \hline
 -1&\mathcal O_B(-b\theta)\oplus\mathcal O_B(-a\theta)\\
 0&\mathcal O_B^{\oplus3}\oplus\mathcal O_B((b-a)\theta)
                          \oplus\mathcal O_B(-(b-a)\theta)\\
 1&\mathcal O_B(b\theta)\oplus\mathcal O_B(a\theta).
 \end{array}
 \tag{E809.4}
\]

For a positive integer \(m\), the cohomology of
\(\mathcal O_B(m\theta)\) is concentrated in degree zero and has
dimension \(m^4\); the negative line bundle has cohomology only in
degree four. Also \(h^q(\mathcal O_B)=\binom4q\). These statements
follow from Kodaira vanishing, Serre duality, and Riemann--Roch with
trivial Todd class.

In the hypercohomology spectral sequence of (E809.4), the sole nonzero
term of total degree two is the degree-two cohomology of the three
trivial summands in degree zero. No higher differential enters or leaves
this term, since the complex has only columns \(-1,0,1\). Thus

\[
 \operatorname{Ext}^2(I_S,I_S)\simeq H^2(\mathcal O_B)^{\oplus3},
 \qquad \dim\operatorname{Ext}^2(I_S,I_S)=18.
 \tag{E809.5}
\]

For total degree one there are the 12 dimensions in
\(H^1(\mathcal O_B)^{\oplus3}\) and the cokernel on global sections of
the map from degree zero to degree one. Its kernel has dimension one:
every global endomorphism of a rank-one ideal with codimension-two
support is scalar, since it is a regular function off that support and
extends over the normal variety \(B\). The cokernel therefore has
dimension

\[
 a^4+b^4-\bigl(3+(b-a)^4\bigr)+1.
\]

There are again no relevant higher differentials. Hence

\[
 \dim\operatorname{Ext}^1(I_S,I_S)
   =a^4+b^4-(b-a)^4+10.
 \tag{E809.6}
\]

\subsubsection{2. The complete-intersection trace is injective in degree
two}\label{the-complete-intersection-trace-is-injective-in-degree-two}

We prove this from the full trace/cap identity, not from a termwise
diagonal trace in a chosen resolution. Put

\[
 W=H^0(B,T_B)\oplus H^1(B,\mathcal O_B),\qquad\dim W=8.
\]

Choose a holomorphic volume form. As in SF-E801, it identifies
\(HH_{-j}(B)\) with \(\bigwedge^{4+j}W\) and makes Hochschild cap
product exterior multiplication. Write \(E=H^1(B,\mathcal O_B)\), and
choose bases \(e_1,\ldots,e_4\) of \(E\) and \(p_1,\ldots,p_4\) of
\(H^0(B,T_B)\) so that contraction with \(\theta\) sends \(p_i\) to
\(e_i\). Up to one fixed nonzero normalization, the character
\(e^{t\theta}\) is the decomposable four-vector

\[
 v_t=\bigwedge_{i=1}^4(p_i-t e_i).
 \tag{E809.7}
\]

It is enough to compute the rank of multiplication by

\[
 v=v_{-a}+v_{-b}-v_{-(a+b)},
\]

because the Koszul resolution gives

\[
 \operatorname{ch}(I_S)=e^{-a\theta}+e^{-b\theta}-e^{-(a+b)\theta}.
 \tag{E809.8}
\]

There is a common 10-dimensional kernel in \(\bigwedge^2W\), namely

\[
 K=\left\langle e_i\wedge p_i,\;
 e_i\wedge p_j+e_j\wedge p_i\ (i<j)\right\rangle.
 \tag{E809.9}
\]

Indeed, on the quotient by the support of \(v_t\) we have \(p_i=t e_i\).
A complementary 18-dimensional space has, for each \(i<j\), the three
basis vectors

\[
 e_i\wedge e_j,\quad
 \frac{e_i\wedge p_j-e_j\wedge p_i}{2},\quad
 p_i\wedge p_j.
\]

Their images in that quotient are respectively \(1,t,t^2\) times
\(e_i\wedge e_j\). Consequently the symmetric bilinear form

\[
 (z,w)\longmapsto z\wedge w\wedge v
\]

on this complement is the tensor product of the nondegenerate exterior
pairing on \(\bigwedge^2E\) with

\[
 H=\sum_{(t,\lambda)=(-a,1),(-b,1),(-(a+b),-1)}
       \lambda\begin{pmatrix}1\\t\\t^2\end{pmatrix}
                    \begin{pmatrix}1&t&t^2\end{pmatrix}.
\]

The graph-to-quotient determinant in this calculation is one,
independent of \(t\). The Vandermonde determinant gives

\[
 \det H=-a^2b^2(a-b)^2\ne0.
 \tag{E809.10}
\]

Thus multiplication by \(v\) from \(\bigwedge^2W\) to \(\bigwedge^6W\)
has rank 18 and kernel exactly \(K\). Its image is

\[
 \mathcal T=K^\perp\subset\bigwedge^6W,
 \qquad\dim\mathcal T=18,
 \tag{E809.11}
\]

where orthogonality uses the perfect top exterior pairing. The
six-dimensional subspace

\[
 H^4(B,\Omega_B^2)=\bigwedge^4E\otimes\bigwedge^2\langle p_1,\ldots,p_4\rangle
\]

lies in \(\mathcal T\): every element of \(K\) has an \(e_i\) factor,
which wedges to zero with \(\bigwedge^4E\).

Let \(\rho^2\) denote Hochschild evaluation on the object. The
Atiyah--Chern identity is

\[
 \sigma_{I_S}^2\rho_{I_S}^2
       =(-)\mathbin{\lrcorner}\operatorname{ch}(I_S).
 \tag{E809.12}
\]

The right-hand side has rank 18, and its middle Ext group has dimension
18 by (E809.5). Therefore \(\sigma_{I_S}^2\) is injective and its image
is exactly \(\mathcal T\).

\subsubsection{3. Disjoint support and the local point
summands}\label{disjoint-support-and-the-local-point-summands}

For any rank-one ideal \(I_T\) with support of codimension at least two,
the unit and trace split its derived endomorphism complex as

\[
 R\mathcal Hom(I_T,I_T)=\mathcal O_B\oplus\mathcal E_T,
 \tag{E809.13}
\]

because trace composed with the unit is multiplication by the rank,
which is one. Moreover \(\mathcal Hom(I_T,I_T)=\mathcal O_B\), so all
cohomology sheaves of \(\mathcal E_T\) have positive degree and are
supported on \(T\).

A coherent derived object supported on a union of disjoint closed
subsets splits into its objects supported on those subsets. For coherent
sheaves this is the support decomposition, or equivalently the Chinese
remainder decomposition on a sufficiently high infinitesimal
neighborhood of the support. The decomposition extends through the
bounded derived category, since there are no Hom or Ext groups between
objects with disjoint supports. It does not require disjoint Zariski
open neighborhoods. Applied to \(S\sqcup Z\), and using that the ideals
coincide near their respective supports, this gives

\[
 \mathcal E_{S\sqcup Z}\simeq\mathcal E_S\oplus
                   \bigoplus_{z\in Z}\mathcal E_z.
 \tag{E809.14}
\]

At a smooth point of a fourfold, the truncated Koszul resolution of the
maximal ideal gives the trace-free local self-Ext spaces

\[
 \mathcal H^j(\mathcal E_z)=\bigwedge^j T_zB,
 \quad j=1,2,3.
 \tag{E809.15}
\]

One can also obtain (E809.15) by applying Hom to
\(0\to I_z\to\mathcal O_B\to\mathcal O_z\to0\) and the full Koszul
resolution of \(\mathcal O_z\). The first three local tangent-space
exterior powers have dimensions 4,6,4. The degree-four map
\(\operatorname{Ext}^4(\mathcal O_z,\mathcal O_B)\to\operatorname{Ext}^4(\mathcal O_z,\mathcal O_z)\)
is an isomorphism in the local Koszul comparison, so no fourth local
self-Ext remains. Since these sheaves have zero-dimensional support,
there is no higher sheaf cohomology. Thus (E809.14) proves the two
self-Ext dimensions in (E809.1).

We now identify the trace images, a point needed for the invariant
statement. The trace on the \(S\)-supported summand agrees with that for
\(I_S\), by restriction and locality of the Atiyah class. More
explicitly, any two extensions of the restricted trace differ by a
morphism supported on \(Z\), and a source supported on \(S\) has no such
morphism. On the shared \(H^2(\mathcal O_B)\) summand, the trace is
multiplication by the Chern character. Replacing \(I_S\) by \(F\)
subtracts \(N[\mathrm{pt}]\), whose product with \(H^2(\mathcal O_B)\)
vanishes for dimensional reasons. Therefore the full 18-dimensional
trace image \(\mathcal T\) of \(I_S\) persists unchanged inside the
trace image of \(F\).

For a single point ideal, evaluation of a bivector in
\(\bigwedge^2H^0(T_B)\) has trace equal to contraction with \(-[z]\), by
(E809.12). This is an isomorphism onto \(H^4(\Omega_B^2)\). Its ordinary
trace component in \(H^2(\mathcal O_B)\) is zero, so evaluation lies in
the trace-free local summand. That summand has dimension six by
(E809.15); its full trace is consequently an isomorphism onto
\(H^4(\Omega_B^2)\), with no other trace components. Locality gives the
same statement for each point in (E809.14).

All the point images are therefore already in \(\mathcal T\). The total
trace has image exactly \(\mathcal T\), rank 18, and kernel dimension
\(6N\), proving (E809.1).

\subsubsection{4. Every finite translation/tensor stabilizer retains a
kernel}\label{every-finite-translationtensor-stabilizer-retains-a-kernel}

The determinant of \(F\) is trivial and its rank is one. If

\[
 t_x^*F\otimes P\simeq F,\qquad P\in\widehat B,
\]

then determinants give \(P\simeq\mathcal O_B\). Thus every finite
stabilizer in \(B\times\widehat B\) consists of translations. The
translation preserves the subscheme \(S\sqcup Z\), hence preserves \(S\)
and permutes \(Z\), since their dimensions differ.

Translations act trivially on Hodge cohomology. Since \(\sigma_{I_S}^2\)
is injective and natural for translations preserving \(S\), they act
trivially on \(\operatorname{Ext}^2(I_S,I_S)\). On the point summands
they permute the points and identify their tangent spaces by the
identity translation derivative. Scalar choices in an equivariant
structure cancel in the conjugation action on self-Ext. The splitting
above is natural for these actions.

It follows that the invariant point contribution has dimension \(6r\),
while the invariant complete-intersection contribution still maps
isomorphically to \(\mathcal T\). For each invariant point vector, its
trace has a unique preimage in that complete-intersection contribution,
and subtracting this preimage gives an invariant trace-kernel vector.
Conversely these are all the kernel vectors, since the
complete-intersection trace is injective. This proves (E809.2).

For \(F\boxtimes F^\vee\), the Kunneth summand

\[
 \operatorname{Ext}^2(F,F)\otimes\operatorname{Hom}(F^\vee,F^\vee)
\]

contains these nonzero vectors tensored with the identity, and their
full trace is zero by the tensor-product trace formula. For
completeness, a finite stabilizer of the product again has trivial
flat-line component by its rank-one trivial determinant. Its two
translation projections preserve \(F\) and \(F^\vee\), as follows by
restricting a product isomorphism to a generic fiber away from the
relevant supports. Such an isomorphism factors up to a scalar into
isomorphisms on the two factors: both factors are simple and have no
negative self-Ext, so the degree-zero Kunneth Hom space is their
one-dimensional tensor product. Hence the invariant vector just
constructed for the first projected group, tensored with the identity on
the second factor, is invariant for the whole group, even when its two
projections are correlated.

An equivalence identifies self-Ext, Hochschild homology, evaluation and
trace; for abelian varieties it also conjugates the translation/tensor
group. Thus Fourier--Mukai equivalences preserve the stated failure for
actual finite stabilizers. Line-bundle twists and shifts are special
cases. No assertion about an arbitrary new equivariant object after an
unrelated isogeny is needed or made. This completes the proof of
SF-E809.

\subsection{SF-E810. A numerical obstruction to every ordinary
same-character
repair}\label{SF-E810}

Let \(F\) be a simple gluable perfect complex on an abelian fourfold.
Here gluable means \(\operatorname{Ext}^{j}(F,F)=0\) for \(j<0\).
Suppose that, after the Fock identification above, its full Chern
character is a nondegenerate two-pure secant

\[
 v=\lambda\operatorname{vol}_L+\mu\operatorname{vol}_M,
 \quad W=L\oplus M,\quad\dim L=\dim M=4,\quad\lambda\mu\ne0.
 \tag{E810.1}
\]

If \(\sigma_F^2\) is injective, then

\[
 \dim\operatorname{Ext}^2(F,F)=12,\qquad
 \operatorname{rank}\sigma_F^1=8,\qquad
 \dim\ker\sigma_F^1=\frac{-\chi(F,F)-2}{2}.
 \tag{E810.2}
\]

In particular, simultaneous injectivity of \(\sigma_F^1\) and
\(\sigma_F^2\) forces

\[
 \chi(F,F)=-2.
 \tag{E810.3}
\]

The ordinary external product \(F\boxtimes F^\vee\) cannot be fully
semiregular unless (E810.3) holds. These statements concern ordinary
full semiregularity; for an arbitrary replacement object they do not
exclude a new finite equivariant action whose invariant Ext groups are
smaller.

\subsubsection{Proof of the character-level
assertion}\label{proof-of-the-character-level-assertion}

For a Calabi--Yau fourfold, the degree-two Hochschild evaluation
\(\rho_F^2\) and full trace \(\sigma_F^2\) are adjoint under Serre
duality and the perfect Hochschild pairing. This is the same
evaluation/trace adjunction used in supplied DEG\_L S-21 and in SF-E802.
Injectivity of \(\sigma_F^2\) therefore makes \(\rho_F^2\) surjective.
By the full trace/cap identity and SF-E801,

\[
 \dim\operatorname{Ext}^2(F,F)=\operatorname{rank}c_v^2=12,
 \qquad\ker\rho_F^2=\ker c_v^2=L\otimes M.
 \tag{E810.4}
\]

Write \(I=L\otimes M\). For \(\xi\in I\) and
\(\alpha\in\operatorname{Ext}^1(F,F)\), compatibility of Hochschild
action with trace gives

\[
 \xi\wedge\sigma_F^1(\alpha)
   =\sigma_F^3\bigl(\rho_F^2(\xi)\circ\alpha\bigr)=0.
 \tag{E810.5}
\]

The common annihilator of \(I\) in \(\bigwedge^5W\) is exactly

\[
 \mathcal A=
 \bigl(\bigwedge^4L\otimes M\bigr)
 \oplus\bigl(L\otimes\bigwedge^4M\bigr),\qquad\dim\mathcal A=8.
 \tag{E810.6}
\]

To see this, decompose a five-vector by bidegrees
\((4,1),(3,2),(2,3),(1,4)\). The extreme components are killed by every
mixed bivector. For either interior component, the collective maps
\(\bigwedge^p L\to\operatorname{Hom}(L,\bigwedge^{p+1}L)\) and
\(\bigwedge^{5-p}M\to\operatorname{Hom}(M,\bigwedge^{6-p}M)\), given by
wedging with every possible additional vector, are injective, because
their source exterior degrees are strictly below four. Their tensor
product is injective as well. Thus a nonzero interior component is
detected by some mixed bivector. Components of different bidegrees
cannot cancel. This proves (E810.6).

Equation (E810.5) gives rank at most eight for \(\sigma_F^1\). On the
other hand \(c_v^1:W\to\bigwedge^5W\), \(x\mapsto x\wedge v\), is
injective: its two summands occupy the distinct bidegrees \((1,4)\),
\((4,1)\) and both coefficients in (E810.1) are nonzero. Since
\(c_v^1=\sigma_F^1\rho_F^1\), the trace has rank exactly eight.

Gluability and Serre duality give no self-Ext outside degrees zero
through four. Simplicity gives \(e_0=e_4=1\); Serre duality gives
\(e_3=e_1\); (E810.4) gives \(e_2=12\). Therefore

\[
 \chi(F,F)=14-2e_1,
 \qquad\dim\ker\sigma_F^1=e_1-8=\frac{-\chi(F,F)-2}{2}.
\]

This proves (E810.2)--(E810.3).

Finally, full semiregularity of \(F\boxtimes F^\vee\) forces injectivity
of \(\sigma_F^2\), by its factor-Ext Kunneth summand. It also forces
injectivity of \(\sigma_F^1\): if a nonzero \(\alpha\) is in the latter
kernel, tensor it with any nonzero element of
\(\operatorname{Ext}^1(F^\vee,F^\vee)\), which exists since \(e_1\ge8\).
The tensor is a nonzero degree-two product-Ext vector whose trace
vanishes. This proves the external-product assertion.

\subsubsection{Application to the tuned complete-intersection
character}\label{application-to-the-tuned-complete-intersection-character}

Put \(p=ab\) and \(s=a+b\). The character of \(I_S\) is

\[
 1-\frac{A}{2}\theta^2+\frac{B}{6}\theta^3
        +\frac{C_{\rm actual}}{24}\theta^4,
 \quad A=2p,\quad B=3ps,\quad
 C_{\rm actual}=-4p(a^2+b^2)-6p^2.
 \tag{E810.7}
\]

A normalized two-exponential moment sequence with rank one and first
Chern class zero must have

\[
 C_{\rm required}=A^2-\frac{B^2}{A}.
\]

The difference is

\[
 C_{\rm actual}-C_{\rm required}=\frac{ab(a-b)^2}{2}=N.
 \tag{E810.8}
\]

This is an integer for every pair of integers \(a,b\): if \(ab\) is odd,
then \(a-b\) is even. Thus no extra parity hypothesis on \(ab\) is
necessary. Since \(\theta^4/24=[\mathrm{pt}]\), adding exactly these
\(N\) disjoint points gives the required character.

Let \(u,v\) be the roots with

\[
 uv=A,\qquad u+v=-B/A=-3s/2.
\]

They are distinct, nonzero, and real, with squared difference

\[
 D=(u-v)^2=\frac{9a^2-14ab+9b^2}{4}
     =ab+\frac94(a-b)^2>0.
 \tag{E810.9}
\]

The weights \(\lambda=-v/(u-v)\), \(\mu=u/(u-v)\) are nonzero, sum to
one and give first moment zero. The recurrence of their moment sequence
gives exactly the coefficients in (E810.7) with \(C_{\rm required}\).
Hence

\[
 \operatorname{ch}(I_{S\sqcup Z})=\lambda e^{u\theta}+\mu e^{v\theta}.
 \tag{E810.10}
\]

The two graph supports in the Fock space are transverse because
\(u-v\ne0\) and \(\theta\) is nondegenerate. Thus (E810.1) applies,
whether or not \(D\) is a rational square. For \((a,b)=(2,4)\),
\(N=16\), \(D=17\), and this is an actual irrational real-quadratic
secant.

Hirzebruch--Riemann--Roch, again with trivial Todd class, gives for
every object with this same character

\[
 \begin{aligned}
 \chi(F,F)
 &=2C_{\rm required}+6A^2
   =8A^2-2B^2/A\\
 &=-4abD
   =-ab\bigl(4ab+9(a-b)^2\bigr).
 \end{aligned}
 \tag{E810.11}
\]

For unequal positive integers this is at most \(-34\), attained at
\((a,b)=(1,2)\). Therefore \textbf{no simple gluable replacement with
the same full Chern character can have both \(\sigma^1\) and
\(\sigma^2\) injective}, and none can yield an ordinarily fully
semiregular external product with its dual. The numerical examples are

\[
 \begin{array}{c|r|r|r|r|r}
 (a,b)&N&\chi(F,F)&\dim\operatorname{Ext}^1(I_{S\sqcup Z})&
 \dim\operatorname{Ext}^2(I_{S\sqcup Z})&\dim\ker\sigma^2\\ \hline
 (1,2)&1&-34&30&24&6\\
 (2,4)&16&-544&330&114&96.
 \end{array}
\]

The first row is conditional on the stated smooth complete intersection;
the second exists for general divisors by basepoint freeness and
Bertini. These numbers also satisfy \(\chi=2-2e_1+e_2\), an independent
consistency check.

\subsection{Exact scope and remaining
directions}\label{exact-scope-and-remaining-directions}

SF-E809 rules out the actual complete-intersection-plus-reduced-points
object, its ordinary external product and every finite
translation/tensor invariant version of those objects. SF-E810 rules out
all \textbf{simple gluable ordinary} same-character repairs of the
external-product strategy. It does not rule out an arbitrary equivariant
replacement with smaller invariant Ext spaces, a nonsimple
indecomposable object, a different full character, or the balanced RM
product character with an even local denominator/group factor. Those are
distinct remaining gates.

\subsection{Primary inputs and proof
provenance}\label{primary-inputs-and-proof-provenance}

\begin{enumerate}
\def\labelenumi{\arabic{enumi}.}
\item
  Buchweitz--Flenner, \emph{A semiregularity map for modules and
  applications to deformations}, Compositio Math. 137 (2003), 135--210,
  §4.1 and Proposition 4.2, pp.166--167, DOI
  \texttt{10.1023/A:1023999012081}: exponential Atiyah trace and
  contraction compatibility. Published PDF:
  \url{https://www.cambridge.org/core/services/aop-cambridge-core/content/view/712BFB5F73E0C38B4B32206699B0277F/S0010437X03100413a.pdf/semiregularity_map_for_modules_and_applications_to_deformations.pdf}.
\item
  Buchweitz--Flenner, \emph{The global decomposition theorem for
  Hochschild (co-)homology of singular spaces via the Atiyah--Chern
  character}, Adv. Math. 217 (2008), 243--281, DOI
  \texttt{10.1016/j.aim.2007.06.013}; author-version Proposition 6.2.1
  and Theorem 6.4.2, author pp.29,31: Hochschild factorization,
  evaluation/trace and module compatibility. The published work, not an
  unrefereed theorem, is being invoked; exact author-copy locators are
  distinguished from journal pagination. Author PDF:
  \url{https://arxiv.org/pdf/math/0606730}.
\item
  The Calabi--Yau evaluation/trace adjunction is the categorical duality
  already proved/used in supplied DEG\_L S-21 and reviewed SF-E802. The
  new common-annihilator argument in degree five and the resulting Euler
  equality are proved here in full.
\item
  Orlov, \emph{Derived categories of coherent sheaves on abelian
  varieties and equivalences between them}, Izv. Math. 66 (2002),
  Theorem 2.10 and Corollary 2.13 (author pp.9--12), supplies
  conjugation of the actual translation/tensor groups under derived
  equivalences; those statements and the corollary proof were read in
  the earlier SF-E803 audit. The failure of the product before any
  equivalence is proved here directly.
\item
  The exact Koszul, Vandermonde, support splitting, invariant-kernel and
  Euler arguments above are new deductions in this work. They invoke no
  Markman or other preprint assertion.
\end{enumerate}

\section{A curvature obstruction to stable balanced-RM bundles in every
rank}\label{a-curvature-obstruction-to-stable-balanced-rm-bundles-in-every-rank}

Author: GPT-6 Astra Pro; inherited effort not exposed; HODGE\_DIVISION
\texttt{/root/sixfolds}. Date: 2026-11-21.

\textbf{Status:} SF-E811, SF-E811-RM and SF-E812 received whole-proof
mathematical PASS reviews from both \texttt{/root} and
\texttt{/root/imprimitive} on 2026-11-21 (GPT-6 Astra Pro; inherited
effort not exposed), after independent reconstruction of the complete
arguments. The root-requested rational-Kähler-class source restriction
is incorporated. The peer used the explicitly identified Siu existence
theorem and previously audited supplied S-21; no independent peer
reopening of Siu is claimed. The exact generic curvature certificate was
actually run and passed. This canonical proof concerns slope-stable or
slope-polystable \textbf{locally free} sheaves in the respective stated
scopes. It makes no assertion about arbitrary perfect complexes and does
not automatically extend to twisted sheaves.

\subsection{SF-E811. Strict mixed discriminant
positivity}\label{SF-E811}

Let \(X=V/\Lambda\) be a complex abelian fourfold with a
translation-invariant Kähler form

\[
 \omega=\omega_1+\omega_2,
\]

where \([\omega]\in H^2(X,\mathbf Q)\) is rational, \(V=V_1\oplus V_2\)
is an orthogonal complex splitting with \(\dim_{\mathbf C}V_i=2\), and
\(\omega_i\) is positive on \(V_i\) and zero on the other summand. Thus
a positive multiple of \([\omega]\) is a polarization, as required by
the exact projective existence source used below. Assume that neither of
the real four-form lines

\[
 \mathbf R\omega_1^2,\qquad \mathbf R\omega_2^2
 \quad\subset H^4(X,\mathbf R)
 \tag{E811.1}
\]

contains a nonzero rational class. For a holomorphic vector bundle \(E\)
of rank \(r>0\), set

\[
 \Delta(E)=c_2(E)-\frac{r-1}{2r}c_1(E)^2
           =\frac{c_1(E)^2}{2r}-\operatorname{ch}_2(E).
 \tag{E811.2}
\]

If \(E\) is slope-stable with respect to \(\omega\), then

\[
 \int_X\Delta(E)\,\omega_1\omega_2\ge0,
\]

and equality holds only when \(E\) is projectively flat. In particular
equality forces \(\Delta(E)=0\).

The proof first treats any Hermitian--Einstein connection. Its only use
of stability, apart from existence of this connection, is to make the
final finite deck action transitive on the parallel blocks. The parallel
algebra bundles used in the proof descend under the full holonomy; they
are not chosen summands that exist only on the universal cover.

\subsubsection{1. The mixed Chern--Weil
identity}\label{the-mixed-chernweil-identity}

The published Donaldson--Uhlenbeck--Yau existence theorem supplies a
Hermitian--Einstein Chern connection on \(E\). A complete projective
proof is in Siu, \emph{Lectures on Hermitian--Einstein metrics for
stable bundles and Kähler--Einstein metrics}, DMV Seminar 8 (1987),
chapter 1, §6.1--§6.4, author-reprint pp.35--39. We use this published
theorem as an input, not a new proof of the nonlinear existence theorem.
The normalization of a Hermitian--Einstein connection and its trace-free
part is also explicit in that source, chapter 1, §1.3--§1.4, author
pp.6--7.

Let \(F\) be the curvature, in the convention that its values on real
tangent vectors are skew-Hermitian, and put

\[
 F^0=F-\frac1r\operatorname{Tr}(F)I.
\]

Then \(F^0\) has type \((1,1)\), \(\Lambda_\omega F^0=0\), and the
Chern--Weil representative of (E811.2) is

\[
 \Delta(E)=\left[\frac1{8\pi^2}\operatorname{Tr}(F^0\wedge F^0)\right].
 \tag{E811.3}
\]

Here is a pointwise computation with all signs fixed. Use real unitary
coordinates \(x_j,y_j\), \(1\le j\le4\), and put

\[
 e_j=dx_j\wedge dy_j,\quad
 \omega_1=e_1+e_2,\quad\omega_2=e_3+e_4,\quad
 \mathrm{vol}=e_1e_2e_3e_4.
\]

A real \((1,1)\)-form can be written

\[
 \alpha=\sum_{j=1}^4 x'_j e_j+
 \sum_{i<j}\left(b_{ij}s_{ij}+c_{ij}t_{ij}\right),
\]

where

\[
 s_{ij}=dx_i\wedge dy_j+dx_j\wedge dy_i,
 \qquad t_{ij}=dx_i\wedge dx_j+dy_i\wedge dy_j.
\]

If \(\Lambda_\omega\alpha=0\), then \(\sum_jx'_j=0\), and direct
exterior multiplication gives

\[
 -\alpha^2\omega_1\omega_2
 =2\left((x'_1+x'_2)^2+
          \sum_{i\le2<j}(b_{ij}^2+c_{ij}^2)\right)\mathrm{vol}.
 \tag{E811.4}
\]

For example \(s_{ij}^2=t_{ij}^2=-2e_ie_j\); distinct off-diagonal basis
terms have zero cross pairing in this expression. The diagonal terms
give \(2(x'_1+x'_2)(x'_3+x'_4)\), and primitivity gives the displayed
sign. This proves (E811.4) directly; the exact script also verifies the
generic polynomial identity and its rank-nine positive semidefinite
matrix on the 15-dimensional primitive space.

Expand \(F^0\) in an orthonormal real basis of \(\mathfrak{su}(r)\) for
the inner product \(-\operatorname{Tr}(AB)\). Equation (E811.4), summed
over that basis, proves the nonnegativity in the theorem. If the
integral is zero, the pointwise nonnegative smooth density is zero
everywhere. Therefore

\[
 F^0=F_1+F_2,
 \qquad F_i\text{ has both real arguments in }V_i,
 \qquad\Lambda_{\omega_i}F_i=0.
 \tag{E811.5}
\]

In particular every mixed real curvature component vanishes. The
distinction between real and complex mixed components causes no gap:
type \((1,1)\) already removes the \((2,0)\) and \((0,2)\) terms, and
(E811.4) removes the remaining mixed terms.

\subsubsection{2. A product gauge and the labelled holonomy
algebras}\label{a-product-gauge-and-the-labelled-holonomy-algebras}

Pull back to the universal cover \(\widetilde X=V_1\times V_2\). The
smooth bundle is trivial there. In a unitary trivialization, subtract
the scalar part of the connection one-form. This produces a trace-free
unitary connection whose curvature is exactly \(F^0\). Choices of
determinant trivialization change the construction only by a scalar
one-form, which has no effect on conjugation on endomorphisms.

We prove the required splitting of this connection. Choose the frame
obtained by parallel transport first on the radial segment in \(V_1\)
and then on the radial segment in \(V_2\). Transport on each rectangle
made of one segment in each factor is independent of its order, because
the pulled-back connection on that rectangle has zero curvature by
(E811.5). Thus the frame is radial in both factors. If \(R_i\) is the
radial vector field of the \(i\)-th factor and \(A=A_1+A_2\) is the
connection one-form in this frame, then

\[
 A(R_1)=A(R_2)=0.
\]

For a constant vector \(u\in V_1\), the equation \(F^0(R_2,u)=0\)
becomes \(R_2(A_1(u))=0\). Smoothness at the origin of \(V_2\) implies
that \(A_1(u)\) is independent of the \(V_2\)-coordinate. The same
argument with the two factors exchanged shows that \(A_2\) depends only
on its own coordinate. The remaining mixed curvature equation then gives

\[
 A=A_1(z_1)+A_2(z_2),\qquad
 [A_1(z_1)(u_1),A_2(z_2)(u_2)]=0
 \tag{E811.6}
\]

for every choice of arguments. All these assertions concern real
coordinates; no holomorphic gauge is required.

Let \(H_i\) be the unitary holonomy subgroup generated by loops in the
\(i\)-th factor, based at a chosen point of \(\widetilde X\). Formula
(E811.6) implies that \(H_1\) and \(H_2\) commute and that the holonomy
of every loop on the universal cover lies in \(H_1H_2\). Let
\(\mathscr A_i\) be the unital complex \(*\)-algebra generated by
\(H_i\), as a subalgebra of the matrix algebra of the fiber. It is
finite-dimensional, and the two algebras commute.

These algebras define intrinsic \textbf{parallel subbundles of
\(\operatorname{End}(E)\) on \(X\)}. We spell out the descent. Changing
the base point along a path in the same factor conjugates its factor
holonomy group by parallel transport. Changing it along the other factor
identifies the group with that of the translated loops, by the
zero-curvature rectangle argument. A general path is handled by
(E811.6), and a loop on the universal cover normalizes both labelled
algebras because its holonomy belongs to the commuting product
\(H_1H_2\). Finally, every deck translation sends a loop tangent to
\(V_i\) to another loop with the same label \(i\); the bundle's deck
identification conjugates their parallel transports, up to the harmless
scalar from removing the determinant. Multiplying a generator by a
nonzero complex scalar does not change the unital complex algebra it
generates. Thus the construction is equivariant under every deck
transformation, and descends without choosing a lift of a projective
factor bundle or a root of the determinant.

The induced connections on \(\mathscr A_i\) are unitary for their
restricted Hilbert--Schmidt metrics. They are the Chern connections of
holomorphic vector bundles, since the induced \((0,2)\) curvature is
zero. Their curvatures are

\[
 F_{\mathscr A_i}=[F_i,-]|_{\mathscr A_i}.
 \tag{E811.7}
\]

Indeed the scalar part of \(F\) commutes with everything, and the other
factor's curvature commutes with \(\mathscr A_i\). The latter also
follows directly by differentiating the commuting factor parallel
transports. Equation (E811.5) makes these induced connections
Hermitian--Einstein with zero Einstein constant.

\subsubsection{3. Rationality kills the two adjoint
curvatures}\label{rationality-kills-the-two-adjoint-curvatures}

The curvature in (E811.7) has both arguments in the real four-plane
underlying \(V_i\). Its Chern--Weil degree-four form therefore belongs
pointwise to the one-dimensional space of real four-forms supported on
that plane. On a torus, averaging differential forms over translations
preserves their de Rham class. Averaging such a form leaves it in the
same constant one-dimensional space. Consequently

\[
 \operatorname{ch}_2(\mathscr A_i)\in
       \mathbf R\omega_i^2\cap H^4(X,\mathbf Q)=0
 \tag{E811.8}
\]

by the hypothesis (E811.1). Rationality is legitimate here:
\(\mathscr A_i\) is an actual finite-rank holomorphic vector bundle on
\(X\), not an unspecified projective or twisted bundle.

For a unitary primitive \((1,1)\) curvature \(G\), the Kähler Hodge-star
identity is

\[
 \operatorname{Tr}(G\wedge G)\wedge\omega^2
     =2\,|G|^2\,\frac{\omega^4}{4!}
 \tag{E811.9}
\]

with the positive inner product \(-\operatorname{Tr}\) and the usual
unitary norm convention. Equivalently, one can use only that the
constant on the right is positive; it follows by expanding in primitive
real \((1,1)\)-forms. Since
\(\operatorname{ch}_2=-\operatorname{Tr}(G^2)/(8\pi^2)\), (E811.8) and
(E811.9) force

\[
 F_{\mathscr A_1}=F_{\mathscr A_2}=0.
 \tag{E811.10}
\]

This conclusion kills the adjoint curvatures, not yet the centers. We
deal with the centers explicitly next.

\subsubsection{4. The central curvatures and a finite deck
orbit}\label{the-central-curvatures-and-a-finite-deck-orbit}

Every factor loop used to define \(H_i\) is contractible on \(X\), since
it is a closed loop on the simply connected universal cover. Its
holonomy on the flat bundle \(\mathscr A_i\) is therefore the identity.
On the matrix algebra this holonomy is conjugation by the corresponding
element of \(H_i\). Thus every element of \(H_i\) commutes with every
element of \(\mathscr A_i\). Since \(H_i\) generates \(\mathscr A_i\),
the algebra \(\mathscr A_i\) is commutative.

The unital \(*\)-algebra \(\mathscr A\) generated by \(\mathscr A_1\)
and \(\mathscr A_2\) is also commutative, because the two factor
algebras commute. It is a parallel bundle of finite-dimensional
commutative \(*\)-algebras. Such a fiber is a direct sum of copies of
\(\mathbf C\): simultaneous unitary diagonalization gives its primitive
orthogonal idempotents. Parallel transport permutes these finitely many
idempotents. A finite-index sublattice of \(\Lambda\) kills this
permutation action; for example, multiplication by \(r!\) suffices since
the number of nonzero idempotent blocks is at most \(r\).

Let \(p:X'\to X\) be the resulting finite étale isogeny. The pullback
splits into parallel holomorphic vector subbundles

\[
 p^*E=\bigoplus_{j=1}^s E_j.
 \tag{E811.11}
\]

The infinitesimal holonomy given by \(F_i\) belongs to \(\mathscr A_i\):
this follows by taking the infinitesimal holonomy of a small factor
rectangle. On each primitive block all elements of \(\mathscr A\) act as
scalars. Hence the full curvature of each \(E_j\) is scalar, including
the original scalar determinant curvature. Each \(E_j\) is therefore
projectively flat. This does not assert that its scalar curvature is
zero.

The finite deck group acts transitively on these blocks. If it had two
or more orbits, the direct sum of one orbit of blocks would descend to a
proper parallel holomorphic subbundle of \(E\). A parallel subbundle of
a Hermitian--Einstein bundle has the same slope as the whole bundle,
because the Einstein constant restricts to the same scalar. This would
contradict slope stability.

Deck transformations of an isogeny of complex tori are translations and
act trivially on rational cohomology. Transitivity therefore gives

\[
 \frac{c_1(E_j)}{\operatorname{rk}E_j}=\alpha
      \quad\text{for one class }\alpha\in H^2(X',\mathbf Q)
      \text{ and every }j.
 \tag{E811.12}
\]

The ranks are equal as well, though equality of normalized classes is
all that is needed. Since each block is projectively flat, Chern--Weil
gives

\[
 \operatorname{ch}(E_j)=\operatorname{rk}(E_j)e^{\alpha},
 \qquad \operatorname{ch}(p^*E)=r e^{\alpha}.
 \tag{E811.13}
\]

Thus \(p^*\Delta(E)=0\), and the injectivity of rational cohomological
pullback under a finite map gives \(\Delta(E)=0\).

One can also see directly that the entire original connection is
projectively flat. Write the scalar curvature of the block \(E_j\) as
\(\beta_j I\). Each \(\beta_j\) is closed of type \((1,1)\), and its
contraction with \(p^*\omega\) is the same constant. The Kähler
identities make it harmonic; this is Siu chapter 1, §1.4(ii), author
p.7. Equation (E811.12) makes these harmonic forms cohomologous, so they
are equal. Hence the curvature of \(p^*E\) is one scalar form times the
identity. The same property descends to \(E\). This handles all central
curvatures and finishes the proof of SF-E811.

\subsection{SF-E811-RM. Consequence for the remaining even-factor bundle
gate}\label{SF-E811-RM}

Suppose that \(X\) has genuine Rosati-invariant real multiplication by a
real quadratic field \(F\), with multiplicities \((2,2)\), and let
\(\theta=\theta_1+\theta_2\) be an \(F\)-compatible polarization. Then
no \(\theta\)-slope-stable vector bundle of positive rank \(N\) can have

\[
 e^B\operatorname{ch}(E)=N
     \left(1-\frac q2\theta_1^2\right)
     \left(1-\frac q2\theta_2^2\right),\qquad q>0,
 \tag{E811.14}
\]

for any rational algebraic divisor \(B\). The assertion holds for every
polarization degree, every rank and every integral order of real
multiplication. The presence of an additional finite translation/tensor
action does not create such an underlying stable bundle.

\textbf{Proof.} The positive invariant forms \(\theta_i\) give the
splitting in SF-E811. To check (E811.1), suppose that a nonzero rational
four-form were a real multiple of \(\theta_i^2\). That four-form is
decomposable, and its contraction kernel is exactly the other real
eigenspace. The kernel of a rational linear map is a rational subspace,
so that eigenspace would contain a nonzero rational vector. Choose
\(u\in F\setminus\mathbf Q\). The rational endomorphism induced by \(u\)
acts on this eigenspace by the irrational real scalar \(\sigma_j(u)\).
For a nonzero rational vector \(z\), the equation \(uz=\sigma_j(u)z\) is
impossible: any nonzero rational coordinate of \(z\) would make
\(\sigma_j(u)\) rational. This proves (E811.1), with no assumption that
\(X\) itself is simple.

The degree-two cohomology part of (E811.14) gives \(B=-c_1(E)/N\). Its
degree-four cohomology part then gives

\[
 \Delta(E)=\frac{Nq}{2}(\theta_1^2+\theta_2^2).
 \tag{E811.15}
\]

Because \(\theta_i^3=0\), this class pairs to zero with
\(\theta_1\theta_2\). By SF-E811 it must be zero. But its pairing with
\(\theta^2\) is

\[
 Nq\int_X\theta_1^2\theta_2^2
       =4Nq\int_X\frac{\theta^4}{24}>0,
\]

a contradiction. In fact only the normalized degree-four cohomology
component of (E811.14) is needed; the rest of the character is
irrelevant to this obstruction.

\subsection{SF-E812. Polystability cannot provide a weakly equivariant
semiregular
repair}\label{SF-E812}

In the balanced-RM setting of SF-E811-RM, let \(E\) be a
slope-polystable vector bundle with character (E811.14). Let
\(G\subset X\times\widehat X\) be any effective finite subgroup with a
specified coherent action and an equivariant structure on \(E\). Then
the full trace on invariant degree-two self-Ext is not injective. In
particular, no such polystable bundle is weakly \(G\)-semiregular,
regardless of its rank, endomorphism dimension, polarization degree or
group order.

This assertion permits actual polystable bundles with the stated Chern
character to exist. It excludes their proposed semiregularity property.
Its equivariant direct summands are constructed on \(X\) itself, not on
an auxiliary isogeny cover.

\subsubsection{1. All stable factors are projectively
flat}\label{all-stable-factors-are-projectively-flat}

Write the actual polystable decomposition

\[
 E=\bigoplus_j E_j,\qquad
 r_j=\operatorname{rk}E_j,\quad
 \alpha_j=c_1(E_j)/r_j,\quad
 \alpha=c_1(E)/N.
\]

The \(E_j\) are slope-stable vector bundles of the same
\(\theta\)-slope. Direct expansion of the discriminant gives

\[
 \Delta(E)=\sum_j\Delta(E_j)
       -\frac12\sum_j r_j(\alpha_j-\alpha)^2.
 \tag{E812.1}
\]

Represent \(\delta_j=\alpha_j-\alpha\) by its invariant real
\((1,1)\)-form. Equality of slopes says \(\Lambda_\theta\delta_j=0\).
Therefore each term

\[
 -\int_X\delta_j^2\theta_1\theta_2
\]

is nonnegative by the elementary identity (E811.4). Each
\(\int_X\Delta(E_j)\theta_1\theta_2\) is nonnegative by SF-E811. The
left side of (E812.1), paired with \(\theta_1\theta_2\), is zero by
(E811.15). Thus every nonnegative summand is zero. SF-E811 applied to
each stable factor shows that every \(E_j\) is projectively flat, and
hence

\[
 \operatorname{ch}(E_j)=r_j e^{\alpha_j}.
 \tag{E812.2}
\]

This step is compatible with the existence of polystable examples having
several different \(\alpha_j\): the negative variance term in (E812.1)
need not vanish as a cohomology class even though its mixed pairing
vanishes.

\subsubsection{2. The slope-class groups are actual equivariant
summands}\label{the-slope-class-groups-are-actual-equivariant-summands}

For each distinct rational class \(\beta\) among the \(\alpha_j\), let
\(E_\beta\) be the direct sum of the factors with \(\alpha_j=\beta\).
Polystable decomposition is unique at the level of stable isotypic
summands: nonisomorphic stable bundles of the same slope have no nonzero
Hom, and multiplicity spaces give the usual canonical isotypic sums.
Grouping these isotypic sums by \(\beta\) is therefore independent of a
choice of decomposition.

Translations and tensor products by degree-zero line bundles preserve
slope stability and the rational class \(c_1/r\). It follows that every
specified isomorphism

\[
 t_x^*E\otimes P\simeq E
\]

preserves each \(E_\beta\). The linearizations restrict, and their
coherence restricts as well. Consequently

\[
 E=\bigoplus_\beta E_\beta
 \tag{E812.3}
\]

is an actual direct-sum decomposition in the given equivariant category.
Each summand is nonzero and has

\[
 \operatorname{ch}(E_\beta)=r_\beta e^\beta,
 \qquad r_\beta>0.
 \tag{E812.4}
\]

\subsubsection{3. Trace/cap incompatibility for any one of these
summands}\label{tracecap-incompatibility-for-any-one-of-these-summands}

Suppose that the full trace on \(\operatorname{Ext}^2(E,E)^G\) were
injective. For the specified effective translation/tensor action,
supplied DEG\_L S-21 identifies the identity Hochschild sector with that
of \(X\), and its invariant Calabi--Yau adjunction makes degree-two
evaluation surjective. Thus the full trace image equals the cap-product
image of the whole character. This is precisely the implication used and
independently reviewed in SF-E802, including the equivariant setting.

Normalize all characters simultaneously by the same \(B\)-field shear in
(E811.14). In the Fock exterior model of SF-E801 the whole character
becomes, up to a nonzero scalar,

\[
 v=\omega'_1\wedge\omega'_2,
 \qquad W=U_1\oplus U_2,\quad\dim U_i=4,
\]

with each \(\omega'_i\in\bigwedge^2U_i\) nondegenerate. Its degree-two
cap kernel is

\[
 K=K_1\oplus K_2,
 \qquad K_i=\ker\left(\bigwedge^2U_i
                 \xrightarrow{\wedge\omega'_i}\bigwedge^4U_i\right),
 \qquad\dim K_i=5.
 \tag{E812.5}
\]

The common annihilator of \(K\) in \(\bigwedge^4W\) is exactly
\(\mathbf C v\), as proved in SF-E801.5. For clarity, the proof is
short: in bidegrees \((0,4),(1,3),(3,1),(4,0)\), multiplication by one
of the primitive spaces detects every nonzero component. In bidegree
\((2,2)\), annihilation by \(K_1\) restricts the first factor to
\(\mathbf C\omega'_1\), and annihilation by \(K_2\) restricts the second
to \(\mathbf C\omega'_2\). The nondegenerate exterior pairing on each
\(\bigwedge^2U_i\) proves these claims.

Because (E812.3) is an actual equivariant direct sum, the trace image of
each \(E_\beta\) is contained in the trace image of \(E\). Evaluation
followed by trace is cap product for each object. Hence

\[
 \operatorname{im}c^2_{e^B\operatorname{ch}(E_\beta)}
        \subset\operatorname{im}c^2_v.
 \tag{E812.6}
\]

The bilinear form \((z,w)\mapsto z\wedge w\wedge v\) on \(\bigwedge^2W\)
is symmetric, and the top exterior pairing is perfect. Taking
annihilators in (E812.6) gives

\[
 K\subset\ker c^2_{e^B\operatorname{ch}(E_\beta)}.
\]

Thus the normalized summand character belongs to the common-annihilator
line \(\mathbf C v\). But (E812.4) makes that character
\(r_\beta e^{B+\beta}\), a nonzero pure exponential. Its Fock
four-vector is decomposable, and its degree-two cap rank is six. The
balanced character has cap rank eighteen. They cannot be proportional.
This contradiction proves SF-E812.

No equivariant simplicity hypothesis was required. In particular, this
argument rules out weakly semiregular polystable repairs that would try
to exploit a large finite group order while leaving the actual
slope-class summands intact.

\subsection{Scope, sources and remaining
gates}\label{scope-sources-and-remaining-gates}

SF-E811 removes the \textbf{stable locally free} realization from the
remaining balanced-RM search even when the polarization/group factor is
even and the rank is at least four. SF-E812 removes every
\textbf{polystable locally free weakly equivariant semiregular}
realization. Neither assertion removes unstable locally free sheaves,
non-polystable semistable extensions, non-locally-free torsion-free
sheaves, arbitrary nonsimple indecomposable objects, twisted objects or
complexes with cohomology in several degrees. No simple perfect complex
is assumed to be slope-stable.

The only deep analytic input is the published Hermitian--Einstein
existence theorem. The exact accessible proof source is Yum-Tong Siu,
\emph{Lectures on Hermitian--Einstein metrics for stable bundles and
Kähler--Einstein metrics}, DMV Seminar 8, Birkhäuser, 1987, chapter 1,
§6.1--§6.4, \textbf{author-reprint pp.35--39}; the same author reprint
has §1.4(ii),(iii) on p.7 and Proposition 1.6 with proof on pp.8--9.
URL:
\url{https://people.math.harvard.edu/~siu/siu_reprints/dmv_book.pdf}.
The source clearly says that its proof presents Donaldson's improved
all-dimensional heat-flow argument. These are author-copy page numbers,
not an asserted identification with another printing.

The original published references are K. Uhlenbeck and S.-T. Yau,
\emph{On the existence of Hermitian--Yang--Mills connections in stable
vector bundles}, Comm. Pure Appl. Math. 39 (1986), S257--S293, DOI
\texttt{10.1002/cpa.3160390714}, with the 1989 note, 42 (1989),
703--707, DOI \texttt{10.1002/cpa.3160420505}; and S. K. Donaldson,
\emph{Infinite determinants, stable bundles and curvature}, Duke Math.
J. 54 (1987), 231--247, DOI \texttt{10.1215/S0012-7094-87-05414-7}.
Their publication metadata was checked. Direct PDFs of these originals
were unavailable in the present retrieval, so no original theorem number
or full original-proof read is claimed.

The mixed sum-of-squares identity, labelled algebra descent,
rational-volume-line argument, treatment of centers, and finite
deck-orbit argument are proved above. They invoke neither a preprint
theorem nor a conjectural geometric decomposition. The exact source
scope and the separate root/peer whole-proof checks are recorded in
SOURCE\_LOG.md and CHECK\_EVENTS.tsv. No formal proof-assistant
verification is claimed.

\section{SF-E813: the balanced-RM obstruction for locally free
Jordan--Hölder
factors}\label{SF-E813}

Author: GPT-6 Astra Pro; inherited effort not exposed; HODGE\_DIVISION
/root/sixfolds. Date: 2026-11-21.

\textbf{Status:} Whole-proof mathematical PASS reviews received from
both /root and /root/imprimitive on 2026-11-21 (GPT-6 Astra Pro;
inherited effort not exposed). Both reviewers reconstructed the complete
argument and checked the exact invoked Mumford statements in the
original book scan. SF-E811--E812 are also independently root/peer
checked. No assertion about singular stable factors or degree-two
derived gluings is made.

\subsection{Statement}\label{SF-E813-statement}

Let \(X\) be a complex abelian fourfold with genuine Rosati-invariant
real multiplication by a real quadratic field, of type \((2,2)\), and
compatible polarization

\[
 \theta=\theta_1+\theta_2.
\]

Let \(E\) be a slope-semistable vector bundle admitting a filtration by
holomorphic subbundles

\[
 0=E^0\subset E^1\subset\cdots\subset E^s=E
 \tag{E813.1}
\]

whose quotients \(Q_j=E^j/E^{j-1}\) are slope-stable \textbf{locally
free} sheaves of the same \(\theta\)-slope as \(E\). Suppose that, for a
rational algebraic divisor \(B\), rank \(N>0\), and \(q>0\),

\[
 e^B\operatorname{ch}(E)=N
    \left(1-\frac q2\theta_1^2\right)
    \left(1-\frac q2\theta_2^2\right).
 \tag{E813.2}
\]

For any effective finite translation/tensor subgroup
\(G\subset X\times\widehat X\), with a specified coherent action and an
equivariant structure on \(E\), the full trace on
\(\operatorname{Ext}^2(E,E)^G\) is not injective. Thus no such bundle is
weakly \(G\)-semiregular.

The filtration condition is essential to the scope: a semistable vector
bundle on a higher-dimensional variety need not have locally free stable
Jordan--Hölder factors. This theorem does not silently replace a
torsion-free factor by its reflexive hull. It also does not assert that
\(E\) itself is polystable.

\subsection{1. All stable factors are projectively
flat}\label{SF-E813-all-stable-factors-are-projectively-flat}

Put

\[
 r_j=\operatorname{rk}Q_j,\qquad
 \alpha_j=c_1(Q_j)/r_j,\qquad
 \alpha=c_1(E)/N,\qquad \delta_j=\alpha_j-\alpha.
\]

Additivity of the Chern character in (E813.1) gives exactly the same
variance identity as for an actual direct sum:

\[
 \Delta(E)=\sum_j\Delta(Q_j)
          -\frac12\sum_jr_j\delta_j^2.
 \tag{E813.3}
\]

Here
\(\Delta(F)=c_1(F)^2/(2\operatorname{rk}F)-\operatorname{ch}_2(F)\). The
normalized degree-four part of (E813.2) gives

\[
 \Delta(E)=\frac{Nq}{2}(\theta_1^2+\theta_2^2),
 \qquad \int_X\Delta(E)\theta_1\theta_2=0.
 \tag{E813.4}
\]

Each \(Q_j\) is an actual stable vector bundle, so SF-E811 applies and
gives nonnegative mixed discriminant pairing. Equality of slopes makes
the invariant representative of each \(\delta_j\) primitive for
\(\theta\). The exact mixed quadratic identity (E811.4) gives

\[
 -\int_X\delta_j^2\theta_1\theta_2\ge0.
\]

Every term on the right of (E813.3), after this pairing, is therefore
nonnegative, and their sum is zero. It follows that every term vanishes.
SF-E811 shows that every \(Q_j\) is projectively flat and

\[
 \operatorname{ch}(Q_j)=r_je^{\alpha_j}.
 \tag{E813.5}
\]

The equality case of the elementary mixed quadratic identity also says
that each invariant \(\delta_j\) is block diagonal for the two
real-multiplication eigenspaces, and primitive on each complex two-plane
separately. Hence every difference

\[
 \alpha_j-\alpha_k=\delta_j-\delta_k
 \tag{E813.6}
\]

has those same two properties.

\subsection{\texorpdfstring{2. Every nonzero difference has Hermitian
signature
\((2,2)\)}{2. Every nonzero difference has Hermitian signature (2,2)}}\label{SF-E813-every-nonzero-difference-has-hermitian-signature-22}

We prove the rationality point explicitly. Let
\(\delta\in H^2(X,\mathbf Q)\cap H^{1,1}(X)\) be nonzero, block diagonal
and primitive on each of the two complex eigenspaces. If one of its
blocks vanished, \(\delta\) would be supported on just one eigenspace,
say the one corresponding to the real embedding \(\sigma_i\). For
\(u\in F\setminus\mathbf Q\) and rational homology vectors \(x,y\), this
would imply

\[
 \delta(ux,y)=\sigma_i(u)\delta(x,y).
 \tag{E813.7}
\]

The left side is rational, because both \(u\) and \(\delta\) are
rational. Nonzero rational \(\delta\) has some rational pair \((x,y)\)
with \(\delta(x,y)\ne0\). Equation (E813.7) would make the irrational
real number \(\sigma_i(u)\) rational, a contradiction. Thus both blocks
are nonzero.

In a unitary basis for either complex two-plane, the Hermitian matrix of
a real primitive \((1,1)\)-form is

\[
 A_i=\begin{pmatrix}a_i&b_i+\mathrm i c_i\\
                      b_i-\mathrm i c_i&-a_i\end{pmatrix},
 \qquad
 \det A_i=-a_i^2-b_i^2-c_i^2<0.
 \tag{E813.8}
\]

Its eigenvalues have opposite signs. The complete four-dimensional
Hermitian form of \(\delta\) is the direct sum of these two nonzero
blocks, hence is nondegenerate of signature \((2,2)\). In particular

\[
 \int_X\frac{\delta^4}{24}>0.
 \tag{E813.9}
\]

This conclusion applies to each unequal pair in (E813.6). No assumption
that the divisor space of \(X\) is minimal has been used.

\subsection{3. Unequal factor classes have no degree-one
extensions}\label{SF-E813-unequal-factor-classes-have-no-degree-one-extensions}

We first give an elementary finite-isogeny splitting argument for the
projectively flat factors, so no classification of semihomogeneous
bundles is being assumed.

Let \(Q\) have rank \(r\), with its projectively flat
Hermitian--Einstein connection. Its scalar curvature is a harmonic
invariant \((1,1)\)-form representing \(c_1(Q)/r\), up to the usual
Chern--Weil scalar. Choose any positive integer \(m\) divisible by
\(r\), and consider multiplication \(\pi=[m]:X\to X\). The line bundle

\[
 L_Q=(\det Q)^{m^2/r}
 \tag{E813.10}
\]

on the source has first Chern class \(m^2c_1(Q)/r\), the normalized
first Chern class of \(\pi^*Q\). Give it a Hermitian metric with that
invariant curvature; this follows from the usual
\(\partial\bar\partial\)-lemma for line-bundle metrics, or directly from
the invariant metric in the Appell--Humbert description. The difference
of the two scalar curvature forms is a harmonic form with zero class,
and is zero. Thus

\[
 \pi^*Q\otimes L_Q^{-1}
\]

is unitary flat. A unitary representation of the abelian fundamental
group of a torus is a direct sum of one-dimensional unitary characters,
by simultaneous diagonalization of commuting normal matrices. Therefore

\[
 \pi^*Q\simeq\bigoplus_{a=1}^r (L_Q\otimes P_a),
 \qquad P_a\in\operatorname{Pic}^0(X).
 \tag{E813.11}
\]

This is a genuine splitting of holomorphic bundles. Neither a projective
action nor a root of a line bundle is substituted for it.

Apply this construction to two factors \(Q_j,Q_k\), taking \(m\)
divisible by both ranks. If \(\alpha_j\ne\alpha_k\), every line-bundle
summand of

\[
 (\pi^*Q_j)^\vee\otimes\pi^*Q_k
\]

has first Chern class \(m^2(\alpha_k-\alpha_j)\), whose Hermitian form
is nondegenerate of signature \((2,2)\). Mumford's vanishing theorem and
its complex index corollary then give cohomology only in degree two. The
exact published locators are \emph{Abelian varieties}, §16, Vanishing
Theorem printed p.150, proof pp.150--153, and the unnumbered complex
Hermitian-index corollary with proof pp.162--163. These statements are
unchanged by a tensor factor in \(\operatorname{Pic}^0\), since that
factor has zero curvature class.

Consequently

\[
 \operatorname{Ext}^{\ell}(\pi^*Q_j,\pi^*Q_k)=0
 \quad\text{for all }\ell\ne2.
 \tag{E813.12}
\]

The pullback map on these Ext groups is injective. Indeed, for a finite
étale map the trace, divided by its degree, splits the unit inclusion
\(\mathcal O_X\to\pi_*\mathcal O_X\). The projection formula makes
\(H^{\ell}(X,Q_j^\vee\otimes Q_k)\) a direct summand of its pullback
cohomology. It follows that

\[
 \operatorname{Ext}^{\ell}(Q_j,Q_k)=0
 \quad(\ell\ne2,\ \alpha_j\ne\alpha_k).
 \tag{E813.13}
\]

In particular both Hom and \(\operatorname{Ext}^1\) vanish across
unequal normalized first Chern classes. Riemann--Roch also identifies
the remaining dimension:

\[
 \dim\operatorname{Ext}^2(Q_j,Q_k)
   =r_jr_k\int_X\frac{(\alpha_k-\alpha_j)^4}{24}>0.
 \tag{E813.14}
\]

The nonzero degree-two extensions are not being declared harmless for
derived constructions. They remain outside this coherent extension
argument.

\subsection{4. Iterated extensions split into actual slope-class
summands}\label{SF-E813-iterated-extensions-split-into-actual-slope-class-summands}

Let \(\mathcal C_\beta\) denote the class of bundles admitting a finite
filtration by projectively flat stable factors of normalized first Chern
class \(\beta\), among the factors in question and their
translation/tensor images. The factor argument above works for all such
images: they have the same normalized class and signature of any unequal
difference.

For \(A\in\mathcal C_\beta\), \(B\in\mathcal C_\gamma\), with
\(\beta\ne\gamma\), induction along their filtrations using the long
exact Hom and Ext sequences gives

\[
 \operatorname{Hom}(A,B)=0,
 \qquad \operatorname{Ext}^1(A,B)=0.
 \tag{E813.15}
\]

Here is the precise induction for the extension assertion. For a short
exact filtration step \(0\to B'\to B\to B''\to0\), the relevant sequence
has \(\operatorname{Ext}^1(A,B')\) on the left of
\(\operatorname{Ext}^1(A,B)\) and \(\operatorname{Ext}^1(A,B'')\) on its
right. Vanishing of both outer terms forces the middle one to vanish,
regardless of any subsequent \(\operatorname{Ext}^2\). Induction in the
first argument uses the analogous contravariant long exact sequence. Hom
is treated the same way. Thus no unproved assertion about an abelian
category of all slope-semistable bundles is needed.

We now induct along (E813.1). Suppose its preceding term has already
split by the distinct classes \(\beta\), and that the next quotient has
class \(\gamma\). The extension class lies in the direct sum of the
groups \(\operatorname{Ext}^1(Q, E_\beta)\). Equation (E813.15) makes
all summands except \(\beta=\gamma\) zero. Therefore the new term is an
extension entirely in the \(\gamma\)-block, direct-summed with all
unchanged other blocks. This proves an actual decomposition

\[
 E=\bigoplus_\beta E_\beta,
 \qquad E_\beta\in\mathcal C_\beta,
 \qquad \operatorname{ch}(E_\beta)=r_\beta e^\beta,
 \quad r_\beta>0.
 \tag{E813.16}
\]

The decomposition is canonical in the sense needed here. Any morphism
from an object in \(\mathcal C_\beta\) to \(E\) has zero projection to
each \(E_\gamma\), \(\gamma\ne\beta\), by (E813.15). Thus it factors
through \(E_\beta\). This characterization makes the summand independent
of any splitting choices in the induction.

Translation and tensoring by a line bundle in \(\operatorname{Pic}^0\)
preserve stability, the given slope, projective flatness and the
normalized first Chern class. In particular they carry \(E_\beta\) into
\(\mathcal C_\beta\). Every specified equivariant isomorphism
\(t_x^*E\otimes P\simeq E\) therefore preserves each summand (E813.16),
by the preceding Hom characterization and its inverse. The
linearizations and their coherence restrict. Consequently (E813.16) is
an actual decomposition in the given equivariant category on \(X\), not
merely after an isogeny.

\subsection{5. The invariant trace
contradiction}\label{SF-E813-the-invariant-trace-contradiction}

If \(E\) were weakly \(G\)-semiregular, supplied DEG\_L S-21 and
invariant Calabi--Yau adjunction would make the trace image equal the
degree-two cap image of the entire character. The trace of each actual
equivariant summand in (E813.16) is contained in that image. Applying
the same \(B\)-field shear to all characters, one obtains

\[
 \operatorname{im}c^2_{r_\beta e^{B+\beta}}
       \subset\operatorname{im}c^2_v,
 \qquad
 v=\left(1-\frac q2\theta_1^2\right)
   \left(1-\frac q2\theta_2^2\right).
 \tag{E813.17}
\]

SF-E801.5 and the full proof in SF-E812 identify the common annihilator
of \(\ker c^2_v\) with the one-dimensional line spanned by \(v\).
Perfect exterior pairing turns the image inclusion (E813.17) into the
corresponding kernel inclusion. It follows that \(r_\beta e^{B+\beta}\)
is proportional to \(v\). This is impossible: the cap rank of a nonzero
pure exponential on a fourfold is six, while the balanced character has
cap rank eighteen. This proves SF-E813.

\subsection{Exact input and scope}\label{SF-E813-exact-input-and-scope}

The only new published input beyond SF-E811--E812 is Mumford's
line-bundle vanishing/index theorem. The primary book scan was retrieved
from

\url{https://wstein.org/edu/Fall2003/252/references/mumford-abvar/Mumford-Abelian_Varieties.pdf}.

The printed pages 150--153 and 162--163 were visually checked; the
intervening index proof pp.154--162 was also OCR-read, but this is not
claimed as an independent audit of every step or of the book's earlier
dependencies. The theorem is a published input. Sug Woo Shin,
\emph{Abelian varieties and Weil representations}, Algebra \& Number
Theory6(8)(2012),1719--1772, DOI10.2140/ant.2012.6.1719, Lemma2.4
printed p.1727, independently confirms the exact nondegenerate vanishing
theorem and cites Mumford §16; that publisher PDF and lemma were opened.
No preprint theorem is imported.

SF-E813 excludes finite equivariant semiregular repairs by nontrivial
coherent extensions of the projectively flat stable factors just
described, including the actual CM polystable examples. It leaves open
semistable sheaves with singular stable factors, unstable bundles,
twisted sheaves, and genuinely derived gluings using the nonzero
degree-two Ext groups. The explicit surviving degree-two groups in
(E813.14) explain why this proof does not assert a no-go result for all
perfect complexes.

\section{SF-E814: an actual simple degree-two cone at the sharp secant
Euler
value}\label{SF-E814}

Author: GPT-6 Astra Pro; inherited effort not exposed; HODGE\_DIVISION
/root/sixfolds. Date: 2026-11-21.

\textbf{Status:} Whole-proof mathematical PASS received from /root on
2026-11-21 (GPT-6 Astra Pro; inherited effort not exposed), after
independent reconstruction of all three long exact sequences and full
trace consequences. This constructs a genuine ordinary derived object of
rational two-pure secant type. It is not a balanced-RM seed and does not
close the split-eightfold gate.

\subsection{Statement and
construction}\label{SF-E814-statement-and-construction}

Let \(X\) be a complex abelian fourfold, and let \(A,B\) be line bundles
for which

\[
 \delta=c_1(B)-c_1(A)
\]

has a nondegenerate Hermitian form of signature \((2,2)\). Put

\[
 m=\chi(A^\vee\otimes B)=\int_X\frac{\delta^4}{24}>0.
\]

The published Mumford vanishing/index theorem, with the exact locators
recorded in SF-E813, gives

\[
 \operatorname{Ext}^j(A,B)=\operatorname{Ext}^j(B,A)=0\quad(j\ne2),
 \qquad \dim\operatorname{Ext}^2(A,B)=m.
 \tag{E814.1}
\]

Choose any nonzero \(e\in\operatorname{Ext}^2(A,B)\), and define \(F\)
by the distinguished triangle

\[
 B[1]\xrightarrow{i}F\xrightarrow{p}A\xrightarrow{e}B[2].
 \tag{E814.2}
\]

This is an actual perfect complex. Its only nonzero cohomology sheaves
are \(\mathcal H^{-1}(F)=B\) and \(\mathcal H^0(F)=A\), and its
Postnikov invariant is \(e\). It has

\[
 \operatorname{ch}(F)=e^{c_1(A)}-e^{c_1(B)}.
 \tag{E814.3}
\]

Then \(F\) is simple and gluable, and its self-Ext dimensions are

\[
 (\operatorname{ext}^0,\operatorname{ext}^1,
   \operatorname{ext}^2,\operatorname{ext}^3,
   \operatorname{ext}^4)
   =(1,m+7,12,m+7,1).
 \tag{E814.4}
\]

Its full degree-two trace is injective. Its full degree-one trace has
rank eight and kernel dimension \(m-1\). Thus when \(m=1\), the full
trace is injective in degrees zero, one and two. In that case
\(F\boxtimes F^\vee\) on \(X\times X\) is also simple, gluable and fully
semiregular in degree two.

\subsection{1. Direct calculation of the lower self-Ext
groups}\label{SF-E814-direct-calculation-of-the-lower-self-ext-groups}

For a line bundle on an abelian fourfold,

\[
 \operatorname{Ext}^j(A,A)=\operatorname{Ext}^j(B,B)
   =H^j(X,\mathcal O_X),\qquad h^j(X,\mathcal O_X)=\binom4j.
 \tag{E814.5}
\]

Apply \(\operatorname{RHom}(A,-)\) to (E814.2). The only cross group is
\(\operatorname{Ext}^2(A,B)\). The connecting map

\[
 \operatorname{Hom}(A,A)\longrightarrow\operatorname{Ext}^2(A,B),
 \qquad 1\longmapsto e
\]

is injective. The resulting long exact sequence gives

\[
 \begin{array}{c|ccccc}
 j&0&1&2&3&4\\ \hline
 \dim\operatorname{Ext}^j(A,F)&0&m+3&6&4&1.
 \end{array}
 \tag{E814.6}
\]

There are no negative groups. In particular the degree-one group fits
into the exact sequence with left term
\(\operatorname{Ext}^2(A,B)/\mathbf Ce\), of dimension \(m-1\), and
right term \(H^1(\mathcal O_X)\), of dimension four.

Apply \(\operatorname{RHom}(B[1],-)\) to the same triangle. Since
\(\operatorname{RHom}(B[1],A)\) has its sole cohomology in degree three,
the maps induced by \(i\) give isomorphisms

\[
 \operatorname{Ext}^j(B,B)
       \xrightarrow{\sim}\operatorname{Ext}^j(B[1],F)
       \quad(j=0,1,2).
 \tag{E814.7}
\]

Again there are no negative groups.

Finally apply \(\operatorname{RHom}(-,F)\). The relevant connecting maps
are

\[
 \partial_j:\operatorname{Ext}^j(B[1],F)
              \longrightarrow\operatorname{Ext}^{j+1}(A,F),
 \tag{E814.8}
\]

obtained by precomposing with the boundary \(e\), with the indicated
shifts. All three maps \(\partial_0,\partial_1,\partial_2\) are zero.
For \(j=0\), the generator of (E814.7) is \(i\), whose composite with
the preceding boundary vanishes by the distinguished-triangle identity.
For \(j=1,2\), (E814.7) represents every class as \(i\) composed with a
class \(a\in\operatorname{Ext}^j(B,B)\). The corresponding connecting
class factors through \(a\circ e\in\operatorname{Ext}^{j+2}(A,B)=0\), by
(E814.1).

Thus the long exact sequence gives

\[
 \operatorname{Hom}(F,F)=\mathbf C,
\]

and short exact sequences with dimensions

\[
 \operatorname{ext}^1(F,F)=(m+3)+4=m+7,
 \qquad
 \operatorname{ext}^2(F,F)=6+6=12.
 \tag{E814.9}
\]

Both outer complexes in the triangle for \(\operatorname{RHom}(-,F)\)
have no negative cohomology, so neither does
\(\operatorname{RHom}(F,F)\). This proves gluability. Serre duality on
the fourfold then supplies all remaining groups in (E814.4), and gives
zero above degree four.

This calculation uses actual distinguished triangles and long exact
sequences. It does not assume formality of the full dg or \(A_\infty\)
category generated by the two line bundles.

\subsection{2. The complete trace
calculation}\label{SF-E814-the-complete-trace-calculation}

The Fock four-vectors of the two exponential terms in (E814.3) are
decomposable. Their supporting four-planes are graphs whose difference
is the Hermitian form of \(\delta\); nondegeneracy of \(\delta\) makes
them transverse. Therefore the complete two-pure secant calculation of
SF-E801 applies, with two nonzero coefficients. In particular

\[
 \operatorname{rank}c^1_{\operatorname{ch}(F)}=8,
 \qquad
 \operatorname{rank}c^2_{\operatorname{ch}(F)}=12.
 \tag{E814.10}
\]

Evaluation followed by full trace is cap product. Its degree-two rank is
twelve, while \(\operatorname{Ext}^2(F,F)\) has dimension twelve by
(E814.9). Consequently the full degree-two trace is injective, and the
degree-two evaluation is surjective.

For completeness, the degree-one trace rank is not being inferred from
the rank of evaluation alone. The argument of SF-E810 gives the needed
upper bound. Let \(L,M\) be the transverse supporting four-planes of the
two pure terms. The degree-two cap kernel is \(I=L\otimes M\). Because
the degree-two trace is injective, this is also the kernel of degree-two
evaluation. The full trace/cap module identity gives, for every
\(z\in\operatorname{Ext}^1(F,F)\) and \(\xi\in I\),

\[
 \xi\wedge\sigma^1(z)=\sigma^3(\rho^2(\xi)\circ z)=0.
\]

The common annihilator of \(I\) in Fock degree five is

\[
 (\bigwedge^4L\otimes M)\oplus(L\otimes\bigwedge^4M),
\]

of dimension eight, by the explicit exterior calculation in SF-E810.
Thus \(\operatorname{rank}\sigma^1\le8\); (E814.10) gives the reverse
inequality. Hence

\[
 \operatorname{rank}\sigma^1=8,
 \qquad \dim\ker\sigma^1=m-1.
 \tag{E814.11}
\]

The degree-zero trace of the identity is the nonzero character (E814.3),
and \(F\) is simple, so degree zero is injective as well. Finally

\[
 \chi(F,F)=2-2(m+7)+12=-2m,
\]

agreeing directly with Riemann--Roch applied to (E814.3). The case
\(m=1\) realizes exactly the necessary value \(\chi=-2\) in SF-E810; it
shows that the exceptional numerical value in that obstruction is
attained by an actual derived object.

\subsection{3. An explicit CM example and its external
product}\label{SF-E814-an-explicit-cm-example-and-its-external-product}

Take \(X=E_{\mathrm i}^{\,4}\), with its product principal polarization.
On each elliptic curve let \(\mathcal O(0)\) be the degree-one line
bundle of its origin, and put

\[
 A=\mathcal O_X,\qquad
 B=\mathcal O(0)\boxtimes\mathcal O(0)
       \boxtimes\mathcal O(-0)\boxtimes\mathcal O(-0).
 \tag{E814.12}
\]

The associated Hermitian matrix is \(S=\operatorname{diag}(1,1,-1,-1)\).
Its signature is \((2,2)\), and \(\chi(B)=\det S=1\). One can see the
unique nonzero cross Ext group without an additional theorem: the two
positive degree-one elliptic factors have only one-dimensional \(H^0\),
the negative factors have only one-dimensional \(H^1\), and Künneth
gives \(H^2(X,B)=\mathbf C\) with all other cohomology zero. Choose its
nonzero class in (E814.2). Equations (E814.4) and (E814.11) then give

\[
 (\operatorname{ext}^0,\operatorname{ext}^1,
   \operatorname{ext}^2,\operatorname{ext}^3,
   \operatorname{ext}^4)=(1,8,12,8,1),
 \quad\chi(F,F)=-2,
 \tag{E814.13}
\]

and full trace injectivity in degrees zero, one and two.

Duality preserves these properties. On \(X\times X\), Künneth and
gluability give

\[
 \operatorname{Ext}^2(F\boxtimes F^\vee,F\boxtimes F^\vee)
 =\bigoplus_{a+b=2}
   \operatorname{Ext}^a(F,F)\otimes
   \operatorname{Ext}^b(F^\vee,F^\vee).
\]

The full trace is the exterior tensor product on these summands. Their
targets have distinct Hochschild bidegrees \((-2,0),(-1,-1),(0,-2)\),
and on each summand both factors are injective. Hence the external
product is fully semiregular. It is simple and gluable by the same
Künneth argument. Its degree-two self-Ext dimension is \(12+8^2+12=88\).

\subsection{Scope relative to the split-eightfold
target}\label{SF-E814-scope-relative-to-the-split-eightfold-target}

The fourfold complex just constructed has the rational two-pure
character \(1-e^X\), with cap rank twelve. The required balanced-RM
character has cap rank eighteen. They are different characters. A fixed
Fourier--Mukai equivalence, isogeny pullback or rational algebraic
\(B\)-field shear preserves the relevant cap rank, so none turns this
fourfold complex into the required balanced-RM input. No claim is made
that its external product solves the eightfold problem either.

The role of this construction is precise: degree-two derived gluing is
genuinely available and can satisfy the full trace conditions, while the
coherent degree-one extensions excluded in SF-E813 cannot. A multivertex
construction with the balanced character still requires a new
calculation, including its finite coherent group action and full
degree-two trace.

\section{SF-E815: a concrete balanced character and the finite-action
obstruction for its line-class cone
model}\label{SF-E815-sf-e815-a-concrete-balanced-character-and-the-finite-action-obstruction-for-its-line-class-cone-model}

Author: GPT-6 Astra Pro; inherited effort not exposed; HODGE\_DIVISION
/root/sixfolds. Date: 2026-11-21. The explicit rank-32 character was
supplied by /root; its complete calculation and the new categorical
obstruction are given below.

\textbf{Status:} Whole-proof mathematical PASS reviews of the complete
base argument were received from /root and /root/imprimitive on
2026-11-21 (GPT-6 Astra Pro; inherited effort not exposed). Both
independently reran the exact seven-character certificate and obtained
PASS. The explicit equivariant Grothendieck-group identity requested by
root is incorporated. Both reviewers also approved the projectively flat
factor extension in §6; its full written adaptation is included below.
The separately written SF-E815R refinement was whole-proof checked by
root and replaces the former numerical bound31 by63; the extended exact
certificate was freshly rerun by both author and root. The categorical
assertions are proved here, not inferred from the computation. The
result concerns the stated cohomology models and does not exclude all
perfect complexes or all nonsimple equivariant objects.

\subsection{1. The explicit actual polystable
bundle}\label{SF-E815-the-explicit-actual-polystable-bundle}

Let \(X=E_{\mathrm i}^{\,4}\), with its product principal polarization
\(\theta\). Identify integral symmetric real matrices with the
corresponding algebraic divisor classes on this CM product. Put

\[
 U=\operatorname{diag}\left(
     \begin{pmatrix}1&1\\1&-1\end{pmatrix},
     \begin{pmatrix}1&1\\1&-1\end{pmatrix}\right),
 \qquad S=\operatorname{diag}(1,1,-1,-1).
 \tag{E815.1}
\]

Then \(U^2=2I\), \(US=SU\), and \(U\) is self-adjoint for \(\theta\). It
defines genuine real multiplication by \(\mathbf Q(\sqrt2)\), of type
\((2,2)\). Write \(\theta=\theta_++\theta_-\) along its two complex
eigenspaces, and let

\[
 X_0=\theta_S=x+y,\qquad Y=\theta_{US}=\sqrt2x-\sqrt2y,
 \qquad Z_+=X_0+Y,\quad Z_-=X_0-Y.
 \tag{E815.2}
\]

Here \(x,y\) are the restrictions of \(X_0\) to the two eigenspaces. The
notation \(X_0\) distinguishes this divisor from the abelian variety
\(X\). On each complex eigenspace the matrix \(S\) has eigenvalues
\(1,-1\). Consequently

\[
 x^2=-\theta_+^2,\qquad y^2=-\theta_-^2,
 \qquad x^3=y^3=0,
 \qquad \int_Xx^2y^2=4.
 \tag{E815.3}
\]

Both \(x\) and \(y\) are primitive on their respective planes. All
rational combinations of \(X_0,Y\) therefore have zero \(\theta\)-slope.
Choose line bundles \(L_\beta\) with the indicated integral divisor
classes and form the actual bundle

\[
 \begin{split}
 E_{\mathrm{poly}}={}&\mathcal O_X^{\oplus24}
 \oplus L_{X_0}^{\oplus2}\oplus L_{-X_0}^{\oplus2}\\
 &\oplus L_{Z_+}\oplus L_{-Z_+}
 \oplus L_{Z_-}\oplus L_{-Z_-}.
 \end{split}
 \tag{E815.4}
\]

Every line summand is stable of slope zero, so this is a polystable
bundle. Its rank is 32; symmetry of the signs makes its odd
Chern-character components zero. Its degree-four cohomology component is

\[
 \operatorname{ch}_2(E_{\mathrm{poly}})
 =2X_0^2+Z_+^2+Z_-^2
 =4X_0^2+2Y^2
 =-8(\theta_+^2+\theta_-^2).
 \tag{E815.5}
\]

For the top component, if \(\beta=ax+by\), then
\(\beta^4=6a^2b^2x^2y^2\). The two coefficients of \(X_0\) have product
one, whereas the two coefficients of either \(Z_+\) or \(Z_-\) have
product \(-1\). Hence

\[
 \begin{split}
 \operatorname{ch}_4(E_{\mathrm{poly}})
 &=\frac{2X_0^4+Z_+^4+Z_-^4}{12}\\
 &=2x^2y^2=2\theta_+^2\theta_-^2.
 \end{split}
 \tag{E815.6}
\]

Thus the \textbf{entire} Chern character is

\[
 C=\operatorname{ch}(E_{\mathrm{poly}})
   =32\left(1-\frac14\theta_+^2\right)
       \left(1-\frac14\theta_-^2\right).
 \tag{E815.7}
\]

This is the balanced character with parameter \(q=1/2\). It is an actual
algebraic realization, not merely a formal class in rational cohomology.
SF-E812 already excludes weak equivariant semiregularity of this
polystable bundle. The question addressed next is whether genuinely
derived gluing within its line-class model can repair that defect.

\subsection{2. Exact independence of the seven exponential
characters}\label{SF-E815-exact-independence-of-the-seven-exponential-characters}

Let

\[
 \mathcal B=\{0,X_0,-X_0,Z_+,-Z_+,Z_-,-Z_-\}.
\]

The seven classes \(e^\beta\), \(\beta\in\mathcal B\), are linearly
independent over \(\mathbf Q\). Here is an exact determinant
certificate. Work over \(\mathbf Q(u)\), \(u^2=2\), in the
nine-dimensional exterior subalgebra

\[
 \mathbf Q(u)[x,y]/(x^3,y^3).
\]

The nine monomials are independent: their \((\deg_x,\deg_y)\) components
are distinct and each nonzero factor \(x^a,y^b\), for \(0\le a,b\le2\),
lies on its own complex two-plane. For \(\beta=aX_0+bY\), the
coefficient of \(x^py^q\) in its exponential is

\[
 \frac{(a+bu)^p(a-bu)^q}{p!q!}.
 \tag{E815.8}
\]

Take the seven columns corresponding, in order, to

\[
 (a,b)=(0,0),(1,0),(-1,0),(1,1),(-1,-1),(1,-1),(-1,1)
\]

and the seven rows

\[
 (p,q)=(0,0),(0,1),(0,2),(1,0),(1,1),(1,2),(2,0).
\]

The determinant of the matrix (E815.8) is exactly \(-512\). This gives
independence even over \(\mathbf Q(u)\), hence over \(\mathbf Q\). The
entries, determinant, full character (E815.7), and Euler numbers below
are supplied in the exact script and JSON certificate.

\subsection{3. The permitted cone model and its actual equivariant class
blocks}\label{SF-E815-the-permitted-cone-model-and-its-actual-equivariant-class-blocks}

Fix an integer \(k\ge1\). Let \(E\) be a perfect complex with

\[
 \operatorname{ch}(E)=kC.
 \tag{E815.9}
\]

Assume that every cohomology sheaf \(\mathcal H^j(E)\) is a vector
bundle admitting a finite filtration by line bundles of the form

\[
 L_\beta\otimes P,\qquad\beta\in\mathcal B,
 \quad P\in\operatorname{Pic}^0(X).
 \tag{E815.10}
\]

This permits arbitrary cohomological degrees, arbitrary extra cancelling
terms, arbitrary multiplicities, and arbitrary extensions within each
class. In particular the aggregate virtual multiplicities need not be
read from the number of vertices of a chosen complex. They will be
determined from its full character.

Let \(G\subset X\times\widehat X\) be any effective finite subgroup with
a specified coherent translation/tensor action, and suppose \(E\) has an
equivariant lift. Individual line factors in (E815.10) are \textbf{not}
assumed to admit linearizations. Nontrivial projective multiplicity
representations are allowed.

For distinct \(\beta,\gamma\in\mathcal B\), their difference is a
nonzero rational combination of \(X_0,Y\), block diagonal and primitive
on both genuine real-multiplication planes. The rationality proof in
SF-E813 shows that both blocks are nonzero and the complete Hermitian
form has signature \((2,2)\). Therefore

\[
 \operatorname{Hom}(L_\beta\otimes P,L_\gamma\otimes Q)=0,
 \qquad
 \operatorname{Ext}^1(L_\beta\otimes P,L_\gamma\otimes Q)=0
 \tag{E815.11}
\]

for all \(P,Q\in\operatorname{Pic}^0(X)\), by the published index
theorem. The exact filtration induction in SF-E813, which only uses
these two vanishings, gives actual canonical decompositions

\[
 \mathcal H^j(E)=\bigoplus_{\beta\in\mathcal B} H^j_\beta.
 \tag{E815.12}
\]

Each \(H^j_\beta\) is filtered entirely by line bundles of class
\(\beta\). Translation/tensor functors preserve that class and are exact
for the usual cohomological structure. Thus the given coherent
linearizations of \(E\) induce coherent structures on its cohomology
sheaves and preserve each canonical summand (E815.12).

Consequently the objects

\[
 U_\beta=\bigoplus_j H^j_\beta[-j]
 \tag{E815.13}
\]

are actual perfect objects of the given equivariant category. They are
not asserted to be direct summands of \(E\). Their characters are

\[
 \operatorname{ch}(U_\beta)=n_\beta e^\beta,
 \qquad n_\beta=\sum_j(-1)^j\operatorname{rk}H^j_\beta\in\mathbf Z.
 \tag{E815.14}
\]

The usual cohomological truncations are canonical and are preserved by
all the exact translation/tensor functors. Their distinguished triangles
therefore exist in the specified equivariant category. Additivity in its
Grothendieck group, followed by the actual cohomology decompositions
(E815.12), gives

\[
 [E]=\sum_j(-1)^j[\mathcal H^j(E)]
      =\sum_{\beta\in\mathcal B}[U_\beta]
      \quad\text{in }K_0(\mathcal C^G).
 \tag{E815.14a}
\]

In particular this is an equivariant Grothendieck-group identity,
although the \(U_\beta\) need not be direct summands of \(E\).
Additivity gives \(\sum_\beta\operatorname{ch}(U_\beta)=kC\). The
independence in §2 therefore forces

\[
 (n_0,n_{X_0},n_{-X_0},n_{Z_+},n_{-Z_+},n_{Z_-},n_{-Z_-})
    =k(24,2,2,1,1,1,1).
 \tag{E815.15}
\]

Thus every possible negative cohomological contribution has been
accounted for before any positivity argument is used.

\subsection{4. Exact cross-Euler scaling for the finite
action}\label{SF-E815-exact-cross-euler-scaling-for-the-finite-action}

We use the identity-sector argument of supplied DEG\_L S-21, with its
specified coherent action and effective-subgroup hypothesis. Its proof
was reread here; only the general categorical argument, not the
secant-specific conclusion, is used. Write
\(F:\mathcal C^G\to\mathcal C\) for forgetting and \(I\) for induction,
where \(\mathcal C=\operatorname{Perf}(X)\). The proof of (S21.6) gives

\[
 \operatorname{ch}_{HH}(IFV)=|G|\operatorname{ch}_{HH}(V)
 \quad\text{for every }V\in\mathcal C^G.
\]

For any two objects \(U,V\in\mathcal C^G\), the right adjunction and
categorical Riemann--Roch now give

\[
 \begin{split}
 \chi_X(FU,FV)
 &=\chi_{\mathcal C^G}(U,IFV)\\
 &=|G|\chi_{\mathcal C^G}(U,V).
 \end{split}
 \tag{E815.16}
\]

This is the cross-pair version of (S21.7). It does not require either
object to be simple or semiregular. It also does not require a model in
which the group acts freely by translations. All possible coherence
cocycles and projective multiplicity spaces are already incorporated in
\(\mathcal C^G\).

The exact inputs in the supplied proof are its vanishing of nonidentity
translation/tensor coefficient sectors, induction/forgetting adjunction,
and the categorical Riemann--Roch pairing. No new unrefereed theorem is
being treated as a published statement here.

Apply (E815.16) to the actual equivariant objects \(U_{X_0},U_{Z_+}\).
By (E815.15) their ordinary Euler pairing is

\[
 \begin{split}
 \chi_X(U_{X_0},U_{Z_+})
 &=2k^2\int_X e^Y\\
 &=8k^2.
 \end{split}
 \tag{E815.17}
\]

The last equality follows from \(Y^4=24x^2y^2\) and (E815.3), or from
the fact that \(Y\) has integral Hermitian matrix \(US\) of determinant
four. Thus

\[
 \frac{8k^2}{|G|}
   =\chi_{\mathcal C^G}(U_{X_0},U_{Z_+})
   \in\mathbf Z_{>0}.
\]

In particular

\[
 |G|\mid8k^2,
 \qquad |G|\le8k^2.
 \tag{E815.18}
\]

This divisibility statement handles arbitrary negative virtual
multiplicities before their aggregate is fixed by (E815.15), and it does
not assume honest linearizations of any individual line factor.

\subsection{5. The full character forces too large an invariant Euler
value}\label{SF-E815-the-full-character-forces-too-large-an-invariant-euler-value}

The character (E815.7) is unchanged by duality. In the ring of §2 it is

\[
 C=32+8(x^2+y^2)+2x^2y^2.
\]

The coefficient of \(x^2y^2\) in \(C^2\) is \(256\); (E815.3) therefore
gives

\[
 \chi_X(E,E)=1024k^2.
 \tag{E815.19}
\]

Equations (E815.16) and (E815.18) imply

\[
 \chi_{\mathcal C^G}(E,E)
       =\frac{1024k^2}{|G|}\ge128.
 \tag{E815.20}
\]

Suppose in addition that the equivariant object is gluable and weakly
semiregular, namely that its invariant self-Ext groups in negative
degrees vanish and its full invariant degree-two trace is injective.
Supplied S-21 invariant Calabi--Yau adjunction and the exact balanced
cap ranks from SF-E801 then give

\[
 e_2^G=18,
 \qquad e_1^G\ge8,
 \qquad e_{4-j}^G=e_j^G.
\]

Consequently

\[
 \chi_{\mathcal C^G}(E,E)
      =2e_0^G-2e_1^G+18\le2e_0^G+2.
 \tag{E815.21}
\]

Combining this with (E815.20) proves the necessary condition

\[
 \boxed{\dim\operatorname{End}_{\mathcal C^G}(E)\ge63.}
 \tag{E815.22}
\]

In particular, \textbf{no complex in this model, for any \(k\ge1\), is
simultaneously gluable, simple in the equivariant category, and weakly
semiregular}. The conclusion allows any effective finite target action,
including mixtures of translation and tensor operations. It also allows
all extra cancelling terms and projective multiplicities permitted in
§3.

The rank-32 polystable bundle (E815.4) is an actual balanced character,
but the proposed multivertex repair within these seven line classes
cannot produce the required simple equivariant seed. A nonsimple weakly
semiregular repair in this model would need at least 63 invariant
endomorphisms; this argument does not exclude it outright.

\subsection{6. General positive aggregate weights in the primitive RM
divisor
space}\label{SF-E815-general-positive-aggregate-weights-in-the-primitive-rm-divisor-space}

There is a useful extension independent of the special seven classes.
Let \(\mathcal B'\) be any finite set of integral algebraic divisor
classes whose pairwise differences are block diagonal and primitive on
both genuine RM planes. Let a perfect equivariant object have
vector-bundle cohomology filtered by line bundles of those classes and
arbitrary degree-zero twists. As above, the unequal-class Hom and
degree-one Ext vanish, so its cohomology admits actual canonical
equivariant class blocks. Form \(U_\beta\) as in (E815.13), with
aggregate virtual ranks \(n_\beta\).

Assume all \(n_\beta\ge0\), and let \(s\) be the number that are
nonzero. For every ordered unequal pair with nonzero ranks, the
signature \((2,2)\) calculation and (E815.16) give

\[
 \chi_{\mathcal C^G}(U_\beta,U_\gamma)
  =\frac{n_\beta n_\gamma}{|G|}
       \int_X\frac{(\gamma-\beta)^4}{24}
  \in\mathbf Z_{>0}.
\]

Self-pairings are zero, since each character is a multiple of a single
exponential. Zero-character blocks have zero Euler pairing with every
object. The equivariant Grothendieck-group identity (E815.14a), valid in
this general class-block setting as well, and bilinearity of the Euler
pairing therefore imply

\[
 \chi_{\mathcal C^G}(E,E)\ge s(s-1).
 \tag{E815.23}
\]

If the full character is balanced with nonzero coefficient, its
degree-two cap rank is 18. Each pure exponential has degree-two cap rank
six, so subadditivity forces \(s\ge3\). Hence (E815.23) gives invariant
Euler at least six. A gluable, equivariantly simple, weakly semiregular
balanced object would instead have invariant Euler at most four, by
(E815.21). This is impossible.

The same theorem permits \textbf{stable projectively flat vector-bundle
factors}, with rational normalized classes
\(\beta=c_1(Q)/\operatorname{rk}Q\), in place of line-bundle factors. We
give the adaptation explicitly. For two factors of unequal classes,
choose an integer \(m\) divisible by their ranks. For a factor \(Q\) of
rank \(r\), the actual line bundle \((\det Q)^{m^2/r}\) has the scalar
curvature of \([m]^*Q\), so the same harmonic-curvature and unitary
abelian-monodromy argument as in SF-E813 gives

\[
 [m]^*Q\simeq
 \bigoplus_{a=1}^{r}(L_Q\otimes P_a),
 \qquad c_1(L_Q)=m^2\beta.
 \tag{E815.24}
\]

For unequal classes, the line differences on this common cover have
signature \((2,2)\). The index theorem makes their Hom and degree-one
Ext zero. Finite étale trace injects the original groups into the
pulled-back groups, so the same two vanishings hold for the stable
factors. The two filtration inductions then give the same canonical
class-block decomposition of each cohomology sheaf, without requiring
that any one stable factor admit a \(G\)-linearization. A block of class
\(\beta\) has character equal to its \textbf{integer rank} times
\(e^\beta\), because the Chern character is additive through its
filtration. Its aggregate virtual coefficient \(n_\beta\) therefore
remains an integer, even though \(\beta\) is rational. Unequal positive
blocks have strictly positive ordinary cross Euler by the signature
calculation; exact scaling makes the invariant cross Euler a positive
integer. The proof of (E815.23) and the cap-rank contradiction is
consequently unchanged. The hypotheses remain locally free stable
projectively flat factors, common primitive differences, and nonnegative
aggregate virtual ranks.

Thus adding more line or projectively flat factor classes with
\textbf{nonnegative aggregate virtual weights}, while keeping all
differences in the common primitive RM space, cannot repair the model. A
construction that escapes this result must change at least one of those
precise hypotheses. Negative aggregate weights in an enlarged class set,
classes outside that primitive space, singular cohomology, twisted
objects, or a sufficiently nonsimple equivariant object are not excluded
by this theorem.

\subsection{Checks and source scope}\label{SF-E815-checks-and-source-scope}

The program rm\_seven\_class\_certificate.py actually computes the
determinant \(-512\), the entire character, self-Euler1024,
opposite-pair Euler16, the matrices \(U^2=2I\), \(US=SU\), and the
unimodularity of \(X_0,Z_+,Z_-\). Its JSON output explicitly disclaims
verification of categorical existence, actions, or trace injectivity.
The proof above gives those logical implications with their stated
hypotheses.

The published Mumford index input and its exact pages are already
recorded in SF-E813. The new categorical scaling (E815.16) is a direct
two-object deduction from the reread induction/forgetting and
Riemann--Roch argument of supplied DEG\_L S-21, equations
(S21.6)--(S21.7). The stable and semistable bundle obstructions
SF-E811--E813 are used only in their previously reviewed scopes. No
Bando--Siu reflexive extension is assumed.

\section{SF-E815R: the exact bound is at least 63 invariant
endomorphisms}\label{SF-E815R-sf-e815r-the-exact-bound-is-at-least-63-invariant-endomorphisms}

Author/checker: GPT-6 Astra Pro; inherited effort not exposed;
HODGE\_DIVISION /root/sixfolds. Date: 2026-11-21. This is a quantitative
refinement of the already reviewed SF-E815, not a new main-gate
exclusion. Root read the entire written refinement and reported
whole-proof mathematical PASS on2026-11-21 (GPT-6 Astra Pro; inherited
effort not exposed), and freshly reran the extended exact certificate.
The author also freshly reran the certificate. This metadata does not
claim a new independent audit of the entirety of the retained S-21
input.

Retain all hypotheses and notation of the seven-class model in SF-E815,
with \(\operatorname{ch}(E)=kC\), \(k\ge1\). This permits all the
cancelling terms, arbitrary cohomological degrees and projective
multiplicity representations allowed there. If \(E\) is gluable and
weakly semiregular for its specified effective finite coherent
translation/tensor action \(G\), then

\[
 |G|\mid8k^2,\qquad
 \chi_G(E,E)\ge128,\qquad
 \dim\operatorname{End}_G(E)\ge63.
 \tag{E815R.1}
\]

In particular the previous lower bound 31 can be replaced by 63.

\textbf{Proof.} SF-E815 constructs actual equivariant perfect objects
\(U_\beta\) by taking alternating direct sums of the canonical
cohomology class blocks. It proves the actual equivariant
Grothendieck-group identity \([E]=\sum_\beta[U_\beta]\), without
asserting that the \(U_\beta\) are summands of \(E\). Its exact
seven-character independence gives

\[
 \operatorname{ch}(U_{X_0})=2k e^{X_0},\qquad
 \operatorname{ch}(U_{Z_+})=k e^{Z_+}.
\]

Their divisor difference is \(Z_+-X_0=Y\). The Hermitian matrix of \(Y\)
is \(US\); its determinant is four, equivalently its real-quadratic norm
is \(-2\) and the fourfold Euler is its square. Hence Riemann--Roch on
the abelian fourfold gives

\[
 \chi_X(U_{X_0},U_{Z_+})
     =2k^2\int_X\frac{Y^4}{24}=8k^2.
 \tag{E815R.2}
\]

The exact two-object induction/forgetting calculation of supplied S-21,
used and proved in SF-E815, gives

\[
 \chi_X(U,V)=|G|\chi_G(U,V)
\]

for these actual equivariant objects. The integer
\(\chi_G(U_{X_0},U_{Z_+})\) is positive, so (E815R.2) gives
\(|G|\mid8k^2\). The entire character has ordinary self-Euler
\(1024k^2\). Therefore

\[
 \chi_G(E,E)=\frac{1024k^2}{|G|}\ge128.
\]

Weak degree-two semiregularity and invariant CY4 adjunction give
\(e_2^G=18\). Hochschild degree-one cap has rank eight, hence
\(e_1^G\ge8\). Gluability and invariant Serre duality give

\[
 \chi_G(E,E)=2e_0^G-2e_1^G+18\le2e_0^G+2.
\]

Combining these inequalities gives \(128\le2e_0^G+2\), or
\(e_0^G\ge63\), proving (E815R.1).

The extended exact script rm\_seven\_class\_certificate.py actually
verifies the unweighted pair Euler four, the weighted value eight at
\(k=1\), and the ratio 128, as well as its previous complete character
and determinant checks. It does not verify categorical existence or
equivariance. Those implications are the proof above and the retained
full scope of SF-E815.

\section{SF-E816: disjoint-class two-cohomology cones cannot be balanced
RM
seeds}\label{SF-E816-sf-e816-disjoint-class-two-cohomology-cones-cannot-be-balanced-rm-seeds}

Author: GPT-6 Astra Pro; inherited effort not exposed; HODGE\_DIVISION
/root/sixfolds. Date: 2026-11-21.

\textbf{Status:} Whole-proof mathematical PASS reviews of §§1--6 were
received from /root and /root/imprimitive on 2026-11-21 (GPT-6 Astra
Pro; inherited effort not exposed). Root additionally read both complete
written corollaries in §§7--8 and reported whole-proof PASS. Imprimitive
independently read and approved the complete written SF-E816a proof, and
verified that SF-E816b faithfully records his own complete derivation.
The latter is contributor verification; root and sixfolds supplied its
independent mathematical checks. The ordinary-gluable overlap argument
in §7 was supplied by /root; the arbitrary-amplitude Postnikov argument
in §8 was supplied by /root/imprimitive. Neither contribution is being
promoted from an unwritten assertion: all proofs and maps are given
below. The results permit negative virtual weights and nonsimple
objects, within their precise factor and class hypotheses.

\subsection{1. Statement and exact
scope}\label{SF-E816-statement-and-exact-scope}

Let \(X\) be a complex abelian fourfold with genuine real multiplication
by a quadratic field \(F\), self-adjoint for a polarization \(\theta\),
with holomorphic multiplicities \((2,2)\). Write

\[
 \theta=\theta_++\theta_-
\]

for its two real eigenspace components. Let
\(G\subset X\times\widehat X\) be an effective finite subgroup with a
specified coherent translation/tensor action of the kind in supplied
DEG\_L S-21. All equivariant objects below belong to this actual
coherent action. No individual stable factor is assumed to admit its own
linearization.

Let \(E\) be a perfect equivariant object whose only nonzero cohomology
sheaves are the vector bundles

\[
 A=\mathcal H^0(E),\qquad B=\mathcal H^{-1}(E).
 \tag{E816.1}
\]

Assume that \(A\) and \(B\) admit filtrations by stable projectively
flat vector bundles. For a factor \(Q\), call
\(\beta(Q)=c_1(Q)/\operatorname{rk}Q\) its normalized divisor class.
Require the following two conditions:

\begin{enumerate}
\def\labelenumi{\arabic{enumi}.}
\item
  For any two unequal normalized classes occurring in either filtration,
  their difference is rational, block diagonal for the two
  real-multiplication planes, and primitive on each plane.
\item
  The sets of normalized classes occurring in \(A\) and in \(B\) are
  disjoint.
\end{enumerate}

Line-bundle factors, with arbitrary degree-zero twists, are included.
Factors and multiplicities may differ between the two sheaves;
extensions within each normalized class are unrestricted.

Then it is impossible that both

\[
 e^D\operatorname{ch}(E)
  =N\left(1-\frac q2\theta_+^2\right)
      \left(1-\frac q2\theta_-^2\right),
 \qquad Nq\ne0,
 \tag{E816.2}
\]

for a rational algebraic divisor \(D\), and that the full invariant
degree-two trace of \(E\) be injective. In particular no such object is
a weakly \(G\)-semiregular balanced RM input, regardless of its
equivariant endomorphism dimension. The parameter \(q\) may be positive,
as in the intended input; the proof only needs it to be nonzero.

The theorem applies to every actual triangle

\[
 B[1]\xrightarrow{i}E\xrightarrow{p}A
       \xrightarrow{e}B[2],
 \qquad e\in\operatorname{Ext}^2_G(A,B),
 \tag{E816.3}
\]

with the stated factors. It does not assume any formality of a dg or
\(A_\infty\) category. In particular it excludes the natural signed
construction with all positive-class vertices in \(A\) and all distinct
negative-class vertices in \(B[1]\), for every matrix of actual
degree-two gluing classes.

\subsection{2. Cross-Ext concentration and canonical class
blocks}\label{SF-E816-cross-ext-concentration-and-canonical-class-blocks}

We give the requisite factor and filtration arguments rather than infer
them from the Chern character. If a rational nonzero divisor difference
\(\delta\) in hypothesis 1 vanished on one real-multiplication plane,
its associated rational alternating form would be supported on the
other. For an irrational field element \(u\in F\), one would have

\[
 \delta(ux,y)=\sigma(u)\delta(x,y)
\]

for all rational first-homology vectors \(x,y\), with \(\sigma(u)\)
irrational. Both the left side and \(\delta(x,y)\) are rational.
Choosing a pair with nonzero value contradicts irrationality. Thus both
blocks of \(\delta\) are nonzero. A nonzero primitive Hermitian form on
a complex two-plane has eigenvalues \(\lambda,-\lambda\), and hence
signature \((1,1)\). The complete form of \(\delta\) is nondegenerate of
signature \((2,2)\).

For stable projectively flat factors \(Q,R\), choose a positive integer
\(m\) divisible by their ranks. We use the actual finite-isogeny
splitting proved in SF-E813. Its construction is as follows. The stable
factor has a Hermitian--Einstein connection for the polarization; its
zero discriminant and the ordinary Chern--Weil energy identity make that
connection projectively flat. For a rank-\(r\) factor \(Q\), the actual
line bundle

\[
 L_Q=(\det Q)^{m^2/r}
\]

has the same harmonic scalar curvature as \([m]^*Q\). Consequently
\([m]^*Q\otimes L_Q^{-1}\) is unitary flat. Its monodromy representation
is a unitary representation of the abelian fundamental group, so it
splits into characters. Therefore

\[
 [m]^*Q\simeq\bigoplus_{a=1}^{r}(L_Q\otimes P_a),
 \qquad P_a\in\operatorname{Pic}^0(X),
 \qquad c_1(L_Q)=m^2\beta(Q).
 \tag{E816.4}
\]

The same holds for \(R\) on this one cover. If their normalized classes
are unequal, every line difference on the cover has signature \((2,2)\).
The published Mumford index theorem gives cohomology only in degree two,
including after every degree-zero twist. Finite étale pullback on Ext is
injective: its composition with trace is multiplication by
\(\deg[m]\ne0\). It follows that

\[
 \operatorname{Ext}^j(Q,R)=0\quad(j\ne2)
 \tag{E816.5}
\]

for unequal normalized classes. Negative Ext is already zero for
sheaves; the assertion above degree four follows as well from the cover,
or from Serre duality. The source is Mumford, \emph{Abelian Varieties},
§16, the vanishing theorem on printed page 150, proof on pages 150--153,
and the complex Hermitian-index corollary and proof on pages 162--163.
The Hermitian--Einstein input is Siu, \emph{Lectures on
Hermitian--Einstein Metrics for Stable Bundles and Kähler--Einstein
Metrics}, Chapter 1, §6.1--§6.4, in the rational polarization scope
already recorded in SF-E811. No reflexive-sheaf extension is used.

The vanishing extends through filtrations in both arguments. Indeed, for
a short exact sequence in one argument, its Ext long exact sequence
shows that concentration in degree two for the two end terms implies the
same concentration for the middle term. Induction first through one
filtration and then through the other proves (E816.5) for any two vector
bundles filtered entirely by unequal classes.

In particular their Hom and degree-one Ext vanish. These two vanishings
give actual canonical class decompositions

\[
 A=\bigoplus_{\beta\in\mathcal A}A_\beta,
 \qquad
 B=\bigoplus_{\gamma\in\mathcal B}B_\gamma,
 \qquad\mathcal A\cap\mathcal B=\varnothing,
 \tag{E816.6}
\]

where each nonzero block is filtered entirely by its indicated
normalized class. To see existence, induct through a filtration. If its
last quotient has class \(\beta\), the extension by every already
separated unequal-class block is zero in Ext degree one. Only the
\(\beta\)-block can extend. This produces the direct sum in (E816.6). To
see its canonicity, maps between two unequal-class blocks are zero. The
projection of any second proposed \(\beta\)-block onto every different
block is therefore zero, so the \(\beta\)-summand is uniquely determined
as a subobject.

Translation and tensoring by \(\operatorname{Pic}^0\) preserve each
rational normalized class, and are exact for the usual cohomological
structure. The coherent action on \(E\) gives coherent actions on
\(A,B\). Canonicity of (E816.6) then gives actual coherent equivariant
structures on every \(A_\beta,B_\gamma\). No equivariant splitting of
their chosen filtrations is asserted or needed.

Since the class sets are disjoint, the filtration argument proves

\[
 \operatorname{Ext}^j_G(A,B)
   =\operatorname{Ext}^j_G(B,A)=0\quad(j\ne2).
 \tag{E816.7}
\]

Here \(\operatorname{Ext}^j_G\) is the invariant Ext group for the
specified coherent action. Over \(\mathbf C\), finite-group invariants
are exact, so the ordinary vanishing implies (E816.7), and all the
ensuing distinguished-triangle long exact sequences hold in the
equivariant category itself.

\subsection{3. Exact survival of both middle self-Ext
groups}\label{SF-E816-exact-survival-of-both-middle-self-ext-groups}

We compute the relevant maps in (E816.3), without any assumption on the
class \(e\). Apply \(\operatorname{RHom}_G(A,-)\). For \(j=2,3\), the
exact segment is

\[
 \operatorname{Ext}^{j+1}_G(A,B)
  \longrightarrow\operatorname{Ext}^{j}_G(A,E)
  \xrightarrow{p_*}\operatorname{Ext}^{j}_G(A,A)
  \longrightarrow\operatorname{Ext}^{j+2}_G(A,B).
\]

Both exterior terms vanish by (E816.7), so

\[
 \operatorname{Ext}^{j}_G(A,E)
       \simeq\operatorname{Ext}^{j}_G(A,A)
       \qquad(j=2,3).
 \tag{E816.8}
\]

Next apply \(\operatorname{RHom}_G(B[1],-)\). Its exact segment is

\[
 \operatorname{Ext}^{j-2}_G(B,A)
  \longrightarrow\operatorname{Ext}^{j}_G(B,B)
  \xrightarrow{i_*}\operatorname{Ext}^{j}_G(B[1],E)
  \longrightarrow\operatorname{Ext}^{j-1}_G(B,A).
\]

For \(j=1,2\), the exterior groups again vanish. Thus

\[
 i_*:\operatorname{Ext}^{j}_G(B,B)
       \xrightarrow{\sim}\operatorname{Ext}^{j}_G(B[1],E)
       \qquad(j=1,2).
 \tag{E816.9}
\]

Finally apply \(\operatorname{RHom}_G(-,E)\). The complete segment
surrounding the middle self-Ext is

\[
 \begin{split}
 \operatorname{Ext}^1_G(B[1],E)
 &\xrightarrow{\partial_1}\operatorname{Ext}^2_G(A,E)
 \longrightarrow\operatorname{Ext}^2_G(E,E)\\
 &\longrightarrow\operatorname{Ext}^2_G(B[1],E)
 \xrightarrow{\partial_2}\operatorname{Ext}^3_G(A,E).
 \end{split}
 \tag{E816.10}
\]

Both displayed connecting maps are zero. More explicitly, (E816.9)
writes any class of degree \(j=1,2\) as

\[
 f=i[j]\circ a[1],
 \qquad a:B\longrightarrow B[j].
\]

Up to the irrelevant conventional sign, its connecting class is

\[
 \partial_j(f)=f[1]\circ e
   =i[j+1]\circ a[2]\circ e.
 \tag{E816.11}
\]

The intermediate composite \(a[2]\circ e\) lies in
\(\operatorname{Ext}^{j+2}_G(A,B)\). This is zero for both \(j=1\) and
\(j=2\), by (E816.7). The calculation takes place among actual
equivariant morphisms and does not replace invariance by a dimension
count.

Equations (E816.8)--(E816.11) prove the short exact sequence

\[
 0\longrightarrow\operatorname{Ext}^2_G(A,A)
  \longrightarrow\operatorname{Ext}^2_G(E,E)
  \longrightarrow\operatorname{Ext}^2_G(B,B)
  \longrightarrow0.
 \tag{E816.12}
\]

There is no claim of a canonical splitting. In particular,

\[
 \operatorname{ext}^2_G(E,E)
   =\operatorname{ext}^2_G(A,A)
     +\operatorname{ext}^2_G(B,B).
 \tag{E816.13}
\]

Thus every actual degree-two gluing class leaves these two middle
self-Ext dimensions intact. The equality also holds for the zero gluing
class. Neither simplicity nor gluability is needed for this equality.

\subsection{4. A balanced four-vector needs at least four pure
terms}\label{SF-E816-a-balanced-four-vector-needs-at-least-four-pure-terms}

We include the exterior-algebra proof in full. Put

\[
 W=H^0(X,T_X)\oplus H^1(X,\mathcal O_X),
 \qquad\dim_{\mathbf C}W=8.
\]

Choose a nonzero holomorphic volume form. The resulting Fock
identification sends a diagonal Hodge character \(v\) to a four-vector
\(\widetilde v\in\bigwedge^4W\), and identifies degree-two cap with

\[
 c_v^2:\bigwedge^2W\longrightarrow\bigwedge^6W,
 \qquad z\longmapsto z\wedge\widetilde v.
 \tag{E816.14}
\]

An exponential \(e^\beta\) becomes a nonzero decomposable four-vector.
Indeed multiplication by the exponential is the invertible divisor shear
applied to the decomposable Fock representative of the unit class. If
\(w\) is a nonzero decomposable four-vector supported on a four-plane
\(L\), choose a complementary four-plane \(M\). Only \(\bigwedge^2M\)
contributes to wedge multiplication by \(w\), and its map is injective.
Hence

\[
 \operatorname{rank}c_w^2=\dim\bigwedge^2M=6.
 \tag{E816.15}
\]

For the nonzero balanced character in (E816.2), the Fock vector has the
form

\[
 v=\omega_+\wedge\omega_-,
 \qquad W=W_+\oplus W_-,\qquad\dim W_\pm=4,
 \tag{E816.16}
\]

where each \(\omega_\pm\in\bigwedge^2W_\pm\) is nondegenerate. Each
two-plane character \(1-q\theta_\pm^2/2\) gives the sum of nonzero
volume bivectors on its translation and tensor two-planes; its square is
nonzero because \(q\ne0\). Overall nonzero scalars do not matter.

Define

\[
 K_\pm=\{z\in\bigwedge^2W_\pm:z\wedge\omega_\pm=0\}.
\]

Each \(K_\pm\) has dimension five. On \(\bigwedge^2W_+\) and
\(\bigwedge^2W_-\), wedge multiplication by \(v\) has rank one in the
distinct bidegrees \((4,2)\) and \((2,4)\). On \(W_+\otimes W_-\) it is
the tensor product of the isomorphisms \(x\mapsto x\wedge\omega_\pm\),
and has rank sixteen in bidegree \((3,3)\). Consequently

\[
 \operatorname{rank}c_v^2=18,
 \qquad\ker c_v^2=K_+\oplus K_-=:K.
 \tag{E816.17}
\]

Its common annihilator in four-vector degree is exactly one-dimensional:

\[
 \{w\in\bigwedge^4W:z\wedge w=0\text{ for every }z\in K\}
       =\mathbf C v.
 \tag{E816.18}
\]

Here is a direct proof. Decompose \(w\) into bidegrees \((p,4-p)\) for
\(W_+,W_-\). The \(K_+\) equations kill the components with \(p=0,1\).
For \(p=0\) this is immediate. For \(p=1\), the bivectors annihilating a
fixed nonzero vector in a four-space form its three-dimensional wedge
subspace; they cannot contain the five-dimensional \(K_+\). Applying
linear functionals to the other tensor factor gives the same assertion
for a tensor-valued component. The \(K_-\) equations likewise kill
\(p=3,4\). For \(p=2\), the wedge pairing on \(\bigwedge^2W_+\) is
nondegenerate, so annihilation by its primitive hyperplane forces the
first factor to lie in \(\mathbf C\omega_+\). The other condition forces
the second factor to lie in \(\mathbf C\omega_-\). This gives precisely
(E816.18).

Now suppose that a nonzero balanced four-vector were a sum of \(s\)
nonzero pure four-vectors, with arbitrary nonzero scalar coefficients.
For \(s\le2\), the rank bound \(18\le6s\) from (E816.15) is impossible.
If \(s=3\), absorb coefficients into the pure terms and write

\[
 v=w_1+w_2+w_3.
\]

Let \(S_i=\operatorname{im}c_{w_i}^2\), each of dimension six. Since

\[
 \operatorname{im}c_v^2\subset S_1+S_2+S_3,
\]

and the left side has dimension eighteen, the sum on the right has
dimension eighteen and is direct. Therefore

\[
 c_v^2(z)=0
 \quad\Longleftrightarrow\quad
 c_{w_1}^2(z)=c_{w_2}^2(z)=c_{w_3}^2(z)=0.
 \tag{E816.19}
\]

Each \(w_i\) is thus annihilated by \(K=\ker c_v^2\). Equation (E816.18)
forces each \(w_i\) to be proportional to \(v\), contradicting their cap
ranks six and eighteen. This direct-sum image argument establishes the
kernel implication explicitly; no unstated image/kernel duality is being
used. We have proved

\[
 \boxed{\text{every nonzero balanced character needs at least four pure exponential terms.}}
 \tag{E816.20}
\]

The conclusion allows signed, rational, real, or complex nonzero
coefficients. It is a statement about the actual number of pure terms in
a presentation, not an independence assumption on those terms. A common
divisor shear is invertible, preserves decomposability and all the
displayed cap ranks, and therefore preserves the conclusion.

\subsection{5. Invariant trace dimensions and the
contradiction}\label{SF-E816-invariant-trace-dimensions-and-the-contradiction}

For any nonzero actual equivariant class block \(A_\beta\), additivity
through its filtration gives

\[
 \operatorname{ch}(A_\beta)=r_\beta e^\beta,
 \qquad r_\beta=\operatorname{rk}A_\beta>0.
\]

Hochschild evaluation is natural for the translation/tensor action. That
action is trivial on Hochschild cohomology: translations and degree-zero
tensor classes act trivially. Evaluation therefore takes values in
\(\operatorname{Ext}^2_G(A_\beta,A_\beta)\). Its composition with the
full trace is cap product with \(r_\beta e^\beta\). Equation (E816.15)
implies

\[
 \operatorname{ext}^2_G(A_\beta,A_\beta)\ge6.
 \tag{E816.21}
\]

This does not require trace injectivity of the block. It is a rank lower
bound from the actual invariant evaluation. The same applies to every
\(B_\gamma\). Because (E816.6) consists of actual coherent equivariant
summands, its diagonal self-Ext groups are direct summands of the
self-Ext of \(A\) and \(B\). Put \(s=|\mathcal A|+|\mathcal B|\).
Equations (E816.13) and (E816.21) give

\[
 \operatorname{ext}^2_G(E,E)\ge6s.
 \tag{E816.22}
\]

On the other hand,

\[
 \operatorname{ch}(E)
   =\sum_{\beta\in\mathcal A}r_\beta e^\beta
       -\sum_{\gamma\in\mathcal B}t_\gamma e^\gamma,
 \qquad r_\beta,t_\gamma>0.
 \tag{E816.23}
\]

The disjoint-class hypothesis ensures that these are \(s\) actual pure
terms with nonzero signed coefficients. Under the common shear \(e^D\),
they are still pure terms. If (E816.2) holds, (E816.20) therefore gives
\(s\ge4\), and so

\[
 \operatorname{ext}^2_G(E,E)\ge24.
 \tag{E816.24}
\]

Finally suppose the invariant full degree-two trace of \(E\) were
injective. The CY4 trace/evaluation adjunction in supplied S-21 makes
the invariant Hochschild evaluation surjective: the two maps are adjoint
for the perfect Serre pairing on invariant degree-two Ext and the
perfect Hochschild pairing. The Serre pairing remains perfect on
invariants because the finite translation/tensor action preserves it and
averaging is exact. The identity-sector normalization may multiply the
trace/cap identity by a nonzero scalar depending on \(|G|\); it changes
no rank.

Thus injectivity implies that the trace image is exactly the image of
cap product with \(\operatorname{ch}(E)\), and gives

\[
 \operatorname{ext}^2_G(E,E)
    =\operatorname{rank}c^2_{\operatorname{ch}(E)}=18,
 \tag{E816.25}
\]

by (E816.17) and invariance under the divisor shear. This contradicts
(E816.24), proving the theorem. There was no equivariant simplicity
assumption and no Euler-divisibility argument in the proof.

\subsection{6. What has and has not been
removed}\label{SF-E816-what-has-and-has-not-been-removed}

SF-E815 excludes equivariantly simple repairs with nonnegative aggregate
weights in the stated primitive class models. The present result
excludes a different proposed escape: negative aggregate weights
implemented by putting disjoint class sets in two adjacent cohomological
degrees and joining them by arbitrary degree-two extensions. It excludes
these objects even when they are nonsimple and their finite action has
projective multiplicity spaces.

If the same normalized class occurs in both \(A\) and \(B\), the cross
groups in degrees one, three or four need not vanish without further
information. The connecting maps in (E816.10) could then be nonzero.
Corollary SF-E816a below shows that ordinary gluability supplies the
needed further information and removes the disjoint-class assumption for
two cohomology sheaves. Objects with three or more nonzero cohomology
sheaves have additional Postnikov data; corollary SF-E816b treats them
when the normalized class sets are globally disjoint. General
overlapping-class complexes with more cohomology sheaves, singular
stable factors, twisted categories, and classes outside the common
primitive real-multiplication space remain outside these results. None
is inferred to exist or to be semiregular.

The proof is an exact long-exact-sequence and exterior-algebra argument.
No new computer calculation is needed or claimed for the universal
theorem. The already recorded exact exterior calculations for SF-E801
independently corroborate the balanced rank eighteen and its
one-dimensional common annihilator; the written derivations above are
the proofs used here.

\subsection{7. SF-E816a: ordinary gluability removes normalized-class
disjointness for two
sheaves}\label{SF-E816-sf-e816a-ordinary-gluability-removes-normalized-class-disjointness-for-two-sheaves}

Retain the setting of §1, including the primitive-difference condition
for unequal normalized classes, but drop condition 2. Instead assume
that \(E\) is \textbf{ordinary gluable}:

\[
 \operatorname{Ext}^j_X(E,E)=0\qquad(j<0).
 \tag{E816.26}
\]

Then the same impossibility theorem holds, even when \(A\) and \(B\)
have overlapping normalized classes and cancelling virtual terms. It
still does not require equivariant simplicity. The ordinary condition
(E816.26), rather than only its invariant counterpart, is the hypothesis
used in this corollary.

\subsubsection{7.1. Gluability detects the extreme cohomology
map}\label{SF-E816-gluability-detects-the-extreme-cohomology-map}

There is a canonical isomorphism

\[
 \operatorname{Hom}_{D^b(X)}(E,E[-1])
       \simeq\operatorname{Hom}_X(A,B).
 \tag{E816.27}
\]

Indeed \(E\) lies in \(D^{\le0}\), whereas \(E[-1]\) lies in
\(D^{\ge0}\). A map between these two objects factors uniquely through
their degree-zero cohomology sheaves. Their degree-zero cohomology
sheaves are \(A\) and \(B\), respectively. This proves (E816.27)
directly from the standard cohomological structure. Ordinary gluability
therefore gives

\[
 \operatorname{Hom}_X(A,B)=0.
 \tag{E816.28}
\]

This conclusion is independent of the Postnikov invariant \(e\).

\subsubsection{7.2. Stable factors of the same normalized
class}\label{SF-E816-stable-factors-of-the-same-normalized-class}

Let \(Q,R\) be stable projectively flat vector bundles of the same
normalized class \(\beta\). Choose the Hermitian--Einstein metrics used
in (E816.4). Their scalar curvature forms are the harmonic
representative of the same class \(\beta\), so the induced connection on
\(Q^\vee\otimes R\) is unitary flat. Since \(\pi_1(X)\) is abelian, its
unitary representation splits into one-dimensional characters. Thus
\(Q^\vee\otimes R\) is an actual direct sum of unitary flat line
bundles.

A nontrivial unitary flat line bundle on a complex torus has no
cohomology. For completeness, equip the torus with its flat Kähler
metric. The Laplacian on a flat line-valued \((0,j)\)-form is the
connection Laplacian, since both the metric curvature and the line
curvature vanish. A harmonic form is therefore parallel by integration
of the squared covariant derivative. The tangent bundle is parallel
trivial, so a nontrivial line character admits no such forms in any
degree. The trivial character contributes precisely
\(H^j(X,\mathcal O_X)\). Consequently the parallel-invariant subspace
gives a natural graded isomorphism

\[
 \operatorname{Ext}^*_X(Q,R)
     \simeq\operatorname{Hom}_X(Q,R)\otimes
                    H^*(X,\mathcal O_X).
 \tag{E816.29}
\]

Stable vector bundles of the same slope have no nonzero Hom unless they
are isomorphic; any nonzero map is an isomorphism by the slope
inequalities for its image. Thus two nonisomorphic stable factors of the
same normalized class have

\[
 \operatorname{Ext}^j_X(Q,R)=0\quad\text{for every }j.
 \tag{E816.30}
\]

No assertion of individual equivariant structures is made here. This is
an ordinary vanishing for the actual stable bundles.

\subsubsection{7.3. Separate stable isomorphism types and exclude a
shared
type}\label{SF-E816-separate-stable-isomorphism-types-and-exclude-a-shared-type}

Within each fixed normalized-class block of \(A\) or \(B\), use
(E816.30) and the same filtration induction as in §2. The block has an
actual canonical decomposition into pieces filtered by one stable
isomorphism type each. Ext degree one gives existence of this
decomposition; Hom vanishing between distinct types gives uniqueness.
Denote the pieces of type \(Q\) by \(A_Q\) and \(B_Q\), when present.

If both pieces were nonzero, a last filtration quotient would give a
surjection \(A_Q\twoheadrightarrow Q\), and a first filtration subobject
would give an injection \(Q\hookrightarrow B_Q\). Their composite is a
nonzero morphism \(A_Q\to B_Q\), and hence a nonzero morphism
\(A\to B\), contradicting (E816.28). Thus no stable isomorphism type
occurs in both sheaves.

For factors of equal normalized class across \(A\) and \(B\), (E816.30)
now gives complete Ext vanishing. For factors of unequal normalized
class, (E816.5) gives concentration in degree two. Two filtration
inductions therefore give

\[
 \operatorname{Ext}^j_X(A,B)
       =\operatorname{Ext}^j_X(B,A)=0\quad(j\ne2),
 \tag{E816.31}
\]

and taking invariants gives (E816.7). The actual equivariant triangle
proof of §3 applies unchanged. Notice that the stable-type
decompositions were used only to prove ordinary vanishing. They need not
be preserved one by one by \(G\); only the normalized-class
decompositions used in the dimension bound need to be equivariant, and
those are canonical as before.

\subsubsection{7.4. The character count still gives the
contradiction}\label{SF-E816-the-character-count-still-gives-the-contradiction}

Let \(s\) now be the number of nonzero normalized-class blocks in \(A\),
plus the number in \(B\), counting a class twice if it occurs in both.
Their characters form the actual signed pure-term presentation
(E816.23), with nonzero positive ranks before the signs, even if some
terms repeat or partially cancel. Multiplication by \(e^D\) preserves
purity. The at-least-four-term result of §4 applies to this presentation
without any independence or distinctness requirement, so a balanced
nonzero character forces \(s\ge4\).

Each of these \(s\) actual equivariant blocks contributes at least six
to the sum in (E816.13), by the invariant evaluation argument in §5.
Thus \(\operatorname{ext}^2_G(E,E)\ge24\). Weak invariant semiregularity
would again give exactly eighteen. This proves SF-E816a, with arbitrary
normalized overlaps and cancellations in the two-cohomology model.

\subsection{8. SF-E816b: arbitrary amplitude with globally disjoint
class
sets}\label{SF-E816-sf-e816b-arbitrary-amplitude-with-globally-disjoint-class-sets}

There is a separate extension with no gluability assumption. Let \(E\)
be any bounded perfect equivariant complex whose cohomology sheaves are
vector bundles filtered by stable projectively flat factors. Assume the
primitive-difference condition from §1 for all unequal normalized
classes occurring anywhere. Assume that the normalized class sets of any
two distinct nonzero cohomology sheaves are disjoint.

Then

\[
 \operatorname{ext}^2_G(E,E)
      =\sum_j\operatorname{ext}^2_G(\mathcal H^j(E),\mathcal H^j(E)).
 \tag{E816.32}
\]

If its character is nonzero balanced after a common divisor shear, this
equality rules out weak invariant semiregularity, regardless of
simplicity or the number of cohomology sheaves. We prove (E816.32) using
actual Postnikov triangles, rather than assume degeneration of a
self-Ext spectral sequence.

Shift the top nonzero cohomological degree to zero, put

\[
 A=\mathcal H^0(E),\qquad F=\tau_{\le-1}E,
\]

and use the canonical equivariant triangle

\[
 F\xrightarrow{i}E\xrightarrow{p}A\xrightarrow{e}F[1].
 \tag{E816.33}
\]

Every successive Postnikov piece of \(F\) is \(\mathcal H^{-k}(E)[k]\)
for some \(k\ge1\). By the disjointness hypothesis and §2, cross Ext
between \(A\) and \(\mathcal H^{-k}(E)\), in either direction, is
concentrated in degree two. It follows that

\[
 \begin{split}
 \operatorname{Ext}^j_G(A,\mathcal H^{-k}(E)[k])
       &=\operatorname{Ext}^{j+k}_G(A,\mathcal H^{-k}(E))=0
              &&(j\ge2),\\
 \operatorname{Ext}^j_G(\mathcal H^{-k}(E)[k],A)
       &=\operatorname{Ext}^{j-k}_G(\mathcal H^{-k}(E),A)=0
              &&(j\le2).
 \end{split}
 \tag{E816.34}
\]

For each fixed \(j\) in the indicated range, vanishing is preserved by
the distinguished-triangle long exact sequence. Induction through the
finite Postnikov tower therefore gives

\[
 \operatorname{Ext}^j_G(A,F)=0\quad(j\ge2),
 \qquad
 \operatorname{Ext}^j_G(F,A)=0\quad(j\le2).
 \tag{E816.35}
\]

Apply \(\operatorname{RHom}_G(A,-)\) to (E816.33). The exterior terms
vanish by (E816.35), giving

\[
 \operatorname{Ext}^j_G(A,E)\simeq\operatorname{Ext}^j_G(A,A)
            \qquad(j=2,3).
 \tag{E816.36}
\]

Apply \(\operatorname{RHom}_G(F,-)\). The two adjacent groups with
target \(A\) vanish, giving

\[
 i_*:\operatorname{Ext}^j_G(F,F)
         \xrightarrow{\sim}\operatorname{Ext}^j_G(F,E)
            \qquad(j=1,2).
 \tag{E816.37}
\]

Finally \(\operatorname{RHom}_G(-,E)\) gives the exact segment

\[
 \operatorname{Ext}^1_G(F,E)\xrightarrow{\partial_1}
 \operatorname{Ext}^2_G(A,E)\longrightarrow
 \operatorname{Ext}^2_G(E,E)\longrightarrow
 \operatorname{Ext}^2_G(F,E)\xrightarrow{\partial_2}
 \operatorname{Ext}^3_G(A,E).
 \tag{E816.38}
\]

Represent any degree-\(j\) class, for \(j=1,2\), by \(f=i[j]\circ a\),
with \(a:F\to F[j]\), using (E816.37). Its connecting image is, up to
sign,

\[
 \partial_j(f)=f[1]\circ e
   =i[j+1]\circ a[1]\circ e.
 \tag{E816.39}
\]

The middle composite belongs to \(\operatorname{Ext}^{j+1}_G(A,F)=0\),
by (E816.35). Hence both connecting maps vanish, and
(E816.36)--(E816.38) give

\[
 0\longrightarrow\operatorname{Ext}^2_G(A,A)
    \longrightarrow\operatorname{Ext}^2_G(E,E)
    \longrightarrow\operatorname{Ext}^2_G(F,F)
    \longrightarrow0.
 \tag{E816.40}
\]

Induction on the number of nonzero cohomology sheaves proves (E816.32).
A shift of a sheaf has the same self-Ext groups, which supplies the base
case.

Each cohomology sheaf has the actual canonical equivariant
normalized-class blocks of §2. Let \(s\) be the total number of such
nonzero blocks. Their signed characters, with signs determined by their
cohomological degrees, present \(\operatorname{ch}(E)\) as \(s\) nonzero
pure terms. Global class disjointness ensures they occur in only one
cohomological degree, but no additional independence claim is needed. If
the full character is balanced, §4 gives \(s\ge4\). Equations (E816.21)
and (E816.32) then give \(\operatorname{ext}^2_G(E,E)\ge6s\ge24\),
contradicting (E816.25).

This proves SF-E816b for arbitrary bounded amplitude and arbitrary
actual higher Postnikov invariants. The general case in which normalized
classes overlap between several different cohomological degrees is not
settled by this argument.

\section{SF-E816c--d: the two-degree gap and a limit to the overlap
argument}\label{SF-E816cd-sf-e816cd-the-two-degree-gap-and-a-limit-to-the-overlap-argument}

Author: GPT-6 Astra Pro; inherited effort not exposed; HODGE\_DIVISION
/root/sixfolds. Date: 2026-11-21. The gap observation in SF-E816c was
supplied by /root. The explicit counterexample in SF-E816d was supplied
by /root/imprimitive. Both are written and independently checked here.
Root read the complete written addendum and reported whole-proof PASS.
Imprimitive independently read and approved SF-E816c; his read of
SF-E816d is contributor verification, while root and sixfolds supply its
independent checks. All reviews are dated2026-11-21, GPT-6 Astra Pro,
inherited effort not exposed.

\subsection{SF-E816c. Any two distinct cohomological
degrees}\label{SF-E816cd-sf-e816c.-any-two-distinct-cohomological-degrees}

The ordinary-gluable obstruction SF-E816a holds when the two nonzero
cohomology sheaves occur in any two distinct degrees, not necessarily
adjacent. Keep all its other factor, rational primitive-difference,
finite coherent action, and balanced-character hypotheses. After an
overall shift, write the two sheaves as

\[
 A=\mathcal H^0(E),\qquad B=\mathcal H^{-k}(E),\qquad k\ge1.
 \tag{E816.41}
\]

Because \(E\in D^{\le0}\) and \(E[-k]\in D^{\ge0}\), their extreme map
is canonically

\[
 \operatorname{Ext}^{-k}_X(E,E)
      =\operatorname{Hom}(E,E[-k])
      \simeq\operatorname{Hom}_X(A,B).
 \tag{E816.42}
\]

Ordinary gluability makes this group zero. The equal-class flat-Hom
argument, stable-type decomposition, and quotient-to-subobject
contradiction of SF-E816a then apply word for word: no stable
isomorphism type can occur in both sheaves. Hence their cross Ext, in
both directions, is concentrated in degree two. This conclusion permits
overlapping normalized classes, since unequal stable types of the same
normalized class are completely Ext-orthogonal.

The canonical equivariant Postnikov triangle is

\[
 B[k]\longrightarrow E\longrightarrow A
      \xrightarrow{e}B[k+1],
 \qquad e\in\operatorname{Ext}^{k+1}_G(A,B).
 \tag{E816.43}
\]

For \(k=1\), SF-E816a already applies. For \(k\ge2\), the last group is
zero, so the triangle splits in the specified equivariant category. Thus

\[
 E\simeq A\oplus B[k].
 \tag{E816.44}
\]

The cross terms in its degree-two self-Ext are
\(\operatorname{Ext}^{2+k}_G(A,B)\) and
\(\operatorname{Ext}^{2-k}_G(B,A)\). Both vanish by concentration in
degree two. Therefore

\[
 \operatorname{ext}^2_G(E,E)
   =\operatorname{ext}^2_G(A,A)+\operatorname{ext}^2_G(B,B).
 \tag{E816.45}
\]

Count the actual nonzero normalized-class blocks of the two sheaves,
counting a repeated class in the two sheaves twice. The full character
is the sum of their pure exponential characters with the signs \(1\) and
\((-1)^k\). The balanced rank/common-annihilator proof in SF-E816 §4
requires at least four terms in this presentation. Invariant evaluation
contributes at least six middle self-Ext directions per actual block.
Equation (E816.45) consequently gives at least twenty-four, whereas weak
invariant semiregularity requires eighteen. This proves the gap
extension.

\subsection{SF-E816d. Three degrees do not force distinct stable types
under
gluability}\label{SF-E816cd-sf-e816d.-three-degrees-do-not-force-distinct-stable-types-under-gluability}

Ordinary gluability alone does not imply that a stable isomorphism type
occurs in at most one cohomological degree of an arbitrary bounded
complex. Here is a concrete actual example within the same primitive
projectively flat factor category.

Take the explicit CM fourfold and index-two line bundle of SF-E814:

\[
 X=E_{\mathrm i}^{\,4},\qquad
 L=\mathcal O(0)\boxtimes\mathcal O(0)
      \boxtimes\mathcal O(-0)\boxtimes\mathcal O(-0).
 \tag{E816.46}
\]

Use the real multiplication matrix \(U\) and divisor matrix \(S\) of
SF-E815. The class \(c_1(L)=\theta_S\) is block diagonal and primitive
on both real-multiplication planes. The elliptic Künneth computation
gives \(H^2(X,L)=\mathbf C\), with every other cohomology group zero.
Choose its nonzero class and the actual triangle

\[
 L[1]\longrightarrow F\longrightarrow\mathcal O_X
            \longrightarrow L[2]
 \tag{E816.47}
\]

of SF-E814. That full proof establishes ordinary gluability of \(F\), as
well as

\[
 \operatorname{Ext}^j_X(\mathcal O_X,F)=0\quad(j\le0).
 \tag{E816.48}
\]

For clarity, the degree-zero vanishing is the injectivity of
\(\operatorname{Hom}(\mathcal O_X,\mathcal O_X)\to
\operatorname{Ext}^2(\mathcal O_X,L)\), which sends the identity to the
chosen nonzero extension class. Negative degrees vanish in the same
triangle since the only cross cohomology of \(L\) is degree two.

Put

\[
 T=\mathcal O_X\oplus F[1].
 \tag{E816.49}
\]

Its nonzero cohomology sheaves are

\[
 \mathcal H^0(T)=\mathcal O_X,\qquad
 \mathcal H^{-1}(T)=\mathcal O_X,\qquad
 \mathcal H^{-2}(T)=L.
 \tag{E816.50}
\]

Thus the identical stable type \(\mathcal O_X\) occurs in two different
degrees. Nevertheless \(T\) is ordinary gluable. The diagonal negative
self-Ext groups vanish for both summands. For a negative integer \(j\),
one cross group is

\[
 \operatorname{Ext}^j_X(\mathcal O_X,F[1])
     =\operatorname{Ext}^{j+1}_X(\mathcal O_X,F)=0
\]

by (E816.48). The other is

\[
 \operatorname{Ext}^j_X(F[1],\mathcal O_X)
     =\operatorname{Ext}^{j-1}_X(F,\mathcal O_X)=0.
\]

Indeed \(F\in D^{\le0}\), while \(\mathcal O_X[j-1]\in D^{\ge1}\) for
these negative \(j\), so the latter vanishing follows from the
cohomological structure. This proves all required negative self-Ext
vanishings directly, with no formality assumption.

The example is not a balanced seed and is not presented as a new
construction for the main target. Its entire character is

\[
 \operatorname{ch}(T)=1-\bigl(1-e^{c_1(L)}\bigr)
                   =e^{c_1(L)},
 \tag{E816.51}
\]

which is pure and has degree-two cap rank six. It is also visibly
decomposable. Its sole role is to refute the blanket inference that
ordinary gluability forces stable-type disjointness across arbitrarily
many cohomological degrees. The three-or-more-degree overlapping case in
SF-E816 therefore needs an additional argument or an actual new
construction.

\section{SF-E817: an actual singular balanced character and the
universal point-defect
obstruction}\label{SF-E817-sf-e817-an-actual-singular-balanced-character-and-the-universal-point-defect-obstruction}

Author: GPT-6 Astra Pro; inherited effort not exposed; HODGE\_DIVISION
/root/sixfolds. Date: 2026-11-21.

\textbf{Status:} Whole-proof mathematical PASS received from /root and
/root/imprimitive on2026-11-21 (GPT-6 Astra Pro; inherited effort not
exposed). Root read the complete written proof and exact certificate and
freshly reran that certificate. Imprimitive read the full proof, freshly
reconstructed the exterior calculation with independent tuple-wedge
code, and freshly reread the supplied S-20/S-21 categorical passages.
The author ran the displayed exact certificate separately. The theorem
concerns torsion-free sheaves with locally free double dual and
zero-dimensional quotient. It does not treat positive-dimensional
defects or singular reflexive hulls.

\subsection{1. Statement}\label{SF-E817-statement}

Let \(X\) be a polarized complex abelian fourfold with a genuine real
quadratic action, of holomorphic type \((2,2)\). Let
\(\theta_+,\theta_-\) be the components of its compatible polarization
\(\theta\), and put

\[
 \Delta=\int_X\frac{\theta^4}{24}>0,
 \qquad z_\pm=\frac{\theta_\pm^2}{2}.
 \tag{E817.1}
\]

Thus \(\int_Xz_+z_-=\Delta\). Let \(E\) be a torsion-free coherent sheaf
of positive rank \(r\), with locally free double dual
\(H=E^{\vee\vee}\), and suppose

\[
 0\longrightarrow E\longrightarrow H\longrightarrow T\longrightarrow0
 \tag{E817.2}
\]

has nonzero zero-dimensional quotient \(T\). Suppose that for a rational
algebraic divisor \(B\) and nonzero real rational \(q\),

\[
 e^B\operatorname{ch}(E)
      =r v_q,
 \qquad
 v_q=(1-qz_+)(1-qz_-).
 \tag{E817.3}
\]

Let \(G\subset X\times\widehat X\) be any effective finite subgroup with
a specified coherent translation/tensor action, and suppose that \(E\)
has an actual equivariant structure. Let \(s\ge1\) be the number of
\(G\)-orbits in the support of \(T\).

Then its full invariant degree-two trace has

\[
 \dim\ker\sigma^2_{E,G}\ge1+3s\ge4.
 \tag{E817.4}
\]

In particular \(E\) is not weakly \(G\)-semiregular. There is no
simplicity, stability, or polystability hypothesis on \(E\) or on \(H\).
Nonreduced finite-length quotients are allowed. The proof uses the
actual canonical double-dual sequence, not an assertion that its
subobject or quotient is a direct summand.

\subsection{2. The equivariant structures on the canonical
sequence}\label{SF-E817-the-equivariant-structures-on-the-canonical-sequence}

Translation and tensoring by a line bundle commute naturally with double
dual. The given coherent isomorphisms on \(E\) therefore induce coherent
isomorphisms on \(H=E^{\vee\vee}\), and the canonical map \(E\to H\) is
equivariant. Its quotient \(T\) inherits the actual coherent structure.
Associativity is preserved by these functorial operations; no new
projective linearization of an individual factor is introduced.

Every finite-length sheaf has a canonical decomposition into its pieces
at the finitely many support points, by the decomposition of its
Artinian support algebra. Grouping these pieces by translation orbits
gives actual equivariant summands

\[
 T=\bigoplus_{\nu=1}^{s}T_\nu.
 \tag{E817.5}
\]

Tensoring by a degree-zero line bundle does not change support. Distinct
support-orbit summands have disjoint support, hence all their cross Ext
groups vanish. Each \(T_\nu\) has positive length \(\ell_\nu\) and
character \(\ell_\nu[\mathrm{pt}]\).

\subsection{3. A canonical exact sequence on middle
self-Ext}\label{SF-E817-a-canonical-exact-sequence-on-middle-self-ext}

We first record two elementary vanishings on a smooth fourfold:

\[
 \operatorname{Ext}^j_X(H,T)=0\quad(j>0),
 \qquad
 \operatorname{Ext}^j_X(T,H)=0\quad(j<4).
 \tag{E817.6}
\]

For the first, \(H\) is locally free and \(H^\vee\otimes T\) has finite
support, so it has no higher cohomology. For the second, the local
Koszul resolution of a point in a regular local ring of dimension four
gives \(\operatorname{Ext}^j(k_p,H)=0\) for \(j<4\). A finite-length
sheaf has a finite filtration with such residue-field factors. The Ext
long exact sequence proves the same vanishing for every factor
extension. The local-to-global passage has no higher-support cohomology,
giving (E817.6). Negative Ext vanishes for sheaves as usual.
Finite-group invariants over \(\mathbf C\) are exact, so these
vanishings also hold for the specified equivariant Ext groups.

Use the actual equivariant triangle

\[
 T[-1]\xrightarrow{i}E\xrightarrow{p}H
       \xrightarrow{a}T.
 \tag{E817.7}
\]

Apply \(\operatorname{RHom}_G(H,-)\). The exact segment for \(j=2,3\) is

\[
 \operatorname{Ext}^{j-1}_G(H,T)
 \longrightarrow\operatorname{Ext}^{j}_G(H,E)
 \xrightarrow{p_*}\operatorname{Ext}^{j}_G(H,H)
 \longrightarrow\operatorname{Ext}^{j}_G(H,T).
\]

The exterior groups vanish, so

\[
 \operatorname{Ext}^{j}_G(H,E)
       \simeq\operatorname{Ext}^{j}_G(H,H)
       \qquad(j=2,3).
 \tag{E817.8}
\]

Next apply \(\operatorname{RHom}_G(T[-1],-)\). Its adjacent groups with
target \(H\) are \(\operatorname{Ext}^{j}_G(T,H)\) and
\(\operatorname{Ext}^{j+1}_G(T,H)\). For \(j=1,2\), both vanish by
(E817.6). Therefore

\[
 i_*:\operatorname{Ext}^{j}_G(T,T)
      \xrightarrow{\sim}\operatorname{Ext}^{j}_G(T[-1],E)
        \qquad(j=1,2).
 \tag{E817.9}
\]

Finally apply \(\operatorname{RHom}_G(-,E)\). The relevant complete
segment is

\[
 \begin{split}
 \operatorname{Ext}^1_G(T[-1],E)
 &\xrightarrow{\partial_1}\operatorname{Ext}^2_G(H,E)
 \longrightarrow\operatorname{Ext}^2_G(E,E)\\
 &\longrightarrow\operatorname{Ext}^2_G(T[-1],E)
 \xrightarrow{\partial_2}\operatorname{Ext}^3_G(H,E).
 \end{split}
 \tag{E817.10}
\]

The two connecting maps are zero. Indeed, using (E817.9), a degree-\(j\)
element, for \(j=1,2\), is represented by

\[
 f=i[j]\circ b[-1],\qquad b:T\longrightarrow T[j].
\]

Its connecting class is, up to sign,

\[
 \partial_j(f)=f[1]\circ a
             =i[j+1]\circ b\circ a.
 \tag{E817.11}
\]

The intermediate composite \(b\circ a\) belongs to
\(\operatorname{Ext}^j_G(H,T)=0\), by (E817.6). These are actual
equivariant morphisms, so the vanishing holds before any dimension
estimate.

Equations (E817.8)--(E817.11) yield the short exact sequence

\[
 0\longrightarrow\operatorname{Ext}^2_G(H,H)
   \longrightarrow\operatorname{Ext}^2_G(E,E)
   \longrightarrow\operatorname{Ext}^2_G(T,T)
   \longrightarrow0.
 \tag{E817.12}
\]

The sequence is not claimed to split canonically. It gives the exact
dimension identity

\[
 e_2^G(E)=e_2^G(H)+e_2^G(T).
 \tag{E817.13}
\]

No degeneration of a spectral sequence, formality assertion, or
stability condition is used.

\subsection{4. The complete cap-rank calculation after adding a point
class}\label{SF-E817-the-complete-cap-rank-calculation-after-adding-a-point-class}

Let \(\ell=\operatorname{length}(T)>0\). Since an exponential fixes the
top point class, (E817.2)--(E817.3) give

\[
 e^B\operatorname{ch}(H)
   =r(1-qz_+)(1-qz_-)+\frac{\ell}{\Delta}z_+z_-.
 \tag{E817.14}
\]

We compute its full degree-two Hochschild cap rank. Let

\[
 W=W_+\oplus W_-,\qquad
 W_\pm=L_\pm\oplus M_\pm,
 \qquad\dim L_\pm=\dim M_\pm=2,
\]

where \(L_\pm\) and \(M_\pm\) are the translation and tensor planes for
the respective real embeddings. Choose nonzero volume bivectors
\(\alpha_\pm\in\bigwedge^2L_\pm\) and \(\beta_\pm\in\bigwedge^2M_\pm\)
so that the Fock representatives of the unit and of \(z_\pm\) are
\(\alpha_\pm\) and \(\beta_\pm\). The full character in (E817.14)
becomes

\[
 w=a\alpha_+\alpha_-+b\beta_+\alpha_-
       +c\alpha_+\beta_-+d\beta_+\beta_-,
 \tag{E817.15}
\]

with

\[
 a=r,\qquad b=c=-rq,\qquad
 d=rq^2+\frac{\ell}{\Delta},\qquad
 ad-bc=\frac{r\ell}{\Delta}\ne0.
 \tag{E817.16}
\]

All four coefficients are nonzero. The multiplication map
\(\bigwedge^2W\to\bigwedge^6W\), \(u\mapsto u\wedge w\), has the
following three independent bidegree blocks.

On \(\bigwedge^2W_+\), only the two directions \(\alpha_+,\beta_+\)
contribute. Their images, in \(\bigwedge^4W_+\otimes\bigwedge^2W_-\),
have coefficient matrix

\[
 \begin{pmatrix}b&a\\d&c\end{pmatrix}.
\]

Its determinant is \(bc-ad\ne0\), so this block has rank two. The
analogous block on \(\bigwedge^2W_-\) has rank two as well.

On \(W_+\otimes W_-\), decompose each factor into its two
translation/tensor planes. The four resulting four-dimensional summands
map isomorphically into distinct summands of
\(\bigwedge^3W_+\otimes\bigwedge^3W_-\), multiplied respectively by
\(a,b,c,d\), up to signs. Hence the mixed block has rank sixteen. Thus

\[
 \operatorname{rank}c^2_{\operatorname{ch}(H)}=2+2+16=20.
 \tag{E817.17}
\]

The common divisor shear does not change this rank. When \(\ell=0\), the
two pure blocks instead have rank one each, recovering the balanced rank
eighteen. For a nonzero point character alone, only \(d\) is nonzero;
the same block calculation gives ranks one, one, and four, totaling six.

The calculation is universal in the coefficients. A nonzero
twenty-by-twenty minor can be chosen with determinant, up to sign,

\[
 (ad-bc)^2(abcd)^4.
 \tag{E817.18}
\]

An exact exterior script checks this generic minor and the stated ranks.
The proof of the universal rank is the three-block argument above.

\subsection{5. The invariant kernel lower
bound}\label{SF-E817-the-invariant-kernel-lower-bound}

For every actual equivariant perfect object, evaluation of Hochschild
degree two takes values in invariant Ext, because translation and
degree-zero tensor operations act trivially on Hochschild cohomology.
Evaluation followed by full trace is cap with its Chern character, up to
the fixed nonzero identity-sector normalization of supplied DEG\_L S-21.
Equation (E817.17) therefore gives

\[
 e_2^G(H)\ge20.
 \tag{E817.19}
\]

For each actual equivariant support-orbit summand \(T_\nu\), its
positive point character has cap rank six. Thus \(e_2^G(T_\nu)\ge6\).
The disjoint-support decomposition (E817.5) gives

\[
 e_2^G(T)\ge6s.
 \tag{E817.20}
\]

Combining these inequalities with (E817.13) yields

\[
 e:=e_2^G(E)\ge20+6s.
 \tag{E817.21}
\]

To obtain an actual trace-kernel bound without assuming semiregularity,
let \(h=\operatorname{rank}\sigma^2_{E,G}\) and \(k=e-h\). CY4
trace/evaluation adjunction says that the evaluation map
\(\rho^2_{E,G}\) also has rank \(h\). The Serre pairing on invariant Ext
is perfect, since the finite translation/tensor action preserves it and
averaging is exact. Cap with the balanced character has rank eighteen.
The elementary rank inequality for the composition gives

\[
 18=\operatorname{rank}(\sigma^2_{E,G}\rho^2_{E,G})
       \ge 2h-e=e-2k.
 \tag{E817.22}
\]

Equations (E817.21)--(E817.22) imply

\[
 k\ge\frac{e-18}{2}\ge1+3s,
\]

which is (E817.4). The bound remains positive for every effective finite
coherent action, including actions that correlate different support
points and carry nontrivial multiplicity representations.

\subsection{6. An explicit actual singular balanced
sheaf}\label{SF-E817-an-explicit-actual-singular-balanced-sheaf}

Use the CM product and integral divisor classes of SF-E815:

\[
 X=E_i^4,\qquad X_0=x+y,\qquad Y=\sqrt2(x-y),
 \qquad x^2=-\theta_+^2,\quad y^2=-\theta_-^2,
 \quad\int_Xx^2y^2=4.
\]

The rationality and integrality of \(X_0,Y\) come from the explicit
Hermitian matrices \(S\) and \(US\) in that theorem. Form the actual
polystable bundle

\[
 H=\bigl(L_{2X_0}\oplus L_{-2X_0}\bigr)^{\oplus3}
       \oplus\bigl(L_Y\oplus L_{-Y}\bigr)^{\oplus6}.
 \tag{E817.23}
\]

Its rank is eighteen and its odd character components vanish. Direct
expansion gives

\[
 \begin{aligned}
 \operatorname{ch}_2(H)
   &=12X_0^2+6Y^2=24(x^2+y^2),\\
 \operatorname{ch}_4(H)
   &=4X_0^4+\tfrac12Y^4=36x^2y^2.
 \end{aligned}
 \tag{E817.24}
\]

Choose sixteen distinct points \(Z\) and a nonzero linear functional on
the fibre of \(H\) at each point. These define an actual surjection
\(H\twoheadrightarrow\mathcal O_Z\). Let \(E\) be its kernel. It is
torsion-free of rank eighteen, and its double dual is \(H\). For
example, dualizing the sequence twice proves this directly because
\(\operatorname{Hom}(T,\mathcal O_X)=\operatorname{Ext}^1(T,\mathcal O_X)=0\).

Locally at a point of \(Z\), an invertible change of basis makes the
surjection \(R^{18}\to k\) the reduction of the first coordinate, so

\[
 E_p\simeq\mathfrak m_p\oplus R^{17}.
\]

The maximal ideal of a regular local ring of dimension four is not free
of rank one. Thus this is a genuinely singular torsion-free sheaf, not a
vector bundle.

Since \([\mathrm{pt}]=x^2y^2/4\), the entire character of its kernel is

\[
 \begin{split}
 \operatorname{ch}(E)
   &=18+24(x^2+y^2)+32x^2y^2\\
   &=18\left(1-\frac43\theta_+^2\right)
          \left(1-\frac43\theta_-^2\right)
     =18v_{8/3}.
 \end{split}
 \tag{E817.25}
\]

All the line summands of \(H\) have zero slope for the product
polarization, because their restrictions to each RM plane are primitive.
The bundle \(H\) is therefore polystable of slope zero. The kernel has
the same first Chern class, and every subsheaf of \(E\) is a subsheaf of
\(H\), so \(E\) is semistable of slope zero. Its singular stable factors
are not covered by the locally free Jordan--Hölder hypothesis in
SF-E813. The surjection can be chosen with all fibre components nonzero;
no direct-sum decomposition of the kernel is assumed.

Nevertheless SF-E817 excludes weak semiregularity of this actual sheaf
under every effective finite coherent action. In the ordinary category,
its sixteen support points give \(s=16\), hence the invariant-free full
trace kernel has dimension at least forty-nine by (E817.4). This lower
bound is not asserted to be sharp.

\subsection{Scope and source record}\label{SF-E817-scope-and-source-record}

This construction genuinely enters the singular-factor route and then
fails by the displayed exact maps. The universal theorem removes all
nonzero finite-length defects of locally free hulls in the balanced
character, without a stability or simplicity premise. It does not remove
defects supported on curves or surfaces, nor a reflexive hull that is
itself singular. Those cases have different local Ext groups and must
not be inferred from (E817.6).

The local vanishings are proved here by the four-variable Koszul complex
and finite-length filtrations. The only categorical input not reproved
is the already reread CY4 trace/evaluation adjunction and coherent
finite-action identity-sector assertion of supplied DEG\_L S-21. The
Fock identification and balanced rank eighteen are also rederived in the
displayed blocks and agree with the previously checked SF-E801. No new
preprint theorem, Hermitian--Einstein extension for singular sheaves, or
dg formality is used.

\section{SF-E818: positive-dimensional defects of a locally free
semistable
hull}\label{SF-E818-sf-e818-positive-dimensional-defects-of-a-locally-free-semistable-hull}

Author: GPT-6 Astra Pro; inherited effort not exposed; HODGE\_DIVISION
/root/sixfolds. Date: 2026-11-21.

\textbf{Status:} Whole-proof mathematical PASS received from /root and
/root/imprimitive on 2026-11-21 (GPT-6 Astra Pro; inherited effort not
exposed). Root reread the precise SF-E811 and SF-E813 hypotheses.
Imprimitive independently reconstructed every proof step and freshly
read the invoked Fulton/Stacks support-cycle statements. The argument
uses only the already reviewed locally free Hermitian--Einstein result
SF-E811, the factor analysis of SF-E813, and the point-defect theorem
SF-E817. It does not use an admissible Hermitian--Einstein metric on a
singular sheaf.

\subsection{Statement}\label{SF-E818-statement}

Let \(X\) be a polarized complex abelian fourfold with genuine
Rosati-invariant real multiplication by a real quadratic field, of
holomorphic type \((2,2)\). Write its compatible polarization as
\(\theta=\theta_++\theta_-\). Suppose that \(E\) is a torsion-free
coherent sheaf of rank \(r>0\), and that its double dual
\(H=E^{\vee\vee}\) is locally free and admits a filtration by subbundles
with stable locally free quotients of the same \(\theta\)-slope. Suppose
also that

\[
 e^B\operatorname{ch}(E)=r(1-q\theta_+^2/2)(1-q\theta_-^2/2),
 \qquad B\in\operatorname{NS}(X)_{\mathbf Q},\qquad q\in\mathbf Q_{>0}.
 \tag{E818.1}
\]

Then \(E\) is not weakly \(G\)-semiregular for any specified coherent
action of an effective finite translation/tensor subgroup
\(G\subset X\times\widehat X\) for which \(E\) is actually equivariant.

In fact the character equation already forces the quotient \(T=H/E\) to
have dimension zero or to be zero. If it is nonzero, SF-E817 gives the
quantitative kernel bound \(1+3s\), where \(s\) is its number of support
orbits. If it is zero, SF-E813 excludes the locally free sheaf itself.

The locally free stable-factor assumption is imposed on the actual hull
\(H\), not on \(E\). Thus the theorem allows singular stable factors of
\(E\), including those arising by elementary modifications. It makes no
assertion when the reflexive hull is singular, or when a locally free
hull has singular stable Jordan--Hölder factors.

\subsection{1. Strict mixed positivity for every effective
surface}\label{SF-E818-strict-mixed-positivity-for-every-effective-surface}

We first prove a small extension of the effectivity argument in supplied
DEG\_L S-22. For every nonzero effective surface cycle \(Z\) on \(X\),

\[
 \int_X[Z]\theta_+\theta_->0.
 \tag{E818.2}
\]

The forms \(\theta_\pm\) have invariant semipositive representatives
supported on the two complex two-dimensional real-multiplication
eigenspaces \(V_\pm\). Hence their mixed wedge restricts nonnegatively
to every complex surface tangent plane. More explicitly, in unitary
bases it is a positive sum of squared norms of the restrictions of the
mixed holomorphic two-forms in \(V_+^*\otimes V_-^*\). Consequently a
zero integral on an effective cycle forces all these mixed forms to pull
back to zero on a projective resolution of each irreducible component.

At a smooth point of a component, write the Plücker bivector of its
tangent plane as \(u_++u_-\), where \(u_\pm\in\bigwedge^2V_\pm\); the
mixed part vanishes by the preceding conclusion. A decomposable bivector
has square zero, whereas this square is \(2u_+\wedge u_-\). Thus either
\(u_+=0\) or \(u_-=0\), and the tangent plane is exactly \(V_-\) or
\(V_+\). The Gauss map of each irreducible component has finite image
and is therefore constant.

For completeness, constant tangent image in \(V_+\), for example, is
impossible here. Let \(Y\) be a smooth projective resolution of the
component and translate its map to \(X\) to send a base point to zero.
All invariant one-forms annihilating \(V_+\) pull back to zero.
Therefore the homomorphism \(\operatorname{Alb}(Y)\to X\) has tangent
image contained in \(V_+\). Its image is a positive-dimensional abelian
subvariety, since the surface map is nonconstant. But the underlying
real eigenspace of \(V_+\) contains no nonzero rational homology vector:
if \(u\in F\setminus\mathbf Q\) and \(z\ne0\) were such a vector, the
equality \(uz=\sigma_+(u)z\), in rational coordinates, would make the
irrational number \(\sigma_+(u)\) rational. A positive-dimensional
abelian subvariety has nonzero rational homology. This contradiction
proves (E818.2). The proof does not assume that \(X\) is simple.

\subsection{2. The surface part of the defect must
vanish}\label{SF-E818-the-surface-part-of-the-defect-must-vanish}

We use the leading-character/support-cycle identity in the following
exact published form. Fulton, \emph{Intersection Theory}, second edition
(1998), Example18.3.11, printed p.363, identifies the
highest-dimensional component of the Riemann--Roch transformation of a
coherent sheaf with its support cycle weighted by generic lengths.
Theorem18.3(3), printed p.353, identifies that transformation on a
smooth variety with
\(\operatorname{ch}(-)\operatorname{td}(T_X)\cap[X]\). Here \(T_X\) is
trivial. Thus a sheaf \(T\) of support dimension at most \(d\) has
\(\operatorname{ch}_j(T)=0\) for \(j<4-d\), and
\(\operatorname{ch}_{4-d}(T)\) is exactly its effective dimension-\(d\)
support cycle. The generic-length convention is also Stacks Project,
Definition42.10.2, tag \texttt{02QV}. These are published/standard
inputs, not claims to reprove the entire Riemann--Roch theorem.

The canonical double-dual map gives

\[
 0\longrightarrow E\longrightarrow H\longrightarrow T\longrightarrow0,
 \qquad \dim\operatorname{Supp}T\le2.
 \tag{E818.3}
\]

Put \(\alpha=c_1(H)/r=c_1(E)/r\). The first Chern component of (E818.1)
gives \(B=-\alpha\). For any positive-rank sheaf define

\[
 \Delta(F)=\frac{c_1(F)^2}{2\operatorname{rk}F}-\operatorname{ch}_2(F).
\]

Since the quotient in (E818.3) has codimension at least two,

\[
 \Delta(E)=\Delta(H)+\operatorname{ch}_2(T),
 \qquad
 \int_X\Delta(E)\theta_+\theta_-=0.
 \tag{E818.4}
\]

Let \(Q_j\) be the stable locally free quotients in the assumed
filtration of \(H\), and set

\[
 r_j=\operatorname{rk}Q_j,\qquad
 \alpha_j=c_1(Q_j)/r_j,\qquad \delta_j=\alpha_j-\alpha.
\]

The exact Chern-character variance identity is

\[
 \Delta(H)=\sum_j\Delta(Q_j)-\frac12\sum_jr_j\delta_j^2.
 \tag{E818.5}
\]

Each factor has the same slope as \(H\), so
\(\int_X\delta_j\theta^3=0\). The invariant representative of
\(\delta_j\) is therefore primitive for \(\theta\). SF-E811 gives

\[
 \int_X\Delta(Q_j)\theta_+\theta_-\ge0,
 \qquad -\int_X\delta_j^2\theta_+\theta_-\ge0.
 \tag{E818.6}
\]

The second inequality is the elementary mixed sum-of-squares identity
(E811.4), and the first uses a smooth Hermitian--Einstein connection
only on the locally free stable factor \(Q_j\). Consequently the mixed
pairing of \(\Delta(H)\) is nonnegative.

The class \(\operatorname{ch}_2(T)\) is the effective fundamental cycle
of the dimension-two part of \(T\): its coefficients are the positive
generic lengths at the irreducible support surfaces. Terms of smaller
support dimension do not contribute. Equations (E818.2), (E818.4), and
(E818.6) therefore force this surface cycle to be zero. Thus

\[
 \dim\operatorname{Supp}T\le1,
 \qquad
 \int_X\Delta(H)\theta_+\theta_-=0.
 \tag{E818.7}
\]

Every nonnegative summand of (E818.5) has zero mixed pairing. By the
equality case of SF-E811, each \(Q_j\) is projectively flat, with
\(\operatorname{ch}(Q_j)=r_je^{\alpha_j}\). By the equality case of
(E811.4), every \(\delta_j\) is block diagonal for \(V_+\oplus V_-\),
and primitive on each complex two-plane separately. This is precisely
the factor conclusion proved in SF-E813; no splitting of \(H\), and no
equivariance of individual Jordan--Hölder factors, is needed in this
step.

\subsection{3. The curve part of the defect must
vanish}\label{SF-E818-the-curve-part-of-the-defect-must-vanish}

Write \(\delta_j=\delta_{j,+}+\delta_{j,-}\). The preceding primitivity
says

\[
 \delta_{j,+}\theta_+=0,\qquad
 \delta_{j,-}\theta_-=0,
 \qquad \delta_{j,+}^3=\delta_{j,-}^3=0.
 \tag{E818.8}
\]

These are identities of invariant forms on the indicated complex
two-dimensional blocks. Expanding the cube gives

\[
 \delta_j^3=3\delta_{j,+}^2\delta_{j,-}
                +3\delta_{j,+}\delta_{j,-}^2.
\]

Multiplication by \(\theta_++\theta_-\) kills every term, either by
excess degree in one block or by one of the first two identities in
(E818.8). Hence

\[
 \theta\delta_j^3=0,
 \qquad
 \int_X\theta\bigl(e^{-\alpha}\operatorname{ch}(H)\bigr)_3
   =\frac16\sum_jr_j\int_X\theta\delta_j^3=0.
 \tag{E818.9}
\]

The normalized character of \(E\) in (E818.1) has zero degree-three
component. By (E818.7), \(T\) has support dimension at most one, so its
degree-three component is unchanged by multiplication by an exponential.
Additivity in (E818.3) therefore gives

\[
 \operatorname{ch}_3(T)
   =\bigl(e^{-\alpha}\operatorname{ch}(H)\bigr)_3,
 \qquad
 \int_X\theta\operatorname{ch}_3(T)=0.
 \tag{E818.10}
\]

For a sheaf with support dimension at most one,
\(\operatorname{ch}_3(T)\) is its effective fundamental curve cycle with
positive generic-length multiplicities. Every nonzero such curve cycle
has strictly positive degree for the ample polarization \(\theta\). Thus
(E818.10) proves that \(T\) has no one-dimensional component. It is
zero-dimensional or zero, as claimed.

\subsection{4. The remaining cases and exact
scope}\label{SF-E818-the-remaining-cases-and-exact-scope}

If \(T\ne0\), the canonical sequence (E818.3) is precisely the sequence
in SF-E817. The specified coherent finite action on \(E\) canonically
induces coherent actions on \(H\) and \(T\); this functorial statement
is proved there. That theorem gives a nonzero trace kernel, indeed
dimension at least \(1+3s\).

If \(T=0\), then \(E=H\) is a locally free sheaf with the stable locally
free filtration assumed in SF-E813 and the same balanced character.
SF-E813 excludes weak equivariant semiregularity. This proves SF-E818 in
all cases.

The new part relative to SF-E817 is the elimination of all
positive-dimensional defects under an explicit condition on the hull.
There is no assertion that every locally free semistable hull
automatically has locally free stable factors. There is also no use of a
Chern--Weil identity for a singular reflexive sheaf, no extension of a
connection across its singularities, and no reduction of a general
derived object to a sheaf.

The smooth Hermitian--Einstein input has exactly the
rational-polarization scope recorded in SF-E811: Siu, \emph{Lectures on
Hermitian--Einstein metrics for stable bundles and Kähler--Einstein
metrics}, DMV Seminar8 (1987), Chapter1, §1.3--§1.4, pp.6--7;
Proposition1.6, pp.8--9; §6.1--§6.4, pp.35--39. Its proof is not
repeated here. The mixed positivity and its equality case are the
independently checked SF-E811 calculation. The strict surface argument
above fully reproves the needed strengthening of supplied DEG\_L S-22.
The categorical trace input is only that already recorded in SF-E813 and
SF-E817.

\section{SF-E819: a local commutator obstruction to signed point
modifications}\label{SF-E819-sf-e819-a-local-commutator-obstruction-to-signed-point-modifications}

Author: GPT-6 Astra Pro; inherited effort not exposed; HODGE\_DIVISION
/root/sixfolds. Date: 2026-11-21.

\textbf{Status:} Complete proof independently reviewed by /root and
/root/g34\_eightfolds on 21 November 2026; final g34 whole-proof review
and independent 256-coordinate exterior certificate PASS. Reviewers:
GPT-6 Astra Pro, inherited effort not exposed. The reviewed
mathematical/source draft has SHA-256
\texttt{10d171d811e27e5eeae0381b00d137b2439fa90c6bfde38046e6cf94baaf2baf}.
The local lemma and the reduced-defect corollary are separate
conclusions. The first-order module argument, lifting of the hull map,
and finite perfect-complex trace are written explicitly. No all-order
quantization theorem, singular Hermitian--Einstein theorem, or formality
assertion is used.

\subsection{1. Statement and new scope}\label{SF-E819-statement-and-new-scope}

Let \(X\) be a complex abelian fourfold with genuine real multiplication
of type \((2,2)\), compatible polarization \(\theta=\theta_++\theta_-\),
and

\[
 v_q=(1-q\theta_+^2/2)(1-q\theta_-^2/2),\qquad q>0.
\]

Let \(E\) be a coherent sheaf, with locally free double dual
\(H=E^{\vee\vee}\), such that the kernel \(T\) and cokernel \(Q\) of its
canonical map to \(H\) have finite support:

\[
 0\longrightarrow T\longrightarrow E\longrightarrow H
       \longrightarrow Q\longrightarrow0.
 \tag{E819.1}
\]

Suppose that \(e^B\operatorname{ch}(E)=r v_q\), where \(r>0\) and \(B\)
is a rational algebraic divisor. If \(E\) is weakly \(G\)-semiregular
for any specified effective coherent finite translation/tensor action,
then at every point \(p\)

\[
 \operatorname{length}Q_p=\operatorname{length}T_p.
 \tag{E819.2}
\]

In particular, a nonzero signed point defect at even one point prevents
weak semiregularity. The theorem allows actual zero-dimensional torsion
in \(E\), which SF-E817 excludes from its hypotheses. It does not settle
modifications with equal lengths at every point, nor singular reflexive
hulls or positive-dimensional defects.

The ordinary Hochschild trace/evaluation identity used here is the
published Buchweitz--Flenner theorem: \emph{The global decomposition
theorem for Hochschild (co-)homology of singular spaces via the
Atiyah--Chern character}, Advances in Mathematics217(2008),243--281,
DOI10.1016/j.aim.2007.06.013, Theorem6.4.2. Exact author-copy locators
are Proposition6.2.1, pp29--30; Proposition6.3.1 and Corollary6.3.2,
pp30--31; Theorem6.4.2 and its proof through Proposition6.4.4, pp31--33,
in arXiv:math/0606730. These identify semiregularity after Hochschild
evaluation with contraction by the full Chern character. Only
compatibility with the specified coherent finite action is taken from
supplied DEG\_L S-21: equation(S21.6), the paragraph following(S21.7),
and the invariant-evaluation paragraph preceding(S21.8). Those are
retained as supplied proved inputs. The obstruction tested is ordinary
degree-two Ext; injectivity on the specified invariant Ext suffices
because Hochschild evaluation takes values in those invariants. No
Serre-adjunction identity is needed here.

\subsection{2. A local first-order
lemma}\label{SF-E819-a-local-first-order-lemma}

Let \(R=\mathbf C[[z_1,z_2,z_3,z_4]]\), let \(M\) be a finite
\(R\)-module, and suppose there is an \(R\)-linear map

\[
 \phi:M\longrightarrow R^r
\]

with finite-length kernel \(T\) and cokernel \(Q\). Let \(\pi\) be a
nonzero constant bivector in these coordinates. Put
\(D=\mathbf C[\varepsilon]/(\varepsilon^2)\), and endow
\(A=R\oplus\varepsilon R\) with the associative first-order product

\[
 f*g=fg+\frac{\varepsilon}{2}\pi(df,dg).
 \tag{E819.3}
\]

Associativity modulo \(\varepsilon^2\) follows from the Leibniz rule for
the biderivation; no higher-order Poisson quantization is required. The
elements \(z_i\) satisfy

\[
 [z_i,z_j]_* =\varepsilon\pi_{ij}.
 \tag{E819.4}
\]

\textbf{Local lemma.} If \(M\) has a \(D\)-flat first-order deformation
to a finite left \(A\)-module, then
\(\operatorname{length}Q=\operatorname{length}T\).

\subsubsection{2.1. The map to the free hull
lifts}\label{SF-E819-the-map-to-the-free-hull-lifts}

Write \(K=\operatorname{im}\phi\). The two exact sequences are

\[
 0\to T\to M\to K\to0,
 \qquad 0\to K\to R^r\to Q\to0.
 \tag{E819.5}
\]

The four-variable Koszul complex and a finite-length filtration give
\(\operatorname{Ext}^j_R(T,R)=\operatorname{Ext}^j_R(Q,R)=0\) for
\(j<4\). Applying \(\operatorname{Hom}_R(-,R^r)\) to (E819.5)
consequently gives

\[
 \operatorname{Ext}^1_R(M,R^r)=0.
 \tag{E819.6}
\]

Let \(M_\varepsilon\) be a flat deformation as in the lemma, and give
\(H_\varepsilon=A^r\) its free left-module structure. The map \(\phi\)
lifts to an \(A\)-linear map

\[
 \phi_\varepsilon:M_\varepsilon\longrightarrow H_\varepsilon.
 \tag{E819.7}
\]

Here is a direct obstruction calculation. Choose a finite free
resolution of \(M\) over the regular local ring \(R\). One can lift it
to a free resolution of \(M_\varepsilon\) with the same finite ranks:
lift generators successively; the kernels remain \(D\)-flat, because
both the free source and its flat quotient are \(D\)-flat. Reduction
modulo \(\varepsilon\) gives the chosen successive kernels. At the last
step a kernel with zero reduction is zero by the nilpotent version of
Nakayama's lemma. Apply \(\operatorname{Hom}_A(-,H_\varepsilon)\). The
reduction map of these Hom complexes is surjective termwise, with kernel
its \(\varepsilon\)-multiple, naturally the corresponding Hom complex
over \(R\). The obstruction to lifting the degree-zero cocycle \(\phi\)
lies in \(\operatorname{Ext}^1_R(M,R^r)\), which is zero by (E819.6).
Correcting a lifted cocycle therefore gives (E819.7).

\subsubsection{2.2. The defect cone is perfect over the dual
numbers}\label{SF-E819-the-defect-cone-is-perfect-over-the-dual-numbers}

Consider the two-term complex, in degrees \(-1,0\),

\[
 C_\varepsilon=[M_\varepsilon\xrightarrow{\phi_\varepsilon}H_\varepsilon].
 \tag{E819.8}
\]

It is a complex of left \(A\)-modules and a complex of flat
\(D\)-modules. Its reduction is \(C_0=[M\to R^r]\), with cohomology
\(T\) in degree \(-1\) and \(Q\) in degree zero. The short exact
sequence of complexes

\[
 0\longrightarrow C_0\xrightarrow{\varepsilon} C_\varepsilon
      \longrightarrow C_0\longrightarrow0
\]

shows that every cohomology group of \(C_\varepsilon\) is
finite-dimensional over \(\mathbf C\). As a complex over the Noetherian
Artinian ring \(D\), it therefore has bounded finite cohomology and is
pseudo-coherent. Its two flat terms give finite Tor amplitude. Thus it
is a perfect \(D\)-complex. The exact general criterion is Stacks
Project, Lemma15.76.2, tag \texttt{0658}; its proof truncates a finite
free resolution using the finite Tor amplitude. In this Noetherian case
the required bounded-above finite free resolution is obtained by
choosing generators of cohomology and successive finite kernels.

Choose a bounded finite free \(D\)-complex \(P^\bullet\) representing
\(C_\varepsilon\). Its Euler rank satisfies

\[
 \sum_j(-1)^j\operatorname{rk}_D P^j
  =\chi(C_0)=\operatorname{length}Q-\operatorname{length}T.
 \tag{E819.9}
\]

\subsubsection{2.3. The commutator trace}\label{SF-E819-the-commutator-trace}

Left multiplication by each \(z_i\) is a \(D\)-linear chain endomorphism
of (E819.8), since \(\phi_\varepsilon\) is \(A\)-linear. Their
commutator is \(\varepsilon\pi_{ij}\) times the identity. Transfer these
endomorphisms to \(P^\bullet\) in the derived category. They can be
represented by chain maps, and the commutator equality holds up to chain
homotopy.

For a chain endomorphism \(f\) of a bounded finite free complex over a
commutative ring, define its supertrace to be
\(\sum_j(-1)^j\operatorname{Tr}(f^j)\). A degree-zero commutator has
zero supertrace by cyclicity of ordinary matrix trace. A null-homotopic
map has zero supertrace because the contributions of \(dh\) and \(hd\)
cancel after shifting the index. Therefore (E819.4) gives

\[
 0=\varepsilon\pi_{ij}
       \bigl(\operatorname{length}Q-\operatorname{length}T\bigr)
       \quad\text{in }D.
 \tag{E819.10}
\]

At least one \(\pi_{ij}\) is nonzero. Since \(\varepsilon\ne0\) in
\(D\), the integer in parentheses must be zero. This proves the local
lemma. It is essential that the trace is taken on the finite perfect
defect complex; no trace of an infinite-dimensional free \(R\)-module is
used.

\subsection{3. Vanishing of Hochschild evaluation gives the local
deformation}\label{SF-E819-vanishing-of-hochschild-evaluation-gives-the-local-deformation}

We recall the elementary first-order identification to fix precisely
what is used. For a \(\mathbf C\)-algebra \(R\), a square-zero
multiplication deformation \(\mu+\varepsilon\mu_1\), and an \(R\)-module
\(M\) with action \(\rho\), an action on \(M\oplus\varepsilon M\) has
the form \(\rho+\varepsilon\rho_1\). The associativity equation is

\[
 \rho(a)\rho_1(b)-\rho_1(ab)+\rho_1(a)\rho(b)
       =\rho\bigl(\mu_1(a,b)\bigr).
 \tag{E819.11}
\]

The left side is the bar differential. The right side is the degree-two
cocycle obtained by applying the Hochschild characteristic morphism to
\(\mu_1\). The bar resolution
\(R\otimes_{\mathbf C}R^{\otimes n}\otimes_{\mathbf C}M\) identifies
this cohomology group with \(\operatorname{Ext}^2_R(M,M)\). Consequently
the obstruction vanishes exactly when (E819.11) has a solution, giving a
flat first-order module deformation. A normalized solution has
\(\rho_1(1)=0\); this can equally be imposed by the normalized bar
resolution. Nilpotent Nakayama shows that the resulting module is finite
over the deformed algebra when \(M\) is finite.

For the completed local ring in §2 one can use a finite free resolution
in place of the bar resolution. Lift its matrices to the product
(E819.3). Their squared differential is \(\varepsilon\) times an
endomorphism cocycle of degree two. Changing the lifted differential by
\(\varepsilon h\) changes this cocycle by \(dh+hd\). Its class is
exactly the characteristic class just described, as comparison with the
bar resolution shows. Thus a zero class gives a finite free deformed
complex. It is flat over \(D\), and reduction has cohomology only in
degree zero. The elementary square-zero cohomology sequence then gives a
flat degree-zero module and no other cohomology. This formulation keeps
all maps continuous for formal power series; it does not assume that an
arbitrary discontinuous action on power series is acceptable.

Restriction of the global characteristic morphism to a small affine
neighborhood and then to its completed local ring is compatible with
this description. In the HKR decomposition

\[
 HH^2(X)=H^2(\mathcal O_X)\oplus H^1(T_X)
                \oplus H^0(\bigwedge^2T_X),
 \tag{E819.12}
\]

the first two components restrict to zero on the affine/formal local
scheme. A global class with bivector component \(\pi\) therefore
restricts to the multiplication obstruction of (E819.3), up to a fixed
nonzero normalization of the HKR map. That normalization does not affect
either vanishing or the trace contradiction.

\subsection{4. The balanced cap kernel supplies a nonzero
bivector}\label{SF-E819-the-balanced-cap-kernel-supplies-a-nonzero-bivector}

For each real-multiplication plane, choose a nonzero invariant
holomorphic bivector \(\pi_\pm\). Nondegeneracy of \(\theta_\pm\)
defines a nonzero element \(\eta_\pm\in H^2(\mathcal O_X)\), supported
on that plane, by

\[
 \pi_\pm\mathbin{\lrcorner}(\theta_\pm^2/2)=\eta_\pm.
\]

Wedge by \(\eta_\pm\) kills the degree-four term of its own
two-dimensional plane. Hence on the factor \(1-q\theta_\pm^2/2\),
contraction by \(\pi_\pm\) gives \(-q\eta_\pm\), while wedge by
\(q\eta_\pm\) gives \(q\eta_\pm\). Their sum is zero. Acting on the
other factor does not change this cancellation, so

\[
 \xi_\pm=\pi_\pm+q\eta_\pm
       \in\ker c^2_{v_q}.
 \tag{E819.13}
\]

For \(\operatorname{ch}(E)=r e^{-B}v_q\), conjugate this class by the
usual divisor shear. The identity
\(e^{-B}\iota_v e^B=\iota_v+(\iota_vB)\wedge\) shows directly that the
shear adds mixed and wedge terms but leaves the projection to
\(\bigwedge^2H^0(T_X)\) unchanged. We obtain a class
\(\xi\in\ker c^2_{\operatorname{ch}(E)}\) whose holomorphic bivector
component is the nonzero constant bivector \(\pi_+\). One could also
take both plane components to obtain a nondegenerate bivector, but only
nonzero is needed.

Suppose that \(\sigma^2_{E,G}\) were injective. Hochschild evaluation
lands in invariant Ext by the stated S-21 compatibility.
Buchweitz--Flenner Theorem6.4.2, together with Proposition6.2.1, gives
the ordinary composition identity; restricting it to these invariant
values gives \(\sigma^2_{E,G}\rho^2_{E,G}=c^2_{\operatorname{ch}(E)}\),
with the consistent HKR convention just fixed. Therefore
\(\rho^2_{E,G}(\xi)=0\). Its image in ordinary Ext is zero, and so is
its restriction at every point. By §3, the completed local module of
\(E\) admits the flat first-order deformation (E819.3).

Complete (E819.1) at any point. In characteristic zero the formal
logarithm of the abelian formal group gives simultaneous formal
coordinates in which all invariant vector fields, and hence both
bivectors \(\pi_\pm\), are constant. Its hull is free, and its kernel
and cokernel have the same lengths as before completion. The local lemma
gives (E819.2). This proves SF-E819 for every specified effective
coherent finite action, without any simplicity assumption.

\subsection{4a. Reduced defects cannot survive even with equal
lengths}\label{SF-E819-a.-reduced-defects-cannot-survive-even-with-equal-lengths}

There is a stronger conclusion when, at a point of the modification,
both finite modules are killed by the maximal ideal. Suppose locally

\[
 Q\simeq k^m,\qquad T\simeq k^n,\qquad k=R/\mathfrak m.
 \tag{E819R.1}
\]

If either module is nonzero, a balanced weakly semiregular sheaf cannot
have this local model. This is stronger than the signed-length
conclusion when \(m=n>0\).

Choose a basis of the free hull so that its quotient to \(Q\) is the
residue map on the first \(m\) coordinates. Then

\[
 K=\mathfrak m\otimes\mathbf C^m\oplus R^{r-m}.
 \tag{E819R.2}
\]

The extension of its free summand by \(T\) splits, since
\(\operatorname{Ext}^1_R(R,T)=0\). For the remaining summand choose
lifts \(e_i(v)\in M\), linear in \(v\in\mathbf C^m\), of \(z_i v\in K\).
The Koszul presentation of \(\mathfrak m\) shows that the extension is
encoded by matrices

\[
 B_{ij}:\mathbf C^m\longrightarrow\mathbf C^n,\qquad
 B_{ij}=-B_{ji},
\]

in the relations

\[
 z_i e_j(v)-z_j e_i(v)=B_{ij}v\in T,
 \qquad z_i T=0.
 \tag{E819R.3}
\]

Indeed the first Koszul syzygies generate all relations among the
\(z_i\), and all next syzygies among (E819R.3) are automatically
respected because \(\mathfrak m T=0\). Equivalently these six matrices
are precisely the class in
\(\operatorname{Ext}^1_R(K,T)=\operatorname{Hom}(\mathbf C^m,\mathbf C^n)\otimes\bigwedge^2\mathbf C^4\).

Assume a flat first-order deformation for a nonzero constant bivector
\(\pi\). Lift the map to the free hull as in §2.1. Lift each basis
vector of \(T\) to \(M_\varepsilon\). Its image under
\(\phi_\varepsilon\) reduces to zero, and hence has the form
\(\varepsilon\) times a vector in the free hull. Reducing that vector at
the origin and projecting to the first \(m\) coordinates defines a
matrix

\[
 C:\mathbf C^n\longrightarrow\mathbf C^m.
\]

Lift the generators \(e_i(v)\) as well. Their images have the form
\(z_i v+\varepsilon a_i(v)\). The lift of the left side of (E819R.3),
minus the lifted right side, is \(\varepsilon\) times an element of
\(M_\varepsilon\), because its reduction is zero. Apply
\(\phi_\varepsilon\), project to the first \(m\) coordinates, and reduce
the coefficient of \(\varepsilon\) at the origin. Terms
\(z_i a_j-z_j a_i\) vanish there. The projected value of the correction
term is zero because \(\phi(M)=K\) maps to zero in \(Q\). The commutator
(E819.4) therefore gives the exact matrix equations

\[
 C B_{ij}=\pi_{ij} I_m\quad\text{for every }i,j.
 \tag{E819R.4}
\]

The local trace lemma has already forced \(m=n\). If \(m=0\), the
modification is trivial. Otherwise some \(\pi_{ij}\ne0\), and (E819R.4)
makes \(C\) invertible. Consequently

\[
 B_{ij}=\pi_{ij}C^{-1}\quad\text{for every }i,j.
 \tag{E819R.5}
\]

Thus the tensor of the six extension matrices lies in the
one-dimensional bivector line \(\mathbf C\pi\), tensored with one fixed
invertible matrix.

Balanced weak semiregularity requires local deformations for both
independent bivectors \(\pi_+\) and \(\pi_-\) of §4. The same fixed
extension matrices would satisfy (E819R.5) for each. This is impossible.
Explicitly, choose a linear functional on \(\bigwedge^2\mathbf C^4\)
that vanishes on \(\pi_+\) and not on \(\pi_-\). Applying it to
(E819R.4) for the second deformation gives zero on the left and a
nonzero scalar identity on the right. This proves the stronger
reduced-defect obstruction.

This proof does not require the two first-order deformations to fit
together into a deformation over a two-dimensional base: separate
vanishing of the two obstruction classes already contradicts the fixed
matrices. For nonreduced finite modules the simple matrix presentation
(E819R.3) is unavailable; equal signed length there remains an open case
in this argument.

\subsection{5. An actual nonpure balanced family that this obstruction
excludes}\label{SF-E819-an-actual-nonpure-balanced-family-that-this-obstruction-excludes}

Use the explicit CM fourfold \(X=E_i^4\), genuine RM matrix \(U\), and
integral divisor classes \(X_0,Y\) from SF-E815:

\[
 U=\operatorname{diag}(A,A),\quad
 A=\begin{pmatrix}1&1\\1&-1\end{pmatrix},\quad
 S=\operatorname{diag}(1,1,-1,-1),\quad
 X_0=\theta_S,\quad Y=\theta_{US}.
\]

Write \(X_0=x+y\), \(Y=\sqrt2(x-y)\), with \(x^2=-\theta_+^2\),
\(y^2=-\theta_-^2\), and \(\int_Xx^2y^2=4\). The following degree-four
isogeny makes a rational half-divisor into an actual integral divisor.
Let

\[
 T_0=\begin{pmatrix}1&0\\1&1+i\end{pmatrix},\qquad
 T=\operatorname{diag}(T_0,T_0),\qquad f=[T]:X\to X.
\]

Then \(\deg f=4\), and direct matrix multiplication gives

\[
 T^*(US)T=\operatorname{diag}(2,-2,-2,2).
 \tag{E819.14}
\]

Thus \(f^*(Y/2)\) is the integral class
\(D_0=\theta_{\operatorname{diag}(1,-1,-1,1)}\). Choose actual line
bundles \(L_{\pm D_0}\) and set

\[
 Q_\pm=f_*L_{\pm D_0}.
\]

These are locally free rank-four sheaves, since \(f\) is finite étale.
The finite étale Grothendieck--Riemann--Roch formula and
\(\operatorname{td}(T_X)=1\) give the entire characters

\[
 \operatorname{ch}(Q_\pm)=4e^{\pm Y/2}.
 \tag{E819.15}
\]

Indeed \(f^*(\pm Y/2)=\pm D_0\), and \(f_*f^*\) multiplies rational
cohomology by four. This invokes the already checked Fulton
Corollary18.3.1(c), printed p.354, or equivalently the finite étale
specialization of GRR; no Fourier--Mukai transform is needed for the
construction.

Form the actual vector bundle

\[
 H=(L_{X_0}\oplus L_{-X_0})^{\oplus6}
            \oplus(Q_+\oplus Q_-)^{\oplus3}.
 \tag{E819.16}
\]

Its rank is thirty-six, its odd character components vanish, and its
entire even character is

\[
 \begin{aligned}
 \operatorname{ch}_2(H)&=6X_0^2+3Y^2=12(x^2+y^2),\\
 \operatorname{ch}_4(H)&=\tfrac12 X_0^4+\tfrac1{16}Y^4
                         =\tfrac92 x^2y^2.
 \end{aligned}
 \tag{E819.17}
\]

Choose a point \(p\), and use the full fibre quotient
\(H\twoheadrightarrow k_p^{36}\). Let \(K\) be its kernel. The exact
sequence and local freeness of \(H\) give

\[
 \operatorname{Ext}^1_X(K,k_p^{34})
     \simeq\operatorname{Ext}^2_X(k_p^{36},k_p^{34})
     \simeq\operatorname{Hom}(\mathbf C^{36},\mathbf C^{34})
                   \otimes\bigwedge^2T_pX.
 \tag{E819.18}
\]

The first isomorphism follows because
\(\operatorname{Ext}^{>0}_X(H,k_p^{34})=0\). The second is the point
Koszul complex. In particular this is a nonzero space of actual
extension classes. For any nonzero class form the actual coherent
extension

\[
 0\longrightarrow k_p^{34}\longrightarrow E\longrightarrow K
          \longrightarrow0.
 \tag{E819.19}
\]

The maximal torsion subsheaf of \(E\) is precisely the displayed point
sheaf, because \(K\) is torsion-free. The double dual is \(H\), and the
canonical map has kernel \(k_p^{34}\) and cokernel \(k_p^{36}\). Thus
its signed local length is two. Since \([p]=x^2y^2/4\), its full
character is

\[
 \begin{split}
 \operatorname{ch}(E)
  &=\operatorname{ch}(H)-2[p]
    =36+12(x^2+y^2)+4x^2y^2\\
  &=36(1-\theta_+^2/3)(1-\theta_-^2/3)
    =36v_{2/3}.
 \end{split}
 \tag{E819.20}
\]

Every such \(E\) is perfect on the smooth fourfold and ordinarily
gluable because it is a coherent sheaf. This is an actual balanced
construction with point torsion and point defect; neither is merely a
virtual class. No indecomposability, simplicity, or finite action is
asserted. For every finite coherent action that might exist, however,
SF-E819 excludes weak semiregularity because \(36-34=2\ne0\).

\subsection{6. Relation to the overlapping three-cohomology
attempt}\label{SF-E819-relation-to-the-overlapping-three-cohomology-attempt}

The signed primitive-divisor class that motivated the construction is

\[
 6(e^{X_0}+e^{-X_0})+3(e^Y+e^{-Y})-2
        =16v_{3/2}.
 \tag{E819.21}
\]

It has five distinct normalized divisor classes and a negative aggregate
coefficient at the zero class. Thus the positive-weight hypothesis of
SF-E815 does not apply. A three-cohomology model can put the four
nonzero classes in degree zero, \(\mathcal O^{36}\) in degree \(-1\),
and \(\mathcal O^{34}\) in degree \(-2\); the repeated zero class also
defeats the globally disjoint-class hypothesis of SF-E816b. The
Fourier--Mukai interpretation suggested this model, but equations
(E819.14)--(E819.20) construct its nonpure coherent counterpart
directly, so no claim about an undeclared inverse-equivalence
deformation is needed for the theorem or its example.

Among nontrivial point modifications with locally free hull, the new
local tests leave open only zero signed length at every point, with at
least one nonreduced defect module at each point where the modification
is nonzero. They do not assert that such modifications are semiregular.
A locally free sheaf itself, with \(T=Q=0\), is not excluded by SF-E819;
its earlier exclusions require the separate SF-E811--E813 hypotheses.
The other outstanding possibilities include singular reflexive hulls,
positive-dimensional torsion or defects, twisted objects outside this
local ordinary-module argument, and derived complexes not equivalent to
the coherent models just analyzed.

\section{SF-E820: common curvilinear point modules cannot repair the
balanced
character}\label{SF-E820-sf-e820-common-curvilinear-point-modules-cannot-repair-the-balanced-character}

Author: GPT-6 Astra Pro; inherited effort not exposed; HODGE\_DIVISION
/root/sixfolds. Date: 2026-11-21.

\textbf{Status:} Complete proof independently whole-reviewed by /root
and /root/g34\_eightfolds on 21 November 2026; both mathematical reviews
PASS. Root freshly reran the author certificate; g34 independently
reconstructed the exterior and full-character calculation. Reviewers:
GPT-6 Astra Pro, inherited effort not exposed. Reviewed draft SHA-256:
\texttt{311c3238ffbd52dccbbd2b4815f8d482378821b0d1bf80052c2ccba2835fe938}.
This is a separate extension of SF-E819R to a specified nonreduced
family. No conclusion for arbitrary Artinian defect modules is asserted.

\subsection{1. Precise local and global
statements}\label{SF-E820-precise-local-and-global-statements}

Let \(R=\mathbf C[[z_1,z_2,z_3,z_4]]\), and let \(\pi_+,\pi_-\) be
constant bivectors that have rank two and complementary two-dimensional
supports. In particular,

\[
 \pi_+\wedge\pi_+=\pi_-\wedge\pi_-=0,
 \qquad \pi_+\wedge\pi_-\ne0.
 \tag{E820.1}
\]

Let \(I\subset R\) be a curvilinear ideal of positive colength \(\ell\).
This means that there are formal parameters \(u_1,u_2,u_3,u_4\) with

\[
 I=(u_1^\ell,u_2,u_3,u_4),\qquad
 A=R/I\simeq\mathbf C[t]/(t^\ell).
 \tag{E820.2}
\]

Suppose that a finite \(R\)-module \(M\) has a map to a free hull

\[
 0\longrightarrow T\longrightarrow M\xrightarrow{\phi}R^r
     \longrightarrow Q\longrightarrow0,
 \qquad T\simeq A^n,\quad Q\simeq A^m.
 \tag{E820.3}
\]

If \(m+n>0\), then \(M\) cannot admit flat first-order deformations for
both multiplication classes \(\pi_+\) and \(\pi_-\).

Consequently, a balanced sheaf in the scope of SF-E819 cannot be weakly
semiregular for any specified effective coherent finite
translation/tensor action if it has a nontrivial local modification of
type(E820.3) at even one point. The formal parameters in(E820.2) need
not be simultaneous constant coordinates for the bivectors. Their
brackets in the new parameters are retained in the proof.

\subsection{2. The extension matrices over the Artinian
quotient}\label{SF-E820-the-extension-matrices-over-the-artinian-quotient}

The local trace lemma SF-E819 gives \(\ell m=\ell n\), hence \(m=n\), as
soon as either nonzero bivector deformation exists. If one of \(m,n\)
were zero while their sum were positive, this already gives a
contradiction. We can therefore assume \(m=n>0\).

Choose a basis of the free hull so that its quotient to \(Q=A^m\) is
reduction modulo \(I\) on the first \(m\) coordinates. Such a basis
exists: lift the quotient matrix to \(R\); an invertible \(m\)-minor
modulo the maximal ideal is a unit minor over \(R\), and elementary
invertible changes reduce the surjection to the displayed projection.
Its kernel is

\[
 K=I^{\oplus m}\oplus R^{r-m}.
 \tag{E820.4}
\]

The free summand contributes no \(\operatorname{Ext}^1(-,T)\), so it
splits off from the extension by \(T\). Put

\[
 f_1=u_1^\ell,\quad f_2=u_2,\quad f_3=u_3,\quad f_4=u_4.
\]

This is a regular sequence. For each \(v\in\mathbf C^m\), choose a lift
\(e_i(v)\in M\) of \(f_i v\in I^m\). The Koszul presentation gives
matrices

\[
 B_{ij}\in\operatorname{Mat}_m(A),\qquad B_{ij}=-B_{ji},
\]

such that

\[
 f_i e_j(v)-f_j e_i(v)=B_{ij}v\in T.
 \tag{E820.5}
\]

These matrices encode the whole extension class. Indeed applying
\(\operatorname{Hom}_R(-,T)\) to the Koszul resolution of \(A^m\) makes
every differential zero, since \(I T=0\). Thus

\[
 \operatorname{Ext}^1_R(K,T)
  \simeq\operatorname{Ext}^2_R(A^m,A^m)
  \simeq\operatorname{Mat}_m(A)\otimes\bigwedge^2\mathbf C^4.
 \tag{E820.6}
\]

There is no condition that \(\mathfrak m T=0\); only \(I T=0\) is used
here. This is the change from the reduced calculation.

\subsection{3. Necessary equations for each bivector
deformation}\label{SF-E820-necessary-equations-for-each-bivector-deformation}

Assume a flat first-order deformation for one of the bivectors \(\pi\).
Lift \(\phi\) to the free deformed hull by SF-E819, section2.1. For an
element of \(T\), the image of a chosen lift is \(\varepsilon\) times a
hull vector. Its class modulo \(K\) is independent of the lift, since
changing that lift by \(\varepsilon M\) changes the vector by an element
of \(K\). It defines an \(R\)-linear map

\[
 C:T\longrightarrow Q,
 \qquad C\in\operatorname{Mat}_m(A).
 \tag{E820.7}
\]

Linearity follows by applying the deformed hull map to a lifted product:
its reduction on \(T\) is zero, so only the ordinary product remains in
the coefficient of \(\varepsilon\).

Lift(E820.5), apply the lifted hull map, and reduce its
\(\varepsilon\)-coefficient modulo \(I\) in the first \(m\) hull
coordinates. Terms \(f_i a_j-f_j a_i\) vanish modulo \(I\). The
correction to the lifted module relation maps into \(K\) and also
vanishes. The first-order commutator is the Poisson bracket of the
actual formal functions \(f_i,f_j\). Consequently

\[
 C B_{ij}=P(\pi)_{ij}I_m,
 \qquad P(\pi)_{ij}=\pi(df_i,df_j)\pmod I.
 \tag{E820.8}
\]

The constant-coordinate star product is evaluated on these possibly
nonlinear functions. No unsupported claim that the brackets stay
constant after straightening the curvilinear scheme is made.

\subsection{4. A unit coefficient makes both bracket forms
proportional}\label{SF-E820-a-unit-coefficient-makes-both-bracket-forms-proportional}

The cotangent vectors \(du_2(0),du_3(0),du_4(0)\) span a
three-dimensional subspace \(W\) of the four-dimensional cotangent
space. The restriction of at least one of \(\pi_+,\pi_-\), viewed as
alternating bilinear forms on that cotangent space, is nonzero on \(W\).

For completeness, a rank-two alternating form on a four-dimensional
space has a two-dimensional kernel. Any three-dimensional isotropic
subspace contains that kernel: otherwise its image in the
two-dimensional symplectic quotient would have dimension two. The
kernels of \(\pi_+\) and \(\pi_-\) are complementary, by(E820.1). A
three-dimensional subspace cannot contain both. Thus \(W\) cannot be
isotropic for both forms.

After interchanging the two signs if necessary, some \(P(\pi_+)_{ij}\)
with \(i,j\in\{2,3,4\}\) has a nonzero constant term. It is therefore a
unit in the local Artinian ring \(A\). Equation(E820.8) makes \(C_+\)
invertible: take determinants of the square matrices, or note that a
product equal to a scalar unit matrix is invertible over a commutative
local ring. It follows that

\[
 B_{ij}=P(\pi_+)_{ij}C_+^{-1}.
 \tag{E820.9}
\]

For the second deformation, substitute this into(E820.8). If
\(D=C_-C_+^{-1}\), then

\[
 D P(\pi_+)_{ij}=P(\pi_-)_{ij}I_m.
\]

The coefficient already chosen to be a unit shows that \(D=a I_m\) for
some \(a\in A\). All the other equations now imply

\[
 P(\pi_-)=a P(\pi_+)
       \quad\text{in }\bigwedge^2 A^4.
 \tag{E820.10}
\]

\subsection{5. The nonzero socle wedge contradicts
proportionality}\label{SF-E820-the-nonzero-socle-wedge-contradicts-proportionality}

Let \(J=(\partial f_i/\partial z_j)\). The bracket matrix is
\(P(\pi)=J\pi J^{\mathsf t}\), reduced modulo \(I\). Exterior powers of
this identity give

\[
 P(\pi_+)\wedge P(\pi_+)=0,
 \qquad
 P(\pi_+)\wedge P(\pi_-)
   =\det(J)(\pi_+\wedge\pi_-).
 \tag{E820.11}
\]

Here both sides are expressed in the ordered four-generator basis of the
Koszul presentation. Since the \(u_i\) are formal parameters,

\[
 \det(J)=\ell u_1^{\ell-1}
      \det\left(\frac{\partial u_i}{\partial z_j}\right).
 \tag{E820.12}
\]

The second determinant is a unit. Its product with \(\ell u_1^{\ell-1}\)
is nonzero in \(A=\mathbf C[u_1]/(u_1^\ell)\), because the ground field
has characteristic zero. By(E820.1), the mixed wedge in(E820.11) is
therefore nonzero. But(E820.10) makes it
\(a P(\pi_+)\wedge P(\pi_+)=0\). This contradiction proves the local
statement.

The global consequence follows from the already reviewed SF-E819
argument: the two balanced cap-kernel classes have bivector components
\(\pi_+,\pi_-\); invariant semiregularity would make both ordinary local
characteristic obstructions vanish. The full published
Buchweitz--Flenner source and the supplied finite-action compatibility
remain exactly those recorded in SF-E819. No additional global
deformation input is used.

\subsection{6. Actual nonreduced balanced candidates covered by the
theorem}\label{SF-E820-actual-nonreduced-balanced-candidates-covered-by-the-theorem}

On the explicit \(E_i^4\) with \(U,X_0,Y\) of SF-E815, put
\(Z_\pm=X_0\pm Y\) and take the actual bundle

\[
 H=\mathcal O^{24}\oplus L_{X_0}^{\oplus2}
       \oplus L_{-X_0}^{\oplus2}
       \oplus L_{Z_+}\oplus L_{-Z_+}
       \oplus L_{Z_-}\oplus L_{-Z_-}.
 \tag{E820.13}
\]

Its complete character, calculated in SF-E815, is

\[
 \operatorname{ch}(H)
  =32-8(\theta_+^2+\theta_-^2)
        +2\theta_+^2\theta_-^2
  =32v_{1/2}.
 \tag{E820.14}
\]

Choose any embedded curvilinear scheme \(Z\) of length \(\ell\ge2\) at a
point, and an integer \(1\le m\le32\). The restriction of \(H\) to the
local Artinian scheme \(Z\) is free, so it has a quotient
\(H\twoheadrightarrow Q=\mathcal O_Z^m\). Let \(K\) be the kernel and
put \(T=\mathcal O_Z^m\). Local freeness of \(H\) and zero-dimensional
support give

\[
 \operatorname{Ext}^1_X(K,T)
   \simeq\operatorname{Ext}^2_X(Q,T)
   \simeq\operatorname{Mat}_m(A)\otimes\bigwedge^2\mathbf C^4.
 \tag{E820.15}
\]

The last space has dimension \(6m^2\ell>0\), by the same Koszul
calculation and the absence of higher cohomology for a point-supported
sheaf. Every nonzero class gives an actual coherent extension

\[
 0\longrightarrow T\longrightarrow E\longrightarrow K
      \longrightarrow0.
 \tag{E820.16}
\]

Its maximal torsion is \(T\), its double dual is \(H\), and its hull-map
cokernel is \(Q\). Its entire character is
\(\operatorname{ch}(E)=\operatorname{ch}(H)+\operatorname{ch}(T)-\operatorname{ch}(Q)=32v_{1/2}\).
Thus these are actual nonpure, ordinarily gluable, nonreduced balanced
sheaves of zero signed local length. They are not excluded merely by the
signed-length statement of SF-E819 or by its reduced corollary. SF-E820
excludes weak semiregularity for every specified coherent finite action
that might exist. No simplicity, indecomposability, or such action is
asserted.

The remaining local case includes defect modules not free over a common
curvilinear Artinian quotient; even equal-length cyclic modules of
embedding dimension at least two are not settled by this proof. Singular
reflexive hulls, positive-dimensional modifications and unrelated
derived objects retain their earlier open status.

\section{SF-E821: local coisotropy excludes ample surface summands in
balanced
cones}\label{SF-E821-sf-e821-local-coisotropy-excludes-ample-surface-summands-in-balanced-cones}

Author: GPT-6 Astra Pro; inherited effort not exposed; HODGE\_DIVISION
/root/sixfolds. Date: 2026-11-21.

\textbf{Status:} Independently whole-reviewed by /root and
/root/imprimitive on 21 November 2026; mathematical reviews PASS, with
fresh peer primary-source checks and independent full exterior
verification for E821. Reviewers: GPT-6 Astra Pro, inherited effort not
exposed. Reviewed draft SHA-256:
\texttt{770138f789de62ebf83087d3f3e5322243727e854e8d458a63003cb19835d749}.
This is a local support criterion, followed by a specific excluded
family. No claim is made for arbitrary extensions whose surface part is
not a local direct summand.

\subsection{1. The support criterion}\label{SF-E821-the-support-criterion}

Let \(X\) be a complex abelian fourfold with genuine real multiplication
of type \((2,2)\), holomorphic tangent decomposition
\(V=V_+\oplus V_-\), and compatible polarization
\(\theta=\theta_++\theta_-\). Put

\[
 v_q=(1-q\theta_+^2/2)(1-q\theta_-^2/2),\qquad q>0.
\]

Let \(E\) be a perfect complex with \(e^B\operatorname{ch}(E)=r v_q\),
where \(r>0\) and \(B\) is an algebraic rational divisor. Suppose that,
near a general smooth point of a reduced irreducible surface
\(S\subset X\), the complex \(E\) has a direct summand \(i_*F[a]\). Here
\(i:S\hookrightarrow X\), \(F\) is a nonzero locally free sheaf on that
smooth part of \(S\), and \(a\) is any integer. Other local summands may
be arbitrary perfect complexes.

If \(E\) is weakly semiregular for a specified effective coherent finite
translation/tensor action, then the conormal plane of \(S\) is isotropic
for both bivectors \(\pi_+\in\bigwedge^2V_+\) and
\(\pi_-\in\bigwedge^2V_-\). Equivalently, at a general smooth point,

\[
 T_pS=\ell_+(p)\oplus\ell_-(p),
 \qquad \ell_\pm(p)\subset V_\pm
 \text{ are complex lines}.
 \tag{E821.1}
\]

Consequently

\[
 \int_S\theta_+^2=\int_S\theta_-^2=0.
 \tag{E821.2}
\]

In particular, such a local summand is impossible on an effective
surface cycle of class \(k^2\theta^2\), for \(k>0\). More generally it
is impossible whenever either pairing in(E821.2) is positive.

This statement does not exclude every surface. In particular, a
real-multiplication abelian surface with tangent type \((1,1)\) has
exactly the mixed tangent plane(E821.1). Such surfaces actually occur in
the explicit \(E_i^4\) example below: either coordinate twofold
corresponding to one \(2\times2\) block of \(U\) has one tangent
eigenline in each RM plane. The criterion imposes no contradiction on
that support class by itself.

\subsection{2. Proof by the two local first-order multiplication
obstructions}\label{SF-E821-proof-by-the-two-local-first-order-multiplication-obstructions}

The full cap-kernel calculation in SF-E819, section4, supplies classes
whose holomorphic bivector components are the two independent
\(\pi_+,\pi_-\). Divisor shear leaves those leading components
unchanged. Published Buchweitz--Flenner Theorem6.4.2 and
Proposition6.2.1 identify evaluation followed by semiregularity with
contraction by the full Chern character. The precise finite-action
compatibility is the same supplied DEG\_L S-21 input as in SF-E819.
Hence injective invariant semiregularity implies that both ordinary
characteristic classes of \(E\) vanish.

Restrict to an affine neighborhood of the point. The
\(H^2(\mathcal O_X)\) and \(H^1(T_X)\) terms of the chosen Hochschild
classes restrict to zero. The remaining local classes are the
first-order multiplication deformations defined by \(\pi_+,\pi_-\).
Naturality of the characteristic morphism with inclusions and
projections shows that a zero characteristic endomorphism on a direct
sum has zero characteristic endomorphism on each direct summand. Shifts
do not change vanishing. Thus each local coherent sheaf \(i_*F\) has
zero first-order characteristic obstruction for both bivectors.

The explicit bar/resolution calculation in SF-E819, section3, now gives
a flat first-order module deformation of \(i_*F\) for each
multiplication class. Let \(I_S\) be its annihilator on the
neighborhood; it equals the reduced ideal of \(S\), because \(F\) is
locally free and nonzero on \(S\). If \(f,g\in I_S\), their actions on a
flat first-order deformation reduce to zero. Each therefore carries the
module into its \(\varepsilon\)-multiple. Their two composites are zero,
since \(\varepsilon^2=0\). The deformed algebra commutator, however,
acts as

\[
 [f,g]_* =\varepsilon\pi(df,dg).
\]

Flatness over the dual numbers identifies the \(\varepsilon\)-multiple
with the original module. We conclude that \(\pi(df,dg)\) annihilates
\(i_*F\), and hence

\[
 \{I_S,I_S\}_{\pi_+}\subset I_S,
 \qquad \{I_S,I_S\}_{\pi_-}\subset I_S.
 \tag{E821.3}
\]

This is the precise first-order coisotropy condition used here. No
all-order quantization or general involutivity theorem is imported.

Let \(N_p^*S\subset V_+^*\oplus V_-^*\) be the two-dimensional conormal
space. Condition(E821.3) says that its projection to either
two-dimensional summand is isotropic for that summand's nondegenerate
alternating form, and thus has dimension at most one. Both projections
must have dimension one: otherwise the two-dimensional conormal space
would project with dimension two to the other summand. Therefore

\[
 N_p^*S=\ell_+^*(p)\oplus\ell_-^*(p)
\]

for lines in the two cotangent planes. Taking annihilators
proves(E821.1). The form \(\theta_+\) has rank one on this surface
tangent plane, and so does \(\theta_-\). Their squares restrict to zero.
Since the same holds on a dense open subset of the smooth locus, it
holds throughout that locus by continuity. The pairing of a smooth form
with the fundamental surface cycle is its integral over that regular
locus; lower-dimensional singular subsets do not contribute to that
cycle integral. This proves(E821.2), including for a surface cycle whose
irreducible components satisfy the local hypothesis.

Finally, \(\theta_+^3=\theta_-^3=0\), while
\(\int_X\theta_+^2\theta_-^2>0\). Thus

\[
 \int_X k^2\theta^2\theta_+^2
   =k^2\int_X\theta_+^2\theta_-^2>0,
\]

and likewise with the two signs interchanged. This contradicts(E821.2)
for a surface cycle of class \(k^2\theta^2\). The proof does not use
gluability or simplicity; it only uses the asserted local summand and
weak semiregularity.

\subsection{3. Why the proposed two-cohomology surface cones have this
local
summand}\label{SF-E821-why-the-proposed-two-cohomology-surface-cones-have-this-local-summand}

Let \(A\) be a vector bundle and let

\[
 B=B_0\oplus C_1\oplus\cdots\oplus C_s,
 \qquad C_j=i_{j*}F_j,
 \tag{E821.4}
\]

where \(B_0\) is coherent and each \(F_j\) is a line bundle, or more
generally a nonzero vector bundle, on the smooth part of a reduced
surface \(S_j\). For any actual class
\(a\in\operatorname{Ext}^2_X(A,B)\), form the actual distinguished
triangle

\[
 B[1]\longrightarrow E_a\longrightarrow A
        \xrightarrow{a}B[2].
 \tag{E821.5}
\]

On an affine neighborhood \(U\) where \(A\) is locally free,
\(\operatorname{Ext}^2_U(A|_U,B|_U)=H^2(U,A^\vee\otimes B)=0\).
Thus(E821.5) splits there as \(A|_U\oplus B|_U[1]\). Near a general
smooth point of a component of any \(S_j\), the corresponding
\(i_{j*}F_j[1]\) is consequently a genuine local direct summand. This is
a local splitting of the actual triangle; the sheaves \(C_j\) are not
being inferred as global summands of an arbitrary extension.

If the full character of \(E_a\) is balanced and any one of these
surface cycles has a positive pairing in(E821.2), SF-E821 excludes weak
semiregularity for every specified coherent finite action that might
exist. This applies to every extension class \(a\), including classes
with nonzero components in the surface Ext groups.

\subsection{4. A signed class outside the primitive-factor
hypotheses}\label{SF-E821-a-signed-class-outside-the-primitive-factor-hypotheses}

Use \(X=E_i^4\) with its standard product principal polarization
\(\theta\), and

\[
 U=\operatorname{diag}\left(
 \begin{pmatrix}1&1\\1&-1\end{pmatrix},
 \begin{pmatrix}1&1\\1&-1\end{pmatrix}\right).
\]

The RM eigenpolarizations satisfy \(\theta=\theta_++\theta_-\) and
\(\theta_U=\sqrt2(\theta_+-\theta_-)\). For an integer \(k\ge1\), put

\[
 D=k\theta,
 \qquad N=k\theta_U.
\]

Direct expansion, using \(\theta_+^3=\theta_-^3=0\), gives the entire
character identity

\[
 \begin{split}
 &18+e^N+e^{-N}
       -e^{D+N}-e^{-D-N}-e^{D-N}-e^{-D+N}\\
 &\qquad=16-4k^2(\theta_+^2+\theta_-^2)
                 +k^4\theta_+^2\theta_-^2
       =16v_{k^2/2}.
 \end{split}
 \tag{E821.6}
\]

The odd terms vanish in the displayed pairs. The degree-two calculation
is \(N^2-(D+N)^2-(D-N)^2=-2D^2-N^2\). The degree-four term is
\((N^4-(D+N)^4-(D-N)^4)/12=k^4\theta_+^2\theta_-^2\). Since
\(\int_X\theta_+^2\theta_-^2=4\), its self-Euler characteristic is
\(256k^4\).

These divisor classes are not primitive on each RM plane. The signed
expression is therefore not covered merely by the primitive-factor
hypotheses of SF-E815 or SF-E816.

\subsection{5. Actual surface sheaves and actual gluable
complexes}\label{SF-E821-actual-surface-sheaves-and-actual-gluable-complexes}

We give an explicit existence construction, rather than leaving(E821.6)
as a virtual class. Write the four elliptic coordinates as
\(x_1,\ldots,x_4\). Choose on each elliptic factor a point \(a_i\) of
exact order \(L=2k+1\). Put

\[
 A_i=\{0,a_i,\ldots,(k-1)a_i\},
 \qquad
 B_i=\{ka_i,\ldots,(2k-1)a_i\}.
\]

Let \(D_0\) be the sum of the coordinate divisors \(x_i=a\) for
\(a\in A_i\), and let \(D_1\) be the analogous sum for \(b\in B_i\).
Each has class \(D=k\theta\), and they have no common irreducible
divisor. Their complete intersection \(S=D_0\cap D_1\) is a reduced
surface. Indeed, locally their equations are products of disjoint
subsets of regular coordinate parameters; the resulting ideal is the
intersection of the prime ideals generated by one parameter from each
subset. Its irreducible components are the smooth coordinate twofolds

\[
 x_i=a\in A_i,\qquad x_j=b\in B_j,\qquad i\ne j.
\]

Their fundamental cycle has class \([S]=k^2\theta^2\).

Choose the standard line bundles with classes \(D\) and \(\pm N\), using
the product polarization and Poincare bundles for the displayed integral
Hermitian matrices. Choose a prime \(p>L\), and a degree-zero line
bundle \(P\) whose restriction to every elliptic factor is a nonzero
point of exact order \(p\). Define

\[
 C_\pm=(L_{\pm N}\otimes L_D\otimes P)|_S,
 \qquad
 B=L_N\oplus L_{-N}\oplus i_*C_+\oplus i_*C_-.
 \tag{E821.7}
\]

The notation \(C_\pm\) in(E821.7) denotes line bundles on \(S\); the
last two terms of \(B\) are their pushforwards.

We claim \(H^0(S,C_+)=H^0(S,C_-)=0\). For the plus sign, the Hermitian
matrix of \(L_N\otimes L_D\) in either two-coordinate block is

\[
 k\begin{pmatrix}2&1\\1&0\end{pmatrix}.
\]

On a component whose two free coordinates belong to that same block,
this is nondegenerate of index one, so it has no sections, even after
any degree-zero twist. This is exactly the already checked Mumford index
theorem, \emph{Abelian Varieties}, section16, printed pp150--153
and162--163. On a component with one free coordinate from each block,
the restriction is an external product of line bundles of degrees \(2k\)
or zero. If either degree is zero, its flat part is an \(L\)-torsion
class plus the specified nonzero \(p\)-torsion class. It is nontrivial
because \(p\) and \(L\) are coprime, so that component has no sections.

Only components with free coordinates \(x_1,x_3\) can remain. The
restriction of a global section to such a component must vanish along
all the divisors \(x_1\in A_1\cup B_1\) and \(x_3\in A_3\cup B_3\): each
is an intersection with a component just shown to have no sections.
Subtracting those \(2k\) points on each factor leaves an external
product of degree-zero line bundles. Again its flat part on either
factor is an \(L\)-torsion class plus a nonzero \(p\)-torsion class, so
it has no sections. The restriction to every irreducible component is
zero. A section of a line bundle on a reduced scheme injects into its
restrictions at the generic points of the irreducible components, so the
section is zero.

For the minus sign the block matrix is
\(k\left(\begin{smallmatrix}0&-1\\-1&2\end{smallmatrix}\right)\). The
identical argument interchanges odd and even coordinate indices; the
possible positive components now have free coordinates \(x_2,x_4\). This
proves the claim.

The bundles \(L_{\pm N}\) themselves have index two and \(h^2=4k^4\), so

\[
 H^0(X,B)=0,
 \qquad H^2(X,B)\supset H^2(X,L_N\oplus L_{-N})
                  \simeq\mathbf C^{8k^4}.
 \tag{E821.8}
\]

Choose any nonzero class
\(a\in\operatorname{Ext}^2_X(\mathcal O_X^{18},B)
=H^2(X,B)\otimes(\mathbf C^{18})^*\), and form(E821.5) with
\(A=\mathcal O_X^{18}\). This is an actual perfect complex, since every
coherent sheaf on the smooth fourfold is perfect. It has cohomology
\(B\) in degree \(-1\) and \(\mathcal O^{18}\) in degree zero.

It is ordinarily gluable. The exact two-cohomology identity proved in
SF-E816a gives
\(\operatorname{Ext}^{-1}(E_a,E_a)=\operatorname{Hom}(\mathcal O^{18},B)=0\).
For degrees at most \(-2\), the standard t-structure gives vanishing
directly from these two cohomological degrees. This argument imposes no
unproved formality or degeneration assertion.

Finally, the Koszul resolution of the actual complete intersection gives

\[
 \operatorname{ch}(i_*C_\pm)
   =e^{\pm N}(e^D-2+e^{-D}).
 \tag{E821.9}
\]

Flat twists do not change this character. Combining(E821.7) and(E821.9),
the class of the actual triangle is exactly(E821.6). Thus these are
actual ordinarily gluable, rank16 balanced perfect complexes, with
\(q=k^2/2\). No simplicity or finite action is asserted. The local
support criterion proves that none can be weakly semiregular for any
specified coherent finite action, because \(S\) has class
\(k^2\theta^2\).

\subsection{6. Remaining scope}\label{SF-E821-remaining-scope}

The local support criterion applies whenever the indicated surface
object is an actual local direct summand, whether or not a global
splitting exists. It does not identify the cohomology pieces of an
arbitrary perfect complex as direct summands. In particular, gluing
through a non-locally-free quotient, through local positive-degree Ext
groups, or through a complex with a genuinely coupled local differential
is outside this proof. Surface support of mixed RM tangent type is not
excluded merely by its existence. These are concrete remaining ways a
surface construction might evade the obstruction.

\section{SF-E821a: a torsion-free quotient still forces the generic
surface
splitting}\label{SF-E821a-sf-e821a-a-torsion-free-quotient-still-forces-the-generic-surface-splitting}

Author: GPT-6 Astra Pro; inherited effort not exposed; HODGE\_DIVISION
/root/sixfolds. Date: 2026-11-21.

\textbf{Status:} Independently whole-reviewed by /root and
/root/imprimitive on 21 November 2026; mathematical reviews PASS, with
fresh peer primary-source checks and independent full exterior
verification for E821. Reviewers: GPT-6 Astra Pro, inherited effort not
exposed. Reviewed draft SHA-256:
\texttt{fc79819f4e29785d58620bfcd32d9c7d878baeab8703bb320458e2fc18e8e7ea}.
It extends only the quotient hypothesis of SF-E821, section3; it does
not identify an arbitrary torsion subsheaf as a direct summand.

In the actual triangle

\[
 B[1]\longrightarrow E\longrightarrow A\xrightarrow{a}B[2],
 \tag{E821a.1}
\]

assume that \(A\) is a torsion-free coherent sheaf on the smooth
integral fourfold \(X\). Suppose that \(B\), near a general smooth point
of a reduced irreducible surface \(S\), has an actual direct summand
\(i_*F\), where \(F\) is nonzero locally free on that smooth part of
\(S\). Then \(E\) has \(i_*F[1]\) as a local direct summand near a
general point of \(S\). In particular the support criterion and
zero-pairing conclusions of SF-E821 apply whenever the full character of
\(E\) is balanced and invariant semiregularity is assumed.

Let \(\eta\) be the generic point of \(S\), and put
\(R=\mathcal O_{X,\eta}\). This is a regular local domain of dimension
two. The finite module \(A_\eta\) is torsion-free. Choose a basis of
\(A_\eta\otimes_R\operatorname{Frac}(R)\). Since \(A_\eta\) is finitely
generated, clearing the finitely many denominators in its generators
gives an injection into a finite free module,

\[
 0\longrightarrow A_\eta\longrightarrow R^s
      \longrightarrow Q\longrightarrow0.
 \tag{E821a.2}
\]

Every finite module over this two-dimensional regular local ring has a
free resolution of length at most two. The exact primary locator is
Stacks Project Proposition10.110.1, tag00O7,
https://stacks.math.columbia.edu/tag/00O7. Its statement and proof,
including the stated finite-resolution dependencies Lemmas10.106.6
and10.104.9, were read for this addendum. Applying
\(\operatorname{Hom}_R(-,N)\) to(E821a.2), for any finite \(N\), gives

\[
 \operatorname{Ext}^2_R(A_\eta,N)
      \simeq\operatorname{Ext}^3_R(Q,N)=0.
 \tag{E821a.3}
\]

Equivalently \(A_\eta\) has projective dimension at most one. This
conclusion can also be seen by truncating a finite free resolution of
\(Q\) and taking its first syzygy; the split free terms do not affect
projective dimension.

A finite free resolution of \(A_\eta\) of length one spreads to a
neighborhood of \(\eta\): lift its finitely many matrix entries after
clearing denominators; the kernel and cokernel of the resulting complex
are coherent and vanish at \(\eta\), so they vanish after shrinking that
neighborhood. The same procedure identifies its last cokernel with \(A\)
there. Shrink further to an affine open \(U\). We obtain

\[
 0\longrightarrow F_1\longrightarrow F_0
     \longrightarrow A|_U\longrightarrow0
\]

with \(F_0,F_1\) finite free. Applying \(\operatorname{Hom}_U(-,B|_U)\)
computes derived Hom by a two-term complex of quasi-coherent sheaves.
Since \(U\) is affine, it has no higher quasi-coherent cohomology, and
hence

\[
 \operatorname{Ext}^2_U(A|_U,B|_U)=0.
 \tag{E821a.4}
\]

The actual connecting class \(a\) in(E821a.1) therefore restricts to
zero on \(U\). The triangle splits there. Taking the assumed direct
summand of \(B|_U\) proves the claim, and the local
characteristic/coisotropy argument of SF-E821 applies at the general
smooth points of \(S\cap U\).

In particular, a singular torsion-free quotient does not by itself evade
the two-cohomology surface obstruction. The result still requires the
specified surface piece to be a genuine local direct summand of \(B\).
It makes no such assertion for a nonsplit local extension of surface
torsion by another sheaf. Quotients with torsion along that surface, or
more complicated local derived couplings, remain outside this addendum.

\section{SF-E822: equal-slope coupled surface modifications require
proportional normal
bivectors}\label{SF-E822-sf-e822-equal-slope-coupled-surface-modifications-require-proportional-normal-bivectors}

Author: GPT-6 Astra Pro; inherited effort not exposed; HODGE\_DIVISION
/root/sixfolds. Date: 2026-11-21.

\textbf{Status:} Complete proof independently whole-reviewed by /root
and /root/g34\_eightfolds on 21 November 2026; both mathematical reviews
PASS. The author and peer independently ran their symbolic resolution
and exterior-norm certificates. Reviewers: GPT-6 Astra Pro, inherited
effort not exposed. Reviewed draft SHA-256:
\texttt{bf9dec6e831599b53db853306b2d0c44867175ed36620c75fd7bed25c0c1b95c}.
Unlike SF-E821, the surface torsion is allowed to be genuinely coupled
to the rest of the local module. The hypotheses on the smooth integral
support, the two vector bundles and their equal degree are essential to
the stated argument.

\subsection{1. Statement}\label{SF-E822-statement}

Let \(X\) be a complex abelian fourfold with the balanced RM data of
SF-E819--E821. Let \(i:S\hookrightarrow X\) be a smooth integral
projective surface, and put \(L=\det N_{S/X}\). Projection of the two
invariant RM bivectors to the normal bundle gives sections

\[
 P_+,P_-\in H^0(S,L).
 \tag{E822.1}
\]

In local equations \(I_S=(f,g)\), the coefficient of \(P_\pm\) in the
corresponding normal determinant frame is \(\pi_\pm(df,dg)\bmod I_S\).

Suppose a coherent sheaf \(M\) has a locally free hull
\(H=M^{\vee\vee}\), and its canonical hull map fits into

\[
 0\longrightarrow i_*T\longrightarrow M\longrightarrow H
       \longrightarrow i_*Q\longrightarrow0,
 \tag{E822.2}
\]

where \(T,Q\) are vector bundles on \(S\) of the same positive rank
\(m\), satisfying

\[
 \int_S c_1(T)\,\theta|_S
     =\int_S c_1(Q)\,\theta|_S.
 \tag{E822.3}
\]

Assume \(e^B\operatorname{ch}(M)=r v_q\), with \(r,q>0\) and \(B\) an
algebraic rational divisor. If \(M\) is weakly semiregular for a
specified effective coherent finite translation/tensor action, then the
two sections in(E822.1) span a vector space of dimension at most one.

Consequently no such weakly semiregular \(M\) exists when
\([S]=k^2\theta^2\), for \(k>0\). The statement applies to actual
nonsplit local extensions; it does not assume that \(i_*T\) is a local
direct summand.

\subsection{2. Intrinsic extension class and the actual local
resolution}\label{SF-E822-intrinsic-extension-class-and-the-actual-local-resolution}

Write \(K=\ker(H\to i_*Q)\). The exact sequence
\(0\to i_*T\to M\to K\to0\) has a local extension class

\[
 \mathcal E\in H^0\bigl(S,\mathcal H om(Q,T)\otimes L\bigr).
 \tag{E822.4}
\]

This is intrinsic: locally freeness of \(H\) identifies
\(\mathcal E xt^1_X(K,i_*T)\) with \(\mathcal E xt^2_X(i_*Q,i_*T)\), and
the two-element Koszul resolution identifies the latter with
\(i_*(\mathcal H om(Q,T)\otimes\det N_{S/X})\). The edge map from the
actual global extension gives(E822.4). Extra global extension data do
not change this local class.

Choose a small affine neighborhood where \(S\) has regular-sequence
equations \(f,g\), where \(H,T,Q\) are trivialized as needed, and where
the quotient \(H\to i_*Q\) is reduction modulo \((f,g)\) on the first
\(m\) coordinates. Then
\(K=(f,g)^{\oplus m}\oplus\mathcal O^{\operatorname{rk}H-m}\). The free
summand splits from the extension locally. The remaining module is the
quotient of the three generator blocks \(R^m\oplus R^m\oplus R^m\) by
the submodule generated by the following relation vectors, for
\(w\in R^m\):

\[
 (-gw,fw,-\widetilde{\mathcal E}w),
 \qquad (0,0,fw),\qquad (0,0,gw).
 \tag{E822.5}
\]

where \(\widetilde{\mathcal E}\) is a matrix lift of(E822.4). In
column-map convention its finite free resolution is

\[
 0\longrightarrow R^m\xrightarrow{d_2}R^{3m}
       \xrightarrow{d_1}R^{3m}\longrightarrow M_{\rm loc}
       \longrightarrow0,
 \tag{E822.6}
\]

with

\[
 d_1=\begin{pmatrix}
 -gI_m&0&0\\
 fI_m&0&0\\
 -\widetilde{\mathcal E}&fI_m&gI_m
 \end{pmatrix},
 \qquad
 d_2=\begin{pmatrix}0\\-gI_m\\fI_m\end{pmatrix}.
 \tag{E822.7}
\]

The first block of columns is the first relation vector in(E822.5). In
particular, the extension matrix maps the quotient-normal relation into
the torsion generators; no transpose convention is left implicit.

Here is the exactness check. If \(d_1(a,b,c)=0\), then \(fa=ga=0\),
hence \(a=0\) since \(f\) is a nonzerodivisor. The remaining equation is
\(fb+gc=0\). The two-element Koszul complex gives \((b,c)=(-gh,fh)\),
which is precisely \(\operatorname{im}d_2\). The same regular-sequence
calculation shows \(d_2\) is injective. The cokernel has the stated
extension by \((R/(f,g))^m\), since the first relation reduces modulo
the torsion generators to the presentation of \((f,g)^m\), and no
nonzero element of the torsion quotient is killed by an additional
relation. Thus(E822.6) resolves the actual module, including at zeros of
\(\det\mathcal E\).

\subsection{3. The characteristic obstruction forces a right inverse up
to the normal
section}\label{SF-E822-the-characteristic-obstruction-forces-a-right-inverse-up-to-the-normal-section}

Let \(\pi\) be either RM bivector. For this matrix calculation use right
modules with column vectors, so maps are left multiplication by the
displayed matrices. This entails no extra vanishing assumption: a
left-module deformation for the opposite product is a right-module
deformation for the stated product, and changing to the opposite
first-order product replaces \(\pi\) by \(-\pi\). Vanishing is linear in
the Hochschild class, so the two balanced vanishings also give these
sign-changed vanishings. Relabel the sign if necessary; it has no effect
on linear dependence of \(P_+,P_-\).

If this characteristic obstruction on \(M_{\rm loc}\) vanishes, the
right-module version of the identical bar/resolution argument of
SF-E819, section3, gives a flat first-order module deformation. Lifting
the finite free resolution, its differentials are the displayed matrices
with the first-order star product, corrected by matrices
\(\varepsilon h_1,\varepsilon h_2\).

The uncorrected star-square is

\[
 d_1*d_2
   =\varepsilon\begin{pmatrix}0\\0\\-\pi(df,dg)I_m\end{pmatrix}.
 \tag{E822.8}
\]

In particular no derivative of \(\widetilde{\mathcal E}\) occurs: its
block is multiplied by the zero block of \(d_2\). The correction
equation is

\[
 d_1h_2+h_1d_2
   =\begin{pmatrix}0\\0\\\pi(df,dg)I_m\end{pmatrix}.
 \tag{E822.9}
\]

Reduce modulo \((f,g)\). The second summand vanishes because \(d_2\)
reduces to zero. Only the \(-\mathcal E\) block of \(d_1\) remains.
Absorbing its sign into the first block of \(h_2\), we obtain a local
morphism

\[
 A_\pi:T\longrightarrow Q,
 \qquad
 \mathcal E A_\pi=P_\pi I_T
       \quad\text{as maps }T\to T\otimes L.
 \tag{E822.10}
\]

The type of \(A_\pi\) follows also from the resolution twists: both the
domain of \(d_2\) and the first column block of \(d_1\) carry the same
factor \(L^{-1}\), leaving \(\mathcal H om(T,Q)\) after reduction.
Changes of normal equations multiply \(\mathcal E\) and \(P_\pi\) by the
same determinant-frame unit; changes of the two bundle bases give the
usual source/target transformations. Thus(E822.10) is intrinsic.

Balanced weak semiregularity supplies vanishing of both ordinary local
characteristic obstructions by exactly the published BF2008 and supplied
invariant-evaluation bridge of SF-E819. Therefore(E822.10) holds locally
for both signs.

\subsection{4. Equal degree makes the local inverse global and
invertible}\label{SF-E822-equal-degree-makes-the-local-inverse-global-and-invertible}

If both sections \(P_+,P_-\) are zero, the asserted conclusion is
already true. Otherwise interchange the signs so that \(P_+\ne0\). On
the dense open set where \(P_+\) is nonzero, equation(E822.10) makes
\(\mathcal E\) an isomorphism, since both bundles have rank \(m\). Its
kernel as a sheaf on the integral surface is zero: it is a subsheaf of a
locally free sheaf with zero generic stalk.

It follows that a local solution \(A_\pi\) of(E822.10) is unique. Two
solutions differ by a map into that zero kernel. The local solutions
therefore glue to global morphisms \(A_\pm:T\to Q\).

The determinant of \(A_+\) is a nonzero section of
\(\det Q\otimes(\det T)^{-1}\). By(E822.3), this line bundle has degree
zero against the ample class \(\theta|_S\). Its zero divisor would be an
effective curve of degree zero, which is impossible unless it is empty.
Hence \(A_+\) is everywhere invertible.

Consequently

\[
 \mathcal E=P_+A_+^{-1},
 \qquad
 P_+ A_+^{-1}A_-=P_-I_T.
 \tag{E822.11}
\]

The trace of the global endomorphism \(A_+^{-1}A_-\) is a global regular
function on the integral projective surface, and is therefore constant.
Taking its trace in(E822.11) shows that \(P_-=cP_+\) for a constant
\(c\in\mathbf C\). This proves the first assertion, without any
assertion that the hull map itself deforms. Such a hull-map lifting
would generally be unjustified for codimension-two defects.

\subsection{5. Independence on an ample complete-intersection
surface}\label{SF-E822-independence-on-an-ample-complete-intersection-surface}

Fix an invariant holomorphic four-form on \(X\). Adjunction identifies
\(L\) with \(K_S\). Under this identification, the normal projection of
an invariant bivector \(\pi\) is the restriction to \(S\) of its
contraction with that four-form. In local coordinates this is exactly
the coefficient \(\pi(df,dg)\) defining(E822.1).

For \(\pi_+,\pi_-\), these contracted forms are nonzero scalar multiples
of the holomorphic volume forms \(\Omega_-,\Omega_+\) on the opposite RM
planes. Suppose \([S]=k^2\theta^2\), and take any nonzero linear
combination \(\Omega=a\Omega_++b\Omega_-\). With the compatible positive
Hermitian forms, each \(\Omega_\pm\wedge\overline{\Omega_\pm}\) is a
positive constant times \(\theta_\pm^2\). The mixed terms have zero
pairing with \(\theta^2\), by their separate holomorphic and
antiholomorphic degrees in the two planes. Therefore

\[
 \int_S\Omega\wedge\overline\Omega
  =k^2\bigl(c_+|a|^2+c_-|b|^2\bigr)
       \int_X\theta_+^2\theta_-^2>0
 \tag{E822.12}
\]

for positive constants \(c_\pm\). No nonzero combination restricts to
zero on \(S\). Thus \(P_+,P_-\) are linearly independent, contradicting
section4.

\subsection{6. Same-character surface repairs satisfy the numerical
hypotheses}\label{SF-E822-same-character-surface-repairs-satisfy-the-numerical-hypotheses}

Suppose in(E822.2) that \(\operatorname{ch}(M)=\operatorname{ch}(H)\),
with \(T,Q\) nonzero vector bundles on the smooth integral surface. Then
\(\operatorname{ch}(i_*T)=\operatorname{ch}(i_*Q)\). Their leading
components are \(\operatorname{rk}(T)[S]\) and
\(\operatorname{rk}(Q)[S]\). Pairing with \(\theta^2\) proves the two
ranks are equal. The next component of the embedding GRR formula is

\[
 \operatorname{ch}_3(i_*T)
   =i_*\left(c_1(T)-\tfrac{\operatorname{rk}T}{2}c_1(N_{S/X})\right),
 \tag{E822.13}
\]

and similarly for \(Q\). The equal ranks cancel the normal term. Pairing
the resulting equality with \(\theta\) gives(E822.3). This is the
smooth-embedding specialization of the already checked Fulton,
\emph{Intersection Theory}, second edition, Corollary18.3.1(c), p354;
expand \(\operatorname{td}(N_{S/X})^{-1}\) through degree one. No
additional Riemann--Roch assertion is assumed.

Thus a balanced locally free hull cannot be repaired, while retaining
its entire character, by equal-character vector-bundle torsion and
defect on such an ample complete-intersection surface. This conclusion
includes choices of \(T,Q\) that are not initially isomorphic.

\subsection{7. Actual nonsplit local balanced examples to which the
obstruction
applies}\label{SF-E822-actual-nonsplit-local-balanced-examples-to-which-the-obstruction-applies}

On the explicit CM fourfold \(E_i^4\), take the actual rank32 bundle
\(H\) of SF-E815, with \(\operatorname{ch}(H)=32v_{1/2}\). Choose a
smooth integral complete intersection \(S=D_0\cap D_1\) of two divisors
in \(|3\theta|\). Such surfaces exist: the previously checked
Bauer--Szemberg1996 Theorem1.1 gives very ampleness of three ample
factors; Bertini is applied twice to that embedding. Exact additional
locator: Stacks Lemma33.47.3 and Remark37.32.4, tags0FD4/0G4G. The
smoothness proof and the smooth-irreducible hyperplane statement were
freshly read for this use; no independent audit of every earlier Bertini
dependency is claimed.

The ideal sheaf has the actual Koszul resolution

\[
 0\to\mathcal O_X(-D_0-D_1)
       \to\mathcal O_X(-D_0)\oplus\mathcal O_X(-D_1)
       \to I_S\to0.
\]

Applying sheaf Hom into \(\mathcal O_S\) makes its differential zero.
Hence each section of \(\mathcal O_S(D_0+D_1)=K_S\) gives an actual
class in \(\operatorname{Ext}^1_X(I_S,\mathcal O_S)\): it is the
degree-one direct summand of the resulting zero-differential Hom
complex. In particular choose the nonzero section \(e=P_+\)
from(E822.1), and form the actual coherent extension

\[
 0\longrightarrow\mathcal O_S\longrightarrow M_e
      \longrightarrow I_S\longrightarrow0.
 \tag{E822.14}
\]

Its class in \(K_0(X)\) equals \([\mathcal O_X]\), since
\([I_S]+[\mathcal O_S]=[\mathcal O_X]\). Set

\[
 E=H\otimes M_e.
 \tag{E822.15}
\]

This is an actual coherent, hence ordinarily gluable, perfect sheaf. Its
hull is \(H\), its torsion and hull-map cokernel are both \(i_*(H|_S)\),
and its complete character is \(32v_{1/2}\). The intrinsic local
extension matrix is \(e I_{32}\). On the dense open subset of \(S\)
where \(e\ne0\), that matrix is invertible. The local extension is
therefore nonsplit: the Koszul class is nonzero, and projecting its
torsion to any positive-rank direct summand still gives a nonzero
extension class. The surface torsion is genuinely coupled there, so the
local-direct-summand hypothesis of SF-E821 is not being used.

SF-E822 excludes weak semiregularity for every specified coherent finite
action that might exist. No such action, simplicity, or a semiregular
object is asserted to have been constructed.

\subsection{8. Remaining scope}\label{SF-E822-remaining-scope}

The proof requires a smooth integral support and vector bundles on that
support, equal ranks, and equal degree as in(E822.3). Unequal degrees
allow determinant zero divisors and destroy the global-inverse step.
Nonreduced support, singular support with non-locally-free modules,
different codimensions, and more general derived local couplings are not
settled. Supports for which the two normal bivector sections are
proportional are also not excluded by this criterion alone.

\section{Unequal-slope surface compensation: actual objects and a
precise remaining
gate}\label{SF-US-unequal-slope-surface-compensation-actual-objects-and-a-precise-remaining-gate}

Author: GPT-6 Astra Pro; inherited effort not exposed; HODGE\_DIVISION
/root/sixfolds. Date: 2026-11-21.

\textbf{Status:} Complete actual construction and local calculation,
whole-reviewed by /root and by its author on 21 November 2026;
mathematical checks PASS. Reviewer and author: GPT-6 Astra Pro,
inherited effort not exposed. Both exact author certificates have
actually passed. This positive-character construction alone is not
counted as a reduction of Q5i. The accompanying SF-E823 proves a global
obstruction to every tensor hull repair, independently of splitting. No
semiregular seed or finite action is claimed.

\subsection{1. A unit extension outside the equal-degree
theorem}\label{SF-US-a-unit-extension-outside-the-equal-degree-theorem}

Let \(X=E_i^4\), with the principal product polarization \(\theta\), and
set \(D=3\theta\). Choose a smooth integral complete intersection
\(S=D_0\cap D_1\), with both \(D_j\in|D|\). Existence uses precisely the
very-ampleness and Bertini inputs already located and read for SF-E822.
Write \(i:S\hookrightarrow X\) and \(L=\det N_{S/X}=\mathcal O_S(2D)\).

Take \(Q=\mathcal O_S\) and \(T=L^{-1}=\mathcal O_S(-2D)\). The sheaf of
local extension classes is

\[
 \mathcal E xt^1_X(I_S,i_*T)
 =i_*\bigl(T\otimes\det N_{S/X}\bigr)=i_*\mathcal O_S.
 \tag{US.1}
\]

The Koszul resolution of \(I_S\), after applying sheaf Hom into
\(i_*T\), has zero differential. Consequently its degree-one term gives
an actual direct summand \(H^0(S,\mathcal O_S)\) of the global Ext
group. The unit section therefore defines an actual coherent extension

\[
 0\longrightarrow i_*\mathcal O_S(-2D)
   \longrightarrow M_0\longrightarrow I_S\longrightarrow0.
 \tag{US.2}
\]

Its torsion is the displayed left term, because the quotient \(I_S\) is
torsion-free. The two-element Koszul resolution gives
\(\mathcal H om(i_*T,\mathcal O_X)=\mathcal E xt^1(i_*T,\mathcal O_X)=0\).
Dualizing (US.2) therefore gives \(M_0^\vee=I_S^\vee=\mathcal O_X\); its
double dual is \(\mathcal O_X\), and the canonical hull-map image is
\(I_S\), so its cokernel is \(i_*\mathcal O_S\). Thus the ranks of the
torsion and defect on \(S\) agree, but their degrees differ by
\(2\int_S D\theta>0\). The equal-degree argument of SF-E822 does not
apply.

More explicitly, with local equations \(f,g\) for \(S\), the unit
extension has the free resolution of SF-E822 with extension matrix
\(1\):

\[
 d_1=\begin{pmatrix}-g&0&0\\f&0&0\\-1&f&g\end{pmatrix},
 \qquad d_2=\begin{pmatrix}0\\-g\\f\end{pmatrix}.
 \tag{US.3}
\]

Eliminating the unit entry gives the actual global locally free
resolution

\[
 0\longrightarrow\mathcal O_X(-4D)
 \xrightarrow{(-g,f)^t}\mathcal O_X(-3D)^{\oplus2}
 \xrightarrow{\left(\begin{smallmatrix}-fg&-g^2\\f^2&fg\end{smallmatrix}\right)}
 \mathcal O_X(-D)^{\oplus2}\longrightarrow M_0\longrightarrow0.
 \tag{US.4}
\]

The entries have the asserted twists, and this is obtained by an
elementary split cancellation from the exact resolution (US.3), not
merely a complex whose Euler character is correct. Its entire character
on a fourfold is

\[
 \operatorname{ch}(M_0)=2e^{-D}-2e^{-3D}+e^{-4D}
       =1-2D^3+4D^4.
 \tag{US.5}
\]

The finite expansion follows directly: the degree-one and degree-two
moments vanish, the degree-three moment is \(-12\), and the degree-four
moment is \(96\), before division by the corresponding factorials.

\subsection{2. Exact local homotopies for every invariant
bivector}\label{SF-US-exact-local-homotopies-for-every-invariant-bivector}

Let \(\pi\) be any constant bivector in formal logarithmic coordinates
on the abelian variety. Put \(P=\pi(df,dg)\), and use the first-order
star product \(a*b=ab+\varepsilon\pi(da,db)/2\), with
\(\varepsilon^2=0\). As in SF-E822, take right modules with column maps;
replacing \(\pi\) by \(-\pi\) translates to the left-module convention.

For (US.3), direct multiplication gives

\[
 d_1*d_2=\varepsilon(0,0,-P)^t.
\]

Define the following explicit correction matrices:

\[
 h_2=(-P,0,0)^t,
 \qquad
 h_1=\begin{pmatrix}0&P&0\\0&0&P\\0&0&0\end{pmatrix}.
 \tag{US.6}
\]

Then

\[
 d_1h_2+h_1d_2=(gP,-fP,P)^t+(-Pg,Pf,0)^t=(0,0,P)^t.
 \tag{US.7}
\]

It follows that \((d_1+\varepsilon h_1)*(d_2+\varepsilon h_2)=0\). The
corrected finite free complex reduces to the exact resolution of
\(M_0\). The elementary lifting criterion gives zero negative homology
and a flat cokernel over the dual numbers: reduce a cycle modulo
\(\varepsilon\), lift the boundary, and repeat on its remaining
\(\varepsilon\)-multiple; the same argument proves that the cokernel's
\(\varepsilon\)-annihilator is its \(\varepsilon\)-image. This
constructs an actual first-order local module deformation. Tensoring
with a locally trivial vector bundle gives the identical calculation in
every rank.

In particular both RM bivectors \(\pi_+\) and \(\pi_-\) have zero local
characteristic obstruction on this unit model. The previous local
obstruction is genuinely absent, not just inconclusive. No globalization
is implied: changes of line-bundle frame change the off-surface
derivatives of \(f,g\), and the displayed local homotopies have not been
glued to a global Hochschild nullhomotopy.

The exact script \texttt{rm\_unequal\_slope\_local\_homotopy.py} checks
the complete first-order calculation with arbitrary functions \(f,g\) in
four variables and all six independent constant bivector coefficients.
It also checks the minimal complex square and the entire character
expansion. Its actual output is
\texttt{rm\_unequal\_slope\_local\_homotopy\_output.json}.

\subsection{3. The required compensating hull
character}\label{SF-US-the-required-compensating-hull-character}

Use the integral symmetric matrices

\[
 U=\operatorname{diag}\!\left(
 \begin{pmatrix}1&1\\1&-1\end{pmatrix},
 \begin{pmatrix}1&1\\1&-1\end{pmatrix}\right),
 \qquad S=\operatorname{diag}(1,1,-1,-1).
\]

They commute, \(U^2=2I\), and each of \(I,U,S,US\) defines an integral
divisor on the product CM fourfold. Write its classes as
\(\theta,\theta_U,\theta_S,\theta_{US}\). Let
\(\omega_1,\ldots,\omega_4\) be the joint eigenline forms, ordered so
that \(\theta_+=\omega_1+\omega_2\) and \(\theta_-=\omega_3+\omega_4\).
Each \(\omega_j^2=0\).

For \(\beta=a\theta+b\theta_U+c\theta_S+d\theta_{US}\), the four
eigenvalues are

\[
 \begin{split}
 \lambda_1&=a+b\sqrt2+c+d\sqrt2,&
 \lambda_2&=a+b\sqrt2-c-d\sqrt2,\\
 \lambda_3&=a-b\sqrt2+c-d\sqrt2,&
 \lambda_4&=a-b\sqrt2-c+d\sqrt2.
 \end{split}
 \tag{US.8}
\]

In the sixteen-dimensional algebra generated by these four line forms,
\(e^\beta=\prod_j(1+\lambda_j\omega_j)\). The balanced character is

\[
 v_q=\left(1-\frac q2\theta_+^2\right)
      \left(1-\frac q2\theta_-^2\right).
 \tag{US.9}
\]

Set \(q=1/2\). Because all products of degree greater than four vanish,
(US.5) is invertible with inverse \(1+2D^3-4D^4\). Thus a hull of rank
\(r\) must satisfy

\[
 \frac{\operatorname{ch}(H)}r
       =v_{1/2}+2D^3-4D^4.
 \tag{US.10}
\]

This equation is for the entire character, not only its degree-two or
top component. Any actual bundle with (US.10) gives an actual coherent
perfect, hence ordinarily gluable, sheaf \(H\otimes M_0\) with character
\(r v_{1/2}\).

\subsection{4. An explicit actual split
hull}\label{SF-US-an-explicit-actual-split-hull}

Choose a line bundle \(L_\beta\) for every integral divisor class in the
following table. Form the direct sum with the listed positive
multiplicities.

\begin{longtable}[]{@{}lr@{}}
\toprule\noalign{}
\((a,b,c,d)\) & Multiplicity \\
\midrule\noalign{}
\endhead
\bottomrule\noalign{}
\endlastfoot
\((-16,0,0,0)\) & 23104075775 \\
\((-8,-6,-8,6)\) & 38474612484 \\
\((-8,-6,8,-6)\) & 38474612484 \\
\((0,-11,0,0)\) & 30281171392 \\
\((-8,6,-8,-6)\) & 38474612484 \\
\((0,0,-16,0)\) & 33081896705 \\
\((0,0,0,-11)\) & 30550999744 \\
\((8,-6,-8,-6)\) & 28751863668 \\
\((-8,6,8,6)\) & 38474612484 \\
\((0,0,0,11)\) & 30550999744 \\
\((0,0,16,0)\) & 33081896705 \\
\((8,-6,8,6)\) & 28751863668 \\
\((0,11,0,0)\) & 30281171392 \\
\((8,6,-8,6)\) & 28751863668 \\
\((8,6,8,-6)\) & 28751863668 \\
\((16,0,0,0)\) & 42549573407 \\
\end{longtable}

Their sum is

\[
 r=522387689472.
 \tag{US.11}
\]

Here is a finite exact way to verify all of (US.10). Index the basis by
the subsets of \(\{1,2,3,4\}\), in binary-mask order from 0 to 15. For a
row \((a,b,c,d)\), its exponential has coefficient
\(\prod_{j\in A}\lambda_j\) on the subset \(A\). After multiplying by
the listed multiplicities, adding and dividing by (US.11), the resulting
sixteen coefficients are

\[
 (1,0,0,-\tfrac12,0,0,0,324,0,0,0,324,-\tfrac12,324,324,
       -\tfrac{31103}{4}).
 \tag{US.12}
\]

They are exactly (US.10), since \(D=3\sum_j\omega_j\): the degree-three
coefficient is \(2\cdot3^3\cdot3!=324\), and the top correction is
\(-4\cdot3^4\cdot4!=-7776\). The balanced top term contributes \(1/4\).
The coefficients of \(\operatorname{ch}(M_0)\) are respectively rank 1,
degree-three coefficients \(-324\), and top coefficient \(7776\).
Multiplication therefore leaves only rank 1, coefficients \(-1/2\) on
subsets \(\{1,2\},\{3,4\}\), and top coefficient \(1/4\), as required.

The script \texttt{rm\_unequal\_slope\_compensation.py} verifies these
equalities in \(\mathbf Q(\sqrt2)\), checks positivity and integrality
of every multiplicity, and multiplies the entire resulting character by
(US.5). Its actual output is
\texttt{rm\_unequal\_slope\_compensation\_output.json}. The failed
scales 4 and 8, at which some interpolation weights were negative, are
retained there; the first tested successful scale was 16.

The method is finite interpolation, not a numerical fit. For
completeness, take the sixteen vertices
\((R\epsilon_1,\ldots,R\epsilon_4)\), invert the linear transformation
(US.8), and round its coefficients to integers. After normalizing rows
of degree \(j\) by \(R^{-j}\), the character matrix tends to the
invertible Boolean-cube evaluation matrix. The normalized target tends
to the vector with empty-subset entry one and all other entries zero.
Its weights consequently tend to \(1/16\), so they are strictly positive
for sufficiently large \(R\). Uniqueness and rationality of the original
divisor characters make the weights rational; the explicit table
verifies this directly at scale 16. Clearing their denominators produces
actual positive integral sums of line bundles. This argument asserts
realizability in this specified divisor subalgebra only.

\subsection{5. What remains missing}\label{SF-US-what-remains-missing}

The sheaf constructed with this hull is

\[
 H_{\mathrm{split}}\otimes M_0
   =\bigoplus_\beta (L_\beta\otimes M_0)^{\oplus n_\beta}.
\]

Its distinct nonzero summand characters are not proportional to the
whole balanced character: their degree-one components already
distinguish them. The ordinary direct-sum argument of SF-E802 therefore
excludes ordinary full semiregularity. The same argument excludes a
finite action whenever the displayed decomposition is an actual
decomposition in that equivariant category. No finite action that avoids
this issue has been constructed; no claim about arbitrary noncanonical
equivariant splittings is being substituted for a proof.

The local calculation does not cure this global problem. On its own,
this construction leaves open whether an actual nonsplit hull or
different global extension data can change the result. The accompanying
SF-E823 settles the first of those possibilities: every actual tensor
product with this unit model is excluded, even for an unstable nonsplit
hull and an arbitrary specified finite action on the tensor product. It
does not settle genuinely different global gluing. A signed K-theory
identity, positivity of interpolation weights, or the two local
nullhomotopies alone gives no semiregularity conclusion. This section is
retained as an actual construction and test case, not counted by itself
as a solution or a narrowing of the remaining Q5i gate.

\section{SF-E823: a global obstruction to every tensor hull repair of
the unit surface
extension}\label{SF-E823-sf-e823-a-global-obstruction-to-every-tensor-hull-repair-of-the-unit-surface-extension}

Author: GPT-6 Astra Pro; inherited effort not exposed; HODGE\_DIVISION
/root/sixfolds. Date: 2026-11-21.

\textbf{Status:} Complete proof independently whole-reviewed by /root
and /root/imprimitive on 21 November 2026; both mathematical/source
reviews PASS. Reviewers: GPT-6 Astra Pro, inherited effort not exposed.
The author and peer independently ran exact exterior certificates; root
also reran the author certificate. Reviewed draft SHA-256:
\texttt{1beaed063137d36e0e879432c5e39006fd95f32eb9a3eb24ff2b084e44696d4e}.
No simplicity, stability, splitting, or equivariant structure on the
hull factor is assumed. The assertion concerns actual tensor products;
it is not a claim about all sheaves with the same local extension
matrix.

\subsection{1. Statement and the actual
input}\label{SF-E823-statement-and-the-actual-input}

Let \(X\) be a complex abelian fourfold with the genuine balanced RM
data of SF-E819: \(\theta=\theta_++\theta_-\), with each positive form
supported on its two-dimensional eigenspace. Put

\[
 v_q=\left(1-\frac q2\theta_+^2\right)
       \left(1-\frac q2\theta_-^2\right),\qquad q>0.
 \tag{E823.1}
\]

Choose \(D=k\theta\), \(k>0\), and a smooth complete intersection
\(S=D_0\cap D_1\), with both divisors in \(|D|\). Let \(M_0\) be the
actual unit extension

\[
 0\to\mathcal O_S(-2D)\to M_0\to I_S\to0,
 \qquad \operatorname{ch}(M_0)=1-2D^3+4D^4,
 \tag{E823.2}
\]

constructed in the unequal-slope compensation note. In particular its
rank is one and its first two positive-degree character components
vanish. Let \(H\) be \textbf{any locally free sheaf of positive rank}
\(r\), and suppose

\[
 e^B\operatorname{ch}(H\otimes M_0)=r v_q,
 \qquad B\in\operatorname{NS}(X)_{\mathbf Q}.
 \tag{E823.3}
\]

Then \(H\otimes M_0\) is not weakly semiregular for any specified
effective coherent finite translation/tensor action in the sense of
supplied S21. The finite action need only be given on the tensor
product; no compatible action on \(H\) or \(M_0\) separately is assumed.

The complete proof below is global. The local corrected resolutions in
the compensation note give zero characteristic obstruction for every
constant bivector on every sufficiently small affine neighborhood. Thus
this result does not reassert the former local obstruction.

\subsection{2. The partial trace and its Atiyah-class
identity}\label{SF-E823-the-partial-trace-and-its-atiyah-class-identity}

Let \(F\) be any perfect complex and \(H\) a vector bundle of rank
\(r\). Tensor duality gives a canonical partial trace

\[
 \operatorname{Tr}_H:
 \operatorname{Ext}^2_X(H\otimes F,H\otimes F)
       \longrightarrow\operatorname{Ext}^2_X(F,F).
 \tag{E823.4}
\]

Explicitly, precompose with the coevaluation
\(F\to H^\vee\otimes H\otimes F\), apply the middle morphism, and
postcompose with evaluation. It is defined in the ordinary derived
category regardless of any finite action. Its value on an endomorphism
supported only on \(F\) is \(r\) times that endomorphism.

Write \(A_H,A_F\) for the Chern-normalized Atiyah classes, with
\(\operatorname{ch}(H)=\operatorname{Tr}\exp(A_H)\). Thus our \(A\) is
the negative of the class denoted \(\mathrm{At}\) in BF2008, whose
Atiyah--Chern character is \(\exp(-\mathrm{At})\). This single sign
convention is used in every formula below. Their tensor identity is

\[
 A_{H\otimes F}=A_H\otimes1+1\otimes A_F.
 \tag{E823.5}
\]

One can see this directly on a finite locally free resolution of \(F\):
choose local connections, and use the tensor connection on each term.
Its commutator with the differential is the sum of the two commutators,
while the differences on overlaps obey the identical Leibniz rule. These
represent the Atiyah classes in the total Čech/endomorphism complex,
giving (E823.5) as a derived identity. The Ext sign and the
differential-form sign cancel when the two Atiyah factors interchange;
this is the usual even total-degree convention that also gives
\(\operatorname{ch}(H\otimes F)=
\operatorname{ch}(H)\operatorname{ch}(F)\).

For a rational algebraic divisor \(B\), put

\[
 A_H^B=A_H+B\operatorname{id}_H,
 \quad c_1^B(H)=\operatorname{Tr}A_H^B,
 \quad \operatorname{ch}_2^B(H)=
       \tfrac12\operatorname{Tr}(A_H^B)^2.
 \tag{E823.6}
\]

These are cohomological definitions; they require no fractional line
bundle. They are the first two components of
\(e^B\operatorname{ch}(H)\). Taking the partial trace of the square of
(E823.5), with \(B\) added to the \(H\) factor, gives the exact identity

\[
 \tfrac12\operatorname{Tr}_H(A_{H\otimes F}+B)^2
  =r\,\tfrac12 A_F^2
    +c_1^B(H)\,A_F
    +\operatorname{ch}_2^B(H)\operatorname{id}_F.
 \tag{E823.7}
\]

This is simply the three terms of the tensor square, with the two mixed
terms combined before dividing by two. In particular, a vanishing first
Chern class removes the whole mixed term in Ext, not merely its trace.

\subsection{3. Hochschild evaluation and the rational
B-field}\label{SF-E823-hochschild-evaluation-and-the-rational-b-field}

For \(\pi\in H^0(X,\bigwedge^2T_X)\) and
\(\eta\in H^2(X,\mathcal O_X)\), the characteristic morphism in the
consistent HKR/Atiyah convention is

\[
 \rho_E(\pi+\eta)=
       \pi\mathbin{\lrcorner}\tfrac12 A_E^2
                         +\eta\operatorname{id}_E.
 \tag{E823.8}
\]

The exact characteristic/Atiyah locator is published BF2008,
Theorem5.1.3, author-copy pp24--25, Proposition5.2.2, p26, and
Corollary5.2.3, p27. The complete invoked proofs and the corollary were
freshly read for this result. The sign convention is displayed in
section4.2, p22. The trace/composition identity remains
Proposition6.2.1, Corollary6.3.2 and Theorem6.4.2, author-copy pp29--33,
including the full Proposition6.4.4 proof already read for SF-E819. The
paper is published in Advances in Mathematics217(2008),243--281,
DOI10.1016/j.aim.2007.06.013; the author copy is
https://arxiv.org/pdf/math/0606730. No unrefereed substitute is used.

For clarity, the B-field version means the following. Conjugate the
Clifford operator of \(\pi+\eta\) on Hodge cohomology by multiplication
by \(e^B\). The resulting class \(\xi_B\in HH^2(X)\) has bivector,
mixed, and \(H^2(\mathcal O_X)\) components. Expanding the square gives

\[
 \rho_E(\xi_B)=
       \pi\mathbin{\lrcorner}\tfrac12(A_E+B)^2
                         +\eta\operatorname{id}_E,
 \qquad
 c_{\operatorname{ch}E}(\xi_B)
    =e^{-B}c_{e^B\operatorname{ch}E}(\pi+\eta).
 \tag{E823.9}
\]

This is the same B-shear used in SF-E819. The equality follows by
expanding \(e^{-B}\iota_\pi e^B\): the constant, one-contraction and
zero-contraction terms match \(A_E^2/2\), \(BA_E\) and \(B^2/2\),
respectively. It is a complex cohomological identity with rational B;
the eigenplane bivectors themselves need not be rational. Neither an
integral B-field nor an isogeny is needed.

Combining (E823.7)--(E823.9) yields

\[
 \begin{split}
 \operatorname{Tr}_H\rho_{H\otimes F}(\xi_B)
   ={}&r\rho_F(\pi)
     +\pi\mathbin{\lrcorner}\bigl(c_1^B(H)A_F\bigr)\\
    &+\bigl(\pi\mathbin{\lrcorner}\operatorname{ch}_2^B(H)
                         +r\eta\bigr)\operatorname{id}_F.
 \end{split}
 \tag{E823.10}
\]

\subsection{4. Balanced vanishing would force an impossible global
vanishing}\label{SF-E823-balanced-vanishing-would-force-an-impossible-global-vanishing}

Take \(F=M_0\). Its first two character components vanish, so comparing
degrees one and two in (E823.3) gives

\[
 c_1^B(H)=0,
 \qquad \operatorname{ch}_2^B(H)
       =-\frac{rq}{2}(\theta_+^2+\theta_-^2).
 \tag{E823.11}
\]

Choose the nonzero invariant bivector \(\pi_+\) on the first plane and
define \(\eta_+=\pi_+\mathbin{\lrcorner}(\theta_+^2/2)\). As checked in
SF-E819, \(\pi_++q\eta_+\) annihilates the entire balanced character
\(v_q\). Thus its B-shear \(\xi_B\) lies in
\(\ker c_{\operatorname{ch}(H\otimes M_0)}\).

Suppose the tensor product were weakly semiregular for the specified
finite action. Supplied S21's invariant-evaluation/forgetting
compatibility and the published BF composition identity imply
\(\rho_{H\otimes M_0}(\xi_B)=0\) in ordinary Ext, exactly as in SF-E819.
Apply (E823.10), taking \(\eta=q\eta_+\). Both its second and third
terms vanish by (E823.11), so

\[
 r\rho_{M_0}(\pi_+)=0.
 \tag{E823.12}
\]

Since \(r>0\), this would give \(\rho_{M_0}(\pi_+)=0\). But the ordinary
BF identity says

\[
 \sigma_{M_0}\rho_{M_0}(\pi_+)
       =\pi_+\mathbin{\lrcorner}\operatorname{ch}(M_0).
 \tag{E823.13}
\]

Its component of Hodge type \((1,3)\), using (E823.2), is

\[
 -2\pi_+\mathbin{\lrcorner}D^3
       =-12k^3\eta_+\theta_-\ne0.
 \tag{E823.14}
\]

Indeed \(\theta_+^3=\theta_-^3=0\), so only \(3k^3\theta_+^2\theta_-\)
contributes to this contraction. The other cubic term has only one
holomorphic slot in the plus plane. The product \(\eta_+\theta_-\) is
nonzero because it has full antiholomorphic degree two in the plus plane
and degree \((1,1)\) in the complementary plane. This contradicts
(E823.12), proving the theorem.

\subsection{5. Exact scope and the next genuinely different
object}\label{SF-E823-exact-scope-and-the-next-genuinely-different-object}

The proof excludes every actual tensor product \(H\otimes M_0\)
satisfying (E823.3), including unstable, indecomposable and nonsimple
vector bundles \(H\). It does not require an equivariant decomposition
or a lift of the action to either tensor factor. Thus replacing the
split interpolated hull by a nonsplit hull while keeping this tensor
construction cannot produce the desired seed.

There may be other global extension classes with the same local unit
matrix and the same hull, torsion and defect, which are not isomorphic
to this tensor product. The partial trace (E823.4) requires the actual
global tensor factor; it is not asserted for arbitrary local tensor
identifications. Such genuinely different global gluing, nonuniform
local extension matrices, other slopes, and more general derived objects
remain open. The local nullhomotopies are retained as evidence that this
new obstruction is global, not a hidden reuse of the equal-degree
hypothesis.

\subsection{6. Exact arithmetic check and limits of that
check}\label{SF-E823-exact-arithmetic-check-and-limits-of-that-check}

\texttt{rm\_tensor\_surface\_certificate.py} has actually passed. Its
eight-generator exterior algebra checks the entire compensated character
for symbolic k,q, both balanced cap kernels, and both nonzero cubic
obstruction formulas. It also checks the B-shear signs on all256
exterior basis elements for each of the two bivectors, with the
displayed exact non-diagonal integral B matrix (512 checks). This is a
finite sign check; the general-B identity and the derived partial trace
are proved in the text, not inferred from that matrix. No computation of
Ext groups, a finite action, or semiregularity is claimed.

\section{Unit surface gluing: the exact local algebra, positive Ext, and
actual nontensor
objects}\label{SF-UG-unit-surface-gluing-the-exact-local-algebra-positive-ext-and-actual-nontensor-objects}

Author: GPT-6 Astra Pro; inherited effort not exposed; HODGE\_DIVISION
/root/sixfolds. Date: 2026-11-21.

\textbf{Status:} Complete local algebra, positive Ext and actual
construction independently whole-reviewed by /root and /root/imprimitive
on 21 November 2026; both mathematical reviews PASS. Author and
reviewers: GPT-6 Astra Pro, inherited effort not exposed. Reviewed draft
SHA-256:
\texttt{16b9e88d5c5c5190628272560e4749daa475da6b585cbc80af56e199abeacc77}.
The author and peer independently ran exact certificates. The final
section states exactly which traces these computations provide; the
accompanying SF-E824 avoids the full-derived-trace issue and excludes
all these unit forms. No semiregular seed or finite action is
constructed.

\subsection{1. The local module and its complete degree-zero
endomorphism
algebra}\label{SF-UG-the-local-module-and-its-complete-degree-zero-endomorphism-algebra}

Work over a regular local ring \(R\) with regular parameters \(x,y\)
cutting out the smooth surface; tangential parameters are unrestricted.
Put \(I=(x,y)\), \(A_0=R/I\), \(v=(-y,x)^t\), \(w=(x,y)\), and

\[
 F=R^2/(I v).
 \tag{UG.1}
\]

This is the actual unit module \(M_0\) of the preceding construction,
with resolution

\[
 P:\quad R\xrightarrow{v}R^2\xrightarrow{v w}R^2\longrightarrow F.
 \tag{UG.2}
\]

The hull map is \(p:F\to I\), \(p(e_1)=x,p(e_2)=y\), and the torsion
\(T=A_0v\) has annihilator exactly \(I\).

Every endomorphism induces a scalar \(a\in R\) on the double dual \(R\).
Its action on \(T\) is multiplication by \(a\bmod I\): compatibility of
the extension with its local unit class identifies the torsion scalar
with the induced quotient scalar. Equivalently, apply the endomorphism
to \(v=-ye_1+xe_2\) and use the hull map. After subtracting
\(a\operatorname{id}\), the endomorphism maps into \(T\) and kills
\(T\), so factors through \(I\to T\). Since the Koszul relation is zero
in \(T\), \(\operatorname{Hom}_R(I,T)=A_0^2\).

Concretely, for a row \(b=(b_1,b_2)\) modulo \(I\), its nilpotent
endomorphism is \(\delta_b=v b\) on the two generators. Thus

\[
 \operatorname{End}_R(F)=\mathcal A=R\oplus J,
 \qquad J=A_0^2,\qquad J^2=0,
 \tag{UG.3}
\]

with multiplication \((a,b)(a',b')=(aa',ab'+a'b)\). In particular this
entire degree-zero algebra is commutative, not just its reduction. The
sum is direct because a scalar with zero hull action is zero in the
domain \(R\).

Globally, for \(S\) cut out by two sections of \(\mathcal O_X(D)\),

\[
 \mathcal E nd_X(M_0)=\mathcal O_X\oplus i_*J_S,
 \qquad J_S=\mathcal O_S(-D)^{\oplus2},\qquad J_S^2=0.
 \tag{UG.4}
\]

Indeed \(J_S=\mathcal H om(I_S,\mathcal O_S(-2D))
=N_{S/X}\otimes\mathcal O_S(-2D)\). Scalar multiplication provides the
global splitting, so no choice of a normal frame is required for (UG.4).

\subsection{2. Positive local Ext and both actions of the nilpotent
ideal}\label{SF-UG-positive-local-ext-and-both-actions-of-the-nilpotent-ideal}

Apply \(\operatorname{Hom}_R(P,F)\). Write \(u=(x,y)^t\) and
\(w_F(a,b)=-ya+xb\). The resulting cochain complex is

\[
 F^2\xrightarrow{u w_F}F^2\xrightarrow{w_F}F.
 \tag{UG.5}
\]

Its last cohomology is \(F/IF=A_0^2\). To compute its middle cohomology,
let \((a,b)\in\ker w_F\). Applying the hull map gives
\(-yp(a)+xp(b)=0\). The regular-sequence relation implies
\(p(a)=xc,p(b)=yc\), for a unique \(c\in R\). Hence
\(a=ce_1+t_1,b=ce_2+t_2\), with \(t_1,t_2\in T\), and \(w_F(a,b)=cv\).
The latter is zero exactly when \(c\in I\). Choose \(m\in F\) with
\(p(m)=c\). Subtracting \(u m\) now leaves an element of \(T^2\). Also
\(uF\cap T^2=0\): if \(u m\in T^2\), then \(xp(m)=yp(m)=0\), so
\(p(m)=0\), and \(m\in T\) gives \(u m=0\). Therefore

\[
 \ker w_F=uF\oplus T^2.
\]

The incoming image is \(u(IF)\). Since \(\ker(u:F\to F^2)=T\) and
\(T\subset IF\), we obtain

\[
 \operatorname{Ext}^1_R(F,F)
       =(F/IF)\oplus T^2=A_0^4,
 \qquad
 \operatorname{Ext}^2_R(F,F)=F/IF=A_0^2.
 \tag{UG.6}
\]

In particular \(I\) annihilates both groups.

Choose any row lift of \(b\), and put \(t_b=bv=-b_1y+b_2x\). A chain
lift of \(\delta_b\) on (UG.2), ordered by degrees \(0,-1,-2\), is

\[
 \bigl(vb,\ t_b I_2,\ t_b\bigr).
 \tag{UG.7}
\]

The two chain identities follow immediately from \(wv=0\). Consequently
the right action of \(J\) on both positive Ext groups is zero: it is
precomposition by \(t_b I\), and \(t_b\in I\). The left action is also
zero. On the last quotient, \(\delta_bF\subset T\subset IF\). On the
middle representatives, it kills \(T^2\), and it sends \(uF\) to
\(uT=0\). This checks both actions, rather than inferring centrality
from commutativity in degree zero.

Restoring the global twists in the resolution gives

\[
 \mathcal E xt^1_X(M_0,M_0)=i_*\mathcal O_S(D)^{\oplus4},
 \qquad
 \mathcal E xt^2_X(M_0,M_0)=i_*\mathcal O_S(3D)^{\oplus2}.
 \tag{UG.8}
\]

The first splitting comes from \(uM_0(2D)\) and
\(\mathcal O_S(-2D)(3D)^2\) in the kernel calculation. The second uses
\((M_0/IM_0)(4D)\), not merely \(M_0/IM_0\). These twists are essential.

\subsection{3. Chain multiplication is not strictly the square-zero
algebra}\label{SF-UG-chain-multiplication-is-not-strictly-the-square-zero-algebra}

Although \(J^2=0\) on \(F\), the chain lifts (UG.7) do not multiply
strictly to zero. There is an explicit nullhomotopy. Put
\(j_b=(b_2,-b_1)^t\). For rows \(b,c\), the product of their chain lifts
is nullhomotopic by

\[
 h^0=j_b c:P^0\to P^{-1},
 \qquad h^{-1}=t_c b:P^{-1}\to P^{-2}.
 \tag{UG.9}
\]

At degree zero, \((vw)j_bc=t_bvc\). At degree minus one, the identity
\(t_bI_2-j_bw=vb\) gives \(v(t_cb)+j_bc(vw)=t_bt_cI_2\). At degree minus
two, \((t_cb)v=t_bt_c\). Thus \(dh+hd\) is precisely the product of the
two chain lifts. This exhibits why a strict trace on degree-zero
matrices does not immediately give a trace on the entire derived descent
data. No dg formality assertion is made.

\subsection{4. Actual nontensor gluing and its classification as a
marked
extension}\label{SF-UG-actual-nontensor-gluing-and-its-classification-as-a-marked-extension}

Let \(H\) be any vector bundle. Consider actual coherent extensions

\[
 0\to i_*(H|_S(-2D))\to M\to H\otimes I_S\to0
 \tag{UG.10}
\]

whose local extension matrix is the identity under the specified
identification of the two rank-\(r\) bundles. Applying Hom to the actual
Koszul resolution of \(H I_S\) again gives a complex with zero
differential. Hence

\[
 \operatorname{Ext}^1_X(H I_S,i_*H|_S(-2D))
  =H^1\bigl(S,\operatorname{End}(H|_S)(-D)\bigr)^{\oplus2}
       \oplus H^0\bigl(S,\operatorname{End}(H|_S)\bigr).
 \tag{UG.11}
\]

Fixing the last component to \(I_r\) leaves the entire first vector
space as actual global extension classes. The zero vector gives
\(H\otimes M_0\).

There is also a direct gluing description, including its absence of a
higher associativity obstruction. Let
\(Y=(X,\mathcal O_X\oplus i_*J_S)\), the split square-zero thickening
defined by (UG.4). Locally (UG.10) is \(M_0^{\oplus r}\), and its
transition automorphisms lie in \(\mathrm{GL}_r(\mathcal A)\). Reduction
modulo \(J_S\) gives the transition functions of \(H\). Starting with
those functions, multiply by \(1+a_{ij}\), where \(a_{ij}\) is a Čech
cocycle in \(\operatorname{End}(H|_S)\otimes J_S\). The cocycle equation
is the ordinary additive one because \(J_S^2=0\). Thus these are true
transition functions, not projective transition data. They define a
vector bundle \(\widetilde H\) on \(Y\), and the actual sheaf
\(M=M_0\otimes_{\mathcal A}\widetilde H\). All constructions can be
checked on a finite affine cover, where evaluation identifies the local
module with \(M_0^r\).

A nonzero cohomology class in the first summand of (UG.11) gives an
object that is not a tensor product with \(M_0\), even allowing an
isomorphism of the hull. Indeed any isomorphism to a tensor product
induces an automorphism \(h\) of the canonical hull \(H\); compatibility
with the unit local class forces the torsion automorphism to be
\(h|_S\). The zero extension parameter is preserved by this pair of
automorphisms, because the tensor construction is functorial in \(H\).
The induced conjugation on the first summand of (UG.11) is invertible,
and cannot send a nonzero class to zero. Changing local splittings adds
only a coboundary. A different tensor hull would have to be isomorphic
to \(H\) by taking double duals, so gives no additional possibility.

\subsection{5. A nonzero parameter in the explicit balanced high-rank
family}\label{SF-UG-a-nonzero-parameter-in-the-explicit-balanced-high-rank-family}

Use the sixteen actual line-bundle classes of the compensation
construction, with the same positive multiplicities. Among their
differences is

\[
 \delta=(8,-6,-8,6)
 =(-8,-6,-8,6)-(-16,0,0,0).
\]

Its four Hermitian eigenvalues, in the previous ordering, are
\(0,16-12\sqrt2,0,16+12\sqrt2\). Put \(D=3\theta\). For \(a=1,2,3\), the
divisor \(\delta-aD\) is nondegenerate with exactly three negative
eigenvalues. The published Mumford index theorem therefore gives
cohomology only in degree three. Exact already checked locator: Mumford,
Abelian Varieties (1970), section16, Vanishing Theorem p150 and its
proof pp150--153, with the complex Hermitian-index corollary and proof
pp162--163. Riemann--Roch on the principal product gives

\[
 h^3(X,L_\delta(-aD))
  =(3a)^2\bigl(288-(16-3a)^2\bigr)
   =\begin{cases}1071&a=1,\\6768&a=2,\\19359&a=3.\end{cases}
 \tag{UG.12}
\]

The restriction Koszul resolution is

\[
 0\to L_\delta(-3D)\to L_\delta(-2D)^2
    \to L_\delta(-D)\to i_*(L_\delta(-D)|_S)\to0.
\]

Only its cohomology row in degree three survives. Thus
\(H^1(S,L_\delta(-D)|_S)\) is the kernel of the first map in that row,
and

\[
 h^1(S,L_\delta(-D)|_S)\ge19359-2\cdot6768=5823>0.
 \tag{UG.13}
\]

This is a dimension inequality valid for every chosen smooth complete
intersection here; no rank claim about the multiplication map is needed.
It proves the existence of a nonzero gluing parameter in (UG.11) for the
explicit compensating hull. The resulting \(M\) is actual coherent,
perfect and ordinarily gluable; its character is unchanged by the
extension parameter and is exactly \(522387689472\,v_{1/2}\). It is
genuinely nontensor by section4. Neither indecomposability, a finite
action nor weak semiregularity is asserted.

\subsection{6. What descends, and why a full derived trace is not
assumed}\label{SF-UG-what-descends-and-why-a-full-derived-trace-is-not-assumed}

Since \(J_S\) acts by zero on both sides of the positive Ext groups,
conjugation by a transition matrix in \(\mathrm{GL}_r(\mathcal A)\) acts
on those groups only through its reduction in
\(\mathrm{GL}_r(\mathcal O_X)\). Consequently the positive self-Ext
sheaves of every (UG.10) are

\[
 \mathcal E xt^q_X(M,M)
 =\operatorname{End}(H|_S)\otimes
       \mathcal E xt^q_X(M_0,M_0),\quad q=1,2.
 \tag{UG.14}
\]

Their ordinary matrix traces therefore give global maps to the
corresponding self-Ext sheaves of \(M_0\). In degree zero, commutativity
of \(\mathcal A\) gives the ordinary bundle trace
\(\mathcal E nd_X(M)=\mathcal E nd_{\mathcal A}(\widetilde H)\to\mathcal A\).

Furthermore \(H^1(S,J_S)=0\). Indeed the restriction resolution of
\(\mathcal O_S(-D)\) has three negative ample line bundles in degrees
\(-2,-1,0\); their cohomology is only in degree four by Mumford's index
theorem, so its first possible hypercohomology degree is two. Thus the
trace of every gluing parameter in (UG.11) is zero in cohomology, and
the exact unit sequence
\(1+J_S\to\mathcal O_Y^\times\to\mathcal O_X^\times\) shows
\(\operatorname{Pic}(Y)\to\operatorname{Pic}(X)\) is injective. The
scalar splitting makes it surjective. In particular \(\det\widetilde H\)
is the pullback of \(\det H\).

These facts do \textbf{not} establish a partial trace on global derived
Ext. One must still check compatibility with the local-to-global
spectral-sequence differentials and the actual Atiyah/characteristic
cocycle. The chain lifts (UG.7) multiply only up to the explicit
homotopies (UG.9); positive cohomology centrality alone does not supply
all their higher compatibilities. The ordinary partial trace of SF-E823
uses an actual \(\mathcal O_X\)-locally free tensor factor and cannot
simply be imported to a bundle over \(Y\). The adjoining SF-E824 avoids
this full-derived-descent issue by proving exactly the trace identity
for the first overlap obstruction. That separate argument excludes all
the global unit forms constructed here. Neither this local algebra
calculation nor that first-overlap argument assumes dg formality or an
asserted deformation of the hull. Other local models remain outside the
stated theorem.

\section{SF-E824: the first overlap obstruction excludes every global
unit
form}\label{SF-E824-sf-e824-the-first-overlap-obstruction-excludes-every-global-unit-form}

Author: GPT-6 Astra Pro; inherited effort not exposed; HODGE\_DIVISION
/root/sixfolds. Date: 2026-11-21.

\textbf{Status:} Complete proof independently whole-reviewed by /root
and /root/imprimitive on 21 November 2026; both mathematical/source
reviews PASS. Author and reviewers: GPT-6 Astra Pro, inherited effort
not exposed. Reviewed draft SHA-256:
\texttt{f42a0232fbbcfbd21a1aa66516888b3d142658d7b3d1aff16c27df7d3f8c1ffe}.
This uses a trace only on the first overlap obstruction, not an unproved
trace on global derived Ext. The complete local algebra and actual
nontensor family are proved in the adjoining unit-gluing analysis,
independently reviewed in full by the same two reviewers.

\subsection{1. Statement}\label{SF-E824-statement}

Let \(X\) be a complex abelian fourfold with the balanced RM data of
SF-E823. Let \(D=k\theta\) be an integral ample divisor, \(k>0\), and
let \(S\) be a smooth complete intersection of two divisors in \(|D|\).
Let \(F=M_0\) be the actual unit extension

\[
 0\to\mathcal O_S(-2D)\to F\to I_S\to0,
 \qquad \operatorname{ch}(F)=1-2D^3+4D^4.
 \tag{E824.1}
\]

Suppose \(M\) is a coherent sheaf locally isomorphic to
\(F^{\oplus r}\), with \(r>0\), and

\[
 e^B\operatorname{ch}(M)=r v_q,
 \qquad B\in\operatorname{NS}(X)_{\mathbf Q},\quad q>0.
 \tag{E824.2}
\]

Then \(M\) is not weakly semiregular for any specified effective
coherent finite translation/tensor action of supplied S21. This covers
all global unit forms, including the actual nontensor extensions
constructed in the adjoining analysis. It assumes no global tensor
factor, simplicity, stability, or equivariant decomposition.

The local algebra already proved there is

\[
 \mathcal A=\mathcal E nd(F)=\mathcal O_X\oplus i_*J,
 \quad J=\mathcal O_S(-D)^2,\quad J^2=0,
 \tag{E824.3}
\]

and \(\mathcal B=\mathcal E xt^1(F,F)=i_*\mathcal O_S(D)^4\). The two
actions of \(\mathcal A\) on \(\mathcal B\) agree and factor through
\(\mathcal O_S\); in particular \(J\) acts by zero on both sides. Also
\(\mathcal E xt^j(F,F)=0\) for \(j>2\).

\subsection{2. The first obstruction class and its exact filtration
meaning}\label{SF-E824-the-first-obstruction-class-and-its-exact-filtration-meaning}

Fix an invariant holomorphic bivector \(\pi\). The first-order sheaf of
algebras

\[
 \mathcal O_{X,\pi}=\mathcal O_X\oplus\varepsilon\mathcal O_X,
 \qquad a*b=ab+\tfrac{\varepsilon}{2}\pi(da,db),
 \quad\varepsilon^2=0,
 \tag{E824.4}
\]

is defined globally. It is associative to this order by the biderivation
identity. The explicit local homotopies of the compensation construction
give flat local module deformations of \(F\), and hence of \(M\).

For \(N=F\) or a local form of \(F^{\oplus r}\), set
\(C_N=R\mathcal H om(N,N)\). Its local characteristic class is zero, so
\(\rho_N(\pi)\in\operatorname{Ext}^2(N,N)\) belongs to the first
filtration piece \(F^1\), the kernel of its map to
\(H^0(\mathcal E xt^2(N,N))\). The truncation triangles

\[
 \tau_{\le1}C_N\to C_N\to\mathcal E xt^2(N,N)[-2],
 \qquad
 \mathcal E nd(N)\to\tau_{\le1}C_N
                   \to\mathcal E xt^1(N,N)[-1]
 \tag{E824.5}
\]

give a canonical map

\[
 F^1\longrightarrow H^1(\mathcal E xt^1(N,N)),
 \tag{E824.6}
\]

whose kernel is the image of \(H^2(\mathcal E nd(N))\). Indeed the first
triangle identifies \(H^2(\tau_{\le1}C_N)\) with \(F^1\), since the
preceding group is \(H^{-1}(\mathcal E xt^2)=0\); the assertion then
follows from the second triangle. Write \(\kappa_N(\pi)\) for the image
of \(\rho_N(\pi)\) under (E824.6).

This class has a direct overlap interpretation. Choose local
deformations \(N_{i,\varepsilon}\) on an affine or sufficiently small
analytic cover. On each intersection, the obstruction to lifting the
original transition isomorphism is an element of the local
\(\operatorname{Ext}^1\) sheaf. The resulting Čech one-cocycle
represents \(\kappa_N(\pi)\). To check this identification, represent
\(\rho_N(\pi)\) in a Čech resolution of the self-Hom complex. The local
nullhomotopies remove its Čech-degree-zero, internal-degree-two
component. The remaining degree-(1,1) component is the difference of
those nullhomotopies after transition; its class is exactly the
morphism-lifting obstruction. The total cocycle equation makes it a Čech
cocycle modulo internal boundaries. This is precisely the map from the
two truncation triangles, and uses no degeneration of a spectral
sequence.

\subsection{3. The unit module has a nonzero first
obstruction}\label{SF-E824-the-unit-module-has-a-nonzero-first-obstruction}

Let \(\pi=\pi_+\), and put
\(\eta_+=\pi_+\mathbin{\lrcorner}(\theta_+^2/2)\). The degree-one
component of the full semiregularity morphism is a derived map

\[
 \sigma_F^1:C_F\longrightarrow\Omega_X^1[1],
\]

inducing \(\operatorname{Ext}^2(F,F)\to H^3(\Omega_X^1)\). Its
restriction along \(\mathcal E nd(F)\to C_F\) is zero. On the scalar
summand \(\mathcal O_X\) of (E824.3), that restriction is
\(\operatorname{ch}_1(F)=0\). On \(i_*J\), the target morphism group is
\(\operatorname{Ext}^1_X(i_*J,\Omega_X^1)=0\): the two-element Koszul
resolution of the regular immersion gives local Hom and local Ext-one
equal to zero against the locally free target, and thus the global group
is also zero. This last assertion follows directly from the low-degree
local-to-global Ext sequence; neither possible total-degree-one term
survives.

It follows that \(\sigma_F^1\) kills the image of
\(H^2(\mathcal E nd(F))\). But published BF gives

\[
 \sigma_F^1\rho_F(\pi_+)
    =\pi_+\mathbin{\lrcorner}\operatorname{ch}_3(F)
    =-12k^3\eta_+\theta_-\ne0,
 \tag{E824.7}
\]

with exactly the normalization checked in SF-E823. Since
\(\rho_F(\pi_+)\) is locally zero, (E824.5)--(E824.6) now show

\[
 \kappa_F(\pi_+)\ne0
       \quad\hbox{in }H^1(X,\mathcal B).
 \tag{E824.8}
\]

This proves nonvanishing of the specific first overlap class, not just
of an unspecified global Ext class.

\subsection{4. The endomorphism-lifting
derivation}\label{SF-E824-the-endomorphism-lifting-derivation}

We give the exact first-order algebra needed for the trace. Work on a
small open set and pass to its affine section modules, or work on
stalks. Choose one local deformation of \(F\), and choose a splitting of
its reduction as a complex vector space; flatness identifies the
underlying dual-number module with \(F\oplus\varepsilon F\). No global
or \(\mathcal O_X\)-linear splitting is being assumed. Write its action
as \(\rho_0+\varepsilon\rho_1\). For \(a\in\mathcal A\), the obstruction
to lifting the endomorphism \(a\) is represented in the bar Ext-one
complex by

\[
 r\longmapsto a\rho_1(r)-\rho_1(r)a.
 \tag{E824.9}
\]

It is a one-cocycle because the two associativity equations have the
same fixed multiplication perturbation (E824.4), which commutes with the
original \(R\)-linear endomorphism. Changing a lift of the endomorphism
adds a bar boundary. Changing the module deformation changes (E824.9) by
\(a\beta-\beta a\), where \(\beta\in\mathcal B\) is the difference class
of the two deformations. This change is zero in \(\mathcal B\), by its
central \(\mathcal A\)-module structure.

Consequently (E824.9) defines a \textbf{global}, choice-independent
derivation of sheaves of complex vector spaces

\[
 \partial_\pi:\mathcal A\longrightarrow\mathcal B,
 \qquad
 \partial_\pi(ab)=a\partial_\pi(b)+\partial_\pi(a)b.
 \tag{E824.10}
\]

The Leibniz equation follows by expanding the commutator with \(ab\).
This derivation need not be \(\mathcal O_X\)-linear. It kills complex
constants, which is the property used below. Formula (E824.9) may be
made using any linear splitting; the boundary and difference-class
calculations show that its Ext class is independent of that choice.

\subsection{5. The normalized matrix overlap
formula}\label{SF-E824-the-normalized-matrix-overlap-formula}

Choose local identifications \(M|_{U_i}=F|_{U_i}^{\oplus r}\), with
transition matrices \(G_{ij}\in\mathrm{GL}_r(\mathcal A)\). Their
reductions define the locally free hull \(H=M^{\vee\vee}\). Choose local
deformations \(F_{i,\varepsilon}\), and use their direct sums for \(M\).
Let \(k_{ij}\in\mathcal B\) be the obstruction to lifting the identity
between the two local deformations of \(F\), with the sign chosen as
\(\rho_{1,i}-\rho_{1,j}\). Use this same Čech sign in (E824.6);
reversing all overlap signs changes neither vanishing nor the trace
argument.

The obstruction to lifting \(G_{ij}\), normalized back into the source
endomorphism group by multiplication by \(G_{ij}^{-1}\), is

\[
 k_{ij} I_r+G_{ij}^{-1}\partial_\pi(G_{ij}).
 \tag{E824.11}
\]

Indeed its bar cochain is \(G\rho_{1,i}-\rho_{1,j}G
=(G\rho_{1,i}-\rho_{1,i}G)+(\rho_{1,i}-\rho_{1,j})G\). The centrality of
\(\mathcal B\) makes conjugation of the second term by \(G\) trivial.
This proves (E824.11) as an Ext-one identity; it is not an assertion
that the matrices lift or satisfy higher descent equations.

Since \(\mathcal A\) is commutative and \(\mathcal B\) is a central
module, the usual polynomial determinant identity is valid for the
derivation (E824.10):

\[
 \operatorname{tr}(G^{-1}\partial_\pi G)
       =\det(G)^{-1}\partial_\pi(\det G).
 \tag{E824.12}
\]

For completeness, differentiate the permutation expansion of \(\det G\).
Each differentiated entry has its signed cofactor; their sum is
\(\operatorname{tr}(\operatorname{adj}(G)\partial_\pi G)\). Multiply by
the inverse determinant. Centrality permits the required reordering of
the undifferentiated entries with the module-valued derivative.

The positive self-Ext sheaf of \(M\) is
\(\operatorname{End}(H|_S)\otimes\mathcal B\), as proved in the
adjoining analysis. Hence its ordinary matrix trace is a globally
defined sheaf map. Taking Čech cohomology of (E824.11)--(E824.12) gives

\[
 \operatorname{tr}\kappa_M(\pi)
     =r\kappa_F(\pi)
       +[\det(G_{ij})^{-1}\partial_\pi(\det G_{ij})]
       \quad\hbox{in }H^1(X,\mathcal B).
 \tag{E824.13}
\]

This is the only trace descent required in the proof.

\subsection{6. The determinant correction vanishes when the hull has
zero first Chern
class}\label{SF-E824-the-determinant-correction-vanishes-when-the-hull-has-zero-first-chern-class}

Let \(Y=(X,\mathcal A)\). The adjoining analysis proves \(H^1(S,J)=0\)
by the negative-ample Koszul resolution and the published Mumford index
theorem. The split unit sequence
\(1+J\to\mathcal A^\times\to\mathcal O_X^\times\), with \(1+J\simeq J\),
therefore gives \(\operatorname{Pic}(Y)\simeq\operatorname{Pic}(X)\).
The determinant of the \(\mathcal A\)-bundle with transition matrices
\(G_{ij}\) is accordingly the pullback of \(\det H\).

If \(c_1(H)=0\), the line bundle \(\det H\in\operatorname{Pic}^0(X)\)
has a flat unitary description, so has locally constant complex
transition functions on a sufficiently fine analytic cover. Here is the
elementary Hodge justification: the exponential sequence identifies
\(\operatorname{Pic}^0\) with \(H^1(\mathcal O_X)/H^1(X,\mathbf Z)\);
projection of purely imaginary real harmonic one-forms onto type
\((0,1)\) is an isomorphism of real vector spaces. A representative is
therefore a flat unitary connection on the corresponding line bundle.
The integral cohomology of an abelian variety is torsion-free, so zero
rational first Chern class gives exactly this \(\operatorname{Pic}^0\)
case.

After changing local determinant frames in \(\mathcal A^\times\), the
cocycle \(\det G_{ij}\) is thus represented by constants. The derivation
\(\partial_\pi\) kills them. A change of frame changes the logarithmic
derivative by a Čech coboundary, by the ordinary Leibniz rule.
Consequently

\[
 [\det(G_{ij})^{-1}\partial_\pi(\det G_{ij})]=0,
 \qquad
 \operatorname{tr}\kappa_M(\pi)=r\kappa_F(\pi).
 \tag{E824.14}
\]

\subsection{7. Contradiction for B zero}\label{SF-E824-contradiction-for-b-zero}

Suppose first that \(B=0\) in (E824.2). The unit form has the same rank
and first Chern class as its hull \(H\): in its hull sequence the
torsion and defect are supported in codimension two, so their first
Chern characters vanish. Thus \(c_1(H)=c_1(M)=0\), and (E824.14)
applies.

The balanced cap kernel contains \(\xi_+=\pi_++q\eta_+\). Weak
semiregularity, published BF composition, and supplied S21
invariant/forgetting compatibility give \(\rho_M(\xi_+)=0\) in ordinary
Ext. The second term \(q\eta_+\operatorname{id}_M\) comes from
\(H^2(\mathcal E nd(M))\), so its first overlap class is zero by
(E824.5)--(E824.6). Therefore \(\kappa_M(\pi_+)=0\). Taking its trace
and using (E824.14) contradicts \(r\kappa_F(\pi_+)\ne0\) from (E824.8).

\subsection{8. Rational B without assuming semiregularity survives an
isogeny}\label{SF-E824-rational-b-without-assuming-semiregularity-survives-an-isogeny}

For general rational algebraic \(B\), choose a positive integer \(m\)
such that \([m]^*B=m^2B\) is the first Chern class of a line bundle
\(L_B\). Write \(f=[m]\), \(X'=X\), and

\[
 M'=f^*M\otimes L_B,
 \qquad F'=f^*F,
 \qquad D'=f^*D,\quad\theta'=f^*\theta.
 \tag{E824.15}
\]

The map is finite étale. The pullback surface is smooth, and its
negative-ample Koszul resolution gives
\(H^0(\mathcal O_{S'})=\mathbf C\), so it is connected and hence
integral. The pullback unit extension is exactly the same unit
construction for \(D'\); tensoring by a line bundle preserves the
property of being locally a direct sum of \(F'\). Its new hull is
\(H'=f^*H\otimes L_B\), with \(c_1(H')=0\), and
\(\operatorname{ch}(M')=r v_q(\theta'_+,\theta'_-)\).

It remains to justify the vanishing needed for section7, rather than
asserting any preservation of semiregularity. Use the étale pullback
bivector \(\pi'_+=m^{-2}\pi_+\); since \(\theta'_\pm=m^2\theta_\pm\),

\[
 \eta'_+=\pi'_+\mathbin{\lrcorner}(\theta_+'{}^2/2)
         =m^2\eta_+=f^*\eta_+.
 \tag{E824.16}
\]

Atiyah classes pull back under an étale map, and tensoring by \(L_B\)
adds \(f^*B\operatorname{id}\). The exact B-shear evaluation formula of
SF-E823 consequently identifies \(\rho_{M'}(\pi'_++q\eta'_+)\) with the
pullback of \(\rho_M(\xi_B)\), under the line-tensor identification of
self-Ext. The latter is zero by the original weak semiregularity
assumption and the balanced kernel. Thus the ordinary vanishing used in
section7 holds on \(X'\), and the same contradiction follows there. No
finite action is transported and no injectivity of a pulled-back
semiregularity map is assumed.

\subsection{9. Scope}\label{SF-E824-scope}

The argument excludes every global form locally isomorphic to the
full-rank unit module \(M_0^r\) on a smooth ample two-divisor complete
intersection. It therefore closes the non-tensor gluing escape for this
particular local model. It does not assert the same result for an added
locally free summand, different torsion or quotient ranks, an extension
matrix with zeros, singular or nonreduced support, different degree
twists, or arbitrary perfect complexes. The genuine remaining seed gate
is preserved in those other scopes.

The only derived filtration used is explicitly (E824.5). Matrix trace on
the first positive Ext sheaf, the endomorphism-lifting derivation and
the determinant identity suffice. A trace on the entire derived
endomorphism complex is neither constructed nor needed.

\subsection{10. Exact source and check
scope}\label{SF-E824-exact-source-and-check-scope}

The ordinary characteristic/trace identities are the directly published
Buchweitz--Flenner2008 inputs located in SF-E823: author-copy
Theorem5.1.3, Proposition5.2.2 and Corollary5.2.3, pp24--27, and
Proposition6.2.1, Corollary6.3.2, Theorem6.4.2 with Proposition6.4.4,
pp29--33. The complete invoked proofs were read in the preceding checks;
no fresh whole-paper audit is asserted here. The first-overlap,
derivation, determinant and truncation arguments are proved completely
above. Supplied DEG\_L S21 equation(S21.6), the paragraph
following(S21.7), and the invariant-evaluation paragraph before (S21.8)
are used only for coherent invariant/forgetting compatibility. The
negative-ample vanishing and explicit parameter construction use the
published Mumford section16 source exactly located in the adjoining
analysis. No new unrefereed theorem or all-order deformation assertion
is used.

\section{SF-UGDC. A nontrivial derived conjugation invisible on all
self-Ext
sheaves}\label{sf-ugdc.-a-nontrivial-derived-conjugation-invisible-on-all-self-ext-sheaves}

Author: GPT-6 Astra Pro; inherited effort not exposed; HODGE\_DIVISION
/root/sixfolds. Date: 2026-11-21.

\textbf{Status:} complete proof independently whole-read by the root
reviewer; PASS. The author exact certificate and the root actual rerun
at19:12UTC passed. It is not a premise of SF-E824, whose proof uses only
the first overlap obstruction. It refutes the proposed inference from
centrality on self-Ext sheaves to an untwisted trace on the full derived
self-Hom complex.

Let \(R=\mathbf C[[x,y,z,w]]\), let \(v=(-y,x)^t\), \(w_0=(x,y)\), and
let \(P\) be the minimal free resolution

\[
 P^{-2}=R\xrightarrow{v}P^{-1}=R^2
             \xrightarrow{v w_0}P^0=R^2
\]

of the unit module \(F\). For \(b=(1,0)\), the square-zero module
endomorphism \(\delta_b\) has the chain lift

\[
 \delta^0=\begin{pmatrix}-y&0\\x&0\end{pmatrix},
 \qquad \delta^{-1}=-yI_2,
 \qquad\delta^{-2}=-y.
 \tag{DC.1}
\]

The chain map \(u=1+\delta\) is invertible: the determinant of its
degree-zero block is \(1-y\), and its remaining blocks are scalar
multiplication by \(1-y\). These are units in \(R\). It lifts the actual
automorphism \(1+\delta_b\) of \(F\), whose inverse on the module is
\(1-\delta_b\).

Put \(C=\mathcal H om_R(P,P)\), with differential
\(d_C(a)=d_Pa-(-1)^{\deg a}ad_P\). Conjugation by \(u\) is an actual
chain automorphism of \(C\). It induces the identity on every cohomology
module. In degree zero this follows from commutativity of
\(\operatorname{End}_R(F)=R\oplus(R/(x,y))^2\); in positive degrees both
left and right nilpotent actions are zero, as proved in the unit-gluing
analysis. Negative Ext is zero because \(F\) is a sheaf.

Nevertheless conjugation by \(u\) is \textbf{not} the identity in
\(D(R)\). Let \(\mathfrak m=(x,y,z,w)\), and consider terms of order one
modulo \(\mathfrak m^2\). The linear part of \(d_P\) is only
\(P^{-2}\xrightarrow{v}P^{-1}\); its map to \(P^0\) is quadratic.
Consequently the linear part \(d_C^{(1)}\) has no arrow into or out of
the block

\[
 V=\operatorname{End}_R(P^0)\subset C^0.
 \tag{DC.2}
\]

For an arbitrary degree-minus-one homotopy \(K\) with power-series
coefficients, the order-one part of \(d_CK+Kd_C\) is
\(d_C^{(1)}K^{(0)}+K^{(0)}d_C^{(1)}\). Its component from \(V\) to \(V\)
is therefore zero. This does not assume that \(K\) preserves the
displayed blocks.

On the other hand, the linear part of \(uau^{-1}-a\) on \(V\) is
\([\delta^0,a]\). At \(a=\operatorname{diag}(1,0)\) this equals

\[
 [\delta^0,a]=\begin{pmatrix}0&0\\x&0\end{pmatrix}\ne0.
 \tag{DC.3}
\]

Thus no such homotopy exists. The bounded complex \(C\) is finite free,
so a null morphism in \(D(R)\) would be nullhomotopic. This proves the
claim.

In particular there cannot be a trace to the fixed complex \(C\) that is
invariant under all these inner changes of local trivialization and is
identity-normalized at rank one. At rank one the required invariance
would make the nontrivial conjugation equal to the identity. The same
contradiction holds if rank-one normalization is any quasiisomorphism to
\(C\), because it is invertible in the derived category.

This conclusion does not rule out a target with its own nontrivial
descent, a cyclic quotient, or a weaker trace on a filtration piece. It
also does not claim that a chosen global gluing has no noncanonical
trace. SF-E824 uses exactly the weaker first-overlap trace, whose
determinant calculation is proved independently. The local
countercalculation is retained to prevent a future replacement of that
proof by a false full-trace inference.

\subsection{Check scope}\label{check-scope}

Root reviewer: GPT-6 Astra Pro; inherited effort not exposed;
HODGE\_DIVISION /root. The entire proof and certificate were read, and
the certificate was actually rerun at2026-11-21 19:12UTC with reported
wall time0.338621487seconds. The author read the complete proof and ran
the exact25-dimensional graded-endomorphism certificate. An initial
purely structural polynomial-equality check failed and was corrected to
compare expanded polynomial differences; that failed attempt remains in
rm\_unit\_derived\_trace\_certificate\_attempts.json. No proof-assistant
verification is claimed. The local End/Ext facts used above are proved
in full in the adjoining unit-gluing analysis; no new outside source is
needed.

\section{SF-E825. The surviving first-overlap obstruction after adding
free
summands}\label{SF-E825:sf-e825.-the-surviving-first-overlap-obstruction-after-adding-free-summands}

Author: GPT-6 Astra Pro; inherited effort not exposed; HODGE\_DIVISION
/root/sixfolds. Date: 2026-11-21.

\textbf{Status:} independently whole-reviewed by /root and
/root/g34\_eightfolds; both mathematical reviews PASS on the frozen
draft
\texttt{34af017d1f0c49a56b9924c68b923052afc003ba953e62e1ff55f6d811079783}.
The author and root actually ran the local certificate; the peer
supplied an independently written exact certificate. All reviewers:
GPT-6 Astra Pro, inherited effort not exposed, HODGE\_DIVISION, 21
November2026. This is a necessary determinant condition, not a universal
exclusion or a construction of a semiregular seed. The free-summand case
is not covered merely by citing E824; the off-diagonal products below
are essential.

\subsection{1. Statement and notation}\label{SF-E825:statement-and-notation}

Let X be the polarized abelian fourfold in the balanced RM setup, with
\(\theta=\theta_++\theta_-\), the two complex rank-two RM planes, and
\(v_q=(1-q\theta_+^2/2)(1-q\theta_-^2/2)\). Let \(L=\mathcal O_X(D)\),
\(D=k\theta\) ample and integral, and let the ordered global sections
f,g of L cut out a smooth integral surface S. Write \(I=I_S\),
\(S_2=(X,\mathcal O_X/I^2)\), and let F be the actual unit sheaf M0
already constructed, with local model

\[
 F=R^2/(Iv),\qquad v=(-g,f)^t,\quad w=(f,g),\quad
 P:\ R\xrightarrow vR^2\xrightarrow{vw}R^2.
 \tag{E825.1}
\]

Suppose M is a coherent sheaf locally isomorphic to
\(F^{\oplus r}\oplus\mathcal O_X^{\oplus s}\), where r is positive and s
is nonnegative. Its hull H is locally free of rank r+s. Its canonical
quotient \(H/(M/\operatorname{tors}M)=i_*Q\) is locally free of rank r
on S. There is a canonically associated vector bundle \(\widetilde Q\)
on S2 restricting to Q. This bundle is constructed explicitly below; no
extension of an arbitrary bundle on S to S2 is assumed.

For an invariant bivector pi define a homomorphism

\[
 \Psi_\pi:\operatorname{Pic}(S_2)\longrightarrow
 H^1(S,\mathcal O_S(D)^{\oplus2})
 \tag{E825.2}
\]

as follows. If ell\_ij are line transition functions on S2, lift them
locally to units near S. In the i-th L-frame the two components of a
representative are \(\{\ell_{ij},f_i\}/\ell_{ij}\) and
\(\{\ell_{ij},g_i\}/\ell_{ij}\), restricted to S. Here
\(\{a,b\}=\pi(da,db)\). This is independent of the lifts because
\(\{I^2,I\}\subset I\). It is a Cech cocycle because the multiplicative
line cocycle holds modulo I squared. Under a normal-frame change
\(f_j=t_{ij}f_i,g_j=t_{ij}g_i\), the extra term has a factor f\_i or
g\_i and vanishes on S. Changing line frames changes this by a Cech
coboundary. The logarithmic product rule proves that (E825.2) is a
homomorphism. We extend it rationally when needed. A class from X is
restricted to S2 first.

If the entire normalized character satisfies
\(e^B\operatorname{ch}(M)=(r+s)v_q\), with B algebraic rational, and M
is weakly G-semi-regular for any specified effective coherent finite
action in supplied S21's scope, then necessarily

\[
 \boxed{\quad
 \Psi_{\pi_\pm}\bigl(\det\widetilde Q+rB\bigr)
       =\frac{7r}{2}\,\Psi_{\pi_\pm}(L).
 \quad}
 \tag{E825.3}
\]

Here pi\_+ and pi\_- are the two invariant bivectors dual to their
rank-two RM planes. In particular, if the rational determinant class
\(\det\widetilde Q+rB\) is the restriction of a flat line class on X,
then M cannot be weakly G-semi-regular. No finite action on Q, its
extension, or a chosen local direct summand is needed.

If det Qtilde extends to a rational line class A on X, (E825.3) is
equivalent to saying that

\[
 C=c_1(A)+rB-\frac{7r}{2}D
 \tag{E825.4}
\]

has zero mixed blocks and is primitive separately in both RM planes.
Consequently \(\theta^3(c_1(A)+rB)=(7r/2)\theta^3D\). This last
reformulation is only asserted under the stated extension-to-X
hypothesis.

\subsection{2. The intrinsic algebra map to the first infinitesimal
neighborhood}\label{SF-E825:the-intrinsic-algebra-map-to-the-first-infinitesimal-neighborhood}

Work first with one F and one free summand. Every FF block has the form
\(aI_2+vb\), with a in R and b a two-entry row modulo I. Its hull scalar
is a, and its action on torsion is a modulo I. Set

\[
 \ell(aI_2+vb)=a+bv\pmod{I^2}.
 \tag{E825.5}
\]

This is well defined: changing b by an I-valued row changes bv by I
squared. For two FF blocks the nilpotent products vanish on F, and
\((bv)(b'v)\in I^2\); hence (E825.5) preserves multiplication.

A map from F to R is c times the hull row w. A map from R to F is a
vector m represented by \((m_1,m_2)^t\). Put \(p(m)=wm\). Then

\[
 m w-p(m)I_2=v(m_2,-m_1),\qquad
 p(m)+(m_2,-m_1)v=0.
 \tag{E825.6}
\]

Thus every FF map factoring through a free module maps to zero under
ell. These identities also show independence of the chosen
representative of m. The free-free and other off-diagonal blocks map to
zero. Taking FF entries and applying (E825.5) therefore gives an actual
unital ring homomorphism

\[
 \mathcal E nd_R(F^r\oplus R^s)\longrightarrow
       M_r(R/I^2).
 \tag{E825.7}
\]

Indeed in each FF block of a product the only extra terms are precisely
factorizations through the free summand, killed by (E825.6). Matrix sums
and the FF multiplicativity complete the proof. A unit maps to a unit.

The fixed ordered global f,g make this compatible with changes of
L-frame. The standard comparison F\_j to F\_i is scalar t\_ij on its
generators and has chain lifts \(t_{ij}^3I_2,t_{ij}^4\) in degrees minus
one and minus two. If a FF block is rewritten using \(v_j=t_{ij}v_i\),
its row b changes by \(t_{ij}^{-1}\); thus bv and (E825.5) are
unchanged. Conjugation by the standard comparisons intertwines (E825.7).
Apply (E825.7) to the actual transition automorphisms of M after these
fixed standard comparisons. Their multiplicative cocycle remains exact
modulo I squared. It defines the vector bundle Qtilde on S2. Modulo I
the map is the action on the canonical quotient Q, proving the claimed
restriction. Changing local identifications of M changes this cocycle by
an ordinary change of bundle frames.

\subsection{3. The first Ext trace that survives off-diagonal
mixing}\label{SF-E825:the-first-ext-trace-that-survives-off-diagonal-mixing}

The full local computation in UG gives

\[
 \operatorname{Ext}^1(F,F)
   =U\oplus V,\quad U=uF/u(IF),\quad V=T^{\oplus2},\quad
 u=(f,g)^t,\quad T=(R/I)v.
 \tag{E825.8}
\]

Both summands have rank two over R/I. Applying Hom to (E825.1) with
target R gives \(\operatorname{Ext}^1(F,R)=R/I\), represented by the
row-pair \(c u\) modulo Iu. Also \(\operatorname{Ext}^{>0}(R,F)=0\) and
\(\operatorname{Ext}^{>0}(R,R)=0\).

Postcomposing the class cu by a map R to F represented by m gives
\(c u m\). As m varies these classes exhaust U. The reversed product is
in Ext1(R,R), hence zero. Therefore every cyclic trace must kill U; one
cannot retain the full unit-module trace from E824 after adding free
terms.

Define a trace on Ext1(F\^{}r plus R\^{}s) by taking, in each diagonal
FF block, its V-component and summing these components. All other blocks
map to zero. This trace is cyclic for every degree-zero endomorphism. On
FF blocks, only the scalar matrix a modulo I acts on positive Ext on
either side; ordinary matrix trace is cyclic. The nilpotent J actions
are zero by UG. An off-diagonal product that returns to FF is exactly a
product just shown to land in U. All remaining products have
off-diagonal or free-free target and have zero trace. This checks all
block possibilities and proves cyclicity, rather than assuming the F
summands are preserved.

Consequently conjugation by any actual automorphism leaves this trace
unchanged. Restoring the global twists, V is \(i_*\mathcal O_S(D)^2\).
The standard frame comparison transforms both V coefficients by t\_ij,
exactly this bundle transition. The trace therefore glues to an
intrinsic sheaf morphism

\[
 \tau_M:\mathcal E xt^1_X(M,M)\longrightarrow
                  i_*\mathcal O_S(D)^{\oplus2}.
 \tag{E825.9}
\]

This is only a trace on a positive Ext sheaf. No trace on the full
derived self-Hom object, decomposition of M, or invariant local summand
is asserted.

\subsection{4. The exact local Poisson deformation and the coefficient
seven
halves}\label{SF-E825:the-exact-local-poisson-deformation-and-the-coefficient-seven-halves}

Use the right-module star product \(a*b=ab+\epsilon\{a,b\}/2\), with
epsilon squared zero. Set P=\{f,g\} and d=vw. A minimal free deformation
of (E825.1) has differentials

\[
 d_\epsilon=d+\frac32\epsilon P I_2,\qquad v_\epsilon=v.
 \tag{E825.10}
\]

Indeed the first-order star product of d and v is \((-3P/2)v\), which
cancels the displayed correction. Its cokernel is flat over the dual
numbers, by the same elementary cycle-lifting argument used in E823:
reduce a cycle modulo epsilon, lift it using the exact original
resolution, and then divide the remaining epsilon cycle and lift once
more. Adding free star modules preserves this assertion.

Let f\_j=t f\_i, g\_j=t g\_i on an overlap, and suppress i. Put
\(\alpha=\{t,f\}\), \(\beta=\{t,g\}\), \(C=f\beta-g\alpha\),
\(z=(-\beta,\alpha)^t\). Then \(P_j=t^2P+tC\). The zeroth-order chain
comparison has matrices \(G_0=tI_2,G_1=t^3I_2,G_2=t^4\). Its first-order
failure to be a chain map in degree minus one to degree zero is

\[
 O_0=t^2\left(\frac32 C I_2+2\{t,d\}\right).
 \tag{E825.11}
\]

This follows by expanding \(G_0*d_{\epsilon,j}-d_{\epsilon,i}*G_1\): the
correction terms give \((3/2)t^2CI_2\), and the two star-product
brackets give \(2t^2\{t,d\}\). The next failure is
\(O_1=(7/2)t^3\{t,v\}\), and \(O_0v_j+dO_1=0\), so the projected cochain
is a self-Ext1 cycle. Normalize on the source by G1 inverse to express
it in the i-frame. Since \(CI_2=-z w+v(\alpha,\beta)\) in F-valued
columns and \(\{t,d\}=z w+v(\alpha,\beta)\), its two components are
exactly

\[
 \frac{O_0}{t^3}
   =\frac{1}{2t}z w+\frac{7}{2t}v(\alpha,\beta).
 \tag{E825.12}
\]

The first term lies in U and is killed by (E825.9). The second is the
surviving V term. The transition t is that of L inverse, because the
coefficients of sections of L obey f\_j=t f\_i. Thus the fixed
unit-module contribution is \(-(7/2)\Psi_\pi(L)\), with precisely the
sign and normalization used in (E825.3).

\subsection{5. Arbitrary gluing contributes the determinant of
Qtilde}\label{SF-E825:arbitrary-gluing-contributes-the-determinant-of-qtilde}

Consider a FF endomorphism \(aI_2+vb\) in a fixed f,g-frame. A chain
lift has higher matrices \((a+bv)I_2,a+bv\). For the local deformation
(E825.10), project its first-order lifting obstruction to V. For a
scalar a, the obstruction is \(\{a,d\}\); writing
\(\alpha_a=\{a,f\},\beta_a=\{a,g\}\) gives
\(\{a,d\}=(-\beta_a,\alpha_a)^tw+v(\alpha_a,\beta_a)\). Its V-component
is therefore \((\alpha_a,\beta_a)\).

For a nilpotent vb, first take b constant. With \(t_b=bv\), direct
multiplication gives the obstruction

\[
 P\bigl(vb-j_bw\bigr),\qquad j_b=(b_2,-b_1)^t.
 \tag{E825.13}
\]

The V-component is P(b1,b2), which equals \((\{t_b,f\},\{t_b,g\})\) on
S. For variable b, the extra terms are explicit boundaries, as follows.
Put

\[
 c=\sum_l\{b_l,v_l\},\qquad
 A_i=v_i c+\sum_l v_l\{b_l,v_i\},\qquad
 B_k=\sum_l v_l\{b_l,w_k\}.
 \tag{E825.13a}
\]

Expanding the two bracket products in the obstruction shows that all
terms differentiating b are exactly \(\frac12 A w+vB\). Each entry of A
lies in I, so Aw represents an element of u(IF), a boundary in (E825.8).
Each entry of B lies in I, so vB is zero in F. This proves vanishing as
an actual Ext class, not just a count of powers of I. Thus in general
the V-component is

\[
 \bigl(\{a+bv,f\},\{a+bv,g\}\bigr)|_S.
 \tag{E825.14}
\]

This also shows directly that it depends only on ell modulo I squared.
For r by r FF matrices the computation is entrywise. After normalization
by an inverse automorphism, its FF scalar matrix modulo I is A inverse,
where A is ell modulo I. Off-diagonal normalization terms are
compositions through a free module and land in U, so contribute no
trace. Nilpotent FF terms act trivially on positive Ext. Hence the
normalized V trace is

\[
 \operatorname{tr}\bigl(A^{-1}\{\ell,h\}\bigr)
   =\frac{\{\det\ell,h\}}{\det\ell}
 \quad\hbox{on }S,\qquad h=f,g.
 \tag{E825.15}
\]

The equality is the ordinary determinant derivative identity, proved by
expanding the determinant in permutations and collecting each
differentiated entry with its cofactor. It is valid here for the
derivation \(R/I^2\to R/I\), \(a\mapsto\{a,f\}|_S\), and likewise for g.

For clarity about choices: a change of a local module deformation by a
class beta changes an endomorphism-lifting obstruction by the commutator
with beta. After normalization and the cyclic trace (E825.9), its change
is zero. Changing local nullhomotopies adds a Cech coboundary. These are
the bar-cocycle computations written in E824, now using cyclicity rather
than falsely assuming the whole Ext1 bimodule is central.

As in E824, REnd(M) has cohomology sheaves only in degrees zero, one and
two. The two truncation triangles give a canonical first-overlap map
\(\kappa_M\) from the kernel of the local Ext2 edge map into
\(H^1(\mathcal E xt^1(M,M))\); its kernel is the image of
\(H^2(\mathcal E nd(M))\). The local deformations (E825.10) make the
Hochschild class rho\_M(pi) belong to this kernel. The Cech description
of kappa is precisely the obstruction to lifting transition
isomorphisms. Combining (E825.12) with (E825.15) therefore proves the
full identity

\[
 \tau_M\kappa_M\bigl(\rho_M(\pi)\bigr)
   =\Psi_\pi(\det\widetilde Q)-\frac{7r}{2}\Psi_\pi(L).
 \tag{E825.16}
\]

It is independent of all the allowed local choices. In particular,
off-diagonal free mixing has been included, not discarded by an assumed
endomorphism decomposition.

\subsection{6. The balanced kernel and the rational B
field}\label{SF-E825:the-balanced-kernel-and-the-rational-b-field}

First B=0. The two classes \(\xi_\pm=\pi_\pm+q\eta_\pm\), with
\(\eta_\pm=\pi_\pm\mathbin\lrcorner\theta_\pm^2/2\), annihilate vq under
the full Hochschild cap action. Weak G-semiregularity, the published BF
trace/evaluation identity, and supplied S21's invariant-to- ordinary
compatibility imply \(\rho_M(\xi_\pm)=0\) in ordinary Ext2. The eta term
is in the image of \(H^2(\mathcal E nd(M))\), so its kappa image is
zero. Applying (E825.16) gives (E825.3) for B=0.

For algebraic rational B choose a positive m such that \([m]^*B=m^2B\)
is the Chern class of a line bundle LB. Pull back M by the finite etale
map {[}m{]} and tensor by LB. Its normalized character is now the
ordinary balanced character. The pulled-back S is smooth and connected:
the negative-ample Koszul resolution has H0 equal to the constants. Thus
it is integral. The local form remains F' to the r plus a free rank-s
bundle; line tensoring does not change a local free form. The canonical
defect bundle becomes \([m]^*\widetilde Q\otimes LB|_{S'_2}\).

Exactly as in E824, etale Atiyah functoriality and line tensoring
transport the ordinary vanishing of the B-sheared characteristic class.
Here \(\pi' =m^{-2}\pi\) and \(\eta'=m^2\eta\). No statement that
semiregularity itself or the finite action pulls back is made. Applying
(E825.16) on the pulled-back space yields the pullback of (E825.3).
Finite etale trace makes the pullback on these coherent cohomology
groups injective in characteristic zero. It follows that (E825.3) holds
on X.

\subsection{7. Nonvanishing and the extension-to-X
consequence}\label{SF-E825:nonvanishing-and-the-extension-to-x-consequence}

For a line class A on X, Psi\_pi(A) is the image of its Atiyah class
under \(\Omega_X^1|_S\xrightarrow{\pi}T_X|_S\to N_{S/X}\). This
description follows immediately by differentiating its line cocycles.
For either normal component the Koszul resolution is

\[
 0\to L^{-1}\to\mathcal O_X^{\oplus2}\to L\to i_*L|_S\to0.
 \tag{E825.17}
\]

Ample and negative-ample vanishing on the abelian fourfold gives
\(H^1(S,L|_S)\simeq H^2(X,\mathcal O_X)^{\oplus2}\). To compute the
class, lift a Cech normal component \(\{\log\ell_{ij},f_i\}\) off S. On
a triple overlap its Cech coboundary is f\_i times the contraction of
\(d\log\ell_{jk}\) and \(d\log t_{ij}\), with a fixed sign depending on
the Cech convention. For the g component it is the same scalar class
multiplied by g\_i. Under (E825.17), the two components therefore map to
\((c,0)\) and \((0,c)\), where c is the contraction by pi of the cup
product of the two line Atiyah classes. Consequently

\[
 \Psi_\pi(A)=0
       \quad\Longleftrightarrow\quad
       \pi\mathbin\lrcorner(c_1(A)D)=0.
 \tag{E825.18}
\]

Only the kernel assertion is used, so the harmless common nonzero Cech/
Atiyah normalization does not enter (E825.3) or change this equivalence.
In particular \(\Psi_{\pi_\pm}(L)\ne0\), because
\(\pi_\pm\mathbin\lrcorner D^2\) is a nonzero multiple of eta\_pm. Flat
classes from X give zero.

For completeness the simultaneous kernel in (E825.18) is computed
directly. Choose a complex unitary basis adapted to the two RM planes,
so D is a positive scalar times the diagonal form and pi\_+ uses
holomorphic indices 1,2, while pi\_- uses 3,4. Write a (1,1) class as
the matrix C\_ab. Contraction of C D by pi\_+ vanishes precisely when
C\_11+C\_22=0 and C\_13,C\_14,C\_23,C\_24=0. The pi\_- equation gives
C\_33+C\_44=0 and C\_31,C\_32,C\_41,C\_42=0. These are exactly the two
separate primitive conditions and absence of mixed blocks. This proves
the reformulation (E825.4), and its stated degree consequence. If the
determinant class plus rB is flat from X, its Psi is zero whereas the
right side of (E825.3) is nonzero, proving the obstruction claimed in
section1.

\subsection{8. Exact scope, source boundary, and remaining construction
question}\label{SF-E825:exact-scope-source-boundary-and-remaining-construction-question}

The certificate rm\_free\_summand\_overlap\_certificate.py was actually
run: PASS,0.094044945seconds internal time. It verifies the corrected
minimal complex, the generic line-frame overlap including its U
coefficient1/2 and V coefficient7/2, nilpotent/scalar obstructions,
free-factorization cancellation, ell multiplicativity and a generic2x2
determinant derivative. It does not compute global Cech groups or prove
a semiregular example. The local/global mathematical arguments are
supplied in full above.

The external inputs are exactly the published BF2008
evaluation/composition proofs already read for E823/E824: Theorem5.1.3,
authorpp24--25; Proposition5.2.2 and Corollary5.2.3, pp26--27;
Proposition6.2.1, pp29--30; Theorem6.4.2 with Proposition6.4.4,
pp31--33. The published article is Advances in
Mathematics217(2008),243--281, DOI10.1016/j.aim.2007.06.013, author copy
https://arxiv.org/pdf/math/0606730 . The vanishing input is Mumford,
Abelian Varieties (1970), section16 printedpp150--153, already read for
E813. Supplied S21 is used solely for coherent invariant evaluation and
ordinary forgetting, at(S21.6) and the compatibility paragraphs
after(S21.7) and before(S21.8). No new preprint theorem or full derived
trace is introduced. E824's truncation/Cech construction is used in its
full proved scope; its centrality hypothesis is replaced here by the
explicit cyclic trace calculation, not silently reused.

When the defect determinant has the required positive degree, (E825.3)
may hold; this proof does not exclude it. It also does not cover
matrices whose rank drops on S, other local torsion/hull models,
singular support, or arbitrary perfect complexes. An actual gluable
object with the full balanced character and this necessary defect
determinant would still require its entire Hochschild evaluation and
semiregularity to be checked. Local first-order deformations alone do
not provide that seed.

\subsection{9. Review and exact scope
record}\label{SF-E825:review-and-exact-scope-record}

Root read the complete original draft, the full variable-b amendment and
final wording, then the complete peer report and independent code; PASS.
The peer read the entire final proof and all author code, and
independently checked all six bivector coefficients with unrestricted
coefficient functions, the explicit derivative boundaries, free-factor
cancellation, a generic3x3 determinant and the full16-variable
primitive/mixed kernel. Its actual computation passed in0.601104seconds.
Report:
../g34\_eightfolds/bireflection/FREE\_SUMMAND\_OVERLAP\_PEER\_AUDIT.md.
The root actual author-certificate rerun passed in0.114255109seconds
internal time,0.44637168seconds tool wall time. These checks preserve
the necessary-condition scope: positive adjusted defect determinants
remain a construction gate, and no universal added-free exclusion,
finite action or semiregular seed is asserted. No PDF compilation or
proof-assistant verification is claimed.

\section{SF-E826. Coupled two-slope extensions: actual indecomposable
objects and a middle-Ext
obstruction}\label{SF-E826:sf-e826.-coupled-two-slope-extensions-actual-indecomposable-objects-and-a-middle-ext-obstruction}

Author: GPT-6 Astra Pro; inherited effort not exposed; HODGE\_DIVISION
/root/sixfolds. Date: 2026-11-21.

\textbf{Status:} independently whole-reviewed by /root and
/root/g34\_eightfolds; both complete mathematical reviews PASS on the frozen
draft
\texttt{58eb0b23279e278860bf305a8038231897d5359ed1441e4f0c9bd797a98727ce}.
The author actually ran the two exact certificates; the peer supplied an
independently written exact certificate. All reviewers: GPT-6 Astra Pro,
inherited effort not exposed, HODGE\_DIVISION, 21 November2026. This
excludes an entire precisely specified extension architecture, including an
actual indecomposable balanced example. It does not exclude all
positive-determinant forms from E825. A connected Ext graph is not used as a
substitute for constructing an extension or checking its endomorphisms. No
formality or vanishing of uncomputed Yoneda products is assumed.

\subsection{1. An actual triangle lemma, with its invariant
form}\label{SF-E826:an-actual-triangle-lemma-with-its-invariant-form}

Let X be a smooth projective fourfold and let

\[
 0\longrightarrow A\xrightarrow{i}E\xrightarrow{p}B\longrightarrow0,
 \qquad e:B\longrightarrow A[1],
 \tag{E826.1}
\]

be an actual coherent extension. Suppose
\(\operatorname{Ext}^1(A,B)=\operatorname{Ext}^3(B,A)=0\). Then

\[
 \operatorname{ext}^2(E,E)\ge
 \operatorname{ext}^2(A,A)+\operatorname{ext}^2(B,B).
 \tag{E826.2}
\]

If a finite group has a coherent action preserving this exact sequence, the
same inequality holds for invariant Ext groups. On an abelian fourfold, the
ordinary Ext1 vanishing follows from the displayed Ext3 vanishing by Serre
duality. The following proof uses actual long exact sequences, not a
degeneration claim for a spectral sequence.

Apply Hom(A,-) to the triangle. The Ext1(A,B) vanishing gives both a
surjection Ext1(A,A) to Ext1(A,E) and an injection
\(i_*:\operatorname{Ext}^2(A,A)\to\operatorname{Ext}^2(A,E)\). Apply
Hom(B,-). The Ext3(B,A) vanishing gives a surjection
\(\operatorname{Ext}^2(B,E)\to\operatorname{Ext}^2(B,B)\). Finally apply
Hom(-,E), obtaining the exact segment

\[
 \operatorname{Ext}^1(A,E)\xrightarrow{\partial_1}
 \operatorname{Ext}^2(B,E)\longrightarrow
 \operatorname{Ext}^2(E,E)\longrightarrow
 \operatorname{Ext}^2(A,E)\xrightarrow{\partial_2}
 \operatorname{Ext}^3(B,E).
 \tag{E826.3}
\]

Every class in Ext1(A,E) is i{[}1{]}a for some a in Ext1(A,A). Its boundary
is i{[}2{]}a{[}1{]}e, so it is in the image of Ext2(B,A). It therefore maps
to zero in Ext2(B,B). Hence the quotient
\(Q=\operatorname{Ext}^2(B,E)/\operatorname{im}\partial_1\), which injects
into Ext2(E,E), still surjects onto Ext2(B,B). For a in Ext2(A,A), the other
boundary of i{[}2{]}a is i{[}3{]}a{[}1{]}e. Its middle composite a{[}1{]}e
is in Ext3(B,A), hence zero. Thus the quotient of Ext2(E,E) by Q contains
the injected copy of Ext2(A,A). This proves(E826.2). All maps are functorial
for the given coherent action; finite invariants are exact in characteristic
zero. The identical argument therefore proves the invariant inequality. No
equivariant splitting of (E826.1) is used.

\subsection{2. The full extension architecture excluded by the
lemma}\label{SF-E826:the-full-extension-architecture-excluded-by-the-lemma}

Work on the balanced RM abelian fourfold, with ample theta, D=4theta, a
smooth ordered complete-intersection surface S, and the actual unit sheaf F
of E823--E825. Put \(\mathscr U=F\otimes O_X(14\theta)\). Consider any
extension(E826.1) with

\[
 A=\mathscr U^{\oplus r}\oplus\bigoplus_{a=1}^{u}L_{\alpha_a}^{\oplus n_a},
 \qquad
 B=\bigoplus_{b=1}^{v}L_{\beta_b}^{\oplus m_b},
 \tag{E826.4}
\]

where r and all multiplicities are positive. Assume:

\begin{itemize}
\tightlist
\item
  all alpha\_a have theta-slope14theta4, are pairwise distinct as Chern
  classes, and differ from14theta;
\item
  all beta\_b have theta-slope minus14theta4 and are pairwise distinct as
  Chern classes;
\item
  Ext3(B,A)=0;
\item
  the entire normalized character is
  \(e^{B_{\mathrm{fld}}}\operatorname{ch}(E)=R v_q\) for q nonzero and
  \(B_{\mathrm{fld}}\) rational algebraic, and u+v is at least4.
\end{itemize}

Then E has \textbf{no weakly semiregular lift for any specified effective
coherent finite translation/tensor action} of the kind in supplied S21. This
is a statement about every extension in this architecture, not just a
selected matrix or an indecomposable example.

Here is the equivariance argument in full. T=tors(E) is canonical. Since B
is locally free, the sequence splits locally, so T=tors(A) and

\[
 A/T=(I_S\otimes O_X(14\theta))^{\oplus r}
       \oplus\bigoplus_a L_{\alpha_a}^{\oplus n_a}.
 \tag{E826.5}
\]

Every factor is torsionfree of rank one and slope14theta4; it is slope
semistable. Their direct sum is semistable. B is semistable of slope
minus14theta4. Thus A/T is the positive part of the two-step HN filtration
of E/T. One can characterize it without invoking any splitting theorem: it
is the sum of the images in E/T of all semistable torsionfree sheaves of
positive slope. Every such image lies in A/T, because a nonzero map to the
semistable negative-slope B is impossible; its image would be both a
quotient of positive slope and a subsheaf of negative slope. Conversely A/T
itself is such a sheaf. This characterization is intrinsic.

Translations and tensoring by a degree-zero line bundle preserve torsion,
semistability and these slopes. Consequently any coherent finite action on E
preserves A, the inverse image of this intrinsic A/T, and induces an actual
coherent action on B. No additional linearization is chosen.

Each displayed line-class block in A and B is also an actual equivariant
summand of its respective quotient. Distinct line bundles of the same slope
with distinct Chern classes have no Hom in either direction: a nonzero
section of their difference would give an effective nonzero divisor of
theta-degree zero. The blocks involving U are orthogonal as well. Indeed U
dual is O\_X(-14theta), so Hom(U,L\_alpha)=H0(L\_alpha-14theta)=0. For the
reverse Hom use the torsion sequence for F: its hull term has no section for
the same reason, and the torsion term on S has class
\(-2D-(\alpha-14\theta)\), whose theta-degree on S is negative. It has no
nonzero section either. End(U)=C, since End(F)=O\_X plus O\_S(-D)\^{}2 and
the latter has no global section. These vanishing statements show that the
displayed projections are the central block projections. An action cannot
mix their different Chern classes, and cannot turn U, which has nonzero
torsion, into a line bundle. It therefore preserves all those projections
and summands. This is an argument on the actual sheaves, not on a
decomposition appearing only after an isogeny.

For each positive-rank line-class block its Chern character is n exp(alpha).
The degree-two Hochschild cap map has rank6: conjugate by exp(alpha), and
its image becomes H2(O\_X), which has dimension6. The published BF
trace/evaluation identity shows that its invariant self-Ext2 dimension is at
least6, whether or not that block itself is semiregular. Evaluation lands in
invariants because these translations/tensor functors act trivially on
Hochschild cohomology. Thus

\[
 \operatorname{ext}^2_\Gamma(A,A)+\operatorname{ext}^2_\Gamma(B,B)
       \ge6(u+v).
 \tag{E826.6}
\]

The triangle lemma applies after invariants and yields the same lower bound
for E. On the other hand, weak semiregularity and the invariant CY4
adjunction in supplied S21, specifically(S21.8), force
\(\operatorname{ext}^2_\Gamma(E,E)=\operatorname{rank}c^2_{Rv_q}=18\). For
clarity, the rank18 calculation is the elementary E801 product calculation:
in the volume-form exterior model the balanced character is omega\_+ wedge
omega\_-, with two nondegenerate two-forms on complementary four-spaces.
Pure exterior-square pieces contribute1 each, and the mixed piece
contributes16 because wedge with each omega is an isomorphism from vectors
to three-vectors. The total is18. A rational B field changes this rank by an
invertible exterior shear, as in the full S21 proof: contraction by a vector
becomes contraction plus wedge by its contraction with B. Thus the
normalized character has the same rank, without pulling back E or its
action. Since6(u+v) is greater than18, there is a contradiction. This proves
the entire architecture assertion.

\subsection{3. Exact character data for an actual coupled
example}\label{SF-E826:exact-character-data-for-an-actual-coupled-example}

We now construct a member with no nontrivial direct-sum decomposition, while
retaining the preceding obstruction. Use the actual Gaussian fourfold and
integral divisor basis(theta,theta\_U,theta\_S,theta\_US) of the
positive-defect construction. Explicitly X=E\_i\^{}4, theta is product
principal, U is block diagonal with two copies of the
matrix{[}{[}1,1{]},{[}1,-1{]}{]}, and S is diag(1,1,-1,-1). Every symmetric
integral matrix gives an actual line bundle by coordinate divisor pullbacks
and the pairwise Poincare bundles (addition pullback minus the two
coordinate divisors). These U,S,US are symmetric integral. D=4theta is very
ample since degree4 separates length-two subschemes on each elliptic factor;
the product embeddings and Segre embedding give the claim. The previously
proved Bertini/Koszul argument gives a smooth integral ordered complete
intersection. For a tuple(a,b,c,d), average its exponential over the four
distinct tuples

\[
 (a,b,c,d),\quad(a,-b,c,-d),\quad(a,b,-c,-d),\quad(a,-b,-c,d).
 \tag{E826.7}
\]

All eight orbits in the following table have four distinct elements. The
last column is the multiplicity of \textbf{each} line bundle in that orbit.
The table is exact rational data; floating linear programming was used only
to discover its support, and supplies no premise of this proof.

\begin{longtable}[]{@{}
  >{\raggedright\arraybackslash}p{(\columnwidth - 4\tabcolsep) * \real{0.2727}}
  >{\raggedleft\arraybackslash}p{(\columnwidth - 4\tabcolsep) * \real{0.3636}}
  >{\raggedleft\arraybackslash}p{(\columnwidth - 4\tabcolsep) * \real{0.3636}}@{}}
\toprule\noalign{}
\begin{minipage}[b]{\linewidth}\raggedright
Orbit representative
\end{minipage} & \begin{minipage}[b]{\linewidth}\raggedleft
Total weight per unit copy
\end{minipage} & \begin{minipage}[b]{\linewidth}\raggedleft
Multiplicity per line
\end{minipage} \\
\midrule\noalign{}
\endhead
\bottomrule\noalign{}
\endlastfoot
(-14, 4, 12, -8) & 9851034439/7111578272 & 118212413268 \\
(-14, 8, 8, -4) & 803399867/888947284 & 77126387232 \\
(-14, 12, 16, -12) & 361040431311/28446313088 & 1083121293933 \\
(-14, 16, 16, -32) & 21876066325/28446313088 & 65628198975 \\
(14, 8, 12, 4) & 154311465847/10667367408 & 1234491726776 \\
(14, 8, 24, -16) & 37616523853/85338939264 & 37616523853 \\
(14, 12, 32, 12) & 2658612777/28446313088 & 7975838331 \\
(-14, 16, 24, 16) & 1/4 & 21334734816 \\
\end{longtable}

Put r=341355757056 and R=32r=10923384225792. The twelve lines with a=14
belong to A, together with U to multiplicity r. The twenty lines with a=-14
belong to B. The free positive lines have total rank15r and the negative
lines total rank16r. Thus both A and B have rank16r.

Here is a compact full-character verification. For an orbit representative
write \(V_0=2b^2+c^2+2d^2\). Its orbit-averaged exponential has seven
independent subset coefficients

\[
 \begin{split}
 &(1,a,\ a^2+2b^2-c^2-2d^2,\ a^2-2b^2+c^2-2d^2,\\
 &\quad a^2-2b^2-c^2+2d^2,\ a^3-aV_0+4bcd,\\
 &\quad a^4-2a^2V_0+16abcd+4b^4+c^4+4d^4
                      -4b^2c^2-8b^2d^2-4c^2d^2).
 \end{split}
 \tag{E826.8}
\]

The three degree-two positions are the pairs12,13,14; the other three pairs
are their orbit partners. All degree-one coefficients agree, as do all
degree-three coefficients. The character of U is
\(2e^{10\theta}-2e^{2\theta}+e^{-2\theta}\), as follows directly from the
actual unit resolution. Weighted by the table, the free part has the exact
seven coefficients

\[
 (31,-14,-212,-196,-196,-1976,-19976).
 \tag{E826.9}
\]

Consequently the sum of A and B is R v\_(1/2) in \textbf{all sixteen} subset
coordinates. The exact certificate independently expands those sixteen
coordinates over Q(sqrt2), so no omitted component or orbit convention is
hidden by the seven-coordinate notation. Every coherent extension of B by A
has that complete character, local form F\^{}r plus a free rank(R-r), and
canonical defect Qtilde=O\_X(14theta)\^{}r on S2. In a triangular transition
matrix the off-diagonal free-to-A extension maps disappear under E825's ring
map ell, so its determinant satisfies the positive requirement exactly.

\subsection{4. Actual Ext1 spaces and absence of the opposite Ext3
direction}\label{SF-E826:actual-ext1-spaces-and-absence-of-the-opposite-ext3-direction}

For a line pair with difference(a,b,c,d), its two Gaussian 2-by-2 blocks
have parameters(a+c,b+d) and(a-c,b-d). Their determinants are
\((a+c)^2-2(b+d)^2\) and \((a-c)^2-2(b-d)^2\). In this table every cross
difference alpha-beta is nondegenerate. Its index is the sum of the numbers
of negative eigenvalues in those two blocks; each contributes one for
negative determinant, zero for positive determinant and positive trace, or
two for positive determinant and negative trace. The full index counts over
the240 line pairs are36 of index0,144 of index1, and60 of index2. None has
index3. The published index theorem therefore gives

\[
 \operatorname{Ext}^3(L_{\beta_b},L_{\alpha_a})=0
       \quad\hbox{for every such pair.}
 \tag{E826.10}
\]

For the unit target U, Ext3(L\_beta,U)=H3(F tensor O(14theta-beta)). In its
length-two resolution only H3 of the first term and H4 of the middle term
could contribute to total degree3. The first term is the line with
class10theta-beta; for every one of the twenty beta's it is nondegenerate of
index0,1 or2, so its H3 vanishes. The middle line has class2theta-beta, of
positive theta-slope16theta4, so its H4 vanishes by Serre duality and the
negative degree of its dual. This proves

\[
 \operatorname{Ext}^3(B,A)=0
 \tag{E826.11}
\]

for every smooth regular choice of the two surface sections. It is not a
numerical rank guess for a differential of the resolution.

For a nondegenerate line pair of index1 its Ext1 dimension is the absolute
product of the two determinants above. For the four beta's in the last orbit
of the input table, the three line terms of the resolution computing
Ext1(L\_beta,U) all have index1, and their dimensions are
4718592,3174400,2467584. Hence

\[
 \dim\operatorname{Ext}^1(L_\beta,U)
        \ge2(4718592-3174400)=3088384.
 \tag{E826.12}
\]

This is again a direct cokernel lower bound, requiring no maximal-rank
assertion. These exact dimensions provide all the arrows used below.

\subsection{5. A uniform explicit matrix construction forcing scalar
pairs}\label{SF-E826:a-uniform-explicit-matrix-construction-forcing-scalar-pairs}

Let a single directed edge have source dimension m and target dimension n.
If it provides at least ceil(max(m/n,n/m))+2 independent Ext1 classes, one
can choose coefficient matrices on those classes such that the only pairs
(T,S) satisfying T E\_j=E\_j S for every chosen matrix are (T,S)=(lambda
I\_n,lambda I\_m).

Here is the full construction. If n is at least m, choose ceil(n/m)
injections F\_j:C\^{}m to C\^{}n whose images span C\^{}n, taking
consecutive blocks of m coordinates and the last block ending at n.~Include
F1 J\_m and F1 J\_m transpose as two further arrow matrices, where J\_m is
the nilpotent single Jordan shift. The relations for F1 and F1 J\_m imply
that S commutes with J\_m; the corresponding transpose relation makes it
commute with J\_m transpose. These two shifts generate the full matrix
algebra: their products give the end-coordinate rank-one projections and
then every matrix unit. Therefore S is scalar. The spanning injections force
T to be the same scalar.

If m is greater than n, use ceil(m/n) coordinate projections P\_j:C\^{}m to
C\^{}n with intersection of kernels zero, and add J\_n P1 and J\_n transpose
P1. Their equations make T commute with both shifts, hence scalar; the joint
injectivity of the projections forces S to be the same scalar. This proves
the claim for all positive m,n, including the one-dimensional cases.

In a bipartite tree apply this construction separately on each edge, using
linearly independent classes in its actual Ext1 space and zero coefficients
on all unused classes. The scalar constraints on adjacent vertices agree.
Connectedness of the tree therefore forces all vertex endomorphisms to be
one common scalar. This is a construction with explicit ordinary matrices,
not a generic-representation or quiver-root theorem.

\subsection{6. The exact tree and the resulting indecomposable
sheaf}\label{SF-E826:the-exact-tree-and-the-resulting-indecomposable-sheaf}

Index A0=U, and list the twelve positive line tuples lexicographically as
A1,\ldots,A12. List the twenty negative line tuples lexicographically as
B0,\ldots,B19. These indices are uniquely determined by the orbit table and
(E826.7). The following tree gives, for every edge, the required number of
matrix arrows and the proved available Ext1 dimension or lower bound.

\begin{longtable}[]{@{}rrrr@{}}
\toprule\noalign{}
A index & B index & Matrices required & Available Ext1 \\
\midrule\noalign{}
\endhead
\bottomrule\noalign{}
\endlastfoot
0 & 0 & 18 & 3088384 \\
0 & 3 & 18 & 3088384 \\
0 & 16 & 18 & 3088384 \\
0 & 19 & 18 & 3088384 \\
2 & 0 & 5 & 7462656 \\
4 & 0 & 60 & 852992 \\
5 & 0 & 60 & 5727232 \\
8 & 0 & 60 & 5669888 \\
9 & 0 & 60 & 9626624 \\
10 & 0 & 4 & 33165056 \\
11 & 0 & 5 & 1990400 \\
12 & 0 & 5 & 7462656 \\
1 & 3 & 5 & 7462656 \\
7 & 3 & 4 & 33165056 \\
6 & 16 & 4 & 33165056 \\
3 & 19 & 4 & 33165056 \\
2 & 1 & 11 & 517888 \\
2 & 5 & 138 & 790272 \\
2 & 7 & 12 & 517888 \\
2 & 9 & 17 & 164864 \\
2 & 11 & 17 & 336896 \\
2 & 13 & 12 & 308992 \\
2 & 15 & 138 & 790272 \\
2 & 18 & 11 & 8070912 \\
4 & 2 & 21 & 435200 \\
4 & 14 & 4 & 156672 \\
4 & 17 & 21 & 2384896 \\
5 & 4 & 4 & 91136 \\
5 & 8 & 13 & 302848 \\
5 & 12 & 19 & 73728 \\
9 & 6 & 19 & 73728 \\
9 & 10 & 13 & 302848 \\
\end{longtable}

This is a tree on all33 vertices, and every displayed capacity is
sufficient. Its multiplicity ratios are the exact ratios from the input
table; the largest number of matrices needed is138. The attached exact
certificate checks the whole tree, all inequalities and every individual Ext
calculation.

The summands A0,\ldots,A12 are pairwise Hom-orthogonal with scalar
individual endomorphism rings, as proved in section2. The same is true of
B0,\ldots,B19. Also Hom(A,B)=0, by their slope comparison (for U use its
hull to a locally free target). Choose the extension class e in Ext1(B,A) by
the matrix construction of section5 on the tree. Since B is locally free,
Ext1(B,A)=H1(B dual tensor A); an actual Cech cocycle glues the local
modules A plus B by upper triangular transition matrices. Their strictly
upper block has square zero, so the usual additive cocycle equation is the
whole gluing condition. There is no additional Maurer--Cartan equation or
assumed dg formality. This gives an actual coherent E in(E826.1), with the
full character and local type already proved in section3. It is ordinary
gluable because it is a coherent sheaf.

We verify indecomposability of this actual object. Every endomorphism of E
preserves A, since its composite from A into B is zero. Its restrictions
(a,b) must satisfy ae=eb. By pairwise Hom orthogonality these are precisely
the vertex matrices, and the tree construction forces all of them to be the
same scalar. Conversely the scalar identity lifts. The kernel of the
restriction map is Hom(B,A), realized by i f p; this is a square-zero ideal
because p i=0. To see that the compatibility condition is sufficient for a
lift, represent e by a Cech cocycle: ae-eb being a coboundary gives local
upper-right entries which complete a,b to an actual global endomorphism.
Therefore

\[
 \operatorname{End}(E)/\operatorname{Hom}(B,A)=\mathbf C,
 \qquad\operatorname{Hom}(B,A)^2=0.
 \tag{E826.13}
\]

An idempotent in this algebra maps to0 or1 in C. If it maps to0 it lies in
the square-zero ideal and is zero; if it maps to1, its complement does so
and is zero. Hence there is no nontrivial idempotent. The coherent category
is idempotent complete, so E is indecomposable. No simplicity is claimed:
Hom(B,A) may be large.

\subsection{7. The invariant middle-Ext obstruction still excludes the whole
family}\label{SF-E826:the-invariant-middle-ext-obstruction-still-excludes-the-whole-family}

All the hypotheses in section2 hold for every extension with these A,B,
including the indecomposable one. There are u+v=12+20=32 canonical
line-class blocks. For any specified coherent effective finite action on its
middle sheaf, the torsion/HN argument provides the required actual actions
on A and B and their line-class blocks. Equation(E826.11), Serre duality and
the exact invariant triangle lemma therefore imply

\[
 \operatorname{ext}^2_\Gamma(E,E)\ge192.
 \tag{E826.14}
\]

The full balanced character has cap rank18. Thus no member is weakly
Gamma-semi-regular. This failure persists after all the specified
free-to-free couplings, even though the selected member is indecomposable
and satisfies the positive defect determinant from E825. The proof of this
failure examines actual invariant Ext and the full Hochschild cap rank; it
does not infer failure merely from the split starting character.

The general scope is the whole two-slope architecture in section2 with
Ext3(B,A)=0 and at least four distinct line-class blocks. To evade this
particular obstruction, one needs a nonzero opposite cross Ext3 direction, a
different intrinsic filtration, or a construction outside that architecture.
No claim is made that those alternatives succeed or that all unstable
sheaves have the stated form.

\subsection{8. Source and check boundary}\label{SF-E826:source-and-check-boundary}

The ordinary abelian index theorem is the previously read original Mumford,
Abelian Varieties(1970), section16 printedpp150--153 and162--163. The unit
sheaf resolution and End description are the fully proved UG/US inputs.
Actual section existence uses the previously checked very-ampleness and
Bertini argument for D=4theta. The general triangle lemma, matrix
construction, Cech extension and endomorphism proof are supplied in full
above.

For the categorical rank consequence, this proof uses the full invariant CY4
adjunction from supplied DEG\_L S21, including the paragraphs preceding
(S21.8), not merely ordinary forgetting compatibility. The entire displayed
S21 proof through(S21.8) was freshly reread for this derivation; its
original source statuses are retained. Published BF2008 Theorem5.1.3,
Proposition5.2.2, Corollary5.2.3 and Proposition6.2.1/Theorem6.4.2 with
Proposition6.4.4, authorpp24--27 and29--33, give evaluation/trace
composition as previously fully read. No supplied preprint premise has been
renamed as a new published theorem. The full balanced rank18 calculation
from E801 was freshly reread and its short proof is repeated in section2.

Exact input is rm\_two\_slope\_exact\_input.json. The complete certificate
rm\_two\_slope\_coupling\_certificate.py verifies the full
sixteen-coordinate character, positive integral multiplicities, all cross
Ext3 vanishings, all148 sufficient edge capacities, and the33-vertex
spanning tree. It imports no floating optimization result as a theorem:
exact rational input is checked from scratch. The exploratory
linear-programming script is retained only as provenance. An initial
certificate output serialization failure was repaired by converting an exact
symbolic integer to a Python integer; the failure is retained in
rm\_two\_slope\_coupling\_attempts.json. No generic multiplication rank,
invariant action or semiregularity is claimed to have been computationally
verified. The latter is excluded by the written invariant argument above.

The separate exact matrix checker rm\_rectangular\_brick\_certificate.py
actually passed all36 pairs1\textless=n,m\textless=6 over the prime1009.
Each relation matrix has rank n²+m²−1 modulo that prime; the scalar solution
supplies the opposite characteristic-zero bound, so those finite cases are
exact. The all-size result is the complete injection/projection and Jordan
proof in section5, not an extrapolation from these cases.

\section{SF-E827. Diagonal Yoneda classes survive arbitrary cross-index
changes}\label{SF-E827:sf-e827.-diagonal-yoneda-classes-survive-arbitrary-cross-index-changes}

Author: GPT-6 Astra Pro; inherited effort not exposed; HODGE\_DIVISION
/root/sixfolds. Date: 2026-11-21.

\textbf{Status:} independently whole-reviewed by /root and
/root/g34\_eightfolds; both mathematical reviews PASS on the final frozen
draft
\texttt{abe42293fb48da62436bfe7fafd7ae19a2546a768ad294f6f89ed4ec451f8a27}.
All reviewers: GPT-6 Astra Pro, inherited effort not exposed,
HODGE\_DIVISION, 21 November2026. No new numerical computation is a premise.
This is a separate strengthening of E826. It removes the global Ext3(B,A)=0
hypothesis in specified line-class settings by following diagonal classes
through the actual triangle maps. It does not claim all positive-determinant
objects or all divisor directions are excluded.

\subsection{1. Exact setting and a diagonal-survival
lemma}\label{SF-E827:exact-setting-and-a-diagonal-survival-lemma}

Let X be an abelian fourfold. Consider an actual coherent extension

\[
 0\to A\xrightarrow{i}E\xrightarrow{p}B\to0,\qquad e:B\to A[1],
 \quad
 A=A_0\oplus\bigoplus_{a=1}^{u}L_a^{\oplus n_a},
 \quad B=\bigoplus_{b=1}^{v}M_b^{\oplus m_b}.
 \tag{E827.1}
\]

Assume a specified coherent finite effective translation/tensor action
preserves this exact sequence and the displayed blocks. These are actual
equivariant summands of A and B, not asserted summands of E. A0 is allowed
to be zero or an arbitrary coherent perfect object in this lemma. In the
application A0 is the unit block and the intrinsic torsion/HN argument of
E826 supplies all the assumed equivariance for every action on E.

Put V=H1(X,O\_X), so H2(O\_X)=Lambda2 V. Let W,Z be vector subspaces of
H2(O\_X) with W intersect Z zero. Assume for every a,b:

\begin{enumerate}
\def\labelenumi{\arabic{enumi}.}
\tightlist
\item
  the Yoneda product of Ext1(L\_a,M\_b) with Ext1(M\_b,L\_a) has image in Z
  inside Ext2(L\_a,L\_a)=H2(O\_X), and its reversed product has image in Z
  inside Ext2(M\_b,M\_b);
\item
  multiplication by W kills Ext1(M\_b,L\_a) into Ext3(M\_b,L\_a), and also
  kills Ext1(L\_a,M\_b) into Ext3(L\_a,M\_b).
\end{enumerate}

Let J be any subset of the B indices such that Ext3(M\_b,A0)=0 for b in J.
Then

\[
 \operatorname{ext}^2_\Gamma(E,E)\ge(\dim W)(u+|J|).
 \tag{E827.2}
\]

If A0=0, J may consist of all v indices. The statement also holds without an
action, by taking the trivial group. No vanishing of all cross Ext3 groups
is assumed.

Here is the complete proof, retaining multiplicity matrices. In Ext2(A,A)
let D\_A be the direct sum over a of W tensor the identity matrix on the
multiplicity space of L\_a. Thus dim D\_A=(dim W)u. Its elements have zero
A0 block and zero off-diagonal blocks. Each is invariant: translations and
degree-zero tensoring act trivially on H2(O\_X), and conjugation fixes the
identity of any multiplicity representation, including a projective one.

In Hom(A,triangle), the kernel of
\(i_*:\operatorname{Ext}^2(A,A)\to\operatorname{Ext}^2(A,E)\) is the image
of Ext1(A,B). For a class gamma in that group, its image is \(e[1]\gamma\).
Its a-th diagonal line block is exactly

\[
 \sum_b e_{ab}[1]\gamma_{ba},
 \tag{E827.3}
\]

with the multiplicity matrices multiplied in this order. Hypothesis1 places
every coefficient in Z, hence the whole block in Z tensor
End(C\^{}\{n\_a\}). This has zero intersection with W tensor its identity.
Thus D\_A injects into Ext2(A,E), regardless of any off-diagonal image or
any component involving A0. The absence of A0 in B is used here.

In Hom(-,E), the next boundary from Ext2(A,E) to Ext3(B,E), on a class
\(i[2]d\) with d in D\_A, is \(i[3]d[1]e\). Its only nonzero possible
components before postcomposition by i are \(d_a[1]e_{ab}\). They are zero
by hypothesis2, also with arbitrary multiplicity matrices, because d\_a is a
scalar identity times a class in W. Consequently the image I of Ext2(E,E) in
Ext2(A,E) contains this injected D\_A.

The B-side contributes independently to the kernel of that map. In
Ext2(B,B), let D\_B be the sum over b in J of W tensor the corresponding
identity matrices. The connecting map from Ext2(B,B) to Ext3(B,A) sends such
a class h to \(e[2]h\). Every line-target component is zero by hypothesis2,
and every A0 component is zero by Ext3(M\_b,A0)=0. Hence D\_B lifts to
Ext2(B,E).

We must check that the preceding boundary in Hom(-,E) cannot kill these
lifted diagonal classes. For any alpha in Ext1(A,E), its boundary is
\(\alpha[1]e\). Projecting to Ext2(B,B) gives

\[
 (p[1]\alpha)[1]e.
 \tag{E827.4}
\]

The b-th diagonal block, b in J, is a sum of line-to-line degree-one
products in Z tensor End(C\^{}\{m\_b\}). The additional possible term
through A0 is zero: by fourfold Serre duality Ext1(A0,M\_b) is dual to
Ext3(M\_b,A0), which is zero. Therefore the diagonal projection of every
such boundary has no W component after passing to H2(O\_X)/Z.

More formally, take diagonal projection, matrix trace divided by m\_b, and
the quotient by Z on each selected block. This is an equivariant linear map
from Ext2(B,B) to a space containing an injected copy of
\(\bigoplus_{b\in J}W\). It kills the projection of all the boundaries
(E827.4), while the lifts of D\_B map onto that copy. Thus

\[
 Q=\operatorname{Ext}^2(B,E)\big/\operatorname{im}
     \bigl(\operatorname{Ext}^1(A,E)\to\operatorname{Ext}^2(B,E)\bigr)
 \tag{E827.5}
\]

has dimension at least(dim W)\textbar J\textbar, and injects into Ext2(E,E).
Its quotient there is I, which already contains D\_A. All constructions
commute with the finite action. Taking invariants is exact; invariant lifts
of D\_B exist by averaging; W and the identity matrices remain fixed. The
same dimension bounds therefore hold for Q invariants and I invariants. This
proves (E827.2). These are the two actual triangle boundary maps, not an
assumed associated-graded decomposition of E.

\subsection{2. Nondegenerate line differences: six classes per positive
block}\label{SF-E827:nondegenerate-line-differences-six-classes-per-positive-block}

Suppose every line bundle L\_a tensor M\_b inverse is nondegenerate. The
abelian index theorem puts all its cohomology in a single degree j, and that
of its inverse in degree4-j. Thus its Ext1 and the reverse Ext1 cannot both
be nonzero. Their products in either diagonal self-Ext2 are zero. Moreover,
whenever its Ext1 is nonzero, j=1 and its Ext3 is zero; the same applies to
the inverse. Hypotheses1--2 of section1 hold with

\[
 Z=0,\qquad W=H^2(X,O_X),\qquad\dim W=6.
 \tag{E827.6}
\]

Consequently ext2\_Gamma(E,E) is at least6(u+\textbar J\textbar). In
particular u at least4 already excludes a weakly semiregular balanced
character of cap rank18, without \textbf{any} condition on Ext3(B,A0) or on
other cross Ext3 groups. Opposite index-three pairs elsewhere do not cancel
the diagonal classes: the extension component and the relevant reverse
degree-one component would have to occur in the same nondegenerate pair,
where they cannot.

For the two-slope architecture in E826, its canonical torsion/HN filtration
and pairwise block Hom vanishings provide the equivariance required in
section1 for every specified finite action on E. Thus this is a theorem for
all extensions in that architecture satisfying the nondegenerate
line-difference condition, not merely for the printed E826 example.

\subsection{3. Degenerate differences in the four-divisor Gaussian
algebra}\label{SF-E827:degenerate-differences-in-the-four-divisor-gaussian-algebra}

There is a stronger conclusion for the actual divisor algebra used in our
constructions, even when some cross differences are degenerate. Write

\[
 X=Y\times Z_0=E_i^2\times E_i^2,
 \qquad V_Y=H^1(O_Y),\quad V_Z=H^1(O_{Z_0}).
 \tag{E827.7}
\]

These are the two S-matrix factors; they are \textbf{not} the irrational RM
embedding planes. In the divisor basis(theta,theta\_U,theta\_S,theta\_US), a
rational difference has two blocks xI+yU2, with U2 squared equal to2. A
nonzero such block is nondegenerate: x2-2y2=0 with rational x,y forces
x=y=0. Hence a nonzero difference is either nondegenerate on X, or is a
pullback of a nondegenerate line class from exactly one of Y,Z0. The
opposite theta slopes in E826 ensure that a cross difference is nonzero.
Arbitrary degree-zero twists do not change the following proof.

Set

\[
 W=V_Y\otimes V_Z,\qquad
 Z=\Lambda^2V_Y\oplus\Lambda^2V_Z
       \ \subset\ H^2(O_X),\qquad\dim W=4.
 \tag{E827.8}
\]

This is a direct decomposition H2(O\_X)=W plus Z. The two projections are
fixed by translations and degree-zero tensor operations on cohomology. We
verify all hypotheses of section1, including every degenerate case. The
nondegenerate case has already been proved in section2.

Suppose the difference line is pulled back from Y. It is an exterior product
L\_Y box P\_Z, up to a degree-zero twist on Y, because Pic0 of a product is
the product of the dual abelian varieties. If P\_Z is nontrivial, all its
cohomology, and all that of its inverse, vanish. Otherwise P\_Z is trivial.
Let j in\{0,1,2\} be the index of the nondegenerate surface line L\_Y.
Kunneth and the surface index theorem give

\[
 H^1(X,L_Y\boxtimes O_{Z_0})=
 \begin{cases}
 H^0(Y,L_Y)\otimes H^1(O_{Z_0}),&j=0,\\
 H^1(Y,L_Y)\otimes H^0(O_{Z_0}),&j=1,\\
 0,&j=2.
 \end{cases}
 \tag{E827.9}
\]

Multiplication by W kills each of these spaces into H3 of the line: it
raises the Y cohomological degree by one, where the nondegenerate line has
no cohomology. The same statement holds for its inverse, whose index is2-j.
Both degree-one spaces can be nonzero only when j=1. In that case their
product lies in H2(O\_Y) tensor H0(O\_Z0), the first pure summand of Z. This
proves hypotheses1 and2, with all multiplicity matrices allowed. The case
pulled back from Z0 is identical with the factors interchanged. This is an
actual cohomological product calculation, not only a count of the dimensions
of the Ext groups.

Therefore for all nonzero cross differences in this rational four-divisor
algebra,

\[
 \operatorname{ext}^2_\Gamma(E,E)\ge4(u+|J|).
 \tag{E827.10}
\]

In particular u at least5 excludes every such two-slope balanced extension,
even with degenerate differences and nonzero opposite Ext3 pairs. If one
uses selected negative blocks in J, the threshold is u+\textbar J\textbar{}
at least5. The defining condition for J is still the explicit Ext3(M\_b,A0)
vanishing; no assertion about the unit block is silently made.

\subsection{4. Balanced consequence and exact remaining
scope}\label{SF-E827:balanced-consequence-and-exact-remaining-scope}

The entire normalized balanced character, with any algebraic rational B
field and q nonzero, has degree-two cap rank18. The exterior proof is
repeated in E826: two pure blocks have rank1 each and the mixed block has
rank16. The B-shear is invertible and preserves this rank. The supplied
invariant CY4 adjunction(S21.8) therefore requires ext2\_Gamma(E,E)=18 under
weak semiregularity. Thus either lower bound exceeding18 is a contradiction
for every specified coherent effective finite action in that framework.

This is a strict enlargement of E826's family exclusion: a merely nonzero
Ext3(B,A) no longer suffices to escape, and in the displayed CM divisor
algebra neither does a degenerate cross pair. The argument still requires
the actual equivariant filtration and blocks, supplied by the proved
two-slope/torsion construction in the applications. It does not assert that
an arbitrary complex or unstable sheaf has such a filtration. When there are
too few positive blocks, selected negative blocks do not meet their
additional vanishing condition, or divisor directions leave the four-divisor
algebra, this proof leaves a genuine gate. No successful seed or general
nonexistence theorem is claimed.

\subsection{5. Source and check boundary}\label{SF-E827:source-and-check-boundary}

The abelian and surface index/vanishing inputs are the original Mumford
Abelian Varieties(1970), section16 printedpp150--153 and162--163, already
read in this run. Nontrivial Pic0 cohomology vanishing is the previously
read published Mukai1981 Theorem2.2 p156 and its base-change consequence, or
the direct flat-character argument already supplied. The Kunneth
factorization, both actual triangle boundary calculations, and all
multiplicity/invariant projections are proved above. The full S21 adjunction
and its stated source statuses are retained exactly as freshly reread for
E826; no new published status is asserted for a supplied preprint input.
BF2008's exact published evaluation/trace locators remain those fully
recorded in E826/E825.

There is no new numerical-rank premise. This proof is a symbolic
Yoneda/Kunneth argument, independently whole-reviewed as recorded above. The
existing E826 exact computation supplies an actual example in its scope but
does not replace the proof for the larger family here.

\section{SF-E828. Equal-hull free blocks are still equivariant after
averaging}\label{SF-E828:sf-e828.-equal-hull-free-blocks-are-still-equivariant-after-averaging}

Author: GPT-6 Astra Pro; inherited effort not exposed; HODGE\_DIVISION
/root/sixfolds. Date: 2026-11-21.

\textbf{Status:} independently whole-reviewed by /root and
/root/imprimitive; both mathematical reviews PASS on the frozen draft
\texttt{d57b98f4adff0567e2ee52c2866ca8d418064f79c815b86b25e14ff149552260}.
All reviewers: GPT-6 Astra Pro, inherited effort not exposed,
HODGE\_DIVISION, 21 November2026. This closes a proposed same-hull exception
to E826/E827. It does not assert that a nonsplit positive HN factor has the
direct-sum form below. No new numerical computation is a premise.

\subsection{1. Statement}\label{SF-E828:statement}

Retain the actual unit sheaf F on the abelian fourfold X, with smooth
integral nonempty S the intersection of two divisors in the ample class D=k
theta used in E826, where k is positive. Its exact sequence and dual are

\[
 0\longrightarrow i_*O_S(-2D)\longrightarrow F\longrightarrow I_S
 \longrightarrow0,\qquad F^\vee=O_X.
 \tag{E828.1}
\]

Let L be a line bundle, and put U=F tensor L. Let

\[
 A=U^{\oplus r}\oplus P,
 \qquad P=\bigoplus_{a=1}^u L_a^{\oplus n_a},
 \tag{E828.2}
\]

where r and n\_a are positive, the c1(L\_a) are pairwise distinct, and all
L\_a have the same theta slope as L. Unlike E826's originally stated
hypothesis, one L\_a is allowed to be isomorphic to L. Also allow a line of
the same Chern class as L with a nontrivial degree-zero difference; because
c1(L\_a) are pairwise distinct there is at most one such class block.

Suppose A carries the actual coherent finite translation/tensor action
induced from the intrinsic two-slope torsion/HN filtration of an extension
as in E826. Then each displayed free block is a canonical invariant
subobject of A. There is an actual equivariant direct-sum decomposition of A
into all those free blocks and a complement whose underlying sheaf is
U\^{}\{r\}. Consequently the conclusions of E826 and E827 apply with their
same Ext and cap hypotheses, without the condition c1(L\_a) different from
c1(L).

The complement is not claimed canonical. The assertion is existence of an
actual coherent equivariant splitting, obtained by averaging an existing
ordinary splitting. No new action or new isogeny is chosen.

\subsection{2. The Hom vanishings and the intrinsic line
subobjects}\label{SF-E828:the-hom-vanishings-and-the-intrinsic-line-subobjects}

First H0(F)=0. In (E828.1), H0(I\_S)=0 because a global regular function on
the connected projective X is constant and S is nonempty. Also
H0(O\_S(-2D))=0 because D restricted to S is ample and its negative has no
nonzero section. The exact sequence proves the claim.

For any L\_a of the stated slope, Hom(L\_a,U)=0. To see this, tensor
(E828.1) by L tensor L\_a inverse. If its first Chern class is nonzero, its
theta-degree on X is zero, so a nonzero section of its hull would give a
nonzero effective divisor of theta-degree zero, impossible. If its class is
zero but the line is nontrivial, it has no section: a section would give an
effective divisor of numerical class zero and would therefore be nowhere
vanishing. If the line is trivial, the ideal factor I\_S has no section as
just proved. In all cases the torsion term has theta-degree
\((-2D+c_1(L)-c_1(L_a))D^2\theta\), which is strictly negative: S has class
D\^{}2, D is the scalar multiple of theta used in E826, and the slope
difference is zero. Thus its space of sections is zero as well.

Distinct L\_a have no Hom between them, by the same theta-degree argument.
It follows that the evaluation image

\[
 \operatorname{im}\bigl(\operatorname{Hom}(L_a,A)\otimes L_a\to A\bigr)
        =L_a^{\oplus n_a}
 \tag{E828.3}
\]

is exactly the displayed block. This description is independent of choices
of basis for Hom and is an actual subsheaf description.

We verify that any functor occurring in the given action preserves the
isomorphism class of each L\_a; preserving its Chern class alone would not
suffice. Let T\_g=t\_x\^{}*(-) tensor Q, with Q degree zero, and put
M=T\_g(L\_a). The given action on A carries the displayed nonzero inclusion
L\_a to A to an inclusion M to A. The line M has the same theta slope as L.
The Hom vanishing proof above applies to every such line M, not just the
displayed L\_a: it gives Hom(M,U)=0, including when c1(M)=c1(L), with any
flat twist. Therefore the nonzero map M to A has a nonzero component into
some L\_b. Since c1(M)=c1(L\_a), the distinct-Chern-class assumption forces
b=a. A nonzero map between two lines of the same first Chern class is an
isomorphism: its effective zero divisor has theta-degree zero, hence is
empty. Thus M is isomorphic to L\_a.

This direct argument also handles c1(L\_a)=c1(L) with a nontrivial flat
difference. It neither uses Krull--Schmidt nor infers isomorphisms of
ambient lines from their restrictions to S. Its only additional datum is the
existing inclusion of each free summand in A.

Thus each functor sends the evaluation image (E828.3) onto itself. The
coherent action on A restricts to an actual coherent action on each free
block and on their sum P.

\subsection{3. Averaging the retraction and the exact
scope}\label{SF-E828:averaging-the-retraction-and-the-exact-scope}

Let j:P to A be the resulting equivariant inclusion. The displayed ordinary
direct sum (E828.2) supplies a retraction q:A to P with qj=Id\_P. The actual
actions on A and P give a linear action on Hom(A,P). Average q over the
finite group:

\[
 q_\Gamma=\frac1{|\Gamma|}\sum_{g\in\Gamma}g(q).
 \tag{E828.4}
\]

For each g, the identity qj=Id and equivariance of j give g(q)j=Id. Hence
q\_Gamma j=Id. The averaged morphism is invariant, so its kernel C is an
actual coherent equivariant complement and A=P direct sum C.

As an ordinary sheaf, C is isomorphic to U\^{}r. Indeed relative to the
original splitting, any retraction has the form (h,Id\_P) for a morphism
h:U\^{}r to P, and its kernel is the graph of -h. This also makes explicit
why possible nonzero Hom(U,L\_a) does not obstruct the splitting after
averaging. The multiplicity actions may be projective after choices of line
trivializations; the actual action on Hom(A,P) and its average are honest
because they are induced from the specified coherent action on the sheaves.

The intrinsic torsion/HN argument of E826 still supplies an action on A:
A/tors A is semistable of its common positive slope, while the quotient B
has strictly smaller negative slope. Allowing a free class equal to the unit
hull class does not change that proof. The invariant direct-sum splitting
just constructed supplies exactly the A-block hypothesis used in E826/E827.
Their actual triangle maps, selected B-block vanishings and rank18
comparison then apply without alteration.

This refines only the explicit direct-sum positive factor (E828.2). If the
positive factor itself is a nonsplit extension of U and line bundles, an
ordinary retraction need not exist and averaging cannot create one. If
additional nonisomorphic line types with the same Chern class are allowed,
the same-slope free evaluation subcategory can be preserved as an orbit, but
no assertion about individual type blocks is made here. Those are separate
remaining constructions, not cases silently covered above.

\subsection{4. Source and check boundary}\label{SF-E828:source-and-check-boundary}

The unit exact sequence and dual are the already whole-reviewed UG/US
proofs. All new subsheaf, line-type and averaging arguments are given above.
The section-vanishing argument uses only positivity of an ample divisor and
the usual global functions on a connected projective variety. E826/E827
retain their exact full supplied S21 adjunction and published BF/Mumford
source boundaries. No new external theorem, numerical rank, constructed
finite group or semiregular object is claimed.

\section{SF-E829. A wedge-pairing gate for arbitrary surface
directions}\label{SF-E829:sf-e829.-a-wedge-pairing-gate-for-arbitrary-surface-directions}

Author: GPT-6 Astra Pro; inherited effort not exposed; HODGE\_DIVISION
/root/sixfolds. Date: 2026-11-21.

\textbf{Status:} independently whole-reviewed by /root and
/root/g34\_eightfolds; both mathematical reviews PASS on the frozen draft
\texttt{4cab77bfcc30b75151e946b3a1da06e29e2461ce59f8787c326c9aa5ab666b28}.
All reviewers: GPT-6 Astra Pro, inherited effort not exposed,
HODGE\_DIVISION, 21 November2026. The peer also supplied an independently
written supplementary exact exterior certificate; the infinite theorem is
the complete proof below, not an extrapolation from finite checks. This
extends E827 beyond its rational four-divisor algebra. It is a necessary
support-incidence gate for the same actual equivariant extension
architecture. It is not a classification of all coherent sheaves, an
existence theorem for a semiregular seed, or an assertion that the remaining
elliptic/surface configurations succeed.

\subsection{1. Exact statement}\label{SF-E829:exact-statement}

Let X be a polarized abelian fourfold and retain the actual coherent
extension and actual equivariant blocks from E827:

\[
 0\longrightarrow A\longrightarrow E\longrightarrow B\longrightarrow0,
 \qquad A=A_0\oplus\bigoplus_{a=1}^u L_a^{n_a},
 \qquad B=\bigoplus_{b=1}^v M_b^{m_b}.
 \tag{E829.1}
\]

Every cross difference line N\_ab=L\_a tensor M\_b inverse is assumed to
have nonzero first Chern class. This is automatic for the opposite slopes in
E826. Write rank(N) for the complex rank of its Hermitian form, relative to
the fixed positive polarization. For every pair, require either rank(N\_ab)
at least2 or \(H^*(X,N_{ab})=H^*(X,N_{ab}^{-1})=0\). Thus a rank-one pair is
allowed only if its flat kernel twist makes it cohomologically irrelevant.
The rank-one cases with nonzero cohomology are deliberately not covered by
the theorem.

Put V=H1(X,O\_X), dim V=4. For each rank-two difference whose nondegenerate
surface part has index one and whose flat kernel part is trivial after the
product isogeny below, let V\_N be the two-dimensional subspace of V coming
from that surface factor. Let z\_N=Lambda2 V\_N, a line in Lambda2 V. Define

\[
 Z=\operatorname{span}_\mathbf C\{z_N\}\subset\Lambda^2V,
 \qquad z=\dim Z.
 \tag{E829.2}
\]

The construction below gives these subspaces canonically after choosing the
fixed polarization. Including an extra acyclic rank-two pair in Z is
permitted and only weakens the bound. Choose a nonzero volume in Lambda4 V
and let q be the nondegenerate symmetric wedge pairing on Lambda2 V. Let rho
be the dimension of the radical of q restricted to Z. Put

\[
 w=6-z-\rho.
 \tag{E829.3}
\]

Let J be any subset of negative-block indices for which Ext3(M\_b,A0)=0. For
every specified coherent effective finite action preserving the actual
sequence and blocks,

\[
 \operatorname{ext}^2_\Gamma(E,E)\ge w(u+|J|).
 \tag{E829.4}
\]

In particular, if the complete normalized character is a nonzero multiple of
the balanced v\_q of E801, with any rational algebraic B field, weak
semiregularity requires w(u+\textbar J\textbar) at most18. The letter q in
the wedge pairing and the parameter in v\_q denote different objects; in the
proof below the pairing is written q\_wedge when needed to avoid ambiguity.

\subsection{2. The product isogeny and the cohomology
calculation}\label{SF-E829:the-product-isogeny-and-the-cohomology-calculation}

We give the needed degenerate-line reduction directly, including the isogeny
and its scope. On the real tangent space of X, the alternating form E\_theta
of the polarization identifies the alternating form E\_N of a line bundle
with a rational linear operator T:

\[
 E_N(x,y)=E_\theta(Tx,y).
 \tag{E829.5}
\]

Both forms are integral on the lattice, and E\_theta is nondegenerate, so T
is rational. Their type(1,1) condition implies that T commutes with the
complex structure. Alternation gives E\_theta(Tx,y)=E\_theta(x,Ty), so T is
self-adjoint for the positive Hermitian form associated with theta. It is
therefore diagonalizable with real eigenvalues. Its kernel K\_R and image
C\_R are rational complex subspaces, and are E\_theta-orthogonal
complements. Intersecting them with the lattice produces complex subtori K
and C. The restriction of theta to each is positive, so these are abelian
subvarieties. Their sum map

\[
 f:C\times K\longrightarrow X
 \tag{E829.6}
\]

is an isogeny: the differential is an isomorphism, and the sum of the two
intersection lattices has finite index in the original lattice. This is a
lattice proof of the needed complement; no generic splitting or rational RM
eigenspace is being treated as an abelian subvariety.

The complex dimension of C is rank(N), and E\_N is nondegenerate on C and
zero on K and on the mixed terms. Thus

\[
 f^*N=N_C\boxtimes P_K,
 \tag{E829.7}
\]

where N\_C is a nondegenerate line bundle and P\_K is a degree-zero line
bundle. To see the line-level assertion, subtract the pullback of the
restriction to C: the remaining line has zero first Chern class, and a flat
character of the product lattice is a product of characters of its two
factors. Absorb its C component into N\_C. This proves(E829.7), not merely
its numerical counterpart.

A nontrivial degree-zero line on an abelian variety has no cohomology. One
may use the previously read published Pic0 vanishing, or the direct
unitary-flat proof: the corresponding flat character is nontrivial, and the
flat Laplacian on its torus has no zero-frequency harmonic mode. Hence if
P\_K is nontrivial, both lines in(E829.7) and their inverses are acyclic. If
P\_K is trivial, the abelian index theorem puts the cohomology of N\_C in
exactly one degree j, and the inverse in degree r-j, where r=dim C. Kunneth
then describes all cohomology and all its products.

The pullback of H1(O\_X) is an isomorphism onto H1(O\_C) direct sum
H1(O\_K), so it identifies V with these two summands. Define V\_N=H1(O\_C)
under this identification. It is independent of lattice indices in the above
construction because C\_R=im T is fixed. The pullback of \(H^*(X,N)\) is
injective by finite-etale trace: trace after pullback is multiplication by
deg f.~Products commute with pullback. Thus every vanishing and every
product-containment proved on C times K descends to X. This use of an
isogeny concerns cohomology and products only; no action, object, or
semiregularity is being pulled back in the conclusion.

\subsection{3. Exhaustive ranks at least
two}\label{SF-E829:exhaustive-ranks-at-least-two}

Let V\_C=H1(O\_C), V\_K=H1(O\_K), with dimensions r and k=4-r. When P\_K is
trivial, the only possible degree-one terms are

\[
 H^1(f^*N)=
 \begin{cases}
 H^0(N_C)\otimes V_K,&j=0,\\
 H^1(N_C)\otimes\mathbf C,&j=1,\\
 0,&j\ge2.
 \end{cases}
 \tag{E829.8}
\]

Multiplication by any class containing a V\_C factor raises the C
cohomological degree away from its sole index and is therefore zero. Only
the component in Lambda2 V\_K can act nontrivially from degree1 to 3. This
observation also applies to N inverse, at its index r-j.

If r=4, the line is nondegenerate on X. Its degree-one group and that of its
inverse cannot both be nonzero, and its own degree-three group is zero
whenever its degree-one group is nonzero. Thus all diagonal degree-one
products vanish and every element of Lambda2 V kills both degree-one groups
into degree3.

If r=3, then k=1 and Lambda2 V\_K=0. Every element again kills degree1 into
3. The indices of N\_C and its inverse sum to3, so both cannot have index at
most1. Thus both degree-one groups cannot occur simultaneously, and all
their diagonal products vanish. This includes the j=0 term with its
one-dimensional kernel factor; no nondegeneracy on the whole X is assumed.

If r=2 and j=0 or2, only one of the two lines can have degree-one
cohomology. The possible group has j=0 and is H0(N\_C) tensor V\_K.
Multiplication by Lambda2 V\_K lands in Lambda3 V\_K=0 because k=2. The
C-containing components already vanish. Thus all diagonal products vanish
and every element of Lambda2 V kills degree1 into3 in these cases.

The sole remaining case with r at least2 is r=2,j=1. Here both lines have

\[
 H^1(f^*N)=H^1(N_C)\otimes\mathbf C,
 \qquad
 H^3(f^*N)=H^1(N_C)\otimes\Lambda^2V_K.
 \tag{E829.9}
\]

Their diagonal degree-one Yoneda products lie in Lambda2 V\_C, the support
line z\_N. Multiplication of a degree-one class by w0 in Lambda2 V depends
only on its projection to Lambda2 V\_K. It is zero whenever that projection
vanishes. Since both V\_C and V\_K have dimension2, this condition is
exactly

\[
 w_0\wedge z_N=0,
 \tag{E829.10}
\]

meaning the wedge with any generator of that line is zero. Equivalently, w0
is orthogonal to z\_N for q\_wedge. This completes the actual cup-product
analysis of every pair allowed in the theorem.

\subsection{4. The surviving diagonal space and invariant exact
sequences}\label{SF-E829:the-surviving-diagonal-space-and-invariant-exact-sequences}

For all cross pairs together, the preceding calculations show that every
incoming diagonal degree-one product lies in Z, and every element of Z
perpendicular kills both relevant degree-one groups into degree3. The
pairing on Lambda2 V is nondegenerate. Thus

\[
 \dim Z^\perp=6-z,
 \qquad Z\cap Z^\perp=\operatorname{rad}(q_{\mathrm{wedge}}|_Z).
 \tag{E829.11}
\]

Choose any vector-space complement W to this radical inside Z perpendicular.
Then W intersects Z trivially, and dim W=6-z-rho=w. All such subspaces are
fixed by the allowed translation/tensor functors on H2(O\_X), since those
functors act trivially on this cohomology. No invariant choice beyond this
trivial action is needed.

Now apply exactly the actual diagonal-survival proof E827.1--5: the
A-diagonal W identity classes inject through Hom(A,triangle); their next
boundary vanishes because W kills Ext1 into Ext3. The selected B-diagonal
classes lift because Ext3(M\_b,A0)=0. Normalized diagonal matrix trace
modulo Z kills the preceding boundary, with its A0 term zero by Serre
duality. Finite invariants are exact and all selected classes are invariant.
Consequently their contributions add to w(u+\textbar J\textbar). This
is(E829.4). There is no degeneration claim for a filtered Ext spectral
sequence and no assumption that different Chern-character terms are summands
of E.

For the full balanced character, its cap rank is18 and the rational B shear
is invertible. The already reviewed full supplied S21 adjunction therefore
forces invariant Ext2 dimension18 under weak semiregularity.
Equation(E829.4) gives the claimed necessary inequality.

\subsection{5. Concrete geometric consequences and the elliptic
limitation}\label{SF-E829:concrete-geometric-consequences-and-the-elliptic-limitation}

With no exceptional rank-two index-one pair, Z=0 and w=6. Thus E827's
six-per-block obstruction extends to every rank-three degeneration and to
every rank-two index-zero or index-two degeneration, without assuming the
four-divisor algebra. For u at least4 the weak balanced seed is excluded
under these hypotheses, even if the unit block has nonzero opposite Ext3.

With a single exceptional support line, z=1 and rho=1 because a decomposable
two-vector squares to zero. Hence w=4. With exactly two independent support
lines generated by z1,z2, there are two alternatives:

\begin{itemize}
\tightlist
\item
  if z1 wedge z2 is nonzero, the restriction pairing has matrix
  {[}{[}0,c{]},{[}c,0{]}{]}, c nonzero, so z=2,rho=0 and w=4;
\item
  if z1 wedge z2 is zero, the restriction is zero, so z=2,rho=2 and w=2.
\end{itemize}

For two-dimensional subspaces of a four-dimensional V, the second case means
the two support planes intersect in a nonzero vector; the first means they
are complementary. These are the actual rational surface support directions
defined from their numerical line forms, not the irrational RM embedding
planes. In particular the two complementary surface directions of E827's
four-divisor algebra give w=4, recovering that earlier bound. Arbitrary
additional directions are controlled by the explicit formula6-dimZ-dimrad,
not by an unstated general-position claim.

The rank-one limitation is genuine for this method. On a product C times K
with C elliptic and K three-dimensional, take a nontrivial positive line
N\_C and trivial P\_K. Then N\_C has index0 and its inverse index1. The two
degree-one spaces are H0(N\_C) tensor V\_K and H1(N\_C inverse) tensor C. By
elliptic Serre duality their diagonal products span

\[
 Z_{\mathrm{ell}}=V_C\wedge V_K.
 \tag{E829.12}
\]

For multiplication to kill the inverse degree-one group into degree3, a
two-vector must have zero Lambda2 V\_K component. Since V\_C is
one-dimensional, the entire allowable annihilator is exactly the same space
V\_C wedge V\_K. It has zero possible complement disjoint from Z\_ell. Thus
the E827 lemma cannot extract any nonzero common W from this single pair.
This is a demonstrated limitation of the diagonal method, not a proof of
semiregularity or existence of a balanced seed. If a rank-one pair is
acyclic because of a nontrivial flat kernel factor, all its products vanish
and it causes no restriction, as permitted in the statement.

\subsection{6. Sources and check boundary}\label{SF-E829:sources-and-check-boundary}

The product isogeny and its numerical/line-level decomposition are proved
above from the rational self-adjoint operator of the two Riemann forms. The
only substantive external cohomology input is the original Mumford Abelian
Varieties(1970), section16 printedpp150--153 and162--163, already fully read
in the invoked scope for E826, together with the previously read Pic0
vanishing or the direct flat proof supplied above. Kunneth, finite trace and
the actual cup products are stated with their full uses. The actual
triangle/invariant argument is the wholly proved E827 lemma; the balanced
rank consequence retains the exact supplied S21 and published BF status
already recorded there. No new numerical experiment, period rank, finite
action or semiregular object is asserted.

\section{SF-E830. The actual extension class controls elliptic escape
directions}\label{SF-E830:sf-e830.-the-actual-extension-class-controls-elliptic-escape-directions}

Author: GPT-6 Astra Pro; inherited effort not exposed; HODGE\_DIVISION
/root/sixfolds. Date: 2026-11-21.

\textbf{Status:} independently whole-reviewed by /root and
/root/g34\_eightfolds; both mathematical reviews PASS on the frozen draft
\texttt{dbe5e856356e49ece65aa9a41cb77851a1eac411bfff5088e4090f3fa9f581a4}.
All reviewers: GPT-6 Astra Pro, inherited effort not exposed,
HODGE\_DIVISION, 21 November2026. The peer also supplied an independently
written supplementary exact exterior certificate; the infinite theorem is
the complete proof below, not an extrapolation from finite checks. This
replaces E827's hypotheses on every degree-one class by hypotheses on the
entries of the actual extension class. It includes the rank-one directions
deliberately outside E829, but gives a necessary rank gate rather than
general nonexistence. An actual full-rank example is constructed at the end
and is explicitly not a balanced RM seed.

\subsection{1. The extension-specific diagonal
lemma}\label{SF-E830:the-extension-specific-diagonal-lemma}

Let X be an abelian fourfold and consider the same actual coherent extension
with actual equivariant block decompositions:

\[
 0\to A\xrightarrow{i}E\xrightarrow{p}B\to0,
 \quad e:B\to A[1],\quad
 A=A_0\oplus\bigoplus_a L_a^{n_a},\quad
 B=\bigoplus_b M_b^{m_b}.
 \tag{E830.1}
\]

Write N\_ab=L\_a tensor M\_b inverse. The matrix of the line-to-line
component e\_ab has entries in H1(N\_ab). Let E\_ab be the span of all those
entries; equivalently it is the image obtained by contracting e\_ab with the
dual of its multiplicity-matrix space. This definition does not depend on
bases of the two multiplicity spaces or on scalar choices identifying the
lines. It records the particular extension, not the whole Ext1 group.

Let W,Z be subspaces of H2(O\_X) with W intersect Z zero. Assume for every
pair:

\begin{itemize}
\tightlist
\item
  the products of E\_ab with all of H1(N\_ab inverse), in either order,
  belong to Z;
\item
  multiplication by W kills E\_ab into H3(N\_ab).
\end{itemize}

There is deliberately no hypothesis about multiplication of W on the whole
reverse group H1(N\_ab inverse). Let J consist of any B indices for which
Ext3(M\_b,A0)=0. Then

\[
 \operatorname{ext}^2_\Gamma(E,E)\ge(\dim W)(u+|J|).
 \tag{E830.2}
\]

Here u is the number of positive line blocks, and the finite coherent action
is required to preserve the actual sequence and blocks, as in E827. No such
decomposition of the middle object is asserted.

We prove the altered hypothesis suffices, retaining every boundary map. The
diagonal identity subspace D\_A=direct sum\_a W Id\_n in Ext2(A,A) is fixed
by the action. The kernel of Ext2(A,A) to Ext2(A,E) comes from Ext1(A,B).
Its a-th diagonal is the sum of e\_ab{[}1{]} gamma\_ba. Every coefficient
contains an entry of the actual e\_ab and a reverse degree-one class, and is
therefore in Z. Thus D\_A injects. Its outgoing boundary in Hom(-,E) is
i{[}3{]}d{[}1{]}e. Every line-target component is multiplication of an entry
of e\_ab by the corresponding W identity class. The second hypothesis kills
exactly these entries; no other reverse Ext1 class occurs in this boundary.
Consequently the image I of Ext2(E,E) in Ext2(A,E) contains D\_A.

For the B-side, take D\_B=direct sum\_b in J W Id\_m inside Ext2(B,B). Its
boundary to Ext3(B,A) is e{[}2{]}h. Again only entries of the actual e\_ab
are multiplied by W. They vanish by the second hypothesis, while the A0
component vanishes by the defining condition on J. Thus D\_B lifts to
Ext2(B,E).

The preceding boundary from Ext1(A,E), projected to Ext2(B,B), is
(p{[}1{]}alpha){[}1{]}e. Its selected b-diagonal is a sum of reverse
degree-one classes multiplied by entries of e\_ab. All such coefficients lie
in Z. The possible A0 term is zero because Ext1(A0,M\_b) is dual to
Ext3(M\_b,A0). Normalized diagonal matrix trace modulo Z therefore kills
this boundary and detects the lifted W identities. It follows that the
quotient Q of Ext2(B,E) by this boundary has dimension at least
dim(W)\textbar J\textbar{} and injects into Ext2(E,E). Its quotient there is
I, which already contains D\_A. These contributions add.

All the actual maps are equivariant. H2(O\_X) is fixed under translations
and degree-zero tensor functors, and conjugation fixes identity matrices and
normalized trace, including projective multiplicity descriptions. Invariant
lifts are obtained by averaging; taking invariants is exact. The same
dimension sum proves(E830.2) in the specified finite-action scope. This is
an actual triangle calculation, with no formality or higher-product
vanishing assumption.

\subsection{2. A support space for the actual entries, in every positive
rank}\label{SF-E830:a-support-space-for-the-actual-entries-in-every-positive-rank}

For each difference with nonzero first Chern class, use the Riemann-form
product isogeny and cohomology identification proved in E829:

\[
 f:C\times K\to X,\qquad f^*N=N_C\boxtimes P_K,
 \quad V=H^1(O_X)=V_C\oplus V_K.
 \tag{E830.3}
\]

The complement is fixed by the polarization. Pullback on the line cohomology
is injective and products commute with pullback. If P\_K is nontrivial, all
the groups vanish; also omit a pair whose actual component e\_ab is zero.
For the remaining cases define a subspace Z\_ab of Lambda2 V as follows.

For rank(N) at least3, put Z\_ab=0. For rank2 with surface index0 or2, also
put Z\_ab=0. The complete calculations in E829 show in these cases that the
products of all degree-one classes with their reverses are zero and all
H2(O\_X) kills degree1 into3. Thus these choices satisfy the new
entry-specific conditions a fortiori.

For rank2 with surface index1, put

\[
 Z_{ab}=\Lambda^2V_C.
 \tag{E830.4}
\]

The degree-one product lies in this line. Any w0 orthogonal to this line
under the four-dimensional wedge pairing has zero Lambda2 V\_K component, so
it kills every pulled-back entry of e\_ab into degree3. This is the same
surface calculation as E829, now ignoring zero actual components.

For rank1 with elliptic index0, the pulled-back entries lie in H0(N\_C)
tensor V\_K, with dim V\_K=3. Define S\_ab to be their coefficient span in
V\_K:

\[
 S_{ab}=\operatorname{im}\left(
   H^0(N_C)^*\otimes E_{ab}\longrightarrow V_K\right),
 \qquad Z_{ab}=V_C\wedge S_{ab}.
 \tag{E830.5}
\]

The displayed contraction includes all matrix entries of e\_ab. If the
pullback is not onto all line cohomology, only the pulled-back entries are
used. Thus no unproved surjectivity of an isogeny pullback is required. The
span is basis independent. Degree-one classes of the inverse line pull back
to H1(N\_C inverse) tensor C. Their products with the entries of e\_ab lie
in V\_C wedge S\_ab, by the elliptic Serre pairing. If w0 is orthogonal to
Z\_ab under wedge pairing, then

\[
 (\operatorname{pr}_{\Lambda^2V_K} w_0)\wedge s=0
       \quad\hbox{for every }s\in S_{ab}.
 \tag{E830.6}
\]

This is precisely the condition that multiplication by w0 kills the
pulled-back entries H0(N\_C) tensor S\_ab into H0(N\_C) tensor Lambda3 V\_K.
All terms of w0 containing V\_C raise the C index and vanish. Finite trace
then gives both the product containment and the required entry vanishing on
X itself.

For rank1 with elliptic index1, put

\[
 Z_{ab}=V_C\wedge V_K.
 \tag{E830.7}
\]

Here the entries lie in H1(N\_C) tensor C, and their reverse degree-one
classes lie in H0(N\_C inverse) tensor V\_K. Products lie in the displayed
three-space. That space is maximal isotropic for the wedge pairing and is
its own perpendicular. Consequently its perpendicular has zero Lambda2 V\_K
component and kills the entries into degree3. This case may contribute no
surviving space, but it still satisfies the theorem. In the opposite-slope
architecture a rank-one difference has positive polarization degree and
therefore elliptic index0; index1 is included only for the stated general
extension lemma.

These computations exhaust every nonzero numerical difference. A numerically
trivial but nontrivial flat difference is acyclic and can also be ignored.
For completeness, if N is the trivial line, put Z\_ab=E\_ab wedge V, where
E\_ab is now a subspace of V. Then the product containment is literal
exterior multiplication. Orthogonality to Z\_ab means w0 wedge xi wedge v=0
for every xi in E\_ab and v in V, equivalently w0 wedge xi=0 in Lambda3 V.
It therefore kills those actual entries. A nonzero trivial-line extension
may again force a zero bound; no opposite-slope application relies on this
case.

\subsection{3. The uniform wedge formula and its elliptic rank
consequence}\label{SF-E830:the-uniform-wedge-formula-and-its-elliptic-rank-consequence}

Set

\[
 Z_e=\sum_{a,b}Z_{ab},\qquad z_e=\dim Z_e,\qquad
 \rho_e=\dim\operatorname{rad}(q_{\mathrm{wedge}}|_{Z_e}),
 \quad w_e=6-z_e-\rho_e.
 \tag{E830.8}
\]

The product calculations show that all incoming products containing the
actual entries lie in Z\_e, and Z\_e perpendicular kills all those entries
into degree3. Choose a complement W to Z\_e intersect Z\_e perpendicular
inside Z\_e perpendicular. Then W intersects Z\_e trivially and dim W=w\_e.
All these cohomological subspaces are fixed under the allowed functors,
which act trivially on H2(O\_X). This does not require a choice of
individual linearizations of the line bundles. Applying section1 gives

\[
 \operatorname{ext}^2_\Gamma(E,E)\ge w_e(u+|J|).
 \tag{E830.9}
\]

Under the full balanced character and the same supplied invariant CY4
adjunction as E826/E827, weak semiregularity therefore requires
w\_e(u+\textbar J\textbar) at most18. The rational algebraic B field acts by
the already proved invertible shear and changes none of these cap
dimensions.

Suppose all the nonzero exceptional components are positive elliptic
components along a single fixed elliptic support direction V\_C; every other
pair is one of the harmless cases above. Let S be the sum of their
coefficient spans in the common three-dimensional V\_K, with r\_e=dim S.
Then Z\_e=V\_C wedge S is totally isotropic, of dimension r\_e, and

\[
 w_e=6-2r_e.
 \tag{E830.10}
\]

Thus ranks0,1,2,3 give respectively6,4,2,0 surviving classes per selected
block. In particular, if u+\textbar J\textbar{} is at least10, weak balanced
semiregularity in this precise situation requires r\_e=3. If
u+\textbar J\textbar{} is between5 and9, it requires r\_e at least2. If
u+\textbar J\textbar{} is4, it at least requires a nonzero elliptic
component. These are necessary statements about the actual extension
matrices; choosing a possible Ext1 space of large dimension does not meet
them by itself.

\subsection{4. An actual full-rank elliptic extension, with its exact
limitation}\label{SF-E830:an-actual-full-rank-elliptic-extension-with-its-exact-limitation}

The last rank condition can be realized by an actual ordinary gluable
object. This example records the boundary of the method and is not a
balanced RM construction. Take X=C times K, with C an elliptic curve and K
any abelian threefold. Let N\_C have degree3 and put N=p\_C\^{}*N\_C.
Riemann--Roch and the elliptic index theorem give h0(N\_C)=3. Choose a basis
s1,s2,s3 of H0(N\_C), and a basis eta1,eta2,eta3 of H1(O\_K). Then

\[
 e=s_1\otimes\eta_1+s_2\otimes\eta_2+s_3\otimes\eta_3
       \in\operatorname{Ext}^1(O_X,N)
 \tag{E830.11}
\]

defines an actual locally free extension

\[
 0\longrightarrow N\longrightarrow E_e\longrightarrow O_X
 \longrightarrow0.
 \tag{E830.12}
\]

Local freeness follows because the quotient is free over every local ring,
so the sequence splits there. Equivalently an actual additive Cech cocycle
for e gives upper triangular transition matrices. Their strictly upper
blocks square to zero, so no additional Maurer--Cartan condition is needed.
The object is coherent and hence has no negative self-Ext. Its coefficient
span S is all of H1(O\_K), so r\_e=3 exactly.

It is indecomposable. Hom(N,O\_X)=0, so every endomorphism preserves the
subline N. Restrictions to N and O\_X are scalars a,b with ae=eb; e nonzero
forces a=b. The kernel is Hom(O\_X,N)=H0(N\_C), a square-zero ideal, and all
common scalar identities lift. Thus End(E\_e) has quotient C by a
square-zero ideal of dimension3 and has no nontrivial idempotent, by the
same direct idempotent calculation used in E826. No simplicity is claimed.

Since c1(N)\^{}2=0, its entire character is

\[
 \operatorname{ch}(E_e)=1+e^{c_1(N)}
       =2e^{c_1(N)/2}.
 \tag{E830.13}
\]

It is a pure exponential character, not the balanced RM character with cap
rank18. No weak semiregularity or finite action is asserted for this
example. Its purpose is only to demonstrate that the required full-rank
elliptic component is attainable by an actual indecomposable extension; the
numerical and categorical balanced seed gate remains open.

\subsection{5. Sources and check boundary}\label{SF-E830:sources-and-check-boundary}

The exact triangle proof in section1 is given in full and differs from E827
only in following the actual entries instead of the whole Ext1 group. The
rational Riemann-form product isogeny and finite-trace product descent are
the full E829 proof. The sole new elliptic ingredients are the original
Mumford index theorem, ordinary curve Serre duality and elliptic
Riemann--Roch, used in the elementary forms written above. Previously read
published BF and supplied S21 retain their exact original scopes. No
uncomputed generic rank, dg formality, constructed finite action or
semiregular seed is assumed. This is a symbolic proof with a direct actual
extension example, not a numerical computation claim.

\section{SF-E831. Every block must meet enough actual exceptional
support}\label{SF-E831:sf-e831.-every-block-must-meet-enough-actual-exceptional-support}

Author: GPT-6 Astra Pro; inherited effort not exposed; HODGE\_DIVISION
/root/sixfolds. Date: 2026-11-21.

\textbf{Status:} independently whole-reviewed by /root and
/root/g34\_eightfolds; both mathematical reviews PASS on the frozen draft
\texttt{870daf7349d51f72253507005b5048dbf6e3e9e3232d9cbc55175d85d50eed45}.
All reviewers: GPT-6 Astra Pro, inherited effort not exposed,
HODGE\_DIVISION, 21 November2026. Both new dependencies E829/E830 have also
received complete independent whole-proof reviews. The peer supplied a
separate exact exterior/incidence supplement. This is a strict blockwise
strengthening of their common-space bounds. It supplies necessary incidence
and rank conditions on an actual balanced extension; it constructs no
balanced semiregular object.

\subsection{1. Separate support spaces at the actual
vertices}\label{SF-E831:separate-support-spaces-at-the-actual-vertices}

Use the actual coherent extension and its specified coherent finite action
from E830, with actual equivariant blocks

\[
 0\to A\xrightarrow{i}E\xrightarrow{p}B\to0,
 \quad A=A_0\oplus\bigoplus_{a=1}^u L_a^{n_a},\quad
 B=\bigoplus_{b=1}^v M_b^{m_b}.
 \tag{E831.1}
\]

For each actual matrix component e\_ab, choose the subspace Z\_ab from the
full E830 support table. All that is required in the proof is the two
properties proved there: its diagonal products with reverse degree-one
classes lie in Z\_ab, and its entries are killed into degree3 by every
two-vector in Z\_ab perpendicular for the wedge pairing. Put Z\_ab=0 when
e\_ab is zero. Let J be any subset with Ext3(M\_b,A0)=0. Define

\[
 Z_a=\sum_bZ_{ab},\qquad Z_b=\sum_aZ_{ab}\ (b\in J),
 \quad w_v=6-\dim Z_v-
       \dim\operatorname{rad}(q_{\mathrm{wedge}}|_{Z_v}).
 \tag{E831.2}
\]

Here v denotes one selected vertex: every positive line block, or a negative
block in J. There is no identification of Z\_a and Z\_b on different
vertices. Then

\[
 \operatorname{ext}^2_\Gamma(E,E)
      \ge\sum_{a=1}^u w_a+\sum_{b\in J}w_b.
 \tag{E831.3}
\]

This statement depends on the actual extension entries, and not merely on
which Ext1 spaces could be nonzero. A component set to zero does not remove
any diagonal classes in this bound.

\subsection{2. Full proof through the two actual boundary
maps}\label{SF-E831:full-proof-through-the-two-actual-boundary-maps}

For each selected vertex choose W\_v inside Z\_v perpendicular complementary
to its intersection with Z\_v. It has dimension w\_v and intersects Z\_v
trivially. Because the allowed functors act trivially on H2(O\_X), these
choices are automatically invariant. Take

\[
 D_A=\bigoplus_a W_a\operatorname{Id}_{n_a}
           \subset\operatorname{Ext}^2(A,A).
 \tag{E831.4}
\]

The a-th diagonal coefficient of the image of gamma in Ext1(A,B) is sum\_b
e\_ab{[}1{]} gamma\_ba. It lies in Z\_a tensor End(C\^{}\{n\_a\}). Thus that
image has zero intersection with D\_A, separately on each diagonal. The map
Ext2(A,A) to Ext2(A,E) is consequently injective on D\_A.

On d in D\_A, the outgoing boundary from Ext2(A,E) in Hom(-,E) is
i{[}3{]}d{[}1{]}e. Its (a,b) line component is multiplication of the actual
entry e\_ab by the scalar class d\_a in W\_a. Since W\_a lies in Z\_a
perpendicular and Z\_ab is contained in Z\_a, this class is in Z\_ab
perpendicular and kills the entry. The A0 target is zero because d has zero
A0 block. Thus all of D\_A lies in the actual image I of Ext2(E,E) in
Ext2(A,E).

For the selected negative vertices, take

\[
 D_B=\bigoplus_{b\in J}W_b\operatorname{Id}_{m_b}
           \subset\operatorname{Ext}^2(B,B).
 \tag{E831.5}
\]

The boundary to Ext3(B,A) is e{[}2{]}h. For b in J, its line-target
components are killed because W\_b lies in every incident Z\_ab
perpendicular. Its A0 component vanishes by Ext3(M\_b,A0)=0. Hence D\_B
lifts into Ext2(B,E).

Project the preceding boundary from Ext1(A,E) to Ext2(B,B). Its b-diagonal
is the sum of reverse degree-one classes multiplied by the actual e\_ab
entries, with the specified multiplicity-matrix compositions. The line terms
belong to Z\_b tensor End(C\^{}\{m\_b\}). The A0 term is zero by Serre
duality from the same selected vanishing. For each b independently,
normalized matrix trace modulo Z\_b kills this boundary and is the identity
on W\_b times the identity matrix. Therefore the quotient

\[
 Q=\operatorname{Ext}^2(B,E)\big/
 \operatorname{im}\bigl(\operatorname{Ext}^1(A,E)
                   \to\operatorname{Ext}^2(B,E)\bigr)
 \tag{E831.6}
\]

has dimension at least sum\_b in J w\_b and injects into Ext2(E,E). Its
quotient there is I, which already contains D\_A. The two bounds add. Every
boundary, projection and trace is equivariant. The identity multiplicity
matrices are fixed even if their chosen descriptions are projective. Finite
invariants are exact, and the invariant diagonal lifts exist by averaging.
Thus the same argument proves(E831.3) after invariants. No spectral-sequence
degeneration or vanishing of uncomputed higher products is used.

\subsection{3. A necessary coverage condition and a weighted edge
budget}\label{SF-E831:a-necessary-coverage-condition-and-a-weighted-edge-budget}

Suppose the full normalized character is a nonzero multiple of balanced
v\_q, with any rational algebraic B field. The reviewed cap rank and full
supplied invariant CY4 adjunction require invariant Ext2 dimension18 under
weak semiregularity. Put

\[
 N=u+|J|,\qquad
 \ell_v=\dim Z_v+
       \dim\operatorname{rad}(q_{\mathrm{wedge}}|_{Z_v})=6-w_v.
 \tag{E831.7}
\]

The exact necessary inequality is

\[
 \sum_v\ell_v\ge6N-18.
 \tag{E831.8}
\]

In particular a selected vertex with Z\_v=0 contributes six surviving
classes. At most three selected vertices can have no nonzero exceptional
component. Thus at least N-3 selected vertices must actually meet a pair
with Z\_ab nonzero. A merely available but unused Ext1 class does not
qualify. This remains true when other vertices have full rank-three elliptic
components: they cannot erase classes at an unrelated vertex.

For an explicit coarse budget put r\_ab=dim Z\_ab. Since rad(q\textbar Z\_v)
is a subspace of Z\_v,

\[
 \ell_v\le2\dim Z_v\le2\sum_{ab\text{ incident to }v}r_{ab}.
 \tag{E831.9}
\]

Each pair is incident to its positive vertex, and also to its negative
vertex exactly when b belongs to J. Consequently

\[
 \sum_{a,b}(1+\mathbf1_{b\in J})r_{ab}\ge3N-9
 \tag{E831.10}
\]

is another necessary condition for weak balance. The precise vertexwise
bound(E831.8) is stronger and remains available whenever the actual
subspaces can be computed.

In the opposite-slope setting, numerical cross differences are nonzero. The
table in E830 then has r\_ab at most3: a surface index-one component has
r\_ab=1, a positive elliptic component has r\_ab equal to its actual kernel
coefficient rank, and all harmless pairs have r\_ab=0. An active pair meets
at most two selected vertices. If N exceeds3, at least ceil((N-3)/2) active
exceptional pairs are necessary, even allowing rank-three elliptic entries.
More precisely the one- or two-endpoint weight in(E831.10) should be used
when some negative vertices are not in J. If all active exceptional pairs
are surface components, or elliptic components of coefficient rank at most1,
then r\_ab at most1 and the number of pairs must be at least ceil(3(N-3)/2).
These bounds are stated only for N greater than3; for smaller N they give no
additional condition.

For the32 distinct line blocks in the E826 example, if all those negative
blocks remain selected, the necessary coarse budgets would be at least15
active maximal elliptic pairs, or at least44 pairs with dim Z\_ab=1. The
original E826 pair differences are nondegenerate and have no such support.
These integers are consequences of the general inequality, not an assertion
that the proposed exceptional pairs can be realized while preserving its
full character or its selected unit vanishings.

\subsection{4. Correlated surface directions needed to remove every local
class}\label{SF-E831:correlated-surface-directions-needed-to-remove-every-local-class}

There is also a useful exact incidence test when dim Z\_v is at most3. Since
its radical has dimension at most dim Z\_v, w\_v=0 implies dim Z\_v at
least3. Under this dimension bound it is therefore equivalent to dim Z\_v=3
and q\_wedge vanishing identically on Z\_v. Such a three-space is maximal
isotropic in the six-dimensional wedge space.

For completeness its elementary classification is as follows. Every nonzero
element z in it is decomposable, because on a four-dimensional space z wedge
z=0 is equivalent to skew rank at most2. Choose two independent elements a
wedge b and a wedge c: their support planes intersect because their mutual
wedge is zero, and independence makes a,b,c linearly independent. If every
support plane contains a, the three-space equals a wedge V. Otherwise choose
a third support plane not containing a. Its intersection lines with
span(a,b) and span(a,c) are distinct and span it, so it lies in
H=span(a,b,c). Its two-vector has a nonzero b wedge c coefficient, and the
whole three-space equals Lambda2 H. Thus the only possibilities are

\[
 Z_v=a\wedge V\quad(\dim\mathbf Ca=1),
 \qquad\hbox{or}\qquad
 Z_v=\Lambda^2H\quad(\dim H=3).
 \tag{E831.11}
\]

Conversely both displayed spaces are three-dimensional and totally
isotropic, so both give w\_v=0. This is a complete linear-algebra proof; no
unproved general-position assertion is used.

In the surface-only case, if at most three independent active support lines
meet a vertex, removing all six diagonal directions therefore requires their
support planes to be of this correlated type: a common one-dimensional
direction, or containment in a common three-dimensional direction. Two
complementary surface planes do not suffice, since their wedge pairing is
nonzero and they leave w\_v=4. If more than three independent support lines
meet the vertex, use(E831.2); the above three-space classification is not
silently extended to that case. These are cohomological support planes from
actual degenerate line forms, not irrational RM eigenspaces. For elliptic
entry coefficients a support subspace may be complex rather than rational;
no extra rationality is inferred from the classification.

\subsection{5. Scope and verification
boundary}\label{SF-E831:scope-and-verification-boundary}

This result applies only when the actual equivariant filtration and line
blocks are present. E826 and E828 explain when their intrinsic HN/evaluation
arguments provide them; a genuinely nonsplit positive factor is still not a
direct sum by assertion. The selected negative vertices continue to require
the exact Ext3(M\_b,A0) vanishing. No numerical count substitutes for that
condition.

The full proof uses the actual-entry support table of E830 and its explicit
finite-isogeny cup-product calculation from E829. The finite-action,
published BF and supplied S21 statuses are exactly those already stated
there. There is no new external classification premise. The exact script
rm\_vertexwise\_wedge\_certificate.py actually passed nine wedge-space
tests, a separately labelled hypothetical six-vertex incidence test, and
all97 closed-form budgets for N=4 through100. Runtime was0.0258956120015s.
The initial output-serialization failure for exact SymPy integers is
retained; only an exact-integer JSON encoder was added. These finite checks
supplement the infinite proof above and do not assert geometric
realizability of the hypothetical incidence data. No semiregular seed,
period-map rank or generic extension-product rank is claimed.

\section{SF-E832. A positive elliptic star cannot compensate the
balanced
character}\label{SF-E832:sf-e832.-a-positive-elliptic-star-cannot-compensate-the-balanced-character}

Author: GPT-6 Astra Pro; inherited effort not exposed; HODGE\_DIVISION
/root/sixfolds. Date: 2026-11-21.

\textbf{Status:} independently whole-reviewed by /root and
/root/g34\_eightfolds; both mathematical reviews PASS on frozen draft
\texttt{0542dce98bcf06689930578a262dd242b43cfbc0a3f2aebaec263222b38d01ad}.
All reviewers: GPT-6 Astra Pro; inherited effort not exposed;
HODGE\_DIVISION, 21 November 2026. The peer supplied a separately
authored exact supplement; its scope is finite algebra, with the
universal proofs reviewed in full. This is a character-level obstruction
and does not require the free terms to be direct summands. Its strongest
version uses the actual necessary determinant condition of E825. Scalar
and general adjusted unit directions are distinguished explicitly. No
numerical search or assumed semiregularity is a premise of the
linear-algebra lemmas.

\subsection{1. Hermitian interpretation of the balanced second
character}\label{SF-E832:hermitian-interpretation-of-the-balanced-second-character}

Let the complex tangent space have the orthogonal rank-two RM splitting
V=V\_+ direct sum V\_-, with positive polarization
theta=theta\_++theta\_-. Normalize theta to the identity Hermitian form.
For any real (1,1) class a, let A be its Hermitian matrix. The Hermitian
operator represented by a²/2 on Lambda2 V is the second compound matrix
Lambda2 A. This follows by expanding a in a unitary eigenbasis: the six
coefficients are the pair products of the four real eigenvalues. The
identification is invariant under unitary changes of basis. In
particular theta²/2 corresponds to the identity, and

\[
 (\theta_+^2+\theta_-^2)/2\quad\longleftrightarrow\quad
 \Pi=\operatorname{pr}_{\Lambda^2V_+}
                +\operatorname{pr}_{\Lambda^2V_-},
 \qquad\operatorname{rank}\Pi=2.
 \tag{E832.1}
\]

All classes here are represented by translation-invariant forms on the
abelian variety. Equality in its exterior-algebra cohomology is equality
of these invariant forms, so this compound-matrix test applies to the
cohomological identities below.

\subsection{2. Scalar-unit lemma, including the zero
scalar}\label{SF-E832:scalar-unit-lemma-including-the-zero-scalar}

No real Hermitian form b can satisfy

\[
 b^2=c\theta^2+d(\theta_+^2+\theta_-^2),
 \qquad c\ge0,\quad d\ne0.
 \tag{E832.2}
\]

If c=0, the compound matrix Lambda2 B would have rank2. If B has rank k,
its compound matrix has rank binomial(k,2), one of0,1,3,6 for k from0
to4. Rank2 is impossible.

If c is positive, let lambda1,\ldots,lambda4 be the real eigenvalues of
B. Four of their six pair products equal c; the other two equal c+d,
which is distinct from c.~Mark the four edges on the four labelled
vertices whose pair products equal c.~A graph with four edges on four
vertices is connected: a disconnected graph on these vertices has at
most three edges. Thus all eigenvalues are nonzero and have the same
sign.

If the graph contains a triangle, its three equations force all three
vertices to have the same eigenvalue, equal to one of plus or minus sqrt
c. Connectedness makes the fourth one equal too. All six products would
be c, a contradiction. If there is no triangle, the four-edge graph is a
four-cycle. Its eigenvalues alternate a,b,a,b, with ab=c positive. The
two missing products are a² and b²; equality to the same c+d forces
abs(a)=abs(b). They have the same sign, so a=b and again all products
are c.~This proves(E832.2) in all cases. It is a spectral argument and
does not assume that B preserves the RM planes.

\subsection{3. A stronger semipositive lemma for a general adjusted unit
direction}\label{SF-E832:a-stronger-semipositive-lemma-for-a-general-adjusted-unit-direction}

Let u be a real Hermitian form preserving V\_+,V\_-, and suppose its
restrictions U,V have strictly positive traces. Let b be a semipositive
Hermitian form. For any positive kappa and any nonzero real d, the
identity

\[
 b^2-\kappa u^2=d(\theta_+^2+\theta_-^2)
 \tag{E832.3}
\]

is impossible. Unlike section2, u need not be scalar or initially
assumed positive.

By(E832.1), this equation is Lambda2 B=kappa Lambda2 diag(U,V)+d Pi.
Since B is semipositive, Lambda2 B is semipositive. Its principal block
on V\_+ tensor V\_- is kappa U tensor V. Both U and V have a positive
eigenvalue, by their positive traces. The semipositivity of their tensor
product then implies that every eigenvalue of each is nonnegative:
multiply any putative negative eigenvalue of one by a positive
eigenvalue of the other. Hence U and V are semipositive and nonzero, and
each has rank1 or2.

Suppose at least one has rank1. The corresponding pure Lambda2 block on
the right side is exactly d.~Semipositivity and d nonzero force
d\textgreater0. The other pure block is kappa times a nonnegative
determinant plus d, so it is positive as well. Thus both lines Lambda2
V\_+ and Lambda2 V\_- are contained in the image of Lambda2 B. But that
image is Lambda2 im B. A nonzero decomposable two-vector belonging to
Lambda2 im B has its whole support plane in im B. Consequently im B
contains both V\_+ and V\_-, and B has rank4. Its compound matrix would
have rank6. In contrast, the right side has rank rank(U)rank(V)+2, equal
to3 or4 in this case. This is a contradiction.

Both U,V must therefore be positive definite. The mixed block has rank4,
so rank(Lambda2 B) is at least4. The possible compound ranks listed in
section2 force rank B=4; semipositivity makes B positive definite. Both
pure lines are preserved nontrivially by Lambda2 B. An invertible linear
map whose compound preserves the line Lambda2 V\_+ must preserve the
support plane V\_+ itself; similarly for V\_-. Hence B is block
diagonal, with positive definite blocks C,D.

On the mixed block the equation gives

\[
 C\otimes D=\kappa U\otimes V.
 \tag{E832.4}
\]

Equality of nonzero simple tensors implies C=tU and D=(kappa/t)V for
some t\textgreater0. For an explicit proof, apply a linear functional
nonzero on U to the first factors: its value on C must be nonzero, and D
is a scalar multiple of V. Substitution gives the matching scalar
multiple for C; positive definiteness forces that scalar to be positive.
On the two pure lines the equation now reads

\[
 (t^2-\kappa)\det U=d,
 \qquad(\kappa^2/t^2-\kappa)\det V=d.
 \tag{E832.5}
\]

For t² different from kappa, these two left sides have opposite signs.
If t² equals kappa, both are zero. Neither possibility equals the same
nonzero d.~This proves the semipositive lemma.

\subsection{4. First two Chern terms of a
star}\label{SF-E832:first-two-chern-terms-of-a-star}

The actual unit sheaf F has the previously proved resolution character

\[
 \operatorname{ch}(F)=2e^{-D}-2e^{-3D}+e^{-4D}.
 \tag{E832.6}
\]

Its rank is1, its first Chern character is zero, and its second Chern
character is zero: the coefficients are respectively2-2+1, -2+6-4,
and1-9+8. Consequently a multiplicity-r unit term F tensor a line with
adjusted real first Chern class u contributes r exp(u) through
cohomological degree4, even after a chosen rational algebraic B shift.

Suppose the remaining free line terms have positive multiplicities n\_j,
with total m\textgreater0, and their adjusted classes are alpha+eta\_j,
where eta\_j²=0. Set R=r+m, with r\textgreater0, and H=sum\_j n\_j
eta\_j. Their full character is

\[
 \sum_jn_je^{\alpha+\eta_j}=e^\alpha(m+H).
 \tag{E832.7}
\]

This identity uses eta\_j²=0 exactly; it is not a truncation. It is
valid for any sheaf or complex with this K0 identity, regardless of
whether the terms occur as actual summands or through nonsplit
extensions.

If its full normalized character were R v\_q with q nonzero, its first
two nonconstant characters would be

\[
 ru+m\alpha+H=0,\qquad
 \operatorname{ch}_2=-\frac{Rq}{2}(\theta_+^2+\theta_-^2).
 \tag{E832.8}
\]

Expanding(E832.6)--(E832.8), eliminating H and completing the square
gives

\[
 b=\alpha+\frac r m u,\qquad
 b^2-\kappa u^2=d(\theta_+^2+\theta_-^2),
 \quad\kappa=\frac{rR}{m^2}>0,\quad d=\frac{Rq}{m}\ne0.
 \tag{E832.9}
\]

For clarity, before division by m the equation is m alpha²+2r alpha u-r
u²=Rq(theta\_+²+theta\_-²); thus no denominator or sign is suppressed in
completing the square.

If u=lambda theta is scalar, section2 applies with c=kappa lambda²,
including lambda=0. Thus \textbf{no square-zero star with scalar
adjusted unit direction can have the balanced character}, even allowing
both signs for the eta\_j. This part requires no positivity of the
individual differences and no use of semiregularity or E825.

\subsection{5. Positive elliptic stars with arbitrary rational B
field}\label{SF-E832:positive-elliptic-stars-with-arbitrary-rational-b-field}

Now require eta\_j=-P\_j, where each P\_j is semipositive of rank at
most1. Zero is allowed, representing a free line at the center alpha.
Equation (E832.8) gives

\[
 b=\alpha+\frac r m u=\frac1m\sum_j n_jP_j\ \ge0.
 \tag{E832.10}
\]

This is the additional positivity absent from an arbitrary square-zero
star. If u preserves the RM planes and has positive trace on both, the
semipositive lemma of section3 contradicts(E832.9). Therefore no such
positive star can have the balanced first two Chern terms under that
adjusted-unit condition.

Here is its exact application to the remaining unit/free architecture.
Take an actual exact sequence with
\(A=(F\otimes L)^{\oplus r}\oplus L_\alpha^{\oplus n_0}\) and
\(B=\bigoplus_j M_j^{\oplus n_j}\), where n0 is nonnegative, m=n0+sum\_j
n\_j is positive, and c1(M\_j)=alpha-P\_j. In the opposite-slope
application A is the positive factor and the M\_j are negative free
lines. Require every nonzero difference P\_j to have rank1 and positive
theta-degree. Rank1 plus positive degree means P\_j is semipositive: it
has one positive eigenvalue and the rest zero. Include the positive free
terms with P\_j=0. No condition on the ranks of the actual extension
matrices is needed for this argument.

Because B is locally free, the exact sequence splits locally. Its local
form is F\^{}r plus a free module. Relative to these local splittings,
the FF diagonal transitions are exactly those of F tensor L; the
additional extension entries go from the free B part into A. After the
fixed F comparisons, the FF transition is the scalar line transition of
L to multiplicity r. E825's actual ring map ell ignores the free
diagonal and all off-diagonal free-to-unit or free-to-free entries,
including their products through a free module. Thus its exact
multiplicative transition cocycle is that of L\^{}r on S2. Thus c1(det
Qtilde)=r c1(L). Under any chosen rational algebraic B field, put
u=c1(L)+B and replace alpha by alpha+B; the differences P\_j are
unchanged. If a coherent effective finite-action weakly semiregular
balanced lift existed, the necessary determinant condition proved in
E825 would give

\[
 u-\frac72D\quad\hbox{has no mixed RM block and is primitive
 on each of the two RM planes.}
 \tag{E832.11}
\]

Here D=k theta, k\textgreater0. Therefore u is block diagonal and both
block traces are7k\textgreater0. All hypotheses of section3 now hold,
while(E832.10) makes b semipositive. The contradiction excludes
\textbf{every actual positive elliptic star in this precise unit/free
architecture, for every rational B field and every specified coherent
effective finite action} in E825's scope. It is stronger than failure of
a chosen extension or a dimension count: its normalized first two Chern
terms and the proved necessary determinant condition already conflict.

This is still a star theorem. More than one center, a negative or
indefinite square-zero difference, a different unit Chern character, or
an object without the stated defect determinant is not covered by this
application. The scalar-unit conclusion of section4 has its separate,
explicitly broader sign allowance. Neither statement classifies
arbitrary unstable sheaves or perfect complexes.

\subsection{6. Why the two positivity hypotheses must be
kept}\label{SF-E832:why-the-two-positivity-hypotheses-must-be-kept}

A spectral proof cannot silently drop both scalarity and semipositivity.
For example, in the orthogonal plane decomposition let

\[
 U_0=\operatorname{diag}(1,1,2,-2),\qquad
 B_0=\operatorname{diag}(2,2,1,-1).
 \tag{E832.12}
\]

The compound difference Lambda2 B0-Lambda2 U0 is zero on the four mixed
basis vectors and equals3 on each of the two pure ones. This is exactly
the general-unit degree-two equation with nonzero right side, at the
level of real Hermitian linear algebra. B0 is indefinite, and U0 also
lacks positive trace in both blocks. Thus it violates the hypotheses
used in section3 and gives no counterexample to it.

This boundary example is not asserted to define a rational line class
for a genuine RM lattice, an integral positive-rank star, the full
balanced Chern character, a coherent object or a semiregular seed. It
only prevents the scalar spectral argument from being overextended
without the actual positivity supplied in section5.

\subsection{7. Source and check
boundary}\label{SF-E832:source-and-check-boundary}

All new compound-matrix, rank, sign and graph arguments are proved
above. The unit resolution and its canonical defect map are the
previously whole-reviewed US/UG and E825 inputs, used exactly as stated.
The only categorical input for the general-B application is E825's
already proved necessary determinant condition, in its preserved
finite-action scope; no new source on stability, Hodge classes or
generic rank is imported. The exact supplementary script
rm\_elliptic\_star\_character\_certificate.py actually passed the unit
expansion, completed square, determinant-product identity, all15
four-edge graphs and the printed boundary matrices in 0.053802978s. It
does not numerically certify the universal Hermitian positivity
argument, the actual E825 interface or the existence of a seed. No
successful seed or general nonexistence outside the stated star
architectures is claimed.

\section{SF-E832R. The uniformly reversed elliptic star is also
excluded}\label{SF-E832R:sf-e832r.-the-uniformly-reversed-elliptic-star-is-also-excluded}

Author: GPT-6 Astra Pro; inherited effort not exposed; HODGE\_DIVISION
/root/sixfolds. Date: 2026-11-21.

\textbf{Status:} independently whole-reviewed by /root and
/root/g34\_eightfolds; both mathematical reviews PASS on frozen draft
\texttt{dc18d6bf237e6133a4f86a74c466d3962380a2b8cbb48f440931012bc8877446}.
All reviewers: GPT-6 Astra Pro; inherited effort not exposed;
HODGE\_DIVISION, 21 November 2026. No new computation is needed or
claimed. The frozen E832 base proof is unchanged. This adds the opposite
uniform sign, not arbitrary mixed-sign or multicenter line
configurations.

\subsection{1. Exact reversed
architecture}\label{SF-E832R:exact-reversed-architecture}

Retain the actual unit F on the smooth complete-intersection surface S
with D=k theta, k\textgreater0. Consider any actual coherent extension

\[
 0\longrightarrow A\longrightarrow E\longrightarrow B\longrightarrow0,
 \qquad A=(F\otimes L)^{\oplus r}
          \oplus\bigoplus_j L_j^{\oplus n_j},
 \qquad B=M^{\oplus n_0},
 \tag{E832R.1}
\]

where r and n0 are positive, each n\_j is positive, and

\[
 c_1(L_j)=\alpha+P_j,\quad\alpha=c_1(M),
 \qquad P_j\ge0,\quad\operatorname{rank}P_j\le1.
 \tag{E832R.2}
\]

The free total rank is m=n0+sum\_j n\_j, and R=r+m. An empty sum is
permitted. In the opposite-slope realization, M is the single negative
free center, and the positive free classes have nonzero rank-one
differences of positive theta-degree from it, so(E832R.2) is automatic.
The sheaf E has no weakly semiregular balanced lift with normalized
character R v\_q, q nonzero, for any algebraic rational B field and any
specified coherent effective finite action in E825's scope.

\subsection{2. Complete proof}\label{SF-E832R:complete-proof}

Include the quotient free terms with P0=0. The actual K0 identity from
(E832R.1) and the fact P\_j²=0 give the free character exp(alpha)(m+H),
with H=sum\_j n\_j P\_j semipositive. After a rational B shift replace
alpha by alpha+B and put u=c1(L)+B. The unit contributes r exp(u)
through its second Chern character, exactly as computed in E832. The
normalized first Chern equation is

\[
 ru+m\alpha+H=0.
 \tag{E832R.3}
\]

Thus the completed-square form in E832.9 is

\[
 b=\alpha+\frac r m u=-\frac{H}{m}\le0.
 \tag{E832R.4}
\]

The same degree-two calculation gives b²-kappa
u²=d(theta\_+²+theta\_-²), with kappa=rR/m² positive and d=Rq/m nonzero.
Put b'=−b. Then b' is semipositive and b'²=b², so it satisfies the
identical compound-square equation.

The exact defect interface must still be checked for this reversed
filtration. Since B is locally free, (E832R.1) splits locally. Relative
to these splittings the FF diagonal transition is that of F tensor L to
multiplicity r. The other diagonal blocks are free, and the extra
extension terms go from the free B part to A. E825's multiplicative ell
map kills every off-diagonal free-to-unit or free-to-free term, as well
as products factoring through a free module. After removing the fixed F
comparisons, its resulting cocycle is precisely the line cocycle of L to
multiplicity r. Hence Qtilde=L\^{}r on S2 and its determinant extends to
r c1(L) on X.

If a balanced weakly semiregular lift existed, E825.3--4 would therefore
force u−7D/2 to have zero mixed RM blocks and zero trace on each
rank-two RM plane. The two traces of u are7k\textgreater0. The
semipositive lemma of E832§3 applies to b',u,kappa,d and gives a
contradiction. This proof uses the intrinsic E825 condition on E; it
does not require an extra action on the chosen extension filtration or
its local direct summands.

\subsection{3. Boundary and checks}\label{SF-E832R:boundary-and-checks}

Together with the base theorem this covers stars with either uniform
semipositive sign. When different leaves have different signs, the sum H
need not be semipositive or seminegative and this proof does not apply.
Multiple centers, different unit characters and arbitrary gluing with an
unproved defect determinant remain outside. No new source, numerical
premise, finite action or successful seed is introduced. The only added
algebra is b'=-b and b'²=b²; the exact reversed A/B interface has been
supplied explicitly above.

\section{SF-E833. The full cubic character excludes every square-zero
star}\label{SF-E833:sf-e833.-the-full-cubic-character-excludes-every-square-zero-star}

Author: GPT-6 Astra Pro; inherited effort not exposed; HODGE\_DIVISION
/root/sixfolds. Date: 2026-11-21.

\textbf{Status:} independently whole-reviewed by /root and
/root/g34\_eightfolds; both mathematical reviews PASS on frozen draft
\texttt{2a74ada9fb1a9894ffc7c1f402e8eeb4d470af5a7d6d63923951055b99f4b58f}.
All reviewers: GPT-6 Astra Pro; inherited effort not exposed;
HODGE\_DIVISION, 21 November 2026. The peer also supplied a separately
authored eight-generator exterior-algebra certificate; the universal
arguments were reviewed in full. This uses the full cubic character and
strictly extends E832's same-center obstruction to arbitrary signs. It
does not treat multiple centers, a different unit character, or an
unproved defect determinant. No assertion about a successful balanced
seed is made.

\subsection{1. The exact character and defect
hypotheses}\label{SF-E833:the-exact-character-and-defect-hypotheses}

Let X be the abelian fourfold with orthogonal rank-two RM tangent planes
V\_+, V\_- and polarization theta=theta\_++theta\_-. Let D=k theta,
k\textgreater0. The actual unit F on the smooth complete-intersection
surface has

\[
 \operatorname{ch}(F)=2e^{-D}-2e^{-3D}+e^{-4D}
                  =1-2D^3+4D^4.
 \tag{E833.1}
\]

Fix positive integers r,m. Consider an actual exact sequence

\[
 0\longrightarrow A\longrightarrow E\longrightarrow B\longrightarrow0,
 \quad A=(F\otimes L)^{\oplus r}\oplus A_{\mathrm f},
 \tag{E833.2}
\]

where A\_f and B are direct sums of line bundles of total rank m. Either
of these two free sums may be empty. Require that the first Chern
classes of all those free lines, counted with positive multiplicities
n\_j, lie on one square-zero star:

\[
 c_1(N_j)=\alpha_0+\eta_j,\qquad \eta_j^2=0,
 \qquad\sum_jn_j=m.
 \tag{E833.3}
\]

The center alpha0 may be any fixed real (1,1) class; it need not occur
as an actual free summand. There is no sign condition on eta\_j and no
slope condition on the given exact sequence. A real Hermitian
square-zero (1,1) form has rank at most one, but that equivalence is not
needed below.

Then E admits no weakly semiregular balanced lift with character R v\_q,
where R=r+m and q is nonzero, for any algebraic rational B field and any
specified coherent effective finite action in E825's scope.

More generally, the same conclusion applies whenever the actual object
has the displayed character identity and E825's defect determinant is
proved to extend from X as r c1(L). This last hypothesis is not assumed
for arbitrary gluing; equation(E833.2) supplies it explicitly in
section5.

\subsection{2. A compound-matrix lemma without a positivity hypothesis
on
b}\label{SF-E833:a-compound-matrix-lemma-without-a-positivity-hypothesis-on-b}

Normalize theta to the identity Hermitian form. Suppose u preserves the
RM planes, and its Hermitian blocks U,V both have the same strictly
positive trace h. Suppose b is any real Hermitian form satisfying

\[
 b^2-\kappa u^2=d(\theta_+^2+\theta_-^2),
 \qquad \kappa>0,\quad d\ne0.
 \tag{E833.4}
\]

Let B be b's matrix and let Pi be the sum of the two pure-plane
projections on Lambda2(V\_+ direct sum V\_-). The invariant-form
identification proved in E832 gives

\[
 \Lambda^2B=\kappa\Lambda^2\operatorname{diag}(U,V)+d\Pi.
 \tag{E833.5}
\]

Neither U nor V is zero because its trace is h. We first show that both
are invertible. Suppose U has rank one.

If V also has rank one, the right side of(E833.5) has rank three: its
two pure-line coefficients are d and d, and its mixed block U tensor V
has rank one. Both pure lines belong to its image. Since im(Lambda2
B)=Lambda2 im(B), their two support planes must both lie in im(B). Thus
B has rank four, whose compound has rank six, a contradiction.

If V has rank two, the pure V\_+ line still has coefficient d and lies
in the image. The mixed image is im(U) tensor V\_-, of dimension two.
Write im(U)=C a with a nonzero in V\_+. The image therefore contains the
bivectors a wedge v for every v in V\_-. A nonzero decomposable bivector
in Lambda2 W has its whole support plane in W: contract by linear
functionals, or write its two independent factors in a basis extending
W. Applying this to W=im(B), the pure V\_+ line first gives V\_+
contained in W; the mixed bivectors then give all V\_- contained in W.
Again B has rank four. But the right side of(E833.5) has mixed rank two
and at most two pure-line ranks, hence rank at most four, a
contradiction. The same argument interchanging the planes treats
rank-one V. Thus U,V are both invertible. No semipositivity or
determinant sign entered this argument.

The mixed block U tensor V now has rank four. A four-dimensional
operator of rank at most three has compound rank at most three, so
(E833.5) forces B invertible. Its compound is invertible too. Each pure
line is preserved, with nonzero eigenvalue, by the right side. Hence B
preserves each support plane. It is block diagonal with invertible
Hermitian blocks C,D1, and

\[
 C\otimes D_1=\kappa U\otimes V.
 \tag{E833.6}
\]

Equality of nonzero simple tensors gives C=tU and D1=(kappa/t)V. The
scalar t is nonzero and real: if two nonzero Hermitian matrices are
complex scalar multiples, taking adjoints makes that scalar equal to its
conjugate. Consequently, as invariant forms,

\[
 u=p+z,\qquad b=t p+(\kappa/t)z,
 \quad p=u|_{V_+},\quad z=u|_{V_-},\quad t\in\mathbb R^*.
 \tag{E833.7}
\]

In particular p and z are nondegenerate. This is the entire conclusion
of the lemma; it does not assert that b or u is semipositive.

\subsection{3. Complete second and third character
calculation}\label{SF-E833:complete-second-and-third-character-calculation}

After an algebraic rational B shift, put u=c1(L)+B and alpha=alpha0+B.
Set H=sum\_j n\_j eta\_j. The actual K0 identity and eta\_j²=0 give

\[
 e^B\operatorname{ch}(E)
  =r e^u(1-2D^3+4D^4)+e^\alpha(m+H).
 \tag{E833.8}
\]

This is the full character, not only its first terms. If it equals R
v\_q, its first, second and third characters are respectively zero,
-Rq(theta\_+²+theta\_-²)/2, and zero. The first equation gives H=-ru-m
alpha. The second equation is exactly E832.9:

\[
 b=\alpha+xu,\quad x=\frac rm>0,\qquad
 b^2-\kappa u^2=d(\theta_+^2+\theta_-^2),
 \quad\kappa=x(1+x)>0,\quad d=(1+x)q\ne0.
 \tag{E833.9}
\]

The cubic character is, with every coefficient displayed,

\[
 0=\frac1m\operatorname{ch}_3(E;B)
   =\frac x6u^3-\frac x2u\alpha^2-\frac13\alpha^3-2xD^3.
 \tag{E833.10}
\]

Indeed the free contribution before eliminating H is m alpha³/6+H
alpha²/2, which becomes -r u alpha²/2-m alpha³/3. The unit contributes r
u³/6-2rD³. Thus there is no missing Chern-class normalization or factor
of six in(E833.10).

Suppose the actual E825 condition has given u block diagonal and both
block traces h=7k. The lemma of section2 applies to(E833.9). With its
notation p,z,t, put

\[
 a=t-x,\qquad c=\kappa/t-x,
 \qquad \alpha=a p+c z.
 \tag{E833.11}
\]

Since p³=z³=0, expand(E833.10) as

\[
 A_+p^2z+A_-pz^2=2xD^3,
 \quad
 A_+=\frac x2(1-a^2-2ac)-a^2c,
 \quad
 A_-=\frac x2(1-c^2-2ac)-ac^2.
 \tag{E833.12}
\]

The two mixed cubic components lie in the complementary summands of
invariant forms with plane bidegrees ((2,2),(1,1)) and ((1,1),(2,2)). On
each plane its normalized polarization is the identity, so p²=(det
U)theta\_+² and z²=(det V)theta\_-². Also
D³=3k³(theta\_+²theta\_-+theta\_+theta\_-²). Thus(E833.12) gives the two
actual equations of (1,1) forms

\[
 A_+\det(U)\,z=6xk^3\theta_-,\qquad
 A_-\det(V)\,p=6xk^3\theta_+.
 \tag{E833.13}
\]

The scalar 6xk³ is nonzero. The first equation forces z to be a nonzero
scalar multiple of theta\_-; the second forces the same for p and
theta\_+. The traces of both blocks are h=7k. Since the trace of each
plane's identity is two, these scalar multiples are both 7k/2. Hence

\[
 u=\frac{7k}{2}\theta.
 \tag{E833.14}
\]

Now(E833.9) contradicts the scalar-unit lemma E832§2, applied with
c0=kappa(7k/2)² and d nonzero. This proves the desired incompatibility.
Notice that if either coefficient A\_+ or A\_- vanishes, the
corresponding equation(E833.13) is already impossible; no division by an
unproved nonzero coefficient has been made.

\subsection{4. Why this is stricter than the uniform-sign
theorem}\label{SF-E833:why-this-is-stricter-than-the-uniform-sign-theorem}

Only the identity eta\_j²=0 was used. Their signs, the sign of their sum
H, and the ranks of the actual extension matrices never entered. The
semipositive lemma E832§3 is not used; the argument replaces its missing
positivity with the nonzero cubic unit term and a separate
singular-block analysis. Even indefinite b and indefinite adjusted u are
accounted for. The earlier real-Hermitian counterexample to a
second-character-only argument does not satisfy these full cubic
equations with the prescribed unit and its defect traces.

The common center is essential to the displayed proof: it supplies the
exact free character e\^{}alpha(m+H) and hence the single
completed-square form b. Several centers produce several such terms, and
the present argument does not combine them into one. Nor does it address
arbitrary perfect complexes with signed virtual free multiplicities or
r/m of unknown sign.

\subsection{5. Actual determinant, finite actions and source
boundary}\label{SF-E833:actual-determinant-finite-actions-and-source-boundary}

For(E833.2), B is locally free, so choose local splittings compatible
with A. Within A the unit transition is precisely that of F tensor L, to
multiplicity r. Every other diagonal block is free; every new extension
entry goes from the free B part to A. After removing the fixed F
comparison, E825's actual multiplicative ring map ell kills all free
blocks and every factorization through a free module. Its FF transition
is exactly the scalar L transition to multiplicity r. Therefore its
intrinsic first-neighborhood defect bundle is Qtilde=L\^{}r, and the
determinant extends as r c1(L) on X.

E825.3--4, for the given coherent effective finite action and the chosen
rational algebraic B field, would force u-7D/2 to have zero mixed RM
blocks and to be primitive separately on both planes. This is exactly
the block-diagonal condition and the equal traces h=7k used above. The
condition is intrinsic to E; the argument does not add an equivariance
assumption on the chosen filtration or its summands. No pullback
preservation of semiregularity is being asserted.

The original unit resolution and E825 condition are the already
whole-reviewed inputs, with their published BF and supplied S21 scope
unchanged. E832§1--2 supplies the explicitly proved invariant-form
identification and scalar spectral lemma. Every new compound-rank,
support-plane, cubic and actual-cocycle step has been proved above. A
separate exact algebra supplement checks the expansions only; it is not
a numerical premise for the universal Hermitian or sheaf arguments.

\section{SF-E834. Mixed discriminant equality for stable reflexive
sheaves}\label{SF-E834:sf-e834.-mixed-discriminant-equality-for-stable-reflexive-sheaves}

Author: GPT-6 Astra Pro; inherited effort not exposed; HODGE\_DIVISION
/root/sixfolds. Date: 2026-11-21.

\textbf{Status:} complete new proof, whole-read by the author and
independently checked in full by /root and /root/g34\_eightfolds.
The invoked published analytic statements and the new arguments pass;
their exact reading scopes are recorded separately. This repairs the
two specific unresolved steps
recorded in reflexive\_curvature\_attempt.md. It uses the published
admissible-metric and general-analytic-set extension theorems at the
exact locators below. It neither assumes a full Chern--Weil current
identity on a singular space nor applies an isolated-point removal
theorem to a curve.

\subsection{1. Statement and precise published
inputs}\label{SF-E834:statement-and-precise-published-inputs}

Let \(X\) be a polarized complex abelian fourfold with a genuine
Rosati-invariant real quadratic action and its orthogonal complex
rank-two tangent eigenspaces \(V_+,V_-\). Let \(\theta=\theta_++\theta_-\) be
the invariant positive polarization form. For a reflexive sheaf \(E\) of
rank \(r>0\) define the rational discriminant class

\[
 \Delta(E)=c_2(E)-\frac{r-1}{2r}c_1(E)^2
          =\frac{c_1(E)^2}{2r}-\operatorname{ch}_2(E).
 \tag{E834.1}
\]

If \(E\) is slope-stable for \(\theta\), then

\[
 \int_X\Delta(E)\theta_+\theta_-\ge0.
 \tag{E834.2}
\]

Equality holds if and only if \(E\) is locally free and projectively flat.
In particular no slope-stable reflexive sheaf has normalized character \(Rv_q\) with \(q>0\), for any rational algebraic \(B\) field, regardless
of an additional coherent finite translation/tensor action.

The analytic inputs are as follows.

\begin{itemize}
\item
  Bando--Siu, \emph{Stable sheaves and Einstein-Hermitian metrics}, in
  \emph{Geometry and Analysis on Complex Manifolds}, World
  Scientific1994, pp39--50, \textbf{Theorem3, p40}: a polystable
  reflexive sheaf on a compact Kähler manifold admits an admissible
  Hermitian--Einstein metric. Admissibility, defined on p39, includes
  square-integrable curvature. The complete invoked §4 proof is on
  pp48--49, with the preceding regularization construction on pp45--48.
  The author's original page images are at
  \url{https://people.math.harvard.edu/~siu/bando_joint_paper/index.html}.
\item
  Biswas, \emph{A criterion for holomorphic extension of principal
  bundles}, Kodai Math. J.34(2011),71--78, \textbf{Theorem3.2 and its
  complete proof, p77} (also Theorem1.2, p72): a holomorphic principal
  bundle for a complex reductive affine group extends uniquely across a
  closed complex analytic subset of codimension at least two if the
  Chern curvature of a smooth maximal-compact reduction has finite
  \(L^n\) norm, where \(n\) is the complex dimension of the ambient manifold.
  We use \(G=\operatorname{GL}_r(\mathbb C)\), \(K=U(r)\), \(n=4\). The primary PDF is
  \url{https://www.jstage.jst.go.jp/article/kodaimath/34/1/34_1_71/_pdf/-char/en}.
  Its vector-bundle input is Harris--Tonegawa, \emph{Analytic
  continuation of vector bundles with \(L^p\)-curvature}, Int. J.
  Math.11(2000),29--40, \textbf{Theorem1, p29},
  DOI10.1142/S0129167X00000040. The original Harris--Tonegawa proof has
  not yet been retrieved here; no claim to have independently audited
  that upstream proof is made. The invoked published Biswas statement
  and its full proof have been read.
\end{itemize}

The final locally free equality implication is exactly the already
whole-reviewed SF-E811 in its rational-polarization scope. All new
averaging, integrability, Bianchi and local-slice arguments are supplied
below. No numerical computation is a premise of this analytic proof.

\subsection{2. Translation averaging supplies the needed Chern--Weil
pairings}\label{SF-E834:translation-averaging-supplies-the-needed-chernweil-pairings}

Let \(S\) be the singular locus of \(E\). A reflexive sheaf on a smooth complex
manifold is locally free in codimension two; thus \(\dim_{\mathbb C}S\) is at most
one. For this familiar fact, locally at codimension at most two a
reflexive module has full depth over the regular local ring, and the
Auslander--Buchsbaum formula gives projective dimension zero. It is also
stated explicitly in Bando--Siu, p46. Put \(U=X\setminus S\).

Choose the admissible Hermitian--Einstein metric \(h\) on \(E\) restricted to \(U\),
with Chern curvature \(F\), and set

\[
 F^0=F-\frac1r\operatorname{Tr}(F)I,
 \qquad
 \beta=\frac1{8\pi^2}\operatorname{Tr}(F^0\wedge F^0).
 \tag{E834.3}
\]

On \(U\) this is the usual smooth Chern--Weil representative of \(\Delta(E)\).
For the fixed flat base metric and the unitary fiber norm,
\(|\beta|\) is at most \(C_r|F|^2\). Therefore
\(\beta\) is integrable on \(X\). We extend its coefficients by arbitrary values
on the measure-zero set \(S\); no closed-current assertion is being made at
this point.

Let \(H_1,H_2\) be integral ample divisor classes. Choose sufficiently large
positive integers \(a,b\) and smooth divisors of classes \(aH_1,bH_2\) whose
intersection \(Y\) is a smooth compact surface. Ordinary Bertini supplies
this after very ampleness; the argument does not require \(Y\) connected.
Let \(t_x\) denote translation by \(x\) in \(X\) and let \(\mu\) be Haar probability.
Except for a measure-zero set of \(x\), the translate \(Y+x\) avoids \(S\). Indeed
the exceptional set is contained in the image \(S-Y\) of \(S\times Y\);
its complex dimension is at most three, while \(X\) has dimension four. One
may equivalently use a finite analytic stratification and the real
Hausdorff-dimension bound at most six.

On each translate disjoint from \(S\), the restriction of \(h\) is an ordinary
smooth Hermitian metric on the vector bundle \(E\) restricted to \(Y+x\). Its
restricted Chern connection has curvature \(F\) restricted to the surface.
Chern--Weil applies without any Einstein condition on the restricted
metric and gives

\[
 \int_{Y+x}\beta
   =\int_{Y+x}\Delta(E|_{Y+x})
   =ab\int_X\Delta(E)H_1H_2.
 \tag{E834.4}
\]

The last identity is topological restriction: \(Y+x\) lies in the locally
free locus, so its derived restriction of \(E\) is the ordinary vector
bundle restriction, and Chern classes restrict. Translation fixes
cohomology, including the class of \(Y\).

Here is absolute Fubini explicitly. Give \(Y\) its induced invariant-base
volume. The norm of restriction of a four-form to a tangent four-plane
is bounded by its ambient norm, up to one harmless fixed convention
constant. Translation is an isometry. Hence

\[
 \int_X\int_Y |t_x^*\beta|\,d\mathrm{vol}_Y\,d\mu(x)
 \le C\operatorname{vol}(Y)\int_X|\beta|\,d\mu
 <\infty.
 \tag{E834.5}
\]

The identity follows by Tonelli applied to \(|\beta(x+y)|\)
and translation invariance of Haar measure. It controls the quadratic
Chern--Weil form by the assumed \(L^2\) curvature; no pointwise bound near \(S\)
was used.

The averaged current \(\gamma_Y=\int_X[Y+x]\,d\mu(x)\) is
translation invariant, closed, and has cohomology class \([Y]=abH_1H_2\).
On a torus, translation-invariant currents have constant coefficients,
hence are smooth forms; a constant form is uniquely determined by its de
Rham class. Consequently \(\gamma_Y\) is exactly \(ab\) times the product of the
invariant representatives of \(H_1\) and \(H_2\). Averaging (E834.4), justified by
(E834.5), gives

\[
 \int_U\beta\,H_1H_2
   =\int_X\Delta(E)H_1H_2.
 \tag{E834.6}
\]

This argument needs only the smooth compact translated surfaces. It does
not say that an arbitrary singular-sheaf Chern current represents
\(\Delta(E)\) against all test forms. Real bilinearity extends (E834.6) to
products of classes in \(\operatorname{NS}(X)_{\mathbb R}\), because the ample cone is open and
spans that vector space.

\subsection{3. The RM mixed form lies in that precise
span}\label{SF-E834:the-rm-mixed-form-lies-in-that-precise-span}

Choose a trace-zero rational real-multiplication operator \(u\) with
\(u^2=d>0\). The \(\theta\)-Rosati symmetry makes \(\lambda(v,w)=\theta(uv,w)\) an invariant rational \((1,1)\) form. After clearing denominators it is
an integral alternating form of type \((1,1)\), hence a divisor class by the
usual line-bundle classification on a complex torus. On the two
eigenspaces it is respectively \(\sqrt d\,\theta_+\) and \(-\sqrt d\,\theta_-\).
Therefore

\[
 \lambda=\sqrt d(\theta_+-\theta_-),\qquad
 \theta_+\theta_-=\frac14\left(\theta^2-\frac1d\lambda^2\right).
 \tag{E834.7}
\]

For an explicit expression by ample products, take a positive rational
\(\epsilon\) with \(\epsilon\sqrt d<1\). Both \(\theta+\epsilon\lambda\) and \(\theta-\epsilon\lambda\) are rational ample classes, and
\(\lambda^2=(\theta^2-(\theta+\epsilon\lambda)(\theta-\epsilon\lambda))/\epsilon^2\).
Thus (E834.6) gives the exact identity actually needed:

\[
 \int_X\Delta(E)\theta_+\theta_-
    =\int_U\beta\,\theta_+\theta_-.
 \tag{E834.8}
\]

This is why only divisor-pair Chern--Weil identities suffice; the
irrational individual RM planes are not treated as abelian subvarieties.

\subsection{4. Equality removes mixed curvature on the regular
locus}\label{SF-E834:equality-removes-mixed-curvature-on-the-regular-locus}

The Hermitian--Einstein equation implies \(\Lambda_\theta F^0=0\), and \(F^0\) has
type \((1,1)\). The pointwise calculation SF-E811.4 applies unchanged at
every point of \(U\). For clarity, in unitary coordinates with
\(\omega_+=e_1+e_2\), \(\omega_-=e_3+e_4\), a real primitive \((1,1)\) form \(\alpha\) has

\[
 -\alpha^2\omega_+\omega_-
 =2\left((x_1+x_2)^2+
       \sum_{i\le2<j}(b_{ij}^2+c_{ij}^2)\right)\mathrm{vol}.
 \tag{E834.9}
\]

Here \(x_i\) are its diagonal entries and \(b_{ij},c_{ij}\) the two real
coefficients of each off-diagonal Hermitian entry. Summing this identity
in an orthonormal skew-Hermitian matrix basis proves that the right side
of (E834.8) is nonnegative. This proves (E834.2).

If it is zero, the nonnegative smooth density is zero everywhere on \(U\).
There are smooth endomorphism-valued real two-forms \(F_+,F_-\) on \(U\) such
that

\[
 F^0=F_++F_-,\qquad
 F_\pm\text{ has both arguments in }V_\pm,\qquad
 \Lambda_{\theta_\pm}F_\pm=0.
 \tag{E834.10}
\]

All real mixed curvature components vanish. The type \((1,1)\) condition
already removed the \((2,0)\) and \((0,2)\) components; equation (E834.9) removes
the remaining mixed ones. This step concerns only \(U\), where all forms and
metrics are smooth.

\subsection{5. Bianchi and clean slices bound both curvature
blocks}\label{SF-E834:bianchi-and-clean-slices-bound-both-curvature-blocks}

Work in a small holomorphic product coordinate box in the universal
cover, using the two parallel complex planes. Take a closed larger
product box contained in the coordinate chart, and a strictly smaller
product box inside it. The intersection of \(S\) with the closed larger box
is compact and has complex dimension at most one. Its projection to
either two-dimensional complex coordinate factor has real measure zero.
Choose a point \(b\) in the smaller second factor outside the projection to
that factor, and a point \(a\) in the smaller first factor outside the other
projection. The horizontal slice through \(b\) and the vertical slice
through \(a\), taken over the closed intermediate factor boxes, are compact
and disjoint from \(S\). They consequently have positive distance from \(S\) and
bounded smooth curvature. This choice is made against the compact larger
box, not only against \(S\) in the smaller open box.

Bianchi gives \(d_\nabla F^0=0\): it holds for \(F\), and the scalar trace
curvature is closed. For constant real vectors \(v,w\) in \(V_+\) and \(z\) in
\(V_-\), its component is

\[
 (\nabla_zF^0)(v,w)+(\nabla_vF^0)(w,z)
                      +(\nabla_wF^0)(z,v)=0.
 \tag{E834.11}
\]

The last two terms vanish because the corresponding mixed components are
identically zero. Thus every coefficient of \(F_+\) is parallel along the
second-factor directions. The corresponding statement holds for \(F_-\)
along the first-factor directions. The connection on \(\operatorname{End}(E)\) is unitary
for the Hilbert--Schmidt norm, so each block norm is constant along
every path in its transverse factor fiber that avoids \(S\).

Every such punctured factor fiber is path connected. Its removed set is
a proper closed analytic subset of a two-dimensional complex box, of
complex dimension at most one. To see connectedness directly, join two
points outside that subset through complex line segments inside the
convex box. On each complex line through an outside point, the analytic
intersection is discrete, since the line is not contained in the subset.
Small detours within the line avoid the finitely many intersection
points on a compact segment. An intermediate outside point and two such
paths therefore connect the chosen points. No simple connectedness or
flatness along that fiber is required.

For a point \((z_1,z_2)\) of the smaller box outside \(S\), connect \(z_2\) to \(b\) within
the second-factor punctured fiber with first coordinate \(z_1\). Parallel
norm transport gives

\[
 |F_+(z_1,z_2)|=|F_+(z_1,b)|.
 \tag{E834.12}
\]

The right side is bounded uniformly by the clean compact horizontal
slice. Similarly
\(|F_-(z_1,z_2)|=|F_-(a,z_2)|\) is uniformly
bounded by the clean vertical slice. Hence \(F^0\) is locally bounded near
every point of \(S\). This proof works also when \(S\) contains curves, and
never identifies the punctured universal cover with a product.

\subsection{6. Smooth determinant normalization and a general-set bundle
extension}\label{SF-E834:smooth-determinant-normalization-and-a-general-set-bundle-extension}

The determinant of a reflexive sheaf on the smooth \(X\) is a holomorphic
line bundle. Choose any smooth Hermitian metric \(h_0\) on that determinant.
On \(U\), let \(g\) be the positive smooth ratio \(h_0/\det h\), and replace \(h\) by

\[
 h'=g^{1/r}h.
 \tag{E834.13}
\]

Then \(\det h'=h_0\). Scalar metric rescaling changes the Chern connection by
a scalar one-form, so it leaves the induced \(\operatorname{End}(E)\) connection and its
trace-free curvature unchanged. It also leaves the norm of every
endomorphism unchanged. Thus the curvature of \(h'\) is

\[
 F'=F^0+\frac1rF_{h_0}I.
 \tag{E834.14}
\]

The first term is locally bounded by section 5, and the second extends
smoothly across \(S\). Hence \(F'\) is locally bounded on \(X\). Compactness gives
finite \(L^4\) norm. The metric \(h'\) need not be Hermitian--Einstein, and
that property is not required in the extension theorem.

Apply Biswas Theorem3.2 with the holomorphic frame \(\operatorname{GL}_r\) bundle of
\(E|_U\) and the \(U(r)\) reduction supplied by \(h'\). The base is the
smooth complex fourfold \(X\), \(S\) is a closed analytic subset of codimension
at least three, and the curvature has finite \(L^4\) norm. Every
hypothesis is now supplied. There is a holomorphic vector-bundle
extension \(\overline E\) on \(X\). Both \(\overline E\) and the original \(E\) are reflexive and
agree outside codimension at least two, so they are uniquely isomorphic:
reflexive sheaves on a smooth normal space equal the pushforward of
their restriction across such a subset. Thus \(E\) is locally free.

Now apply the already proved stable-bundle theorem SF-E811 to \(E\). The
same discriminant class has zero mixed pairing, so \(E\) is projectively
flat. Conversely a locally free projectively flat \(E\) has \(\Delta(E)=0\),
which proves the equality characterization.

\subsection{7. The balanced consequence and exact
boundary}\label{SF-E834:the-balanced-consequence-and-exact-boundary}

For \(e^B\operatorname{ch}(E)=Rv_q\) with \(q>0\), the first two terms imply
\(\Delta(E)=Rq(\theta_+^2+\theta_-^2)/2\). It pairs to zero with
\(\theta_+\theta_-\), but its pairing with \(\theta^2\) is positive. The equality
theorem would force \(\Delta(E)=0\), a contradiction. A finite
translation/tensor action cannot alter this obstruction on the
underlying stable reflexive sheaf.

This theorem does not directly assert the result for an unstable sheaf,
a semistable sheaf with an unexamined filtration, or a perfect complex.
Nor does it claim that the singular Hermitian--Einstein connection
itself has been smoothly extended by the new argument: only the
holomorphic bundle extension is needed, followed by the existing locally
free result. The translated-surface averaging and clean-slice bounds are
the new steps; the published analytic existence and extension theorems
are stated inputs with their source status explicit.

\section{SF-E835. Every semistable torsion-free balanced sheaf reduces
to the checked hull
gate}\label{SF-E835:sf-e835.-every-semistable-torsion-free-balanced-sheaf-reduces-to-the-checked-hull-gate}

Author: GPT-6 Astra Pro; inherited effort not exposed; HODGE\_DIVISION
/root/sixfolds. Date: 2026-11-21.

\textbf{Status:} complete new proof, whole-read by the author and
independently checked in full by /root and /root/g34\_eightfolds.
Its analytic premise SF-E834 has separately passed both whole-proof
reviews. The remaining algebraic steps are proved in full below. The statement concerns actual semistable
torsion-free sheaves; it does not extend to unstable sheaves or
arbitrary perfect complexes.

\subsection{1. Statement}\label{SF-E835:statement}

Retain the genuine polarized RM abelian fourfold \(X\) and
\(\theta=\theta_++\theta_-\) of SF-E834. Let \(E\) be any \(\theta\)-slope-semistable
torsion-free coherent sheaf of positive rank \(R\), and suppose

\[
 e^B\operatorname{ch}(E)
   =R\left(1-\frac q2\theta_+^2\right)
      \left(1-\frac q2\theta_-^2\right),
 \qquad B\in\operatorname{NS}(X)_{\mathbb Q},\quad q\in\mathbb Q_{>0}.
 \tag{E835.1}
\]

Then the actual hull \(H=E^{\vee\vee}\) is locally free and has a
filtration by subbundles whose quotients are stable projectively flat
vector bundles of the same \(\theta\)-slope. The quotient \(H/E\) is supported on
finitely many points or is zero. Consequently \(E\) is not weakly
\(G\)-semiregular for any specified coherent effective finite
translation/tensor action on \(E\) in the scope of SF-E818.

This supplies the hypotheses on the hull which were explicitly assumed
in E818. In particular a singular reflexive hull or a singular stable
factor in a locally free hull is no longer an unchecked semistable
escape, by the separately checked E834. No chosen stable-factor filtration is
assumed to be preserved by \(G\).

\subsection{2. Stable torsion-free factors and stability of their
hulls}\label{SF-E835:stable-torsion-free-factors-and-stability-of-their-hulls}

Choose a finite slope-stable-factor filtration

\[
 0=E_0\subset E_1\subset\cdots\subset E_s=E,
 \qquad Q_i=E_i/E_{i-1},\qquad r_i=\operatorname{rk}Q_i>0,
 \tag{E835.2}
\]

with torsion-free quotients of slope \(\mu_\theta(E)\). Here is a direct
existence argument adequate for this use. If a semistable torsion-free
sheaf is not stable, choose an equal-slope proper subsheaf of smallest
positive rank and saturate it. Saturation adds a torsion quotient whose
codimension-one cycle is effective. Any increase of slope would
contradict semistability of the ambient sheaf; thus the saturated sheaf
still has equal slope. Minimality of its rank makes it stable. The
saturated quotient is torsion-free and semistable: a larger-slope
subsheaf of the quotient would lift to a larger-slope subsheaf of the
ambient sheaf. Induction on rank gives (E835.2). This does not assume
local freeness of its factors.

Set

\[
 H_i=Q_i^{**},\qquad T_i=H_i/Q_i,
 \quad \alpha_i=\frac{c_1(Q_i)}{r_i}=\frac{c_1(H_i)}{r_i},
 \quad \alpha=\frac{c_1(E)}R,\quad\delta_i=\alpha_i-\alpha.
 \tag{E835.3}
\]

Every \(H_i\) is stable reflexive. Indeed, for any proper positive-rank
subsheaf \(M\) of \(H_i\) with \(0<\operatorname{rk}M<r_i\), the
intersection \(M\) with \(Q_i\) has the same rank and first Chern class as \(M\):
the quotient is supported in codimension at least two, since \(H_i/Q_i\)
is. Stability of \(Q_i\) then gives \(\mu(M)=\mu(M\cap Q_i)<\mu(Q_i)=\mu(H_i)\). Each \(T_i\) has support dimension at
most two.

The first Chern component of (E835.1) gives \(B=-\alpha\). Equality of the
factor slopes implies \(\int_X\delta_i\theta^3=0\). Hence each
invariant representative \(\delta_i\) is primitive for \(\theta\).

\subsection{3. Mixed discriminant equality eliminates all surface
defects}\label{SF-E835:mixed-discriminant-equality-eliminates-all-surface-defects}

Let \(\Delta\) be the discriminant convention (E834.1). Additivity of Chern
characters in (E835.2) gives the exact variance identity

\[
 \Delta(E)=\sum_i\Delta(Q_i)-\frac12\sum_i r_i\delta_i^2
          =\sum_i\Delta(H_i)+\sum_i\operatorname{ch}_2(T_i)
                         -\frac12\sum_i r_i\delta_i^2.
 \tag{E835.4}
\]

For example, subtracting the first expression from the definition of
\(\Delta(E)\) reduces to \(\sum_i r_i(\alpha_i-\alpha)^2=\sum_i r_i\alpha_i^2-R\alpha^2\), using \(\sum_i r_i\alpha_i=R\alpha\). The second expression uses
\(c_1(Q_i)=c_1(H_i)\) and \(\operatorname{ch}_2(Q_i)=\operatorname{ch}_2(H_i)-\operatorname{ch}_2(T_i)\).

Every mixed pairing on the right side is nonnegative:

\begin{itemize}
\item
  E834 applies to the stable reflexive \(H_i\).
\item
  The leading-character/support-cycle identity makes \(\operatorname{ch}_2(T_i)\) the
  effective surface cycle with positive generic lengths. E818\S1 proves
  that every nonzero effective surface has strictly positive pairing
  with \(\theta_+\theta_-\) on a genuine RM fourfold.
\item
  The elementary primitive mixed sum-of-squares formula E811.4 gives
  \(-\int_X\delta_i^2\theta_+\theta_-\) nonnegative.
\end{itemize}

The left side has zero mixed pairing, by (E835.1). Each nonnegative
summand therefore vanishes. E834's equality case makes every \(H_i\) a
locally free projectively flat stable bundle. Strict mixed positivity
eliminates every two-dimensional component of \(T_i\). The equality case of
E811.4 makes every \(\delta_i\) block diagonal for \(V_+\oplus V_-\) and
primitive on each plane separately. Thus

\[
 \operatorname{ch}(H_i)=r_i e^{\alpha_i},\qquad
 \dim\operatorname{Supp}T_i\le1,
 \qquad \delta_i=\delta_{i,+}+\delta_{i,-},\quad
 \delta_{i,\pm}\theta_\pm=0.
 \tag{E835.5}
\]

The forms in the last identities are invariant forms on the indicated
planes. No claim that those planes are rational subvarieties is made.
The complete strict-surface proof is E818\S1, including its Pluecker and
Albanese argument, not a supposition of general-position positivity.

\subsection{4. The cubic character eliminates all curve
defects}\label{SF-E835:the-cubic-character-eliminates-all-curve-defects}

On each rank-two complex plane, \(\delta_{i,+}^3=\delta_{i,-}^3=0\).
Therefore

\[
 \delta_i^3=3\delta_{i,+}^2\delta_{i,-}
                 +3\delta_{i,+}\delta_{i,-}^2,
 \qquad \theta\delta_i^3=0.
 \tag{E835.6}
\]

Multiplication by \(\theta\) kills the two summands either by excess plane
degree or by the separate primitivity in (E835.5). The actual \(K_0\) identity \([E]=\sum_i[H_i]-\sum_i[T_i]\) and \(B=-\alpha\) give

\[
 0=\left(e^{-\alpha}\operatorname{ch}(E)\right)_3
     =\frac16\sum_i r_i\delta_i^3-\sum_i\operatorname{ch}_3(T_i).
 \tag{E835.7}
\]

Here \(e^{-\alpha}\) leaves the third character of each \(T_i\) unchanged
because \(T_i\) already has support dimension at most one. Multiply by
\(\theta\) and integrate. The first term is zero by (E835.6). Each \(\operatorname{ch}_3(T_i)\)
is its effective curve cycle, with positive generic-length weights. The
ample \(\theta\) has strictly positive degree on every nonzero effective
curve. Thus every \(T_i\) has support dimension zero or is zero. Let \(P\) be
the finite union of their supports and put \(U=X\setminus P\).

The exact leading-character source is the same as in E818: Fulton,
\emph{Intersection Theory}, second edition, Example18.3.11 p363 together
with Theorem18.3(3) p353, since the tangent bundle of \(X\) is trivial. No
sign is lost when replacing a sheaf by its effective leading support
cycle. The actual short exact sequences \(Q_i\) into \(H_i\) supply the
positive weights; this argument would not apply to virtual defect
differences.

\subsection{5. Extensions of vector bundles extend across these
points}\label{SF-E835:extensions-of-vector-bundles-extend-across-these-points}

We give the required algebraic extension step rather than assuming that
a reflexive hull preserves exact sequences. For any two vector bundles
\(V,W\) on the smooth projective fourfold \(X\), restriction gives an
isomorphism

\[
 \operatorname{Ext}^1_X(V,W)
       \xrightarrow{\ \sim\ }
 \operatorname{Ext}^1_U(V|_U,W|_U).
 \tag{E835.8}
\]

Indeed both groups are \(H^1\) of the corresponding bundle \(V^\vee\otimes W\),
since the first argument is locally free. The reduced finite set \(P\) is
locally cut out by four regular parameters. For a locally free sheaf \(K\),
its local-cohomology sheaves with support \(P\) vanish in degrees below four.
Explicitly, the augmented Cech complex on those four parameters computes
local cohomology; regularity of the sequence makes its cohomology vanish
except in degree four, and tensoring with the free stalk of \(K\) preserves this
statement. This is precisely Stacks Lemma36.6.9, tag0G7P, whose complete
proof was read. The local-to-global support spectral sequence therefore
gives \(H_P^j(X,K)=0\) for \(j<4\). The support triangle and its
long exact sequence, Stacks §20.21 tag0A39, now give \(H^1(X,K)\) isomorphic to \(H^1(U,K|_U)\). This proves (E835.8). No analytic/algebraic
comparison on the nonproper \(U\) is used: the sheaves and the actual
extensions being restricted are algebraic; holomorphic sheaves on the
original projective \(X\) may first be algebraized by GAGA.

On \(U\), each \(Q_i\) agrees with \(H_i\). The restricted filtration (E835.2) is
therefore a filtration by vector bundles with vector-bundle quotients.
Inductively construct vector bundles \(K_i\) on \(X\) extending \(E_i|_U\)
with a filtration of the prescribed kind. Set \(K_1=H_1\). Given \(K_{i-1}\),
the actual sequence on \(U\) is

\[
 0\longrightarrow K_{i-1}|_U\longrightarrow E_i|_U
             \longrightarrow H_i|_U\longrightarrow0.
 \tag{E835.9}
\]

Its extension class lifts by (E835.8) to an actual sequence on \(X\)

\[
 0\longrightarrow K_{i-1}\longrightarrow K_i
                    \longrightarrow H_i\longrightarrow0.
 \tag{E835.10}
\]

The middle term is locally free: on an affine local ring the free
quotient \(H_i\) splits the sequence. The new inclusion is a subbundle
inclusion, again because its quotient is locally free. Restriction to \(U\)
identifies (E835.10) with (E835.9), including the indicated subobject and
quotient. Induction produces \(K_s\) with stable locally free quotients
\(H_i\) of the same \(\theta\)-slope.

The original hull \(H\) and \(K_s\) are reflexive and agree on \(U\). The
complement has codimension four, so the usual reflexive Hartogs extension
property identifies them uniquely. Thus \(H\) is \(K_s\) and has the claimed
actual subbundle filtration. The canonical injection \(E\) into \(H\) is an
isomorphism on \(U\), so \(H/E\) is supported on the finite set \(P\).

\subsection{6. Applying the existing categorical
obstruction}\label{SF-E835:applying-the-existing-categorical-obstruction}

All hypotheses of SF-E818 have now been proved for \(E\): its actual hull is
locally free with stable locally free same-slope factors, and it has the
same balanced character. If \(H/E\) is nonzero, the canonical double-dual
sequence and the given coherent action give the point-defect kernel of
SF-E817, of dimension at least \(1+3s\) for \(s\) support orbits. If \(H/E\) is zero,
SF-E813 excludes the locally free sheaf with this filtration.
Equivalently one may invoke the complete theorem E818 directly.

Neither the chosen stable-factor filtration (E835.2) nor its hull factors
need carry separate \(G\) actions. The only action used at the final step is
the specified coherent action on \(E\) and its canonical functorial actions
on \(H\) and \(H/E\), as proved in E817/E818. The scalar twisting and Serre
adjunction inputs are inherited in exactly their reviewed scopes.

\subsection{7. Source and review
boundary}\label{SF-E835:source-and-review-boundary}

The analytic premise E834 remains a separate theorem and has passed its
own whole-proof reviews. The strict surface positivity, leading support
cycle, point-defect and locally free filtration results are the reviewed
E818/E817/E813 statements. All stable-hull, variance, curve, extension-
lifting and actual-filtration steps above are new complete arguments.
The exact local-cohomology inputs were freshly read in full at
\url{https://stacks.math.columbia.edu/tag/0G7P} and
\url{https://stacks.math.columbia.edu/tag/0A39}; the depth form was also
checked at tag0AVZ. No finite numerical certificate is needed or
claimed. An unstable sheaf or a perfect complex outside these actual
sheaf hypotheses remains an open construction route.

\section{SF-E836. Complete-intersection point defects with a commuting
generator
triple}\label{SF-E836:sf-e836.-complete-intersection-point-defects-with-a-commuting-generator-triple}

Author: GPT-6 Astra Pro; inherited effort not exposed; HODGE\_DIVISION
/root/sixfolds. Date: 2026-11-21.

\textbf{Status:} complete new proof checked in full by the author,
/root and /root/imprimitive. The root has also read the entire canonical
conversion; the peer audit concerns the complete frozen draft.
All twenty-one display blocks are unchanged by conversion. It extends the reviewed reduced and curvilinear
point-modification obstructions in a precise common
complete-intersection family. No assertion for arbitrary Artinian
modules or all embedding-dimension-four ideals is made. The necessary
companion equations below were previously retained as an unaudited
calculation in \texttt{rm\_artin\_local\_attempts.md}; they are now
proved in full.

\subsection{1. The local statement and its two concrete
corollaries}\label{SF-E836:the-local-statement-and-its-two-concrete-corollaries}

Let \(R=\mathbb C[[x,y,z,w]]\), with maximal ideal \(\mathfrak m\). Let
\(\pi_+,\pi_-\) be two constant rank-two bivectors with complementary
supports. After independent linear choices inside their two planes and
nonzero scalar normalizations, write

\[
 \pi_+=\partial_x\wedge\partial_y,
 \qquad \pi_-=\partial_z\wedge\partial_w.
 \tag{E836.1}
\]

Let \(I\subset R\) be a proper ideal of positive finite colength
generated by a regular sequence \(f_1,f_2,f_3,f_4\). Put \(A=R/I\) and
\(\ell=\dim_{\mathbb C}A>0\). Suppose that

\[
 \{f_j,f_k\}_+\in I\qquad(2\le j<k\le4).
 \tag{E836.2}
\]

Thus a particular triple of the actual ideal generators Poisson-commutes
modulo the ideal for one of the two bivectors. Consider any finite
\(R\)-module with an exact hull-modification sequence

\[
 0\longrightarrow A^n\longrightarrow M\xrightarrow{\phi}R^r
      \longrightarrow A^m\longrightarrow0,
 \qquad m+n>0.
 \tag{E836.3}
\]

Then \(M\) cannot have flat first-order module deformations for both
multiplication classes \(\pi_+\) and \(\pi_-\).

The hypothesis applies in each of the following cases.

\begin{enumerate}
\def\labelenumi{\arabic{enumi}.}
\item
  \textbf{Every common complete-intersection quotient of embedding
  dimension at most three.} The ideal contains a function with a nonzero
  linear part. Its linear part may mix both RM cotangent planes; no
  pure-plane conormal assumption is imposed. Section 5 proves the
  required generator reduction without changing the constant Poisson
  coordinates nonlinearly.
\item
  \textbf{Every presentation with three generators independent of one
  coordinate in one RM plane.} For example, \(I=(x^a,y^b,z^c,w^d)\),
  with arbitrary positive integers \(a,b,c,d\), satisfies (E836.2),
  including the embedding-dimension-four case \(a,b,c,d\ge2\). The first
  generator may depend on all four coordinates; only the indicated three
  need be independent of \(x\).
\end{enumerate}

The second case is a sufficient presentation condition, not a
classification of complete intersections. The remaining
full-embedding-dimension-four complete-intersection gate is not silently
identified with this case.

\subsection{2. Extension matrices and the map to the free deformed
hull}\label{SF-E836:extension-matrices-and-the-map-to-the-free-deformed-hull}

For a bivector \(\pi\), let \(D=\mathbb C[\varepsilon]/(\varepsilon^2)\)
and let \(R_\pi=R\oplus\varepsilon R\) have product

\[
 a*b=ab+\frac{\varepsilon}{2}\pi(da,db).
 \tag{E836.4}
\]

This is the same continuous formal first-order algebra used in E819. In
particular \([a,b]_* =\varepsilon\pi(da,db)\).

If either nonzero bivector deformation exists, E819's
finite-perfect-cone supertrace lemma gives \(\ell m=\ell n\). If either
multiplicity vanishes, this contradicts \(m+n>0\). We therefore assume
\(m=n>0\). The remainder of the proof is independent of the magnitude of
\(m\).

Choose a free-hull basis that identifies its quotient to \(A^m\) with
reduction of the first \(m\) coordinates. Indeed, a surjective matrix
over the local ring \(A\) has a unit \(m\)-minor; lift elementary
operations and that unit minor to \(R\). The image of \(\phi\) is then

\[
 K=I^{\oplus m}\oplus R^{r-m}.
 \tag{E836.5}
\]

The free summand has no \(\operatorname{Ext}^1_R(-,A^m)\), so it splits
in the ordinary extension of \(K\) by \(T=A^m\). No assertion that this
chosen splitting survives a deformation is needed.

Choose elements \(e_i(v)\in M\), depending \(R\)-linearly on
\(v\in R^m\), that lift \(f_i v\) in the first \(m\) hull coordinates.
For every pair there are matrices
\(B_{ij}=-B_{ji}\in\operatorname{Mat}_m(A)\) such that

\[
 f_i e_j(v)-f_j e_i(v)=B_{ij}\bar v\in T.
 \tag{E836.6}
\]

These are exactly the extension matrices: the Koszul resolution of
\(A\), followed by applying \(\operatorname{Hom}_R(-,A^m)\), has all
differentials zero because \(I A^m=0\). Thus

\[
 \operatorname{Ext}^1_R(K,A^m)
   \simeq\operatorname{Ext}^2_R(A^m,A^m)
   \simeq\operatorname{Mat}_m(A)\otimes\bigwedge^2\mathbb C^4.
 \tag{E836.7}
\]

Now assume a flat first-order deformation \(M_\varepsilon\) for \(\pi\).
The map \(\phi\) lifts to an \(R_\pi\)-linear map
\(\phi_\varepsilon:M_\varepsilon\to R_\pi^r\). We recall why the
reviewed E819 lifting applies: finite-length modules over the regular
local fourfold have \(\operatorname{Ext}^j(-,R)=0\) for \(j<4\), by a
finite-length filtration and the four-variable Koszul resolution. The
two exact sequences in (E836.3) therefore give
\(\operatorname{Ext}^1_R(M,R^r)=0\). Lift a finite free resolution and
the degree-zero map. The map-lifting obstruction is precisely this zero
group. Flatness permits the lifted resolution and all kernel
identifications over the dual numbers, as in E819 section 2.1. This uses
actual continuous finite free complexes, not discontinuous operations on
formal power series.

For \(t\in T\), choose a lift \(\widetilde t\in M_\varepsilon\). Its
hull image has the form \(\varepsilon c\), for \(c\in R^r\). Changing
the lift changes \(c\) by an element of \(K\), so its class in \(Q=A^m\)
is well-defined. This defines an \(A\)-linear map

\[
 C_\pi:T\longrightarrow Q,
 \qquad C_\pi\in\operatorname{Mat}_m(A).
 \tag{E836.8}
\]

For linearity, apply the deformed hull map to \(a*\widetilde t\). The
zero reduction of the hull image leaves \(\varepsilon ac\), so the
ordinary scalar action is the induced one. Because \(I\) annihilates
both modules, this factors through \(A\).

\subsection{3. Both companion equations are
necessary}\label{SF-E836:both-companion-equations-are-necessary}

Set

\[
 P^\pi_{ij}=\pi(df_i,df_j)\pmod I\in A.
\]

The previously reviewed hull relation gives

\[
 C_\pi B_{ij}=P^\pi_{ij}I_m.
 \tag{E836.9}
\]

For clarity, lift (E836.6) to the deformed module and apply
\(\phi_\varepsilon\). The hull images of the lifted \(e_i\) are
\(f_i v+\varepsilon a_i\). Their commutator contribution is
\(\varepsilon P^\pi_{ij}v\). The terms \(f_i a_j-f_j a_i\) vanish in the
first \(m\) coordinates modulo \(I\). The correction to the lifted
relation reduces to an element of \(M\), whose hull image is in \(K\),
and vanishes in the same quotient. This proves (E836.9), with the stated
sign convention.

There is also the companion equation in the opposite matrix order:

\[
 B_{ij}C_\pi=P^\pi_{ij}I_m.
 \tag{E836.10}
\]

Here is its full derivation. Since \(f_i t=0\), there exists
\(\widetilde m_i\in M_\varepsilon\) such that
\(f_i*\widetilde t=\varepsilon\widetilde m_i\). Write \(m_i\) for its
reduction. Flatness identifies \(\varepsilon M_\varepsilon\) with \(M\),
so \(m_i\) is determined by this equation. Applying the hull map gives

\[
 \phi(m_i)=f_i c.
 \tag{E836.11}
\]

Split \(c=(c_1,c_2)\) according to the first \(m\) and the remaining
free coordinates. Choose an ordinary splitting of the free part in
(E836.5), denoted \(s:R^{r-m}\to M\). Equation (E836.11) implies

\[
 m_i=e_i(c_1)+s(f_i c_2)+t_i,
 \qquad t_i\in T.
 \tag{E836.12}
\]

Compute the commutator of \(f_i,f_j\) on \(\widetilde t\). On one hand
it is \(\varepsilon P^\pi_{ij}t\). On the other hand it is
\(\varepsilon(f_i m_j-f_j m_i)\). The terms from \(t_i\) vanish because
\(I T=0\), and those from \(s\) cancel by ordinary commutativity. The
remaining terms, by (E836.6), equal \(B_{ij}\bar c_1=B_{ij}C_\pi(t)\).
This proves (E836.10) for every \(t\). No derivative of a chosen lift of
\(c_1\) is omitted: this calculation is already the coefficient of
\(\varepsilon\), where the module action is the ordinary commutative
one.

Apply the two companion equations separately to \(\pi_+\) and \(\pi_-\),
writing \(C_+,C_-\) and \(P^+,P^-\). Associativity gives

\[
 \begin{aligned}
 P^+_{ij}P^-_{kl}I_m
   &=(C_+B_{ij})(C_-B_{kl})\\
   &=C_+(B_{ij}C_-)B_{kl}
     =P^-_{ij}P^+_{kl}I_m.
 \end{aligned}
 \tag{E836.13}
\]

Since \(m>0\), equality of scalar matrices implies the scalar equalities

\[
 P^+_{ij}P^-_{kl}-P^-_{ij}P^+_{kl}=0
          \quad\text{in }A
          \quad\text{for every }i,j,k,l.
 \tag{E836.14}
\]

No coefficient has been divided out in this argument. In particular it
is valid over the nonreduced Artinian ring even when every bracket
coefficient is a zero divisor.

\subsection{4. The commuting triple makes the mixed wedge vanish term by
term}\label{SF-E836:the-commuting-triple-makes-the-mixed-wedge-vanish-term-by-term}

Regard \(P^\pm\) as alternating two-forms on a free \(A\)-module with
basis indexed by the four ideal generators. Hypothesis (E836.2) gives
\(P^+_{jk}=0\) for \(j,k\ge2\). Equation (E836.14), used for the pairs
\((1,i)\) and \((j,k)\) with \(\{i,j,k\}=\{2,3,4\}\), consequently gives

\[
 P^+_{1i}P^-_{jk}=0.
 \tag{E836.15}
\]

Every term of the coefficient of \(P^+\wedge P^-\) is either one of
(E836.15), with a sign, or contains \(P^+_{jk}=0\). Hence that mixed
wedge is zero.

Let \(J=(\partial f_i/\partial z_j)\), with coordinates ordered as
\((x,y,z,w)\). The bracket forms satisfy
\(P^\pm=J\pi_\pm J^{\mathsf t}\) modulo \(I\). Functoriality of the
fourth exterior power gives

\[
 P^+\wedge P^-
    =\det(J)(\pi_+\wedge\pi_-).
 \tag{E836.16}
\]

The right side is nonzero. The two constant bivectors have complementary
supports, and the Jacobian determinant is nonzero in \(A\): for a
zero-dimensional complete intersection over a characteristic-zero field,
the Jacobian ideal is its one-dimensional socle. The exact published
source is Eisenbud--Ulrich, \emph{Duality and socle generators for
residual intersections}, \textbf{J. reine angew. Math. 756 (2019),
183--226, Theorem 7.1, p.210}, with the complete-intersection proof in
\textbf{Appendix A, Theorem A.5, pp.224--225}, using Theorem A.1,
pp.222--223, and Theorem A.3, p.224. The complete invoked statements and
proofs were read in the published author-hosted copy. Theorem 7.1
explicitly uses complete local power-series algebras, so the present
formal setting is included. The quotient is nonzero and finite, and
characteristic zero does not divide its positive length. These are
precisely the hypotheses needed here.

Equations (E836.15)--(E836.16) contradict one another. This proves the
local theorem. Unlike a proposed unrestricted argument from coefficient
minors alone, this proof uses the commuting triple to kill the mixed
wedge term by term; it does not infer proportionality of two forms over
an arbitrary nonreduced ring.

\subsection{5. Every embedding-dimension-at-most-three complete
intersection has the required
presentation}\label{SF-E836:every-embedding-dimension-at-most-three-complete-intersection-has-the-required-presentation}

For \(A=R/I\),
\(\operatorname{edim}A=\dim_{\mathbb C}\mathfrak m/(\mathfrak m^2+I)\).
If this is at most three, there is \(f\in I\) with \(df(0)\ne0\). At
least one of its two RM-plane components is nonzero. Interchange signs
if needed and make a linear change only within that plane so that
\(\partial f/\partial x(0)\ne0\). Rescale the corresponding constant
bivector if necessary. The other plane is unchanged, and both bivectors
still have the form (E836.1). The linear part of \(f\) is allowed to
have components in the other plane.

The element \(f\) is not in \(\mathfrak m I\), since
\(I\subset\mathfrak m\) and \(f\notin\mathfrak m^2\). Extend its nonzero
class to a basis of \(I/\mathfrak m I\), and lift to generators
\(f,h_2,h_3,h_4\) of \(I\). This is an invertible generator change from
any minimal four-generator presentation. These generators are a regular
sequence: in the regular local ring, four generators of the height-four
ideal form a system of parameters and hence a regular sequence.

The formal implicit-function construction gives a unique
\(\varphi\in\mathbb C[[y,z,w]]\), with zero constant term, such that

\[
 f=u\bigl(x-\varphi(y,z,w)\bigr),\qquad u\in R^\times.
 \tag{E836.17}
\]

One may obtain the solution successively by total degree, since the
coefficient of \(x\) in the linear part is nonzero. Substitution makes
\(f\) vanish; formal division by the monic degree-one expression gives
(E836.17), with a unit quotient. In particular \(\varphi\) may have a
nonzero linear \(z\)- or \(w\)-term. We replace the first ideal
generator by \(f_1=x-\varphi\), using the unit \(u\); \textbf{we do not
use \(f_1\) as a new Poisson coordinate}.

For \(j\ge2\), substitute \(x=\varphi\) in \(h_j\), obtaining
\(g_j\in\mathbb C[[y,z,w]]\). The difference \(h_j-g_j\) is divisible by
\(x-\varphi\); subtract this multiple of the first generator. These are
invertible triangular generator operations. We have obtained

\[
 I=(x-\varphi(y,z,w),g_2(y,z,w),g_3(y,z,w),g_4(y,z,w)).
 \tag{E836.18}
\]

Because each \(g_j\) is independent of the unchanged coordinate \(x\),
its bracket with every other \(g_k\) for
\(\pi_+=\partial_x\wedge\partial_y\) is zero as a formal function, not
merely modulo \(I\). Thus (E836.2) holds even for a mixed linear part of
\(f\).

All equations in sections 2--4 may simply be derived for this new
regular sequence. Equivalently their covariance is explicit. For a
generator change \(f'=Uf\), \(U\in\operatorname{GL}_4(R)\),
differentiation modulo \(I\) gives \(df'=Udf\), so
\(P'=UPU^{\mathsf t}\) in \(A\). Taking lifts \(e'_i=\sum_a U_{ia}e_a\)
in the ordinary commutative module gives
\(B'_{ij}=\sum_{a,b}U_{ia}U_{jb}B_{ab}\) with the alternating
convention. This is the second exterior-power change in the Koszul
extension matrices. The map \(C:T\to Q\) is unchanged. Thus both
companion equations are preserved. The Jacobian determinant changes by
\(\det(U)\) modulo \(I\), a unit. No untracked derivative terms survive
modulo \(I\).

\subsection{6. Global balanced consequence and actual noncurvilinear
examples}\label{SF-E836:global-balanced-consequence-and-actual-noncurvilinear-examples}

Let a balanced coherent sheaf on the genuine RM fourfold have a locally
free hull and finite kernel and cokernel of its canonical hull map. At
any point, suppose those two defect modules are free over a common
quotient satisfying the theorem, with positive total multiplicity. The
reviewed E819 characteristic-map argument supplies two global cap-kernel
classes whose bivector components are \(\pi_+,\pi_-\). If the sheaf were
weakly semiregular for a specified coherent effective finite
translation/tensor action, invariant evaluation would vanish on these
classes. Forgetting equivariance and restricting to the completed local
ring gives flat first-order deformations for both bivectors, as proved
there. The local theorem contradicts this. The rational algebraic \(B\)
shear leaves the bivector projections unchanged. No action on a chosen
presentation or on individual local defect bases is assumed.

There are actual noncurvilinear balanced sheaves to which this new
statement applies. Take the rank-32 vector bundle \(H\) constructed and
fully checked in E815/E820, with \(\operatorname{ch}(H)=32v_{1/2}\). At
one point choose a finite subscheme whose completed ideal, in constant
RM-plane coordinates, is

\[
 I_3=(x+z,y^2,z^2,w^2)
       \quad\text{or}\quad
 I_4=(x^2,y^2,z^2,w^2).
 \tag{E836.19}
\]

The first quotient has length eight and embedding dimension three, with
a mixed linear conormal; the second has length sixteen and embedding
dimension four and has the three-generator presentation from case 2.
These are Artinian complete intersections. They define actual algebraic
finite subschemes: each contains a sufficiently high power of the
maximal ideal, and the local ring and its completion have identical
finite quotients. Thus the finitely many required formal jets are
algebraic local functions. No algebraicity of the whole formal logarithm
is asserted.

For either subscheme \(Z\), choose \(1\le m\le32\), a quotient
\(H\twoheadrightarrow Q=\mathcal O_Z^m\), its kernel \(K\), and
\(T=\mathcal O_Z^m\). Such a quotient exists because a vector bundle
over the local Artinian scheme is free. Since \(H\) is locally free and
\(T\) is point-supported, all its positive Ext groups into \(T\) vanish.
The exact sequence for \(K\) and the local Koszul resolution give

\[
 \operatorname{Ext}^1_X(K,T)
   \simeq\operatorname{Ext}^2_X(Q,T),
 \qquad
 \dim_{\mathbb C}\operatorname{Ext}^1_X(K,T)=6m^2\ell>0.
 \tag{E836.20}
\]

Every extension class gives an actual coherent sequence
\(0\to T\to E\to K\to0\). Its torsion is exactly \(T\), because \(K\) is
torsion-free. Its double dual is \(H\), its hull cokernel is \(Q\), and
its entire character is
\(\operatorname{ch}(H)+\operatorname{ch}(T)-\operatorname{ch}(Q)=32v_{1/2}\).
It is ordinarily gluable as a coherent sheaf. The new theorem excludes
weak semiregularity for every specified coherent finite action, if such
an action exists. No simplicity, indecomposability, or action is
asserted.

These examples have equal signed point defect and are not curvilinear.
The first-order obstruction therefore supplies a strict extension of
E819's signed-length statement and E820's common-curvilinear family.
Arbitrary common complete intersections without a commuting triple,
unequal Artinian module types, positive-dimensional singularities,
unstable torsion-free sheaves, and unrelated perfect complexes remain
outside the theorem.

\subsection{7. Source and certificate
boundary}\label{SF-E836:source-and-certificate-boundary}

The new published source is the 2019 Eisenbud--Ulrich paper at the exact
locators above, DOI \texttt{10.1515/crelle-2017-0045}, published primary
PDF:
\url{https://www.math.purdue.edu/~bulrich/papers/published\%20version\%20of\%20Duality\%20for\%20Residual\%20Intersections.pdf}.
Theorem 7.1 and the entire invoked Appendix A proofs were freshly read.
This is an invoked published Jacobian-socle theorem, not an independent
reproof of all upstream different theory.

The finite exact supplement checks the generic twelve-variable
eliminated Jacobian/bracket identity, the noncommutative associative
expansion underlying (E836.13), and the two explicit Artinian socles. It
is not a substitute for the module-lifting, covariance,
formal-generator, global-action, or published-socle arguments. The older
failed matrix solution and the nonintegrable Jacobian counterexample
remain in \texttt{rm\_artin\_local\_attempts.md}; no unrestricted
inference from coefficient minors has been revived.

\section{An integrable boundary for the complete-intersection argument}
\label{sec:ci-cross-minor-boundary}

\paragraph{Status and scope.}
The full derivation and exact supplement have passed the root check and
an independent whole-proof check by /root/sixfolds. This is a
counterexample to an overstrong scalar inference. It supplies neither a
module deformation nor a weakly semiregular object and does not reduce
the main supplier gate.

\begin{proposition}[HD-CI-B1: scalar cross minors do not force a zero Jacobian]
\label{HD-CI-B1}
There is a length-sixteen complete intersection
$A=\mathbf C[[x,y,z,w]]/(f_1,f_2,f_3,f_4)$ of embedding dimension four
with complementary constant bivectors
$\pi_+=\partial_x\wedge\partial_y$ and
$\pi_-=\partial_z\wedge\partial_w$ such that all scalar minors
\[
 P^+_{ij}P^-_{kl}-P^-_{ij}P^+_{kl}=0\quad\hbox{in }A,
 \qquad P^{\pm}_{ij}=\{f_i,f_j\}_{\pi_\pm},
\]
while $\det\operatorname{Jac}(f)$ is a nonzero socle element.
No regular generating sequence for this ideal has the commuting-triple
property of SF-E836 with these same two bivectors.
\end{proposition}

\begin{proof}
Let $J$ be the rational matrix with successive columns
\[
 u=(1,0,-3/5,-8/5)^t,\quad v=(0,1,1,1)^t,\quad
 u'=(1,0,-9/5,-16/5)^t,\quad v'=(0,1,2,3)^t.
\]
In the coordinate order $(12,13,14,23,24,34)$, the two column wedges are
\[
 p=u\wedge v=(1,1,1,3/5,8/5,1),\qquad
 q=u'\wedge v'=(1,2,3,9/5,16/5,1).
\]
Both are decomposable. The ratios $q_{ij}/p_{ij}$ agree within each
complementary pair, with values respectively $1,2,3$. Hence every
complementary scalar coefficient minor vanishes. Nevertheless
\[
 p\wedge q=2(1-16/5+9/5)=-4/5=\det J\ne0.
\]
Put $\ell=J(x,y,z,w)^t$ and $f_i=\ell_i^2$. Invertibility gives
\[
 A\simeq\mathbf C[[\ell_1,\ell_2,\ell_3,\ell_4]]/
                 (\ell_1^2,\ell_2^2,\ell_3^2,\ell_4^2).
\]
The sixteen squarefree monomials are a basis, so the asserted length
and embedding dimension follow. These are actual power series, whose
actual gradient matrix is
\[
 \operatorname{Jac}(f)=2\operatorname{diag}(\ell_1,\ell_2,\ell_3,\ell_4)J.
\]
It follows that
$P^+_{ij}=4\ell_i\ell_jp_{ij}$ and
$P^-_{ij}=4\ell_i\ell_jq_{ij}$.
If the pairs $(ij)$ and $(kl)$ share an index, both terms in the scalar
minor contain a square $\ell_a^2$ and vanish in $A$.
If they are disjoint, the scalar coefficient minor vanishes by the
explicit complementary-pair calculation. This proves all the stated
minor identities, including coincident and ordered pairs.

The actual Jacobian determinant is
\[
 \det\operatorname{Jac}(f)=-\frac{64}{5}\ell_1\ell_2\ell_3\ell_4.
\]
The top squarefree monomial is nonzero and is annihilated by the maximal
ideal, so it is a nonzero socle element. The coefficient of the mixed
bracket wedge is the same nonzero element.

For completeness, replacing a minimal regular generating sequence by
another multiplies its column of generators by an invertible matrix
$U$ modulo the ideal. Modulo that ideal the derivative terms involving
$U$ disappear, so both bracket arrays transform by $\bigwedge^2U$.
The vanishing of all scalar minors is preserved, while their mixed
wedge is multiplied by $\det U$, a unit. If a new generator triple
commuted under $\pi_+$, each surviving mixed-wedge term would be one
of those vanishing scalar minors. This would contradict its invariant
nonvanishing. Thus no such generator presentation exists.
\end{proof}

\paragraph{Logical boundary and exact checks.}
The two companion matrix orders of SF-E836 imply the displayed scalar
minor identities, but HD-CI-B1 does not solve the companion equations.
It does not construct a flat deformation of any coherent hull-modification
module. The root and separate peer programs check all thirty-six ordered
minor equations, all actual-gradient entries, and the nonzero socle;
these checks supplement the complete two-case proof above.
In the common Grothendieck framework this is a local Artin scheme and
$K$-theoretic deformation boundary for a proposed condition-A supplier
(cote16, scan050; SIGN\_NEXT SN-C03/C04), not a condition-I production
of cycles (scan049, cote44, Theorem4). No change to the main detector
gate or to the conditional status of Markman, Severn or032 follows.

\section{SF-E837. Every formal square-thickening point modification is
excluded}\label{SF-E837:sf-e837.-every-formal-square-thickening-point-modification-is-excluded}

Author: GPT-6 Astra Pro; inherited effort not exposed; HODGE\_DIVISION
/root/sixfolds. Date: 2026-11-21.

\textbf{Status:} complete new proof checked in full by the author,
/root and /root/g34\_eightfolds. The root also read the complete
canonical conversion, with all seventeen display blocks unchanged.
It uses the opposite-order companion equations proved in E836
sections 2--3, whose separate whole-proof reviews have passed. It does not use E836's commuting-triple hypothesis. The
integrable scalar-cross-minor boundary remains valid; the new argument
uses a strictly stronger actual matrix constraint.

\subsection{1. Statement}\label{SF-E837:statement}

Let \(R=\mathbb C[[z_1,z_2,z_3,z_4]]\), and let \(\pi_+,\pi_-\) be
constant rank-two bivectors with complementary supports. Let
\(u_1,u_2,u_3,u_4\) be any formal regular parameters, possibly nonlinear
and mixing both planes. Set

\[
 I=(u_1^2,u_2^2,u_3^2,u_4^2),
 \qquad A=R/I,
 \qquad\mathfrak n=(\bar u_1,\bar u_2,\bar u_3,\bar u_4)\subset A.
 \tag{E837.1}
\]

Consider a finite module with an exact sequence

\[
 0\longrightarrow A^n\longrightarrow M\longrightarrow R^r
     \longrightarrow A^m\longrightarrow0,
 \qquad m+n>0.
 \tag{E837.2}
\]

Then \(M\) has no pair of flat first-order module deformations for the
two multiplication classes \(\pi_+,\pi_-\). The assertion holds for
every positive multiplicity and every formal orientation of the square
thickening. It consequently excludes weak balanced semiregularity for a
coherent sheaf with such a nontrivial local point modification and any
specified coherent effective finite translation/tensor action in E819's
scope.

The proof has two ingredients: a universal matrix lemma over the
length-sixteen square algebra and the constant terms of the two
transformed bivectors. It does not identify a vanishing coefficient
cross minor with a vanishing Jacobian.

\subsection{2. Four monomial annihilators force a square matrix to be
invertible}\label{SF-E837:four-monomial-annihilators-force-a-square-matrix-to-be-invertible}

We write \(t_i\) for \(\bar u_i\), identifying
\(A=\mathbb C[t_1,t_2,t_3,t_4]/(t_1^2,t_2^2,t_3^2,t_4^2)\). Its
degree-one and degree-two pieces are

\[
 V=\operatorname{span}_{\mathbb C}(t_1,t_2,t_3,t_4),
 \qquad W=\operatorname{span}_{\mathbb C}(t_it_j: i<j),
 \qquad \dim V=4,\quad\dim W=6.
 \tag{E837.3}
\]

\textbf{Matrix lemma.} Let \(C\in\operatorname{Mat}_m(A)\). If the
cokernel of \(C:A^m\to A^m\) is annihilated by at least four distinct
monomials \(t_it_j\), then \(C\) is invertible.

First split off all unit pivots by invertible row and column operations
over the local ring. If a unit entry exists, scale it to one, clear its
row and column, and repeat on the remaining square block. The process
terminates in a block decomposition with an identity block and a square
block \(C'\) of size \(k\), all of whose entries belong to
\(\mathfrak n\). The cokernel is unchanged up to isomorphism. If
\(k=0\), the claim holds. Assume \(k>0\); the four monomials still
annihilate \(\operatorname{coker}C'\).

Work modulo \(\mathfrak n^3\) and write the matrix as \(C_1+C_2\), where
the entries are homogeneous of degrees one and two. Let

\[
 L:\mathbb C^k\longrightarrow V\otimes\mathbb C^k,
 \qquad L(v)=C_1v,
 \qquad a=\operatorname{rank}_{\mathbb C}L\le k.
 \tag{E837.4}
\]

Any input modulo \(\mathfrak n^2\) is a sum of a constant vector \(v_0\)
and a linear vector \(v_1\). Its image modulo \(\mathfrak n^3\) is
\(C_1v_0+C_2v_0+C_1v_1\). Such an image has no degree-one term precisely
when \(v_0\in\ker L\). Therefore the intersection of the image with the
pure degree-two target is exactly

\[
 \operatorname{im}(C'\bmod\mathfrak n^3)
       \cap(W\otimes\mathbb C^k)
     = C_2(\ker L)+ V\cdot\operatorname{im}L.
 \tag{E837.5}
\]

The second summand has dimension at most \(4a\), because it is the image
of \(V\otimes\operatorname{im}L\). The first has dimension at most
\(k-a\). Thus the dimension in (E837.5) is at most

\[
 k+3a\le4k.
 \tag{E837.6}
\]

Choose four distinct pair monomials that annihilate the cokernel and let
\(J\subset W\) be their four-dimensional span. Annihilation means that
each of these monomials times every constant target vector belongs to
the image of \(C'\). Hence \(J\otimes\mathbb C^k\), of dimension \(4k\),
is contained in (E837.5). Equality must hold throughout (E837.6):
\(a=k\), \(\ker L=0\), and

\[
 V\cdot\operatorname{im}L=J\otimes\mathbb C^k.
 \tag{E837.7}
\]

Write \(C_1=\sum_i t_i A_i\), with constant \(k\)-by-\(k\) matrices
\(A_i\). For any pair \(\{i,j\}\) whose monomial is not in the chosen
four, multiply \(C_1v\) by \(t_i\). Its coefficient at \(t_it_j\) is
\(A_jv\), which must be zero by (E837.7). Multiplication by \(t_j\)
similarly gives \(A_i=0\). There are exactly two omitted distinct edges
among the six pairs of four vertices. Their endpoints cover at least
three vertices. Consequently at most one of the matrices \(A_i\) can be
nonzero.

If the one possible index is \(s\), then \(C_1=t_sA_s\), and the
dimension of \(V\cdot\operatorname{im}L\) is at most
\(3\operatorname{rank}A_s\le3k\), since \(t_s^2=0\). This contradicts
(E837.7). Thus \(k>0\) was impossible, proving the matrix lemma. The
argument applies in every size; no simultaneous diagonalization or
matrix commutativity was assumed.

\subsection{3. The actual companion equations in the square
algebra}\label{SF-E837:the-actual-companion-equations-in-the-square-algebra}

If a deformation exists for either nonzero bivector, the E819
finite-cone trace gives \(16n=16m\). We may therefore assume \(n=m>0\).
For the regular generators \(f_i=u_i^2\), the actual module construction
and both matrix-order calculations of E836 give fixed matrices
\(B_{ij}\) and two matrices \(C_+,C_-\) in \(\operatorname{Mat}_m(A)\)
with

\[
 C_\pm B_{ij}=B_{ij}C_\pm=P^\pm_{ij}I_m,
 \qquad
 P^\pm_{ij}=4t_it_j\,a^\pm_{ij}(t),
 \tag{E837.8}
\]

where \(a^\pm_{ij}\) is the image of the actual formal function
\(\pi_\pm(du_i,du_j)\) in \(A\). These functions need not be constant.
Their values at the origin define alternating forms

\[
 p_{ij}=a^+_{ij}(0),\qquad q_{ij}=a^-_{ij}(0).
 \tag{E837.9}
\]

The derivative matrix of the formal parameter change is invertible.
Hence \(p\) and \(q\) are decomposable rank-two forms with complementary
supports, so

\[
 p\wedge p=q\wedge q=0,
 \qquad p\wedge q\ne0.
 \tag{E837.10}
\]

The scalar cross-minor consequences of the two matrix orders are
\(P^+_{ij}P^-_{kl}=P^-_{ij}P^+_{kl}\). For disjoint pairs, all higher
terms of the functions \(a^\pm\) disappear after multiplying by the top
squarefree monomial. That monomial is nonzero in \(A\). Thus

\[
 p_{ij}q_{kl}=q_{ij}p_{kl}
       \quad\text{whenever }\{i,j,k,l\}=\{1,2,3,4\}.
 \tag{E837.11}
\]

This reduction is legitimate even for a nonlinear formal parameter
change; it uses the explicit square-algebra socle, not a linearity
assumption on that change.

\subsection{4. The small-support case already contradicts scalar
minors}\label{SF-E837:the-small-support-case-already-contradicts-scalar-minors}

Let \(\mathcal E\) be the set of unordered pairs \(\{i,j\}\) for which
at least one of \(p_{ij},q_{ij}\) is nonzero. Suppose first that
\(|\mathcal E|\le3\). Among at most three edges there is at most one
pair of disjoint complementary edges: two distinct complementary pairs
require four different edges. If there is no complementary pair, the
mixed wedge is zero, contradicting (E837.10).

For the possible unique complementary pair, write the two coefficients
of \(p\) as \(a,b\) and those of \(q\) as \(c,d\), with a fixed
ordering. The decomposable identities (E837.10) give \(ab=cd=0\), since
this is the only possible term in each Pfaffian. Equation (E837.11)
gives \(ad=cb\). If their common value is \(z\), then

\[
 z^2=(ad)(cb)=(ab)(cd)=0.
 \tag{E837.12}
\]

These are complex numbers, so \(z=0\). The only possible mixed-wedge
coefficient is, up to sign, \(ad+cb=0\), again a contradiction.
Therefore

\[
 |\mathcal E|\ge4.
 \tag{E837.13}
\]

No analogous step is made for arbitrary coefficients in a nonreduced
ring; only the constant complex coefficients are used here.

\subsection{5. The large-support case forces an impossible scalar
ratio}\label{SF-E837:the-large-support-case-forces-an-impossible-scalar-ratio}

Choose a nonzero complex number \(\lambda\) outside the finite set of
roots of \(p_{ij}+\lambda q_{ij}\) for \(\{i,j\}\in\mathcal E\). Define
the actual matrix pencil

\[
 C_\lambda=C_++\lambda C_-,
 \qquad P^\lambda_{ij}=P^+_{ij}+\lambda P^-_{ij}.
 \tag{E837.14}
\]

Linearity of (E837.8) gives
\(C_\lambda B_{ij}=B_{ij}C_\lambda=P^\lambda_{ij}I_m\). This is an
algebraic matrix identity; the proof does not need to assume an
additional deformation for every parameter value.

For every edge in \(\mathcal E\),
\(P^\lambda_{ij}=4t_it_j(a^+_{ij}+\lambda a^-_{ij})\), and the factor in
parentheses is a unit of the local Artinian ring by the choice of
\(\lambda\). Consequently each \(t_it_j\) in at least four distinct
pairs annihilates \(\operatorname{coker}C_\lambda\). The matrix lemma
proves that \(C_\lambda\) is invertible.

Put \(D=C_+C_\lambda^{-1}\). Solving the first pencil equation for
\(B_{ij}\), and substituting in its plus equation, gives

\[
 D P^\lambda_{ij}=P^+_{ij}I_m.
 \tag{E837.15}
\]

Let \(D_0\) be the constant matrix modulo \(\mathfrak n\). The
degree-two component of (E837.15), for every edge in \(\mathcal E\), is

\[
 (p_{ij}+\lambda q_{ij})D_0=p_{ij}I_m.
 \tag{E837.16}
\]

One such nonzero coefficient forces \(D_0=d I_m\) for a complex scalar
\(d\). All remaining equations imply \(p=d(p+\lambda q)\) as complete
alternating forms, with zero coefficients off \(\mathcal E\) as well.
Since \(p\ne0\), \(d\ne0\). Since \(\lambda\ne0\), this makes \(p\) and
\(q\) linearly dependent, contradicting their complementary supports in
(E837.10). This completes the local theorem.

Both parts of the proof are necessary: the dense-support argument uses
the actual square matrices and their cokernels, while the small-support
argument uses only constant Pluecker and cross-minor identities. The
unrestricted scalar-cross-minor implication refuted by the integrable
boundary example is never invoked.

\subsection{6. Global consequence and actual excluded
objects}\label{SF-E837:global-consequence-and-actual-excluded-objects}

The global passage is precisely the reviewed E819 one. A balanced weakly
semiregular sheaf, with its specified coherent finite translation/tensor
action, makes the invariant evaluation vanish on the two balanced
cap-kernel classes. Forgetting the action and restricting to the
completed local ring gives flat first-order deformations for the two
constant bivectors. A rational algebraic \(B\) shear preserves those
bivector components. Thus a single nontrivial point modification of type
(E837.2) is impossible. The chosen formal parameters and individual
defect bases do not need equivariant structures.

There are actual ordinarily gluable balanced coherent sheaves in this
family before semiregularity is imposed. Take the rank-32 bundle \(H\)
of E815/E820, with full character \(32v_{1/2}\), and any finite
subscheme \(Z\) with the completed ideal (E837.1). It is an actual
algebraic subscheme: the ideal contains a power of the maximal ideal,
and finite quotients of the local ring and its completion agree. Choose
a quotient \(H\twoheadrightarrow\mathcal O_Z^m\), with \(1\le m\le32\),
let \(K\) be its kernel, and choose any extension

\[
 0\longrightarrow\mathcal O_Z^m\longrightarrow E\longrightarrow K
       \longrightarrow0.
 \tag{E837.17}
\]

As in E836, the extension space has dimension \(96m^2\), because the
quotient has length sixteen and a four-generator Koszul resolution. The
torsion of \(E\) is the displayed point module, its double dual is
\(H\), and its hull cokernel is the same point module. Therefore its
entire character is \(32v_{1/2}\). It is ordinarily gluable because it
is a coherent sheaf. Every such object is excluded from weak
semiregularity by the theorem, without any assertion that a prescribed
finite action actually exists.

In particular the parent's checked integrable boundary ideal, defined by
four squared mixed rational linear forms and having all scalar cross
minors zero but a nonzero mixed wedge, is now excluded by the stronger
matrix argument. Its counterexample to the scalar inference is
unchanged. This is a strict new family beyond the E836 commuting-triple
criterion, not a reclassification of that boundary example as a
deformation.

Arbitrary Artinian defect modules, common complete-intersection algebras
not isomorphic to this formal square thickening and not covered by E836,
unstable torsion-free objects, and unrelated perfect complexes remain
outside the theorem. The new proof uses no additional external analytic
or algebraic classification theorem. Its finite supplement only checks
the elementary edge combinatorics and selected exact coefficient
identities; the all-size filtered-matrix argument is the proof in
section 2.

\section{SF-E838. A leading-quadratic criterion for common point
defects}\label{SF-E838:sf-e838.-a-leading-quadratic-criterion-for-common-point-defects}

Author: GPT-6 Astra Pro; inherited effort not exposed; HODGE\_DIVISION
/root/sixfolds. Date: 2026-11-21.

\textbf{Status:} complete new proof checked in full by the author, /root and
/root/imprimitive. Both reviewers also read the entire canonical
conversion; all twenty-one displays are unchanged. The actual companion equations and local-to-global
characteristic interface are the reviewed SF-E836 and SF-E819 inputs.
The new matrix lemma is proved for arbitrary commutative local Artinian algebras.
The application excludes a nonempty open family of quadratic complete
intersections, and every higher-order perturbation with those quadratic
jets. Its residual polynomial locus is a necessary condition only.

\subsection{1. An arbitrary-Artin square-matrix
lemma}\label{SF-E838:an-arbitrary-artin-square-matrix-lemma}

Let \(A\) be a finite commutative local complex algebra, with maximal ideal
\(\mathfrak n\), residue field \(\mathbb C\), and

\[
 V=\mathfrak n/\mathfrak n^2,\qquad W=\mathfrak n^2/\mathfrak n^3,
 \qquad e=\dim V,\quad h_2=\dim W.
 \tag{E838.1}
\]

\textbf{Lemma.} If \(h_2>e\), a square matrix
\(C\in\operatorname{Mat}_m(A)\) whose cokernel is annihilated by
\(\mathfrak n^2\) is invertible.

Split off unit pivots by actual invertible row and column operations
over the local ring. This preserves the cokernel up to isomorphism and
leaves an identity block together with a square block \(C'\) of size
\(k\), all of whose entries lie in \(\mathfrak n\). If \(k=0\), the
assertion holds. Suppose \(k>0\).

Choose complex vector-space splittings of the two successive filtration
quotients in \(A/\mathfrak n^3\). They are not algebra splittings. Write
\(C'=C_1+C_2\) modulo \(\mathfrak n^3\), with coefficients in the chosen
lifts of \(V\) and \(W\), respectively. Set

\[
 L:\mathbb C^k\longrightarrow V\otimes\mathbb C^k,\qquad L(v)=C_1v,
 \qquad a=\operatorname{rank}_{\mathbb C}L\le k.
 \tag{E838.2}
\]

An input has a constant part \(v_0\), a linear part \(v_1\), and a part
in \(\mathfrak n^2\), which does not contribute modulo
\(\mathfrak n^3\). Its image is \(C_1v_0+C_2v_0+C_1v_1\). This lies in
the pure second filtration quotient exactly when \(v_0\in\ker L\).
Consequently

\[
 \operatorname{im}(C'\bmod\mathfrak n^3)\cap(W\otimes\mathbb C^k)
   =C_2(\ker L)+V\cdot\operatorname{im}L.
 \tag{E838.3}
\]

The multiplication on the right is the actual induced map
\(V\otimes V\to W\). The two summands have dimensions at most \(k-a\)
and \(ea\). Since \(h_2>e\) entails \(e\ge1\), their sum has dimension
at most

\[
 k+(e-1)a\le ek<h_2k.
 \tag{E838.4}
\]

On the other hand, annihilation of the cokernel by \(\mathfrak n^2\)
puts every element of \(W\otimes\mathbb C^k\) in (E838.3), a
contradiction. Thus the residual block is absent. This proves the lemma
in every matrix size, without any grading, matrix commutativity, or
simultaneous diagonalization hypothesis.

\subsection{2. A general necessary companion
criterion}\label{SF-E838:a-general-necessary-companion-criterion}

Let \(R=\mathbb C[[z_1,z_2,z_3,z_4]]\), and let \(\pi_+,\pi_-\) be
constant rank-two bivectors with complementary supports. Let
\(I=(f_1,f_2,f_3,f_4)\) be an Artinian complete-intersection ideal, and
put \(A=R/I\). Consider an actual finite module fitting into

\[
 0\longrightarrow A^n\longrightarrow M\longrightarrow R^r
       \longrightarrow A^m\longrightarrow0,\qquad m+n>0.
 \tag{E838.5}
\]

Suppose that \(M\) has flat first-order module deformations for both
multiplication classes \(\pi_+,\pi_-\). The reviewed finite-cone trace
in E819 gives \(n=m>0\). The full actual lifting and opposite-order
calculation of E836, sections 2--3, gives matrices
\(B_{ij},C_+,C_-\in\operatorname{Mat}_m(A)\) satisfying

\[
 C_\pm B_{ij}=B_{ij}C_\pm=P^\pm_{ij}I_m,
 \qquad P^\pm_{ij}=\{f_i,f_j\}_{\pi_\pm}\bmod I.
 \tag{E838.6}
\]

Assume the following three algebraic conditions: all \(P^\pm_{ij}\)
belong to \(\mathfrak n^2\); \(h_2>e\); and some constant linear
combination \(P^0_{ij}=aP^+_{ij}+bP^-_{ij}\) spans
\(W=\mathfrak n^2/\mathfrak n^3\) as \(i<j\) varies. Let \(J\) be the
ideal generated by these six \(P^0_{ij}\). Spanning says

\[
 \mathfrak n^2=J+\mathfrak n^3.
 \tag{E838.7}
\]

In the finite \(A\)-module \(\mathfrak n^2/J\), this identity is exactly
\(\mathfrak n^2/J=\mathfrak n(\mathfrak n^2/J)\). Iterating it through
the nilpotent maximal ideal, or applying Nakayama's lemma, gives
\(J=\mathfrak n^2\). This uses the ideal generated by the brackets, not
only their complex linear span.

The actual matrix \(C_0=aC_++bC_-\) obeys \(C_0B_{ij}=P^0_{ij}I_m\).
Therefore \(J\), and hence \(\mathfrak n^2\), annihilates its cokernel.
The lemma proves that \(C_0\) is invertible. No extra deformation for
the pencil bivector is assumed; this is a linear combination of two
existing matrix identities.

Put \(D_\pm=C_\pm C_0^{-1}\). Solving for \(B_{ij}\) gives

\[
 D_\pm P^0_{ij}=P^\pm_{ij}I_m.
 \tag{E838.8}
\]

Let \(\overline P\) denote the class in \(W\), and let \(D_{\pm,0}\) be
the constant matrix modulo \(\mathfrak n\). Some
\(\overline P^0_{ij}\ne0\). Applying a complex linear functional on
\(W\) that is nonzero on that element to (E838.8) modulo
\(\mathfrak n^3\) forces \(D_{\pm,0}=d_\pm I_m\). All six equations then
give

\[
 (\overline P^+_{ij})_{i<j}=d_+(\overline P^0_{ij})_{i<j},
 \qquad
 (\overline P^-_{ij})_{i<j}=d_-(\overline P^0_{ij})_{i<j}.
 \tag{E838.9}
\]

Thus the two leading bracket tuples must be linearly dependent in
\(W^6\). In particular, the three conditions together with independence
of these two tuples exclude the two deformations. This conclusion
concerns an actual module over the original algebra; no module
deformation or Ext group has been base changed to a smaller quotient.

\subsection{3. Regular quadratic jets make the tuples
independent}\label{SF-E838:regular-quadratic-jets-make-the-tuples-independent}

Let \(S=\mathbb C[z_1,z_2,z_3,z_4]\). Suppose that
\(F_1,\ldots,F_4\in S_2\) form a homogeneous regular sequence with
Artinian quotient

\[
 A_0=S/(F_1,F_2,F_3,F_4),\qquad
 Q=\operatorname{span}_{\mathbb C}(F_1,F_2,F_3,F_4).
 \tag{E838.10}
\]

The four elementary exact sequences obtained by successively multiplying
by a nonzerodivisor of degree two give

\[
 \operatorname{Hilb}_{A_0}(t)
       =\frac{(1-t^2)^4}{(1-t)^4}=(1+t)^4.
 \tag{E838.11}
\]

In particular \(\dim(A_0)_1=4\), \(\dim(A_0)_2=6\), and the length is
sixteen.

Take arbitrary formal perturbations \(f_i=F_i+O(\mathfrak m^3)\), where
\(\mathfrak m=(z_1,z_2,z_3,z_4)\). Their quotient is again an Artinian
complete intersection, and its associated graded algebra is \(A_0\).
Here is the filtered justification. If the initial form of an element
\(f\) in a separated filtered quotient is a nonzerodivisor, then for
every nonzero \(h\) the initial form of \(fh\) is the nonzero product of
the two initial forms. Hence the initial ideal of the principal ideal
\((f)\) is generated by that initial form, and the associated graded
quotient is its graded quotient. Apply this successively to the four
regular initial forms \(F_i\). All intermediate quotients are separated
Noetherian complete local rings. The resulting graded quotient is
(E838.11); alternatively its vanishing in degree five gives
\(\mathfrak n^5=\mathfrak n^6\), and Nakayama gives \(\mathfrak n^5=0\).
The same initial-form nonzerodivisor argument proves regularity of the
\(f_i\).

Thus for the actual perturbed algebra, \(e=4\), \(h_2=6>4\), and

\[
 W=\mathfrak n^2/\mathfrak n^3=S_2/Q,\qquad
 \overline P^\pm_{ij}=\{F_i,F_j\}_{\pi_\pm}\bmod Q.
 \tag{E838.12}
\]

The leading bracket tuples in (E838.12) are independent. To prove this,
work in the separate homogeneous complete intersection \(A_0\), not with
an unjustified conclusion about the perturbed Jacobian. Let \(J_F\) be
the actual Jacobian matrix of the four quadrics, and write

\[
 \mathcal P_\pm=J_F\pi_\pm J_F^{\mathsf t}
       \in\Lambda^2\mathbb C^4\otimes(A_0)_2.
 \tag{E838.13}
\]

Functoriality of exterior powers gives
\(\mathcal P_\pm\wedge\mathcal P_\pm=0\), since each original bivector
has rank two, and

\[
 \mathcal P_+\wedge\mathcal P_-
       =\det(J_F)\,(\pi_+\wedge\pi_-)\ne0
       \quad\text{in }\Lambda^4\mathbb C^4\otimes(A_0)_4.
 \tag{E838.14}
\]

The nonvanishing is precisely the published complete-intersection
Jacobian-socle theorem: Eisenbud--Ulrich, \emph{J. reine angew. Math.}
756 (2019), Theorem 7.1, p.~210, with Appendix A.1, A.3 and A.5,
pp.~222--225.\cite{SFEisenbudUlrich2019} Apply it to the completion of the homogeneous Artinian
algebra; the finite quotient itself is unchanged. The characteristic is
zero, the four generators are a regular sequence, and the Jacobian
generates its nonzero socle. This exact published statement and the
invoked Appendix proofs were fully read for E836.

If the two six-tuples were dependent over \(\mathbb C\), one would be a
scalar multiple of the other, or one would vanish. Either possibility,
together with their individual exterior squares being zero, would make
(E838.14) zero. Thus they are independent. Section 2 now proves the
following theorem.

\textbf{Quadratic-jet theorem.} If some pencil of leading brackets spans
\(S_2/Q\), every formal perturbation with these regular quadratic jets
is excluded from the two-deformation property of (E838.5), in every
positive common multiplicity. Consequently no nontrivial point
modification of this type can occur in the balanced weak-semiregular
coherent-sheaf setting of E819, with its specified coherent finite
translation/tensor action and arbitrary rational algebraic \(B\) shear.

\subsection{4. An exact seven-coefficient residual
test}\label{SF-E838:an-exact-seven-coefficient-residual-test}

Choose any ordered monomial basis of the ten-dimensional space \(S_2\).
For the ordered four quadrics, form the ten-by-ten coefficient matrix
whose first four columns are the \(F_i\) and whose last six columns are
their brackets for \(a\pi_++b\pi_-\), in the order
\(12,13,14,23,24,34\). Define

\[
 \mathcal D_F(a,b)
   =\det[\,F_1\ F_2\ F_3\ F_4\mid
       (\{F_i,F_j\}_{a\pi_++b\pi_-})_{i<j}\,]
   =\sum_{j=0}^{6}c_j(F)a^{6-j}b^j.
 \tag{E838.15}
\]

The pencil spans \(S_2/Q\) for some \((a,b)\) exactly when this
polynomial is not identically zero. Each \(c_j\) is a homogeneous
polynomial of degree sixteen in the forty coefficients of the ordered
four quadrics: the first four columns have degree one, and the last six
have degree two. If the chosen ordered basis of \(Q\) is changed by
\(T\in\mathrm{GL}_4(\mathbb C)\), the first columns transform by \(T\)
and the bracket columns by \(\Lambda^2T\), so

\[
 \mathcal D_{TF}(a,b)=(\det T)^4\mathcal D_F(a,b),
 \tag{E838.16}
\]

using \(\det(\Lambda^2T)=(\det T)^3\). The latter identity follows for
diagonal matrices by counting the three occurrences of each eigenvalue,
and then for every matrix by polynomial identity. A change of the
ambient monomial basis contributes a common nonzero scalar as well.

Therefore the condition \(c_0=\cdots=c_6=0\) is basis independent on the
Grassmannian \(\operatorname{Gr}(4,S_2)\). On each usual affine
Grassmannian chart the seven coefficients are regular polynomial
functions, of degree at most sixteen in its twenty-four coordinates.
This is a closed residual locus relative to the regular-quadratic
complete-intersection locus. The theorem proves only necessity of lying
in it. It does not say the seven equations are independent, that this
locus has codimension seven, or that any point in it admits a companion
matrix or a module deformation.

\subsection{5. Exact nonemptiness and strict new
scope}\label{SF-E838:exact-nonemptiness-and-strict-new-scope}

Normalize \(\pi_+=\partial_{z_1}\wedge\partial_{z_2}\) and
\(\pi_-=\partial_{z_3}\wedge\partial_{z_4}\). Let \(\ell=Jz\), with

\[
 J=\begin{pmatrix}
 1&0&1&0\\
 0&1&0&1\\
 -3/5&1&-9/5&2\\
 -8/5&1&-16/5&3
 \end{pmatrix},\qquad
 \det J=-4/5,\qquad F_i=\ell_i^2.
 \tag{E838.17}
\]

Use the monomial basis
\((z_1^2,z_2^2,z_3^2,z_4^2,z_1z_2,z_1z_3,z_1z_4,z_2z_3,z_2z_4,z_3z_4)\).
The two alternating coefficient forms in the \(\ell\)-coordinates have
pair coefficients

\[
 p=(1,1,1,3/5,8/5,1),\qquad
 q=(1,2,3,9/5,16/5,1).
 \tag{E838.18}
\]

In the corresponding \(\ell\)-monomial basis, the first four columns are
the four square basis vectors, and the six bracket columns are diagonal
with entries \(4(ap_{ij}+bq_{ij})\). Their determinant is
\(4^6(24/25)(a+b)^2(a+2b)^2(a+3b)^2\). Passing to the \(z\)-monomial
basis multiplies this by
\(\det(\operatorname{Sym}^2J)=(\det J)^5=(-4/5)^5\). The exponent five
is obtained on diagonal matrices by counting each variable twice in its
square and once in each of three pair monomials; polynomial identity
gives the general case. Thus

\[
 \mathcal D_F(a,b)
   =-\frac{100663296}{78125}
       (a+b)^2(a+2b)^2(a+3b)^2.
 \tag{E838.19}
\]

This is an exact rank certificate, also reconstructed directly by the
adjacent symbolic checker. In particular both pure endpoints
\(\mathcal D_F(1,0)\) and \(\mathcal D_F(0,1)\) are nonzero.

There is an equally explicit nonsquare witness. In these same mixed
linear parameters put

\[
 F^\sharp_1=\ell_1^2+\ell_2\ell_3,\qquad
 F^\sharp_2=\ell_2^2,\quad F^\sharp_3=\ell_3^2,\quad
 F^\sharp_4=\ell_4^2.
 \tag{E838.20}
\]

Its quotient is a free module with basis \(1,\ell_1\) over
\(\mathbb C[\ell_2,\ell_3,\ell_4]/(\ell_2^2,\ell_3^2,\ell_4^2)\),
because the remaining relation is monic of degree two in \(\ell_1\).
Thus its length is sixteen, and it is a quadratic complete intersection.
In its quadratic quotient use the six pair monomials as basis. For any
constant alternating coefficients \(P_{ij}\) in the
\(\ell\)-coordinates, its first three brackets reduce to

\[
 \{F^\sharp_1,F^\sharp_2\}=4P_{12}\ell_1\ell_2,\qquad
 \{F^\sharp_1,F^\sharp_3\}=4P_{13}\ell_1\ell_3,\qquad
 \{F^\sharp_1,F^\sharp_4\}
   =4P_{14}\ell_1\ell_4
      +2P_{24}\ell_3\ell_4+2P_{34}\ell_2\ell_4.
 \tag{E838.21}
\]

The last three brackets are \(4P_{ij}\ell_i\ell_j\), for
\(2\le i<j\le4\). The resulting six-by-six matrix is triangular, with
exactly the same six diagonal entries as for the four squares.
Eliminating the first four columns of the ten-by-ten matrix contributes
determinant one. Therefore \(\mathcal D_{F^\sharp}\) is exactly
(E838.19).

The quadratic relation space of (E838.20) has only three projective
rank-one quadrics. Indeed, a linear combination with nonzero coefficient
\(c\) at \(F^\sharp_1\) has an isolated nonzero \(\ell_1^2\) entry and a
nonzero \(c\ell_2\ell_3\) entry. If it were the square of a linear form,
the nonzero \(\ell_1\) coefficient and the vanishing of all
\(\ell_1\ell_j\) coefficients would force the other three coefficients
of that linear form to vanish, contradicting the \(\ell_2\ell_3\) term.
With \(c=0\), the diagonal quadratic has rank one only when exactly one
of its three coefficients is nonzero. In contrast a coordinate-square
space contains four independent rank-one quadrics. Hence (E838.20) is
not in that orbit, and no formal parameter change or local-algebra
isomorphism can turn its algebra into the square algebra: either map
preserves the cotangent space and its quadratic relation space. This is
an exact nonsquare example, independent of the dimension argument that
follows.

For completeness, the regular-quadratic locus is open. It is exactly the
locus on which the degree-five multiplication map
\(S_3^{\oplus4}\to S_5\), \((g_i)\mapsto\sum F_i g_i\), is surjective.
For a regular sequence this follows from (E838.11). Conversely
surjectivity puts every monomial of degree five in the ideal, so its
quotient is Artinian. A four-element parameter ideal in the
four-dimensional regular local ring is a regular sequence; the precise
standard source is Stacks Lemmas 10.106.3, tag 00NQ, and 10.104.2, tag
02JN, with the full proof in Proposition 10.103.4, tag 00N6\cite{SFStacksQuadraticJets}. Homogeneity
transfers regularity back to the graded polynomial algebra: a nonzero
homogeneous kernel element would survive localization at the homogeneous
maximal ideal. Surjectivity is the nonvanishing of a maximal matrix
minor, hence is open. The witness (E838.17) lies in this locus.

The Grassmannian is irreducible of dimension \(4(10-4)=24\). Its
standard matrix charts are affine spaces with twenty-four coordinates;
the space of full-rank four-by-ten frames is an irreducible open subset
of affine forty-space and maps surjectively to the Grassmannian, which
proves irreducibility. These explicit chart constructions are also
Stacks Section 27.22, tag 089R, Lemma 27.22.1\cite{SFStacksQuadraticJets}. Consequently the simultaneous conditions
of regularity and nonvanishing of both pure endpoints define a nonempty
open subset \(\mathcal U\) of dimension twenty-four.

The orbit of \(\operatorname{span}(z_1^2,z_2^2,z_3^2,z_4^2)\) under
linear changes of the four variables has closure of dimension at most
fifteen. One does not need an orbit-stabilizer classification: scalar
matrices act trivially on this subspace, so on each of the sixteen
charts in which a chosen matrix entry is nonzero, normalize that entry
to one. The image is parametrized by at most fifteen algebraic
coordinates. Its closure has dimension at most fifteen, since the
coordinate functions of the image generate a field of transcendence
degree at most fifteen. The finite union of these image closures has the
same bound. It cannot contain the twenty-four-dimensional open set
\(\mathcal U\).

It follows that \(\mathcal U\) contains quadratic spaces outside the
closure of the coordinate-square orbit. For such a space, neither its
homogeneous ideal nor any higher-order perturbation can be a formal
coordinate-square ideal: the linear terms of four formal parameters
would put its four-dimensional quadratic relation space in that orbit.
The same observation rules out an abstract local-algebra isomorphism to
the square algebra, since such an isomorphism preserves the cotangent
space and its quadratic relations. Thus the theorem has strict scope
beyond E837.

It also has strict scope beyond the E836 commuting-triple criterion for
either original bivector. At either endpoint, nonvanishing of the
determinant means that the bracket map \(\Lambda^2Q\to S_2/Q\) is an
isomorphism. Any other minimal regular generator set has leading
quadrics related by an invertible constant matrix. Three independent
generators have a three-dimensional exterior-square space, which cannot
be killed by this isomorphism. Hence no such presentation has its last
three generators pairwise commuting modulo the ideal for either original
bivector.

\subsection{6. Global application, actual excluded objects, and
boundaries}\label{SF-E838:global-application-actual-excluded-objects-and-boundaries}

The global implication is exactly E819's reviewed one: invariant
evaluation on each balanced cap-kernel class vanishes under weak
semiregularity; forgetting the finite action and restricting to the
completed local ring gives the two ordinary module deformations. An
algebraic rational \(B\) shear preserves their bivector components. No
equivariance of the individual quadratic generators or module bases is
required.

The excluded family contains actual balanced coherent sheaves, before
imposing semiregularity. Choose any homogeneous quadratic complete
intersection in the nonempty open subset outside the square-orbit
closure, form its length-sixteen subscheme \(Z\) at a point of the
abelian fourfold, and take the rank-thirty-two bundle \(H\) of E815/E820
with entire character \(32v_{1/2}\). A finite formal quotient is an
actual algebraic subscheme because the local ring and its completion
have the same quotients by powers of the maximal ideal. For
\(1\le m\le32\), choose \(H\twoheadrightarrow\mathcal O_Z^m\), write
\(K\) for its kernel, and take any extension
\(0\to\mathcal O_Z^m\to E\to K\to0\). The four-generator Koszul
resolution gives
\(\dim\operatorname{Ext}^1(K,\mathcal O_Z^m)=6m^2\operatorname{length}Z=96m^2\).
Its torsion, hull and hull quotient are the displayed ones, and its full
character is \(32v_{1/2}\). It is ordinarily gluable as a coherent
sheaf. The theorem excludes its weak balanced semiregularity; no finite
action, simplicity, or successful seed is asserted.

Common complete intersections whose regular quadratic jets lie in the
seven-coefficient residual locus are not decided by this criterion,
although E836 or E837 can decide some of them. Ideals with higher
initial degrees, unequal Artinian defect algebras, arbitrary Artinian
defect modules, unstable torsion-free objects, and unrelated perfect
complexes remain outside the new assertion. In particular no passage to
a square quotient of a weighted power algebra is used. The exact
supplement checks the determinant witness and finite algebra identities;
it is not a finite verification of the all-size matrix lemma or of any
geometric module deformation.

\section{SF-E839. Five leading brackets already obstruct a common point
defect}\label{SF-E839:sf-e839.-five-leading-brackets-already-obstruct-a-common-point-defect}

Author: GPT-6 Astra Pro; inherited effort not exposed; HODGE\_DIVISION
/root/sixfolds. Date: 2026-11-21.

\textbf{Status:} complete new refinement checked in full by the author, /root
and /root/g34\_eightfolds. The root also read the entire canonical
conversion and both exact certificate reports. E838 is retained unchanged. The new necessary
condition is that every leading bracket pencil have rank at most four.
An exact nonsquare rank-five example, with identically zero determinant
pencil, proves strict new scope relative to E838 as well as E836 and
E837.

\subsection{1. The stronger filtered-matrix
lemma}\label{SF-E839:the-stronger-filtered-matrix-lemma}

Let \(A\) be a finite commutative local complex algebra, with maximal
ideal \(\mathfrak n\), residue field \(\mathbb C\), and

\[
 V=\mathfrak n/\mathfrak n^2,\qquad W=\mathfrak n^2/\mathfrak n^3,
 \qquad e=\dim V.
 \tag{E839.1}
\]

Suppose that a square matrix \(C\in\operatorname{Mat}_m(A)\) has
cokernel annihilated by elements \(p_1,\ldots,p_s\in\mathfrak n^2\)
whose classes in \(W\) have span of dimension \(\rho>e\). Then \(C\) is
invertible.

Here is the complete dimension argument, so the stronger hypothesis is
not hidden in a reference to E838. Split all unit pivots by invertible
row and column operations. If a nonempty residual square block \(C'\) of
size \(k\) remains, all its coefficients lie in \(\mathfrak n\). Choose
complex vector-space splittings modulo \(\mathfrak n^3\), and write
\(C'=C_1+C_2\), with coefficients in the respective lifts of \(V,W\).
Let

\[
 L:\mathbb C^k\longrightarrow V\otimes\mathbb C^k,\qquad
 L(v)=C_1v,\qquad a=\operatorname{rank}L\le k.
 \tag{E839.2}
\]

A constant input has no first-order image exactly when it belongs to
\(\ker L\). Linear inputs contribute through \(C_1\), and quadratic
inputs contribute nothing modulo \(\mathfrak n^3\). Therefore the pure
second-order part of the image is exactly

\[
 \operatorname{im}(C'\bmod\mathfrak n^3)\cap(W\otimes\mathbb C^k)
   =C_2(\ker L)+V\cdot\operatorname{im}L.
 \tag{E839.3}
\]

The hypothesis forces \(e\ge1\): if \(e=0\), nilpotence gives
\(\mathfrak n=\mathfrak n^2=0\) and hence \(W=0\), contrary to
\(\rho>e\). Its dimension is at most \((k-a)+ea\le ek\). Each of the \(\rho\)
independent leading annihilators times every constant target vector lies
in that image, providing a subspace of dimension \(\rho k>ek\). This is
a contradiction. Thus no residual block remains. Only the actual linear
span of these leading annihilators was needed; their ideal need not be
all of \(\mathfrak n^2\).

The strict inequality cannot simply be replaced by equality in this
abstract lemma. In the four-variable square algebra, multiplication by
\(l=t_1+t_2+t_3+t_4\) is not invertible, and its cokernel is annihilated
by the four quadrics \(lt_i\). These are independent: a relation has
coefficients \(a_i+a_j=0\) at every \(t_it_j\); using three distinct
indices gives \(a_i=0\) in characteristic zero. This is only a boundary
for the matrix lemma, not a simultaneous module-deformation example.

\subsection{2. The rank-at-most-four necessary
condition}\label{SF-E839:the-rank-at-most-four-necessary-condition}

Keep the formal geometric setting of E838. Namely
\(R=\mathbb C[[z_1,z_2,z_3,z_4]]\), the two constant rank-two bivectors
\(\pi_+,\pi_-\) have complementary supports, and the leading quadrics
\(F_1,\ldots,F_4\) of the regular generators
\(f_i=F_i+O(\mathfrak m^3)\) form an Artinian homogeneous regular
sequence. Put

\[
 Q=\operatorname{span}(F_1,F_2,F_3,F_4),\qquad
 A=R/(f_1,f_2,f_3,f_4),\qquad W=S_2/Q.
 \tag{E839.4}
\]

The complete filtered initial-form proof of E838 gives
\(\operatorname{gr}A=S/(F)\), \(e=4\), \(\dim W=6\), and length sixteen.
It also proves that the two leading bracket tuples
\((\{F_i,F_j\}_{\pi_\pm}\bmod Q)_{i<j}\) are independent, using the
nonzero Jacobian socle in this homogeneous complete intersection.

Consider an actual common point modification

\[
 0\longrightarrow A^n\longrightarrow M\longrightarrow R^r
       \longrightarrow A^m\longrightarrow0,\qquad m+n>0.
 \tag{E839.5}
\]

If \(M\) has flat first-order deformations for both multiplication
classes, E819 gives \(n=m>0\), and E836 gives the actual common
companion matrices

\[
 C_\pm B_{ij}=B_{ij}C_\pm=P^\pm_{ij}I_m,\qquad
 P^\pm_{ij}=\{f_i,f_j\}_{\pi_\pm}\bmod I.
 \tag{E839.6}
\]

For any complex \(a,b\), form \(C_0=aC_++bC_-\). Its cokernel is
annihilated by \(P^0_{ij}=aP^+_{ij}+bP^-_{ij}\). If their classes in
\(W\) span a space of dimension at least five, the lemma makes \(C_0\)
invertible. Then

\[
 (C_\pm C_0^{-1})P^0_{ij}=P^\pm_{ij}I_m.
 \tag{E839.7}
\]

Taking the second filtration quotient and applying a linear functional
nonzero on one nonzero \(\overline P^0_{ij}\) forces the constant part
of each \(C_\pm C_0^{-1}\) to be scalar. All six equations therefore
make both leading tuples scalar multiples of the same one, contrary to
their independence. This proves:

\textbf{Rank-five obstruction.} For every such common point modification
admitting both deformations, the linear map

\[
 T_{a,b}:\Lambda^2Q\longrightarrow S_2/Q,\qquad
 F_i\wedge F_j\longmapsto
       \{F_i,F_j\}_{a\pi_++b\pi_-}\bmod Q
 \tag{E839.8}
\]

has rank at most four for every \(a,b\in\mathbb C\). Hence a single
rank-five or rank-six pencil value excludes weak balanced semiregularity
in the actual coherent-sheaf setting of E819, with its specified
coherent finite translation/tensor action and rational algebraic \(B\)
shear.

There is no extra deformation for a pencil parameter and no base change
of an Ext group or module deformation to a square quotient. Only the two
necessary identities for the original module were combined.

\subsection{3. The exact residual rank
equations}\label{SF-E839:the-exact-residual-rank-equations}

On a standard chart of \(\operatorname{Gr}(4,S_2)\), choose the ordered
generator matrix with a fixed four-by-four identity pivot block. After
reducing the six bracket columns by those four generators, \(T_{a,b}\)
is an explicit six-by-six matrix with entries linear in \(a,b\). The
rank condition is equivalent to vanishing of all its thirty-six
five-by-five minors, identically in the pencil variables. Each such
minor has degree five in \(a,b\), so its six coefficients give at most

\[
 36\cdot6=216
 \tag{E839.9}
\]

regular polynomial equations on each chart.

The chart polynomials have degree at most fourteen in the quadratic
coefficients. Indeed each five-by-five quotient minor is a nine-by-nine
augmented minor: retain all four pivot rows, the chosen five other rows,
all four generator columns and the chosen five bracket columns. Its
determinant is unchanged by eliminating the pivot columns. The four
generator columns have coefficient degree one, and the five bracket
columns have coefficient degree two, giving total degree fourteen before
normalization; substituting the identity pivot block does not increase
that degree.

These rank conditions are independent of the chosen generator basis and
quotient basis because rank is invariant under invertible changes of
both. They therefore define an intrinsic closed residual locus within
the regular-quadratic locus. No independence of the 216 equations,
codimension count, sufficiency, or existence of a module on the residual
locus is asserted.

\subsection{4. An exact rank-five nonsquare
example}\label{SF-E839:an-exact-rank-five-nonsquare-example}

Take the two standard bivectors in the \(z\)-coordinates and put
\(\ell=Jz\), where

\[
 J=\begin{pmatrix}
 0&1&0&1\\
 1&1&3&2\\
 1&0&1&0\\
 1&0&2&0
 \end{pmatrix},\qquad \det J=-1.
 \tag{E839.10}
\]

Use the actual quadrics

\[
 F_1=\ell_1^2+\ell_2\ell_3,\qquad
 F_2=\ell_2^2,\qquad F_3=\ell_3^2,\qquad F_4=\ell_4^2.
 \tag{E839.11}
\]

The quotient is free with basis \(1,\ell_1\) over the length-eight
square algebra in \(\ell_2,\ell_3,\ell_4\); hence it has length sixteen
and is an Artinian quadratic complete intersection. The transformed
bivector coefficients, in pair order \(12,13,14,23,24,34\), are

\[
 p=(-1,-1,-1,-1,-1,0),\qquad
 q=(-3,-1,-2,-2,-4,0).
 \tag{E839.12}
\]

Both are decomposable and their mixed wedge equals \(\det J=-1\), as
follows directly by applying \(J\) to the two standard decomposable
forms.

Use the six pair monomials as the basis of \(S_2/Q\). The complete
unrestricted-coefficient calculation (E838.21) applies to the same four
relation shapes. With \(P_{ij}=ap_{ij}+bq_{ij}\), its six-by-six bracket
matrix is

\[
 \begin{pmatrix}
 4P_{12}&0&0&0&0&0\\
 0&4P_{13}&0&0&0&0\\
 0&0&4P_{14}&0&0&0\\
 0&0&0&4P_{23}&0&0\\
 0&0&2P_{34}&0&4P_{24}&0\\
 0&0&2P_{24}&0&0&4P_{34}
 \end{pmatrix}.
 \tag{E839.13}
\]

Here \(P_{34}=0\) identically, so the last column vanishes and every
pencil value has rank at most five. Deleting its last row and last
column gives the exact minor

\[
 -1024(a+3b)(a+b)(a+2b)^2(a+4b).
 \tag{E839.14}
\]

It is nonzero at both pure endpoints. Thus both maps \(T_{1,0}\),
\(T_{0,1}\) have rank exactly five, and the generic pencil rank is five.
In particular the entire determinant pencil of E838 is identically zero,
while the new criterion excludes this example and every higher-order
perturbation having these quadratic jets.

The example lies outside E837. Its relation space has only three
projective rank-one quadrics: a nonzero coefficient at \(F_1\) produces
an isolated \(\ell_1^2\) term and a nonzero \(\ell_2\ell_3\) term, which
cannot be a square of a linear form because all three \(\ell_1\ell_j\)
coefficients vanish. With that coefficient zero, only the three
coordinate-square lines are rank one. A formal square relation space has
four independent rank-one quadrics. Formal parameter changes and
abstract local-algebra isomorphisms preserve this cotangent-space
relation invariant.

It also lies outside E836's commuting-triple criterion for every minimal
regular generator presentation, for either original bivector. Any such
presentation changes the leading quadrics by an invertible constant
matrix. The kernel of either six-dimensional bracket map has dimension
one. Three independent leading generators have a three-dimensional
exterior-square subspace; it cannot be contained in a one-dimensional
kernel. Their three brackets therefore cannot all vanish modulo the
ideal. The same argument applies to every higher-order perturbation with
these jets.

\subsection{5. Actual geometric scope and source
boundary}\label{SF-E839:actual-geometric-scope-and-source-boundary}

The global passage remains the precise reviewed E819 one: invariant
vanishing for the specified coherent finite action implies ordinary
vanishing after forgetting the action, and completion supplies the two
local deformations. An algebraic rational \(B\) shear preserves their
bivector components. No chosen local generator basis needs an
equivariant structure.

There are actual balanced coherent objects with the new excluded local
algebra. Realize its length-sixteen finite quotient at a point of the
abelian fourfold, take the rank-thirty-two bundle \(H\) with full
character \(32v_{1/2}\), choose a quotient
\(H\twoheadrightarrow\mathcal O_Z^m\) for \(1\le m\le32\), and let \(K\)
be its kernel. Any extension \(0\to\mathcal O_Z^m\to E\to K\to0\) has
precisely this torsion, locally free hull \(H\), and identical hull
quotient. Its full character is \(32v_{1/2}\), and it is ordinarily
gluable as a coherent sheaf. The Koszul calculation gives
extension-space dimension \(96m^2\). The theorem excludes weak
semiregularity; it asserts no action, simplicity, or successful seed.

No new outside theorem is used. E838's complete associated-graded and
Jacobian argument, with the precisely located published Eisenbud--Ulrich
Theorem 7.1 and Appendix proofs, is retained as a checked input. The new
all-size matrix proof, actual rank-five witness, residual-minor count
and strict-scope statements are proved above. The finite supplement
checks exact matrices and minors only. Rank-at-most-four residual jets,
higher initial degrees, other Artinian defect modules and unrelated
perfect complexes remain open in this result.

\section{SF-E840: a principal-image condition at the rank-four
boundary}\label{SF-E840:sf-e840-a-principal-image-condition-at-the-rank-four-boundary}

Status: complete new author derivation; whole root and independent
/root/g34\_eightfolds proof, source and code audits passed. This
strengthens the residual necessary condition of E839. The example below
survives the scalar cross-minor equations and every rank-five pencil
minor. No semiregular object is constructed.

\subsection{1. Quadratic multiplication on a regular four-variable
intersection}\label{SF-E840:quadratic-multiplication-on-a-regular-four-variable-intersection}

Let \(V\) be a four-dimensional complex vector space and let \(Q\) be a
four-dimensional subspace of \(\operatorname{Sym}^2(V)\) generated by a
regular sequence of quadrics. Put

\[
 W=\operatorname{Sym}^2(V)/Q,
 \qquad \dim W=6,
 \qquad \mu:V\otimes V\longrightarrow W.
 \tag{E840.1}
\]

For every subspace \(L\) of \(V\) of dimension at least two,

\[
                     \dim\mu(VL)\ge5.
 \tag{E840.2}
\]

It suffices to take \(\dim L=2\). The degree-two ideal \(VL\) has
dimension seven: its quotient is \(\operatorname{Sym}^2(V/L)\), of
dimension three. If its image in \(W\) had dimension at most four, then
\(\dim(Q\cap VL)\) would be at least three. Thus the four quadrics in
\(Q\) restrict to a space of dimension at most one on the
two-dimensional quotient \(V/L\). Such a collection has a common
projective zero over \(\mathbf C\). Indeed, either every restriction is
zero or they are multiples of one homogeneous binary quadratic, which
has a zero in \(\mathbf P^1\). Lifting the corresponding quotient point
gives a nonzero common zero of \(Q\) in \(V^*\), contradicting the
regular-sequence, finite-length hypothesis. This proves (E840.2).

For a four-dimensional subspace \(J\) of \(W\) define the intrinsic
vector space

\[
 K(J)=\{\ell\in V:\mu(\ell V)\subseteq J\}
      =\ker\bigl(V\longrightarrow\operatorname{Hom}(V,W/J)\bigr).
 \tag{E840.3}
\]

Equation (E840.2) shows that \(\dim K(J)\) is at most one.

\subsection{2. The equality case of the square-matrix
bound}\label{SF-E840:the-equality-case-of-the-square-matrix-bound}

Let \(A\) be a commutative local complex Artinian algebra with maximal
ideal \(\mathfrak n\). Assume its multiplication through degree two
identifies \(\mathfrak n/\mathfrak n^2\) with \(V\) and
\(\mathfrak n^2/\mathfrak n^3\) with \(W\) as in (E840.1). Let \(C\) be
a square matrix over \(A\), and suppose elements of \(\mathfrak n^2\)
annihilate \(\operatorname{coker}C\) and have leading classes spanning a
four-dimensional \(J\) in \(W\).

If \(C\) is not invertible, then

\[
                   J=\mu(\ell V)
       \quad\text{for a unique line }\mathbf C\ell\subset V,
       \qquad \dim K(J)=1.
 \tag{E840.4}
\]

In particular the multiplication map \(V\) to \(W\) given by \(\ell\) is
injective.

Here is the complete equality argument. Split unit pivots by invertible
row and column operations. The cokernel is unchanged, and a nonempty
residual square block \(C'\) of size \(k\) has entries in
\(\mathfrak n\). Choose vector-space splittings modulo \(\mathfrak n^3\)
and write \(C'=C_1+C_2\) there. Let

\[
 L_1:\mathbf C^k\longrightarrow V\otimes\mathbf C^k,
 \qquad U=\operatorname{im}L_1,\qquad a=\dim U.
 \tag{E840.5}
\]

The pure degree-two part of the image of \(C'\) modulo \(\mathfrak n^3\)
is exactly

\[
             C_2(\ker L_1)+\mu(VU).
 \tag{E840.6}
\]

To verify this, write an input as a constant vector plus a vector of
linear terms plus quadratic terms. The linear output is \(C_1\) applied
to the constant vector, so it vanishes precisely when that constant lies
in \(\ker L_1\). The quadratic output is then \(C_2\) on that constant
plus multiplication of \(C_1\) with the linear input. Terms of degree at
least two in the input contribute nothing modulo \(\mathfrak n^3\). This
uses a vector-space splitting only, not a graded or algebra splitting of
\(A\).

The annihilator hypothesis puts \(J\otimes\mathbf C^k\) inside (E840.6),
because every generating annihilator times each constant target vector
is an actual column image. Consequently

\[
 4k\le\dim\bigl(C_2(\ker L_1)+\mu(VU)\bigr)
       \le(k-a)+4a\le4k.
 \tag{E840.7}
\]

Every inequality is equality. Thus \(a=k\), \(\ker L_1=0\), and

\[
                     \mu(VU)=J\otimes\mathbf C^k.
 \tag{E840.8}
\]

Let \(L\) be the span in \(V\) of all contractions of \(U\) by linear
functionals on \(\mathbf C^k\). Contracting (E840.8) shows that
\(\mu(VL)\) is contained in \(J\). The space \(L\) is nonzero because
\(U\) has dimension \(k>0\). By (E840.2), \(L\) cannot have dimension at
least two. Write \(L=\mathbf C\ell\). Then \(U\) is contained in
\(\ell\otimes\mathbf C^k\), and both spaces have dimension \(k\), so
they are equal. Equation (E840.8) gives \(J=\mu(\ell V)\). Its dimension
is four, proving injectivity. Equation (E840.2) also proves uniqueness
of the line and \(\dim K(J)=1\).

All higher-order entries of \(C\) are allowed. No quotient of a
deformation or base change of an Ext group has been made.

\subsection{3. Necessary condition for actual balanced point
modifications}\label{SF-E840:necessary-condition-for-actual-balanced-point-modifications}

Use the precise geometric setup of E819 and E836: a nontrivial finite
modification on a genuine RM fourfold with locally free hull, whose
local torsion and hull quotient are respectively \(A^m\) and \(A^n\) for
the same complete-intersection Artin algebra. Weak balanced middle
semiregularity, with the allowed coherent finite translation/tensor
action and rational algebraic B-field, gives the ordinary vanishings for
both RM bivectors after the already proved invariant-to-ordinary and
B-shear steps. The finite-perfect cone and supertrace in E819 force
\(m=n>0\).

Suppose the ideal has regular quadratic initial forms \(Q\); arbitrary
higher terms are permitted. E838 identifies the associated graded
algebra with the homogeneous regular intersection. E836 supplies common
extension matrices \(B_{ij}\) and square companion matrices \(C_+,C_-\)
satisfying both orders of the actual equations

\[
 C_\pm B_{ij}=B_{ij}C_\pm=P_{\pm,ij}I_m.
 \tag{E840.9}
\]

Here \(P\) are the brackets of the actual ideal generators. Their
degree-two initial tuples \(P_+^0\) and \(P_-^0\) are linearly
independent in \(W^6\). For completeness, the homogeneous Jacobian
calculation gives their two self-wedges zero, while their mixed wedge is
\(\det(\operatorname{Jac}Q)\) times the nonzero transverse constant
bivector wedge. The Jacobian generates the nonzero socle of the
homogeneous characteristic-zero complete intersection. Thus the tuples
cannot be proportional. This is exactly the published Eisenbud--Ulrich
socle input and full proof scope already read in E836/E838, not a new
claim about arbitrary nonhomogeneous initial forms.

For a nonzero pair \((a,b)\), form the linear combinations
\(C_0=aC_++bC_-\) and \(P_0=aP_++bP_-\). The two already existing
companion equations give

\[
                         C_0B_{ij}=P_{0,ij}I_m.
 \tag{E840.10}
\]

This does not assume a separate deformation in a third pencil direction.

If the degree-two tuple \(P_0\) has rank four, let \(J\) be its span in
\(W\). The matrix \(C_0\) cannot be a unit. Otherwise
\(B_{ij}=P_{0,ij}C_0^{-1}\), and the other companion equations give
\(D_\pm P_{0,ij}=P_{\pm,ij}I\) with \(D_\pm=C_\pm C_0^{-1}\). Passing to
degree two uses only the constant matrices \(D_\pm(0)\). A functional
nonzero on some \(P_{0,ij}\) forces each constant matrix to be scalar,
and then all six equations make both complete initial tuples
proportional to \(P_0\), contrary to their independence.

Equation (E840.10) annihilates \(\operatorname{coker}C_0\) by the six
\(P_{0,ij}\). Applying the equality lemma therefore proves the new
necessary condition:

\[
 \boxed{\text{For every rank-four initial pencil direction,}
        \quad J=\ell V\text{ in }W,
        \quad\dim K(J)=1.}
 \tag{E840.11}
\]

If \(K(J)=0\) in even one rank-four direction, weak balanced middle
semiregularity is impossible for every multiplicity \(m\). This
condition is intrinsic under changes of the relation basis and local
coordinates. It can be tested by the ordinary linear map in (E840.3); no
choice of a matrix companion is required for the test. No codimension or
sufficiency statement is made for the remaining locus.

\subsection{4. An exact rank-four witness surviving the previous scalar
equations}\label{SF-E840:an-exact-rank-four-witness-surviving-the-previous-scalar-equations}

Use formal parameters \(x\),\(y\),\(z\),\(w\) and the regular quadratic
intersection

\[
 A=\mathbf C[x,y,z,w]/(x^2+yw,y^2,z^2,w^2).
 \tag{E840.12}
\]

It has length sixteen: it is free with basis \(1,x\) over
\(\mathbf C[y,z,w]/(y^2,z^2,w^2)\). In particular \(xyzw\) is nonzero.
The six squarefree degree-two monomials form a basis of \(W\).

In these parameter coordinates choose the two constant decomposable
transverse bivectors, in the order 12,13,14,23,24,34,

\[
 p=(1,1,0,0,1,1),\qquad
 q=(1,-1,0,0,-1,1),\qquad p\wedge q=4.
 \tag{E840.13}
\]

These are actual transforms of the two standard block bivectors. The
four columns

\[
 (1,0,0,-1),\quad(0,1,1,0),\quad
 (1,0,0,1),\quad(0,1,-1,0)
 \tag{E840.14}
\]

form an invertible matrix of determinant four; the first and second
column wedges are \(p\), and the third and fourth column wedges are
\(q\). Thus the construction is compatible with genuine transverse RM
Poisson coordinates and does not use nonintegrable prescribed gradients.

For the actual four generators in (E840.12), direct differentiation and
reduction give the entire pencil

\[
 P(a,b)=\bigl(4(a+b)xy,\;
 4(a-b)xz-2(a+b)yz,\;0,\;0,\;
 4(a-b)yw,\;4(a+b)zw\bigr).
 \tag{E840.15}
\]

Its rank is four when \((a+b)(a-b)\) is nonzero. It is three when
\(a=b\) is nonzero and two when \(a=-b\) is nonzero. Hence every
five-minor and the full determinant vanish identically. The pure
endpoints both have rank four.

All scalar cross minors

\[
             P_{+,ij}P_{-,kl}-P_{-,ij}P_{+,kl}
 \tag{E840.16}
\]

vanish in \(A\). Indeed, the only nonzero products among
\(xy\),\(xz\),\(yw\),\(zw\) are \(xy\) times \(zw\) and \(xz\) times
\(yw\), both equal to \(xyzw\). The coefficients of each complementary
pair have the same plus-to-minus ratio. The extra \(yz\) term
annihilates each of those four monomials and itself. This verifies every
ordered pair, including the zero tuple entries. Nevertheless the mixed
wedge is \(64xyzw\), and \(\det(\operatorname{Jac}Q)=16xyzw\), both
nonzero.

At the plus endpoint the image is

\[
                   J=\langle xy,\;2xz-yz,\;yw,\;zw\rangle.
 \tag{E840.17}
\]

For \(\ell=\alpha x+\beta y+\gamma z+\delta w\), require \(\ell V\) to
be in \(J\). The quotient \(W/J\) is detected by the \(xw\) coefficient
and by the functional taking \(xz\) coefficient plus twice the \(yz\)
coefficient. The products \(\ell x\) and \(\ell w\) force
\(\delta=\alpha=0\). The second functional on \(\ell x\) forces
\(\gamma=0\), and its value on \(\ell z\) then forces \(\beta=0\). Thus
\(K(J)=0\). The new criterion excludes this example and every
higher-order perturbation having these quadratic initial forms.

\subsection{5. Strictness and actual
examples}\label{SF-E840:strictness-and-actual-examples}

This witness is outside the earlier families for precise invariant
reasons.

First, the relation space has exactly three projective rank-one
quadrics. A general relation is
\(\alpha(x^2+yw)+\beta y^2+\gamma z^2+\delta w^2\). A two-by-two minor
involving the \(x^2\) and \(yw\) entries forces \(\alpha=0\) in rank
one; the remaining diagonal expression has precisely the three
coordinate-square lines. A coordinate-square regular intersection has
four such lines, so no formal coordinate change identifies this algebra
with E837's family.

Second, each pure bracket map from the six-dimensional exterior square
of the relation space has rank four and kernel dimension two. The
exterior square of any independent triple of minimal generators has
dimension three, and hence cannot lie in that kernel. Thus no minimal
presentation has the commuting triple required in E836, for either pure
bivector.

Third, (E840.15) satisfies the identically zero determinant condition of
E838 and all rank-five minors of E839. Equation (E840.16) also satisfies
the scalar companion consequences, so the new obstruction is not merely
a re-use of a failed cross-minor test.

Actual coherent examples with these defects exist by the already
explicit balanced-character hull construction in E838/E839. Choose the
local length-sixteen subscheme with (E840.12) at a point, a locally free
balanced-character hull \(H\), and a surjection \(H\) to \(A^m\) with
\(m\) at most rank \(H\). Let \(K\) be its kernel. Extensions

\[
                      0\longrightarrow A^m
                      \longrightarrow E\longrightarrow K
                      \longrightarrow0
 \tag{E840.18}
\]

have character \(\operatorname{ch}(E)=\operatorname{ch}(H)\), torsion
\(A^m\), double dual \(H\), and hull quotient \(A^m\). The extension
space has dimension \(6\cdot16\cdot m^2\), by the complete-intersection
Koszul resolution and the vanishing of positive Ext from the locally
free hull to the point-supported module. These are actual coherent,
ordinarily gluable objects. The theorem excludes weak balanced middle
semiregularity for every such extension; no simplicity or successful
finite action is asserted.

\subsection{6. Sources, checks, and
boundary}\label{SF-E840:sources-checks-and-boundary}

The new multiplication and equality arguments are proved in full above.
Their inherited actual-companion, invariant/B-field, associated-graded,
and socle inputs are exactly the already checked E819/E836/E838 ones.
The published Jacobian-socle locator is Eisenbud--Ulrich, J. reine
angew. Math. 756 (2019), Theorem 7.1 p.210 and Appendix A.1/A.3/A.5
pp.222--225.\cite{SFEisenbudUlrich2019} The invoked complete proofs were read in the preceding
E836/E838 source audit; no new unrefereed premise is used.

The separate exact certificate reconstructs the complete length-sixteen
quotient and actual bracket matrices; it supplements the universal
proof. The complete independent review has passed. Rank-at-most-three pencils,
rank-four directions with principal multiplication image, higher initial
degrees, unequal or nonfree Artinian defects, and unrelated perfect
complexes remain open. This theorem is not an unrestricted
complete-intersection or arbitrary-singular-sheaf exclusion.

\section{SF-E842--843. A socle-pair obstruction and arbitrary parameter-power point defects}
\label{SF-E842:section}

\paragraph{Status and provenance.}
Author: \texttt{/root/sixfolds}, GPT-6 Astra Pro; inherited effort is not exposed.
The complete frozen Markdown proof was prepared on 22 November 2026 during
an execution-environment outage and preserved in runtime memory. The author,
root, and the independent \texttt{/root/g34\_eightfolds} reviewer have read the
complete proof; all three mathematical reviews passed. The peer also checked
the complete author program and produced an independent exact certificate.
The durable draft preserves its original pending-review status and all
outage-time metadata. SF-E841 remains a separate unpromoted derived draft.

\subsection{Local setup and the previously proved deformation interface}
\label{SF-E842:interface}

Let \(R=\mathbb C[[x_1,x_2,x_3,x_4]]\), let
\(I=(f_1,f_2,f_3,f_4)\) be a proper finite-colength complete-intersection
ideal, and set \(A=R/I\). Let \(\pi_+,\pi_-\) be constant rank-two
bivectors with complementary supports. Thus \(\pi_+\wedge\pi_-\) is a
nonzero scalar multiple of the coordinate volume. Write \(P_+,P_-\)
for the alternating four-by-four matrices, or equivalently their six
coefficient tuples,
\[
 P_{\pm,ij}=\pi_\pm(df_i,df_j)\pmod I,\qquad i<j.
 \tag{E842.1}
\]

Consider an actual finite \(R\)-module with the hull-modification sequence
\[
 0\longrightarrow A^n\longrightarrow M\longrightarrow R^r
       \longrightarrow A^m\longrightarrow0,\qquad m+n>0.
 \tag{E842.2}
\]

The exact E819/E836 interface is used, with no new deformation premise.
If \(M\) has flat first-order deformations for both multiplication classes
\(\pi_+\) and \(\pi_-\), the finite-perfect-cone supertrace argument gives
\(m=n>0\). The lifted map to the free hull and the ordinary Koszul extension
description then give actual square matrices
\(C_+,C_-,B_{ij}\in\operatorname{Mat}_m(A)\) with both orders
\[
 C_\pm B_{ij}=B_{ij}C_\pm=P_{\pm,ij}I_m.
 \tag{E842.3}
\]

Here the same \(B_{ij}\) are the extension matrices for both bivectors.
A chosen free summand in the hull image need not remain split under
deformation.

For precision, this interface uses
\(\operatorname{Ext}_R^j(A,R)=0\) for \(j<4\), which lifts the hull map.
Its first-order value on torsion defines \(C_\pi:T\to Q\).
Lifting the relation \(f_i e_j-f_j e_i=B_{ij}\) in \(M\) and applying that
hull map gives \(C_\pi B_{ij}=P^\pi_{ij}I\).
Applying the deformed commutator of \(f_i,f_j\) to a lifted torsion element
gives \(B_{ij}C_\pi=P^\pi_{ij}I\). These are the two complete actual-module
calculations proved in E836, Sections 2--3. No commutativity among the
\(B_{ij}\), or between \(C_+\) and \(C_-\), is assumed here.

The complete-intersection Jacobian-socle theorem is also retained with its
precise characteristic-zero hypothesis:
\[
 \Delta_f=\det(\partial f_i/\partial x_j)\pmod I
       \quad\hbox{generates the nonzero one-dimensional socle of }A.
 \tag{E842.4}
\]

The exact published source is Eisenbud--Ulrich, \emph{Duality and socle
generators for residual intersections}, \emph{J. reine angew. Math.}
756 (2019), 183--226, Theorem 7.1, p.~210, with the complete invoked proofs
in Appendix A.1, A.3 and A.5, pp.~222--225.\cite{SFEisenbudUlrich2019}
Those complete proofs were read in the earlier E836/E838/E839 reviews.
The formal-local theorem statement and the complete invoked Appendix proofs
were freshly read again during the present author and independent peer
audits. No preprint is substituted for this input.

\subsection{The elementary square-matrix socle-pair lemma}
\label{SF-E842:matrix}

\paragraph{Lemma SF-E842a.}
Let \(A\) be a commutative local complex Artinian algebra. Suppose
\(p,q\in A\) satisfy
\[
 0\ne pq\in\operatorname{Soc}(A),\qquad
                         pA\cap qA=\mathbf C\,pq.
 \tag{E842.5}
\]

If \(C,B_p,B_q\) are square matrices of positive size over \(A\) and
\[
 CB_p=B_pC=pI,\qquad CB_q=B_qC=qI,
 \tag{E842.6}
\]
then \(C\) is invertible.

\paragraph{Proof.}
The two orders, not any assumed commutation of \(B_p\) and \(B_q\), give
\[
 pB_q=B_pCB_q=qB_p.
 \tag{E842.7}
\]

Each entry of \(qB_p\) therefore belongs to \(pA\cap qA\). Consequently
\[
 qB_p=pqD
 \tag{E842.8}
\]
for a constant complex matrix \(D\). Apply \(C\) to this equality.
Since \(pq\) is a nonzero socle element, every entry of the maximal ideal
kills it, so
\[
 pqI=qCB_p=C(pqD)=pq\,C(0)D.
 \tag{E842.9}
\]

The nonzero scalar vector \(pq\) gives \(C(0)D=I\) over \(\mathbb C\).
The square complex matrix \(C(0)\) is invertible, and therefore \(C\) is
invertible over the local ring \(A\). This proves the lemma for every
matrix size and with arbitrary higher entries. No quotient deformation
or freeness under base change has been used.

\subsection{An invertible symplectic pencil companion is impossible}
\label{SF-E842:pencil}

\paragraph{Lemma SF-E842b.}
In the complete-intersection setup (E842.1)--(E842.3), no matrix
\[
                  C=aC_++bC_-,\qquad ab\ne0,
 \tag{E842.10}
\]
is invertible.

\paragraph{Proof.}
Put \(P=aP_++bP_-\). The two existing companion equations imply
\(CB_{ij}=B_{ij}C=P_{ij}I\); a deformation in a third direction is not
being assumed.

Each \(P_\pm\) is a decomposable two-form over \(A\), since it is obtained
by applying the Jacobian matrix to the decomposable constant bivector
\(\pi_\pm\). Hence its Pfaffian is zero. The Pfaffian transformation
identity gives
\[
 \operatorname{Pf}(P_\pm)=0,\qquad
 \operatorname{Pf}(P)
   =\Delta_f\,\operatorname{Pf}(a\Pi_++b\Pi_-)
   =\lambda\Delta_f,\quad \lambda\ne0.
 \tag{E842.11}
\]

Here \(\Pi\) denotes the constant skew matrix in the original \(x\)
coordinates. The coefficient \(\lambda\) is nonzero because \(ab\ne0\)
and the two planes are complementary. By (E842.4),
\(\operatorname{Pf}(P)\) is a nonzero socle element.

If \(C\) were invertible, set \(D_+=aC_+C^{-1}\),
\(D_-=bC_-C^{-1}\), and \(d_\pm=\operatorname{tr}(D_\pm)/m\in A\).
These are defined over \(A\) in characteristic zero.
Substitution of \(B_{ij}=C^{-1}P_{ij}\) gives
\[
 D_+P_{ij}=aP_{+,ij}I,\quad
 D_-P_{ij}=bP_{-,ij}I,\quad
 D_++D_-=I.
 \tag{E842.12}
\]

Taking traces yields the exact identities
\[
             aP_+=d_+P,\qquad bP_-=d_-P,\qquad d_++d_-=1.
 \tag{E842.13}
\]

Their Pfaffians give
\(d_+^2\operatorname{Pf}(P)=d_-^2\operatorname{Pf}(P)=0\).
Multiplication on the nonzero one-dimensional socle only sees the constant
term. Thus \(d_+(0)^2=d_-(0)^2=0\), so both constants are zero.
This contradicts \(d_+(0)+d_-(0)=1\). The lemma follows.
Notice that no homogeneous initial form, matrix diagonalization,
or proportionality assumption was needed.

\paragraph{Theorem SF-E842.}
Suppose that for some \(ab\ne0\), two entries \(p,q\) of
\(P=aP_++bP_-\) satisfy the socle-pair condition (E842.5).
Then no module (E842.2) has flat first-order deformations for both
\(\pi_+\) and \(\pi_-\).

\paragraph{Proof.}
If a multiplicity is zero, the already proved supertrace equality
contradicts \(m+n>0\). Otherwise use the common positive square size and
\(C=aC_++bC_-\). Lemma E842a makes \(C\) invertible, whereas E842b
forbids this. This proves the theorem. The socle-pair condition is a
sufficient, explicitly testable criterion; it is not claimed for every
complete intersection.

\subsection{All parameter-power complete intersections, with nonlinear mixed parameters}
\label{SF-E843:powers}

\paragraph{Theorem SF-E843.}
Let \(u_1,\ldots,u_4\) be any regular system of parameters of \(R\).
They may be nonlinear formal functions mixing the two Poisson planes.
For arbitrary positive integers \(a_i\), put
\[
                  I=(u_1^{a_1},u_2^{a_2},u_3^{a_3},u_4^{a_4}).
 \tag{E843.1}
\]

Then the conclusion of SF-E842 holds for every nontrivial module (E842.2).

\paragraph{Proof.}
Set \(v_i=u_i^{a_i-1}\in A\), and \(\tau=v_1v_2v_3v_4\).
The monomial basis \(0\le r_i<a_i\) proves that \(A\) has length
\(\prod_i a_i\) and socle \(\mathbb C\tau\), with \(\tau\ne0\).
Put \(E_{ij}=v_iv_j\).

Use the already existing pencil \(\pi=\pi_++\pi_-\), which is
nondegenerate. Its coefficient matrix on the cotangent vectors \(du_i\)
at the closed point is nondegenerate because those vectors form a basis.
Its nonzero Pfaffian implies that at least one complementary pair of entries
\(h_{ij}(0),h_{kl}(0)\) is nonzero, where
\(\{i,j,k,l\}=\{1,2,3,4\}\). The actual entries of \(P\) are
\[
 P_{ij}=a_i a_j E_{ij}h_{ij},\qquad
                   h_{ij}=\pi(du_i,du_j)\pmod I.
 \tag{E843.2}
\]

For that complementary pair, both \(h_{ij}\) and \(h_{kl}\) are units.
Hence the two principal ideals are \(E_{ij}A\) and \(E_{kl}A\).
In the monomial basis, the first requires maximal exponent on \(i,j\)
and the second on \(k,l\). Their intersection requires maximal exponents
on every variable and is exactly \(\mathbb C\tau\).
Their product is a nonzero constant multiple of \(\tau\).
Thus these two actual entries satisfy (E842.5), and E842 proves the assertion.

This proof handles every exponent and every nonlinear parameter change
directly in \(A\). In particular, \(A\) is never replaced by its square
quotient; no deformation, Ext group, or free module is improperly base changed.
If some \(a_i=1\), the assertion includes cases already excluded by E836.
The new scope includes all \(a_i\ge2\) in arbitrary mixed parameter coordinates.

\subsection{Exact strict cubic witnesses}
\label{SF-E843:witnesses}

Take constant Poisson coordinates \(X=(x_1,x_2,y_1,y_2)\), with
\(\pi_+=\partial_{x_1}\wedge\partial_{x_2}\) and
\(\pi_-=\partial_{y_1}\wedge\partial_{y_2}\).
Define \(w=JX\) by the integral matrix
\[
 J=\begin{pmatrix}
 1&0&1&0\\
 0&1&0&1\\
 -3&5&-9&10\\
 -8&5&-16&15
 \end{pmatrix},\qquad \det J=-20.
 \tag{E843.3}
\]

The bivector coefficient tuples in the order \(12,13,14,23,24,34\) are
\[
 p=(1,5,5,3,8,25),\qquad q=(1,10,15,9,16,25).
 \tag{E843.4}
\]

Each tuple is decomposable: its Pfaffian is respectively
\(25-40+15=0\) and \(25-160+135=0\).
Their mixed wedge coefficient is \(-20\), so the planes are transverse.
This is the earlier dense rational boundary after scaling the last two
parameters by five; the present certificate is entirely integral.

There are two explicit variants. In the first set \(u_i=w_i\).
In the second set
\[
 u_1=w_1+w_3^2,\quad u_2=w_2+w_4^2,\quad u_3=w_3,\quad u_4=w_4.
 \tag{E843.5}
\]

The second change is a polynomial automorphism with determinant one
and triangular inverse
\(w_1=u_1-u_3^2,\ w_2=u_2-u_4^2,\ w_3=u_3,\ w_4=u_4\).
In either case set \(I=(u_1^3,\ldots,u_4^3)\).
The algebra has length \(81\), embedding dimension four and socle degree
eight; \(\tau=u_1^2u_2^2u_3^2u_4^2\ne0\).

For the linear variant, all six brackets are exactly
\(P_{+,ij}=9p_{ij}E_{ij}\) and \(P_{-,ij}=9q_{ij}E_{ij}\).
In the nonlinear variant the same formula has coefficient units whose
constants are \(p_{ij},q_{ij}\).
Every scalar cross-minor vanishes in both variants:
\[
                  P_{+,e}P_{-,f}-P_{-,e}P_{+,f}=0.
 \tag{E843.6}
\]

If \(e,f\) share an index, the product contains \(u_i^4\) and is zero.
If they are complementary, the product of the \(E\) entries is \(\tau\),
which kills every positive-degree coefficient.
The relevant ratios \(q_e/p_e\) are equal on each complementary pair:
\(1\) on \(12/34\), \(2\) on \(13/24\), and \(3\) on \(14/23\).
This proves all scalar cross-minors, including all ordered pairs.

The actual generator Jacobian and the mixed bracket wedge are
\[
 \det\bigl(\partial(u_i^3)/\partial X_j\bigr)
       =-1620\,\tau,\qquad
                P_+\wedge P_-=-1620\,\tau
 \tag{E843.7}
\]
with the mixed-wedge coefficient convention of (E843.4).
Indeed \(3^4\det J=81(-20)\), and the nonlinear triangular change
has determinant one.

For \(P=P_++P_-\), choose its \(12\) and \(34\) entries.
Their constant coefficients relative to \(E_{12},E_{34}\) are \(18\)
and \(450\). Their ideals each have dimension nine, their intersection
has dimension one, and their product is \(8100\tau\).
This is the socle pair used by the new theorem.

These examples are outside the earlier local exclusions for exact reasons.
There is no linear generator and the embedding dimension is four.
The lowest ideal order is three, invariant under formal coordinate
and minimal-generator changes, so no regular quadratic initial presentation
exists; the square-thickening and quadratic-jet E837--840 families do not apply.
For each pure bivector the leading quartic bracket map from the exterior
square of the four-dimensional minimal relation space has rank six:
all six coefficients in (E843.4) are nonzero and the six \(E_{ij}\)
are independent degree-four monomials.
A triple of independent minimal generators has a three-dimensional exterior
square, which cannot lie in that zero kernel.
The exact covariance under an arbitrary invertible generator change is
the one proved in E836. Thus no minimal presentation has the commuting
triple required by E836, for either pure bivector.
Finally, (E843.6) shows that this is not an example already rejected
by the scalar companion minors.

\subsection{Actual coherent balanced examples and the equivariant consequence}
\label{SF-E843:sheaves}

Use the actual locally free balanced-character hull \(H\) from E815/E820,
or a direct sum of copies, on the explicit genuine RM abelian fourfold.
Choose a point and the finite subscheme \(Z\) with one of the preceding
completed ideals. It is an algebraic finite subscheme: the ideal contains
the ninth power of the maximal ideal, so its finite quotient is already
a quotient of the algebraic local ring modulo that power.
Thus no algebraization of an infinite formal family is asserted.

For any positive \(m\), take enough copies of \(H\) to have rank at least
\(m\), choose a surjection \(H\to\mathcal O_Z^m\) using a local basis
on this finite subscheme, and put \(K\) equal to its kernel.
There are actual extensions
\[
 0\longrightarrow O_Z^m\longrightarrow E\longrightarrow K\longrightarrow0,
 \qquad
                \operatorname{ch}(E)=\operatorname{ch}(H).
 \tag{E843.8}
\]

The torsion is precisely \(\mathcal O_Z^m\), the double dual is \(H\),
and the hull quotient is \(\mathcal O_Z^m\): \(K\) is torsion-free and
agrees with \(H\) away from codimension four.
Smoothness makes these coherent sheaves perfect, and coherence gives
ordinary negative-Ext vanishing.
Because positive Ext from \(H\) to a point-supported sheaf vanishes,
the complete-intersection Koszul resolution gives
\[
 \operatorname{Ext}^1(K,O_Z^m)\simeq
 \operatorname{Ext}^2(O_Z^m,O_Z^m),\qquad
                         \dim=6\cdot81\,m^2=486m^2.
 \tag{E843.9}
\]

Thus the examples form actual extension families, not formal characters.

The already proved E819 invariant-to-ordinary characteristic-map argument
says that weak balanced middle semiregularity for any specified coherent
effective finite translation/tensor action forces the two local
deformations for \(\pi_+,\pi_-\).
Its rational algebraic \(B\)-field version leaves these bivector projections
unchanged. Restricting at the chosen point and applying E843 therefore
excludes weak semiregularity for every extension in (E843.8).
No action on a chosen local basis, no simplicity, and no equivariant
splitting of the extension is assumed.
This does not assert a universal obstruction for unrelated Artinian modules
or arbitrary perfect complexes.

\subsection{Exact computation and check status}
\label{SF-E843:checks}

A new exact certificate ran in the available JavaScript runtime at
\texttt{2026-11-22T00:56:11.438Z}, with recorded duration \(0.006\) seconds.
The first run passed.
Every polynomial integer addition and product is guarded by
\texttt{Number.isSafeInteger}.
It reconstructs both actual gradient matrices, all \(6561\) basis products
of the length-\(81\) algebra, both pure Pfaffians, the nonzero mixed
and actual Jacobian, all \(36\) ordered scalar cross-minors in each variant,
both initial ranks, the two nine-dimensional principal ideals and their
one-dimensional socle intersection.
This finite check supplements the universal all-size proofs.

The complete runtime source and original first result were preserved
during the outage and saved after restoration as
\texttt{rm\_socle\_pair\_power\_certificate.js} and
\texttt{rm\_socle\_pair\_power\_certificate\_first\_result.json}.
The frozen v1 draft is
\texttt{rm\_socle\_pair\_power\_obstruction\_draft.md}.
A restored-runtime rerun at \texttt{2026-11-22T01:09:37.894Z} passed in
\(0.018\) seconds and had exactly the same invariant output.
The independent peer's separately authored exact reconstruction also passed;
its original failure-and-repair history concerns only its own initial
\texttt{Array.map} callback and is retained in the peer artifacts.
Root has now read the complete canonical TeX proof, the entire author
program and its first complete output, and the complete independent peer
report. The separate peer also bound its full review to this canonical
file and checked that all22 displays match the frozen draft. No fresh
root execution of the author or peer programs is claimed here.
No successful semiregular seed has been constructed.

\section{SF-E844. A non-power complete-intersection family satisfying the socle-pair criterion}\label{SF-E844:sf-e844.-a-non-power-complete-intersection-family-satisfying-the-socle-pair-criterion}

Author: /root/sixfolds, GPT-6 Astra Pro; inherited effort not exposed.
Date: 2026-11-22.

\textbf{Status:} the complete author, root and separate independent peer whole-proof readings have passed. This is an explicit derived application of the already accepted SF-E842 criterion, not a separate reduction of the residual gate defined by that criterion.
The theorem SF-E842 is an already whole-reviewed input; this note does not
modify that frozen theorem or claim its socle criterion for every complete
intersection. The new point is an explicit family outside the parameter-power
scope of E843 and all the previously listed scalar/commuting-triple/quadratic
tests. No semiregular seed is constructed.

\subsection{1. The actual complete intersections and their finite algebra}\label{SF-E844:the-actual-complete-intersections-and-their-finite-algebra}

Use the same constant Poisson coordinates and integral mixing matrix as E843:
\(X=(x_1,x_2,y_1,y_2)\), with
\(\pi_+=\partial_{x_1}\wedge\partial_{x_2}\) and
\(\pi_-=\partial_{y_1}\wedge\partial_{y_2}\), and put \(u=JX\), where

\[
 J=\begin{pmatrix}
 1&0&1&0\\
 0&1&0&1\\
 -3&5&-9&10\\
 -8&5&-16&15
 \end{pmatrix},\qquad \det J=-20.
 \tag{E844.1}
\]

For any nonzero complex number \(\lambda\), define

\[
 f_1=u_1^3+\lambda u_2^2u_3^2,\quad
 f_2=u_2^3,\quad f_3=u_3^3,\quad f_4=u_4^3,\qquad
 I_\lambda=(f_1,f_2,f_3,f_4),\quad
 A_\lambda=\mathbb C[[u_1,u_2,u_3,u_4]]/I_\lambda.
 \tag{E844.2}
\]

The algebra is finite and is a complete intersection. To verify finiteness
and its exact basis without a flatness assumption, set
\(B=\mathbb C[u_2,u_3,u_4]/(u_2^3,u_3^3,u_4^3)\).
The polynomial quotient
\(B[u_1]/(u_1^3+\lambda u_2^2u_3^2)\) is free over \(B\), with basis
\(1,u_1,u_1^2\), because the polynomial in \(u_1\) is monic.
Its generators are all nilpotent:
\(u_1^6=\lambda^2u_2^4u_3^4=0\) and \(u_j^3=0\) for \(j=2,3,4\).
It therefore equals its formal completion and is \(A_\lambda\).
Consequently

\[
 \{u_1^{r_1}u_2^{r_2}u_3^{r_3}u_4^{r_4}:0\le r_i\le2\}
 \text{ is a basis},\qquad \dim_{\mathbb C}A_\lambda=81.
 \tag{E844.3}
\]

The ideal is generated by four elements and has height four in the regular
local ring of dimension four; hence they form a regular sequence.
Alternatively, the first three coordinate powers \(u_2^3,u_3^3,u_4^3\)
form a regular sequence and the remaining monic relation is a nonzerodivisor,
which can be checked by successive powers of the nilpotent ideal
\((u_2,u_3,u_4)\).

The only nonzero reduction of a forbidden exponent is

\[
                   u_1^3=-\lambda u_2^2u_3^2.
 \tag{E844.4}
\]

This reduction increases total degree by one. A monomial of total degree
at least nine either vanishes during reduction or would have to reduce
to a basis monomial of degree at least nine, whereas all basis monomials
have degree at most eight. Thus the ninth maximal-ideal power is zero.
The degree filtration of the displayed basis has associated graded algebra
\(\mathbb C[u_1,u_2,u_3,u_4]/(u_1^3,u_2^3,u_3^3,u_4^3)\).
In particular, the lowest nonzero ideal order is three and the embedding
dimension is four.

The element
\(\tau=u_1^2u_2^2u_3^2u_4^2\) is nonzero by the basis and is killed by
every \(u_i\), including \(u_1\) by (E844.4).
The complete-intersection Jacobian-socle theorem, in its precise
characteristic-zero form already used in E842, gives
\(\operatorname{Soc}(A_\lambda)=\mathbb C\tau\). The exact published locator is Eisenbud--Ulrich, Theorem 7.1, p.~210, with Appendix A.1/A.3/A.5, pp.~222--225.\cite{SFEisenbudUlrich2019}

\subsection{2. The two actual bracket entries form a socle pair}\label{SF-E844:the-two-actual-bracket-entries-form-a-socle-pair}

The constant coefficient tuples of \(\pi_+,\pi_-\) in the \(u\)-coordinates,
in the order \(12,13,14,23,24,34\), are

\[
 p=(1,5,5,3,8,25),\qquad q=(1,10,15,9,16,25).
 \tag{E844.5}
\]

Put \(P_\pm=(\pi_\pm(df_i,df_j))_{i<j}\) in \(A_\lambda^6\), and
\(P=P_++P_-\). The derivative of the additional term in \(f_1\) is

\[
 d(\lambda u_2^2u_3^2)
       =2\lambda u_2u_3^2\,du_2+2\lambda u_2^2u_3\,du_3.
 \tag{E844.6}
\]

In the bracket with \(f_2=u_2^3\), the first term pairs trivially with
\(du_2\), while the second contains \(u_2^4\) and vanishes.
The unmodified pair \(f_3,f_4\) is immediate. Therefore the two exact
pencil entries, including all higher terms, are

\[
             P_{12}=18u_1^2u_2^2,\qquad
             P_{34}=450u_3^2u_4^2.
 \tag{E844.7}
\]

Let \(E_{12}=u_1^2u_2^2\), \(E_{34}=u_3^2u_4^2\).
The ideal \(E_{12}A_\lambda\) has the nine basis monomials
\(E_{12}u_3^au_4^b\), \(0\le a,b\le2\).
Indeed multiplication by \(u_2\) kills it, and multiplication by \(u_1\)
also kills it because \(u_1E_{12}=-\lambda u_2^4u_3^2=0\).
Likewise, \(E_{34}A_\lambda\) has basis
\(E_{34}u_1^au_2^b\), \(0\le a,b\le2\).
The relation \(u_1^3=-\lambda u_2^2u_3^2\) contributes nothing after
multiplication by \(E_{34}\), because it then contains \(u_3^4\).
The two indicated basis subsets meet exactly in \(\tau\). Hence

\[
 P_{12}A_\lambda\cap P_{34}A_\lambda=\mathbb C\tau,\qquad
                        P_{12}P_{34}=8100\tau\ne0.
 \tag{E844.8}
\]

This is precisely the sufficient condition in SF-E842. No comparison with a
different quotient, and no base change of a deformation, has been used.

The Jacobian provides a further exact check. In the \(u\)-coordinates the
derivative matrix is triangular apart from the two entries in the first row,
so its determinant is \(81\tau\). Returning to \(X\) gives

\[
 \det\bigl(\partial f_i/\partial X_j\bigr)
       =81\det(J)\tau=-1620\tau,\qquad
                    P_+\wedge P_-=-1620\tau.
 \tag{E844.9}
\]

The wedge convention is the one fixed in E843.4.

\subsection{3. These are not parameter-power ideals in any formal coordinates}\label{SF-E844:these-are-not-parameter-power-ideals-in-any-formal-coordinates}

Suppose, for contradiction, that there are formal parameters \(v_1,\ldots,v_4\)
and positive integers \(a_i\) such that
\(I_\lambda=(v_1^{a_1},\ldots,v_4^{a_4})\).
There is no element of ideal order less than three, so every \(a_i\ge3\).
The degree-three initial relation space of \(I_\lambda\) has dimension four.
Only generators with \(a_i=3\) can contribute to this space, so all four
exponents must equal three.

Write \(\ell_i\) for the linear part of \(v_i\).
Since \(v_i^3\in I_\lambda\), its initial cubic belongs to

\[
                 \operatorname{span}_{\mathbb C}
                           (u_1^3,u_2^3,u_3^3,u_4^3).
 \tag{E844.10}
\]

A nonzero linear form whose cube belongs to this span has only one nonzero
coordinate: otherwise its cube has a nonzero coefficient \(3a^2b\)
on a mixed monomial \(u_j^2u_k\).
The four \(\ell_i\) are independent, so after reordering and rescaling the
\(v_i\) we may assume \(\ell_i=u_i\).

Write \(v_1=u_1+Q_1+O(\mathfrak m^3)\), with \(Q_1\) homogeneous quadratic.
Because \(v_1^3\) belongs to the ideal generated by the four \(f_j\), write
\(v_1^3=\sum_j a_{1j}(u)f_j\).
Comparison of degree-three terms gives \(a_{11}(0)=1\) and
\(a_{1j}(0)=0\) for \(j\ne1\).
Comparison of degree-four terms therefore gives, for certain linear forms
\(L_j\),

\[
              3u_1^2Q_1
                 =\lambda u_2^2u_3^2+\sum_{j=1}^4L_j u_j^3.
 \tag{E844.11}
\]

This is impossible. The monomial \(u_2^2u_3^2\) is divisible neither by
\(u_1^2\) nor by any \(u_j^3\), so its coefficient is zero on the left
and \(\lambda\ne0\) on the right. This excludes every formal parameter
change, without assuming that it preserves the Poisson planes.

\subsection{4. The old local tests still do not exclude this family}\label{SF-E844:the-old-local-tests-still-do-not-exclude-this-family}

Every bracket has order at least four. The additional terms from (E844.6)
have order at least five. The product of a higher bracket term with any
bracket term has order at least nine and is therefore zero in \(A_\lambda\).
Thus every scalar cross-minor is determined by the original quartic terms:

\[
                P_{+,e}P_{-,f}-P_{-,e}P_{+,f}=0
                       \quad\text{for every ordered }e,f.
 \tag{E844.12}
\]

For a direct verification, if the two edge indices overlap outside
coordinate \(1\), the product contains \(u_j^4\) with \(j\ge2\) and vanishes.
If the shared index is \(1\), replacing \(u_1^3\) by (E844.4) raises the
degree of the degree-eight product to at least nine, so it vanishes too.
For complementary edges the product is \(\tau\); the coefficient vanishes
because the ratios \(q_e/p_e\) agree on each complementary pair, as in E843.
This proves all scalar cross-minors, rather than only selected minors.

The initial pure bracket maps have rank six, since all coefficients in
(E844.5) are nonzero and the six monomials \(u_i^2u_j^2\) are independent
in the degree-four associated graded quotient.
By the exact covariance under arbitrary minimal-generator changes,
no independent generator triple can commute for either pure bivector.
Thus E836's commuting-triple criterion does not apply in any presentation.

The ideal has lowest order three and embedding dimension four, so it is
outside the square and regular quadratic-jet families E837--840, as well as
the lower-embedding-dimension part of E836. Section 3 excludes E843's full
parameter-power family, even allowing nonlinear mixed coordinates.
The assertion here is precisely an explicit additional application of E842;
it is not a claim that any unclassified complete intersection satisfies E842.

\subsection{5. Actual coherent objects excluded by the criterion}\label{SF-E844:actual-coherent-objects-excluded-by-the-criterion}

On the genuine RM abelian fourfold with the reviewed balanced locally free
hull \(H\), realize this quotient at a point as a finite subscheme \(Z_\lambda\).
This is algebraic because \(I_\lambda\) contains the ninth maximal-ideal
power. Take sufficiently many copies of \(H\), choose a surjection
\(H\to\mathcal O_{Z_\lambda}^m\), and let \(K\) be its kernel.
Every extension

\[
 0\longrightarrow\mathcal O_{Z_\lambda}^m
   \longrightarrow E\longrightarrow K\longrightarrow0
 \tag{E844.13}
\]

is an actual coherent, hence perfect and ordinarily gluable, sheaf.
Its torsion is exactly \(\mathcal O_{Z_\lambda}^m\), its double dual is \(H\),
its hull quotient is \(\mathcal O_{Z_\lambda}^m\), and
\(\operatorname{ch}(E)=\operatorname{ch}(H)\).
The same exact sequence and complete-intersection Koszul calculation as in
E843 give the actual extension-space dimension \(486m^2\).

For any specified allowed coherent effective finite translation/tensor action,
weak balanced middle semiregularity gives the two local ordinary bivector
deformations by the reviewed E819 interface, also with a rational algebraic
\(B\)-field. The socle pair (E844.8) and E842 exclude those deformations.
Thus every extension (E844.13) fails the required semiregularity.
No equivariant splitting or simplicity is assumed.

\subsection{6. Exact certificate and review boundary}\label{SF-E844:exact-certificate-and-review-boundary}

The companion criterion and universal all-multiplicity obstruction are the
written E842 proof. A separate exact quotient-ring program reconstructs this
family, its multiplication and actual gradients; its numerical scope is
recorded in the adjacent JSON. The non-power assertion is the formal
degree-three/four argument in Section 3, not an inference from finitely many
coordinate changes. The separate independent peer has read the entire draft and canonical proof, the whole author checker and its complete successful output, and has executed a distinct exact symbolic reconstruction. Root read the entire draft, canonical proof and peer report and compared all thirteen displays. No fresh root code execution or whole-code reading is claimed for this example.

\paragraph{Certificate provenance.}
The exact BigInt checker passed at \texttt{2026-11-22T01:17:56.329Z} in \(0.060296346\) seconds. It checks the \(\lambda=1\) witness, including all \(6561\) basis products, \(531441\) associativity triples, \(220\) degree-nine monomials, \(36\) scalar cross-minors, the actual Jacobian, and both principal-ideal dimensions. The first attempt stopped with a Boolean--BigInt precedence error; the original source, empty standard output, error and one-parenthesis repair are preserved. No mathematical revision was required. The universal \(\lambda\ne0\), formal-coordinate and all-multiplicity assertions rest on the written proof.

\section{SF-E846: homogeneous quadratic point modifications}\label{SF-E846:section}

\textbf{Checked mathematical status, 2026-11-22.} The author, root,
and independent reviewer \texttt{/root/g34\_eightfolds} have read the complete
frozen proof; all three mathematical reviews passed, with no repair.
The independent full proof/source/code report was completed by
\texttt{2026-11-22 02:30:14 UTC}. The frozen Markdown original retains its
author-review status. The author and root have also read all 604 lines of the canonical conversion; all 19 frozen display bodies are byte-identical. No strict comparison with all previous local Artin exclusions, and no independent main-gate movement, is claimed.
Author: \texttt{/root/sixfolds}, GPT-6 Astra Pro; inherited reasoning effort is not exposed.
The all-order noncommutative semiregularity implication is not assumed here.
The global argument below instead constructs a class in ordinary Ext.
Its exact scope is homogeneous scalar local extensions; arbitrary filtered
perturbations and arbitrary matrix-valued extension classes are not included.

\subsection{1. The exact local module and its resolution}\label{SF-E846:s1}

Put \(R=\mathbb C[[x_1,x_2,x_3,x_4]]\). Let \(f_1,\ldots,f_4\) be homogeneous quadrics forming a
regular sequence, \(I=(f_1,\ldots,f_4)\), and \(A=R/I\). Let \(U=\mathbb C^4\) with ordered basis
\(e_1,\ldots,e_4\). Let \(\iota_i\) denote contraction by the dual basis element and let
\(e_i\wedge\) denote exterior multiplication. Write

\[
 K_p=R\otimes\Lambda^pU,\quad d=\sum_i f_i\iota_i,
 \qquad
 \beta=\sum_{i<j}b_{ij}\iota_j\iota_i,
 \tag{HK.1}
\]

where the six \(b_{ij}\) are linear forms, not all zero. Scalars commute
throughout. No spanning hypothesis on the \(b_{ij}\) is imposed. The convention makes
\(\beta(e_i\wedge e_j)=b_{ij}\) for \(i<j\). The even contraction \(\beta\) commutes with \(d\).
There is an actual extension \(M_b\) with graded exact sequence

\[
 0\longrightarrow A(-3)\longrightarrow M_b\longrightarrow I
 \longrightarrow0.
 \tag{HK.2}
\]

Here an \(R(-j)\) generator has degree \(j\). A resolution has
\(L_p=K_p(-3)\) direct-sum \(K_{p+1}\) for \(0\le p\le3\), with \(L_4=K_4(-3)\).
In this notation the first \(K_p\) generator has degree \(3+2p\) and the second
\(K_{p+1}\) generator has degree \(2+2p\). Its differential is

\[
 D=\begin{pmatrix}d&\beta\\0&-d\end{pmatrix}.
 \tag{HK.3}
\]

More explicitly the bottom complex is the resolution of \(I\), ending in \(K_1\);
the top complex resolves \(A(-3)\). Their triangular extension is exact in
negative cohomological degrees by the cohomology sequence of the two
complexes. Its degree-zero cohomology is (HK.2). The canonical hull map
\(\phi:M_b\to R\) is the map onto \(I\) followed by inclusion. Its kernel is \(A(-3)\)
and its cokernel is \(A\). The module is rank one and is locally free away
from the closed point. All displayed maps are continuous for the maximal
ideal topology. The notation \(K_p(-3)\) is shorthand for the stated generator
degrees, not for a second independent shift of an already shifted \(K_p\).

\subsection{2. First-order corrections, with every lowest weight exhausted}\label{SF-E846:s2}

Let \(\pi\) be a constant nondegenerate bivector. Use right modules and the
Moyal product

\[
 a\star b=ab+\frac h2\{a,b\}_\pi+
 \frac{h^2}{8}\sum\pi^{ij}\pi^{kl}
       (\partial_i\partial_k a)(\partial_j\partial_l b)
        \pmod {h^3}.
 \tag{HK.4}
\]

The opposite convention is obtained by replacing \(\pi\) by \(-\pi\); existence
and nonexistence statements are unchanged. A matrix acts on a column
of a free right module by left multiplication. Composition of matrices
therefore uses star products in the displayed order.

Give \(h\) degree two. Assume \(M_b\) admits a flat first-order deformation.
The finite free resolution lifts with the same ranks, by successive
kernel lifting and nilpotent Nakayama exactly as in E819 section 2.1.
Write its differential \(D+hH\). A homogeneous endomorphism of cohomological
degree one and internal degree below -3 is zero by the generator shifts.
The degree -3 part can only be a constant lower-left block \(c_p:K_p\to K_p\)
(\(1\le p\le4\)). The equation \(dc-cd=0\) implies \(c_p=C_0\) identity on the exterior
factor. To see this start with \(K_4\), which has exterior rank one, then
compare the coefficients of the four independent quadrics successively
on \(K_4\),\(K_3\),\(K_2\). Contractions of \(K_4\) span \(K_3\), of \(K_3\) span \(K_2\), and of \(K_2\) span \(K_1\).
The diagonal equations also give \(\beta C_0=C_0\beta=0\). Since some \(b_{ij}\) is
nonzero, \(C_0=0\). Thus the entire degree -3 correction vanishes.

The homogeneous part \(H\) of degree -2 has the form

\[
 H=\begin{pmatrix}a&0\\C&e\end{pmatrix},
 \tag{HK.5}
\]

where \(a,e\) have constant coefficients and exterior degree -1, while \(C\)
is a linear scalar acting identically on each exterior degree. Indeed
the same chain argument now compares linear multiples of the \(f_i\).
There are no nonzero linear syzygies among a regular sequence of quadrics:
its Koszul first syzygies start in total degree four. This forces the
lower-left operators to be \(C\) times exterior identity. Upper-right degree
-2 corrections would need a coefficient of degree -1 and hence are zero.

Put \(P_{ij}=\{f_i,f_j\}_\pi\). The diagonal equations are

\[
 \{d,a\}=\iota_{P-Cb},\qquad
 \{d,e\}=-\iota_{P-Cb}.
 \tag{HK.6}
\]

In particular there are unique constant two-forms \(A_l\) on \(U\) such that
\(P-Cb=\sum_l f_l A_l\). Define \(h_A=\sum_l(e_l\wedge)\iota_{A_l}\). Then
\(\{d,h_A\}=\iota_{P-Cb}\). All constant solutions of the first equation in (HK.6)
are \(h_A+\iota_v\), for \(v\) in \(U^*\). The second equation gives \(-h_A+\iota_w\) on its
truncated domain \(K_2\),\(K_3\),\(K_4\). Here is the complete uniqueness argument.
After subtracting \(h_A\), comparison of the coefficients \(f_i\) says the
operator anticommutes with every \(\iota_i\). Its value on \(K_4\) determines a
unique contraction \(\iota_v\). After subtracting that contraction, the
equations on \(K_4\) force the operator to vanish on \(K_3\), those on \(K_3\) force
vanishing on \(K_2\), and those on \(K_2\) force vanishing on \(K_1\). For the bottom
truncation only \(K_4\) and \(K_3\) are needed to determine \(K_3\) and \(K_2\); no equation
on an absent \(K_1\to K_0\) bottom map is used. Thus

\[
 a=h_A+\iota_v,\qquad e=-h_A+\iota_w.
 \tag{HK.7}
\]

For reference the remaining first-order equation is the genuine
off-diagonal identity

\[
 G_\pi+[h_A,\beta]+\iota_{v+w}\beta=0,
 \qquad
 G_\pi=\sum\pi^{ij}(\partial_i d)(\partial_j\beta).
 \tag{HK.8}
\]

It is not replaced by scalar cross minors. No existence of a solution
of (HK.8) for a previously unexcluded two-direction example is claimed.

\subsection{3. A nonzero second-order obstruction of internal degree -4}\label{SF-E846:s3}

First \(C\) is nonzero. If \(C=0\), (HK.6) says every \(P_{ij}\) vanishes in \(A\). The
Pfaffian identity for the nondegenerate \(\pi\) says \(\operatorname{Pf}(P)\) is a nonzero
constant times the Jacobian determinant of \(f\). The latter is nonzero
in \(A\) by the characteristic-zero complete-intersection socle theorem,
contradiction. This is the already checked Eisenbud--Ulrich Theorem 7.1,
p.~210, with full Appendix A.1/A.3/A.5 proofs, pp.~222--225.\cite{SFEisenbudUlrich2019}

Take the degree -2 solution \(H\) furnished by a first-order lift.
The coefficient of \(h^2\) in (\(D+hH\)) star (\(D+hH\)) is a cocycle \(O\) of
cohomological degree two and internal degree -4. The lower-left block
is particularly simple:

\[
 O_{21}=-\iota_{\{f,C\}_\pi}+C\iota_{v+w}.
 \tag{HK.9}
\]

Indeed the second-order term involving \(D\) twice has zero lower-left
block. The two first-order star terms with \(D\),\(H\) give
\(-\sum\pi^{ij}(\partial_i d)(\partial_j C)\), while the ordinary square \(H^2\)
has lower-left block \(Ca+eC=C\iota_{v+w}\). This proves (HK.9) with the
right-module convention of (HK.4).

Here is the cocycle verification. Associativity gives
(\(D+hH\)) star ((\(D+hH\)) star (\(D+hH\))) equal to
((\(D+hH\)) star (\(D+hH\))) star (\(D+hH\)). The constant and first-order
curvatures are zero, so the coefficient of \(h^2\) is \(DO-OD=0\). The
constant-Poisson product (HK.4) is associative through this order:
\(\exp((h/2)\pi^{ij}\partial_i\otimes\partial_j)\) has commuting differential
operators, and the Leibniz rule makes the two triple products equal.
No global quantized category is being asserted.

There is no endomorphism of cohomological degree one and internal
degree -4: its four block coefficient degrees would be -2,-3,-1,-2.
Consequently any nonzero \(O\) represents a nonzero class in
\(\operatorname{Ext}^2_R(M_b,M_b)\) of internal degree -4.

If \(O\) vanished, (HK.9), evaluated on \(K_2\to K_1\), would give

\[
 \{f_i,C\}_\pi=z_i C\qquad (z_i\in\mathbf C).
 \tag{HK.10}
\]

Because \(\pi\) is nondegenerate and \(C\) is a nonzero linear form, its
Hamiltonian vector field is a nonzero constant vector \(v_C\) with
\(v_C(C)=0\). Choose linear coordinates with \(C=x\) and \(v_C=\partial_y\).
Equation (HK.10) says every homogeneous quadric \(f_i\) has \(y\)-dependence
only in the monomial \(xy\). All four quadrics therefore vanish at the
nonzero projective point on the \(y\)-axis, contradicting regularity.
Hence \(O\) is nonzero. This argument does not require (and does not
assert) the existence of any second-order deformation.

Other internal degrees of a first-order correction cannot affect
the weight -4 equation: degree -3 is zero, all remaining degrees
are at least -2, and the coefficient of the first Moyal correction
lowers internal degree by two. The second-order correction has no
degree -4 part. This exhausts all corrections for the exactly
homogeneous module in (HK.2).

\subsection{4. The ordinary local-to-global injection}\label{SF-E846:s4}

Let \(E\) be a coherent sheaf on a smooth proper fourfold \(X\) with locally
free hull \(H\) and finite kernel and cokernel of \(\phi:E\to H\). Put
\(\mathcal C=\operatorname{cone}(\phi)\), in the distinguished triangle

\[
 E\xrightarrow\phi H\longrightarrow\mathcal C
      \xrightarrow\delta E[1].
 \tag{HK.11}
\]

All cohomology of \(\mathcal C\) is supported on a finite set \(Z\). Since \(E\)
is perfect, \(R\operatorname{Hom}(E,\mathcal C)\) also has finite coherent cohomology
supported on \(Z\). Its support spectral sequence has no rows of
positive sheaf cohomology. Restriction and completion consequently
give canonical isomorphisms

\[
 \operatorname{Ext}^j_X(E,\mathcal C)
   \simeq\bigoplus_{p\in Z}
       \operatorname{Ext}^j_{\widehat{\mathcal O}_{X,p}}
            (M_p,C_p).
 \tag{HK.12}
\]

For any \(M_p\) with a free hull and finite kernel/cokernel, the two
short exact sequences through the hull image, together with
\(\operatorname{Ext}^j_R(\text{finite length},R)=0\) for \(j<4\), give
\(\operatorname{Ext}^1_R(M_p,R^r)=\operatorname{Ext}^2_R(M_p,R^r)=0\). Therefore the local triangle gives
\(\operatorname{Ext}^1_R(M_p,C_p)\simeq\operatorname{Ext}^2_R(M_p,M_p)\).

The actual preceding global arrow \(\operatorname{Ext}^1_X(E,H)\to\operatorname{Ext}^1_X(E,\mathcal C)\)
is zero. This is not inferred from a spectral-sequence edge alone.
The commutative square has top row that arrow, bottom row
the map from the direct sum of \(\operatorname{Ext}^1_R(M_p,R^r)\) to the direct sum
of \(\operatorname{Ext}^1_R(M_p,C_p)\), and vertical arrows restriction followed by
completion. The bottom source is zero, while the right vertical
arrow is the isomorphism (HK.12). Hence the top arrow is zero.
The global triangle then gives
a canonical injection

\[
 j:\bigoplus_{p\in Z}\operatorname{Ext}^2_R(M_p,M_p)
       \hookrightarrow\operatorname{Ext}^2_X(E,E).
 \tag{HK.13}
\]

It is the connecting map after the identifications just made and is
a section of ordinary restriction to the local Ext groups. It is
natural for isomorphisms of the entire hull triangle, in particular
for every specified coherent translation/tensor action.

\subsection{5. The trace-weight lemma, including the canonical volume}\label{SF-E846:s5}

Suppose \(X\) is an abelian fourfold, and at \(p\) the hull map is the graded
map \(M_b\to R\) of section 1, with finite direct-sum multiplicities and
possibly a free summand. Let \(u\) be a homogeneous local \(\operatorname{Ext}^2\) class of
degree \(d\), and let \(a\in\operatorname{Ext}^1_R(M_b,C_p)\) be its preimage. The connecting
map has internal degree zero. For a global degree-two Hochschild
class \(\gamma\), naturality of composition gives

\[
 \operatorname{Tr}_E\bigl(j(u)\,\rho_E(\gamma)\bigr)
   =\operatorname{Tr}_E\bigl(\delta\,[a\,\rho_E(\gamma)]\bigr),
 \tag{HK.14}
\]

up to the single harmless cyclic sign. The bracketed term belongs
to \(\operatorname{Ext}^3_X(E,\mathcal C)\), so (HK.12) makes it depend only on the actual
completed restriction of \(\rho_E(\gamma)\). In particular the \(H^1(T_X)\)
and \(H^2(\mathcal O_X)\) HKR components restrict to zero; their possible global
hypercohomology representatives contribute no extra boundary term
to this equality in the derived category. The remaining constant
bivector evaluation is homogeneous of internal degree -2.

Fix a nonzero invariant volume \(\Omega\) on \(X\). It identifies the
Calabi--Yau trace used on the left of (HK.14) with the Serre trace
after tensoring by \(\Omega\). Rotate this latter trace from \(E\) to the
finite-supported perfect object \(\mathcal C\). This is the usual cyclic
identity for evaluation and coevaluation of bounded locally free
complexes; its sign depends only on the two cohomological degrees.
Thus the result is, up to sign, the Serre trace on \(\mathcal C\) of
the composition of \(a\rho\) with \(\delta\otimes\Omega\).

The finite-support comparison can be made directly at the complex
level. The complex \(R\operatorname{Hom}(C_p,C_p\otimes\omega_R)\) has finite-supported
cohomology, so applying \(R\Gamma_{\mathfrak m}\) to it is an isomorphism. Apply \(R\Gamma_{\mathfrak m}\) to
its perfect-complex evaluation/supertrace map to \(\omega_R\). On degree
four this gives

\[
 \operatorname{Ext}^4_R(C_p,C_p\otimes\omega_R)
   \longrightarrow H^4_{\mathfrak m}(\omega_R)
   \xrightarrow{\operatorname{Res}_p}\mathbf C.
 \tag{HK.15}
\]

For the corresponding global perfect complex the same evaluation
map, followed by forgetting supports, lands in \(H^4(X,\omega_X)\).
These maps commute with restriction and completion. Lipman's
Theorem 9.1 and Corollary 9.2 identify the composite with the global
Serre trace with the last residue in (HK.15): \emph{Dualizing sheaves,
differentials and residues on algebraic varieties}, Astérisque 117
(1984), pp.~82--85.\cite{SFLipman1984} The full invoked proofs, including the comparison
diagrams, were read in this run. The trace/evaluation interpretation
for perfect complexes is also Stacks, Section 48.27 tag 0FVU,
Lemma 48.27.1(5),(7), Remarks 48.27.2--3 and Lemma 48.27.4;\cite{SFStacksE846Duality} its use
here follows by taking the bounded locally free representatives
in the same evaluation diagram, not by extending a scalar formula
without a map of complexes.

For precision, \(\omega_R=R\,dx_1\wedge\cdots\wedge dx_4\). The local-cohomology
group is the span of inverse monomials \(x_1^{-a_1}\cdots x_4^{-a_4}\,dx_1\wedge\cdots\wedge dx_4\),
all \(a_i\ge1\). The residue extracts the coefficient with every \(a_i=1\).
This description is independent of regular formal coordinates;
Lipman §7, pp.~64--69, Lemma 7.3 and Theorem 7.4 give the formula and
its complete coordinate-independence and local-duality proofs.
Consequently the residue is equivariant for simultaneous dilation
of the coordinates and their differentials: it vanishes on every
homogeneous element of nonzero total internal degree. This can
also be checked term by term in the displayed inverse-monomial basis.

The restriction of \(\Omega\) in invariant formal-logarithm coordinates
is a nonzero constant multiple of \(dx_1\wedge\cdots\wedge dx_4\), hence
has degree +4. The connecting map \(\delta\) has degree zero with the
graded hull map of §1. For \(u\) of degree \(d=-4\), the first factor \(a\rho\)
has degree -6, while \(\delta\otimes\Omega\) has degree +4. The value
in (HK.15) therefore has total degree -2, so its residue is zero.
The same argument works for any \(d\) with \(d+2\) nonzero.

Thus the image \(j(u)\) pairs to zero with every ordinary Hochschild
evaluation in degree two. The fourfold Serre trace/evaluation
adjunction supplied and checked in DEG\_L S20/S21 identifies this
annihilator with the kernel of the full semiregularity map \(\sigma^2\).
This invokes precisely the adjunction, as well as the nondegenerate
Hochschild pairing on an abelian fourfold; the first-order identity
\(\sigma\rho=\operatorname{contr}(\operatorname{ch})\) alone would not suffice.

For multiple copies of \(M_b\), take \(u\) tensor the identity on the
multiplicity space and zero on a free summand. Naturality of
Hochschild evaluation makes its restriction block diagonal under
any actual module isomorphism. The trace on these copies is the
sum of the same zero pairings. The free summand's hull cone is
acyclic, so it contributes nothing to (HK.15).

\subsection{6. Actual finite actions, with the stabilizer and scalar cocycle}\label{SF-E846:s6}

Assume the completed local module is actually isomorphic to
\(M_b\otimes\mathbb C^m\) direct-sum \(R^s\), with \(m>0\). This is an isomorphism of
modules, not merely of their associated graded modules or \(K\)-classes.
The hull map is functorial under it, since double duals are functorial.
We prove the existence of an invariant nonzero class without
assuming the original action preserves a grading or chosen resolution.

First consider only the point stabilizer \(G_p\) in the specified finite
translation/tensor action. If a translation \(t_x\) fixes \(p\), then \(x=0\).
Therefore \(G_p\) acts trivially on the completed base ring; the action
at this stage is \(R\)-linear, not semilinear. Its Picard line bundles
are torsion and have flat local trivializations. Coherent global
product isomorphisms between these flat line bundles are constant
in such trivializations: their ratios are global holomorphic
functions on the connected proper abelian variety. Hence the
local operators satisfy \(gh=\alpha(g,h)(gh)\) for one constant scalar
cocycle \(\alpha\) in \(Z^2(G_p,\mathbb C^*)\). Scalar ambiguities disappear on Ext.

The twisted group algebra \(\mathbb C_\alpha[G_p]\) is semisimple. Indeed if a
subspace has a linear projection onto it, averaging its conjugates
over \(G_p\) makes an invariant projection; the scalar factors in
conjugation cancel. Its finite-dimensional simple representations
\(S_\lambda\) therefore give an \(R\)-linear isotypic decomposition

\[
 M_p\simeq\bigoplus_\lambda S_\lambda\otimes N_\lambda,
 \qquad N_\lambda=\operatorname{Hom}_{\mathbf C_\alpha[G_p]}
                    (S_\lambda,M_p).
 \tag{HK.16}
\]

The modules \(N_\lambda\) are finite over \(R\), since they are images of
idempotents on finite direct sums of \(M_p\). The action on the right
is on \(S_\lambda\) alone.

We spell out the required local decomposition fact. Every
endomorphism of \(M_b\) induces multiplication by a scalar \(a\in R\) on
its rank-one double dual and multiplication by \(t\) in \(A\) on its torsion
\(A(-3)\). The extension-class condition is
\((t-a\bmod I)b_{ij}=0\) in \(A\) for all \(i\),\(j\). This follows by applying the
two induced endomorphisms to the actual class in \(\operatorname{Ext}^1_R(I,A)\),
which the Koszul resolution identifies with \(A\otimes\Lambda^2U^*\).
Since one \(b_{ij}\) is a nonzero linear class, this condition forces
\(t(0)=a(0)\). If \(a(0)\) is nonzero, both the induced torsion and hull-image
maps are invertible, hence so is the original endomorphism by the
short exact sequence (HK.2). If \(a(0)=0\), its difference from the
identity is invertible by the same argument. Thus \(\operatorname{End}_R(M_b)\) is
local, with residue \(\mathbb C\), and \(M_b\) is indecomposable. It is not \(R\):
its torsion \(A\) is nonzero. The ring \(\operatorname{End}_R(R)=R\) is local as well.

For completeness, the Krull--Schmidt assertion used here holds
for finite modules over a complete Noetherian local ring. Their
endomorphism algebras are finite and complete over \(R\). Modulo the
maximal ideal they are finite-dimensional algebras; primitive
idempotents lift through their nilpotent radicals and then through
the powers of the maximal ideal. At the complete step, the usual
Newton correction to \(e^2-e\) gives a convergent lift, since \(2e-1\) is
invertible near an idempotent and commutes with every polynomial
in \(e\). This gives a finite decomposition into summands with local
endomorphism rings. Uniqueness follows by writing the identity
of any indecomposable summand as the sum of its composites through
another decomposition. At least one composite is a unit in its
local endomorphism ring. It identifies that summand with one on
the other side; split it off and repeat. In particular every direct
summand of \(M_b^m\) direct-sum \(R^s\) is a sum of copies of \(M_b\) and \(R\).

Apply this fact to all \(N_\lambda\) in (HK.16). There is an actual
\(G_p\)-equivariant isomorphism

\[
 M_p\simeq(M_b\otimes V)\oplus(R\otimes W),
 \tag{HK.17}
\]

where \(V\) and \(W\) carry projective representations with the same
cocycle \(\alpha\), and \(\dim V=m>0\). This isomorphism need not preserve
the original grading; the homogeneous grading is used on its
\(M_b\) factor after the isomorphism, and all characteristic maps
and pairings are transported naturally. Thus \(u\otimes\operatorname{id}_V\)
is a nonzero \(G_p\)-invariant local \(\operatorname{Ext}^2\) class, and §5 proves that
its global image has zero semiregularity.

Finally translate this class over the \(G\)-orbit of \(p\). Choose one
representative for each coset of \(G_p\) and transport by the actual
coherent functor. The result is independent of the representatives,
because the starting class is \(G_p\)-invariant and the action on Ext
is honest. Its direct sum over the orbit is nonzero and \(G\)-invariant.
The injection (HK.13) is natural for the hull triangle, so the
global image remains nonzero and invariant. Each translated
trace pairing is zero by naturality, or by the same formal-log
calculation at that point. Hence it lies in \(\ker\sigma^2\) on invariant Ext.

\subsection{7. The balanced-character consequence and its exact hypotheses}\label{SF-E846:s7}

Let \(X\) be an abelian fourfold of genuine real-multiplication type
(2,2), with polarization \(\theta=\theta_++\theta_-\). Let \(E\) be a coherent
sheaf whose canonical hull map has locally free hull and finite
kernel/cokernel. Suppose at one point its completed module, in
invariant formal-logarithm coordinates, has the exact type
\(M_b^m\) direct-sum \(R^s\) of §6. Suppose moreover

\[
 e^B\operatorname{ch}(E)=r v_q,
 \qquad
 v_q=(1-q\theta_+^2/2)(1-q\theta_-^2/2),
 \qquad r>0,\quad q>0,
 \tag{HK.18}
\]

where \(B\) is any rational algebraic divisor. Then \(E\) is not weakly
\(G\)-semiregular for any specified effective coherent finite
translation/tensor action as in supplied S21.

Suppose to the contrary that it is. The two cap-kernel classes
\(\xi_\pm=\pi_\pm+q\eta_\pm\) of E819§4, conjugated by the full \(B\)-field shear,
retain complementary constant bivector projections \(\pi_\pm\). Their
sum \(\xi\) has nondegenerate projection \(\pi=\pi_++\pi_-\). Evaluation lies
in invariant Ext by supplied S21, and the published
Buchweitz--Flenner composition identity gives \(\sigma^2\rho^2(\xi)=0\).
Injectivity of \(\sigma^2\) on that invariant Ext forces \(\rho^2(\xi)=0\).
Restrict and complete. The \(H^1(T)\) and \(H^2(\mathcal O)\) components vanish,
so the local module has a flat first-order deformation for \(\pi\),
by the explicit bar/free-resolution argument of E819§3.

Under any actual module splitting (HK.17), the characteristic
class restricts diagonally to each \(M_b\) copy. Thus \(M_b\) itself has
zero first-order characteristic class and a first-order lift.
Sections 2--3 construct a nonzero local class \(u\) of degree -4.
Sections 4--6 give a nonzero invariant global \(\operatorname{Ext}^2\) class in \(\ker\sigma^2\),
contradiction. The second-order local obstruction was used only
as an explicit cocycle construction in ordinary Ext. At no point
is a second-order noncommutative lift deduced from semiregularity.

The published first-order input is Buchweitz--Flenner,
\emph{The global decomposition theorem for Hochschild (co-)homology of
singular spaces via the Atiyah--Chern character}, Adv. Math. 217 (2008),
Theorem 6.4.2, with Proposition 6.2.1, Proposition 6.3.1, Corollary 6.3.2
and the proof through Proposition 6.4.4 (author-copy pp.~29--33).\cite{SF-BF2008}
The invariant-evaluation and CY4-adjunction interfaces are the
explicit supplied proved inputs DEG\_L S20/S21: S20's first proof
paragraph (the Serre-transpose identity), S21.6 and the paragraphs
immediately preceding S21.8. Those exact passages were reread in
this run; no publication status is added to the supplied results.

\subsection{8. Actual global objects in the stated local scope}\label{SF-E846:s8}

The local type is an actual coherent-sheaf type, not only a
formal character. For any regular homogeneous four-quadric
sequence the Hilbert series of \(A\) is \((1+t)^4\), so \(\mathfrak m^5\) is contained
in \(I\). The quotient \(R\to A\) therefore factors through \(R/\mathfrak m^5\). In
invariant formal-logarithm coordinates at a point of an abelian
fourfold, the canonical isomorphism of finite jets with the
algebraic local ring realizes it as a closed algebraic subscheme
\(Z\) of length 16. Its completed ideal is exactly \(I\).

There is a global extension
\(0\to\mathcal O_Z\to M\to\mathcal I_Z\to0\) with any prescribed local class \(b\). Indeed
\(R\operatorname{Hom}(\mathcal I_Z,\mathcal O_Z)\) has finite-supported cohomology, and its support
spectral sequence collapses, so global \(\operatorname{Ext}^1\) identifies with the
completed local \(\operatorname{Ext}^1\). The Koszul calculation identifies the latter
with \(A\otimes\Lambda^2U^*\), in which the six specified linear classes
are legitimate coordinates. The resulting sheaf has local module
\(M_b\), canonical hull \(\mathcal O_X\), and identical length 16 kernel and
cokernel. Additivity gives \(\operatorname{ch}(M)=\operatorname{ch}(\mathcal O_X)=1\).

For a concrete balanced ambient character use the actual rank 32
polystable bundle \(H\) of E815/E820 on the CM RM fourfold \(E_i^4\).
Its decomposition into actual line bundles has multiplicities
\(24,2,2,1,1,1,1\) and \(\operatorname{ch}(H)=32v_{1/2}\). Distinct line types have no
Hom between them: their normalized differences are the primitive
index-two classes checked in E813/E815. Thus every global hull
endomorphism is a constant block matrix on the multiplicity
spaces, with algebra \(\operatorname{Mat}_{24}(\mathbb C)\) direct-sum \(\operatorname{Mat}_2(\mathbb C)\) direct-sum
\(\operatorname{Mat}_2(\mathbb C)\) direct-sum \(\mathbb C^4\).

Choose 33 distinct points \(p_1,\ldots,p_{33}\). At each choose the same
quadratic CI in invariant logarithm coordinates and its length 16
algebraic subscheme \(Z_j\). In local frames of the 32 line summands,
define a surjection \(H\to Q=\bigoplus_j\mathcal O_{Z_j}\). At \(p_j\) for \(j\le32\)
use the row \(e_j\); at \(p_{33}\) use the row \((1,\ldots,1)\). Each is a surjection,
and arbitrary such jets define actual morphisms to the finite
sheaf \(Q\). Let \(K\) be its kernel and \(T=Q\). At each point the hull
sequence has \(K_p\simeq I\) direct-sum \(R^{31}\). The support spectral sequence
again identifies \(\operatorname{Ext}^1_X(K,T)\) with the direct sum of its completed
local groups. Choose the nonzero scalar linear class \(b\) on the
\(I\) factor at every point. All \(\operatorname{Ext}^1\) from the free local factors
vanish, so this constructs an actual coherent extension

\[
 0\longrightarrow T\longrightarrow E\longrightarrow K
       \longrightarrow0,
 \qquad E_p\simeq M_b\oplus R^{31}\quad(p\in\{p_j\}).
 \tag{HK.19}
\]

Its maximal torsion is \(T\), its double dual is \(H\), and its hull
cokernel is \(Q\). Hence its full character is \(\operatorname{ch}(H)\). It is an
ordinary gluable object because it is a coherent sheaf.

This \(E\) is indecomposable. Let \(e\) be an idempotent endomorphism.
Its induced hull endomorphism \(a\) preserves \(K\) and hence induces
an endomorphism on \(Q\). At \(p_j\), \(j\le32\), compatibility with the
quotient row \(e_j\) forces the \(j\)-th row of the constant block
matrix \(a\) to be a scalar multiple of \(e_j\). Thus \(a\) is diagonal.
The last row forces all 32 diagonal entries to agree, so
\(a=\lambda\operatorname{id}\). Since \(e\) is idempotent, \(\lambda\) is 0 or 1.
If \(\lambda=0\), the induced endomorphism of each torsion algebra \(A\)
is multiplication by an idempotent \(t_j\). Compatibility with \(b\)
gives \(t_jb=0\), so \(t_j\) has zero residue, as in §6. A local Artin
algebra has no nonzero idempotent of zero residue, hence \(t_j=0\).
Now \(e\) has image in \(T\) and vanishes on \(T\), so \(e^2=0\) and therefore \(e=0\).
If \(\lambda=1\), apply this argument to \(1-e\) to get \(e=1\). There is no
nontrivial idempotent, which proves indecomposability.

The identity action is an actual finite coherent action on these
examples; no extra action is presumed to exist. The theorem
excludes weak semiregularity for it and for every further specified
action if one exists. This is actual-family evidence, not a strict
comparison with all earlier local Artin criteria. No particular
quadratic CI surviving those earlier criteria is asserted here.

\subsection{9. Current boundary and certificate scope}\label{SF-E846:s9}

There is no claim here for arbitrary filtered quadratic jets, arbitrary
matrix-valued \(b\), nonfree Artin defects, or all complete intersections.
No second-order noncommutative semiregularity theorem is imported.
The failed broad bridge through Perry remains separate: his 2026
preprint Theorem 4.1 assumes geometric origin. Pridham's published
2024 cyclic argument has explicit additional hypotheses in section 1.5
and Remark 1.21; these have not been verified for the present Moyal-plus-
gerbe deformation, and are not used in the proposed ordinary trace route.

The entire local calculation, support/trace interfaces and finite-action
step have passed independent whole-proof review. A strict scope comparison
with all previous Artin exclusions has not been established; therefore no
independent main-gate movement is claimed.

The exact author checker rm\_homogeneous\_koszul\_certificate.py
checks the complete finite exterior operator spaces: the lower
chain space has dimension 1 (69 variables, rank 68), the top
anticommutant dimension 4 (56 variables, rank 52), and the truncated
bottom anticommutant dimension 4 (52 variables, rank 48). It also
checks the degree bound, the Hamiltonian-axis implication, and
the entire 31-by-31 second-order matrix identity at three rational
inputs, retaining the first-order Bianchi remainder. The first
actual run passed 34 checks on 2026-11-22 at 02:13:46--02:13:51UTC,
elapsed 5.900344789 seconds. These arithmetic checks neither produce
a first-order existence witness nor verify the global argument.

\section{Placement in the Grothendieck programme}

\subsection{Cycle conditions and the corrected critical map}
The new detectors provide integral witnesses in the cycle spaces that
enter condition A; the uniform bound supplies their finite-type
parameter space. This is stronger data than the bare assertion that
Lefschetz maps preserve rational algebraic classes. We retain the
correction on cote16, scan050 (archive folio49), in SIGN\_NEXT
SN-C03/C04: on a $2m$-fold $T$, the critical cycle map is $L_T^2$, and
its hyperplane factorization is $j_*L_Wj^*$. The weighted Gysin
condition $L_Tj_*$ cannot be replaced by ordinary lower-critical Gysin.
The historical $C_{\rm G}$ on that page is the algebraicity of the
Lefschetz inverse (modern B), not the modern Kunneth condition.
These distinctions apply to every detector and transfer in this paper;
no repaired handwritten condition is used as a proof of the Hodge
conjecture. The exact modern weak-Lefschetz locator is SGA5, ExposeVII,
Theorem7.1 and its Poincare-dual form; the project correction is the
cited SIGN\_NEXT pair.

The norm bound and the positive Gram matrices in HD-WQ, HD-G05 and
SF-06/07 are the use of condition I on cote16 scan049, and of cote44,
scan002/footer1, Theorem4. They invert a nondegenerate pairing once
geometric detectors exist. The resultant projector in the supplied
Theorem W is the polynomial, Cayley--Hamilton inversion mechanism of
cote16 scans048--049. Our shorter Gaussian and Eisenstein projectors
are explicit integral character-polynomial instances of that mechanism.
The extra cancellation in HD-WQ is a new arithmetic bound on the
multiplier, not a new source of cycles.

\subsection{Schemes over a finitely generated ring over the integers}
All fixed curves, finite covers, group actions, polarizations,
correspondences and effective cycles in a finite portion of the
construction descend to a finitely generated ring over $\Z$ and extend
after shrinking it. This is finite-presentation descent, EGA IV$_3$,
8.8.2; it is not a claim of a smooth model over all of $\operatorname{Spec}\Z$.
For the tame cover constructions invert the group order and the
finitely many level, order and denominator integers involved. The
finite-cover comparison has its standard complex-algebraic placement
in SGA1, ExposeXII, Theorem5.1. Bounded degrees lead to finitely many
cycle parameters as in the accepted HECKE\_DIV R21 and attack R1/D.
The FGA locators used by those proofs are Bourbaki expose221,
Lemme2.4, p.255, and Theoremes3.1--3.2, pp.259--260; one must retain
the bounded-degree argument, rather than cite only a single
fixed-Hilbert-polynomial scheme.

The kernel computations and conductor changes are actual homomorphisms
of abelian schemes after such a base change. Divisor translation,
proper intersection, pullback by a finite etale isogeny and proper
pushforward preserve integral multiplicities in the ways explicitly
proved above. The constants refer to the stated polarization; they
are not assertions independent of an unbounded conductor or lattice.

\subsection{Lambda-rings, K-theory, and the sheaf alternative}
The representation identities for finite groups are identities in
their integral representation rings; exterior powers give the
determinant characters. Their $h^1$ realizations and rational group
projectors are algebraic correspondences. This observation by itself
does not create a middle-dimensional effective cycle. The supplier
and contraction proofs do that additional work.

For a perfect complex the Chern character and exterior operations live
in the K-theoretic intersection framework of SGA6, Expose0, section1,
equations(1.1)--(1.3), p.1, and ExposeV, 2.1--2.2, especially
Lemma2.2.1. Newton and Adams operations have the exact further
locators ExposeV, Appendix7.1, pp.360--361, and Proposition7.5,
pp.362--364. The semiregularity gate of DEG\_L S-14 asks for an actual
perfect complex, a coherent finite translation/tensor action, the
entire equation $e^B\operatorname{ch}(E)=\lambda t+P(\theta)$, and
ordinary negative self-Ext vanishing. A virtual K-class, one Chern
component, or a formal exterior identity does not supply those
hypotheses. Our separate checks keep each hypothesis visible.

\subsection{Etale realizations, weights, and motives}
Smooth proper specialization is SGA4, ExposeXVI, Corollary2.2.
The rational cycle-class compatibility used by the supplied arithmetic
reduction is SGA6, ExposeX, Appendix section7, equation(7.13.10),
pp.578--579. Trace and correspondence operations have the exact
locators SGA5, ExposeIII, section3 and 5.1--5.3. Thus, after the
specified spreading and shrinking, the finite group projectors,
quotient maps and isogeny comparisons act on the lisse sheaves
$R^1f_*\Q_\ell$ and their exterior powers. The Weil summand belongs
to $R^{2n}f_*\Q_\ell(n)$; its untwisted weight is $2n$ and its Tate
twist has weight zero. These are realization statements, not a
deduction of algebraicity from a Frobenius eigenvalue.

In homological motives the rational Weil projector cuts out the
two-dimensional determinant summand. A detecting cycle and its
$K$-endomorphism translate give two independent Tate morphisms into
that summand, with the integral denominators controlled by HD-WQ.
The bounded degree is extra information on representatives and is
not contained in the isomorphism class of the rational motive.
For data over a number field, once these two cycles are defined over
a finite extension, the corresponding good local factor on the
untwisted summand is $(1-q^n u)^{-2}$; before making that extension
there can be a finite Artin action on its Tate twist. This follows
from the fixed algebraic classes and does not estimate any Rallis
value, conductor, Chow degree, or multiplier from an L-function.
No RH or GRH assertion is used. In particular DEG\_L B1's Rallis
point $s=3/2$ remains distinct from a geometric detector estimate.

\subsection{The remaining ghost statement}
The source exclusions D-BI-01--24, D-SPOR, D-SPW, D-G2T,
D-ESC, D-A2, L34, D-HES-F, D-HES-01--06, D-C3-01--03, D-A7-01--03, D-HOS-01--06, D-UNW, D-LNW, D-OC2, D-OC2-F, D-EXS, D-OSM, HD-SPIN, HD-ALT, D-AST, D-EX-01--03, D-UP-01--03,
D-R5-01--04 and HD-E8IF1 concern finite-monodromy coefficient
systems and the dimension of their weight-one period maps. Their
algebraic-geometric objects are the finite covers and rational group
projectors described above; their integral representation-ring and
exterior-power identities belong to the stated lambda-ring framework.
Their semistable graph arguments identify a square-zero monodromy
image in the original coefficient space, and the maximal isotropic
image excludes a nonsplit polarization. The explicit scalar
ramification arguments use finite etale cyclotomic algebras after
inverting the conductor and contrast them with ramification of the
specified coefficient field. These restrictions apply to potential
suppliers for condition A at cote16, scan050, and do not produce
cycles by positivity or extract a bound from a motivic L-function.

SF-US, SF-UG and SF-E811--833 concern actual bundles, coherent sheaves and perfect complexes on a
fourfold carrying genuine real multiplication. The Chern-character
identities are identities in rational cohomology induced by actual
K-classes, in the SGA6 framework specified above. Equivariant
Grothendieck-group additivity, trace and evaluation have the exact
project locator DEG\_L S-21, equations(S21.6)--(S21.7), while the
deformation-input requirements are DEG\_L S-14. The new obstructions
restrict those possible inputs; they neither prove a general negative
algebraicity statement nor identify a ghost. SF-E814 gives an actual
derived example with a different, two-pure character and retains that
scope. Spreading any fixed algebraic construction uses EGA IV$_3$,
8.8.2 as above; the analytic Hermitian--Einstein existence theorem is
not asserted to be a theorem of EGA or SGA. Once an algebraic cycle
exists its etale realization and weights have the stated SGA4/SGA5
placement, but no new cycle or L-function conclusion follows merely
from a Chern--Weil inequality or a self-Ext dimension.

The new bounds close the ghost gate for the Weil planes in the
specified families. Failure of the cyclic or imprimitive construction
is only failure of that source. A counterexample to algebraicity
would give the ghost conclusion only under the exact arithmetic and
CM hypotheses of GHOSTS B02/B03 and HECKE\_DIV P6, with 032 still
conditional as required. We do not use 032 to manufacture the
effective cycles in this paper. Likewise the unbounded core index
in the supplied Theorem E is compatible with our bounded detectors:
fixed, explicitly controlled isogenies in a proof are not a uniform
estimate along an unbounded Hecke transport problem.


\part{Conditional deductions and limitations of alternative routes}
% Checked auxiliary and conditional CM results. See cm_transport.md and CHECK_EVENTS.tsv.
\section{Auxiliary CM transport and a conditional secant construction}
\label{sec:cm-transport}

The linear and descent results in this section are independent of a
conjectural cycle construction. A separately labelled theorem uses
Markman's unrefereed source sheaf conditionally. The explicit
rank-zero construction elsewhere in this paper supersedes that
conditional route for the balanced Gaussian CM example; the conditional
route is retained with its precise scope and its unsuccessful boundary
specialization recorded.

\begin{proposition}[SF-CM01: semiregularity under a fixed isogeny]
\label{SF-CM01}
Let $f:A\to Y$ be an isogeny of complex abelian varieties and let
$\overline F$ be an ordinary semiregular coherent sheaf on $Y$.
Put $G=\ker f$ and $F=f^*\overline F$. Then $F$ is a perfect
coherent sheaf with no negative self-Ext and is $G$-semiregular for
the canonical coherent action of the actual finite subgroup
$G\subset A\times\{0\}$ by translations. If its rank is $r\ne0$,
then
\[
 e^{-c_1(F)/r}\operatorname{ch}(F)
 =f^*\bigl(e^{-c_1(\overline F)/r}\operatorname{ch}(\overline F)\bigr).
\]
\end{proposition}
\begin{proof}
The isogeny is finite etale. Its translations satisfy $f\circ t_g=f$,
so functorial pullback gives canonical isomorphisms $t_g^*F\simeq F$
satisfying the cocycle identity. This is an honest coherent action,
with no projective ambiguity or invisible central kernel. Finite
etale descent, and exactness of taking invariants in characteristic
zero, give
\[
 \operatorname{Ext}^2_A(F,F)^G
 \simeq\operatorname{Ext}^2_Y(\overline F,\overline F).
\]
The Atiyah class and its trace commute with etale pullback. Thus the
full semiregularity map restricted to this invariant subspace is
$f^*\circ\sigma_{\overline F}$. On every Hodge-cohomology summand
of $\bigoplus_qH^{q+2}(Y,\Omega_Y^q)$ the pullback is injective,
since $f_*f^*=(\deg f)\operatorname{id}$. The restricted map is
therefore injective. This is strong $G$-semiregularity, and hence
weak $G$-semiregularity in the convention of DEG\_L S-14; see also
\cite[Definition~2.6 and Lemma~2.7, pp.~5--6]{SF-Perry2604}.
Every coherent sheaf on the smooth $A$ is perfect, and negative
self-Ext vanishes for sheaves in degree zero. Naturality of rank and
Chern character proves the final equality.
\end{proof}
This statement uses the full Atiyah-trace map. Injectivity of a
single middle component is not being substituted for semiregularity.

\begin{proposition}[SF-CM02: the entire invariant character]
\label{SF-CM02}
Let $(A,K,h)$ be a polarized Weil sixfold of signature $(3,3)$.
If a rational even cohomology class $\kappa$ remains Hodge under
all local deformations in its full nine-dimensional Weil period
domain, then
\[
 \kappa=P(h)+t,\qquad P\in\mathbf Q[x],\quad t\in W_K\subset H^6(A,\mathbf Q).
\]
If $\kappa_3$ is not proportional to $h^3$, then $t\ne0$.
\end{proposition}
\begin{proof}
The Hodge circle of a balanced Weil point lies in $\operatorname{SU}(3,3)$:
on either six-dimensional $K$-eigenspace its determinant is
$z^3z^{-3}=1$. Its conjugates are the Hodge circles in the period
domain. Analyticity extends the assumed Hodge condition from an open
neighborhood to the connected period domain. The infinitesimal
generators of these circles generate a nonzero ideal in the simple
real Lie algebra $\mathfrak{su}(3,3)$, hence all of that Lie algebra.
The class $\kappa$ is consequently $\operatorname{SU}(3,3)$-invariant.

After scalar extension to $\mathbf C$, write $H^1=U\oplus U^*$,
where $U$ is the standard six-dimensional representation of
$\operatorname{SL}_6$. The decomposition
\[
 \bigwedge^k(U\oplus U^*)
 =\bigoplus_{a+b=k}\bigwedge^aU\otimes\bigwedge^bU^*
\]
contains an invariant exactly when $a=b$, or when
$(a,b)=(6,0),(0,6)$. Indeed the exterior powers of the standard
representation are irreducible with distinct highest weights, except
that $\bigwedge^0U$ and $\bigwedge^6U$ are both trivial. For $a=b$
the invariant line is generated by the $a$th power of the invariant
pairing $h$. The two exceptional lines are the determinant lines
$\bigwedge^6U$ and $\bigwedge^6U^*$ in degree six. Thus the entire
invariant subspace is
\[
 \mathbf Q[h]/(h^7)\oplus W_K.
\]
The polynomial summand already includes $\mathbf Qh^3$. The
decomposition is rational, since $h$ and the sum of the conjugate
determinant lines are rational. The asserted formula and nonvanishing
criterion follow.
\end{proof}
This supplies the full-character conclusion needed by S-14, rather
than a conclusion about the middle character alone.

\begin{proposition}[SF-CM04: the fixed hyperbolic spin model forces splitness]
\label{SF-CM04}
Let $(V,q)$ be a hyperbolic rational quadratic space of dimension
$4n$. Suppose $f^2=-d$, $d>0$, and
$q(fx,y)=-q(x,fy)$. Endow $V$ with its
$K=\mathbf Q(\sqrt{-d})$ structure through $f$. The Hermitian form
\[
 H(x,y)=d\,q(x,y)+\sqrt{-d}\,q(fx,y)
\]
(up to interchanging its linear argument convention) is split of
$K$-rank $2n$. Its determinant class is $[(-1)^n]$.
\end{proposition}
\begin{proof}
Skew-adjointness gives $q(fx,x)=0$, hence $H(x,x)=dq(x,x)$.
A nonzero rational isotropic vector for the hyperbolic $q$ is
therefore $K$-isotropic for $H$. Nondegeneracy splits off a
Hermitian hyperbolic plane containing it. The rational quadratic
form $\operatorname{Re}H/d$ on this plane is hyperbolic of rank
four: the two isotropic $K$-lines are two perfectly paired isotropic
rational planes. The Hermitian orthogonal complement is also
$q$-orthogonal and $K$-stable. Rational Witt cancellation makes
$q$ on that complement hyperbolic. Repeating gives $n$ Hermitian
hyperbolic planes, with determinant $(-1)^n$ modulo norms.
\end{proof}
In Markman's construction the ambient lattice is
$H^1(X)\oplus H^1(\widehat X)$ with the evaluation pairing, a sum
of $2n$ hyperbolic planes. His oriented $K$-secant endomorphism
and form satisfy these identities
\cite[\S3.1, equation~(3.1.2) and Lemma~3.1.2, pp.~22--23]{SF-Markman2502}.
Consequently the split determinant is not restricted to the explicit
ample secants $\exp(\pm\sqrt{-d}\Theta)$. Another secant in this
same quadratic model cannot repair the nonsplit determinant gap.
The linear-algebra proposition is unconditional; geometric uses of
the unrefereed construction remain conditional.

\begin{proposition}[SF-CM05: the literal elliptic-chain boundary fails]
\label{SF-CM05}
On $A_3=E_1\times E_2\times E_3$, for $b\in E_2\setminus\{0\}$,
put
\[
 C_b=(E_1\times\{0\}\times\{0\})
 \cup(\{0\}\times E_2\times\{0\})
 \cup(\{0\}\times\{b\}\times E_3).
\]
Its ideal sheaf is not semiregular. If a finite translation subgroup
$H\subset A_3$ gives pairwise disjoint translates of $C_b$, the
ideal of their union with its canonical $H$-linearization is not
weakly $H$-semiregular. The defect persists for the natural product
action on exterior products of two such orbit ideals, and under line
bundle twists or derived equivalences.
\end{proposition}
\begin{proof}
The curve is a connected genus-three chain with two ordinary nodes.
At either node its ideal has local form $(z,xy)$ in the smooth
threefold, so it is a local complete intersection. Its normal sheaf
has at least four independent global sections: the three translations
of $A_3$, and the variation of $b$. To verify independence, restrict
to the smooth portions of the three components. A translation vector
$(v_1,v_2,v_3)$ has normal components
\[
 (v_2,v_3),\qquad(v_1,v_3),\qquad(v_1,v_2).
\]
The $b$-variation is zero on the first two components and nonzero
only in the $E_2$-normal direction of the third. A translation
matching its first two restrictions must be zero and cannot match
the third. The same calculation shows that the three translation
directions themselves are independent. These are Hilbert tangent
vectors: $C_b$ is a flat family locally in $b$, with constant nodal
local equations, component genera and incidence relations.

Write $I=\mathcal I_{C_b}$. In the exact sequence obtained by applying
$\operatorname{Hom}(I,-)$ to
$0\to I\to\mathcal O_{A_3}\to\mathcal O_{C_b}\to0$,
the map $\operatorname{Hom}(I,I)\to\operatorname{Hom}(I,\mathcal O_{A_3})$
is the isomorphism $\mathbf C\to\mathbf C$. Consequently
\[
 H^0(N_{C_b/A_3})=\operatorname{Hom}(I,\mathcal O_{C_b})
 \hookrightarrow\operatorname{Ext}^1(I,I).
\]
Tensoring by $\operatorname{Pic}^0(A_3)$ gives another
three-dimensional subspace $H^1(\mathcal O_{A_3})$. Its trace is
the identity because $I$ has rank one, whereas the embedded ideal
deformations have trivial determinant. Thus the two subspaces are
independent and $\operatorname{ext}^1(I,I)\ge7$.
Serre duality on the abelian threefold gives
$\operatorname{ext}^2(I,I)\ge7$. The full semiregularity target has
dimension
\[
 h^2(\mathcal O_{A_3})+h^3(\Omega^1_{A_3})=3+3=6,
\]
so the map is not injective.

For the orbit union $Z_b=\bigcup_{h\in H}(C_b+h)$, disjointness
persists locally in $b$. Its invariant normal directions contain the
same four-dimensional diagonal subspace. The three trace directions
from $H^1(\mathcal O_{A_3})$ are invariant. Translation preserves
the volume form, so equivariant Serre duality yields
$\dim\operatorname{Ext}^2(\mathcal I_{Z_b},\mathcal I_{Z_b})^H\ge7$.
More decisively, descent identifies the invariant object with the
ideal of the embedded image of $C_b$ in the abelian threefold
$A_3/H$. That object's ordinary semiregularity target again has
dimension six. It is therefore not semiregular in the invariant
category, which is exactly failure of weak $H$-semiregularity.
The automorphisms of the ideal are scalars; other linearizations
only tensor its descent by a torsion line bundle and do not repair
the kernel.

On the exterior product, take a nonzero kernel element in one
factor's $\operatorname{Ext}^2$ and tensor it with the other
factor's identity. The Kunneth summand is nonzero but its full trace
is zero. Twisting by line bundles multiplies the trace by an
invertible exponential; derived equivalence transports the
categorical trace and its kernel. Thus the defect persists under
the stated operations.
\end{proof}
This proposition concerns the literal boundary chain and its natural
orbit repair. It does not obstruct smooth CM Jacobians or arbitrary
complexes on $E^6$. In particular, the chain used in an openness
argument cannot itself be substituted for the required semiregular
CM seed without an additional proof.

\subsection{Conditional subgroup: Markman's source sheaf}
\label{sec:conditional-cm-source}

\begin{theorem}[SF-CM03: a conditional balanced CM seed]
\label{SF-CM03}
Assume the sheaf construction and semiregularity assertions of
Markman \cite{SF-Markman2502}, specifically Theorem~1.4.1(2)--(4),
Lemmas~9.1.4, 9.3.1, 9.3.5, 9.3.11 and Remark~9.3.7.
Let $K=\mathbf Q(\sqrt{-m})$, $m>0$ squarefree, and let $E$ be
any elliptic curve with $\operatorname{End}_{\mathbf Q}(E)=K$.
On $A=E^6$ with the standard balanced action
$\operatorname{diag}(a,a,a,\bar a,\bar a,\bar a)$, there exist:
an integral ample compatible polarization $\theta$ of determinant
$[-1]$, a reflexive coherent sheaf $F$ of rank $32m$, and an honest
finite translation group $G\subset A\times\{0\}$ for which $F$
is $G$-semiregular and has no negative self-Ext, such that
\[
 e^{B_0}\operatorname{ch}(F)=P(\theta)+t,
 \quad B_0=-c_1(F)/(32m),\quad P\in\mathbf Q[x],\quad0\ne t\in W_K.
\]
The equality holds in the full even cohomology. This statement has no
class-number hypothesis. It does not assert a prescribed principal
polarization, a prescribed Weil direction, or a numerical detector
bound.
\end{theorem}
\begin{proof}
We specify the finite choices and identify precisely where the
conditional source assertions are used.

\emph{The radicand.} Set $d=4m$, $N=d+1$, and $r=8d=32m$.
Then $d\ge4$ is even, $K=\mathbf Q(\sqrt{-d})$, and
$\gcd(r,N^2)=1$. Thus the construction satisfies the source
hypothesis $d\ge3$ and admits its untwisted determinant-power
normalization. In particular the inadmissible $d=1$ Gaussian
specialization is avoided.

\emph{A good smooth CM Jacobian.} Let $\tau_0\in\mathfrak H_3$
be the period matrix of the product principally polarized $E^3$.
The orbit $\operatorname{Sp}_6(\mathbf Q)\tau_0$ is analytically
dense. Indeed the Cayley map
$X\mapsto(1-X)(1+X)^{-1}$ on the symplectic Lie algebra is rational
and parametrizes a neighborhood of the identity. Rational
Lie-algebra points are dense there. As
$\operatorname{Sp}_6(\mathbf R)$ is connected, its rational points
are dense throughout, and it acts transitively on $\mathfrak H_3$.
A rational symplectic transform changes the marked lattice to a
commensurable one, so every orbit point is a principally polarized
abelian variety isogenous to $E^3$.

The locus of nonhyperelliptic genus-three Jacobians contains an
analytic open subset of $\mathcal A_3$. To check the needed rank,
a smooth plane quartic has $K_C=\mathcal O_C(1)$, and
\[
 \operatorname{Sym}^2H^0(K_C)\longrightarrow H^0(K_C^2)
\]
is an isomorphism between six-dimensional spaces: no nonzero
quadric contains the quartic. This is the dual period differential,
so the differential has rank six. The good orbit and intersection
hypotheses in the conditional source Lemmas~9.1.4 and~9.3.1 hold
on a nonempty open subset of the finite moduli cover recording the
$N$-torsion groups. A finite etale cover is open, so the image gives
a nonempty analytic open subset of good Jacobians in $\mathcal A_3$.
Density supplies $q\in\operatorname{Sp}_6(\mathbf Q)$ in that
open subset. Choose its smooth curve $C$, put $X=J(C)$, and choose
the good cyclic groups $G_1,G_2$ of order $N$. Then $X$ is isogenous
to $E^3$.

This choice establishes existence in a nonempty open subset. No
specific rational matrix, plane quartic equation or torsion point
coordinates are asserted to have been computed for this conditional
construction.

\emph{An ordinary semiregular quotient.} The source theorem gives
a simple reflexive sheaf $\mathcal E$ of rank $r$ on
$M=X\times\widehat X$, together with finite translation/tensor
symmetries. Lemma~9.3.5, p.~84, uses even $d$ and
$\gcd(r,N^2)=1$ to cancel the tensor part by a determinant power.
One may take an integer satisfying $ar\equiv-1\pmod{N^2}$,
with the sign adjusted to the source convention for the intermediate
derived object. The result descends along a translation quotient
$q_M:M\to Y$ to an ordinary reflexive sheaf $\overline F$ of
rank $r$. Remark~9.3.7, p.~86, identifies the normalized twisted
sheaf $\mathcal B$ with the normalization of this ordinary sheaf;
Lemma~9.3.11, p.~88, states that $\mathcal B$ is semiregular.

This also proves ordinary semiregularity of $\overline F$.
Replacing its Atiyah class by
$\operatorname{at}_{\overline F}-(c_1(\overline F)/r)\operatorname{id}$
multiplies the entire Atiyah-trace map by
$\exp(-c_1(\overline F)/r)$. On the full target
$\bigoplus_qH^{q+2}(\Omega_Y^q)$ this is an invertible triangular
operator. The normalized local scalar twists cancel in the
endomorphism complex, so the two self-Ext spaces agree. Thus
injectivity of the normalized twisted map is equivalent to
injectivity of the ordinary map. Moreover
\[
 q_M^*\kappa(\overline F)=\kappa(\mathcal E),\qquad
 \kappa(F')=e^{-c_1(F')/\operatorname{rank}(F')}\operatorname{ch}(F').
\]

\emph{The balanced action on the actual power.} Since $Y$ is
isogenous to $X\times\widehat X$, it is isogenous to $E^6$ and
$\operatorname{End}_{\mathbf Q}(Y)\simeq M_6(K)$. Here the central
copy of $K$ initially comes from the elliptic CM action. Markman's
Weil embedding $\eta:K\to M_6(K)$ commutes with this center.
Over the central $K$ its generator has two eigenvalues
$\pm\sqrt{-d}$. The two eigenspaces have $K$-dimension three,
because their holomorphic dimensions are the Weil multiplicities
$(3,3)$. A matrix in $\operatorname{GL}_6(K)$ therefore conjugates
$\eta$ to the specified balanced diagonal embedding. Compose with
an isogeny identification and clear denominators to obtain a
$K$-linear isogeny $f:E^6_{\rm balanced}\to Y$. Equality of the
elliptic CM field and the secant field is essential to this
argument; it is not claimed for an unrelated elliptic field.

\emph{Pullback and the entire character.} Set $F=f^*\overline F$
and $G=\ker f$. SF-CM01 gives the honest translation action,
strong $G$-semiregularity and negative-Ext vanishing. The positive
rational Weil polarization $h_Y$ pulls back to a compatible
positive class. Choose $c>0$ such that $\theta=cf^*h_Y$ is an
integral ample line-bundle class. The source determinant is $[-1]$
by its Lemma~3.1.3, pp.~23--24, or by SF-CM04. Quotient and
$K$-isogeny transport preserve it, and scaling by $c$ changes the
rank-six determinant by $c^6=\operatorname{Nm}(c^3)$, leaving its
class unchanged.

Theorem~1.4.1(3) states that the entire $\kappa(\mathcal E)$ stays
Hodge on the full Weil period domain. Fixed isogenies identify the
local period domains, so the same is true of $\kappa(F)$.
SF-CM02 gives $\kappa(F)=P(\theta)+t$. By
Theorem~1.4.1(4), the source $\kappa_3(\mathcal E)$ has nonzero
Weil projection: its $K$-translates together with $h^3$ span three
dimensions. Isogeny pullback is injective and preserves the Weil
plane, so $t\ne0$. Finally $B_0=-c_1(F)/r$ is algebraic and
$e^{B_0}\operatorname{ch}(F)=\kappa(F)$, as required. Replacing
$\theta$ by a very ample positive multiple merely rescales $P$.
\end{proof}

\paragraph{What the conditional theorem changes.}
It constructs a genuine coherent object at the balanced CM point,
subject to the stated preprint input. The first-order assertion of
Theorem~1.4.1 alone would not suffice: the finite quotient and
Lemma~9.3.11 are essential here. Choosing even $d$ removes the gerbe
coherence issue entirely. SF-CM02 supplies the whole character, not
just its middle term. One nonzero Weil direction is enough for the
detector programme; its rational $K$-translates span the plane.
There is no assertion that arbitrary direct sums preserve
semiregularity while prescribing another direction.

Applying the supplied DEG\_L S-14/Perry theorem gives bounded
division on the associated split determinant family conditionally
on these source assertions. The resulting bound remains
existential until the chosen CM Jacobian, sheaf presentation and
degree data are made explicit. It is not presented as a numerical
detector constant, and the chosen polarization on $E^6$ is not
identified with its product principal polarization. Nonsplit
determinants remain outside this fixed hyperbolic secant model.

In the programme's language, this is a conditional $K$-theory input
to S-14, with a full $\lambda$-ring Chern character and an object
whose deformation theory is controlled. The subsequent uniform
integral-cycle step uses the supplied HECKE\_DIV/GHOSTS reduction,
condition A at cote~16 scan~050. Neither condition I (scan~049 and
cote~44 Theorem~4) nor a bound on the CM isogeny or core index $J$
produces the required object. No additional archive or EGA/SGA
locator is claimed for the new representation and descent arguments.


\part{Open gates and unchecked work}
\section{Current state of the remaining work}
This is an active checkpoint draft, not a final claim of completion.
The rank-zero perfect-complex construction D-SR-01--06 is checked and
included above. The coherent variant SF-QC01--05 has now passed a full
proof reading and exact arithmetic checks, giving a numerical degree.
HD-CM6 explicitly charges its source-to-target lattice map. The remaining
main sixfold gate is the nonsplit determinant outside the Eisenstein
case, rather than numerical control on the split families.
The $G_{24}$ and $F_4$-extension proofs are checked, including their full
limiting-period argument, integral field action, cycle transfer, and
fixed-lattice cost. They reach the split determinant only. Other sixfold
determinants, except the Eisenstein cases settled by HD-E6 and the
Gaussian signed class three settled by HD-E8Q, remain open in this
run. The universal local determinant formula D-DET-01 and
the normal-closure and monomial barriers are now checked. They exclude
every dominant monomial nonsplit source, every reflection-containing
source over the seven fields other than the Gaussian and Eisenstein
fields, and the natural Gaussian groups $G_{29},G_{31}$ in the
precise scopes stated above. The complete Gaussian reflection-component
reduction D-GAUSS-01--03 now leaves only the scalar $E_8$ extension.
An exact twelve-branch tuple generates that full group, has signature
$(3,3)$ and signed discriminant $3$, and gives a connected
nine-parameter family. D-E8-GRR-01--06 and D-E8R-01--06 now independently
prove actual period rank nine on a connected component. HD-FS1--2 and
SF-F1201--04 supply its rational Weil classes algebraically; SF-SKQ01--05
gives an explicit effective degree on every required product of
Fermat-twelve hypersurfaces. D-E8-XFER01--03 and HD-E8Q1--6 now
charge the entire cyclic/Brauer correspondence cost; D-E8-LAT01--02
charges the nonsplit target-lattice comparison. The resulting
$D_3(f^{12}\Delta)$ and quadratic division bound are explicit.
The explicit two-pair tuple has graph rank two and an anisotropic
residual plane; it is not asserted to lie in either dominant component.

The checked bireflection reduction D-BI-01--07 further shows that a
remaining reflection-free nonsplit source must have coefficient degree
at least fifteen and cannot be induced from a representation of dimension
one, two, or three. D-BI-08 also excludes scalar alternating and
symmetric deleted-permutation inducing representations in every degree.
D-BI-09--10 forces every primitive essentially self-dual inducing
representation of dimension at least nine to have dimension at least
twenty-four; the original primitive exact-$K$ representation has even
dimension outside the Gaussian and Eisenstein fields. D-EX-01--03
excludes every scalar exterior-square deleted-permutation inducing
representation in arbitrary degree and with arbitrary inducing field.
D-UP-01--03 likewise excludes every scalar unordered-pair Specht block
$S^{(n-2,2)}$. D-BI-11 rules out inducing dimension four over
$\Q(\sqrt{-d})$ for $d=7,11,19,43,67,163$. For the Gaussian and
$\Q(\sqrt{-2})$ fields, it reduces that dimension to the explicitly
listed biquadratic inducing fields, with ten reflection inertias and
two moving-rank-two inertias on a rational intermediate base.
D-BI-12 further removes every real reflection normal closure in
inducing dimension four, by constructing a rank-three monodromy
logarithm over the original quadratic field. Only the stated
$G_{29},G_{31},G_{32}$ rows remain in that dimension.
D-BI-13 further shows that over $\Q(\sqrt{-p})$, for
$p=7,11,19,43,67,163$, every inducing image is reflection-free,
regardless of its character field. D-BI-14 forces an inducing
degree of at least $7(p-1)$ whenever its image is scalar times a
finite model defined over a field unramified at $p$; the bound is
$12(p-1)$ for a primitive tensor-indecomposable inducing image.
D-BI-14a strengthens the latter bound to
$(p-1)(25(p-1)-22)/4$. D-BI-15 bounds the period rank of a
larger-field inducing cover with a nontrivial compatible base
automorphism by $\lfloor R/2\rfloor-1$, where $R$ is its number
of genuine branch points. Thus the remaining twelve-branch
larger-field cases require a nonnormal inducing subgroup and
more than one imprimitivity system of the given block dimension.
D-R5-01--03 leaves only $\mu_4G_{33}$ over $\Q(\zeta_{12})$
for a non-Eisenstein inducing block of dimension five, with
original field $\Q(i)$. HD-E8IF1 gives the two exact larger-field
$E_8$ branch profiles. Existence, original-field descent and
dominance are separate remaining requirements for those profiles.
D-BI-16 reduces inducing dimension six to the Gaussian
$\mu_4G_{34}$ model over $\Q(\zeta_{12})$, with the exact
twelve-to-fourteen-branch moving-rank profiles stated above.
D-BI-17 excludes inducing dimension seven throughout the
non-Eisenstein class-number-one scope. D-R5-04 supplies an actual
natural $\Q(\zeta_{12})$ cover with the rank-five profile and
excludes its specific Nielsen component by a rank-three boundary.
Its local signed determinant is a norm over the larger CM field;
this does not prove splitness over the original Gaussian field or
identify the other Nielsen components.
D-BI-18--20 now give a structural normal-core gate. Every primitive,
factorwise tensor-indecomposable inducing image of degree at least five
has a unique nonabelian simple-power core modulo scalars, whose
quasisimple central-product lift acts irreducibly and has transitive
factors. Extraspecial alternatives are excluded. In inducing degree
at least ten, a nonalternating Lie core has field size at most twenty-seven;
essentially self-dual cores obey the stronger stated field bounds.
D-SPOR-01--03 excludes every sporadic core. D-BI-21 excludes all
rank-one Lie cores, including the small exceptional alternating
isomorphisms and covers in its scope. D-BI-22 excludes all Suzuki,
Ree and Tits cores. D-BI-23 excludes every transitive linear-Weil
core of $\operatorname{SL}_n(q)$ for $n\ge3$, in all field sizes.
D-G2T-01--03 excludes all $G_2(q)$ with $q\ge3$ and all
${}^3D_4(q)$ with $q\ge2$, including mixed factors and all compatible
outer automorphisms. D-SPW-01--03 reduces odd symplectic Weil cores
to one degree-four or degree-five factor of $\operatorname{Sp}_4(3)$;
D-SPW-04 gives the same final small endpoint for all its irreducible
representations in the nonsplit non-Eisenstein problem.
D-BI-24 reduces all transitive ordinary unitary-Weil cores, in all
field sizes, to the single $\operatorname{PSU}_4(2)$ degrees five
and six; only five remains in that nonsplit scope.
D-ESC-01--03 excludes all even-characteristic $\operatorname{Sp}_4(q)$
and $\operatorname{Sp}_6(q)$ cores in the main nonsplit problem.
For $q\ge4$ its theorem also permits mixed ranks and field sizes;
the separate transitive $\operatorname{Sp}_6(2)$ endpoint uses the
precise non-Eisenstein degree-seven exclusion. D-A2-01--03 excludes
all ordinary quasisimple $\operatorname{PSL}_3(q)$ and
$\operatorname{PSU}_3(q)$ cores for $q\ge3$, including mixed fields
and types. L34-1--3 supplies the exceptional $\operatorname{PSL}_3(4)$
normalizer coverage by the published thirty-six bicyclic classes;
the computed universal table is only a redundant check.
The successful $\operatorname{PSL}_3(2)=\operatorname{PSL}_2(7)$
case is explicitly outside that exclusion. D-HES-F excludes arbitrary
mixed products of $\operatorname{Sp}_8(2)$, $\operatorname{Sp}_{10}(2)$
and $\operatorname{Sp}_{12}(2)$, with all normalizing scalars and factor
permutations. D-HES-01--03 then excludes every
$\operatorname{Sp}_{2n}(2)$ core with $n\ge4$, and arbitrary mixed products.
Its explicit elementary parabolic subgroup replaces the previously
unverified semisimple estimate. The only character-ratio premise is the
verified published unipotent bound; the essential finite endpoint inputs
are the ATLAS $\operatorname{Sp}_8(2)$ and $\operatorname{Sp}_{10}(2)$
tables. No missing Gluck--Magaard (1998) premise or separately computed
$\operatorname{Sp}_{12}(2)$ table is needed for this all-rank theorem.
D-HES-04--06 then excludes every higher symplectic core
$\operatorname{Sp}_{2n}(q)$ with $n\ge4$ and $q$ even, allowing mixed
ranks and fields, all compatible outer operators, scalars and factor
permutations. Its uniform normalizer support $7e/64$, together with the
nontrivial bilinear multiplier, forces $e\le64$; the actual factor
degrees exclude every case. The $q=4$ field-involution endpoint is
proved by explicit symplectic descent and dual-weight pairs.
D-C3-01--03 excludes every ordinary quasisimple
$\operatorname{PSp}_6(q)$ core for odd $q\ge3$, including mixed
fields and all compatible outer, scalar and factor-permutation actions.
The complete similitude table, careful Clifford restriction and
projective semilinear descent at $q=9$ give uniform normalizer support
$3e/16$, forcing $e\le24$ and excluding every endpoint. No selfduality
premise is needed. The generic full-table premise is explicitly
L\"ubeck's dissertation and its official mathematical implementation,
not a separately refereed journal table. D-OC2-01--04 now excludes every
transitive ordinary $\operatorname{PSp}_4(q)$ core with odd $q\ge5$
in the nonselfdual-original source problem, allowing all compatible
normalizers, larger inducing fields and either determinant. Its ordinary
unipotent restrictions, elementary shear and Heisenberg estimates, and
explicit field descent include the characteristic-equals-field-order case.
The complete universal $q=5$ endpoint is supplied by D-OC2-F; the separate
finite $q=7,9$ tables have only their stated center-positive coverage.
D-EXS-01--03 excludes arbitrary mixed products of the four simple types
$F_4(2)$, $F_4(3)$, $E_6(2)$ and ${}^2E_6(2)$, including all covers and
compatible normalizers. Complete exact spectra give common normalizer
support $e/5$ and force $e\le15$, below the minimum factor degree $52$.
D-HOS-03--06 excludes every higher symplectic core
$\operatorname{Sp}_{2n}(q)$ with $n\ge4$ and odd $q\ge7$,
allowing arbitrary mixed fields and ranks, all compatible normalizers,
scalar actions and factor permutations. Its explicit elementary
normalizer estimate $9234/55223>1/6$ forces total degree
$e\le165669/181<916$, below every factor degree $1029$ or more.
D-HOS-07--10 now excludes all higher $q=5$ cores, with arbitrary
mixed ranks, full compatible normalizers and tensor permutations.
Opposite radicals give a trivial-isotypic proportion at most one half;
the exact Heisenberg carrier argument includes defining-characteristic
order. The resulting normalizer support is $31e/250$. Every ordinary
core factor is self-dual, so the original imaginary character field
forces a nontrivial form multiplier. Two such inertias and the other
ten force $e\le25<250$, the elementary factor-degree floor.
D-HOS-11 supplies full-normalizer support $13e/108$ for every
ordinary higher $q=3$ factor. D-HOS-12 excludes all self-dual such
factors, also allowing arbitrary mixtures with the $q=5$ factors:
the nontrivial multiplier gives $e\le324/11<30$, and exact integer
costs exclude all three remaining degrees $27,28,29$. D-HOS-20--22 now excludes every remaining nonselfdual higher $q=3$
source. For every odd prime at least five, symmetric-form weight
orbits have fixed fraction at most $1/9$. Stable isotropic primary
pairs and the published restriction bound give support $4a/15$;
when no such pair exists the full centralizer is a norm-one torus
of order at most $10^{n/2}$, and column orthogonality is stronger.
Combining all projective-prime cases gives support $2a/9$ for even
half-rank and $19a/104$ for odd half-rank. Scalar-inner compatible
normalizers, tensor convolution and factor cycles preserve these
bounds. Twelve genuine inertias force degree at most $9$ or $156/5$,
contradicting respectively $27$ or $81$. Together with D-HOS-12,
this closes all ordinary transitive higher $q=3$ cores, $n\ge4$;
the prior rank and inertia restrictions are now superseded.
Small ranks retain their separately stated earlier results. No new
arbitrary mixture of $q=5$ with nonselfdual higher-field cases is
asserted here.
D-UNW-01--03 and D-LNW-01 give the explicit natural-dimension cutoff
$n<10^{61}$ for transitive primitive factorwise tensor-indecomposable
unitary and linear cores. Above that cutoff all non-Weil factors have
full-normalizer support at least $11e/60$, forcing $e\le30$;
degree separation and the already proved Weil exclusions treat the
other alternative. In the nonsplit main-source scope the existing
field bound $q\le27$ also applies. This makes the ranges of simple factor types finite. It does not
enumerate tensor multiplicities, scalar extensions, or actual sources.
D-OSM-01--04 excludes arbitrary mixed cores from the eight types
$\Omega_7(3)$, $\Omega_8^-(2)$, $\mathrm{P}\Omega_8^\pm(3)$,
$\Omega_{10}^\pm(2)$ and $\Omega_{12}^\pm(2)$, with all covers and
compatible normalizers. Fifteen complete tables give common support
$e/5$ and force $e\le15<27$. The ninth type $\Omega_9(3)$ is
excluded under transitivity and the nonselfdual-original hypothesis:
its two nonselfdual characters have quarter normalizer support, while
the other case uses the nontrivial form multiplier and gives
$e\le41/4<780$. This ninth conclusion is not enlarged to arbitrary
mixtures. The small $\Omega_8^+(2)$ representations retain their
separate gates. Other small-field higher-rank Lie representations
remain outside these exclusions.
D-OL-01--03 gives a separate effective orthogonal reduction at natural
dimension $d\ge10$. In the transitive setting every factor degree is
less than $q^{(2\cdot10^{29})d}$, and every descended ordinary character
has level $\ell<10^{30}$. A genuinely spin factor, or in odd
characteristic a character which does not extend to the special
orthogonal group, must have $d<2\cdot10^{30}$. The mixed-core
degree obstruction uses full-normalizer support $3e/16$ and $e\le24$.
At that stage the remaining ordinary low-level natural dimension is unbounded;
no theta irreducibility is used. D-OL-04--05 then proves that a level
$\ell$ ordinary character extends irreducibly to the full orthogonal
group, and is self-dual, whenever $d\ge4\ell+2$. The proof compares
actual ordered-tuple orbits and their commutant algebras.
Thus a nonselfdual ordinary factor has $d\le4\cdot10^{30}-3$.
For $d\ge4\cdot10^{30}-2$ every factor is ordinary and self-dual;
the actual inducing image has a nontrivial bilinear-form multiplier.
The published ratio estimate at $q\ge16$ then forces $e\le64$,
contradicting the factor-degree bound. The unbounded natural-dimension
case is therefore confined to $q\in\{2,3,4,5,7,8,9,11,13\}$ and
ordinary self-dual characters of level below $10^{30}$ at that stage.
D-OL-06--08 now closes this unbounded-dimension case. For every
$d\ge22$ and these nine fields, an ordinary irreducible character
extending to $O$ has full compatible-normalizer support at least
$7a/64$. Opposite abelian radicals, quadratic weight counts, reflection
commutators and finite-field descent give the complete bound, including
the possible character discrepancy between extensions to $O$.
The actual nontrivial multiplier budget forces $e\le64$, below the
factor degree. Together with D-OL-01--05 this leaves only
$d\le4\cdot10^{30}-3$ for every transitive orthogonal core.
This is a finite simple-factor parameter range, not an enumeration
of tensor powers, scalar extensions or actual sources.
D-OL-09--11 sharpens two ordinary-character subranges: for all $q\ge3$,
a remaining selfdual ordinary factor has $d\le21$, without an
extension-to-$O$ assumption. For all even $q$, every remaining ordinary
transitive factor has $d\le1134$, without a selfduality assumption.
D-OL-12--13 further reduces ordinary odd-dimensional factors over odd
$q\ge5$ to $d\le1133$. In the nonselfdual case, compatible linear
normalizers are scalar-inner by full-$O$ reality and Clifford restriction.
Signed line-radical actions, including moving rank one, have fixed weight
fraction below $1/4$; the remaining inner cases have restricted natural
support at least $567$. The published spin-group bound then gives
$\nu>\{1-(19/20)^9\}a/2$, and twelve genuine inertias force degree
below $28$, contrary to the factor floor. Arbitrary mixed field components
are included. At $q=3$ the proposed quarter fixed-weight shortcut is
explicitly false. D-OL-14--15 nevertheless closes the large ordinary
$q=3$ odd-natural endpoint by a different argument. Only actual negative
reflections can have projective support below the stronger constant;
higher roots are covered by the full restricted-support character bound.
The exact budget leaves twelve bad inertias, each active on one factor.
Their generation would make at most twelve negative reflections generate
the perfect natural orthogonal factor. Their common nonzero perpendicular
instead carries a nontrivial scalar sign character, a contradiction.
Thus all ordinary odd-natural transitive cores over every odd field
satisfy $d\le1133$. D-OL-16--18 closes the remaining ordinary
even-natural odd-field case at $d\ge1136$. The nonsquare normalizer
uses an invariant isotropic two-plane and its full Heisenberg radical;
the square and field cases retain all compatible normalizers. At
$q=3$, a determinant argument leaves only projective reflections,
and twelve such actual generators contradict perfectness on their
common natural fixed space. Consequently every surviving ordinary
transitive factor has $d\le1134$, with $d\le1133$ in odd dimension.
D-OL-19--21 now removes the genuine-Spin boundary in odd
characteristic and even natural dimension: every nontrivial
irreducible module of every represented Spin central quotient
is covered, giving $d\le1134$ in this parity as well.
The even-Clifford lift and its odd bimodule prove actual
kernel preservation, while unique odd-primary radical lifts
retain all multiplicities and normalizer actions. D-OL-22--24 supplies the separate odd-dimensional genuine-Spin
argument. Its order-four Heisenberg carrier covers both natural
$SO$ cosets, while higher two-power roots are handled by their
squares and the published restricted-support estimate. Thus
every odd-characteristic transitive Spin factor, with arbitrary
represented central kernel and any nontrivial complex irreducible
module, satisfies $d\le1133$ in odd dimension and $d\le1134$
in even dimension. The same source contradiction excludes mixed
large-dimensional factors of both parities.
D-OL-27--29 now sharpens this odd-characteristic bound to
$d\le185$ in odd dimension and $d\le186$ in even dimension,
for every represented central kernel and every nontrivial factor
module. Its new trace theorem needs restricted natural support
only $93$. A weighted count of actual trivial-factor spaces,
the regular-seven orthogonal centralizer, and an exact packing
of good blocks replace the earlier coarse restriction estimate.
The actual-source contradiction permits arbitrary mixed factors
above these thresholds. Neither numerical exponent of the
withdrawn D-OL-25--26 candidate is used; that candidate remains
withdrawn. Characteristic two keeps its separately proved bounds.

For alternating cores, HD-SPIN-02 excludes all faithful double-cover
spin factors in degree $n\ge8$. HD-ALT-02 and D-AST-01--03 together
exclude the entire ordinary range $8\le n\le14$. HD-ALT-03 excludes
all ordinary factors for $n\ge15$ with
$\max(\lambda_1,\lambda'_1)\le\lceil n/2\rceil$, including both
nonreal and real split constituents. Its remaining large ordinary
factors extend to real Specht representations and their full inducing
image is essentially self-dual. HD-ALT-04 further forces every
transitive ordinary factor label to have a row or column longer than
$3n/4$: the complete normalizer bound gives $e\le64$, whereas
each factor has degree at least $125$. HD-ALT-05 then strengthens
this to a row or column longer than $6n/7$. Its tail-content
estimate gives $e\le144$, against an unsplit factor degree
at least $250$. HD-ALT-06a/b keeps the exact first-row square-content
contribution and strengthens the remaining row or column bound to
$87n/100$. Its quantitative necessary inequality is
$\delta(n,k)\le1/10+3/(5e)$ for the explicitly given tail function.
The excluded annulus gives $e<2000$ for $n\ge30$, against the
precise cubic degree floor, and $e<54$ at the fifteen smaller endpoints.
HD-ALT-07 now removes the fixed-point-free involution loss with the
published Fomin--Lulov rim-hook estimate. Its common full-normalizer
support fraction is exactly $k(n-1-k)/((n-1)(n-2))$, giving the stronger
necessary inequality with this fraction in place of $\delta(n,k)$.
Every remaining row or column must consequently exceed
$\tfrac12(1+\sqrt{3/5})n$. The label $(88,12)$ at $n=100$ passes
the old finite-degree gate but fails this new one, proving strictness.
HD-AST-04--06 retains the actual original field in the rational
tensor-wreath scalar coordinate. Quadratic conductors force two inertia
phases of specified orders. Rational eigenspaces and the fourth-root
trace bound exclude standard tensors through $n=19$ over every
non-Eisenstein field, through $20$ when $8$ divides the discriminant,
and through $5(p-1)$ for odd discriminant primes $p\ge7$.
HD-ALT-08 then treats all remaining long labels: the necessary tail
inequality is $\beta(n,k)\le1/(5h)+3/(5e)$, with $h=4$ in the
conductor-eight case and $h=p-1$ in the odd-prime case. The required
row or column exceeds $\tfrac12(1+\sqrt{1-4/(5h)})n$.
For $\mathbf Q(\sqrt{-2})$ this is over $0.9472135954n$;
the other six non-Gaussian/non-Eisenstein class-number-one fields
have stronger thresholds. The nonstandard label $(94,6)$ at $n=100$
shows strictness beyond ALT07. The Gaussian general-shape gate is
unchanged by this conductor argument.
D-ATT-01--03 gives stronger exact support and trace estimates for both
two-box-tail cores. For every tensor count at least two, both labels
are excluded through $n=39$ over every non-Eisenstein field. The
conductor-eight exterior case also excludes $n=40$; for an odd
discriminant prime $p\ge7$, the unordered-pair and exterior ranges
extend respectively through $10(p-1)-1$ and $10(p-1)$.
The degree-$39$ profiles pass the preceding shape gate, giving a
separate strict reduction. D-ATT-04--06 further excludes both
two-box labels for every $n\ge15$ with exactly two transitive factors,
over every non-Eisenstein original field. Exactly two marked factor
swaps and an actual degree-$2n$ quotient give a Riemann--Hurwitz
contradiction. D-ATT-07--09 removes every larger transitive
tensor count for these same two labels over non-Eisenstein fields.
Its exact weighted-cycle error includes all short word orbits, and
the factor-fixing index bound keeps its full constant term.
HD-AST-07--09 independently excludes the standard two-factor case
in every degree $n\ge15$, over all original fields, using a new
uniform tensor index estimate and the actual degree-$2n$ quotient.
HD-AST-10--12 then excludes every larger standard tensor count
over non-Eisenstein fields by pair selection and the same exact
cycle error. These all-count results supersede the finite cutoffs
for the three labels. Every remaining ordinary long label in the
main non-Eisenstein problem now has at least four boxes off its long
row or column and satisfies the exact ALT07/ALT08 shape bounds.
HD-AT301--03 further removes the entire two-row label $(n-3,3)$,
for $n\ge15$: a single factor is excluded over every original
field, and every transitive tensor count over non-Eisenstein fields.
Its new all-eigenvalue character-defect estimate has a full positive
constant term, so the actual natural quotient Riemann--Hurwitz
identity exceeds the permitted moving-rank budget.
D-HK301--03 and HD-MT301--03 exclude the other two three-box
labels, $(n-3,1,1,1)$ and $(n-3,2,1)$, in the same single-factor
and all-count scopes. Thus all three-box tails disappear from
the remaining non-Eisenstein ordinary transitive gate.
The derived HD-ALT-10 gives $n\ge20h-2$: respectively
$38,78,118,198,358,838,1318,3238$ in the eight non-Eisenstein
class-number-one fields. These are necessary source inequalities,
not existence or dominance statements.
D-AFT01--03 gives an effective threshold $N_\mu$ for every
fixed partition tail $\mu$ of size at least three. For
$n\ge N_\mu$ its character-polynomial coefficient bounds yield
$\nu(W_nx)\ge A\operatorname{ind}(x)+B$ with a full positive
constant. The actual natural quotient then excludes one factor
over every original imaginary quadratic field, and every
transitive tensor count over non-Eisenstein original fields.
HD-TW01 now aggregates all weighted factor cycles while retaining
both actual quotient-cover indices and the full additive support
constant. It extends the same fixed-tail thresholds to every tensor
count over every imaginary quadratic field. The published effective
mixing theorem and a uniform small-tail coefficient bound then give
HD-TW03: with $M_*=1+241\cdot240!\cdot2^{239}$,
$F_*=240M_*+2$, and
$\mathcal N=[100(1+M_*)F_*^{240}]^2$, every remaining ordinary
alternating factor in the main non-Eisenstein transitive-core gate
satisfies $n<\mathcal N$. The former growing-tail range is finite.
This retains the primitive factorwise tensor-indecomposable hypothesis
of D-BI19; primitive factorwise tensor-decomposable images remain outside
that structural theorem.
The scalar single-factor
standard, exterior-square and unordered-pair families remain excluded
in all degrees by D-BI-08, D-EX and D-UP. D-A7-01--03 excludes
every transitive $A_7$-cover core in the nonsplit problem, including
the double, triple and sixfold covers and all compatible normalizers.
The sole standard-factor exception is excluded by its actual split
polarization in the single-factor case and a scalar-multiplier budget
in the tensor case. The finitely bounded long-shape candidates still require their actual
field, rank and cycle arguments. No surviving group profile has been promoted to a source
without original-field descent and an actual period-rank computation.

HD-TC01--03 now bounds the previously uncontrolled tensor and
scalar parameters for any fixed finite factor/module catalogue.
For a transitive core with factor degree $a$ and total degree
$e=a^t$, the actual factor quotient and cyclic all-power traces give
$t\le\max\{4,\lceil5a/4\rceil-1\}$.
If $P$ is the represented core modulo its center and
$M=|\operatorname{Aut}(P)|$, the actual inducing image $J$ has
scalar kernel of order dividing $eM$ and $|J|\le eM^2$.
This makes the group, tuple and complex Nielsen-component rank
enumeration finite and effective whenever the factor catalogue
is finite. It does not execute that enumeration, settle original
quadratic-field descent or polarization compatibility, or bound
an arbitrarily larger original inducing overgroup. Those
interfaces remain part of the actual supplier problem.
The derived HD-TC04 can use a stronger local support fraction directly;
it is not separately counted as movement.

D-OG01--03 proves that arbitrary original overgroups and induction
indices really can be unbounded with a fixed inducing image: wreath
padding of the actual Gaussian E8 source preserves rank nine and the
signed determinant class three, with a fully identified polarization
scale. Proper-field descent requires the actual rational isogeny-power
locus; a complex rank calculation alone does not supply it.
HD-PD01--02 gives a complementary restriction without bounding those
original groups. For any integer $Q$, set
$M_Q=\operatorname{lcm}(8,|\operatorname{disc}K|,1,\ldots,Q)$.
Every prime $p\equiv-1\pmod{M_Q}$ is excluded as a signed determinant
class for a naturally polarized exact-$K$ primitive sixfold with a
faithful inducing image of order at most $Q$.
The integral local proof permits $p$ to divide the original group or
quotient-transfer order. It also excludes every $K$-isogenous target
with $\operatorname{End}^0=K$, regardless of its initially specified
polarization. The resulting infinitely many determinant exclusions
concern this supplier route, not arbitrary algebraic cycles or a lower
bound for $\delta_W$. No global exhaustion of factorwise tensor-decomposable
images is assumed.

For the genuine balanced-real-multiplication eightfold alternative,
SF-E811 excludes every stable locally free normalized balanced character,
and SF-E812--813 excludes weakly equivariant semiregular semistable
bundles with locally free stable factors. SF-E818 extends this to
balanced torsion-free sheaves whose locally free hull has such a
stable-factor filtration: mixed positivity eliminates surface and curve
defects, and SF-E817 excludes every remaining nonzero point defect.
The point-defect theorem itself does not assume stability or simplicity
and gives trace-kernel dimension at least $1+3s$, for $s$ support orbits.
SF-E834 now extends mixed discriminant nonnegativity and its equality
characterization to stable reflexive sheaves. Translation averaging
over complete-intersection surfaces supplies precisely the required
Chern--Weil pairings. Equality removes mixed curvature; Bianchi and
compact transverse slices bound both pure blocks across the entire
singular set. A smooth determinant normalization and the published
Biswas extension theorem make the sheaf locally free, after which
SF-E811 gives projective flatness. Thus no stable reflexive sheaf has
the positive balanced normalized character.
SF-E835 applies this to the reflexive hulls of all stable factors of
an arbitrary semistable torsion-free sheaf. The exact discriminant
variance and effective support cycles remove surface defects; the
normalized cubic character removes curve defects. Actual extension
classes lift across the remaining finite set by local cohomology in
dimension four. Consequently the actual reflexive hull has the precise
locally free stable-factor filtration required by SF-E818.
Every semistable torsion-free balanced sheaf is therefore excluded
from weak equivariant semiregularity, without an equivariant choice
of stable-factor filtration.
SF-E836 independently restricts the actual point-torsion/point-cokernel
class: if both defects are free over one common finite complete
intersection admitting a commuting generator triple for either
complementary bivector, the two first-order deformations are impossible.
Both orders of the actual companion-matrix equations yield the scalar
cross-minor identities; the commuting triple then contradicts the
published Jacobian-socle theorem. This includes every such complete
intersection of embedding dimension at most three, and all coordinate
power ideals. HD-CI-B1 records an actual length-sixteen integrable
counterexample to extending the scalar inference without that hypothesis.
It does not solve the stronger matrix equations or produce a supplier.
SF-E837 now uses those stronger equations to exclude every common
formal square thickening, with arbitrary nonlinear mixed parameters.
A square matrix whose cokernel is killed by four distinct pair
monomials must be invertible, by its filtered degree-two image and
the two omitted-edge constraints. The bracket-support alternatives
then force proportional complementary bivectors, a contradiction.
This includes the HD-CI-B1 ideal without altering its counterexample
to the scalar inference. It does not extend to arbitrary Artin modules.
SF-E838 next excludes a nonempty open family of regular quadratic
complete intersections and every higher-order perturbation of those
jets. A full-rank bracket pencil forces its actual companion matrix
invertible, using the strict Hilbert inequality $6>4$; the resulting
proportional leading tuples contradict the nonzero Jacobian socle.
Seven degree-six-pencil coefficients, each of degree sixteen in the
quadratic data, cut out a necessary residual locus. The explicit
nonsquare example has nonzero pure endpoints and no commuting-triple
presentation, so this is strict progress beyond both E836 and E837.
It asserts no module deformation on the residual locus.
SF-E839 strengthens that locus to rank at most four for every pencil.
Five independent quadratic annihilators already contradict the same
filtered image bound. Its thirty-six five-minors yield at most 216
coefficient equations, each of degree at most fourteen on a normalized
Grassmannian chart. An explicit nonsquare rank-five example has an
identically zero determinant pencil and no commuting-triple presentation,
proving strictness relative to all three preceding criteria. The equality
case for the abstract matrix lemma is explicitly false, with its own
small algebraic boundary example. No residual-locus existence follows.
SF-E840 determines the relevant equality constraint: every rank-four
pencil image must be $\ell V$ for a unique linear factor. The
regular-quadratic multiplication lemma and the exact filtered matrix
equalities prove this for every multiplicity. An explicit nonsquare
intersection has rank-four pencils and all preceding scalar and
rank-five tests zero, yet no common linear factor. Thus the new
condition strictly narrows the remaining quadratic locus.
SF-E842--843 adds an all-size socle-pair obstruction in the actual
common complete-intersection point-defect module. If two entries of
a symplectic bracket pencil have nonzero socle product and intersect
only in its line, the actual companion matrix must be invertible;
the nonzero Jacobian Pfaffian forbids that invertibility.
Consequently every parameter-power complete intersection is excluded,
for arbitrary positive exponents and arbitrary nonlinear mixed regular
parameters. Explicit length81 cubic examples pass all scalar companion
minors and lie outside the earlier commuting-triple and quadratic
families, so this strictly reduces the local gate. Actual coherent
balanced extensions of dimension $486m^2$ are constructed and excluded
from weak semiregularity. Unrelated Artin modules and complete
intersections outside the criterion remain open; SF-E841 stays an
unpromoted derived draft.

SF-E819 now includes actual point torsion as well as point cokernel.
For any balanced coherent sheaf with locally free hull and finite
kernel/cokernel, weak equivariant semiregularity forces equal lengths
at each point. SF-E819R excludes every nontrivial local modification
when both modules are reduced. Hence only equal-length modifications
with a nonreduced module at every nontrivial point remain in that local
class. SF-E820 further excludes every nontrivial point at which both
modules are free over the same curvilinear Artinian quotient, of any
length. It constructs actual rank-thirty-two balanced nonreduced
examples of equal signed local length and excludes them by two
first-order local quantization obstructions. The actual rank-thirty-six balanced nonpure sheaf is constructed
and excluded; it is not a successful seed. A locally free sheaf with
zero defect is outside this local conclusion and still requires the
separate stability hypotheses of the earlier results.

For derived constructions, SF-E815/R requires at least sixty-three
invariant endomorphisms in the seven-class model and excludes the
stated simple nonnegative-aggregate projectively flat models.
SF-E816/a/b/c excludes ordinary-gluable primitive projectively flat
models in any two cohomology degrees, including overlaps and arbitrary
gaps, and arbitrary bounded amplitude with globally disjoint classes.
The actual three-degree example SF-E816d refutes automatic disjointness,
but has a pure character and is not a balanced seed. Remaining routes
include unstable torsion-free objects, the nonreduced equal-length
modifications just described,
positive-dimensional torsion or defects outside the following surface
restrictions, and overlapping class sets in at least three cohomology
degrees. SF-E821/E821a proves that a genuine local surface summand in
a balanced complex has mixed RM tangent type and zero pairing with
each squared eigenpolarization. This excludes the entire rank-sixteen
ample-surface cone family, even with an arbitrary torsion-free
quotient. SF-E822 also excludes genuinely coupled coherent surface
modifications when the smooth integral support has class
$k^2\theta^2$ and its torsion and defect bundles have equal rank and
equal polarization degree. In particular it excludes all such repairs
preserving the entire hull character, including actual rank-thirty-two
nonsplit examples. SF-US constructs an actual unequal-degree unit
surface extension and a positive integral split hull with the entire
balanced character; all constant bivectors have explicit local
nullhomotopies. This construction is not a semiregular seed.
SF-E823 excludes every actual global tensor product $H\otimes M_0$
with that unit extension, arbitrary locally free positive-rank $H$,
and rational $B$ giving the balanced character. A derived partial trace
forces a nonzero cubic character obstruction to vanish. This includes
unstable and nonsplit hulls, without requiring an action on either
factor separately. SF-UG then constructs actual nontensor forms of
the unit module, including at least $5823$ independent parameters in
one explicit summand for the same positive integral balanced character.
SF-E824 excludes every global form locally isomorphic to $M_0^r$,
including all this nontensor gluing. Its canonical first overlap
obstruction has a nonzero trace; the determinant correction vanishes
after the justified algebraic $B$ normalization. No trace on the full
derived self-Hom complex and no preservation of semiregularity under
isogeny is assumed. SF-E825 treats arbitrary added locally free summands
and their off-diagonal mixing. It constructs the intrinsic defect bundle
$\widetilde Q$ on the first infinitesimal neighborhood and proves
$\Psi_{\pi_\pm}(\det\widetilde Q+rB)
=(7r/2)\Psi_{\pi_\pm}(\mathcal O_X(D))$ as a necessary condition.
An adjusted flat determinant extending from $X$ is therefore excluded.
Positive adjusted determinants satisfying this condition remain possible;
there is no universal added-free exclusion or balanced semiregular seed.
SF-E826 excludes every actual two-slope extension in its specified
unit-plus-line architecture with $\operatorname{Ext}^3(B,A)=0$
and at least four distinct line-class blocks. The actual invariant
triangle maps give $\operatorname{ext}^2_\Gamma(E,E)\ge6(u+v)>18$.
It also constructs an actual indecomposable object of full balanced
character and positive defect determinant, using a proved rectangular
matrix construction and a checked 33-vertex tree. That object still
has invariant middle Ext at least $192$ and is not a seed.
SF-E827 then removes the global cross-$\operatorname{Ext}^3$
vanishing requirement in two precise settings. Actual diagonal Yoneda
classes give $\operatorname{ext}^2_\Gamma(E,E)\ge6u$ when every
cross-line difference is nondegenerate, and at least $4u$ for every
nonzero cross difference in the rational four-divisor Gaussian algebra,
including all degenerate cases. Thus $u\ge4$ or respectively $u\ge5$
is excluded; selected negative blocks contribute under their explicit
unit-block vanishing condition. The actual equivariant filtration and
blocks are essential hypotheses. Too few such blocks and more general
divisor or filtration types remain open.
SF-E828 closes the apparent equal-hull exception: same-slope free
line blocks remain actual invariant subobjects even when one has the
unit hull's Chern class, and averaging the existing splitting gives
an actual equivariant unit complement. Thus the E826/E827 conclusions
retain their Ext and cap hypotheses but no longer require those two
Chern classes to differ. A nonsplit positive HN factor and several
distinct flat isomorphism types within one Chern-class block remain
outside this statement.
SF-E829 treats every nonzero cross-line difference of Hermitian rank
at least two. Only rank-two index-one surface directions contribute
exceptional product lines. If their span is $Z\subset\Lambda^2H^1(O_X)$,
then $w=6-\dim Z-\dim\operatorname{rad}(q_\wedge|_Z)$ diagonal
classes survive per selected block; the actual invariant middle Ext
is at least $w(u+|J|)$. The product-isogeny and all cup-product cases
are proved. Nonacyclic elliptic directions are outside that common
annihilator theorem.
SF-E830 then follows the actual extension-matrix entries through both
triangle boundaries. Its entry-specific support $Z_e$ obeys the same
wedge formula and includes elliptic differences. Along one positive
elliptic direction, actual coefficient rank $r_e$ leaves $6-2r_e$
classes per selected block. Ten or more selected blocks require
$r_e=3$; five through nine require $r_e\ge2$. The constructed
full-rank indecomposable elliptic extension has pure exponential
character and is explicitly not a balanced seed. These are necessary
support and rank conditions, with the actual filtration, action and
selected unit-block vanishings retained.
SF-E831 strengthens this to a separate support space $Z_v$ at each
actual selected block: invariant middle Ext is at least
$\sum_v(6-\dim Z_v-\dim\operatorname{rad}(q_\wedge|_{Z_v}))$.
At most three selected blocks can be untouched by actual exceptional
components. For $N>3$ selected blocks, at least
$\lceil(N-3)/2\rceil$ active exceptional pairs are required, or
$\lceil3(N-3)/2\rceil$ if every pair support has dimension at most one.
Its weighted inequality retains the distinction between one and two
selected endpoints. The zero-survivor three-spaces are classified
explicitly as $a\wedge V$ or $\Lambda^2H$, $\dim H=3$.
These incidence conditions do not assert geometric realizability.
SF-E832 and E832R exclude actual elliptic stars with either uniform
semipositive sign, for arbitrary rational algebraic $B$ and every
specified coherent effective finite action in E825's scope.
The actual locally split extension identifies the defect bundle as
$(L|_{S_2})^r$; its necessary determinant condition gives two positive
RM block traces. The completed-square second-character equation then
contradicts a universal compound-matrix lemma. For scalar adjusted
unit direction, the character obstruction already permits mixed signs.
SF-E833 then removes the uniform-sign restriction: every actual
single-center square-zero star in its unit/free extension architecture
is excluded, even for a real center which is not an actual line class.
The compound-image argument handles singular adjusted-unit blocks
without positivity. The full cubic character forces both surviving
blocks to be scalar; their equal defect traces give
$u=(7k/2)\theta$, contradicting the scalar-unit lemma.
Multiple centers, different unit characters, signed virtual
multiplicities and unproved defect determinants remain outside
this conclusion.
Rank-deficient unit matrices, other unequal-degree extensions, other
support classes, nonreduced or singular surface modules, and more general
derived couplings remain open. Twisted objects and
sufficiently nonsimple constructions retain their stated separate gates.
The full trace--evaluation composition used in SF-E819 is directly the
published Buchweitz--Flenner Theorem~6.4.2, not a theorem of 032.

The cyclic-source failure for the other quadratic fields does not refute
their Hodge conjecture. The imprimitive rank obstruction likewise
does not exclude other nonmonomial curve sources or higher-dimensional
complexes. D-NR-01--04 supplies a terminating exact procedure for the generic
rank of every specified Nielsen component, and a finite tabulation
procedure for each fixed group degree. Only its stated small matrix
example has been executed; no unbounded-degree table is claimed. The task's remaining scope is
recorded row by row in \texttt{QUESTIONS.tsv} and \texttt{LEDGER.md}.

\appendix
\section{Reconstructed eightfold matrices}
\subsection*{Exact matrices for the G34 certificate}

These are outputs of the separate exact reconstruction, not inputs copied from a printed monodromy table. The conventions and all numerical inputs are in D-G34-02.

\[ S=\begin{pmatrix}
0 & 0 & 0 & 0 & 3 & 0 & 0 & 0 \\
0 & 0 & 0 & 0 & 0 & 1 & 0 & 0 \\
0 & 0 & 0 & 0 & 0 & 0 & 1 & 0 \\
0 & 0 & 0 & 0 & 0 & 0 & 0 & 3 \\
-3 & 0 & 0 & 0 & -3-6\omega & -3-3\omega & 4+5\omega & -6-6\omega \\
0 & -1 & 0 & 0 & -3\omega & -1-2\omega & 1+4\omega & -2-4\omega \\
0 & 0 & -1 & 0 & 1+5\omega & 3+4\omega & -2-4\omega & 3+6\omega \\
0 & 0 & 0 & -3 & -6\omega & -2-4\omega & 3+6\omega & -3-6\omega
\end{pmatrix}.\]

\[ X_{1}=\begin{pmatrix}
2 & \frac{1}{3}-\frac{1}{3}\omega & \frac{1}{3}+\frac{2}{3}\omega & -\omega \\
2+\omega & 2 & 0 & 2+\omega \\
-1-2\omega & 0 & 2 & -1+\omega \\
1+\omega & \frac{1}{3}-\frac{1}{3}\omega & -\frac{2}{3}-\frac{1}{3}\omega & 2
\end{pmatrix}.\]

\[ X_{2}=\begin{pmatrix}
2 & -2\omega & -\frac{2}{3}+\frac{2}{3}\omega & 2-2\omega \\
6+6\omega & 6 & -4-2\omega & 12+6\omega \\
-4-2\omega & -2+2\omega & 2 & -6 \\
4+2\omega & 2-2\omega & -2 & 6
\end{pmatrix}.\]

\[ X_{3}=\begin{pmatrix}
0 & 0 & 0 & 0 \\
0 & 3 & -\omega & 3+3\omega \\
0 & 1+\omega & 3 & 3\omega \\
0 & -\omega & -1-\omega & 3
\end{pmatrix}.\]

\[ X_{4}=\begin{pmatrix}
0 & 0 & 0 & 0 \\
0 & 2 & 0 & 2+\omega \\
0 & 0 & 2 & -1+\omega \\
0 & \frac{1}{3}-\frac{1}{3}\omega & -\frac{2}{3}-\frac{1}{3}\omega & 2
\end{pmatrix}.\]

\[ X_{5}=\begin{pmatrix}
0 & 0 & 0 & 0 \\
0 & 0 & 0 & 0 \\
0 & 0 & 4 & -2+2\omega \\
0 & 0 & -\frac{4}{3}-\frac{2}{3}\omega & 3
\end{pmatrix}.\]

\[ X_{6}=\begin{pmatrix}
0 & 0 & 0 & 0 \\
0 & 0 & 0 & 0 \\
0 & 0 & 2 & 0 \\
0 & 0 & 0 & 0
\end{pmatrix}.\]

\[ X_{7}=\begin{pmatrix}
0 & 0 & 0 & 0 \\
0 & 0 & 0 & 0 \\
0 & 0 & 2 & -1+\omega \\
0 & 0 & -\frac{2}{3}-\frac{1}{3}\omega & 2
\end{pmatrix}.\]

\[ X_{8}=\begin{pmatrix}
0 & 0 & 0 & 0 \\
0 & 0 & 0 & 0 \\
0 & 0 & 0 & 0 \\
0 & 0 & 0 & 2
\end{pmatrix}.\]

\[ X_{9}=\begin{pmatrix}
2 & 2 & -\frac{4}{3}-\frac{2}{3}\omega & 2 \\
6 & 6 & -4-2\omega & 6 \\
-2+2\omega & -2+2\omega & 2 & -2+2\omega \\
2 & 2 & -\frac{4}{3}-\frac{2}{3}\omega & 2
\end{pmatrix}.\]

\[ X_{10}=\begin{pmatrix}
3 & 2-\omega & -2-\omega & 3 \\
9+3\omega & 8 & -5-4\omega & 9+3\omega \\
-3+3\omega & -1+4\omega & 3 & -3+3\omega \\
3 & 2-\omega & -2-\omega & 3
\end{pmatrix}.\]

\[ X_{11}=\begin{pmatrix}
5 & \frac{4}{3}-\frac{7}{3}\omega & -\frac{4}{3}+\frac{1}{3}\omega & 3-2\omega \\
11+7\omega & 8 & -5-3\omega & 11+4\omega \\
-5-\omega & -2+3\omega & 4 & -5+2\omega \\
5+2\omega & \frac{7}{3}-\frac{4}{3}\omega & -\frac{7}{3}-\frac{2}{3}\omega & 5
\end{pmatrix}.\]

\[ X_{12}=\begin{pmatrix}
7 & \frac{2}{3}-\frac{11}{3}\omega & -\frac{2}{3}+\frac{5}{3}\omega & 3-4\omega \\
13+11\omega & 9 & -5-2\omega & 13+5\omega \\
-7-5\omega & -3+2\omega & 5 & -7+\omega \\
7+4\omega & \frac{8}{3}-\frac{5}{3}\omega & -\frac{8}{3}-\frac{1}{3}\omega & 7
\end{pmatrix}.\]

\[ X_{13}=\begin{pmatrix}
5 & \frac{4}{3}-\frac{7}{3}\omega & -\frac{1}{3}+\frac{4}{3}\omega & 3-2\omega \\
11+7\omega & 8 & -3 & 11+4\omega \\
-5-4\omega & -3 & 3 & -5-\omega \\
5+2\omega & \frac{7}{3}-\frac{4}{3}\omega & -\frac{4}{3}+\frac{1}{3}\omega & 5
\end{pmatrix}.\]

\[ X_{14}=\begin{pmatrix}
2 & \frac{4}{3}-\frac{4}{3}\omega & -\frac{4}{3}-\frac{2}{3}\omega & 2 \\
8+4\omega & 8 & -4-4\omega & 8+4\omega \\
-2+2\omega & 4\omega & 2 & -2+2\omega \\
2 & \frac{4}{3}-\frac{4}{3}\omega & -\frac{4}{3}-\frac{2}{3}\omega & 2
\end{pmatrix}.\]

\[ X_{15}=\begin{pmatrix}
0 & 0 & 0 & 0 \\
0 & 2 & -2-2\omega & 4+2\omega \\
0 & 2\omega & 2 & -2+2\omega \\
0 & \frac{2}{3}-\frac{2}{3}\omega & -\frac{4}{3}-\frac{2}{3}\omega & 2
\end{pmatrix}.\]

\[ X_{16}=\begin{pmatrix}
2 & -\omega & \frac{1}{3}+\frac{2}{3}\omega & 1-\omega \\
3+3\omega & 3 & -2-\omega & 6+3\omega \\
-1-2\omega & -1+\omega & 2 & -3 \\
2+\omega & 1-\omega & -1 & 3
\end{pmatrix}.\]

\[ X_{17}=\begin{pmatrix}
2 & \frac{2}{3}-\frac{2}{3}\omega & 0 & 2 \\
4+2\omega & 2 & 0 & 4+2\omega \\
0 & 0 & 0 & 0 \\
2 & \frac{2}{3}-\frac{2}{3}\omega & 0 & 2
\end{pmatrix}.\]

The sixteen row-flattened blocks with indices $1,2,3,4,5,6,7,9,10,11,12,13,14,15,16,17$, in that order, have determinant $1728$. All exact principal minors and the auxiliary rational minor are in the accompanying JSON certificate.

\section{Reconstructed sixfold matrices}
\subsection*{Exact matrices for the G24 certificate}

These are computed outputs from the six displayed roots and the nine twist words, with exact arithmetic in $\mathbb Q(\tau)$, $\tau^2+\tau+2=0$.

\[S=\begin{pmatrix}
0 & 0 & 0 & -1-\frac{1}{2}\tau & 1-\frac{1}{2}\tau & -1-\tau \\
0 & 0 & 0 & 1 & 1+\tau & 1 \\
0 & 0 & 0 & -1-\frac{1}{2}\tau & -1-\frac{1}{2}\tau & -1 \\
\frac{1}{2}-\frac{1}{2}\tau & -1 & \frac{1}{2}-\frac{1}{2}\tau & -\frac{1}{2}-\tau & \frac{1}{2} & \frac{1}{2}-\frac{1}{2}\tau \\
-\frac{3}{2}-\frac{1}{2}\tau & \tau & \frac{1}{2}-\frac{1}{2}\tau & -\frac{1}{2} & \frac{1}{2}+\tau & \frac{1}{2}-\frac{1}{2}\tau \\
-\tau & -1 & 1 & -1-\frac{1}{2}\tau & -1-\frac{1}{2}\tau & 0
\end{pmatrix}.\]

\[X_{1}=\begin{pmatrix}
-2-\tau & 2-\tau & -2-2\tau \\
0 & 0 & 0 \\
0 & 0 & 0
\end{pmatrix}.\]

\[X_{2}=\begin{pmatrix}
0 & 0 & 0 \\
2 & 2+2\tau & 2 \\
0 & 0 & 0
\end{pmatrix}.\]

\[X_{3}=\begin{pmatrix}
0 & 0 & 0 \\
0 & 0 & 0 \\
-2-\tau & -2-\tau & -2
\end{pmatrix}.\]

\[X_{4}=\begin{pmatrix}
\tau & 0 & 0 \\
2 & 0 & 0 \\
-\tau & 0 & 0
\end{pmatrix}.\]

\[X_{5}=\begin{pmatrix}
0 & \tau & 0 \\
0 & 2+2\tau & 0 \\
0 & \tau & 0
\end{pmatrix}.\]

\[X_{6}=\begin{pmatrix}
0 & 0 & -2-2\tau \\
0 & 0 & -4 \\
0 & 0 & -2
\end{pmatrix}.\]

\[X_{7}=\begin{pmatrix}
-1 & 1 & -1-\tau \\
1+\frac{1}{2}\tau & 1+\frac{1}{2}\tau & 0 \\
0 & 0 & 0
\end{pmatrix}.\]

\[X_{8}=\begin{pmatrix}
-\frac{1}{2}+\frac{1}{4}\tau & -\frac{1}{2}-\frac{3}{4}\tau & -1-\frac{1}{2}\tau \\
1+\frac{1}{2}\tau & -1+\frac{1}{2}\tau & \tau \\
-\frac{1}{2}-\frac{1}{4}\tau & -\frac{3}{2}-\frac{1}{4}\tau & -1+\frac{1}{2}\tau
\end{pmatrix}.\]

\[X_{9}=\begin{pmatrix}
0 & 1 & -1-\frac{3}{2}\tau \\
0 & 0 & -2 \\
0 & -1 & -1+\frac{1}{2}\tau
\end{pmatrix}.\]


\subsection*{Exact matrices for the F4 extension certificate}

These matrices are independently reconstructed from the five displayed reflection roots, the rank-two involution $z$, and the nine interval twists over $\mathbb Q(t)$, $t^2=-2$.

\[S=\begin{pmatrix}
0 & 0 & 0 & -1 & 1+\frac{1}{2}t & \frac{1}{2}t \\
0 & 0 & 0 & 1 & -2 & -1-t \\
0 & 0 & 0 & 2 & -1 & -1-t \\
1 & -1 & -2 & 0 & -1-\frac{1}{2}t & -\frac{1}{2}t \\
-1+\frac{1}{2}t & 2 & 1 & 1-\frac{1}{2}t & t & -1-\frac{1}{2}t \\
\frac{1}{2}t & 1-t & 1-t & -\frac{1}{2}t & 1-\frac{1}{2}t & 0
\end{pmatrix}.\]

\[X_{1}=\begin{pmatrix}
-2 & 2+t & t \\
0 & 0 & 0 \\
0 & 0 & 0
\end{pmatrix}.\]

\[X_{2}=\begin{pmatrix}
0 & 0 & 0 \\
2 & -4 & -2-2t \\
0 & 0 & 0
\end{pmatrix}.\]

\[X_{3}=\begin{pmatrix}
0 & 0 & 0 \\
0 & 0 & 0 \\
2 & -1 & -1-t
\end{pmatrix}.\]

\[X_{4}=\begin{pmatrix}
0 & 2t & -2 \\
0 & -\frac{4}{3}+\frac{2}{3}t & -\frac{2}{3}-\frac{2}{3}t \\
0 & \frac{2}{3}+\frac{2}{3}t & -\frac{2}{3}+\frac{1}{3}t
\end{pmatrix}.\]

\[X_{5}=\begin{pmatrix}
-2 & 0 & 0 \\
-\frac{4}{3}+\frac{2}{3}t & 0 & 0 \\
\frac{2}{3}-\frac{1}{3}t & 0 & 0
\end{pmatrix}.\]

\[X_{6}=\begin{pmatrix}
0 & t & -2-t \\
0 & -\frac{2}{3}+\frac{1}{3}t & -\frac{4}{3}-\frac{1}{3}t \\
0 & \frac{1}{3}+\frac{1}{3}t & -\frac{1}{3}-\frac{1}{3}t
\end{pmatrix}.\]

\[X_{7}=\begin{pmatrix}
-1 & t & -1 \\
1 & -3+\frac{1}{2}t & -2-\frac{3}{2}t \\
0 & 0 & 0
\end{pmatrix}.\]

\[X_{8}=\begin{pmatrix}
-t & \frac{5}{2}t & -3+\frac{1}{2}t \\
\frac{2}{3}-\frac{1}{3}t & -\frac{5}{3}+\frac{5}{6}t & -\frac{4}{3}-\frac{5}{6}t \\
\frac{5}{3}-\frac{1}{3}t & -\frac{7}{6}+\frac{5}{6}t & -\frac{11}{6}-\frac{5}{6}t
\end{pmatrix}.\]

\[X_{9}=\begin{pmatrix}
-1 & \frac{3}{2}t & -2-\frac{1}{2}t \\
-1 & -\frac{1}{3}+\frac{2}{3}t & -1 \\
1 & \frac{1}{6}+\frac{1}{6}t & -\frac{1}{2}-\frac{1}{2}t
\end{pmatrix}.\]


\section{Reproduction and check status}
The scripts in \texttt{scripts/} and the agent directories contain every
numerical input, finite-field minor and exact group computation used
above. Source logs give primary URLs and locators, with hashes for
downloaded witnesses. Third-party full texts are not part of this
deliverable. The run did not compile the supplied Alp\"oge Lean imports;
no RH or GRH assertion is needed by any new proof here. No numbered
Masullo general-GRH theorem is invented from the supplied index.

\begin{thebibliography}{99}
% Include these inside the root thebibliography; keys are namespace-prefixed.
\bibitem{SF-Schoen1988} C.~Schoen, \emph{Hodge classes on self-products of a variety with an automorphism}, Compositio Math.~65 (1988), 3--32. Unramified Theorem~2.0, pp.~12--13; Lemma~2.6; \S3, pp.~24--25. \url{https://www.numdam.org/article/CM_1988__65_1_3_0.pdf}.
\bibitem{SF-Schoen1998} C.~Schoen, \emph{Addendum to: Hodge classes on self-products of a variety with an automorphism}, Compositio Math.~114 (1998), 329--336. \S\S6--7, p.~332; Lemma~13 and equations~(13.1)--(13.2), pp.~334--335. \url{https://doi.org/10.1023/A:1000566205021}.
\bibitem{SF-Koike2004} K.~Koike, \emph{Algebraicity of some Weil Hodge Classes}, Canad. Math. Bull.~47(4) (2004), 566--572. Theorem~2.1, Proposition~3.1, Corollary~2.1; equations~(1)--(2), p.~568, and (6)--(9), pp.~571--572. \url{https://doi.org/10.4153/CMB-2004-055-x}.
\bibitem{SF-vanGeemen1996} B.~van Geemen, \emph{Theta functions and cycles on some abelian fourfolds}, Math. Z.~222 (1996), 617--631. Theorem~5.3 and Remark~5.4, pp.~629--630.
\bibitem{SF-LangeOrtega2011} H.~Lange and A.~Ortega, \emph{Prym varieties of cyclic coverings}, Geom. Dedicata~150 (2011), 391--403. \S4, induced polarization-type paragraph; author version p.~6, citing Birkenhake--Lange Corollary~12.1.4 and Lemma~12.3.1. \url{https://arxiv.org/abs/0805.1020}.
\bibitem{SF-Mumford1974} D.~Mumford, \emph{Prym Varieties I}, in \emph{Contributions to Analysis}, Academic Press (1974), 325--350. \S3, Corollary~2 immediately after Corollary~1 on p.~332.
\bibitem{SF-Fulton} W.~Fulton, \emph{Intersection Theory}, second edition, Springer (1998), Example~8.4.6, component Bezout inequality.
\bibitem{SF-FarbMargalit} B.~Farb and D.~Margalit, \emph{A Primer on Mapping Class Groups}, Princeton University Press (2012), Theorem~6.4.
\bibitem{SF-Stein} W.~Stein, \emph{Algebraic Number Theory}, author-hosted notes, subsection ``Class Number 1,'' explicit list of the nine negative fundamental discriminants. \url{https://wstein.org/books/ant/ant/node46.html}.
\bibitem{SF-Markman2502} E.~Markman, \emph{Cycles on abelian $2n$-folds of Weil type from secant sheaves on abelian $n$-folds}, arXiv:2502.03415v2, revision 8 June 2025, \textbf{preprint; uses here are conditional}. Theorem~1.4.1, pp.~7--8; Theorem~1.5.1, p.~9; Lemma~3.1.3, pp.~23--24; Lemmas~9.3.1, 9.3.5, 9.3.11, pp.~81, 84, 88; Remark~9.3.7, p.~86. \url{https://arxiv.org/pdf/2502.03415v2}.
\bibitem{SF-Perry2604} A.~Perry, \emph{The semiregularity theorem for equivariant noncommutative varieties}, arXiv:2604.00511v2, revision 24 June 2026, \textbf{preprint}. Theorem~1.2, p.~3; Definition~2.6 and Lemma~2.7, pp.~5--6; \S1.1, pp.~3--4. Used through the authorized DEG\_L S-14 input. \url{https://arxiv.org/pdf/2604.00511v2}.

\bibitem{G34:Davis}
D.~M. Davis, \emph{Homotopy type and \(v_1\)-periodic homotopy groups of
\(p\)-compact groups}, Topology and its Applications \textbf{156} (2008),
300--321, DOI:10.1016/j.topol.2008.07.008.
Exact locator: Theorem~5.1 and its proof, pp.~312--313, equation~(5.2),
and the \(270+486\) root description on p.~313.
Author's published copy: \url{https://www.lehigh.edu/~dmd1/TOPOL3486.pdf}.

\bibitem{G34:FGP}
P. Frediani, A. Ghigi, M. Penegini, \emph{Shimura varieties in the Torelli
locus via Galois coverings}, International Mathematics Research Notices
\textbf{2015}(20), 10595--10623, DOI:10.1093/imrn/rnu272.
Exact locator: Theorem~2.10 (Chevalley--Weil); author version p.~8.
\url{https://arxiv.org/pdf/1402.0973}.

\bibitem{G34:DWmon}
M. Dettweiler, S. Wewers, \emph{Variation of local systems and parabolic
cohomology}, Israel Journal of Mathematics \textbf{156} (2006), 157--185.
Exact locators: Proposition~1.1, Lemma~2.2, \S2.2, Theorem~2.5;
author version pp.~4--7 and 8--12.
\url{https://arxiv.org/pdf/math/0310139}.

\bibitem{G34:DWdual}
M. Dettweiler, S. Wewers, \emph{Variation of parabolic cohomology and
Poincar\'e duality}, S\'eminaires et Congr\`es \textbf{13} (2006), 145--164.
Exact locator: Theorem~2.5 (author version p.~12).
\url{https://arxiv.org/pdf/math/0411119};
publisher record \url{https://smf.emath.fr/node/43624}.

\bibitem{G34:Schmid}
W. Schmid, \emph{Variation of Hodge Structure: The Singularities of the
Period Mapping}, Inventiones Mathematicae \textbf{22} (1973), 211--319.
Exact locators: Theorem~(4.12), p.~233; Theorem~(6.16), p.~255.
\url{https://webhomes.maths.ed.ac.uk/~v1ranick/papers/schmid.pdf}.

\bibitem{G34:Schoen}
C. Schoen, \emph{Hodge classes on self-products of a variety with an
automorphism}, Compositio Mathematica \textbf{65} (1988), 3--32.
Exact locators: simple-exponent definition and Theorem~2.0, p.~11;
Lemmas~2.2--2.6, pp.~13--17; Lemma~2.8 and Sublemma~2.9, pp.~18--19;
Corollary~3.1, pp.~24--25.  The required prime-quotient step in the printed
proof of Lemma~2.8 is replaced by D-G34-04 above.
\url{https://www.numdam.org/item/CM_1988__65_1_3_0.pdf}.

\bibitem{G34:Serre}
J.-P. Serre, \emph{Linear Representations of Finite Groups}, Graduate Texts
in Mathematics~42, Springer, 1977.  Exact locator: \S10.5, Theorems~19--20,
pp.~78--79.  DOI:10.1007/978-1-4684-9458-7.

\bibitem{G34:Milne}
J.~S. Milne, \emph{Lefschetz classes on abelian varieties}, Duke Mathematical
Journal \textbf{96} (1999), 639--675.  Exact locators: Lemma~3.1,
Theorem~3.2, Lemma~3.5, Proposition~3.6 and Remark~3.7;
author version pp.~14--18.
\url{https://www.jmilne.org/math/articles/1999aP.pdf}.

\bibitem{G34:SGA4}
M. Artin, A. Grothendieck, J.-L. Verdier (eds.), \emph{Th\'eorie des topos
et cohomologie \'etale des sch\'emas (SGA4)}, Tome~3, Lecture Notes in
Mathematics~305, Springer, 1973.  Exact locators: Expos\'e~XII,
Th\'eor\`eme~5.1(iii) (proper base change), and Expos\'e~XVI,
Th\'eor\`eme~4.1 (classical comparison).  Locators also cross-checked
against L.~Illusie's author-hosted review, Theorem~1.5,
\url{https://www.imo.universite-paris-saclay.fr/~luc.illusie/wolfson2.pdf},
and A.~Beilinson, J.~Bernstein, P.~Deligne, \emph{Faisceaux pervers},
Ast\'erisque~100 (1982), p.~150.

\bibitem{G34:Severn}
B. Severn, \emph{Weil classes on abelian sixfolds and eightfolds of Weil type
for \(\mathbb Q(\sqrt{-3})\)}, Hexagon~2610.00387v2, internal version~1.5,
21 October 2026, \textbf{unrefereed}.  Numerical-data provenance only:
Appendix~A.  No theorem of this manuscript is a premise in the new proof.
Its Theorems~A--C, if used as inputs, remain conditional.  The portal heading
has a different, sixfold-only title; the actual downloaded PDF and matching
TeX source have the title given here.
\url{https://hexagonmath.org/pdf/2610.00387v2};
\url{https://hexagonmath.org/src/2610.00387v2/files}.

\bibitem{G34:Hecke}
\emph{HECKE\_DIV}, supplied project run, final package
\texttt{RETURN\_HECKE\_DIV\_final\_20261121T0204Z},
file \texttt{HECKE\_DIV.tex}, statements SQ1--SQ7.
Full quantitative proofs rechecked in the present HODGE\_DIVISION run;
their previously conditional geometric premises are replaced here by
D-G34-01--06.  Check events record the actual readers and check types.

\bibitem{G24:Deraux}
Martin Deraux, \emph{Non-arithmetic lattices and the Klein quartic},
Journal f\"ur die reine und angewandte Mathematik \textbf{754} (2019),
253--279. DOI: \url{https://doi.org/10.1515/crelle-2017-0005}.
Author manuscript dated 17 January 2017:
\url{https://www-fourier.univ-grenoble-alpes.fr/~deraux/recherche/Klein.pdf}.
Locators: \S3, pp.~4--6, especially Proposition 3.1 and Remark 3.1;
equation (13), p.~24, specialized to \(p=2\).

\bibitem{GF4:Cahn}
Patricia Cahn, Ruth Haas, Aloysius G. Helminck, Juan Li, and Jeremy Schwartz,
\emph{Permutation notations for the exceptional Weyl group \(F_4\)},
Involve \textbf{5}, no.~1 (2012), 81--90.
\url{https://msp.org/involve/2012/5-1/involve-v5-n1-p09-s.pdf}.
Locator: \S3, p.~83, standard simple roots and the full 48-root system.

\bibitem{FIELD:Cohen}
A.~M.~Cohen, \emph{Finite complex reflection groups},
Ann. Sci. \(\acute E\)cole Norm. Sup. (4) \textbf{9} (1976), 379--436.
\url{https://www.numdam.org/article/ASENS_1976_4_9_3_379_0.pdf}.

\bibitem{FIELD:CohenErr}
A.~M.~Cohen, \emph{Erratum: ``Finite complex reflection groups''},
Ann. Sci. \(\acute E\)cole Norm. Sup. (4) \textbf{11} (1978), 613.
\url{https://www.numdam.org/article/ASENS_1978_4_11_4_613_0.pdf}.

\bibitem{FIELD:Malle}
G.~Malle, \emph{On the rationality and fake degrees of characters of
cyclotomic algebras}, J. Math. Sci. Univ. Tokyo \textbf{6} (1999), 647--677.
\url{https://www.ms.u-tokyo.ac.jp/journal/pdf/jms060402.pdf}.

\bibitem{WG:Carter}
R.~W. Carter, \emph{Conjugacy classes in the Weyl group}, Compositio
Mathematica \textbf{25} (1972), 1--59.
Exact locators: Table~1, p.~48 (root systems); Table~9, p.~50 (\(E_6\));
Table~10, pp.~51--52 (\(E_7\)); Table~11, pp.~52--54 (\(E_8\)).
\url{https://www.numdam.org/item/CM_1972__25_1_1_0.pdf}.

\bibitem{WG:CHEVIE}
J. Michel, \emph{The development version of the CHEVIE package of GAP3},
Journal of Algebra \textbf{435} (2015), 308--336.
Author-maintained data by M. Geck, J. Michel and the CHEVIE team:
\texttt{pkg/chevie/tbl/weyle6.g}, \texttt{weyle7.g}, \texttt{weyle8.g},
\url{https://github.com/jmichel7/gap3-jm}.
Exact documentation: GAP3 manual \S87.1, \texttt{ChevieClassInfo},
\url{https://webusers.imj-prg.fr/~jean.michel/gap3/htm/chap087.htm}.
The source hashes and compact mathematical witness data are retained
with the exact certificates.

% Bibitems used by the imprimitive and semiregular CM sections.
% Include inside the main paper's thebibliography environment.
% Native URLs and exact consulted locators are also in SOURCES.md.

\bibitem{IM-LehrerTaylor}
G.~I. Lehrer and D.~E. Taylor,
\emph{Unitary Reflection Groups}, Cambridge University Press, 2009,
Chapter~2, \S\S1--3, pp.~23--27, especially Proposition~2.10;
authors' errata, 5 February 2025,
\url{https://www.maths.usyd.edu.au/u/don/papers/errata-uggr.pdf}.

\bibitem{IM-LewisWang}
J.~B. Lewis and J.~Wang,
The Hurwitz action in complex reflection groups,
\emph{Combinatorial Theory} \textbf{2}(1) (2022), \S2.1,
\url{https://escholarship.org/uc/item/4xg2m191}.

\bibitem{IM-LangeOrtega}
H.~Lange and A.~Ortega,
Prym varieties of cyclic coverings,
\emph{Geometriae Dedicata} \textbf{150} (2011), 391--403,
doi:10.1007/s10711-010-9512-9.
Proposition~4.1 and proof, author version pp.~7--10,
\url{https://arxiv.org/abs/0805.1020}.

\bibitem{IM-FarbMargalit}
B.~Farb and D.~Margalit,
\emph{A Primer on Mapping Class Groups}, Princeton University Press, 2012.
Proposition~6.2, p.~166; Theorem~6.4, p.~170;
Theorem~4.6 and Fact~4.7, pp.~97--98.
Authors' publication page:
\url{https://www.math.uchicago.edu/~farb/books.html}.

\bibitem{IM-StacksDeform}
The Stacks Project authors,
\emph{The Stacks Project}, Deformation Theory, \S93.10,
Lemma~93.10.3, tag~0E3V,
\url{https://stacks.math.columbia.edu/tag/0E3V}.
Used as background only; the morphism-deformation argument is proved in
D-IM-07.

\bibitem{SR-BuchweitzFlenner2003}
R.-O. Buchweitz and H.~Flenner,
A semiregularity map for modules and applications to deformations,
\emph{Compositio Mathematica} \textbf{137} (2003), 135--210,
doi:10.1023/A:1023999012081.
\S4.1 and Proposition~4.2, pp.~166--167.

\bibitem{SR-BuchweitzFlenner2008}
R.-O. Buchweitz and H.~Flenner,
The global decomposition theorem for Hochschild (co-)homology of singular
spaces via the Atiyah--Chern character,
\emph{Advances in Mathematics} \textbf{217} (2008), 243--281,
doi:10.1016/j.aim.2007.06.013.
Proposition~6.2.1 and Theorem~6.4.2 in the author version,
pp.~29 and~31,
\url{https://arxiv.org/abs/math/0606730}.

\bibitem{SR-Mukai1981}
S.~Mukai,
Duality between $D(X)$ and $D(\widehat X)$ with its application to
Picard sheaves,
\emph{Nagoya Mathematical Journal} \textbf{81} (1981), 153--175,
doi:10.1017/S002776300001922X.
Theorem~2.2, p.~156, and equation~(3.1), p.~158.

\bibitem{SR-Beauville2014}
A.~Beauville,
Some surfaces with maximal Picard number,
\emph{Journal de l'\'Ecole polytechnique---Math\'ematiques}
\textbf{1} (2014), 101--116,
doi:10.5802/jep.5.
\S4, Lemma~2, p.~106, and Example~3, p.~107,
\url{https://www.numdam.org/article/JEP_2014__1__101_0.pdf}.

\bibitem{SR-Perry2026}
A.~Perry,
\emph{The semiregularity theorem for equivariant noncommutative varieties},
arXiv:2604.00511v2, 24 June 2026 (preprint).
Theorem~1.2, p.~3; Remarks~2.3 and~2.5, p.~5;
Definition~2.6 and Lemma~2.7, pp.~5--6,
\url{https://arxiv.org/abs/2604.00511}.

\bibitem{SR-BirkenhakeLange}
C.~Birkenhake and H.~Lange,
\emph{Complex Abelian Varieties}, second edition, Springer, 2004,
Proposition~11.2.1 (Poincar\'e formula),
\url{https://link.springer.com/book/10.1007/978-3-662-06307-1}.
The proposition is a classical input; the original book text was not
directly retrieved in this run. An independently consulted primary use
is Schoen~\cite{SR-Schoen1988}, p.~25.

\bibitem{SR-BauerSzemberg}
T.~Bauer and T.~Szemberg,
On tensor products of ample line bundles on abelian varieties,
\emph{Mathematische Zeitschrift} \textbf{223} (1996), 79--85,
Theorem~1.1(b).
The author paper's theorem and proof were consulted; three equal ample
factors give the very ampleness of $3\Theta$ used in D-SR-01.

\bibitem{SR-Schoen1988}
C.~Schoen,
Hodge classes on self-products of a variety with an automorphism,
\emph{Compositio Mathematica} \textbf{65} (1988), 3--32.
The simple-branch definition and Theorem~2.0, p.~11;
Corollary~3.1 and proof, pp.~24--25,
\url{https://www.numdam.org/item/CM_1988__65_1_3_0.pdf}.

\bibitem{IM:Looijenga1997}
E.~Looijenga,
Prym representations of mapping class groups,
\emph{Geometriae Dedicata} \textbf{64} (1997), 69--83,
doi:10.1023/A:1004909416648.
Proposition~2.2 and Theorem~2.4, author version p.~5;
the hyperbolic-basis proof, pp.~8--9, and monodromy proof for
genus at least three, pp.~10--12.
Author's PostScript: \url{https://webspace.science.uu.nl/~looij101/pr2.ps}.

\bibitem{NR:RW}
M.~Romagny and S.~Wewers,
\emph{Hurwitz spaces}, S\'eminaires et Congr\`es \textbf{13} (2006),
313--341.  Theorem~3.1, Propositions~3.2 and~3.4, pp.~319--321;
Theorem~4.3, pp.~323--327; Theorem~4.11 and Remark~4.15(i),
pp.~328--331.  Published full text: \url{https://emis.de/ft/47206}.


\bibitem{IM:IlyashenkoYakovenko}
Y.~Ilyashenko and S.~Yakovenko,
\emph{Lectures on Analytic Differential Equations},
Graduate Studies in Mathematics 86, AMS, 2008.
Author-posted extract of Section 16, pp.~265--274;
Theorems 16.15--16.16 and Corollaries 16.19--16.20.
\url{https://yakovenko.wordpress.com/wp-content/uploads/2008/03/sect-16-from-thebook.pdf}.

\bibitem{IM:FredianiGhigiPenegini}
P.~Frediani, A.~Ghigi and M.~Penegini,
Shimura varieties in the Torelli locus via Galois coverings,
\emph{Int. Math. Res. Not. IMRN} 2015, 10595--10623.
Theorem 2.10; author manuscript arXiv:1402.0973v3, p.~8.
\url{https://arxiv.org/pdf/1402.0973}.

\bibitem{IM:JohanssonArb}
F.~Johansson,
Arb: efficient arbitrary-precision midpoint-radius interval arithmetic,
\emph{IEEE Transactions on Computers} 66 (2017), no.~8, 1281--1292.
DOI: 10.1109/TC.2017.2690633.
\url{https://arxiv.org/abs/1611.02831}.

\bibitem{E8GRR:Tseng}
H.-H.~Tseng, \emph{Orbifold quantum Riemann--Roch, Lefschetz and Serre},
Geometry \& Topology \textbf{14} (2010), 1--81,
doi:10.2140/gt.2010.14.1.
\url{https://msp.org/gt/2010/14-1/gt-v14-n1-p01-p.pdf}.
The degree-one specialization used here is from \S7.3,
(7.3.6.1), pp.~60--61, and (7.3.7.2), p.~64;
marking-gerbe conventions are Remark2.4.2, p.~19, and \S2.5.1, p.~21.

\bibitem{E8GRR:BBT}
B.~Bakker, Y.~Brunebarbe and J.~Tsimerman,
\emph{o-minimal GAGA and a conjecture of Griffiths},
Inventiones Mathematicae \textbf{232} (2023), 163--228,
doi:10.1007/s00222-022-01166-1.
The 48-page author version has Lemmas6.15 and6.17 and their
proofs on pp.~42--44:
\url{https://benjamin-bakker.github.io/griffiths.pdf}.

\bibitem{E8GRR:CD}
F.~Catanese and M.~Dettweiler,
\emph{The direct image of the relative dualizing sheaf needs not be semiample},
C.~R.~Acad.~Sci.~Paris, S\'er.~I \textbf{352} (2014), 241--244,
doi:10.1016/j.crma.2013.12.015.
The canonical-extension identification is stated in the proof of
Theorem2.1, p.~242.
\url{https://www.numdam.org/item/CRMATH_2014__352_3_241_0.pdf}.

\bibitem{E8GRR:Michel}
J.~Michel, \emph{The development version of the CHEVIE package of GAP3},
Journal of Algebra \textbf{435} (2015), 308--336,
doi:10.1016/j.jalgebra.2015.03.031.

\bibitem{E8GRR:CHEVIE}
J.~Michel and the CHEVIE authors,
\emph{GAP3 CHEVIE manual}, Chapter81 (package and library scope),
Chapter87 (class data), and the author-maintained table
\texttt{pkg/chevie/tbl/weyle8.g}.
\url{https://webusers.imj-prg.fr/~jean.michel/gap3/htm/chap081.htm};
\url{https://webusers.imj-prg.fr/~jean.michel/gap3/htm/chap087.htm};
\url{https://raw.githubusercontent.com/jmichel7/gap3-jm/master/pkg/chevie/tbl/weyle8.g}.
The retrieved table's full hash is recorded in D-E8-GRR-04 and in
\texttt{character\_input.json}.

\bibitem{E8GRR:Carter}
R.~W.~Carter, \emph{Conjugacy classes in the Weyl group},
Compositio Mathematica \textbf{25} (1972), 1--59;
Table11, pp.~52--54.
\url{https://www.numdam.org/item/CM_1972__25_1_1_0.pdf}.

\bibitem{E8GRR:lead66}
cm\_symmetry\_scout and lelong\_ridge,
\emph{Branch-exponent determinant sources and a dominant nonsplit
Gaussian sixfold family}, rStation paper66, version1 (28 October2026),
\S4.  Unrefereed AI-authored manuscript; used for attribution of
the proposed route and numerical target only, not as a theorem premise.
\url{https://rstation.ai/problems/hodge_conjecture/papers/66}.

\bibitem{KoblitzRohrlich1978}
N.~Koblitz and D.~Rohrlich,
\emph{Simple Factors in the Jacobian of a Fermat Curve},
Canadian Journal of Mathematics \textbf{30} (1978), no.~6, 1183--1205.
doi:10.4153/CJM-1978-099-6.
Relevant published locator: introductory lattice decomposition and CM-type
discussion, pp.~1183--1184. The later prime-to-six theorems are not used at
level twelve.

\bibitem{ShiodaKatsura1979}
T.~Shioda and T.~Katsura,
\emph{On Fermat varieties},
Tohoku Mathematical Journal, Second Series \textbf{31} (1979), no.~1, 97--115.
doi:10.2748/tmj/1178229881.
Relevant locators: Lemmas~1.1--1.2, pp.~98--99;
Theorem~1.7, pp.~101--102; Lemma~2.1 and Proposition~2.4, pp.~103--104.

\bibitem{Truong2020}
T.~T.~Truong,
\emph{Relative dynamical degrees of correspondences over a field of
arbitrary characteristic},
Journal f\"ur die reine und angewandte Mathematik
\textbf{758} (2020), 139--182.
doi:10.1515/crelle-2017-0052.
Quantitative intersection input: Lemma~2.2 and its full proof,
author version arXiv:1605.05049, p.~13.
The exact printed lemma pagination was not independently verified.

\bibitem{Siu1987}
Y.-T. Siu,
\emph{Lectures on Hermitian--Einstein metrics for stable bundles and
K\"ahler--Einstein metrics}, DMV Seminar~8, Birkh\"auser, 1987.
Author reprint: \url{https://people.math.harvard.edu/~siu/siu_reprints/dmv_book.pdf}.
Chapter~1, \S1.4(ii),(iii), author-reprint p.~7;
Proposition~1.6, pp.~8--9; \S6.1--\S6.4, pp.~35--39.
The existence input is used only for rational polarization classes.

\bibitem{MumfordAbelian1970}
D. Mumford, \emph{Abelian varieties}, Tata Institute of Fundamental
Research Studies in Mathematics~5, Oxford University Press, 1970.
Section~16, Vanishing Theorem, p.~150, proof pp.~150--153;
complex Hermitian-index corollary and proof, pp.~162--163.

\bibitem{SF-Fulton1998Support}
W.~Fulton, \emph{Intersection Theory}, second edition,
Ergebnisse der Mathematik und ihrer Grenzgebiete, 3.~Folge, vol.~2,
Springer, 1998, Theorem~18.3(3), p.~353, and Example~18.3.11, p.~363.
The top-term proof is Step~8 on p.~359.
\bibitem{SF-StacksSupportCycle}
The Stacks Project Authors, \emph{The Stacks Project},
Definition~42.10.2 and Lemma~42.10.4,
\url{https://stacks.math.columbia.edu/tag/02QV}, accessed 21 November 2026.

% New primary inputs for SF-E834; upstream/source-read scope is recorded in
% rm_reflexive_mixed_curvature_draft.md and SOURCE_LOG.md.
\bibitem{SFBS1994}
S. Bando and Y.-T. Siu,
\emph{Stable sheaves and Einstein-Hermitian metrics},
in \emph{Geometry and Analysis on Complex Manifolds},
World Scientific, 1994, pp.~39--50.
Theorem~3, p.~40; admissibility, p.~39; regularization and complete
invoked \S4 proof, pp.~45--49.
Author's original reprint images:
\url{https://people.math.harvard.edu/~siu/bando_joint_paper/index.html}.

\bibitem{SFBiswas2011}
I. Biswas,
\emph{A criterion for holomorphic extension of principal bundles},
Kodai Math. J. \textbf{34} (2011), 71--78.
Theorem~3.2 and proof, p.~77; also Theorem~1.2, p.~72.
\url{https://www.jstage.jst.go.jp/article/kodaimath/34/1/34_1_71/_pdf/-char/en}.

\bibitem{SFHarrisTonegawa2000}
A. Harris and Y. Tonegawa,
\emph{Analytic continuation of vector bundles with $L^p$-curvature},
Internat. J. Math. \textbf{11} (2000), no.~1, 29--40.
Theorem~1, p.~29, identified in the complete proof of
Biswas, Theorem~3.2. DOI: \url{https://doi.org/10.1142/S0129167X00000040}.
The original Harris--Tonegawa proof was not independently retrieved in
this run; the invoked published extension theorem is Biswas's precise
statement above.

\bibitem{SFStacksPointExtension}
The Stacks Project Authors,
Lemma~36.6.9, tag~0G7P (local cohomology for a regular sequence),
\S20.21, tag~0A39 (support triangle and local-to-global support sequence),
and Lemma~47.11.1, tag~0AVZ (depth and local cohomology).
Complete invoked proofs read on 21 November 2026.
\url{https://stacks.math.columbia.edu/tag/0G7P};
\url{https://stacks.math.columbia.edu/tag/0A39};
\url{https://stacks.math.columbia.edu/tag/0AVZ}.

\bibitem{SFEisenbudUlrich2019}
D. Eisenbud and B. Ulrich,
\emph{Duality and socle generators for residual intersections},
J. reine angew. Math. \textbf{756} (2019), 183--226.
DOI: \url{https://doi.org/10.1515/crelle-2017-0045}.
Theorem~7.1, p.~210, for complete local algebras in characteristic zero;
Appendix~A, Theorem~A.5 and proof, pp.~224--225, with Theorem~A.1,
pp.~222--223, and Theorem~A.3, p.~224.
Published author-hosted copy:
\url{https://www.math.purdue.edu/~bulrich/papers/published%20version%20of%20Duality%20for%20Residual%20Intersections.pdf}.

\bibitem{SFStacksQuadraticJets}
The Stacks Project Authors,
\emph{The Stacks Project}, accessed 21 November 2026.
Lemma 10.106.3, tag \texttt{00NQ}, regular local rings are Cohen--Macaulay;
Lemma 10.104.2, tag \texttt{02JN}, and its complete proof in
Proposition 10.103.4, tag \texttt{00N6}, on parameter ideals and regular
sequences; Section 27.22, tag \texttt{089R}, Lemma 27.22.1, complete Grassmannian
chart construction.
\url{https://stacks.math.columbia.edu/tag/00NQ};
\url{https://stacks.math.columbia.edu/tag/02JN};
\url{https://stacks.math.columbia.edu/tag/00N6};
\url{https://stacks.math.columbia.edu/tag/089R}.

\bibitem{SF-BF2008}
R.-O.~Buchweitz and H.~Flenner,
\emph{The global decomposition theorem for Hochschild (co-)homology
of singular spaces via the Atiyah--Chern character},
Advances in Mathematics \textbf{217} (2008), no.~1, 243--281.
DOI: 10.1016/j.aim.2007.06.013.
Invoked: Proposition~6.2.1, Proposition~6.3.1, Corollary~6.3.2,
Theorem~6.4.2 and Proposition~6.4.4; author version
\url{https://arxiv.org/pdf/math/0606730}, pp.~29--33.

\bibitem{SF-StacksPerfect}
The Stacks Project Authors, \emph{The Stacks Project},
Lemma~15.76.2, tag~0658, and Lemma~15.66.5, tag~064U,
\url{https://stacks.math.columbia.edu/tag/0658} and
\url{https://stacks.math.columbia.edu/tag/064U}.
Accessed 21 November 2026.

\bibitem{SFLipman1984}
J.~Lipman,
\emph{Dualizing sheaves, differentials and residues on algebraic varieties},
Ast\'erisque \textbf{117} (1984), Soci\'et\'e Math\'ematique de France.
Invoked: Lemma~7.3 and Theorem~7.4, with complete proofs, pp.~64--69;
Theorem~9.1 and Corollary~9.2, with complete proofs and comparison
diagrams, pp.~82--85.
Published primary PDF:
\url{https://smf.emath.fr/system/files/filepdf/AS-117.pdf}.

\bibitem{SFStacksE846Duality}
The Stacks Project Authors, \emph{The Stacks Project}.
Section~48.27, tag~0FVU: Lemma~48.27.1(5),(7),
Remarks~48.27.2--3 and Lemma~48.27.4;
\url{https://stacks.math.columbia.edu/tag/0FVU}.
For the stated local-duality and trace interfaces also see
Theorem~47.18.3, tag~0A81,
\url{https://stacks.math.columbia.edu/tag/0A81}, and
Section~48.7, tag~0AWG,
\url{https://stacks.math.columbia.edu/tag/0AWG}.
The complete invoked statements and proofs were read on 22 November 2026.

\bibitem{SF-StacksRegularDimension}
The Stacks Project Authors, \emph{The Stacks Project},
Proposition~10.110.1, tag~00O7, and its displayed finite-resolution
dependencies Lemmas~10.106.6 and~10.104.9,
\url{https://stacks.math.columbia.edu/tag/00O7},
\url{https://stacks.math.columbia.edu/tag/00NT},
\url{https://stacks.math.columbia.edu/tag/00NG}.
Accessed 21 November 2026.

\bibitem{SF-StacksBertini}
The Stacks Project Authors, \emph{The Stacks Project},
Lemma~33.47.3 and Remark~37.32.4,
\url{https://stacks.math.columbia.edu/tag/0FD4},
\url{https://stacks.math.columbia.edu/tag/0G4G}.
Accessed 21 November 2026.

\bibitem{SF-BF2008Evaluation}
R.-O. Buchweitz and H. Flenner,
\emph{The global decomposition theorem for Hochschild (co-)homology
of singular spaces via the Atiyah--Chern character},
Advances in Mathematics \textbf{217} (2008), 243--281,
doi:10.1016/j.aim.2007.06.013.
For the evaluation formula: Theorem~5.1.3, Proposition~5.2.2 and
Corollary~5.2.3, author-copy pp.~24--27; the signed exponential is
in Section~4.2, p.~22.
Author copy: \url{https://arxiv.org/pdf/math/0606730}.

\bibitem{E8LAT:Landherr}
W.~Landherr, \emph{\"Aquivalenz Hermitescher Formen \"uber einem
beliebigen algebraischen Zahlk\"orper},
Abhandlungen aus dem Mathematischen Seminar der Universit\"at
Hamburg \textbf{11} (1935), 245--248,
doi:10.1007/BF02940728.
The invariants are given on p.~246; their completeness is proved
in the first paragraph of p.~247.
\url{https://webhomes.maths.ed.ac.uk/~v1ranick/papers/landherr.pdf}.

\bibitem{BI:GS} R.~M. Guralnick and J.~Saxl,
\emph{Generation of finite almost simple groups by conjugates},
J. Algebra \textbf{268} (2003), 519--571,
doi:10.1016/S0021-8693(03)00182-0.
Theorem~7.1 and Remark~(3), pp.~548--549; the tensor convention is
stated immediately before Theorem~7.1.

\bibitem{BI:GT} R.~M. Guralnick and P.~H. Tiep,
\emph{Decompositions of small tensor powers and Larsen's conjecture},
Represent. Theory \textbf{9} (2005), 138--208,
doi:10.1090/S1088-4165-05-00192-5.
Table~II and Proposition~7.10.  Author version
\texttt{arXiv:math/0502080}, Table~II p.~7.

\bibitem{BI:AtlasU52} R.~A. Wilson and collaborators,
\emph{ATLAS of Finite Group Representations},
$U_5(2)$ page, characteristic-zero representation
\texttt{U52G1-Ar10B0.M}; explicit generators over $\mathbf Z[i]$
and invariant alternating and Hermitian forms.
\texttt{https://brauer.maths.qmul.ac.uk/Atlas/clas/U52/}.

\bibitem{BI:LR} J.~B. Lewis and V.~Reiner,
\emph{Circuits and Hurwitz action in finite root systems},
New York J. Math. \textbf{22} (2016), 1457--1486.
Corollary~1.4, p.~1459 and proof \S4.3, pp.~1480--1481;
Proposition~4.1, p.~1473.
Published PDF: \url{https://nyjm.albany.edu/j/2016/22-63v.pdf}.

\bibitem{BI:DHP} H.~Duyan, Z.~Halasi, and K.~Podoski,
\emph{Random bases for coprime linear groups},
J. Group Theory \textbf{23} (2020), no.~1, 133--157,
doi:10.1515/jgth-2019-0043.
Definition~2.3 and Lemma~2.4, pp.~137--138;
Lemma~3.1, p.~139; Theorem~4.1(2) and Remark~4.2, p.~141.
Publisher PDF:
\url{https://www.degruyterbrill.com/document/doi/10.1515/jgth-2019-0043/pdf}.
The updated eighteen-page author version
\url{https://zhalasi.web.elte.hu/papers/16_CoprimeRandomBase.pdf}
has Lemma~3.1 on pp.~5--6 and Theorem~4.1(2) on p.~7.
The seventeen-page arXiv v1, \texttt{arXiv:1903.00692v1}, instead
numbers the character bound Theorem~3.1(2), p.~5.

\bibitem{BI:Gluck1995} D.~Gluck,
\emph{Sharper character value estimates for groups of Lie type},
J. Algebra \textbf{174} (1995), no.~1, 229--266,
doi:10.1006/jabr.1995.1127.
Attribution of the character estimate; the precise bound used here
is the published restatement in \cite[Theorem~4.1(2)]{BI:DHP}.

\bibitem{BI:Hoffman}
C.~Hoffman, \emph{Cross characteristic projective representations
for some classical groups}, Journal of Algebra \textbf{229} (2000),
666--677, doi:10.1006/jabr.2000.8300.
The complete classical lower-bound table was read in the author's
14-page version, Table~1 on numbered p.~14, available at
\url{https://www.academia.edu/31207891/Cross_Characteristic_Projective_Representations_for_Some_Classical_Groups}.
The exceptional $E$-type rows there attributed to an unpublished
manuscript are not used.

\bibitem{BI:Luebeck}
F.~L\"ubeck, \emph{Smallest degrees of representations of exceptional
groups of Lie type}, Communications in Algebra \textbf{29} (2001),
2147--2169, doi:10.1081/AGB-100002175.
Author version: \url{https://www.math.rwth-aachen.de/~Frank.Luebeck/preprints/smdegexc.pdf}.
Theorem~2.1, numbered p.~4; first-degree tables in \S5.1--5.10,
numbered pp.~8--22.  The author's notation for the Suzuki and
Ree groups uses a square root of the field size, which has been
converted explicitly in the certificate.

\bibitem{BI:HHM}
G.~Hiss, W.~J.~Husen and K.~Magaard,
\emph{Imprimitive irreducible modules for finite quasisimple groups},
Memoirs of the American Mathematical Society \textbf{234} (2015),
no.~1104, doi:10.1090/memo/1104.
Author version: \url{https://arxiv.org/pdf/1211.6350v2}.
Introduction, numbered p.~4, and Chapter~5, Table~1, numbered p.~35,
give the ordinary-cover distinction and the complete exceptional
Schur-multiplier list used here.  No imprimitive-character
classification theorem from the monograph is needed for D-BI-20.

\bibitem{BI:WhiteSL2}
D.~L.~White, \emph{Character degrees of extensions of
$\operatorname{PSL}_2(q)$ and $\operatorname{SL}_2(q)$},
Journal of Group Theory \textbf{16} (2013), 1--33,
doi:10.1515/jgt-2012-0026.
Published PDF: \url{https://www.degruyterbrill.com/document/doi/10.1515/jgt-2012-0026/pdf}.
Theorem~B, p.~4; automorphism descriptions, pp.~2--3 and~6;
Lemma~3.1, p.~7.  A separately numbered author version is retained
at \url{https://www.math.kent.edu/~white/papers/sl2.pdf}; all
page locators in D-BI-20 refer to the published 33-page article.

\bibitem{BI:MalleRee}
G.~Malle, \emph{Die unipotenten Charaktere von ${}^2F_4(q^2)$},
Communications in Algebra \textbf{18} (1990), no.~7, 2361--2381,
doi:10.1080/00927879008824026.
The complete character table supplied by Malle is retained from
CTblLib's \texttt{2F4(8)} record; its CHEVIE provenance is also
explicitly documented in \cite[\S10.2, numbered p.~131]{BI:HHM}.
The unipotent-character paper is credited as the original source;
its complete proof was not freshly reconstructed in this run.

\bibitem{BI:GerardinLinear}
P.~G\'erardin, \emph{Weil representations associated to finite fields},
Journal of Algebra \textbf{46} (1977), 54--101,
doi:10.1016/0021-8693(77)90394-5.
Precise linear-family input: Proposition~4.2 and its complete proof,
pp.~82--83; the natural function model is also given in Proposition~1.4
and Corollary~1.4, pp.~59--60. The full published article copy read
was \url{https://www.academia.edu/126991842/Weil_representations_associated_to_finite_fields}.

\bibitem{BI:AtlasL34}
R.~A. Wilson, R.~A. Parker and J.~N. Bray,
\emph{ATLAS of Finite Group Representations: Linear group $L_3(4)$},
\url{https://brauer.maths.qmul.ac.uk/Atlas/lin/L34/},
accessed 21 November2026.
Precise inputs: displayed order, multiplier and outer group;
the natural three-dimensional $3.L_3(4)$ representations;
the list of cyclic outer extensions and automorphism descriptions.
The six complete numerical character tables used in D-BI-23b
are the published ATLAS records in CTblLib1.3.11, not computed
generic tables. Their complete source strings and normal fusions
are retained in \texttt{linear\_weil\_inputs.json}.

\bibitem{BI:BreuerL34}
T.~Breuer, \emph{Using Table Automorphisms for Constructing
Character Tables in GAP}, official CTblLib computational documentation,
\S2.6-4 (June2010), outer group and its three involution types;
\S2.7-4 (September2005), complete construction of the Schur cover
of $L_3(4)$ from its multiplier $C_4^2\times C_3$, with explicit
reference to the printed \emph{ATLAS of Finite Groups}, p.~23.
\url{https://www.math.rwth-aachen.de/~Thomas.Breuer/ctbllib/doc2/chap2.html}.
The diagonal, field and graph descriptions are also given in his
\emph{GAP Computations Concerning Probabilistic Generation of Finite
Simple Groups}, \S11.5-7, official computational supplement to
Breuer--Guralnick--Kantor's published generation work:
\url{https://www.math.rwth-aachen.de/~Thomas.Breuer/ctbllib/doc2/chap11_mj.html}.
Both invoked structural passages were freshly read on21 November2026.

\bibitem{BI:GMST2002}
R.~M. Guralnick, K.~Magaard, J.~Saxl and P.~H. Tiep,
\emph{Cross characteristic representations of symplectic and unitary groups},
Journal of Algebra \textbf{257} (2002), 291--347,
doi:10.1016/S0021-8693(02)00527-6.
Precise input: the ordinary unitary Weil degrees and extensions,
p.~293; the full published page was freshly read. The author's
earlier version is at
\url{https://www.academia.edu/14851735/Cross_characteristic_representations_of_symplectic_and_unitary_groups}.

\bibitem{BI:GLOST2018}
R.~M. Guralnick, M.~W. Liebeck, E.~A. O'Brien, A.~Shalev and P.~H. Tiep,
\emph{Surjective word maps and Burnside's $p^aq^b$ theorem},
Inventiones Mathematicae \textbf{213} (2018), 589--695,
doi:10.1007/s00222-018-0795-z.
Precise input: \S4.3, equations(4.3)--(4.4), unitary Weil
projector characters and their extensions; author-version p.~16,
\url{https://www.ma.imperial.ac.uk/~mwl/burnside.pdf}.
The complete invoked subsection was freshly read. Its pagination
is explicitly the author version, not an asserted journal page.

\bibitem{BI:YangWuRevinVdovin2023}
N.~Yang, Zh.~Wu, D.~O. Revin and E.~P. Vdovin,
\emph{On the sharp Baer--Suzuki theorem for the $\pi$-radical
of a finite group}, Sbornik: Mathematics \textbf{214} (2023),
108--147, doi:10.4213/sm9698e.
Precise input: \S2.2, definitions and Lemma6, p.~114,
giving the full inner-diagonal and field automorphisms of
simple unitary groups. This is the published restatement
of Gorenstein--Lyons--Solomon, Theorem2.5.12, actually read
in this run; no fresh proof of that classification is claimed.
\url{https://www.mathnet.ru/php/archive.phtml?jrnid=sm&option_lang=eng&paperid=9698&wshow=paper}.

\bibitem{SPW:Gerardin}
P. G\'erardin, \emph{Weil representations associated to finite fields},
J. Algebra \textbf{46} (1977), 54--101,
doi:10.1016/0021-8693(77)90394-5.
Precise input: Theorem~2.4(a), pp.~63--64, complete proof
pp.~64--69; Schr\"odinger formulas (2.8)--(2.9), p.~65;
Theorem~4.4(c), pp.~84--85, and Corollary~4.4(a), p.~85.

\bibitem{SPW:Taylor}
J. Taylor, \emph{Action of automorphisms on irreducible characters
of symplectic groups}, J. Algebra \textbf{505} (2018), 211--246,
doi:10.1016/j.jalgebra.2018.03.008.
Precise automorphism input: \S1.5 and \S12.5,
author PDF numbered pp.~2 and~24.

\bibitem{G2T:EY}
H.~Enomoto and H.~Yamada,
The characters of $G_2(2^n)$,
\emph{Japan. J. Math. (N.S.)} \textbf{12} (1986), 325--377.
doi:10.4099/math1924.12.325.
Published tables: Table~IV-1, p.~364; Table~IV-2, pp.~365--370;
and Table~IV-3, p.~371.
\url{https://www.jstage.jst.go.jp/article/math1924/12/2/12_2_325/_pdf/-char/en}.

\bibitem{G2T:Enomoto}
H.~Enomoto,
The characters of the finite Chevalley group $G_2(q)$, $q=3^f$,
\emph{Japan. J. Math. (N.S.)} \textbf{2} (1976), 191--248.
doi:10.4099/math1924.2.191.
Published tables: Table~VII-1, p.~239, and Table~VII-2, pp.~240--247.
\url{https://www.jstage.jst.go.jp/article/math1924/2/2/2_2_191/_pdf/-char/en}.

\bibitem{G2T:ChangRee}
B.~Chang and R.~Ree,
The characters of $G_2(q)$,
in \emph{Symposia Mathematica}, Vol.~XIII,
Academic Press, London--New York, 1974, pp.~395--413.
The complete ordinary table is accessed here through the official
generic entries \texttt{G2.111}, \texttt{G2.113},
\texttt{G2.121}, \texttt{G2.123} in \cite{G2T:GCT}.

\bibitem{G2T:DM}
D.~I.~Deriziotis and G.~O.~Michler,
Character table and blocks of finite simple triality groups ${}^3D_4(q)$,
\emph{Trans. Amer. Math. Soc.} \textbf{303} (1987), 39--70.
doi:10.1090/S0002-9947-1987-0896007-9.
Relevant locations: \S4, Theorem~4.3, Tables~3.6 and~4.4;
the complete mathematical table input is supplied by the three
\texttt{3D4} entries of \cite{G2T:GCT}.

\bibitem{G2T:CHEVIE}
M.~Geck, G.~Hiss, F.~L\"ubeck, G.~Malle, and G.~Pfeiffer,
CHEVIE---a system for computing and processing generic character tables,
\emph{Appl. Algebra Engrg. Comm. Comput.} \textbf{7} (1996), 175--210.

\bibitem{G2T:GCT}
The OSCAR/GenericCharacterTables developers,
\emph{GenericCharacterTables.jl}, official mathematical table library,
commit \texttt{87231738271ed548f62e7f1c98ff2bbfc1ed30d3},
retrieved 21 November 2026.
Complete inventories and published source attributions:
\url{https://oscar-system.github.io/GenericCharacterTables.jl/stable/tables_list/}.
Frozen source directory:
\url{https://github.com/oscar-system/GenericCharacterTables.jl/tree/87231738271ed548f62e7f1c98ff2bbfc1ed30d3/data/Tables}.

\bibitem{ESC:Enomoto}
H.~Enomoto,
The characters of the finite symplectic group $\operatorname{Sp}(4,q)$,
$q=2^f$,
\emph{Osaka J. Math.} \textbf{9} (1972), 75--94.
doi:10.18910/8596.
Precise complete-character construction: \S4, pp.~80--85;
complete group tables IV-1 and IV-2, pp.~92--94.
\url{https://ir.library.osaka-u.ac.jp/repo/ouka/all/8596/ojm09_01_07.pdf}.
The actual mathematical input used is the corrected official
\texttt{Sp4.0} table; its documented recomputation by F.~L\"ubeck
is distinguished from an unchecked transcription of the original print.

\bibitem{ESC:Luebeck}
F.~L\"ubeck,
\emph{Charaktertafeln f\"ur die Gruppen $\operatorname{CSp}_6(q)$
mit ungeradem $q$ und $\operatorname{Sp}_6(q)$ mit geradem $q$},
Dissertation, Universit\"at Heidelberg, 1993.
Full author text:
\url{https://www.math.rwth-aachen.de/~Frank.Luebeck/preprints/CharTabCSp6.pdf}.
Complete-table statement and supplied CHEVIE files: Introduction,
pp.~3--4; construction overview \S2.2, pp.~18--21;
computational mathematical input and checks \S9, pp.~78--79.
The implementation used here is the full \texttt{Sp6.0} table,
not its \texttt{uniSp6.0} subtable, in \cite{G2T:GCT}.

\bibitem{D:BMO2017}
T.~Breuer, G.~Malle, and E.~A.~O'Brien,
\emph{Reliability and reproducibility of Atlas information},
in \emph{Finite Simple Groups: Thirty Years of the Atlas and Beyond},
Contemporary Mathematics \textbf{694}, American Mathematical Society,
2017, pp.~21--31, DOI: 10.1090/conm/694/13960.
Invoked: \S3.3.1 and Lemma~3.1; the complete argument was read on
pp.~5--6 of the author version
\url{https://agag-malle.math.rptu.de/~malle/download/atlas.pdf}.
The publisher metadata confirms the published page range; author
pagination is stated separately.  The local centralization-label
discrepancy is recorded explicitly in the proof and is not used.

\bibitem{A2:SF73}
W.~A.~Simpson and J.~S.~Frame,
\emph{The character tables for $\operatorname{SL}(3,q)$,
$\operatorname{SU}(3,q^2)$, $\operatorname{PSL}(3,q)$,
$\operatorname{PSU}(3,q^2)$}, Canadian Journal of Mathematics
\textbf{25} (1973), 486--494, doi:10.4153/CJM-1973-049-7.
Complete class and character Tables~1a, 1b and~2, as corrected
in the official implementation \cite{G2T:GCT}; its frozen source
metadata records the exact affected entries and parameter headers.
The paper's unitary notation $q^2$ denotes the quadratic field
extension; the present text uses $\operatorname{SU}_3(q)$.

\bibitem{A2:Geck90}
M.~Geck, \emph{Irreducible Brauer characters of the 3-dimensional
special unitary groups in non-defining characteristic},
Communications in Algebra \textbf{18} (1990), 563--584,
doi:10.1080/00927879008823932.
\S3, corrections to the ordinary character table used by the
official complete \texttt{SU3.2} implementation.  The present
run accesses this corrected mathematical input through that
official table; it does not claim a new reading of the whole
original article or use an unverified modular-character theorem.

\bibitem{HES:GM2006}
D.~Gluck and K.~Magaard,
\emph{Absolute Fourth Moments and Finiteness of Linear Groups},
Communications in Algebra \textbf{34} (2006), no.~9, 3387--3407,
doi:10.1080/00927870600788448.
Exact input: the unconditional multiplier and outer-automorphism statements
for $\operatorname{Sp}_{2n}(2)$, $n\ge4$, in the author manuscript, p.~25.
Author manuscript read at
\url{https://www.academia.edu/29518749/Absolute_Fourth_Moments_and_Finiteness_of_Linear_Groups}.
The paper's additional small-fourth-moment hypotheses are not imported.

\bibitem{HES:S12Data}
T.~Breuer and contributors, \emph{GAP Character Table Library},
official complete ordinary table \texttt{S12(2)}, 477 conjugacy classes.
Exact table provenance: computed by Christoph K\"ohler using
group-theoretic and character-theoretic methods.
Official table record:
\url{https://www.math.rwth-aachen.de/~Thomas.Breuer/ctbllib/ctbltoc/data/S12\%282\%29.html}.
Library source \texttt{data/ctosymp3.tbl}; its revision record dates the
addition to 1 October 2001. Accessed 21 November 2026.
No separately verified refereed-publication claim is made for this
construction. The full numerical input, provenance, and exact derived-bound
certificate are retained with D-HES-F.

\bibitem{HES:Liebeck2017}
M.~W.~Liebeck,
\emph{Character ratios for finite groups of Lie type, and applications},
Contemporary Mathematics \textbf{694} (2017), 193--208.
Exact input: Theorem~2.4(i), author pp.~3--4, with the simply connected
quasisimple-group conventions on p.~2: for a noncentral unipotent element
and a nontrivial ordinary irreducible character the ratio is at most $3/4$
when $q\le4$, and at most $1/(\sqrt q-1)$ when $q>4$.
The subtraction is outside the radical. This is an explicit published
restatement of Gluck (1995).
Author PDF: \url{https://www.ma.imperial.ac.uk/~mwl/princeton-survey-final.pdf}.
The nonunipotent bound in part~(ii) is not substituted for part~(i).

\bibitem{HES:Gow1981}
R.~Gow,
\emph{Products of two involutions in classical groups of characteristic~2},
Journal of Algebra \textbf{71} (1981), 583--591.
Exact input: theorem on p.~583, symplectic proof in Section~3, pp.~589--591.
Original published scan read at
\url{https://www.academia.edu/115642614/Products_of_two_involutions_in_classical_groups_of_characteristic_2}.

\bibitem{HES:LiXia2022}
C.~H.~Li and B.~Xia,
\emph{Factorizations of Almost Simple Groups with a Solvable Factor,
and Cayley Graphs of Solvable Groups},
Memoirs of the American Mathematical Society \textbf{279} (2022),
no.~1375, v+99 pp., doi:10.1090/memo/1375.
Exact input: Section~2.2.2, Table~2.1, numbered p.~14 of the author
version \url{https://arxiv.org/pdf/1408.0350v3}, giving
$\operatorname{Out}(\operatorname{PSp}_{2m}(q))=C_d\times C_f$,
$q=p^f$, $d=(2,q-1)$, $(m,p)\ne(2,2)$.
The author's publication list links this version to the published memoir:
\url{https://sites.google.com/site/binzhouxia/research}.
The full relevant table and its preceding hypotheses were read and
rendered. The publisher PDF returned HTTP~403; no assertion that its
pagination was checked is made.

\bibitem{C3:Luebeck}
F.~L\"ubeck,
\emph{Charaktertafeln f\"ur die Gruppen $\operatorname{CSp}_6(q)$
mit ungeradem $q$ und $\operatorname{Sp}_6(q)$ mit geradem $q$},
Dissertation, Universit\"at Heidelberg, 1993.
\url{https://www.math.rwth-aachen.de/~Frank.Luebeck/preprints/CharTabCSp6.pdf}.
Complete-table statement and mathematical files: Introduction,
pp.~3--4; complete construction overview \S2.2, pp.~18--21;
CHEVIE data and checks \S9, pp.~78--79. The full generic
$\operatorname{CSp}_6(q)$ table is used, not only its printed
unipotent-character appendix. The full dissertation was retained;
the invoked passages were freshly read. Its complete construction
was not independently reconstructed in this run.

\bibitem{C3:GCT}
F.~L\"ubeck and the CHEVIE/GenericCharacterTables contributors,
\emph{GenericCharacterTables.jl}, official complete table
\texttt{CSp6.1}.
Inventory and exact scope:
\url{https://oscar-system.github.io/GenericCharacterTables.jl/stable/tables_list/}.
The immutable mathematical input is
\texttt{data/Tables/C3/CSp6.1.jl} at repository commit
\texttt{87231738271ed548f62e7f1c98ff2bbfc1ed30d3},
\url{https://github.com/oscar-system/GenericCharacterTables.jl}.
Original source SHA-256:
\texttt{f5e68e9429054370df7376f174a095bc548cf1ecf7ce44f5446f1eaebba250dc}.
Retrieved 21 November 2026. The input is parsed as mathematical
data without executing the external Julia program.

\bibitem{C3:Small}
T.~Breuer and contributors, \emph{GAP Character Table Library},
version 1.3.11, complete ordinary tables
\texttt{2.S6(3)} and \texttt{2.S6(3).2}, both with exact
source string ``origin: ATLAS of finite groups''.
The complete original values, class orders and sizes, power maps,
fusion and provenance are retained in \texttt{small\_inputs.json}.
Official table records:
\url{https://www.math.rwth-aachen.de/~Thomas.Breuer/ctbllib/ctbltoc/data/2.S6\%283\%29.html}
and
\url{https://www.math.rwth-aachen.de/~Thomas.Breuer/ctbllib/ctbltoc/data/2.S6\%283\%29.2.html}.
The GAP and separate exact polynomial-Fourier checks are retained;
no independent construction of the original ATLAS tables is claimed.

% Additional existing bibliography entries used verbatim:
% BI:HHM; HES:LiXia2022; HES:Liebeck2017.

\bibitem{HD:A7Tables}
J.~H.~Conway, R.~T.~Curtis, S.~P.~Norton, R.~A.~Parker and
R.~A.~Wilson, \emph{Atlas of Finite Groups}, Oxford, 1985;
complete ordinary tables as retained in T.~Breuer and contributors,
\emph{GAP Character Table Library}, version~1.3.11, exact identifiers
\texttt{6.A7} and \texttt{6.A7.2}, both with provenance
``origin: ATLAS of finite groups''.
Official records:
\url{https://www.math.rwth-aachen.de/~Thomas.Breuer/ctbllib/ctbltoc/data/6.A7.html},
\url{https://www.math.rwth-aachen.de/~Thomas.Breuer/ctbllib/ctbltoc/data/6.A7.2.html}.
The simple-group metadata
\texttt{ExtensionInfoCharacterTable(CharacterTable("A7"))}
give multiplier and outer-group orders six and two. The exact
\texttt{A7.2} record identifies $S_7=\operatorname{Aut}(A_7)$:
\url{https://www.math.rwth-aachen.de/~Thomas.Breuer/ctbllib/ctbltoc/data/A7.2.html}.
These web records identify version~1.3.10 and their May~2025
generation dates; the actual retained export is from~1.3.11.
The complete values, power maps and core fusion are mathematical
inputs to the new exact spectrum inequalities, not tables
independently constructed in this run. Retrieved 21 November 2026.

\bibitem{OC2:CTbl}
T.~Breuer and contributors, \emph{GAP Character Table Library},
version 1.3.11. Complete tables \texttt{2.S4(5)},
\texttt{2.S4(5).2}, \texttt{S4(7)}, \texttt{S4(7).2},
\texttt{S4(9)} and \texttt{S4(9).2\_i} for $i=1,2,3$.
The first two have ATLAS provenance; the latter tables were computed
by Dixon's algorithm and character-theoretic methods, with the
original attribution to Thomas Breuer and the 2001/2004 dates retained.
Official table pages, freshly read:
\url{https://www.math.rwth-aachen.de/~Thomas.Breuer/ctbllib/ctbltoc/data/S4(7).html},
\url{https://www.math.rwth-aachen.de/~Thomas.Breuer/ctbllib/ctbltoc/data/S4(9).html},
and the adjacent three \texttt{S4(9).2\_i.html} pages, explicitly
identifying respectively field, diagonal and diagonal-field extensions.
The web inventory identifies version 1.3.10, while the actual retained
export is from version 1.3.11. Complete mathematical data and exact
source strings are retained in \texttt{odd\_c2/character\_inputs.json}.

\bibitem{OC2:GCT}
F.~L\"ubeck and the CHEVIE/GenericCharacterTables contributors,
\emph{GenericCharacterTables.jl}, official table inventory,
type $C_2$, entry \texttt{uniCSp4.1}.
\url{https://oscar-system.github.io/GenericCharacterTables.jl/stable/tables_list/}.
Freshly read 21 November 2026. The text explicitly limits this entry
to unipotent characters and warns that a preliminary full table
contains errors. It is not used as a complete ordinary table here.

% Existing HES:LiXia2022 supplies the full automorphism table.

\bibitem{OC2:Dabbaghian2006}
V.~Dabbaghian-Abdoly,
\emph{Constructing representations of the finite symplectic group
$\operatorname{Sp}(4,q)$}, Journal of Algebra \textbf{303} (2006),
618--625, DOI \href{https://doi.org/10.1016/j.jalgebra.2005.04.020}
{10.1016/j.jalgebra.2005.04.020}.
Complete author-uploaded published text:
\url{https://www.researchgate.net/publication/222038219_Constructing_representations_of_the_finite_symplectic_group_Sp4q}.
Also indexed at
\url{https://www.academia.edu/78281746/Constructing_representations_of_the_finite_symplectic_group}.
The invoked inputs are Table~2, p.~621, its full specification on
p.~620, the complete restriction formulas pp.~622--623, and the
basis averaging identities pp.~623--624. They concern all ordinary
characters on unipotent classes. The four inconsistent indexed
restriction identities are not used; the direct Table~2 rows and
the failed and successful exact transcription checks are retained.
The binary-PDF retrieval attempts returned HTTP403; the complete
indexed published author text was read. No unobserved PDF image
inspection is claimed.

\bibitem{OC2:Vinroot2004}
C.~R.~Vinroot, \emph{A factorization in $\operatorname{GSp}(V)$},
Linear and Multilinear Algebra \textbf{52} (2004), 385--403,
DOI \href{https://doi.org/10.1080/03081080412331316079}
{10.1080/03081080412331316079}.
Author full text:
\url{https://www.math.wm.edu/~vinroot/gspfactor.pdf}.
Theorem~1 on author p.~1 restates Wonenburger's factorization into
two skew-symplectic involutions for arbitrary fields of characteristic
different from two. The theorem and full stated hypotheses were
freshly read; the original Wonenburger proof is cited, not newly
reconstructed. The resulting class-inversion corollary is proved
explicitly in D-OC2-03.

% Existing HES:LiXia2022, BI:HHM, and the reviewed SPW inputs are used.

\bibitem{HD:ExceptionalSmallTables}
T.~Breuer and contributors, \emph{GAP Character Table Library},
version~1.3.11. Exact complete ordinary identifiers
\texttt{2.F4(2)}, \texttt{2.F4(2).2}, \texttt{F4(3)},
\texttt{E6(2)} and \texttt{(2\^{}2x3).2E6(2)}.
The first two have ATLAS provenance (Conway et al., Oxford, 1985);
the extension's pre-1.3 square-map error is corrected in the retained
version. Official records:
\url{https://www.math.rwth-aachen.de/~Thomas.Breuer/ctbllib/ctbltoc/data/2.F4(2).html},
\url{https://www.math.rwth-aachen.de/~Thomas.Breuer/ctbllib/ctbltoc/data/2.F4(2).2.html}.
The official \texttt{F4(2).2} record also identifies the full automorphism group.
For \texttt{E6(2)}, the official record identifies B.~Fischer's table
with W.~Unger's 2016 Magma corrections to order-91 power maps and
irrationalities:
\url{https://www.math.rwth-aachen.de/~Thomas.Breuer/ctbllib/ctbltoc/data/E6(2).html}.
The universal \texttt{(2\^{}2x3).2E6(2)} record specifies
\texttt{PossibleCharacterTablesOfTypeV4G} construction:
\url{https://www.math.rwth-aachen.de/~Thomas.Breuer/ctbllib/ctbltoc/data/(2^2x3).2E6(2).html}.
The $F_4(3)$ complete table's original data string is
``computed by Eamonn O'Brien using Magma, May/June 2024'';
see \texttt{ctoexce1.tbl}, the block beginning \texttt{MOT("F4(3)")},
and the official \emph{Changes in the GAP Character Table Library},
p.~15, release~1.3.10 (May~2025), $F_4(3)$ entry:
\url{https://www.math.rwth-aachen.de/homes/Thomas.Breuer/ctbllib/doc/ctbldiff.pdf}.
These source records and exact values are cited mathematical inputs,
not independent constructions in this run. Retrieved 21 November 2026.

\bibitem{HD:OrthogonalSmallTables}
T.~Breuer and contributors, \emph{GAP Character Table Library},
version~1.3.11, complete ordinary tables listed in (OSM.1).
The exact identifiers, full values, indicators, normal fusions and
original provenance strings are retained in
\texttt{research/orthogonal\_small\_inputs.json}.
The meaning of multiplier/outer metadata and the universal-cover
naming convention is the official manual, \S3.7-3,
\emph{ExtensionInfoCharacterTable}; source descriptions are \S3.7-5,
\emph{InfoText}:
\url{https://www.math.rwth-aachen.de/~Thomas.Breuer/ctbllib/doc/chap3.html}.
The sixfold $\Omega_7(3)$ core/outer pair, the $\Omega_8^-(2)$ pair,
and both $\Omega_{10}^{\pm}(2)$ pairs have ATLAS provenance
(Conway et al., Oxford, 1985). The official outer records explicitly
identify the full automorphism groups for the centerless even-field
types. Exact source pages are obtained by appending the displayed
identifier and \texttt{.html} to
\url{https://www.math.rwth-aachen.de/~Thomas.Breuer/ctbllib/ctbltoc/data/}.
All fifteen records were freshly read on 21 November 2026.

The $\operatorname{Spin}_8^+(3)$ table records the
\texttt{PossibleCharacterTablesOfTypeV4G} construction; the accepted
specific official table is the premise, not a new uniqueness proof
for that general construction. The $2.O_8^-(3)$ and $2.O_9(3)$
tables are attributed to Max Neunh\"offer, April~2008.
The two $O_{12}^{\pm}(2)$ core tables record Magma V2.23-9
computations of 4 September 2018. Their outer tables have distinct
provenance: J.~S.~Frame~(1986), with Aachen technical support,
for the plus type; J.~S.~Frame and H.~Pahlings~(1987), CAS library,
for the minus type. These are not collectively labeled ATLAS tables.
The automatically generated web pages show version~1.3.10;
the actual retained mathematical export uses version~1.3.11.

\bibitem{UNW:LT2024}
M.~Larsen and P.~H.~Tiep,
\emph{Uniform character bounds for finite classical groups},
Annals of Mathematics \textbf{200} (2024), 1--70,
doi:10.4007/annals.2024.200.1.1.
Publisher record: \url{https://annals.math.princeton.edu/2024/200-1/p01}.
Complete accepted author version read:
\url{https://arxiv.org/pdf/2403.09046}, dated 14 March 2024;
Theorem~4.4 and proof, numbered pp.~19--20;
Theorem~5.5 and complete proof, numbered pp.~25--26.
The last display on p.~26 gives
$\sigma=1/(6507\cdot2^{21}\cdot10^{18})>7\cdot10^{-29}$.
The publication was accepted 12 March 2024 and published online
3 July 2024 according to the publisher record.  Page locators here
refer to the author version, not the 70-page journal typesetting.

\bibitem{UNW:GLT2020}
R.~M.~Guralnick, M.~Larsen and P.~H.~Tiep,
\emph{Character levels and character bounds},
Forum of Mathematics, Pi \textbf{8} (2020), e2, 81 pp.,
doi:10.1017/fmp.2019.9.
Complete published PDF:
\url{https://www.cambridge.org/core/services/aop-cambridge-core/content/view/77A6CCCB36F94337FD1DAA7F67B40379/S205050861900009Xa.pdf/character_levels_and_character_bounds.pdf}.
Definition~2 and Lemma~7.1, pp.~51--52; Theorem~7.3, p.~52;
complete unitary dual-pair argument, \S8.3, pp.~65--67,
especially equations (8.17)--(8.18), pp.~65--66;
Corollary~8.6 with complete proof, pp.~67--68.
Complete general-linear dual-pair argument, \S8.1, pp.~55--61,
including Lemma~8.2 and equations (8.9)--(8.11), pp.~59--60;
level-one Weil identification, p.~61.

\bibitem{OL:GLT2024}
R.~M.~Guralnick, M.~Larsen and P.~H.~Tiep,
\emph{Character levels and character bounds for finite classical groups},
Inventiones mathematicae \textbf{235} (2024), 151--210,
doi:10.1007/s00222-023-01221-5.
Full published PDF:
\url{https://link.springer.com/content/pdf/10.1007/s00222-023-01221-5.pdf}.
Exact inputs, with their complete proofs read:
Definition~3.2(ii) and Lemmas~3.3--3.4, pp.~160--161;
Theorem~3.7, pp.~165--167;
Theorem~1.5(iii), p.~154, with its proof on p.~167;
Proposition~6.3, pp.~193--194;
Proposition~7.3, pp.~198--200.
Corollary~4.7, p.~174, is distinguished from an irreducibility theorem.
The online publication date is 29 September 2023; the volume year is 2024.


\bibitem{OL:KTV2023}
J.~Kim, S.~Trefethen and C.~R.~Vinroot,
\emph{Reality properties of finite simple orthogonal groups},
Advances in Group Theory and Applications \textbf{16} (2023), 39--66,
doi:10.32037/agta-2023-007.
Full published PDF:
\url{https://www.advgrouptheory.com/journal/Volumes/16/KimTrefethenVinroot.pdf}.
Exact inputs read: full orthogonal groups are strongly real in odd
and even characteristic, pp.~39--40 and Theorem~1, p.~44;
the determinant/spinor-norm quotient in odd characteristic and
index-two derived subgroup in even characteristic, Section~2,
pp.~41--43.  The invoked statements explicitly attribute their
original proofs to Wonenburger, Ellers--Nolte and Gow; no new reading
of those original proofs is claimed for this section.

\bibitem{OS:Taylor1992}
D.~E.~Taylor, \emph{The Geometry of the Classical Groups},
Sigma Series in Pure Mathematics, vol.~9, Heldermann Verlag, 1992.
Author's corrected full text (April 2021 corrections):
\url{https://www.maths.usyd.edu.au/u/don/papers/gcg.pdf}.
Publisher record:
\url{https://www.heldermann.de/SSPM/SSPM09/sspm09.htm}.
Exact inputs read: Witt's theorem and complete proof,
Theorem~7.4, pp.~57--58; finite-field quadratic decomposition,
Lemma~11.1, Theorem~11.2 and ensuing discussion, pp.~138--140;
reflection generation, Theorems~11.39, 11.41 and
Corollary~11.42, pp.~157--159, including the full characteristic-two
qualification and proofs; perfectness and simplicity,
Theorems~11.47--11.48, pp.~161--162; reflection spinor norm, p.~163.
The author version has 237 PDF pages; these are its numbered
mathematical page locators, not PDF indices.

\bibitem{OS:TaylorVinroot2018}
G.~K.~Taylor and C.~R.~Vinroot,
\emph{On involutions and indicators of finite orthogonal groups},
Journal of the Australian Mathematical Society \textbf{105} (2018),
380--416, doi:10.1017/S144678871800006X.
Exact input: Lemma~2.1, p.~383, the even-characteristic equivalence
$g\in\Omega\Longleftrightarrow\operatorname{rank}(g-I)$ is even.
Full published PDF:
\url{https://www.cambridge.org/core/services/aop-cambridge-core/content/view/7E9E87830BB65CBA396A3BD82A2863C1/S144678871800006Xa.pdf/on_involutions_and_indicators_of_finite_orthogonal_groups.pdf}.

% Previously defined OL:KTV2023, OL:GLT2024, HES:LiXia2022 and
% HES:Liebeck2017 retain their exact published statements and locators.
% No unavailable GM1998 input is used.

\bibitem{SF:Wonenburger1962}
Maria J. Wonenburger,
\emph{The Clifford Algebra and the Group of Similitudes},
Canadian Journal of Mathematics \textbf{14} (1962), 45--59,
\url{https://doi.org/10.4153/CJM-1962-004-1}.
The exact input is Part~II, pp.~48--49: the associated semilinear
automorphism of the even Clifford algebra and the complete defining-ideal
calculation. The full-algebra extension in Theorem~3 is not invoked.
The original published PDF was freshly read; the Clifford quotient and
reversal definitions on pp.~46--47 were also read.

\bibitem{SF:LOST2011}
Martin W. Liebeck, E. A. O'Brien, Aner Shalev, and Pham Huu Tiep,
\emph{Commutators in finite quasisimple groups},
Bulletin of the London Mathematical Society \textbf{43} (2011), 1079--1092,
\url{https://doi.org/10.1112/blms/bdr043}.
Section~4 defines the norm-one even Clifford group and records the
Spin-to-Omega exact sequence. The freshly read author version places
this passage on pp.~10--11:
\url{https://wwwf.imperial.ac.uk/~mwl/ORE-PLUS2.PDF}.
Publisher pagination for this passage was not independently matched.
The full proof of Lemma~4.1 on author pp.~11--12 was read, but the
new argument uses a full inverse-image Levi and does not require that lemma.

% Existing HOS:LST2011 is used at proof of Theorem4.3.6 Step1 pp1918--1919
% for finite Spin quasisimplicity and the exact double-cover sequence,
% and Theorem4.3.6/Remark4.3.7 for the character bound.
% Existing HES:LiXia2022 remains the complete outer-order table input
% already fully incorporated in the reviewed natural-lift proof D-OL07.

\bibitem{HOS:Vinroot2004}
C.~R.~Vinroot, \emph{A factorization in $\operatorname{GSp}(V)$},
Linear and Multilinear Algebra \textbf{52} (2004), 385--403,
DOI \href{https://doi.org/10.1080/03081080412331316079}
{10.1080/03081080412331316079}.
Author full text: \url{https://www.math.wm.edu/~vinroot/gspfactor.pdf}.
Theorem~1 on author p.~1 states Wonenburger's factorization into
two anti-symplectic involutions over every field of characteristic
different from two. The complete statement and domain were freshly
read for D-HOS-10. The original factorization proof is cited as a
published input, not claimed as a new reconstruction here.

% Existing BI:HHM and HES:LiXia2022 supply precisely the complete
% universal-cover list (Chapter5 Table1 p35) and type-C automorphisms
% (Section2.2.2 Table2.1, author p14), respectively.

\bibitem{HOS:LST2011}
M.~Larsen, A.~Shalev and P.~H.~Tiep,
\emph{The Waring problem for finite simple groups},
Annals of Mathematics \textbf{174} (2011), 1885--1950.
doi:10.4007/annals.2011.174.3.10.
Published full text:
\url{https://annals.math.princeton.edu/wp-content/uploads/annals-v174-n3-p10-p.pdf}.
The inputs here are Definition~4.1.1 and (4.1.1), p.~1907;
Theorem~4.3.6, pp.~1917--1918, especially (4.3.2);
the complete proof, pp.~1918--1923, especially its part~4 on p.~1922;
and Remark~4.3.7, p.~1923.

\bibitem{HOS:Taylor1992}
D. E. Taylor, \emph{The Geometry of the Classical Groups},
Heldermann Verlag, Berlin, 1992. Corrected author version,
unitary classification and complete order calculation pp.~115--118;
finite orthogonal classification and complete order calculation
pp.~138--141.
\url{https://www.maths.usyd.edu.au/u/don/papers/gcg.pdf}.
The Larsen--Shalev--Tiep reference is the existing
\texttt{HOS:LST2011} entry, here Proposition~4.2.3 and its complete
proof pp.~1909--1912.

\bibitem{SPOR:Atlas} J.~H.~Conway, R.~T.~Curtis, S.~P.~Norton,
R.~A.~Parker, and R.~A.~Wilson, \emph{Atlas of Finite Groups},
Oxford University Press, 1985, xxxiv+252 pp.
The sporadic group character tables, their Schur covers, and outer
extensions; exact electronic table identifiers and row/column data
are retained in the accompanying finite certificate.

\bibitem{SPOR:CTblLib} T.~Breuer and contributors,
\emph{The GAP Character Table Library}, version 1.3.11 (2025).
\url{https://www.math.rwth-aachen.de/~Thomas.Breuer/ctbllib/}.
Manual \S2.1 (tables and class fusions), \S2.2 (identifiers),
\S3.7 entry \texttt{ExtensionInfoCharacterTable} (Schur multipliers
and outer automorphism groups), and the thirty-eight exact table
identifiers in D-SPOR-01.  Author archive SHA-256:
\texttt{f201e9da571c681a7bbfb1e4683bdb2c80a23b005099ed0453521f59cb443434}.
Every used table's \texttt{InfoText} identifies the published ATLAS
as its origin.  The library is an input, not a new proof of those tables.

\bibitem{SPOR:GAP} The GAP Group,
\emph{GAP---Groups, Algorithms, and Programming}, version 4.16.1 (2026),
\url{https://www.gap-system.org/}.
Reference manual entry \texttt{EigenvaluesChar}; the cyclic Fourier
identity is independently proved and implemented in D-SPOR-01.

\bibitem{Root:Mumford}
D. Mumford, \emph{Prym Varieties I}, in Contributions to Analysis,
Academic Press, 1974, pp.325--350. Exact locator: section3,
Corollaries1--2, p.332.
\url{https://www.dam.brown.edu/people/mumford/alg_geom/papers/1974a--PrymVar-I-NC.pdf}.
\bibitem{Root:vanGeemen}
B. van Geemen, \emph{Theta functions and cycles on some abelian fourfolds},
Mathematische Zeitschrift 221 (1996), 617--631.
Exact locator: Theorem5.3, pp.629--630; Remark5.4, p.630.
\url{https://sites.unimi.it/vangeemen/ThetafctsCycles4folds_MathZ1996.pdf}.
\bibitem{Root:Fulton}
W. Fulton, \emph{Intersection Theory}, second edition, Springer, 1998.
Exact locator: Example8.4.6, the degree bound for components of
projective intersections used in the Bezout estimate.
\bibitem{Root:Attack}
\emph{HODGE\_DIVISION\_ATTACK\_20261121.md}, supplied attack by
Claude Opus5.5, effort high. Exact locators: R1--R4/Theorem W,
Theorem D, Lemma Dom, Propositions O/O$'$, Corollary O$''$,
and Theorem E. These are accepted inputs of the present task.
\bibitem{Root:Project}
Supplied project returns: HECKE\_DIV R21--R23, P0, P6, SQ1--SQ7;
DEG\_L B1, B19, B21, DIV1--DIV3, S-14;
SIGN\_NEXT SN-C03/C04 and SN-14;
GHOSTS B02/B03/B06 and R12/R20--R24 via the supplied intake.
The GHOSTS full paper is absent from the supplied directory.
\bibitem{Root:DeligneIII}
P. Deligne, \emph{Th\'eorie de Hodge III}, Publications
Math\'ematiques de l'IH\'ES \textbf{44} (1974), 5--77.
Exact locators: Theorem8.2.4(iii), pp.38--39;
Proposition8.2.5 and equation8.2.5.2, p.39.
\url{https://www.numdam.org/item/10.1007/BF02685881.pdf}.
\bibitem{Root:SK}
T. Shioda and T. Katsura, \emph{On Fermat varieties}, Tohoku
Mathematical Journal (2) \textbf{31} (1979), 97--115.
Exact locators: Lemmas1.1--1.6, pp.98--101; Theorem1.7 and
Remark1.8, pp.101--102; Lemmas2.1--2.3 and Proposition2.4,
pp.103--104. DOI:10.2748/tmj/1178229881.
\url{https://www.jstage.jst.go.jp/article/tmj1949/31/1/31_1_97/_pdf/-char/en}.
\bibitem{Root:Gale}
OpenAI, \emph{Abelian covers, Gale correspondences, and the Hodge
conjecture for powers}, 6 October2026, \textbf{unrefereed preprint}.
Exact locator: Lemma3.4, pp.10--14, in the public deposit at commit
\texttt{1581d4cbdd69a0faf34b2089614fe2a814316a6a}.
\url{https://github.com/openai/math/tree/1581d4cbdd69a0faf34b2089614fe2a814316a6a/preprints}.
Credited geometric lead; its theorem is not a premise of HD-FS1.
\bibitem{Root:Rstation66}
\texttt{cm\_symmetry\_scout} and \texttt{lelong\_ridge},
\emph{Branch-exponent determinant sources and a dominant nonsplit
Gaussian sixfold family}, rStation paper66, version1,
28 October2026, \textbf{unrefereed manuscript}.
Exact locators: sections2 (inertia-exponent supply) and4
(Hurwitz--Hodge volume). \url{https://rstation.ai/problems/hodge_conjecture/papers/66}.
Used as a lead and numerical comparison, not as an unverified theorem.
\bibitem{Root:Groth}
A. Grothendieck, Montpellier archive, cote16, scans048--050
(condition I, polynomial inverse, and corrected critical map);
cote44, scan002/footer1, Theorem4. Exact mathematical source corrections
are those of the supplied SIGN\_NEXT SN-C03/C04. No new transcription
of the original handwritten scans is claimed in this run.
\bibitem{Root:Frameworks}
EGA IV$_3$, 8.8.2; FGA, Bourbaki expose221, Lemme2.4 and
Theoremes3.1--3.2; SGA1 XII5.1; SGA4 XVI2.2;
SGA5 III3, 5.1--5.3, and VII7.1;
SGA6 Expose0 section1, (1.1)--(1.3), ExposeV2.2.1,
Appendix7.1 and Proposition7.5, and ExposeX Appendix7, (7.13.10).
These exact framework locators are carried from the supplied
SIGN\_NEXT and DEG\_L source checks, with the scope explained in the text.
\bibitem{HD:GT2012}
R.~M.~Guralnick and P.~H.~Tiep,
\emph{A problem of Koll\'ar and Larsen on finite linear groups
and crepant resolutions},
Journal of the European Mathematical Society \textbf{14}
(2012), 605--657, doi:10.4171/JEMS/313.
Published PDF: \url{https://ems.press/content/serial-article-files/31824?nt=1}.
Lemma~2.23 and its complete proof, pp.~621--622;
Proposition~3.4 and its proof, pp.~626--627.
HD-SPIN-01 uses the explicit local $2.A6$ input and the recorded
spin-degree lower bound, with a new moving-rank deduction.

\bibitem{HD:CMN2022} B. Collins, S. Matsumoto and J. Novak,
\emph{The Weingarten calculus}, Notices Amer. Math. Soc. 69(5)
(2022), 734--745, doi:10.1090/noti2474.
Theorem4.2, read in author version arXiv:2109.14890v1, p.9;
published and author page numbers are not identified.
\bibitem{HD:BNOT2015} C. Bessenrodt, H. N. Nguyen, J. B. Olsson
and H. P. Tong-Viet, \emph{Complex group algebras of the double
covers of the symmetric and alternating groups}, Algebra Number
Theory9(3) (2015), 601--628, doi:10.2140/ant.2015.9.601.
Section2A p.604; Lemmas3.1--3.2 p.607.
\bibitem{HD:TongViet2012} H. P. Tong-Viet, \emph{Alternating and
sporadic simple groups are determined by their character degrees},
Algebras Represent. Theory15(2) (2012), 379--389,
doi:10.1007/s10468-010-9247-1. Lemma4(a) and Corollary5,
read in the author full text, pp.2--3; Rasala Result3 is cited there.
\bibitem{HD:GKT2019} R. M. Guralnick, N. M. Katz and P. H. Tiep,
\emph{Rigid local systems and alternating groups}, Tunisian J. Math.
1(3) (2019), 295--320, doi:10.2140/tunis.2019.1.295.
Proof of Theorem9.7, p.312: precise low-degree Specht consequence
of Rasala1977 Result3; no semisimple-regular hypothesis is imported
into the cited degree-classification statement.

\bibitem{HD:FominLulov1995}
S.~Fomin and N.~Lulov, \emph{On the number of rim hook tableaux},
Zapiski Nauchnykh Seminarov POMI \textbf{223} (1995), 219--226.
Theorem1.1 and the character equality immediately following it,
p.~220; Corollary2.2 and Theorem2.3, p.~221; the complete
invoked proof in Sections1--2, pp.~219--223.
English translation: Journal of Mathematical Sciences \textbf{87}
(1997), 4118--4123, DOI\ 10.1007/BF02355806.
\url{https://www.mathnet.ru/eng/znsl4388}.

\bibitem{HD:MilneCFT2020}
J.~S.~Milne, \emph{Class Field Theory}, version4.03,
6 August 2020, author course notes. Chapter V, Example3.2, printed p.~157,
Remark3.8, printed p.~158,
and Examples3.10--3.11, printed p.~159: the ray-class-field
conductor criterion and the quadratic conductor computation.
\url{https://www.jmilne.org/math/CourseNotes/CFT.pdf}.

\bibitem{HD:JamesLiebeck2001}
Gordon James and Martin Liebeck,
\emph{Representations and Characters of Groups}, second edition,
Cambridge University Press, 2001,
\url{https://doi.org/10.1017/CBO9780511814532}.
Theorem~29.13, pp.~344--345, gives the difference of consecutive
subset permutation characters for the two-row irreducible.
Its complete proof and Proposition~29.12 were freshly read in
the institutional scan of pp.~343--345:
\url{https://www.math.colostate.edu/~hulpke/lectures/m667/JamesLiebeckSnRep.pdf}.
The publisher record confirms the second edition and print date;
no fresh proof of all preceding representation theory is claimed.

\bibitem{HD:GarsiaGoupil2009}
A.~M. Garsia and A. Goupil,
\emph{Character Polynomials, their q-Analogs and the Kronecker Product},
Electronic Journal of Combinatorics \textbf{16}(2) (2009), R19,
\url{https://doi.org/10.37236/85}.
The exact inputs are Proposition~I.1, p.~2, its complete proof
in Section~2, equations (2.1)--(2.6), pp.~9--11, and the
ordinary symmetric-group rows of Table~5.1, p.~31.
Primary published PDF:
\url{https://www.combinatorics.org/ojs/index.php/eljc/article/download/v16i2r19/pdf}.
The invoked proof and specific table formulas were freshly read;
the Hecke-algebra deformation results are not used.

\bibitem{HK:FultonHarris1991}
William Fulton and Joe Harris,
\emph{Representation Theory: A First Course},
Graduate Texts in Mathematics 129, Springer, 1991.
Proposition~3.12 and its full proof, pp.~31--32;
Exercise~4.6, p.~48, and the branching discussion, p.~59.
Institutional copy:
\url{https://websites.math.leidenuniv.nl/edixhoven/teaching/2019-2020/CIMPA/Fulton-Harris.pdf}.


\bibitem{HD:MSP2007}
T.~W. M\"uller and J.-C. Schlage-Puchta,
\emph{Character theory of symmetric groups, subgroup growth of Fuchsian
groups, and random walks}, Advances in Mathematics \textbf{213}(2)
(2007), 919--982. Theorem2 and Lemma8(ii).
DOI: \url{https://doi.org/10.1016/j.aim.2007.01.016}.
Author version: \url{https://www.mathematik.uni-rostock.de/storages/uni-rostock/Alle_MNF/Mathematik/Struktur/Lehrstuehle/Algebra/papers/Tri.pdf}.
Theorem2 is on author p.~8; Lemma8(ii) on p.~13; the explicit
4000 discussion is on pp.~16--17 and is a numerical proof sketch.
Theorem1 is not asserted with the threshold4000 in this run.

\bibitem{HD:MullerSLC2006}
T.~W. M\"uller, \emph{Character theory of symmetric groups, analysis of
long relators, and random walks}, S\'eminaire Lotharingien de
Combinatoire \textbf{56} (2006), Article B56a.
Section5.1, equation(24), p.~25; Theorem9, p.~26;
Lemma7 and equation(30), pp.~28--30.
\url{https://www.mat.univie.ac.at/~slc/wpapers/s56mueller.pdf}.
The exact squared-relative-norm normalization and its full invoked
character-orthogonality proof were read.

\bibitem{OG:Serre1977}
J.-P. Serre, \emph{Linear Representations of Finite Groups},
Graduate Texts in Mathematics 42, Springer, New York, 1977.
Chapter 12, \emph{Rationality questions}, pp.~90--101,
in particular Section 12.2, pp.~92--93, and Section 12.3,
Theorem 24 and its proof, p.~94; also Section 7.2,
Theorem 13 and Proposition 21 with proofs, pp.~56--57.
Chapter DOI: \url{https://doi.org/10.1007/978-1-4684-9458-7_12}.
Institutional full scan:
\url{https://www.sas.rochester.edu/mth/sites/doug-ravenel/otherpapers/Serre-Reps.pdf}.

\bibitem{OG:Zheng2025}
W. Zheng, \emph{Lectures on Algebra}, author lecture notes,
10 January 2025, 176 PDF pages.
Proposition 4.3.2 and proof, printed p.~145;
Proposition 4.4.2 and explicit adjunction maps, pp.~150--152;
Corollary 4.6.12 and proof, pp.~164--165;
Proposition 4.6.14 and proof, p.~165.
\url{https://server.mcm.ac.cn/~zheng/algebra.pdf}.
This citation is to author notes, not claimed to be a refereed publication.

\bibitem{PD:MilneANT}
J. S. Milne, \emph{Algebraic Number Theory},
author course notes, version 3.08, 19 July 2020, 166 pages.
Statements 8.15--8.18 and their proofs,
printed pp.~142--143, on unramified Frobenius in towers and
the cyclotomic formula.
\url{https://www.jmilne.org/math/CourseNotes/ANT.pdf}.

\bibitem{PD:Harrison2009}
J. Harrison, \emph{A formalized proof of Dirichlet's theorem
on primes in arithmetic progression},
Journal of Formalized Reasoning \textbf{2}(1) (2009), 63--83.
The exact coprime-residue infinitude statement is in
Section 1, p.~63.
DOI: \url{https://doi.org/10.6092/issn.1972-5787/1558}.
Publisher PDF: \url{https://jfr.unibo.it/article/download/1558/1078}.

\bibitem{FQ:Weibel1994}
Charles A. Weibel,
\emph{An Introduction to Homological Algebra},
Cambridge Studies in Advanced Mathematics 38,
Cambridge University Press, 1994.
Chapter~6, especially \S6.5, pp.~177--180, and \S6.9, pp.~198--201.
Chapter DOI: \texttt{10.1017/CBO9781139644136.007}.
Complete published-text academic mirror:
\url{https://www.sas.rochester.edu/mth/sites/doug-ravenel/otherpapers/weibel-homv2.pdf}.

\bibitem{FQ:HRT2001}
Robert B. Howlett, Leonard J. Rylands and Donald E. Taylor,
Matrix generators for exceptional groups of Lie type,
\emph{Journal of Symbolic Computation} 31 (2001), no.~4, 429--445.
DOI: \texttt{10.1006/jsco.2000.0431}.
Complete author version, \S3 pp.~4--12 and \S4 pp.~12--17:
\url{https://www.maths.usyd.edu.au/u/bobh/hrt.pdf}.

\bibitem{FQ:Wilson2008}
Robert A. Wilson,
\emph{Finite simple groups}, Lecture~1, author lecture notes, 2008,
p.~1 (classification list and sporadic-order range), p.~2 (alternating order).
This is identified as an author lecture source, not as a page of the
later published textbook.
\url{https://webspace.maths.qmul.ac.uk/r.a.wilson/FSG2008/lect1.pdf}.

\bibitem{FQ:AtlasMonster}
ATLAS of Finite Group Representations,
\emph{Sporadic group M (Monster)}, exact group order and multiplier data.
\url{https://brauer.maths.qmul.ac.uk/Atlas/spor/M/}.

\bibitem{HD:MondererNeftin2022}
T. Monderer and D. Neftin,
\emph{Symmetric Galois groups under specialization},
Israel Journal of Mathematics \textbf{248}(1) (2022),201--227.
Theorem1.2 and Remark2.8; the full function-field proof is Theorem5.1.
DOI: \url{https://doi.org/10.1007/s11856-022-2302-x}.
Author revised version,18 December2020:
\url{https://neftin.net.technion.ac.il/files/2020/12/symmetric_181220.pdf}.
Author pp.~2,7--8,16--17. The constant3500000 uses the classification
of finite simple groups; the stated classification-free constant
is60000000000. The authors' numerical inspection behind those
constants is used as published input and was not recomputed here.

\bibitem{FTS:Balcerzyk1997}
S. Balcerzyk,
\emph{Representation ring of the sequence of alternating groups},
Colloquium Mathematicum \textbf{74}(1) (1997),147--155.
DOI: \url{https://doi.org/10.4064/cm-74-1-147-155}.
Published institutional PDF:
\url{https://matwbn.icm.edu.pl/ksiazki/cm/cm74/cm74113.pdf}.
Exact inputs: formulas(3)--(4), pp.~148--149, and the
diagonal-hook character value in the proof of Lemma1, p.~151.
The full indexed primary text was read; subsequent direct fetch
and screenshot attempts returned403.

\bibitem{FTS:YangStaroletov2021}
N. Yang and A. Staroletov,
\emph{The minimal polynomials of powers of cycles in the ordinary
representations of symmetric and alternating groups},
Journal of Algebra and Its Applications \textbf{20}(11) (2021),2150209.
DOI: \url{https://doi.org/10.1142/S0219498821502091}.
Author version: \url{https://arxiv.org/abs/1911.06049},
version3,3 May2020, Propositions2--5 on pp.~5--6.
This is an explicitly identified author version of the published work,
not a claimed comparison with inaccessible typeset pages.

\bibitem{LQ:Collins2008}
M.~J. Collins,
\emph{Bounds for finite primitive complex linear groups},
Journal of Algebra \textbf{319}(2) (2008),759--776.
DOI: \url{https://doi.org/10.1016/j.jalgebra.2005.11.042}.
Exact input: Theorem A and its degree table, p.~761;
the complete final proof is Section5, pp.~771--774.
Published typeset PDF in the institutional Oxford repository:
\url{https://ora.ox.ac.uk/objects/uuid:e1afc89d-0fed-4be7-a361-fb53f4698959/files/mc5139af1b119028c01c68527aa980659}.

\bibitem{TD:MilneEC2021}
J.~S. Milne, \emph{Elliptic Curves}, second edition,
World Scientific,2021, x+308pp.
ISBN9789811221859 and9789811221835.
Author-hosted second-edition PDF:
\url{https://jmilne.org/math/Books/EC2.pdf}.
ChapterIV, Proposition7.13 with preceding proof, p.~140,
and Remark7.18 with complete proof, p.~141.

% This section reuses the existing keys OG:Serre1977 and BI:DHP.
% Serre: Linear Representations of Finite Groups, Springer1977,
% sections2.3--2.4,6.2 Proposition10(p48),6.4--6.5 Proposition15(p52),12.1,
% and12.3 Theorem24(p94).
% Duyan--Halasi--Podoski: Random bases for coprime linear groups,
% J.GroupTheory23(2020)133--157, Definition2.3(p137),
% Theorem4.1(2) and its complete published proof(p141).

\end{thebibliography}
\end{document}
